\documentclass[lang=cn,11pt]{elegantbook}

\title{工科数学分析}
\subtitle{学习笔记}

\author{夏同}
\institute{华南理工大学}
\date{\today}
\version{1.0}
\bioinfo{自定义}{信息}

\extrainfo{靡不有初，鲜克有终}

\setcounter{tocdepth}{3}

\logo{logo-blue.png}
\cover{cover.jpg}

% 本文档命令
\usepackage{array}
\usepackage{mathtools}
% `\oiint` is provided by `esint`; fall back to `\iint` if the package is unavailable.
\IfFileExists{esint.sty}{
  \usepackage{esint}
}{
  \providecommand{\oiint}{\iint}
}
\ExplSyntaxOn
\tl_gset:Nn \CJKttdefault { ebtt }
\ExplSyntaxOff
\setCJKmonofont[BoldFont={FangSong},ItalicFont={FangSong},BoldItalicFont={FangSong}]{FangSong}
\setCJKfamilyfont{ebfs}[BoldFont={FangSong},ItalicFont={FangSong},BoldItalicFont={FangSong}]{FangSong}
\setCJKfamilyfont{ebkai}[BoldFont={KaiTi},ItalicFont={KaiTi},BoldItalicFont={KaiTi}]{KaiTi}
\setlength{\emergencystretch}{3em}
\renewcommand*{\fangsong}{\CJKfamily{ebfs}}
\renewcommand*{\kaishu}{\CJKfamily{ebkai}}
\newcommand{\ccr}[1]{\makecell{{\color{#1}\rule{1cm}{1cm}}}}
\newenvironment{shortexenum}
  {\par\noindent\setlength{\tabcolsep}{0pt}\renewcommand{\arraystretch}{1.00}%
   \begin{tabular}{@{}>{\raggedright\arraybackslash}p{0.485\linewidth}@{\hspace{0.03\linewidth}}>{\raggedright\arraybackslash}p{0.485\linewidth}@{}}}
  {\end{tabular}\par}
\newenvironment{shorttrienum}
  {\par\noindent\setlength{\tabcolsep}{0pt}\renewcommand{\arraystretch}{1.00}%
   \begin{tabular}{@{}>{\raggedright\arraybackslash}p{0.315\linewidth}@{\hspace{0.0275\linewidth}}>{\raggedright\arraybackslash}p{0.315\linewidth}@{\hspace{0.0275\linewidth}}>{\raggedright\arraybackslash}p{0.315\linewidth}@{}}}
  {\end{tabular}\par}
\newcommand{\shortitem}[2]{#1.\enspace #2}
\newcommand{\shortchoice}[2]{#1\enspace #2}
\newenvironment{shortsingleenum}
  {\par\noindent\setlength{\tabcolsep}{0pt}\renewcommand{\arraystretch}{0.98}%
   \begin{tabular}{@{}>{\raggedright\arraybackslash}p{\linewidth}@{}}}
  {\end{tabular}\par}
\newcommand{\shortsingleitem}[2]{#1.\enspace #2\\}
\newcommand{\shortsinglechoice}[2]{#1\enspace #2\\}
\newenvironment{shortanswerenum}
  {\par\noindent\setlength{\tabcolsep}{0pt}\renewcommand{\arraystretch}{1.00}%
   \begin{tabular}{@{}>{\raggedright\arraybackslash}p{\linewidth}@{}}}
  {\end{tabular}\par}
\newcommand{\shortansweritem}[2]{#1.\enspace #2\\}
\newenvironment{compactanswerenum}
  {\begin{enumerate}\setlength{\topsep}{0pt}\setlength{\partopsep}{0pt}\setlength{\itemsep}{0pt}\setlength{\parsep}{0pt}\setlength{\parskip}{0pt}}
  {\end{enumerate}}
\newenvironment{compactenum}
  {\begin{enumerate}\setlength{\topsep}{0pt}\setlength{\partopsep}{0pt}\setlength{\itemsep}{0pt}\setlength{\parsep}{0pt}\setlength{\parskip}{0pt}}
  {\end{enumerate}}
\begin{document}

\maketitle
\frontmatter

\tableofcontents
\makeenvindex
\mainmatter

\chapter{不定积分}

\section{知识讲解}

\subsection{模法：有理分式与多项式逆元}

\begin{note}[什么是"对多项式取模"]
你已经熟悉整数取模：$17\bmod 5=2$ 意思是 $17$ 除以 $5$ 余 $2$。多项式取模完全类似：$P(x)\bmod A(x)$ 就是用 $A(x)$ 去除 $P(x)$ 后的余式。特别地，"令 $A(x)=0$"是一种快速取模的方法——把 $A(x)$ 的根代入，或者用 $A(x)=0$ 这个关系消去高次项，效果等同于做多项式除法取余。

所谓"$B(x)$ 在模 $A(x)$ 下的逆元"，就是找一个多项式 $B^{-1}(x)$ 使得 $B(x)\cdot B^{-1}(x)\equiv 1\pmod{A(x)}$（即乘积除以 $A(x)$ 余 $1$）。只要 $A(x)$ 和 $B(x)$ 互素（没有公因式），这样的逆元一定存在，可以用辗转相除法求得。
\end{note}

\par\medskip\noindent 模法处理有理分式时，关键是把系数比较改写成模运算。\par

题目要求把有理分式拆成部分分式，且分母已经分解成若干互素因式或重因式；如果直接比较系数会列出很大的方程组，就立刻改走 mod 法。

若 \(\displaystyle \frac{P(x)}{A(x)B(x)}=\frac{U(x)}{A(x)}+\frac{V(x)}{B(x)},\ (A,B)=1\)，则 \(\displaystyle U(x)B(x)+V(x)A(x)=P(x)\)，且 \(\displaystyle U(x)B(x)\equiv P(x)\pmod{A(x)}\)。所以本质是在模 $A$ 下求 $B$ 的逆元。若有重因式 $(x-a)^2$，则
\[
B=\left[(x-a)^2F(x)\right]_{x=a},
\qquad
A=\left[\frac{\mathrm d}{\mathrm{d}x}(x-a)^2F(x)\right]_{x=a}.
\]

分母需要完全因式分解，并按因式类型写出正确的部分分式形式：一次因式配常数分子，不可约二次因式配一次分子，重因式要把每一层都写全。简单线性因式可用掩盖法；重线性因式可先取最高层系数，再求导取下一层。若分母是两个互素大块 $A,B$，就对 $A$ 取模求 $U$，对 $B$ 取模求 $V$，避免硬比较所有系数。若出现 $A^mB^n$，则从一次情形的裴蜀等式出发，再升幂或连续带余除法。每个分式还要做分子降次，保证分子次数严格低于对应分母块。

标准落笔是：先写“设 $\dfrac{P}{AB}=\dfrac{U}{A}+\dfrac{V}{B}$，则 $UB+VA=P$”。然后写“模 $A$”或“令 $x=a$”，直接拿到一块分子。若有重根，再写“对 $(x-a)^2F(x)$ 求导并代入 $x=a$”。求出局部块以后，作差相消，把剩余项继续降维。

“原分母可拆成互素两块，比较系数太长，改用 mod 法。对 $A$ 取模得 $U\equiv PB^{-1}\pmod A$，从而 $U=\cdots$；同理 $V=\cdots$，故原式分解为 $\cdots$。”

若分母尚未完全因式分解，先分解再谈 mod。若两个分母块不互素，不能直接求逆元，先把公因式拆开。若题目只是单个线性因式或单个二次因式，直接凑微分往往更短。

检查三件事：部分分式写法有没有漏层；每一项分子次数是否合格；把关键根代回去时，左右两边是否真的相等。
\par\medskip


有理分式分解 $\frac{P(x)}{A(x)B(x)}=\frac{U(x)}{A(x)} + \frac{V(x)}{B(x)}$。通分后，在等式两边对 $A(x)$ 取模（也就是令 $A(x)=0$），那么 $V(x)A(x)$ 这一项消失：\(\displaystyle U(x)B(x)+V(x)A(x)=P(x)\)，于是 \(\displaystyle U(x)B(x)\equiv P(x)\pmod{A(x)}\)，也就是 \(\displaystyle U(x)\equiv P(x)\,B^{-1}(x)\pmod{A(x)}\)。
问题转化为求 $B(x)$ 在模 $A(x)$ 下的“乘法逆元” $B^{-1}(x)$（也就是多项式的倒数）。
\begin{itemize}
\item 只要 $A(x)$ 和 $B(x)$ 没公因式，必存在多项式 $u(x)$ 和 $v(x)$使${u(x)A(x) + v(x)B(x) = 1}$。
\item 遇到重因式 $(x^2+1)^2$ 时，类比把十进制数字转成二进制一样：只要把分子连续除以底座 $(x^2+1)$ 取余数，就能将其逐步拆解。
\item 二项式升幂法：遇到混合重因式 $A^m B^n$ 时，可先找到一次方情形的裴蜀等式，再将两边提升到 $(m+n-1)$ 次方展开。
\end{itemize}

\par\medskip\noindent 双元法把实圆和虚圆底座连成一条线：先认底座，再认微元。\par

\begin{note}[从一个熟悉的例子看双元法的动机]
考虑 $\displaystyle\int\frac{\mathrm{d}x}{\sqrt{1-x^2}}$。标准做法是令 $x=\sin\theta$，但换一个视角：设 $y=\sqrt{1-x^2}$，则 $x^2+y^2=1$（实圆底座）。对两边微分得 $x\,\mathrm{d}x+y\,\mathrm{d}y=0$，即 $\mathrm{d}x/y=-\mathrm{d}y/x$。而原积分恰好就是 $\int\mathrm{d}x/y$，它等于角度的累积 $-\theta+C=\arcsin x+C$。

关键观察：引入第二个变量 $y$ 并用约束 $x^2+y^2=1$ 把它们绑在一起后，微分关系自动产生，不需要再做"换元后回代"。这就是双元法的核心思想。
\end{note}

下面这些母式、微分公式和典型递推，都是同一底座的不同写法。

\par\medskip\noindent 本章内容概览。\par

第1章的题目看起来很多，其实都在做同一件事：把被积式送回已经熟悉的导数结构，或者送进少数几条基本母式。真正容易失分的地方不是算不出来，而是入口选错，例如该先凑微分时跑去分部，该先配方时直接硬换元，该用恒等变形时一上来就套万能代换。因此，下面的内容按入口类型展开；入口一旦确定，同一道题中应保持同一路径推进，不要来回试法。
\begin{itemize}
\item 入口类型包括：
\begin{compactenum}
\item 若能一眼看成 $g(\varphi(x))\varphi'(x)$，先凑微分；最常见的是 \(\displaystyle \int \frac{f'(x)}{f(x)}\,\mathrm{d}x=\ln|f(x)|+C\)，以及 \(\displaystyle \int (ax+b)^n\,\mathrm{d}x=\frac{(ax+b)^{n+1}}{a(n+1)}+C\ (n\neq -1)\)。
凡是只出现一个整体 $ax+b$、$\sqrt{ax+b}$、$\ln(ax+b)$、$\mathrm e^{ax+b}$，第一反应都是把这个整体当成新元。
\item 若是纯有理函数 $\dfrac{P(x)}{Q(x)}$，先看次数；分子次数不低于分母时先做带余除法，再做部分分式。一次因式配常数分子，不可约二次因式配一次分子，重因式每一层都要写全。
\item 若是乘积题，尤其是多项式与指数函数、三角函数、对数函数、反三角函数的乘积，先看分部积分。通常把 $\ln x$、$\arcsin x$、$\arctan x$ 这类求导会变简单的部分放进 $u$，把指数函数、三角函数或多项式放进 $\mathrm{d}v$。
\item 若含二次式或二次根式，先配方，再把它归到三种标准底座之一：\(\displaystyle a^2-x^2,\ x^2+a^2,\ x^2-a^2\)。
后面所有“含 $ax^2+bx+c$”与“含 $\sqrt{ax^2+bx+c}$”的题，都是这一步之后的继续计算。
\item 若题目反复出现 $\sqrt{a^2-x^2}$、$\sqrt{x^2+a^2}$、$\sqrt{x^2-a^2}$，或者分子分母本来就由 $x,y,\mathrm{d}x,\mathrm{d}y$ 组合出来，优先考虑双元法，不要急着在三角、双曲和代数形式之间来回折腾。
\item 若是三角函数积分，先分清它是幂次题、乘积题、不同频率题，还是 $R(\sin x,\cos x)$ 型有理式；顺序始终是“先恒等变形，再考虑万能代换”。
\item 若外层是反三角函数、对数函数或幂次对数，先考虑分部积分；若分母里正好带一个 $x$，再看能否令 $t=\ln x$。
\item 若题目是定积分，不要先求原函数，先查奇偶性、周期性、对称性、端点映射与正交性；能直接判零、化半区间或写递推，就不要多走一步。
\end{compactenum}

\item 有理函数先走“除法 $\to$ 分式 $\to$ 真分式积分”这条线。若
\[
\frac{P(x)}{A(x)B(x)}=\frac{U(x)}{A(x)}+\frac{V(x)}{B(x)},\qquad (A,B)=1,
\]
则通分后应满足 \(\displaystyle U(x)B(x)+V(x)A(x)=P(x)\)。
比较系数太长时就改用 mod 法：对 $A(x)$ 取模可得 $U(x)B(x)\equiv P(x)\pmod{A(x)}$，对 $B(x)$ 取模可得 $V(x)A(x)\equiv P(x)\pmod{B(x)}$。因此第1章里凡是分母已经拆成若干互素块、但比较系数会列出大方程组的题，都应优先想到 mod 法；若有重线性因式，就先取最高层系数，再求导取下一层。做完分解以后，每一项还要再检查一次：分子次数必须严格低于对应分母。

\item 分部积分只在“做一次会变简单”时使用。母式始终是 \(\displaystyle \int u\,\mathrm{d}v=uv-\int v\,\mathrm{d}u\)。
常见稳定选法有三组：
\begin{itemize}
\item $u=\ln x,\arcsin x,\arctan x$ 一类求导会变简单的外层函数；
\item $\mathrm{d}v=x^n\,\mathrm{d}x,\ \mathrm{e}^{ax}\,\mathrm{d}x,\ \sin bx\,\mathrm{d}x,\ \cos bx\,\mathrm{d}x$ 一类容易积分的部分；
\item 遇到 $x^n\mathrm e^{ax}$、$x^n\sin bx$、$x^n\cos bx$ 时，把多项式放进 $u$，每做一次就降一次幂。
\end{itemize}
指数与三角的乘积常用两条闭式：
\begin{align*}
\int \mathrm e^{ax}\cos bx\,\mathrm{d}x
&=\frac{\mathrm e^{ax}}{a^2+b^2}\bigl(a\cos bx+b\sin bx\bigr)+C,\\
\int \mathrm e^{ax}\sin bx\,\mathrm{d}x
&=\frac{\mathrm e^{ax}}{a^2+b^2}\bigl(a\sin bx-b\cos bx\bigr)+C.
\end{align*}
若题目不是乘积，而是纯有理函数或一眼可凑微分，就不要硬套分部积分。

\item 二次式与二次根式先配方，再按母式处理。第1章所有这类题最后都要落回下面几条基础公式：
\begin{align*}
\int \frac{\mathrm{d}x}{x^2+a^2}&=\frac{1}{a}\arctan\frac{x}{a}+C,\qquad
\int \frac{\mathrm{d}x}{x^2-a^2}=\frac{1}{2a}\ln\left|\frac{x-a}{x+a}\right|+C,\\
\int \frac{\mathrm{d}x}{\sqrt{a^2-x^2}}&=\arcsin\frac{x}{a}+C,\qquad
\int \frac{\mathrm{d}x}{\sqrt{x^2+a^2}}=\ln\bigl|x+\sqrt{x^2+a^2}\bigr|+C,\\
\int \frac{\mathrm{d}x}{\sqrt{x^2-a^2}}&=\ln\bigl|x+\sqrt{x^2-a^2}\bigr|+C.
\end{align*}
若题目是 $ax^2+b$、$ax^2+bx+c$ 或对应根式，先提出系数、完成平方，再决定它究竟属于哪一类；分类一旦判错，后面对数、反三角、双曲函数都会全错。

\item 双元法不是零散方法，而是整章的一条脉络。看到实圆底座，就设
\[
y=\sqrt{a^2-x^2},\qquad x^2+y^2=a^2,\qquad x\,\mathrm{d}x+y\,\mathrm{d}y=0.
\]
看到虚圆底座，就设
\[
y=\sqrt{x^2+a^2}\ \text{或}\ \sqrt{x^2-a^2},
\qquad
y^2-x^2=a^2\ \text{或}\ x^2-y^2=a^2,
\qquad
x\,\mathrm{d}x=y\,\mathrm{d}y.
\]
随后只做三件事：
\begin{itemize}
\item 先认母式
\[
\frac{\mathrm{d}x}{y},\qquad \frac{\mathrm{d}y}{x},\qquad x\,\mathrm{d}y-y\,\mathrm{d}x,\qquad y\,\mathrm{d}x-x\,\mathrm{d}y,
\]
不要急着做大计算；
\item 再认参数微元
\[
\frac{\mathrm{d}x}{y}=-\mathrm{d}\theta,\qquad \frac{\mathrm{d}y}{x}=\mathrm{d}\theta\quad(\text{实圆}),\qquad
\frac{\mathrm{d}x}{y}=\mathrm{d}\phi,\qquad \frac{\mathrm{d}y}{x}=\mathrm{d}\phi\quad(\text{虚圆})
\]
\item 最后再认比值微元
\[
\mathrm{d}\!\left(\frac{x}{y}\right)=\frac{y\,\mathrm{d}x-x\,\mathrm{d}y}{y^2},
\qquad
\int \frac{\mathrm{d}x}{y}=
\begin{cases}
\arctan\dfrac{x}{y}+C,& x^2+y^2=a^2,\\[4pt]
\ln(x+y)+C,& y^2-x^2=a^2\ \text{或}\ x^2-y^2=a^2.
\end{cases}
\]
\end{itemize}
第1章里大量“经典组合”、高次根式分母、纯根式题与某些三角题，本质上都是把式子压回这些微元后完成的。若引入 $y$ 后既不能化成角度微元，也不能化成面积微元或比值微元，就说明双元入口选早了，应改走别的路。

\item 遇到“经典组合”不要当新题做。像
\[
\frac{1}{(1+x^2)^n},\qquad \frac{1}{(1-x^2)^n},\qquad
\frac{x^{2n}}{\sqrt{1-x^2}},\qquad
\frac{1}{x^{2n-1}\sqrt{1-x^2}}
\]
以及与 $\sqrt{x^2-1}$ 相关的同类题，都不是新方法，它们只是前面二次根式和双元法的递推展开。做法可表述为：设 $I_n$，利用底座恒等式把高次幂拆低，必要时配合一次分部积分，把它递推回 $\int \dfrac{\mathrm{d}x}{y}$、$\int \dfrac{\mathrm{d}x}{y^3}$ 或 $\int y\,\mathrm{d}x$ 这几条母式。

\item 三角函数积分先判结构，不要上来就万能代换。本章三角题的判断结构如下：
\begin{compactenum}
\item 若是 $\sin^m x\cos^n x$，先看奇偶：有一个奇幂就拆一个出来，剩余部分用 $1-\sin^2x$ 或 $1-\cos^2x$；两者都偶时再用半角公式。
\item 若是 $\sec^m x\tan^n x$，优先找 $\sec x\tan x$ 或 $\sec^2x$；若是 $\csc^m x\cot^n x$，优先找 $\csc x\cot x$ 或 $\csc^2x$。
\item 若是 $\sin ax\sin bx$、$\sin ax\cos bx$、$\cos ax\cos bx$，先积化和差；频率不同而你还在硬乘幂，通常就已经走偏了。
\item 只有在式子整体就是 $R(\sin x,\cos x)$ 且恒等变形不顺时，才用
\[
t=\tan\frac{x}{2},\qquad
\sin x=\frac{2t}{1+t^2},\qquad
\cos x=\frac{1-t^2}{1+t^2},\qquad
\mathrm{d}x=\frac{2}{1+t^2}\,\mathrm{d}t.
\]
\end{compactenum}
幂次题、递推题、正交题和 Wallis 递推题，本质上都在这一条线上，不要把它们拆成互不相干的零碎方法。

\item 反三角、指数、对数、双曲函数都各有常用入口。
\begin{itemize}
\item 反三角函数积分先分部，再把余项送回实圆根式或二次母式；例如 $\arcsin,\arccos$ 的导数把题目送回 $\sqrt{a^2-x^2}$，$\arctan$ 的导数把题目送回 $x^2+a^2$。
\item 指数函数积分先分清是纯指数、指数乘多项式，还是指数乘三角；纯指数直接积，指数乘多项式用分部降幂，指数乘三角直接用前面的闭式或两次分部积分。
\item 对数函数积分要盯住 \(\displaystyle \int \ln x\,\mathrm{d}x=x\ln x-x+C,\quad \int \frac{\mathrm{d}x}{x\ln x}=\ln|\ln x|+C\)。
若 \(\displaystyle I_{m,n}=\int x^m(\ln x)^n\,\mathrm{d}x,\ m\neq -1\)，则直接用
\(\displaystyle I_{m,n}=\frac{x^{m+1}}{m+1}(\ln x)^n-\frac{n}{m+1}I_{m,n-1}\)。
若 $m=-1$，立刻令 $t=\ln x$，不要再分部。
\item 双曲函数积分直接把 \(\displaystyle p=\cosh x,\ q=\sinh x,\ p^2-q^2=1\) 当成虚圆底座，记住 \(\displaystyle (\sinh x)'=\cosh x,\ (\cosh x)'=\sinh x,\ \int \tanh x\,\mathrm{d}x=\ln\cosh x+C\)。
若只是低次幂，先凑导数；若是平方，先半角；若天然长成 $\cosh,\sinh$ 的组合，就按虚圆底座处理。
\end{itemize}

\item 定积分先做“结构检查”，后做计算。对称区间先判奇偶性；整周期先判周期性与正交性；三角乘积先看是否能直接用积化和差或复指数。要做 Wallis 递推时，先写 \(\displaystyle p=\sin x,\ q=\cos x,\ p^2+q^2=1\)，
再用端点 $x=0,\frac{\pi}{2}$ 消边界项。定积分里双元法的价值主要在于暴露端点关系和递推关系，而不是机械地“先求原函数再代上下限”。

\item \textbf{常见遗漏。} 不定积分一定补上 $+C$；定积分不能补 $+C$。若中途除过某个因子、取过倒数、开过方、设过新元，就要回头补查零解、常值解或特殊支；若题目来自根式或对数，还要核对定义域、绝对值以及回代后的变量是否已经全部换回原来的 $x$。这一层收束的重点，是把所有可能被代换掩盖掉的分支重新带回原题。
\end{itemize}
\par\medskip

\par\medskip\noindent\textbf{双元微分公式}\par
\begin{theorem}[微分公式]
\begin{itemize}
\item 实圆：$x\mathrm{d}x+y\mathrm{d}y=0$；虚圆：$x\mathrm{d}x=y\mathrm{d}y$；
\item 参数微元，实圆：$\frac{\mathrm{d}x}{y}=-\mathrm{d}\theta, \quad \frac{\mathrm{d}y}{x}=\mathrm{d}\theta$；
\item 参数微元，虚圆：$\frac{\mathrm{d}x}{y}=\mathrm{d}\phi,\quad \frac{\mathrm{d}y}{x}=\mathrm{d}\phi$；
\item 面积2倍的微元，$y\mathrm{d}x-x\mathrm{d}y=(x^2\pm y^2)\mathrm{d}\theta$，注意顺序和微分的顺序一致；
\item 比值微元，$\mathrm{d}\left(\frac{x}{y}\right)=\frac{y\mathrm{d}x-x\mathrm{d}y}{y^2}$；
\end{itemize}
\end{theorem}
\begin{enumerate}
    \item 实圆：$x\mathrm{d}x+y\mathrm{d}y=0$；虚圆：$x\mathrm{d}x=y\mathrm{d}y$；几何解释如下：
\begin{itemize}
\item 实圆：瞬时位移向量 $(\mathrm{d}x, \mathrm{d}y)$ 必须垂直于位置向量$(x,y)$，二者内积为零（两个向量 $(a,b)$ 与 $(c,d)$ 的内积定义为 $ac+bd$，内积为零即互相垂直）
\item 虚圆：等价于：$\frac{\mathrm{d}x}{x} = \frac{\mathrm{d}y}{y}  $，即 $x$ 与 $y$ 的相对变化率（对数微分）必须严格相等。
\end{itemize}
\item 参数微元，实圆：$\frac{\mathrm{d}x}{y}=-\mathrm{d}\theta,\quad \frac{\mathrm{d}y}{x}=\mathrm{d}\theta$；虚圆：$\frac{\mathrm{d}x}{y}=\mathrm{d}\phi,\quad \frac{\mathrm{d}y}{x}=\mathrm{d}\phi$；几何解释如下：
\begin{itemize}
\item 实圆:$ \mathrm{d}\theta  $ 是位置向量绕原点逆时针转过的有向小角（数学正方向）。此时指针尖端的位移大小等于弧长 $  a\mathrm{d}\theta  $，方向正好是 $  \vec{r}  $ 逆时针旋转 90° 得到的向量 $  (-y, x)  $（将向量 $(a,b)$ 逆时针转 90° 就是交换分量并取反第一个：$(-b,a)$），模也是$a$。所以 \(\displaystyle (\mathrm{d}x, \mathrm{d}y)=\mathrm{d}\theta(-y,x)\)。
\item 虚圆：$\mathrm{d}\phi$ 是位置向量绕原点逆时针转过的有向小角，此时指针尖端的位移大小等于 $  a\mathrm{d}\phi  $，方向正好与 $\vec{r}$ 沿$y=x$对称的向量$(y,x)$一致。所以 \(\displaystyle (\mathrm{d}x, \mathrm{d}y)=\mathrm{d}\phi(y,x)\)。
\end{itemize}
\item 面积2倍的微元，利用扇形面积两倍等于半径平方乘以角度（弧度）即得，注意方向。面积微元可以等比例地转化为角度微元，角度微元可以等比例地转化为面积微元。
\begin{itemize}
\item $x\mathrm{d}y-y\mathrm{d}x=(x^2\pm y^2)\mathrm{d}\theta$（逆向面积），对应逆时针角度微元$\mathrm{d}\theta$，即实圆关系的$\frac{\mathrm{d}y}{x}$，虚圆关系的$\frac{\mathrm{d}y}{x},\frac{\mathrm{d}x}y$
\item $y\mathrm{d}x-x\mathrm{d}y=(x^2\pm y^2)(-\mathrm{d}\theta)$（顺向面积）对应顺时针角度微元$-\mathrm{d}\theta$，即实圆关系的$\frac{\mathrm{d}x}{y}$.
\end{itemize}
这些公式并不难记忆。实际运算时，可以先补齐分母中的 $x^2\pm y^2$，再观察分子与分母的对应关系：
\begin{itemize}
\item 对于$\frac{x\mathrm{d}y-y\mathrm{d}x}{x^2+y^2}$，观察到分子左侧的$x\mathrm{d}y$和分母的$x^2$，相除就是$\frac{\mathrm{d}y}{x}$，所以这是$\frac{\mathrm{d}y}{x},-\frac{\mathrm{d}x}{y}$的合分比之后的结果。
\item 对于$\frac{y\mathrm{d}x-x\mathrm{d}y}{x^2+y^2}$，观察到分子左侧的$y\mathrm{d}x$和分母的$y^2$（无关哪一侧，加法满足交换律），相除就是$\frac{\mathrm{d}x}{y}$，所以这是$\frac{\mathrm{d}x}{y},-\frac{\mathrm{d}y}{x}$的合分比之后的结果。
\item 对于$\frac{x\mathrm{d}y-y\mathrm{d}x}{x^2-y^2}$，观察到分子左侧的$x\mathrm{d}y$和分母的$x^2$，相除就是$\frac{\mathrm{d}y}{x}$，所以这是$\frac{\mathrm{d}y}{x},\frac{\mathrm{d}x}{y}$的合分比之后的结果。
\item 对于$\frac{y\mathrm{d}x-x\mathrm{d}y}{x^2-y^2}$，观察到分子左侧的$y\mathrm{d}x$和分母的$-y^2$，相除就是$-\frac{\mathrm{d}x}{y}$，所以这是$-\frac{\mathrm{d}x}{y},-\frac{\mathrm{d}y}{x}$的合分比之后的结果。
\end{itemize}
\item 比值微元，$\mathrm{d}\left(\frac{x}{y}\right)=\frac{y\mathrm{d}x-x\mathrm{d}y}{y^2}$；这个式子经常反过来用作积分降次。它不仅可以由代数商法则得到，也可以从\textbf{中心投影与面积的相似缩放}作几何解释：
\begin{itemize}
\item \textbf{比值即投影坐标：} 比值 $\frac{x}{y}$ 在几何上代表什么？想象平面上有一条高度归一化的水平参考线 $y=1$。将原点与动点 $P(x,y)$ 连成一条射线，交该辅助线于投影点 $P'$。根据相似三角形可知，投影点 $P'$ 的坐标恰为 $\left(\frac{x}{y}, 1\right)$。也就是说，比值 $\frac{x}{y}$ 本质上是动点向 $y=1$ 直线作中心投影后的\textbf{横坐标}。
\item \textbf{投影扫过的面积：} 当投影点 $P'$ 在直线上发生纯水平的微小位移 $\mathrm{d}\left(\frac{x}{y}\right)$ 时，它绕原点构成的微小投影三角形，底为 $\mathrm{d}\left(\frac{x}{y}\right)$，高恒为 1。用二维外积计算其逆时针 2 倍有向面积，恰为横纵交叉相乘：$\frac{x}{y}\cdot 0 - 1\cdot\mathrm{d}\left(\frac{x}{y}\right) = -\mathrm{d}\left(\frac{x}{y}\right)$。
\item \textbf{相似缩放原理：} 实际动点 $P(x,y)$ 和投影点 $P'$ 始终在同一条射线上，两者的向径距离缩放比例刚好是纵坐标之比（$y$ 倍）。根据平面几何基本定理\textbf{“相似图形的面积比等于相似比的平方”}，实际动点扫过的面积必定等于投影点扫过面积的 $y^2$ 倍！
\item \textbf{公式推导：} \textbf{实际的 2 倍面积} $= y^2 \times$ \textbf{投影的 2 倍面积}。即
\(\displaystyle x\mathrm{d}y - y\mathrm{d}x = y^2 \left[ -\mathrm{d}\left(\frac{x}{y}\right) \right]\)，
将负号移入括号并整理，即得 $\mathrm{d}\left(\frac{x}{y}\right)=\frac{y\mathrm{d}x-x\mathrm{d}y}{y^2}$。
\end{itemize}
\end{enumerate}

\par\medskip\noindent\textbf{双元法：底座与微元}\par

题面反复出现 $\sqrt{a^2-x^2}$、$\sqrt{x^2+a^2}$、$\sqrt{x^2-a^2}$、三角函数、双曲函数，或者被积式已经天然长成 $x,y,\mathrm{d}x,\mathrm{d}y$ 的组合。尤其是分子里出现 $x\,\mathrm{d}y$、$y\,\mathrm{d}x$、$\mathrm{d}x/y$、$\mathrm{d}y/x$ 时，优先考虑双元法。

实圆底座：
\begin{align*}
y&=\sqrt{a^2-x^2},\qquad x^2+y^2=a^2,\qquad x\,\mathrm{d}x+y\,\mathrm{d}y=0,\\
\frac{\mathrm{d}x}{y}&=-\mathrm{d}\theta,\qquad
\frac{\mathrm{d}y}{x}=\mathrm{d}\theta,\qquad
x\,\mathrm{d}y-y\,\mathrm{d}x=(x^2+y^2)\,\mathrm{d}\theta.
\end{align*}
虚圆底座：
\begin{align*}
y&=\sqrt{x^2+a^2}\text{ 或 }\sqrt{x^2-a^2},\qquad y^2-x^2=a^2\text{ 或 }x^2-y^2=a^2,\\
x\,\mathrm{d}x&=y\,\mathrm{d}y,\qquad
\frac{\mathrm{d}x}{y}=\mathrm{d}\phi,\qquad
\frac{\mathrm{d}y}{x}=\mathrm{d}\phi,\qquad
\mathrm{d}\!\left(\frac{x}{y}\right)=\frac{y\,\mathrm{d}x-x\,\mathrm{d}y}{y^2}.
\end{align*}

如果根式是 $a^2-x^2$，就建实圆底座；如果根式是 $x^2+a^2$ 或 $x^2-a^2$，就建虚圆底座。建完底座后不要急着积分，先判分子分母更像哪一种微元：像 $\mathrm{d}x/y$、$\mathrm{d}y/x$ 就走角度微元；像 $x\,\mathrm{d}y-y\,\mathrm{d}x$ 就走面积微元；像 $y\,\mathrm{d}x-x\,\mathrm{d}y$ 除以平方项，就走比值微元。最后再回代到 $x$。

常见推导是“设 $y=\sqrt{\cdots}$，写出底座恒等式，微分，选微元，积分，回代”。例如先写 \(\displaystyle y=\sqrt{\cdots},\ x^2\pm y^2=a^2\)，
再写微分关系，随后把原积分强行改写成 $\mathrm{d}\theta$、$\mathrm{d}\phi$、$\mathrm{d}(x/y)$ 或 $\mathrm{d}(y/x)$。

“设 $y=\sqrt{\cdots}$，则有 $x^2\pm y^2=a^2$，并且 $x\,\mathrm{d}x\pm y\,\mathrm{d}y=0$。原式化为 $\cdots=\int \mathrm{d}\theta$（或 $\int \mathrm{d}\phi,\ \int \mathrm{d}(x/y)$），故结果为 $\cdots$。”

若题目一眼就是标准三角代换或部分分式，别为用双元而用双元。若引入 $y$ 后分子分母仍然无法落到角度微元、面积微元、比值微元之一，就换方法。若是定积分，别忘了端点也要一起映射到双元坐标。

检查三件事：底座恒等式有没有写对号；微分关系有没有丢负号；回代后答案是否只剩原变量 $x$。
\par\medskip

 
\par\medskip\noindent\textbf{双元母式}\par
\begin{theorem}[$\int\frac{\mathrm{d}x}{y}$型]
\begin{itemize}
    \item 实圆：$\int\frac{\mathrm{d}x}{y}=\arctan\frac{x}{y}$；
    \item 虚圆：$\int \frac{\mathrm{d}x}{y}=\ln(x+y)$.
\end{itemize}
\end{theorem}
\par\medskip\noindent\textbf{虚圆参数坐标}\par

需要处理 $\sqrt{x^2+a^2}$、$\sqrt{x^2-a^2}$、$\ln(x+\sqrt{x^2+a^2})$、$\ln(x+\sqrt{x^2-a^2})$，或者双元法已经走到虚圆底座时，就直接调用这张表。

\begin{align*}
\cosh t&=\frac{e^t+e^{-t}}{2},\qquad
\sinh t=\frac{e^t-e^{-t}}{2},\qquad
\tanh t=\frac{\sinh t}{\cosh t},\\
\cosh^2 t-\sinh^2 t&=1,\qquad
(\cosh t)'=\sinh t,\qquad
(\sinh t)'=\cosh t,\\
\operatorname{arsinh}u&=\ln\!\left(u+\sqrt{u^2+1}\right),\qquad
\operatorname{arccosh}u=\ln\!\left(u+\sqrt{u^2-1}\right),\\
x&=a\cosh t,\ y=a\sinh t\Rightarrow x+y=a\,e^t,\qquad
t=\ln\frac{x+y}{a}.
\end{align*}

如果根式是 $x^2+a^2$，优先把答案记成 $\ln(x+\sqrt{x^2+a^2})$ 或 $\operatorname{arsinh}(x/a)$。如果根式是 $x^2-a^2$，优先记成 $\ln(x+\sqrt{x^2-a^2})$ 或 $\operatorname{arccosh}(x/a)$。若题目已经走到虚圆底座，就把参数角直接看成 $t$，不要再额外造新字母。

先写 $x=a\cosh t,\ y=a\sinh t$，再把 $\mathrm{d}x=y\,\mathrm{d}t$ 和 $x+y=a e^t$ 一起用上。若只需要最终结果，也可以直接把 $\int \mathrm{d}x/y$ 写成 $\ln(x+y)+C$。

“取 $x=a\cosh t,\ y=a\sinh t$，则 $\mathrm{d}x=y\,\mathrm{d}t$，故原式化为 $\int \mathrm{d}t=t+C=\ln(x+y)+C$。”

若题目本身更适合三角代换或部分分式，不要强行双曲化。若出现的是 $a^2-x^2$，这是实圆不是虚圆，不能套双曲函数表。

检查对数里是否写成了 $x+y$；检查常数 $a$ 是否被正确吸收到积分常数里。
\par\medskip


对于虚圆 \(x^2 - y^2 = 1\)（或一般常数），参数方程可写为 \(x = \cosh \phi\)，\(y = \sinh \phi\)，那么
\(\displaystyle x + y = \cosh \phi + \sinh \phi = e^\phi,\quad\phi = \ln(x + y)\)。
而微分 \(\frac{\mathrm{d}x}{y} = \mathrm{d}\phi\)（因为 \(\mathrm{d}x = \sinh \phi \, \mathrm{d}\phi = y \, \mathrm{d}\phi\)），因此
\(\displaystyle \int \frac{\mathrm{d}x}{y} = \phi + C = \ln(x + y) + C\)。
在双曲线上，\(\phi\) 是“双曲角度”（类似圆中的角度），而 \(x+y\) 恰好是双曲线上的一个自然指数坐标，其对数就是双曲角。想象双曲线 \(x^2 - y^2 = 1\)，它由两支组成，通常考虑右支。对于右支上的点，\(x > 0\)，且 \(x+y > 0\)，随着 \(\phi\) 增大，点沿双曲线向右上方移动，\(x+y\) 指数增长。所以 \(\ln(x+y)\) 度量了沿双曲线移动的“路程”或“参数增量”。

对于实圆 \(x^2 + y^2 = 1\)，最自然的参数是角度 \(\theta\)，例如取 \(x = \cos\theta\)，\(y = \sin\theta\)。此时
\(\displaystyle \frac{\mathrm{d}x}{y} = -\mathrm{d}\theta\)，而
\(\displaystyle \arctan\frac{x}{y} = \arctan(\cot\theta) = \frac{\pi}{2} - \theta\)
（在适当范围内），故
\(\displaystyle \int \frac{\mathrm{d}x}{y} = -\theta + C = \arctan\frac{x}{y} + \left(C - \frac{\pi}{2}\right)\)。
即两者相差常数，所以不定积分公式成立。

\begin{note}[母式来源]
$\displaystyle \int \frac{\mathrm{d}x}{y}$ 是整个双元法里最像“母式”的一条。实圆底座里，它本质上是在积累角度，所以答案自然落到反三角函数；虚圆底座里，它本质上是在积累双曲角，而 $x+y$ 恰好扮演指数坐标，所以答案自然落到对数。

因此记忆时抓住“圆看角，虚圆看对数坐标”就够了。后面大量公式其实只是把更复杂的被积式重新压回这一条母式或它的比值变体。
\end{note}

\begin{theorem}[$\int \frac{\mathrm{d}x}{y^3}$型]

\begin{align*}
\int \frac{\mathrm{d}x}{y^3} &=\int\frac{1}{y^2}\frac{\mathrm{d}x}{y}=\int\frac{1}{y^2}\frac{y\mathrm{d}x-x\mathrm{d}y}{y^2\pm x^2} =\frac{1}{y^2\pm x^2}\int\frac{y\mathrm{d}x-x\mathrm{d}y}{y^2}=\frac{1}{y^2\pm x^2}\int\mathrm{d}\left(\frac{x}{y}\right)\\
\int \frac{\mathrm{d}x}{xy^2}&=\int\frac1{xy}\frac{\mathrm{d}x}{y}=\int\frac1{xy}\frac{y\mathrm{d}x-x\mathrm{d}y}{y^2\pm x^2}=\frac{1}{y^2\pm x^2}\int\left(\frac{\mathrm{d}x}{x}-\frac{\mathrm{d}y}{y}\right)\\
\int \frac{\mathrm{d}x}{x^2y}&=\int\frac{1}{x^2}\frac{\mathrm{d}x}{y}=\int\frac{1}{x^2}\frac{y\mathrm{d}x-x\mathrm{d}y}{y^2\pm x^2}=\frac{-1}{y^2\pm x^2}\int\mathrm{d}\left(\frac{y}{x}\right)
\end{align*}
\end{theorem}
这里的关键在于：先判断基本公式不能直接适用，再借助微分关系把积分转化为可处理的类型。

\begin{theorem}[$\int x\mathrm{d}y$型]

\begin{align*}
\int x\mathrm{d}y&=\int\frac12\left(x\mathrm{d}y+y\mathrm{d}x+x\mathrm{d}y-y\mathrm{d}x\right)=\frac12\int\mathrm{d}(xy)+\frac{x^2\pm y^2}2\int\mathrm{d}\frac{x}{y}\\
\int y\mathrm{d}x&=\int\frac12\left(y\mathrm{d}x+x\mathrm{d}y+y\mathrm{d}x-x\mathrm{d}y\right)=\frac12\int\mathrm{d}(xy)+\frac{x^2\pm y^2}2\int\mathrm{d}\frac{y}{x}
\end{align*}
\end{theorem}
这里的处理要点是先把被积式拆成两部分，再分别构造成全微分。

\begin{theorem}[$\int\frac{x\mathrm{d}y}{y}$型]

\[
\int \frac{\mathrm{d}x}{xy}=\frac{1}{C}\int\frac{x^2\pm y^2}{xy}\mathrm{d}x
=\frac1{C}\left(\int\frac{x\mathrm{d}x}{y}\pm \int\frac{y\mathrm{d}x}{x}\right),\qquad
\int\frac{x\mathrm{d}y}{y}=\int\frac{xy\mathrm{d}y}{y^2}.
\]
再利用关系式 $y\mathrm{d}y=\mp x\mathrm{d}x$，即可把分子中的 $y\mathrm{d}y$ 与分母中的 $y^2$ 一并转化，从而化为前面已经处理过的基本类型。
\end{theorem}

\par\medskip\noindent\textbf{三次分母与对称拆分}\par

看到 $\dfrac{\mathrm{d}x}{y^3}$、$\dfrac{\mathrm{d}x}{x^2y}$、$\dfrac{\mathrm{d}x}{xy^2}$、$x\,\mathrm{d}y$、$y\,\mathrm{d}x$ 这类既不像母式、又明显带平方或立方分母的题，就直接调用这一套。

若底座常数记为 $C=x^2\pm y^2$ 或 $C=y^2\pm x^2$，则优先凑
\[
\int \frac{\mathrm{d}x}{y^3}
=\frac{1}{C}\int \mathrm{d}\!\left(\frac{x}{y}\right),\qquad
\int \frac{\mathrm{d}x}{x^2y}
=-\frac{1}{C}\int \mathrm{d}\!\left(\frac{y}{x}\right).
\]
若遇到 $x\,\mathrm{d}y$ 或 $y\,\mathrm{d}x$，先拆
\[
x\,\mathrm{d}y=\frac12\,\mathrm{d}(xy)+\frac12(x\,\mathrm{d}y-y\,\mathrm{d}x),\qquad
y\,\mathrm{d}x=\frac12\,\mathrm{d}(xy)+\frac12(y\,\mathrm{d}x-x\,\mathrm{d}y).
\]

如果分母是三次根式，就别硬分部，先查能否凑成 $\mathrm{d}(x/y)$ 或 $\mathrm{d}(y/x)$。如果分子是 $x\,\mathrm{d}y$ 或 $y\,\mathrm{d}x$，先做“对称部分 $+\,$反对称部分”拆分，再把反对称部分送回面积微元。

先把底座常数补出来，再把分子凑成比值微元或面积微元。例如先写
\[
\frac{1}{y^3}\,\mathrm{d}x
=\frac{1}{C}\frac{y\,\mathrm{d}x-x\,\mathrm{d}y}{y^2},
\]
也可先写
\[
x\,\mathrm{d}y
=\frac12\,\mathrm{d}(xy)+\frac12(x\,\mathrm{d}y-y\,\mathrm{d}x).
\]

“由底座常数关系 $x^2\pm y^2=C$，原式可化为 $\dfrac1C\int \mathrm{d}(x/y)$（或 $\dfrac{-1}{C}\int \mathrm{d}(y/x)$），故积分立刻化为初等形式。”

若分母幂次不高、母式已经能直接积分，就不要过度拆分。若凑出的比值微元仍然带多余变量，说明还没补够底座常数，先回头补常数再说。

检查常数 $C$ 来自哪条底座恒等式；检查对称拆分后的两个部分是否都已处理完。
\par\medskip


\begin{theorem}[$\int\dfrac{\mathrm{d}x}{y^{2n}}$型]
对于负偶次幂，可结合常数关系代换、微分代换与分部积分逐步降次。
\end{theorem}
例如：
\begin{align*}
\int\dfrac{\mathrm{d}x}{y^4}&=\frac{1}{C}\int\frac{x^2\pm y^2}{y^4}\mathrm{d}x=\frac{1}{C}\left(\int\frac{x^2\mathrm{d}x}{y^4}\pm\int\frac{\mathrm{d}x}{y^2}\right)\\
&=\begin{cases}\frac{1}{C}\left(\int\frac{xy\mathrm{d}y}{y^4}\pm\int\frac{\mathrm{d}x}{y^2}\right)=\frac{1}{C}\left(\int\frac{x\mathrm{d}y}{y^3}\pm\int\frac{\mathrm{d}x}{y^2}\right)\\\frac{1}{C}\left(\int\frac{-xy\mathrm{d}y}{y^4}\pm\int\frac{\mathrm{d}x}{y^2}\right)=\frac{1}{C}\left(\int\frac{-x\mathrm{d}y}{y^3}\pm\int\frac{\mathrm{d}x}{y^2}\right)\end{cases}
\end{align*}
对于 $\displaystyle\int\frac{x\mathrm{d}y}{y^3}$ 这类积分，还可以再做一次分部积分：
\[
\int\frac{x\mathrm{d}y}{y^3}
=\int x\mathrm{d}\left(\frac1{-2}\frac{1}{y^2}\right)
=-\frac12\int x\mathrm{d}\left(\frac{1}{y^2}\right)
=-\frac12\frac{x}{y^2}+\frac12\int\frac{\mathrm{d}x}{y^2}.
\]
当分母的幂次更高时，也可沿着这一思路反复降次。

\par\medskip\noindent\textbf{负偶幂降次}\par

出现 $\dfrac{\mathrm{d}x}{y^{2n}}$、$\dfrac{x\,\mathrm{d}y}{y^{2n-1}}$、$\dfrac{y\,\mathrm{d}x}{y^{2n-1}}$ 这类高次负幂时，直接沿这一思路处理。

降次核心不是背最终答案，而是背这一步：
\[
\frac{1}{y^{2n}}
=\frac{x^2\pm y^2}{C\,y^{2n}}
=\frac{x^2}{C\,y^{2n}}\pm \frac{1}{C\,y^{2n-2}},
\]
其中 $C=x^2\pm y^2$ 或 $C=y^2\pm x^2$ 是底座常数。随后再用 \(\displaystyle x\,\mathrm{d}x=\mp y\,\mathrm{d}y\quad\text{或}\quad x\,\mathrm{d}x=y\,\mathrm{d}y\)
把高次幂继续往下拉。

降次时要补入底座常数，把一部分压回低一阶的 $\dfrac{1}{y^{2n-2}}$，另一部分变成 $x\,\mathrm{d}y$、$y\,\mathrm{d}x$ 或比值微元；这个过程反复进行，直到回到 $\int \mathrm{d}x/y$、$\int \mathrm{d}x/y^3$ 这类已经会做的母式。

常见推导是“补常数 $\to$ 分拆 $\to$ 降次 $\to$ 回母式”。不要一上来就分部积分，因为高次负偶幂真正值钱的是底座恒等式。

“由底座常数 $x^2\pm y^2=C$，先将 $1/y^{2n}$ 拆成 $x^2/(C y^{2n})$ 与 $1/(C y^{2n-2})$ 两部分；前者借微分关系降次，后者直接回到较低幂次积分。”

若幂次本来就低到可以直接母式处理，就不要再机械降次。若补常数后没有出现可用的微分结构，说明底座选错了。

检查每次降次后幂次确实下降了；检查最后是否真的回到了已经会做的母式。
\par\medskip


\begin{theorem}[$\int y^{2n-1}\mathrm{d}x$型]

设 $C=y^2\pm x^2$，则
\begin{align*}
\int \frac{\mathrm{d}x}{y}&=\int \frac{\mathrm{d}x}{y}\\
\int y\mathrm{d}x&=\frac12 xy+\frac12 C\int \frac{\mathrm{d}x}{y}\\
\int y^3\mathrm{d}x&=\frac14 xy^3+\frac14\frac32 Cxy+\frac14\frac32 C^2\int\frac{\mathrm{d}x}{y}\\
\int y^5\mathrm{d}x&=\frac16 xy^5+\frac16\frac54 Cxy^3+\frac16\frac54\frac32 C^2xy +\frac16\frac54\frac32 C^3\int\frac{\mathrm{d}x}{y}\\
\int y^7\mathrm{d}x&=\frac18 xy^7+\frac18\frac76 Cxy^5+\frac18\frac76\frac54 C^2xy^3+\frac18\frac76\frac54\frac32 C^3xy +\frac18\frac76\frac54\frac32 C^4\int\frac{\mathrm{d}x}{y}
\end{align*}
由此可以归纳出一般公式
\begin{align*}
\int y^{2n-1}\mathrm{d}x
&=\sum_{k=0}^{n-1}\frac{(2n-1)!!\left(2n-2k-2\right)!!}{(2n)!!\left(2n-2k-1\right)!!}\left(y^2\pm x^2\right)^kxy^{2n-1-2k}\\
&\quad+\frac{(2n-1)!!}{(2n)!!}\left(y^2\pm x^2\right)^n\int\frac{\mathrm{d}x}{y}
\end{align*}
\end{theorem}
设双元底座常数为 $C = y^2 \pm x^2$。对该恒等式两边求微分，得
\(\displaystyle 2y \mathrm{d}y \pm 2x \mathrm{d}x = 0 \implies y \mathrm{d}y = \mp x \mathrm{d}x\)。
再由 $C=y^2\pm x^2$ 可得 \(\displaystyle \mp x^2 = y^2 - C\)。

现设目标积分为 $\displaystyle I_n = \int y^{2n-1} \mathrm{d}x$（其中 $n \ge 1$）。

\textbf{第一步：分部积分。}
对 $I_n$ 作分部积分，取 $u=y^{2n-1},\ \mathrm{d}v=\mathrm{d}x$：
\[
I_n = \int y^{2n-1} \mathrm{d}x = x y^{2n-1} - \int x \mathrm{d}\left(y^{2n-1}\right) = x y^{2n-1} - (2n-1) \int x y^{2n-2} \mathrm{d}y.
\]

\textbf{第二步：利用微分关系降次。}
由 $y\,\mathrm{d}y=\mp x\,\mathrm{d}x$，有
\[
\int x y^{2n-2} \mathrm{d}y
= \int x y^{2n-3} (y \mathrm{d}y)
= \int x y^{2n-3} (\mp x \mathrm{d}x)
= \int y^{2n-3} (\mp x^2) \mathrm{d}x.
\]
再由 $\mp x^2 = y^2 - C$ 消去 $x^2$，得
\[
\int y^{2n-3} (\mp x^2) \mathrm{d}x
= \int y^{2n-3} (y^2 - C) \mathrm{d}x
= \int y^{2n-1} \mathrm{d}x - C \int y^{2n-3} \mathrm{d}x
= \mathbf{I_n - C I_{n-1}}.
\]

\textbf{第三步：建立递推关系。}
将上式代回第一步的分部积分公式，得
\(\displaystyle I_n=x y^{2n-1}-(2n-1)(I_n-CI_{n-1})
=x y^{2n-1}-(2n-1)I_n+(2n-1)CI_{n-1}\)。
整理得
\(\displaystyle 2n I_n=x y^{2n-1}+(2n-1)C I_{n-1}\)，即
\(\displaystyle \mathbf{I_n=\frac{1}{2n}x y^{2n-1}+\frac{2n-1}{2n}C I_{n-1}}\)。
这就是所需的递推公式。

\textbf{第四步：递推展开与双阶乘表示。}
将递推式逐次展开，可得：
\begin{align*}
I_n &= \frac{1}{2n} x y^{2n-1} + \frac{2n-1}{2n} C \left( \frac{1}{2n-2} x y^{2n-3} + \frac{2n-3}{2n-2} C I_{n-2} \right) \\
&= \frac{1}{2n} x y^{2n-1} + \frac{2n-1}{2n(2n-2)} C x y^{2n-3} + \frac{(2n-1)(2n-3)}{2n(2n-2)} C^2 I_{n-2} \\
&= \dots
\end{align*}
观察第 $k$ 次降次（$k=0,1,\dots,n-1$）对应项的系数，分子是奇数连乘，分母是偶数连乘。将其写成双阶乘（$!!$）形式时，只需补齐缺失因子：
\[
\text{第 } k \text{ 项系数}
=\frac{(2n-1)(2n-3)\cdots(2n-2k+1)}{(2n)(2n-2)\cdots(2n-2k)}
=\mathbf{\frac{(2n-1)!!(2n-2k-2)!!}{(2n)!!(2n-2k-1)!!}}.
\]
\begin{itemize}
    \item 分子补齐：$\displaystyle (2n-1)(2n-3)\cdots(2n-2k+1) = \frac{(2n-1)!!}{\mathbf{(2n-2k-1)!!}} $
    \item 分母补齐：$\displaystyle (2n)(2n-2)\cdots(2n-2k) = \frac{(2n)!!}{\mathbf{(2n-2k-2)!!}} $
\end{itemize}
对应的代数项为 $C^k x y^{2n-1-2k}$，将 $C$ 替换回 $y^2 \pm x^2$，即 $\mathbf{\left(y^2 \pm x^2\right)^k x y^{2n-1-2k}}$。

\textbf{第五步：写出末项。}
当递推进行到 $k=n$ 时，剩余积分与累积系数为
\(\displaystyle I_0=\int y^{-1}\,\mathrm{d}x=\int\frac{\mathrm{d}x}{y}\)，
\(\displaystyle \frac{(2n-1)(2n-3)\cdots 1}{(2n)(2n-2)\cdots 2}
=\mathbf{\frac{(2n-1)!!}{(2n)!!}}\)。
同时伴随因子 $C^n = \left(y^2 \pm x^2\right)^n$。

综上，得到通式
\begin{align*}
\int y^{2n-1}\mathrm{d}x =& \sum_{k=0}^{n-1}\frac{(2n-1)!!\left(2n-2k-2\right)!!}{(2n)!!\left(2n-2k-1\right)!!}\left(y^2\pm x^2\right)^k x y^{2n-1-2k} \\
&+ \frac{(2n-1)!!}{(2n)!!}\left(y^2\pm x^2\right)^n\int\frac{\mathrm{d}x}{y}
\end{align*}

\begin{note}[双阶乘系数从哪来]
这条通式看上去长，但结构并不复杂。每做一次递推，都会固定地乘上一个“奇数比偶数”的系数，并把幂次降两阶，所以双阶乘并不是突然冒出来的，而只是把这一串重复乘积压缩写成了记号。

因此真正该记的不是整条公式，而是递推骨架 \(\displaystyle I_n=\frac{1}{2n}xy^{2n-1}+\frac{2n-1}{2n}CI_{n-1}\)。一旦知道每一步都长这个样子，就知道最后为什么一定会出现奇数连乘、偶数连乘和母式 $\int \frac{\mathrm{d}x}{y}$。
\end{note}

\begin{theorem}[$\int\sec^n x\mathrm{d}x$型]
\begin{align*}
\int\sec x\,\mathrm{d}x &= 1\ln|\sec x+\tan x|,\\
\int\sec^{3}x\,\mathrm{d}x &= \frac{1}{2}\tan x\sec x+\frac{1}{2}1\ln|\sec x+\tan x|,\\
\int\sec^{5}x\,\mathrm{d}x &= \frac{1}{4}\tan x\sec^{3}x+\frac{1}{4}\frac{3}{2}\tan x\sec x+\frac{1}{2}\frac{3}{2}1\ln|\sec x+\tan x|,\\
\int\sec^{7}x\,\mathrm{d}x &= \frac{1}{6}\tan x\sec^{5}x+\frac{1}{6}\frac{5}{4}\tan x\sec^{3}x+\frac{1}{6}\frac{5}{4}\frac{3}{2}\tan x\sec x+\frac{1}{6}\frac{5}{4}\frac{3}{2}1\ln|\sec x+\tan x|.
\end{align*}
\end{theorem}
置 $p=\sec x,\ q=\tan x,\ p^{2}=q^{2}+1$，且 $\mathrm{d}q=p^{2}\mathrm{d}x$，
则 \(\displaystyle\int\sec^{2n+1}x\,\mathrm{d}x=\int p^{2n-1}\,\mathrm{d}q\)。

\begin{theorem}[$\int\csc^{2n+1}x\mathrm{d}x$型]
\begin{align*}
\int\csc x\,\mathrm{d}x &= -1\ln|\csc x+\cot x|,\\
\int\csc^3x\,\mathrm{d}x &= -\frac{1}{2}\cot x\csc x-\frac{1}{2}1\ln|\csc x+\cot x|,\\
\int\csc^5x\,\mathrm{d}x &= -\frac{1}{4}\cot x\csc^3x-\frac{1}{4}\frac{3}{2}\cot x\csc x-\frac{1}{4}\frac{3}{2}1\ln|\csc x+\cot x|,\\
\int\csc^7x\,\mathrm{d}x &= -\frac{1}{6}\cot x\csc^5x-\frac{1}{6}\frac{5}{4}\cot x\csc^3x-\frac{1}{6}\frac{5}{4}\frac{3}{2}\cot x\csc x-\frac{1}{6}\frac{5}{4}\frac{3}{2}1\ln|\csc x+\cot x|.
\end{align*}
\end{theorem}
留意到恒等式
\[
\csc^{2}x=\frac{1}{\sin^{2}x}=\frac{\sin^{2}x+\cos^{2}x}{\sin^{2}x}=1+\cot^{2}x.
\]
置 $p=\csc x,\ q=\cot x,\ p^{2}=q^{2}+1$，并且
\[
\mathrm{d}q=\mathrm{d}\left(\frac{\cos x}{\sin x}\right)
=\frac{-\sin^{2}x-\cos^{2}x}{\sin^{2}x}\,\mathrm{d}x
=-p^{2}\,\mathrm{d}x,\qquad
\int\csc^{2n+1}x\,\mathrm{d}x=-\int p^{2n-1}\,\mathrm{d}q.
\]

\begin{theorem}[$\cos^{2n}x$型]

\begin{align*}
\int \mathrm{d}x &= 1x,\\
\int\cos^2x\,\mathrm{d}x &= \frac{1}{2}\sin x\cos x+\frac{1}{2}1x,\\
\int\cos^4x\,\mathrm{d}x &= \frac{1}{4}\sin x\cos^{3}x+\frac{1}{4}\frac{3}{2}\sin x\cos x+\frac{1}{4}\frac{3}{2}1x,\\
\int\cos^6x\,\mathrm{d}x &= \frac{1}{6}\sin x\cos^5x+\frac{1}{6}\frac{5}{4}\sin x\cos^3x+\frac{1}{6}\frac{5}{4}\frac{3}{2}\sin x\cos x+\frac{1}{6}\frac{5}{4}\frac{3}{2}1x,\\
\int\cos^8x\,\mathrm{d}x &= \frac{1}{8}\sin x\cos^7x+\frac{1}{8}\frac{7}{6}\sin x\cos^5x+\frac{1}{8}\frac{7}{6}\frac{5}{4}\sin x\cos^3x+\frac{1}{8}\frac{7}{6}\frac{5}{4}\frac{3}{2}\sin x\cos x\\
&\qquad+\frac{1}{8}\frac{7}{6}\frac{5}{4}\frac{3}{2}1x.
\end{align*}
\end{theorem}
置 $p=\cos x,\ q=\sin x$，则 $\mathrm{d}q=p\,\mathrm{d}x$，
\(\displaystyle \int\cos^{2n}x\,\mathrm{d}x=\int p^{2n-1}\,\mathrm{d}q\)。

\begin{theorem}[$\sin^{2n}x$型]
    
    \begin{align*}
        \int\mathrm{d}x=&1x\\
        \int\sin^2x\mathrm{d}x=&-\frac{1}{2}\sin x\cos x+\frac{1}{2}1x\\
        \int\sin^4x\operatorname{d}x=&-\frac{1}{4}\sin^3x\cos x-\frac{1}{4}\frac{3}{2}\sin x\cos x+\frac{1}{4}\frac{3}{2}1x\\
        \int\sin^6x\operatorname{d}x=&-\frac{1}{6}\sin^5x\cos x-\frac{1}{6}\frac{5}{4}\sin^3x\cos x-\frac{1}{6}\frac{5}{4}\frac{3}{2}\sin x\cos x+\frac{1}{6}\frac{5}{4}\frac{3}{2}1x\\
        \int\sin^8x\operatorname{d}x=&-\frac{1}{8}\sin^7x\cos x-\frac{1}{8}\frac{7}{6}\sin^5x\cos x-\frac{1}{8}\frac{7}{6}\frac{5}{4}\sin^3x\cos x-\frac{1}{8}\frac{7}{6}\frac{5}{4}\frac{3}{2}\sin x\cos x\\
        &+\frac{1}{8}\frac{7}{6}\frac{5}{4}\frac{3}{2}1x
    \end{align*}
\end{theorem}
置$p=\sin x,q=\cos x$，则$\mathrm{d}q=p\mathrm{d}x$，$\int\sin^{2n}x\mathrm{d}x=\int p^{2n-1}\mathrm{d}q$.

\begin{theorem}[$\frac{1}{(1+x^2)^n}$型]
    
    \begin{align*}
        \int\frac{1}{1+x^2}\mathrm{d}x=&\arctan x\\
        \int\frac{1}{\left(1+x^2\right)^2}\mathrm{d}x=&\frac{1}{2}\frac{x}{1+x^2}+\frac{1}{2}1\arctan x\\
        \int\frac{1}{\left(1+x^2\right)^3}\mathrm{d}x=&\frac{1}{4}\frac{x}{\left(1+x^2\right)^2}+\frac{1}{4}\frac{3}{2}\frac{x}{1+x^2}+\frac{1}{4}\frac{3}{2}1\arctan x\\
        \int\frac{1}{\left(1+x^{2}\right)^{4}}\mathrm{d}x=&\frac{1}{6}\frac{x}{\left(1+x^{2}\right)^{3}}+\frac{1}{6}\frac{5}{4}\frac{x}{\left(1+x^{2}\right)^{2}}+\frac{1}{6}\frac{5}{4}\frac{3}{2}\frac{x}{1+x^{2}}+\frac{1}{6}\frac{5}{4}\frac{3}{2}1\arctan x\\
        \int\frac{1}{\left(1+x^{2}\right)^{5}}\mathrm{d}x=&\frac{1}{8}\frac{x}{\left(1+x^{2}\right)^{4}}+\frac{1}{8}\frac{7}{6}\frac{x}{\left(1+x^{2}\right)^{3}}+\frac{1}{8}\frac{7}{6}\frac{5}{4}\frac{x}{\left(1+x^{2}\right)^{2}}+\frac{1}{8}\frac{7}{6}\frac{5}{4}\frac{3}{2}\frac{x}{1+x^{2}}\\
        &+\frac{1}{8}\frac{7}{6}\frac{5}{4}\frac{3}{2}1\arctan x
    \end{align*}
\end{theorem}
记$y=\sqrt{1+x^{2}}$，则$y^2=x^2+1,\mathrm{d}y=\frac{x}{\sqrt{1+x^2}}\mathrm{d}x=\frac{x}{y}\mathrm{d}x$
\begin{align*}
\int\frac{\mathrm{d}x}{(1+x^{2})^{2}}&=\int\frac{\mathrm{d}x}{y^{4}}=\int\frac{y^{2}-x^{2}}{y^{4}}\mathrm{d}x=\int\frac{\mathrm{d}x}{y^{2}}-\int\frac{x\mathrm{d}y}{y^{3}}\\
&=\int\frac{\mathrm{d}x}{1+x^{2}}-\int x\mathrm{d}\left(y^{-2}\cdot\frac{1}{-2}\right) \quad\text{（下一步用分部积分：}\textstyle\int u\,\mathrm{d}v=uv-\int v\,\mathrm{d}u\text{）}\\
&=\arctan x+\frac{1}{2}\cdot\left[\frac{x}{y^{2}}-\int\frac{\mathrm{d}x}{y^{2}}\right]
\end{align*}

\begin{theorem}[$\frac{1}{(1-x^2)^n}$型]

\begin{align*}
\int\frac{1}{1-x^2}\mathrm{d}x=&\operatorname{artanh}x\quad(\text{注:}\operatorname{artanh}x=\frac{1}{2}\ln\left|\frac{1+x}{1-x}\right|)\\
\int\frac{1}{\left(1-x^2\right)^2}\mathrm{d}x=&\frac{1}{2}\frac{x}{1-x^2}+\frac{1}{2}1\operatorname{artanh}x\\
\int\frac{1}{\left(1-x^2\right)^3}\mathrm{d}x=&\frac{1}{4}\frac{x}{\left(1-x^2\right)^2}+\frac{1}{4}\frac{3}{2}\frac{x}{1-x^2}+\frac{1}{4}\frac{3}{2}1\operatorname{artanh}x\\
\int\frac{1}{(1-x^{2})^{4}}\mathrm{d}x=&\frac{1}{6}\frac{x}{(1-x^2)^3}+\frac{1}{6}\frac{5}{4}\frac{x}{(1-x^2)^2}+\frac{1}{6}\frac{5}{4}\frac{3}{2}\frac{x}{1-x^2}+\frac{1}{6}\frac{5}{4}\frac{3}{2}1\operatorname{artanh}x\\
\int\frac{1}{\left(1-x^{2}\right)^{5}}\mathrm{d}x=&\frac{1}{8}\frac{x}{\left(1-x^{2}\right)^{4}}+\frac{1}{8}\frac{7}{6}\frac{x}{\left(1-x^{2}\right)^{3}}+\frac{1}{8}\frac{7}{6}\frac{5}{4}\frac{x}{\left(1-x^{2}\right)^{2}}+\frac{1}{8}\frac{7}{6}\frac{5}{4}\frac{3}{2}\frac{x}{1-x^{2}}\\
&+\frac{1}{8}\frac{7}{6}\frac{5}{4}\frac{3}{2}1\mathrm{artanh}x\end{align*}
\end{theorem}
记$y=\sqrt{1-x^{2}}$，则$x^2+y^2=1$，则
\begin{align*}
    \int\frac{1}{(1-x^{2})^{2}}\mathrm{d}x&=\int\frac{\mathrm{d}x}{y^{4}}=\int\frac{x^{2}+y^{2}}{y^{4}}\mathrm{d}x=\int\frac{x^{2}}{y^{4}}\mathrm{d}x+\int\frac{\mathrm{d}x}{y^{2}}\\
&=-\int\frac{x\mathrm{d}y}{y^3}+\int\frac{\mathrm{d}x}{1-x^{2}}=\text{arctanh}~x+\int\frac{\mathrm{d}x}{1-x^{2}}\\
\end{align*}

\begin{theorem}[$\frac{x^{2n}}{\sqrt{1-x^2}}\mathrm{d}x$型]
    
    \begin{align*}
        \int\frac{1}{\sqrt{1-x^{2}}}\mathrm{d}x=&-\arccos x\\
        \int\frac{x^2}{\sqrt{1-x^2}}\mathrm{d}x=&-\frac{1}{2}x\sqrt{1-x^2}-\frac{1}{2}1\arccos x\\
        \int\frac{x^4}{\sqrt{1-x^2}}\mathrm{d}x=&-\frac{1}{4}x^3\sqrt{1-x^2}-\frac{1}{4}\frac{3}{2}x\sqrt{1-x^2}-\frac{1}{4}\frac{3}{2}1\arccos x\\
        \int\frac{x^{6}}{\sqrt{1-x^{2}}}\mathrm{d}x=&-\frac{1}{6}x^5\sqrt{1-x^2}-\frac{1}{6}\frac{5}{4}x^3\sqrt{1-x^2}-\frac{1}{6}\frac{5}{4}\frac{3}{2}x\sqrt{1-x^2}-\frac{1}{6}\frac{5}{4}\frac{3}{2}1\arccos x\\
        \int\frac{x^{8}}{\sqrt{1-x^{2}}}\mathrm{d}x=&-\frac{1}{8}x^{7}\sqrt{1-x^{2}}-\frac{1}{8}\frac{7}{6}x^{5}\sqrt{1-x^{2}}-\frac{1}{8}\frac{7}{6}\frac{5}{4}x^{3}\sqrt{1-x^{2}}-\frac{1}{8}\frac{7}{6}\frac{5}{4}\frac{3}{2}x\sqrt{1-x^{2}}\\
        &-\frac{1}{8}\frac{7}{6}\frac{5}{4}\frac{3}{2}1\arccos x
    \end{align*}
\end{theorem}
记$y=\sqrt{1-x^{2}}$，$y^{2}+x^{2}=1\Rightarrow x\mathrm{d}x=-y\mathrm{d}y$
\[\int\frac{x^{2n}}{\sqrt{1-x^{2}}}\mathrm{d}x=\int\frac{x^{2n}\mathrm{d}x}{y}=-\int x^{2n-1}\mathrm{d}y\]
特别地,当$n=1$时
\begin{align*}
\int\frac{x^{2}}{\sqrt{1-x^{2}}}\mathrm{d}x=&-\int x\mathrm{d}y=-\frac{1}{2}\times y-\frac{1}{2}\cdot\arctan\frac{y}{x}\\=&-\frac{1}{2}\times\sqrt{1-x^{2}}-\frac{1}{2}\arctan\frac{\sqrt{1-x^{2}}}{x}=-\frac{1}{2}\times\sqrt{1-x^{2}}-\frac{1}{2}\arccos x
\end{align*}

\begin{theorem}[$\int\frac{1}{x^{2n-1}\sqrt{1-x^2}}\mathrm{d}x$型]
    
    \begin{align*}
    \int\frac{1}{x\sqrt{1-x^2}}\mathrm{d}x=&-\ln\frac{1+\sqrt{1-x^2}}{x}\\
    \int\frac{1}{x^{3}\sqrt{1-x^{2}}}\mathrm{d}x=&-\frac{1}{2}\frac{\sqrt{1-x^{2}}}{x^{2}}-\frac{1}{2}1\ln\frac{1+\sqrt{1-x^{2}}}{x}\\
    \int\frac{1}{x^5\sqrt{1-x^2}}\mathrm{d}x=&-\frac{1}{4}\frac{\sqrt{1-x^2}}{x^4}-\frac{1}{4}\frac{3}{2}\frac{\sqrt{1-x^2}}{x^2}-\frac{1}{4}\frac{3}{2}1\ln\frac{1+\sqrt{1-x^2}}{x}\\
    \int\frac{1}{x^7\sqrt{1-x^2}}\mathrm{d}x=&-\frac{1}{6}\frac{\sqrt{1-x^2}}{x^6}-\frac{1}{6}\frac{5}{4}\frac{\sqrt{1-x^2}}{x^4}-\frac{1}{6}\frac{5}{4}\frac{3}{2}\frac{\sqrt{1-x^2}}{x^2}-\frac{1}{6}\frac{5}{4}\frac{3}{2}1\ln\frac{1+\sqrt{1-x^2}}{x}\\
    \int\frac{1}{x^9\sqrt{1-x^2}}\mathrm{d}x=&-\frac{1}{8}\frac{\sqrt{1-x^{2}}}{x^{8}}-\frac{1}{8}\frac{7}{6}\frac{\sqrt{1-x^{2}}}{x^{6}}-\frac{1}{8}\frac{7}{6}\frac{5}{4}\frac{\sqrt{1-x^{2}}}{x^{4}}-\frac{1}{8}\frac{7}{6}\frac{5}{4}\frac{3}{2}\frac{\sqrt{1-x^{2}}}{x^{2}}\\
    &-\frac{1}{8}\frac{7}{6}\frac{5}{4}\frac{3}{2}1\ln\frac{1+\sqrt{1-x^{2}}}{x}
    \end{align*}
\end{theorem}
记$y=\sqrt{1-x^{2}}\Rightarrow y^{2}+x^{2}=1\Rightarrow x\mathrm{d}x=-y\mathrm{d}y$
\[\int\frac{\mathrm{d}x}{x^{2n+1}\cdot\sqrt{1-x^{2}}}=\int\frac{\mathrm{d}x}{x^{2n+1}\cdot y}=-\int\frac{\mathrm{d}y}{x^{2n+2}}\]

\begin{theorem}[$\int\frac{x^{2n}}{\sqrt{x^{2}-1}}\mathrm{d}x$型]
    \begin{align*}
    \int\frac{1}{\sqrt{x^{2}-1}}\mathrm{d}x=&\operatorname{arccos}x=\ln\left|x+\sqrt{x^2-1}\right|\\
    \int\frac{x^{2}}{\sqrt{x^{2}-1}}\mathrm{d}x=&\frac{1}{2}x\sqrt{x^2-1}+\frac{1}{2}1\operatorname{arccos}x\\
    \int\frac{x^4}{\sqrt{x^2-1}}\mathrm{d}x=&\frac{1}{4}x^3\sqrt{x^2-1}+\frac{1}{4}\frac{3}{2}x\sqrt{x^2-1}+\frac{1}{4}\frac{3}{2}1\operatorname{arccos}x\\
    \int\frac{x^{6}}{\sqrt{x^{2}-1}}\mathrm{d}x=&\frac{1}{6}x^5\sqrt{x^2-1}+\frac{1}{6}\frac{5}{4}x^3\sqrt{x^2-1}+\frac{1}{6}\frac{5}{4}\frac{3}{2}x\sqrt{x^2-1}+\frac{1}{6}\frac{5}{4}\frac{3}{2}1\mathrm{arccos}x\\
    \int\frac{x^8}{\sqrt{x^2-1}}\mathrm{d}x=&\frac{1}{8}x^7\sqrt{x^2-1}+\frac{1}{8}\frac{7}{6}x^5\sqrt{x^2-1}+\frac{1}{8}\frac{7}{6}\frac{5}{4}x^3\sqrt{x^2-1}+\frac{1}{8}\frac{7}{6}\frac{5}{4}\frac{3}{2}x\sqrt{x^2-1}\\
    &+\frac{1}{8}\frac{7}{6}\frac{5}{4}\frac{3}{2}1\mathrm{arccos}~x
    \end{align*}
\end{theorem}
置$y=\sqrt{x^{2}-1}\Rightarrow x^{2}-y^{2}=1$，则
\[\int\frac{x^{2n}}{\sqrt{x^{2}-1}}\mathrm{d}x=\int\frac{x^{2n}}{y}\mathrm{d}x=\int x^{2n-1}\mathrm{d}y\]
特别地，当$n=1$时，有
\[\int\frac{x^{2}}{\sqrt{x^{2}-1}}\mathrm{d}x=\int x\mathrm{d}y=\frac{1}{2}xy+\frac{1}{2}\cdot\ln(x+y)\]

\begin{theorem}[$\int\frac{\mathrm{d}x}{x^{2n+1}\sqrt{x^{2}-1}}$型]
\begin{align*}
    \int\frac{1}{x\sqrt{x^{2}-1}}\mathrm{d}x=&\arccos\frac{1}{|x|}\\
    \int\frac{1}{x^{3}\sqrt{x^{2}-1}}\mathrm{d}x=&\frac{1}{2}\frac{\sqrt{x^{2}-1}}{x^{2}}+\frac{1}{2}1\arccos\frac{1}{|x|}\\
    \int\frac{1}{x^5\sqrt{x^2-1}}\mathrm{d}x=&\frac{1}{4}\frac{\sqrt{x^2-1}}{x^4}+\frac{1}{4}\frac{3}{2}\frac{\sqrt{x^2-1}}{x^2}+\frac{1}{4}\frac{3}{2}1\arccos\frac{1}{|x|}\\
    \int\frac{1}{x^7\sqrt{x^2-1}}\mathrm{d}x=&\frac{1}{6}\frac{\sqrt{x^2-1}}{x^6}+\frac{1}{6}\frac{5}{4}\frac{\sqrt{x^2-1}}{x^4}+\frac{1}{6}\frac{5}{4}\frac{3}{2}\frac{\sqrt{x^2-1}}{x^2}+\frac{1}{6}\frac{5}{4}\frac{3}{2}1\arccos\frac{1}{|x|}\\
    \int\frac{1}{x^{9}\sqrt{x^{2}-1}}\mathrm{d}x=&\frac{1}{8}\frac{\sqrt{x^{2}-1}}{x^{8}}+\frac{1}{8}\frac{7}{6}\frac{\sqrt{x^{2}-1}}{x^{6}}+\frac{1}{8}\frac{7}{6}\frac{5}{4}\frac{\sqrt{x^{2}-1}}{x^{4}}+\frac{1}{8}\frac{7}{6}\frac{5}{4}\frac{3}{2}\frac{\sqrt{x^{2}-1}}{x^{2}}\\
    &+\frac{1}{8}\frac{7}{6}\frac{5}{4}\frac{3}{2}1\arccos\frac{1}{|x|}
\end{align*}
\end{theorem}
置$y=\sqrt{x^{2}-1}$，则$x^2-y^2=1$，则
\[\int\frac{\mathrm{d}x}{x^{2n+1}\sqrt{x^{2}-1}}=\int\frac{\mathrm{d}x}{x^{2n+1}\cdot y}=\int\frac{\mathrm{d}y}{x^{2n+2}}\]
当$n=0$时\[\int\frac{\mathrm{d}x}{x\sqrt{x^2-1}}=\int\frac{\mathrm{d}y}{x^2}=\int\frac{\mathrm{d}y}{x^2}=\int\frac{\mathrm{d}y}{1+y^2}=\arctan y\]

\begin{theorem}[$\int\frac{x^{2n}}{\sqrt{x^{2}+1}}\mathrm{d}x$型]
    \begin{align*}
        \int\frac{1}{\sqrt{x^2+1}}\mathrm{d}x=&\operatorname{arcsinh}x\\
        \int\frac{x^2}{\sqrt{x^2+1}}\mathrm{d}x=&\frac{1}{2}x\sqrt{x^2+1}-\frac{1}{2}1\operatorname{arcsinh}x\\
        \int\frac{x^4}{\sqrt{x^2+1}}\mathrm{d}x=&\frac{1}{4}x^{3}\sqrt{x^{2}+1}-\frac{1}{4}\frac{3}{2}x\sqrt{x^{2}+1}+\frac{1}{4}\frac{3}{2}1\operatorname{arcsinh}x\\
        \int\frac{x^6}{\sqrt{x^2+1}}\mathrm{d}x=&\frac{1}{6}x^5\sqrt{x^2+1}-\frac{1}{6}\frac{5}{4}x^3\sqrt{x^2+1}+\frac{1}{6}\frac{5}{4}\frac{3}{2}x\sqrt{x^2+1}-\frac{1}{6}\frac{5}{4}\frac{3}{2}1\operatorname{arcsinh}x\\
        \int\frac{x^8}{\sqrt{x^2+1}}\mathrm{d}x=&\frac{1}{8}x^{7}\sqrt{x^{2}+1}-\frac{1}{8}\frac{7}{6}x^{5}\sqrt{x^{2}+1}+\frac{1}{8}\frac{7}{6}\frac{5}{4}x^{3}\sqrt{x^{2}+1}-\frac{1}{8}\frac{7}{6}\frac{5}{4}\frac{3}{2}x\sqrt{x^{2}+1}\\
        &+\frac{1}{8}\frac{7}{6}\frac{5}{4}\frac{3}{2}1\mathrm{arcsinh}x
    \end{align*}
\end{theorem}
置$y=\sqrt{x^{2}+1}$，则$y^2-x^2=1$，于是：
\[\int\frac{x^{2n}}{\sqrt{x^{2}+1}}\mathrm{d}x=\int\frac{x^{2n-1}\cdot x\mathrm{d}x}{y}=\int x^{2n-1}\mathrm{d}y\]
当n=1时
\begin{align*}
\int\frac{x^{2}}{\sqrt{x^{2}+1}}\mathrm{d}x&=\int x\mathrm{d}y=\frac{1}{2}xy-\frac{1}{2}\ln(x+y) =\frac{1}{2}\times\sqrt{x^{2}+1}-\frac{1}{2}\ln(x+\sqrt{x^{2}+1})
\end{align*}

\begin{theorem}[$\int\frac{\mathrm{d}x}{x^{2n+1}\cdot\sqrt{x^{2}+1}}$型]
    \begin{align*}
        \int\frac{1}{x\sqrt{x^2+1}}\mathrm{d}x&=-\ln\frac{1+\sqrt{x^{2}+1}}{x}\\
        \int\frac{1}{x^3\sqrt{x^2+1}}\mathrm{d}x=&-\frac{1}{2}\frac{\sqrt{x^2+1}}{x^2}+\frac{1}{2}1\ln\frac{1+\sqrt{x^2+1}}{x}\\
        \int\frac{1}{x^{5}\sqrt{x^{2}+1}}\mathrm{d}x=&-\frac{1}{4}\frac{\sqrt{x^2+1}}{x^4}+\frac{1}{4}\frac{3}{2}\frac{\sqrt{x^2+1}}{x^2}-\frac{1}{4}\frac{3}{2}1\ln\frac{1+\sqrt{x^2+1}}{x}\\
        \int\frac{1}{x^7\sqrt{x^2+1}}\mathrm{d}x=&-\frac{1}{6}\frac{\sqrt{x^2+1}}{x^6}+\frac{1}{6}\frac{5}{4}\frac{\sqrt{x^2+1}}{x^4}-\frac{1}{6}\frac{5}{4}\frac{3}{2}\frac{\sqrt{x^2+1}}{x^2}+\frac{1}{6}\frac{5}{4}\frac{3}{2}1\ln\frac{1+\sqrt{x^2+1}}{x}\\
        \int\frac{1}{x^{9}\sqrt{x^{2}+1}}\mathrm{d}x=&-\frac{1}{8}\frac{\sqrt{x^{2}+1}}{x^{8}}+\frac{1}{8}\frac{7}{6}\frac{\sqrt{x^{2}+1}}{x^{6}}-\frac{1}{8}\frac{7}{6}\frac{5}{4}\frac{\sqrt{x^{2}+1}}{x^{4}}+\frac{1}{8}\frac{7}{6}\frac{5}{4}\frac{3}{2}\frac{\sqrt{x^{2}+1}}{x^{2}}\\
        &-\frac{1}{8}\frac{7}{6}\frac{5}{4}\frac{3}{2}1\ln\frac{1+\sqrt{x^{2}+1}}{x}
    \end{align*}
\end{theorem}
置$y=\sqrt{x^{2}+1}$，则$y^{2}-x^{2}=1$，
\[\int\frac{\mathrm{d}x}{x^{2n+1}\cdot\sqrt{x^{2}+1}}=\int\frac{\mathrm{d}x}{x^{2n+1}\cdot y}=\int\frac{\mathrm{d}y}{x^{2n+2}}\]
当$n=0$时
\begin{align*}
\int\frac{\mathrm{d}x}{x\cdot\sqrt{x^{2}+1}}=&\int\frac{\mathrm{d}y}{x^{2}}=\int\frac{y^{2}-x^{2}}{x^{2}}\mathrm{d}y=\int\frac{y^{2}\mathrm{d}y}{x^{2}}-\int \mathrm{d}y\\=&\int\frac{y^{2}}{y^{2}-1}\mathrm{d}y-\int \mathrm{d}y=\int\frac{\mathrm{d}y}{y^{2}-1}=-\text{arctanh} y
\end{align*}

\par\medskip\noindent\textbf{经典组合与递推}\par

题目里同时出现两个根式、两个三角量、两个双曲量，或者某个复杂代数式明显可以拆成“平方和 / 平方差 = 常数”。这时不要先算，先试图造出一组 $p,q$ 或 $p,q,r$。

先找题面里最自然的两个量，把它们记成 $p,q$。若能得到 $p^2+q^2=\text{常数}$，就按实圆处理；若能得到 $p^2-q^2=\text{常数}$，就按虚圆处理。若题目还有第三个根式或第三个频率，就再补一个 $r$，把三者压进同一个底座。只有在你先把底座写出来以后，后面的 $\mathrm{d}p,\mathrm{d}q,\mathrm{d}r$ 才会自己长出路。

\textbf{经典双元公式。}
\noindent 经典双元
\begin{itemize}
\item $p=\sin x,q=\cos x;~p=\sec x,q=\tan x$.
\item $p=\sqrt{x^2-1},q=\sqrt{1-x^2},r=\sqrt{x^2+1}$.
\item $p=\sqrt{1+a^2x^2},q=ax$.
\item $p=\sqrt{2+x^3},q=x^{\frac32}.$和$p=\sqrt{a+x^k},q=x^{\frac{k}2}$
\item $p=\frac{\sqrt2\sqrt{x^2+1}}{x+1},q=\frac{x-1}{x+1}$，则$p^2=1+q^2$.
\item $p=\frac{x}{\sqrt{1-x^2}},q=\frac{\sqrt{2x^2-1}}{\sqrt{1-x^2}}\Rightarrow p^2-1=q^2$.
\item $p=x+\frac{a}{x},q=x-\frac{a}{x}$.
\item $p=\sqrt{x-a}+\sqrt{b-x},q=\sqrt{x-a}-\sqrt{b-x}$.
\item $p=\cos x+\sin x,q=\cos x-\sin x\Rightarrow p^2+q^2=2,\mathrm{d}p=q\mathrm{d}x,\mathrm{d}q=-p\mathrm{d}x$.
\item $p=\sin x+2\cos x,q=\cos x-2\sin x$.
\item $p=\sqrt{\sec^4x+\tan^4x},q=\sqrt2\sec x\tan x\Rightarrow p^2-q^2=1$.
\item $p=\sqrt{\sec x+1},q=\sqrt{\sec x-1}\Rightarrow p^2-q^2=2,pq=\tan x$.
\item $p=\sqrt{1+\sin x}=\sin \frac{x}2+\cos\frac{x}2,q=\cos\frac{x}2-\sin\frac{x}2$.
\end{itemize}
经典三元
\begin{itemize}
\item $p=\sin x,q=\cos x,r=\sqrt{1+\sin^2 x}$.
\item $p=\sqrt{x}+\sqrt{\frac{2}{x}},q=\sqrt{x}-\sqrt{\frac{2}{x}},r=\sqrt{x+\frac{2}{x}+2}$.
\item $p=\frac{\sqrt2\sqrt{1+x+x^2}}{\sqrt{x^2+1}},q=\frac{x+1}{\sqrt{x^2+1}},r=\frac{1-x}{\sqrt{x^2+1}}\implies q^2=p^2-1=2-r^2$.
\item $p=x+\frac{a}{x},q=x-\frac{a}{x},r=\sqrt{x^2+\frac{a^2}{x^2}+b}$.
\item $r=\sqrt{x+\frac1x+1},s=\sqrt{x+\frac1x-1},p=\sqrt{x}+\frac1{\sqrt{x}},q=\sqrt{x}-\frac1{\sqrt{x}}$.
\item $p=\cos x+\sin x,q=\cos x-\sin x,r=\sqrt{2\sin x\cos x}\implies r^2=p^2-1=1-q^2,\mathrm{d}p=q\mathrm{d}x,\mathrm{d}q=-p\mathrm{d}x$.
\item $p=\sin\frac{x}{2}+\cos\frac{x}{2},q=\cos\frac{x}{2}-\sin\frac{x}{2}=\sqrt{1-\sin x},r=\sqrt{\sin x}$.
\item $p=\sin x+\cos x,q=\cos x-\sin x,r=\sqrt{2+\sin 2x}\implies r^2-p^2=1,p^2+q^2=2$.
\item $p=x+\sqrt{1-x^2},q=x-\sqrt{1-x^2},r=\sqrt{2+2x\sqrt{1-x^2}}\implies r^2-p^2=1,p^2+q^2=2$.
\item $p,q=\frac{\sqrt{\sqrt{1+x^4}+1}\pm\sqrt{\sqrt{1+x^4}-1}}{\sqrt2x},r=\frac{\sqrt[4]{1+x^4}}{x},p^2-q^2=2,p^2-r^2=1$.
\end{itemize}
经典恒等式
\begin{itemize}
\item$\sin^2x+\cos^2x=1\Leftrightarrow (\sqrt{1-x^2})^2+x^2=1$
\item $(x+1)^2+(x-1)^2=2(x^2+1)\Leftrightarrow \left(\frac{\sqrt2\sqrt{x^2+1}}{x+1}\right)^2=1+\left(\frac{x-1}{x+1}\right)^2$.
\item $x^2-(1-x^2)=2x^2-1\Leftrightarrow \left(\frac{x}{\sqrt{1-x^2}}\right)^2-\left(\frac{\sqrt{2x^2-1}}{\sqrt{1-x^2}}\right)^2=1$.
\item 偶数次多项式配方也能出恒等式。比如$1+x^2+3x^4=\left(\frac{x^2}{2}+1\right)^2+\left(\frac{\sqrt{11}}{2}x^2\right)^2$.
\item $2(1+x+x^2)=(x+1)^2+(x^2+1)$.
\item $5(x^2+4)=4(x-1)^2+(x+4)^2$.
\item $\left(x-\frac{a}{x}\right)^2+4a=\left(x+\frac{a}{x}\right)^2$
\end{itemize}

落笔时可写成“设 $p=\cdots,\ q=\cdots$，则 $p^2\pm q^2=\text{常数}$，且 $\mathrm{d}p=\cdots,\ \mathrm{d}q=\cdots$，故原积分化为 $\cdots$。” 看到题目和上述关系长得像时，优先使用这一表达，不要临场重新发明。

若只差一步凑微分就能结束，不要强行造三元。若造出的 $p,q$ 之间没有清晰的平方恒等式，说明这组搭配不值钱，应换另一组。

检查你造出的 $p,q,r$ 是否真的能把原式全部改写；检查微分关系是否和底座符号一致，特别是实圆的负号与虚圆的正号。
\par\medskip


\subsection{\texorpdfstring{含有 $ax+b$ 与 $\sqrt{ax+b}$ 的积分}{一次式与一次根式的积分}}

\par\medskip\noindent\textbf{一次式积分}\par

分母或根式里反复出现 $ax+b$，分子却是 $1,x,x^2$ 或者 $x$ 的低次多项式；或者题目长成 $\dfrac{1}{x(ax+b)}$、$\dfrac{x}{(ax+b)^2}$ 这类“一个一次式配一个 $x$”的结构。令 \(\displaystyle u=ax+b,\ \mathrm{d}u=a\,\mathrm{d}x,\ x=\frac{u-b}{a}\)。
\[
\int \frac{\mathrm{d}x}{ax+b}=\frac{1}{a}\ln|ax+b|+C,\qquad
\int (ax+b)^\alpha \,\mathrm{d}x=\frac{(ax+b)^{\alpha+1}}{a(\alpha+1)}+C.
\]
并记 \(\displaystyle x=\frac{(ax+b)-b}{a},\ x^k=\frac{\bigl((ax+b)-b\bigr)^k}{a^k},\ \frac{1}{x(ax+b)}=\frac{1}{b}\left(\frac{1}{x}-\frac{a}{ax+b}\right)\)。

如果整题只围绕一个 $ax+b$，先令 $u=ax+b$。若分子里有 $x$ 或 $x^2$，一律先把 $x$ 改写成 $\dfrac{u-b}{a}$，把原式压成 $u$ 的有理函数。若分母还带一个 $x$，优先做拆项，把 $\dfrac{1}{x(ax+b)}$、$\dfrac{1}{x^2(ax+b)}$ 这类结构拆回母式。若出现 $(ax+b)^2$，仍然先改写分子，再把整题看成关于 $u$ 的幂函数或有理分式。
常见开头是“令 $u=ax+b$，则 $x=\dfrac{u-b}{a}$，$\mathrm{d}x=\dfrac{1}{a}\mathrm{d}u$”。若分子含 $x^k$，直接写成 $\dfrac{(u-b)^k}{a^k}$；若分母还有 $x$，先拆成 $\dfrac{1}{x}$ 母式与 $\dfrac{1}{u}$ 母式的线性组合。
“原式围绕 $ax+b$ 展开，令 $u=ax+b$。再用 $x=\dfrac{u-b}{a}$ 消去 $x$，原积分化为 $u$ 的初等积分，故结果为 $\cdots$。”
若题目分母已经是两个不同一次因式的乘积，直接部分分式可能更短。若 $a=0$，这节写法失效，要先把题目还原成常数情形。
\par\medskip


\par\medskip\noindent\textbf{一次根式积分}\par

题目中反复出现 $\sqrt{ax+b}$、$\dfrac{1}{\sqrt{ax+b}}$、$\dfrac{\sqrt{ax+b}}{x}$、$\dfrac{1}{x\sqrt{ax+b}}$ 这类一次根式。
\begin{align*}
y&=\sqrt{ax+b},\qquad y^2=ax+b,\qquad x=\frac{y^2-b}{a},\qquad \mathrm{d}x=\frac{2y}{a}\,\mathrm{d}y,\\
\int \frac{\mathrm{d}x}{x\sqrt{ax+b}}
&=\int \frac{2}{y^2-b}\,\mathrm{d}y,\qquad
\int \frac{\sqrt{ax+b}}{x}\,\mathrm{d}x=\int \frac{2y^2}{y^2-b}\,\mathrm{d}y.
\end{align*}

做题时先设 $y=\sqrt{ax+b}$，于是 $y^2=ax+b$、$x=\dfrac{y^2-b}{a}$、$\mathrm{d}x=\dfrac{2y}{a}\mathrm{d}y$；原式很快就会变成关于 $y$ 的有理式。若只剩 $y$ 的幂，就直接幂积分；若出现 $\dfrac{1}{y^2-b}$，则按 $b>0$ 用对数、$b<0$ 用反正切；若出现 $\dfrac{y^2}{y^2-b}$ 或 $\dfrac{1}{(y^2-b)^2}$，先拆成常数项加母式，不要硬算。

卷面上通常先写“设 $y=\sqrt{ax+b}$，则 $y^2=ax+b,\ x=\dfrac{y^2-b}{a},\ \mathrm{d}x=\dfrac{2y}{a}\mathrm{d}y$”，然后把原式化成 $y$ 的有理式，再按幂函数、$\dfrac{1}{y^2-b}$ 或其递推变体收尾。若题目绕着的不止一个一次根式，或者根号外层还套着一般二次式，就应转去二次根式或纯根式双元；若 $a=0$，先把题目退化成常数根式。最后核对是否把 $y$ 全换回 $\sqrt{ax+b}$，以及 $b>0$ 与 $b<0$ 的收尾公式有没有混用、根号前的系数 $2/a$ 是否漏掉。
\par\medskip


\subsection{二次式及其根式}

\par\medskip\noindent\textbf{$x^2\pm a^2$ 基本母式}\par

分母直接长成 $x^2+a^2$、$x^2-a^2$ 或 $(x^2+a^2)^n$，没有线性项，也还没进入根式。

\[
\int \frac{\mathrm{d}x}{x^2+a^2}=\frac{1}{a}\arctan\frac{x}{a}+C,
\qquad
\int \frac{\mathrm{d}x}{x^2-a^2}=\frac{1}{2a}\ln\left|\frac{x-a}{x+a}\right|+C.
\]
对 \(\displaystyle I_n=\int \frac{\mathrm{d}x}{(x^2+a^2)^n}\)，必背递推：
\[
I_n=
\frac{x}{2(n-1)a^2(x^2+a^2)^{n-1}}
+\frac{2n-3}{2(n-1)a^2}I_{n-1}.
\]

若是 $x^2+a^2$，先提出 $a^2$，再换成标准母式 $1+u^2$。若是 $x^2-a^2$，先因式分解成 $(x-a)(x+a)$，不要硬背奇怪变形。若是 $(x^2+a^2)^n$，直接调用递推，把 $n$ 一路降到 $1$。这一节的核心不是乱换元，而是先认出“反正切母式、对数母式、递推母式”三条主路。
对 $x^2+a^2$ 型，写“令 $u=x/a$”；对 $x^2-a^2$ 型，写“先分解并部分分式”；对高次幂型，先写 $I_n$，再把递推公式直接代上去。
“原式属于 $x^2+a^2$ 母式，令 $u=x/a$ 后得到反正切。” 或 “原式属于 $x^2-a^2$ 母式，分解为 $(x-a)(x+a)$ 后作部分分式。” 或 “原式为 $(x^2+a^2)^{-n}$，直接按递推降次。”
若题目其实是 $ax^2+b$ 或 $ax^2+bx+c$，不要套本节公式，要先提出系数或配方。若分母伴随根号，这一节也不是终点，要跳去根式章的底座写法。
检查 $a$ 有没有提出干净；检查 $x^2-a^2$ 的对数答案分子分母顺序是否写对；检查递推里 $2n-3$ 和 $2(n-1)a^2$ 有没有抄错。
\par\medskip


\par\medskip\noindent\textbf{$ax^2+b$ 型积分的处理}\par

题面中心是 $ax^2+b$，分子只到 $x^2$ 或者分母里还挂着 $x,x^2,x^3$ 这类幂。

\[
\int \frac{\mathrm{d}x}{ax^2+b}
=
\begin{cases}
\dfrac{1}{\sqrt{ab}}\arctan\!\left(\sqrt{\dfrac{a}{b}}\,x\right)+C, & b>0,\\[0.9em]
\dfrac{1}{2\sqrt{-ab}}
\ln\left|\dfrac{\sqrt{a}x-\sqrt{-b}}{\sqrt{a}x+\sqrt{-b}}\right|+C, & b<0.
\end{cases}
\]
\[
\int \frac{x}{ax^2+b}\,\mathrm{d}x=\frac{1}{2a}\ln|ax^2+b|+C,\qquad
\frac{1}{x(ax^2+b)}=\frac{1}{2b}\left(\frac{2}{x}-\frac{2ax}{ax^2+b}\right).
\]

先判 $b$ 的符号：$b>0$ 走反正切，$b<0$ 走对数。若分子正好带一个 $x$，立刻凑 $\mathrm{d}(ax^2+b)=2ax\,\mathrm{d}x$。若分子次数和分母同阶，就用 $ax^2=(ax^2+b)-b$ 拆项。若分母外面还挂着 $x^k$，优先把常数 $b$ 改写成 $(ax^2+b)-ax^2$，把题目拆回 $\dfrac{1}{x^m}$ 和基础母式。若分母平方出现 $(ax^2+b)^2$，先分部积分或先把 $ax^2$ 拆开，再回到基础母式。
开头可写“按 $b$ 的符号分类”。然后如果有 $x$，就写“凑 $\mathrm{d}(ax^2+b)$”；如果有 $x^{-1}$、$x^{-2}$、$x^{-3}$，就写“用 $b=(ax^2+b)-ax^2$ 拆项”。
“原式围绕 $ax^2+b$ 展开，先按 $b$ 的符号确定反正切或对数母式；其余部分通过 $ax^2=(ax^2+b)-b$ 或 $b=(ax^2+b)-ax^2$ 拆回基础母式。”
若中间已经出现 $bx$，说明应跳去一般二次式章节。若根号已经罩在外面，则应跳去 $\sqrt{x^2\pm a^2}$ 或一般二次根式章节。
检查 $b>0,b<0$ 的答案族是否写反；检查对数前面的 $\sqrt a,\sqrt{-b}$ 是否漏掉；检查拆项后每一项都确实简化了难度。
\par\medskip


\par\medskip\noindent\textbf{一般二次式积分的处理}\par

分母是 $ax^2+bx+c$，分子至多一次，题目核心障碍是中间的 $bx$ 项。

先配方：\(\displaystyle ax^2+bx+c=a\left(x+\frac{b}{2a}\right)^2+\frac{4ac-b^2}{4a}\)。
记判别量 $\Delta=b^2-4ac$。则
\[
\int \frac{\mathrm{d}x}{ax^2+bx+c}
=
\begin{cases}
\dfrac{1}{\sqrt{\Delta}}
\ln\left|\dfrac{2ax+b-\sqrt{\Delta}}{2ax+b+\sqrt{\Delta}}\right|+C, & \Delta>0,\\[1.1em]
\dfrac{2}{\sqrt{-\Delta}}
\arctan\dfrac{2ax+b}{\sqrt{-\Delta}}+C, & \Delta<0.
\end{cases}
\]
\[
\int \frac{x}{ax^2+bx+c}\,\mathrm{d}x
=\frac{1}{2a}\ln|ax^2+bx+c|
-\frac{b}{2a}\int \frac{\mathrm{d}x}{ax^2+bx+c}.
\]

处理这类积分时，配方与判别式决定结果类型：$\Delta>0$ 对应对数母式，$\Delta<0$ 对应反正切母式。若分子是 $x$，就把 $2ax+b$ 凑出来，让一部分变成对数，另一部分回到基础母式。若题目再往上叠根号，则配方后的结果应转入二次根式章节，不在这里硬打。
常见处理是把二次式先配方，再按判别式分情况。如果分子含 $x$，就补一个 $b$，写成 \(\displaystyle x=\frac{(2ax+b)-b}{2a}\)，然后一半变对数，一半使用上一条基础公式。
“将 $ax^2+bx+c$ 配方为 $a(x+\frac{b}{2a})^2+\cdots$。由判别式 $\Delta$ 的符号，原式分别化为对数母式或反正切母式。”
若 $b=0$，应退回上一节 $ax^2+b$。若题目是 $\sqrt{ax^2+bx+c}$ 而不是倒数，应跳去二次根式章节。
检查配方常数是否写成 $\dfrac{4ac-b^2}{4a}$；检查 $\Delta$ 的正负是否和最终函数族一致；检查分子拆成 $(2ax+b)-b$ 时前面的 $\dfrac{1}{2a}$ 是否保留。
\par\medskip


\par\medskip\noindent\textbf{$\sqrt{x^2+a^2}$ 型积分的处理}\par

题目出现 $\sqrt{x^2+a^2}$、$\dfrac{1}{\sqrt{x^2+a^2}}$、$\dfrac{1}{(x^2+a^2)^{3/2}}$、$\ln(x+\sqrt{x^2+a^2})$，或者分子分母都只由 $x$ 与 $\sqrt{x^2+a^2}$ 组成。

\begin{align*}
y&=\sqrt{x^2+a^2},\qquad y^2-x^2=a^2,\qquad x\,\mathrm{d}x=y\,\mathrm{d}y,\\
\int \frac{\mathrm{d}x}{y}
&=\ln(x+y)+C
=\operatorname{arsinh}\frac{x}{a}+C,\\
\int \frac{\mathrm{d}x}{y^3}&=\frac{x}{a^2y}+C,\qquad
\int \frac{\mathrm{d}x}{x^2y}=-\frac{y}{a^2x}+C.
\end{align*}

若题目只要求一条标准母式，可直接使用 $\int \mathrm{d}x/\sqrt{x^2+a^2}=\ln(x+\sqrt{x^2+a^2})+C$。若题目里高频出现 $x,y,\mathrm{d}x,\mathrm{d}y$，优先双元法；若题目明显更适合 $\tan t,\sec t$，可改用 $x=a\tan t$。分子里有 $x\,\mathrm{d}x$ 时，立刻换成 $y\,\mathrm{d}y$；分母是三次根式时，立刻查比值微元。
先设 $y=\sqrt{x^2+a^2}$，把所有根式统一写成 $y$。然后用 $y^2-x^2=a^2$ 改写平方项，用 $x\,\mathrm{d}x=y\,\mathrm{d}y$ 转移微分，把题目压回母式、比值微元或对数微分。
“设 $y=\sqrt{x^2+a^2}$，则 $y^2-x^2=a^2$，且 $x\,\mathrm{d}x=y\,\mathrm{d}y$。原式化为 $\cdots$，最后回到 $\int \mathrm{d}x/y$ 或 $\int \mathrm{d}(x/y)$，故结果为 $\cdots$。”
若题目已经天然适合 $x=a\tan t$，三角代换同样合法，但要少走回代弯路。若根式不是 $x^2+a^2$ 而是 $a^2-x^2$ 或 $x^2-a^2$，要换成对应章节的底座。
检查对数里是否写成 $x+\sqrt{x^2+a^2}$；检查回代后是否完全消掉了临时变量 $y$。
\par\medskip


\begin{note}[双元法里的约束别丢]
\begin{itemize}
  \item 设 $y=\sqrt{x^2+a^2}$ 后，$x$ 与 $y$ 仍受约束 $y^2-x^2=a^2$，不能把它们当成彼此独立的两个变量。
  \item 看到 $\int \frac{\mathrm{d}x}{y}$、$\int \frac{\mathrm{d}x}{y^3}$、$\int \frac{\mathrm{d}x}{x^2y}$ 这类形式时，不要急着硬算，先检查是否能凑出 $\mathrm{d}\ln(x+y)$、$\mathrm{d}\!\left(\frac{x}{y}\right)$ 或 $\mathrm{d}\!\left(\frac{y}{x}\right)$。
\end{itemize}
\end{note}

\begin{note}[这一组为什么是虚圆版本]
令 $y=\sqrt{x^2-a^2}$ 后，立刻得到虚圆底座 $x^2-y^2=a^2$。这意味着本节最先该检查的不是三角代换，而是四类虚圆母动作：$\dfrac{\mathrm{d}x}{y}$、$\mathrm{d}\!\left(\dfrac{x}{y}\right)$、$\mathrm{d}\!\left(\dfrac{y}{x}\right)$、$\mathrm{d}\ln(x+y)$。

如果题目主要由 $x,y$ 及其商组成，那么双元法往往比设 $x=a\sec t$ 更短，因为它始终保留了代数底座 $x^2-y^2=a^2$，不用来回在三角和代数形式之间切换。
\end{note}

\par\medskip\noindent\textbf{虚圆根式的处理}\par

题目含 $\sqrt{x^2-a^2}$、$\dfrac{1}{\sqrt{x^2-a^2}}$、$\dfrac{1}{x\sqrt{x^2-a^2}}$、$\dfrac{\sqrt{x^2-a^2}}{x^2}$ 等结构，答案常出现对数或反三角函数。

\begin{align*}
y&=\sqrt{x^2-a^2},\qquad x^2-y^2=a^2,\qquad x\,\mathrm{d}x=y\,\mathrm{d}y,\\
\int \frac{\mathrm{d}x}{y}&=\ln|x+y|+C,\\
\mathrm{d}\!\left(\frac{x}{y}\right)&=-\frac{a^2}{y^3}\,\mathrm{d}x,\qquad
\mathrm{d}\!\left(\frac{y}{x}\right)=\frac{a^2}{x^2y}\,\mathrm{d}x,\\
\int y\,\mathrm{d}x&=\frac{1}{2}xy-\frac{a^2}{2}\ln|x+y|+C.
\end{align*}

先设 $y=\sqrt{x^2-a^2}$。若分母只有一个 $y$，优先调用 $\ln|x+y|$。若分母是 $y^3$，优先看 $\dfrac{x}{y}$ 的微分；若分母是 $x^2y$，优先看 $\dfrac{y}{x}$ 的微分。若分子含 $x$，就用 $x\,\mathrm{d}x=y\,\mathrm{d}y$；若分子含 $x^2$，就换成 $y^2+a^2$。只有当双元改写不顺时，才退回 $x=a\sec t$。

落笔可写为“设 $y=\sqrt{x^2-a^2}$，则 $x^2-y^2=a^2,\ x\,\mathrm{d}x=y\,\mathrm{d}y$”。随后把原式压回四个母动作：$\dfrac{\mathrm{d}x}{y}$、$\dfrac{\mathrm{d}x}{y^3}$、$\dfrac{\mathrm{d}x}{x^2y}$、$y\,\mathrm{d}x$。

“设 $y=\sqrt{x^2-a^2}$。由 $x^2-y^2=a^2$ 及 $x\,\mathrm{d}x=y\,\mathrm{d}y$，原式化为虚圆母式 $\cdots$，故积分为 $\cdots$。”

若根式其实是 $\sqrt{x^2+a^2}$ 或 $\sqrt{a^2-x^2}$，不要乱套符号。若题目只有 $\dfrac{1}{x^2-a^2}$ 而没有根号，应回到上一节的部分分式写法。

检查对数里是 $x+\sqrt{x^2-a^2}$ 而不是别的组合；检查 $\dfrac{x}{y}$ 与 $\dfrac{y}{x}$ 的符号是否写对；检查回代后是否仍满足 $x^2-a^2>0$ 的定义域。
\par\medskip


\begin{note}[这一组为什么是实圆版本]
令 $y=\sqrt{a^2-x^2}$ 后，得到的是实圆底座 $x^2+y^2=a^2$。和上一节相比，最关键的变化是符号：这里有 $x\,\mathrm{d}x=-y\,\mathrm{d}y$，所以参数角、面积微元、比值微元都会自动带上实圆的符号结构。

因此凡是答案里自然出现 $\arcsin$、$\arccos$、面积微元或实圆比值微元的题，都应优先想到这一套底座，而不是把它只当成“又一类根式积分写法”。
\end{note}

\par\medskip\noindent\textbf{实圆根式的处理}\par

题目含 $\sqrt{a^2-x^2}$、$\dfrac{1}{\sqrt{a^2-x^2}}$、$\dfrac{1}{x\sqrt{a^2-x^2}}$、$\dfrac{\sqrt{a^2-x^2}}{x}$ 等结构，答案常出现 $\arcsin$、$\arccos$。

\begin{align*}
y&=\sqrt{a^2-x^2},\qquad x^2+y^2=a^2,\qquad x\,\mathrm{d}x=-y\,\mathrm{d}y,\\
\int \frac{\mathrm{d}x}{y}&=\arcsin\frac{x}{a}+C,\\
\mathrm{d}\!\left(\frac{x}{y}\right)&=\frac{a^2}{y^3}\,\mathrm{d}x,\qquad
\mathrm{d}\!\left(\frac{y}{x}\right)=-\frac{a^2}{x^2y}\,\mathrm{d}x,\\
\int y\,\mathrm{d}x&=\frac{1}{2}xy+\frac{a^2}{2}\arcsin\frac{x}{a}+C.
\end{align*}

先设 $y=\sqrt{a^2-x^2}$。若分母只有一个 $y$，直接出 $\arcsin(x/a)$。若分母是 $y^3$，优先看 $\dfrac{x}{y}$ 的微分；若分母是 $x^2y$，优先看 $\dfrac{y}{x}$ 的微分。若分子含 $x$，就用 $x\,\mathrm{d}x=-y\,\mathrm{d}y$；若分子含 $x^2$，先改成 $a^2-y^2$。若题目还伴随 $\arcsin,\arccos$，就把这里当成它们的基础母式库。

常见处理是“设 $y=\sqrt{a^2-x^2}$，则 $x^2+y^2=a^2,\ x\,\mathrm{d}x=-y\,\mathrm{d}y$”。然后把原式收缩到 $\dfrac{\mathrm{d}x}{y}$、$\dfrac{\mathrm{d}x}{y^3}$、$\dfrac{\mathrm{d}x}{x^2y}$ 或 $y\,\mathrm{d}x$ 这些基本母式。

“设 $y=\sqrt{a^2-x^2}$。由实圆底座 $x^2+y^2=a^2$ 及 $x\,\mathrm{d}x=-y\,\mathrm{d}y$，原式化为 $\cdots$，故结果为 $\cdots$。”

若根号改成 $\sqrt{x^2-a^2}$，符号体系全部要换。若题目只是 $\dfrac{1}{x^2+a^2}$，那是反正切母式，不在本节。

检查 $\arcsin$ 与 $\arccos$ 有没有写混；检查实圆里负号是否保留；检查定义域 $|x|\le a$ 是否与题目相符。
\par\medskip


\begin{note}[配方后二次根式为什么还能回到双元]
这一组表面上比 $\sqrt{x^2\pm a^2}$ 更复杂，但关键只是先把二次式配方，再用
\(\displaystyle u=2ax\pm b,\ v=2\sqrt a\,y\) 把它改造成标准的实圆或虚圆底座。换句话说，这一步不是额外发明新方法，而是在制造前面已经熟悉的双元母场景。

所以读题时若看到根号里是一般二次式，不必急着分别记很多公式。先问自己：能否把它压成 $u^2+v^2=\Delta$ 或 $v^2-u^2=\Delta$？一旦能，后面就只是调用前面的 $\int \frac{\mathrm{d}u}{v}$、面积微元和比值微元。
\end{note}

\par\medskip\noindent\textbf{一般二次根式的处理}\par

根号里一旦是一般二次式，最稳的路径不是单独记新公式，而是先配方，把它压成前面已经熟悉的实圆或虚圆底座。若是 $ax^2+bx+c$，就对应虚圆情形处理；若是 $c+bx-ax^2$，就对应实圆情形处理。

若 \(\displaystyle y=\sqrt{ax^2+bx+c},\ u=2ax+b,\ v=2\sqrt a\,y\)，则 \(\displaystyle v^2-u^2=4ac-b^2\)。
于是
\[
\int \frac{\mathrm{d}x}{\sqrt{ax^2+bx+c}}
=\frac{1}{\sqrt a}\ln|u+v|+C
=\frac{1}{\sqrt a}\ln|2ax+b+2\sqrt a\,y|+C.
\]
若 \(\displaystyle y=\sqrt{c+bx-ax^2},\ u=2ax-b,\ v=2\sqrt a\,y\)，则 \(\displaystyle u^2+v^2=b^2+4ac\)，从而
\[
\int \frac{\mathrm{d}x}{\sqrt{c+bx-ax^2}}
=\frac{1}{\sqrt a}\arcsin\frac{u}{\sqrt{b^2+4ac}}+C.
\]

到这一步后，后面的动作已经很明确：分子若含 $x$，就用 $x=\dfrac{u\mp b}{2a}$ 消去；若要算 $\int y\,\mathrm{d}x$，直接使用面积微元，不必重做分部积分。卷面上可以写成“设 $y=\sqrt{\cdots}$，再令 $u=2ax\pm b,\ v=2\sqrt a\,y$，则得到实圆 / 虚圆底座”，这样后面所有积分都尽量改成 $\dfrac{\mathrm{d}u}{v}$、$v\,\mathrm{d}u$、$\dfrac{u\,\mathrm{d}u}{v}$ 这三类。若根号内其实已是 $x^2\pm a^2$，就不用多此一举；若根号外面只是倒数二次式，也该回到一般二次式积分章节。最后核对 $u=2ax+b$ 还是 $2ax-b$、$\sqrt a$ 有没有漏乘，以及结果的函数族是否和底座类型一致。
\par\medskip


\begin{note}[纯根式双元的想法]
这一组最容易误以为只能靠乘共轭或继续套换元。其实更系统的做法，是把两个根号本身直接当成双元：一个承担 $p$，一个承担 $q$，然后看它们满足的是 $p^2-q^2=\text{常数}$ 还是 $p^2+q^2=\text{常数}$。

这样做的好处是，题目里原本零散的根式差、根式和、分式结构，会立刻回收到标准实圆或虚圆底座里。换句话说，这里不是“根式方法”孤立出现，而是双元法向纯根式情形的自然延伸。
\end{note}

\par\medskip\noindent\textbf{纯根式双元处理}\par

题目出现 $\sqrt{\dfrac{x-a}{x-b}}$、$\sqrt{\dfrac{x-a}{b-x}}$、$\sqrt{(x-a)(b-x)}$ 这类“两个一次根式互相缠绕”的结构时，最自然的做法就是把两个根号本身直接当成双元 \(p,q\)。

若设
\[
p=\sqrt{x-a},\qquad q=\sqrt{x-b},\qquad
p^2-q^2=b-a,\qquad \mathrm{d}x=2p\,\mathrm{d}p=2q\,\mathrm{d}q.
\]
若设
\[
p=\sqrt{x-a},\qquad q=\sqrt{b-x},\qquad
p^2+q^2=b-a,\qquad \mathrm{d}x=2p\,\mathrm{d}p=-2q\,\mathrm{d}q.
\]

先看两个根号是“同向增长”还是“一个增一个减”。若同向，就优先造虚圆 $p^2-q^2=\text{常数}$；若一增一减，就优先造实圆 $p^2+q^2=\text{常数}$。然后把原式中的根式商、根式积全部改成 $p/q$、$pq$、$p^2-q^2$ 或 $p^2+q^2$，能直接变成 $\dfrac{\mathrm{d}p}{q}$、$\dfrac{\mathrm{d}q}{p}$、$p\,\mathrm{d}q$、$q\,\mathrm{d}p$ 时，题目就已经被收回到标准母式里了。

卷面上可以直接写“设 $p=\sqrt{x-a}, q=\sqrt{x-b}$”或“设 $p=\sqrt{x-a}, q=\sqrt{b-x}$”，再把底座恒等式和 $\mathrm{d}x$ 的表达带上。若根号内部不是一次式，而是二次式或三角式，就不要硬套这一节；若已经能单次换元化成 $\sqrt{t^2\pm c}$，也可以直接转去前面的根式章节。最后只需核对 $q$ 的定义、底座的平方和还是平方差，以及收尾函数到底该是对数还是反三角。
\par\medskip


\subsection{含有三角函数的积分}

\begin{note}[三角函数其实就是实圆双元的参数化]
把 \(\displaystyle p=\sin x,\ q=\cos x\) 写出来以后，立刻回到实圆底座 $p^2+q^2=1$；而把 $p=\sec x,q=\tan x$ 写出来，则回到另一种常见的代数底座 $p^2-q^2=1$。所以三角函数积分里很多看似分散的做法，本质上仍是在调用“底座 + 微元 + 递推/对数微元”这同一套结构。

因此，当题目可以明显分离出一对彼此满足平方恒等式的三角量时，就可以把它当成双元题看，而不是只靠积化和差、半角或万能代换硬推。
\end{note}

\par\medskip\noindent\textbf{三角函数积分的处理}\par

题面是三角函数积分，尤其是幂次、乘积、倒幂、混合频率、$R(\sin x,\cos x)$ 这几类。

\begin{align*}
\int \sin x\,\mathrm{d}x&=-\cos x+C,\quad
\int \cos x\,\mathrm{d}x=\sin x+C,\\
\int \tan x\,\mathrm{d}x&=-\ln|\cos x|+C,\quad
\int \cot x\,\mathrm{d}x=\ln|\sin x|+C,\\
\int \sec x\,\mathrm{d}x&=\ln|\sec x+\tan x|+C,\quad
\int \csc x\,\mathrm{d}x=\ln|\csc x-\cot x|+C.
\end{align*}
\[
\sin x=\frac{2t}{1+t^2},\quad
\cos x=\frac{1-t^2}{1+t^2},\quad
\mathrm{d}x=\frac{2}{1+t^2}\,\mathrm{d}t,
\quad t=\tan\frac{x}{2}.
\]
\begin{align*}
\int \sin^n x\,\mathrm{d}x
&=-\frac{1}{n}\sin^{n-1}x\cos x+\frac{n-1}{n}\int \sin^{n-2}x\,\mathrm{d}x,\\
\int \cos^n x\,\mathrm{d}x
&=\frac{1}{n}\cos^{n-1}x\sin x+\frac{n-1}{n}\int \cos^{n-2}x\,\mathrm{d}x.
\end{align*}
\begin{align*}
\int \sec^n x\,\mathrm{d}x
&=\frac{\sin x}{(n-1)\cos^{n-1}x}
+\frac{n-2}{n-1}\int \sec^{n-2}x\,\mathrm{d}x,\\
\int \csc^n x\,\mathrm{d}x
&=-\frac{\cos x}{(n-1)\sin^{n-1}x}
+\frac{n-2}{n-1}\int \csc^{n-2}x\,\mathrm{d}x.
\end{align*}

单个基本函数可直接使用结论；若是 $\sin^m x\cos^n x$，先看奇偶：有一个奇幂时，拆一个出来，剩余部分用 $1-\sin^2 x$ 或 $1-\cos^2 x$；两者都偶时，用半角。若是 $\sec,\tan$ 型，优先找 $\sec x\tan x$ 或 $\sec^2x$；若是 $\csc,\cot$ 型，优先找 $\csc x\cot x$ 或 $\csc^2x$。若是不同频率的 $\sin ax\cos bx$、$\sin ax\sin bx$、$\cos ax\cos bx$，先积化和差。若是 $R(\sin x,\cos x)$ 怎么变都不顺，最后再上万能代换 $t=\tan(x/2)$。

做题时先写“判奇偶”或“判频率”。如果是幂次题，就先写“拆一因子 + 余项降次”或“半角降次”；如果是混合频率，就先写积化和差；如果是有理式，就先写万能代换公式。能用积化和差解决的，不要直接上万能代换；能拆奇幂的，不要先半角。最后只需检查奇偶判定、半角系数 $\frac12$ 和 $\frac14$、万能代换有没有把 $x$ 全换掉，以及结果是否回到原变量。
\par\medskip


\subsection{\texorpdfstring{含有反三角函数的积分 $(a>0)$}{含有反三角函数的积分 (a>0)}}

\par\medskip\noindent\textbf{反三角函数积分的处理}\par

被积函数外层常是 $\arcsin\dfrac{x}{a}$、$\arccos\dfrac{x}{a}$ 或 $\arctan\dfrac{x}{a}$。
前面通常再乘 $1,x,x^2$ 等低次多项式。

\begin{compactenum}
\item \(\displaystyle \left(\arcsin\frac{x}{a}\right)'=\frac{1}{\sqrt{a^2-x^2}}\)。
\item \(\displaystyle \left(\arccos\frac{x}{a}\right)'=-\frac{1}{\sqrt{a^2-x^2}}\)。
\item \(\displaystyle \left(\arctan\frac{x}{a}\right)'=\frac{a}{x^2+a^2}\)。
\end{compactenum}
因此本节余项一定回流到前面两类母式：\(\displaystyle \sqrt{a^2-x^2}\) 或 \(\displaystyle x^2+a^2\)。

这类题通常先分部积分，把反三角函数放进 $u$；再对导数项分类判断，若是 $\arcsin,\arccos$，余项回到 \(\displaystyle \frac{x^k}{\sqrt{a^2-x^2}}\)，若是 $\arctan$，余项回到 \(\displaystyle \frac{x^k}{x^2+a^2}\)；分子若是 $x^2$，就优先拆成 $a^2-(a^2-x^2)$ 或 $(x^2+a^2)-a^2$，不要做无意义长除。
\begin{compactenum}
\item 先分部积分，把反三角函数放进 $u$。
\item 再对导数项分类判断。若是 $\arcsin,\arccos$，余项回到 \(\displaystyle \frac{x^k}{\sqrt{a^2-x^2}}\)；若是 $\arctan$，余项回到 \(\displaystyle \frac{x^k}{x^2+a^2}\)。
\item 分子若是 $x^2$，优先拆成 $a^2-(a^2-x^2)$ 或 $(x^2+a^2)-a^2$，不要做无意义长除。
\end{compactenum}

常见开头是“取 $u=\arcsin\frac{x}{a}$（或 $\arccos,\arctan$），$\mathrm{d}v=x^m\,\mathrm{d}x$”。分部以后，剩下的积分应明确回到 $\sqrt{a^2-x^2}$ 或 $x^2+a^2$ 的母式。不要一上来对 $\arcsin,\arccos,\arctan$ 自己再设新元，通常比直接分部更长；若题目本身就是导数母式，如 $\dfrac{1}{\sqrt{a^2-x^2}}$，应回前面章节，不要留在本节硬转。最后核对分部积分的正负号，尤其是 $\arccos$ 的导数自带负号，以及最后留下的母式是否确实属于实圆根式或反正切母式。
\par\medskip


\subsection{含有指数函数的积分}

\par\medskip\noindent\textbf{指数函数积分的处理}\par

题目含 $e^{ax}$、$a^x$，并与多项式、三角函数、三角幂相乘。

\[
\begin{aligned}
\int e^{ax}\,\mathrm{d}x&=\frac{1}{a}e^{ax}+C,\\
\int a^x\,\mathrm{d}x&=\frac{1}{\ln a}a^x+C.
\end{aligned}
\]
\[
\begin{aligned}
I_n&=\int x^n e^{ax}\,\mathrm{d}x
=\frac{x^n e^{ax}}{a}-\frac{n}{a}I_{n-1},\\
J_n&=\int x^n a^x\,\mathrm{d}x
=\frac{x^n a^x}{\ln a}-\frac{n}{\ln a}J_{n-1}.
\end{aligned}
\]
\begin{compactenum}
\item \(\displaystyle \int e^{ax}\sin bx\,\mathrm{d}x=\frac{e^{ax}}{a^2+b^2}(a\sin bx-b\cos bx)+C\)。
\item \(\displaystyle \int e^{ax}\cos bx\,\mathrm{d}x=\frac{e^{ax}}{a^2+b^2}(a\cos bx+b\sin bx)+C\)。
\end{compactenum}
\[
\int e^{ax}\sin^n bx\,\mathrm{d}x
=\frac{e^{ax}\sin^{n-1}bx}{a^2+n^2b^2}(a\sin bx-nb\cos bx)
+\frac{n(n-1)b^2}{a^2+n^2b^2}I_{n-2}.
\]
\[
\int e^{ax}\cos^n bx\,\mathrm{d}x
=\frac{e^{ax}\cos^{n-1}bx}{a^2+n^2b^2}(a\cos bx+nb\sin bx)
+\frac{n(n-1)b^2}{a^2+n^2b^2}J_{n-2}.
\]

纯指数直接积；若是多项式乘指数，反复分部积分，永远让多项式降次；若是指数乘 $\sin bx$ 或 $\cos bx$，可直接使用闭式，不必临场重推两次分部；若是指数乘三角幂，就按上面的递推把幂次降下去。若底数不是 $e$，就先看成 $a^x=e^{x\ln a}$，只是把 $a$ 换成 $\ln a$。

开头先判“三类对象”：纯指数、指数乘多项式、指数乘振荡项。多项式类统一写递推；振荡类统一写闭式或递推；不要在一种题里混好几套方法。若只是 $a^x$，无需先写成 $e^{x\ln a}$ 再大算，直接套母式更快。最后核对 $a^2+b^2$ 分母是否漏写、三角项前的符号是否正确，以及递推是否真的降了一阶而不是写回原题。
\par\medskip


\subsection{含有对数函数的积分}

\par\medskip\noindent\textbf{对数函数积分的处理}\par

题面含 $\ln x$、$(\ln x)^n$、$x^m(\ln x)^n$、$\dfrac{1}{x\ln x}$ 这几类。

\(\displaystyle \int \ln x\,\mathrm{d}x=x\ln x-x+C\)，以及 \(\displaystyle \int \frac{\mathrm{d}x}{x\ln x}=\ln|\ln x|+C\)。
对 \(\displaystyle I_{m,n}=\int x^m(\ln x)^n\,\mathrm{d}x\)，当 $m\neq -1$ 时，
有 \(\displaystyle I_{m,n}=\frac{x^{m+1}}{m+1}(\ln x)^n-\frac{n}{m+1}I_{m,n-1}\)。
当 $m=-1$ 时，令 $t=\ln x$，则 \(\displaystyle \int \frac{(\ln x)^n}{x}\,\mathrm{d}x=\int t^n\,\mathrm{d}t\)。

若分母正好含一个 $x$，先检查能否把 $\ln x$ 看成新元 $t$。若外面是 $\ln x$ 或 $(\ln x)^n$，优先分部积分，并且通常让 $u=(\ln x)^n$。若前面还乘了 $x^m$，只要 $m\neq -1$，照递推直接降 $n$；若 $m=-1$，立刻改用 $t=\ln x$。对数题的关键不是花哨换元，而是认出“分部降幂”和“$x^{-1}\mathrm{d}x$ 变 $\mathrm{d}t$”这两条路。

常见开头是“取 $u=(\ln x)^n,\ \mathrm{d}v=x^m\mathrm{d}x$”或“令 $t=\ln x$，则 $\mathrm{d}t=\dfrac{\mathrm{d}x}{x}$”。前者负责 $m\neq -1$，后者负责 $m=-1$。若没有 $\dfrac{\mathrm{d}x}{x}$，不要硬做 $t=\ln x$；若只是 $\dfrac{1}{x\ln x}$，不要分部积分，直接凑微分最快。最后核对 $m=-1$ 是否单独处理、分部积分里前面的 $\dfrac{1}{m+1}$，以及 $\ln|\ln x|$ 的绝对值有没有漏。
\par\medskip


\subsection{含有双曲函数的积分}

\begin{note}[双曲函数就是虚圆双元的参数化]
若记 \(\displaystyle p=\cosh x,\qquad q=\sinh x\)，则立刻得到虚圆底座 $p^2-q^2=1$，并伴随微分关系 $\mathrm{d}p=q\,\mathrm{d}x,\ \mathrm{d}q=p\,\mathrm{d}x$。因此本节的对数微元、降次递推、面积型拆分，都不是新花样，而是前面虚圆母式在双曲参数下的直接改写。

所以遇到双曲函数积分时，最稳的想法往往不是先展开成指数，而是先看它是否已经天然给出了一个虚圆双元坐标系。
\end{note}

\par\medskip\noindent\textbf{双曲函数积分的处理}\par

题目含 $\sinh x,\cosh x,\tanh x$ 及其低次幂，或者明显能写成双曲恒等式。

\begin{align*}
(\sinh x)'&=\cosh x,\qquad (\cosh x)'=\sinh x,\qquad \cosh^2x-\sinh^2x=1,\\
\int \sinh x\,\mathrm{d}x&=\cosh x+C,\qquad
\int \cosh x\,\mathrm{d}x=\sinh x+C,\qquad
\int \tanh x\,\mathrm{d}x=\ln\cosh x+C,\\
\sinh^2x&=\frac{\cosh 2x-1}{2},
\qquad
\cosh^2x=\frac{\cosh 2x+1}{2}.
\end{align*}

若是单个 $\sinh,\cosh,\tanh$，可直接使用结论；若是 $\sinh^m x\cosh^n x$，先看能否拆出一个导数因子：$\sinh$ 配 $\mathrm{d}(\cosh x)$，$\cosh$ 配 $\mathrm{d}(\sinh x)$；若是平方，优先用半角公式。若题目里天然出现 \(\displaystyle p=\cosh x,\ q=\sinh x\)，就把它当成虚圆底座 $p^2-q^2=1$ 来看，不必急着拆指数。

开头先判“直接导数 / 半角降次 / 双元底座”。基本题直接凑微分，平方题直接半角，复杂题再写 $p=\cosh x,\ q=\sinh x$。若题目只是 $e^x$ 与 $e^{-x}$ 的组合，直接回指数函数章节也可以；若双曲幂次很高但没有明显导数因子，先试半角再决定是否双元化。最后核对 $\cosh^2x-\sinh^2x=1$ 的符号、$\ln\cosh x$ 是否需要绝对值，以及半角公式里的 $\pm 1$ 有没有写反。
\par\medskip


\subsection{定积分}

\par\medskip\noindent\textbf{定积分中的双元法思路}\par

定积分里出现圆、半圆、三角函数整周期、Wallis 递推、正交性，或者换成双元坐标后端点和对称性能明显简化时，就可以沿这一思路处理。

实圆端点常映到 $(0,1)$、$(1,0)$、$(-1,0)$ 等标准点；复指数法常用
\[
\mathrm{e}^{ix}=\cos x+i\sin x,\qquad
\cos x=\frac{\mathrm{e}^{ix}+\mathrm{e}^{-ix}}{2},\qquad
\sin x=\frac{\mathrm{e}^{ix}-\mathrm{e}^{-ix}}{2i}.
\]
若区间是整周期，先查奇偶性与正交性，再决定是否真的展开。

先判奇偶性、周期性、对称性，能秒杀就别换元。若需要双元法，不是为了求原函数，而是为了把端点、边界项、正交项暴露出来；若是三角函数乘积，优先比较“积化和差”和“复指数双元”谁更短。常见推导可以概括成“端点映射 $\to$ 对称判零或递推 $\to$ 必要时再积分”，例如 Wallis 递推先把 $p=\sin x,\ q=\cos x$ 写出来，再用端点 $(0,1)$、$(1,0)$ 消边界项。若直接用牛顿-莱布尼茨就能结束，不要为了统一风格强行双元化；若区间不是整周期，复指数展开后未必更短。最后核对端点映射是否正确，以及正交性使用条件是否真是整数频率和完整周期。
\par\medskip


\section{习题}

\subsection{第一节习题：模法与多项式逆元}
\begin{exercise}[混合分解 $\frac{x^3+2}{(x-1)^2(x^2+1)}$]
    分解 $\displaystyle F(x) = \frac{x^3+2}{(x-1)^2(x^2+1)}$
\end{exercise}
\begin{solution}
原分母包含重根线性因式 $(x-1)^2$ 与二次因式 $(x^2+1)$。优先使用赋值与求导法处理 $(x-1)^2$。设分解形式中包含 $\frac{A}{x-1} + \frac{B}{(x-1)^2}$，令剩余部分函数为 $G(x) = \frac{x^3+2}{x^2+1}$。
\begin{itemize}
\item \textbf{求最高次系数 $B$（掩盖赋值）：}
直接将 $x=1$ 代入 $G(x)$ 中：\(\displaystyle B = G(1) = \frac{1^3+2}{1^2+1} = \frac{3}{2}\)。
\item \textbf{求次高次系数 $A$（求导）：}
对 $G(x)$ 求一阶导数，并代入 $x=1$：
\[
G'(x)=\frac{3x^2(x^2+1)-(x^3+2)(2x)}{(x^2+1)^2},\qquad A=G'(1)=0.
\]
因此，关于 $(x-1)$ 的部分分式仅有一项 $\frac{3/2}{(x-1)^2}$。
\item \textbf{作差相消与整体降维：}
将已求得的部分分式从原分式中减去，设剩余部分为 $H(x)$：
\[
H(x) = \frac{x^3+2}{(x-1)^2(x^2+1)} - \frac{3/2}{(x-1)^2} = \frac{x^3+2 - \frac{3}{2}(x^2+1)}{(x-1)^2(x^2+1)} = \frac{x^3 - \frac{3}{2}x^2 + \frac{1}{2}}{(x-1)^2(x^2+1)}.
\]
根据代数基本定理，分子必然含有 $(x-1)^2$ 因式。通过多项式综合除法或观察法对其进行因式分解，并代回 $H(x)$ 直接约去 $(x-1)^2$：
\[
x^3 - \frac{3}{2}x^2 + \frac{1}{2} = (x-1)^2 \left(x + \frac{1}{2}\right),\qquad
H(x) = \frac{(x-1)^2 \left(x + \frac{1}{2}\right)}{(x-1)^2(x^2+1)}
= \frac{x + \frac{1}{2}}{x^2+1}.
\]
\item \textbf{合并结果：}
此时所有因式已被完全分离，直接组合即可得到最终分解式 \(\displaystyle F(x) = \frac{3}{2(x-1)^2} + \frac{2x+1}{2(x^2+1)}\)。
\end{itemize}
（注：若 $H(x)$ 的分母中仍包含多个高次不可约多项式，此时再对其应用带余除法的 mod 法，即可规避复数运算与庞大的方程组。）
\end{solution}

\begin{exercise}[分解$\frac{x+2}{(x^2+2x+2)(x^2+1)}$]
分解$\dfrac{x+2}{(x^2+2x+2)(x^2+1)}$
\end{exercise}
\begin{solution}
设大分母 $A = x^2+2x+2$，小分母 $B = x^2+1$。我们要求它们组合出 1。
\begin{itemize}
\item 先大除以小（列竖式即可）：得到$A \equiv 2x+1 \pmod{x^2+1}$
\item 然后用除数$B= x^2+1$除以余式$2x+1$，继续列竖式，可得商 $\frac{2x-1}{4}$，余数 $\frac{5}{4}$
\item 写带余除法：$B - (2x+1)\left(\frac{2x-1}{4}\right) = \frac{5}{4}$，将原本余式$2x+1$替换为$A$和$B$的带余除法式$2x+1=A-B$:\[\frac{5}{4} = B - (A - B)\left(\frac{2x-1}{4}\right)\Leftrightarrow 5 = B\left(2x+3\right) - A\left(2x-1\right)\]
\item 所以原式可以如下分解：
\begin{align*}
\frac{x+2}{AB}=&\frac{(x+2)5}{5AB}=\frac{(x+2)\left( B\left(2x+3\right) - A\left(2x-1\right)\right)}{5AB}\\
=&\frac{(x+2)(2x+3)}{5(x^2+2x+2)} - \frac{(x+2)(2x-1)}{5(x^2+1)}=\frac{2x^2+7x+6}{5(x^2+2x+2)} - \frac{2x^2+3x-2}{5(x^2+1)}\\
=&\frac{2x^2+7x+6}{5(x^2+2x+2)}-\frac25 - \left(\frac{2x^2+3x-2}{5(x^2+1)}-\frac25\right)\\
=&\frac{4-3x}{5(x^2+1)}+ \frac{3x+2}{5(x^2+2x+2)}
\end{align*}
\end{itemize}
\end{solution}

\begin{exercise}[提取 $\frac{x^3+3x}{(x^2+1)^2}$]
提取 $\dfrac{x^3+3x}{(x^2+1)^2}$
\end{exercise}
\begin{solution}
先把分子拆成“一个完整的 $x(x^2+1)$ 再加余项”的形式：
\(\displaystyle x^3+3x=x(x^2+1)+2x\)。
于是原式立刻分成两项：
\(\displaystyle \frac{x^3+3x}{(x^2+1)^2}
= \frac{x(x^2+1) + 2x}{(x^2+1)^2}
= \frac{x}{x^2+1} + \frac{2x}{(x^2+1)^2}\)。
\end{solution}

\begin{exercise}[分解$\frac{1}{(x-1)^2(x^2+1)^2}$]
分解$\dfrac{1}{(x-1)^2(x^2+1)^2}$
\end{exercise}
\begin{solution}
设 $A = x-1, B = x^2+1$。
\begin{itemize}
\item 竖式计算得$B$除以$A$的商$x+1$，余式$2$，即$B - A(x+1) = (x^2+1) - (x^2-1) = 2$
\item 分母 $A^2 B^2$ 指数之和减一 $2+2-1=3$，把上面的等式两边三次方，并按 $B^2$ 和 $A^2$ 提取公因式分组：
\begin{align*}
\left(B-(x+1)A\right)^3&=8,\\
B^3-3(x+1)AB^2+3(x+1)^2A^2B-(x+1)^3A^3&=8,\\
8&=B^2 \underbrace{\left[ B -3(x+1)A \right]}_{\text{这是} A^2 \text{的分子} U(x)}\\
&\quad +A^2 \underbrace{\left[ {3(x+1)^2}B - (x+1)^3A \right]}_{\text{这是} B^2\text{ 的分子} V(x)}.
\end{align*}
\item 代入$A = x-1, B = x^2+1$展开括号：$8=B^2(4-2x^2)+A^2(2x^4+4x^3+6x^2+8x+4)$
\item 继续竖式除法：\begin{align*}
U(x) \div (x-1)^2 &\implies 4-2x^2 =-2(x-1)^2 + 6-4x\\
V(x) \div (x^2+1)^2 &\implies 2x^4+4x^3+6x^2+8x+4= 2(x^2+1)^2 +4x^3+2x^2+8x+2
\end{align*}
\item 分离出的常数抵消。结果：$ \frac{3-2x}{4(x-1)^2} + \frac{2x^3+x^2+4x+1}{4(x^2+1)^2}$
\end{itemize}
\end{solution}

\begin{exercise}[分解 $\frac{2x^2 + 3}{(x-1)(x^2+x+1)}$]
    分解 $\displaystyle F(x) = \frac{2x^2 + 3}{(x-1)(x^2+x+1)}$
\end{exercise}
\begin{solution}
设二次因式 $A = x^2+x+1$，一次因式 $B = x-1$。
\begin{itemize}
\item 利用带余除法，多项式 $A$ 除以 $B$ 可得商为 $x+2$，余数为 $3$，即 $A = B(x+2) + 3$。
\item 移项并两边同除以 $3$，构造出关于常数 $1$ 的恒等式：\(\displaystyle 1 = \frac{1}{3}A - \frac{x+2}{3}B\)。
\item 将原分式的分子 $P(x) = 2x^2+3$ 乘以该恒等式，并除以总分母 $AB$ 进行裂项：
\[
\frac{P(x)}{AB} = \frac{P(x) \left(\frac{1}{3}A - \frac{x+2}{3}B\right)}{AB} = \frac{2x^2+3}{3B} - \frac{(2x^2+3)(x+2)}{3A}.
\]
\item 对拆分后的两部分分别通过加减项进行分子降次，化为真分式：
\begin{align*}
\frac{2x^2+3}{3(x-1)} &= \frac{2(x^2-1)+5}{3(x-1)} = \frac{2(x+1)}{3} + \frac{5}{3(x-1)} \\
\frac{(2x^2+3)(x+2)}{3(x^2+x+1)} &= \frac{2x^3+4x^2+3x+6}{3(x^2+x+1)} = \frac{2x(x^2+x+1) + 2(x^2+x+1) - x + 4}{3(x^2+x+1)} \\
&= \frac{2x+2}{3} + \frac{-x+4}{3(x^2+x+1)}
\end{align*}
\item 将上述两式相减，多项式部分相互抵消，得到最终结果
\[
F(x) = \left(\frac{2x+2}{3} + \frac{5}{3(x-1)}\right) - \left(\frac{2x+2}{3} + \frac{-x+4}{3(x^2+x+1)}\right) = \frac{5}{3(x-1)} + \frac{x-4}{3(x^2+x+1)}.
\]
\end{itemize}
\end{solution}

\begin{exercise}[分解 $\frac{x^5}{(x^2+1)^3}$]
    分解 $\displaystyle F(x) = \frac{x^5}{(x^2+1)^3}$
\end{exercise}
\begin{solution}
\begin{itemize}
\item 观察分母结构，令 $t = x^2+1$，则隐含关系为 $x^2 = t-1$。
\item 将分子 $x^5$ 写成以 $t$ 为基底的多项式展开：\(\displaystyle x^5=x\cdot(x^2)^2=x(t-1)^2=x(t^2-2t+1)\)。再把 $t$ 替换回 $x^2+1$ 并按项约分：
\[
F(x)=\frac{x(x^2+1)^2-2x(x^2+1)+x}{(x^2+1)^3}
=\frac{x}{x^2+1}-\frac{2x}{(x^2+1)^2}+\frac{x}{(x^2+1)^3}.
\]
\end{itemize}
\end{solution}

\begin{exercise}[分解 $\frac{x}{(x-1)^2(x+1)^2}$]
    分解 $\displaystyle F(x) = \frac{x}{(x-1)^2(x+1)^2}$
\end{exercise}
\begin{solution}
设 $A = x+1$，$B = x-1$。
\begin{itemize}
\item 观察因式关系易得 $A - B = 2$，即 $1 = \frac{A - B}{2}$。
\item 原分母包含 $A^2 B^2$，总指数为 $4$。将上述关系式两端同时升至 $3$ 次方，再按 $A^2$ 和 $B^2$ 分组并代回化简：
\begin{align*}
1&=\left(\frac{A-B}{2}\right)^3=\frac{1}{8}(A^3-3A^2B+3AB^2-B^3),\\
1&=A^2\left(\frac{A-3B}{8}\right)+B^2\left(\frac{3A-B}{8}\right),\\
1&=A^2\left(\frac{2-x}{4}\right)+B^2\left(\frac{x+2}{4}\right).
\end{align*}
\item 将该等式代入原分式分子，并与分母 $A^2B^2$ 约分；再对两项分别通过加减项化为真分式：
\begin{align*}
F(x)&=\frac{x \cdot 1}{A^2 B^2}
=\frac{x(2-x)}{4B^2}+\frac{x(x+2)}{4A^2}
=\frac{2x-x^2}{4(x-1)^2}+\frac{x^2+2x}{4(x+1)^2},\\
\frac{2x-x^2}{4(x-1)^2} &= \frac{-(x^2-2x+1)+1}{4(x-1)^2} = -\frac{1}{4} + \frac{1}{4(x-1)^2},\\
\frac{x^2+2x}{4(x+1)^2} &= \frac{(x^2+2x+1)-1}{4(x+1)^2} = \frac{1}{4} - \frac{1}{4(x+1)^2},\\
\therefore\quad
F(x) &= \frac{1}{4(x-1)^2} - \frac{1}{4(x+1)^2}.
\end{align*}
\end{itemize}
\end{solution}

\begin{exercise}[分解 $\frac{x^2}{(x^2+1)^2(x^2+x+1)}$]
    分解 $\displaystyle F(x) = \frac{x^2}{(x^2+1)^2(x^2+x+1)}$
\end{exercise}
\begin{solution}
设 $A = x^2+x+1$，$B = x^2+1$。原式分母结构为 $A B^2$。
\begin{itemize}
\item 观察分子的 $x^2$，易知其等价于 $B - 1$。因此，原分式可以优先实现结构降次：\(\displaystyle \frac{x^2}{A B^2} = \frac{B - 1}{A B^2} = \frac{1}{AB} - \frac{1}{A B^2}\)。
此时问题转化为求 $\frac{1}{AB}$ 与 $\frac{1}{A B^2}$ 的分解，从而规避高次多项式的展开与除法。
\item 对 $\frac{1}{AB}$ 进行分解：由 $A - B = x$，利用带余除法 $B \div x$ 得到 $B = x \cdot x + 1$。
将 $x = A-B$ 代入余式关系得 $1 = B - x(A-B) = (x+1)B - xA$。除以 $AB$ 得到 \(\displaystyle \frac{1}{AB} = \frac{(x+1)B - xA}{AB} = \frac{x+1}{A} - \frac{x}{B}\)。
\item 对 $\frac{1}{A B^2}$ 进行分解：由上述结果两边同除以 $B$ 可得 $\frac{1}{AB^2} = \frac{x+1}{AB} - \frac{x}{B^2}$。
为求 $\frac{x+1}{AB}$，将恒等式 $1 = (x+1)B - xA$ 两端同乘 $(x+1)$，再除以 $AB$，得
\[
x+1=(x+1)^2B-x(x+1)A,\qquad
\frac{1}{AB^2}=\frac{(x+1)^2}{A}-\frac{x^2+x}{B}-\frac{x}{B^2}.
\]
\item 通过加减项，对上述项进行分子降次化简：
\(\displaystyle \frac{(x+1)^2}{A}=1+\frac{x}{A}\)，
\(\displaystyle \frac{x^2+x}{B}=1+\frac{x-1}{B}\)。
代回得 $\frac{1}{A B^2} = \left(1 + \frac{x}{A}\right) - \left(1 + \frac{x-1}{B}\right) - \frac{x}{B^2} = \frac{x}{A} - \frac{x-1}{B} - \frac{x}{B^2}$。
\item 最终两式相减合并：
\begin{align*}
F(x) &= \frac{1}{AB} - \frac{1}{AB^2} = \left(\frac{x+1}{A} - \frac{x}{B}\right) - \left(\frac{x}{A} - \frac{x-1}{B} - \frac{x}{B^2}\right) \\
&= \frac{x+1-x}{A} + \frac{-x+(x-1)}{B} + \frac{x}{B^2} \\
&= \frac{1}{x^2+x+1} - \frac{1}{x^2+1} + \frac{x}{(x^2+1)^2}
\end{align*}
\end{itemize}
\end{solution}

\begin{exercise}[分解 $F(x) = \frac{1}{x(x-1)(x^2+1)}$]
    分解 $\displaystyle F(x) = \frac{1}{x(x-1)(x^2+1)}$
\end{exercise}
\begin{solution}
\begin{itemize}
\item 首先分离因式 $x$。设 $A = x$，其余部分 $M = (x-1)(x^2+1) = x^3 - x^2 + x - 1$。
利用带余除法，$M \div A$ 可得商为 $x^2-x+1$，余数为 $-1$，即 $M = A(x^2-x+1) - 1$。
\item 移项得常数代换 $1 = A(x^2-x+1) - M$，代入原分式分子并裂项：\(\displaystyle \frac{1}{AM} = \frac{A(x^2-x+1) - M}{AM}= \frac{x^2-x+1}{(x-1)(x^2+1)} - \frac{1}{x}\)。
由此完成第一步拆分。
\item 接下来分解新分式 $\frac{x^2-x+1}{(x-1)(x^2+1)}$。设 $B = x^2+1$，$C = x-1$。
带余除法 $B \div C$ 得商为 $x+1$，余数为 $2$，即 $B = C(x+1) + 2$。
\item 构造恒等式 $1 = \frac{1}{2}B - \frac{x+1}{2}C$，乘入分子 $P(x) = x^2-x+1$ 并除以 $BC$：
\[
\frac{P(x)}{BC} = \frac{P(x)\left(\frac{1}{2}B - \frac{x+1}{2}C\right)}{BC} = \frac{x^2-x+1}{2(x-1)} - \frac{(x^2-x+1)(x+1)}{2(x^2+1)}.
\]
\item 对两部分分别通过加减项降次：
\(\displaystyle \frac{x^2-x+1}{2(x-1)}=\frac{x}{2}+\frac{1}{2(x-1)}\)，
\(\displaystyle \frac{x^3+1}{2(x^2+1)}=\frac{x}{2}+\frac{1-x}{2(x^2+1)}\)。
\item 将上述两式相减，一次项相互抵消，并得到原分式分解：
\begin{align*}
\left(\frac{x}{2} + \frac{1}{2(x-1)}\right) - \left(\frac{x}{2} + \frac{1-x}{2(x^2+1)}\right)
&= \frac{1}{2(x-1)} + \frac{x-1}{2(x^2+1)},\\
F(x) &= -\frac{1}{x} + \frac{1}{2(x-1)} + \frac{x-1}{2(x^2+1)}.
\end{align*}
\end{itemize}
\end{solution}

\begin{exercise}[分解 $\frac{x^3}{(x-1)(x+1)(x^2+x+1)}$]
    分解 $\displaystyle F(x) = \frac{x^3}{(x-1)(x+1)(x^2+x+1)}$
\end{exercise}
\begin{solution}
\begin{itemize}
\item 对于一次因式 $(x-1)$ 和 $(x+1)$，可直接利用赋值法求出其对应的部分分式分子：
令 $x=1$，掩去分母中的 $(x-1)$，代入剩余部分得 $\frac{1^3}{(1+1)(1+1+1)} = \frac{1}{6}$，即含 $\frac{1/6}{x-1}$。
令 $x=-1$，掩去分母中的 $(x+1)$，代入剩余部分得 $\frac{(-1)^3}{(-1-1)(1-1+1)} = \frac{1}{2}$，即含 $\frac{1/2}{x+1}$。
\item 将这两项合并通分：
\[
\frac{1}{6(x-1)} + \frac{1}{2(x+1)} = \frac{1(x+1) + 3(x-1)}{6(x^2-1)} = \frac{2x-1}{3(x^2-1)}.
\]
\item 将原分式减去已分离的部分，通分相减并约分：
\begin{align*}
\text{剩余项}
&= \frac{x^3}{(x^2-1)(x^2+x+1)} - \frac{2x-1}{3(x^2-1)},\\
3x^3 - (2x-1)(x^2+x+1)
&= x^3 - x^2 - x + 1=(x^2-1)(x-1),\\
\text{剩余项}
&= \frac{(x^2-1)(x-1)}{3(x^2-1)(x^2+x+1)} = \frac{x-1}{3(x^2+x+1)}.
\end{align*}
\item 综上所述，三部分合并得到最终结果：\(\displaystyle F(x) = \frac{1}{6(x-1)} + \frac{1}{2(x+1)} + \frac{x-1}{3(x^2+x+1)}\)。
\end{itemize}
\end{solution}

\subsection{第二节习题：\texorpdfstring{含有 $ax+b$ 的积分}{含有 ax+b 的积分}}
\begin{exercise}[含有 $ax+b$ 的积分 - 1]
$\displaystyle \int \frac{\mathrm{d}x}{ax+b}$
\end{exercise}
\begin{solution}
令 \(\displaystyle u=ax+b,\ \mathrm{d}u=a\,\mathrm{d}x\)。则 \(\displaystyle \int \frac{\mathrm{d}x}{ax+b}= \frac{1}{a}\int \frac{\mathrm{d}u}{u}= \frac{1}{a}\ln|u|+C= \frac{1}{a}\ln|ax+b|+C\)。
\end{solution}

\begin{exercise}[含有 $ax+b$ 的积分 - 2]
$\displaystyle \int (ax+b)^{\alpha}\mathrm{d}x$
\end{exercise}
\begin{solution}
令 \(\displaystyle u=ax+b,\ \mathrm{d}u=a\,\mathrm{d}x\)。
则
\[
\int (ax+b)^{\alpha}\mathrm{d}x
= \frac{1}{a}\int u^{\alpha}\,\mathrm{d}u
= \frac{1}{a(\alpha+1)}u^{\alpha+1}+C
= \frac{1}{a(\alpha+1)}(ax+b)^{\alpha+1}+C\qquad (\alpha\neq -1).
\]
\end{solution}

\begin{exercise}[含有 $ax+b$ 的积分 - 3]
$\displaystyle \int \frac{x}{ax+b}\mathrm{d}x$
\end{exercise}
\begin{solution}
先把分子改写为 \(\displaystyle x=\frac{(ax+b)-b}{a}\)。
于是
\[
\int \frac{x}{ax+b}\mathrm{d}x
= \frac{1}{a}\int \frac{(ax+b)-b}{ax+b}\,\mathrm{d}x
= \frac{1}{a^2}\int \left(1-\frac{b}{ax+b}\right)\mathrm{d}(ax+b)
= \frac{1}{a^2}\bigl(ax+b-b\ln|ax+b|\bigr)+C.
\]
\end{solution}

\begin{exercise}[含有 $ax+b$ 的积分 - 4]
$\displaystyle \int \frac{x^2}{ax+b}\mathrm{d}x$
\end{exercise}
\begin{solution}
先把 \(\displaystyle x^2=\frac{[(ax+b)-b]^2}{a^2}\) 代入原积分。于是
\begin{align*}
\int \frac{x^2}{ax+b}\mathrm{d}x
&= \frac{1}{a^2}\int \frac{[(ax+b)-b]^2}{ax+b}\,\mathrm{d}x \\
&= \frac{1}{a^3}\int \frac{(ax+b)^2-2b(ax+b)+b^2}{ax+b}\,\mathrm{d}(ax+b) \\
&= \frac{1}{a^3}\int \left(ax+b-2b+\frac{b^2}{ax+b}\right)\mathrm{d}(ax+b) \\
&= \frac{1}{a^3}\left[\frac{1}{2}(ax+b)^2-2b(ax+b)+b^2\ln|ax+b|\right]+C.
\end{align*}
\end{solution}

\begin{exercise}[含有 $ax+b$ 的积分 - 5]
$\displaystyle \int \frac{\mathrm{d}x}{x(ax+b)}$
\end{exercise}
\begin{solution}
先做拆项：\(\displaystyle \frac{1}{x(ax+b)}
=\frac{1}{b}\cdot\frac{ax+b-ax}{x(ax+b)}
=\frac{1}{b}\left(\frac{1}{x}-\frac{a}{ax+b}\right)\)。
因此
\[
\int \frac{\mathrm{d}x}{x(ax+b)}
= \frac{1}{b}\int \left(\frac{1}{x}-\frac{a}{ax+b}\right)\mathrm{d}x
= \frac{1}{b}\bigl(\ln|x|-\ln|ax+b|\bigr)+C
= -\frac{1}{b}\ln\left|\frac{ax+b}{x}\right|+C.
\]
\end{solution}

\begin{exercise}[含有 $ax+b$ 的积分 - 6]
$\displaystyle \int \frac{\mathrm{d}x}{x^2(ax+b)}$
\end{exercise}
\begin{solution}
先拆成 \(\displaystyle \frac{1}{x^2(ax+b)}
=\frac{1}{b}\left(\frac{1}{x^2}-\frac{a}{x(ax+b)}\right)\)，这样就能调用上一题的结果。于是
\begin{align*}
\int \frac{\mathrm{d}x}{x^2(ax+b)}
&= \frac{1}{b}\int \left(\frac{1}{x^2}-\frac{a}{x(ax+b)}\right)\mathrm{d}x \\
&= -\frac{1}{bx}-\frac{a}{b}\int \frac{\mathrm{d}x}{x(ax+b)} \\
&= -\frac{1}{bx}+\frac{a}{b^2}\ln\left|\frac{ax+b}{x}\right|+C.
\end{align*}
\end{solution}

\begin{exercise}[含有 $ax+b$ 的积分 - 7]
$\displaystyle \int \frac{x}{(ax+b)^2}\mathrm{d}x$
\end{exercise}
\begin{solution}
先改写分子为 \(\displaystyle x=\frac{(ax+b)-b}{a}\)。于是
\begin{align*}
\int \frac{x}{(ax+b)^2}\mathrm{d}x
&= \frac{1}{a^2}\int \frac{ax+b-b}{(ax+b)^2}\,\mathrm{d}(ax+b) \\
&= \frac{1}{a^2}\int \left(\frac{1}{ax+b}-\frac{b}{(ax+b)^2}\right)\mathrm{d}(ax+b) \\
&= \frac{1}{a^2}\left[\ln|ax+b|+\frac{b}{ax+b}\right]+C.
\end{align*}
\end{solution}

\begin{exercise}[含有 $ax+b$ 的积分 - 8]
$\displaystyle \int \frac{x^2}{(ax+b)^2}\mathrm{d}x$
\end{exercise}
\begin{solution}
仍令 \(\displaystyle x^2=\frac{[(ax+b)-b]^2}{a^2}\)，
则
\begin{align*}
\int \frac{x^2}{(ax+b)^2}\mathrm{d}x
&= \frac{1}{a^3}\int \frac{(ax+b-b)^2}{(ax+b)^2}\mathrm{d}(ax+b) \\
&= \frac{1}{a^3}\int \left(1-\frac{2b}{ax+b}+\frac{b^2}{(ax+b)^2}\right)\mathrm{d}(ax+b) \\
&= \frac{1}{a^3}\left[ax+b-2b\ln|ax+b|-\frac{b^2}{ax+b}\right]+C.
\end{align*}
\end{solution}

\begin{exercise}[含有 $ax+b$ 的积分 - 9]
$\displaystyle \int \frac{\mathrm{d}x}{x(ax+b)^2}$
\end{exercise}
\begin{solution}
先拆成 \(\displaystyle \frac{1}{x(ax+b)^2}
=\frac{1}{b}\left(\frac{1}{x(ax+b)}-\frac{a}{(ax+b)^2}\right)\)，然后分别积分：
\[
\int \frac{\mathrm{d}x}{x(ax+b)^2}
= \frac{1}{b}\int \left(\frac{1}{x(ax+b)}-\frac{a}{(ax+b)^2}\right)\mathrm{d}x
= -\frac{1}{b^2}\ln\left|\frac{ax+b}{x}\right|+\frac{1}{b(ax+b)}+C.
\]
\end{solution}

\subsection{第三节习题：\texorpdfstring{含有 $\sqrt{ax+b}$ 的积分}{含有 sqrt(ax+b) 的积分}}
\begin{exercise}[含有 $\sqrt{ax+b}$ 的积分 - 10]
$\displaystyle \int \sqrt{ax+b}\mathrm{d}x$
\end{exercise}
\begin{solution}
记 $y=\sqrt{ax+b}$，则 \(\displaystyle y^2=ax+b,\ x=\frac{y^2-b}{a},\ \mathrm{d}x=\frac{2y}{a}\mathrm{d}y\)。
原积分中的根式正好就是 $y$，因此可直接化为幂函数积分
\[
\int \sqrt{ax+b}\mathrm{d}x = \int y\cdot\frac{2y}{a}\mathrm{d}y = \frac{2}{a}\int y^2\mathrm{d}y = \frac{2}{3a}y^3+C = \frac{2}{3a}\sqrt{(ax+b)^3}+C.
\]
\end{solution}

\begin{exercise}[含有 $\sqrt{ax+b}$ 的积分 - 11]
$\displaystyle \int x\sqrt{ax+b}\mathrm{d}x$
\end{exercise}
\begin{solution}
记 $y=\sqrt{ax+b}$，则 \(\displaystyle y^2=ax+b,\ x=\frac{y^2-b}{a},\ \mathrm{d}x=\frac{2y}{a}\mathrm{d}y\)。
这样分子里的 $x$ 与根式都能化成 $y$ 的多项式：
\begin{align*}
\int x\sqrt{ax+b}\mathrm{d}x &= \int \frac{y^2-b}{a}\cdot y\cdot\frac{2y}{a}\mathrm{d}y = \frac{2}{a^2}\int (y^4-by^2)\mathrm{d}y \\
&= \frac{2}{a^2}\left(\frac{1}{5}y^5-\frac{b}{3}y^3\right)+C = \frac{2}{15a^2}y^3(3y^2-5b)+C \\
&= \frac{2}{15a^2}(3ax-2b)\sqrt{(ax+b)^3}+C.
\end{align*}
\end{solution}

\begin{exercise}[含有 $\sqrt{ax+b}$ 的积分 - 12]
$\displaystyle \int x^2\sqrt{ax+b}\mathrm{d}x$
\end{exercise}
\begin{solution}
记 $y=\sqrt{ax+b}$，则 \(\displaystyle y^2=ax+b,\ x=\frac{y^2-b}{a},\ \mathrm{d}x=\frac{2y}{a}\mathrm{d}y\)。
把 $x^2$ 先写成 $\left(\dfrac{y^2-b}{a}\right)^2$，原积分就退化为 $y$ 的多项式积分：
\begin{align*}
\int x^2\sqrt{ax+b}\mathrm{d}x &= \int \left(\frac{y^2-b}{a}\right)^2\cdot y\cdot\frac{2y}{a}\mathrm{d}y = \frac{2}{a^3}\int (y^6-2by^4+b^2y^2)\mathrm{d}y \\
&= \frac{2}{a^3}\left(\frac{1}{7}y^7-\frac{2b}{5}y^5+\frac{b^2}{3}y^3\right)+C = \frac{2}{105a^3}y^3(15y^4-42by^2+35b^2)+C \\
&= \frac{2}{105a^3}(15(ax+b)^2-42b(ax+b)+35b^2)\sqrt{(ax+b)^3}+C \\
&= \frac{2}{105a^3}(15a^2x^2-12abx+8b^2)\sqrt{(ax+b)^3}+C.
\end{align*}
\end{solution}

\begin{exercise}[含有 $\sqrt{ax+b}$ 的积分 - 13]
$\displaystyle \int \frac{x}{\sqrt{ax+b}}\mathrm{d}x$
\end{exercise}
\begin{solution}
记 $y=\sqrt{ax+b}$，则 \(\displaystyle y^2=ax+b,\ x=\frac{y^2-b}{a},\ \mathrm{d}x=\frac{2y}{a}\mathrm{d}y\)。
由于分母中已有一个 $y$，代换后正好可以约去：
\begin{align*}
\int \frac{x}{\sqrt{ax+b}}\mathrm{d}x &= \int \frac{\frac{y^2-b}{a}}{y}\cdot \frac{2y}{a}\mathrm{d}y = \frac{2}{a^2}\int (y^2-b)\mathrm{d}y \\
&= \frac{2}{a^2}\left(\frac{1}{3}y^3-by\right)+C = \frac{2}{3a^2}y(y^2-3b)+C \\
&= \frac{2}{3a^2}(ax-2b)\sqrt{ax+b}+C.
\end{align*}
\end{solution}

\begin{exercise}[含有 $\sqrt{ax+b}$ 的积分 - 14]
$\displaystyle \int \frac{x^2}{\sqrt{ax+b}}\mathrm{d}x$
\end{exercise}
\begin{solution}
记 $y=\sqrt{ax+b}$，则 \(\displaystyle y^2=ax+b,\ x=\frac{y^2-b}{a},\ \mathrm{d}x=\frac{2y}{a}\mathrm{d}y\)。
代换后分母中的根式与 $\mathrm{d}x$ 里的 $y$ 抵消，只剩普通多项式：
\begin{align*}
\int \frac{x^2}{\sqrt{ax+b}}\mathrm{d}x &= \int \frac{\left(\frac{y^2-b}{a}\right)^2}{y}\cdot\frac{2y}{a}\mathrm{d}y = \frac{2}{a^3}\int (y^4-2by^2+b^2)\mathrm{d}y \\
&= \frac{2}{a^3}\left(\frac{1}{5}y^5-\frac{2b}{3}y^3+b^2y\right)+C = \frac{2}{15a^3}y(3y^4-10by^2+15b^2)+C \\
&= \frac{2}{15a^3}(3a^2x^2-4abx+8b^2)\sqrt{ax+b}+C.
\end{align*}
\end{solution}

\begin{exercise}[含有 $\sqrt{ax+b}$ 的积分 - 15]
$\displaystyle \int \frac{\mathrm{d}x}{x\sqrt{ax+b}}$
\end{exercise}
\begin{solution}
记 $y=\sqrt{ax+b}$，则 $y^2=ax+b \implies a\mathrm{d}x=2y\mathrm{d}y$，$x=\frac{y^2-b}{a}$，于是：
\begin{align*}
\int \frac{\mathrm{d}x}{x\sqrt{ax+b}} &= \int \frac{\frac{2y}{a}\mathrm{d}y}{\frac{y^2-b}{a}\cdot y} = \int \frac{2}{y^2-b}\mathrm{d}y \\
&= \begin{cases}
\frac{1}{\sqrt{b}}\ln\left|\frac{y-\sqrt{b}}{y+\sqrt{b}}\right|+C & (b>0) \\
\frac{2}{\sqrt{-b}}\arctan\frac{y}{\sqrt{-b}}+C & (b<0)
\end{cases} = \begin{cases}
\frac{1}{\sqrt{b}}\ln\left|\frac{\sqrt{ax+b}-\sqrt{b}}{\sqrt{ax+b}+\sqrt{b}}\right|+C & (b>0), \\
\frac{2}{\sqrt{-b}}\arctan\sqrt{\frac{ax+b}{-b}}+C & (b<0).
\end{cases}
\end{align*}
\end{solution}

\begin{exercise}[含有 $\sqrt{ax+b}$ 的积分 - 16]
$\displaystyle \int \frac{\mathrm{d}x}{x^2\sqrt{ax+b}}$
\end{exercise}
\begin{solution}
记 $y=\sqrt{ax+b}$，则 \(\displaystyle y^2=ax+b,\ x=\frac{y^2-b}{a},\ \mathrm{d}x=\frac{2y}{a}\mathrm{d}y\)。
代换后会出现 $\dfrac{1}{(y^2-b)^2}$，此时把分子凑成 $2y^2-(y^2-b)$，就能拆成一个恰好可分部的项与第 15 题的母式：
\begin{align*}
\int \frac{\mathrm{d}x}{x^2\sqrt{ax+b}} &= \int \frac{\frac{2y}{a}\mathrm{d}y}{\left(\frac{y^2-b}{a}\right)^2 \cdot y} = \int \frac{2a}{(y^2-b)^2}\mathrm{d}y = \frac{a}{b}\int \frac{y^2+b-(y^2-b)}{(y^2-b)^2}\mathrm{d}y \\
&= \frac{a}{b}\int \frac{2y^2-(y^2-b)}{(y^2-b)^2}\mathrm{d}y = \frac{a}{b}\int \frac{2y^2}{(y^2-b)^2}\mathrm{d}y - \frac{a}{b}\int \frac{1}{y^2-b}\mathrm{d}y \\
&= \frac{a}{b}\int y\mathrm{d}\left(\frac{-1}{y^2-b}\right) - \frac{a}{2b}\int \frac{2}{y^2-b}\mathrm{d}y \\
&= -\frac{ay}{b(y^2-b)} + \frac{a}{b}\int \frac{\mathrm{d}y}{y^2-b} - \frac{a}{2b}\int \frac{2}{y^2-b}\mathrm{d}y \\
&= -\frac{\sqrt{ax+b}}{bx} - \frac{a}{2b}\int \frac{\mathrm{d}x}{x\sqrt{ax+b}} + C.
\end{align*}
\end{solution}

\begin{exercise}[含有 $\sqrt{ax+b}$ 的积分 - 17]
$\displaystyle \int \frac{\sqrt{ax+b}}{x}\mathrm{d}x$
\end{exercise}
\begin{solution}
记 $y=\sqrt{ax+b}$，则 \(\displaystyle y^2=ax+b,\ x=\frac{y^2-b}{a},\ \mathrm{d}x=\frac{2y}{a}\mathrm{d}y\)。
代换后把分式先拆成常数项加第 15 题的母式：
\begin{align*}
\int \frac{\sqrt{ax+b}}{x}\mathrm{d}x &= \int \frac{y}{\frac{y^2-b}{a}}\cdot\frac{2y}{a}\mathrm{d}y = \int \frac{2y^2}{y^2-b}\mathrm{d}y = \int \left(2 + \frac{2b}{y^2-b}\right)\mathrm{d}y \\
&= 2y + b\int \frac{2}{y^2-b}\mathrm{d}y = 2\sqrt{ax+b} + b\int \frac{\mathrm{d}x}{x\sqrt{ax+b}} + C.
\end{align*}
\end{solution}

\begin{exercise}[含有 $\sqrt{ax+b}$ 的积分 - 18]
$\displaystyle \int \frac{\sqrt{ax+b}}{x^2}\mathrm{d}x$
\end{exercise}
\begin{solution}
记 $y=\sqrt{ax+b}$，则 \(\displaystyle y^2=ax+b,\ x=\frac{y^2-b}{a},\ \mathrm{d}x=\frac{2y}{a}\mathrm{d}y\)。
代换后出现 $\dfrac{2ay^2}{(y^2-b)^2}$，把它写成 $ay\cdot\dfrac{2y}{(y^2-b)^2}$，即可直接凑出微分：
\begin{align*}
\int \frac{\sqrt{ax+b}}{x^2}\mathrm{d}x &= \int \frac{y}{\left(\frac{y^2-b}{a}\right)^2}\cdot\frac{2y}{a}\mathrm{d}y = \int \frac{2ay^2}{(y^2-b)^2}\mathrm{d}y \\
&= \int ay\cdot\frac{2y}{(y^2-b)^2}\mathrm{d}y = a\int y\mathrm{d}\left(\frac{-1}{y^2-b}\right) = -\frac{ay}{y^2-b} + \int \frac{a}{y^2-b}\mathrm{d}y \\
&= -\frac{\sqrt{ax+b}}{x} + \frac{a}{2}\int \frac{\mathrm{d}x}{x\sqrt{ax+b}} + C.
\end{align*}
\end{solution}

\subsection{第四节习题：\texorpdfstring{含有 $x^2\pm a^2$ 的积分}{平方项积分}}
\begin{exercise}[含有 $x^2\pm a^2$ 的积分 - 19]
$\displaystyle \int \frac{\mathrm{d}x}{x^2+a^2}$
\end{exercise}
\begin{solution}
提出分母中的常数因子 $a^2$，并令 \(\displaystyle u=\frac{x}{a},\ \mathrm{d}u=\frac{1}{a}\mathrm{d}x\)。
这样分母就化成熟悉的 $1+u^2$ 型：
\begin{align*}
\int \frac{\mathrm{d}x}{x^2+a^2} &= \int \frac{\mathrm{d}x}{a^2\left[\left(\frac{x}{a}\right)^2+1\right]} = \frac{1}{a}\int \frac{\frac{1}{a}\mathrm{d}x}{\left(\frac{x}{a}\right)^2+1} = \frac{1}{a}\int \frac{\mathrm{d}\left(\frac{x}{a}\right)}{\left(\frac{x}{a}\right)^2+1} \\
&= \frac{1}{a}\arctan\frac{x}{a}+C\ (a\neq 0).
\end{align*}
\end{solution}

\begin{exercise}[含有 $x^2\pm a^2$ 的积分 - 20]
$\displaystyle \int \frac{\mathrm{d}x}{(x^2+a^2)^n}$
\end{exercise}
\begin{solution}
记 $I_n = \int \frac{\mathrm{d}x}{(x^2+a^2)^n}$。为了建立降次递推关系，我们从 $I_{n-1}$ 出发，利用分部积分法展开：
\begin{align*}
I_{n-1} &= \int \frac{1}{(x^2+a^2)^{n-1}}\mathrm{d}x = \frac{x}{(x^2+a^2)^{n-1}} - \int x\mathrm{d}\left[\frac{1}{(x^2+a^2)^{n-1}}\right] \\
&= \frac{x}{(x^2+a^2)^{n-1}} - \int x\left[-(n-1)(x^2+a^2)^{-n}\cdot 2x\right]\mathrm{d}x \\
&= \frac{x}{(x^2+a^2)^{n-1}} + 2(n-1)\int \frac{x^2}{(x^2+a^2)^n}\mathrm{d}x \\
&= \frac{x}{(x^2+a^2)^{n-1}} + 2(n-1)\int \frac{(x^2+a^2)-a^2}{(x^2+a^2)^n}\mathrm{d}x \\
&= \frac{x}{(x^2+a^2)^{n-1}} + 2(n-1)\int \frac{1}{(x^2+a^2)^{n-1}}\mathrm{d}x - 2a^2(n-1)\int \frac{1}{(x^2+a^2)^n}\mathrm{d}x \\
&= \frac{x}{(x^2+a^2)^{n-1}} + 2(n-1)I_{n-1} - 2a^2(n-1)I_n.
\end{align*}
将上述等式两边关于目标项 $I_n$ 进行移项整理：
\[
2a^2(n-1)I_n = \frac{x}{(x^2+a^2)^{n-1}} + 2(n-1)I_{n-1} - I_{n-1} = \frac{x}{(x^2+a^2)^{n-1}} + (2n-3)I_{n-1}.
\]
等式两边同除以 $2a^2(n-1)$，即得最终的递推公式：
\begin{align*}
I_n &= \frac{x}{2(n-1)a^2(x^2+a^2)^{n-1}} + \frac{2n-3}{2(n-1)a^2}I_{n-1} \\
&= \frac{x}{2(n-1)a^2(x^2+a^2)^{n-1}} + \frac{2n-3}{2(n-1)a^2}\int \frac{\mathrm{d}x}{(x^2+a^2)^{n-1}}.
\end{align*}
\end{solution}

\begin{exercise}[含有 $x^2\pm a^2$ 的积分 - 21]
$\displaystyle \int \frac{\mathrm{d}x}{x^2-a^2}$
\end{exercise}
\begin{solution}
先把分母因式分解为 $(x-a)(x+a)$，再设 \(\displaystyle \frac{1}{x^2-a^2}=\frac{A}{x-a}+\frac{B}{x+a}\)。
比较系数得 $A=\dfrac{1}{2a},\ B=-\dfrac{1}{2a}$，于是：
\begin{align*}
\int \frac{\mathrm{d}x}{x^2-a^2} &= \int \frac{\mathrm{d}x}{(x-a)(x+a)} = \int \frac{1}{2a}\left(\frac{1}{x-a} - \frac{1}{x+a}\right)\mathrm{d}x \\
&= \frac{1}{2a}\left( \int \frac{\mathrm{d}x}{x-a} - \int \frac{\mathrm{d}x}{x+a} \right) = \frac{1}{2a}(\ln|x-a| - \ln|x+a|) + C \\
&= \frac{1}{2a}\ln\left|\frac{x-a}{x+a}\right|+C.
\end{align*}
\end{solution}

\subsection{第五节习题：\texorpdfstring{含有 $ax^2+b\ (a>0)$ 的积分}{二次项加常数积分}}
\begin{exercise}[含有 $ax^2+b$ 的积分 - 22]
$\displaystyle \int \frac{\mathrm{d}x}{ax^2+b}$
\end{exercise}
\begin{solution}
由于 $a>0$，根据常数 $b$ 的符号分两种情况讨论：

当 $b>0$ 时，利用基础积分公式 $\int \frac{\mathrm{d}u}{u^2+A^2} = \frac{1}{A}\arctan\frac{u}{A}+C$：
\begin{align*}
\int \frac{\mathrm{d}x}{ax^2+b} &= \int \frac{\mathrm{d}x}{(\sqrt{a}x)^2+(\sqrt{b})^2} = \frac{1}{\sqrt{a}}\int \frac{\mathrm{d}(\sqrt{a}x)}{(\sqrt{a}x)^2+(\sqrt{b})^2} \\
&= \frac{1}{\sqrt{a}}\cdot\frac{1}{\sqrt{b}}\arctan\frac{\sqrt{a}x}{\sqrt{b}}+C = \frac{1}{\sqrt{ab}}\arctan\left(\sqrt{\frac{a}{b}}x\right)+C.
\end{align*}

当 $b<0$ 时，记 $-b>0$，利用基础积分公式 $\int \frac{\mathrm{d}u}{u^2-A^2} = \frac{1}{2A}\ln\left|\frac{u-A}{u+A}\right|+C$：
\begin{align*}
\int \frac{\mathrm{d}x}{ax^2+b} &= \int \frac{\mathrm{d}x}{(\sqrt{a}x)^2-(\sqrt{-b})^2} = \frac{1}{\sqrt{a}}\int \frac{\mathrm{d}(\sqrt{a}x)}{(\sqrt{a}x)^2-(\sqrt{-b})^2} \\
&= \frac{1}{\sqrt{a}}\cdot\frac{1}{2\sqrt{-b}}\ln\left|\frac{\sqrt{a}x-\sqrt{-b}}{\sqrt{a}x+\sqrt{-b}}\right|+C = \frac{1}{2\sqrt{-ab}}\ln\left|\frac{\sqrt{a}x-\sqrt{-b}}{\sqrt{a}x+\sqrt{-b}}\right|+C.
\end{align*}
\end{solution}

\begin{exercise}[含有 $ax^2+b$ 的积分 - 23]
$\displaystyle \int \frac{x}{ax^2+b}\mathrm{d}x$
\end{exercise}
\begin{solution}
这是最典型的“分子差一个常数倍就是分母导数”的情形。因为
\(\displaystyle \mathrm{d}(ax^2+b)=2ax\,\mathrm{d}x\)，
所以先把分子 $x$ 改写成 $\dfrac{1}{2a}(2ax)$：
\(\displaystyle
\int \frac{x}{ax^2+b}\mathrm{d}x
= \frac{1}{2a}\int \frac{2ax}{ax^2+b}\mathrm{d}x
= \frac{1}{2a}\int \frac{\mathrm{d}(ax^2+b)}{ax^2+b}
= \frac{1}{2a}\ln|ax^2+b|+C\)。
\end{solution}

\begin{exercise}[含有 $ax^2+b$ 的积分 - 24]
$\displaystyle \int \frac{x^2}{ax^2+b}\mathrm{d}x$
\end{exercise}
\begin{solution}
分子次数与分母相同，先做一次代数拆项最省力。把
\(\displaystyle ax^2=(ax^2+b)-b\)
代回分子，就能把原式拆成“常数项 + 基本母式”：
\begin{align*}
\int \frac{x^2}{ax^2+b}\mathrm{d}x &= \frac{1}{a}\int \frac{ax^2}{ax^2+b}\mathrm{d}x = \frac{1}{a}\int \frac{(ax^2+b)-b}{ax^2+b}\mathrm{d}x = \frac{1}{a}\int \left(1 - \frac{b}{ax^2+b}\right)\mathrm{d}x \\
&= \frac{1}{a}\int 1\,\mathrm{d}x - \frac{b}{a}\int \frac{\mathrm{d}x}{ax^2+b} = \frac{x}{a} - \frac{b}{a}\int \frac{\mathrm{d}x}{ax^2+b}.
\end{align*}
\end{solution}

\begin{exercise}[含有 $ax^2+b$ 的积分 - 25]
$\displaystyle \int \frac{\mathrm{d}x}{x(ax^2+b)}$
\end{exercise}
\begin{solution}
这一题若直接对 $x$ 做部分分式会比较绕，更方便的是先凑出 $\mathrm{d}(x^2)$。因为
\(\displaystyle \mathrm{d}(x^2)=2x\,\mathrm{d}x\)，
故可先把原积分改写成关于 $x^2$ 的有理积分，再做拆分：
\begin{align*}
\int \frac{\mathrm{d}x}{x(ax^2+b)} &= \int \frac{x}{x^2(ax^2+b)}\mathrm{d}x = \frac{1}{2}\int \frac{\mathrm{d}(x^2)}{x^2(ax^2+b)} = \frac{1}{2}\int \frac{1}{b}\left(\frac{1}{x^2} - \frac{a}{ax^2+b}\right)\mathrm{d}(x^2) \\
&= \frac{1}{2b}\left( \int \frac{\mathrm{d}(x^2)}{x^2} - \int \frac{a\,\mathrm{d}(x^2)}{ax^2+b} \right) = \frac{1}{2b}\left( \ln|x^2| - \int \frac{\mathrm{d}(ax^2+b)}{ax^2+b} \right) \\
&= \frac{1}{2b}\left( \ln(x^2) - \ln|ax^2+b| \right) + C = \frac{1}{2b}\ln\frac{x^2}{|ax^2+b|}+C.
\end{align*}
\end{solution}

\begin{exercise}[含有 $ax^2+b$ 的积分 - 26]
$\displaystyle \int \frac{\mathrm{d}x}{x^2(ax^2+b)}$
\end{exercise}
\begin{solution}
这一题的关键仍是用
\(\displaystyle b=(ax^2+b)-ax^2\)
把分式拆成两项，其中一项是 $x^{-2}$，另一项回到第 22 题的基本母式：
\begin{align*}
\int \frac{\mathrm{d}x}{x^2(ax^2+b)} &= \frac{1}{b}\int \frac{b}{x^2(ax^2+b)}\mathrm{d}x = \frac{1}{b}\int \frac{(ax^2+b)-ax^2}{x^2(ax^2+b)}\mathrm{d}x \\
&= \frac{1}{b}\int \left(\frac{1}{x^2} - \frac{a}{ax^2+b}\right)\mathrm{d}x = \frac{1}{b}\int x^{-2}\mathrm{d}x - \frac{a}{b}\int \frac{\mathrm{d}x}{ax^2+b} \\
&= -\frac{1}{bx} - \frac{a}{b}\int \frac{\mathrm{d}x}{ax^2+b}.
\end{align*}
\end{solution}

\begin{exercise}[含有 $ax^2+b$ 的积分 - 27]
$\displaystyle \int \frac{\mathrm{d}x}{x^3(ax^2+b)}$
\end{exercise}
\begin{solution}
与上题代数方法类似，在分子上凑出 $b = (ax^2+b)-ax^2$ 进行分离：
\begin{align*}
\int \frac{\mathrm{d}x}{x^3(ax^2+b)} &= \frac{1}{b}\int \frac{b}{x^3(ax^2+b)}\mathrm{d}x = \frac{1}{b}\int \frac{(ax^2+b)-ax^2}{x^3(ax^2+b)}\mathrm{d}x \\
&= \frac{1}{b}\int \left(\frac{1}{x^3} - \frac{a}{x(ax^2+b)}\right)\mathrm{d}x = \frac{1}{b}\int x^{-3}\mathrm{d}x - \frac{a}{b}\int \frac{\mathrm{d}x}{x(ax^2+b)}
\end{align*}
此时，代入我们在第 25 题中得到的结论 $\displaystyle \int \frac{\mathrm{d}x}{x(ax^2+b)} = \frac{1}{2b}\ln\frac{x^2}{|ax^2+b|}$：
\begin{align*}
\int \frac{\mathrm{d}x}{x^3(ax^2+b)} &= \frac{1}{b}\left(-\frac{1}{2x^2}\right) - \frac{a}{b}\left(\frac{1}{2b}\ln\frac{x^2}{|ax^2+b|}\right) + C \\
&= -\frac{1}{2bx^2} - \frac{a}{2b^2}\ln\frac{x^2}{|ax^2+b|} + C = \frac{a}{2b^2}\ln\frac{|ax^2+b|}{x^2} - \frac{1}{2bx^2} + C.
\end{align*}
\end{solution}

\begin{exercise}[含有 $ax^2+b$ 的积分 - 28]
$\displaystyle \int \frac{\mathrm{d}x}{(ax^2+b)^2}$
\end{exercise}
\begin{solution}
为了得到降次递推公式，我们从低次积分 $\int \frac{\mathrm{d}x}{ax^2+b}$ 出发，运用分部积分法：
\begin{align*}
\int \frac{\mathrm{d}x}{ax^2+b} &= \frac{x}{ax^2+b} - \int x\,\mathrm{d}\left(\frac{1}{ax^2+b}\right) = \frac{x}{ax^2+b} - \int x\left(-\frac{2ax}{(ax^2+b)^2}\right)\mathrm{d}x \\
&= \frac{x}{ax^2+b} + 2a\int \frac{x^2}{(ax^2+b)^2}\mathrm{d}x = \frac{x}{ax^2+b} + \frac{2a}{a}\int \frac{ax^2}{(ax^2+b)^2}\mathrm{d}x \\
&= \frac{x}{ax^2+b} + 2\int \frac{(ax^2+b)-b}{(ax^2+b)^2}\mathrm{d}x \\
&= \frac{x}{ax^2+b} + 2\int \frac{1}{ax^2+b}\mathrm{d}x - 2b\int \frac{\mathrm{d}x}{(ax^2+b)^2}.
\end{align*}
将包含目标积分 $\int \frac{\mathrm{d}x}{(ax^2+b)^2}$ 的项移至等式左边，其余项移至右边合并：
\begin{align*}
2b\int \frac{\mathrm{d}x}{(ax^2+b)^2}
&= \frac{x}{ax^2+b} + 2\int \frac{\mathrm{d}x}{ax^2+b} - \int \frac{\mathrm{d}x}{ax^2+b}
= \frac{x}{ax^2+b} + \int \frac{\mathrm{d}x}{ax^2+b},\\
\int \frac{\mathrm{d}x}{(ax^2+b)^2}
&= \frac{x}{2b(ax^2+b)} + \frac{1}{2b}\int \frac{\mathrm{d}x}{ax^2+b}.
\end{align*}
\end{solution}

\subsection{第六节习题：\texorpdfstring{含有 $ax^2+bx+c\ (a>0)$ 的积分}{二次三项式积分}}
\begin{exercise}[含有 $ax^2+bx+c$ 的积分 - 29]
$\displaystyle \int \frac{\mathrm{d}x}{ax^2+bx+c}$
\end{exercise}
\begin{solution}
首先对二次三项式进行配方，并提取常数 $a$：
\begin{align*}
\int \frac{\mathrm{d}x}{ax^2+bx+c} &= \int \frac{\mathrm{d}x}{a\left(x^2+\frac{b}{a}x\right)+c} = \int \frac{\mathrm{d}x}{a\left(x^2+\frac{b}{a}x+\frac{b^2}{4a^2}\right)+c-\frac{b^2}{4a}} \\
&= \int \frac{\mathrm{d}x}{a\left(x+\frac{b}{2a}\right)^2+\frac{4ac-b^2}{4a}} = \frac{1}{a}\int \frac{\mathrm{d}\left(x+\frac{b}{2a}\right)}{\left(x+\frac{b}{2a}\right)^2+\frac{4ac-b^2}{4a^2}}.
\end{align*}

此时，根据判别式 $\Delta = b^2-4ac$ 的符号分两种情况讨论：
\begin{itemize}
\item 当 $b^2 > 4ac$ 时，$\frac{4ac-b^2}{4a^2} < 0$，记为负平方，并使用基础积分公式
\[
\frac{4ac-b^2}{4a^2}=-\left(\frac{\sqrt{b^2-4ac}}{2a}\right)^2,\qquad
\int \frac{\mathrm{d}u}{u^2-A^2} = \frac{1}{2A}\ln\left|\frac{u-A}{u+A}\right|+C,
\]
有
\begin{align*}
\int \frac{\mathrm{d}x}{ax^2+bx+c}
&= \frac{1}{a}\int \frac{\mathrm{d}\left(x+\frac{b}{2a}\right)}{\left(x+\frac{b}{2a}\right)^2-\left(\frac{\sqrt{b^2-4ac}}{2a}\right)^2} \\
&= \frac{1}{a} \cdot \frac{1}{2\left(\frac{\sqrt{b^2-4ac}}{2a}\right)}
   \ln\left|\frac{x+\frac{b}{2a}-\frac{\sqrt{b^2-4ac}}{2a}}{x+\frac{b}{2a}+\frac{\sqrt{b^2-4ac}}{2a}}\right| + C \\
&= \frac{1}{\sqrt{b^2-4ac}}
   \ln\left|\frac{\frac{2ax+b-\sqrt{b^2-4ac}}{2a}}{\frac{2ax+b+\sqrt{b^2-4ac}}{2a}}\right| + C \\
&= \frac{1}{\sqrt{b^2-4ac}}
   \ln\left|\frac{2ax+b-\sqrt{b^2-4ac}}{2ax+b+\sqrt{b^2-4ac}}\right| + C.
\end{align*}
\item 当 $b^2 < 4ac$ 时，$\frac{4ac-b^2}{4a^2} > 0$，记为正平方，并使用基础积分公式
\[
\frac{4ac-b^2}{4a^2}=\left(\frac{\sqrt{4ac-b^2}}{2a}\right)^2,\qquad
\int \frac{\mathrm{d}u}{u^2+A^2} = \frac{1}{A}\arctan\frac{u}{A}+C,
\]
有
\begin{align*}
\int \frac{\mathrm{d}x}{ax^2+bx+c}
&= \frac{1}{a}\int \frac{\mathrm{d}\left(x+\frac{b}{2a}\right)}{\left(x+\frac{b}{2a}\right)^2+\left(\frac{\sqrt{4ac-b^2}}{2a}\right)^2} \\
&= \frac{1}{a} \cdot \frac{1}{\frac{\sqrt{4ac-b^2}}{2a}}
   \arctan\frac{x+\frac{b}{2a}}{\frac{\sqrt{4ac-b^2}}{2a}} + C \\
&= \frac{2}{\sqrt{4ac-b^2}}
   \arctan\frac{\frac{2ax+b}{2a}}{\frac{\sqrt{4ac-b^2}}{2a}} + C \\
&= \frac{2}{\sqrt{4ac-b^2}}
   \arctan\frac{2ax+b}{\sqrt{4ac-b^2}} + C.
\end{align*}
\end{itemize}
\end{solution}

\begin{exercise}[含有 $ax^2+bx+c$ 的积分 - 30]
$\displaystyle \int \frac{x}{ax^2+bx+c}\mathrm{d}x$
\end{exercise}
\begin{solution}
为了让分子与分母导数对上，先注意 \(\displaystyle \mathrm{d}(ax^2+bx+c)=(2ax+b)\mathrm{d}x\)。
于是把分子中的 $x$ 改造成 $(2ax+b)-b$ 的形式，原积分就能拆成一个对数项加一个基础母式：
\begin{align*}
\int \frac{x}{ax^2+bx+c}\mathrm{d}x &= \frac{1}{2a}\int \frac{2ax}{ax^2+bx+c}\mathrm{d}x = \frac{1}{2a}\int \frac{(2ax+b)-b}{ax^2+bx+c}\mathrm{d}x \\
&= \frac{1}{2a}\int \frac{2ax+b}{ax^2+bx+c}\mathrm{d}x - \frac{b}{2a}\int \frac{\mathrm{d}x}{ax^2+bx+c} \\
&= \frac{1}{2a}\int \frac{\mathrm{d}(ax^2+bx+c)}{ax^2+bx+c} - \frac{b}{2a}\int \frac{\mathrm{d}x}{ax^2+bx+c} \\
&= \frac{1}{2a}\ln|ax^2+bx+c|-\frac{b}{2a}\int \frac{\mathrm{d}x}{ax^2+bx+c}.
\end{align*}
\end{solution}

\subsection{第七节习题：\texorpdfstring{含有 $\sqrt{x^2+a^2}\ (a>0)$ 的积分}{根式积分（加号）}}
\begin{exercise}[含有 $\sqrt{x^2+a^2}$ 的积分 - 31]
$\displaystyle \int \frac{\mathrm{d}x}{\sqrt{x^2+a^2}}$
\end{exercise}
\begin{solution}
\textbf{思路：}这是整组题的基石公式。方法一把根号三角化；方法二直接寻找对数微分 $\mathrm{d}\ln(x+y)$，后面很多题都会反复调用这个结果。
\begin{enumerate}
\item 方法一：令 $x = a\tan t$，$\mathrm{d}x = a\sec^2 t\,\mathrm{d}t$，则 \(\displaystyle \int \frac{\mathrm{d}x}{\sqrt{x^2+a^2}} = \int \sec t\,\mathrm{d}t = \ln|\sec t+\tan t|+C = \ln(x+\sqrt{x^2+a^2})+C\)。
\item 方法二：记 $y=\sqrt{x^2+a^2}$。由 \(\displaystyle \mathrm{d}y=\frac{x}{y}\,\mathrm{d}x\)，得 \(\displaystyle \int \frac{\mathrm{d}x}{y} = \int \frac{\mathrm{d}x+\mathrm{d}y}{x+y} = \ln(x+y)+C = \ln(x+\sqrt{x^2+a^2})+C\)。
\end{enumerate}
这个例子也提示：以后只要分母里恰好出现一次 $\sqrt{x^2+a^2}$，都要先留意 $\mathrm{d}\ln(x+\sqrt{x^2+a^2})$ 这一母式。
\end{solution}

\begin{exercise}[含有 $\sqrt{x^2+a^2}$ 的积分 - 32]
$\displaystyle \int \frac{\mathrm{d}x}{\sqrt{(x^2+a^2)^3}}$
\end{exercise}
\begin{solution}
\textbf{思路：}分母里出现 $y^3$ 时，最自然的双元微分不是 $\ln(x+y)$，而是比值 $\dfrac{x}{y}$ 的微分。
\begin{enumerate}
\item 方法一：令 $x = a\tan t$，$\mathrm{d}x = a\sec^2 t\,\mathrm{d}t$，则 \(\displaystyle \int \frac{\mathrm{d}x}{(x^2+a^2)^{\frac{3}{2}}} = \frac{1}{a^2}\int \cos t\,\mathrm{d}t = \frac{x}{a^2\sqrt{x^2+a^2}}+C\)。
\item 方法二：记 $y=\sqrt{x^2+a^2}$。由 \(y^2-x^2=a^2\) 得 \(\displaystyle \mathrm{d}\!\left(\frac{x}{y}\right)=\frac{y\,\mathrm{d}x-x\,\mathrm{d}y}{y^2}=\frac{a^2}{y^3}\,\mathrm{d}x\)，所以 \(\displaystyle \int \frac{\mathrm{d}x}{y^3}=\frac{1}{a^2}\int \mathrm{d}\!\left(\frac{x}{y}\right)=\frac{x}{a^2y}+C\)。
\end{enumerate}
这个例子也提示：凡是看到 $\dfrac{\mathrm{d}x}{(x^2+a^2)^{3/2}}$ 或更高次分母，都可以优先检查 $\dfrac{x}{y}$、$\dfrac{y}{x}$ 这样的比值微分。
\end{solution}

\begin{exercise}[含有 $\sqrt{x^2+a^2}$ 的积分 - 33]
$\displaystyle \int \frac{x}{\sqrt{x^2+a^2}}\mathrm{d}x$
\end{exercise}
\begin{solution}
\textbf{思路：}分子里带一个 $x$，最直接的信号就是凑 $\mathrm{d}(x^2+a^2)$；在双元法里则对应于把 $x\,\mathrm{d}x$ 换成 $y\,\mathrm{d}y$。
\begin{enumerate}
\item 方法一：凑微分法：\(\displaystyle \int \frac{x}{\sqrt{x^2+a^2}}\mathrm{d}x=\frac{1}{2}\int (x^2+a^2)^{-\frac{1}{2}}\mathrm{d}(x^2+a^2)=\sqrt{x^2+a^2}+C\)。
\item 方法二：记 $y=\sqrt{x^2+a^2}$。由于 \(x\,\mathrm{d}x=y\,\mathrm{d}y\)，所以原积分立即变成 \(\displaystyle \int \frac{x}{y}\mathrm{d}x=\int \frac{y\mathrm{d}y}{y}=\int \mathrm{d}y=y+C=\sqrt{x^2+a^2}+C\)。
\end{enumerate}
\end{solution}

\begin{exercise}[含有 $\sqrt{x^2+a^2}$ 的积分 - 34]
$\displaystyle \int \frac{x}{\sqrt{(x^2+a^2)^3}}\mathrm{d}x$
\end{exercise}
\begin{solution}
\textbf{思路：}这一题仍然由 $x\,\mathrm{d}x$ 触发，只是多了一个 $y^2$ 的幂次，因此换到 $y$ 后会变成最普通的幂函数积分。
\begin{enumerate}
\item 方法一：凑微分法：\(\displaystyle \int \frac{x}{(x^2+a^2)^{\frac{3}{2}}}\mathrm{d}x
= \frac{1}{2}\int (x^2+a^2)^{-\frac{3}{2}}\mathrm{d}(x^2+a^2)
=-\frac{1}{\sqrt{x^2+a^2}}+C\)。
\item 方法二：记 $y=\sqrt{x^2+a^2}$，再用 $x\,\mathrm{d}x=y\,\mathrm{d}y$ 转换变量：\(\displaystyle \int \frac{x}{y^3}\mathrm{d}x=\int \frac{y\mathrm{d}y}{y^3}=\int y^{-2}\mathrm{d}y=-\frac{1}{y}+C=-\frac{1}{\sqrt{x^2+a^2}}+C\)。
\end{enumerate}
\end{solution}

\begin{exercise}[含有 $\sqrt{x^2+a^2}$ 的积分 - 35]
$\displaystyle \int \frac{x^2}{\sqrt{x^2+a^2}}\mathrm{d}x$
\end{exercise}
\begin{solution}
\textbf{思路：}出现 $x^2$ 时，双元法最重要的操作是先用 $x^2=y^2-a^2$ 把分子降成 $y$ 与常数两部分，再调用第 31 题的基石公式。
\begin{enumerate}
\item 方法一：分部积分法：
\begin{align*}
\int \frac{x^2}{\sqrt{x^2+a^2}}\mathrm{d}x &= \int x\mathrm{d}\left(\sqrt{x^2+a^2}\right) = x\sqrt{x^2+a^2} - \int \sqrt{x^2+a^2}\mathrm{d}x \\
&= x\sqrt{x^2+a^2} - \int \frac{x^2+a^2}{\sqrt{x^2+a^2}}\mathrm{d}x \\
&= x\sqrt{x^2+a^2} - \int \frac{x^2}{\sqrt{x^2+a^2}}\mathrm{d}x - a^2\int \frac{\mathrm{d}x}{\sqrt{x^2+a^2}} \\
\implies 2\int \frac{x^2}{\sqrt{x^2+a^2}}\mathrm{d}x &= x\sqrt{x^2+a^2} - a^2\ln(x+\sqrt{x^2+a^2}) .
\end{align*}
\item 方法二：记 $y=\sqrt{x^2+a^2}$。因为 $y^2-x^2=a^2$，所以 \(x^2=y^2-a^2\)。把它代回原式，原积分就会自动拆成“一个 $\int y\,\mathrm{d}x$”与“一个基石对数积分”：
\begin{align*}
\int \frac{x^2}{y}\mathrm{d}x &= \int \frac{y^2-a^2}{y}\mathrm{d}x = \int y\mathrm{d}x - a^2\int \frac{\mathrm{d}x}{y} \\
&= \left(\frac{1}{2}xy + \frac{a^2}{2}\ln(x+y)\right) - a^2\ln(x+y) + C \\
&= \frac{1}{2}xy - \frac{a^2}{2}\ln(x+y)+C = \frac{x}{2}\sqrt{x^2+a^2}-\frac{a^2}{2}\ln(x+\sqrt{x^2+a^2})+C.
\end{align*}
\end{enumerate}
这个例子也提示：以后只要分子里是 $x^2$、$x^4$ 一类偶次幂，都应先问自己能否用 $x^2=y^2-a^2$ 降阶。
\end{solution}

\begin{exercise}[含有 $\sqrt{x^2+a^2}$ 的积分 - 36]
$\displaystyle \int \frac{x^2}{\sqrt{(x^2+a^2)^3}}\mathrm{d}x$
\end{exercise}
\begin{solution}
\textbf{思路：}与上一题同理，先把 $x^2$ 改写成 $y^2-a^2$，便会拆成第 31 题和第 32 题两种已知母式。
\begin{enumerate}
\item 方法一：分部积分法：
\begin{align*}
\int x\cdot\frac{x}{(x^2+a^2)^{\frac{3}{2}}}\mathrm{d}x &= \int x\mathrm{d}\left(-\frac{1}{\sqrt{x^2+a^2}}\right) = -\frac{x}{\sqrt{x^2+a^2}} + \int \frac{\mathrm{d}x}{\sqrt{x^2+a^2}} \\
&= -\frac{x}{\sqrt{x^2+a^2}} + \ln(x+\sqrt{x^2+a^2})+C.
\end{align*}
\item 方法二：记 $y=\sqrt{x^2+a^2}$，并用 $x^2=y^2-a^2$ 降阶：
\begin{align*}
\int \frac{x^2}{y^3}\mathrm{d}x &= \int \frac{y^2-a^2}{y^3}\mathrm{d}x = \int \frac{\mathrm{d}x}{y} - a^2\int \frac{\mathrm{d}x}{y^3} \\
&= \ln(x+y) - a^2\left(\frac{x}{a^2y}\right)+C = -\frac{x}{\sqrt{x^2+a^2}} + \ln(x+\sqrt{x^2+a^2})+C.
\end{align*}
\end{enumerate}
\end{solution}

\begin{exercise}[含有 $\sqrt{x^2+a^2}$ 的积分 - 37]
$\displaystyle \int \frac{\mathrm{d}x}{x\sqrt{x^2+a^2}}$
\end{exercise}
\begin{solution}
\textbf{思路：}这里最大的障碍是额外多出的 $x$。双元法的关键是先用 $\dfrac{\mathrm{d}x}{y}=\dfrac{\mathrm{d}y}{x}$ 把微分转移到 $y$ 上，再利用 $x^2=y^2-a^2$ 完成一元分式积分。
\begin{enumerate}
\item 方法一：令 $x=a\tan t$，则 $\mathrm{d}x=a\sec^2 t\mathrm{d}t$：
\begin{align*}
\int \frac{a\sec^2 t}{a\tan t\cdot a\sec t}\mathrm{d}t &= \frac{1}{a}\int \csc t\mathrm{d}t = \frac{1}{a}\ln|\csc t-\cot t|+C \\
&= \frac{1}{a}\ln\left|\frac{\sqrt{x^2+a^2}}{x}-\frac{a}{x}\right|+C = \frac{1}{a}\ln\frac{\sqrt{x^2+a^2}-a}{|x|}+C.
\end{align*}
\item 方法二：记 $y=\sqrt{x^2+a^2}$。由 \(\displaystyle x\,\mathrm{d}x=y\,\mathrm{d}y,\ \frac{\mathrm{d}x}{y}=\frac{\mathrm{d}y}{x}\)，可把原题中“不好处理的 $\mathrm{d}x/y$”改写成“对 $y$ 的普通积分”：\(\displaystyle \int \frac{\mathrm{d}x}{xy} = \int \frac{1}{x}\left(\frac{\mathrm{d}y}{x}\right) = \int \frac{\mathrm{d}y}{x^2} = \int \frac{\mathrm{d}y}{y^2-a^2}\)。继续积分得 \(\displaystyle \frac{1}{2a}\ln\left|\frac{y-a}{y+a}\right|+C = \frac{1}{2a}\ln\frac{(y-a)^2}{y^2-a^2}+C = \frac{1}{a}\ln\frac{y-a}{|x|}+C\)。
\end{enumerate}
这个例子也提示：当根式外面额外多出 $x^{-1}$、$x^{-2}$ 之类的因子时，优先考虑把 $\mathrm{d}x/y$ 改写成 $\mathrm{d}y/x$。
\end{solution}

\begin{exercise}[含有 $\sqrt{x^2+a^2}$ 的积分 - 38]
$\displaystyle \int \frac{\mathrm{d}x}{x^2\sqrt{x^2+a^2}}$
\end{exercise}
\begin{solution}
\textbf{思路：}与第 32 题相呼应，这一题最合适的微分对象是 $\dfrac{y}{x}$，因为它的导数会自然生出 $\dfrac{1}{x^2y}$。
\begin{enumerate}
\item 方法一：令 $x=a\tan t$，则 $\mathrm{d}x=a\sec^2 t\mathrm{d}t$：
\begin{align*}
\int \frac{a\sec^2 t}{a^2\tan^2 t\cdot a\sec t}\mathrm{d}t &= \frac{1}{a^2}\int \frac{\cos t}{\sin^2 t}\mathrm{d}t = \frac{1}{a^2}\int \csc t\cot t\mathrm{d}t = -\frac{1}{a^2\sin t}+C \\
&= -\frac{\sqrt{x^2+a^2}}{a^2x}+C.
\end{align*}
\item 方法二：记 $y=\sqrt{x^2+a^2}$。为了得到 $\dfrac{\mathrm{d}x}{x^2y}$，考虑反过来的比值 $\dfrac{y}{x}$：\(\displaystyle \mathrm{d}\left(\frac{y}{x}\right)=\frac{x\,\mathrm{d}y-y\,\mathrm{d}x}{x^2}\)。再代入 $\mathrm{d}y=\dfrac{x}{y}\mathrm{d}x$，便有 \(\displaystyle \mathrm{d}\left(\frac{y}{x}\right)=\frac{x^2-y^2}{x^2y}\,\mathrm{d}x=-\frac{a^2}{x^2y}\,\mathrm{d}x\)。因此 \(\displaystyle \int \frac{\mathrm{d}x}{x^2y}=-\frac{1}{a^2}\int \mathrm{d}\left(\frac{y}{x}\right)=-\frac{y}{a^2x}+C=-\frac{\sqrt{x^2+a^2}}{a^2x}+C\)。
\end{enumerate}
\end{solution}

\begin{exercise}[含有 $\sqrt{x^2+a^2}$ 的积分 - 39]
$\displaystyle \int \sqrt{x^2+a^2}\mathrm{d}x$
\end{exercise}
\begin{solution}
\textbf{思路：}这是一道“回到基石公式”的整合题。无论采用分部积分还是双元法，最后都要落回第 31 题的对数母式。
\begin{enumerate}
\item 方法一：分部积分法：
\begin{align*}
\int \sqrt{x^2+a^2}\mathrm{d}x &= x\sqrt{x^2+a^2} - \int x\cdot\frac{x}{\sqrt{x^2+a^2}}\mathrm{d}x = x\sqrt{x^2+a^2} - \int \frac{x^2+a^2-a^2}{\sqrt{x^2+a^2}}\mathrm{d}x \\
&= x\sqrt{x^2+a^2} - \int \sqrt{x^2+a^2}\mathrm{d}x + a^2\int \frac{\mathrm{d}x}{\sqrt{x^2+a^2}}.
\end{align*}
移项得
移项后得到
\begin{align*}
2\int \sqrt{x^2+a^2}\,\mathrm{d}x
&=x\sqrt{x^2+a^2}+a^2\ln(x+\sqrt{x^2+a^2})+C,\\
\int \sqrt{x^2+a^2}\,\mathrm{d}x
&=\frac{x}{2}\sqrt{x^2+a^2}
+\frac{a^2}{2}\ln(x+\sqrt{x^2+a^2})+C.
\end{align*}
\item 方法二：记 $y=\sqrt{x^2+a^2}$。为了把 $\int y\,\mathrm{d}x$ 改写成熟悉的对象，先把它拆成
\[
2y\,\mathrm{d}x=(y\,\mathrm{d}x+x\,\mathrm{d}y)+(y\,\mathrm{d}x-x\,\mathrm{d}y).
\]
其中第一组是 $\mathrm{d}(xy)$，第二组则可借助 $x\,\mathrm{d}y=\dfrac{x^2}{y}\mathrm{d}x$ 化简，因此
\begin{align*}
\int y\mathrm{d}x
&= \frac{1}{2}\int (y\mathrm{d}x+x\mathrm{d}y+y\mathrm{d}x-x\mathrm{d}y)\\
&= \frac{1}{2}xy + \frac{1}{2}\int \frac{y^2-x^2}{y}\mathrm{d}x\\
&= \frac{1}{2}xy + \frac{a^2}{2}\int \frac{\mathrm{d}x}{y}\\
&= \frac{1}{2}x\sqrt{x^2+a^2}+\frac{a^2}{2}\ln(x+\sqrt{x^2+a^2})+C.
\end{align*}
\end{enumerate}
这个例子也提示：以后遇到 $\int (x^2+a^2)^{m+\frac12}\mathrm{d}x$，往往都需要“分部积分 + 降回基石公式”或“先设 $y=\sqrt{x^2+a^2}$ 再建递推”这两条脉络。
\end{solution}

\begin{exercise}[含有 $\sqrt{x^2+a^2}$ 的积分 - 40]
$\displaystyle \int \sqrt{(x^2+a^2)^3}\mathrm{d}x$
\end{exercise}
\begin{solution}
\textbf{思路：}把 $y=\sqrt{x^2+a^2}$ 记成统一符号后，原积分就是 $\int y^3\,\mathrm{d}x$。它比第 39 题只高了一个幂次，因此最有效的办法是建立递推，把它降回 $\int y\,\mathrm{d}x$。
\begin{enumerate}
\item 方法一：分部积分法。记 \(\displaystyle I_2=\int (x^2+a^2)^{\frac{3}{2}}\mathrm{d}x,\qquad y=\sqrt{x^2+a^2}\)。
取 $u=y^3,\ \mathrm{d}v=\mathrm{d}x$，则 \(\displaystyle I_2=x y^3-3\int x^2y\,\mathrm{d}x\)。
再用 \(\displaystyle x^2=y^2-a^2\) 化简：
\[
I_2=x y^3-3\int (y^3-a^2y)\,\mathrm{d}x=x y^3-3I_2+3a^2\int y\,\mathrm{d}x,\qquad
4I_2=x y^3+3a^2\int y\,\mathrm{d}x.
\]
代入第 39 题结果 \(\displaystyle \int y\,\mathrm{d}x=\frac{1}{2}xy+\frac{a^2}{2}\ln(x+y)+C\)，得到
\[
I_2=\frac{1}{4}x y^3+\frac{3a^2}{8}xy+\frac{3a^4}{8}\ln(x+y)+C
=\frac{x}{8}(2x^2+5a^2)\sqrt{x^2+a^2}+\frac{3a^4}{8}\ln(x+\sqrt{x^2+a^2})+C.
\]
\item 方法二：直接在双元框架内建立递推。设 \(\displaystyle I_1=\int y\,\mathrm{d}x,\qquad I_2=\int y^3\,\mathrm{d}x\)。
由 $y\,\mathrm{d}x=\mathrm{d}(xy)-x\,\mathrm{d}y$ 可得
\[
I_2=\int y^2(y\,\mathrm{d}x)=\int y^2\bigl(\mathrm{d}(xy)-x\,\mathrm{d}y\bigr)
=x y^3-3\int x y^2\,\mathrm{d}y.
\]
而
\begin{align*}
x\,\mathrm{d}y&=\frac{x^2}{y}\,\mathrm{d}x=\frac{y^2-a^2}{y}\,\mathrm{d}x,\\
\int x y^2\,\mathrm{d}y&=\int (y^2-a^2)y\,\mathrm{d}x=I_2-a^2I_1,\\
I_2&=x y^3-3(I_2-a^2I_1),\\
4I_2&=x y^3+3a^2I_1.
\end{align*}
再代入 $I_1=\dfrac12xy+\dfrac{a^2}{2}\ln(x+y)$，得到与方法一相同的结果。
\end{enumerate}
\end{solution}

\begin{exercise}[含有 $\sqrt{x^2+a^2}$ 的积分 - 41]
$\displaystyle \int x\sqrt{x^2+a^2}\mathrm{d}x$
\end{exercise}
\begin{solution}
\textbf{思路：}这四题给出的都是平面闭曲线上的第二类曲线积分，而且区域内部没有奇点，所以先看 Green 公式比先参数化边界更稳。真正决定快慢的，是 $\frac{\partial Q}{\partial x}-\frac{\partial P}{\partial y}$ 能否直接为零，或者化成面积分后能否再借对称性消掉。
\begin{enumerate}
\item 方法一：凑微分法：\(\displaystyle \int x\sqrt{x^2+a^2}\mathrm{d}x=\frac{1}{2}\int (x^2+a^2)^{\frac{1}{2}}\mathrm{d}(x^2+a^2)=\frac{1}{3}(x^2+a^2)^{\frac{3}{2}}+C\)。
\item 方法二：记 $y=\sqrt{x^2+a^2}$，利用 $x\mathrm{d}x=y\mathrm{d}y$ 直接转化：\(\displaystyle \int xy\mathrm{d}x=\int y(y\mathrm{d}y)=\int y^2\mathrm{d}y=\frac{1}{3}y^3+C=\frac{1}{3}\sqrt{(x^2+a^2)^3}+C\)。
\end{enumerate}
\end{solution}

\begin{exercise}[含有 $\sqrt{x^2+a^2}$ 的积分 - 42]
$\displaystyle \int x^2\sqrt{x^2+a^2}\mathrm{d}x$
\end{exercise}
\begin{solution}
\textbf{思路：}这一题最怕把问题越做越重。更稳妥的路线是先把它降到“第 40 题减去 $a^2$ 倍第 39 题”，这样层次最清楚。
\begin{enumerate}
\item 方法一：分部积分法。令 \(\displaystyle u=x,\qquad \mathrm{d}v=x\sqrt{x^2+a^2}\,\mathrm{d}x\)。由第 41 题，\(\displaystyle v=\int x\sqrt{x^2+a^2}\,\mathrm{d}x=\frac13(x^2+a^2)^{3/2}\)。
于是
\[
\int x^2\sqrt{x^2+a^2}\mathrm{d}x = x\left(\frac{1}{3}(x^2+a^2)^{\frac{3}{2}}\right) - \frac{1}{3}\int (x^2+a^2)^{\frac{3}{2}}\mathrm{d}x.
\]
再代入第 40 题结果：
\[
\int (x^2+a^2)^{3/2}\mathrm{d}x=\frac{x}{8}(2x^2+5a^2)\sqrt{x^2+a^2}+\frac{3a^4}{8}\ln(x+\sqrt{x^2+a^2})+C,
\]
便得
\[
\int x^2\sqrt{x^2+a^2}\,\mathrm{d}x
=\frac{x}{8}(2x^2+a^2)\sqrt{x^2+a^2}-\frac{a^4}{8}\ln(x+\sqrt{x^2+a^2})+C.
\]
\item 方法二：记 $y=\sqrt{x^2+a^2}$，利用 $x^2 = y^2-a^2$ 对代数多项式进行结构展开：
\begin{align*}
\int x^2y\mathrm{d}x &= \int (y^2-a^2)y\mathrm{d}x = \int y^3\mathrm{d}x - a^2\int y\mathrm{d}x = I_2 - a^2I_1 \\
&= \left(\frac{1}{4}xy^3+\frac{3a^2}{4}I_1\right) - a^2I_1 = \frac{1}{4}xy^3 - \frac{a^2}{4}\left(\frac{1}{2}xy+\frac{a^2}{2}\ln(x+y)\right) \\
&= \frac{xy}{8}(2y^2-a^2) - \frac{a^4}{8}\ln(x+y) = \frac{x}{8}(2x^2+a^2)\sqrt{x^2+a^2}-\frac{a^4}{8}\ln(x+\sqrt{x^2+a^2})+C.
\end{align*}
\end{enumerate}
这个例子也提示：只要分子里出现 $x^2y$ 这类“偶次幂乘根式”的结构，都可以优先尝试“高一阶母式减低一阶母式”的组合，而不要从头硬做。
\end{solution}

\begin{exercise}[含有 $\sqrt{x^2+a^2}$ 的积分 - 43]
$\displaystyle \int \frac{\sqrt{x^2+a^2}}{x}\mathrm{d}x$
\end{exercise}
\begin{solution}
\textbf{思路：}题目问的是面积，因此最自然的入口不是把曲线分段积分到底，而是先翻译成 Green 面积公式；这样边界一旦参数化，后面就只剩代入与化简。
\begin{enumerate}
\item 方法一：令 $y = \sqrt{x^2+a^2} \implies x^2=y^2-a^2, x\mathrm{d}x=y\mathrm{d}y$：
\begin{align*}
\int \frac{\sqrt{x^2+a^2}}{x^2}x\mathrm{d}x &= \int \frac{y^2}{y^2-a^2}\mathrm{d}y = \int \left(1+\frac{a^2}{y^2-a^2}\right)\mathrm{d}y \\
&= y + \frac{a}{2}\ln\left|\frac{y-a}{y+a}\right|+C = \sqrt{x^2+a^2}+a\ln\frac{\sqrt{x^2+a^2}-a}{|x|}+C.
\end{align*}
\item 方法二：记 $y=\sqrt{x^2+a^2}$，直接对被积函数分子进行底座展开，规避所有换元：
\begin{align*}
\int \frac{y}{x}\mathrm{d}x &= \int \frac{y^2}{xy}\mathrm{d}x = \int \frac{x^2+a^2}{xy}\mathrm{d}x = \int \frac{x\mathrm{d}x}{y} + a^2\int \frac{\mathrm{d}x}{xy} \\
&= y + a^2\left(\frac{1}{a}\ln\frac{y-a}{|x|}\right)+C = \sqrt{x^2+a^2}+a\ln\frac{\sqrt{x^2+a^2}-a}{|x|}+C.
\end{align*}
\end{enumerate}
\end{solution}

\begin{exercise}[含有 $\sqrt{x^2+a^2}$ 的积分 - 44]
$\displaystyle \int \frac{\sqrt{x^2+a^2}}{x^2}\mathrm{d}x$
\end{exercise}
\begin{solution}
\begin{enumerate}
\item 方法一：分部积分法：
\begin{align*}
\int \sqrt{x^2+a^2}\mathrm{d}\left(-\frac{1}{x}\right) &= -\frac{\sqrt{x^2+a^2}}{x} + \int \frac{1}{x}\cdot\frac{x}{\sqrt{x^2+a^2}}\mathrm{d}x \\
&= -\frac{\sqrt{x^2+a^2}}{x} + \int \frac{\mathrm{d}x}{\sqrt{x^2+a^2}} = -\frac{\sqrt{x^2+a^2}}{x}+\ln(x+\sqrt{x^2+a^2})+C.
\end{align*}
\item 方法二：记 $y=\sqrt{x^2+a^2}$，配合简单的分部积分，将复杂根式直接消解于对数基石之中：
\begin{align*}
\int \frac{y}{x^2}\mathrm{d}x &= \int y\mathrm{d}\left(-\frac{1}{x}\right) = -\frac{y}{x} + \int \frac{1}{x}\mathrm{d}y = -\frac{y}{x} + \int \frac{1}{x}\left(\frac{x\mathrm{d}x}{y}\right) \\
&= -\frac{y}{x} + \int \frac{\mathrm{d}x}{y} = -\frac{\sqrt{x^2+a^2}}{x}+\ln(x+\sqrt{x^2+a^2})+C.
\end{align*}
\end{enumerate}
\end{solution}

\subsection{第八节习题：\texorpdfstring{含有 $\sqrt{x^2-a^2}\ (a>0)$ 的积分}{根式积分（减号）}}
\begin{exercise}[含有 $\sqrt{x^2-a^2}$ 的积分 - 45]
$\displaystyle \int \frac{\mathrm{d}x}{\sqrt{x^2-a^2}}$
\end{exercise}
\begin{solution}
\begin{enumerate}
\item 方法一：令 $x = a\sec t$，$\mathrm{d}x = a\sec t\tan t\mathrm{d}t$，则：
\begin{align*}
\int \frac{\mathrm{d}x}{\sqrt{x^2-a^2}} &= \int \frac{a\sec t\tan t}{a\tan t}\mathrm{d}t = \int \sec t\mathrm{d}t = \ln|\sec t+\tan t|+C \\
&= \ln\left|\frac{x}{a}+\frac{\sqrt{x^2-a^2}}{a}\right|+C = \ln|x+\sqrt{x^2-a^2}|+C_1.
\end{align*}
\item 方法二：记 $y=\sqrt{x^2-a^2}$，则 $x^2-y^2=a^2 \implies x\mathrm{d}x=y\mathrm{d}y$。提取虚圆对数微元：
\begin{align*}
\int \frac{\mathrm{d}x}{y} &= \int \frac{\mathrm{d}x+\frac{x}{y}\mathrm{d}x}{x+y} = \int \frac{\mathrm{d}x+\mathrm{d}y}{x+y} = \int \mathrm{d}\ln|x+y| \\
&= \ln|x+y|+C = \ln|x+\sqrt{x^2-a^2}|+C.
\end{align*}
\end{enumerate}
\end{solution}

\begin{exercise}[含有 $\sqrt{x^2-a^2}$ 的积分 - 46]
$\displaystyle \int \frac{\mathrm{d}x}{\sqrt{(x^2-a^2)^3}}$
\end{exercise}
\begin{solution}
\begin{enumerate}
\item 方法一：令 $x = a\sec t$，$\mathrm{d}x = a\sec t\tan t\mathrm{d}t$，则：
\begin{align*}
\int \frac{\mathrm{d}x}{(x^2-a^2)^{\frac{3}{2}}} &= \int \frac{a\sec t\tan t}{a^3\tan^3 t}\mathrm{d}t = \frac{1}{a^2}\int \frac{\cos t}{\sin^2 t}\mathrm{d}t = -\frac{1}{a^2\sin t}+C \\
&= -\frac{x}{a^2\sqrt{x^2-a^2}}+C.
\end{align*}
\item 方法二：记 $y=\sqrt{x^2-a^2}$，利用双元比值微元 \(\displaystyle \mathrm{d}\left(\frac{x}{y}\right) = \frac{y\mathrm{d}x-x\mathrm{d}y}{y^2} = \frac{y^2-x^2}{y^3}\mathrm{d}x = -\frac{a^2}{y^3}\mathrm{d}x\)，得 \(\displaystyle \int \frac{\mathrm{d}x}{y^3}=-\frac{1}{a^2}\int \mathrm{d}\left(\frac{x}{y}\right)=-\frac{x}{a^2y}+C=-\frac{x}{a^2\sqrt{x^2-a^2}}+C\)。
\end{enumerate}
\end{solution}

\begin{exercise}[含有 $\sqrt{x^2-a^2}$ 的积分 - 47]
$\displaystyle \int \frac{x}{\sqrt{x^2-a^2}}\mathrm{d}x$
\end{exercise}
\begin{solution}
\begin{enumerate}
\item 方法一：凑微分法：\(\displaystyle \int \frac{x}{\sqrt{x^2-a^2}}\mathrm{d}x=\frac{1}{2}\int (x^2-a^2)^{-\frac{1}{2}}\mathrm{d}(x^2-a^2)=\sqrt{x^2-a^2}+C\)。
\item 方法二：记 $y=\sqrt{x^2-a^2}$，利用底层微分关系 $x\mathrm{d}x=y\mathrm{d}y$：\(\displaystyle \int \frac{x}{y}\mathrm{d}x=\int \frac{y\mathrm{d}y}{y}=\int \mathrm{d}y=y+C=\sqrt{x^2-a^2}+C\)。
\end{enumerate}
\end{solution}

\begin{exercise}[含有 $\sqrt{x^2-a^2}$ 的积分 - 48]
$\displaystyle \int \frac{x}{\sqrt{(x^2-a^2)^3}}\mathrm{d}x$
\end{exercise}
\begin{solution}
\begin{enumerate}
\item 方法一：凑微分法：\(\displaystyle \int \frac{x}{(x^2-a^2)^{\frac{3}{2}}}\mathrm{d}x
= \frac{1}{2}\int (x^2-a^2)^{-\frac{3}{2}}\mathrm{d}(x^2-a^2)
=-\frac{1}{\sqrt{x^2-a^2}}+C\)。
\item 方法二：记 $y=\sqrt{x^2-a^2}$，利用 $x\mathrm{d}x=y\mathrm{d}y$，有 \(\displaystyle \int \frac{x}{y^3}\mathrm{d}x=\int \frac{y\mathrm{d}y}{y^3}=\int y^{-2}\mathrm{d}y=-\frac{1}{y}+C=-\frac{1}{\sqrt{x^2-a^2}}+C\)。
\end{enumerate}
\end{solution}

\begin{exercise}[含有 $\sqrt{x^2-a^2}$ 的积分 - 49]
$\displaystyle \int \frac{x^2}{\sqrt{x^2-a^2}}\mathrm{d}x$
\end{exercise}
\begin{solution}
\begin{enumerate}
\item 方法一：分部积分法：
\begin{align*}
\int \frac{x^2}{\sqrt{x^2-a^2}}\mathrm{d}x &= \int x\mathrm{d}\left(\sqrt{x^2-a^2}\right) = x\sqrt{x^2-a^2} - \int \sqrt{x^2-a^2}\mathrm{d}x \\
&= x\sqrt{x^2-a^2} - \int \frac{x^2-a^2}{\sqrt{x^2-a^2}}\mathrm{d}x \\
&= x\sqrt{x^2-a^2} - \int \frac{x^2}{\sqrt{x^2-a^2}}\mathrm{d}x + a^2\int \frac{\mathrm{d}x}{\sqrt{x^2-a^2}} \\
\implies 2\int \frac{x^2}{\sqrt{x^2-a^2}}\mathrm{d}x &= x\sqrt{x^2-a^2} + a^2\ln|x+\sqrt{x^2-a^2}| .
\end{align*}
\item 方法二：记 $y=\sqrt{x^2-a^2}$，利用底座代换式 $x^2 = y^2+a^2$ 进行多项式拆项：
\begin{align*}
\int \frac{x^2}{y}\mathrm{d}x &= \int \frac{y^2+a^2}{y}\mathrm{d}x = \int y\mathrm{d}x + a^2\int \frac{\mathrm{d}x}{y} \\
&= \left(\frac{1}{2}xy - \frac{a^2}{2}\ln|x+y|\right) + a^2\ln|x+y| + C \\
&= \frac{1}{2}xy + \frac{a^2}{2}\ln|x+y|+C = \frac{x}{2}\sqrt{x^2-a^2}+\frac{a^2}{2}\ln|x+\sqrt{x^2-a^2}|+C.
\end{align*}
\end{enumerate}
\end{solution}

\begin{exercise}[含有 $\sqrt{x^2-a^2}$ 的积分 - 50]
$\displaystyle \int \frac{x^2}{\sqrt{(x^2-a^2)^3}}\mathrm{d}x$
\end{exercise}
\begin{solution}
\begin{enumerate}
\item 方法一：分部积分法：
\begin{align*}
\int x\cdot\frac{x}{(x^2-a^2)^{\frac{3}{2}}}\mathrm{d}x &= \int x\mathrm{d}\left(-\frac{1}{\sqrt{x^2-a^2}}\right) = -\frac{x}{\sqrt{x^2-a^2}} + \int \frac{\mathrm{d}x}{\sqrt{x^2-a^2}} \\
&= -\frac{x}{\sqrt{x^2-a^2}} + \ln|x+\sqrt{x^2-a^2}|+C.
\end{align*}
\item 方法二：记 $y=\sqrt{x^2-a^2}$，利用恒等式 $x^2 = y^2+a^2$ 进行拆分：
\begin{align*}
\int \frac{x^2}{y^3}\mathrm{d}x &= \int \frac{y^2+a^2}{y^3}\mathrm{d}x = \int \frac{\mathrm{d}x}{y} + a^2\int \frac{\mathrm{d}x}{y^3} \\
&= \ln|x+y| + a^2\left(-\frac{x}{a^2y}\right)+C = -\frac{x}{\sqrt{x^2-a^2}} + \ln|x+\sqrt{x^2-a^2}|+C.
\end{align*}
\end{enumerate}
\end{solution}

\begin{exercise}[含有 $\sqrt{x^2-a^2}$ 的积分 - 51]
$\displaystyle \int \frac{\mathrm{d}x}{x\sqrt{x^2-a^2}}$
\end{exercise}
\begin{solution}
\begin{enumerate}
\item 方法一：令 $x=a\sec t$，则 $\mathrm{d}x=a\sec t\tan t\mathrm{d}t$，于是 \(\displaystyle \int \frac{a\sec t\tan t}{a\sec t\cdot a\tan t}\mathrm{d}t = \frac{1}{a}\int \mathrm{d}t = \frac{1}{a}t+C = \frac{1}{a}\arccos\frac{a}{|x|}+C\)。
\item 方法二：记 $y=\sqrt{x^2-a^2}$。利用双元微分对称性 $x\mathrm{d}x=y\mathrm{d}y \implies \frac{\mathrm{d}x}{y} = \frac{\mathrm{d}y}{x}$：
\begin{align*}
\int \frac{\mathrm{d}x}{xy} &= \int \frac{1}{x}\left(\frac{\mathrm{d}y}{x}\right) = \int \frac{\mathrm{d}y}{x^2} = \int \frac{\mathrm{d}y}{y^2+a^2} \\
&= \frac{1}{a}\arctan\frac{y}{a}+C = \frac{1}{a}\arccos\frac{a}{|x|}+C.
\end{align*}
\end{enumerate}
\end{solution}

\begin{exercise}[含有 $\sqrt{x^2-a^2}$ 的积分 - 52]
$\displaystyle \int \frac{\mathrm{d}x}{x^2\sqrt{x^2-a^2}}$
\end{exercise}
\begin{solution}
\begin{enumerate}
\item 方法一：令 $x=a\sec t$，则 $\mathrm{d}x=a\sec t\tan t\mathrm{d}t$，于是 \(\displaystyle \int \frac{a\sec t\tan t}{a^2\sec^2 t\cdot a\tan t}\mathrm{d}t = \frac{1}{a^2}\int \cos t\mathrm{d}t = \frac{1}{a^2}\sin t+C = \frac{\sqrt{x^2-a^2}}{a^2x}+C\)。
\item 方法二：记 $y=\sqrt{x^2-a^2}$，利用双元比值微元 \(\displaystyle \mathrm{d}\left(\frac{y}{x}\right) = \frac{x\mathrm{d}y-y\mathrm{d}x}{x^2} = \frac{x^2-y^2}{x^2y}\mathrm{d}x = \frac{a^2}{x^2y}\mathrm{d}x\)，得 \(\displaystyle \int \frac{\mathrm{d}x}{x^2y}=\frac{1}{a^2}\int \mathrm{d}\left(\frac{y}{x}\right)=\frac{y}{a^2x}+C=\frac{\sqrt{x^2-a^2}}{a^2x}+C\)。
\end{enumerate}
\end{solution}

\begin{exercise}[含有 $\sqrt{x^2-a^2}$ 的积分 - 53]
$\displaystyle \int \sqrt{x^2-a^2}\mathrm{d}x$
\end{exercise}
\begin{solution}
\begin{enumerate}
\item 方法一：分部积分法：
\begin{align*}
\int \sqrt{x^2-a^2}\mathrm{d}x &= x\sqrt{x^2-a^2} - \int x\cdot\frac{x}{\sqrt{x^2-a^2}}\mathrm{d}x = x\sqrt{x^2-a^2} - \int \frac{x^2-a^2+a^2}{\sqrt{x^2-a^2}}\mathrm{d}x \\
&= x\sqrt{x^2-a^2} - \int \sqrt{x^2-a^2}\mathrm{d}x - a^2\int \frac{\mathrm{d}x}{\sqrt{x^2-a^2}}.
\end{align*}
由此可得所求结果。
\item 方法二：记 $y=\sqrt{x^2-a^2}$，由面积微元 $y\mathrm{d}x - x\mathrm{d}y = (y^2-x^2)\frac{\mathrm{d}x}{y} = -a^2\frac{\mathrm{d}x}{y}$ 直接降阶：
\begin{align*}
\int y\mathrm{d}x &= \frac{1}{2}\int (y\mathrm{d}x+x\mathrm{d}y+y\mathrm{d}x-x\mathrm{d}y) = \frac{1}{2}xy - \frac{a^2}{2}\int \frac{\mathrm{d}x}{y} \\
&= \frac{1}{2}x\sqrt{x^2-a^2}-\frac{a^2}{2}\ln|x+\sqrt{x^2-a^2}|+C.
\end{align*}
\end{enumerate}
\end{solution}

\begin{exercise}[含有 $\sqrt{x^2-a^2}$ 的积分 - 54]
$\displaystyle \int \sqrt{(x^2-a^2)^3}\mathrm{d}x$
\end{exercise}
\begin{solution}
\begin{enumerate}
\item 方法一：分部积分法 $u=(x^2-a^2)^{\frac{3}{2}}, \mathrm{d}v=\mathrm{d}x$：
\[
\int (x^2-a^2)^{\frac{3}{2}}\mathrm{d}x = x(x^2-a^2)^{\frac{3}{2}} - 3\int x^2\sqrt{x^2-a^2}\mathrm{d}x \dots
\]
\item 方法二：记 $y=\sqrt{x^2-a^2}$，结合恒等式 $x^2=y^2+a^2$ 递推：
\begin{align*}
I_2 &= \int y^3\mathrm{d}x = \int y^2(y\mathrm{d}x) = \int y^2(\mathrm{d}(xy)-x\mathrm{d}y) = xy^3 - \int 2y(xy\mathrm{d}y) \\
&= xy^3 - 2\int y(x^2\mathrm{d}x) = xy^3 - 2\int y(y^2+a^2)\mathrm{d}x = xy^3 - 2I_2 - 2a^2 I_1 \\
\implies I_2 &= \frac{1}{3}xy^3 - \frac{2a^2}{3}\left(\frac{1}{2}xy-\frac{a^2}{2}\ln|x+y|\right) = \frac{x}{6}y(2y^2-2a^2) + \frac{a^4}{3}\ln|x+y| \\
&= \frac{x}{8}(2x^2-5a^2)\sqrt{x^2-a^2} + \frac{3a^4}{8}\ln|x+\sqrt{x^2-a^2}|+C.
\end{align*}
\end{enumerate}
\end{solution}

\begin{exercise}[含有 $\sqrt{x^2-a^2}$ 的积分 - 55]
$\displaystyle \int x\sqrt{x^2-a^2}\mathrm{d}x$
\end{exercise}
\begin{solution}
\begin{enumerate}
\item 方法一：凑微分法：\(\displaystyle \int x\sqrt{x^2-a^2}\mathrm{d}x=\frac{1}{2}\int (x^2-a^2)^{\frac{1}{2}}\mathrm{d}(x^2-a^2)=\frac{1}{3}(x^2-a^2)^{\frac{3}{2}}+C\)。
\item 方法二：记 $y=\sqrt{x^2-a^2}$，利用 $x\mathrm{d}x=y\mathrm{d}y$：\(\displaystyle \int xy\mathrm{d}x=\int y(y\mathrm{d}y)=\int y^2\mathrm{d}y=\frac{1}{3}y^3+C=\frac{1}{3}\sqrt{(x^2-a^2)^3}+C\)。
\end{enumerate}
\end{solution}

\begin{exercise}[含有 $\sqrt{x^2-a^2}$ 的积分 - 56]
$\displaystyle \int x^2\sqrt{x^2-a^2}\mathrm{d}x$
\end{exercise}
\begin{solution}
\begin{enumerate}
\item 方法一：分部积分法：\(\displaystyle \int x^2\sqrt{x^2-a^2}\mathrm{d}x = x\left(\frac{1}{3}(x^2-a^2)^{\frac{3}{2}}\right) - \frac{1}{3}\int (x^2-a^2)^{\frac{3}{2}}\mathrm{d}x\)。由此可得所求结果。
\item 方法二：记 $y=\sqrt{x^2-a^2}$，利用 $x^2 = y^2+a^2$ 展开拆解：
\begin{align*}
\int x^2y\mathrm{d}x &= \int (y^2+a^2)y\mathrm{d}x = \int y^3\mathrm{d}x + a^2\int y\mathrm{d}x = I_2 + a^2I_1 \\
&= \left(\frac{1}{4}xy^3-\frac{3a^2}{4}I_1\right) + a^2I_1 = \frac{1}{4}xy^3 + \frac{a^2}{4}\left(\frac{1}{2}xy-\frac{a^2}{2}\ln|x+y|\right) \\
&= \frac{x}{8}(2x^2-a^2)\sqrt{x^2-a^2}-\frac{a^4}{8}\ln|x+\sqrt{x^2-a^2}|+C.
\end{align*}
\end{enumerate}
\end{solution}

\begin{exercise}[含有 $\sqrt{x^2-a^2}$ 的积分 - 57]
$\displaystyle \int \frac{\sqrt{x^2-a^2}}{x}\mathrm{d}x$
\end{exercise}
\begin{solution}
\begin{enumerate}
\item 方法一：令 $y = \sqrt{x^2-a^2} \implies x^2=y^2+a^2, x\mathrm{d}x=y\mathrm{d}y$：
\begin{align*}
\int \frac{\sqrt{x^2-a^2}}{x^2}x\mathrm{d}x &= \int \frac{y^2}{y^2+a^2}\mathrm{d}y = \int \left(1-\frac{a^2}{y^2+a^2}\right)\mathrm{d}y \\
&= y - a\arctan\frac{y}{a}+C = \sqrt{x^2-a^2}-a\arccos\frac{a}{|x|}+C.
\end{align*}
\item 方法二：记 $y=\sqrt{x^2-a^2}$，代数重构分离基础微元：
\begin{align*}
\int \frac{y}{x}\mathrm{d}x &= \int \frac{y^2}{xy}\mathrm{d}x = \int \frac{x^2-a^2}{xy}\mathrm{d}x = \int \frac{x\mathrm{d}x}{y} - a^2\int \frac{\mathrm{d}x}{xy} \\
&= y - a^2\left(\frac{1}{a}\arccos\frac{a}{|x|}\right)+C = \sqrt{x^2-a^2}-a\arccos\frac{a}{|x|}+C.
\end{align*}
\end{enumerate}
\end{solution}

\begin{exercise}[含有 $\sqrt{x^2-a^2}$ 的积分 - 58]
$\displaystyle \int \frac{\sqrt{x^2-a^2}}{x^2}\mathrm{d}x$
\end{exercise}
\begin{solution}
\begin{enumerate}
\item 方法一：分部积分法：
\begin{align*}
\int \sqrt{x^2-a^2}\mathrm{d}\left(-\frac{1}{x}\right) &= -\frac{\sqrt{x^2-a^2}}{x} + \int \frac{1}{x}\cdot\frac{x}{\sqrt{x^2-a^2}}\mathrm{d}x \\
&= -\frac{\sqrt{x^2-a^2}}{x}+\ln|x+\sqrt{x^2-a^2}|+C.
\end{align*}
\item 方法二：记 $y=\sqrt{x^2-a^2}$，分部引入对数微元：
\begin{align*}
\int \frac{y}{x^2}\mathrm{d}x &= \int y\mathrm{d}\left(-\frac{1}{x}\right) = -\frac{y}{x} + \int \frac{1}{x}\mathrm{d}y = -\frac{y}{x} + \int \frac{1}{x}\left(\frac{x\mathrm{d}x}{y}\right) \\
&= -\frac{y}{x} + \int \frac{\mathrm{d}x}{y} = -\frac{\sqrt{x^2-a^2}}{x}+\ln|x+\sqrt{x^2-a^2}|+C.
\end{align*}
\end{enumerate}
\end{solution}

\subsection{第九节习题：\texorpdfstring{含有 $\sqrt{a^2-x^2}\ (a>0)$ 的积分}{根式积分（圆弧型）}}
\begin{exercise}[含有 $\sqrt{a^2-x^2}$ 的积分 - 59]
$\displaystyle \int \frac{\mathrm{d}x}{\sqrt{a^2-x^2}}$
\end{exercise}
\begin{solution}
\begin{enumerate}
\item 方法一：标准积分公式：
\(\displaystyle \int \frac{\mathrm{d}x}{\sqrt{a^2-x^2}} = \arcsin\frac{x}{a}+C\)。
\item 方法二：记 $y=\sqrt{a^2-x^2}$，此即为实圆底座的最基础微分关系：
\(\displaystyle \int \frac{\mathrm{d}x}{y} = \arcsin\frac{x}{a}+C\)。
\end{enumerate}
\end{solution}

\begin{exercise}[含有 $\sqrt{a^2-x^2}$ 的积分 - 60]
$\displaystyle \int \frac{\mathrm{d}x}{\sqrt{(a^2-x^2)^3}}$
\end{exercise}
\begin{solution}
\begin{enumerate}
\item 方法一：令 $x=a\sin t$，则 $\mathrm{d}x=a\cos t\mathrm{d}t$，于是 \(\displaystyle \int \frac{a\cos t}{a^3\cos^3 t}\mathrm{d}t = \frac{1}{a^2}\int \sec^2 t\mathrm{d}t = \frac{1}{a^2}\tan t+C = \frac{x}{a^2\sqrt{a^2-x^2}}+C\)。
\item 方法二：记 $y=\sqrt{a^2-x^2}$，由于 $x^2+y^2=a^2$，提取双元比值微元 \(\displaystyle \mathrm{d}\left(\frac{x}{y}\right) = \frac{y\mathrm{d}x-x\mathrm{d}y}{y^2} = \frac{y^2+x^2}{y^3}\mathrm{d}x = \frac{a^2}{y^3}\mathrm{d}x\)，得 \(\displaystyle \int \frac{\mathrm{d}x}{y^3}=\frac{1}{a^2}\int \mathrm{d}\left(\frac{x}{y}\right)=\frac{x}{a^2y}+C=\frac{x}{a^2\sqrt{a^2-x^2}}+C\)。
\end{enumerate}
\end{solution}

\begin{exercise}[含有 $\sqrt{a^2-x^2}$ 的积分 - 61]
$\displaystyle \int \frac{x}{\sqrt{a^2-x^2}}\mathrm{d}x$
\end{exercise}
\begin{solution}
\begin{enumerate}
\item 方法一：凑微分法：\(\displaystyle \int \frac{x}{\sqrt{a^2-x^2}}\mathrm{d}x=-\frac{1}{2}\int (a^2-x^2)^{-\frac{1}{2}}\mathrm{d}(a^2-x^2)=-\sqrt{a^2-x^2}+C\)。
\item 方法二：记 $y=\sqrt{a^2-x^2}$，底座关系 $x^2+y^2=a^2 \implies x\mathrm{d}x = -y\mathrm{d}y$：\(\displaystyle \int \frac{x}{y}\mathrm{d}x=\int \frac{-y\mathrm{d}y}{y}=-\int \mathrm{d}y=-y+C=-\sqrt{a^2-x^2}+C\)。
\end{enumerate}
\end{solution}

\begin{exercise}[含有 $\sqrt{a^2-x^2}$ 的积分 - 62]
$\displaystyle \int \frac{x}{\sqrt{(a^2-x^2)^3}}\mathrm{d}x$
\end{exercise}
\begin{solution}
\begin{enumerate}
\item 方法一：凑微分法：\(\displaystyle \int \frac{x}{(a^2-x^2)^{\frac{3}{2}}}\mathrm{d}x
= -\frac{1}{2}\int (a^2-x^2)^{-\frac{3}{2}}\mathrm{d}(a^2-x^2)
=\frac{1}{\sqrt{a^2-x^2}}+C\)。
\item 方法二：记 $y=\sqrt{a^2-x^2}$，利用 $x\mathrm{d}x = -y\mathrm{d}y$，有 \(\displaystyle \int \frac{x}{y^3}\mathrm{d}x=\int \frac{-y\mathrm{d}y}{y^3}=\int -y^{-2}\mathrm{d}y=\frac{1}{y}+C=\frac{1}{\sqrt{a^2-x^2}}+C\)。
\end{enumerate}
\end{solution}

\begin{exercise}[含有 $\sqrt{a^2-x^2}$ 的积分 - 63]
$\displaystyle \int \frac{x^2}{\sqrt{a^2-x^2}}\mathrm{d}x$
\end{exercise}
\begin{solution}
\begin{enumerate}
\item 方法一：令 $x=a\sin t$，则 $\mathrm{d}x=a\cos t\mathrm{d}t$，于是
\begin{align*}
\int \frac{a^2\sin^2 t}{a\cos t}a\cos t\mathrm{d}t
&=a^2\int \sin^2 t\mathrm{d}t\\
&=-\frac{x}{2}\sqrt{a^2-x^2}
+\frac{a^2}{2}\arcsin\frac{x}{a}+C.
\end{align*}
\item 方法二：记 $y=\sqrt{a^2-x^2}$，利用底座关系 $x^2 = a^2-y^2$ 进行代数拆解：
\begin{align*}
\int \frac{x^2}{y}\mathrm{d}x
&= \int \frac{a^2-y^2}{y}\mathrm{d}x\\
&= a^2\int \frac{\mathrm{d}x}{y} - \int y\mathrm{d}x \\
&= a^2\arcsin\frac{x}{a}
- \left(\frac{1}{2}xy + \frac{a^2}{2}\arcsin\frac{x}{a}\right)+C\\
&= -\frac{x}{2}\sqrt{a^2-x^2}+\frac{a^2}{2}\arcsin\frac{x}{a}+C.
\end{align*}
\end{enumerate}
\end{solution}

\begin{exercise}[含有 $\sqrt{a^2-x^2}$ 的积分 - 64]
$\displaystyle \int \frac{x^2}{\sqrt{(a^2-x^2)^3}}\mathrm{d}x$
\end{exercise}
\begin{solution}
\begin{enumerate}
\item 方法一：分部积分法：
\begin{align*}
\int x\cdot\frac{x}{(a^2-x^2)^{\frac{3}{2}}}\mathrm{d}x &= \int x\mathrm{d}\left(\frac{1}{\sqrt{a^2-x^2}}\right) = \frac{x}{\sqrt{a^2-x^2}} - \int \frac{\mathrm{d}x}{\sqrt{a^2-x^2}} \\
&= \frac{x}{\sqrt{a^2-x^2}} - \arcsin\frac{x}{a}+C.
\end{align*}
\item 方法二：记 $y=\sqrt{a^2-x^2}$，利用 $x^2 = a^2-y^2$：
\begin{align*}
\int \frac{x^2}{y^3}\mathrm{d}x &= \int \frac{a^2-y^2}{y^3}\mathrm{d}x = a^2\int \frac{\mathrm{d}x}{y^3} - \int \frac{\mathrm{d}x}{y} \\
&= a^2\left(\frac{x}{a^2y}\right) - \arcsin\frac{x}{a}+C = \frac{x}{\sqrt{a^2-x^2}}-\arcsin\frac{x}{a}+C.
\end{align*}
\end{enumerate}
\end{solution}

\begin{exercise}[含有 $\sqrt{a^2-x^2}$ 的积分 - 65]
$\displaystyle \int \frac{\mathrm{d}x}{x\sqrt{a^2-x^2}}$
\end{exercise}
\begin{solution}
\begin{enumerate}
\item 方法一：令 $x=a\sin t$ 进行换元，化简为对数形式。
\item 方法二：记 $y=\sqrt{a^2-x^2}$。利用双元微分对称性 $x\mathrm{d}x=-y\mathrm{d}y \implies \frac{\mathrm{d}x}{y} = -\frac{\mathrm{d}y}{x}$：
\begin{align*}
\int \frac{\mathrm{d}x}{xy} &= \int \frac{1}{x}\left(-\frac{\mathrm{d}y}{x}\right) = \int \frac{-\mathrm{d}y}{x^2} = \int \frac{-\mathrm{d}y}{a^2-y^2} \\
&= -\frac{1}{2a}\ln\left|\frac{a+y}{a-y}\right|+C = \frac{1}{2a}\ln\frac{(a-y)^2}{a^2-y^2}+C = \frac{1}{a}\ln\frac{a-\sqrt{a^2-x^2}}{|x|}+C.
\end{align*}
\end{enumerate}
\end{solution}

\begin{exercise}[含有 $\sqrt{a^2-x^2}$ 的积分 - 66]
$\displaystyle \int \frac{\mathrm{d}x}{x^2\sqrt{a^2-x^2}}$
\end{exercise}
\begin{solution}
\begin{enumerate}
\item 方法一：令 $x=a\sin t$，则 $\mathrm{d}x=a\cos t\mathrm{d}t$，于是 \(\displaystyle \int \frac{a\cos t}{a^2\sin^2 t\cdot a\cos t}\mathrm{d}t = \frac{1}{a^2}\int \csc^2 t\mathrm{d}t = -\frac{1}{a^2}\cot t+C = -\frac{\sqrt{a^2-x^2}}{a^2x}+C\)。
\item 方法二：记 $y=\sqrt{a^2-x^2}$，利用双元比值微元 \(\displaystyle \mathrm{d}\left(\frac{y}{x}\right) = \frac{x\mathrm{d}y-y\mathrm{d}x}{x^2} = \frac{-x^2-y^2}{x^2y}\mathrm{d}x = -\frac{a^2}{x^2y}\mathrm{d}x\)，得 \(\displaystyle \int \frac{\mathrm{d}x}{x^2y}=-\frac{1}{a^2}\int \mathrm{d}\left(\frac{y}{x}\right)=-\frac{y}{a^2x}+C=-\frac{\sqrt{a^2-x^2}}{a^2x}+C\)。
\end{enumerate}
\end{solution}

\begin{exercise}[含有 $\sqrt{a^2-x^2}$ 的积分 - 67]
$\displaystyle \int \sqrt{a^2-x^2}\mathrm{d}x$
\end{exercise}
\begin{solution}
\begin{enumerate}
\item 方法一：令 $x=a\sin t$ 换元计算。
\item 方法二：记 $y=\sqrt{a^2-x^2}$，利用实圆面积微元 $y\mathrm{d}x - x\mathrm{d}y = (x^2+y^2)\frac{\mathrm{d}x}{y} = a^2\frac{\mathrm{d}x}{y}$ 直接降阶：
\begin{align*}
\int y\mathrm{d}x &= \frac{1}{2}\int (y\mathrm{d}x+x\mathrm{d}y+y\mathrm{d}x-x\mathrm{d}y) = \frac{1}{2}xy + \frac{a^2}{2}\int \frac{\mathrm{d}x}{y} \\
&= \frac{1}{2}x\sqrt{a^2-x^2}+\frac{a^2}{2}\arcsin\frac{x}{a}+C.
\end{align*}
\end{enumerate}
\end{solution}

\begin{exercise}[含有 $\sqrt{a^2-x^2}$ 的积分 - 68]
$\displaystyle \int \sqrt{(a^2-x^2)^3}\mathrm{d}x$
\end{exercise}
\begin{solution}
\begin{enumerate}
\item 方法一：分部积分法或三角换元。
\item 方法二：记 $y=\sqrt{a^2-x^2}$，结合恒等式 $x^2=a^2-y^2$ 递推：
\begin{align*}
I_2 &= \int y^3\mathrm{d}x = \int y^2(\mathrm{d}(xy)-x\mathrm{d}y) = xy^3 - \int 2y(xy\mathrm{d}y) = xy^3 + 2\int y(x^2\mathrm{d}x) \\
&= xy^3 + 2\int y(a^2-y^2)\mathrm{d}x = xy^3 - 2I_2 + 2a^2 I_1 \\
\implies I_2 &= \frac{1}{3}xy^3 + \frac{2a^2}{3}\left(\frac{1}{2}xy+\frac{a^2}{2}\arcsin\frac{x}{a}\right) = \frac{x}{6}y(2y^2+2a^2) + \frac{a^4}{3}\arcsin\frac{x}{a} \\
&= \frac{x}{8}(5a^2-2x^2)\sqrt{a^2-x^2} + \frac{3a^4}{8}\arcsin\frac{x}{a}+C.
\end{align*}
\end{enumerate}
\end{solution}

\begin{exercise}[含有 $\sqrt{a^2-x^2}$ 的积分 - 69]
$\displaystyle \int x\sqrt{a^2-x^2}\mathrm{d}x$
\end{exercise}
\begin{solution}
\begin{enumerate}
\item 方法一：凑微分法：\(\displaystyle \int x\sqrt{a^2-x^2}\mathrm{d}x = -\frac{1}{2}\int (a^2-x^2)^{\frac{1}{2}}\mathrm{d}(a^2-x^2) = -\frac{1}{3}(a^2-x^2)^{\frac{3}{2}}+C\)。
\item 方法二：记 $y=\sqrt{a^2-x^2}$，利用 $x\mathrm{d}x=-y\mathrm{d}y$，有 \(\displaystyle \int xy\mathrm{d}x=\int y(-y\mathrm{d}y)=-\int y^2\mathrm{d}y=-\frac{1}{3}y^3+C=-\frac{1}{3}\sqrt{(a^2-x^2)^3}+C\)。
\end{enumerate}
\end{solution}

\begin{exercise}[含有 $\sqrt{a^2-x^2}$ 的积分 - 70]
$\displaystyle \int x^2\sqrt{a^2-x^2}\mathrm{d}x$
\end{exercise}
\begin{solution}
\begin{enumerate}
\item 方法一：分部积分法或三角换元：\(\displaystyle \int x^2\sqrt{a^2-x^2}\mathrm{d}x = x\left(-\frac{1}{3}(a^2-x^2)^{\frac{3}{2}}\right) + \frac{1}{3}\int (a^2-x^2)^{\frac{3}{2}}\mathrm{d}x\)。由此可得所求结果。
\item 方法二：记 $y=\sqrt{a^2-x^2}$，利用 $x^2 = a^2-y^2$ 展开拆解：
\begin{align*}
\int x^2y\mathrm{d}x &= \int (a^2-y^2)y\mathrm{d}x = a^2\int y\mathrm{d}x - \int y^3\mathrm{d}x = a^2I_1 - I_2 \\
&= a^2I_1 - \left(\frac{1}{4}xy^3+\frac{3a^2}{4}I_1\right) = -\frac{1}{4}xy^3 + \frac{a^2}{4}\left(\frac{1}{2}xy+\frac{a^2}{2}\arcsin\frac{x}{a}\right) \\
&= \frac{x}{8}(2x^2-a^2)\sqrt{a^2-x^2}+\frac{a^4}{8}\arcsin\frac{x}{a}+C.
\end{align*}
\end{enumerate}
\end{solution}

\begin{exercise}[含有 $\sqrt{a^2-x^2}$ 的积分 - 71]
$\displaystyle \int \frac{\sqrt{a^2-x^2}}{x}\mathrm{d}x$
\end{exercise}
\begin{solution}
\begin{enumerate}
\item 方法一：令 $y = \sqrt{a^2-x^2} \implies x^2=a^2-y^2, x\mathrm{d}x=-y\mathrm{d}y$：
\begin{align*}
\int \frac{\sqrt{a^2-x^2}}{x^2}x\mathrm{d}x &= \int \frac{-y^2}{a^2-y^2}\mathrm{d}y = \int \left(1-\frac{a^2}{a^2-y^2}\right)\mathrm{d}y \\
&= y + \frac{a}{2}\ln\left|\frac{y-a}{y+a}\right|+C = \sqrt{a^2-x^2}+a\ln\frac{a-\sqrt{a^2-x^2}}{|x|}+C.
\end{align*}
\item 方法二：记 $y=\sqrt{a^2-x^2}$，代数重构分离对数微元：
\begin{align*}
\int \frac{y}{x}\mathrm{d}x &= \int \frac{y^2}{xy}\mathrm{d}x = \int \frac{a^2-x^2}{xy}\mathrm{d}x = a^2\int \frac{\mathrm{d}x}{xy} - \int \frac{x\mathrm{d}x}{y} \\
&= a^2\left(\frac{1}{a}\ln\frac{a-y}{|x|}\right) - (-y)+C = \sqrt{a^2-x^2}+a\ln\frac{a-\sqrt{a^2-x^2}}{|x|}+C.
\end{align*}
\end{enumerate}
\end{solution}

\begin{exercise}[含有 $\sqrt{a^2-x^2}$ 的积分 - 72]
$\displaystyle \int \frac{\sqrt{a^2-x^2}}{x^2}\mathrm{d}x$
\end{exercise}
\begin{solution}
\begin{enumerate}
\item 方法一：分部积分法：
\begin{align*}
\int \sqrt{a^2-x^2}\mathrm{d}\left(-\frac{1}{x}\right) &= -\frac{\sqrt{a^2-x^2}}{x} + \int \frac{1}{x}\cdot\frac{-x}{\sqrt{a^2-x^2}}\mathrm{d}x \\
&= -\frac{\sqrt{a^2-x^2}}{x}-\arcsin\frac{x}{a}+C.
\end{align*}
\item 方法二：记 $y=\sqrt{a^2-x^2}$，分部引入面积基础微元：
\begin{align*}
\int \frac{y}{x^2}\mathrm{d}x &= \int y\mathrm{d}\left(-\frac{1}{x}\right) = -\frac{y}{x} + \int \frac{1}{x}\mathrm{d}y = -\frac{y}{x} - \int \frac{1}{x}\left(\frac{x\mathrm{d}x}{y}\right) \\
&= -\frac{y}{x} - \int \frac{\mathrm{d}x}{y} = -\frac{\sqrt{a^2-x^2}}{x}-\arcsin\frac{x}{a}+C.
\end{align*}
\end{enumerate}
\end{solution}

\subsection{第十节习题：\texorpdfstring{含有 $\sqrt{ax^2+bx+c}\ (a>0)$ 的积分}{二次根式积分}}
\begin{exercise}[含有 $\sqrt{ax^2+bx+c}$ 的积分 - 73]
$\displaystyle \int \frac{\mathrm{d}x}{\sqrt{ax^2+bx+c}}$
\end{exercise}
\begin{solution}
\begin{enumerate}
\item 方法一：配方法：
\begin{align*}
\int \frac{\mathrm{d}x}{\sqrt{a\left(x+\frac{b}{2a}\right)^2+\frac{4ac-b^2}{4a}}} &= \frac{1}{\sqrt{a}}\ln\left|x+\frac{b}{2a}+\sqrt{\left(x+\frac{b}{2a}\right)^2+\frac{4ac-b^2}{4a^2}}\right|+C_1 \\
&= \frac{1}{\sqrt{a}}\ln\left|2ax+b+2\sqrt{a}\sqrt{ax^2+bx+c}\right|+C.
\end{align*}
\item 方法二：记 $y=\sqrt{ax^2+bx+c}$。令 $u=2ax+b, v=2\sqrt{a}y$，则 $v^2-u^2=4ac-b^2=\Delta$（虚圆底座）。
由 $v\mathrm{d}v=u\mathrm{d}u \implies \frac{\mathrm{d}u}{v}=\frac{\mathrm{d}v}{u}=\frac{\mathrm{d}(u+v)}{u+v}=\mathrm{d}\ln|u+v|$，且 $\mathrm{d}x=\frac{\mathrm{d}u}{2a}$：
\begin{align*}
\int \frac{\mathrm{d}x}{y} &= \int \frac{\mathrm{d}u/(2a)}{v/(2\sqrt{a})} = \frac{1}{\sqrt{a}}\int \frac{\mathrm{d}u}{v} = \frac{1}{\sqrt{a}}\ln|u+v|+C \\
&= \frac{1}{\sqrt{a}}\ln|2ax+b+2\sqrt{a}\sqrt{ax^2+bx+c}|+C.
\end{align*}
\end{enumerate}
\end{solution}

\begin{exercise}[含有 $\sqrt{ax^2+bx+c}$ 的积分 - 74]
$\displaystyle \int \sqrt{ax^2+bx+c}\mathrm{d}x$
\end{exercise}
\begin{solution}
\begin{enumerate}
\item 方法一：配方后利用分部积分法或三角换元法计算。
\item 方法二：同上，设 $u=2ax+b, v=2\sqrt{a}y$，利用虚圆面积微元 $\int v\mathrm{d}u = \frac{1}{2}uv + \frac{\Delta}{2}\int \frac{\mathrm{d}u}{v}$：
\begin{align*}
\int y\mathrm{d}x &= \int \frac{v}{2\sqrt{a}}\frac{\mathrm{d}u}{2a} = \frac{1}{4a\sqrt{a}}\int v\mathrm{d}u = \frac{1}{4a\sqrt{a}}\left(\frac{1}{2}uv + \frac{4ac-b^2}{2}\int \frac{\mathrm{d}u}{v}\right) \\
&= \frac{uv}{8a\sqrt{a}} + \frac{4ac-b^2}{8a\sqrt{a}}\ln|u+v|+C \\
&= \frac{2ax+b}{4a}\sqrt{ax^2+bx+c} + \frac{4ac-b^2}{8\sqrt{a^3}}\ln|2ax+b+2\sqrt{a}\sqrt{ax^2+bx+c}|+C.
\end{align*}
\end{enumerate}
\end{solution}

\begin{exercise}[含有 $\sqrt{ax^2+bx+c}$ 的积分 - 75]
$\displaystyle \int \frac{x}{\sqrt{ax^2+bx+c}}\mathrm{d}x$
\end{exercise}
\begin{solution}
\begin{enumerate}
\item 方法一：凑微分法：
\begin{align*}
\int \frac{x}{\sqrt{ax^2+bx+c}}\mathrm{d}x &= \frac{1}{2a}\int \frac{2ax+b-b}{\sqrt{ax^2+bx+c}}\mathrm{d}x = \frac{1}{a}\sqrt{ax^2+bx+c} - \frac{b}{2a}\int \frac{\mathrm{d}x}{\sqrt{ax^2+bx+c}}.
\end{align*}
由此可得所求结果。
\item 方法二：设 $u=2ax+b, v=2\sqrt{a}y$，代入 $x=\frac{u-b}{2a}$，在双元框架内代数拆解：
\begin{align*}
\int \frac{x}{y}\mathrm{d}x &= \int \frac{(u-b)/(2a)}{v/(2\sqrt{a})}\frac{\mathrm{d}u}{2a} = \frac{1}{2a\sqrt{a}}\int \frac{u-b}{v}\mathrm{d}u \\
&= \frac{1}{2a\sqrt{a}}\left(\int \frac{u\mathrm{d}u}{v} - b\int \frac{\mathrm{d}u}{v}\right) = \frac{1}{2a\sqrt{a}}\left(\int \mathrm{d}v - b\int \frac{\mathrm{d}u}{v}\right) \\
&= \frac{v}{2a\sqrt{a}} - \frac{b}{2a\sqrt{a}}\ln|u+v|+C = \frac{1}{a}y - \frac{b}{2\sqrt{a^3}}\ln|2ax+b+2\sqrt{a}y|+C.
\end{align*}
\end{enumerate}
\end{solution}

\begin{exercise}[含有 $\sqrt{ax^2+bx+c}$ 的积分 - 76]
$\displaystyle \int \frac{\mathrm{d}x}{\sqrt{c+bx-ax^2}}$
\end{exercise}
\begin{solution}
\begin{enumerate}
\item 方法一：配方法：
\begin{align*}
\int \frac{\mathrm{d}x}{\sqrt{\frac{b^2+4ac}{4a}-a\left(x-\frac{b}{2a}\right)^2}} &= \frac{1}{\sqrt{a}}\arcsin\frac{x-\frac{b}{2a}}{\sqrt{\frac{b^2+4ac}{4a^2}}}+C = \frac{1}{\sqrt{a}}\arcsin\frac{2ax-b}{\sqrt{b^2+4ac}}+C.
\end{align*}
\item 方法二：记 $y=\sqrt{c+bx-ax^2}$。令 $u=2ax-b, v=2\sqrt{a}y$，则 $u^2+v^2=b^2+4ac=\Delta$（实圆底座）。
由 $u\mathrm{d}u+v\mathrm{d}v=0 \implies \frac{\mathrm{d}u}{v}=-\frac{\mathrm{d}v}{u}=\mathrm{d}\arcsin\frac{u}{\sqrt{\Delta}}$：
\begin{align*}
\int \frac{\mathrm{d}x}{y} &= \int \frac{\mathrm{d}u/(2a)}{v/(2\sqrt{a})} = \frac{1}{\sqrt{a}}\int \frac{\mathrm{d}u}{v} = \frac{1}{\sqrt{a}}\arcsin\frac{u}{\sqrt{\Delta}}+C \\
&= \frac{1}{\sqrt{a}}\arcsin\frac{2ax-b}{\sqrt{b^2+4ac}}+C.
\end{align*}
（注：原题答案的负号为 $\arccos$ 变体或 $u=b-2ax$ 引起的等价符号差。）
\end{enumerate}
\end{solution}

\begin{exercise}[含有 $\sqrt{ax^2+bx+c}$ 的积分 - 77]
$\displaystyle \int \sqrt{c+bx-ax^2}\mathrm{d}x$
\end{exercise}
\begin{solution}
\begin{enumerate}
\item 方法一：配方后利用分部积分法或实圆三角换元法计算。
\item 方法二：同上，设 $u=2ax-b, v=2\sqrt{a}y$，利用实圆面积微元 $\int v\mathrm{d}u = \frac{1}{2}uv + \frac{\Delta}{2}\int \frac{\mathrm{d}u}{v}$：
\begin{align*}
\int y\mathrm{d}x &= \int \frac{v}{2\sqrt{a}}\frac{\mathrm{d}u}{2a} = \frac{1}{4a\sqrt{a}}\int v\mathrm{d}u = \frac{1}{4a\sqrt{a}}\left(\frac{1}{2}uv + \frac{b^2+4ac}{2}\int \frac{\mathrm{d}u}{v}\right) \\
&= \frac{uv}{8a\sqrt{a}} + \frac{b^2+4ac}{8a\sqrt{a}}\arcsin\frac{u}{\sqrt{\Delta}}+C \\
&= \frac{2ax-b}{4a}\sqrt{c+bx-ax^2} + \frac{b^2+4ac}{8\sqrt{a^3}}\arcsin\frac{2ax-b}{\sqrt{b^2+4ac}}+C.
\end{align*}
\end{enumerate}
\end{solution}

\begin{exercise}[含有 $\sqrt{ax^2+bx+c}$ 的积分 - 78]
$\displaystyle \int \frac{x}{\sqrt{c+bx-ax^2}}\mathrm{d}x$
\end{exercise}
\begin{solution}
\begin{enumerate}
\item 方法一：凑微分法：
\begin{align*}
\int \frac{x}{\sqrt{c+bx-ax^2}}\mathrm{d}x &= -\frac{1}{2a}\int \frac{b-2ax-b}{\sqrt{c+bx-ax^2}}\mathrm{d}x = -\frac{1}{a}\sqrt{c+bx-ax^2} + \frac{b}{2a}\int \frac{\mathrm{d}x}{\sqrt{c+bx-ax^2}}.
\end{align*}
由此可得所求结果。
\item 方法二：设 $u=2ax-b, v=2\sqrt{a}y$，代入 $x=\frac{u+b}{2a}$ 进行双元降维：
\begin{align*}
\int \frac{x}{y}\mathrm{d}x &= \int \frac{(u+b)/(2a)}{v/(2\sqrt{a})}\frac{\mathrm{d}u}{2a} = \frac{1}{2a\sqrt{a}}\int \frac{u+b}{v}\mathrm{d}u \\
&= \frac{1}{2a\sqrt{a}}\left(\int \frac{u\mathrm{d}u}{v} + b\int \frac{\mathrm{d}u}{v}\right) = \frac{1}{2a\sqrt{a}}\left(-\int \mathrm{d}v + b\int \frac{\mathrm{d}u}{v}\right) \\
&= -\frac{v}{2a\sqrt{a}} + \frac{b}{2a\sqrt{a}}\arcsin\frac{u}{\sqrt{\Delta}}+C = -\frac{1}{a}y + \frac{b}{2\sqrt{a^3}}\arcsin\frac{2ax-b}{\sqrt{b^2+4ac}}+C.
\end{align*}
\end{enumerate}
\end{solution}

\subsection{第十一节习题：\texorpdfstring{含有 $\sqrt{\frac{x-a}{x-b}}$ 或 $\sqrt{(x-a)(b-x)}$ 的积分}{含有 sqrt((x-a)/(x-b)) 或 sqrt((x-a)(b-x)) 的积分}}
\begin{exercise}[含有 $\sqrt{\frac{x-a}{x-b}}$ 的积分 - 79]
$\displaystyle \int \sqrt{\frac{x-a}{x-b}}\mathrm{d}x$
\end{exercise}
\begin{solution}
\begin{enumerate}
\item 方法一：令 $t = \sqrt{x-b}$，则 $x = t^2+b, \mathrm{d}x = 2t\mathrm{d}t$：
\begin{align*}
\int \sqrt{\frac{t^2+b-a}{t^2}}2t\mathrm{d}t &= 2\int \sqrt{t^2+(b-a)}\mathrm{d}t = 2\left[\frac{t}{2}\sqrt{t^2+b-a} + \frac{b-a}{2}\ln|t+\sqrt{t^2+b-a}|\right]+C \\
&= (x-b)\sqrt{\frac{x-a}{x-b}}+(b-a)\ln(\sqrt{|x-b|}+\sqrt{|x-a|})+C.
\end{align*}
\item 方法二：引入纯根式双元 $p=\sqrt{|x-a|}, q=\sqrt{|x-b|}$。假定 $x>a,b$，则 $p^2-q^2=b-a$（虚圆）。
$\mathrm{d}x = 2p\mathrm{d}p = 2q\mathrm{d}q$，且 $p\mathrm{d}p=q\mathrm{d}q \implies \frac{\mathrm{d}p}{q}=\frac{\mathrm{d}q}{p}$：
\begin{align*}
\int \frac{p}{q}\mathrm{d}x &= \int \frac{p}{q}(2q\mathrm{d}q) = 2\int p\mathrm{d}q = \int (p\mathrm{d}q+q\mathrm{d}p) + \int (p\mathrm{d}q-q\mathrm{d}p) \\
&= pq + \int (p^2-q^2)\frac{\mathrm{d}q}{p} = pq + (b-a)\int \frac{\mathrm{d}q}{p} = pq + (b-a)\ln|p+q|+C \\
&= \sqrt{x-a}\sqrt{x-b} + (b-a)\ln(\sqrt{x-a}+\sqrt{x-b})+C.
\end{align*}
\end{enumerate}
\end{solution}

\begin{exercise}[含有 $\sqrt{\frac{x-a}{b-x}}$ 的积分 - 80]
$\displaystyle \int \sqrt{\frac{x-a}{b-x}}\mathrm{d}x$
\end{exercise}
\begin{solution}
\begin{enumerate}
\item 方法一：令 $t = \sqrt{b-x}$，则 $x = b-t^2, \mathrm{d}x = -2t\mathrm{d}t$：
\begin{align*}
\int \sqrt{\frac{b-a-t^2}{t^2}}(-2t)\mathrm{d}t &= -2\int \sqrt{b-a-t^2}\mathrm{d}t = -2\left[\frac{t}{2}\sqrt{b-a-t^2} + \frac{b-a}{2}\arcsin\frac{t}{\sqrt{b-a}}\right]+C \\
&= -t\sqrt{b-a-t^2} - (b-a)\arcsin\frac{\sqrt{b-x}}{\sqrt{b-a}}+C \\
&= (x-b)\sqrt{\frac{x-a}{b-x}} + (b-a)\arcsin\sqrt{\frac{x-a}{b-a}}+C_1.
\end{align*}
\item 方法二：设 $p=\sqrt{x-a}, q=\sqrt{b-x}$。则 $p^2+q^2=b-a$（实圆）。
$\mathrm{d}x = 2p\mathrm{d}p = -2q\mathrm{d}q$，且 $p\mathrm{d}p=-q\mathrm{d}q \implies \frac{\mathrm{d}p}{q}=-\frac{\mathrm{d}q}{p}$：
\begin{align*}
\int \frac{p}{q}\mathrm{d}x &= \int \frac{p}{q}(-2q\mathrm{d}q) = -2\int p\mathrm{d}q = -\int (p\mathrm{d}q+q\mathrm{d}p) - \int (p\mathrm{d}q-q\mathrm{d}p) \\
&= -pq - \int (-p^2-q^2)\frac{\mathrm{d}p}{q} = -pq + (b-a)\int \frac{\mathrm{d}p}{q} \\
&= -pq + (b-a)\arcsin\frac{p}{\sqrt{b-a}}+C = -\sqrt{x-a}\sqrt{b-x} + (b-a)\arcsin\sqrt{\frac{x-a}{b-a}}+C.
\end{align*}
\end{enumerate}
\end{solution}

\begin{exercise}[含有 $\sqrt{(x-a)(b-x)}$ 的积分 - 81]
$\displaystyle \int \frac{\mathrm{d}x}{\sqrt{(x-a)(b-x)}}$
\end{exercise}
\begin{solution}
\begin{enumerate}
\item 方法一：令 $x-a = t^2$，则 $\mathrm{d}x = 2t\mathrm{d}t$：
\[
\int \frac{2t\mathrm{d}t}{t\sqrt{b-t^2-a}} = 2\int \frac{\mathrm{d}t}{\sqrt{b-a-t^2}} = 2\arcsin\frac{t}{\sqrt{b-a}}+C = 2\arcsin\sqrt{\frac{x-a}{b-a}}+C.
\]
\item 方法二：引入正交基底 $p=\sqrt{x-a}, q=\sqrt{b-x}$，则 $p^2+q^2=b-a,\ \mathrm{d}x=2p\mathrm{d}p$。于是
\[
\int \frac{\mathrm{d}x}{pq} = \int \frac{2p\mathrm{d}p}{pq} = 2\int \frac{\mathrm{d}p}{q} = 2\arcsin\frac{p}{\sqrt{b-a}}+C = 2\arcsin\sqrt{\frac{x-a}{b-a}}+C.
\]
\end{enumerate}
\end{solution}

\begin{exercise}[含有 $\sqrt{(x-a)(b-x)}$ 的积分 - 82]
$\displaystyle \int \sqrt{(x-a)(b-x)}\mathrm{d}x$
\end{exercise}
\begin{solution}
\begin{enumerate}
\item 方法一：令 $x-a = t^2$，转化为 $\int 2t^2\sqrt{b-a-t^2}\mathrm{d}t$，再通过 $t = \sqrt{b-a}\sin\theta$ 进行三角换元。
\item 方法二：设 $p=\sqrt{x-a}, q=\sqrt{b-x}$。利用 $p, q$ 的乘积与平方差构造辅助变量：
令 $u = p^2-q^2 = 2x-a-b$，$v = 2pq = 2\sqrt{(x-a)(b-x)}$。
此时 $u^2+v^2 = (p^2+q^2)^2 = (b-a)^2$（半径为 $b-a$ 的实圆），且 $\mathrm{d}u = 2\mathrm{d}x$：
\begin{align*}
\int pq\mathrm{d}x &= \int \frac{v}{2}\frac{\mathrm{d}u}{2} = \frac{1}{4}\int v\mathrm{d}u = \frac{1}{4}\left(\frac{1}{2}uv + \frac{(b-a)^2}{2}\arcsin\frac{u}{b-a}\right) \\
&= \frac{uv}{8} + \frac{(b-a)^2}{8}\arcsin\frac{2x-a-b}{b-a}+C \\
&= \frac{2x-a-b}{4}\sqrt{(x-a)(b-x)} + \frac{(b-a)^2}{4}\arcsin\sqrt{\frac{x-a}{b-a}}+C_1.
\end{align*}
（注：$\arcsin(2t-1) \equiv 2\arcsin\sqrt{t} - \frac{\pi}{2}$，其中常数项可并入积分常数。）
\end{enumerate}
\end{solution}

\subsection{第十二节习题：含有三角函数的积分}
\begin{exercise}[含有三角函数的积分 - 83]
$\displaystyle \int \sin x\mathrm{d}x$
\end{exercise}
\begin{solution}
\begin{enumerate}
\item 方法一：利用基本导数公式直接计算：\(\displaystyle \int \sin x\mathrm{d}x = \int \mathrm{d}(-\cos x) = -\cos x+C\)。
\item 方法二：记 $p=\sin x, q=\cos x$，根据复合函数求导法则，建立实圆微分约束 $\mathrm{d}q=-\sin x\mathrm{d}x=-p\mathrm{d}x$。代入积分中有 \(\displaystyle \int \sin x\mathrm{d}x=\int p\mathrm{d}x=\int (-\mathrm{d}q)\)，故 \(\displaystyle \int \sin x\mathrm{d}x=-q+C=-\cos x+C\)。
\end{enumerate}
\end{solution}

\begin{exercise}[含有三角函数的积分 - 84]
$\displaystyle \int \cos x\mathrm{d}x$
\end{exercise}
\begin{solution}
\begin{enumerate}
\item 方法一：利用基本导数公式直接计算：\(\displaystyle \int \cos x\mathrm{d}x = \int \mathrm{d}(\sin x) = \sin x+C\)。
\item 方法二：记 $p=\sin x, q=\cos x$，建立实圆微分约束 $\mathrm{d}p=\cos x\mathrm{d}x=q\mathrm{d}x$。代入积分中有 \(\displaystyle \int \cos x\mathrm{d}x=\int q\mathrm{d}x=\int \mathrm{d}p\)，故 \(\displaystyle \int \cos x\mathrm{d}x=p+C=\sin x+C\)。
\end{enumerate}
\end{solution}

\begin{exercise}[含有三角函数的积分 - 85]
$\displaystyle \int \tan x\mathrm{d}x$
\end{exercise}
\begin{solution}
\begin{enumerate}
\item 方法一：将正切函数展开为正弦与余弦的比值，利用凑微分法将分子纳入微分号内。先写
\[
\int \tan x\,\mathrm{d}x
=\int \frac{\sin x}{\cos x}\,\mathrm{d}x
=\int \frac{1}{\cos x}\bigl(-\mathrm{d}(\cos x)\bigr),
\]
于是
\(\displaystyle \int \tan x\,\mathrm{d}x=-\int \frac{\mathrm{d}(\cos x)}{\cos x}=-\ln|\cos x|+C\)。
\item 方法二：记 $p=\sin x, q=\cos x$，根据底层微元关系 $\mathrm{d}q = -p\mathrm{d}x$，即 $\mathrm{d}x = -\frac{\mathrm{d}q}{p}$，提取实圆对数微元：
\[
\int \tan x\,\mathrm{d}x=\int \frac{p}{q}\left(-\frac{\mathrm{d}q}{p}\right)=-\int \frac{\mathrm{d}q}{q},
\]
所以 \(\displaystyle \int \tan x\,\mathrm{d}x=-\ln|q|+C=-\ln|\cos x|+C\)。
\end{enumerate}
\end{solution}

\begin{exercise}[含有三角函数的积分 - 86]
$\displaystyle \int \cot x\mathrm{d}x$
\end{exercise}
\begin{solution}
\begin{enumerate}
\item 方法一：将余切函数展开为余弦与正弦的比值，凑微分法计算：
\[
\int \cot x\mathrm{d}x
=\int \frac{\cos x}{\sin x}\mathrm{d}x
=\int \frac{1}{\sin x}\mathrm{d}(\sin x).
\]
所以结果为 \(\displaystyle \ln|\sin x|+C\)。
\item 方法二：记 $p=\sin x, q=\cos x$，由 $\mathrm{d}p = q\mathrm{d}x$ 得 $\mathrm{d}x = \dfrac{\mathrm{d}p}{q}$，于是 \(\displaystyle \int \cot x\mathrm{d}x=\int \frac{q}{p}\frac{\mathrm{d}p}{q}=\int \frac{\mathrm{d}p}{p}\)。所以结果为 \(\displaystyle \ln|p|+C=\ln|\sin x|+C\)。
\end{enumerate}
\end{solution}

\begin{exercise}[含有三角函数的积分 - 87]
$\displaystyle \int \sec x\mathrm{d}x$
\end{exercise}
\begin{solution}
\begin{enumerate}
\item 方法一：分子分母同乘 $(\sec x+\tan x)$ 构造导数结构。此时分子化为 $\sec^2x+\sec x\tan x$，正好是 $\tan x+\sec x$ 的导数，所以 \(\displaystyle \int \sec x\mathrm{d}x=\int \frac{\mathrm{d}(\tan x+\sec x)}{\tan x+\sec x}\)，结果为 \(\displaystyle \ln|\sec x+\tan x|+C\)。
\item 方法二：记双元 $p=\sec x, q=\tan x$。利用微分关系 $\mathrm{d}p=pq\mathrm{d}x$ 与 $\mathrm{d}q=p^2\mathrm{d}x$，两式相加得 $\mathrm{d}(p+q) = p(p+q)\mathrm{d}x$，即 \(p\mathrm{d}x = \dfrac{\mathrm{d}(p+q)}{p+q}\)。于是 \(\displaystyle \int \sec x\mathrm{d}x=\int \frac{\mathrm{d}(p+q)}{p+q}\)，故结果为 \(\displaystyle \ln|p+q|+C=\ln|\sec x+\tan x|+C\)。
\end{enumerate}
\end{solution}

\begin{exercise}[含有三角函数的积分 - 88]
$\displaystyle \int \csc x\mathrm{d}x$
\end{exercise}
\begin{solution}
\begin{enumerate}
\item 方法一：分子分母同乘 $(\csc x-\cot x)$ 构造导数结构。此时分子化为 $\csc^2x-\csc x\cot x$，正好是 $-\cot x-\csc x$ 的导数，所以 \(\displaystyle \int \csc x\mathrm{d}x=\int \frac{\mathrm{d}(-\cot x-\csc x)}{\csc x-\cot x}\)，结果为 \(\displaystyle \ln|\csc x-\cot x|+C\)。
\item 方法二：记双元 $p=\csc x, q=\cot x$。利用微分关系 $\mathrm{d}p=-pq\mathrm{d}x$ 与 $\mathrm{d}q=-p^2\mathrm{d}x$，两式相减得 $\mathrm{d}(p-q) = p(p-q)\mathrm{d}x$，即 \(p\mathrm{d}x = \dfrac{\mathrm{d}(p-q)}{p-q}\)。于是 \(\displaystyle \int \csc x\mathrm{d}x=\int \frac{\mathrm{d}(p-q)}{p-q}\)，故结果为 \(\displaystyle \ln|p-q|+C=\ln|\csc x-\cot x|+C\)。
\end{enumerate}
\end{solution}

\begin{exercise}[含有三角函数的积分 - 89]
$\displaystyle \int \sec^2 x\mathrm{d}x$
\end{exercise}
\begin{solution}
\begin{enumerate}
\item 方法一：直接根据导数公式 $\frac{\mathrm{d}}{\mathrm{d}x}(\tan x) = \sec^2 x$ 积分：\(\displaystyle \int \sec^2 x\mathrm{d}x=\tan x+C\)。
\item 方法二：记 $p=\sec x, q=\tan x$，直接应用微元约束 $\mathrm{d}q=p^2\mathrm{d}x$：\(\displaystyle \int \sec^2 x\mathrm{d}x=\int p^2\mathrm{d}x=\int \mathrm{d}q=q+C=\tan x+C\)。
\end{enumerate}
\end{solution}

\begin{exercise}[含有三角函数的积分 - 90]
$\displaystyle \int \csc^2 x\mathrm{d}x$
\end{exercise}
\begin{solution}
\begin{enumerate}
\item 方法一：直接根据导数公式 $\frac{\mathrm{d}}{\mathrm{d}x}(\cot x) = -\csc^2 x$ 积分：\(\displaystyle \int \csc^2 x\mathrm{d}x=-\cot x+C\)。
\item 方法二：记 $p=\csc x, q=\cot x$，直接应用微元约束 $\mathrm{d}q=-p^2\mathrm{d}x$：\(\displaystyle \int \csc^2 x\mathrm{d}x=\int p^2\mathrm{d}x=\int (-\mathrm{d}q)=-q+C=-\cot x+C\)。
\end{enumerate}
\end{solution}

\begin{exercise}[含有三角函数的积分 - 91]
$\displaystyle \int \sec x\tan x\mathrm{d}x$
\end{exercise}
\begin{solution}
\begin{enumerate}
\item 方法一：直接根据导数公式 $\frac{\mathrm{d}}{\mathrm{d}x}(\sec x) = \sec x\tan x$ 积分：\(\displaystyle \int \sec x\tan x\mathrm{d}x=\sec x+C\)。
\item 方法二：记 $p=\sec x, q=\tan x$，利用交叉微分约束 $\mathrm{d}p=pq\mathrm{d}x$：\(\displaystyle \int \sec x\tan x\mathrm{d}x=\int pq\mathrm{d}x=\int \mathrm{d}p=p+C=\sec x+C\)。
\end{enumerate}
\end{solution}

\begin{exercise}[含有三角函数的积分 - 92]
$\displaystyle \int \csc x\cot x\mathrm{d}x$
\end{exercise}
\begin{solution}
\begin{enumerate}
\item 方法一：直接根据导数公式 $\frac{\mathrm{d}}{\mathrm{d}x}(\csc x) = -\csc x\cot x$ 积分：\(\displaystyle \int \csc x\cot x\mathrm{d}x=-\csc x+C\)。
\item 方法二：记 $p=\csc x, q=\cot x$，利用交叉微分约束 $\mathrm{d}p=-pq\mathrm{d}x$：\(\displaystyle \int \csc x\cot x\mathrm{d}x=\int pq\mathrm{d}x=\int (-\mathrm{d}p)=-p+C=-\csc x+C\)。
\end{enumerate}
\end{solution}

\begin{exercise}[含有三角函数的积分 - 93]
$\displaystyle \int \sin^2 x\mathrm{d}x$
\end{exercise}
\begin{solution}
\begin{enumerate}
\item 方法一：利用半角公式 $\sin^2 x = \frac{1-\cos 2x}{2}$ 进行线性化降次，可得 \(\displaystyle \int \sin^2 x\mathrm{d}x=\int \frac{1-\cos 2x}{2}\mathrm{d}x\)，故结果为 \(\displaystyle \frac{x}{2}-\frac{1}{4}\sin 2x+C\)。
\item 方法二：记 $p=\sin x, q=\cos x$ 且 $p^2+q^2=1$。通过实圆面积的分部展开，完全规避倍角公式：
\begin{align*}
\int p^2\mathrm{d}x &= \int p(-\mathrm{d}q) = -pq - \int (-q)\mathrm{d}p = -pq + \int q^2\mathrm{d}x \\
2\int p^2\mathrm{d}x &= \int p^2\mathrm{d}x + \left( -pq + \int q^2\mathrm{d}x \right) = -pq + \int (p^2+q^2)\mathrm{d}x \\
2\int p^2\mathrm{d}x &= -pq + \int 1\mathrm{d}x = x - pq \\
\implies \int \sin^2 x\mathrm{d}x &= \frac{x}{2} - \frac{1}{2}pq+C = \frac{x}{2} - \frac{1}{4}\sin 2x+C.
\end{align*}
\end{enumerate}
\end{solution}

\begin{exercise}[含有三角函数的积分 - 94]
$\displaystyle \int \cos^2 x\mathrm{d}x$
\end{exercise}
\begin{solution}
\begin{enumerate}
\item 方法一：利用半角公式 $\cos^2 x = \frac{1+\cos 2x}{2}$ 降次，可得 \(\displaystyle \int \cos^2 x\mathrm{d}x=\int \frac{1+\cos 2x}{2}\mathrm{d}x\)，故结果为 \(\displaystyle \frac{x}{2}+\frac{1}{4}\sin 2x+C\)。
\item 方法二：记 $p=\sin x, q=\cos x$。通过对称面积代换推导：
\begin{align*}
\int q^2\mathrm{d}x &= \int q(\mathrm{d}p) = pq - \int p\mathrm{d}q = pq - \int p(-p\mathrm{d}x) = pq + \int p^2\mathrm{d}x \\
2\int q^2\mathrm{d}x &= \int q^2\mathrm{d}x + \left( pq + \int p^2\mathrm{d}x \right) = pq + \int (p^2+q^2)\mathrm{d}x \\
2\int q^2\mathrm{d}x &= pq + \int 1\mathrm{d}x = x + pq \\
\implies \int \cos^2 x\mathrm{d}x &= \frac{x}{2} + \frac{1}{2}pq+C = \frac{x}{2} + \frac{1}{4}\sin 2x+C.
\end{align*}
\end{enumerate}
\end{solution}

\begin{exercise}[含有三角函数的积分 - 95]
$\displaystyle \int \sin^n x\mathrm{d}x$
\end{exercise}
\begin{solution}
\begin{enumerate}
\item 方法一：分离一个 $\sin x$ 进入微分符号，利用分部积分法 $u=\sin^{n-1}x, \mathrm{d}v=\sin x\mathrm{d}x$：
\begin{align*}
I_n &= \int \sin^{n-1}x (-\mathrm{d}(\cos x)) = -\sin^{n-1}x\cos x + \int \cos x \mathrm{d}(\sin^{n-1}x) \\
&= -\sin^{n-1}x\cos x + \int \cos x (n-1)\sin^{n-2}x \cos x \mathrm{d}x \\
&= -\sin^{n-1}x\cos x + (n-1)\int \sin^{n-2}x (1-\sin^2 x) \mathrm{d}x \\
&= -\sin^{n-1}x\cos x + (n-1)I_{n-2} - (n-1)I_n
\end{align*}
移项合并得 $n I_n = -\sin^{n-1}x\cos x + (n-1)I_{n-2}$，两边同除以 $n$ 即可得证。
\item 方法二：记 $p=\sin x, q=\cos x$，在双元底座 $p^2+q^2=1$ 约束下直接纯代数递推：
\begin{align*}
I_n &= \int p^n\mathrm{d}x = \int p^{n-1}(-\mathrm{d}q) = -p^{n-1}q + \int q \mathrm{d}(p^{n-1}) \\
&= -p^{n-1}q + \int q(n-1)p^{n-2}\mathrm{d}p = -p^{n-1}q + (n-1)\int q^2 p^{n-2}\mathrm{d}x \\
&= -p^{n-1}q + (n-1)\int (1-p^2)p^{n-2}\mathrm{d}x = -p^{n-1}q + (n-1)I_{n-2} - (n-1)I_n \\
n I_n &= -p^{n-1}q + (n-1)I_{n-2} \implies I_n = -\frac{1}{n}\sin^{n-1}x\cos x+\frac{n-1}{n}\int \sin^{n-2}x\mathrm{d}x.
\end{align*}
\end{enumerate}
\end{solution}

\begin{exercise}[含有三角函数的积分 - 96]
$\displaystyle \int \cos^n x\mathrm{d}x$
\end{exercise}
\begin{solution}
\begin{enumerate}
\item 方法一：分离一个 $\cos x$，利用分部积分法 $u=\cos^{n-1}x, \mathrm{d}v=\cos x\mathrm{d}x$：
\begin{align*}
I_n &= \int \cos^{n-1}x \mathrm{d}(\sin x) = \cos^{n-1}x\sin x - \int \sin x \mathrm{d}(\cos^{n-1}x) \\
&= \cos^{n-1}x\sin x - \int \sin x (n-1)\cos^{n-2}x (-\sin x)\mathrm{d}x \\
&= \cos^{n-1}x\sin x + (n-1)\int \cos^{n-2}x (1-\cos^2 x)\mathrm{d}x \\
&= \cos^{n-1}x\sin x + (n-1)I_{n-2} - (n-1)I_n
\end{align*}
移项合并得 $n I_n = \cos^{n-1}x\sin x + (n-1)I_{n-2}$，两边同除以 $n$ 得证。
\item 方法二：记 $p=\sin x, q=\cos x$，在实圆底座下完全对称展开：
\begin{align*}
I_n &= \int q^n\mathrm{d}x = \int q^{n-1}(\mathrm{d}p) = q^{n-1}p - \int p\mathrm{d}(q^{n-1}) \\
&= q^{n-1}p - \int p(n-1)q^{n-2}(-\mathrm{d}x p) = q^{n-1}p + (n-1)\int p^2 q^{n-2}\mathrm{d}x \\
&= q^{n-1}p + (n-1)\int (1-q^2)q^{n-2}\mathrm{d}x = q^{n-1}p + (n-1)I_{n-2} - (n-1)I_n \\
n I_n &= q^{n-1}p + (n-1)I_{n-2} \implies I_n = \frac{1}{n}\cos^{n-1}x\sin x+\frac{n-1}{n}\int \cos^{n-2}x\mathrm{d}x.
\end{align*}
\end{enumerate}
\end{solution}

\begin{exercise}[含有三角函数的积分 - 97]
$\displaystyle \int \frac{\mathrm{d}x}{\sin^n x}$
\end{exercise}
\begin{solution}
\begin{enumerate}
\item 方法一：重写为 $I_n = \int \csc^n x \mathrm{d}x$，分离平方项进行分部积分：
\begin{align*}
I_n &= \int \csc^{n-2}x \csc^2 x \mathrm{d}x = \int \csc^{n-2}x (-\mathrm{d}(\cot x)) \\
&= -\csc^{n-2}x \cot x + \int \cot x \mathrm{d}(\csc^{n-2}x) \\
&= -\csc^{n-2}x \cot x + \int \cot x (n-2)\csc^{n-3}x (-\csc x \cot x)\mathrm{d}x \\
&= -\csc^{n-2}x \cot x - (n-2)\int \csc^{n-2}x (\csc^2 x - 1)\mathrm{d}x \\
&= -\csc^{n-2}x \cot x - (n-2)I_n + (n-2)I_{n-2} \\
(n-1)I_n &= -\csc^{n-2}x \cot x + (n-2)I_{n-2} \implies I_n = -\frac{1}{n-1}\frac{\cos x}{\sin^{n-1}x} + \frac{n-2}{n-1}I_{n-2}.
\end{align*}
\item 方法二：记 $p=\sin x, q=\cos x$，本题目标为 $I_{-n} = \int p^{-n}\mathrm{d}x$。
利用第 95 题已证得的代数递推式 $k I_k = -p^{k-1}q + (k-1)I_{k-2}$，将其指标参数替换为 $k = 2-n$，得 \(\displaystyle (2-n)I_{2-n} = -p^{(2-n)-1}q + (2-n-1)I_{(2-n)-2} = -p^{1-n}q + (1-n)I_{-n}\)。
重排方程，将 $I_{-n}$ 孤立在一侧：
\begin{align*}
(n-1)I_{-n} &= -p^{1-n}q + (n-2)I_{2-n} \\
I_{-n} &= -\frac{1}{n-1}\frac{q}{p^{n-1}} + \frac{n-2}{n-1}I_{-(n-2)} \\
\int \frac{\mathrm{d}x}{\sin^n x} &= -\frac{1}{n-1}\frac{\cos x}{\sin^{n-1}x}+\frac{n-2}{n-1}\int \frac{\mathrm{d}x}{\sin^{n-2}x}.
\end{align*}
\end{enumerate}
\end{solution}

\begin{exercise}[含有三角函数的积分 - 98]
$\displaystyle \int \frac{\mathrm{d}x}{\cos^n x}$
\end{exercise}
\begin{solution}
\begin{enumerate}
\item 方法一：重写为 $I_n = \int \sec^n x \mathrm{d}x$，分离平方项进行分部积分：
\begin{align*}
I_n &= \int \sec^{n-2}x \sec^2 x \mathrm{d}x = \int \sec^{n-2}x \mathrm{d}(\tan x) \\
&= \sec^{n-2}x \tan x - \int \tan x \mathrm{d}(\sec^{n-2}x) \\
&= \sec^{n-2}x \tan x - \int \tan x (n-2)\sec^{n-3}x (\sec x \tan x)\mathrm{d}x \\
&= \sec^{n-2}x \tan x - (n-2)\int \sec^{n-2}x (\sec^2 x - 1)\mathrm{d}x \\
&= \sec^{n-2}x \tan x - (n-2)I_n + (n-2)I_{n-2} \\
(n-1)I_n &= \sec^{n-2}x \tan x + (n-2)I_{n-2} \implies I_n = \frac{1}{n-1}\frac{\sin x}{\cos^{n-1}x} + \frac{n-2}{n-1}I_{n-2}.
\end{align*}
\item 方法二：记 $p=\sin x, q=\cos x$，目标为 $I_{-n} = \int q^{-n}\mathrm{d}x$。
利用第 96 题已证得的代数递推式 $k I_k = q^{k-1}p + (k-1)I_{k-2}$，将指标替换为 $k = 2-n$，得 \(\displaystyle (2-n)I_{2-n} = q^{(2-n)-1}p + (2-n-1)I_{(2-n)-2} = q^{1-n}p + (1-n)I_{-n}\)。
重排方程提取 $I_{-n}$：
\begin{align*}
(n-1)I_{-n} &= q^{1-n}p + (n-2)I_{2-n} \\
I_{-n} &= \frac{1}{n-1}\frac{p}{q^{n-1}} + \frac{n-2}{n-1}I_{-(n-2)} \\
\int \frac{\mathrm{d}x}{\cos^n x} &= \frac{1}{n-1}\frac{\sin x}{\cos^{n-1}x}+\frac{n-2}{n-1}\int \frac{\mathrm{d}x}{\cos^{n-2}x}.
\end{align*}
\end{enumerate}
\end{solution}

\begin{exercise}[含有三角函数的积分 - 99]
$\displaystyle \int \cos^m x\sin^n x\mathrm{d}x$
\end{exercise}
\begin{solution}
\begin{enumerate}
\item 方法一：常规分部积分，令 $I_{m,n} = \int \cos^m x \sin^n x \mathrm{d}x$。结合微元 $\sin^n x \cos x \mathrm{d}x = \frac{1}{n+1}\mathrm{d}(\sin^{n+1}x)$：
\begin{align*}
I_{m,n} &= \int \cos^{m-1}x \cdot \cos x \sin^n x \mathrm{d}x = \frac{1}{n+1}\int \cos^{m-1}x \mathrm{d}(\sin^{n+1}x) \\
&= \frac{1}{n+1}\cos^{m-1}x \sin^{n+1}x - \frac{1}{n+1}\int \sin^{n+1}x (m-1)\cos^{m-2}x (-\sin x)\mathrm{d}x \\
&= \frac{1}{n+1}\cos^{m-1}x \sin^{n+1}x + \frac{m-1}{n+1}\int \cos^{m-2}x \sin^n x (1-\cos^2 x)\mathrm{d}x \\
&= \frac{1}{n+1}\cos^{m-1}x \sin^{n+1}x + \frac{m-1}{n+1}I_{m-2,n} - \frac{m-1}{n+1}I_{m,n}
\end{align*}
两边同乘 $(n+1)$ 并提取 $I_{m,n}$ 得 $(m+n)I_{m,n} = \cos^{m-1}x \sin^{n+1}x + (m-1)I_{m-2,n}$，除以 $m+n$ 即证第一式。另一重递推式同理。
\item 方法二：记 $p=\sin x, q=\cos x$。在实圆双元体系内分部积分，并利用底座 $p^2+q^2=1$ 展开：
\begin{align*}
\int q^m p^n\mathrm{d}x &= \int q^{m-1}p^n (\mathrm{d}p) = \frac{1}{n+1}\int q^{m-1}\mathrm{d}(p^{n+1}) \\
&= \frac{1}{n+1}q^{m-1}p^{n+1} - \frac{1}{n+1}\int p^{n+1}\mathrm{d}(q^{m-1}) \\
&= \frac{1}{n+1}q^{m-1}p^{n+1} - \frac{1}{n+1}\int p^{n+1}(m-1)q^{m-2}(-\mathrm{d}x p) \\
&= \frac{1}{n+1}q^{m-1}p^{n+1} + \frac{m-1}{n+1}\int q^{m-2}p^n (p^2)\mathrm{d}x \\
&= \frac{1}{n+1}q^{m-1}p^{n+1} + \frac{m-1}{n+1}\int q^{m-2}p^n (1-q^2)\mathrm{d}x \\
I_{m,n} &= \frac{1}{n+1}q^{m-1}p^{n+1} + \frac{m-1}{n+1}I_{m-2,n} - \frac{m-1}{n+1}I_{m,n} \\
\frac{m+n}{n+1} I_{m,n} &= \frac{1}{n+1}q^{m-1}p^{n+1} + \frac{m-1}{n+1}I_{m-2,n} \\
\implies I_{m,n} &= \frac{1}{m+n}\cos^{m-1}x\sin^{n+1}x+\frac{m-1}{m+n}\int \cos^{m-2}x\sin^n x\mathrm{d}x.
\end{align*}
\end{enumerate}
\end{solution}

\begin{exercise}[含有三角函数的积分 - 100]
$\displaystyle \int \sin ax\cos bx\mathrm{d}x$
\end{exercise}
\begin{note}[复指数双元什么时候最值钱]
复指数法不是另起炉灶，而是把三角函数看成复指数的实部、虚部，再把原来“乘积里有两个频率”的结构打包成若干单指数模态。这样一来，积分的难点就从三角恒等变成了指数的代数展开。

所以当题目里同时出现不同频率的三角函数，或者后面还会叠加指数阻尼项时，优先想到复指数双元法往往最稳，因为它能把“振荡结构”一次性拉直到复平面的单指数上。
\end{note}
\begin{solution}
\begin{enumerate}
\item 方法一：先设 $a\neq \pm b$。利用积化和差公式 \(\displaystyle \sin ax\cos bx=\frac12[\sin(a+b)x+\sin(a-b)x]\)，可得结果为 \(\displaystyle -\frac{1}{2(a+b)}\cos(a+b)x-\frac{1}{2(a-b)}\cos(a-b)x+C\)。
当 $b=\pm a$ 且 $a\neq 0$ 时，\(\displaystyle \int \sin ax\cos bx\,\mathrm{d}x=\frac{1}{2a}\sin^2 ax+C\)。
\item 方法二：复指数双元法。在 $a\neq \pm b$ 的条件下，由欧拉公式展开可得：
\begin{align*}
\int \sin ax\cos bx\,\mathrm{d}x
&=\int \frac{e^{iax}-e^{-iax}}{2i}\frac{e^{ibx}+e^{-ibx}}{2}\,\mathrm{d}x \\
&= \frac{1}{4i}\int\!\bigl(e^{i(a+b)x}+e^{i(a-b)x}-e^{-i(a-b)x}-e^{-i(a+b)x}\bigr)\,\mathrm{d}x \\
&= -\frac{1}{2(a+b)}\cos(a+b)x-\frac{1}{2(a-b)}\cos(a-b)x+C.
\end{align*}
\end{enumerate}
\end{solution}

\begin{exercise}[含有三角函数的积分 - 101]
$\displaystyle \int \sin ax\sin bx\mathrm{d}x$
\end{exercise}
\begin{solution}
\begin{enumerate}
\item 方法一：先设 $a\neq \pm b$。利用积化和差公式 \(\displaystyle \sin ax\sin bx=-\frac12[\cos(a+b)x-\cos(a-b)x]\)，可得结果为 \(\displaystyle -\frac{1}{2(a+b)}\sin(a+b)x+\frac{1}{2(a-b)}\sin(a-b)x+C\)。
特殊情形为：当 $b=a$ 且 $a\neq 0$ 时，\(\displaystyle \int \sin^2 ax\,\mathrm{d}x=\frac{x}{2}-\frac{\sin 2ax}{4a}+C\)；当 $b=-a$ 且 $a\neq 0$ 时，\(\displaystyle \int \sin ax\sin(-ax)\,\mathrm{d}x=-\frac{x}{2}+\frac{\sin 2ax}{4a}+C\)。
\item 方法二：复指数双元法：
\begin{align*}
\int \sin ax\sin bx\,\mathrm{d}x
&=\int \frac{e^{iax}-e^{-iax}}{2i}\frac{e^{ibx}-e^{-ibx}}{2i}\,\mathrm{d}x \\
&= -\frac{1}{4}\int\!\bigl(e^{i(a+b)x}-e^{i(a-b)x}-e^{-i(a-b)x}+e^{-i(a+b)x}\bigr)\,\mathrm{d}x \\
&= -\frac{1}{2(a+b)}\sin(a+b)x+\frac{1}{2(a-b)}\sin(a-b)x+C.
\end{align*}
\end{enumerate}
\end{solution}

\begin{exercise}[含有三角函数的积分 - 102]
$\displaystyle \int \cos ax\cos bx\mathrm{d}x$
\end{exercise}
\begin{solution}
\begin{enumerate}
\item 方法一：先设 $a\neq \pm b$。利用积化和差公式 \(\displaystyle \cos ax\cos bx=\frac12[\cos(a+b)x+\cos(a-b)x]\)，可得结果为 \(\displaystyle \frac{1}{2(a+b)}\sin(a+b)x+\frac{1}{2(a-b)}\sin(a-b)x+C\)。当 $b=\pm a$ 且 $a\neq 0$ 时，\(\displaystyle \int \cos ax\cos bx\,\mathrm{d}x=\frac{x}{2}+\frac{\sin 2ax}{4a}+C\)。
\item 方法二：复指数双元法：
\begin{align*}
\int \cos ax\cos bx\,\mathrm{d}x
&=\int \frac{e^{iax}+e^{-iax}}{2}\frac{e^{ibx}+e^{-ibx}}{2}\,\mathrm{d}x \\
&= \frac{1}{4}\int\!\bigl(e^{i(a+b)x}+e^{i(a-b)x}+e^{-i(a-b)x}+e^{-i(a+b)x}\bigr)\,\mathrm{d}x \\
&= \frac{1}{2(a+b)}\sin(a+b)x+\frac{1}{2(a-b)}\sin(a-b)x+C.
\end{align*}
\end{enumerate}
\end{solution}

\begin{exercise}[含有三角函数的积分 - 103]
$\displaystyle \int \frac{\mathrm{d}x}{a+b\sin x}$
\end{exercise}
\begin{solution}
\begin{enumerate}
\item 方法一：若 $a=0$，则原式化为 $\dfrac1b\int \csc x\,\mathrm{d}x$，可直接积分。以下设 $a\neq 0$。令 $t=\tan\frac{x}{2}$，则 $\mathrm{d}x=\frac{2\mathrm{d}t}{1+t^2}, \sin x=\frac{2t}{1+t^2}$：
\begin{align*}
\int \frac{\mathrm{d}x}{a+b\sin x} &= \int \frac{\frac{2\mathrm{d}t}{1+t^2}}{a+b\frac{2t}{1+t^2}} = \int \frac{2\mathrm{d}t}{a(1+t^2)+2bt} \\
&= 2\int \frac{\mathrm{d}t}{at^2+2bt+a} = \frac{2}{a} \int \frac{\mathrm{d}t}{\left(t+\frac{b}{a}\right)^2 + \frac{a^2-b^2}{a^2}}.
\end{align*}
若 $a^2>b^2$：
\begin{align*}
\text{原式} &= \frac{2}{a} \frac{a}{\sqrt{a^2-b^2}} \arctan \left( \frac{t+b/a}{\sqrt{a^2-b^2}/a} \right) = \frac{2}{\sqrt{a^2-b^2}}\arctan\frac{a\tan\frac{x}{2}+b}{\sqrt{a^2-b^2}}+C.
\end{align*}
若 $a^2<b^2$：
\begin{align*}
\text{原式} &= \frac{2}{a} \int \frac{\mathrm{d}t}{\left(t+\frac{b}{a}\right)^2 - \frac{b^2-a^2}{a^2}} = \frac{2}{a} \frac{a}{2\sqrt{b^2-a^2}} \ln\left| \frac{t+\frac{b}{a}-\frac{\sqrt{b^2-a^2}}{a}}{t+\frac{b}{a}+\frac{\sqrt{b^2-a^2}}{a}} \right| \\
&= \frac{1}{\sqrt{b^2-a^2}}\ln\left|\frac{a\tan\frac{x}{2}+b-\sqrt{b^2-a^2}}{a\tan\frac{x}{2}+b+\sqrt{b^2-a^2}}\right|+C.
\end{align*}
若 $a^2=b^2$，即 $b=\pm a$，则可直接化简为
\begin{align*}
\int\frac{\mathrm{d}x}{a+a\sin x}
&=\frac1a\int\frac{1-\sin x}{\cos^2 x}\,\mathrm{d}x
=\frac1a(\tan x-\sec x)+C,\\
\int\frac{\mathrm{d}x}{a-a\sin x}
&=\frac1a\int\frac{1+\sin x}{\cos^2 x}\,\mathrm{d}x
=\frac1a(\tan x+\sec x)+C.
\end{align*}
\item 方法二：作同样的代换 $p=\tan\frac{x}{2}$，得到
\[
\int \frac{\mathrm{d}x}{a+b\sin x}
= \int \frac{\frac{2\mathrm{d}p}{1+p^2}}{a+b\frac{2p}{1+p^2}}
= \int \frac{2\mathrm{d}p}{ap^2+2bp+a}.
\]
后续只需按二次型判别式 $\Delta=4(b^2-a^2)$ 的符号分类讨论，便回到上面的三种情形。
\end{enumerate}
\end{solution}

\begin{exercise}[含有三角函数的积分 - 104]
$\displaystyle \int \frac{\mathrm{d}x}{a+b\cos x}$
\end{exercise}
\begin{solution}
\begin{enumerate}
\item 方法一：若 $a=0$，则原式化为 $\dfrac1b\int \sec x\,\mathrm{d}x$，可直接积分。以下设 $a\neq 0$。令 $t=\tan\frac{x}{2}$，则 $\mathrm{d}x=\frac{2\mathrm{d}t}{1+t^2}, \cos x=\frac{1-t^2}{1+t^2}$：
\begin{align*}
\int \frac{\mathrm{d}x}{a+b\cos x} &= \int \frac{\frac{2\mathrm{d}t}{1+t^2}}{a+b\frac{1-t^2}{1+t^2}} = \int \frac{2\mathrm{d}t}{a(1+t^2)+b(1-t^2)} = \int \frac{2\mathrm{d}t}{(a-b)t^2+(a+b)}.
\end{align*}
若 $a^2>b^2$ (设 $a>b$)：
\begin{align*}
\text{原式} &= \frac{2}{a-b} \int \frac{\mathrm{d}t}{t^2 + \frac{a+b}{a-b}} = \frac{2}{a-b} \sqrt{\frac{a-b}{a+b}} \arctan\left( t\sqrt{\frac{a-b}{a+b}} \right) \\
&= \frac{2}{\sqrt{a^2-b^2}}\arctan\left(\sqrt{\frac{a-b}{a+b}}\tan\frac{x}{2}\right)+C.
\end{align*}
若 $a^2<b^2$：
\begin{align*}
\text{原式} &= \int \frac{2\mathrm{d}t}{-(b-a)t^2+(b+a)} = \frac{2}{b-a}\int \frac{\mathrm{d}t}{\frac{b+a}{b-a}-t^2} \\
&= \frac{2}{b-a}\frac{1}{2\sqrt{\frac{b+a}{b-a}}} \ln\left| \frac{\sqrt{\frac{b+a}{b-a}}+t}{\sqrt{\frac{b+a}{b-a}}-t} \right| = \frac{1}{b+a}\sqrt{\frac{b+a}{b-a}}\ln\left|\frac{\tan\frac{x}{2}+\sqrt{\frac{b+a}{b-a}}}{\tan\frac{x}{2}-\sqrt{\frac{b+a}{b-a}}}\right|+C.
\end{align*}
若 $a^2=b^2$，即 $b=\pm a$，则
\begin{align*}
\int\frac{\mathrm{d}x}{a+a\cos x}
&=\frac{1}{a}\int\frac{\mathrm{d}x}{1+\cos x}
=\frac{1}{a}\tan\frac{x}{2}+C,\\
\int\frac{\mathrm{d}x}{a-a\cos x}
&=\frac{1}{a}\int\frac{\mathrm{d}x}{1-\cos x}
=-\frac{1}{a}\cot\frac{x}{2}+C.
\end{align*}
\item 方法二：作同样的代换 $p=\tan\frac{x}{2}$，得到 \(\displaystyle \int \frac{\mathrm{d}x}{a+b\cos x} = \int \frac{2\mathrm{d}p}{(a-b)p^2+(a+b)}\)。
再按二次型的符号分情况积分，即得到与上面一致的结果。
\end{enumerate}
\end{solution}

\begin{exercise}[含有三角函数的积分 - 105]
$\displaystyle \int \frac{\mathrm{d}x}{a^2\cos^2 x+b^2\sin^2 x}$
\end{exercise}
\begin{solution}
\begin{enumerate}
\item 方法一：分子分母同除以 $\cos^2 x$，利用 $t=\tan x$ 代换：
\begin{align*}
\int \frac{\mathrm{d}x}{a^2\cos^2 x+b^2\sin^2 x} &= \int \frac{\sec^2 x \mathrm{d}x}{a^2+b^2\tan^2 x} = \int \frac{\mathrm{d}(\tan x)}{a^2+b^2\tan^2 x} \\
&= \frac{1}{b^2}\int \frac{\mathrm{d}(\tan x)}{(a/b)^2+\tan^2 x} = \frac{1}{b^2}\frac{b}{a}\arctan\left(\frac{b}{a}\tan x\right)+C \\
&= \frac{1}{ab}\arctan\left(\frac{b}{a}\tan x\right)+C.
\end{align*}
\item 方法二：记 $p=\tan x, q=\sec x$。利用底座流形恒等式 $q^2=1+p^2$ 结合微元关系 $\mathrm{d}p=q^2\mathrm{d}x \implies \mathrm{d}x=\frac{\mathrm{d}p}{q^2}$ 实现非暴力解构：
\begin{align*}
\int \frac{\mathrm{d}x}{a^2\cos^2 x+b^2\sin^2 x} &= \int \frac{\frac{\mathrm{d}p}{q^2}}{a^2\frac{1}{q^2} + b^2\frac{p^2}{q^2}} = \int \frac{\mathrm{d}p}{a^2+b^2 p^2} \\
&= \frac{1}{ab}\arctan\left(\frac{b}{a}p\right)+C = \frac{1}{ab}\arctan\left(\frac{b}{a}\tan x\right)+C.
\end{align*}
\end{enumerate}
\end{solution}

\begin{exercise}[含有三角函数的积分 - 106]
$\displaystyle \int \frac{\mathrm{d}x}{a^2\cos^2 x-b^2\sin^2 x}$
\end{exercise}
\begin{solution}
\begin{enumerate}
\item 方法一：分子分母同除以 $\cos^2 x$，利用 $t=\tan x$ 代换：
\begin{align*}
\int \frac{\mathrm{d}x}{a^2\cos^2 x-b^2\sin^2 x} &= \int \frac{\sec^2 x \mathrm{d}x}{a^2-b^2\tan^2 x} = \int \frac{\mathrm{d}(\tan x)}{a^2-b^2\tan^2 x} \\
&= \frac{1}{b^2}\int \frac{\mathrm{d}(\tan x)}{(a/b)^2-\tan^2 x} = \frac{1}{b^2}\frac{b}{2a}\ln\left|\frac{\frac{a}{b}+\tan x}{\frac{a}{b}-\tan x}\right|+C \\
&= \frac{1}{2ab}\ln\left|\frac{a+b\tan x}{a-b\tan x}\right|+C.
\end{align*}
\item 方法二：记 $p=\tan x, q=\sec x$ 且 $\mathrm{d}x=\frac{\mathrm{d}p}{q^2}$。同法直接将根基还原至有理状态下：
\begin{align*}
\int \frac{\frac{\mathrm{d}p}{q^2}}{a^2\frac{1}{q^2} - b^2\frac{p^2}{q^2}} &= \int \frac{\mathrm{d}p}{a^2-b^2 p^2} = \frac{1}{2a}\int \left( \frac{1}{a-bp} + \frac{1}{a+bp} \right)\frac{\mathrm{d}p}{b} \\
&= \frac{1}{2ab}\ln\left|\frac{a+bp}{a-bp}\right|+C = \frac{1}{2ab}\ln\left|\frac{b\tan x+a}{-b\tan x+a}\right|+C.
\end{align*}
\end{enumerate}
\end{solution}

\begin{exercise}[含有三角函数的积分 - 110]
$\displaystyle \int x\sin ax\mathrm{d}x$
\end{exercise}
\begin{solution}
\begin{enumerate}
\item 方法一：分部积分法 $u=x, \mathrm{d}v=\sin ax\mathrm{d}x$：
\begin{align*}
\int x\sin ax\mathrm{d}x &= \int x\mathrm{d}\left(-\frac{1}{a}\cos ax\right) = -\frac{x}{a}\cos ax - \int \left(-\frac{1}{a}\cos ax\right)\mathrm{d}x \\
&= -\frac{x}{a}\cos ax + \frac{1}{a}\int \cos ax \mathrm{d}x = -\frac{x}{a}\cos ax + \frac{1}{a^2}\sin ax+C.
\end{align*}
\item 方法二：严格算子代数法 (Operator Algebra)。令微分算子 $D = \frac{\mathrm{d}}{\mathrm{d}x}$，规避所有积分过程中的“猜测”与反向试错组合。利用位移定理 $D^{-1}(e^{iax} V(x)) = e^{iax}(D+ia)^{-1}V(x)$ 处理含欧拉指数的积分 $I = \text{Im}[D^{-1}(x e^{iax})]$：
\begin{align*}
(D+ia)^{-1}x &= \frac{1}{ia\left(1 + \frac{\mathrm{d}}{ia}\right)}x = \left(\frac{1}{ia} - \frac{\mathrm{d}}{(ia)^2} + \dots\right)x = \left(-i\frac{1}{a} + \frac{\mathrm{d}}{a^2}\right)x = -i\frac{x}{a} + \frac{1}{a^2} \\
D^{-1}(x e^{iax}) &= e^{iax}\left(\frac{1}{a^2} - i\frac{x}{a}\right) = (\cos ax + i\sin ax)\left(\frac{1}{a^2} - i\frac{x}{a}\right) \\
&= \left(\frac{1}{a^2}\cos ax + \frac{x}{a}\sin ax\right) + i\left(\frac{1}{a^2}\sin ax - \frac{x}{a}\cos ax\right) \\
\text{取虚部得} \quad I &= -\frac{x}{a}\cos ax + \frac{1}{a^2}\sin ax+C.
\end{align*}
\end{enumerate}
\end{solution}

\begin{exercise}[含有三角函数的积分 - 111]
$\displaystyle \int x\cos ax\mathrm{d}x$
\end{exercise}
\begin{solution}
\begin{enumerate}
\item 方法一：分部积分法 $u=x, \mathrm{d}v=\cos ax\mathrm{d}x$：
\[
\int x\cos ax\mathrm{d}x
= \int x\mathrm{d}\left(\frac{1}{a}\sin ax\right)
= \frac{x}{a}\sin ax + \frac{1}{a^2}\cos ax+C.
\]
\item 方法二：严格算子代数法。直接取上一题微积分算子逆运算结果实部结构：
\begin{align*}
\int x\cos ax\mathrm{d}x &= \text{Re}[D^{-1}(x e^{iax})] = \text{Re}\left[(\cos ax + i\sin ax)\left(\frac{1}{a^2} - i\frac{x}{a}\right)\right] \\
&= \frac{x}{a}\sin ax + \frac{1}{a^2}\cos ax+C.
\end{align*}
\end{enumerate}
\end{solution}

\begin{exercise}[含有三角函数的积分 - 112]
$\displaystyle \int x^n\sin ax\mathrm{d}x$
\end{exercise}
\begin{solution}
\begin{enumerate}
\item 方法一：分部积分法递推降解幂次：
\begin{align*}
\int x^n\sin ax\mathrm{d}x &= \int x^n \mathrm{d}\left(-\frac{1}{a}\cos ax\right) \\
&= -\frac{1}{a}x^n\cos ax - \int \left(-\frac{1}{a}\cos ax\right)n x^{n-1}\mathrm{d}x \\
&= -\frac{1}{a}x^n\cos ax + \frac{n}{a}\int x^{n-1}\cos ax\mathrm{d}x.
\end{align*}
\item 方法二：记 $p=\sin ax, q=\cos ax$。通过交叉微元约束 $\mathrm{d}q = -a\sin ax\mathrm{d}x = -ap\mathrm{d}x$，得出代换关系 $p\mathrm{d}x = -\frac{\mathrm{d}q}{a}$，纯粹用代数结构平推而避免记忆公式模型：
\begin{align*}
\int x^n p\mathrm{d}x &= \int x^n \left( -\frac{\mathrm{d}q}{a} \right) = -\frac{1}{a}\int x^n \mathrm{d}q \\
&= -\frac{1}{a}\left(x^n q - \int q\mathrm{d}(x^n)\right) = -\frac{1}{a}x^n q + \frac{n}{a}\int x^{n-1}q\mathrm{d}x \\
&= -\frac{1}{a}x^n\cos ax+\frac{n}{a}\int x^{n-1}\cos ax\mathrm{d}x.
\end{align*}
\end{enumerate}
\end{solution}

\begin{exercise}[含有三角函数的积分 - 113]
$\displaystyle \int x^n\cos ax\mathrm{d}x$
\end{exercise}
\begin{solution}
\begin{enumerate}
\item 方法一：分部积分法递推：
\begin{align*}
\int x^n\cos ax\mathrm{d}x &= \int x^n \mathrm{d}\left(\frac{1}{a}\sin ax\right) \\
&= \frac{1}{a}x^n\sin ax - \int \frac{1}{a}\sin ax\cdot n x^{n-1}\mathrm{d}x \\
&= \frac{1}{a}x^n\sin ax - \frac{n}{a}\int x^{n-1}\sin ax\mathrm{d}x.
\end{align*}
\item 方法二：记 $p=\sin ax, q=\cos ax$。利用交叉微元约束 $\mathrm{d}p = a\cos ax\mathrm{d}x = aq\mathrm{d}x$ 推出 $q\mathrm{d}x = \frac{\mathrm{d}p}{a}$，以对称操作重构该过程：
\begin{align*}
\int x^n q\mathrm{d}x &= \int x^n \left(\frac{\mathrm{d}p}{a}\right) = \frac{1}{a}\int x^n \mathrm{d}p \\
&= \frac{1}{a}\left(x^n p - \int p\mathrm{d}(x^n)\right) = \frac{1}{a}x^n p - \frac{n}{a}\int x^{n-1}p\mathrm{d}x \\
&= \frac{1}{a}x^n\sin ax-\frac{n}{a}\int x^{n-1}\sin ax\mathrm{d}x.
\end{align*}
\end{enumerate}
\end{solution}

\subsection{第十三节习题：\texorpdfstring{含有反三角函数的积分 $(a>0)$}{含有反三角函数的积分 (a>0)}}
\begin{exercise}[含有反三角函数的积分 - 114]
$\displaystyle \int \arcsin\frac{x}{a}\mathrm{d}x$
\end{exercise}
\begin{solution}
\begin{enumerate}
\item 方法一：利用分部积分法，将 $\mathrm{d}x$ 看作 $\mathrm{d}(x)$：
\begin{align*}
\int \arcsin\frac{x}{a}\mathrm{d}x &= x\arcsin\frac{x}{a} - \int x\mathrm{d}\left(\arcsin\frac{x}{a}\right) \\
&= x\arcsin\frac{x}{a} - \int x\cdot\frac{1}{\sqrt{1-(x/a)^2}}\cdot\frac{1}{a}\mathrm{d}x \\
&= x\arcsin\frac{x}{a} - \int \frac{x}{\sqrt{a^2-x^2}}\mathrm{d}x \\
&= x\arcsin\frac{x}{a} + \frac{1}{2}\int (a^2-x^2)^{-\frac{1}{2}}\mathrm{d}(a^2-x^2) \\
&= x\arcsin\frac{x}{a} + \sqrt{a^2-x^2}+C.
\end{align*}
\item 方法二：记反三角函数的底层代数流形为 $y=\sqrt{a^2-x^2}$。此时构成实圆底座 $x^2+y^2=a^2$，蕴含微分关系 $x\mathrm{d}x=-y\mathrm{d}y$。
我们提取反三角微元 $\mathrm{d}\left(\arcsin\frac{x}{a}\right) = \frac{\mathrm{d}x}{\sqrt{a^2-x^2}} = \frac{\mathrm{d}x}{y}$。
在双元结构下直接实施分部积分与降维：
\begin{align*}
\int \arcsin\frac{x}{a}\mathrm{d}x &= x\arcsin\frac{x}{a} - \int x\left(\frac{\mathrm{d}x}{y}\right) = x\arcsin\frac{x}{a} - \int \frac{x\mathrm{d}x}{y} \\
&= x\arcsin\frac{x}{a} - \int \frac{-y\mathrm{d}y}{y} = x\arcsin\frac{x}{a} + \int \mathrm{d}y \\
&= x\arcsin\frac{x}{a} + y + C = x\arcsin\frac{x}{a} + \sqrt{a^2-x^2}+C.
\end{align*}
\end{enumerate}
\end{solution}

\begin{exercise}[含有反三角函数的积分 - 115]
$\displaystyle \int x\arcsin\frac{x}{a}\mathrm{d}x$
\end{exercise}
\begin{solution}
\begin{enumerate}
\item 方法一：分部积分法 $u=\arcsin\frac{x}{a}, \mathrm{d}v=x\mathrm{d}x$：
\begin{align*}
\int x\arcsin\frac{x}{a}\mathrm{d}x &= \int \arcsin\frac{x}{a}\mathrm{d}\left(\frac{x^2}{2}\right) = \frac{x^2}{2}\arcsin\frac{x}{a} - \frac{1}{2}\int \frac{x^2}{\sqrt{a^2-x^2}}\mathrm{d}x \\
&= \frac{x^2}{2}\arcsin\frac{x}{a} - \frac{1}{2}\int \frac{a^2-(a^2-x^2)}{\sqrt{a^2-x^2}}\mathrm{d}x \\
&= \frac{x^2}{2}\arcsin\frac{x}{a} - \frac{a^2}{2}\int \frac{\mathrm{d}x}{\sqrt{a^2-x^2}} + \frac{1}{2}\int \sqrt{a^2-x^2}\mathrm{d}x \\
&= \frac{x^2}{2}\arcsin\frac{x}{a} - \frac{a^2}{2}\arcsin\frac{x}{a} + \frac{1}{2}\left( \frac{x}{2}\sqrt{a^2-x^2} + \frac{a^2}{2}\arcsin\frac{x}{a} \right) + C \\
&= \left(\frac{x^2}{2}-\frac{a^2}{4}\right)\arcsin\frac{x}{a} + \frac{x}{4}\sqrt{a^2-x^2}+C.
\end{align*}
\item 方法二：记 $y=\sqrt{a^2-x^2}$。利用底层关系 $x^2=a^2-y^2$ 以及反三角微元 $\mathrm{d}\left(\arcsin\frac{x}{a}\right) = \frac{\mathrm{d}x}{y}$：
\begin{align*}
\int x\arcsin\frac{x}{a}\mathrm{d}x &= \frac{x^2}{2}\arcsin\frac{x}{a} - \frac{1}{2}\int x^2\left(\frac{\mathrm{d}x}{y}\right) = \frac{x^2}{2}\arcsin\frac{x}{a} - \frac{1}{2}\int \frac{a^2-y^2}{y}\mathrm{d}x \\
&= \frac{x^2}{2}\arcsin\frac{x}{a} - \frac{a^2}{2}\int \frac{\mathrm{d}x}{y} + \frac{1}{2}\int y\mathrm{d}x \\
&= \frac{x^2}{2}\arcsin\frac{x}{a} - \frac{a^2}{2}\arcsin\frac{x}{a} + \frac{1}{2}\left( \frac{1}{2}xy + \frac{a^2}{2}\int \frac{\mathrm{d}x}{y} \right) \\
&= \left(\frac{x^2}{2}-\frac{a^2}{2}+\frac{a^2}{4}\right)\arcsin\frac{x}{a} + \frac{1}{4}xy+C \\
&= \left(\frac{x^2}{2}-\frac{a^2}{4}\right)\arcsin\frac{x}{a} + \frac{x}{4}\sqrt{a^2-x^2}+C.
\end{align*}
\end{enumerate}
\end{solution}

\begin{exercise}[含有反三角函数的积分 - 116]
$\displaystyle \int x^2\arcsin\frac{x}{a}\mathrm{d}x$
\end{exercise}
\begin{solution}
\begin{enumerate}
\item 方法一：分部积分法 $u=\arcsin\frac{x}{a}, \mathrm{d}v=x^2\mathrm{d}x$：
\begin{align*}
\int x^2\arcsin\frac{x}{a}\mathrm{d}x &= \int \arcsin\frac{x}{a}\mathrm{d}\left(\frac{x^3}{3}\right) = \frac{x^3}{3}\arcsin\frac{x}{a} - \frac{1}{3}\int \frac{x^3}{\sqrt{a^2-x^2}}\mathrm{d}x \\
&= \frac{x^3}{3}\arcsin\frac{x}{a} + \frac{1}{6}\int x^2 (a^2-x^2)^{-\frac{1}{2}}\mathrm{d}(a^2-x^2) \\
&= \frac{x^3}{3}\arcsin\frac{x}{a} + \frac{1}{3}\int x^2\mathrm{d}\left(\sqrt{a^2-x^2}\right) \\
&= \frac{x^3}{3}\arcsin\frac{x}{a} + \frac{1}{3}\left( x^2\sqrt{a^2-x^2} - \int \sqrt{a^2-x^2}(2x)\mathrm{d}x \right) \\
&= \frac{x^3}{3}\arcsin\frac{x}{a} + \frac{1}{3}x^2\sqrt{a^2-x^2} + \frac{1}{3}\int (a^2-x^2)^{\frac{1}{2}}\mathrm{d}(a^2-x^2) \\
&= \frac{x^3}{3}\arcsin\frac{x}{a} + \frac{1}{3}x^2\sqrt{a^2-x^2} + \frac{2}{9}(a^2-x^2)^{\frac{3}{2}}+C \\
&= \frac{x^3}{3}\arcsin\frac{x}{a} + \frac{1}{9}\sqrt{a^2-x^2}\left(3x^2+2(a^2-x^2)\right)+C \\
&= \frac{x^3}{3}\arcsin\frac{x}{a} + \frac{1}{9}(x^2+2a^2)\sqrt{a^2-x^2}+C.
\end{align*}
\item 方法二：记 $y=\sqrt{a^2-x^2}$。利用 $x\mathrm{d}x=-y\mathrm{d}y$，即 $\frac{x}{y}\mathrm{d}x=-\mathrm{d}y$。在双元空间内直接积分重构，化简步骤降维为纯多项式计算：
\begin{align*}
\int x^2\arcsin\frac{x}{a}\mathrm{d}x &= \frac{x^3}{3}\arcsin\frac{x}{a} - \frac{1}{3}\int x^3\left(\frac{\mathrm{d}x}{y}\right) = \frac{x^3}{3}\arcsin\frac{x}{a} - \frac{1}{3}\int x^2\left(\frac{x\mathrm{d}x}{y}\right) \\
&= \frac{x^3}{3}\arcsin\frac{x}{a} - \frac{1}{3}\int (a^2-y^2)(-\mathrm{d}y) = \frac{x^3}{3}\arcsin\frac{x}{a} + \frac{1}{3}\int (a^2-y^2)\mathrm{d}y \\
&= \frac{x^3}{3}\arcsin\frac{x}{a} + \frac{1}{3}\left( a^2y - \frac{1}{3}y^3 \right)+C = \frac{x^3}{3}\arcsin\frac{x}{a} + \frac{1}{9}y(3a^2-y^2)+C \\
&= \frac{x^3}{3}\arcsin\frac{x}{a} + \frac{1}{9}\sqrt{a^2-x^2}\left(3a^2-(a^2-x^2)\right)+C \\
&= \frac{x^3}{3}\arcsin\frac{x}{a} + \frac{1}{9}(x^2+2a^2)\sqrt{a^2-x^2}+C.
\end{align*}
\end{enumerate}
\end{solution}

\begin{exercise}[含有反三角函数的积分 - 117]
$\displaystyle \int \arccos\frac{x}{a}\mathrm{d}x$
\end{exercise}
\begin{solution}
\begin{enumerate}
\item 方法一：利用分部积分法。注意 $\mathrm{d}\left(\arccos\frac{x}{a}\right) = -\frac{1}{\sqrt{a^2-x^2}}\mathrm{d}x$：
\begin{align*}
\int \arccos\frac{x}{a}\mathrm{d}x &= x\arccos\frac{x}{a} - \int x\mathrm{d}\left(\arccos\frac{x}{a}\right) \\
&= x\arccos\frac{x}{a} - \int x\left( -\frac{1}{\sqrt{a^2-x^2}} \right)\mathrm{d}x = x\arccos\frac{x}{a} + \int \frac{x}{\sqrt{a^2-x^2}}\mathrm{d}x \\
&= x\arccos\frac{x}{a} - \frac{1}{2}\int (a^2-x^2)^{-\frac{1}{2}}\mathrm{d}(a^2-x^2) \\
&= x\arccos\frac{x}{a} - \sqrt{a^2-x^2}+C.
\end{align*}
\item 方法二：记底座流形 $y=\sqrt{a^2-x^2}$。微分关系 $x\mathrm{d}x=-y\mathrm{d}y$。
提取微元 $\mathrm{d}\left(\arccos\frac{x}{a}\right) = -\frac{\mathrm{d}x}{y}$。在双元结构下操作：
\begin{align*}
\int \arccos\frac{x}{a}\mathrm{d}x &= x\arccos\frac{x}{a} - \int x\left(-\frac{\mathrm{d}x}{y}\right) = x\arccos\frac{x}{a} + \int \frac{x\mathrm{d}x}{y} \\
&= x\arccos\frac{x}{a} + \int \frac{-y\mathrm{d}y}{y} = x\arccos\frac{x}{a} - \int \mathrm{d}y \\
&= x\arccos\frac{x}{a} - y + C = x\arccos\frac{x}{a} - \sqrt{a^2-x^2}+C.
\end{align*}
\end{enumerate}
\end{solution}

\begin{exercise}[含有反三角函数的积分 - 118]
$\displaystyle \int x\arccos\frac{x}{a}\mathrm{d}x$
\end{exercise}
\begin{solution}
\begin{enumerate}
\item 方法一：分部积分法：
\begin{align*}
\int x\arccos\frac{x}{a}\mathrm{d}x &= \frac{x^2}{2}\arccos\frac{x}{a} - \frac{1}{2}\int x^2\left(-\frac{1}{\sqrt{a^2-x^2}}\right)\mathrm{d}x = \frac{x^2}{2}\arccos\frac{x}{a} + \frac{1}{2}\int \frac{x^2}{\sqrt{a^2-x^2}}\mathrm{d}x \\
&= \frac{x^2}{2}\arccos\frac{x}{a} + \frac{1}{2}\int \frac{a^2-(a^2-x^2)}{\sqrt{a^2-x^2}}\mathrm{d}x \\
&= \frac{x^2}{2}\arccos\frac{x}{a} + \frac{a^2}{2}\int \frac{\mathrm{d}x}{\sqrt{a^2-x^2}} - \frac{1}{2}\int \sqrt{a^2-x^2}\mathrm{d}x \\
&= \frac{x^2}{2}\arccos\frac{x}{a} - \frac{a^2}{2}\arccos\frac{x}{a} - \frac{1}{2}\left( \frac{x}{2}\sqrt{a^2-x^2} - \frac{a^2}{2}\arccos\frac{x}{a} \right) + C \\
&= \left(\frac{x^2}{2}-\frac{a^2}{4}\right)\arccos\frac{x}{a} - \frac{x}{4}\sqrt{a^2-x^2}+C.
\end{align*}
*(注：这里使用了基本公式 $\int \frac{\mathrm{d}x}{\sqrt{a^2-x^2}} = -\arccos\frac{x}{a}$，保持函数族一致性)*
\item 方法二：记 $y=\sqrt{a^2-x^2}$。反三角微元 $\mathrm{d}\left(\arccos\frac{x}{a}\right) = -\frac{\mathrm{d}x}{y}$。已知 $\int \frac{\mathrm{d}x}{y} = -\arccos\frac{x}{a}$ 及面积微元 $\int y\mathrm{d}x = \frac{1}{2}xy + \frac{a^2}{2}\int\frac{\mathrm{d}x}{y}$。
\begin{align*}
\int x\arccos\frac{x}{a}\mathrm{d}x &= \frac{x^2}{2}\arccos\frac{x}{a} - \frac{1}{2}\int x^2\left(-\frac{\mathrm{d}x}{y}\right) = \frac{x^2}{2}\arccos\frac{x}{a} + \frac{1}{2}\int \frac{a^2-y^2}{y}\mathrm{d}x \\
&= \frac{x^2}{2}\arccos\frac{x}{a} + \frac{a^2}{2}\int \frac{\mathrm{d}x}{y} - \frac{1}{2}\int y\mathrm{d}x \\
&= \frac{x^2}{2}\arccos\frac{x}{a} + \frac{a^2}{2}\left(-\arccos\frac{x}{a}\right) - \frac{1}{2}\left( \frac{1}{2}xy + \frac{a^2}{2}\left(-\arccos\frac{x}{a}\right) \right) + C \\
&= \left(\frac{x^2}{2}-\frac{a^2}{2}+\frac{a^2}{4}\right)\arccos\frac{x}{a} - \frac{1}{4}xy+C \\
&= \left(\frac{x^2}{2}-\frac{a^2}{4}\right)\arccos\frac{x}{a} - \frac{x}{4}\sqrt{a^2-x^2}+C.
\end{align*}
\end{enumerate}
\end{solution}

\begin{exercise}[含有反三角函数的积分 - 119]
$\displaystyle \int x^2\arccos\frac{x}{a}\mathrm{d}x$
\end{exercise}
\begin{solution}
\begin{enumerate}
\item 方法一：分部积分法：
\begin{align*}
\int x^2\arccos\frac{x}{a}\mathrm{d}x &= \frac{x^3}{3}\arccos\frac{x}{a} - \frac{1}{3}\int x^3\left(-\frac{1}{\sqrt{a^2-x^2}}\right)\mathrm{d}x = \frac{x^3}{3}\arccos\frac{x}{a} + \frac{1}{3}\int \frac{x^3}{\sqrt{a^2-x^2}}\mathrm{d}x \\
&= \frac{x^3}{3}\arccos\frac{x}{a} - \frac{1}{6}\int x^2 (a^2-x^2)^{-\frac{1}{2}}\mathrm{d}(a^2-x^2) \\
&= \frac{x^3}{3}\arccos\frac{x}{a} - \frac{1}{3}\int x^2\mathrm{d}\left(\sqrt{a^2-x^2}\right) \\
&= \frac{x^3}{3}\arccos\frac{x}{a} - \frac{1}{3}\left( x^2\sqrt{a^2-x^2} - \int \sqrt{a^2-x^2}(2x)\mathrm{d}x \right) \\
&= \frac{x^3}{3}\arccos\frac{x}{a} - \frac{1}{3}x^2\sqrt{a^2-x^2} - \frac{1}{3}\int (a^2-x^2)^{\frac{1}{2}}\mathrm{d}(a^2-x^2) \\
&= \frac{x^3}{3}\arccos\frac{x}{a} - \frac{1}{3}x^2\sqrt{a^2-x^2} - \frac{2}{9}(a^2-x^2)^{\frac{3}{2}}+C \\
&= \frac{x^3}{3}\arccos\frac{x}{a} - \frac{1}{9}\sqrt{a^2-x^2}\left(3x^2+2(a^2-x^2)\right)+C \\
&= \frac{x^3}{3}\arccos\frac{x}{a} - \frac{1}{9}(x^2+2a^2)\sqrt{a^2-x^2}+C.
\end{align*}
\item 方法二：记 $y=\sqrt{a^2-x^2}$。利用反向微元 $\mathrm{d}\left(\arccos\frac{x}{a}\right) = -\frac{\mathrm{d}x}{y}$ 及转换规则 $\frac{x}{y}\mathrm{d}x=-\mathrm{d}y$：
\begin{align*}
\int x^2\arccos\frac{x}{a}\mathrm{d}x &= \frac{x^3}{3}\arccos\frac{x}{a} - \frac{1}{3}\int x^3\left(-\frac{\mathrm{d}x}{y}\right) = \frac{x^3}{3}\arccos\frac{x}{a} + \frac{1}{3}\int x^2\left(\frac{x\mathrm{d}x}{y}\right) \\
&= \frac{x^3}{3}\arccos\frac{x}{a} + \frac{1}{3}\int (a^2-y^2)(-\mathrm{d}y) = \frac{x^3}{3}\arccos\frac{x}{a} - \frac{1}{3}\int (a^2-y^2)\mathrm{d}y \\
&= \frac{x^3}{3}\arccos\frac{x}{a} - \frac{1}{3}\left( a^2y - \frac{1}{3}y^3 \right)+C = \frac{x^3}{3}\arccos\frac{x}{a} - \frac{1}{9}y(3a^2-y^2)+C \\
&= \frac{x^3}{3}\arccos\frac{x}{a} - \frac{1}{9}\sqrt{a^2-x^2}\left(3a^2-(a^2-x^2)\right)+C \\
&= \frac{x^3}{3}\arccos\frac{x}{a} - \frac{1}{9}(x^2+2a^2)\sqrt{a^2-x^2}+C.
\end{align*}
\end{enumerate}
\end{solution}

\begin{exercise}[含有反三角函数的积分 - 120]
$\displaystyle \int \arctan\frac{x}{a}\mathrm{d}x$
\end{exercise}
\begin{solution}
\begin{enumerate}
\item 方法一：利用分部积分法：
\begin{align*}
\int \arctan\frac{x}{a}\mathrm{d}x &= x\arctan\frac{x}{a} - \int x\mathrm{d}\left(\arctan\frac{x}{a}\right) \\
&= x\arctan\frac{x}{a} - \int x \frac{1}{1+(x/a)^2}\cdot\frac{1}{a}\mathrm{d}x \\
&= x\arctan\frac{x}{a} - \int \frac{ax}{x^2+a^2}\mathrm{d}x \\
&= x\arctan\frac{x}{a} - \frac{a}{2}\int \frac{\mathrm{d}(x^2+a^2)}{x^2+a^2} \\
&= x\arctan\frac{x}{a} - \frac{a}{2}\ln(x^2+a^2)+C.
\end{align*}
\item 方法二：记反正切函数的双元底座为 $y^2=x^2+a^2$。在此结构下 $2x\mathrm{d}x = 2y\mathrm{d}y \implies x\mathrm{d}x = y\mathrm{d}y$。
提取反三角微元 $\mathrm{d}\left(\arctan\frac{x}{a}\right) = \frac{a}{x^2+a^2}\mathrm{d}x = \frac{a}{y^2}\mathrm{d}x$。将其带入分部体系内实现降次消元：
\begin{align*}
\int \arctan\frac{x}{a}\mathrm{d}x &= x\arctan\frac{x}{a} - \int x\left(\frac{a}{y^2}\mathrm{d}x\right) = x\arctan\frac{x}{a} - a\int \frac{x\mathrm{d}x}{y^2} \\
&= x\arctan\frac{x}{a} - a\int \frac{y\mathrm{d}y}{y^2} = x\arctan\frac{x}{a} - a\int \frac{\mathrm{d}y}{y} \\
&= x\arctan\frac{x}{a} - a\ln|y|+C = x\arctan\frac{x}{a} - a\ln\sqrt{x^2+a^2}+C \\
&= x\arctan\frac{x}{a} - \frac{a}{2}\ln(a^2+x^2)+C.
\end{align*}
\end{enumerate}
\end{solution}

\begin{exercise}[含有反三角函数的积分 - 121]
$\displaystyle \int x\arctan\frac{x}{a}\mathrm{d}x$
\end{exercise}
\begin{solution}
\begin{enumerate}
\item 方法一：分部积分法 $u=\arctan\frac{x}{a}, \mathrm{d}v=x\mathrm{d}x$：
\begin{align*}
\int x\arctan\frac{x}{a}\mathrm{d}x &= \frac{x^2}{2}\arctan\frac{x}{a} - \frac{1}{2}\int x^2\left(\frac{a}{x^2+a^2}\right)\mathrm{d}x = \frac{x^2}{2}\arctan\frac{x}{a} - \frac{a}{2}\int \frac{x^2}{x^2+a^2}\mathrm{d}x \\
&= \frac{x^2}{2}\arctan\frac{x}{a} - \frac{a}{2}\int \frac{(x^2+a^2)-a^2}{x^2+a^2}\mathrm{d}x \\
&= \frac{x^2}{2}\arctan\frac{x}{a} - \frac{a}{2}\int \left(1 - \frac{a^2}{x^2+a^2}\right)\mathrm{d}x \\
&= \frac{x^2}{2}\arctan\frac{x}{a} - \frac{a}{2}\left( x - a\arctan\frac{x}{a} \right)+C \\
&= \left(\frac{x^2}{2}+\frac{a^2}{2}\right)\arctan\frac{x}{a} - \frac{ax}{2}+C \\
&= \frac{1}{2}(a^2+x^2)\arctan\frac{x}{a} - \frac{ax}{2}+C.
\end{align*}
\item 方法二：记底座流形 $y^2=x^2+a^2$。微元 $\mathrm{d}\left(\arctan\frac{x}{a}\right) = \frac{a}{y^2}\mathrm{d}x$。利用 $\int \frac{\mathrm{d}x}{y^2} = \frac{1}{a}\arctan\frac{x}{a}$ 与分式拆分重构：
\begin{align*}
\int x\arctan\frac{x}{a}\mathrm{d}x &= \frac{x^2}{2}\arctan\frac{x}{a} - \frac{1}{2}\int x^2\left(\frac{a}{y^2}\mathrm{d}x\right) = \frac{x^2}{2}\arctan\frac{x}{a} - \frac{a}{2}\int \frac{y^2-a^2}{y^2}\mathrm{d}x \\
&= \frac{x^2}{2}\arctan\frac{x}{a} - \frac{a}{2}\int \left(1 - \frac{a^2}{y^2}\right)\mathrm{d}x \\
&= \frac{x^2}{2}\arctan\frac{x}{a} - \frac{a}{2}x + \frac{a^3}{2}\int \frac{\mathrm{d}x}{y^2} \\
&= \frac{x^2}{2}\arctan\frac{x}{a} - \frac{ax}{2} + \frac{a^3}{2}\left(\frac{1}{a}\arctan\frac{x}{a}\right)+C \\
&= \frac{x^2}{2}\arctan\frac{x}{a} - \frac{ax}{2} + \frac{a^2}{2}\arctan\frac{x}{a}+C \\
&= \frac{1}{2}(a^2+x^2)\arctan\frac{x}{a} - \frac{ax}{2}+C.
\end{align*}
\end{enumerate}
\end{solution}

\begin{exercise}[含有反三角函数的积分 - 122]
$\displaystyle \int x^2\arctan\frac{x}{a}\mathrm{d}x$
\end{exercise}
\begin{solution}
\begin{enumerate}
\item 方法一：分部积分法 $u=\arctan\frac{x}{a}, \mathrm{d}v=x^2\mathrm{d}x$：
\begin{align*}
\int x^2\arctan\frac{x}{a}\mathrm{d}x &= \frac{x^3}{3}\arctan\frac{x}{a} - \frac{1}{3}\int x^3\left(\frac{a}{x^2+a^2}\right)\mathrm{d}x = \frac{x^3}{3}\arctan\frac{x}{a} - \frac{a}{3}\int \frac{x^3}{x^2+a^2}\mathrm{d}x \\
&= \frac{x^3}{3}\arctan\frac{x}{a} - \frac{a}{3}\int \frac{x(x^2+a^2)-a^2x}{x^2+a^2}\mathrm{d}x \\
&= \frac{x^3}{3}\arctan\frac{x}{a} - \frac{a}{3}\int \left( x - \frac{a^2x}{x^2+a^2} \right)\mathrm{d}x \\
&= \frac{x^3}{3}\arctan\frac{x}{a} - \frac{a}{3}\left( \frac{x^2}{2} - \frac{a^2}{2}\int \frac{\mathrm{d}(x^2+a^2)}{x^2+a^2} \right) \\
&= \frac{x^3}{3}\arctan\frac{x}{a} - \frac{ax^2}{6} + \frac{a^3}{6}\ln(x^2+a^2)+C.
\end{align*}
\item 方法二：建立底座 $y^2=x^2+a^2$，蕴含 $x\mathrm{d}x=y\mathrm{d}y$。在双元框架下可将多项式长除法转化为单变元幂次运算：
\begin{align*}
\int x^2\arctan\frac{x}{a}\mathrm{d}x &= \frac{x^3}{3}\arctan\frac{x}{a} - \frac{1}{3}\int x^3\left(\frac{a}{y^2}\mathrm{d}x\right) = \frac{x^3}{3}\arctan\frac{x}{a} - \frac{a}{3}\int x^2\frac{x\mathrm{d}x}{y^2} \\
&= \frac{x^3}{3}\arctan\frac{x}{a} - \frac{a}{3}\int (y^2-a^2)\frac{y\mathrm{d}y}{y^2} = \frac{x^3}{3}\arctan\frac{x}{a} - \frac{a}{3}\int \frac{y^2-a^2}{y}\mathrm{d}y \\
&= \frac{x^3}{3}\arctan\frac{x}{a} - \frac{a}{3}\int \left( y - \frac{a^2}{y} \right)\mathrm{d}y \\
&= \frac{x^3}{3}\arctan\frac{x}{a} - \frac{a}{3}\left( \frac{y^2}{2} - a^2\ln|y| \right)+C \\
&= \frac{x^3}{3}\arctan\frac{x}{a} - \frac{a(x^2+a^2)}{6} + \frac{a^3}{3}\ln\sqrt{x^2+a^2}+C \\
&= \frac{x^3}{3}\arctan\frac{x}{a} - \frac{ax^2}{6} + \frac{a^3}{6}\ln(x^2+a^2)+C.
\end{align*}
其中常数项已吸收入积分常数 $C$ 中。
\end{enumerate}
\end{solution}

\subsection{第十四节习题：含有指数函数的积分}
\begin{exercise}[含有指数函数的积分 - 123]
$\displaystyle \int a^x\mathrm{d}x$
\end{exercise}
\begin{solution}
\begin{enumerate}
\item 方法一：利用指数函数的基本积分公式直接求解：\(\displaystyle \int a^x\mathrm{d}x=\frac{1}{\ln a}a^x+C\)。
\item 方法二：算子代数法 (Operator Algebra)。记微分算子 $D = \frac{\mathrm{d}}{\mathrm{d}x}$，其逆运算可形式地记作积分算子 $D^{-1} = \int \cdot \mathrm{d}x$。
对于指数函数 $e^{kx}$，求导是乘以 $k$，那么积分就是除以 $k$，即 $D^{-1}(e^{kx}) = \frac{1}{k}e^{kx}$。
利用底数转换 $a^x = e^{x\ln a}$，可得 \(\displaystyle \int a^x\mathrm{d}x=D^{-1}(e^{x\ln a})=\frac{1}{\ln a}e^{x\ln a}+C=\frac{1}{\ln a}a^x+C\)。
\end{enumerate}
\end{solution}

\begin{exercise}[含有指数函数的积分 - 124]
$\displaystyle \int e^{ax}\mathrm{d}x$
\end{exercise}
\begin{solution}
\begin{enumerate}
\item 方法一：利用基本公式凑微分：\(\displaystyle \int e^{ax}\mathrm{d}x=\frac{1}{a}\int e^{ax}\mathrm{d}(ax)=\frac{1}{a}e^{ax}+C\)。
\item 方法二：算子代数法。如同上一题所述，把求积分视为求导算子 $D$ 的逆运算。因为 $D(e^{ax}) = a e^{ax}$，所以算子方程可以直接写出：\(\displaystyle \int e^{ax}\mathrm{d}x=D^{-1}(e^{ax})=\frac{1}{a}e^{ax}+C\)。
\end{enumerate}
\end{solution}

\begin{exercise}[含有指数函数的积分 - 125]
$\displaystyle \int xe^{ax}\mathrm{d}x$
\end{exercise}
\begin{solution}
\begin{enumerate}
\item 方法一：分部积分法 $u=x, \mathrm{d}v=e^{ax}\mathrm{d}x$：
\begin{align*}
\int xe^{ax}\mathrm{d}x &= \int x\mathrm{d}\left(\frac{1}{a}e^{ax}\right) = \frac{x}{a}e^{ax} - \int \frac{1}{a}e^{ax}\mathrm{d}x \\
&= \frac{x}{a}e^{ax} - \frac{1}{a^2}e^{ax}+C = \frac{1}{a^2}(ax-1)e^{ax}+C.
\end{align*}
\item 方法二：算子位移定理 (Shift Theorem)。这是算子代数法的核心灵魂，彻底规避繁琐的分部积分。
定理指出：$D^{-1}[e^{ax} V(x)] = e^{ax}(D+a)^{-1}V(x)$。（即将指数项提出算子外，代价是算子 $D$ 加上常数 $a$）。
对于 $(D+a)^{-1}$，我们可以像无穷等比数列一样利用泰勒展开：$(D+a)^{-1} = \frac{1}{a(1 + D/a)} = \frac{1}{a} - \frac{\mathrm{d}}{a^2} + \frac{D^2}{a^3} - \dots$
将其作用于多项式 $x$ 上（注意 $D(x)=1, D^2(x)=0$），过程变为纯粹的代数展开：
\begin{align*}
\int xe^{ax}\mathrm{d}x &= D^{-1}(e^{ax} \cdot x) = e^{ax}(D+a)^{-1}x \\
&= e^{ax}\left(\frac{1}{a} - \frac{\mathrm{d}}{a^2}\right)x = e^{ax}\left(\frac{x}{a} - \frac{1}{a^2}\right) + C \\
&= \frac{1}{a^2}(ax-1)e^{ax}+C.
\end{align*}
\end{enumerate}
\end{solution}

\begin{exercise}[含有指数函数的积分 - 126]
$\displaystyle \int x^n e^{ax}\mathrm{d}x$
\end{exercise}
\begin{solution}
\begin{enumerate}
\item 方法一：分部积分法降次：
\begin{align*}
\int x^n e^{ax}\mathrm{d}x &= \frac{1}{a}\int x^n \mathrm{d}(e^{ax}) = \frac{1}{a}x^n e^{ax} - \frac{1}{a}\int e^{ax}nx^{n-1}\mathrm{d}x \\
&= \frac{1}{a}x^n e^{ax}-\frac{n}{a}\int x^{n-1}e^{ax}\mathrm{d}x.
\end{align*}
\item 方法二：算子代数法。利用上一题讲授的“位移定理”，无需去凑分部积分的符号，直接写出递推的代数本质：
\begin{align*}
I_n &= D^{-1}(e^{ax} x^n) = e^{ax}(D+a)^{-1}x^n \\
&= e^{ax}\frac{1}{a}\left(1 + \frac{\mathrm{d}}{a}\right)^{-1}x^n = \frac{e^{ax}}{a}\left(1 - \frac{\mathrm{d}}{a+D}\right)x^n \\
&= \frac{e^{ax}}{a}x^n - \frac{1}{a}D\left[ e^{-ax} \cdot e^{ax}(D+a)^{-1}x^n \right] \quad \text{(这里反向使用位移定理)} \\
&= \frac{1}{a}x^n e^{ax} - \frac{1}{a} D(D^{-1}(x^n e^{ax})) \dots \text{实际上更简单的视角是直接展开算子：} \\
&= e^{ax}\left(\frac{1}{a} - \frac{\mathrm{d}}{a^2} + \frac{D^2}{a^3} - \dots\right)x^n = \frac{e^{ax}}{a}x^n - \frac{n}{a}\left[ e^{ax}\left(\frac{1}{a} - \frac{\mathrm{d}}{a^2} + \dots\right)x^{n-1} \right] \\
&= \frac{1}{a}x^n e^{ax}-\frac{n}{a}\int x^{n-1}e^{ax}\mathrm{d}x.
\end{align*}
\end{enumerate}
\end{solution}

\begin{exercise}[含有指数函数的积分 - 127]
$\displaystyle \int x a^x\mathrm{d}x$
\end{exercise}
\begin{solution}
\begin{enumerate}
\item 方法一：分部积分法 $u=x, \mathrm{d}v=a^x\mathrm{d}x$：
\[
\int x a^x\mathrm{d}x
= \int x\mathrm{d}\left(\frac{a^x}{\ln a}\right)
= \frac{x a^x}{\ln a} - \frac{a^x}{(\ln a)^2}+C.
\]
\item 方法二：算子位移定理。将 $a^x$ 视为 $e^{x\ln a}$，常数 $a$ 在算子中的角色就是 $\ln a$。完全套用算子的泰勒展开公式：
\begin{align*}
\int x e^{x\ln a}\mathrm{d}x &= D^{-1}(e^{x\ln a} \cdot x) = e^{x\ln a}(D+\ln a)^{-1}x \\
&= a^x\left(\frac{1}{\ln a} - \frac{\mathrm{d}}{(\ln a)^2}\right)x \\
&= a^x\left(\frac{x}{\ln a} - \frac{1}{\ln^2 a}\right)+C = \frac{x a^x}{\ln a}-\frac{a^x}{\ln^2 a}+C.
\end{align*}
\end{enumerate}
\end{solution}

\begin{exercise}[含有指数函数的积分 - 128]
$\displaystyle \int x^n a^x\mathrm{d}x$
\end{exercise}
\begin{solution}
\begin{enumerate}
\item 方法一：分部积分法降次：
\begin{align*}
\int x^n a^x\mathrm{d}x &= \frac{1}{\ln a}\int x^n \mathrm{d}(a^x) = \frac{x^n a^x}{\ln a} - \frac{1}{\ln a}\int a^x n x^{n-1}\mathrm{d}x \\
&= \frac{x^n a^x}{\ln a}-\frac{n}{\ln a}\int x^{n-1}a^x\mathrm{d}x.
\end{align*}
\item 方法二：算子代数法。同理将 $a^x$ 化为 $e^{x\ln a}$，这在高级积分理论中是处理非自然底数指数函数的统一方式。
算子推演直接给出多项式下降的本质结构：
\begin{align*}
\int x^n e^{x\ln a}\mathrm{d}x &= D^{-1}(e^{x\ln a}x^n) = e^{x\ln a}\left(\frac{1}{\ln a} - \frac{\mathrm{d}}{\ln^2 a} + \dots\right)x^n \\
&= \frac{x^n a^x}{\ln a} - \frac{n}{\ln a} D^{-1}(e^{x\ln a}x^{n-1}) \\
&= \frac{x^n a^x}{\ln a}-\frac{n}{\ln a}\int x^{n-1}a^x\mathrm{d}x.
\end{align*}
\end{enumerate}
\end{solution}

\begin{exercise}[含有指数函数的积分 - 129]
$\displaystyle \int e^{ax}\sin bx\mathrm{d}x$
\end{exercise}
\begin{solution}
\begin{enumerate}
\item 方法一：两次分部积分法，构造循环方程后求解（略去冗长的移项过程）。
\item 方法二：复指数双元法 (Complex Exponential Dual-Variable)。遇到指数与三角的乘积，高级做法是直接将其打包为复数域上的纯指数函数。
根据欧拉公式 $e^{i\theta} = \cos\theta + i\sin\theta$，显然 $e^{ax}\sin bx$ 是 $e^{(a+ib)x}$ 的虚部（$\text{Im}$）。
在复数域中，积分依然遵循最简单的算子规则 $D^{-1}(e^{kx}) = \frac{1}{k}e^{kx}$，根本不需要分部积分：
\begin{align*}
\int e^{ax}\sin bx\mathrm{d}x &= \text{Im}\left[ \int e^{(a+ib)x}\mathrm{d}x \right] = \text{Im}\left[ \frac{1}{a+ib}e^{(a+ib)x} \right] \\
&= \text{Im}\left[ \frac{a-ib}{a^2+b^2} e^{ax}(\cos bx + i\sin bx) \right]
\end{align*}
只需要将虚部（即含 $i$ 的项相乘的部分）提取出来：
\(\displaystyle \text{虚部}=\frac{e^{ax}}{a^2+b^2}\left(a\sin bx-b\cos bx\right)+C\)，
因此
\(\displaystyle \int e^{ax}\sin bx\mathrm{d}x=\frac{e^{ax}}{a^2+b^2}(a\sin bx-b\cos bx)+C\)。
\end{enumerate}
\end{solution}

\begin{exercise}[含有指数函数的积分 - 130]
$\displaystyle \int e^{ax}\cos bx\mathrm{d}x$
\end{exercise}
\begin{solution}
\begin{enumerate}
\item 方法一：两次分部积分法。
\item 方法二：复指数双元法。同上题，$e^{ax}\cos bx$ 恰好是复指数函数 $e^{(a+ib)x}$ 的实部（$\text{Re}$）。
直接复用上一题的代数展开结果，再提取实部即可：
\[
\int e^{ax}\cos bx\mathrm{d}x = \text{Re}\left[ \int e^{(a+ib)x}\mathrm{d}x \right] = \text{Re}\left[ \frac{a-ib}{a^2+b^2} e^{ax}(\cos bx + i\sin bx) \right].
\]
实部为 \(\displaystyle \frac{e^{ax}}{a^2+b^2}\left[a\cos bx-(-b)\sin bx\right]+C\)，因此
\(\displaystyle \int e^{ax}\cos bx\mathrm{d}x=\frac{e^{ax}}{a^2+b^2}(a\cos bx+b\sin bx)+C\)。
\end{enumerate}
\end{solution}

\begin{exercise}[含有指数函数的积分 - 131]
$\displaystyle \int e^{ax}\sin^n bx\mathrm{d}x$
\end{exercise}
\begin{solution}
\begin{enumerate}
\item 方法一：多次分部积分法并构造递推方程（过程极易出错）。
\item 方法二：微分代数系构造法。该方法的关键在于先对目标函数求导，在导数层面建立封闭的线性关系，再逆向得到积分递推式。
令目标结构 $U = e^{ax}\sin^n bx$，以及伴随结构 $V = e^{ax}\sin^{n-1}bx \cos bx$。
分别对它们求导（纯代数运算）：
1) $U' = a e^{ax}\sin^n bx + nb e^{ax}\sin^{n-1}bx \cos bx = aU + nbV$
2) $V' = a e^{ax}\sin^{n-1}bx \cos bx + b(n-1)e^{ax}\sin^{n-2}bx \cos^2 bx - be^{ax}\sin^n bx$
将 $\cos^2 bx = 1 - \sin^2 bx$ 代入 (2) 中化简：
$V' = aV + b(n-1)e^{ax}\sin^{n-2}bx - nb e^{ax}\sin^n bx = aV - nbU + b(n-1)e^{ax}\sin^{n-2}bx$

我们要解的是 $\int U \mathrm{d}x = I_n$。为了消去等式右边烦人的伴随项 $V$，我们将 (1) 式乘 $a$，(2) 式乘 $-nb$，然后相加：
$aU' - nbV' = (a^2 U + anb V) - (anb V - n^2b^2 U + nb^2(n-1)e^{ax}\sin^{n-2}bx)$
$aU' - nbV' = (a^2+n^2b^2)U - n(n-1)b^2 e^{ax}\sin^{n-2}bx$
现在，对这个极为干净的等式两边同时积分（将导数符号抹去，生成 $I_n$ 和 $I_{n-2}$）：
$aU - nbV = (a^2+n^2b^2)I_n - n(n-1)b^2 I_{n-2}$
移项并除以 $(a^2+n^2b^2)$，即可得到：
\begin{align*}
I_n &= \frac{aU - nbV}{a^2+n^2b^2} + \frac{n(n-1)b^2}{a^2+n^2b^2}I_{n-2} \\
&= \frac{e^{ax}\sin^{n-1}bx}{a^2+b^2n^2}(a\sin bx-nb\cos bx)+\frac{n(n-1)b^2}{a^2+b^2n^2}\int e^{ax}\sin^{n-2}bx\mathrm{d}x.
\end{align*}
\end{enumerate}
\end{solution}

\begin{exercise}[含有指数函数的积分 - 132]
$\displaystyle \int e^{ax}\cos^n bx\mathrm{d}x$
\end{exercise}
\begin{solution}
\begin{enumerate}
\item 方法一：两次分部积分构造递推。
\item 方法二：微分代数系构造法。与上题逻辑完全对称，构建封闭的求导空间：
令目标结构 $U = e^{ax}\cos^n bx$，伴随结构 $V = e^{ax}\cos^{n-1}bx \sin bx$。
求导并使用 $\sin^2 bx = 1 - \cos^2 bx$ 展开：
1) $U' = aU - nbV$
2) $V' = aV + b(n-1)e^{ax}\cos^{n-2}bx \sin^2 bx + be^{ax}\cos^n bx = aV - nbU + b(n-1)e^{ax}\cos^{n-2}bx$

消去伴随项 $V$：将 (1) 式乘 $a$，(2) 式乘 $nb$，然后相加：
$aU' + nbV' = (a^2 U - anb V) + (anb V - n^2b^2 U + nb^2(n-1)e^{ax}\cos^{n-2}bx)$
$aU' + nbV' = (a^2+n^2b^2)U + n(n-1)b^2 e^{ax}\cos^{n-2}bx$
对等式两边积分，剥离导数算子：
$aU + nbV = (a^2+n^2b^2)I_n + n(n-1)b^2 I_{n-2}$
整理即得递推公式：
\begin{align*}
I_n &= \frac{aU + nbV}{a^2+n^2b^2} + \frac{n(n-1)b^2}{a^2+n^2b^2}I_{n-2} \\
&= \frac{e^{ax}\cos^{n-1}bx}{a^2+b^2n^2}(a\cos bx+nb\sin bx)+\frac{n(n-1)b^2}{a^2+b^2n^2}\int e^{ax}\cos^{n-2}bx\mathrm{d}x.
\end{align*}
\end{enumerate}
\end{solution}

\subsection{第十五节习题：含有对数函数的积分}
\begin{exercise}[含有对数函数的积分 - 132(b)]
$\displaystyle \int \ln x\mathrm{d}x$
\end{exercise}
\begin{solution}
\begin{enumerate}
\item 方法一：分部积分法，令 $u=\ln x, \mathrm{d}v=\mathrm{d}x$：\(\displaystyle \int \ln x\mathrm{d}x=x\ln x-\int x\mathrm{d}(\ln x)=x\ln x-\int x\cdot \frac{1}{x}\mathrm{d}x=x\ln x-x+C\)。
\item 方法二：对数-指数空间映射法。令 $x = e^t$，则 $t = \ln x$，且 $\mathrm{d}x = e^t \mathrm{d}t$。原积分化为
\(\displaystyle \int \ln x\mathrm{d}x = \int t e^t \mathrm{d}t\)。
在 $t$ 空间作分部积分：\(\displaystyle \int te^t\,\mathrm{d}t=te^t-\int e^t\,\mathrm{d}t=e^t(t-1)+C\)。
将 $x = e^t,\ t = \ln x$ 代回，得到 \(\displaystyle x(\ln x - 1)+C = x\ln x - x+C\)。
\end{enumerate}
\end{solution}

\begin{exercise}[含有对数函数的积分 - 133]
$\displaystyle \int \frac{\mathrm{d}x}{x\ln x}$
\end{exercise}
\begin{solution}
\begin{enumerate}
\item 方法一：凑微分法：
\(\displaystyle \int \frac{\mathrm{d}x}{x\ln x} = \int \frac{1}{\ln x}\mathrm{d}(\ln x) = \ln|\ln x|+C\)。
\item 方法二：对数-指数空间映射法。令 $x=e^t$，则 $\mathrm{d}x = e^t\mathrm{d}t, \ln x = t$：
\(\displaystyle \int \frac{\mathrm{d}x}{x\ln x} = \int \frac{e^t \mathrm{d}t}{e^t \cdot t} = \int \frac{\mathrm{d}t}{t} = \ln|t|+C\)。
映射回原变量，得到 \(\displaystyle \ln|\ln x|+C\)。
这种变换虽然在本题中并非唯一选择，但它为处理含对数的积分提供了统一的变量转换框架。
\end{enumerate}
\end{solution}

\begin{exercise}[含有对数函数的积分 - 134]
$\displaystyle \int x^n\ln x\mathrm{d}x$
\end{exercise}
\begin{solution}
\begin{enumerate}
\item 方法一：先设 $n\neq -1$。取 $u=\ln x,\ \mathrm{d}v=x^n\mathrm{d}x$，由分部积分得
\begin{align*}
\int x^n\ln x\mathrm{d}x &= \frac{x^{n+1}}{n+1}\ln x - \int \frac{x^{n+1}}{n+1}\frac{1}{x}\mathrm{d}x = \frac{x^{n+1}}{n+1}\ln x - \frac{x^{n+1}}{(n+1)^2}+C \\
&= \frac{1}{n+1}x^{n+1}\left(\ln x-\frac{1}{n+1}\right)+C.
\end{align*}
若 \(n=-1\)，则 \(\displaystyle \int \frac{\ln x}{x}\,\mathrm{d}x=\frac12(\ln x)^2+C\)。
\item 方法二：对数-指数空间映射法。令 $x=e^t$，则 $\mathrm{d}x=e^t\mathrm{d}t,\ \ln x=t$。原积分化为 \(\displaystyle \int x^n\ln x\mathrm{d}x = \int e^{nt} t e^t \mathrm{d}t = \int t e^{(n+1)t}\mathrm{d}t\)。
当 $n\neq -1$ 时，在 $t$ 空间作分部积分，并在第二行代回 $t=\ln x,\ e^{(n+1)t}=x^{n+1}$：
\begin{align*}
\int t e^{(n+1)t}\mathrm{d}t
&= t\cdot \frac{e^{(n+1)t}}{n+1}-\int \frac{e^{(n+1)t}}{n+1}\,\mathrm{d}t
= \frac{t\,e^{(n+1)t}}{n+1}-\frac{e^{(n+1)t}}{(n+1)^2}+C,\\
\int x^n\ln x\mathrm{d}x
&= x^{n+1}\left( \frac{\ln x}{n+1} - \frac{1}{(n+1)^2} \right)+C
= \frac{1}{n+1}x^{n+1}\left(\ln x-\frac{1}{n+1}\right)+C.
\end{align*}
\end{enumerate}
\end{solution}

\begin{exercise}[含有对数函数的积分 - 135]
$\displaystyle \int (\ln x)^n\mathrm{d}x$
\end{exercise}
\begin{solution}
\begin{enumerate}
\item 方法一：分部积分法递推，令 $u=(\ln x)^n, \mathrm{d}v=\mathrm{d}x$，得
\[
\int (\ln x)^n\mathrm{d}x = x(\ln x)^n - \int x \cdot n(\ln x)^{n-1}\frac{1}{x}\mathrm{d}x = x(\ln x)^n-n\int (\ln x)^{n-1}\mathrm{d}x+C.
\]
\item 方法二：对数-指数空间映射法。令 $x=e^t,\ \mathrm{d}x=e^t\mathrm{d}t$，则原积分化为 \(\displaystyle I_n = \int t^n e^t \mathrm{d}t\)。
在 $t$ 空间作分部积分，取 $u=t^n,\ \mathrm{d}v=e^t\mathrm{d}t$，得 \(\displaystyle I_n=t^n e^t-n\int t^{n-1}e^t\,\mathrm{d}t\)。
再换回原变量，得到 \(\displaystyle I_n = x(\ln x)^n-n\int (\ln x)^{n-1}\mathrm{d}x+C\)。
\end{enumerate}
\end{solution}

\begin{exercise}[含有对数函数的积分 - 136]
$\displaystyle \int x^m(\ln x)^n\mathrm{d}x$
\end{exercise}
\begin{solution}
\begin{enumerate}
\item 方法一：先设 $m\neq -1$。取 $u=(\ln x)^n,\ \mathrm{d}v=x^m\mathrm{d}x$，则
\begin{align*}
\int x^m(\ln x)^n\mathrm{d}x &= \frac{x^{m+1}}{m+1}(\ln x)^n - \int \frac{x^{m+1}}{m+1}n(\ln x)^{n-1}\frac{1}{x}\mathrm{d}x \\
&= \frac{1}{m+1}x^{m+1}(\ln x)^n-\frac{n}{m+1}\int x^m(\ln x)^{n-1}\mathrm{d}x+C.
\end{align*}
若 \(m=-1\)，则 \(\displaystyle \int \frac{(\ln x)^n}{x}\,\mathrm{d}x=\int t^n\,\mathrm{d}t\)。因此，当 $n\neq -1$ 时，\(\displaystyle \int \frac{(\ln x)^n}{x}\,\mathrm{d}x=\frac{(\ln x)^{n+1}}{n+1}+C\)；当 $n=-1$ 时，\(\displaystyle \int \frac{\mathrm{d}x}{x\ln x}=\ln|\ln x|+C\)。
\item 方法二：在 $m\neq -1$ 的前提下，令 $x=e^t$，则原积分转化为指数函数与多项式的乘积：\(\displaystyle I_{m,n} = \int e^{mt} t^n e^t \mathrm{d}t = \int t^n e^{(m+1)t} \mathrm{d}t\)。
在 $t$ 空间作分部积分，得
\[
\int t^n e^{(m+1)t}\mathrm{d}t = \frac{e^{(m+1)t}}{m+1}t^n-\frac{n}{m+1}\int t^{n-1}e^{(m+1)t}\mathrm{d}t.
\]
将结果换回 $x$ 变量（注意 $e^{(m+1)t}=x^{m+1},\ t=\ln x$），便得到
\(\displaystyle I_{m,n} = \frac{1}{m+1}x^{m+1}(\ln x)^n-\frac{n}{m+1}\int x^m(\ln x)^{n-1}\mathrm{d}x+C\)。
\end{enumerate}
\end{solution}

\subsection{第十六节习题：含有双曲函数的积分}
\begin{exercise}[含有双曲函数的积分 - 137]
$\displaystyle \int \mathrm{sinh}x\mathrm{d}x$
\end{exercise}
\begin{solution}
\begin{enumerate}
\item 方法一：利用定义 $\sinh x = \frac{e^x-e^{-x}}{2}$ 积分即可得 $\cosh x+C$。
\item 方法二：双曲双元法 (Hyperbolic Dual-Variable)。
记 \(p=\cosh x,\ q=\sinh x\)。则 $p^2-q^2=1$，并且
\(\displaystyle \mathrm{d}p=q\,\mathrm{d}x,\qquad \mathrm{d}q=p\,\mathrm{d}x\)，因此
\(\displaystyle \int \mathrm{sinh}x\mathrm{d}x=\int q\mathrm{d}x=\int \mathrm{d}p=p+C=\mathrm{cosh}x+C\)。
\end{enumerate}
\end{solution}

\begin{exercise}[含有双曲函数的积分 - 138]
$\displaystyle \int \mathrm{cosh}x\mathrm{d}x$
\end{exercise}
\begin{solution}
\begin{enumerate}
\item 方法一：利用定义
\(\displaystyle \cosh x = \frac{e^x+e^{-x}}{2}\)，逐项积分便得
\(\displaystyle \int \cosh x\,\mathrm{d}x=\frac{1}{2}\int (e^x+e^{-x})\,\mathrm{d}x=\frac{1}{2}(e^x-e^{-x})+C=\sinh x+C\)。
\item 方法二：双曲双元法。记 $p = \cosh x, q = \sinh x$，调用微分约束 $\mathrm{d}q = p\mathrm{d}x$：
\(\displaystyle \int \mathrm{cosh}x\mathrm{d}x=\int p\mathrm{d}x=\int \mathrm{d}q=q+C=\mathrm{sinh}x+C\)。
\end{enumerate}
\end{solution}

\begin{exercise}[含有双曲函数的积分 - 139]
$\displaystyle \int \mathrm{tanh}x\mathrm{d}x$
\end{exercise}
\begin{solution}
\begin{enumerate}
\item 方法一：利用 $\tanh x = \frac{\sinh x}{\cosh x}$ 凑微分法积分：
\[
\int \frac{\mathrm{sinh}x}{\mathrm{cosh}x}\mathrm{d}x = \int \frac{\mathrm{d}(\mathrm{cosh}x)}{\mathrm{cosh}x} = \ln\mathrm{cosh}x+C.
\]
\item 方法二：双曲双元法。记 $p = \cosh x, q = \sinh x$。提取虚圆对数微元关系 $q\mathrm{d}x = \mathrm{d}p$：
\[
\int \mathrm{tanh}x\mathrm{d}x
= \int \frac{q}{p}\mathrm{d}x
= \int \frac{1}{p}(q\mathrm{d}x)
= \int \frac{\mathrm{d}p}{p}
= \ln|p|+C = \ln\mathrm{cosh}x+C.
\]
这里最后一步用到 $\cosh x \ge 1$ 恒正。
\end{enumerate}
\end{solution}

\begin{exercise}[含有双曲函数的积分 - 140]
$\displaystyle \int \mathrm{sinh}^2 x\mathrm{d}x$
\end{exercise}
\begin{solution}
\begin{enumerate}
\item 方法一：利用双曲半角公式 $\sinh^2 x = \frac{\cosh 2x - 1}{2}$ 展开后积分。
\item 方法二：双曲双元法。记 $p = \cosh x,\ q = \sinh x$，则 $p^2 - q^2 = 1$，且 $\mathrm{d}p = q\,\mathrm{d}x,\ \mathrm{d}q = p\,\mathrm{d}x$。对 $\int q^2\,\mathrm{d}x$ 作分部处理：
\[
\int q^2\mathrm{d}x = \int q(q\mathrm{d}x) = \int q \mathrm{d}p = pq - \int p\mathrm{d}q = pq - \int p(p\mathrm{d}x) = pq - \int p^2\mathrm{d}x.
\]
移项得 \(\displaystyle \int (p^2+q^2)\mathrm{d}x = pq\)。
再利用 $p^2 = 1 + q^2$，便有
\[
\int (1+2q^2)\mathrm{d}x = pq \implies \int q^2\mathrm{d}x
= \frac{1}{2}(pq - x) = -\frac{x}{2}+\frac{1}{4}\mathrm{sinh}2x+C.
\]
\end{enumerate}
\end{solution}

\begin{exercise}[含有双曲函数的积分 - 141]
$\displaystyle \int \mathrm{cosh}^2 x\mathrm{d}x$
\end{exercise}
\begin{solution}
\begin{enumerate}
\item 方法一：利用双曲半角公式 $\cosh^2 x = \frac{\cosh 2x + 1}{2}$ 展开后积分。
\item 方法二：沿用上题记号 $p=\cosh x,\ q=\sinh x$。由上题已经得到
\(\displaystyle \int (p^2+q^2)\mathrm{d}x = pq\)。
又因为 $q^2 = p^2 - 1$，所以
\[
\int (2p^2 - 1)\mathrm{d}x = pq \implies \int p^2\mathrm{d}x
= \frac{1}{2}(pq + x) = \frac{x}{2}+\frac{1}{4}\mathrm{sinh}2x+C.
\]
\end{enumerate}
\end{solution}

\subsection{第十七节习题：定积分}
\begin{exercise}[定积分 - 142]
$\displaystyle \int_{-\pi}^{\pi} \cos nx\mathrm{d}x = \int_{-\pi}^{\pi} \sin nx\mathrm{d}x$
\end{exercise}
\begin{solution}
\begin{enumerate}
\item 方法一：直接利用牛顿-莱布尼茨公式积分：
\begin{align*}
\int_{-\pi}^{\pi} \cos nx\mathrm{d}x &= \left[ \frac{1}{n}\sin nx \right]_{-\pi}^{\pi} = \frac{1}{n}(\sin n\pi - \sin(-n\pi)) = 0. \\
\int_{-\pi}^{\pi} \sin nx\mathrm{d}x &= \left[ -\frac{1}{n}\cos nx \right]_{-\pi}^{\pi} = -\frac{1}{n}(\cos n\pi - \cos(-n\pi)) = 0.
\end{align*}
\item 方法二：复指数双元法 (Complex Exponential Dual-Variable)。
\textbf{教学提示：}在复平面上，三角函数可以看作复指数函数的投影。根据欧拉公式 $e^{ix} = \cos x + i\sin x$，我们可以将 $\cos nx$ 和 $\sin nx$ 合并为一个复指数积分 $\int_{-\pi}^{\pi} e^{inx}\mathrm{d}x$。
对任何非零整数 $n$，复指数在完整周期 $[-\pi, \pi]$ 上的积分为 0：
\[
\int_{-\pi}^{\pi} e^{inx}\mathrm{d}x = \left[ \frac{1}{in}e^{inx} \right]_{-\pi}^{\pi} = \frac{1}{in}(e^{in\pi} - e^{-in\pi}) = \frac{1}{in}((-1)^n - (-1)^n) = 0.
\]
既然复数积分本身等于 0，它的实部与虚部也都等于 0，因此 \(\displaystyle \int_{-\pi}^{\pi} \cos nx\,\mathrm{d}x=0,\ \int_{-\pi}^{\pi} \sin nx\,\mathrm{d}x=0\)。
\end{enumerate}
\end{solution}

\begin{exercise}[定积分 - 143]
$\displaystyle \int_{-\pi}^{\pi} \cos mx\sin nx\mathrm{d}x$
\end{exercise}
\begin{solution}
\begin{enumerate}
\item 方法一：利用奇偶性。因为 $\cos mx$ 是偶函数，$\sin nx$ 是奇函数，偶函数乘奇函数等于奇函数。奇函数在对称区间 $[-\pi, \pi]$ 上的积分为 0。
\item 方法二：复指数双元法。
\textbf{教学提示：}如果遇到非对称区间或者更复杂的乘积，奇偶性失效时，可以改用复指数展开：
$\cos \theta = \frac{e^{i\theta} + e^{-i\theta}}{2}$，$\sin \theta = \frac{e^{i\theta} - e^{-i\theta}}{2i}$。
将三角乘法转化为复指数的代数展开：
\begin{align*}
\int_{-\pi}^{\pi} \cos mx\sin nx\mathrm{d}x &= \int_{-\pi}^{\pi} \left( \frac{e^{imx} + e^{-imx}}{2} \right) \left( \frac{e^{inx} - e^{-inx}}{2i} \right) \mathrm{d}x \\
&= \frac{1}{4i} \int_{-\pi}^{\pi} \left( e^{i(m+n)x} - e^{i(m-n)x} + e^{-i(m-n)x} - e^{-i(m+n)x} \right) \mathrm{d}x
\end{align*}
利用上题结论：只要指数上的整数 $k \neq 0$，$\int_{-\pi}^{\pi} e^{ikx}\mathrm{d}x = 0$。
由于 $m, n$ 为正整数，$m+n \neq 0$，前后两项积分都为 0。
即使 $m=n$，中间两项也相互抵消。
因此无论 $m, n$ 取何值，积分都为 0。
\end{enumerate}
\end{solution}

\begin{exercise}[定积分 - 144]
$\displaystyle \int_{-\pi}^{\pi} \cos mx\cos nx\mathrm{d}x$
\end{exercise}
\begin{solution}
\begin{enumerate}
\item 方法一：利用积化和差公式 $\cos \alpha \cos \beta = \frac{1}{2}[\cos(\alpha+\beta) + \cos(\alpha-\beta)]$ 展开：\(\displaystyle \int_{-\pi}^{\pi} \cos mx\cos nx\mathrm{d}x=\frac{1}{2} \int_{-\pi}^{\pi} [\cos(m+n)x + \cos(m-n)x] \mathrm{d}x\)。
若 $m \neq n$，两项积分均为 0；若 $m = n$，后一项变为 $\cos 0 = 1$，积分为 $\frac{1}{2} \cdot 2\pi = \pi$。
\item 方法二：复指数双元法。直接代入欧拉结构式，像拆分多项式一样展开：
\begin{align*}
\int_{-\pi}^{\pi} \cos mx\cos nx\mathrm{d}x &= \int_{-\pi}^{\pi} \left( \frac{e^{imx} + e^{-imx}}{2} \right) \left( \frac{e^{inx} + e^{-inx}}{2} \right) \mathrm{d}x \\
&= \frac{1}{4} \int_{-\pi}^{\pi} \left( e^{i(m+n)x} + e^{i(m-n)x} + e^{-i(m-n)x} + e^{-i(m+n)x} \right) \mathrm{d}x
\end{align*}
\textbf{教学提示：}引导学生观察这四项。
包含 $(m+n)$ 的项必定为 0。
对于包含 $(m-n)$ 的项：
如果 $m \neq n$，它们的积分也为 0。
如果 $m = n$，则 $e^{i(m-n)x}=1$，中间两项合并为 $2$。
此时积分化为常数积分，结果为 $\pi$。
这就是傅里叶级数中正交性的来源。
\end{enumerate}
\end{solution}

\begin{exercise}[定积分 - 145]
$\displaystyle \int_{-\pi}^{\pi} \sin mx\sin nx\mathrm{d}x$
\end{exercise}
\begin{solution}
\begin{enumerate}
\item 方法一：利用积化和差公式 $\sin \alpha \sin \beta = -\frac{1}{2}[\cos(\alpha+\beta) - \cos(\alpha-\beta)]$ 展开讨论（过程同上）。
\item 方法二：复指数双元法。再次利用复指数展开的代数性质（注意分母是 $2i$，相乘是 $4i^2 = -4$）：
\begin{align*}
\int_{-\pi}^{\pi} \sin mx\sin nx\mathrm{d}x &= \int_{-\pi}^{\pi} \left( \frac{e^{imx} - e^{-imx}}{2i} \right) \left( \frac{e^{inx} - e^{-inx}}{2i} \right) \mathrm{d}x \\
&= -\frac{1}{4} \int_{-\pi}^{\pi} \left( e^{i(m+n)x} - e^{i(m-n)x} - e^{-i(m-n)x} + e^{-i(m+n)x} \right) \mathrm{d}x
\end{align*}
同样利用指数项的正交性：当 $m \neq n$ 时，所有振荡项的积分都为 0。
当 $m = n$ 时，中间包含 $(m-n)$ 的两项变为 $-e^0 - e^0 = -2$。
积分变为：$-\frac{1}{4} \int_{-\pi}^{\pi} (-2) \mathrm{d}x = -\frac{1}{4} \times (-2) \times 2\pi = \pi$。
\end{enumerate}
\end{solution}

\begin{exercise}[定积分 - 146]
$\displaystyle \int_{0}^{\pi} \sin mx\sin nx\mathrm{d}x = \int_{0}^{\pi} \cos mx\cos nx\mathrm{d}x$
\end{exercise}
\begin{solution}
\begin{enumerate}
\item 方法一：积化和差，因为区间是 $[0, \pi]$，结果恰好是 $[-\pi, \pi]$ 的一半（由于被积函数都是偶函数）。
\item 方法二：复指数双元法。
\textbf{教学提示：}可以直接告诉学生，傅里叶正交性的性质不仅存在于全周期 $[-\pi, \pi]$，在半周期 $[0, \pi]$ 上同样有效。
拿 $\cos mx \cos nx$ 举例，展开后依然是 $\frac{1}{4} ( e^{i(m+n)x} + e^{i(m-n)x} + \dots )$。
由于 $m+n$ 和 $m-n$ 都是整数 $k$，在 $[0, \pi]$ 上 $\int_0^\pi e^{ikx} \mathrm{d}x = \frac{e^{ik\pi} - 1}{ik} = \frac{(-1)^k - 1}{ik}$。
虽然单个指数项在 $[0,\pi]$ 上未必分别积分为零，但它们总是以共轭成对的形式出现，例如 $e^{ikx}+e^{-ikx}=2\cos kx$。而当 $k\neq 0$ 时，\(\displaystyle \int_0^\pi \cos kx\,\mathrm{d}x=0\)。
所以，只有当 $m=n$ 产生常数项 $1$ 时才会留下数值：
$\int_{0}^{\pi} \cos mx\cos mx\mathrm{d}x = \frac{1}{4} \int_{0}^{\pi} (e^{2imx} + 1 + 1 + e^{-2imx}) \mathrm{d}x = \frac{1}{4} \int_{0}^{\pi} 2 \mathrm{d}x = \frac{\pi}{2}$。
因此得到所求结果。
\end{enumerate}
\end{solution}

\begin{exercise}[定积分 - 147]
$\displaystyle I_n = \int_{0}^{\frac{\pi}{2}} \sin^n x\mathrm{d}x = \int_{0}^{\frac{\pi}{2}} \cos^n x\mathrm{d}x$
\end{exercise}
\begin{note}[Wallis 递推的双元视角]
这题最值得看的，不是最后的递推式，而是“为什么端点会自动干净”。把 $p=\sin x,\ q=\cos x$ 代入后，区间端点刚好对应实圆第一象限的两个顶点 $(0,1)$ 与 $(1,0)$，所以分部积分里冒出的边界项天然消失。

因此这道题可以看成“实圆双元 + 端点信息 + 递推”的定积分版本：先在双元坐标里看清边界，再做代数降次，整个结构就会非常透明。
\end{note}
\begin{solution}
\begin{enumerate}
\item 方法一：分部积分法也可以得到同样的递推公式。
\item 方法二：实圆双元代数法。
\textbf{第一步：建立记号。}
记 \(\displaystyle p=\sin x,\qquad q=\cos x\)。则 $p^2+q^2=1$，并且 \(\displaystyle \mathrm{d}p=q\,\mathrm{d}x,\ \mathrm{d}q=-p\,\mathrm{d}x\)。
\textbf{第二步：考察端点。}
当 $x=0$ 时，$(p,q)=(0,1)$；当 $x=\frac{\pi}{2}$ 时，$(p,q)=(1,0)$，因此 \(\displaystyle [p^{n-1}q]_0^{\pi/2}=0\)。
\textbf{第三步：推导递推公式。}
我们要计算 $I_n = \int_0^{\pi/2} p^n \mathrm{d}x$。
写成微分形式为
$I_n = \int_0^{\pi/2} p^{n-1} \cdot p \mathrm{d}x = \int_0^{\pi/2} p^{n-1} (-\mathrm{d}q)$
由乘积求导公式 $\mathrm{d}(p^{n-1}q) = q\mathrm{d}(p^{n-1}) + p^{n-1}\mathrm{d}q$，可得
$I_n = -[p^{n-1}q]_0^{\pi/2} + \int_0^{\pi/2} q \mathrm{d}(p^{n-1})$
于是
$I_n = \int_0^{\pi/2} q \cdot (n-1)p^{n-2}\mathrm{d}p$
再由 $\mathrm{d}p=q\,\mathrm{d}x$，得
$I_n = (n-1)\int_0^{\pi/2} p^{n-2} q^2 \mathrm{d}x$
利用 $q^2 = 1 - p^2$，便有
$I_n = (n-1)\int_0^{\pi/2} p^{n-2} (1-p^2) \mathrm{d}x = (n-1)I_{n-2} - (n-1)I_n$
整理得
$n I_n = (n-1) I_{n-2} \implies I_n = \frac{n-1}{n}I_{n-2}$
这就得到了 Wallis 递推公式。再结合初值 $I_0 = \frac{\pi}{2},\ I_1 = 1$，即可继续求出各项。
\end{enumerate}
\end{solution}

\chapter{微分方程}

\section{知识讲解}

\subsection{概念与分类判断：解、通解、特解与反求}

\begin{definition}[通解] \label{def:ode-general-solution}
若微分方程的解中所含互相独立的任意常数的个数与方程的阶数相同，则称这样的解为该方程的\textbf{通解}（General Solution）。\\
更严格地，对于 $n$ 阶方程，若解 $y=y(x;c_1,\dots,c_n)$ 满足在某邻域内 Wronski 行列式不为零，则 $c_1,\dots,c_n$ 独立且该解为通解。\\
\textbf{反例}：$y=(C_1+2C_2)\sin t$ 中两个常数可合并为一个，不是通解。
\end{definition}

\begin{definition}[特解与定解条件] \label{def:ode-initial-value}
为确定通解中的任意常数而附加的条件称为\textbf{定解条件}。定解条件和微分方程共同构成一个\textbf{定解问题}。满足定解条件的解称为\textbf{特解}。\\
反映系统初始状态的条件称为\textbf{初始条件}。初始条件与微分方程构成的定解问题称为\textbf{初值问题}（Cauchy Problem/IVP），标准形式：\\
一阶：$\begin{cases} y'=f(x,y), \\ y(x_0)=y_0; \end{cases}$~~~~~二阶：$\begin{cases} y''=f(x,y,y'), \\ y(x_0)=y_0,\; y'(x_0)=y_0'. \end{cases}$
\end{definition}

\begin{definition}[积分曲线] \label{def:ode-integral-curve}
微分方程解的图形（即解曲线 $y=y(x)$ 在平面上的图像）称为该方程的\textbf{积分曲线}（Integral Curve）。通解对应于积分曲线族；初值问题的几何意义则是从曲线族中找出经过定点 $(x_0, y_0)$（及指定切线方向）的那条特定曲线。
\end{definition}

已知含任意常数的函数族（通解），可通过足够多次求导后联立消去常数来反推原方程。

\begin{exercise}
已知通解 $y=c_1\mathrm{e}^{-x}+c_2\mathrm{e}^{2x}$，求其所满足的微分方程。
\end{exercise}

\begin{solution}
对 $y$ 求导，得 \(\displaystyle y'=-c_1\mathrm{e}^{-x}+2c_2\mathrm{e}^{2x},\qquad
y''=c_1\mathrm{e}^{-x}+4c_2\mathrm{e}^{2x}\)。
$y+y''=5c_2\mathrm{e}^{2x}$ 与 $y'-y=-3c_2\mathrm{e}^{2x}$ 联立消去 $c_2\mathrm{e}^{2x}$：$3(y+y'')+5(y'-y)=0$，整理得 \(\displaystyle \boxed{y''-y'-2y=0}\)。
这是二阶常系数线性齐次方程，特征根恰为 $r_1=-1,\; r_2=2$，与出发时的通解结构吻合。
\end{solution}

\subsection{不要因为不能套公式就以为自己不会做，看看是否可分离、齐次与可化齐次}

\begin{remark}[丢失的特解]
形如\(\displaystyle \frac{\mathrm{d}y}{\mathrm{d}x}=f(x)\,\phi(y)\)，若存在常数 $y_0$ 使得 $\phi(y_0)=0$，则直接代入验证可知 $y\equiv y_0$ 也是原方程的一个特解。\textbf{这类特解往往在分离变量过程中丢失，且有时不包含在通解表达式中，必须另行补上才能得到全部解}。
\end{remark}

\begin{exercise} \label{ex:sep-circle}
求 $\displaystyle\frac{\mathrm{d}y}{\mathrm{d}x}=-\frac{x}{y}$ 的通解。
\end{exercise}
\begin{solution}
将方程变形为 $y\,\mathrm{d}y=-x\,\mathrm{d}x$，两边积分得 \(\displaystyle \int y\,\mathrm{d}y=-\int x\,\mathrm{d}x\)，即 \(\displaystyle \frac{1}{2}y^{2}=-\frac{1}{2}x^{2}+C\)。故 \(\displaystyle \boxed{x^{2}+y^{2}=C}\)，其中 $C>0$ 为任意常数。这是一族以原点为圆心的同心圆。
\end{solution}

\begin{exercise}
求方程 $x^{2}y\,\mathrm{d}y+\sqrt{1-y^{2}}\,\mathrm{d}x=0$ 的全部解。
\end{exercise}
\begin{solution}
当 $1-y^{2}\neq 0$ 时，方程分离为 \(\displaystyle -\frac{y\,\mathrm{d}y}{\sqrt{1-y^{2}}}=\frac{\mathrm{d}x}{x^{2}}\)，两边积分得 \(\displaystyle \sqrt{1-y^{2}}=-\frac{1}{x}+C\)，即 \(\displaystyle \sqrt{1-y^{2}}+\frac{1}{x}=C\)。

此外，易见 $y=\pm 1$ 也满足原方程（此时 $\sqrt{1-y^2}=0$），
且它们不包含在上面的通解中。因此方程的\textbf{全部解}为
\(\displaystyle \boxed{\sqrt{1-y^{2}}+\frac{1}{x}=C \quad\text{及}\quad y=\pm 1}\)。
\end{solution}

\begin{exercise}
求 $\displaystyle\frac{\mathrm{d}y}{\mathrm{d}x}=y^{2}\cos x$ 满足初始条件 $y(0)=1$ 的特解。
\end{exercise}
\begin{solution}
当 $y\neq 0$ 时分离变量：$\displaystyle\frac{\mathrm{d}y}{y^{2}}=\cos x\,\mathrm{d}x$，
积分得 \(\displaystyle -\frac{1}{y}=\sin x + C\)，从而
\(\displaystyle y=-\frac{1}{\sin x+C}\)。
另有特解 $y\equiv 0$，但它不满足初值条件 $y(0)=1$，故本题中应舍去。

代入 $y(0)=1$ 得 $-1=C$，故所求特解为
\(\displaystyle \boxed{y=\frac{1}{1-\sin x}},\ x\neq\frac{\pi}{2}+2k\pi\)。
\end{solution}

\begin{exercise} \label{ex:sep-linear-homog}
求 $\displaystyle\frac{\mathrm{d}y}{\mathrm{d}x}=p(x)\,y$ 的通解，其中 $p(x)$ 是连续函数。
\end{exercise}
\begin{solution}
当 $y\neq 0$ 时分离变量并积分：$\int\frac{\mathrm{d}y}{y}=\int p(x)\,\mathrm{d}x \;\implies\; \ln|y|=\int p(x)\,\mathrm{d}x+\bar{C}$。
令 $C=\pm\mathrm{e}^{\bar{C}}$ 得
\(\displaystyle \boxed{y=C\exp\!\left(\int p(x)\,\mathrm{d}x\right)}\)。
此外 $y\equiv 0$ 显然也是解，且已包含在上述通解中（取 $C=0$）。
\end{solution}

此例的结果非常重要，它正是后续一阶线性方程求解公式的基础。
同时要记住：通解中的积分 $\int p(x)\,\mathrm{d}x$ 只需取一个原函数即可，积分常数已经吸收到 $C$ 中。

\subsubsection{齐次方程}

对于齐次方程，作代换 $u=\dfrac{y}{x}$（即 $y=ux$），由乘积法则
$\displaystyle\frac{\mathrm{d}y}{\mathrm{d}x}=u+x\frac{\mathrm{d}u}{\mathrm{d}x}$，
代入原方程得
\begin{equation}
u+x\frac{\mathrm{d}u}{\mathrm{d}x}=\varphi(u)
\quad\implies\quad
\frac{\mathrm{d}u}{\mathrm{d}x}=\frac{\varphi(u)-u}{x}. \tag{5.2}
\end{equation}

\begin{property}[齐次方程的完整解法]
\begin{enumerate}
  \item 若 $\varphi(u)-u\neq 0$，则 (5.2) 分离为
        $\displaystyle\frac{\mathrm{d}u}{\varphi(u)-u}=\frac{\mathrm{d}x}{x}$，
        积分后回代 $u=\dfrac{y}{x}$ 即得通解；
  \item 若 $\varphi(u)-u=0$ 有实根 $u_1,u_2,\dots,u_m$，
        则直线族 $y=u_i\,x$（$i=1,\dots,m$）也是解，一般不包含在通解中；
  \item 若 $\varphi(u)-u\equiv 0$，则 $\dfrac{\mathrm{d}u}{\mathrm{d}x}\equiv 0$，
        解为一族过原点的直线 $y=cx$。
\end{enumerate}
\end{property}

\begin{exercise}
求 $\displaystyle\frac{\mathrm{d}y}{\mathrm{d}x}=\frac{y}{x}+\tan\frac{y}{x}$ 的通解。
\end{exercise}
\begin{solution}
\begin{enumerate}
    \item \textbf{判断方程类型并代换}：
    观察到方程右端的函数完全由变量比例 $\frac{y}{x}$ 构成，这说明它是一个齐次微分方程。
    令 $u=\frac{y}{x}$，即 $y=ux$。两边对 $x$ 求导得到 $\frac{\mathrm{d}y}{\mathrm{d}x} = u + x\frac{\mathrm{d}u}{\mathrm{d}x}$。
    \item \textbf{代入并化简方程}：
    将 $y=ux$ 及导数表达式代入原方程，得 $u + x\frac{\mathrm{d}u}{\mathrm{d}x} = u + \tan u$。
    消去两边的 $u$，得到一个可分离变量的方程 $x\frac{\mathrm{d}u}{\mathrm{d}x} = \tan u$。
    \item \textbf{分离变量并积分}：
    将含有 $u$ 的项移到左边，含有 $x$ 的项移到右边（假设 $\tan u \neq 0$ 且 $x \neq 0$）：\(\displaystyle \frac{\mathrm{d}u}{\tan u} = \frac{\mathrm{d}x}{x} \implies \frac{\cos u}{\sin u}\mathrm{d}u = \frac{\mathrm{d}x}{x}\)。
    对等式两边同时进行不定积分，左边凑微分为 $\mathrm{d}(\sin u)$，于是 \(\displaystyle \int \frac{\mathrm{d}(\sin u)}{\sin u} = \int \frac{\mathrm{d}x}{x} \implies \ln|\sin u| = \ln|x| + C_1\)。
    \item \textbf{处理常数并回代}：
    利用对数性质化简上式：$|\sin u| = \mathrm{e}^{C_1}|x| \implies \sin u = \pm\mathrm{e}^{C_1}x$。
    令 $C = \pm\mathrm{e}^{C_1}$，则有 $\sin u = Cx$。
    同时，我们需要检验分离变量时遗漏的解 $\tan u = 0$，即 $\sin u = 0$。此时对应的正是 $C=0$ 的情况，因此常数 $C$ 可以是任意实数。
    最后，将 $u = \frac{y}{x}$ 代回，得到原方程的隐式通解
    \(\displaystyle \boxed{\sin\frac{y}{x} = Cx}\)。
\end{enumerate}
\end{solution}

\begin{exercise}
求解 $(x+y)\,\mathrm{d}x-(y-x)\,\mathrm{d}y=0$。
\end{exercise}
\begin{solution}
\begin{enumerate}
    \item \textbf{方程变形与判断}：
    将方程写成导数形式：$\frac{\mathrm{d}y}{\mathrm{d}x} = \frac{x+y}{y-x} = \frac{1 + \frac{y}{x}}{\frac{y}{x} - 1}$。
    右端是 $\frac{y}{x}$ 的函数，故为齐次微分方程。
    \item \textbf{变量代换}：
    令 $y=ux$，则其微分为 $\mathrm{d}y=u\,\mathrm{d}x+x\,\mathrm{d}u$。直接将它们代入原微分形式中：\(\displaystyle (x+ux)\,\mathrm{d}x - (ux-x)(u\,\mathrm{d}x+x\,\mathrm{d}u) = 0\)。
    \item \textbf{展开与化简}：
    提取公因子 $x$（假设 $x \neq 0$）：\(\displaystyle x\left[ (1+u)\,\mathrm{d}x - (u-1)(u\,\mathrm{d}x+x\,\mathrm{d}u) \right] = 0\)。
    化简括号内的项，把 $\mathrm{d}x$ 和 $\mathrm{d}u$ 的系数归拢：\(\displaystyle (1+u)\,\mathrm{d}x - (u^2-u)\,\mathrm{d}x - x(u-1)\,\mathrm{d}u = 0 \;\Longrightarrow\; (1 + 2u - u^2)\,\mathrm{d}x - x(u-1)\,\mathrm{d}u = 0\)。
    \item \textbf{分离变量与积分}：
    将 $\mathrm{d}x$ 和 $\mathrm{d}u$ 分开，得
    \(\displaystyle \frac{\mathrm{d}x}{x} = \frac{u-1}{1+2u-u^2}\,\mathrm{d}u\)。
    观察到右边分母的导数为 $\frac{\mathrm{d}}{\mathrm{d}u}(1+2u-u^2) = 2-2u = -2(u-1)$，因此
    \(\displaystyle \int \frac{\mathrm{d}x}{x}
    = -\frac{1}{2} \int \frac{\mathrm{d}(1+2u-u^2)}{1+2u-u^2}\)。
    积分得到 $\ln|x| = -\frac{1}{2}\ln|1+2u-u^2| + C_1$。
    \item \textbf{化简并回代}：
    方程两边同乘 $2$ 得 \(2\ln|x| + \ln|1+2u-u^2| = 2C_1\)，合并对数为 \(\displaystyle \ln\left| x^2(1+2u-u^2) \right| = 2C_1\)。去对数，令 $C = \pm\mathrm{e}^{2C_1}$，得 $x^2(1+2u-u^2) = C$。
    （注：当分离变量时令 $1+2u-u^2=0$ 即 $u=1\pm\sqrt{2}$ 得到的特解，正好对应 $C=0$ 的情形，因此 $C$ 为任意常数。）
    将 $u=\frac{y}{x}$ 代回上式：
    \(\displaystyle x^2\left(1 + 2\frac{y}{x} - \frac{y^2}{x^2}\right) = C\)。
    展开得到最终的通解：
    \(\displaystyle \boxed{x^2+2xy-y^2=C}\)。
\end{enumerate}
\end{solution}

\begin{exercise}
求解 $x\,\mathrm{d}y-y\,\mathrm{d}x=\sqrt{x^{2}+y^{2}}\,\mathrm{d}x$。
\end{exercise}
\begin{solution}
\begin{enumerate}
    \item \textbf{方程化简与代换}：
    先合并同类项，把 $\mathrm{d}x$ 移到一起：
    \(\displaystyle x\,\mathrm{d}y = \left( y + \sqrt{x^2+y^2} \right)\,\mathrm{d}x\)。
    假设 $x>0$（对于 $x<0$ 的情况推导完全类似），两边同除以 $x\mathrm{d}x$，将根号内的 $x^2$ 提出来，得到齐次方程的形式：\(\displaystyle \frac{\mathrm{d}y}{\mathrm{d}x} = \frac{y}{x} + \frac{\sqrt{x^2+y^2}}{x} = \frac{y}{x} + \sqrt{1+\left(\frac{y}{x}\right)^2}\)。
    \item \textbf{代入并分离变量}：
    令 $u=\frac{y}{x}$，则 $\frac{\mathrm{d}y}{\mathrm{d}x} = u + x\frac{\mathrm{d}u}{\mathrm{d}x}$，代入方程得 $u + x\frac{\mathrm{d}u}{\mathrm{d}x} = u + \sqrt{1+u^2} \implies x\frac{\mathrm{d}u}{\mathrm{d}x} = \sqrt{1+u^2}$。
    分离变量，将 $u$ 的项移到左边，$x$ 的项移到右边：$\frac{\mathrm{d}u}{\sqrt{1+u^2}} = \frac{\mathrm{d}x}{x}$。
    \item \textbf{积分与化简}：
    利用基本积分公式 \(\displaystyle \int \frac{\mathrm{d}u}{\sqrt{1+u^2}}
    = \ln(u+\sqrt{1+u^2})+C\)。对两边进行积分，得
    \(\displaystyle \ln\left(u+\sqrt{1+u^2}\right) = \ln x + \ln C_1 = \ln(C_1 x)\)。
    去对数（因为 $\sqrt{1+u^2}>|u|$，对数真数恒正，并取 $C>0$），得 $u + \sqrt{1+u^2} = Cx$。
    \item \textbf{回代并消除根号}：
    将 $u=\frac{y}{x}$ 代回方程：
    \(\displaystyle \frac{y}{x} + \sqrt{1+\left(\frac{y}{x}\right)^2} = Cx
    \implies \sqrt{x^2+y^2} = Cx^2 - y\)。
    为了得到不含根号的显式解，将方程两边平方：$x^2 + y^2 = (Cx^2 - y)^2 = C^2x^4 - 2Cx^2y + y^2$。
    两边消去 $y^2$，然后提取 $x^2$（已知 $x \neq 0$），得 $x^2 = C^2x^4 - 2Cx^2y \implies 1 = C^2x^2 - 2Cy$。
    解出 $y$，得到通解
    \(\displaystyle 2Cy = C^2x^2 - 1 \implies \boxed{y = \frac{1}{2}\left(Cx^2 - \frac{1}{C}\right)}\)。
\end{enumerate}
\end{solution}

\subsubsection{可化为齐次方程的类型}

\begin{exercise} \label{ex:ode-quasi-homog}
求解方程 $\displaystyle\frac{\mathrm{d}y}{\mathrm{d}x}=\frac{a_1x+b_1y+c_1}{a_2x+b_2y+c_2}$（系数均为常数）。
\end{exercise}
\begin{solution}
这类方程虽然不齐次，但可以通过几何平移或代数变换转化为齐次方程。根据系数的情况，我们分为三种情形进行详细讨论：
\begin{enumerate}
    \item \textbf{情形一：$c_1=c_2=0$}\\
    方程直接就是 $\frac{\mathrm{d}y}{\mathrm{d}x}=\frac{a_1x+b_1y}{a_2x+b_2y}$。分子分母同除以 $x$，立即可见其为标准的零次齐次方程。只需令 $u = \frac{y}{x}$ 即可化为可分离变量的方程。

    \item \textbf{情形二：行列式 $\begin{vmatrix}a_1&b_1\\a_2&b_2\end{vmatrix}=0$}\\
    这意味着直线 $a_1x+b_1y+c_1=0$ 和 $a_2x+b_2y+c_2=0$ 平行。因此系数成比例，即存在常数 $k$ 使得 $\frac{a_1}{a_2} = \frac{b_1}{b_2} = k$。
    此时原方程分子部分可写为 $k(a_2x+b_2y)+c_1$。
    方程变为 $\frac{\mathrm{d}y}{\mathrm{d}x} = \frac{k(a_2x+b_2y)+c_1}{a_2x+b_2y+c_2}$。
    通过作变量代换 $u = a_2x+b_2y$，则 $\frac{\mathrm{d}u}{\mathrm{d}x} = a_2 + b_2\frac{\mathrm{d}y}{\mathrm{d}x}$，代入即可将方程彻底化为一个仅包含 $u$ 和 $x$ 的可分离变量方程。

    \item \textbf{情形三：行列式 $\begin{vmatrix}a_1&b_1\\a_2&b_2\end{vmatrix}\neq0$ 且 $c_1,c_2$ 不全为零}\\
    这表明平面上两条直线 $a_1x+b_1y+c_1=0$ 和 $a_2x+b_2y+c_2=0$ 必相交于唯一一点。
    设交点坐标为 $(\alpha, \beta)$。为了消去常数项 $c_1, c_2$，我们通过坐标平移将原点移动到交点处。
    作代换：$X = x - \alpha, \quad Y = y - \beta$。此时 $\mathrm{d}X = \mathrm{d}x, \mathrm{d}Y = \mathrm{d}y$。
代入原方程得到 \(\displaystyle \frac{\mathrm{d}Y}{\mathrm{d}X} = \frac{a_1(X+\alpha)+b_1(Y+\beta)+c_1}{a_2(X+\alpha)+b_2(Y+\beta)+c_2} = \frac{a_1X+b_1Y + (a_1\alpha+b_1\beta+c_1)}{a_2X+b_2Y + (a_2\alpha+b_2\beta+c_2)}\)。
因为 $(\alpha, \beta)$ 是交点，满足 $a_1\alpha+b_1\beta+c_1=0$ 和 $a_2\alpha+b_2\beta+c_2=0$，所以常数项消失，方程化为 \(\displaystyle \frac{\mathrm{d}Y}{\mathrm{d}X} = \frac{a_1X+b_1Y}{a_2X+b_2Y}\)。
    这成功地将原方程化为情形一的齐次方程，令 $u = \frac{Y}{X}$ 继续求解即可。
\end{enumerate}
\end{solution}

\begin{exercise} \label{ex:ode-intersect-line}
求解 $\displaystyle\frac{\mathrm{d}y}{\mathrm{d}x}=\frac{x-y+1}{x+y-3}$。
\end{exercise}
\begin{solution}
\begin{enumerate}
    \item \textbf{求交点进行平移}：
    计算系数行列式 $\begin{vmatrix}1 & -1 \\ 1 & 1\end{vmatrix} = 2 \neq 0$，属于相交直线的类型。
    联立方程组求交点：
    联立 \(\displaystyle x-y+1=0,\ x+y-3=0\)，
    两式相加得 $2x-2=0 \implies x=1$；相减得 $-2y+4=0 \implies y=2$，即交点为 $(\alpha, \beta) = (1,2)$。

    \item \textbf{坐标平移化为齐次方程}：
    作代换 $X=x-1,\;Y=y-2$，则 $\mathrm{d}X=\mathrm{d}x,\;\mathrm{d}Y=\mathrm{d}y$。原方程平移后常数项消失，变为：
    $\frac{\mathrm{d}Y}{\mathrm{d}X}=\frac{X-Y}{X+Y}$。

    \item \textbf{分离变量与积分}：
    令 $Y=uX$，则 $\frac{\mathrm{d}Y}{\mathrm{d}X} = u+X\frac{\mathrm{d}u}{\mathrm{d}X}$，代入化简：
    $u+X\frac{\mathrm{d}u}{\mathrm{d}X} = \frac{X-uX}{X+uX} = \frac{1-u}{1+u}$。
    移项并通分，\(\displaystyle X\frac{\mathrm{d}u}{\mathrm{d}X} = \frac{1-u}{1+u} - u = \frac{1-2u-u^2}{1+u}
    \implies \frac{1+u}{1-2u-u^2}\mathrm{d}u = \frac{\mathrm{d}X}{X}\)。
    左侧分母的导数为 $-2-2u = -2(1+u)$，据此凑微分，\(\displaystyle -\frac{1}{2} \int \frac{\mathrm{d}(1-2u-u^2)}{1-2u-u^2} = \int \frac{\mathrm{d}X}{X}\)。
    积分得到：
    $-\frac{1}{2}\ln|1-2u-u^2| = \ln|X| + C_1$。

    \item \textbf{整理结果并回代}：
    两边同乘 $-2$，得 $\ln|1-2u-u^2| = -2\ln|X| - 2C_1 = \ln(X^{-2}) + C_2$。去对数得：
    $1-2u-u^2 = \frac{C_3}{X^2} \implies X^2(1-2u-u^2) = C_3$。
    把 $u = \frac{Y}{X}$ 代回，化简得 $\displaystyle X^2 - 2XY - Y^2 = C_3$。
    将原来的变量 $X = x-1$ 和 $Y = y-2$ 替换回来，\(\displaystyle (x-1)^2 - 2(x-1)(y-2) - (y-2)^2 = C_3\)。
    展开并合并同类项，再把常数吸收到右边，可写成
    \(\displaystyle \boxed{x^{2}-y^{2}+2x+6y-2xy=C}\)。
\end{enumerate}
\end{solution}

\begin{exercise}
求解初值问题
$\begin{cases}(x^2+2xy-y^2)\mathrm{d}x+(y^2+2xy-x^2)\mathrm{d}y=0,\\ y(1)=1.\end{cases}$
\end{exercise}
\begin{solution}
\begin{enumerate}
    \item \textbf{转换为齐次方程标准形式}：
    把方程重写为导数的形式：
    $\frac{\mathrm{d}y}{\mathrm{d}x} = -\frac{x^2+2xy-y^2}{y^2+2xy-x^2}$。
    分子分母同除以 $x^2$，显然它是一个齐次方程：
    $\frac{\mathrm{d}y}{\mathrm{d}x} = \frac{(y/x)^2 - 2(y/x) - 1}{(y/x)^2 + 2(y/x) - 1}$。
    \item \textbf{令 $y=ux$ 进行代换}：
    代入 $\frac{\mathrm{d}y}{\mathrm{d}x} = u + x\frac{\mathrm{d}u}{\mathrm{d}x}$ 到方程中：
    $u + x\frac{\mathrm{d}u}{\mathrm{d}x} = \frac{u^2-2u-1}{u^2+2u-1}$。
    将 $u$ 移到右边并通分化简，可得 \(\displaystyle x\frac{\mathrm{d}u}{\mathrm{d}x} = \frac{u^2-2u-1 - u(u^2+2u-1)}{u^2+2u-1} = \frac{u^2-2u-1 - u^3-2u^2+u}{u^2+2u-1} = \frac{-u^3-u^2-u-1}{u^2+2u-1}\)。
    对分子进行因式分解：$-u^3-u^2-u-1 = -u^2(u+1) - (u+1) = -(u+1)(u^2+1)$。
    于是方程化为：
    $x\frac{\mathrm{d}u}{\mathrm{d}x} = -\frac{(u+1)(u^2+1)}{u^2+2u-1}$。
    \item \textbf{分离变量与有理分式积分}：
    将含有 $u$ 的移动到一边，得到：
    $\frac{u^2+2u-1}{(u+1)(u^2+1)}\mathrm{d}u = -\frac{\mathrm{d}x}{x}$。
    利用待定系数法对左侧进行有理分式分解。设：
    $\frac{u^2+2u-1}{(u+1)(u^2+1)} = \frac{A}{u+1} + \frac{Bu+C}{u^2+1}$。
    通分合并分子：$A(u^2+1) + (Bu+C)(u+1) = u^2+2u-1$。
    令 $u=-1$，得 $2A = 1-2-1 = -2 \implies A=-1$。
    比较 $u^2$ 的系数：$A+B=1 \implies B=2$。
    比较常数项：$A+C=-1 \implies C=0$。
    分解结果为：
    $\left(\frac{2u}{u^2+1} - \frac{1}{u+1}\right)\mathrm{d}u = -\frac{\mathrm{d}x}{x}$。
    \item \textbf{积分回代与定常数}：
    对等式两边进行积分：\(\displaystyle \ln(u^2+1) - \ln|u+1| = -\ln|x| + \ln|C| \implies \ln\frac{u^2+1}{|u+1|} = \ln\frac{|C|}{|x|}\)。
    去掉对数，取常数 $C_1$：
    $\frac{u^2+1}{u+1} = \frac{C_1}{x}$。
    将 $u = \frac{y}{x}$ 替换回去：\(\displaystyle \frac{\frac{y^2}{x^2}+1}{\frac{y}{x}+1} = \frac{C_1}{x} \implies \frac{\frac{y^2+x^2}{x^2}}{\frac{y+x}{x}} = \frac{C_1}{x} \implies \frac{y^2+x^2}{x(y+x)} = \frac{C_1}{x}\)。
    两边乘以 $x$ 得到通解：
    $\frac{x^2+y^2}{x+y} = C_1$。
    应用初始条件 $y(1)=1$，代入 $x=1, y=1$：
    $C_1 = \frac{1^2+1^2}{1+1} = \frac{2}{2} = 1$。
    整理得最终特解为：
    $\boxed{x^{2}+y^{2}=x+y}$。
\end{enumerate}
\end{solution}

\begin{exercise}
求解：(1) $f(xy)y\,\mathrm{d}x+g(xy)x\,\mathrm{d}y=0$；
\quad (2) $\displaystyle\frac{\mathrm{d}y}{\mathrm{d}x}=\frac{1}{x+y}$。
\end{exercise}
\begin{solution}
\begin{enumerate}
    \item[\textbf{(1)}] \textbf{整体代换法解决耦合变量}：
    观察到方程中函数的变量均以 $xy$ 的形式结合在一起，自然想到令 $u = xy$ 作为整体代换。则有 $y = \frac{u}{x}$。
    根据乘积求导法则计算其全微分：$\mathrm{d}u = y\,\mathrm{d}x + x\,\mathrm{d}y$，从中可解出 $x\,\mathrm{d}y = \mathrm{d}u - y\,\mathrm{d}x = \mathrm{d}u - \frac{u}{x}\,\mathrm{d}x$。
    将这些替换带入原方程 $f(xy)y\,\mathrm{d}x + g(xy)x\,\mathrm{d}y = 0$ 中：\(\displaystyle f(u)\frac{u}{x}\,\mathrm{d}x + g(u)\left(\mathrm{d}u - \frac{u}{x}\,\mathrm{d}x\right) = 0\)。
    将 $\mathrm{d}x$ 和 $\mathrm{d}u$ 的各项合并归类：\(\displaystyle \left[ u f(u) - u g(u) \right] \frac{\mathrm{d}x}{x} + g(u)\mathrm{d}u = 0\)。
    分离变量，把所有含有 $u$ 的项移到方程的另一边：\(\displaystyle \frac{\mathrm{d}x}{x} = - \frac{g(u)}{u[f(u)-g(u)]}\mathrm{d}u\)。
    对两边进行积分，即得隐式通解表达式：\(\displaystyle \ln|x| + \int \frac{g(u)}{u[f(u)-g(u)]}\,\mathrm{d}u = C\)。

    \item[\textbf{(2)}] \textbf{线性组合代换法}：
    方程右侧分母含有 $x+y$ 的组合，令 $u = x+y$ 可化简方程。
    两边对 $x$ 求导得到 $\frac{\mathrm{d}u}{\mathrm{d}x} = 1 + \frac{\mathrm{d}y}{\mathrm{d}x}$，即 $\frac{\mathrm{d}y}{\mathrm{d}x} = \frac{\mathrm{d}u}{\mathrm{d}x} - 1$。
    代入原方程中：\(\displaystyle \frac{\mathrm{d}u}{\mathrm{d}x} - 1 = \frac{1}{u} \implies \frac{\mathrm{d}u}{\mathrm{d}x} = 1 + \frac{1}{u} = \frac{u+1}{u}\)。
    这就是一个可分离变量的方程。将 $u$ 的变量放到左边：\(\displaystyle \frac{u}{u+1}\mathrm{d}u = \mathrm{d}x\)。
    为了方便积分，将分子改写为 $(u+1)-1$：\(\displaystyle \int \left( 1 - \frac{1}{u+1} \right)\mathrm{d}u = \int \mathrm{d}x\)，从而 \(\displaystyle u-\ln|u+1|=x+C_1 \xrightarrow{u=x+y} (x+y)-\ln|x+y+1|=x+C_1 \implies \boxed{y-\ln|x+y+1|=C_1}\)。
    如果对该式进一步求解关于 $x$ 的表达式：$\ln|x+y+1| = y - C_1 \implies x+y+1 = \pm\mathrm{e}^{-C_1}\mathrm{e}^y$。设 $C = \pm\mathrm{e}^{-C_1}$，则有另一种形式：\(\displaystyle \boxed{x=C\mathrm{e}^{y}-y-1}\)。
\end{enumerate}
\end{solution}

\begin{property}[复合结构的整体代换]
方程中出现 $f(xy),\,f(x\pm y),\,f(x^2\pm y^2),\,f(y/x)$ 等复合结构时，
通常应尝试相应代换 $u=xy,\,x\pm y,\,x^2\pm y^2,\,y/x$ 等，
将方程化为可分离变量型。
\end{property}

\subsection{线性、伯努利、黎卡提与结构代换}

\subsubsection{线性齐次方程的求解，别忘记换元或者对调x,y}

\begin{property}[通解的结构]
将通解公式展开，可以明显地看到它由两部分叠加而成：
\[
y = \underbrace{C_0\exp\!\left(-\int P(x)\,\mathrm{d}x\right)}_{\text{对应齐次方程的通解}} \;+\; \underbrace{\exp\!\left(-\int P(x)\,\mathrm{d}x\right)\int Q(x)\exp\!\left(\int P(x)\,\mathrm{d}x\right)\mathrm{d}x}_{\text{原非齐次方程的一个特解}}
\]
即：\textbf{非齐次方程的通解 = 对应齐次方程的通解 + 非齐次方程的任意一个特解}。
\end{property}

\begin{theorem}[积分因子法 / 凑微分法] \label{thm:integrating-factor}
将原非齐次方程 $\frac{\mathrm{d}y}{\mathrm{d}x}+P(x)y=Q(x)$ 的两边同时乘以积分因子 $\mu(x)=\exp\!\left(\int P(x)\,\mathrm{d}x\right)$：\(\displaystyle \mu(x)\frac{\mathrm{d}y}{\mathrm{d}x}+P(x)\mu(x)y=Q(x)\mu(x)\)。

观察方程左边。利用乘积求导法则，
\(\displaystyle \frac{\mathrm{d}}{\mathrm{d}x}\left[ y \cdot \exp\!\left(\int P(x)\,\mathrm{d}x\right) \right]
= \mu(x)\frac{\mathrm{d}y}{\mathrm{d}x} + y\mu'(x)
= Q(x)\mu(x)\)。

对上式两边同时关于 $x$ 积分（左边积分即消去导数符号，右边加上任意常数 $C_0$），得
\(\displaystyle y \cdot \exp\!\left(\int P(x)\,\mathrm{d}x\right) = \int Q(x)\exp\!\left(\int P(x)\,\mathrm{d}x\right)\mathrm{d}x + C_0\)。

两边同乘 $\exp\!\left(-\int P(x)\,\mathrm{d}x\right)$，即可解得 $y$，从而得到通解公式：
\[
\boxed{ y=\exp\!\left(-\int P(x)\,\mathrm{d}x\right)\left[\int Q(x)\exp\!\left(\int P(x)\,\mathrm{d}x\right)\mathrm{d}x+C_0\right] }
\]
\end{theorem}

\begin{exercise}
求方程 $\displaystyle\frac{\mathrm{d}y}{\mathrm{d}x}-\frac{y}{x}=x^{2}$ 的通解。
\end{exercise}
\begin{solution}
\begin{enumerate}
    \item \textbf{确认标准形式与提取系数}：
    方程已是一阶线性微分方程的标准形式 $\frac{\mathrm{d}y}{\mathrm{d}x}+P(x)y=Q(x)$，其中 $P(x)=-\frac{1}{x},\ Q(x)=x^{2}$。

    \item \textbf{计算积分因子}：
    构造积分因子 $\mu(x) = \exp\!\left(\int P(x)\,\mathrm{d}x\right)$，于是 \(\displaystyle \mu(x)=\frac{1}{x}\)（因为 \(\displaystyle \exp(-\ln x)=\frac{1}{x}\)）。
    （注：此处为寻找一个特解作为积分因子，故可省略绝对值和积分常数。）

    \item \textbf{方程同乘积分因子并凑微分}：
    将方程两边同时乘以 $\frac{1}{x}$，根据乘积求导法则 $(uv)' = u'v + uv'$ 逆推，可化为
    \(\displaystyle \frac{1}{x}\frac{\mathrm{d}y}{\mathrm{d}x} - \frac{1}{x^{2}}y = x
    \implies \frac{\mathrm{d}}{\mathrm{d}x}\left( \frac{y}{x} \right) = x\)。

    \item \textbf{直接积分求通解}：
    对上式两边同时关于 $x$ 积分，得 \(\displaystyle \frac{y}{x} = \int x\,\mathrm{d}x \implies \frac{y}{x} = \frac{x^{2}}{2}+C_0\)。
    两边同乘 $x$，展开即得通解 $\boxed{y=C_0x+\frac{x^{3}}{2}}$。
\end{enumerate}
\end{solution}

\begin{exercise}
解初值问题
$\begin{cases}(x^{2}-1)y'+2xy=\cos x,\\ y(0)=1.\end{cases}$
\end{exercise}
\begin{solution}
\begin{enumerate}
    \item \textbf{直接观察并凑全微分}：
    观察原方程左边 $(x^{2}-1)y'+2xy$。
    注意到 $2x$ 恰好是 $(x^2-1)$ 的导数，即 $(x^2-1)' = 2x$。
    根据乘积求导法则 $(uv)'=u'v+uv'$，原方程左边恰可写成导数展开式 $(x^{2}-1)y' + (x^{2}-1)'y = \cos x$，因此可将其写成全导数形式 $[(x^{2}-1)y]' = \cos x$（这一步与乘积分因子的做法等价）。

    \item \textbf{两边直接积分并代入初值}：
    \(\displaystyle (x^{2}-1)y=\int \cos x\,\mathrm{d}x=\sin x+C_0,\qquad y=\frac{\sin x+C_0}{x^{2}-1}\)。
    由 $y(0)=1$ 得 $C_0=-1$，故 $\boxed{y=\frac{\sin x - 1}{x^{2}-1} = \frac{1-\sin x}{1-x^{2}}}$。
\end{enumerate}
\end{solution}

\begin{exercise}
解方程 $y\ln y\,\mathrm{d}x+(x-\ln y)\,\mathrm{d}y=0$。
\end{exercise}
\begin{solution}
\begin{enumerate}
    \item \textbf{转换自变量与因变量}：
    若按常规写成 $\frac{\mathrm{d}y}{\mathrm{d}x}=-\frac{y\ln y}{x-\ln y}$，方程既不是线性方程也不是可分离变量方程。
    \textbf{关键观察}：我们尝试交换自变量和因变量的角色，将 $x$ 视为 $y$ 的函数，将原方程除以 $\mathrm{d}y$，得到
    \(\displaystyle y\ln y\frac{\mathrm{d}x}{\mathrm{d}y}+(x-\ln y)=0 \implies \frac{\mathrm{d}x}{\mathrm{d}y}+\frac{1}{y\ln y}x=\frac{1}{y}\)。

    \item \textbf{计算积分因子}：
    此处自变量为 $y$，$P(y)=\frac{1}{y\ln y}$。构造积分因子 $\mu(y) = \exp\!\left(\int P(y)\,\mathrm{d}y\right)$：
    \(\displaystyle \mu(y) = \exp\!\left(\int \frac{1}{y\ln y}\,\mathrm{d}y\right) = \exp\!\left(\int \frac{1}{\ln y}\,\mathrm{d}(\ln y)\right) = \exp(\ln(\ln y)) = \ln y\)。

    \item \textbf{方程同乘积分因子并凑微分}：
    将标准形式方程两边同时乘以 $\ln y$，得
    \(\displaystyle \ln y \cdot \frac{\mathrm{d}x}{\mathrm{d}y} + \frac{1}{y} \cdot x = \frac{\ln y}{y}
    \implies \frac{\mathrm{d}}{\mathrm{d}y}(x \ln y) = \frac{\ln y}{y}\)。

    \item \textbf{积分求通解}：
    两边同时关于 $y$ 积分（右侧凑微分 $\mathrm{d}(\ln y)=\frac{1}{y}\mathrm{d}y$），得
    \(\displaystyle x \ln y = \int \frac{\ln y}{y}\,\mathrm{d}y = \int \ln y\,\mathrm{d}(\ln y) = \frac{1}{2}(\ln y)^{2}+C_0\)。
    等式两边同除以 $\ln y$，得到原方程的通解
    \(\displaystyle \boxed{x=\frac{1}{2}\ln y+\frac{C_0}{\ln y}}\)
    （另外，观察原方程可知，使分母 $\ln y=0$ 的函数 $y \equiv 1$ 导致退化，代回原微分形式恒成立，故 $y \equiv 1$ 也是一个奇解）。
\end{enumerate}
\end{solution}
\begin{property}[参数化曲面下的有向投影面积元]
    在实际解题中，出题人往往会化成含有$y'$的形式，此时我们要先化成同时含有$\mathrm{d}y,\mathrm{d}x$的形式，这样更容易看出来选择谁做自变量。
\end{property}

\begin{exercise}[质点在驱动力与阻力下的运动]
质量为 $m$ 的质点从静止开始运动，受到与时间成正比的驱动力 $k_1 t$（与速度同向），以及与速度成正比的阻力 $k_2 v$。求速度 $v(t)$。
\end{exercise}
\begin{solution}
\begin{enumerate}
    \item \textbf{建立微分方程并化为标准形式}：
    根据牛顿第二定律 $F=ma$，合外力为驱动力与阻力之差，故
    \(\displaystyle m\frac{\mathrm{d}v}{\mathrm{d}t}=k_1t-k_2v,\qquad
    v(0)=0,\qquad \frac{\mathrm{d}v}{\mathrm{d}t}+\frac{k_2}{m}v=\frac{k_1}{m}t\)。

    \item \textbf{求解对应齐次方程}：
    齐次方程为 $\frac{\mathrm{d}v}{\mathrm{d}t}+\frac{k_2}{m}v=0$。分离变量积分：
    \(\displaystyle \frac{\mathrm{d}v}{v}=-\frac{k_2}{m}\,\mathrm{d}t \implies \ln|v|=-\frac{k_2}{m}t+C_1 \implies v_h=C\mathrm{e}^{-\frac{k_2}{m}t}\)

    \item \textbf{常数变易法求特解并代入初值}：
    \(\displaystyle v=C(t)\mathrm{e}^{-\frac{k_2}{m}t}
    \implies C'(t)\mathrm{e}^{-\frac{k_2}{m}t}=\frac{k_1}{m}t
    \implies C'(t)=\frac{k_1}{m}t\mathrm{e}^{\frac{k_2}{m}t}\)。
    积分得
    \(\displaystyle C(t)=\frac{k_1}{k_2}t\mathrm{e}^{\frac{k_2}{m}t}-\frac{k_1m}{k_2^{2}}\mathrm{e}^{\frac{k_2}{m}t}+C_0\)，
    从而
    \(\displaystyle v(t)=\frac{k_1}{k_2}t-\frac{k_1m}{k_2^{2}}+C_0\mathrm{e}^{-\frac{k_2}{m}t}\)。
    再由 $v(0)=0$ 得 $C_0=\dfrac{k_1m}{k_2^{2}}$，故
    \(\displaystyle \boxed{v(t)=\frac{k_1}{k_2}t-\frac{k_1m}{k_2^{2}}\left(1-\mathrm{e}^{-\frac{k_2}{m}t}\right)}\)
\end{enumerate}
\end{solution}

\subsubsection{伯努利方程}

\begin{definition}[伯努利方程] \label{def:bernoulli-eq}
形如
\(\displaystyle \frac{\mathrm{d}y}{\mathrm{d}x}+P(x)\,y=Q(x)\,y^{n},\quad n \neq 0,1\)
的方程称为\textbf{伯努利方程}，其中 $n$ 称为伯努利指数。
\end{definition}

\begin{theorem}[伯努利方程的标准化方法]
伯努利方程可以通过特定的变量代换转化为一阶线性微分方程。具体推导步骤如下：

\begin{enumerate}
    \item \textbf{消除等式右侧的 $y$ 项}：假设 $y \neq 0$，将原方程两边同时除以 $y^{n}$，得到 \(\displaystyle y^{-n}\frac{\mathrm{d}y}{\mathrm{d}x}+P(x)\,y^{1-n}=Q(x)\)。
    \item 由 \( \displaystyle y^ny'=\frac{1}{n+1}(y^{n+1})' \) 联想到 \( \displaystyle y^{-n}\frac{\mathrm{d}y}{\mathrm{d}x}=\frac{1}{1-n}\dfrac{\mathrm{d}y^{1-n}}{\mathrm{d}x} \)，于是化为 \(\displaystyle \frac{\mathrm{d}y^{1-n}}{\mathrm{d}x}+(1-n)P(x)y^{1-n}=(1-n)Q(x)\)。
\end{enumerate}
即可得到原方程的通解。需要注意的是，当 $n>0$ 时，除以 $y^n$ 的操作预设了 $y \neq 0$，但显然 $y\equiv 0$ 也是原方程的一个解（此解可能包含在通解中，也可能是奇解，需单独说明）。
\end{theorem}

\begin{exercise}
解方程 $\displaystyle\frac{\mathrm{d}y}{\mathrm{d}x}=\frac{4}{x}y+x\sqrt{y}\;(y>0,x \neq 0)$。
\end{exercise}
\begin{solution}
\begin{enumerate}
    \item \textbf{识别方程类型}：原方程可改写为 $\displaystyle \frac{\mathrm{d}y}{\mathrm{d}x}-\frac{4}{x}y=xy^{\frac{1}{2}}$。这是一个伯努利指数为 $n=\frac{1}{2}$ 的伯努利方程。

    \item \textbf{变量代换}：方程两边同除以 $y^{\frac{1}{2}}$，得
    \[
    y^{-\frac{1}{2}}\frac{\mathrm{d}y}{\mathrm{d}x}
    -\frac{4}{x}y^{\frac{1}{2}}=x,\qquad
    z=y^{\frac{1}{2}},\qquad
    \frac{\mathrm{d}z}{\mathrm{d}x}
    =\frac{1}{2}y^{-\frac{1}{2}}\frac{\mathrm{d}y}{\mathrm{d}x}.
    \]

    \item \textbf{化为线性方程}：代入原式，得到
    \(\displaystyle 2\frac{\mathrm{d}z}{\mathrm{d}x}-\frac{4}{x}z=x
    \quad\Longrightarrow\quad
    \frac{\mathrm{d}z}{\mathrm{d}x}-\frac{2}{x}z=\frac{x}{2}\)。
    这已经是一个标准的一阶线性非齐次方程。

    \item \textbf{求解线性方程}：寻找积分因子 $\mu(x) = \exp\!\left(\int -\frac{2}{x}\,\mathrm{d}x\right) = x^{-2}$。方程两边同乘 $x^{-2}$ 后，
    \begin{align*}
    x^{-2}\frac{\mathrm{d}z}{\mathrm{d}x}-2x^{-3}z=\frac{1}{2}x^{-1},
    \frac{\mathrm{d}}{\mathrm{d}x}\left(zx^{-2}\right)=\frac{1}{2x},
    z=x^{2}\left(\frac{1}{2}\ln|x|+C\right).
    \end{align*}

    \item \textbf{回代求原通解}：将 $z = \sqrt{y}$ 代回，得
    \(\displaystyle \sqrt{y} = x^{2}\left(\frac{1}{2}\ln|x| + C\right)\)，两边平方即得原方程的通解 \(\displaystyle \boxed{y=x^{4}\left(C+\frac{1}{2}\ln|x|\right)^{2}}\)。
\end{enumerate}
\end{solution}

\subsubsection{黎卡提方程与刘维尔方法}

\begin{theorem}[黎卡提方程的刘维尔解法] \label{thm:riccati-liouville}
设已知黎卡提方程 $\frac{\mathrm{d}y}{\mathrm{d}x}=P(x)\,y^{2}+Q(x)\,y+R(x)$ 的一个特解 $y=\phi(x)$。通过以下严密的代换过程，可以求得其通解：

\begin{enumerate}
    \item \textbf{引入扰动函数}：作代换 $y=\phi(x)+u$，其中 $u=u(x)$ 为新的未知函数。对其求导得 $\frac{\mathrm{d}y}{\mathrm{d}x}=\phi'(x)+\frac{\mathrm{d}u}{\mathrm{d}x}$。

    \item \textbf{代入原方程并展开}：将 $y$ 和 $\frac{\mathrm{d}y}{\mathrm{d}x}$ 代入原黎卡提方程中：
    \(\displaystyle \phi'(x)+\frac{\mathrm{d}u}{\mathrm{d}x} = P(x)[\phi(x)+u]^{2} + Q(x)[\phi(x)+u] + R(x)\)
    展开等式右侧的括号：
    \(\displaystyle \phi'(x)+\frac{\mathrm{d}u}{\mathrm{d}x} = P(x)[\phi(x)^{2} + 2\phi(x)u + u^{2}] + Q(x)\phi(x) + Q(x)u + R(x)\)
    重新整理同类项，把不含 $u$ 的项、含 $u$ 的一次项和二次项分别组合：\(\displaystyle \phi'(x)+\frac{\mathrm{d}u}{\mathrm{d}x} = \left[P(x)\phi(x)^{2} + Q(x)\phi(x) + R(x)\right] + \left[2P(x)\phi(x) + Q(x)\right]u + P(x)u^{2}\)。

    \item \textbf{利用特解性质消元}：因为 $y=\phi(x)$ 是原方程的特解，所以它必定满足原方程，即恒有 $\phi'(x) = P(x)\phi(x)^{2} + Q(x)\phi(x) + R(x)$。
    利用这一性质，上式左右两边的第一项恰好完全抵消！等式大幅简化为：\(\displaystyle \frac{\mathrm{d}u}{\mathrm{d}x} = \underbrace{[2P(x)\phi(x)+Q(x)]}_{\text{记为 }S(x)}u + P(x)u^{2}\)。

    \item \textbf{化为伯努利方程并求解}：记 $S(x) = 2P(x)\phi(x)+Q(x)$，上式变为 \(\displaystyle \frac{\mathrm{d}u}{\mathrm{d}x} - S(x)u = P(x)u^{2}\)。这是一个指数为 $n=2$ 的\textbf{伯努利方程}！同除以 $u^2$ 得 \(\displaystyle u^{-2}\frac{\mathrm{d}u}{\mathrm{d}x} - S(x)u^{-1} = P(x)\)。
    令 $w = u^{-1}$，则 $\frac{\mathrm{d}w}{\mathrm{d}x} = -u^{-2}\frac{\mathrm{d}u}{\mathrm{d}x}$。代入后化为关于 $w$ 的线性方程：
    \(\displaystyle -\frac{\mathrm{d}w}{\mathrm{d}x} - S(x)w = P(x)
    \quad\Longrightarrow\quad
    \frac{\mathrm{d}w}{\mathrm{d}x} + S(x)w = -P(x)\)。

    \item \textbf{得出最终解的结构}：利用一阶线性方程的通解公式解出
    \[
    w(x)=\exp\!\left(-\int S(x)\,\mathrm{d}x\right)
    \left[
    C-\int P(x)\exp\!\left(\int S(x)\,\mathrm{d}x\right)\mathrm{d}x
    \right].
    \]
    由于 $u = \frac{1}{w}$ 且 $y = \phi(x) + u$，最终得到黎卡提方程的通解公式：
    \[
    \boxed{y=\phi(x)+\frac{\exp\!\left(\displaystyle\int S(x)\,\mathrm{d}x\right)}{C-\displaystyle\int P(x)\exp\!\left(\int S(x)\,\mathrm{d}x\right)\mathrm{d}x}}
    \]
\end{enumerate}
\end{theorem}

\begin{conclusion}[黎卡提方程的通解求法]
$\frac{\mathrm{d}y}{\mathrm{d}x}=P(x)\,y^{2}+Q(x)\,y+R(x)$若有特解$y=\phi(x)$，那么设$S(x)$是$S(x) = 2P(x)\phi(x)+Q(x)$，通解：
    \[
    \boxed{y=\phi(x)+\frac{\exp\!\left(\displaystyle\int S(x)\,\mathrm{d}x\right)}{C-\displaystyle\int P(x)\exp\!\left(\int S(x)\,\mathrm{d}x\right)\mathrm{d}x}}
    \]
\end{conclusion}

\subsection{恰当方程与积分因子}

许多一阶微分方程经过移项、去分母后，都可写成 \(\displaystyle M(x,y)\mathrm{d}x + N(x,y)\mathrm{d}y = 0\)。如果存在函数 $u(x,y)$ 使得 $M=\frac{\partial u}{\partial x}$、$N=\frac{\partial u}{\partial y}$，那么原方程便等价于 $\mathrm{d}u=0$，从而得到通解 $u(x,y)=C$。这样的方程称为\textbf{恰当微分方程}。

随意给出一个 $M(x,y)\mathrm{d}x + N(x,y)\mathrm{d}y = 0$，假设方程是恰当的，即确实存在这样一个 $u(x,y)$，使得 \(\displaystyle M=\frac{\partial u}{\partial x},\quad N=\frac{\partial u}{\partial y}\)。根据多元微积分中的克莱罗定理（又称二阶混合偏导数相等定理），如果函数 $u(x,y)$ 的二阶混合偏导数连续，那么求导的顺序可以交换：\(\displaystyle \frac{\partial^2 u}{\partial y \partial x}=\frac{\partial^2 u}{\partial x \partial y}\)。将 $M$ 和 $N$ 代入上式，可得
    \[
    \frac{\partial^2 u}{\partial y \partial x}
    =\frac{\partial}{\partial y}\left(\frac{\partial u}{\partial x}\right)
    =\frac{\partial M}{\partial y}=M_y,
    \qquad
    \frac{\partial^2 u}{\partial x \partial y}
    =\frac{\partial}{\partial x}\left(\frac{\partial u}{\partial y}\right)
    =\frac{\partial N}{\partial x}=N_x.
    \]
由此，我们得出判定条件 \(\displaystyle \frac{\partial M}{\partial y} = \frac{\partial N}{\partial x}\)。如果 $\frac{\partial M}{\partial y} \neq \frac{\partial N}{\partial x}$，方程不是恰当的，怎么办？能不能给整个方程乘上一个非零函数$\mu(x,y)$，使它变成恰当方程？若设$\mu(x,y)=X(x)Y(y)$，则新方程的恰当性条件
\begin{align*}
\displaystyle \frac{\partial \mu M}{\partial y} = \frac{\partial \mu N}{\partial x}&\Leftrightarrow \mu_yM+\mu M_y=\mu_xN+\mu N_x\Leftrightarrow\mu_xN-\mu_yM=\mu M_y-\mu N_x\\
&\Leftrightarrow NX'(x)Y(y)-MX(x)Y'(y)=X(x)Y(y)(M_y-N_x)\\
&\Leftrightarrow N\dfrac{X'(x)}{X(x)}-M\dfrac{Y'(y)}{Y(y)}=M_y-N_x
\end{align*}
等价于说恰当微分要求：如果我们能找到两个单变量函数 \(a(x),b(y)\)，使得
\[
M_y-N_x=N(x,y)a(x)-M(x,y)b(y)\implies (\ln X)'=a(x),(\ln Y)'=b(y)
\]
分别积分：
\[
\ln|X|=\int a(x)\,\mathrm{d}x,\ln|Y|=\int b(y)\,\mathrm{d}y\implies
X(x)=\exp\!\left(\int a(x)\,\mathrm{d}x\right),
Y(y)=\exp\!\left(\int b(y)\,\mathrm{d}y\right),
\]所以\[
\boxed{\mu(x,y)=
\exp\!\left(\int a(x)\,\mathrm{d}x\right)
\exp\!\left(\int b(y)\,\mathrm{d}y\right)}.
\]
积分时出现的非零常数因子可以省略，因为积分因子整体乘上一个非零常数，所得恰当方程的解集不变。


\begin{exercise}[用 $X(x)Y(y)$ 积分因子看换元题]
\[
y'=\frac{1}{x}+\mathrm e^y,\qquad
y'+\sin y+x\cos y+x=0.
\]
\end{exercise}

\begin{solution}
\begin{enumerate}
    \item 先化为标准形式$\left(-\frac1x-\mathrm e^y\right)\mathrm{d}x+\mathrm{d}y=0$，
    因此$M=-\frac1x-\mathrm e^y,N=1$，由$M_y-N_x=-\mathrm e^y$可知原方程不是恰当方程，设积分因子为$\mu=X(x)Y(y)$，代入 \(M=-\dfrac1x-\mathrm e^y,\ N=1\)，得
    \[
    1\cdot\frac{X'}{X}-\left(-\frac1x-\mathrm e^y\right)\dfrac{Y'}{Y}=-\mathrm e^y.
    \]
    右边只有一项 \(-\mathrm e^y\)。左边里唯一能产生 \(\mathrm e^y\) 的地方是$\mathrm e^y \dfrac{Y'}{Y}$，所以$\dfrac{Y'}{Y}=-1,Y=c\mathrm{e}^{-y}$，代入得
    \[\frac{X'}{X}+\left(-\frac1x-\mathrm e^y\right)=-\mathrm e^y\implies (\ln X)'=\frac1{x}\]
    那么$\ln X=\ln x+C,X=cx$，所以积分因子可取$\mu=x\mathrm e^{-y}$，它乘回原方程：
    \[\mu M=x\mathrm e^{-y}\left(-\frac1x-\mathrm e^y\right)=-\mathrm e^{-y}-x,\mu N=x\mathrm e^{-y}\implies\left(-\mathrm e^{-y}-x\right)\mathrm{d}x+x\mathrm e^{-y}\mathrm{d}y=0\]
    再验一次恰当性：
    \[
    \frac{\partial}{\partial y}\left(-\mathrm e^{-y}-x\right)=\mathrm e^{-y},
    \qquad
    \frac{\partial}{\partial x}\left(x\mathrm e^{-y}\right)=\mathrm e^{-y}.
    \]
    两者相等，说明积分因子找对了。现在求势函数 \(u(x,y)\)。因为
    \[
    u_y=x\mathrm e^{-y},
    \]
    所以先对 \(y\) 积分。积分时把 \(x\) 看成常数：
    \[
    u=\int x\mathrm e^{-y}\,\mathrm{d}y+\psi(x)
    =-x\mathrm e^{-y}+\psi(x).
    \]
    这里的 \(\psi(x)\) 不能省略，对 \(u\) 求 \(x\) 偏导：
    \[
    u_x=-\mathrm e^{-y}+\psi'(x)=-\mathrm e^{-y}-x\implies 
    \psi(x)=-\frac{x^2}{2}.
    \]
    所以势函数可取
    \[
    u(x,y)=-x\mathrm e^{-y}-\frac{x^2}{2}.
    \]
    恰当方程的通解是 \(u=C\).
    \item 标准形式：$\bigl(\sin y+x(1+\cos y)\bigr)\mathrm{d}x+\mathrm{d}y=0$，因此$M=\sin y+x(1+\cos y),N=1$，$M_y-N_x=\cos y-x\sin y$，不是恰当方程。设$\mu=X(x)Y(y)$，由
    \[N\frac{X'}{X}-M\frac{Y'}{Y}=\frac{X'}{X}-(\sin y+x(1+\cos y))\frac{Y'}{Y}=M_y-N_x=\cos y-x\sin y\]不妨$(1+\cos y)\frac{Y'}{Y}=\sin y$，
\[\ln Y=\int\frac{\sin y}{1+\cos y}\,\mathrm{d}y=-\int\frac{1}{t}\,\mathrm{d}t=-\ln|t|+C=-\ln|1+\cos y|+C\]
   可取$Y(y)=\frac{1}{1+\cos y}$，重新代入：
   \[\frac{X'}{X}-(\sin y+x(1+\cos y))\frac{\sin y}{1+\cos y}=\cos y-x\sin y\Rightarrow\frac{X'}{X}-\frac{\sin^2 y}{1+\cos y}=\cos y\]
即$\frac{X'}{X}-\frac{1-\cos^2 y}{1+\cos y}=\frac{X'}{X}-(1-\cos y)=\cos y$，所以$(\ln X)'=1\Rightarrow X=c\mathrm e^{x}$，可取积分因子$\mu=\frac{\mathrm e^x}{1+\cos y}$，把积分因子乘回原方程。先算新的 \(M\)：
    \[
    \mu M
    =\frac{\mathrm e^x}{1+\cos y}\bigl(\sin y+x(1+\cos y)\bigr)=\mathrm e^x\frac{\sin y}{1+\cos y}+x\mathrm e^x
    =\mathrm e^x\left(\tan\frac y2+x\right).
    \]
    再算新的$\mu N=\frac{\mathrm e^x}{1+\cos y}$，所以新方程是
    \[
    \mathrm e^x\left(\tan\frac y2+x\right)\mathrm{d}x
    +\frac{\mathrm e^x}{1+\cos y}\mathrm{d}y=0.
    \]
    求势函数 \(u\)。因为
    \[
    u_y=\mu N=\frac{\mathrm e^x}{1+\cos y},
    \]
    所以对 \(y\) 积分。积分时把 \(\mathrm e^x\) 看成常数：
    \[
    u=\mathrm e^x\int\frac{1}{1+\cos y}\,\mathrm{d}y+\psi(x)=\mathrm e^{x}\int\sec^2\frac y2\mathrm{d}\frac{y}{2}+\psi(x)=\mathrm e^{x}\tan\frac y2+\psi(x).
    \]
    再对 \(u\) 求 \(x\) 偏导：
    \[
    u_x=\mathrm e^x\tan\frac y2+\psi'(x)=\mu M=\mathrm e^x\left(\tan\frac y2+x\right)
    =\mathrm e^x\tan\frac y2+x\mathrm e^x\implies
    \psi'(x)=x\mathrm e^x.
    \]
    于是势函数为
    \[
    u(x,y)=\mathrm e^x\tan\frac y2+\mathrm e^x(x-1)
    =\mathrm e^x\left(\tan\frac y2+x-1\right).
    \]
    通解为
    \[
    \mathrm e^x\left(\tan\frac y2+x-1\right)=C\Leftrightarrow
    \boxed{\tan\frac y2=C\mathrm e^{-x}+1-x}.
    \]
    最后补充分支。上面 \(1+\cos y\)，因此必须单独检查$\cos y=-1\Leftrightarrow y=(2k+1)\pi,k\in\mathbb Z$，
    若 \(y\) 是这样的常值函数，则 \(y'=0,\ \sin y=0,\ \cos y=-1\)。代回原方程：
    \[
    y'+\sin y+x\cos y+x
    =0+0-x+x=0.
    \]
    所以$y=(2k+1)\pi$也是原方程的常值解，不能被积分因子分支漏掉。
\end{enumerate}
\end{solution}

\begin{exercise}[换元一例]
求微分方程初值问题
$\displaystyle y'\sec^{2}y+\frac{x}{1+x^{2}}\tan y=x,\quad y\big|_{x=0}=0$
的解。
\end{exercise}
\begin{solution}
先化为标准形式$\left(\frac{x}{1+x^2}\tan y-x\right)\mathrm{d}x+\sec^2 y\,\mathrm{d}y=0$，
因此$M=\frac{x}{1+x^2}\tan y-x,N=\sec^2 y$，由$M_y-N_x=\frac{x}{1+x^2}\sec^2 y$可知原方程不是恰当方程。设积分因子为$\mu=X(x)Y(y)$，代入得
\[
\sec^2 y\cdot\frac{X'}{X}-\left(\frac{x}{1+x^2}\tan y-x\right)\dfrac{Y'}{Y}=\frac{x}{1+x^2}\sec^2 y.
\]
右边含有 $\sec^2 y$。左边能产生 $\sec^2 y$ 的地方是 $\sec^2 y \frac{X'}{X}$，因此可取 $\dfrac{Y'}{Y}=0$，即 $Y(y)=1$。代入上式简化为
\[
\sec^2 y\cdot\frac{X'}{X}=\frac{x}{1+x^2}\sec^2 y\implies (\ln X)'=\frac{x}{1+x^2}.
\]
积分得 $\ln X = \frac{1}{2}\ln(1+x^2) = \ln\sqrt{1+x^2}$，可取 $X(x)=\sqrt{1+x^2}$。因此积分因子为 $\mu=\sqrt{1+x^2}$。将其乘回原方程：
\[
\left(\frac{x}{\sqrt{1+x^2}}\tan y-x\sqrt{1+x^2}\right)\mathrm{d}x+\sqrt{1+x^2}\sec^2 y\,\mathrm{d}y=0.
\]
现在求势函数 $u(x,y)$。因为
\[
u_y=\sqrt{1+x^2}\sec^2 y,~~u=\sqrt{1+x^2}\tan y+\psi(x).
\]
对 $u$ 求 $x$ 偏导：
\[
u_x=\frac{x}{\sqrt{1+x^2}}\tan y+\psi'(x)=\frac{x}{\sqrt{1+x^2}}\tan y-x\sqrt{1+x^2}\implies \psi'(x)=-x\sqrt{1+x^2}.
\]
积分得 $\psi(x)=-\frac{1}{3}(1+x^2)^{3/2}$，所以势函数可取
\[
u(x,y)=\sqrt{1+x^2}\tan y-\frac{1}{3}(1+x^2)^{3/2}.
\]
原方程的通解为
\[
\sqrt{1+x^2}\tan y-\frac{1}{3}(1+x^2)^{3/2}=C.
\]
代入初值条件 $y(0)=0$，得 $0-\frac{1}{3}=C \implies C=-\frac{1}{3}$。
因此特解为
\[
\sqrt{1+x^2}\tan y-\frac{1}{3}(1+x^2)^{3/2}=-\frac{1}{3}\implies \sqrt{1+x^2}\tan y=\frac{1}{3}\left((1+x^2)^{3/2}-1\right).
\]
两边同除以 $\sqrt{1+x^2}$ 即得显式特解：
\[
\boxed{y=\arctan\!\left[\frac13\!\left(1+x^{2}-\frac{1}{\sqrt{1+x^{2}}}\right)\right]}.
\]
\end{solution}


\begin{exercise} \label{ex:ode-rc-circuit}
如图所示的 $R$--$C$ 电路，电容 $C$ 初始无电荷。
合上开关 K 后电池 $E$ 对电容充电，求电压 $u_C$ 随时间的变化规律。
\end{exercise}

\begin{solution}
\begin{enumerate}
    \item \textbf{建立微分方程：}
    根据基尔霍夫电压定律（KVL），闭合回路的电压降之和等于电源电动势：\(u_C + u_R = E\)。
    根据欧姆定律，电阻上的电压为 $u_R = RI$。同时，电容器的电荷量 $Q$ 与电压 $u_C$ 的关系为 $Q = Cu_C$。因为充电电流即为电荷量的变化率，所以 $I = \frac{\mathrm{d}Q}{\mathrm{d}t} = C\frac{\mathrm{d}u_C}{\mathrm{d}t}$。
    将 $u_R$ 替换后代入 KVL 方程，得到关于 $u_C$ 的一阶常微分方程：\(\displaystyle RC\frac{\mathrm{d}u_C}{\mathrm{d}t} + u_C = E\)。

    \item \textbf{分离变量求解通解：} 由 \(\displaystyle RC\frac{\mathrm{d}u_C}{\mathrm{d}t}=E-u_C\)，得 \(\displaystyle
    \frac{\mathrm{d}u_C}{E-u_C}=\frac{1}{RC}\mathrm{d}t\)。
    两边积分得 \(\displaystyle -\ln|E-u_C|=\frac{t}{RC}+C_1\)，即 \(\displaystyle E-u_C=\mathrm{e}^{-C_1}\mathrm{e}^{-\frac{t}{RC}}\)，从而 \(\displaystyle u_C(t)=E-K\mathrm{e}^{-\frac{t}{RC}}\)。

    \item \textbf{代入初始条件确定特解：}
    已知电容初始无电荷，即在 $t=0$ 时 $u_C(0) = 0$。代入上式：\(\displaystyle 0 = E - K\mathrm{e}^0 \implies K = E\)。
    将 $K=E$ 代回通解，得到 \(\displaystyle u_C(t)=E\left(1-\mathrm{e}^{-\frac{t}{RC}}\right)\)。
    定义时间常数 $\tau=RC$：当 $t=3\tau$ 时 $u_C \approx 0.95E$，工程上通常认为充电基本完成。
\end{enumerate}
\end{solution}

\begin{exercise} \label{ex:ode-searchlight-mirror}
制造探照灯反射镜时要求点光源发出的光经镜面反射后平行射出，
试求反射镜面的几何形状。
\end{exercise}

\begin{solution}
\begin{enumerate}
    \item \textbf{导出微分方程：}
    取光源为坐标原点，$x$ 轴平行于反射方向。设镜面由曲线 $y=f(x)$ 绕 $x$ 轴旋转而成。
    根据光学反射定律和平面几何关系，可以推导出曲线上任一点 $(x, y)$ 处的切线斜率满足微分方程 \(\displaystyle \frac{\mathrm{d}y}{\mathrm{d}x} = \frac{y}{x+\sqrt{x^{2}+y^{2}}}\)（5.4）。

    \item \textbf{利用反函数视角进行齐次代换：}
    考虑到原方程含有根号且 $x$ 的结构复杂，我们将其转换为关于 $\frac{\mathrm{d}x}{\mathrm{d}y}$ 的方程：\(\displaystyle \frac{\mathrm{d}x}{\mathrm{d}y} = \frac{x+\sqrt{x^2+y^2}}{y} = \frac{x}{y} + \sqrt{\left(\frac{x}{y}\right)^2 + 1}\)。
    这是一个标准的齐次方程。令 $v = \frac{x}{y}$，则 $x = vy$。两边对 $y$ 求导得：
    \(\displaystyle \frac{\mathrm{d}x}{\mathrm{d}y} = v + y\frac{\mathrm{d}v}{\mathrm{d}y}\)。

    \item \textbf{代入并分离变量：}
    \(\displaystyle v+y\frac{\mathrm{d}v}{\mathrm{d}y}=v+\sqrt{v^2+1} \implies y\frac{\mathrm{d}v}{\mathrm{d}y}=\sqrt{v^2+1} \implies \frac{\mathrm{d}v}{\sqrt{v^2+1}}=\frac{\mathrm{d}y}{y}\)。

    \item \textbf{积分与化简回代：}
    \(\displaystyle \ln(v+\sqrt{v^2+1})=\ln|y|+\ln C_1 \implies v+\sqrt{v^2+1}=C_1y \quad (C_1>0)\)。
    回代并平方整理，得 \(\displaystyle \frac{x}{y}+\sqrt{\left(\frac{x}{y}\right)^2+1}=C_1y
    \implies \sqrt{x^2+y^2}=C_1y^2-x
    \implies y^2=\frac{2}{C_1}x+\frac{1}{C_1^2}\)。
    令新的常数 $C = \frac{1}{C_1}$，则方程化简为
    \(\displaystyle \boxed{y^{2}=C(C+2x)} \quad (C>0)\)。
    这是一条抛物线！因此反射镜面为旋转抛物面 $y^{2}+z^{2}=C(C+2x)$。
\end{enumerate}
\end{solution}

\subsection{\texorpdfstring{逐次积分型，含 $x$ 缺 $y$ 型，含 $y$ 缺 $x$ 型}{逐次积分型，含 x 缺 y 型，含 y 缺 x 型}}

\begin{remark}[两种降阶视角的选择]
设 $p=y'$。同一个二阶导数可以从两种角度理解：\(\displaystyle
y''=\frac{\mathrm{d}p}{\mathrm{d}x},\qquad
y''=\frac{\mathrm{d}p}{\mathrm{d}y}\cdot\frac{\mathrm{d}y}{\mathrm{d}x}
=p\,\frac{\mathrm{d}p}{\mathrm{d}y}\)。
到底选哪一种，关键看方程中还显含哪个量。
\begin{itemize}
    \item 若方程\textbf{不显含 $y$}，或整理后只含 $x,y',y''$，即 $y''=F(x,y')$，则令 $p=p(x)$ 最自然，用 \(\displaystyle y''=\frac{\mathrm{d}p}{\mathrm{d}x}\)。若能化为 \(\displaystyle \frac{\mathrm{d}p}{\mathrm{d}x}=X(x)P(p),\quad \frac{\mathrm{d}p}{P(p)}=X(x)\,\mathrm{d}x\)，便可直接分离变量。
    \item 若方程\textbf{不显含 $x$}，但显含 $y$，即 $y''=F(y,y')$，则令 $p=p(y)$ 最自然，用 \(\displaystyle y''=p\,\frac{\mathrm{d}p}{\mathrm{d}y}\)。
    若能化为 \(\displaystyle p\,\frac{\mathrm{d}p}{\mathrm{d}y}=Y(y)P(p),\quad \frac{p\,\mathrm{d}p}{P(p)}=Y(y)\,\mathrm{d}y\)，便可分离变量。
    \item 若方程只含 $y'$ 与 $y''$，即 $y''=F(y')$，两种方法通常都能用。若目标是先求 $p(x)$，用 $y''=\dfrac{\mathrm{d}p}{\mathrm{d}x}$ 更直接；若后续更关心 $p$ 与 $y$ 的关系，则用 $y''=p\dfrac{\mathrm{d}p}{\mathrm{d}y}$。
\end{itemize}
一句话记忆：缺 $y$ 看 $x$，用 $\dfrac{\mathrm{d}p}{\mathrm{d}x}$；缺 $x$ 看 $y$，用 $p\dfrac{\mathrm{d}p}{\mathrm{d}y}$。如果用了 $\dfrac{\mathrm{d}p}{\mathrm{d}x}$ 后右端仍含未知的 $y$，通常说明应改用 $p=p(y)$。
\end{remark}

\begin{exercise}[逐次积分]
求 $y'''=\mathrm{e}^{3x}-\cos x$ 的通解。
\end{exercise}
\begin{solution}
积分一次得 \(\displaystyle y''=\dfrac13\mathrm{e}^{3x}-\sin x+C_1\)，再积一次得 \(\displaystyle y'=\dfrac19\mathrm{e}^{3x}+\cos x+C_1x+C_2\)。
第三次积分即得通解：
\(\displaystyle \boxed{y=\frac{1}{27}\mathrm{e}^{3x}+\sin x+\frac{C_1}{2}x^{2}+C_2x+C_3}\)。
\end{solution}

\begin{exercise}
解初值问题 \(\displaystyle y''=\dfrac{3x^{2}}{1+x^{3}}\,y',\quad y|_{x=0}=1,\quad y'|_{x=0}=4\)。
\end{exercise}
\begin{solution}
方程不显含 $y$，属于 $y''=f(x,y')$ 型。令 $y'=p(x)$，则 $y''=p'$，代入得 \(\displaystyle p'=\frac{3x^{2}}{1+x^{3}}\,p\)。

分离变量得 \(\displaystyle \frac{\mathrm{d}p}{p}=\frac{3x^{2}\,\mathrm{d}x}{1+x^{3}}\)，积分得 \(\displaystyle \ln|p|=\ln(1+x^{3})+\ln C_1 \implies p=C_{1}(1+x^{3})\)。

由 $y'(0)=4$ 得 $C_1=4$，故 \(\displaystyle y'=4(1+x^{3})=4+4x^{3}\)。

再次积分得 \(\displaystyle y=\int(4+4x^{3})\,\mathrm{d}x=x^{4}+4x+C_{2}\)。由 $y(0)=1$ 得 $C_2=1$，故特解为 \(\displaystyle \boxed{y=x^{4}+4x+1.}\)
\end{solution}

\begin{exercise}
求 $y^{(5)}-\dfrac{1}{x}y^{(4)}=0$ 的通解。
\end{exercise}
\begin{solution}
方程不显含 $y$ 且最高阶导数为 $5$ 阶，缺 $y,y',y'',y'''$。
令 $p=y^{(4)}$，降为关于 $p$ 的一阶线性齐次方程 \(\displaystyle p'-\frac{1}{x}\,p=0\)。

分离变量得 \(\displaystyle \frac{\mathrm{d}p}{p}=\frac{\mathrm{d}x}{x}\)，故 $p=C_{1}x$。
于是 $y^{(4)}=C_{1}x$，对 $x$ 连续积分四次，得到
\begin{align*}
y&=\frac{C_1}{120}x^{5}+\frac{C_2}{24}x^{4}
+\frac{C_3}{6}x^{3}+\frac{C_4}{2}x^{2}+C_5x+C_6,\\
\text{常数重新吸并后}\qquad
&\boxed{y=C_1 x^{5}+C_2 x^{4}+C_3 x^{3}+C_4 x^{2}+C_5 x+C_{6}}.
\end{align*}
\end{solution}


\begin{exercise}
求 $y''=\dfrac{1+y'^{2}}{2y}$ 的通解。
\end{exercise}
\begin{solution}
方程不显含 $x$，属于 $y''=f(y,y')$ 型。令 $y'=p(y)$，则 $y''=p\,\dfrac{\mathrm{d}p}{\mathrm{d}y}$，代入原方程得 \(\displaystyle p\,\frac{\mathrm{d}p}{\mathrm{d}y}=\frac{1+p^{2}}{2y}\)。
分离变量并积分，得 \(\displaystyle \frac{2p\,\mathrm{d}p}{1+p^{2}}=\frac{\mathrm{d}y}{y}\)，即 \(\displaystyle \ln(1+p^{2})=\ln|y|+\ln C \quad (C>0)\)。
故 $1+p^{2}=Cy$，即 \(\displaystyle p=\pm\sqrt{Cy-1},\qquad \frac{\mathrm{d}y}{\mathrm{d}x}=\pm\sqrt{Cy-1}\)。
分离变量并整理得隐式通解：
\(\displaystyle \boxed{4(Cy-1)=C^{2}(x+C_{1})^{2}.}\)
\end{solution}

\begin{exercise}
求 $yy''-y'^{2}=0$ 的通解。
\end{exercise}
\begin{solution}
令 $y'=p(y)$，则 $y''=p\,\dfrac{\mathrm{d}p}{\mathrm{d}y}$。代入得
由 \(\displaystyle yp\,\frac{\mathrm{d}p}{\mathrm{d}y}-p^{2}=0\) 可化为
\(\displaystyle p\!\left(y\,\frac{\mathrm{d}p}{\mathrm{d}y}-p\right)=0\)。

由 $p=0$ 得 $y=C$；由 $y\,\dfrac{\mathrm{d}p}{\mathrm{d}y}-p=0$ 得 \(\displaystyle \frac{\mathrm{d}p}{p}=\frac{\mathrm{d}y}{y},\qquad p=C_{1}y\)。
即 $y'=C_{1}y$。此为可分离变量方程，通解为
\(\displaystyle \boxed{y=C_{2}\,\mathrm{e}^{C_{1}x}.}\)
\end{solution}


\begin{exercise}[导弹追击舰船] \label{ex:ode-pursuit}
如图所示，位于坐标原点 $O$ 的甲舰向正 $x$ 轴上点 $A(1,0)$ 处的乙舰发射制导导弹。导弹头始终对准乙舰。
已知乙舰以最大恒定速度 $v_{0}$ 沿平行于 $y$ 轴的直线行驶，导弹的速度为 $5v_{0}$。
求导弹轨迹方程，以及乙舰被击中时的位置。
\end{exercise}

\begin{solution}
建立坐标系如图所示。设导弹轨迹为 $y=y(x)$，经过时间 $t$ 导弹位于 $P(x,y)$，
此时乙舰位于 $Q(1,v_{0}t)$。

由于导弹始终瞄准目标，直线 $PQ$ 就是轨迹曲线在 $P$ 点的切线，斜率为
\(\displaystyle y'=\tan\theta=\frac{v_{0}t-y}{1-x}\)，即 \(\displaystyle v_{0}t=(1-x)y'+y\)。记作 (†)。

另一方面，弧长 $\widehat{OP}$ 的长度是距离 $|AQ|$ 的 5 倍（速度比 $5:1$），即 \(\displaystyle \int_{0}^{x}\sqrt{1+y'^{2}}\;\mathrm{d}x=5v_{0}t\)。记作 (‡)。

将 (†) 代入 (‡) 并消去 $t$，得 \(\displaystyle (1-x)y'+y=\frac15\int_{0}^{x}\sqrt{1+y'^{2}}\;\mathrm{d}x\)。记作 (§)。

对 $x$ 求导（利用初始条件 $y(0)=0,\;y'(0)=0$），得 \(\displaystyle -y'+(1-x)y''+y'=\frac15\sqrt{1+y'^{2}}\)，即 \(\displaystyle (1-x)y'' =\frac15\sqrt{1+y'^{2}}\)。

令 $y'=p$，则 \(\displaystyle
(1-x)p'=\frac15\sqrt{1+p^{2}},\qquad
\frac{\mathrm{d}p}{\sqrt{1+p^{2}}}=\frac{1}{5}\frac{\mathrm{d}x}{1-x}\)。
于是 \(\displaystyle \ln(p+\sqrt{1+p^{2}})=-\frac15\ln(1-x)+C\)。由 $p(0)=0$ 知 $C=0$，故
$p+\sqrt{1+p^{2}}=(1-x)^{-1/5}$。取倒数有 $\sqrt{1+p^{2}}-p=(1-x)^{1/5}$，两式相减并除以 2 得
\(\displaystyle y'=p=\frac12\!\left[(1-x)^{-1/5}-(1-x)^{1/5}\right]\)。

最后积分（利用 $y(0)=0$ 确定常数）得
\(\displaystyle \boxed{
y=-\frac58(1-x)^{4/5}+\frac5{12}(1-x)^{6/5}+\frac{5}{24}
}\)。

当导弹到达 $x=1$ 时，$y=\dfrac5{24}$。即乙舰航行到 $\bigl(1,\dfrac{5}{24}\bigr)$ 时被击中
\end{solution}


\begin{exercise}[悬链线的形状] \label{ex:ode-catenary}
如图所示，一条均匀、柔软且不可拉伸的细弦两端悬挂。当载荷沿弦均匀分布时，求弦在平衡状态下的形状曲线方程。
\end{exercise}

\begin{solution}
取弦最低点为坐标原点 $O(0,0)$。设弦的（常值）线密度为 $\rho$，最低点水平张力为 $H$，重力加速度为 $g$。平衡时弦的方程为 $y=y(x)$。

在弦上任一点 $B(x,y)$ 取微元弧段 $\widehat{OB}$ 做受力分析。设该处弦切线与 $x$ 轴夹角为 $\theta$，则
\(\displaystyle T\sin\theta = \rho g s,\ T\cos\theta = H\)，其中 $s$ 为 $\widehat{OB}$ 段的弧长，$T$ 为拉力。

两式相除得 \( \displaystyle H\tan\theta=\rho g\int_{0}^{x}\sqrt{1+y'^{2}}\;\mathrm{d}x \)。由于 $y'=\tan\theta$，上式化为 \(\displaystyle y'=\frac{\rho g}{H}\int_{0}^{x}\sqrt{1+y'^{2}}\;\mathrm{d}x\)。

记常数 $a=\dfrac{\rho g}{H}>0$，则 \(\displaystyle y''=a\sqrt{1+y'^{2}}\)。记作 (††)。
此方程不显含 $y$，属 $y''=f(x,y')$ 型。令 $y'=p$，则 \(\displaystyle \frac{\mathrm{d}p}{\mathrm{d}x}=a\sqrt{1+p^{2}}\)，从而 \(\displaystyle \frac{\mathrm{d}p}{\sqrt{1+p^{2}}} = a\mathrm{d}x\)，积分得 \(\displaystyle \ln(p+\sqrt{1+p^{2}})=ax+C\)。

由 $y'(0)=0$ 知 $C=0$，故 $p+\sqrt{1+p^{2}}=\mathrm{e}^{ax}$，解出 \(\displaystyle p=\dfrac12(\mathrm{e}^{ax}-\mathrm{e}^{-ax})\)。
，积分并利用 $y(0)=0$，得 \(\displaystyle y=\frac1{2a}\bigl(\mathrm{e}^{ax}+\mathrm{e}^{-ax}\bigr)+C_2=\frac1a\cosh(ax)-\frac{1}{a}\)，即
\(\displaystyle y=\frac{H}{\rho g}\cosh\!\left(\frac{\rho g}{H}x\right)-\frac{H}{\rho g}\)。
\end{solution}

\subsection{线性理论，特征根，比对系数}

\subsubsection{线性算子，叠加原理，线性相关/无关，朗斯基行列式，通解=齐次通解+非齐次特解}

\begin{definition}[二阶线性微分方程及初值问题]
称形如 $y''+P(x) y'+Q(x) y=f(x)$ 的方程为\textbf{二阶线性微分方程}（由于方程关于未知函数 $y$ 及其导数 $y', y''$ 都是一次的）。
\begin{itemize}
    \item 若 $f(x) \equiv 0$，则方程称为\textbf{二阶线性齐次方程}；
    \item 若 $f(x) \not\equiv 0$，则称为\textbf{二阶线性非齐次方程}；
    \item 若 $P(x), Q(x)$ 均为常数，则称方程为\textbf{二阶常系数线性方程}。
\end{itemize}
附加初始条件的定解问题：
\[
\begin{cases}
y''+P(x) y'+Q(x) y=f(x) \\
y(x_{0})=y_{0} \\
y'(x_{0})=y_{0}'
\end{cases}
\]
称为二阶线性微分方程的\textbf{初值问题}（Cauchy 问题）。
\end{definition}

首先，下面的定理保证了解的合理性与适定性。

\begin{theorem}[解的存在唯一性定理]
对于二阶线性初值问题，若函数 $P(x), Q(x)$ 及 $f(x)$ 在区间 $(a, b)$ 上连续，且 $x_0 \in (a, b)$，则存在唯一的一个函数 $y(x)$ 在 $(a, b)$ 上满足该初值问题。
\end{theorem}

该定理有一个重要的推论：对于对应的齐次方程 $y''+P(x) y'+Q(x) y=0$，如果满足初始条件 $y(x_{0})=0$ 且 $y'(x_{0})=0$，那么它的解\textbf{必定恒为零}（即平凡解 $y \equiv 0$）。

\begin{definition}[二阶线性微分算子]
对具有二阶连续导数的函数 $y \in C^{2}(a, b)$，由齐次方程左端的部分，我们定义一个映射 \(\displaystyle L: y \mapsto y''+P(x) y'+Q(x) y\)，即记为 $L[y]=y''+P(x) y'+Q(x) y$。该映射 $L$ 称为微分算子。
\end{definition}

\begin{exercise}[算子的计算]
设 $P(x)=0, Q(x)=x$，即算子为 $L[y]=y''+x y$。计算 $L[\cos x]$ 与 $L[x^3]$。
\end{exercise}
\begin{solution}
\begin{enumerate}
    \item 对于 $y=\cos x$，求二阶导数得 $y'' = -\cos x$。代入算子得 \(\displaystyle L[\cos x] = -\cos x + x \cos x = (x-1) \cos x\)。
    \item 对于 $y=x^{3}$，求二阶导数得 $y'' = 6x$。代入算子得 \(\displaystyle L[x^{3}] = 6 x + x \cdot x^{3} = x^{4} + 6 x\)。
\end{enumerate}
\end{solution}

\begin{property}[微分算子的线性性质]
微分算子 $L$ 满足以下两条核心性质（其中 $c, c_1, c_2$ 为常数， $y, y_1, y_2 \in C^2(a,b)$）：
齐次性 \(L[cy] = cL[y]\)，可加性 \(L[y_1+y_2] = L[y_1] + L[y_2]\)。
综合可得 \textbf{线性性}：$L[c_1 y_1 + c_2 y_2] = c_1 L[y_1] + c_2 L[y_2]$。
\end{property}

\begin{proof}
\begin{enumerate}
    \item 证明齐次性：
    \(\displaystyle L[c y] = (c y)''+P(x)(c y)'+Q(x)(c y) = c\left[y''+P(x) y'+Q(x) y\right] = c L[y]\)。
    \item 证明可加性：
    \(\displaystyle L[y_{1}+y_{2}] = (y_{1}+y_{2})''+P(x)(y_{1}+y_{2})'+Q(x)(y_{1}+y_{2}) = y_{1}''+y_{2}''+P(x) y_{1}'+P(x) y_{2}'+Q(x) y_{1}+Q(x) y_{2} = L[y_{1}]+L[y_{2}]\)。
\end{enumerate}
\end{proof}

\begin{theorem}[线性叠加原理]
设 $y_{1}$ 与 $y_{2}$ 为齐次方程 $L[y]=0$ 的两个解，则它们的任意线性组合 $y=c_{1} y_{1}+c_{2} y_{2}$ 也为该方程的解。
\end{theorem}

\begin{proof}
因为 $y_{1}, y_{2}$ 均为 $L[y]=0$ 的解，所以 $L[y_{1}]=0$ 且 $L[y_{2}]=0$。利用算子的线性性：
\(\displaystyle L[c_{1}y_{1}+c_{2}y_{2}] = c_{1}L[y_{1}]+c_{2}L[y_{2}] = 0 + 0 = 0\)。
所以 $y=c_{1} y_{1}+c_{2} y_{2}$ 也是 $L[y]=0$ 的解。
\end{proof}

现在产生一个问题：若 $L[y_{1}]=L[y_{2}]=0$，那么 $y=c_{1} y_{1}+c_{2} y_{2}$ 是否一定是 $L[y]=0$ 的\textbf{通解}？
结论：\textbf{不一定}。关键在于 $y_{1}$ 和 $y_{2}$ 是否足够“独立”。

例如考虑方程 $y''-y=0$。易见 $y_{1}=e^{x}$ 和 $y_{2}=2 e^{x}$ 均为解，但它们的线性组合 $y = c_{1}e^{x}+2c_{2}e^{x} = (c_{1}+2c_{2})e^{x}$ 实际上只包含一个独立常数，无法涵盖如 $e^{-x}$ 这样的解，故不是通解。
为了严格刻画解的独立性，引入以下概念：

\begin{definition}[线性相关与线性无关]
考虑定义在 $(a, b)$ 上的函数 $y_{1}(x)$ 和 $y_{2}(x)$。若存在不全为零的常数 $c_{1}, c_{2}$，使得对任意 $x \in (a, b)$ 均有 \(\displaystyle c_{1}y_{1}(x)+c_{2}y_{2}(x) \equiv 0\)，则称 $y_{1}(x)$ 与 $y_{2}(x)$ 在该区间上\textbf{线性相关}。否则，称为\textbf{线性无关}。也就是说，对于两个函数而言：
\begin{itemize}
    \item 线性相关 $\iff$ 其中一个函数是另一个函数的常数倍（即 $y_1(x) \equiv k y_2(x)$）。
    \item 线性无关 $\iff$ 它们的比值不是常数。
\end{itemize}
\end{definition}

\begin{definition}[朗斯基行列式 (Wronskian)]
对定义在 $(a, b)$ 上的可导函数 $y_{1}(x)$ 和 $y_{2}(x)$，我们称
\(\displaystyle W(x) := \begin{vmatrix}
y_{1}(x) & y_{2}(x) \\
y_{1}'(x) & y_{2}'(x)
\end{vmatrix} = y_{1}(x)y_{2}'(x) - y_{1}'(x)y_{2}(x)\)
为这两个函数的\textbf{朗斯基行列式}（Wronskian）。
\end{definition}

\begin{exercise}[线性相关性的判别]
利用朗斯基行列式或定义，判别下列各组函数在指定区间上的线性相关性：
(1) $1$ 与 $x$ 在 $(0, 1)$ 上；
(2) $\sin x$ 与 $\cos x$ 在 $(0, \pi)$ 上；
(3) $\cos^2 x$ 与 $\sin^2 x - 1$ 在 $(0, 1)$ 上。
\end{exercise}

\begin{solution}
\begin{enumerate}
    \item \textbf{证明 $1, x$ 线性无关}：
    若有常数 $c_{1}, c_{2}$ 使得 $c_{1} \cdot 1 + c_{2} x \equiv 0$，两边对 $x$ 求导得 $c_{2} \equiv 0$。代回原式得 $c_{1} = 0$。
    或者通过朗斯基行列式计算：\(\displaystyle W(x) = \begin{vmatrix} 1 & x \\ 0 & 1 \end{vmatrix} = 1 \neq 0\)。
    故 $1$ 和 $x$ 线性无关。

    \item \textbf{证明 $\sin x, \cos x$ 线性无关}：
    计算它们的朗斯基行列式：\(\displaystyle W(x) = \begin{vmatrix} \sin x & \cos x \\ \cos x & -\sin x \end{vmatrix} = -\sin^2 x - \cos^2 x = -1 \neq 0\)。
    因为 $W(x) \neq 0$，所以 $\sin x$ 和 $\cos x$ 在 $(0, \pi)$ 上线性无关。

    \item \textbf{证明 $\cos^2 x, \sin^2 x - 1$ 线性相关}：
    观察两函数关系，由三角恒等式知 $\sin^2 x - 1 = -\cos^2 x$。
    即存在不全为零的常数组合 $1 \cdot (\cos^2 x) + 1 \cdot (\sin^2 x - 1) \equiv 0$。故它们线性相关。
\end{enumerate}
\end{solution}

\begin{remark}[朗斯基行列式与线性相关性的关系]
对于一般的函数组，若线性相关，则必有 $W(x) \equiv 0$。但反之不一定成立（即 $W(x) \equiv 0$ 不能直接推导线性相关，存在分段函数反例）。
但是，如果这两个函数是某个\textbf{二阶齐次线性微分方程的解}，那么 $W(x)$ 具有更强的性质，由以下定理保证。
\end{remark}

\begin{theorem}[解的朗斯基行列式判别法]
如果 $y_{1}(x)$ 和 $y_{2}(x)$ 是二阶线性齐次方程 $L[y]=0$ 的两个解，则它们在区间 $(a,b)$ 上\textbf{线性无关}的充分必要条件是：对于任意的 $x \in (a, b)$，均有 $W(x) \neq 0$。
\end{theorem}

\begin{proof}
这里仅证明必要性（即若线性无关，则 $W(x) \neq 0$）。采用反证法。
\begin{enumerate}
    \item 假设存在某个点 $x_{0} \in (a, b)$，使得 $W(x_{0})=0$。
    \item 考虑关于未知数 $c_{1}, c_{2}$ 的齐次线性代数方程组：\(\displaystyle c_{1} y_{1}(x_{0}) + c_{2} y_{2}(x_{0}) = 0,\quad c_{1} y_{1}'(x_{0}) + c_{2} y_{2}'(x_{0}) = 0\)。
    \item 因为其系数矩阵的行列式恰为 $W(x_{0}) = 0$，根据齐次线性方程组非零解判别法，该方程组必然存在\textbf{非零解} $(c_{1}^*, c_{2}^*)$。
    \item 利用这组非零常数构造函数 $Y(x) = c_{1}^* y_{1}(x) + c_{2}^* y_{2}(x)$。由叠加原理，$Y(x)$ 是 $L[y]=0$ 的解。
    \item 结合第 2 步的方程组，该解满足初始条件：$Y(x_{0})=0$ 且 $Y'(x_{0})=0$。
    \item 根据解的存在唯一性定理推论，必有 \( \displaystyle Y(x) \equiv 0\)，即 \(\displaystyle c_{1}^* y_{1}(x) + c_{2}^* y_{2}(x) \equiv 0 \quad (\forall x \in (a, b))\)。
    这与 $y_1, y_2$ 线性无关的前提矛盾！故假设不成立，必定处处有 $W(x) \neq 0$。
\end{enumerate}
\end{proof}

\begin{theorem}[二阶线性齐次方程通解的结构]
设 $y_{1}(x)$ 和 $y_{2}(x)$ 是齐次方程 $L[y]=0$ 的两个\textbf{线性无关解}，则：
\(\displaystyle y = c_{1} y_{1}(x) + c_{2} y_{2}(x)\)
为 $L[y]=0$ 的通解（其中 $c_1, c_2$ 为任意常数）。称 $\{y_{1}, y_{2}\}$ 为方程的一个\textbf{基础解系}。该通解包含了方程的所有解。
\end{theorem}

\begin{proof}
先由前一定理知道：既然 $y_1,y_2$ 线性无关，那么它们的朗斯基行列式在所讨论区间上处处非零。任取一点 $x_0$，便有 \(\displaystyle W(x_0)=y_1(x_0)y_2'(x_0)-y_1'(x_0)y_2(x_0)\neq 0\)。
现在设 $y$ 是方程 $L[y]=0$ 的任意一个解。我们希望选取常数 $c_1,c_2$，使 \(\displaystyle c_1y_1(x_0)+c_2y_2(x_0)=y(x_0),\quad c_1y_1'(x_0)+c_2y_2'(x_0)=y'(x_0)\)。
由于系数行列式正是 $W(x_0)\neq0$，这个二元一次方程组有唯一解 $c_1,c_2$。令 \(\displaystyle z(x)=c_1y_1(x)+c_2y_2(x)\)。由叠加原理，$z$ 也是 $L[y]=0$ 的解，并且按照上面的构造，\(\displaystyle z(x_0)=y(x_0),\qquad z'(x_0)=y'(x_0)\)。于是 $y$ 与 $z$ 满足同一个二阶线性齐次方程，又具有相同初值。由初值问题解的唯一性，必有 $z(x)\equiv y(x)$。这说明任意解都可写成 \(\displaystyle y=c_1y_1+c_2y_2\)。
反过来，任意这样的线性组合当然仍是齐次方程的解，因此它正好给出全部解，也就是该方程的通解。
\end{proof}

\begin{theorem}[非齐次方程通解的结构]
设 $y_{1}, y_{2}$ 是对应齐次方程 $L[y]=0$ 的基础解系，而 $y^{*}$ 是非齐次方程 $L[y]=f(x)$ 的\textbf{一个特解}，则：
\(\displaystyle y = c_{1} y_{1} + c_{2} y_{2} + y^{*}\)
为非齐次方程 $L[y]=f(x)$ 的通解。
\end{theorem}

\begin{proof}
\begin{enumerate}
    \item \textbf{验证是解}：
    利用算子 $L$ 的线性性：
    \(\displaystyle L[y] = L[c_{1} y_{1} + c_{2} y_{2} + y^{*}] = c_{1} L[y_{1}] + c_{2} L[y_{2}] + L[y^{*}] = 0 + 0 + f(x) = f(x)\)。
    故它是解。
    \item \textbf{验证完备性（包含所有解）}：
    设 $\bar{y}(x)$ 为 $L[y]=f(x)$ 的任意一个解。考虑差值函数 $\bar{y} - y^{*}$：
    \(\displaystyle L[\bar{y} - y^{*}] = L[\bar{y}] - L[y^{*}] = f(x) - f(x) = 0\)。
    这说明 $\bar{y} - y^{*}$ 是对应齐次方程 $L[y]=0$ 的解。根据齐次通解结构定理，必然存在常数 $c_1, c_2$ 使得：
    \(\displaystyle \bar{y} - y^{*} = c_{1} y_{1} + c_{2} y_{2} \implies \bar{y} = c_{1} y_{1} + c_{2} y_{2} + y^{*}\)。
    证明完毕。
\end{enumerate}
\end{proof}

\begin{exercise}[非齐次通解的构造]
已知 $y''-y=x$，求其通解。
\end{exercise}
\begin{solution}
\begin{enumerate}
    \item 先看对应的齐次方程 \(\displaystyle y''-y=0\)。其特征方程为 $\lambda^2-1=0$，解得 $\lambda=\pm1$，故基础解系可取为 $y_1=e^x,\ y_2=e^{-x}$。
    \item 再求非齐次方程的一个特解。由于右端是一次多项式，试猜 \(\displaystyle y^*=Ax+B,\qquad 0-(Ax+B)=x\)。
    比较系数可得 $A=-1,\ B=0$，所以 $y^*=-x$。
    \item 因而原方程通解为 \(\displaystyle \boxed{y = c_{1} e^{x} + c_{2} e^{-x} - x}\)。
\end{enumerate}
\end{solution}

\begin{theorem}[非齐次方程的线性叠加原理]
设 $y_{1}^{*}$ 和 $y_{2}^{*}$ 分别是方程 $L[y]=f_{1}(x)$ 和 $L[y]=f_{2}(x)$ 的解，则 $y^{*} = y_{1}^{*}+y_{2}^{*}$ 是叠加方程 $L[y] = f_{1}(x) + f_{2}(x)$ 的特解。
\end{theorem}

\begin{proof}
这条结论直接来自线性算子 $L$ 的线性性。因为 $L[y_1^*]=f_1(x)$、$L[y_2^*]=f_2(x)$，且
\(\displaystyle L[y_1^*+y_2^*]=L[y_1^*]+L[y_2^*]=f_1(x)+f_2(x)\)。
因此 $y^*=y_1^*+y_2^*$ 确实是方程 $L[y]=f_1(x)+f_2(x)$ 的一个特解。
\end{proof}

\begin{exercise}[非齐次叠加原理应用]
求解方程 $y''+y=x+e^{x}$。
\end{exercise}
\begin{solution}
\begin{enumerate}
    \item 求对应齐次方程 $y''+y=0$ 的通解为 $y_h = c_{1} \sin x+c_{2} \cos x$。
    \item 将右端项拆分为 $f_1(x)=x$ 和 $f_2(x)=e^x$。
    \item 显然 $y_1^* = x$ 是 $y''+y=x$ 的特解。
    \item 将 $y_2^* = A e^x$ 代入 $y''+y=e^x$ 得 $2A e^x = e^x \implies A=1/2$，故 $y_2^* = \frac{1}{2} e^{x}$ 为第二个特解。
    \item 组合得到总的通解：\(\displaystyle y = c_{1} \sin x + c_{2} \cos x + x + \frac{1}{2} e^{x}\)。
\end{enumerate}
\end{solution}

\subsubsection{特征根公式做题法}

\begin{definition}[特征方程与特征根] \label{def:characteristic-eq}
称关于 $\lambda$ 的二次代数方程 \(\displaystyle \boxed{\lambda^{2}+p\lambda+q=0}\)为齐次方程 $y''+py'+qy=0$ 的\textbf{特征方程}，其根称为\textbf{特征根}或\textbf{特征值}。
\begin{itemize}
\item 若特征方程有两互异实根 $\lambda_1,\lambda_2$，则 $y''+py'+qy=0$ 的通解为
\(\displaystyle \boxed{y=c_1\,\mathrm{e}^{\lambda_1x}+c_2\,\mathrm{e}^{\lambda_2x}}\)。
\item 若特征方程有一对共轭单复根 $\lambda_{1,2}=\alpha\pm i\beta$ ($\beta \neq 0$)，则 $y''+py'+qy=0$ 的实值通解为
\(\displaystyle \boxed{y=\mathrm{e}^{\alpha x}\bigl(c_1\cos\beta x+c_2\sin\beta x\bigr)}\)。
\item 若特征方程有二重实根 $\lambda$（即 $\lambda^{2}+p\lambda+q=(\lambda-\lambda)^2=0$），则 $y''+py'+qy=0$ 的通解为
\(\displaystyle \boxed{y=(c_1+c_2x)\,\mathrm{e}^{\lambda x}}\)。
\end{itemize}
\end{definition}


\begin{property}[$n$ 阶方程的通解结构] \label{property:norder-homog-solution}
根据代数基本定理，$n$ 次特征方程在复数域中有 $n$ 个根（计重数）。按根的类型分类：
\begin{enumerate}
  \item[(i)] 每个\textbf{$k$ 重实根} $\lambda_i$ 贡献 $k$ 个线性无关解：
        $\mathrm{e}^{\lambda_ix},\;x\mathrm{e}^{\lambda_ix},\;\dots,\;x^{k-1}\mathrm{e}^{\lambda_ix}$；
  \item[(ii)] 每对\textbf{$k$ 重共轭复根} $\alpha_i\pm i\beta_i$ 贡献 $2k$ 个实值线性无关解：
  \begin{align*}
  \text{余弦族：}\quad&
  \mathrm{e}^{\alpha_i x}\cos\beta_i x,\quad x\mathrm{e}^{\alpha_i x}\cos\beta_i x,\quad \dots,\quad x^{k-1}\mathrm{e}^{\alpha_i x}\cos\beta_i x;\\
  \text{正弦族：}\quad&
  \mathrm{e}^{\alpha_i x}\sin\beta_i x,\quad x\mathrm{e}^{\alpha_i x}\sin\beta_i x,\quad \dots,\quad x^{k-1}\mathrm{e}^{\alpha_i x}\sin\beta_i x.
  \end{align*}
\end{enumerate}
\end{property}

\begin{exercise}
求解 $y''+y=0$。
\end{exercise}
\begin{solution}
先设试探解 $y=e^{\lambda x}$，代入方程得特征方程
\(\displaystyle \lambda^{2}+1=0\)，从而 \(\displaystyle \lambda_{1,2}=\pm i\)，故
\(\displaystyle y=c_1\cos x+c_2\sin x\)。
\end{solution}

\begin{exercise}
求解 $y''-y=0$。
\end{exercise}
\begin{solution}
同样令 $y=e^{\lambda x}$，可得特征方程
\(\displaystyle \lambda^{2}-1=0\)，从而 \(\displaystyle \lambda_1=1,\ \lambda_2=-1\)，故
\(\displaystyle y=c_1\mathrm{e}^{x}+c_2\mathrm{e}^{-x}\)。
\end{solution}

\begin{exercise}
解初值问题
$\begin{cases}
16y''-24y'+9y=0,\\
y(0)=4,\;y'(0)=2.
\end{cases}$
\end{exercise}
\begin{solution}
特征方程 $16\lambda^{2}-24\lambda+9=0$，即 $(4\lambda-3)^{2}=0$，
得二重根 $\lambda_{1,2}=3/4$。由情形三，通解为 $y=(c_1+c_2x)\,\mathrm{e}^{3x/4}$。

求导：$y'=c_2\,\mathrm{e}^{3x/4}+\dfrac34(c_1+c_2x)\,\mathrm{e}^{3x/4}$。
代入初值得 \(\displaystyle y(0)=c_1=4,\ y'(0)=c_2+\tfrac34 c_1=2\)，故 $c_2=-1$，从而
\(\displaystyle y=(4-x)\,\mathrm{e}^{3x/4}\)。
\end{solution}

\begin{exercise}
求解 $y^{(5)}+2y'''+y'=0$。
\end{exercise}
\begin{solution}
特征方程 $\lambda^5+2\lambda^3+\lambda=\lambda(\lambda^4+2\lambda^2+1)
=\lambda(\lambda^2+1)^2=0$。

根的情况：
\begin{itemize}
  \item $\lambda_1=0$（单实根）$\to$ 解 $1$；
  \item $\lambda_{2,3}=\pm i$（二重复根）$\to$ 解 $\cos x,\;x\cos x,\;\sin x,\;x\sin x$。
\end{itemize}
共 5 个线性无关解，通解为
\(\displaystyle \boxed{y=c_1+(c_2+c_3x)\cos x+(c_4+c_5x)\sin x}\)。
\end{solution}


\begin{theorem}[比较系数法] \label{thm:undetermined-coefficients-type2}
对于二阶常系数线性非齐次微分方程 $y'' + py' + qy = f(x)$，，若右端形如
\[
f(x) = \mathrm{e}^{\lambda x}\bigl[P_m(x)\cos\omega x+Q_n(x)\sin\omega x\bigr]
\]
（其中 $P_m(x), Q_n(x)$ 分别为 $m$ 次和 $n$ 次实多项式），则方程存在形如
\[
y^*=x^k\,\mathrm{e}^{\lambda x}\bigl[T_M(x)\cos\omega x+G_M(x)\sin\omega x\bigr]
\]
的特解。其中待定多项式 $T_M(x)$ 和 $G_M(x)$ 的次数均为 $M = \max(m, n)$。
比例因子 $x^k$ 中的 $k$ 仅取决于复数 $\sigma = \lambda + \mathrm{i}\omega$ 是否为特征方程 $r^2+pr+q=0$ 的根：
\begin{itemize}
    \item 若 $\sigma$ 不是特征根（即 $\phi(\sigma) \neq 0$），则取 $k=0$；
    \item 若 $\sigma$ 是单特征根，则取 $k=1$（此时若 $\omega \neq 0$，$\sigma$ 必为单根）；
    \item 若 $\sigma$ 是重特征根，则取 $k=2$（此时必有 $\omega=0$，且 $\lambda$ 为特征方程的二重实根）。
\end{itemize}
\end{theorem}


\begin{exercise}
求微分方程 $y''-2y'+2y=x\mathrm{e}^x\cos x$ 的一个特解。
\end{exercise}
\begin{solution}
特征方程为 $r^2-2r+2=0$，其特征根为 $r=1\pm \mathrm{i}$。
这里 $\lambda=1, \omega=1$，由于 $\sigma = 1+\mathrm{i}$ 是特征方程的单根，故取 $k=1$。
又因 $P_1(x)=x, Q_0(x)=0$，待定多项式最高次数 $M = \max(1, 0) = 1$。
设特解形式为：
\[
y^* = x\mathrm{e}^x\bigl[(Ax+B)\cos x + (Cx+D)\sin x\bigr] = \mathrm{e}^x\bigl[(Ax^2+Bx)\cos x + (Cx^2+Dx)\sin x\bigr]
\]
代入原方程左端，化简整理得：
\[
(y^*)'' - 2(y^*)' + 2y^* = \mathrm{e}^x\bigl[(4Cx + 2A + 2D)\cos x + (-4Ax + 2C - 2B)\sin x\bigr]
\]
与右端项 $x\mathrm{e}^x\cos x$ 比较系数，得：
\[
\begin{cases}
4C = 1 \\
2A + 2D = 0 \\
-4A = 0 \\
2C - 2B = 0
\end{cases}
\implies
A = 0,\ B = \frac{1}{4},\ C = \frac{1}{4},\ D = 0
\]
故所求的一个特解为：
\[
y^* = \frac{1}{4}x\mathrm{e}^x(\cos x + x\sin x)
\]
\end{solution}

\begin{exercise}
求 $y''-y=2x\,\mathrm{e}^{x}$ 的通解。
\end{exercise}
\begin{solution}
齐次通解 $y_h=c_1\mathrm{e}^x+c_2\mathrm{e}^{-x}$。
$\lambda=1$ 是特征方程的单根（$k=1$），故设特解形式为：
\(\displaystyle y^*=x\,\mathrm{e}^x(b_0x+b_1)=\mathrm{e}^x(b_0 x^2+b_1 x)\)。
比较系数得 \(\displaystyle 4b_0x+2(b_0+b_1)=2x\)，即
\(\displaystyle 4b_0=2,\ 2(b_0+b_1)=0\)，故
\(\displaystyle b_0=\frac12,\ b_1=-\frac12\)。
于是特解为 \(\displaystyle y^*=\frac12 x\mathrm{e}^x(x-1)\)，最终通解为
\(\displaystyle y=c_1\mathrm{e}^x+c_2\mathrm{e}^{-x}+\frac12 x\mathrm{e}^x(x-1)\)。
\end{solution}

\begin{exercise}
求 $y''+4y'+4y=\cos 2x$ 的一个特解。
\end{exercise}
\begin{solution}
特征根为 $\lambda_{1,2}=-2$（二重实根）。此处 $\lambda+i\omega=2i$ 不是特征根（$k=0$）。
设特解为 $y^*=A\cos 2x+B\sin 2x$。代入方程并比较 $\cos 2x,\sin 2x$ 的系数，
得 \(\displaystyle 8B=1,\ -8A=0\)，故 \(\displaystyle A=0,\ B=\frac18\)，从而
\(\displaystyle y^*=\frac18\sin 2x\)。
\end{solution}

\subsection{常数变易法和几个刘维尔公式}

\begin{conclusion}[黎卡提方程的通解求法]
$\frac{\mathrm{d}y}{\mathrm{d}x}=P(x)\,y^{2}+Q(x)\,y+R(x)$若有特解$y=\phi(x)$，那么设$S(x)$是$S(x) = 2P(x)\phi(x)+Q(x)$，通解：
    \[
    \boxed{y=\phi(x)+\frac{\exp\!\left(\displaystyle\int S(x)\,\mathrm{d}x\right)}{C-\displaystyle\int P(x)\exp\!\left(\int S(x)\,\mathrm{d}x\right)\mathrm{d}x}}
    \]
\end{conclusion}

除此以外，本节继续介绍几个比较难背的公式。
\begin{theorem}[常数变易法公式推导]
对于已经化为首项系数为 $1$ 的非齐次方程 $y''+P(x)y'+Q(x)y=f(x)$，设 $y_{1}(x), y_{2}(x)$ 为对应齐次方程 $y''+P(x)y'+Q(x)y=0$ 的基础解系，记
\(\displaystyle W(x)=\begin{vmatrix} y_1 & y_2 \\ y_1' & y_2' \end{vmatrix}=y_1y_2'-y_2y_1'\)，
则非齐次方程的一个特解为
\[
\boxed{
y^{*}(x) = -y_{1}(x) \int \frac{y_{2}(x) f(x)}{W(x)} \mathrm{d}x
+ y_{2}(x) \int \frac{y_{1}(x) f(x)}{W(x)} \mathrm{d}x
}.
\]
\end{theorem}
\begin{proof}
\begin{enumerate}
    \item \textbf{变易常数}：
    已知齐次通解为 $y = c_{1} y_{1} + c_{2} y_{2}$。设特解为 $y^{*}=c_{1}(x)y_{1}+c_{2}(x)y_{2}$。求一阶导数：
    \[
    (y^{*})' = c_{1}(x) y_{1}' + c_{2}(x) y_{2}' + \left[ c_{1}'(x) y_{1} + c_{2}'(x) y_{2} \right]
    \]
    为避免继续求导产生二阶导数，令 $c_{1}'(x) y_{1} + c_{2}'(x) y_{2} = 0$，从而 $(y^{*})' = c_{1}(x) y_{1}' + c_{2}(x) y_{2}'$。
    \[(y^{*})'' = c_{1}(x) y_{1}'' + c_{2}(x) y_{2}'' + c_{1}'(x) y_{1}' + c_{2}'(x) y_{2}'\]
    将 $y^*, (y^*)', (y^*)''$ 代入原非齐次方程 $y''+P(x) y'+Q(x) y=f(x)$ 中，按 $c_1(x), c_2(x)$ 以及剩余项重新分组：
    \[
    c_{1}(x)[y_{1}''+P(x) y_{1}'+Q(x) y_{1}] + c_{2}(x)[y_{2}''+P(x) y_{2}'+Q(x) y_{2}] + c_{1}'(x) y_{1}' + c_{2}'(x) y_{2}' = f(x).
    \]
    因为 $y_1, y_2$ 是齐次方程的解，方括号内各项恒等于 0，于是方程化为
    \(\displaystyle c_{1}'(x) y_{1}' + c_{2}'(x) y_{2}' = f(x)\)
    联立方程：
    \[
    \begin{cases}
    y_{1}(x) c_{1}'(x) + y_{2}(x) c_{2}'(x) = 0 \\
    y_{1}'(x) c_{1}'(x) + y_{2}'(x) c_{2}'(x) = f(x)
    \end{cases}
    \]
    该方程组的系数矩阵的行列式恰好是朗斯基行列式 $W(x) = y_{1}y_{2}' - y_{2}y_{1}'$。因为 $y_1, y_2$ 是基础解系，线性无关，故 $W(x) \neq 0$。
    根据克拉默法则（Cramer's Rule），解得
    \[
    c_{1}'(x) = \frac{\begin{vmatrix} 0 & y_{2} \\ f(x) & y_{2}' \end{vmatrix}}{W(x)} = -\frac{y_{2}(x) f(x)}{W(x)},\qquad
    c_{2}'(x) = \frac{\begin{vmatrix} y_{1} & 0 \\ y_{1}' & f(x) \end{vmatrix}}{W(x)} = \frac{y_{1}(x) f(x)}{W(x)}.
    \]
    后续显然。
\end{enumerate}
\end{proof}

\begin{exercise}[常数变易法的实际应用]
求方程 $y''+y=\sec x$ 的通解。
\end{exercise}

\begin{solution}
\begin{enumerate}
    \item \textbf{求齐次基础解系}：
    对应齐次方程 $y''+y=0$ 的特征方程为 $r^2+1=0$，得 $r = \pm i$。
    基础解系为 $y_{1}=\sin x, y_{2}=\cos x$。

    \item \textbf{设特解形式并列方程组}：
    令非齐次特解为 $y^{*} = c_{1}(x) \sin x + c_{2}(x) \cos x$。
    根据常数变易法推导出的标准方程组结构，直接写出关于 $c_1', c_2'$ 的方程组：
    \[
    \begin{cases}
    c_{1}'(x) \sin x + c_{2}'(x) \cos x = 0 \quad &(1) \\
    c_{1}'(x) \cos x - c_{2}'(x) \sin x = \sec x \quad &(2)
    \end{cases}
    \]

    \item \textbf{求解方程组中的导函数}：
    使用加减消元法。将方程 (1) 两边同乘 $\sin x$，方程 (2) 两边同乘 $\cos x$：
    \[
    \begin{cases}
    c_{1}'(x) \sin^2 x + c_{2}'(x) \sin x \cos x = 0 \\
    c_{1}'(x) \cos^2 x - c_{2}'(x) \sin x \cos x = \sec x \cdot \cos x = 1
    \end{cases}
    \]
    两式相加，消去 $c_2'$ 项，得 $c_{1}'(x)(\sin^2 x + \cos^2 x) = 1 \implies c_{1}'(x) = 1$；再代回方程 (1)，得 $\sin x + c_{2}'(x) \cos x = 0 \implies c_{2}'(x) = -\frac{\sin x}{\cos x} = -\tan x$。

    \item \textbf{积分求特解系数}：
    对导函数积分以求出系数函数（求特解时可令积分常数为0）：$c_{1}(x) = \int 1 \, \mathrm{d}x = x,\ c_{2}(x) = \int (-\tan x) \, \mathrm{d}x = \int \frac{-\sin x}{\cos x} \, \mathrm{d}x = \ln |\cos x|$。

    \item \textbf{写出最终通解}：
    代回特解假设式，并将齐次通解与特解相加，得到 \(\displaystyle y^{*} = x \sin x + \cos x \ln |\cos x|\)，从而 \(\displaystyle \boxed{y} = \boxed{C_{1} \sin x + C_{2} \cos x + x \sin x + \cos x \ln |\cos x|}\)。
\end{enumerate}
\end{solution}


\begin{theorem}[告知一个通解的情形]
设 $y_1$ 是标准型齐次方程 $y''+P(x)y'+Q(x)y=0$ 在所讨论区间上的一个不为零的解，可得
\[\displaystyle y_2(x)=y_1(x)\int \frac{1}{y_1^2(x)}e^{-\int P(x)\mathrm{d}x}\,\mathrm{d}x\]
\end{theorem}

\begin{proof}
由前文定理，设 $y_1(x), y_2(x)$ 为齐次方程 $y''+P(x)y'+Q(x)y=0$ 的基础解系，其朗斯基行列式为：
\[
W(x) = y_1 y_2' - y_2 y_1',~~W'(x) = y_1' y_2' + y_1 y_2'' - y_2' y_1' - y_2 y_1'' = y_1 y_2'' - y_2 y_1''
\]
将 $y_1'' = -P(x)y_1' - Q(x)y_1$ 和 $y_2'' = -P(x)y_2' - Q(x)y_2$ 代入上式：
\[
W'(x) = y_1 \bigl[-P(x)y_2' - Q(x)y_2\bigr] - y_2 \bigl[-P(x)y_1' - Q(x)y_1\bigr]= -P(x)(y_1 y_2' - y_2 y_1') = -P(x)W(x)
\]
这是一个关于 $W(x)$ 的一阶可分离变量微分方程，解得：
\[
W(x) = C \mathrm{e}^{-\int P(x)\mathrm{d}x}
\]
不失一般性，取常数 $C = 1$，即 $W(x) = \mathrm{e}^{-\int P(x)\mathrm{d}x}$。

另一方面，观察 $W(x)$ 的结构，由于 $y_1(x) \neq 0$，可将其变形为：
\[
W(x) = y_1^2 \left( \frac{y_2}{y_1} \right)'
\]
两边同除以 $y_1^2$，得：
\[
\left( \frac{y_2}{y_1} \right)' = \frac{W(x)}{y_1^2} = \frac{1}{y_1^2} \mathrm{e}^{-\int P(x)\mathrm{d}x}
\]
两边对 $x$ 积分，得：
\[
\frac{y_2}{y_1} = \int \frac{1}{y_1^2} \mathrm{e}^{-\int P(x)\mathrm{d}x} \,\mathrm{d}x\implies
y_2(x) = y_1(x) \int \frac{1}{y_1^2} \mathrm{e}^{-\int P(x)\mathrm{d}x} \,\mathrm{d}x
\]
\end{proof}

\begin{theorem}[已知一个齐次解时的非齐次通解]
设 $P(x),Q(x),f(x)$ 在区间 $I$ 上连续。考虑标准型非齐次方程 $y''+P(x)y'+Q(x)y=f(x)$。若已知其对应齐次方程 $y''+P(x)y'+Q(x)y=0$ 有一个在 $I$ 上不为零的解 $y_1(x)$，则先由刘维尔公式\[\displaystyle y_2(x)=y_1(x)\int \frac{1}{y_1^2(x)}e^{-\int P(x)\mathrm{d}x}\,\mathrm{d}x\]
此时 $y_1,y_2$ 构成对应齐次方程的一组基础解系，并且可取 $W(x)=y_1y_2'-y_2y_1'=e^{-\int P(x)\mathrm{d}x}$。
于是由常数变易法，非齐次方程的通解为
\[
\boxed{
y=C_1y_1+C_2y_2-y_1\int\frac{y_2f(x)}{W(x)}\,\mathrm{d}x
+y_2\int\frac{y_1f(x)}{W(x)}\,\mathrm{d}x
}
\]
也就是
\[
\boxed{
y=C_1y_1+C_2y_2-y_1\int y_2f(x)e^{\int P(x)\mathrm{d}x}\,\mathrm{d}x
+y_2\int y_1f(x)e^{\int P(x)\mathrm{d}x}\,\mathrm{d}x
}.
\]
\end{theorem}
\begin{proof}
首先，证明 $y_1(x), y_2(x)$ 构成对应齐次方程的基础解系。
由已知定理，齐次方程的第二个解为：
\[
y_2(x) = y_1(x) \int \frac{1}{y_1^2(x)} \mathrm{e}^{-\int P(x)\mathrm{d}x} \,\mathrm{d}x
\]
由于 $y_1(x) \neq 0$，两边同除以 $y_1(x)$ 得：
\[
\frac{y_2(x)}{y_1(x)} = \int \frac{1}{y_1^2(x)} \mathrm{e}^{-\int P(x)\mathrm{d}x} \,\mathrm{d}x
\]
因为被积函数 $\frac{1}{y_1^2(x)} \mathrm{e}^{-\int P(x)\mathrm{d}x} > 0$，所以其积分不为常数，即比值 $\frac{y_2(x)}{y_1(x)}$ 不是常数。
因此，$y_1(x)$ 与 $y_2(x)$ 线性无关，构成对应齐次方程的一组基础解系。齐次方程的通解为 $Y(x) = C_1 y_1(x) + C_2 y_2(x)$。

其次，求非齐次方程的通解。
根据二阶线性微分方程解的结构，非齐次方程的通解为齐次通解与非齐次特解之和，即 $y = C_1 y_1 + C_2 y_2 + y^*$。
由常数变易法公式定理，非齐次方程的一个特解为：
\[
y^{*}(x) = -y_{1}(x) \int \frac{y_{2}(x) f(x)}{W(x)} \mathrm{d}x
+ y_{2}(x) \int \frac{y_{1}(x) f(x)}{W(x)} \mathrm{d}x
\]
其中 $W(x)$ 为 $y_1, y_2$ 的朗斯基行列式。由前文推导可知：
\[
W(x) = y_1 y_2' - y_2 y_1' = \mathrm{e}^{-\int P(x)\mathrm{d}x}
\]
由此可得：
\[
\frac{1}{W(x)} = \mathrm{e}^{\int P(x)\mathrm{d}x}
\]
将 $\frac{1}{W(x)}$ 代入特解 $y^*(x)$ 的表达式中，得：
\[
y^*(x) = -y_1\int y_2 f(x)\mathrm{e}^{\int P(x)\mathrm{d}x}\,\mathrm{d}x
+ y_2\int y_1 f(x)\mathrm{e}^{\int P(x)\mathrm{d}x}\,\mathrm{d}x
\]
将特解 $y^*(x)$ 与齐次通解相加，即得非齐次方程的通解：
\[
y = C_1y_1+C_2y_2-y_1\int\frac{y_2f(x)}{W(x)}\,\mathrm{d}x
+y_2\int\frac{y_1f(x)}{W(x)}\,\mathrm{d}x
\]
也就是：
\[
y = C_1y_1+C_2y_2-y_1\int y_2f(x)\mathrm{e}^{\int P(x)\mathrm{d}x}\,\mathrm{d}x
+y_2\int y_1 f(x)\mathrm{e}^{\int P(x)\mathrm{d}x}\,\mathrm{d}x
\]
\end{proof}

\begin{exercise}[降阶法的应用]
已知方程 $(1-x^{2}) y''-2 x y'+2 y=0$ 有一个解为 $y_{1}=x$，求其通解。
\end{exercise}

\begin{solution}
\begin{enumerate}
    \item \textbf{化为首项系数为 1 的标准形式}：
    在应用公式前，必须将方程两边同除以 $(1-x^2)$，提取标准的 $P(x)$：
    \(\displaystyle y'' - \frac{2 x}{1-x^{2}} y' + \frac{2}{1-x^{2}} y = 0\)。
    此时得到 $P(x) = -\frac{2 x}{1-x^{2}}$。
    严格来说，下面的公式应先在不含 $x=\pm1$ 且不含 $y_1=x$ 零点的子区间上使用，最后所得第二解可以延拓到 $x=0$ 附近。

    \item \textbf{计算刘维尔公式中的指数部分}：
    利用凑微分法，令 $v = 1-x^2, \mathrm{d}v = -2x\mathrm{d}x$，则 \(\displaystyle -\int P(x)\,\mathrm{d}x=\int \frac{2x}{1-x^2}\,\mathrm{d}x=\int \frac{-\mathrm{d}v}{v}=-\ln|v|=-\ln|1-x^2|\)。
    在固定区间上，绝对值只差一个常数因子；该常数因子会被第二解的任意常数吸收。为简化计算，取
    \(\displaystyle e^{-\int P(x) \mathrm{d}x}=e^{-\ln|1-x^2|}\sim \frac{1}{1-x^2}\)。

    \item \textbf{计算外层积分求 $u(x)$}：
    将 $y_1^2 = x^2$ 和刚才求出的指数项代入公式，得 \(\displaystyle u(x)=\int \frac{1}{x^2}\cdot\frac{1}{1-x^2}\,\mathrm{d}x\)，且 \(\displaystyle \frac{1}{x^2(1-x^2)}=\frac{1}{x^2}+\frac{1}{1-x^2}\)。
    从而 \(\displaystyle u(x)=\int\left(\frac{1}{x^2}+\frac{1}{1-x^2}\right)\mathrm{d}x=-\frac{1}{x}+\frac{1}{2}\ln\left|\frac{1+x}{1-x}\right|\)。

    \item \textbf{写出第二解与最终通解}：
    根据降阶设定 $y_2 = x \cdot u(x)$，乘回 $y_1 = x$，得
    \begin{align*}
    y_{2}(x)
    &=x \left( -\frac{1}{x} + \frac{1}{2} \ln \left| \frac{1+x}{1-x} \right| \right)\\
    &=-1 + \frac{x}{2} \ln \left| \frac{1+x}{1-x} \right|.
    \end{align*}
    故该微分方程的通解为
    \[
    \boxed{y=C_{1} x + C_{2} \left( -1 + \frac{x}{2} \ln \left| \frac{1+x}{1-x} \right| \right)}.
    \]
\end{enumerate}
\end{solution}

\subsection{Euler 方程:变系数化为常系数}

Euler 方程虽不满足 ``常系数'' 条件，但可通过变量代换\textbf{精确地}化为常系数线性方程的特殊类型

\begin{definition}[Euler 方程 / Cauchy--Euler Equation] \label{def:euler-eq}
形如 \(\displaystyle x^n y^{(n)}+p_1 x^{n-1}y^{(n-1)}+\cdots+p_{n-1}xy'+p_n y=f(x)\)
的方程称为 \textbf{Euler 方程}（$n$ 阶），其中 $p_1,\dots,p_n$ 为常数。

其特点是：$y^{(k)}$ 前的系数恰好是 $x^k$ 的常数倍——各阶导数的 ``权重''
与自变量的幂次相匹配。
\end{definition}

\begin{theorem}[Euler 方程的标准化变换]
作自变量替换 $x=\mathrm{e}^t$（即 $t=\ln x$，要求 $x>0$；$x<0$ 时类似处理），则
\[
\frac{\mathrm{d}y}{\mathrm{d}x}=\frac{\mathrm{d}y}{\mathrm{d}t}\cdot\frac{\mathrm{d}t}{\mathrm{d}x}
=\mathrm{e}^{-t}\frac{\mathrm{d}y}{\mathrm{d}t},
\qquad
\frac{\mathrm{d}^2y}{\mathrm{d}x^2}=x^{-2}\left(\frac{\mathrm{d}^2y}{\mathrm{d}t^2}-\frac{\mathrm{d}y}{\mathrm{d}t}\right),
\]
更一般地，\(\displaystyle x^k\frac{\mathrm{d}^k y}{\mathrm{d}x^k}
=\frac{\mathrm{d}^k y}{\mathrm{d}t^k}
+\beta_1\frac{\mathrm{d}^{k-1}y}{\mathrm{d}t^{k-1}}
+\cdots+\beta_{k-1}\frac{\mathrm{d}y}{\mathrm{d}t}\)，
其中 $\beta_1,\dots,\beta_{k-1}$ 为仅依赖于 $k$ 的确定常数。
\end{theorem}
特别地，对\textbf{二阶 Euler 方程}
$x^2y''+p_1xy'+p_2y=f(x)$，经 $x=\mathrm{e}^t$ 变换后化为
\(\displaystyle \frac{\mathrm{d}^2 y}{\mathrm{d}t^2}+(p_1-1)\frac{\mathrm{d}y}{\mathrm{d}t}+p_2\,y=f(\mathrm{e}^t)\)，
这是一个关于自变量 $t$ 的\textbf{二阶常系数}线性方程！

\begin{remark}[变换的本质]
从算子的角度看，Euler 方程之所以能够化为常系数方程，是因为
\(\displaystyle D:=x\frac{\mathrm{d}}{\mathrm{d}x}\)
在变量替换 $x=\mathrm{e}^t$ 下恰好对应于对 $t$ 的普通求导算子。
\begin{enumerate}
    \item 设 $y=y(x)=y(\mathrm{e}^t)$，则由链式法则
    \(\displaystyle D y=x\frac{\mathrm{d}y}{\mathrm{d}x}
    =\frac{\mathrm{d}y}{\mathrm{d}t}\)。
    这说明在 $t$ 变量下，算子 $D$ 与 $\dfrac{\mathrm{d}}{\mathrm{d}t}$ 具有相同的作用。

    \item 需要特别注意：\textbf{$D^k$ 并不等于 $x^k\dfrac{\mathrm{d}^k}{\mathrm{d}x^k}$}。例如
    \(\displaystyle D^2y=D(Dy)=x\frac{\mathrm{d}}{\mathrm{d}x}\!\left(x\frac{\mathrm{d}y}{\mathrm{d}x}\right)
    =x^2y''+xy'\)。
    因而 $D^k$ 与 $x^k y^{(k)}$ 之间总会相差若干低阶导数项。

    \item 另一方面，由
    \(\displaystyle \frac{\mathrm{d}}{\mathrm{d}x}=x^{-1}D\)
    出发，可以直接计算
    \(\displaystyle x^2\frac{\mathrm{d}^2y}{\mathrm{d}x^2}=x^2(x^{-1}D)(x^{-1}D)y=D(D-1)y\)，
    \(\displaystyle x^3\frac{\mathrm{d}^3y}{\mathrm{d}x^3}=x^3(x^{-1}D)(x^{-1}D)(x^{-1}D)y=D(D-1)(D-2)y\)。
    一般地，有 \(\displaystyle \boxed{x^k\frac{\mathrm{d}^k y}{\mathrm{d}x^k}=D(D-1)\cdots(D-k+1)y}\)。

    \item 因此，Euler 方程左端各项都可以写成关于 $D$ 的多项式作用在 $y$ 上；而在替换 $x=\mathrm{e}^t$ 后，$D$ 又化为 $\dfrac{\mathrm{d}}{\mathrm{d}t}$，于是原方程便转化为关于 $t$ 的常系数线性方程。
\end{enumerate}
\end{remark}

\begin{exercise}
求解 $x^2y''-2y=x\;(x>0)$。
\end{exercise}
\begin{solution}
令 $t=\ln x$（即 $x=\mathrm{e}^t$），则方程化为
\(\displaystyle \frac{\mathrm{d}^2 y}{\mathrm{d}t^2}-\frac{\mathrm{d}y}{\mathrm{d}t}-2y=\mathrm{e}^t\)。
特征方程 $\lambda^2-\lambda-2=0$，根为 $\lambda_1=2,\;\lambda_2=-1$。
齐次通解 $y_h=c_1\mathrm{e}^{2t}+c_2\mathrm{e}^{-t}$。

$\lambda=1$ 不是特征根，设特解 $y^*=A\mathrm{e}^t$。代入得
\(\displaystyle (A-A-2A)\mathrm{e}^t=\mathrm{e}^t\Rightarrow -2A=1\Rightarrow A=-\dfrac12\)。

$t$-空间中的通解：$y=c_1\mathrm{e}^{2t}+c_2\mathrm{e}^{-t}-\dfrac12\mathrm{e}^t$。
代回 $x=\mathrm{e}^t$ 得 \(\displaystyle y(x)=c_1x^2+\frac{c_2}{x}-\frac{x}{2}\)。
\end{solution}

\begin{exercise}[积分方程转化为微分方程]
已知对所有实数 $x$ 有
$f'(x)=x^2+\displaystyle\int_0^x f(t)\,\mathrm{d}t$，且 $f(0)=2$。求 $f(x)$。
\end{exercise}
\begin{solution}
对给定式两边关于 $x$ 求导（注意右端是变上限积分，需用微积分基本定理），得
\(\displaystyle f''(x)=2x+f(x)\)，
即 $f''-f=2x$。特征方程 $\lambda^2-1=0$，根为 $\pm 1$。

观察可知 $y^*=-2x$ 是一个特解（代入验证：$0-(-2x)=2x$ ）。
通解 $f=c_1\mathrm{e}^x+c_2\mathrm{e}^{-x}-2x$。

由条件 $f(0)=2$，又 $f'(x)=c_1\mathrm{e}^x-c_2\mathrm{e}^{-x}-2$，故 $f'(0)=c_1-c_2-2=0$。
联立 \(\displaystyle c_1+c_2=2,\ c_1-c_2=2\) 得 $c_1=2,\ c_2=0$，从而
\(\displaystyle f(x)=2(\mathrm{e}^x-x)\)。
\end{solution}

\subsection{方程组专题：分类、二阶系统与消元}

\begin{definition}[微分方程组] \label{def:ode-system}
含有 $n$ 个未知函数 $y_1(t),y_2(t),\dots,y_n(t)$ 的 $n$ 个微分方程联立的系统称为\textbf{微分方程组}，
要求方程的个数等于未知函数的个数，且联立方程对每个未知函数及其出现的各阶导数都是\textbf{线性}的（即一次方）。
\end{definition}


为便于叙述，下面以两个未知函数的情形展开讨论（三个及以上情形完全类似）。

设自变量为 $t$，因变量（未知函数）为 $x=x(t)$ 和 $y=y(t)$。

\begin{definition}[一阶微分方程组的一般形式] \label{def:first-order-system-general}
\textbf{隐式形式}（Implicit Form）：\(\displaystyle F_1\!\left(t,x,y,\dfrac{\mathrm{d}x}{\mathrm{d}t},\dfrac{\mathrm{d}y}{\mathrm{d}t}\right)=0,\qquad F_2\!\left(t,x,y,\dfrac{\mathrm{d}x}{\mathrm{d}t},\dfrac{\mathrm{d}y}{\mathrm{d}t}\right)=0\)。
其中 $F_1,F_2$ 是五个变量的已知连续函数。\\
\textbf{显式/典则形式}（Explicit / Canonical Form）：\(\displaystyle \dfrac{\mathrm{d}x}{\mathrm{d}t}=f_1(t,x,y),\qquad \dfrac{\mathrm{d}y}{\mathrm{d}t}=f_2(t,x,y)\) \textnormal{($\star$)}。
其中 $f_1,f_2$ 是在某区域 $D\subset\mathbb R^3$ 上定义的已知函数。
\end{definition}

\begin{note}[提醒：先判能否由隐式化为显式]
并非所有的隐式方程组都能写成典则形式 ($\star$)。根据隐函数定理，能够从中局部解出 $\dfrac{\mathrm{d}x}{\mathrm{d}t}$ 和 $\dfrac{\mathrm{d}y}{\mathrm{d}t}$ 的前提是雅可比行列式不为零，即 \(\displaystyle \frac{\partial(F_1,F_2)}{\partial(x',y')} \neq 0\)。
在以下的讨论中，我们总假设方程组已经化简为典则形式。
\end{note}

\begin{definition}[一阶方程组的初值问题] \label{def:ivp-sys}
在典则形式 ($\star$) 上附加初始条件 \(\displaystyle x\big|_{t=t_0}=x_0,\qquad y\big|_{t=t_0}=y_0\)，
称为一阶微分方程组的\textbf{初值问题}（或 Cauchy 问题），其中 $(t_0;x_0,y_0)\in D$ 为给定点。
\end{definition}

初值问题的几何意义在于：在由 $(t,x,y)$ 构成的相空间中，指定了一个确定的起始点，要求找出恰好经过该点的积分曲线。

\begin{definition}[一阶方程组的通解] \label{def:general-solution-sys}
若含有两个独立任意常数 $C_1,C_2$ 的函数组 \(\displaystyle x=\varphi_1(t;C_1,C_2),\qquad y=\varphi_2(t;C_1,C_2)\) \textnormal{($\star\star$)}
恒满足方程组 ($\star$)，且对任意适当的初值 $(t_0;x_0,y_0)$，总存在特定的常数 $C_1^*,C_2^*$ 使得该函数组满足初始条件，则称 ($\star\star$) 为该方程组的\textbf{通解}（General Solution）。
\end{definition}

\begin{remark}[通解的独立性]
``独立任意常数'' 的含义是：$C_1,C_2$ 无法被合并或消去——它们确实代表了二维解空间的两个自由度。
这在几何上对应于：通解是一族两参数曲线族，覆盖了相平面上的每一点（在存在唯一性区域内）。
\end{remark}

有时通解以隐式形式给出更为方便：

\begin{definition}[一阶方程组的通积分] \label{def:first-integral-sys}
若含参数 $C_1,C_2$ 的隐式方程组 \(\displaystyle \Phi_1(t,x,y;C_1,C_2)=0,\qquad \Phi_2(t,x,y;C_1,C_2)=0\) \textnormal{($\star\star\star$)}
所确定的每一组隐函数解 $(x(t),y(t))$ 均构成原方程组的解，则称 ($\star\star\star$) 为该方程组的\textbf{通积分}（General Integral）或\textbf{隐式通解}。
\end{definition}

\subsubsection{二阶微分方程组}

将上述概念自然推广到二阶情形：

\begin{definition}[二阶微分方程组的一般形式]
\textbf{隐式形式}：\(\displaystyle F_1\!\left(t,x,y,x',y',x'',y''\right)=0,\qquad F_2\!\left(t,x,y,x',y',x'',y''\right)=0\)。
\textbf{典则形式}：\(\displaystyle \dfrac{\mathrm{d}^{2}x}{\mathrm{d}t^{2}}=f_1\!\left(t,x,y,x',y'\right),\qquad \dfrac{\mathrm{d}^{2}y}{\mathrm{d}t^{2}}=f_2\!\left(t,x,y,x',y'\right)\)。
\end{definition}

\begin{definition}[二阶方程组的初值问题] \label{def:ivp-sys-2nd}
对于二阶方程组，初值问题需要附加初始位置和速度条件：
\(\displaystyle x(t_0)=x_0,\ y(t_0)=y_0,\ x'(t_0)=x'_0,\ y'(t_0)=y'_0\)。
共需四个条件，对应于确定特解中的四个任意常数。
\end{definition}

\begin{definition}[二阶方程组的通解与通积分]
记 $\mathbf C=(C_1,C_2,C_3,C_4)$。含有\textbf{四个}独立任意常数的解表达式
\(\displaystyle x=\psi_1(t;\mathbf C),\quad y=\psi_2(t;\mathbf C)\)，
或对应的隐式形式（含四常数的两个独立关系式）称为二阶方程组的\textbf{通解}（或\textbf{通积分}）。
\end{definition}

方程组的``阶''定义为其中所含最高阶导数的阶数。事实上，任何高阶方程组都可以通过引入新的中间变量，降维转化为一个维度更高的一阶方程组；例如将 $2$ 元二阶组通过设 $u=x,\ v=y,\ w=x',\ z=y'$ 转化为 $4$ 维一阶组。因此 $n$ 阶 $m$ 元微分方程组的通解应含有 $n\times m$ 个独立的任意常数。

\subsubsection{常系数线性微分方程组的消元法}

\begin{exercise}[消元法 I：重根情形] \label{ex:sys-elim-ex1}
求解常系数线性齐次方程组 \(\displaystyle \dfrac{\mathrm{d}x}{\mathrm{d}t}=-3x-y,\qquad \dfrac{\mathrm{d}y}{\mathrm{d}t}=x-y\)。
\end{exercise}
\begin{solution}
由第一式消去 $y$，可将系统化为关于 $x$ 的二阶方程：
\begin{enumerate}
    \item \textbf{选择：} 由第一式中可以最轻易地解出 $y$：\(\displaystyle y=-\frac{\mathrm{d}x}{\mathrm{d}t}-3x\)，再求导得
    \(\displaystyle \frac{\mathrm{d}y}{\mathrm{d}t}=-\frac{\mathrm{d}^{2}x}{\mathrm{d}t^{2}}-3\frac{\mathrm{d}x}{\mathrm{d}t}\)。
    \item \textbf{代入：} 将上面两式代入原系统的第二式 $y'=x-y$ 中，得
    \(\displaystyle -\frac{\mathrm{d}^{2}x}{\mathrm{d}t^{2}}-3\frac{\mathrm{d}x}{\mathrm{d}t}=x-\left(-\frac{\mathrm{d}x}{\mathrm{d}t}-3x\right)\)，整理为
    \(\displaystyle \frac{\mathrm{d}^{2}x}{\mathrm{d}t^{2}}+4\frac{\mathrm{d}x}{\mathrm{d}t}+4x=0\)。
    \item \textbf{求解单方程：} 这时得到的二阶方程的特征方程为 $\lambda^{2}+4\lambda+4=(\lambda+2)^{2}=0$，存在二重实根 $\lambda_{1,2}=-2$。
    故 $x$ 的通解为 \(\displaystyle x(t)=(C_1+C_2t)\,\mathrm{e}^{-2t}\)。
    \item \textbf{回代：} 为了求出 $y$，先计算
    \(\displaystyle x'(t)=C_2\mathrm{e}^{-2t}-2(C_1+C_2t)\mathrm{e}^{-2t}=\mathrm{e}^{-2t}(C_2-2C_1-2C_2t)\)。
    代回前面的 $y$ 表达式，得
    \(\displaystyle y(t)=-\mathrm{e}^{-2t}(C_2-2C_1-2C_2t)-3(C_1+C_2t)\mathrm{e}^{-2t}=-\mathrm{e}^{-2t}(C_1+C_2+C_2t)\)。
\end{enumerate}
最终通解为 \(\displaystyle \boxed{x(t)=(C_1+C_2t)\,\mathrm{e}^{-2t},\qquad y(t)=-(C_1+C_2+C_2t)\,\mathrm{e}^{-2t}}\)。
\end{solution}

\begin{exercise}[消元法 II：带初值条件的非齐次型转化] \label{ex:sys-elim-ex2}
求解定解问题 \(\displaystyle \dfrac{\mathrm{d}x}{\mathrm{d}t}=3x-2y,\qquad \dfrac{\mathrm{d}y}{\mathrm{d}t}=2x-y,\qquad x(0)=1,\ y(0)=0\)。
\end{exercise}
\begin{solution}
本次我们尝试从第二式出发消去 $x$（体现路线选择的灵活性）。
\begin{enumerate}
    \item \textbf{选择：} 由第二式解出 $x$：\(\displaystyle x=\frac12\left(\frac{\mathrm{d}y}{\mathrm{d}t}+y\right)\)，再求导得
    \(\displaystyle \frac{\mathrm{d}x}{\mathrm{d}t}=\frac12\left(\frac{\mathrm{d}^{2}y}{\mathrm{d}t^{2}}+\frac{\mathrm{d}y}{\mathrm{d}t}\right)\)。
    \item \textbf{代入：} 将上面两式代入第一式 $x'=3x-2y$ 中，得
    \(\displaystyle \frac12\left(y''+y'\right)=\frac32(y'+y)-2y=\frac32y'-\frac12y\)，整理为
    \(\displaystyle y''-2y'+y=0\)。
    \item \textbf{求解单方程：} 这时得到的二阶方程的特征方程为 $\lambda^{2}-2\lambda+1=(\lambda-1)^{2}=0$，得二重根 $\lambda_{1,2}=1$，故
    \(\displaystyle y(t)=(C_1+C_2t)\,\mathrm{e}^{t}\)。
    \item \textbf{回代：} 为了计算 $x$，先求
    \(\displaystyle y'(t)=C_2\mathrm{e}^t+(C_1+C_2t)\mathrm{e}^t=\mathrm{e}^t(C_1+C_2+C_2t)\)。
    代回前面的 $x$ 表达式，得
    \(\displaystyle x(t)=\frac12\left[\mathrm{e}^{t}(C_1+C_2+C_2t)+\mathrm{e}^{t}(C_1+C_2t)\right]=\mathrm{e}^{t}\left(C_1+\frac{1}{2}C_2+C_2t\right)\)。
    \item \textbf{定初值：} 代入 $t=0$，有
    \(\displaystyle y(0)=C_1=0,\qquad x(0)=C_1+\frac{1}{2}C_2=1 \implies \frac{1}{2}C_2=1 \implies C_2=2\)。
\end{enumerate}
将 $C_1=0, C_2=2$ 代入通解，得特解 \(\displaystyle \boxed{x(t)=(1+2t)\,\mathrm{e}^{t},\qquad y(t)=2t\,\mathrm{e}^{t}}\)。
\end{solution}

\section{习题}

\subsection{第一节习题：微分方程的基本概念}

\begin{exercise}[相关概念]
回答下列问题：（1）什么叫微分方程？（2）微分方程的阶如何定义？（3）什么叫微分方程的通解？通解中含多少个任意常数？（4）什么叫初值问题？什么叫微分方程的特解？
\end{exercise}
\begin{solution}
（1）含有未知函数及其导数（或微分）的方程称为微分方程，它刻画的是“函数值”和“变化率”之间的关系，而不是只给出某个孤立数值。（2）微分方程中出现的未知函数的最高阶导数的阶数，称为该微分方程的阶；例如含有 $\dfrac{\mathrm{d}^2y}{\mathrm{d}x^2}$ 且不含更高阶导数时，就是二阶微分方程。（3）如果微分方程的解中含有任意常数，且这些常数相互独立，并且个数恰好等于方程的阶数，就称这样的解为通解；这里“个数等于阶数”是因为 $n$ 阶方程通常需要 $n$ 个条件来唯一确定一个特解。（4）初值问题是指附加了初始条件的微分方程求解问题；特解则是把通解中的任意常数用这些条件确定以后得到的具体解，因此不再含自由常数。
\end{solution}

\begin{exercise}[验证解]
证明对任意常数 $P_0$，函数 $P(t)=P_0\mathrm{e}^t$ 满足微分方程 $\displaystyle \frac{\mathrm{d}P}{\mathrm{d}t}=P$。
\end{exercise}
\begin{solution}
对 $P(t)=P_0\mathrm{e}^t$ 求导，得 \(\displaystyle \frac{\mathrm{d}P}{\mathrm{d}t}=P_0\mathrm{e}^t=P(t)\)。
所以该函数确为微分方程的解。
\end{solution}

\begin{exercise}[求参数]
设 $Q=C\mathrm{e}^{kt}$ 满足微分方程 $\displaystyle \frac{\mathrm{d}Q}{\mathrm{d}t}=-0.03Q$，问  $k$ 为何值？
\end{exercise}
\begin{solution}
对 $Q=C\mathrm{e}^{kt}$ 求导，得 \(\displaystyle \frac{\mathrm{d}Q}{\mathrm{d}t}=Ck\mathrm{e}^{kt}=kQ\)，
与 $\frac{\mathrm{d}Q}{\mathrm{d}t}=-0.03Q$ 对比即得 $k=-0.03$。
\end{solution}

\begin{exercise}[求参数]
设 $y=\cos\omega t$，求 $\omega$ 的值，使 $y$ 满足方程 $\displaystyle \frac{\mathrm{d}^2y}{\mathrm{d}t^2}+9y=0$。
\end{exercise}
\begin{solution}
由 $y=\cos\omega t$ 可得 \(\displaystyle \frac{\mathrm{d}y}{\mathrm{d}t}=-\omega\sin\omega t,\qquad
\frac{\mathrm{d}^2y}{\mathrm{d}t^2}=-\omega^2\cos\omega t\)。
代入方程 $\dfrac{\mathrm{d}^2y}{\mathrm{d}t^2}+9y=0$，得 $(9-\omega^2)\cos\omega t=0$。
要使它对任意 $t$ 恒成立，只能有$9-\omega^2=0\implies\omega=\pm 3$.
\end{solution}

\begin{exercise}[求参数]
在下列各题中, 利用所给的初始条件确定函数关系中的常数:
\[
\begin{array}{ll}
(1)&\displaystyle x^2-y^2=C,\quad y|_{x=0}=5;\\
(2)&\displaystyle y=\mathrm{e}^{2x}(C_1+C_2x),\quad y|_{x=0}=0,\ y'|_{x=0}=1;\\
(3)&\displaystyle y=C_1\sin(x-C_2),\quad y|_{x=\pi}=1,\ y'|_{x=\pi}=0.
\end{array}
\]
\end{exercise}
\begin{solution}
\begin{enumerate}
    \item[(1)] 将 $x=0$ 和 $y=5$ 代入方程：
    $0^2 - 5^2 = C \implies C = -25$。
    \item[(2)] 首先利用 $y(0)=0$：
    $0 = \mathrm{e}^0(C_1 + C_2 \cdot 0) \implies C_1 = 0$。
    将 $C_1=0$ 代入原函数，得 $y = C_2 x \mathrm{e}^{2x}$。对其求导：
    $y' = C_2(\mathrm{e}^{2x} + 2x\mathrm{e}^{2x})$。
    利用 $y'(0)=1$：
    $1 = C_2(\mathrm{e}^0 + 0) \implies C_2 = 1$。
    所以常数为 $C_1 = 0$, $C_2 = 1$。（特解为 $y=x\mathrm{e}^{2x}$）
    \item[(3)] 对原函数求导：
    $y' = C_1\cos(x-C_2)$。
    代入初始条件：\(\displaystyle y(\pi)=C_1\sin(\pi-C_2)=C_1\sin C_2=1\)（式1），
    \(\displaystyle y'(\pi)=C_1\cos(\pi-C_2)=-C_1\cos C_2=0\)（式2）。
    由（式1）知 $C_1 \neq 0$，因此由（式2）必有 $\cos C_2 = 0$。解得 $C_2 = \frac{\pi}{2} + k\pi$ ($k \in \mathbb{Z}$)。
    分两种情形：当 $k$ 为偶数时，即 $C_2 = \frac{\pi}{2} + 2m\pi$ ($m \in \mathbb{Z}$) 时，$\sin C_2 = 1$，代入（式1）得 $C_1 = 1$；当 $k$ 为奇数时，即 $C_2 = -\frac{\pi}{2} + 2m\pi$ ($m \in \mathbb{Z}$) 时，$\sin C_2 = -1$，代入（式1）得 $C_1 = -1$。
    所以常数可取为 $C_1=1, C_2=\frac{\pi}{2}+2k\pi$，或 $C_1=-1, C_2=-\frac{\pi}{2}+2k\pi$ ($k \in \mathbb{Z}$)。
\end{enumerate}
\end{solution}

\begin{exercise}[函数对应微分方程]
指出下列给出的函数满足哪个微分方程:
\par\noindent
\begin{tabular}{@{}lll@{}}
\toprule
 \textbf{函数} & \textbf{微分方程} \\
\midrule
(1)  $y=2\sin x$ & (a) $\displaystyle \frac{\mathrm{d}y}{\mathrm{d}x}=-2y$ \\
(2) $y=\sin2x$ & (b) $\displaystyle \frac{\mathrm{d}y}{\mathrm{d}x}=2y$ \\
(3)  $y=\mathrm{e}^{2x}$ & (c) $\displaystyle \frac{\mathrm{d}^2y}{\mathrm{d}x^2}=-y$ \\
(4)  $y=\mathrm{e}^{-2x}$ & (d) $\displaystyle \frac{\mathrm{d}^2y}{\mathrm{d}x^2}=-4y$ \\
\bottomrule
\end{tabular}
\par
\end{exercise}
\begin{solution}
\begin{enumerate}
    \item[(1)] 对于 $y=2\sin x$，有 $y' = 2\cos x,\ y'' = -2\sin x = -y$，因此函数 (1) 满足微分方程 (c)。
    \item[(2)] 对于 $y=\sin 2x$，有 $y' = 2\cos 2x,\ y'' = -4\sin 2x = -4y$，因此函数 (2) 满足微分方程 (d)。
    \item[(3)] 对于 $y=\mathrm{e}^{2x}$，有 $y' = 2\mathrm{e}^{2x} = 2y$，因此函数 (3) 满足微分方程 (b)。
    \item[(4)] 对于 $y=\mathrm{e}^{-2x}$，有 $y' = -2\mathrm{e}^{-2x} = -2y$，因此函数 (4) 满足微分方程 (a)。
\end{enumerate}
\end{solution}

\begin{exercise}[函数对应微分方程]
指出下列哪个函数是哪个微分方程的解?
\par\noindent
\begin{tabular}{@{}lll@{}}
\toprule
\textbf{函数} & \textbf{微分方程} \\
\midrule
(1)  $y=\mathrm{e}^x$ & (a) $y''-y=0$ \\
(2)  $y=x^3$ & (b) $x^2y''+2xy'-2y=0$ \\
(3)  $y=\mathrm{e}^{-x}$ & (c) $x^2y''-6y=0$ \\
(4)  $y=x^{-2}$ & \\
\bottomrule
\end{tabular}
\par
\end{exercise}
\begin{solution}
逐一验证左侧给出的函数：
\begin{enumerate}
\item 对 $y=\mathrm{e}^x$，有 $y'=\mathrm{e}^x,\ y''=\mathrm{e}^x$，代入 (a) 得 $y''-y=\mathrm{e}^x-\mathrm{e}^x=0$，故对应 (a)。
\item 对 $y=x^3$，有 $y'=3x^2,\ y''=6x$，代入 (c) 得 $x^2y''-6y=x^2\cdot 6x-6x^3=0$，故对应 (c)。
\item 对 $y=\mathrm{e}^{-x}$，有 $y'=-\mathrm{e}^{-x},\ y''=\mathrm{e}^{-x}$，代入 (a) 仍得 $y''-y=\mathrm{e}^{-x}-\mathrm{e}^{-x}=0$，故它同样对应 (a)。
这也说明 (1) 与 (3) 可以共同组成方程 (a) 的一组基础解。
\item 对 $y=x^{-2}$，有 $y'=-2x^{-3},\ y''=6x^{-4}$。代入 (b)，得
\(\displaystyle x^2y''+2xy'-2y=x^2(6x^{-4})+2x(-2x^{-3})-2x^{-2}=6x^{-2}-4x^{-2}-2x^{-2}=0\)，故对应 (b)。
\end{enumerate}
\end{solution}

\begin{exercise}[曲线族对应微分方程]
曲线族常常作为微分方程的解出现. 试将下列曲线与微分方程对应起来:
\par\noindent
\begin{tabular}{@{}lll@{}}
\toprule
 \textbf{曲线族} & \textbf{微分方程} \\
\midrule
(1) $y=x\mathrm{e}^{kx}$ & (a) $\displaystyle \frac{\mathrm{d}y}{\mathrm{d}x}=\frac{y}{x}$ \\
(2) $y=x^p$ & (b) $\displaystyle \frac{\mathrm{d}y}{\mathrm{d}x}=\frac{y\ln y}{x}$ \\
(3) $y=\mathrm{e}^{kx}$ & (c) $\displaystyle \frac{\mathrm{d}y}{\mathrm{d}x}=\frac{y}{x}\left(1+\ln\frac{y}{x}\right)$ \\
(4) $y=mx$ & (d) $\displaystyle \frac{\mathrm{d}y}{\mathrm{d}x}=\frac{y\ln y}{x\ln x}$ \\
\bottomrule
\end{tabular}
\par
\end{exercise}
\begin{solution}
我们需要通过消去曲线族中的参数（$k, p, m$）来得到微分方程：
\begin{enumerate}
    \item[(1)] 对于 $y=x\mathrm{e}^{kx}$，两边对 $x$ 求导得
    \(\displaystyle \frac{\mathrm{d}y}{\mathrm{d}x}=\mathrm{e}^{kx}+kx\mathrm{e}^{kx}\)。
    由原方程可得 $\mathrm{e}^{kx}=\dfrac{y}{x}$，两边取对数得 $kx=\ln\dfrac{y}{x}$。将其代入导数表达式中，得
    \(\displaystyle \frac{\mathrm{d}y}{\mathrm{d}x}=\frac{y}{x}+\left(\ln\frac{y}{x}\right)\frac{y}{x}
    =\frac{y}{x}\left(1+\ln\frac{y}{x}\right)\)。
    因此，(1) 对应 (c)。
    \item[(2)] 对于 $y=x^p$，两边对 $x$ 求导得
    \(\displaystyle \frac{\mathrm{d}y}{\mathrm{d}x}=px^{p-1}=p\frac{x^p}{x}=p\frac{y}{x}\)。
    由原方程 $y=x^p$ 两边取自然对数得 $\ln y=p\ln x$，从而 $p=\dfrac{\ln y}{\ln x}$。将 $p$ 代入导数表达式中，得
    \(\displaystyle \frac{\mathrm{d}y}{\mathrm{d}x}=\frac{\ln y}{\ln x}\cdot\frac{y}{x}
    =\frac{y\ln y}{x\ln x}\)。
    因此，(2) 对应 (d)。
    \item[(3)] 对于 $y=\mathrm{e}^{kx}$，两边对 $x$ 求导得
    \(\displaystyle \frac{\mathrm{d}y}{\mathrm{d}x}=k\mathrm{e}^{kx}=ky\)。
    由原方程两边取自然对数得 $\ln y=kx$，从而 $k=\dfrac{\ln y}{x}$。将 $k$ 代入导数表达式中，得
    \(\displaystyle \frac{\mathrm{d}y}{\mathrm{d}x}=\frac{\ln y}{x}\cdot y=\frac{y\ln y}{x}\)。
    因此，(3) 对应 (b)。
    \item[(4)] 对于 $y=mx$，两边对 $x$ 求导得
    \(\displaystyle \frac{\mathrm{d}y}{\mathrm{d}x}=m=\frac{y}{x}\)。
    因此，(4) 对应 (a)。
\end{enumerate}
综上所述，对应关系为：(1) --- (c)，(2) --- (d)，(3) --- (b)，(4) --- (a)。
\end{solution}

\begin{exercise}[曲线族对应微分方程]
求抛物线族 $\displaystyle y=C_1(x-C_2)^2$ 的微分方程.
\end{exercise}
\begin{solution}
曲线族 $y=C_1(x-C_2)^2$
含有两个独立常数 $C_1,C_2$，因此所求微分方程应为二阶方程。对原式连续求导，得
\(\displaystyle y'=2C_1(x-C_2),\qquad y''=2C_1\)。
由第二式可得 $C_1=\dfrac{y''}{2}$，又由原方程知 $C_1(x-C_2)^2=y$，因此
\(\displaystyle (y')^2=4C_1^2(x-C_2)^2=4C_1\bigl[C_1(x-C_2)^2\bigr]=4C_1y=2yy''\)。
于是所求微分方程为 \(\displaystyle \boxed{2yy''-(y')^2=0}\)。
\end{solution}

\subsection{第二节习题：变量可分离方程及齐次方程}

\subsubsection{A组}

\begin{exercise}[相关概念]
回答下列问题:
\begin{shorttrienum}
    \shortitem{1}{变量可分离方程具有什么形状?} &
    \shortitem{2}{什么叫齐次方程?} &
    \shortitem{3}{用什么办法将齐次方程化为变量可分离方程?}
\end{shorttrienum}
\end{exercise}
\begin{solution}
\begin{enumerate}
    \item[(1)] 变量可分离方程是指可以表示为 \(\displaystyle \frac{\mathrm{d}y}{\mathrm{d}x} = f(x)g(y)\) 或 \(\displaystyle M(x)\mathrm{d}x + N(y)\mathrm{d}y = 0\) 形式的一阶微分方程，其中方程的右端可以分解为一个仅含 $x$ 的函数与一个仅含 $y$ 的函数的乘积。

    \item[(2)] 若一阶方程 $\displaystyle\frac{\mathrm{d}y}{\mathrm{d}x} = f(x,y)$ 中的函数 $f(x,y)$ 满足对任意非零实数 $t$ 都有 \(\displaystyle f(tx,ty) = f(x,y)\)，或者等价地，它可以写成仅依赖于比值 $\frac{y}{x}$ 的函数形式 $\varphi\big(\frac{y}{x}\big)$，则称该微分方程为齐次方程。

    \item[(3)] 通常采用变量代换的方法，令 $u = \frac{y}{x}$（即 $y = ux$），则 \(\displaystyle \frac{\mathrm{d}y}{\mathrm{d}x} = u + x\frac{\mathrm{d}u}{\mathrm{d}x}\)。
    将此式代入原齐次方程 $\displaystyle \frac{\mathrm{d}y}{\mathrm{d}x} = \varphi\big(\frac{y}{x}\big)$ 中，可得 \(\displaystyle u + x\frac{\mathrm{d}u}{\mathrm{d}x} = \varphi(u),\quad \frac{\mathrm{d}u}{\mathrm{d}x} = \frac{\varphi(u)-u}{x}\)，从而转化为关于未知函数 $u(x)$ 的变量可分离方程。
\end{enumerate}
\end{solution}

\begin{exercise}[求解方程]
对 $k>0$ 求微分方程 $\displaystyle \frac{\mathrm{d}H}{\mathrm{d}t}=-k(H-20)$ 的通解.
\end{exercise}
\begin{solution}
对方程进行变量分离。当 $H \neq 20$ 时，\(\displaystyle \frac{\mathrm{d}H}{H-20} = -k\,\mathrm{d}t,\quad \ln|H-20| = -kt + C_1\)。故 \(H-20 = \pm\mathrm{e}^{C_1}\mathrm{e}^{-kt}\)。令 $C = \pm\mathrm{e}^{C_1}$，则通解为 \(\displaystyle \mathbf{H(t) = 20 + C\mathrm{e}^{-kt}}\)。
（注意：当 $C=0$ 时，特解 $H \equiv 20$ 也包含在上述通解之中。）
\end{solution}

\begin{exercise}[求解方程]
求初值问题 $\displaystyle \frac{\mathrm{d}P}{\mathrm{d}t}=2P-2Pt,\,P\big|_{t=0}=5$ 的解, 并画出其图形.
\end{exercise}
\begin{solution}
提取公因式，将微分方程化为变量可分离形式 \(\displaystyle \frac{\mathrm{d}P}{\mathrm{d}t} = 2P(1-t)\)。分离变量并积分得 \(\displaystyle \frac{\mathrm{d}P}{P}=2(1-t)\,\mathrm{d}t,\ \ln|P|=2t-t^2+C_1,\ P=C\mathrm{e}^{2t-t^2}\)。
再由初始条件 $P(0)=5$ 得 $C=5$。
故所求初值问题的解为 \(\displaystyle \mathbf{P(t) = 5\mathrm{e}^{2t-t^2}}\)。
\textbf{图形描述：} 此函数图象为一族类似高斯钟形曲线。对其求导发现极值点在 $t=1$ 处取得，最大值为 $P(1) = 5\mathrm{e}$。曲线在 $t=1$ 两侧对称，随着 $t \to \pm\infty$，$P(t)$ 迅速衰减趋于 $0$。
\end{solution}

\begin{exercise}[求解方程]
求解下列微分方程：
\begin{shortexenum}
  \shortitem{1}{$\displaystyle \frac{\mathrm{d}P}{\mathrm{d}t}=0.02P,\ P(0)=20$；} &
  \shortitem{2}{$\displaystyle \frac{\mathrm{d}y}{\mathrm{d}x}+\frac{y}{3}=0,\ y(0)=10$；}\\
  \shortitem{3}{$\displaystyle \frac{\mathrm{d}y}{\mathrm{d}x}=2y-4$，过点 $(2,5)$；} &
  \shortitem{4}{$\displaystyle \frac{\mathrm{d}z}{\mathrm{d}t}=t\mathrm{e}^z$，过原点；}\\
  \shortitem{5}{$\displaystyle \frac{\mathrm{d}y}{\mathrm{d}x}=\frac{5y}{x},\ y\big|_{x=1}=3$；} &
  \shortitem{6}{$\displaystyle \frac{\mathrm{d}u}{\mathrm{d}t}=u+ut^2,\ u(0)=5$；}\\
  \shortitem{7}{$\displaystyle \frac{\mathrm{d}\omega}{\mathrm{d}\theta}=\theta\omega^2\sin\theta^2,\ \omega(0)=1$；} &
  \shortitem{8}{$\displaystyle x(x+1)\frac{\mathrm{d}y}{\mathrm{d}x}=y^2,\ y(1)=1$。}
\end{shortexenum}
\end{exercise}
\begin{solution}
\begin{enumerate}
    \item[(1)] 分离变量并积分，\(\displaystyle \frac{\mathrm{d}P}{P}=0.02\,\mathrm{d}t \implies \ln|P|=0.02t+C \implies P(t)=C_1\mathrm{e}^{0.02t}\)。由 $P(0)=20$ 得 $C_1=20$，故 \(\displaystyle \mathbf{P(t)=20\mathrm{e}^{0.02t}}\)。

    \item[(2)] 分离变量并积分，\(\displaystyle \frac{\mathrm{d}y}{y}=-\frac{1}{3}\,\mathrm{d}x \implies y=C\mathrm{e}^{-\frac{x}{3}}\)。由 $y(0)=10$ 得 $C=10$，故 \(\displaystyle \mathbf{y=10\mathrm{e}^{-\frac{x}{3}}}\)。

    \item[(3)] 分离变量并积分，\(\displaystyle \frac{\mathrm{d}y}{y-2}=2\,\mathrm{d}x \implies \ln|y-2|=2x+C \implies y=2+C_1\mathrm{e}^{2x}\)。再由点 $(2,5)$ 得 $C_1=3\mathrm{e}^{-4}$，故 \(\displaystyle \mathbf{y=2+3\mathrm{e}^{2x-4}}\)。

    \item[(4)] 分离变量并积分，\(\displaystyle \mathrm{e}^{-z}\mathrm{d}z=t\,\mathrm{d}t \implies -\mathrm{e}^{-z}=\frac{1}{2}t^2+C\)。由 $z(0)=0$ 得 $C=-1$，于是 \(\displaystyle \mathrm{e}^{-z}=1-\frac{1}{2}t^2 \implies \mathbf{z(t)=-\ln\left(1-\frac{1}{2}t^2\right)}\)。

    \item[(5)] 分离变量并积分，\(\displaystyle \frac{\mathrm{d}y}{y}=\frac{5}{x}\mathrm{d}x \implies \ln|y|=5\ln|x|+C \implies y=C_1x^5\)。由 $y|_{x=1}=3$ 得 $C_1=3$，故 \(\displaystyle \mathbf{y=3x^5}\)。

    \item[(6)] 提取公因式 \(\displaystyle \frac{\mathrm{d}u}{\mathrm{d}t} = u(1+t^2)\)，分离并积分，\(\displaystyle \frac{\mathrm{d}u}{u}=(1+t^2)\mathrm{d}t \implies \ln|u|=t+\frac{1}{3}t^3+C \implies u=C_1\mathrm{e}^{t+\frac{t^3}{3}}\)。由 $u(0)=5$ 得 $C_1=5$，故 \(\displaystyle \mathbf{u(t)=5\mathrm{e}^{t+\frac{t^3}{3}}}\)。

    \item[(7)] 分离变量得 \(\displaystyle \frac{\mathrm{d}\omega}{\omega^2}=\theta\sin(\theta^2)\mathrm{d}\theta \implies -\frac{1}{\omega}=-\frac{1}{2}\cos(\theta^2)+C\)。
    由 $\omega(0)=1$ 得 $C=-\dfrac{1}{2}$，于是 \(\displaystyle -\frac{1}{\omega}=-\frac{1}{2}\bigl(\cos(\theta^2)+1\bigr)=-\cos^2\left(\frac{\theta^2}{2}\right)\implies \mathbf{\omega=\sec^2\left(\frac{\theta^2}{2}\right)}\)。

    \item[(8)] 分离变量得
    \[
    \frac{\mathrm{d}y}{y^2}
    =\frac{\mathrm{d}x}{x(x+1)}
    =\left(\frac{1}{x}-\frac{1}{x+1}\right)\mathrm{d}x
    \implies -\frac{1}{y}=\ln\left|\frac{x}{x+1}\right|+C.
    \]
    由 $y(1)=1$ 得 $C=\ln 2-1$，于是 \(\displaystyle -\frac{1}{y}=\ln\left|\frac{2x}{x+1}\right|-1\implies \mathbf{y=\frac{1}{1-\ln\left|\frac{2x}{x+1}\right|}}\)。
\end{enumerate}
\end{solution}

\begin{exercise}[求解方程]
求微分方程 $\displaystyle \frac{\mathrm{d}y}{\mathrm{d}t}=1-y$ 的通解, 并求适合初始条件 $y\big|_{t=0}=0$ 的特解.
\end{exercise}
\begin{solution}
分离变量可得 $\frac{\mathrm{d}y}{1-y}=\mathrm{d}t$，从而 $-\ln|1-y|=t+C$，即 $y=1-C_1\mathrm{e}^{-t}$。
由 $y(0)=0$ 得 $C_1=1$，故 $\mathbf{y=1-\mathrm{e}^{-t}}$。
其中 $C_1$ 为任意常数，前面的关系式给出的是\textbf{通解}，最后一行给出的是满足初始条件的\textbf{特解}。
\end{solution}

\begin{exercise}[求解齐次方程]
求下列齐次方程的通解:\\
    (1) $\displaystyle x\frac{\mathrm{d}y}{\mathrm{d}x}=y\ln\frac{y}{x}$; ~~
    (2) $\displaystyle y^2+x^2\frac{\mathrm{d}y}{\mathrm{d}x}=xy\frac{\mathrm{d}y}{\mathrm{d}x}$; ~~
    (3) $\displaystyle \frac{\mathrm{d}y}{\mathrm{d}x}=2\sqrt{\frac{y}{x}}+\frac{y}{x}$; ~~
    (4) $\displaystyle (x^2+y^2)\mathrm{d}x-xy\mathrm{d}y=0$.
\end{exercise}
\begin{solution}
针对齐次方程，统一令 $y=ux$，则 $y' = u + xu'$：
\begin{enumerate}
    \item[(1)] 原方程可化为 $\displaystyle\frac{\mathrm{d}y}{\mathrm{d}x} = \frac{y}{x}\ln\frac{y}{x}$。代入得
    \[
    u + x\frac{\mathrm{d}u}{\mathrm{d}x} = u\ln u
    \implies
    x\frac{\mathrm{d}u}{\mathrm{d}x} = u(\ln u - 1),
    \qquad
    \frac{\mathrm{d}u}{u(\ln u - 1)} = \frac{\mathrm{d}x}{x}.
    \]
    积分得 $\ln|\ln u - 1| = \ln|x| + C \implies \ln u - 1 = Cx$，即 $\ln\frac{y}{x} = Cx+1$，故通解为 $\mathbf{y = x\mathrm{e}^{Cx+1}}$。

    \item[(2)] 原方程整理为 $\displaystyle\frac{\mathrm{d}y}{\mathrm{d}x} = \frac{y^2}{xy-x^2}$。代入得
    \[
    u + xu' = \frac{u^2}{u-1}
    \implies
    xu' = \frac{u^2}{u-1} - u = \frac{u}{u-1},
    \qquad
    \frac{u-1}{u}\mathrm{d}u = \frac{\mathrm{d}x}{x}.
    \]
    积分得 \(\displaystyle u - \ln|u| = \ln|x| + C \implies \frac{y}{x} = \ln|x| + \ln\left|\frac{y}{x}\right| + C = \ln|y| + C\)。通解为（以及特解 $y=0$）$\mathbf{y = C_1\mathrm{e}^{y/x}}$。

    \item[(3)] 代入得 \(\displaystyle u + xu' = 2\sqrt{u} + u \implies xu' = 2\sqrt{u},\quad \frac{\mathrm{d}u}{2\sqrt{u}} = \frac{\mathrm{d}x}{x}\)。积分得 $\sqrt{u} = \ln|x| + C$。将 $u=y/x$ 回代得通解 $\mathbf{\sqrt{\frac{y}{x}} = \ln|x| + C}$，另有特解 $y=0$。

    \item[(4)] 原方程化为 $\displaystyle\frac{\mathrm{d}y}{\mathrm{d}x} = \frac{x^2+y^2}{xy} = \frac{x}{y} + \frac{y}{x}$。代入得 \(\displaystyle u + xu' = \frac{1}{u} + u \implies xu' = \frac{1}{u},\quad u\,\mathrm{d}u = \frac{\mathrm{d}x}{x}\)。积分得 $\frac{1}{2}u^2 = \ln|x| + C_1$。即 $\frac{y^2}{2x^2} = \ln|x| + C_1$，通解为 $\mathbf{y^2 = 2x^2(\ln|x| + C_1)=x^2(\ln x^2+C)}$。
\end{enumerate}
\end{solution}

\begin{exercise}[求解齐次方程]
求下列初值问题的解:\\
(1) $\displaystyle \frac{\mathrm{d}y}{\mathrm{d}x}=\frac{x}{y}+\frac{y}{x},\,y\big|_{x=1}=2$;
~~(2)$\displaystyle (x^2+2xy-y^2)\mathrm{d}x+(y^2+2xy-x^2)\mathrm{d}y=0,\,y\big|_{x=1}=1$.
\end{exercise}
\begin{solution}
\begin{enumerate}
    \item[(1)] 由上一题结论可知该方程通解为 $y^2=2x^2(\ln|x|+C)$。由 $y(1)=2$ 得 $C=2$，故 $y^2=2x^2(\ln x+2)$。
    考虑到当 $x=1$ 时 $y=2>0$，开方取正支，得到特解 $\mathbf{y = x\sqrt{2\ln x + 4}}$。

    \item[(2)] 将方程显式化为 $\displaystyle\frac{\mathrm{d}y}{\mathrm{d}x} = \frac{y^2-2xy-x^2}{y^2+2xy-x^2}$。令 $y=ux$，得
    \(\displaystyle u+xu'=\frac{u^2-2u-1}{u^2+2u-1}
    \implies xu'=-\frac{(u+1)(u^2+1)}{u^2+2u-1}
    \implies \left(\frac{2u}{u^2+1}-\frac{1}{u+1}\right)\mathrm{d}u=-\frac{\mathrm{d}x}{x}\)。
    积分得 $\ln(u^2+1)-\ln|u+1|=-\ln|x|+C \implies \frac{u^2+1}{|u+1|}|x|=C_1$。回代 $u=\frac{y}{x}$，整理得 $\frac{x^2+y^2}{x+y} = C_1$。代入初始条件 $y(1)=1$ 得 $\frac{1+1}{1+1} = C_1 \implies C_1 = 1$，故所求初值问题的解为 $\mathbf{x^2+y^2=x+y}$。
\end{enumerate}
\end{solution}

\subsubsection{B组}

\begin{exercise}[求解方程]
化下列方程为齐次方程, 并求出其通解:\\
(1) $\displaystyle \frac{\mathrm{d}y}{\mathrm{d}x}=\frac{y-x-2}{x+y+4}$;~~(2)$\displaystyle (x-y-1)\mathrm{d}x+(4y+x-1)\mathrm{d}y=0$.
\end{exercise}
\begin{solution}
\begin{enumerate}
    \item[(1)] 联立方程组解两条直线的交点，并作平移变换 $X=x+3, Y=y+1$，有
    联立 $y-x-2=0,\ x+y+4=0$ 得交点 $(-3,-1)$；再令 $Y=uX$，则
    \(\displaystyle u+Xu'=\frac{u-1}{1+u}\implies Xu'=\frac{-1-u^2}{1+u}\)。
    因而 \(\displaystyle
    \frac{1+u}{1+u^2}\mathrm{d}u=-\frac{\mathrm{d}X}{X}
    \implies \arctan u+\frac{1}{2}\ln(1+u^2)=-\ln|X|+C\)。
    合并对数项并回代 $u=Y/X$ 及原变量，得 \(\displaystyle \arctan u + \ln|X\sqrt{1+u^2}| = C\)，即
    \(\displaystyle \mathbf{\arctan\frac{y+1}{x+3} + \frac{1}{2}\ln\big((x+3)^2+(y+1)^2\big) = C}\)。

    \item[(2)] 化为标准形 $y' = -\frac{x-y-1}{x+4y-1}$。联立 $x-y-1=0,\ x+4y-1=0$ 得交点 $(1,0)$；令 $X=x-1,\ Y=y$，再取 $Y=uX$，有
    \(\displaystyle u+Xu'=\frac{u-1}{1+4u}\implies Xu'=\frac{-1-4u^2}{1+4u}\)。
    因而 \(\displaystyle \frac{1+4u}{1+4u^2}\mathrm{d}u=-\frac{\mathrm{d}X}{X}
    \implies \frac{1}{2}\arctan(2u)+\frac{1}{2}\ln(1+4u^2)=-\ln|X|+C_1\)。
    化简得 \(\displaystyle \arctan(2u) + \ln\big(X^2(1+4u^2)\big) = C\)。
    回代后得到通解 \(\displaystyle \mathbf{\arctan\frac{2y}{x-1} + \ln\big((x-1)^2+4y^2\big) = C}\)。
\end{enumerate}
\end{solution}

\begin{exercise}[求解方程]
求下列方程的通解:(1)$\displaystyle \sqrt{1-y^2}=3x^2yy'$;~~(2)$\displaystyle x\mathrm{d}y-y\mathrm{d}x=\sqrt{x^2+y^2}\mathrm{d}x$.
\end{exercise}
\begin{solution}
\begin{enumerate}
    \item[(1)] 原方程可直接分离变量：\(\displaystyle \frac{y}{\sqrt{1-y^2}}\mathrm{d}y=\frac{1}{3x^2}\mathrm{d}x\)，积分得 \(\displaystyle -\sqrt{1-y^2}=-\frac{1}{3x}+C\)，即 \(\displaystyle \mathbf{\sqrt{1-y^2}}=\mathbf{\frac{1}{3x}-C}\)。
    另外显然 $y=\pm 1$ 也是解。

    \item[(2)] 对方程同除以 $x^2$（假设 $x>0$），并令 $u = \frac{y}{x}$，得 \(\displaystyle \mathrm{d}\left(\frac{y}{x}\right)=\frac{1}{x}\sqrt{1+\left(\frac{y}{x}\right)^2}\mathrm{d}x\)，即
    \(\displaystyle \frac{\mathrm{d}u}{\sqrt{1+u^2}}=\frac{\mathrm{d}x}{x}\)。积分后 \(\displaystyle \ln(u+\sqrt{1+u^2})=\ln x+C_1\)，从而 \(\displaystyle y+\sqrt{x^2+y^2}=Cx^2\)。
    移项平方消去根号，可化简得通解：\(\displaystyle \mathbf{y = \frac{1}{2}\left(Cx^2 - \frac{1}{C}\right)},\ x=0\)。
    对于 $x<0$ 推导同理，结论一致。
\end{enumerate}
\end{solution}

\begin{exercise}[求解方程]
由曲线上任意一点引法线, 它在纵轴上截得线段的长度等于该点到坐标原点的距离, 求此曲线的方程.
\end{exercise}
\begin{solution}
\begin{enumerate}
    \item \textbf{建立微分方程}：
    设曲线方程为 $y=y(x)$。曲线上任意一点 $(x,y)$ 处的法线方程为 $Y - y = -\frac{1}{y'}(X - x)$。令 $X=0$，得到法线在纵轴上的截距 \(\displaystyle Y_0 = y + \frac{x}{y'} = y + x\frac{\mathrm{d}x}{\mathrm{d}y}\)。由题意，该截距的绝对值等于点 $(x,y)$ 到坐标原点的距离 $\sqrt{x^2+y^2}$，即 \(\displaystyle y + x\frac{\mathrm{d}x}{\mathrm{d}y} = \pm\sqrt{x^2+y^2}\)。

    \item \textbf{凑全微分求解}：
    将方程两边同乘 $\mathrm{d}y$，得到 $y\,\mathrm{d}y + x\,\mathrm{d}x = \pm\sqrt{x^2+y^2}\,\mathrm{d}y$。注意到左端恰好是 $\frac{1}{2}\mathrm{d}(x^2+y^2)$，方程可写为
    \(\displaystyle \frac{\mathrm{d}(x^2+y^2)}{2\sqrt{x^2+y^2}} = \pm\mathrm{d}y \implies \mathrm{d}\left(\sqrt{x^2+y^2}\right) = \pm\mathrm{d}y\)
    两边直接积分得 $\sqrt{x^2+y^2} = \pm y + C$。

    \item \textbf{代数化简与常数吸收}：
    将等式两边平方，消除根号，得 $x^2+y^2 = y^2 \pm 2Cy + C^2 \implies x^2 = \pm 2Cy + C^2$。
    由于 $C$ 是任意常数（可正可负），表达式 $\pm 2C$ 中的符号可以完全被常数 $C$ 本身吸收（即若取负号，可视作将常数重新记为 $-C$，且其平方项保持不变）。因此，两种情况实际上代表同一族曲线：\(\displaystyle x^2=2Cy+C^2=C(2y+C)\)，即 \(\displaystyle \boxed{x^2=C(2y+C)}\)。
\end{enumerate}
\end{solution}

\begin{exercise}[实际应用题]
质量为 $1\ \mathrm{kg}$ 的质点受外力的作用作直线运动, 该力和时间成正比, 和质点运动的速度成反比. 在 $t=10\ \mathrm{s}$ 时速度为 $5\ \mathrm{m/s}$, 力为 $4\ \mathrm{N}$. 问从运动开始经过 $20\ \mathrm{s}$ 后的速度是多少?
\end{exercise}
\begin{solution}
设外力为 $F$，运动速度为 $v$。根据题意有 $F = k\frac{t}{v}$。
由 $t=10, v=5, F=4$，代入并结合牛顿第二定律 $F = ma$（其中 $m=1\ \mathrm{kg}$），得
\(\displaystyle 4=k\frac{10}{5}\implies k=2,\quad
\frac{\mathrm{d}v}{\mathrm{d}t}=\frac{2t}{v}
\implies \frac{1}{2}v^2=t^2+C
\implies v^2=2t^2+C_1\)。
再由 $t=10,\ v=5$ 得 $25=200+C_1$，所以 $C_1=-175$。于是 $v^2=2(20^2)-175=625$，即 $\mathbf{v=25\ \mathrm{m/s}}$。
（由于速度方向与外力一致，取正）。
\end{solution}

\begin{exercise}[实际应用题]
我们所使用的可耕地面积随着世界人口的增加而增加. 令 $A(t)$ 表示在 $t$ 年所使用的可耕地的总公顷数 ($1\ \mathrm{hm}=1\times10^4\ \mathrm{m}^2$).
\begin{shortsingleenum}
    \shortsingleitem{1}{试解释为什么我们期望 $A(t)$ 满足微分方程 $\displaystyle \frac{\mathrm{d}A}{\mathrm{d}t}=kA$ 是合理的? 关于世界人口你作了些什么假设? 它与可耕地面积有什么关系?}
    \shortsingleitem{2}{在 1950 年, 约有 $1\times10^9\ \mathrm{hm}$ 的可耕地在使用, 到 1980 年则有 $2\times10^9\ \mathrm{hm}$. 如果可供使用的可耕地的总公顷数为 $3.2\times10^9$, 试问, 这些耕地何时被开垦并耕种完 ($t=0$ 指 1950 年)?}
\end{shortsingleenum}
\end{exercise}
\begin{solution}
\begin{enumerate}
    \item[(1)] \textbf{合理性解释：} 假设世界人口 $P(t)$ 在此时期内不受资源极度匮乏的影响，按照马尔萨斯模型呈指数增长，即 $\frac{\mathrm{d}P}{\mathrm{d}t} = \lambda P$。同时假设人均需要使用的可耕地面积维持一个常数 $\alpha$，即 $A(t) = \alpha P(t)$。因此，$A(t)$ 的变化率也必与当前可耕地面积成正比：\(\displaystyle \frac{\mathrm{d}A}{\mathrm{d}t}=\alpha \frac{\mathrm{d}P}{\mathrm{d}t}=\alpha(\lambda P)=\lambda A\)。若令 $k=\lambda$，即得所期望的微分方程。

    \item[(2)] 该微分方程的解为 $A(t) = A_0\mathrm{e}^{kt}$。已知 $t=0$ 时 $A(0) = 1\times 10^9$，故 $A(t)=10^9\mathrm{e}^{kt}$。再由 $A(30)=2\times 10^9$ 得 \(\displaystyle 2\times 10^9=10^9\mathrm{e}^{30k} \implies \mathrm{e}^{30k}=2 \implies k=\frac{\ln 2}{30}\)。
    当 $A(t)=3.2\times 10^9$ 时，
    \(\displaystyle \mathrm{e}^{\frac{\ln 2}{30}t}=3.2\)，从而
    \(\displaystyle t=30\frac{\ln 3.2}{\ln 2}\approx 50.34\ (\text{年})\)。
    即大约需要 $50.34$ 年，对应于 \textbf{2000年}左右这些耕地将被开垦并耕种完。
\end{enumerate}
\end{solution}

\begin{exercise}[实际应用题]
将一只白薯放进 $200^\circ\mathrm{C}$ 的烤炉内, 其加热情况服从牛顿加热(或冷却)定律.
\begin{shortexenum}
    \shortitem{1}{写出白薯的温度函数 $H(t)$ 所满足的微分方程;} &
    \shortitem{2}{若白薯放进炉内时的温度是 $20^\circ\mathrm{C}$, 求解所得的微分方程;} \\
    \shortitem{3}{假设 30 分钟后白薯的温度为 $120^\circ\mathrm{C}$, 试确定牛顿加热定律中的比例常数 $k$.} &
\end{shortexenum}
\end{exercise}
\begin{solution}
\begin{enumerate}
    \item[(1)] 设比例常数为 $k$ (其中 $k<0$)。由牛顿冷却（加热）定律，白薯温度变化率与它同周围环境（烤炉）的温差成正比。微分方程为
    \(\displaystyle \frac{\mathrm{d}H}{\mathrm{d}t} = k(H - 200)\)。

    \item[(2)] 这是一个可分离变量方程，\(\displaystyle \frac{\mathrm{d}H}{H-200} = k\,\mathrm{d}t\)。两边积分并代入初始条件 $H(0) = 20$，得
    \(\displaystyle \ln|H-200|=kt+C_1\)，即
    \(\displaystyle H-200=C\mathrm{e}^{kt}\)。
    代入 $H(0)=20$ 得 $C=-180$，所以
    \(\displaystyle H(t)=200-180\mathrm{e}^{kt}\)。

    \item[(3)] 代入 $t=30$ 时 $H=120$，得
    \(\displaystyle 120 = 200 - 180\mathrm{e}^{30k}\)，即
    \(\displaystyle \mathrm{e}^{30k} = \frac{4}{9}\)，从而
    \(\displaystyle k = \frac{1}{30}\ln\frac{4}{9} \approx -0.027\)。
\end{enumerate}
\end{solution}

\begin{exercise}[实际应用题]
放射性元素由于原子中不断放出微观粒子, 它的含量不断减少, 称为衰变. 由实验知道, 放射性元素镭的衰变率与这时镭的存量成正比(比例系数为 $k$). 设在开始时镭的存量为 $R_0$ 克, 求任意时刻 $t$ 的含量 $R(t)$. 经验材料断定, 镭经 1600 年后, 只余原质量之半, 试由此确定比例常数 $k$.
\end{exercise}
\begin{solution}
根据题意列出微分方程，并利用初始条件与半衰期条件：
\(\displaystyle \frac{\mathrm{d}R}{\mathrm{d}t}=-kR,\qquad R(t)=C\mathrm{e}^{-kt}\)。
再由 \(R(0)=R_0\) 得 \(\displaystyle R(t)=R_0\mathrm{e}^{-kt}\)。
由半衰期条件
\(\displaystyle \frac{1}{2}R_0=R_0\mathrm{e}^{-1600k}\)，即
\(\displaystyle \mathrm{e}^{-1600k}=\frac{1}{2}\)，从而
\(\displaystyle k=\frac{\ln 2}{1600}\approx 4.33\times 10^{-4}\ (\text{年}^{-1})\)。
\end{solution}

\begin{exercise}[实际应用题]
设在某一溶液中有两种物质, 在反应开始时, 两种物质的量分别为 $a$ 与 $b$; 又设在时刻 $t$, 两种物质已经起反应的量相等, 记为 $x$, 利用化学反应的基本定律: 反应进行的速率与尚未起反应的量的乘积成正比(比例系数为 $k$). 求变量 $x$ 随时间 $t$ 的变化规律.
\end{exercise}
\begin{solution}
由化学反应基本定律（质量作用定律），反应速率 $\frac{\mathrm{d}x}{\mathrm{d}t}$ 与两种物质尚未起反应的量 $(a-x)$ 和 $(b-x)$ 之积成正比，微分方程为：
$\frac{\mathrm{d}x}{\mathrm{d}t} = k(a-x)(b-x), \quad x(0)=0$。
分离变量得：$\frac{\mathrm{d}x}{(a-x)(b-x)} = k\,\mathrm{d}t$。积分过程需根据 $a$ 与 $b$ 的大小关系分类讨论：

\textbf{情形一：$a \neq b$}
利用部分分式分解，有 $\frac{1}{b-a}\left(\frac{1}{a-x} - \frac{1}{b-x}\right)\mathrm{d}x = k\,\mathrm{d}t$。两边积分得 $\frac{1}{b-a}\big(-\ln|a-x| + \ln|b-x|\big) = kt + C \implies \frac{1}{b-a}\ln\left|\frac{b-x}{a-x}\right| = kt + C$。
代入 $x(0)=0$，得 $C = \frac{1}{b-a}\ln\left|\frac{b}{a}\right|$，回代化简为 $\ln\frac{b-x}{a-x} = (b-a)kt + \ln\frac{b}{a}$。
进一步解出
\(\displaystyle x(t) = \frac{ab\big(1-\mathrm{e}^{(b-a)kt}\big)}{a - b\mathrm{e}^{(b-a)kt}}\)。

\textbf{情形二：$a = b$}
方程变为 $\frac{\mathrm{d}x}{(a-x)^2} = k\,\mathrm{d}t$，积分得 $\frac{1}{a-x} = kt + C$。
代入 $x(0)=0$，得 $C = \frac{1}{a}$，从而 $\frac{1}{a-x} = kt + \frac{1}{a} = \frac{akt+1}{a}$。
解得
\(\displaystyle x(t) = \frac{a^2kt}{akt+1}\)。
\end{solution}

\begin{exercise}[实际应用题]
设质量为 $m$ 的炸弹在具有水平初速 $v_0$ 而不具有垂直初速的情况下在空气中降落. 现在, 我们不研究炸弹在水平方向的位移, 而只研究炸弹在重力方向的降落. 由实验知, 空气的阻力与炸弹下降的速度平方成正比: $R=\mu v^2$, 求炸弹下降的速度与时间的关系.
\end{exercise}
\begin{solution}
取竖直向下为正方向。根据牛顿第二定律并令极限速度 $v_T = \sqrt{\frac{mg}{\mu}}$，有：
$m\frac{\mathrm{d}v}{\mathrm{d}t}=mg-\mu v^2,\ \frac{\mathrm{d}v}{v_T^2-v^2}=\frac{g}{v_T^2}\,\mathrm{d}t$。
积分得 $\frac{1}{2v_T}\ln\left|\frac{v_T+v}{v_T-v}\right|=\frac{g}{v_T^2}t+C$。
再由 $v(0)=0$ 得 $C=0$，于是 $\frac{v_T+v}{v_T-v}=\mathrm{e}^{\frac{2g}{v_T}t}$，从而
\[
v(t)=v_T\frac{\mathrm{e}^{\frac{2g}{v_T}t}-1}{\mathrm{e}^{\frac{2g}{v_T}t}+1}
=v_T\tanh\left(\frac{g}{v_T}t\right)
=\sqrt{\frac{mg}{\mu}}\tanh\left(\sqrt{\frac{\mu g}{m}}t\right).
\]
\end{solution}

\begin{exercise}[求解曲线]
一曲线通过点 $(2,0)$, 且其切点和 $y$ 轴间的切线段有定长 2, 求这曲线.
\end{exercise}
\begin{solution}
设所求曲线的方程为 $y=y(x)$。在任意一点 $(x,y)$ 处的切线方程为 $Y-y = y'(X-x)$。
令 $X=0$ 得该切线在 $y$ 轴上的截距 $Y_0 = y - xy'$。
从切点 $(x,y)$ 到 $y$ 轴交点 $(0, Y_0)$ 的切线段长度为 \(\displaystyle d = \sqrt{x^2 + (y-Y_0)^2}= \sqrt{x^2 + (xy')^2}= |x|\sqrt{1+(y')^2}\)。由题意 $d=2$，又由于曲线过点 $(2,0)$，局部有 $x>0$，故 \(\displaystyle x\sqrt{1+(y')^2} = 2 \implies (y')^2 = \frac{4}{x^2}-1 \implies y' = \pm\frac{\sqrt{4-x^2}}{x}\)，
从而
\(\displaystyle y = \pm\int \frac{\sqrt{4-x^2}}{x}\mathrm{d}x\)。
作三角代换令 $x = 2\sin\theta$，则 $\mathrm{d}x = 2\cos\theta\,\mathrm{d}\theta$，积分变为
\[
\int \frac{2\cos\theta}{2\sin\theta} 2\cos\theta\,\mathrm{d}\theta
= 2\int \frac{\cos^2\theta}{\sin\theta}\mathrm{d}\theta
= 2\int \frac{1-\sin^2\theta}{\sin\theta}\mathrm{d}\theta
= 2\ln|\csc\theta - \cot\theta| - 2\cos\theta + C.
\]
将 $\sin\theta = x/2$ 回代得 $\cos\theta = \frac{\sqrt{4-x^2}}{2}$，再利用曲线过点 $(2,0)$，得 $y=\pm \left( 2\ln\left|\frac{2-\sqrt{4-x^2}}{x}\right| + \sqrt{4-x^2} \right)+C$。
由点 $(2,0)$ 得 $C=0$，故
\(\displaystyle y = \pm \left( \sqrt{4-x^2} - 2\ln\left|\frac{2+\sqrt{4-x^2}}{x} \right| \right)\)。
\end{solution}

\begin{exercise}[求解曲线]
一曲线通过点 $(2,3)$, 它在两坐标轴间的任意切线段均被切点所平分, 求这曲线方程.
\end{exercise}
\begin{solution}
切点为 $(x,y)$ 的切线方程为 $Y-y = y'(X-x)$。
令 $Y=0$，得 $x$ 轴上的截距 $X_A = x - \frac{y}{y'}$；
令 $X=0$，得 $y$ 轴上的截距 $Y_B = y - x y'$。
依题意切点平分切线段，故点 $(x,y)$ 是 $(X_A, 0)$ 和 $(0, Y_B)$ 的中点，即 $x=\frac{X_A}{2}=\frac{x-y/y'}{2}$，从而 $\frac{\mathrm{d}y}{\mathrm{d}x}=-\frac{y}{x}$，再积分得 $\frac{\mathrm{d}y}{y}=-\frac{\mathrm{d}x}{x} \implies \ln|y|=-\ln|x|+C_1 \implies xy=C$。
再由曲线过点 $(2,3)$，得 $C=6$，所以
\(\displaystyle xy=6\)。
\end{solution}

\begin{exercise}[求解曲线]
求曲线 $y=y(x)$, 使它正交于圆心在 $x$ 轴上且过原点的任何圆(注: 两曲线正交是指在交点处两曲线的切线互相垂直).
\end{exercise}
\begin{solution}
所给已知圆族圆心在 $(a,0)$ 并过原点，故 \(\displaystyle (x-a)^2+y^2=a^2
\implies x^2+y^2=2ax
\implies y'=\frac{y^2-x^2}{2xy}\)。
因此正交轨线满足 $\frac{\mathrm{d}y}{\mathrm{d}x}=\frac{2xy}{x^2-y^2}$。这是一个齐次微分方程。令 $y = ux$，则
\[
u+xu'=\frac{2u}{1-u^2}
\implies
\frac{1-u^2}{u(1+u^2)}\mathrm{d}u=\frac{\mathrm{d}x}{x}
\implies
\left(\frac{1}{u}-\frac{2u}{1+u^2}\right)\mathrm{d}u=\frac{\mathrm{d}x}{x}.
\]
积分得 $\ln|u|-\ln(1+u^2)=\ln|x|+\ln|C| \implies \frac{u}{1+u^2}=Cx$，回代即得 $\frac{y/x}{1+y^2/x^2}=Cx \implies x^2+y^2=Cy$。
结果即为正交轨线族方程 $\mathbf{x^2+y^2=Cy}$，它恰是圆心在 $y$ 轴上且过原点的所有圆。
\end{solution}

\begin{exercise}[落体问题]
    设跳伞运动员从跳伞塔下落后, 所受空气的阻力与速度成正比, 运动员离塔时的速度为零, 求运动员下落过程中速度和时间的函数关系.
\end{exercise}
\begin{solution}
设阻力系数为 $k(k>0)$，运动员质量为 $m$。
由牛顿第二定律，运动方程为 $m\frac{\mathrm{d}v}{\mathrm{d}t} = mg - kv$，这可化为 $\frac{\mathrm{d}v}{\mathrm{d}t} = g - \frac{k}{m}v = -\frac{k}{m}\left(v - \frac{mg}{k}\right)$。分离变量并积分得
\(\displaystyle \frac{\mathrm{d}v}{v-\frac{mg}{k}} = -\frac{k}{m}\mathrm{d}t\)，即
\(\displaystyle \ln\left|v-\frac{mg}{k}\right| = -\frac{k}{m}t + C_1\)。
从而 $v = \frac{mg}{k} + C\mathrm{e}^{-\frac{k}{m}t}$。
代入初值 $v(0)=0$ 得 $0 = \frac{mg}{k} + C \implies C = -\frac{mg}{k}$。
因此，速度和时间关系为 $\mathbf{v(t) = \frac{mg}{k}\left(1 - \mathrm{e}^{-\frac{k}{m}t}\right)}$。
\end{solution}

\begin{exercise}[实际应用题]
一电动机开动后, 每分钟温度升高 $10^\circ\mathrm{C}$, 同时将按冷却定律不断散发热量. 设电动机安置在一保持 $15^\circ\mathrm{C}$ 恒温的房间里, 求电机温度 $\theta$ 与时间 $t$ 的函数关系.
\end{exercise}
\begin{solution}
设电机温度为 $\theta(t)$，散热服从牛顿冷却定律。综合起来有微分方程
\(\displaystyle \frac{\mathrm{d}\theta}{\mathrm{d}t}=10+k(\theta-15)\quad (k<0)\)，即
\(\displaystyle \frac{\mathrm{d}\theta}{\theta-\left(15-\frac{10}{k}\right)}=k\,\mathrm{d}t\)。从而
\(\displaystyle \ln\left|\theta-\left(15-\frac{10}{k}\right)\right|=kt+C_1\)，即
\(\displaystyle \theta=15-\frac{10}{k}+C\mathrm{e}^{kt}\)。
再由 $\theta(0)=15$ 得 $C=\dfrac{10}{k}$，所以 $\mathbf{\theta(t)=15-\frac{10}{k}\bigl(1-\mathrm{e}^{kt}\bigr)}$。
（若初值未给定为室温，则通解包含常数项 $C$）。
\end{solution}

\subsection{第三节习题：一阶线性微分方程}

\begin{exercise}[基础：直接套用公式或常数变易法]
求解下列一阶线性方程：
\begin{shortexenum}
  \shortitem{1}{$y'+y=\mathrm{e}^{-x},\; y(0)=1$} &
  \shortitem{2}{$xy'-2y=x^{2}\ln x$} \\
  \shortitem{3}{$y'\cos x+y\sin x=\cos^{2}x,\; y(\pi)=\pi$} &
\end{shortexenum}
\end{exercise}
\begin{solution}
\begin{enumerate}
  \item[(1)] \textbf{求解 $y'+y=\mathrm{e}^{-x},\; y(0)=1$}
  \begin{enumerate}
      \item 此为标准形式的一阶线性微分方程，其中 $P(x)=1, Q(x)=\mathrm{e}^{-x}$。
      \item 计算积分因子：$\mu(x) = \exp\!\left(\int 1\,\mathrm{d}x\right) = \mathrm{e}^{x}$。
      \item 方程两边同乘积分因子 $\mathrm{e}^{x}$，左边化为完全导数：\(\displaystyle \mathrm{e}^{x}y' + \mathrm{e}^{x}y = 1 \implies (y\mathrm{e}^{x})' = 1\)。
      \item 两边同时对 $x$ 积分，得 \(\displaystyle \int (y\mathrm{e}^{x})'\,\mathrm{d}x = \int 1\,\mathrm{d}x \implies y\mathrm{e}^{x} = x + C \implies y = (x+C)\mathrm{e}^{-x}\)。
      \item 代入初值条件 $y(0)=1$，得 \(\displaystyle 1 = (0+C)\mathrm{e}^{0} \implies C=1\)。
      \item 故初值问题的特解为：$\mathbf{y = (x+1)\mathrm{e}^{-x}}$。
  \end{enumerate}

  \item[(2)] \textbf{求解 $xy'-2y=x^{2}\ln x$}
  \begin{enumerate}
      \item 将方程两边同除以 $x$，化为标准形式：\(\displaystyle y' - \frac{2}{x}y = x\ln x\)。
      \item 计算积分因子（注意 $\mathrm{e}^{\ln a} = a$）：\(\displaystyle \mu(x) = \exp\!\left(\int-\frac{2}{x}\,\mathrm{d}x\right) = \mathrm{e}^{-2\ln|x|} = x^{-2}\)。
      \item 方程两边同乘 $x^{-2}$，左边凑成完全导数形式：\(\displaystyle x^{-2}y' - 2x^{-3}y = x^{-1}\ln x \implies (yx^{-2})' = \frac{\ln x}{x}\)。
      \item 两边同时积分，右侧利用凑微分法 $\mathrm{d}(\ln x) = \frac{1}{x}\mathrm{d}x$，得 \(\displaystyle yx^{-2} = \int \frac{\ln x}{x}\,\mathrm{d}x = \int \ln x\,\mathrm{d}(\ln x) = \frac{1}{2}(\ln x)^2 + C\)。
      \item 两边同乘 $x^2$，整理得通解：$\mathbf{y = x^2\left[ \frac{1}{2}(\ln x)^2 + C \right]}$。
  \end{enumerate}

  \item[(3)] \textbf{求解 $y'\cos x+y\sin x=\cos^{2}x,\; y(\pi)=\pi$}
  \begin{enumerate}
      \item 方程两边同除以 $\cos x$，化为标准型：\(\displaystyle y' + y\tan x = \cos x\)。
      \item 计算积分因子：\(\displaystyle \mu(x) = \exp\!\left(\int\tan x\,\mathrm{d}x\right) = \exp(-\ln|\cos x|) = \sec x\)。
      \item 方程两边同乘 $\sec x$（即 $\frac{1}{\cos x}$），得 \(\displaystyle y'\sec x + y\sec x\tan x = 1 \implies (y\sec x)' = 1\)。
      \item 两边积分得到 \(\displaystyle y\sec x = \int 1\,\mathrm{d}x \implies y\sec x = x + C \implies y = (x+C)\cos x\)。
      \item 代入初值条件 $y(\pi)=\pi$，得 \(\displaystyle \pi = (\pi+C)\cos\pi \implies \pi = -(\pi+C) \implies \pi = -\pi - C \implies C = -2\pi\)。
      \item 故特解为：$\mathbf{y = (x-2\pi)\cos x}$。
  \end{enumerate}
\end{enumerate}
\end{solution}

\begin{exercise}[基础：伯努利方程]
求解下列伯努利方程（先判断指数 $n$）：\\
（1） $y'-\dfrac{1}{x}y=x^{2}y^{3}$~~（2）$y'+\dfrac{2}{x}y=\dfrac{1}{x^{2}}y^{2}\sin x$~~（3）$3(1+x^{2})y'-2xy=(1+x^{2})^{2}y^{4}$
\end{exercise}
下面三题均先化为伯努利方程的线性化形式：同除以 $y^n$ 后，把 $y'$ 与 $y$ 的幂次凑成可微结构，再换元化成线性方程。
\begin{solution}
\begin{enumerate}
  \item[(1)] \textbf{求解 $y'-\frac{1}{x}y=x^{2}y^{3}$}
  \begin{enumerate}
      \item 此为 $n=3$ 的伯努利方程。两边同除以 $y^3$，得 \(\displaystyle y^{-3}y' - \frac{1}{x}y^{-2} = x^2\)。
      \item 作变量代换，令 $z = y^{1-3} = y^{-2}$。对其求导得 $z' = -2y^{-3}y'$，即 $y^{-3}y' = -\frac{1}{2}z'$，故代回原方程可得 \(\displaystyle -\frac{1}{2}z' - \frac{1}{x}z = x^2 \implies z' + \frac{2}{x}z = -2x^2\)。
      \item 计算积分因子：\(\displaystyle \mu(x) = \exp\!\left(\int \frac{2}{x}\,\mathrm{d}x\right)= \mathrm{e}^{2\ln|x|} = x^2\)。两边同乘 $x^2$，化为 \((zx^2)' = -2x^4\)，两边积分得 \(\displaystyle zx^2 = \int -2x^4\,\mathrm{d}x = -\frac{2}{5}x^5 + C \implies z = Cx^{-2} - \frac{2}{5}x^3\)。
      \item 回代 $z=y^{-2}$，得通解：\(\displaystyle y^{-2} = Cx^{-2} - \frac{2}{5}x^3 \implies \mathbf{y^2 = \frac{5x^2}{5C - 2x^5}} \quad (\text{另外 } y\equiv0 \text{ 也是解})\)。
  \end{enumerate}

  \item[(2)] \textbf{求解 $y'+\frac{2}{x}y=\frac{1}{x^{2}}y^{2}\sin x$}
  \begin{enumerate}
      \item 此为 $n=2$ 的伯努利方程。两边同除以 $y^2$，得 \(\displaystyle y^{-2}y' + \frac{2}{x}y^{-1} = \frac{\sin x}{x^2}\)。
      \item 作代换 $z = y^{1-2} = y^{-1}$，则 $z' = -y^{-2}y'$。代入化简得 \(\displaystyle -z' + \frac{2}{x}z = \frac{\sin x}{x^2} \implies z' - \frac{2}{x}z = -\frac{\sin x}{x^2}\)。
      \item 计算积分因子：\(\displaystyle \mu(x) = \exp\!\left(\int-\frac{2}{x}\,\mathrm{d}x\right) = x^{-2}\)。方程化为 \(\displaystyle (zx^{-2})' = -\frac{\sin x}{x^4}\)。
      两边积分（右侧为非初等积分，保留积分号）得 \(\displaystyle zx^{-2} = \int -\frac{\sin x}{x^4}\,\mathrm{d}x + C \implies z = x^2\left( C - \int \frac{\sin x}{x^4}\,\mathrm{d}x \right)\)。
      \item 回代 $z=y^{-1}$，得通解：\(\displaystyle \mathbf{y = \frac{1}{x^2\left( C - \displaystyle\int \frac{\sin x}{x^4}\,\mathrm{d}x \right)}} \quad (\text{另外 } y\equiv0 \text{ 也是解})\)。
  \end{enumerate}

  \item[(3)] \textbf{求解 $3(1+x^{2})y'-2xy=(1+x^{2})^{2}y^{4}$}
  \begin{enumerate}
      \item 先将方程化为标准形式：\(\displaystyle y' - \frac{2x}{3(1+x^2)}y = \frac{1+x^2}{3}y^4\)。这是 $n=4$ 的伯努利方程。两边同除以 $y^4$，得 \(\displaystyle y^{-4}y' - \frac{2x}{3(1+x^2)}y^{-3} = \frac{1+x^2}{3}\)。
      \item 作代换 $z = y^{1-4} = y^{-3}$，则 $z' = -3y^{-4}y'$，即 $y^{-4}y' = -\frac{1}{3}z'$。代入得 \(\displaystyle -\frac{1}{3}z' - \frac{2x}{3(1+x^2)}z = \frac{1+x^2}{3} \implies z' + \frac{2x}{1+x^2}z = -(1+x^2)\)。
      \item 计算积分因子：\(\displaystyle \mu(x) = \exp\!\left(\int \frac{2x}{1+x^2}\,\mathrm{d}x\right)= \exp(\ln(1+x^2)) = 1+x^2\)。方程化为 \(\displaystyle (z(1+x^2))' = -(1+x^2)^2 = -(1 + 2x^2 + x^4)\)，
      两边积分得 \(\displaystyle z(1+x^2)= \int -(1 + 2x^2 + x^4)\,\mathrm{d}x = -x - \frac{2}{3}x^3 - \frac{1}{5}x^5 + C\)。
      \item 回代 $z=y^{-3}$，整理得通解：\(\displaystyle \mathbf{y^{-3} = \frac{C - x - \frac{2}{3}x^3 - \frac{1}{5}x^5}{1+x^2}} \quad (\text{另外 } y\equiv0 \text{ 也是解})\)。
  \end{enumerate}
\end{enumerate}
\end{solution}

\begin{exercise}[提高：角色互换与全微分]
解方程 $y'\,(x\ln y-1)=y\ln y$。
\end{exercise}
\begin{solution}
\begin{enumerate}
    \item \textbf{化为微分形式}：
    先将导数 $y'$ 替换为微商形式 $\frac{\mathrm{d}y}{\mathrm{d}x}$，得 \(\displaystyle \frac{\mathrm{d}y}{\mathrm{d}x}(x\ln y-1)=y\ln y\)。
    方程两边同乘 $\mathrm{d}x$，将变导数形式转化为纯微分形式，并展开：\(\displaystyle x\ln y\,\mathrm{d}y - \mathrm{d}y = y\ln y\,\mathrm{d}x\)。

    \item \textbf{项的重组与凑微分}：
    移项，将含有 $x, y$ 乘积的项放在等式同一侧，提取公因式 $\ln y$：\(\displaystyle x\ln y\,\mathrm{d}y - y\ln y\,\mathrm{d}x = \mathrm{d}y \implies \ln y\,(x\,\mathrm{d}y - y\,\mathrm{d}x) = \mathrm{d}y\)。
这里出现了常见的微分结构 $(x\,\mathrm{d}y - y\,\mathrm{d}x)$。结合商的微分公式 \(\displaystyle \mathrm{d}\left(\frac{x}{y}\right) = \frac{y\,\mathrm{d}x - x\,\mathrm{d}y}{y^2}\)，
在方程两边同除以 $y^2$，得 \(\displaystyle \ln y\,\frac{x\,\mathrm{d}y - y\,\mathrm{d}x}{y^2}= \frac{\mathrm{d}y}{y^2}\)。
由于 $\frac{x\,\mathrm{d}y - y\,\mathrm{d}x}{y^2} = -\mathrm{d}\left(\frac{x}{y}\right)$，原方程可化为 \(\displaystyle -\ln y\,\mathrm{d}\left(\frac{x}{y}\right) = \frac{\mathrm{d}y}{y^2} \implies \mathrm{d}\left(\frac{x}{y}\right) = -\frac{\mathrm{d}y}{y^2\ln y}\)。

    \item \textbf{直接积分求通解}：
    等式两边已经完全分离，直接进行积分：\(\displaystyle \frac{x}{y} = \int -\frac{1}{y^2\ln y}\,\mathrm{d}y + C \implies x = Cy - y \int \frac{1}{y^2\ln y}\,\mathrm{d}y\)。
注：右侧的积分为超越积分，无法用初等函数有限形式表示。保留积分号即为最终解：\(\displaystyle \mathbf{x = Cy - y \int \frac{\mathrm{d}y}{y^2\ln y}}\)。
\end{enumerate}
\end{solution}

\begin{exercise}[提高：三角变形与全微分]
解方程 $y'\sec x-y\tan x=\sec x\tan x$。
\end{exercise}
\begin{solution}
\begin{enumerate}
    \item \textbf{化为微分形式}：
    先将 $y'$ 替换为 $\frac{\mathrm{d}y}{\mathrm{d}x}$：
    \(\displaystyle \frac{\mathrm{d}y}{\mathrm{d}x}\sec x-y\tan x=\sec x\tan x\)
    方程两边同乘 $\mathrm{d}x$，转化为纯微分形式：
    \(\displaystyle \sec x\,\mathrm{d}y - y\tan x\,\mathrm{d}x = \sec x\tan x\,\mathrm{d}x\)

    \item \textbf{化简三角函数与凑微分}：
    为消除分母并暴露出更基本的微分关系，方程两边同时乘以 $\cos^2 x$，得
    \(\displaystyle \cos x\,\mathrm{d}y - y\sin x\,\mathrm{d}x = \sin x\,\mathrm{d}x\)。
    又因为 $\mathrm{d}(\cos x) = -\sin x\,\mathrm{d}x$，
    所以 \(\displaystyle \cos x\,\mathrm{d}y + y\,\mathrm{d}(\cos x) = -\mathrm{d}(\cos x)\)，
    进而 \(\displaystyle \mathrm{d}(y\cos x) = \mathrm{d}(-\cos x)\)。

    \item \textbf{直接积分求解}：
    对两边直接去微分符号（积分），并加上任意常数 $C$：
    \(\displaystyle y\cos x = -\cos x + C\)
    两边同除以 $\cos x$，解出通解：
    \(\displaystyle \mathbf{y = C\sec x - 1}\)
\end{enumerate}
\end{solution}

\begin{exercise}[黎卡提方程]
验证 $y=\dfrac{1}{x}$ 是黎卡提方程 $y'=-y^{2}-\dfrac{1}{x}y+\dfrac{1}{x^2}$ 的一个特解，并由此求出通解。
\end{exercise}
\begin{solution}
\begin{enumerate}
    \item \textbf{验证特解}：将 $y = \frac{1}{x}$ 及其导数 $y' = -\frac{1}{x^2}$ 代入原方程：
    左端 $= -\frac{1}{x^2}$；
    右端 $= - \left(\frac{1}{x}\right)^2 - \frac{1}{x}\left(\frac{1}{x}\right) + \frac{1}{x^2} = -\frac{1}{x^2} - \frac{1}{x^2} + \frac{1}{x^2} = -\frac{1}{x^2}$。
    左右两端相等，故 $y = \frac{1}{x}$ 确为方程的特解。

    \item \textbf{刘维尔代换}：利用找到的特解，设原方程的解为 $y = \frac{1}{x} + u(x)$，对其求导得 $y' = -\frac{1}{x^2} + u'$。代入原方程可得 \(\displaystyle -\frac{1}{x^2}+u'=-\left(\frac{1}{x}+u\right)^2-\frac{1}{x}\left(\frac{1}{x}+u\right)+\frac{1}{x^2}\)，即 \(\displaystyle -\frac{1}{x^2}+u'=-\left(\frac{1}{x^2}+\frac{2u}{x}+u^2\right)-\left(\frac{1}{x^2}+\frac{u}{x}\right)+\frac{1}{x^2}\)。

    \item \textbf{展开并消元}：去括号并合并同类项，注意常数项 $-\frac{1}{x^2}$ 在左右两边恰好抵消，可得 \(\displaystyle u' = -\frac{2}{x}u - u^2 - \frac{1}{x}u \implies u' = -\frac{3}{x}u - u^2\)。

    \item \textbf{化为伯努利方程求解}：整理得 $u' + \frac{3}{x}u = -u^2$，这是 $n=2$ 的伯努利方程。
    两边同除以 $u^2$，得 $u^{-2}u' + \frac{3}{x}u^{-1} = -1$。令 $w = u^{-1}$，则 $w' = -u^{-2}u'$，代入方程得
    \(\displaystyle -w' + \frac{3}{x}w = -1 \implies w' - \frac{3}{x}w = 1\)。计算积分因子：\(\displaystyle \mu(x) = \exp\!\left(\int -\frac{3}{x}\,\mathrm{d}x\right) = x^{-3}\)。
    方程化为 \((wx^{-3})' = x^{-3}\)，两边积分得 \(\displaystyle wx^{-3} = \int x^{-3}\,\mathrm{d}x = -\frac{1}{2}x^{-2} + C \implies w = Cx^3 - \frac{x}{2}\)。

    \item \textbf{回代求最终通解}：由于 $u = w^{-1}$，有 $u = \frac{1}{Cx^3 - \frac{x}{2}}$。代回 $y = \frac{1}{x} + u$ 并通分化简得
    \(\displaystyle y = \frac{1}{x} + \frac{1}{Cx^3 - \frac{x}{2}} = \frac{1}{x} + \frac{2}{2Cx^3 - x} = \mathbf{\frac{2Cx^2 + 1}{2Cx^3 - x}}\)。
\end{enumerate}
\end{solution}

\subsubsection{A组}

\begin{exercise}[基本概念]
回答下列问题:
\begin{shorttrienum}
  \shortitem{1}{什么叫一阶线性微分方程?} &
  \shortitem{2}{常数变易法的基本思想是什么?} &
  \shortitem{3}{如何将伯努利方程化为线性方程?}
\end{shorttrienum}
\end{exercise}
\begin{solution}
\begin{enumerate}
  \item[(1)] 若一个一阶微分方程中，关于未知函数 $y$ 及其一阶导数 $y'$ 都是一次的（即没有平方、高次方，也没有将 $y$ 放在诸如 $\sin y$ 等非线性函数内部的情形），则称之为一阶线性微分方程。其一般标准形式可写作 $\frac{\mathrm{d}y}{\mathrm{d}x} + P(x)y = Q(x)$。

  \item[(2)] 常数变易法的基本思想是：首先求解非齐次方程所对应的齐次方程 $y' + P(x)y = 0$，得到其通解形式 $y_h = C \mathrm{e}^{-\int P(x)\mathrm{d}x}$。然后，我们保留该解的整体代数结构，但假设通解中的任意常数 $C$ 不是固定的常数，而是一个关于自变量 $x$ 的待定函数 $C(x)$，即令 $y = C(x) \mathrm{e}^{-\int P(x)\mathrm{d}x}$。将其代入原非齐次方程中，方程中的 $C(x)$ 项会必然相消，只剩下 $C'(x)$ 项，通过一次积分求解出 $C(x)$，最终得到非齐次方程的通解。

  \item[(3)] 对于标准形式的伯努利方程 $y' + P(x)y = Q(x)y^n$（其中 $n \neq 0, 1$），通过以下两步核心代换可将其化为线性方程：
      第一步：等式两边同除以非线性因子 $y^n$，得到 $y^{-n}y' + P(x)y^{1-n} = Q(x)$。
      第二步：引入新的未知函数进行变量代换，令 $z = y^{1-n}$。利用复合函数求导法则，其导数为 $z' = (1-n)y^{-n}y'$。将该关系代入上一步骤的方程中，即可把原方程化为关于 $z$ 的一阶线性微分方程：\(\displaystyle z' + (1-n)P(x)z = (1-n)Q(x)\)。
\end{enumerate}
\end{solution}

\begin{exercise}[一阶线性微分方程求通解]
求下列微分方程的通解:\\
\begin{tabular}{ll}
  (1) $\displaystyle \frac{\mathrm{d}y}{\mathrm{d}x} + y\tan x = \cos x$; &
  (2) $\displaystyle \frac{\mathrm{d}y}{\mathrm{d}x} - y\cot x = 2x\sin x$; \\[2ex]
  (3) $\displaystyle (y^2 - 6x)\frac{\mathrm{d}y}{\mathrm{d}x} + 2y = 0$; &
  (4) $\displaystyle y' + y\cos x = \mathrm{e}^{-\sin x}$.
\end{tabular}
\end{exercise}
\begin{solution}
\begin{enumerate}
    \item[(1)] \textbf{求解 $\frac{\mathrm{d}y}{\mathrm{d}x} + y\tan x = \cos x$}
    \begin{enumerate}
        \item 积分因子 $\mu(x) = \exp\!\left(\int \tan x\,\mathrm{d}x\right) = \mathrm{e}^{-\ln|\cos x|} = \sec x$。
        \item 方程两边同乘 $\sec x$ 得到 $(y\sec x)' = \cos x \cdot \sec x = 1$。
        \item 积分得 $y\sec x = \int 1\,\mathrm{d}x = x + C$。
        \item 故通解为：$\mathbf{y = (x+C)\cos x}$。
    \end{enumerate}

    \item[(2)] \textbf{求解 $\frac{\mathrm{d}y}{\mathrm{d}x} - y\cot x = 2x\sin x$}
    \begin{enumerate}
        \item 积分因子 $\mu(x) = \exp\!\left(\int -\cot x\,\mathrm{d}x\right) = \mathrm{e}^{-\ln|\sin x|} = \csc x$。
        \item 两边同乘 $\csc x$ 得 $(y\csc x)' = 2x\sin x \cdot \csc x = 2x$。
        \item 积分得 $y\csc x = \int 2x\,\mathrm{d}x = x^2 + C$。
        \item 故通解为：$\mathbf{y = (x^2+C)\sin x}$。
    \end{enumerate}

    \item[(3)] \textbf{求解 $(y^2 - 6x)\frac{\mathrm{d}y}{\mathrm{d}x} + 2y = 0$}
    \begin{enumerate}
        \item 方程关于 $y$ 不是线性的，采取角色互换策略，视 $x$ 为 $y$ 的函数。将原方程转化为
        \(\displaystyle 2y\frac{\mathrm{d}x}{\mathrm{d}y} = 6x - y^2,\qquad
        \frac{\mathrm{d}x}{\mathrm{d}y} - \frac{3}{y}x = -\frac{y}{2}\)。
        \item 这是一阶线性微分方程。积分因子 $\mu(y) = \exp\!\left(\int -\frac{3}{y}\,\mathrm{d}y\right) = y^{-3}$。
        \item 方程化为 $(x y^{-3})' = -\frac{y}{2}y^{-3} = -\frac{1}{2}y^{-2}$。
        \item 积分得 $xy^{-3} = \int -\frac{1}{2}y^{-2}\,\mathrm{d}y = \frac{1}{2}y^{-1} + C$。
        \item 故通解为：$\mathbf{x = Cy^3 + \frac{1}{2}y^2} \quad (\text{另外 } y\equiv0 \text{ 也是解})$。
    \end{enumerate}

    \item[(4)] \textbf{求解 $y' + y\cos x = \mathrm{e}^{-\sin x}$}
    \begin{enumerate}
        \item 积分因子 $\mu(x) = \exp\!\left(\int \cos x\,\mathrm{d}x\right) = \mathrm{e}^{\sin x}$。
        \item 两边同乘 $\mathrm{e}^{\sin x}$ 得 $(y\mathrm{e}^{\sin x})' = \mathrm{e}^{-\sin x} \cdot \mathrm{e}^{\sin x} = 1$。
        \item 积分得 $y\mathrm{e}^{\sin x} = \int 1\,\mathrm{d}x = x + C$。
        \item 故通解为：$\mathbf{y = (x+C)\mathrm{e}^{-\sin x}}$。
    \end{enumerate}
\end{enumerate}
\end{solution}

\begin{exercise}[一阶线性微分方程初值问题]
求下列初值问题的解:\\
\begin{tabular}{ll}
  (1) $\displaystyle (x+1)y' + y = 2\mathrm{e}^{-x},\,y(1)=0$; &
  (2) $\displaystyle \frac{\mathrm{d}y}{\mathrm{d}x} - y\tan x = \sec x,\,y(0)=0$; \\[2ex]
  (3) $\displaystyle \frac{\mathrm{d}y}{\mathrm{d}x} + \frac{3}{x}y = \frac{2}{x^3},\,y(1)=1$; &
  (4) $\displaystyle y' + \frac{y}{x} + \mathrm{e}^x = 0,\,y(1)=0$.
\end{tabular}
\end{exercise}
\begin{solution}
\begin{enumerate}
  \item[(1)] \textbf{求解 $(x+1)y' + y = 2\mathrm{e}^{-x},\,y(1)=0$}
  \begin{enumerate}
      \item 观察方程左端，利用乘积导数法则 $(uv)'=u'v+uv'$，可知它恰好是 $((x+1)y)'$，于是 $((x+1)y)' = 2\mathrm{e}^{-x}$，从而 $(x+1)y = -2\mathrm{e}^{-x} + C$。
      \item 代入初值条件 $y(1)=0$，得 $C = 2\mathrm{e}^{-1}$，故 \(\displaystyle \mathbf{y} = \mathbf{\frac{2\mathrm{e}^{-1} - 2\mathrm{e}^{-x}}{x+1}}\)。
  \end{enumerate}

  \item[(2)] \textbf{求解 $\frac{\mathrm{d}y}{\mathrm{d}x} - y\tan x = \sec x,\,y(0)=0$}
  \begin{enumerate}
      \item 计算积分因子：$\mu(x) = \exp\!\left(\int -\tan x\,\mathrm{d}x\right) = \mathrm{e}^{\ln|\cos x|} = \cos x$。
      \item 方程两边同乘 $\cos x$ 并积分，得 $y'\cos x - y\sin x = 1 \implies (y\cos x)' = 1$，从而 $y\cos x = x + C$。
      \item 代入初值条件 $y(0)=0$，得 $C = 0$，故初值问题的解为 \(\displaystyle \mathbf{y} = \mathbf{\frac{x}{\cos x} = x\sec x}\)。
  \end{enumerate}

  \item[(3)] \textbf{求解 $\frac{\mathrm{d}y}{\mathrm{d}x} + \frac{3}{x}y = \frac{2}{x^3},\,y(1)=1$}
  \begin{enumerate}
      \item 计算积分因子：$\mu(x) = \exp\!\left(\int \frac{3}{x}\,\mathrm{d}x\right) = \mathrm{e}^{3\ln|x|} = x^3$。
      \item 两边同乘 $x^3$ 并积分，得 $(yx^3)' = 2$，从而 $yx^3 = 2x + C$。
      \item 代入初值条件 $y(1)=1$，得 $C = -1$，故初值问题的解为 \(\displaystyle \mathbf{y} = \mathbf{\frac{2x - 1}{x^3} = 2x^{-2} - x^{-3}}\)。
  \end{enumerate}

  \item[(4)] \textbf{求解 $y' + \frac{y}{x} + \mathrm{e}^x = 0,\,y(1)=0$}
  \begin{enumerate}
      \item 整理为标准形式 $y' + \frac{1}{x}y = -\mathrm{e}^x$。积分因子 $\mu(x) = \exp\!\left(\int \frac{1}{x}\,\mathrm{d}x\right) = x$。
      \item 两边同乘 $x$ 得到 $(xy)' = -x\mathrm{e}^x$。
      \item 利用分部积分法对右端进行积分 $\int u\,\mathrm{d}v = uv - \int v\,\mathrm{d}u$（令 $u=x, \mathrm{d}v=\mathrm{e}^x\mathrm{d}x$），得 \(\displaystyle xy = \int -x\mathrm{e}^x\,\mathrm{d}x = -\left(x\mathrm{e}^x - \int \mathrm{e}^x\,\mathrm{d}x\right) = (1-x)\mathrm{e}^x + C\)。
      \item 代入初值条件 $y(1)=0$，得 $C = 0$，故初值问题的解为 \(\displaystyle \mathbf{y} = \mathbf{\frac{1-x}{x}\mathrm{e}^x}\)。
  \end{enumerate}
\end{enumerate}
\end{solution}

\begin{exercise}[求曲线方程]
设一曲线通过原点, 且它在点$(x,y)$处的切线斜率等于$2x+y$. 求这曲线方程.
\end{exercise}
\begin{solution}
\begin{enumerate}
    \item \textbf{建立微分方程}：由导数的几何意义可知，曲线上任意一点 $(x,y)$ 的切线斜率即为其导数值 $y'$。根据题意建立微分方程 \(\displaystyle y' = 2x + y \implies y' - y = 2x\)。
    \item \textbf{求解一阶线性方程}：这是一阶线性常微分方程，积分因子为 $\mu(x) = \exp\!\left(\int -1\,\mathrm{d}x\right) = \mathrm{e}^{-x}$。方程两边同乘 $\mathrm{e}^{-x}$，再对两边关于 $x$ 积分，右侧运用分部积分法（令 $u=2x, \mathrm{d}v=\mathrm{e}^{-x}\mathrm{d}x$），得 \(\displaystyle (y\mathrm{e}^{-x})' = 2x\mathrm{e}^{-x}\)。积分并化简得 \(\displaystyle y\mathrm{e}^{-x} = -2x\mathrm{e}^{-x} - 2\mathrm{e}^{-x} + C,\quad y = -2(x+1) + C\mathrm{e}^x\)。
    \item \textbf{应用初值条件并得出最终方程}：由于曲线通过原点 $(0,0)$，代入通解可得 \(\displaystyle 0 = -2(0+1) + C\mathrm{e}^0 \implies C = 2,\qquad \mathbf{y = 2\mathrm{e}^x - 2x - 2}\)。
\end{enumerate}
\end{solution}

\begin{exercise}[应用题: 电路模型]
设有一个由电阻$R=10\ \Omega$、电感$L=2\ \mathrm{H}$和电源电压$E=20\sin5t$ (V) 串联组成的电路. 开关$\mathrm{S}$合上后, 电路中有电流通过. 求电流$I$与时间$t$的函数关系. (利用回路电压定律$\displaystyle E=RI+L\frac{\mathrm{d}I}{\mathrm{d}t}$.)
\end{exercise}
\begin{solution}
\begin{enumerate}
    \item \textbf{建立方程模型}：根据题干给出的基尔霍夫电压定律（回路电压定律），代入具体的物理参数 $R=10, L=2, E=20\sin 5t$，得
    \(\displaystyle 2\frac{\mathrm{d}I}{\mathrm{d}t} + 10I = 20\sin 5t
    \quad\Longrightarrow\quad
    \frac{\mathrm{d}I}{\mathrm{d}t} + 5I = 10\sin 5t\)。
    这是一个关于电流 $I(t)$ 的一阶线性非齐次常微分方程。
    \item \textbf{寻找积分因子}：计算积分因子 $\mu(t) = \exp\!\left(\int 5\,\mathrm{d}t\right) = \mathrm{e}^{5t}$。方程两边同乘 $\mathrm{e}^{5t}$，得 \(\displaystyle (I\mathrm{e}^{5t})' = 10\mathrm{e}^{5t}\sin 5t\)。
    \item \textbf{积分求解}：利用标准积分公式 \(\displaystyle \int \mathrm{e}^{at}\sin(bt)\,\mathrm{d}t = \frac{\mathrm{e}^{at}}{a^2+b^2}(a\sin bt - b\cos bt)\)，
    此处 $a=5, b=5$，有
    \[
    \int \mathrm{e}^{5t}\sin 5t\,\mathrm{d}t
    =\frac{\mathrm{e}^{5t}}{50}\cdot 5(\sin 5t-\cos 5t)
    =\frac{\mathrm{e}^{5t}}{10}(\sin 5t-\cos 5t).
    \]
    将积分结果代回，可得
    \(\displaystyle I\mathrm{e}^{5t}
    =10\left[\frac{\mathrm{e}^{5t}}{10}(\sin 5t-\cos 5t)\right]+C
    =\mathrm{e}^{5t}(\sin 5t-\cos 5t)+C\)，
    两边同乘 $\mathrm{e}^{-5t}$ 得到通解 \(\displaystyle I(t)=\sin 5t-\cos 5t+C\mathrm{e}^{-5t}\)。
    \item \textbf{利用物理初值条件求特解}：对于含有电感 $L$ 的 RL 串联电路，依据电磁学原理，合上开关的瞬间电流无法突变，即满足初值条件 $I(0) = 0$。代入求 $C$，得 \(\displaystyle 0=\sin0-\cos0+C\mathrm{e}^{0}= -1+C,\quad C=1\)。
    \item \textbf{得出结论}：电流与时间的函数关系为 \(\displaystyle \mathbf{I(t)=\sin 5t-\cos 5t+\mathrm{e}^{-5t}}\)。
\end{enumerate}
\end{solution}

\begin{exercise}[溶液混合问题]
一容器内盛有$50\ \mathrm{L}$的盐水溶液, 其中含有$10\ \mathrm{g}$的盐. 现将每升含$2\ \mathrm{g}$盐的溶液以每分钟$5\ \mathrm{L}$的速率注入容器, 并不断进行搅拌, 使混合液迅速达到均匀, 同时混合液以每分钟$3\ \mathrm{L}$的速率流出容器. 问在任一时刻$t$容器中含盐量是多少?
\end{exercise}
\begin{solution}
\begin{enumerate}
    \item \textbf{建立状态函数}：设时刻 $t$ (分钟) 时，容器中的含盐量为 $y(t)$ (克)。
    首先计算时刻 $t$ 容器内溶液的总体积 $V(t)$。初始体积为 $50\ \mathrm{L}$，每分钟流入 $5\ \mathrm{L}$，流出 $3\ \mathrm{L}$，体积净增加率为 $2\ \mathrm{L}/\mathrm{min}$，故 \(\displaystyle V(t) = 50 + (5-3)t = 50 + 2t \ (\mathrm{L})\)。
    \item \textbf{建立导数方程}：根据质量守恒定律，盐的浓度变化率 $\frac{\mathrm{d}y}{\mathrm{d}t} = \text{流入盐的速率} - \text{流出盐的速率}$：
    流入速率 = 注入溶液浓度 $\times$ 注入速率 = $2\ \mathrm{g/L} \times 5\ \mathrm{L/min} = 10\ \mathrm{g/min}$。
    流出速率 = 容器中瞬时浓度 $\times$ 流出速率 = $\frac{y(t)}{V(t)} \times 3\ \mathrm{L/min} = \frac{3y}{50+2t}\ \mathrm{g/min}$。
    由此建立一阶线性非齐次微分方程 \(\displaystyle \frac{\mathrm{d}y}{\mathrm{d}t} = 10 - \frac{3y}{50+2t} \implies y' + \frac{3}{50+2t}y = 10\)。
    \item \textbf{求解线性方程}：积分因子为 $\mu(t) = (50+2t)^{3/2}$，从而 \(\displaystyle \left(y(50+2t)^{3/2}\right)' = 10(50+2t)^{3/2},\qquad y(t) = 100 + 4t + C(50+2t)^{-3/2}\)。
    \item \textbf{应用初值条件求特解并得出最终函数}：题干指出初始 $t=0$ 时，容器内含 $10\ \mathrm{g}$ 盐，即初值条件 $y(0) = 10$。代入得 $C = -90(50)^{3/2}$，因此
    \(\displaystyle \mathbf{y(t)} = \mathbf{100 + 4t - 90\left(\frac{50}{50+2t}\right)^{\!3/2}}\)。
\end{enumerate}
\end{solution}

\subsubsection{B组}

\begin{exercise}[一阶线性微分方程初值问题]
求微分方程的初值问题 $\displaystyle y'\sec^2 y + \frac{x}{1+x^2}\tan y = x,\,y\big|_{x=0}=0$ 的解。
\end{exercise}
\begin{solution}
\begin{enumerate}
    \item \textbf{化为微分形式并提取基本微分}：
    先将导数 $y'$ 替换为微商形式 $\frac{\mathrm{d}y}{\mathrm{d}x}$，再将方程两边同乘 $\mathrm{d}x$，可得 \(\displaystyle \sec^2 y\,\mathrm{d}y + \frac{x}{1+x^2}\tan y\,\mathrm{d}x = x\,\mathrm{d}x\)。
    敏锐地观察到第一项 $\sec^2 y\,\mathrm{d}y$ 恰好是 $\tan y$ 的全微分，即 $\mathrm{d}(\tan y) = \sec^2 y\,\mathrm{d}y$，代入方程得 \(\displaystyle \mathrm{d}(\tan y) + \frac{x}{1+x^2}\tan y\,\mathrm{d}x = x\,\mathrm{d}x\)。

    \item \textbf{适当变形并凑全微分}：
    为了将左边凑成乘积的微分 $\mathrm{d}(uv) = u\,\mathrm{d}v + v\,\mathrm{d}u$，需要先处理中间项的分母 $(1+x^2)$。
    注意到 $\mathrm{d}(1+x^2) = 2x\,\mathrm{d}x$，因此可将方程两边同时乘以 $\sqrt{1+x^2}$：
    \(\displaystyle \mathrm{d}\left(\tan y \cdot \sqrt{1+x^2}\right)= x\sqrt{1+x^2}\,\mathrm{d}x\)。

    \item \textbf{两边直接积分}：
    对两边同时去掉微分符号（进行积分），右侧通过将 $x\mathrm{d}x$ 凑成 $\frac{1}{2}\mathrm{d}(1+x^2)$ 来计算，得 \(\displaystyle \tan y \cdot \sqrt{1+x^2}= \int (1+x^2)^{1/2} \cdot \frac{1}{2}\,\mathrm{d}(1+x^2)= \frac{1}{3}(1+x^2)^{3/2} + C\)。

    \item \textbf{代入初值求解特解}：
    利用初值条件 $y(0)=0$，此时 $x=0$，由 $\tan 0 \cdot \sqrt{1+0} = \frac13(1+0)^{3/2} + C$ 得 $C = -\frac13$，故 \(\displaystyle \tan y = \frac{1}{3}(1+x^2) - \frac{1}{3\sqrt{1+x^2}}\)。
    故 \(\displaystyle \mathbf{y} = \mathbf{\arctan\!\left[ \frac{1}{3}\left(1+x^2 - \frac{1}{\sqrt{1+x^2}}\right) \right]}\)。
\end{enumerate}
\end{solution}

\begin{exercise}[求通解]
求下列微分方程的通解:(1)$\displaystyle xy' = y(2y\ln x - 1)$;~~(2) $\displaystyle y' = x^3y^3 - xy$.
\end{exercise}
\begin{solution}
\begin{enumerate}
  \item[(1)] \textbf{求解 $xy' = y(2y\ln x - 1)$}
  \begin{enumerate}
      \item 展开右侧并整理为伯努利方程的标准形式：
      即 \(\displaystyle y' + \frac{1}{x}y = \frac{2\ln x}{x}y^2\)。
      \item 此为 $n=2$ 的伯努利方程。作代换 $z = y^{1-2} = y^{-1}$，则有 $z' = -y^{-2}y'$。
      方程两边同除以 $y^2$ 并乘以 $-1$：
      即 \(\displaystyle z' - \frac{1}{x}z = -\frac{2\ln x}{x}\)。
      \item 积分因子 $\mu(x) = \exp\!\left(\int -\frac{1}{x}\,\mathrm{d}x\right) = x^{-1}$，故转化为 $(zx^{-1})' = -\frac{2\ln x}{x^2}$。
      \item 两边积分，利用分部积分法求解右侧（令 $u=\ln x, \mathrm{d}v=-2x^{-2}\mathrm{d}x$）：
      \(\displaystyle zx^{-1} = \int -\frac{2\ln x}{x^2}\,\mathrm{d}x\)。
      分部积分后得 \(\displaystyle zx^{-1}=\frac{2\ln x}{x}-\int 2x^{-2}\,\mathrm{d}x=\frac{2\ln x}{x}+\frac{2}{x}+C\)。
      \item 两边同乘 $x$ 得通解 $z = 2\ln x + 2 + Cx$。回代 $z = y^{-1}$ 解得 \(\displaystyle \mathbf{y = \frac{1}{Cx + 2\ln x + 2}}\)，另外 $y\equiv0$ 也是解。
  \end{enumerate}

  \item[(2)] \textbf{求解 $y' = x^3y^3 - xy$}
  \begin{enumerate}
      \item 整理为标准形式 $y' + xy = x^3y^3$。此为 $n=3$ 的伯努利方程。
      \item 作代换 $z = y^{1-3} = y^{-2}$，则 $z' = -2y^{-3}y'$。方程两边同除以 $y^3$ 并乘以 $-2$：
      即 \(\displaystyle z' - 2xz = -2x^3\)。
      \item 积分因子 $\mu(x) = \exp\!\left(\int -2x\,\mathrm{d}x\right) = \mathrm{e}^{-x^2}$，故转化为 $(z\mathrm{e}^{-x^2})' = -2x^3\mathrm{e}^{-x^2}$。
      \item 两边积分。通过换元积分法结合分部积分法（令 \(u = x^2,\ \mathrm{d}u = 2x\mathrm{d}x\)），有
      \[
      z\mathrm{e}^{-x^2}=\int -2x^3\mathrm{e}^{-x^2}\,\mathrm{d}x
      =\int -u\mathrm{e}^{-u}\,\mathrm{d}u
      =(x^2+1)\mathrm{e}^{-x^2}+C .
      \]
      \item 得到通解 $z = x^2 + 1 + C\mathrm{e}^{x^2}$。回代 $z = y^{-2}$ 得
      \(\displaystyle \mathbf{y^2 = \frac{1}{x^2 + 1 + C\mathrm{e}^{x^2}}}\)，另外 $y\equiv0$ 也是解。
  \end{enumerate}
\end{enumerate}
\end{solution}

\begin{exercise}[积分方程求未知函数]
已知连续函数$f(x)$满足条件$\displaystyle f(x) = \int_{0}^{3x} f\left(\frac{t}{3}\right)\mathrm{d}t + \mathrm{e}^{2x}$, 求$f(x)$.
\end{exercise}
\begin{solution}
\begin{enumerate}
    \item \textbf{消除积分限与被积函数的参数耦合}：对积分项作换元处理。令 $u = \frac{t}{3}$，则 $\mathrm{d}t = 3\mathrm{d}u$。当积分下限 $t=0$ 时，新的积分下限 $u=0$；当积分上限 $t=3x$ 时，新的积分上限 $u=x$。
    将换元结果代入原方程，可得 \(\displaystyle f(x) = \int_{0}^{x} f(u) \cdot 3\mathrm{d}u + \mathrm{e}^{2x} = 3\int_{0}^{x} f(t)\mathrm{d}t + \mathrm{e}^{2x}\)。
    \item \textbf{化为微分方程}：这是一个含有未知函数变上限积分的积分方程。利用微积分基本定理，等式两边同时关于 $x$ 求导，化解积分号，得 \(\displaystyle f'(x) = 3f(x) + 2\mathrm{e}^{2x} \implies f'(x) - 3f(x) = 2\mathrm{e}^{2x}\)。
    这转化为一阶线性非齐次常微分方程。
    \item \textbf{求解线性方程并确定常数}：求出积分因子 $\mu(x) = \exp\!\left(\int -3\,\mathrm{d}x\right) = \mathrm{e}^{-3x}$。方程两边同乘 $\mathrm{e}^{-3x}$ 并积分，再利用 $x=0$ 时的隐含初值条件，得 \(\displaystyle (f(x)\mathrm{e}^{-3x})' = 2\mathrm{e}^{-x},\qquad f(x)\mathrm{e}^{-3x} = -2\mathrm{e}^{-x} + C\)。
    从而 $f(x) = -2\mathrm{e}^{2x} + C\mathrm{e}^{3x}$。又由 $f(0)=1$ 得 $-2 + C = 1$，所以 $C = 3$，于是 \(\displaystyle \mathbf{f(x)} = \mathbf{3\mathrm{e}^{3x} - 2\mathrm{e}^{2x}}\)。
\end{enumerate}
\end{solution}

\begin{exercise}[积分与变量无关]
设$f(x)$连续, 积分$\displaystyle \int_{0}^{1} [f(x) + xf(xt)]\mathrm{d}t$与$x$无关, 求$f(x)$.
\end{exercise}
\begin{solution}
\begin{enumerate}
    \item \textbf{积分项的拆分与化简}：将题干积分基于加法拆分为两项分别处理。对于第一项，积分变量是 $t$，因此 $f(x)$ 相当于常数，可以直接提出来；对于第二项，令 $u = xt$，则 $\mathrm{d}u = x\mathrm{d}t$，可得 \(\displaystyle I(x) = \int_{0}^{1} f(x)\mathrm{d}t + \int_{0}^{1} xf(xt)\mathrm{d}t = f(x) + \int_{0}^{x} f(u)\mathrm{d}u\)。
    \item \textbf{利用题意建立微分方程}：题干指出，该积分值与变量 $x$ 无关。也就是说，$I(x)$ 是一个常数。因此，它关于 $x$ 的导数必定恒等于 0。
    对 $I(x)$ 两边关于 $x$ 求导，得 \(\displaystyle I'(x) = f'(x) + f(x) = 0\)。
    \item \textbf{求解齐次微分方程}：这是一个极基础的齐次一阶线性（且可分离变量）微分方程。由 $f'(x) = -f(x)$ 得 $\frac{f'(x)}{f(x)} = -1$，积分可得 \(\displaystyle \ln|f(x)| = -x + C_1,\qquad f(x) = C\mathrm{e}^{-x}\)。
    （为了严谨可代回原式验算：$I(x) = C\mathrm{e}^{-x} + \int_0^x C\mathrm{e}^{-u}\,\mathrm{d}u = C\mathrm{e}^{-x} - C\mathrm{e}^{-x} + C = C$，确实与 $x$ 无关，恒为常数）。
    \item \textbf{得出最终结论}：满足条件的函数族为 $\mathbf{f(x) = C\mathrm{e}^{-x}}$（其中 $C$ 为任意实数常数）。
\end{enumerate}
\end{solution}

\subsection{第四节习题：可降阶的高阶方程}

\begin{exercise}[缺 $y$ 型与缺 $x$ 型]
求解下列方程：
\begin{shortexenum}
  \shortitem{1}{$(1+x^{2})y''+xy'=1,\;y(0)=0,\;y'(0)=-1$} &
  \shortitem{2}{$yy''+(y')^{2}=0,\;y(0)=1,\;y'(0)=1$} \\
  \shortitem{3}{$y''=y'^{3}/\sqrt{y}$} &
\end{shortexenum}
\end{exercise}
\begin{solution}
\begin{enumerate}
  \item[(1)] \textbf{思路}：方程缺 $y$，令 $p=y'$，则 $y''=p'$。
  代入原方程得 \(\displaystyle (1+x^2)p' + xp = 1 \Rightarrow p' + \frac{x}{1+x^2}p = \frac{1}{1+x^2}\)。这是一阶线性常微分方程，积分因子为 \(\displaystyle \mu(x)=\exp\!\left(\int \frac{x}{1+x^2}\,\mathrm{d}x\right)=\sqrt{1+x^2}\)。两边同乘 $\sqrt{1+x^2}$ 并积分得 \(\displaystyle (p\sqrt{1+x^2})' = \frac{1}{\sqrt{1+x^2}},\quad p\sqrt{1+x^2} = \ln(x+\sqrt{1+x^2}) + C_1\)。
  代入初值 $y'(0)=p(0)=-1$，得 $C_1=-1$，于是
  \(\displaystyle y' = p = \frac{\ln(x+\sqrt{1+x^2})-1}{\sqrt{1+x^2}},\qquad
  y = \frac{1}{2}\ln^2(x+\sqrt{1+x^2}) - \ln(x+\sqrt{1+x^2}) + C_2\)。
  由 $y(0)=0$ 得 $0 = 0 - 0 + C_2 \implies C_2=0$。特解为 \(\displaystyle \boxed{y = \frac{1}{2}\ln^2(x+\sqrt{1+x^2}) - \ln(x+\sqrt{1+x^2})}\)。

  \item[(2)] \textbf{思路}：利用“凑全微分法”，观察到左侧恰好是乘积求导法则的展开：\(\displaystyle (yy')' = yy'' + (y')^2 = 0,\qquad yy' = C_1,\qquad \frac{1}{2}y^2 = x + C_2\)。
  由初值 $y(0)=1, y'(0)=1$ 得 $C_1=1$，再由 $y(0)=1$ 得 $C_2 = \frac{1}{2}$，故
  \(\displaystyle y^2 = 2x + 1 \implies \boxed{y = \sqrt{2x+1}}\)。
  （因 $y(0)=1>0$ 故取正号）。

  \item[(3)] \textbf{思路}：方程缺 $x$，令 $p=y'$，则 $y''=p\frac{\mathrm{d}p}{\mathrm{d}y}$。
  代入方程得 \(\displaystyle p\frac{\mathrm{d}p}{\mathrm{d}y} = \frac{p^3}{\sqrt{y}}\)。
  若 $p \neq 0$，两边同除以 $p^3$ 并分离变量，得
  \(\displaystyle \frac{\mathrm{d}p}{p^2} = \frac{\mathrm{d}y}{\sqrt{y}},\qquad
  -\frac{1}{p} = 2\sqrt{y} + C_1
  \implies y' = p = \frac{-1}{2\sqrt{y}+C_1}\)。
  再次分离变量并积分得
  \(\displaystyle (2\sqrt{y}+C_1)\mathrm{d}y = -\mathrm{d}x,\qquad
  \boxed{\frac{4}{3}y^{3/2} + C_1y = -x + C_2}\)。
  此即为原方程的隐式通解。
\end{enumerate}
\end{solution}

\subsubsection{A组}

\begin{exercise}[求通解]
求下列各微分方程的通解:\\
\begin{tabular}{lll}
  (1) $y'' = x + \cos x$; &
  (2) $y'' = \dfrac{1}{1+x^2}$; &
  (3) $y'' - y' = x$; \\
  (4) $y'' + y'^2 = 1$; &
  (5) $yy'' + y'^2 = 0$; &
  (6) $y'' = y'^3 + y'$.
\end{tabular}
\end{exercise}
\begin{solution}
\begin{enumerate}
    \item[(1)] \textbf{直接积分法}。原方程 $y'' = x + \cos x$，积分两次得
    \(\displaystyle y' = \frac{1}{2}x^2 + \sin x + C_1,\qquad
    y = \frac{1}{6}x^3 - \cos x + C_1 x + C_2\)。

    \item[(2)] \textbf{直接积分法}。原方程 $y'' = \frac{1}{1+x^2}$，先得
    \(\displaystyle y' = \arctan x + C_1\)。再利用 $\int \arctan x \mathrm{d}x = x\arctan x - \int \frac{x}{1+x^2}\mathrm{d}x$，得通解
    \(\displaystyle y = x\arctan x - \frac{1}{2}\ln(1+x^2) + C_1 x + C_2\)。

    \item[(3)] \textbf{缺 $y$ 型，令 $p=y'$}。原方程化为 \(p' - p = x\)。这是一阶线性微分方程。积分因子为 $\mu(x) = \mathrm{e}^{-x}$，故
    \[
    (\mathrm{e}^{-x}p)' = x\mathrm{e}^{-x},
    \qquad
    \mathrm{e}^{-x}p = -x\mathrm{e}^{-x} - \mathrm{e}^{-x} + C_1.
    \]
    从而 \(p = y' = C_1\mathrm{e}^x - x - 1\)，再积分得通解
    \(\displaystyle y = C_1\mathrm{e}^x - \frac{1}{2}x^2 - x + C_2\)。

    \item[(4)] \textbf{缺 $y$ 型，令 $p=y'$}。方程化为 \(p' + p^2 = 1\)，即 $p' = 1-p^2$。分离变量（注意可能遗漏解 $p=\pm1 \implies y=\pm x+C$）得
    \(\displaystyle \frac{\mathrm{d}p}{1-p^2} = \mathrm{d}x,\quad
    \frac{1}{2}\ln\left|\frac{1+p}{1-p}\right| = x + C_1\)。
    因而先得 $p$，再积分得通解：
    \begin{align*}
    p = y' &= \frac{C\mathrm{e}^{2x}-1}{C\mathrm{e}^{2x}+1}
    = 1 - \frac{2}{C\mathrm{e}^{2x}+1},\\
    y &= x - \int \frac{2\mathrm{e}^{-2x}}{C+\mathrm{e}^{-2x}}\mathrm{d}x
    = x + \frac{1}{C}\ln(C+\mathrm{e}^{-2x}) + C_2.
    \end{align*}

    \item[(5)] \textbf{凑全微分法}。原方程 $yy'' + (y')^2 = 0$，注意到左端恰好是 $(yy')'$ 的展开式，因此 \(\displaystyle (yy')' = 0 \implies yy' = C_1 \implies \frac{1}{2}y^2 = C_1 x + C_2,\ y^2 = C_1 x + C_2\)。

    \item[(6)] \textbf{缺 $x$ 型，令 $p=y'$，则 $y''=p\frac{\mathrm{d}p}{\mathrm{d}y}$}。方程化为
    \(\displaystyle p\frac{\mathrm{d}p}{\mathrm{d}y} = p^3 + p \implies p\left(\frac{\mathrm{d}p}{\mathrm{d}y} - p^2 - 1\right) = 0\)。若 $p=0$，则 $y=C$。若 $p \neq 0$，则
    \(\displaystyle \frac{\mathrm{d}p}{\mathrm{d}y} = p^2 + 1 \implies \frac{\mathrm{d}p}{p^2+1} = \mathrm{d}y,
    \qquad
    \arctan p = y + C_1 \implies p = y' = \tan(y+C_1)\)。
    分离变量并积分得 \(\displaystyle \cot(y+C_1)\mathrm{d}y = \mathrm{d}x,\qquad
    \ln|\sin(y+C_1)| = x + C_2\)。
\end{enumerate}
\end{solution}

\begin{exercise}[求初值问题]
求解下列初值问题:\\
\begin{tabular}{ll}
  (1) $y^3 y'' + 1 = 0,\ y(1)=1,\ y'(1)=0$; &
  (2) $y''' = e^{2x},\ y(1)=y'(1)=y''(1)=0$; \\
  (3) $y'' = 3\sqrt{y},\ y(0)=1,\ y'(0)=2$; &
  (4) $y'' - e^{2y} = 0,\ y(0)=0,\ y'(0)=1$; \\
  (5) $yy'' = 2(y'^2 - y'),\ y(0)=1,\ y'(0)=2$; &
  (6) $y'' = \dfrac{3x^2}{1+x^3}y',\ y(0)=1,\ y'(0)=4$.
\end{tabular}
\end{exercise}
\begin{solution}
\begin{enumerate}
    \item[(1)] \textbf{缺 $x$ 型，令 $p=y'$}，则 $y''=p\frac{\mathrm{d}p}{\mathrm{d}y}$。代入得：
    \(\displaystyle y^3 p \frac{\mathrm{d}p}{\mathrm{d}y} = -1 \implies p\mathrm{d}p = -y^{-3}\mathrm{d}y\)，积分得
    \(\displaystyle \frac{1}{2}p^2 = \frac{1}{2}y^{-2} + C_1\)。
    代入初始条件：当 $x=1$ 时，$y=1, p=y'=0$。即 $0 = \frac{1}{2}(1) + C_1 \implies C_1 = -\frac{1}{2}$。
    于是
    \(\displaystyle p^2 = y^{-2} - 1 = \frac{1-y^2}{y^2} \implies y' = \pm\frac{\sqrt{1-y^2}}{y},
    \qquad
    \frac{y\mathrm{d}y}{\sqrt{1-y^2}} = \pm\mathrm{d}x \implies -\sqrt{1-y^2} = \pm x + C_2\)。
    当 $x=1$ 时 $y=1$，得 $0 = \pm 1 + C_2 \implies C_2 = \mp 1$。
    代回得 $-\sqrt{1-y^2} = \pm(x-1)$。平方得 $1-y^2 = (x-1)^2$。
    故特解为圆的方程（的局部）：
    \(\displaystyle \boxed{(x-1)^2 + y^2 = 1}\)。

    \item[(2)] \textbf{直接积分法}。$y''' = \mathrm{e}^{2x}$。
    积分一次得 \(\displaystyle y'' = \frac{1}{2}\mathrm{e}^{2x} + C_1 \xrightarrow{y''(1)=0} C_1=-\frac{1}{2}\mathrm{e}^2\)；
    再积一次得 \(\displaystyle y' = \frac{1}{4}\mathrm{e}^{2x} - \frac{1}{2}\mathrm{e}^2 x + C_2 \xrightarrow{y'(1)=0} C_2=\frac{1}{4}\mathrm{e}^2\)；
    第三次积分得 \(\displaystyle y = \frac{1}{8}\mathrm{e}^{2x} - \frac{1}{4}\mathrm{e}^2 x^2 + \frac{1}{4}\mathrm{e}^2 x + C_3 \xrightarrow{y(1)=0} C_3=-\frac{1}{8}\mathrm{e}^2\)。
    特解：
    \(\displaystyle \boxed{y = \frac{1}{8}\mathrm{e}^{2x} - \frac{1}{4}\mathrm{e}^2 x^2 + \frac{1}{4}\mathrm{e}^2 x - \frac{1}{8}\mathrm{e}^2}\)。

    \item[(3)] \textbf{缺 $x$ 型，令 $p=y'$}，得：
    \(\displaystyle p\frac{\mathrm{d}p}{\mathrm{d}y} = 3y^{1/2}\)，积分得 \(\displaystyle \frac{1}{2}p^2 = 2y^{3/2} + C_1\)。
    由初值 $y(0)=1, p=2$，得 $2 = 2(1) + C_1 \implies C_1=0$。所以
    \(\displaystyle p^2 = 4y^{3/2} \implies y' = 2y^{3/4} \quad (\text{因 } y'(0)>0 \text{ 取正}),
    \qquad
    y^{-3/4}\mathrm{d}y = 2\mathrm{d}x \implies 4y^{1/4} = 2x + C_2\)。
    由初值 $y(0)=1$ 得 $4(1) = 0 + C_2 \implies C_2=4$。
    特解：
    \(\displaystyle 4y^{1/4} = 2x + 4 \implies y^{1/4} = \frac{x}{2} + 1 \implies \boxed{y = \left(\frac{x}{2} + 1\right)^4}\)。

    \item[(4)] \textbf{缺 $x$ 型，令 $p=y'$}，得
    \(\displaystyle p\frac{\mathrm{d}p}{\mathrm{d}y} = \mathrm{e}^{2y},\qquad
    \frac{1}{2}p^2 = \frac{1}{2}\mathrm{e}^{2y} + C_1\)。
    初值 $y(0)=0, p=1$，得 $\frac{1}{2} = \frac{1}{2} + C_1 \implies C_1=0$。故 \( \displaystyle p = \mathrm{e}^y \implies \mathrm{e}^{-y}\mathrm{d}y = \mathrm{d}x\)，积分得 \(\displaystyle -\mathrm{e}^{-y} = x + C_2\)。
    初值 $y(0)=0 \implies -1 = 0 + C_2 \implies C_2=-1$。特解：
    \(\displaystyle -\mathrm{e}^{-y} = x - 1 \implies \mathrm{e}^{-y} = 1-x \implies \boxed{y = -\ln(1-x)}\)。

    \item[(5)] \textbf{缺 $x$ 型，令 $p=y'$}。方程为
    \(\displaystyle yp\frac{\mathrm{d}p}{\mathrm{d}y} = 2(p^2-p)\)。分离变量（已知 $p=2 \neq 0,1$）得
    \(\displaystyle \frac{p\mathrm{d}p}{p(p-1)} = \frac{2\mathrm{d}y}{y} \implies \frac{\mathrm{d}p}{p-1} = \frac{2\mathrm{d}y}{y},
    \qquad
    \ln|p-1| = 2\ln y + \ln C_1 \implies p-1 = C_1 y^2\)。
    初值 $y(0)=1, p=2$，得 $1 = C_1(1) \implies C_1=1$。于是 $p = y' = y^2 + 1$。
    分离变量得 \(\displaystyle \frac{\mathrm{d}y}{y^2+1} = \mathrm{d}x \implies \arctan y = x + C_2\)。
    初值 $y(0)=1 \implies \frac{\pi}{4} = 0 + C_2 \implies C_2=\frac{\pi}{4}$。特解：
    \(\displaystyle \boxed{y = \tan\left(x + \frac{\pi}{4}\right)}\)。

    \item[(6)] \textbf{缺 $y$ 型}。已经在正文例题中给出详尽过程。
    特解为 \(\displaystyle \boxed{y = x^4 + 4x + 1}\)。
\end{enumerate}
\end{solution}

\begin{exercise}[二阶方程的初值问题]
求方程 $yy'' + y'^2 = 1$ 经过点 $(0,1)$ 且在这一点与直线 $x+y=1$ 相切的积分曲线.
\end{exercise}
\begin{solution}
\textbf{思路}：“经过 $(0,1)$”即 $y(0)=1$；“与直线 $x+y=1$（即 $y=-x+1$）相切”即该点处的导数相等，故 $y'(0)=-1$。这就是变相给出了初值条件。

观察原方程 $yy'' + (y')^2 = 1$，左边可以凑成全微分，因此 \(\displaystyle (yy')' = 1,\ yy' = x + C_1\)。
代入初值 $x=0, y=1, y'=-1$，得 $1 \cdot (-1) = 0 + C_1 \implies C_1 = -1$。于是
\(\displaystyle yy' = x - 1 \implies y\mathrm{d}y = (x-1)\mathrm{d}x,\ \frac{1}{2}y^2 = \frac{1}{2}x^2 - x + C_2\)。
代入初值 $x=0, y=1$，得 $\frac{1}{2} = 0 - 0 + C_2 \implies C_2 = \frac{1}{2}$。
方程化为 \(\displaystyle y^2 = x^2 - 2x + 1 = (x-1)^2\)。
开方得 $y = \pm(x-1)$。由于 $y(0)=1$，我们必须取负号侧以满足条件（如果取 $+$, 则 $-1=1$ 不符）。
故满足条件的积分曲线方程为 \(\displaystyle \boxed{y = 1-x}\)。
\end{solution}

\begin{exercise}[可降阶的一阶齐次型综合应用]
求方程 $xy'' = y' \ln \dfrac{y'}{x}$ 的通解.
\end{exercise}
\begin{solution}
\textbf{思路}：方程缺 $y$，首先令 $p=y'$，则 $y''=p'$。
代入化简得 \(\displaystyle xp' = p \ln\left(\frac{p}{x}\right) \implies p' = \frac{p}{x} \ln\left(\frac{p}{x}\right)\)。
我们发现降阶后，它是一个\textbf{一阶齐次方程}！

令 $u = \frac{p}{x}$，则 $p = ux$，$p' = u'x + u$。代入方程得 \(\displaystyle u'x + u = u\ln u \implies x\frac{\mathrm{d}u}{\mathrm{d}x} = u(\ln u - 1),\ \frac{\mathrm{d}u}{u(\ln u - 1)} = \frac{\mathrm{d}x}{x}\)。
两边积分（左边运用凑微分 $\mathrm{d}(\ln u)$）得
\(\displaystyle \int \frac{\mathrm{d}(\ln u)}{\ln u - 1} = \int \frac{\mathrm{d}x}{x} \implies \ln|\ln u - 1| = \ln|x| + \ln C_1\)，
从而 \(\displaystyle \ln u - 1 = C_1 x \implies u = \mathrm{e}^{C_1 x + 1},\qquad y' = x \mathrm{e}^{C_1 x + 1}\)。
最后，通过分部积分求解 $y$：
    \(\displaystyle y = \int x \mathrm{e}^{C_1 x + 1} \mathrm{d}x
    = \mathrm{e} \left( \frac{x}{C_1}\mathrm{e}^{C_1 x} - \frac{1}{C_1^2}\mathrm{e}^{C_1 x} \right) + C_2\)。
此即为原方程的通解。
\end{solution}

\begin{exercise}[缺 $x$ 型二阶方程求解]
求方程 $y''(1+y^2) = 2y y'^2$ 的通解.
\end{exercise}
\begin{solution}
\textbf{思路}：缺 $x$ 型，令 $p=y'$，代入 $y'' = p\frac{\mathrm{d}p}{\mathrm{d}y}$。原方程化为 \(\displaystyle p\frac{\mathrm{d}p}{\mathrm{d}y}(1+y^2) = 2yp^2\)。排除 $p=0$（对应常数解 $y=C$）后，分离变量并积分，得
\(\displaystyle \frac{\mathrm{d}p}{p} = \frac{2y}{1+y^2}\mathrm{d}y,\quad
\ln|p| = \ln(1+y^2) + \ln C_1
\implies p = C_1(1+y^2)\)。
还原变量并再次分离变量，得
\(\displaystyle \frac{\mathrm{d}y}{\mathrm{d}x} = C_1(1+y^2),\quad
\frac{\mathrm{d}y}{1+y^2} = C_1 \mathrm{d}x,\quad
\arctan y = C_1 x + C_2\)。
或写作显式形式：$y = \tan(C_1 x + C_2)$。
\end{solution}

\subsubsection{B组}

\begin{exercise}[建立运动微分方程：追击问题]
设物体 $A$ 从点 $(0,1)$ 出发, 以速度大小为常数 $v$ 沿 $y$ 轴正向运动. 物体 $B$ 从点 $(-1,0)$ 与 $A$ 同时出发, 其速度大小为 $2v$, 方向始终指向 $A$. 求物体 $B$ 的运动微分方程及初始条件.
\end{exercise}
\begin{solution}
\textbf{思路}：
根据题意，$A$ 在时刻 $t$ 的坐标为 $(0, 1+vt)$。
设 $B$ 的轨迹线方程为 $y=y(x)$。在任意时刻 $t$，$B$ 位于坐标点 $(x,y)$。
因为 $B$ 的运动方向始终指向 $A$，因此轨迹曲线在 $(x,y)$ 处的切线必须通过点 $A(0, 1+vt)$。
切线斜率为 \(\displaystyle y' = \frac{(1+vt) - y}{0 - x} = \frac{y - 1 - vt}{x}\)，从中反解得 \(\displaystyle vt = y - x y' - 1\)。等式两边对自变量 $x$ 求导，可得 \(\displaystyle v \frac{\mathrm{d}t}{\mathrm{d}x} = y' - (y' + x y'') = -xy''\)。

另一方面，我们考察 $B$ 的运动速度。$B$ 的速度大小为 $\frac{\mathrm{d}s}{\mathrm{d}t} = 2v$。
由于 $B$ 从 $(-1,0)$ 出发向 $y$ 轴靠近，它的 $x$ 坐标是递增的，故 $\mathrm{d}x > 0$，所以弧长微分公式为 $\frac{\mathrm{d}s}{\mathrm{d}x} = \sqrt{1+(y')^2}$。
应用链式法则得 \(\displaystyle \frac{\mathrm{d}t}{\mathrm{d}x} = \frac{\mathrm{d}t}{\mathrm{d}s} \cdot \frac{\mathrm{d}s}{\mathrm{d}x} = \frac{1}{2v} \sqrt{1+(y')^2}\)。

将两个关于 $\frac{\mathrm{d}t}{\mathrm{d}x}$ 的表达式等同起来，得 \(\displaystyle v \left( \frac{1}{2v} \sqrt{1+(y')^2} \right) = -xy'' \implies xy'' + \frac{1}{2}\sqrt{1+(y')^2} = 0\)。

最后确认初始条件：
$t=0$ 时，$B$ 在 $(-1,0)$，所以 $y(-1) = 0$。
初始时刻 $A$ 在 $(0,1)$，由于 $B$ 指向 $A$，初始斜率 $y'(-1) = \frac{1-0}{0-(-1)} = 1$。
综合即得所求微分方程模型：\(\displaystyle xy'' + \frac{1}{2}\sqrt{1+(y')^2} = 0,\ y(-1) = 0,\ y'(-1) = 1\)。
\end{solution}

\begin{exercise}[积分方程求导转化法]
设第一象限内的曲线 $y=y(x)$ 对应于 $0\leqslant x\leqslant a$ 一段的弧长等于曲边梯形 $D=\{0\leqslant y\leqslant y(x),0\leqslant x\leqslant a\}$ 的面积, 其中 $a>0$ 是任意给定的, $y(0)=1$, 求 $y(x)$.
\end{exercise}
\begin{solution}
\textbf{思路}：将几何条件翻译为积分表达式：
弧长 $s = \int_0^a \sqrt{1+(y')^2}\mathrm{d}x$，面积 $S = \int_0^a y(x)\mathrm{d}x$。
由于对任意 $a>0$ 均有此等式成立，由变上限积分求导定理，直接对变上限积分等式两边对 $x$ 求导，可得 \(\displaystyle \int_0^x \sqrt{1+(y')^2}\mathrm{d}t = \int_0^x y(t)\mathrm{d}t \implies \sqrt{1+(y')^2} = y\)。
两边平方并移项，结合曲线位于第一象限且 $y(0)=1$，积分面积随 $x$ 增加而增加，弧长必在延伸，故 $y > 1$（随后）且斜率为正 $y'>0$，于是
\(\displaystyle (y')^2 = y^2 - 1,\ y' = \sqrt{y^2-1},\ \frac{\mathrm{d}y}{\sqrt{y^2-1}} = \mathrm{d}x\)。
积分并代入初值 $y(0)=1$，得
\(\displaystyle \ln|y+\sqrt{y^2-1}| = x + C,\ C = 0 \implies \operatorname{arccosh}(y) = x \implies y = \cosh x = \frac{\mathrm{e}^x + \mathrm{e}^{-x}}{2}\)。
\end{solution}

\begin{exercise}[物理建模与微分方程的推导]
一条均匀柔软不可拉伸的细弦, 两端悬起, 当载荷沿弧均匀分布时, 求弦的平衡状态时的形状.
\end{exercise}
\begin{solution}
\textbf{思路}：此问题为经典的悬链线模型。
取细弦最低点为坐标原点，切线方向水平。设弧段 $\widehat{OB}$ 的水平拉力恒为 $H$，微元弧长的重力为 $\rho g s$（其中 $s$ 为从原点算起的弧长）。
由平衡方程，绳索在 $B(x,y)$ 处的切向张力 $T$ 必须与重力和水平拉力平衡，即 \(\displaystyle T \cos\theta = H,\ T \sin\theta = \rho g s\)。两式相除，并利用导数的几何意义 $y' = \tan\theta$，得 \(\displaystyle y' = \frac{\rho g}{H} s = \frac{\rho g}{H} \int_0^x \sqrt{1+(y')^2}\mathrm{d}x,\ y'' = \frac{\rho g}{H} \sqrt{1+(y')^2}\)。
这是一个缺 $y$ 型微分方程。令 $p=y'$，$a=\frac{\rho g}{H}$，则
\(\displaystyle \frac{\mathrm{d}p}{\mathrm{d}x} = a\sqrt{1+p^2} \implies \frac{\mathrm{d}p}{\sqrt{1+p^2}} = a\mathrm{d}x,
\qquad
\ln(p+\sqrt{1+p^2}) = ax \implies p = \sinh(ax)\)。
再积分一次并利用 $y(0)=0$ 解出 \(\displaystyle \boxed{y = \frac{1}{a}(\cosh(ax)-1)}\)。
此即著名的双曲余弦方程。
\end{solution}

\begin{exercise}[基于牛顿第二定律降阶求解]
一物体只受地球引力作用, 自无穷高处落下, 求这个物体落向地面的速度.
\end{exercise}
\begin{solution}
\textbf{思路}：此问题利用万有引力定律与动力学方程建模。
设物体质量为 $m$，地球质量为 $M$，物体距地心距离为 $r$。
根据牛顿第二定律，\(\displaystyle m\frac{\mathrm{d}^2r}{\mathrm{d}t^2} = -G\frac{Mm}{r^2},\ v\frac{\mathrm{d}v}{\mathrm{d}r} = -\frac{GM}{r^2}\)。
这是一个缺 $t$ 型二阶方程。令速度 $v = \frac{\mathrm{d}r}{\mathrm{d}t}$，则加速度 $\frac{\mathrm{d}^2r}{\mathrm{d}t^2} = v\frac{\mathrm{d}v}{\mathrm{d}r}$。
分离变量并积分，得 \(\displaystyle \int v\mathrm{d}v = -GM \int r^{-2}\mathrm{d}r \implies \frac{1}{2}v^2 = \frac{GM}{r} + C\)。
由初始条件“自无穷高处落下”，即当 $r \to \infty$ 时，$v = 0$。代入得 $C=0$。
所以 \(\displaystyle v^2 = \frac{2GM}{r}\)。
当物体到达地面时，距地心距离为地球半径 $R$。考虑到地表重力加速度 $g = \frac{GM}{R^2}$，即 $GM = gR^2$。
落向地面的瞬间速度大小为 \(\displaystyle \boxed{v = \sqrt{\frac{2GM}{R}} = \sqrt{2gR}}\)。
（这恰好也是从地球表面发射物体的第二宇宙速度 / 逃逸速度的倒推过程）。
\end{solution}

\begin{exercise}[变加速直线运动的积分]
重量为 $300\ \mathrm{kg}$ 的摩托艇以 $66\ \mathrm{m/s}$ 的初速度直线前进, 如果水的阻力与速度成正比, 且当速度为 $1\ \mathrm{m/s}$ 时, 阻力为 $10\ \mathrm{kg}$. 问经过多少时间艇的速度降为 $8\ \mathrm{m/s}$?
\end{exercise}
\begin{solution}
\textbf{思路}：此处的“重量 300kg”实质为质量 $m = 300 \mathrm{kg}$，阻力 10kg 指的是 $10\ \mathrm{kgf}$（公斤力），转化为牛顿（$\mathrm{N}$）应乘重力加速度 $g$（通常取 $9.8\ \mathrm{m/s^2}$），即阻力 $f = 10g$。
因为阻力与速度成正比，设 $f = kv$。由题意当 $v=1$ 时 $f=10g$，故比例系数 $k=10g$。
根据牛顿第二定律建立动力学方程：\(\displaystyle m\frac{\mathrm{d}v}{\mathrm{d}t} = -kv \implies 300\frac{\mathrm{d}v}{\mathrm{d}t} = -10g v,\ \frac{\mathrm{d}v}{v} = -\frac{g}{30}\mathrm{d}t,\ \ln v = -\frac{g}{30}t + C_1\)。
初始状态 $t=0$ 时，$v=66$，所以 $C_1 = \ln 66$。
得到速度随时间的变化关系：\(\displaystyle \ln\left(\frac{v}{66}\right) = -\frac{g}{30}t \implies t = \frac{30}{g} \ln\left(\frac{66}{v}\right)\)。
当艇的速度降为 $v=8\ \mathrm{m/s}$ 时，代入上式计算所需时间：
\(\displaystyle t = \frac{30}{9.8} \ln\left(\frac{66}{8}\right) \approx 3.06 \times \ln(8.25) \approx 3.06 \times 2.11 = 6.46\ \mathrm{s}\)
（若物理计算中取近似 $g=10\ \mathrm{m/s^2}$，则 $t = 3 \ln(8.25) \approx 6.33\ \mathrm{s}$）。
\end{solution}

% ============================================================
% §5.5 二阶线性微分方程
% ============================================================

\subsection{第五节习题：二阶微分方程}

\begin{exercise}[基础：验证与判定]
\begin{shortexenum}
  \shortitem{1}{验证 $y_1=\mathrm{e}^{2x}$ 和 $y_2=x\mathrm{e}^{2x}$ 都是
        $y''-4y'+4y=0$ 的解，判断它们是否构成基础解系；} &
  \shortitem{2}{设 $y_1,y_2$ 是某二阶齐次线性方程的线性无关解，
        证明 $z_1=y_1+y_2$ 和 $z_2=y_1-y_2$ 也构成该方程的一组基础解系；} \\
  \shortitem{3}{已知 $y''-2xy'+4y=0$ 的一个解为 $y_1=1-2x^2$，求其通解。} &
\end{shortexenum}
\end{exercise}
\begin{solution}
\begin{enumerate}
  \item[(1)] \textbf{解的验证与基础解系判定：}
  \begin{enumerate}
      \item 验证 $y_1 = \mathrm{e}^{2x}$：
      计算其导数：$y_1' = 2\mathrm{e}^{2x}$，$y_1'' = 4\mathrm{e}^{2x}$。代入原方程左端：
      \(\displaystyle y_1'' - 4y_1' + 4y_1 = 4\mathrm{e}^{2x} - 4(2\mathrm{e}^{2x}) + 4\mathrm{e}^{2x} = 4\mathrm{e}^{2x} - 8\mathrm{e}^{2x} + 4\mathrm{e}^{2x} = 0\)
      等式成立，故 $y_1$ 是方程的解。
      \item 验证 $y_2 = x\mathrm{e}^{2x}$：
      利用乘积求导法则：$y_2' = \mathrm{e}^{2x} + 2x\mathrm{e}^{2x} = (1+2x)\mathrm{e}^{2x}$。
      继续求二阶导：$y_2'' = 2\mathrm{e}^{2x} + 2(1+2x)\mathrm{e}^{2x} = (4+4x)\mathrm{e}^{2x}$。代入原方程左端：
      \(\displaystyle y_2'' - 4y_2' + 4y_2 = (4+4x)\mathrm{e}^{2x} - 4(1+2x)\mathrm{e}^{2x} + 4x\mathrm{e}^{2x} = (4 + 4x - 4 - 8x + 4x)\mathrm{e}^{2x} = 0\)
      等式成立，故 $y_2$ 也是方程的解。
      \item 判断基础解系：
      计算两个解的朗斯基行列式，得 \(\displaystyle W(x)=y_1y_2'-y_1'y_2
      =\mathrm{e}^{2x}(1+2x)\mathrm{e}^{2x}-(2\mathrm{e}^{2x})(x\mathrm{e}^{2x})
      =\mathrm{e}^{4x}\)。
      因为对于任意 $x$ 都有 $W(x) = \mathrm{e}^{4x} \neq 0$，所以 $y_1, y_2$ 线性无关，\textbf{它们构成方程的一组基础解系}。
  \end{enumerate}

  \item[(2)] \textbf{证明：}
  \begin{enumerate}
      \item 根据齐次线性微分方程的叠加原理，由于 $y_1, y_2$ 是解，它们的任意线性组合 $z_1 = y_1 + y_2$ 和 $z_2 = y_1 - y_2$ 必然也是该方程的解。
      \item 要证明 $z_1, z_2$ 构成基础解系，只需证明它们线性无关，即计算其朗斯基行列式 $W(z_1, z_2)$ 是否非零：
      \(\displaystyle W(z_1, z_2) = z_1 z_2' - z_1' z_2
      = (y_1+y_2)(y_1'-y_2') - (y_1'+y_2')(y_1-y_2)
      = -2W(y_1, y_2)\)。
      \item 因为已知 $y_1, y_2$ 线性无关，所以 $W(y_1, y_2) \neq 0$。进而推导出 $W(z_1, z_2) = -2 W(y_1, y_2) \neq 0$。
      因此 $z_1, z_2$ 线性无关，\textbf{确实构成该方程的一组基础解系}。
  \end{enumerate}

  \item[(3)] \textbf{求解通解：}
  \begin{enumerate}
      \item 验证特解：令 $y_1 = 1-2x^2$，则 $y_1' = -4x$, $y_1'' = -4$。代入 $y'' - 2xy' + 4y = -4 - 2x(-4x) + 4(1-2x^2) = -4 + 8x^2 + 4 - 8x^2 = 0$，确实为方程的解。
      \item 利用刘维尔（Liouville）公式求第二解。原方程已经是标准形式 $y'' + P(x)y' + Q(x)y = 0$，其中 $P(x) = -2x$。
      \item 计算刘维尔公式中的积分因子：\(\displaystyle \exp\left(-\int P(x)\,\mathrm{d}x\right)=\exp\left(\int 2x\,\mathrm{d}x\right)=\mathrm{e}^{x^2}\)。
      \item 代入降阶公式；由于右侧积分无法用初等函数形式表达，可取 \(\displaystyle y_2(x)=(1-2x^2)I(x)\)，其中 \(\displaystyle I(x)=\int \frac{\mathrm{e}^{x^2}}{(1-2x^2)^2}\,\mathrm{d}x\)。
      因而通解为
      \[
      \mathbf{y=C_1(1-2x^2)+C_2(1-2x^2)\int_{x_0}^x
      \frac{\mathrm{e}^{t^2}}{(1-2t^2)^2}\,\mathrm{d}t}.
      \]
  \end{enumerate}
\end{enumerate}
\end{solution}

\begin{exercise}[基础：常数变易法]
用常数变易法求下列方程的通解：
\begin{shortexenum}
  \shortitem{1}{$y''+y=\cot x$；} &
  \shortitem{2}{$y''-3y'+2y=\dfrac{\mathrm{e}^x}{1+\mathrm{e}^x}$；} \\
  \shortitem{3}{$xy''-y'=x^2\;(x>0)$。} &
\end{shortexenum}
\end{exercise}
\noindent\textbf{本题可直接按二阶线性非齐次方程的标准形式处理。}

先化为标准型 $y''+P(x)y'+Q(x)y=f(x)$，再用
\[
\boxed{y^{*}(x) = -y_{1}(x) \int \frac{y_{2}(x) f(x)}{W(x)} \mathrm{d}x + y_{2}(x) \int \frac{y_{1}(x) f(x)}{W(x)} \mathrm{d}x}
\]
其中 $W(x)=y_1y_2'-y_2y_1'$。
\begin{solution}
\begin{enumerate}
\item[(1)] \textbf{求解 $y''+y=\cot x$}
\begin{enumerate}
\item \textbf{求齐次通解与基础解系}：特征方程 $r^2+1=0$，解得根为 $r=\pm \mathrm{i}$。齐次方程的基础解系为 $y_1 = \sin x$, $y_2 = \cos x$。
\item \textbf{计算朗斯基行列式}：
\(\displaystyle W(x) = \begin{vmatrix} y_1 & y_2 \\ y_1' & y_2' \end{vmatrix} = \begin{vmatrix} \sin x & \cos x \\ \cos x & -\sin x \end{vmatrix} = -1\)
原方程已是标准型，非齐次项 $f(x) = \cot x$。
\item \textbf{套用常数变易法公式求特解}：
代入公式 \(\displaystyle y^{*}(x) = -y_{1} \int \frac{y_{2} f(x)}{W(x)} \mathrm{d}x + y_{2} \int \frac{y_{1} f(x)}{W(x)} \mathrm{d}x\)，得 \(\displaystyle y^* = -\sin x \int \frac{\cos x \cdot \cot x}{-1} \mathrm{d}x + \cos x \int \frac{\sin x \cdot \cot x}{-1} \mathrm{d}x\)。
分别计算两个积分：
第一个积分：
\begin{align*}
\int \cos x \cdot \cot x \,\mathrm{d}x
&= \int \frac{\cos^2 x}{\sin x} \,\mathrm{d}x\\
&= \int \frac{1-\sin^2 x}{\sin x} \,\mathrm{d}x\\
&= \int (\csc x - \sin x)\,\mathrm{d}x\\
&= \ln|\csc x - \cot x| + \cos x.
\end{align*}
第二个积分为 $-\int \sin x \cdot \cot x \mathrm{d}x = -\int \cos x \mathrm{d}x = -\sin x$。
\item \textbf{写出特解与通解}：
将积分结果代回：\(\displaystyle y^* = \sin x (\ln|\csc x - \cot x| + \cos x)+ \cos x (-\sin x)= \sin x \ln|\csc x - \cot x|\)。
最终通解为：$\mathbf{y = C_1 \sin x + C_2 \cos x + \sin x \ln|\csc x - \cot x|}$。
\end{enumerate}
\item[(2)] \textbf{求解 $y''-3y'+2y=\frac{\mathrm{e}^x}{1+\mathrm{e}^x}$}
\begin{enumerate}
\item \textbf{求齐次通解与基础解系}：特征方程 $r^2 - 3r + 2 = 0$，解得 $r_1=1, r_2=2$。基础解系为 $y_1=\mathrm{e}^x, y_2=\mathrm{e}^{2x}$。
\item \textbf{计算朗斯基行列式}：\(\displaystyle W(x) = \begin{vmatrix} \mathrm{e}^x & \mathrm{e}^{2x} \\ \mathrm{e}^x & 2\mathrm{e}^{2x} \end{vmatrix}= \mathrm{e}^{3x}\)。
方程标准，非齐次项 $f(x) = \frac{\mathrm{e}^x}{1+\mathrm{e}^x}$。
    \item \textbf{套用常数变易法公式求特解}：
    \begin{align*}
    y^*
    &= -\mathrm{e}^x
    \int \frac{\mathrm{e}^{2x} \cdot \frac{\mathrm{e}^x}{1+\mathrm{e}^x}}{\mathrm{e}^{3x}} \mathrm{d}x
    + \mathrm{e}^{2x}
    \int \frac{\mathrm{e}^x \cdot \frac{\mathrm{e}^x}{1+\mathrm{e}^x}}{\mathrm{e}^{3x}} \mathrm{d}x\\
    &= -\mathrm{e}^x \int \frac{1}{1+\mathrm{e}^x} \mathrm{d}x
    + \mathrm{e}^{2x} \int \frac{\mathrm{e}^{-x}}{1+\mathrm{e}^x} \mathrm{d}x.
    \end{align*}
    分别计算两个积分：\(\displaystyle \int \frac{1}{1+\mathrm{e}^x}\mathrm{d}x=x-\ln(1+\mathrm{e}^x)\)。
    令 $t=\mathrm e^x$，则 \(\displaystyle \int \frac{\mathrm{e}^{-x}}{1+\mathrm{e}^x}\mathrm{d}x=\int\frac{1}{t^2(1+t)}\,\mathrm{d}t\)，计算得 \(\displaystyle -x-\mathrm e^{-x}+\ln(1+\mathrm e^x)\)。
\item \textbf{写出特解与通解}：
代回 $y^*$ 表达式：
\begin{align*}
y^*
&= -\mathrm{e}^x\bigl(x-\ln(1+\mathrm{e}^x)\bigr)
+\mathrm{e}^{2x}\bigl(-x-\mathrm{e}^{-x}+\ln(1+\mathrm{e}^x)\bigr)\\
&=(\mathrm{e}^x+\mathrm{e}^{2x})\ln(1+\mathrm{e}^x)
-x(\mathrm{e}^x+\mathrm{e}^{2x})-\mathrm{e}^x.
\end{align*}
注意到末尾的 $-\mathrm{e}^x$ 与齐次通解 $C_1\mathrm{e}^x$ 结构重合，可被合并吸收。
最终通解为
\[
y = C_1\mathrm{e}^x + C_2\mathrm{e}^{2x}
+(\mathrm{e}^x+\mathrm{e}^{2x})\ln(1+\mathrm{e}^x)
-x(\mathrm{e}^x+\mathrm{e}^{2x}).
\]
\end{enumerate}
\item[(3)] \textbf{求解 $xy''-y'=x^2\;(x>0)$}
\begin{enumerate}
\item \textbf{化为标准型并提取 $f(x)$}：将原方程两边同除以 $x$，化为首项系数为 1 的标准形式 $y'' - \frac{1}{x}y' = x$。此时非齐次项 $f(x) = x$。
\item \textbf{求齐次通解与基础解系}：解齐次方程 $xy'' - y' = 0$。分离变量得 $\frac{y''}{y'} = \frac{1}{x}$，积分两次得到 $y = \frac{1}{2}C_1 x^2 + C_0$。取基础解系为 $y_1 = 1, y_2 = x^2$。
\item \textbf{计算朗斯基行列式并求特解}：
\(\displaystyle W(x) = \begin{vmatrix} 1 & x^2 \\ 0 & 2x \end{vmatrix} = 2x\)。
    \(\displaystyle y^* = -1 \cdot \int \frac{x^2 \cdot x}{2x} \mathrm{d}x + x^2 \int \frac{1 \cdot x}{2x} \mathrm{d}x
    = -\int \frac{x^2}{2} \mathrm{d}x + x^2 \int \frac{1}{2} \mathrm{d}x
    = -\frac{x^3}{6} + x^2 \left(\frac{x}{2}\right) = \frac{x^3}{3}\)。
\item \textbf{写出通解}：
最终通解为：$\mathbf{y = C_1 + C_2 x^2 + \frac{1}{3}x^3}$。
\end{enumerate}
\end{enumerate}
\end{solution}

\begin{exercise}[提高：叠加原理的应用]
设 $y_1$ 是 $y''+py'+qy=f_1(x)$ 的解，$y_2$ 是 $y''+py'+qy=f_2(x)$ 的解。
问：$y_1+y_2$ 满足什么方程？
若要求 $y''+py'+qy=f_1(x)+f_2(x)$ 的一个特解，
你能否仅通过 $y_1$ 和 $y_2$ 将它 ``拼出来''？
\end{exercise}
\begin{solution}
\textbf{解析：}
\begin{enumerate}
    \item 定义线性微分算子 $L[y] = y'' + py' + qy$。根据题意已知：$L[y_1] = f_1(x)$ 且 $L[y_2] = f_2(x)$。
    \item 考察 $y_1 + y_2$ 作用于算子 $L$ 的结果。根据导数和运算法则的线性性质，将其展开并代入已知条件：
    \(\displaystyle L[y_1 + y_2] = (y_1 + y_2)'' + p(y_1 + y_2)' + q(y_1 + y_2) = (y_1'' + py_1' + qy_1) + (y_2'' + py_2' + qy_2) = L[y_1] + L[y_2] = f_1(x) + f_2(x)\)。
    这说明 \textbf{$y_1 + y_2$ 满足的方程正是 $y''+py'+qy=f_1(x)+f_2(x)$}。
    \item \textbf{结论}：若要求 $y''+py'+qy=f_1(x)+f_2(x)$ 的一个特解，我们完全可以仅通过将 $y_1$ 和 $y_2$ “拼出来”（直接相加）即可。这就是线性常微分方程中极为重要的\textbf{非齐次叠加原理}。
\end{enumerate}
\end{solution}

\begin{exercise}[刘维尔公式的推广]
将刘维尔公式推广到三阶情形：设 $y_1$ 是
$y'''+p(x)y''+q(x)y'+r(x)y=0$ 的一个非零特解，
试推导求另两个线性无关解的降阶公式。
（提示：依次令 $y_2=u\,y_1$, $y_3=v\,y_1$，逐级降阶）
\end{exercise}
\begin{solution}
\textbf{解析推导：}
已知 $y_1$ 是三阶齐次方程 $L[y] = y''' + p(x)y'' + q(x)y' + r(x)y = 0$ 的非零特解。我们通过变量代换将其逐级降阶：
\begin{enumerate}
    \item \textbf{第一次降阶（从三阶降为二阶）}：
    令 $y = u y_1$，利用乘积求导法则计算各阶导数：\(\displaystyle y' = u'y_1 + u y_1',\quad
    y'' = u''y_1 + 2u'y_1' + u y_1'',\quad
    y''' = u'''y_1 + 3u''y_1' + 3u'y_1'' + u y_1'''\)。
    将上述导数代入原方程，并按 $u$ 的各阶导数提取公因式：\(\displaystyle y_1 u''' + (3y_1' + py_1) u'' + (3y_1'' + 2py_1' + qy_1) u' + u\underbrace{(y_1''' + py_1'' + qy_1' + ry_1)}_{L[y_1]=0} = 0\)。
    因为 $y_1$ 是原方程的解，故含有 $u$ 的项恰好消去。
    令 $w = u'$，则 $w' = u'', w'' = u'''$。代入上式并在两边同除以 $y_1$（因为 $y_1 \neq 0$），得到关于 $w$ 的标准二阶齐次微分方程 \(\displaystyle w'' + \left(3\frac{y_1'}{y_1} + p(x)\right)w' + \left(3\frac{y_1''}{y_1} + 2p(x)\frac{y_1'}{y_1} + q(x)\right)w = 0\)，记为 (*)。

    \item \textbf{确立第一级降阶的解关系}：
    假设我们能找到二阶方程 (*) 的一个特解 $w_1$。由于 $w_1 = u'$，对 $w_1$ 积分即可得到 $u$，从而得到原三阶方程的第二个线性无关解：
    \(\displaystyle y_2 = y_1 \int w_1\,\mathrm{d}x\)。
    反过来，如果已知第二个解 $y_2$，我们也就知道了 $w_1$ 的表达式：
    \(\displaystyle w_1 = u' = \left(\frac{y_2}{y_1}\right)' = \frac{y_1 y_2' - y_1' y_2}{y_1^2} = \frac{W(y_1, y_2)}{y_1^2}\)。
    其中 $W(y_1, y_2)$ 是 $y_1, y_2$ 的朗斯基行列式。

    \item \textbf{第二次降阶（利用二阶刘维尔公式求 $w_2$）}：
    现在问题转化为：已知二阶方程 (*) 的一个特解 $w_1$，如何求其第二个线性无关解 $w_2$？
    其中 $P(x)=3\frac{y_1'}{y_1}+p(x)$。直接套用二阶的刘维尔公式（降阶法）：\(\displaystyle w_2= w_1 \int \frac{1}{w_1^2}\exp\left(-\int P(x)\,\mathrm{d}x\right) \mathrm{d}x\)。
    指数部分为
    \[
    -\int P(x)\,\mathrm{d}x
    =-\int \left(3\frac{y_1'}{y_1} + p(x)\right)\mathrm{d}x
    =-3\ln|y_1|-\int p(x)\,\mathrm{d}x,
    \]
    所以积分因子为 \(\displaystyle \exp\left(-\int P(x)\,\mathrm{d}x\right)= \frac{1}{y_1^3}\exp\left(-\int p(x)\,\mathrm{d}x\right)\)。

    \item \textbf{化简并得到 $y_3$ 的最终公式}：
    将积分因子和 $w_1 = \frac{W(y_1, y_2)}{y_1^2}$ 代入 $w_2$ 的公式中。注意到 $\frac{1}{w_1^2} = \frac{y_1^4}{W(y_1, y_2)^2}$，于是
    \[
    w_2
    = w_1 \int \frac{y_1^4}{W(y_1, y_2)^2}\cdot \frac{1}{y_1^3}
    \exp\left(-\int p(x)\,\mathrm{d}x\right)\mathrm{d}x,
    \]
    即
    \[
    w_2=\left(\frac{y_2}{y_1}\right)'
    \int \frac{y_1}{W(y_1, y_2)^2}
    \exp\left(-\int p(x)\,\mathrm{d}x\right)\mathrm{d}x.
    \]

    因为题设提示第三个解为 $y_3 = v y_1$，且同理 $w_2 = v'$。我们对 $w_2$ 积分得到 $v$，再乘上 $y_1$，即得到第三个线性无关解的最终降阶积分公式：
    \[
    \boxed{ y_3 = y_1 \int \left[ \left(\frac{y_2}{y_1}\right)' \int \frac{y_1}{W(y_1, y_2)^2} \exp\left(-\int p(x)\,\mathrm{d}x\right) \mathrm{d}x \right] \mathrm{d}x }.
    \]
    （注：此公式表明，只要已知原方程的两个线性无关解 $y_1, y_2$，便可以通过积分运算求出第三个线性无关解 $y_3$，这可以视为刘维尔公式在三阶情形下的推广。）
\end{enumerate}
\end{solution}

\subsubsection{A组}

\begin{exercise}[基础：概念理解]
回答下列问题
\begin{shortexenum}
    \shortitem{1}{“算子”的含义是什么？算子 $L$ 把什么函数变成什么函数？} &
    \shortitem{2}{$L$ 的线性运算性质是什么？} \\
    \shortitem{3}{什么叫两个函数 $y_1(x)$ 与 $y_2(x)$ 线性相关、线性无关？方程 $L[y]=0$ 的基础解系是什么？} &
    \shortitem{4}{方程 $L[y]=0$ 的通解结构是怎样的？何谓线性叠加原理 1？} \\
    \shortitem{5}{方程 $L[y]=f(x)$ 的通解结构是怎样的？何谓线性叠加原理 2？} &
\end{shortexenum}
\end{exercise}
\begin{solution}
\begin{enumerate}
    \item[(1)] 算子是一种“输入一个对象，输出另一个对象”的映射。
    在这里，二阶线性微分算子 $L$ 的输入是足够可导的函数 $y(x)$，输出是
    \(\displaystyle L[y]=y'' + P(x)y' + Q(x)y\)。
    \item[(2)] $L$ 的线性性质体现在两点：
    \(\displaystyle L[cy]=cL[y],\qquad L[y_1+y_2]=L[y_1]+L[y_2]\)。因而综合起来就有 \(\displaystyle L[c_1y_1+c_2y_2]=c_1L[y_1]+c_2L[y_2]\)。
    \item[(3)] 若存在不全为零的常数 $c_1,c_2$ 使 \(\displaystyle c_1y_1(x)+c_2y_2(x)\equiv0\)，
    则称 $y_1,y_2$ 线性相关；若只能由 $c_1=c_2=0$ 才成立，则称线性无关。
    对二阶齐次方程 $L[y]=0$，任取两个线性无关解组成的集合 $\{y_1,y_2\}$ 就叫基础解系。
    \item[(4)] 齐次方程 $L[y]=0$ 的通解是基础解系的全体线性组合：\(\displaystyle y=c_1y_1+c_2y_2\)。
    线性叠加原理 1 指的是：齐次方程的两个解做任意线性组合后，仍然还是齐次方程的解。
    \item[(5)] 非齐次方程 $L[y]=f(x)$ 的通解结构是 \(\displaystyle y=c_1y_1+c_2y_2+y^*\)，
    其中 $y^*$ 是任意一个特解。
    线性叠加原理 2 指的是：若 $y_1^*$、$y_2^*$ 分别对应右端 $f_1,f_2$，则它们的和 $y_1^*+y_2^*$ 对应右端 $f_1+f_2$。
\end{enumerate}
\end{solution}

\begin{exercise}[解的验证]
验证 $y=2\cos x+3\sin x$ 是方程 $y''+y=0$ 的解.
\end{exercise}
\begin{solution}
解析验证：
计算该函数的一阶和二阶导数，得 $y' = -2\sin x + 3\cos x,\ y'' = -2\cos x - 3\sin x$。代入原微分方程左端，
\(\displaystyle y'' + y = (-2\cos x - 3\sin x) + (2\cos x + 3\sin x) = 0\)。
因为等式恒成立，所以验证了 $y=2\cos x+3\sin x$ 确实是方程 $y''+y=0$ 的解。
\end{solution}

\begin{exercise}[含参方程的求解]
设 $y=A\cos\alpha x$ 是方程 $y''+5y=0$ 满足条件 $y'(1)=3$ 的解, 试求 $A$ 与 $\alpha$ 的值.
\end{exercise}
\begin{solution}
解析推导：
首先保证 $y=A\cos\alpha x$ 满足微分方程。求其两阶导数并代入方程 $y''+5y=0$ 中：
$y' = -A\alpha\sin\alpha x,\ y'' = -A\alpha^2\cos\alpha x$，故
\(\displaystyle -A\alpha^2\cos\alpha x + 5A\cos\alpha x = 0 \implies A(5-\alpha^2)\cos\alpha x = 0 \implies \alpha = \pm\sqrt{5}\)。
为了使其对所有 $x$ 成立（且 $A \neq 0$ 否则 $y \equiv 0$ 无法满足 $y'(1)=3$），必须满足上式。
接着利用初始条件 $y'(1)=3$：
代入 $x=1$ 到 $y'$ 表达式中，得 $-A\alpha\sin\alpha = 3 \implies -A(\pm\sqrt{5})\sin(\pm\sqrt{5}) = 3$。
因为正弦函数是奇函数 $\sin(-\sqrt{5}) = -\sin(\sqrt{5})$，无论 $\alpha$ 取正负，乘积 $-A\alpha\sin\alpha$ 都等于 $-A\sqrt{5}\sin\sqrt{5}$。
因此 $-A\sqrt{5}\sin\sqrt{5} = 3 \implies A = -\frac{3}{\sqrt{5}\sin\sqrt{5}}$。
综上所述，参数的值为 $\alpha = \pm\sqrt{5}$ 且 $A = -\frac{3}{\sqrt{5}\sin\sqrt{5}}$。
\end{solution}

\begin{exercise}[边值问题求解]
求 $A,B$ 及 $\omega$ 的值, 使 $y=A\cos\omega x+B\sin\omega x$ 成为边值问题
\(\displaystyle y''+16y=0,\quad y(0)=2,\quad y\left(\frac{\pi}{8}\right)=3\)
的解.
\end{exercise}
\begin{solution}
解析推导：
步骤一：利用微分方程确定 $\omega$。
求函数二阶导数并代入方程 $y''+16y=0$，有 $y'' = -A\omega^2\cos\omega x - B\omega^2\sin\omega x = -\omega^2 y$，故
\(\displaystyle (-\omega^2 + 16)y = 0 \implies \omega^2 = 16 \implies \omega = \pm 4\)。
由于余弦是偶函数，正弦是奇函数，正负号的作用可以被常数 $A, B$ 的重新分配所吸收，一般我们取正频率 $\omega = 4$。
步骤二：利用边界条件确定 $A$ 和 $B$。
将 $\omega = 4$ 代入原式得 $y = A\cos 4x + B\sin 4x$。
利用两个边值条件，\(\displaystyle A\cos(0) + B\sin(0) = 2 \implies A = 2\)，以及 \(\displaystyle A\cos\left(4 \cdot \frac{\pi}{8}\right) + B\sin\left(4 \cdot \frac{\pi}{8}\right) = 3 \implies A\cos\left(\frac{\pi}{2}\right) + B\sin\left(\frac{\pi}{2}\right) = 3\)。
即 $A(0) + B(1) = 3 \implies B = 3$。
若取 $\omega = -4$，同理可得 $A=2$，但由于 $\sin(-4 \cdot \pi/8) = -1$，则会得出 $B=-3$。
因此解为：$\omega = \pm 4, A = 2, B = \pm 3$ （符号同号对应）。
\end{solution}

\begin{exercise}[线性算子的计算]
设 $L[y]=y''-3xy'+3y$, 试计算（1）$L[\mathrm{e}^x]$;~~（2） $L[\cos\sqrt{3}x]$;~~（3） $L[x^2+3x]$.
\end{exercise}
\begin{solution}
\begin{enumerate}
    \item[(1)] 对于 $y = \mathrm{e}^x$，有 $y' = \mathrm{e}^x$，$y'' = \mathrm{e}^x$。代入算子：
    $L[\mathrm{e}^x] = \mathrm{e}^x - 3x\mathrm{e}^x + 3\mathrm{e}^x = (4 - 3x)\mathrm{e}^x$。
    \item[(2)] 对于 $y = \cos\sqrt{3}x$，有 $y' = -\sqrt{3}\sin\sqrt{3}x$，$y'' = -3\cos\sqrt{3}x$。代入算子：
    $L[\cos\sqrt{3}x] = -3\cos\sqrt{3}x - 3x(-\sqrt{3}\sin\sqrt{3}x) + 3\cos\sqrt{3}x = 3\sqrt{3}x\sin\sqrt{3}x$。
    \item[(3)] 对于 $y = x^2+3x$，有 $y' = 2x+3$，$y'' = 2$。代入算子：
    $L[x^2+3x] = 2 - 3x(2x+3) + 3(x^2+3x) = 2 - 6x^2 - 9x + 3x^2 + 9x = 2 - 3x^2$。
\end{enumerate}
\end{solution}

\begin{exercise}[函数组的线性相关性]
下列函数组在其定义区间内哪些是线性无关的?\\(1)$x,x^2$;(2)$\cos2x,\sin2x$;(3)$x,2x$;(4)$\mathrm{e}^{\alpha x},\mathrm{e}^{\beta x}\ (\alpha\neq\beta)$;(5)$\sin2x,\cos x\sin x$.
\end{exercise}
\begin{solution}
判断方法：若两函数成比例（即比值为常数），则线性相关；否则线性无关。我们也可以用朗斯基行列式辅助判断。
\begin{enumerate}
    \item[(1)] 线性无关。$x$ 与 $x^2$ 的比值为 $x$，不是常数。其 $W(x) = x(2x) - 1(x^2) = x^2 \neq 0$（在避开原点的区间上）。
    \item[(2)] 线性无关。两者的比值为 $\tan 2x$，不是常数。
    \item[(3)] 线性相关。因为 $2x = 2 \cdot x$，其中一个是另一个的常数倍。
    \item[(4)] 线性无关。比值 $\mathrm{e}^{(\alpha-\beta)x}$ 不是常数（因 $\alpha \neq \beta$）。
    \item[(5)] 线性相关。根据倍角公式 $\sin 2x = 2\sin x\cos x$，可知 $\sin 2x = 2 \cdot (\cos x\sin x)$，成比例关系。
\end{enumerate}
\end{solution}

\begin{exercise}[非齐次叠加原理与猜测法]
求方程 $y''+y=3\mathrm{e}^{-x}+2x$ 的通解.
\end{exercise}
\begin{solution}
\begin{enumerate}
    \item 求对应的齐次方程 $y''+y=0$ 的通解。
特征方程为 $r^2+1=0 \implies r=\pm i$。
基础解系为 $\sin x, \cos x$，所以 $y_h = C_1\sin x + C_2\cos x$。
\item 利用非齐次叠加原理，分别猜测右侧两部分的特解。
\begin{itemize}
\item 对部分一：求方程 $y''+y = 3\mathrm{e}^{-x}$ 的特解 $y_1^*$。
由于 $-1$ 不是特征根，合理猜测 $y_1^* = A\mathrm{e}^{-x}$。
代入方程：$A\mathrm{e}^{-x} + A\mathrm{e}^{-x} = 3\mathrm{e}^{-x} \implies 2A = 3 \implies A = \frac{3}{2}$。
故 $y_1^* = \frac{3}{2}\mathrm{e}^{-x}$。
\item 对部分二：求方程 $y''+y = 2x$ 的特解 $y_2^*$。
右端为一阶多项式，合理猜测 $y_2^* = Bx + C$。
代入方程：$0 + (Bx + C) = 2x \implies B=2, C=0$。
故 $y_2^* = 2x$。
\end{itemize}
\item 组合通解。
由叠加原理 2，原非齐次方程的一个特解为 $y^* = y_1^* + y_2^* = \frac{3}{2}\mathrm{e}^{-x} + 2x$。
原方程的通解为
\(\displaystyle \boxed{y = y_h + y^* = C_1\sin x + C_2\cos x + \frac{3}{2}\mathrm{e}^{-x} + 2x}\)。
\end{enumerate}
\end{solution}

\begin{exercise}[解的验证与通解构造]
验证 $y_1=\mathrm{e}^{x^2}$ 和 $y_2=x\mathrm{e}^{x^2}$ 都是方程 $y''-4xy'+(4x^2-2)y=0$ 的解, 并写出该方程的通解.
\end{exercise}
\begin{solution}
计算 $y_1$ 的导数：$y_1' = 2x\mathrm{e}^{x^2}$，$y_1'' = 2\mathrm{e}^{x^2} + 4x^2\mathrm{e}^{x^2}$。
代入原方程左侧：
\(\displaystyle (2\mathrm{e}^{x^2} + 4x^2\mathrm{e}^{x^2}) - 4x(2x\mathrm{e}^{x^2}) + (4x^2-2)\mathrm{e}^{x^2} = \mathrm{e}^{x^2}(2 + 4x^2 - 8x^2 + 4x^2 - 2) = 0\)。
所以 $y_1$ 是方程的解。

计算 $y_2$ 的导数：$y_2' = \mathrm{e}^{x^2} + 2x^2\mathrm{e}^{x^2}$，$y_2'' = 2x\mathrm{e}^{x^2} + 4x\mathrm{e}^{x^2} + 4x^3\mathrm{e}^{x^2} = 6x\mathrm{e}^{x^2} + 4x^3\mathrm{e}^{x^2}$。
代入原方程左侧：
\(\displaystyle (6x\mathrm{e}^{x^2} + 4x^3\mathrm{e}^{x^2}) - 4x(\mathrm{e}^{x^2} + 2x^2\mathrm{e}^{x^2}) + (4x^2-2)(x\mathrm{e}^{x^2})
= \mathrm{e}^{x^2}(6x + 4x^3 - 4x - 8x^3 + 4x^3 - 2x) = 0\)。
所以 $y_2$ 也是方程的解。

验证线性无关性：
其朗斯基行列式 $W(x) = y_1 y_2' - y_1' y_2 = \mathrm{e}^{x^2}(\mathrm{e}^{x^2}+2x^2\mathrm{e}^{x^2}) - 2x\mathrm{e}^{x^2}(x\mathrm{e}^{x^2}) = \mathrm{e}^{2x^2} \neq 0$。
因为两者线性无关，它们构成了该二阶齐次方程的基础解系。
故方程的通解为 \(\displaystyle \boxed{y = C_1\mathrm{e}^{x^2} + C_2 x\mathrm{e}^{x^2}}\)。
\end{solution}

\begin{exercise}[积分算子的线性性证明]
试证明: 由 $T[y]=\int_a^x t^2 y(t)\mathrm{d}t$ 定义的算子 $T$ 是线性的, 即要证 \(\displaystyle T[Cy] = CT[y],\quad T[y_1+y_2] = T[y_1]+T[y_2]\)。
\end{exercise}
\begin{solution}
（1）验证齐次性：$T[Cy] = \int_a^x t^2 (C y(t)) \mathrm{d}t = C \int_a^x t^2 y(t) \mathrm{d}t = C T[y]$。
（2）验证可加性：$T[y_1+y_2] = \int_a^x t^2 (y_1(t) + y_2(t)) \mathrm{d}t = \int_a^x \left(t^2 y_1(t) + t^2 y_2(t)\right) \mathrm{d}t = \int_a^x t^2 y_1(t) \mathrm{d}t + \int_a^x t^2 y_2(t) \mathrm{d}t = T[y_1] + T[y_2]$。
由此证明，积分算子 $T$ 完全满足线性算子的定义。
\end{solution}

\subsubsection{B组}

\begin{exercise}[物理建模：LRC电路微分方程]
设有一个由电阻 $R$、自感 $L$、电容 $C$ 和电源 $E$ 串联组成的电路, 其中 $R$、$L$ 及 $C$ 为常数, 电源电动势是时间 $t$ 的函数: $E=E_m\sin\omega t$, 这里 $E_m$ 及 $\omega$ 也是常数，设电路中的电流为 $I(t)$, 电容器极板上的电量为 $Q(t)$, 两极板间的电压为 $V_C$, 自感电动势为 $E_L$, 由电学知 \(\displaystyle I = \frac{\mathrm{d}Q}{\mathrm{d}t},\qquad
V_C = \frac{Q}{C},\qquad
E_L = -L\frac{\mathrm{d}I}{\mathrm{d}t}\)，
试根据回路电压定律, 导出串联电路的振荡方程(即关于 $V_C$ 的微分方程).
\end{exercise}
\begin{solution}
根据基尔霍夫回路电压定律 (KVL)，闭合回路中的各部分电压降之和等于电源提供的电动势，即 $V_R + (-E_L) + V_C = E(t)$，也就是 \(\displaystyle IR + L\frac{\mathrm{d}I}{\mathrm{d}t} + V_C = E_m\sin\omega t\)。
我们的目标是推导出一个全部由待求函数 $V_C$ 及其导数组成的方程。由定义 $Q = C V_C$ 可知，
\[
I = \frac{\mathrm{d}Q}{\mathrm{d}t}
= \frac{\mathrm{d}}{\mathrm{d}t}(C V_C)
= C\frac{\mathrm{d}V_C}{\mathrm{d}t},
\qquad
\frac{\mathrm{d}I}{\mathrm{d}t}
= \frac{\mathrm{d}}{\mathrm{d}t}\left(C\frac{\mathrm{d}V_C}{\mathrm{d}t}\right)
= C\frac{\mathrm{d}^2V_C}{\mathrm{d}t^2}.
\]
将上面得到的 $I$ 和 $\frac{\mathrm{d}I}{\mathrm{d}t}$ 代入到回路电压方程中，并将方程两边同时除以 $LC$，按标准形式重排：
\[
\boxed{\frac{\mathrm{d}^2V_C}{\mathrm{d}t^2}
+ \frac{R}{L}\frac{\mathrm{d}V_C}{\mathrm{d}t}
+ \frac{1}{LC}V_C
= \frac{E_m}{LC}\sin\omega t}.
\]
这就是熟知的串联二阶 LRC 强迫振荡电路方程。
\end{solution}

\begin{exercise}[物理建模：弹簧振子动力学]
设两个相同的重物一起悬挂在弹簧的一端, 如果其中一个重物坠下, 试求另一个重物的运动方程, 并用合理猜测法求出方程的通解及特解.
\end{exercise}
\begin{solution}
\begin{enumerate}
    \item 解析推导：设单个重物的质量为 $m$，弹簧的劲度系数为 $k$。
起初挂着 $2m$ 的重物，处于静止状态。此时弹簧的伸长量（相对于原长）为 $y_{initial} = \frac{2mg}{k}$。
在 $t=0$ 时刻，其中一个重物坠落，系统只剩下质量为 $m$ 的物体。
对于剩下的这一个重物 $m$，它的新的静力平衡位置应该是 $y_{eq} = \frac{mg}{k}$。
我们以“这个新的平衡位置”为坐标原点（取向下为正，坐标为 $y$），则系统的微分方程直接写为无阻尼自由振动标准形式 $m y'' + ky = 0 \implies y'' + \frac{k}{m}y = 0$。
\item 求通解（合理猜测法）：猜测解的形式为 $y = \mathrm{e}^{rt}$，代入方程得特征方程 $r^2 + \frac{k}{m} = 0$，解得纯虚根 $r = \pm i\sqrt{\frac{k}{m}}$。
由欧拉公式转化，一般实数通解形式为
\(\displaystyle y(t) = C_1\cos\left(\sqrt{\frac{k}{m}}t\right) + C_2\sin\left(\sqrt{\frac{k}{m}}t\right)\)。
    \item 求特解：在 $t=0$ 的瞬间，物体的位置还处于原来挂两个重物时的位置，即在新的平衡位置下方 $\frac{mg}{k}$ 处。并且坠落瞬间没有初速度，因此 $y(0) = \frac{2mg}{k} - \frac{mg}{k} = \frac{mg}{k}, \ y'(0) = 0$。
代入通解：
由 $y(0) = C_1 = \frac{mg}{k}$。
由 $y'(t) = -C_1\sqrt{\frac{k}{m}}\sin\left(\sqrt{\frac{k}{m}}t\right) + C_2\sqrt{\frac{k}{m}}\cos\left(\sqrt{\frac{k}{m}}t\right)$，代入 $t=0$ 得 $y'(0) = C_2\sqrt{\frac{k}{m}} = 0 \implies C_2=0$。
于是特解为 \(\displaystyle \boxed{y(t) = \frac{mg}{k}\cos\left(\sqrt{\frac{k}{m}}t\right)}\)。
\end{enumerate}
\end{solution}

\begin{exercise}[根据基础解系反求微分方程]
试构造线性齐次微分方程, 已知它的基础解系如下:\\(1)$y_1=\sin x,y_2=\cos x$;(2)$y_1=x,y_2=x^2$.
\end{exercise}
\begin{solution}
我们要反向利用方程解的性质，从解构造回微分方程。
\begin{enumerate}
    \item[(1)] 已知 $y_1=\sin x, y_2=\cos x$。
    它们都是三角函数，求两阶导数极容易发现循环性质：$y'' = -\sin x = -y$。
    移项后可得所求的二阶常系数线性方程：
    \(\displaystyle \boxed{y'' + y = 0}\)。
    \item[(2)] 已知 $y_1=x, y_2=x^2$。通解形式为 $y = C_1 x + C_2 x^2$。
    我们通过连续求导来消去任意常数 $C_1, C_2$。
    \(\displaystyle y' = C_1 + 2C_2 x\)。
    为了消去 $C_1$，将第一式乘以 $x$，与 $y$ 的表达式做运算，可得 \(\displaystyle xy' = C_1 x + 2C_2 x^2\)。
    用 $xy'$ 减去 $y$，再对 $y$ 再次求导，得 \(\displaystyle xy' - y = (C_1 x + 2C_2 x^2) - (C_1 x + C_2 x^2) = C_2 x^2,\qquad y'' = 2C_2\)。
    因而 \(\displaystyle \frac{1}{2}x^2 y'' = C_2 x^2\)，于是
    \(\displaystyle \boxed{x^2 y'' - 2x y' + 2y = 0}\)。
\end{enumerate}
\end{solution}

\begin{exercise}[刘维尔公式应用]
已知 $y_1=\mathrm{e}^x$ 是线性齐次方程 $(2x-1)y''-(2x+1)y'+2y=0$ 的一个解, 求此方程的通解.
\end{exercise}
\begin{solution}
这是求已知一解求第二解的典型问题。先化为首一标准形式
\(\displaystyle y'' + P(x)y' + Q(x)y = 0,\qquad P(x) = -\frac{2x+1}{2x-1}\)。
原方程两端同除以 $(2x-1)$ 后，计算得 \(\displaystyle \int P(x)\mathrm{d}x=-x-\ln|2x-1|\)，故 \(\displaystyle \exp\left(-\int P(x)\mathrm{d}x\right)=\mathrm{e}^x(2x-1)\)。
利用刘维尔公式求第二解：\(\displaystyle y_2 = y_1 \int \frac{\exp\left(-\int P(x)\mathrm{d}x\right)}{y_1^2} \mathrm{d}x
= \mathrm{e}^x \int \mathrm{e}^{-x}(2x-1) \mathrm{d}x
= \mathrm{e}^x \bigl[-\mathrm{e}^{-x}(2x+1)\bigr]
= -(2x+1)\)。
为书写简便，且由于是在求解基础解系可以忽略常数倍，我们直接取 $y_2 = 2x+1$ 即可。最终的齐次方程通解为
\(\displaystyle \boxed{y = C_1\mathrm{e}^x + C_2(2x+1)}\)。
\end{solution}

\begin{exercise}[常数变易法的综合应用]
已知 $y_1=x$ 是线性齐次方程 $x^2y''-2xy'+2y=0$ 的一个解, 求线性非齐次方程 $x^2y''-2xy'+2y=2x^3$ 的通解.
\end{exercise}
\begin{solution}
本题分为两步走：先利用刘维尔公式找出对应的齐次方程的另一个基础解，再使用常数变易法求出非齐次特解。
\begin{enumerate}
\item 求齐次方程基础解系。
将 $x^2y''-2xy'+2y=0$ 化为标准形 $y'' - \frac{2}{x}y' + \frac{2}{x^2}y = 0$。$P(x) = -\frac{2}{x}$。
利用刘维尔公式计算另一解 $y_2$：
\(\displaystyle \exp\left(-\int P(x)\mathrm{d}x\right) = \exp\left(\int \frac{2}{x}\mathrm{d}x\right) = \mathrm{e}^{2\ln x} = x^2,\qquad
y_2 = x \int \frac{x^2}{x^2} \mathrm{d}x = x^2\)。
此时齐次方程的基础解系为 $y_1=x, y_2=x^2$，其朗斯基行列式 $W(x) = x(2x) - 1(x^2) = x^2$。
\item 用常数变易法求非齐次方程特解。
将非齐次方程同样化为标准形
\(\displaystyle y'' - \frac{2}{x}y' + \frac{2}{x^2}y = 2x\)。
此时自由项为 $f(x) = 2x$。设特解 $y^* = c_1(x)x + c_2(x)x^2$，则
\(\displaystyle c_1' x + c_2' x^2 = 0,\qquad c_1' + 2xc_2' = 2x
\implies c_2' = 2,\qquad c_1' = -2x\)。
积分得 $c_2 = 2x,\ c_1 = -x^2$，从而 $y^* = (-x^2)(x) + (2x)(x^2) = x^3$，故
\(\displaystyle \boxed{y = C_1 x + C_2 x^2 + x^3}\)。
\end{enumerate}
\end{solution}

\subsection{第六节习题：二阶常系数微分方程}

\subsubsection{A组}

\begin{exercise}[基本概念]
回答下列问题:
\begin{shortexenum}
    \shortitem{1}{方程 $L[y]=y''+py'+qy=0$ 的特征方程是什么?} &
    \shortitem{2}{$k$ 个函数线性无关是怎样定义的?}
\end{shortexenum}
\end{exercise}
\begin{solution}
\begin{enumerate}
    \item[(1)] \textbf{特征方程的推导}：利用欧拉待定指数法，假设齐次线性微分方程有形如 $y=\mathrm{e}^{\lambda x}$ 的解。
    对其求导可得 $y' = \lambda \mathrm{e}^{\lambda x}$， $y'' = \lambda^2 \mathrm{e}^{\lambda x}$。将其代入原方程 $y''+py'+qy=0$ 中，得到 \(\displaystyle (\lambda^2 + p\lambda + q)\mathrm{e}^{\lambda x} = 0 \implies \lambda^2 + p\lambda + q = 0\)。由于指数函数 $\mathrm{e}^{\lambda x}$ 恒不为零，消去该因子后得到关于常数 $\lambda$ 的一元二次代数方程，该代数方程即为原微分方程的\textbf{特征方程}。

    \item[(2)] \textbf{线性无关}：对定义在区间 $I$ 上的 $k$ 个函数 $y_1(x), y_2(x), \dots, y_k(x)$，若存在一组常数 $c_1, c_2, \dots, c_k$，使得恒等式 \(\displaystyle c_1y_1(x) + c_2y_2(x) + \dots + c_ky_k(x) \equiv 0 \quad (\forall x \in I)\) 成立，且这组常数中\textbf{至少有一个不为零}，则称这 $k$ 个函数\textbf{线性相关}；若\textbf{只有当} $c_1 = c_2 = \dots = c_k = 0$ 时，上述恒等式才成立，则称它们\textbf{线性无关}。
\end{enumerate}
\end{solution}

\begin{exercise}[齐次方程通解]
求下列微分方程的通解:\\
\begin{tabular}{lll}
  (1) $y''-6y'+9y=0$; &
  (2) $4y''-12y'+9y=0$; &
  (3) $y''+6y'+13y=0$; \\
  (4) $y''-2y'+y=0$; &
  (5) $y^{(4)}-y=0$; &
  (6) $y'''-6y''+3y'+10y=0$.
\end{tabular}
\end{exercise}
\begin{solution}
\begin{enumerate}
    \item[(1)] \textbf{求解思路}：写出特征方程 $\lambda^2 - 6\lambda + 9 = 0$。
    因式分解得 $(\lambda-3)^2 = 0$，解得二重实根 $\lambda_{1,2} = 3$。根据二重实根的通解公式，通解为 \(\displaystyle \mathbf{y = (C_1 + C_2 x)\mathrm{e}^{3x}}\)。

    \item[(2)] \textbf{求解思路}：写出特征方程 $4\lambda^2 - 12\lambda + 9 = 0$。
    完全平方式配方得 $(2\lambda-3)^2 = 0$，解得二重实根 $\lambda_{1,2} = \frac{3}{2}$。因此，通解为 \(\displaystyle \mathbf{y = (C_1 + C_2 x)\mathrm{e}^{\frac{3}{2}x}}\)。

    \item[(3)] \textbf{求解思路}：写出特征方程 $\lambda^2 + 6\lambda + 13 = 0$。
    利用求根公式 $\lambda = \frac{-6 \pm \sqrt{36-52}}{2} = \frac{-6 \pm \sqrt{-16}}{2} = -3 \pm 2i$，得到一对共轭复根 $\alpha = -3, \beta = 2$。因此，通解为 \(\displaystyle \mathbf{y = \mathrm{e}^{-3x}(C_1\cos 2x + C_2\sin 2x)}\)。

    \item[(4)] \textbf{求解思路}：写出特征方程 $\lambda^2 - 2\lambda + 1 = 0$。
    因式分解得 $(\lambda-1)^2 = 0$，解得二重实根 $\lambda_{1,2} = 1$。因此，通解为 \(\displaystyle \mathbf{y = (C_1 + C_2 x)\mathrm{e}^x}\)。

    \item[(5)] \textbf{求解思路}：这是一道四阶常系数线性方程。写出特征方程 $\lambda^4 - 1 = 0$。
    因式分解：$(\lambda^2 - 1)(\lambda^2 + 1) = (\lambda-1)(\lambda+1)(\lambda^2+1) = 0$。解得四个互异的特征根：$\lambda_1 = 1$, $\lambda_2 = -1$, $\lambda_{3,4} = \pm i$。分别对应指数解和三角函数解，通解为 \(\displaystyle \mathbf{y = C_1\mathrm{e}^x + C_2\mathrm{e}^{-x} + C_3\cos x + C_4\sin x}\)。

    \item[(6)] \textbf{求解思路}：这是三阶常系数齐次方程。特征方程为 $\lambda^3 - 6\lambda^2 + 3\lambda + 10 = 0$。
    利用试根法，当 $\lambda = -1$ 时，$(-1)^3 - 6(-1)^2 + 3(-1) + 10 = -1 - 6 - 3 + 10 = 0$，故含有因式 $(\lambda+1)$。
    利用多项式长除法或综合除法：
    $(\lambda^3 - 6\lambda^2 + 3\lambda + 10) \div (\lambda+1) = \lambda^2 - 7\lambda + 10 = (\lambda-2)(\lambda-5)$。
    解得三个互异的实根：$\lambda_1 = -1, \lambda_2 = 2, \lambda_3 = 5$。通解为 \(\displaystyle \mathbf{y = C_1\mathrm{e}^{-x} + C_2\mathrm{e}^{2x} + C_3\mathrm{e}^{5x}}\)。
\end{enumerate}
\end{solution}

\begin{exercise}[齐次方程初值问题]
求下列初值问题的解:\\
\begin{tabular}{ll}
  (1) $y''+4y'+29y=0,\ y(0)=0,\ y'(0)=15$; &
  (2) $4y''+4y'+y=0,\ y(0)=2,\ y'(0)=0$; \\
  (3) $y''-4y'+3y=0,\ y(0)=6,\ y'(0)=10$; &
  (4) $y''+25y=0,\ y(0)=2,\ y'(0)=5$.
\end{tabular}
\end{exercise}
\begin{solution}
\begin{enumerate}
    \item[(1)] \textbf{求解通解与初值}：
    \begin{enumerate}
        \item 特征方程 $\lambda^2 + 4\lambda + 29 = 0 \implies \lambda = \frac{-4 \pm \sqrt{16-116}}{2} = -2 \pm 5i$。
        \item 通解形式为 $y = \mathrm{e}^{-2x}(C_1\cos 5x + C_2\sin 5x)$。
        \item 代入初始条件 $y(0) = 0 \implies \mathrm{e}^0(C_1\cos 0 + C_2\sin 0) = C_1 = 0$。故 $y = C_2\mathrm{e}^{-2x}\sin 5x$。
        \item 对 $y$ 求导：$y' = C_2(-2\mathrm{e}^{-2x}\sin 5x + 5\mathrm{e}^{-2x}\cos 5x)$。
        \item 代入初始条件 $y'(0) = 15 \implies C_2(0 + 5\cos 0) = 5C_2 = 15 \implies C_2 = 3$。
        \item 满足条件的特解为 \(\displaystyle \mathbf{y = 3\mathrm{e}^{-2x}\sin 5x}\)。
    \end{enumerate}

    \item[(2)] \textbf{求解通解与初值}：
    \begin{enumerate}
        \item 特征方程 $4\lambda^2 + 4\lambda + 1 = 0 \implies (2\lambda+1)^2 = 0 \implies \lambda_{1,2} = -\frac{1}{2}$。
        \item 通解形式为 $y = (C_1 + C_2 x)\mathrm{e}^{-\frac{1}{2}x}$。
        \item 代入 $y(0) = 2 \implies (C_1 + 0)\mathrm{e}^0 = C_1 = 2$。故 $y = (2 + C_2 x)\mathrm{e}^{-\frac{1}{2}x}$。
        \item 求导：$y' = C_2\mathrm{e}^{-\frac{1}{2}x} - \frac{1}{2}(2 + C_2 x)\mathrm{e}^{-\frac{1}{2}x}$。
        \item 代入 $y'(0) = 0 \implies C_2 - \frac{1}{2}(2 + 0) = 0 \implies C_2 - 1 = 0 \implies C_2 = 1$。
        \item 满足条件的特解为 \(\displaystyle \mathbf{y = (2 + x)\mathrm{e}^{-\frac{x}{2}}}\)。
    \end{enumerate}

    \item[(3)] \textbf{求解通解与初值}：
    \begin{enumerate}
        \item 特征方程 $\lambda^2 - 4\lambda + 3 = 0 \implies (\lambda-1)(\lambda-3) = 0 \implies \lambda_1 = 1, \lambda_2 = 3$。
        \item 通解形式为 $y = C_1\mathrm{e}^x + C_2\mathrm{e}^{3x}$。
        \item 求导得 $y' = C_1\mathrm{e}^x + 3C_2\mathrm{e}^{3x}$。
        \item 代入初始条件联立方程组，得
        \(\displaystyle y(0) = C_1 + C_2 = 6,\qquad y'(0) = C_1 + 3C_2 = 10\)。
        \item 下式减去上式得 $2C_2 = 4 \implies C_2 = 2$。代回第一式得 $C_1 = 4$。
        \item 满足条件的特解为 \(\displaystyle \mathbf{y = 4\mathrm{e}^x + 2\mathrm{e}^{3x}}\)。
    \end{enumerate}

    \item[(4)] \textbf{求解通解与初值}：
    \begin{enumerate}
        \item 特征方程 $\lambda^2 + 25 = 0 \implies \lambda = \pm 5i$。
        \item 通解形式为 $y = C_1\cos 5x + C_2\sin 5x$。
        \item 代入 $y(0) = 2 \implies C_1\cos 0 + C_2\sin 0 = C_1 = 2$。
        \item 求导得 $y' = -5C_1\sin 5x + 5C_2\cos 5x$。
        \item 代入 $y'(0) = 5 \implies -0 + 5C_2\cos 0 = 5C_2 = 5 \implies C_2 = 1$。
        \item 满足条件的特解为 \(\displaystyle \mathbf{y = 2\cos 5x + \sin 5x}\)。
    \end{enumerate}
\end{enumerate}
\end{solution}

\begin{exercise}[非齐次方程通解]
求下列微分方程的通解:\\
\begin{tabular}{lll}
  (1) $y''+y'-2y=-2\sin x$; &
  (2) $y''+y=4\sin x$; &
  (3) $y''+4y=17\mathrm{e}^{-x}\cos 2x$; \\
  (4) $y''+a^2y=\mathrm{e}^x$; &
  (5) $y''+5y'+4y=3-2x$; &
  (6) $y''-6y'+9y=\mathrm{e}^{2x}(x+1)$; \\
  (7) $y''+y=3x\mathrm{e}^x+5\sin x$; &
  (8) $y''-y=\sin^2 x$. &
\end{tabular}
\end{exercise}
\begin{solution}
\begin{enumerate}
    \item[(1)] \textbf{求齐次通解}：$\lambda^2+\lambda-2=0 \implies (\lambda-1)(\lambda+2)=0 \implies \lambda_1=1, \lambda_2=-2$。齐次通解为 $y_h = C_1\mathrm{e}^x + C_2\mathrm{e}^{-2x}$。\\
    \textbf{求特解}：非齐次项 $-2\sin x$ 对应的等效根为 $\lambda = i$。由于 $i$ 不是特征根（不共振），设特解 $y^* = A\cos x + B\sin x$。
    求导：$(y^*)' = -A\sin x + B\cos x$, $(y^*)'' = -A\cos x - B\sin x$。代入原方程 $y''+y'-2y=-2\sin x$，得 \(\displaystyle (-A\cos x - B\sin x) + (-A\sin x + B\cos x) - 2(A\cos x + B\sin x) = -2\sin x\)，即 \(\displaystyle (-3A+B)\cos x + (-A-3B)\sin x = -2\sin x\)。比较系数得到 \(\displaystyle -3A + B = 0,\qquad -A - 3B = -2\)。
    将上式代入下式：$-A - 9A = -10A = -2 \implies A = \frac{1}{5}$，进而 $B = \frac{3}{5}$。
    \textbf{通解}：\(\displaystyle \mathbf{y = C_1\mathrm{e}^x + C_2\mathrm{e}^{-2x} + \frac{1}{5}\cos x + \frac{3}{5}\sin x}\)。

    \item[(2)] \textbf{求齐次通解}：$\lambda^2+1=0 \implies \lambda = \pm i$。齐次通解为 $y_h = C_1\cos x + C_2\sin x$。\\
    \textbf{求特解}：非齐次项 $4\sin x$ 对应等效根 $\lambda = i$。该根是特征单根（发生共振，需乘 $x$），设特解 $y^* = x(A\cos x + B\sin x)$。
    求导：
    $(y^*)' = (A\cos x + B\sin x) + x(-A\sin x + B\cos x)$
    $(y^*)'' = (-A\sin x + B\cos x) + (-A\sin x + B\cos x) + x(-A\cos x - B\sin x) = -2A\sin x + 2B\cos x - x(A\cos x + B\sin x)$
    代入原方程 $y''+y = 4\sin x$，得 \(\displaystyle [-2A\sin x + 2B\cos x - x(A\cos x + B\sin x)] + x(A\cos x + B\sin x) = 4\sin x\)。含有 $x$ 的项恰好抵消，于是 \(\displaystyle -2A\sin x + 2B\cos x = 4\sin x \implies -2A = 4,\ 2B = 0 \implies A = -2,\ B = 0\)。
    特解 $y^* = -2x\cos x$。
    \textbf{通解}：$\mathbf{y = C_1\cos x + C_2\sin x - 2x\cos x}$。

    \item[(3)] \textbf{求齐次通解}：$\lambda^2+4=0 \implies \lambda=\pm 2i$。齐次通解为 $y_h = C_1\cos 2x + C_2\sin 2x$。\\
    \textbf{求特解}：非齐次项 $17\mathrm{e}^{-x}\cos 2x$ 对应根 $\lambda = -1 + 2i$（不是特征根）。设 $y^* = \mathrm{e}^{-x}(A\cos 2x + B\sin 2x)$。
    利用微分算子法（位移定理）简化计算 $(D^2+4)y^* = 17\mathrm{e}^{-x}\cos 2x$：\(\displaystyle \mathrm{e}^{-x}((D-1)^2+4)[A\cos 2x + B\sin 2x] = 17\mathrm{e}^{-x}\cos 2x\)。
    消去 $\mathrm{e}^{-x}$，展开算子 $(D^2-2D+5)[A\cos 2x + B\sin 2x] = 17\cos 2x$。
    因为 $D^2(\cos 2x) = -4\cos 2x$ 以及 $D^2(\sin 2x) = -4\sin 2x$，故算子可简化为 $(-4-2D+5) = (1-2D)$，即 \(\displaystyle (1-2D)[A\cos 2x + B\sin 2x] = A\cos 2x + B\sin 2x - 2(-2A\sin 2x + 2B\cos 2x) = 17\cos 2x\)，从而 \(\displaystyle (A-4B)\cos 2x + (4A+B)\sin 2x = 17\cos 2x\)。
    对比系数：$4A+B=0 \implies B=-4A$。代入另一式：$A-4(-4A) = 17A = 17 \implies A=1, B=-4$。
    \textbf{通解}：$\mathbf{y = C_1\cos 2x + C_2\sin 2x + \mathrm{e}^{-x}(\cos 2x - 4\sin 2x)}$。

    \item[(4)] \textbf{求齐次通解}：$\lambda^2+a^2=0 \implies \lambda = \pm ai$。齐次通解为 $y_h = C_1\cos ax + C_2\sin ax$。\\
    \textbf{求特解}：非齐次项 $\mathrm{e}^x$ 对应 $\lambda=1$（不是根）。设 $y^* = A\mathrm{e}^x$。
    代入方程 $y''+a^2y = \mathrm{e}^x$，得 \(\displaystyle A\mathrm{e}^x + a^2 A\mathrm{e}^x = \mathrm{e}^x \implies A(1+a^2)=1 \implies A=\frac{1}{1+a^2}\)。
    \textbf{通解}：$\mathbf{y = C_1\cos ax + C_2\sin ax + \frac{\mathrm{e}^x}{1+a^2}}$。

    \item[(5)] \textbf{求齐次通解}：$\lambda^2+5\lambda+4=0 \implies \lambda = -1, -4$。齐次通解为 $y_h = C_1\mathrm{e}^{-x} + C_2\mathrm{e}^{-4x}$。\\
    \textbf{求特解}：非齐次项 $3-2x$ 是多项式，对应 $\lambda=0$（不是根）。设 $y^* = Ax + B$。
    求导得 $(y^*)'=A, (y^*)''=0$。代入方程 $y''+5y'+4y=3-2x$，得 \(\displaystyle 0 + 5A + 4(Ax+B) = 3 - 2x \implies 4Ax + (5A+4B) = -2x + 3\)。
    比较同次幂系数：
    $4A = -2 \implies A = -\frac{1}{2}$；
    $5A + 4B = 3 \implies 5(-\frac{1}{2}) + 4B = 3 \implies 4B = \frac{11}{2} \implies B = \frac{11}{8}$。
    \textbf{通解}：$\mathbf{y = C_1\mathrm{e}^{-x} + C_2\mathrm{e}^{-4x} - \frac{1}{2}x + \frac{11}{8}}$。

    \item[(6)] \textbf{求齐次通解}：$\lambda^2-6\lambda+9=0 \implies (\lambda-3)^2=0 \implies \lambda_{1,2}=3$。齐次通解为 $y_h = (C_1+C_2x)\mathrm{e}^{3x}$。\\
    \textbf{求特解}：非齐次项 $\mathrm{e}^{2x}(x+1)$ 对应 $\lambda=2$（不是根）。设 $y^* = \mathrm{e}^{2x}(Ax+B)$。
    使用微分算子法 $(D-3)^2 y^* = \mathrm{e}^{2x}(x+1)$：\(\displaystyle \mathrm{e}^{2x}(D+2-3)^2[Ax+B] = \mathrm{e}^{2x}(D-1)^2[Ax+B] = \mathrm{e}^{2x}(D^2-2D+1)[Ax+B] = \mathrm{e}^{2x}(0 - 2A + Ax+B) = \mathrm{e}^{2x}(Ax + B - 2A) = \mathrm{e}^{2x}(x+1)\)。
    比较系数：$A = 1$，$B - 2A = 1 \implies B - 2 = 1 \implies B = 3$。特解 $y^* = \mathrm{e}^{2x}(x+3)$。
    \textbf{通解}：$\mathbf{y = (C_1 + C_2 x)\mathrm{e}^{3x} + \mathrm{e}^{2x}(x+3)}$。

    \item[(7)] \textbf{求齐次通解}：$\lambda^2+1=0 \implies \lambda=\pm i$。齐次通解为 $y_h = C_1\cos x + C_2\sin x$。\\
    \textbf{求特解}：运用叠加原理分离非齐次项 $3x\mathrm{e}^x$ 和 $5\sin x$。
    对于 $3x\mathrm{e}^x$，对应 $\lambda=1$（不共振），设 $y_1^* = \mathrm{e}^x(Ax+B)$。
    代入方程 $(y_1^*)'' + y_1^*$：
    \(\displaystyle \mathrm{e}^x(Ax+B+2A) + \mathrm{e}^x(Ax+B) = \mathrm{e}^x(2Ax+2A+2B) = 3x\mathrm{e}^x\)。
    比较系数得 $2A=3 \implies A=\frac{3}{2}$；$2A+2B=0 \implies B=-\frac{3}{2}$。$y_1^* = \mathrm{e}^x\left(\frac{3}{2}x - \frac{3}{2}\right)$。
    对于 $5\sin x$，对应 $\lambda=i$（单根共振，需乘 $x$），设 $y_2^* = x(C\cos x + D\sin x)$。
    参考第(2)题的求导结果，代入 $(y_2^*)'' + y_2^* = 5\sin x$ 得 \(\displaystyle -2C\sin x + 2D\cos x = 5\sin x \implies -2C=5,\ 2D=0 \implies C=-\frac{5}{2}, D=0\)。
    得到 $y_2^* = -\frac{5}{2}x\cos x$。将各项叠加：
    \textbf{通解}：$\mathbf{y = C_1\cos x + C_2\sin x + \mathrm{e}^x\left(\frac{3}{2}x - \frac{3}{2}\right) - \frac{5}{2}x\cos x}$。

    \item[(8)] \textbf{求齐次通解}：$\lambda^2-1=0 \implies \lambda=\pm 1$。齐次通解为 $y_h = C_1\mathrm{e}^x + C_2\mathrm{e}^{-x}$。\\
    \textbf{求特解}：右端为 $\sin^2 x$，利用半角公式降幂得 $f(x) = \frac{1}{2} - \frac{1}{2}\cos 2x$。分为两部分。
    对常数项 $\frac{1}{2}$，不共振，设 $y_1^* = A$。代入 $y''-y$ 得 $0 - A = \frac{1}{2} \implies y_1^* = -\frac{1}{2}$。
    对 $-\frac{1}{2}\cos 2x$，不共振，设 $y_2^* = B\cos 2x$。
    代入方程 $(y_2^*)''-y_2^*$，得 \(\displaystyle -4B\cos 2x - B\cos 2x = -5B\cos 2x = -\frac{1}{2}\cos 2x \implies B = \frac{1}{10}\)。
    得到 $y_2^* = \frac{1}{10}\cos 2x$。
    \textbf{通解}：$\mathbf{y = C_1\mathrm{e}^x + C_2\mathrm{e}^{-x} - \frac{1}{2} + \frac{1}{10}\cos 2x}$。
\end{enumerate}
\end{solution}

\begin{exercise}[非齐次方程初值问题]
求下列初值问题的解:
\begin{shortexenum}
    \shortitem{1}{$y''+y+\sin 2x=0$, $y(\pi)=1$, $y'(\pi)=1$;} &
    \shortitem{2}{$y''-3y'+2y=5$, $y(0)=1$, $y'(0)=2$;} \\
    \shortitem{3}{$y''-y=4x\mathrm{e}^x$, $y(0)=0$, $y'(0)=1$;} &
    \shortitem{4}{$y''-4y'=5$, $y(0)=1$, $y'(0)=0$.}
\end{shortexenum}
\end{exercise}
\begin{solution}
\begin{enumerate}
    \item[(1)] \textbf{求通解}：方程移项化为 $y''+y=-\sin 2x$。特征根 $\lambda=\pm i$，齐次通解 $y_h = C_1\cos x + C_2\sin x$。
    由于非齐次项 $-\sin 2x$ 的频率为 2（不共振），设特解 $y^* = A\sin 2x$。
    代入得 $(y^*)''+y^* = -4A\sin 2x + A\sin 2x = -3A\sin 2x = -\sin 2x \implies A = \frac{1}{3}$。
    通解为：$y = C_1\cos x + C_2\sin x + \frac{1}{3}\sin 2x$。\\
    \textbf{定初值}：代入 $x=\pi$：
    代入得 \(\displaystyle y(\pi) = C_1\cos\pi + C_2\sin\pi + \frac{1}{3}\sin 2\pi = -C_1 = 1\)，
    故 \(\displaystyle C_1 = -1\)。
    求导：$y' = -C_1\sin x + C_2\cos x + \frac{2}{3}\cos 2x$。
    代入 $x=\pi$：
    代入得 \(\displaystyle y'(\pi) = -C_1\sin\pi + C_2\cos\pi + \frac{2}{3}\cos 2\pi = -C_2 + \frac{2}{3} = 1\)，
    故 \(\displaystyle C_2 = -\frac{1}{3}\)。
    \textbf{特解}：$\mathbf{y = -\cos x - \frac{1}{3}\sin x + \frac{1}{3}\sin 2x}$。

    \item[(2)] \textbf{求通解}：特征根 $\lambda^2-3\lambda+2=0 \implies \lambda=1, 2$。齐次通解 $y_h = C_1\mathrm{e}^x + C_2\mathrm{e}^{2x}$。
    非齐次项是常数 $5$，因 $\lambda=0$ 不共振，设特解为 $y^* = A$。代入方程得 $2A = 5 \implies y^* = \frac{5}{2}$。
    通解为：$y = C_1\mathrm{e}^x + C_2\mathrm{e}^{2x} + \frac{5}{2}$。\\
    \textbf{定初值}：
    代入 $y(0) = C_1 + C_2 + \frac{5}{2} = 1 \implies C_1 + C_2 = -\frac{3}{2}$。
    求导得 $y' = C_1\mathrm{e}^x + 2C_2\mathrm{e}^{2x}$，代入 $y'(0) = C_1 + 2C_2 = 2$。
    解方程组：下式减上式得 $C_2 = 2 - (-\frac{3}{2}) = \frac{7}{2}$。代回求得 $C_1 = -5$。
    \textbf{特解}：$\mathbf{y = -5\mathrm{e}^x + \frac{7}{2}\mathrm{e}^{2x} + \frac{5}{2}}$。

    \item[(3)] \textbf{求通解}：特征根 $\lambda^2-1=0 \implies \lambda=\pm 1$。齐次通解 $y_h = C_1\mathrm{e}^x + C_2\mathrm{e}^{-x}$。
    非齐次项为 $4x\mathrm{e}^x$，因 $\lambda=1$ 为特征单根（共振，需乘 $x$），设特解 $y^* = x(Ax+B)\mathrm{e}^x = (Ax^2+Bx)\mathrm{e}^x$。
    对其求导：
    $(y^*)' = (2Ax+B)\mathrm{e}^x + (Ax^2+Bx)\mathrm{e}^x = [Ax^2+(2A+B)x+B]\mathrm{e}^x$
    $(y^*)'' = (2Ax+2A+B)\mathrm{e}^x + [Ax^2+(2A+B)x+B]\mathrm{e}^x = [Ax^2+(4A+B)x+(2A+2B)]\mathrm{e}^x$
    代入原方程 $y''-y=4x\mathrm{e}^x$：
    即 \(\displaystyle [Ax^2+(4A+B)x+(2A+2B)]\mathrm{e}^x - [Ax^2+Bx]\mathrm{e}^x = (4Ax + 2A+2B)\mathrm{e}^x = 4x\mathrm{e}^x\)。
    比较同次幂系数：$4A=4 \implies A=1$；$2A+2B=0 \implies B=-1$。特解 $y^* = (x^2-x)\mathrm{e}^x$。
    通解为：$y = C_1\mathrm{e}^x + C_2\mathrm{e}^{-x} + (x^2-x)\mathrm{e}^x$。\\
    \textbf{定初值}：
    代入 $y(0) = C_1 + C_2 = 0 \implies C_2 = -C_1$。
    求导：$y' = C_1\mathrm{e}^x - C_2\mathrm{e}^{-x} + (2x-1)\mathrm{e}^x + (x^2-x)\mathrm{e}^x = C_1\mathrm{e}^x - C_2\mathrm{e}^{-x} + (x^2+x-1)\mathrm{e}^x$。
    代入 $y'(0) = C_1 - C_2 - 1 = 1 \implies C_1 - C_2 = 2$。
    由于 $C_2 = -C_1$，代入得 $2C_1 = 2 \implies C_1 = 1, C_2 = -1$。
    \textbf{特解}：$\mathbf{y = \mathrm{e}^x(x^2-x+1) - \mathrm{e}^{-x}}$。

    \item[(4)] \textbf{求通解}：特征方程 $\lambda^2-4\lambda=0 \implies \lambda=0, 4$。齐次通解 $y_h = C_1 + C_2\mathrm{e}^{4x}$。
    非齐次项是常数 $5$，因对应 $\lambda=0$ 为单根（共振，乘 $x$），设特解 $y^* = Ax$。
    代入原方程 $y''-4y'=5$ 得 $0 - 4A = 5 \implies A = -\frac{5}{4}$。
    通解为：$y = C_1 + C_2\mathrm{e}^{4x} - \frac{5}{4}x$。\\
    \textbf{定初值}：
    代入 $y(0) = C_1 + C_2 = 1$。
    求导：$y' = 4C_2\mathrm{e}^{4x} - \frac{5}{4}$。代入 $y'(0) = 4C_2 - \frac{5}{4} = 0 \implies C_2 = \frac{5}{16}$。
    从而 $C_1 = 1 - C_2 = 1 - \frac{5}{16} = \frac{11}{16}$。
    \textbf{特解}：$\mathbf{y = \frac{11}{16} + \frac{5}{16}\mathrm{e}^{4x} - \frac{5}{4}x}$。
\end{enumerate}
\end{solution}

\begin{exercise}[含参方程的通解]
求通解：(1)$y''+4y'+4y=\mathrm{e}^{\alpha x}$;~~(2)$y''+a^2y=\sin x$ ($a>0$).
\end{exercise}
\begin{solution}
\begin{enumerate}
    \item[(1)] \textbf{求解思路}：特征方程 $\lambda^2+4\lambda+4=0 \implies (\lambda+2)^2=0 \implies \lambda=-2$（二重实根）。齐次通解为 $y_h = (C_1 + C_2 x)\mathrm{e}^{-2x}$。
    非齐次项为指数函数 $\mathrm{e}^{\alpha x}$，其指数部分的系数为 $\alpha$。由于待定系数法的结构依赖于 $\alpha$ 是否与特征根重合，需分类讨论：
    \begin{itemize}
        \item \textbf{情形 1（非共振）：若 $\alpha \neq -2$}。不发生共振，设特解为 $y^* = A\mathrm{e}^{\alpha x}$，代入原方程：
        \(\displaystyle A\alpha^2\mathrm{e}^{\alpha x} + 4A\alpha\mathrm{e}^{\alpha x} + 4A\mathrm{e}^{\alpha x} = \mathrm{e}^{\alpha x} \implies A(\alpha^2+4\alpha+4) = 1 \implies A = \frac{1}{(\alpha+2)^2}\)。
        此时通解为 $\mathbf{y = (C_1 + C_2 x)\mathrm{e}^{-2x} + \frac{\mathrm{e}^{\alpha x}}{(\alpha+2)^2}}$。
        \item \textbf{情形 2（二重共振）：若 $\alpha = -2$}。$\lambda=-2$ 为二重根，发生深度共振。设特解为 $y^* = Ax^2\mathrm{e}^{-2x}$。
        求导：$(y^*)' = (2Ax - 2Ax^2)\mathrm{e}^{-2x}$，$(y^*)'' = (2A - 8Ax + 4Ax^2)\mathrm{e}^{-2x}$。代入方程 $y''+4y'+4y=\mathrm{e}^{-2x}$：
        \(\displaystyle (2A - 8Ax + 4Ax^2)\mathrm{e}^{-2x} + 4(2Ax - 2Ax^2)\mathrm{e}^{-2x} + 4Ax^2\mathrm{e}^{-2x} = 2A\mathrm{e}^{-2x} = \mathrm{e}^{-2x} \implies A = \frac{1}{2}\)。
        此时通解为 $\mathbf{y = (C_1 + C_2 x)\mathrm{e}^{-2x} + \frac{1}{2}x^2\mathrm{e}^{-2x}}$。
    \end{itemize}

    \item[(2)] \textbf{求解思路}：特征方程 $\lambda^2+a^2=0 \implies \lambda=\pm ai$。齐次通解为 $y_h = C_1\cos ax + C_2\sin ax$。
    非齐次项 $\sin x$ 对应的等价指数根为 $\pm i$。同样需要按是否发生共振分类讨论：
    \begin{itemize}
        \item \textbf{情形 1（非共振）：若 $a \neq 1$}（因 $a>0$，即特征根 $\pm ai \neq \pm i$）。设 $y^* = A\sin x$（因方程只含偶次导数，无需设余弦项）。
        代入得 $(y^*)''+a^2y^* = -A\sin x + a^2 A\sin x = (a^2-1)A\sin x = \sin x \implies A = \frac{1}{a^2-1}$。
        此时通解为 $\mathbf{y = C_1\cos ax + C_2\sin ax + \frac{\sin x}{a^2-1}}$。
        \item \textbf{情形 2（共振）：若 $a = 1$}。此时特征根为 $\pm i$，对应齐次解已含 $\sin x$ 与 $\cos x$。发生单根共振，设 $y^* = x(A\cos x + B\sin x)$。
        求二次导数得 $(y^*)'' = -2A\sin x + 2B\cos x - x(A\cos x + B\sin x)$。代入方程 $y''+y = \sin x$：
        \(\displaystyle (-2A\sin x + 2B\cos x - y^*) + y^* = -2A\sin x + 2B\cos x = \sin x \implies A = -\frac{1}{2}, B = 0\)。
        特解为 $y^* = -\frac{1}{2}x\cos x$。
        此时通解为 $\mathbf{y = C_1\cos x + C_2\sin x - \frac{1}{2}x\cos x}$。
    \end{itemize}
\end{enumerate}
\end{solution}

\begin{exercise}[Euler方程求解]
求下列欧拉方程的通解:\\
\begin{tabular}{lll}
  (1) $x^2y''-xy'+y=0$; &
  (2) $x^2y''+xy'+y=x$; &
  (3) $x^2y''+2xy'-n(n+1)y=0$; \\
  (4) $xy''+2y'=12\ln x$; &
  (5) $x^2y''-xy'+4y=x\sin(\ln x)$; &
  (6) $x^2y''-3xy'+4y=0$.
\end{tabular}
\end{exercise}
\begin{solution}
对于欧拉方程（Euler Equation），我们统一采用自变量代换 $x = \mathrm{e}^t$（即 $t = \ln x$）将其转化为以 $t$ 为自变量的常系数常微分方程。
\noindent\textbf{欧拉代换算子规律。}记微分算子 $D_t = \frac{\mathrm{d}}{\mathrm{d}t}$。利用链式法则可推导 \(\displaystyle x y' = D_t y,\qquad x^2 y'' = D_t(D_t-1)y = D_t^2 y - D_t y,\qquad x^3 y''' = D_t(D_t-1)(D_t-2)y\)。
利用上述规律，对各题代入替换求解：
\begin{enumerate}
    \item[(1)] 代入原方程 $x^2y''-xy'+y=0$ 得 \(\displaystyle (D_t^2 - D_t)y - D_t y + y = 0 \implies y_{tt}'' - 2y_t' + y = 0\)。
    常系数方程的特征根 $\lambda^2 - 2\lambda + 1 = 0 \implies (\lambda-1)^2 = 0$。
    在 $t$ 空间解得：$y = (C_1 + C_2 t)\mathrm{e}^t$。
    反代回 $x$ 空间（利用 $t=\ln x, \mathrm{e}^t=x$）：\(\displaystyle \mathbf{y = x(C_1 + C_2\ln x)}\)。

    \item[(2)] 代入原方程 $x^2y''+xy'+y=x$ 得 \(\displaystyle (D_t^2 - D_t)y + D_t y + y = \mathrm{e}^t \implies y_{tt}'' + y = \mathrm{e}^t\)。
    特征根 $\lambda^2 + 1 = 0 \implies \lambda=\pm i$。齐次通解 $y_h = C_1\cos t + C_2\sin t$。
    非齐次特解不共振，设为 $y^* = A\mathrm{e}^t \implies (y^*)'' = A\mathrm{e}^t$。代入得 $A\mathrm{e}^t + A\mathrm{e}^t = 2A\mathrm{e}^t = \mathrm{e}^t \implies A=\frac{1}{2}$。
    $t$ 空间通解：$y = C_1\cos t + C_2\sin t + \frac{1}{2}\mathrm{e}^t$。
    反代回 $x$ 空间：
    \(\displaystyle \mathbf{y = C_1\cos(\ln x) + C_2\sin(\ln x) + \frac{1}{2}x}\)。

    \item[(3)] 代入原方程 $x^2y''+2xy'-n(n+1)y=0$ 得：
    即 \(\displaystyle (D_t^2 - D_t)y + 2D_t y - n(n+1)y = 0 \implies y_{tt}'' + y_t' - n(n+1)y = 0\)。
    特征方程 $\lambda^2 + \lambda - n(n+1) = 0$。因式分解得 $(\lambda-n)(\lambda+n+1) = 0$，故 $\lambda_1=n, \lambda_2=-(n+1)$。
    在 $t$ 空间解得：$y = C_1\mathrm{e}^{nt} + C_2\mathrm{e}^{-(n+1)t}$。
    反代回 $x$ 空间：
    \(\displaystyle \mathbf{y = C_1 x^n + C_2 x^{-(n+1)}}\)。

    \item[(4)] 先将原方程同乘 $x$ 凑成标准形式：$x^2 y'' + 2xy' = 12x\ln x$。
    代入算子关系得 \(\displaystyle (D_t^2 - D_t)y + 2D_t y = 12t\mathrm{e}^t \implies y_{tt}'' + y_t' = 12t\mathrm{e}^t\)。
    特征方程 $\lambda^2+\lambda=0 \implies \lambda=0, -1$。齐次通解 $y_h = C_1 + C_2\mathrm{e}^{-t}$。
    设特解 $y^* = \mathrm{e}^t(At+B)$，求导得 $(y^*)'=\mathrm{e}^t(At+A+B)$, $(y^*)''=\mathrm{e}^t(At+2A+B)$。
    代入得 $(y^*)''+(y^*)' = \mathrm{e}^t(2At+3A+2B) = 12t\mathrm{e}^t$。
    比较系数得 $2A=12 \implies A=6$；$3A+2B=0 \implies B=-9$。
    $t$ 空间通解：$y = C_1 + C_2\mathrm{e}^{-t} + \mathrm{e}^t(6t-9)$。
    反代回 $x$ 空间：
    \(\displaystyle \mathbf{y = C_1 + C_2 x^{-1} + x(6\ln x - 9)}\)。

    \item[(5)] 代入原方程 $x^2y''-xy'+4y=x\sin(\ln x)$ 得 \(\displaystyle (D_t^2 - D_t)y - D_t y + 4y = \mathrm{e}^t\sin t \implies y_{tt}'' - 2y_t' + 4y = \mathrm{e}^t\sin t\)。
    特征根 $\lambda^2-2\lambda+4=0 \implies \lambda = 1 \pm \sqrt{3}i$。齐次解 $y_h = \mathrm{e}^t(C_1\cos\sqrt{3}t + C_2\sin\sqrt{3}t)$。
    非齐次项对应的指数为 $1 \pm i$，与特征根不同，不发生共振。设特解 $y^* = \mathrm{e}^t(A\cos t + B\sin t)$。
    连续求导得：
    \[
    (y^*)' = \mathrm{e}^t[(A+B)\cos t + (B-A)\sin t],\qquad
    (y^*)'' = \mathrm{e}^t[2B\cos t - 2A\sin t].
    \]
    代入方程合并后化简为
    \begin{align*}
    &\mathrm{e}^t[2B\cos t - 2A\sin t]
    - 2\mathrm{e}^t[(A+B)\cos t + (B-A)\sin t]
    + 4\mathrm{e}^t[A\cos t + B\sin t]\\
    &= \mathrm{e}^t[(2B - 2A - 2B + 4A)\cos t
    + (-2A - 2B + 2A + 4B)\sin t]\\
    &= \mathrm{e}^t(2A\cos t + 2B\sin t)
    = \mathrm{e}^t\sin t.
    \end{align*}
    解得：$2A=0, 2B=1 \implies A=0, B=\frac{1}{2}$。
    $t$ 空间解得：$y = \mathrm{e}^t(C_1\cos\sqrt{3}t + C_2\sin\sqrt{3}t) + \frac{1}{2}\mathrm{e}^t\sin t$。
    反代回 $x$ 空间：\(\displaystyle \mathbf{y = x\left(C_1\cos(\sqrt{3}\ln x) + C_2\sin(\sqrt{3}\ln x) + \frac{1}{2}\sin(\ln x)\right)}\)。

    \item[(6)] 代入原方程 $x^2y''-3xy'+4y=0$ 得 \(\displaystyle (D_t^2 - D_t)y - 3D_t y + 4y = 0 \implies y_{tt}'' - 4y_t' + 4y = 0\)。
    特征方程 $\lambda^2-4\lambda+4=0 \implies (\lambda-2)^2=0$。
    在 $t$ 空间解得：$y = (C_1 + C_2 t)\mathrm{e}^{2t}$。
    反代回 $x$ 空间（注意 $\mathrm{e}^{2t} = (\mathrm{e}^t)^2 = x^2$）：
    \(\displaystyle \mathbf{y = x^2(C_1 + C_2\ln x)}\)。
\end{enumerate}
\end{solution}

\begin{exercise}[Euler方程初值问题]
求方程 $x^2y''-xy'+y=2x$ 满足条件 $y(1)=0$, $y'(1)=1$ 的特解.
\end{exercise}
\begin{solution}
\begin{enumerate}
    \item \textbf{运用代换法求解通解：}\\
    令 $x = \mathrm{e}^t$，将原方程化为以 $t$ 为自变量的常系数常微分方程。
    代入算子 $x y' = y_t'$, $x^2 y'' = y_{tt}'' - y_t'$，得
    \(\displaystyle (y_{tt}'' - y_t') - y_t' + y = 2\mathrm{e}^t \implies y_{tt}'' - 2y_t' + y = 2\mathrm{e}^t\)。
    求齐次通解，特征方程 $\lambda^2-2\lambda+1=0 \implies (\lambda-1)^2=0$。齐次特征根为 $\lambda_{1,2} = 1$（二重根），所以齐次通解为 $y_h = (C_1 + C_2 t)\mathrm{e}^t$。
    对于右侧的非齐次项 $2\mathrm{e}^t$，它的指数 $\lambda=1$ 是特征二重根，发生深度共振（需乘 $t^2$）。设特解 $y^* = A t^2 \mathrm{e}^t$，求导：
    $(y^*)' = A(t^2+2t)\mathrm{e}^t$, $(y^*)'' = A(t^2+4t+2)\mathrm{e}^t$。
    代入方程得
    \(\displaystyle (y^*)'' - 2(y^*)' + y^* = A(t^2+4t+2 - 2t^2-4t + t^2)\mathrm{e}^t = 2A\mathrm{e}^t = 2\mathrm{e}^t\)，
    故 \( \displaystyle A=1\)。
    因此 $y$ 在 $t$ 空间的通解为：$y(t) = (C_1 + C_2 t + t^2)\mathrm{e}^t$。
    反代回自变量 $x = \mathrm{e}^t, t=\ln x$ 得到原始通解形式：
    \(\displaystyle y(x)=x(C_1 + C_2\ln x + \ln^2 x)\)。

    \item \textbf{利用初始条件定常数：}\\
    首先运用初值条件 $y(1)=0$，代入 $x=1$ 得
    $y(1) = 1 \cdot (C_1 + C_2 \ln 1 + \ln^2 1) = C_1 = 0$。
    于是通解化简为 $y(x) = x(C_2\ln x + \ln^2 x)$。
    对其利用乘积法则求导，并代入第二个初值条件 $y'(1)=1$，得
    \(\displaystyle y'(x)=1 \cdot (C_2\ln x + \ln^2 x) + x\left(C_2\frac{1}{x} + 2\ln x \cdot \frac{1}{x}\right)=C_2\ln x + \ln^2 x + C_2 + 2\ln x\)，
    从而 \(\displaystyle y'(1)=C_2\ln 1 + \ln^2 1 + C_2 + 2\ln 1 = C_2 = 1\)，故
    \(\displaystyle \boxed{y = x(\ln x + \ln^2 x)}\)。
\end{enumerate}
\end{solution}

\subsubsection{B组}

\begin{exercise}[无阻尼自由振动物理建模]
一个重 $p(\mathrm{kg})$ 的物体挂在弹簧下, 把弹簧拉长 $a(\mathrm{cm})$, 再用手把弹簧拉长 $A(\mathrm{cm})$ 后无初速地松开. 求弹簧的振动规律(不计介质阻力).
\end{exercise}
\begin{solution}
\begin{enumerate}
    \item \textbf{受力分析与参数推导：}\\
    （注：在经典力学习题中，物体重 $p(\mathrm{kg})$ 通常指千克力 $\mathrm{kgf}$，其质量为 $m = \frac{p}{g}$。此重物挂上后弹簧静力伸长了 $a$，故弹簧劲度系数满足 $ka = mg = p$，即 $k = \frac{p}{a}$）。
    \item \textbf{建立微分方程：}\\
    以重物静止时的平衡位置为坐标原点，向下为正方向建立 $y$ 轴。在该位置重力与初始静弹性力抵消，仅剩偏离平衡点带来的弹性恢复力起作用。
    根据牛顿第二定律，系统的动力学方程（无阻尼）为
    \(\displaystyle m y'' = -k y\)，即 \(\displaystyle \frac{p}{g} y'' = -\frac{p}{a} y\)，整理得
    \(\displaystyle y'' + \frac{g}{a} y = 0\)。
    这是一个标准的无阻尼简谐振动方程。其特征方程为 $\lambda^2 + \frac{g}{a} = 0$，解得固有圆频率 $\omega = \sqrt{\frac{g}{a}}$。
    方程的通解为 \(\displaystyle y(t)=C_1\cos\left(\sqrt{\frac{g}{a}}t\right)+C_2\sin\left(\sqrt{\frac{g}{a}}t\right)\)。

    \item \textbf{根据初始条件定常数：}\\
    题设指出“再用手把弹簧拉长 $A$ 后无初速地松开”。这意味着在 $t=0$ 瞬间，物体偏离静力平衡位置的距离为正 $A$，初速度为零。
    得到初始条件：$y(0) = A$，$y'(0) = 0$。
    将 $t=0$ 代入通解：
    $y(0) = C_1 = A$
    求导得 $y'(t) = -C_1\sqrt{\frac{g}{a}}\sin\left(\sqrt{\frac{g}{a}}t\right) + C_2\sqrt{\frac{g}{a}}\cos\left(\sqrt{\frac{g}{a}}t\right)$，代入 $t=0$：
    $y'(0) = C_2\sqrt{\frac{g}{a}} = 0 \implies C_2=0$
    最终振动规律为：
    \(\displaystyle \boxed{y(t) = A\cos\left(\sqrt{\frac{g}{a}}t\right)}\)。
\end{enumerate}
\end{solution}

\begin{exercise}[阻尼自由振动物理建模]
一个质量为 $m$ 的质点在一个与距离成正比的外力作用下离开中心 $O$ 点, 介质的阻力与速度成正比, 求质点的运动规律.
\end{exercise}
\begin{solution}
\begin{enumerate}
    \item \textbf{物理建模与微分方程的建立：}\\
    题设指出外力“与距离成正比”，这是一种指向平衡位置的典型弹性恢复力特征，根据胡克定律设为 $F_{ext} = -k x$（$k>0$ 为比例系数）。
    介质的阻力与速度成正比，且总是与运动方向相反阻碍运动，设为 $F_{fric} = -\mu x'$（$\mu>0$ 为阻尼系数）。
    由牛顿第二定律 $m x'' = \Sigma F$，得到完整的运动微分方程
    \(\displaystyle m x'' = -k x - \mu x' \implies m x'' + \mu x' + kx = 0\)。

    \item \textbf{方程求解与物理分类讨论：}\\
    这是经典的常系数二阶齐次方程，其特征方程为 $m\lambda^2 + \mu\lambda + k = 0$。
    通过求根公式解得 \(\displaystyle \lambda_{1,2} = \frac{-\mu \pm \sqrt{\mu^2 - 4mk}}{2m}\)。质点的运动规律将完全取决于根的性质，即阻尼判别式 $\Delta = \mu^2 - 4mk$ 的符号：
    \begin{enumerate}
        \item[(1)] \textbf{过阻尼情形}（大阻尼，当 $\mu^2 > 4mk$ 时）：
        特征根为两个不相等的负实数 $\lambda_1, \lambda_2$。运动规律为非周期性的纯衰减过程：\(\displaystyle \mathbf{x(t) = C_1\mathrm{e}^{\lambda_1 t} + C_2\mathrm{e}^{\lambda_2 t}}\)。
        质点由于阻力过大，只会缓慢地滑回平衡位置而不会发生来回越位的振荡。

        \item[(2)] \textbf{临界阻尼情形}（当 $\mu^2 = 4mk$ 时）：
        特征根为二重等负实根 \(\lambda = -\frac{\mu}{2m}\)。运动规律为 \(\displaystyle x(t) = (C_1 + C_2 t)\mathrm{e}^{-\frac{\mu}{2m}t}\)。
        这是介于振荡与不振荡之间的临界态，质点能以最快的速度返回平衡位置。

        \item[(3)] \textbf{欠阻尼情形}（小阻尼，当 $\mu^2 < 4mk$ 时）：
        特征根为一对共轭复数 $\lambda = -\frac{\mu}{2m} \pm i\omega_d$，其中带阻尼圆频率 $\omega_d = \frac{\sqrt{4mk - \mu^2}}{2m}$。运动规律为衰减振荡：\(\displaystyle \mathbf{x(t) = \mathrm{e}^{-\frac{\mu}{2m}t}(C_1\cos\omega_d t + C_2\sin\omega_d t)}\)。
        质点作振幅随指数律不断缩小的往复振荡，直到最终停留在平衡位置。
    \end{enumerate}
\end{enumerate}
\end{solution}

\begin{exercise}[利用切线条件的非齐次方程初值问题]
设函数 $y=y(x)$ 满足微分方程
\(\displaystyle y''-3y'+2y=2\mathrm{e}^x\)，
且其图形在点 $(0,1)$ 处的切线与曲线 $y=x^2-x+1$ 在该点的切线重合。求函数 $y=y(x)$。
\end{exercise}
\begin{solution}
\begin{enumerate}
    \item \textbf{将几何条件转化为初值条件：}\\
    题设明确指出解曲线经过点 $(0,1)$，因此得到初值条件 $y(0) = 1$。
    同时，解曲线在该点处的切线与抛物线 $y_0 = x^2-x+1$ 的切线重合，这意味着它们在该点的斜率（一阶导数值）相等。
    求抛物线的导数 $y_0'(x) = 2x - 1$，在 $x=0$ 处的斜率为 $y_0'(0) = -1$。
    故得到完整的初值条件：$y(0) = 1, y'(0) = -1$。

    \item \textbf{求解微分方程的通解：}\\
    原方程对应的齐次特征方程为 $\lambda^2 - 3\lambda + 2 = 0 \implies \lambda_1=1, \lambda_2=2$。
    非齐次项 $2\mathrm{e}^x$ 的指数 $\lambda=1$ 恰是特征单根，发生共振。设特解结构为 $y^* = Ax\mathrm{e}^x$。
    对其求导代入：$(y^*)' = A(x+1)\mathrm{e}^x$, $(y^*)'' = A(x+2)\mathrm{e}^x$。
    代入原方程 $y''-3y'+2y$，得 \(\displaystyle A(x+2)\mathrm{e}^x - 3A(x+1)\mathrm{e}^x + 2Ax\mathrm{e}^x = A\mathrm{e}^x(x+2 - 3x-3 + 2x) = -A\mathrm{e}^x = 2\mathrm{e}^x\)。
    解得 $A=-2$。因此特解为 $y^* = -2x\mathrm{e}^x$。
    方程通解为 \(\displaystyle y = C_1\mathrm{e}^x + C_2\mathrm{e}^{2x} - 2x\mathrm{e}^x\)。

    \item \textbf{利用初值确定常数：}\\
    将 $x=0$ 代入通解中，得 \(\displaystyle y(0) = C_1 + C_2 - 0 = 1 \implies C_1 + C_2 = 1\)。
    对通解求导，\(\displaystyle y'(x) = C_1\mathrm{e}^x + 2C_2\mathrm{e}^{2x} - 2(x+1)\mathrm{e}^x\)；
    再代入 $x=0$ 和 $y'(0) = -1$，得 \(\displaystyle y'(0) = C_1 + 2C_2 - 2 = -1 \implies C_1 + 2C_2 = 1\)。
    联立解方程组得 $C_2 = 0, C_1 = 1$。
    因此，满足全部条件的函数为 \(\displaystyle \mathbf{y(x) = \mathrm{e}^x - 2x\mathrm{e}^x = (1-2x)\mathrm{e}^x}\)。
\end{enumerate}
\end{solution}

\begin{exercise}[变加速力学建模：链条滑落]
一质量均匀的链条挂在一无摩擦的钉子上, 运动开始时, 链条的一边下垂 $8\ \mathrm{m}$, 另一边下垂 $10\ \mathrm{m}$. 试问整个链条滑过钉子需多少时间?
\end{exercise}
\begin{solution}
\begin{enumerate}
    \item \textbf{建立力学方程：}\\
    设链条的线密度为 $\rho$，总长度固定为 $L = 8 + 10 = 18\ \mathrm{m}$。
    为了描述系统的运动，设时刻 $t$ 时，较长那一端（初始下垂 $10\ \mathrm{m}$）下垂的长度为 $x(t)$。由于总长不变，另一边下垂的长度即为 $18 - x(t)$。
    由于钉子绝对光滑，整个链条受到的驱动力仅来自于两侧下垂链条的重力差。根据牛顿第二定律对系统整体（总质量 $M = 18\rho$）列方程，化简得到一个常系数二阶非齐次微分方程：
    \(\displaystyle F_{net} = \rho g x - \rho g (18-x) = \rho g(2x - 18)\)，
    且 \(\displaystyle M x'' = F_{net} \implies 18\rho x'' = \rho g(2x - 18) \implies x'' - \frac{g}{9}x = -g\)。

    \item \textbf{求解运动规律通解：}\\
    齐次特征方程 $\lambda^2 - \frac{g}{9} = 0 \implies \lambda = \pm \frac{\sqrt{g}}{3}$。令角频率记号 $\omega = \frac{\sqrt{g}}{3}$。
    齐次通解为 $x_h(t) = C_1\mathrm{e}^{\omega t} + C_2\mathrm{e}^{-\omega t}$，通常物理上写为双曲函数形式 $x_h(t) = C_1\cosh(\omega t) + C_2\sinh(\omega t)$ 会更方便。
    观察右侧常数非齐次项 $-g$，由于不共振，直接猜设特解 $x^* = A$。代入得 $-\frac{g}{9}A = -g \implies A = 9$。
    通解为 \(\displaystyle x(t) = C_1\cosh(\omega t) + C_2\sinh(\omega t) + 9\)。

    \item \textbf{代入初始状态定常数：}\\
    静止释放的初始时刻，状态为 $x(0) = 10, x'(0) = 0$。
    代入通解：
    $x(0) = C_1\cosh 0 + C_2\sinh 0 + 9 = C_1 + 9 = 10 \implies C_1 = 1$。
    求导：$x'(t) = \omega(C_1\sinh\omega t + C_2\cosh\omega t)$。
    代入得 $x'(0) = \omega(0 + C_2) = 0 \implies C_2 = 0$。
    所以链条滑落的运动方程确切为 $x(t) = 9 + \cosh(\omega t)$。

    \item \textbf{求解目标滑落时间：}\\
    当链条完全滑过钉子的瞬间，较长端的长度恰好占据了整个链条，即 $x(t) = 18\ \mathrm{m}$。代入运动方程，并利用反双曲余弦公式 $\operatorname{arccosh} y = \ln(y + \sqrt{y^2-1})$，得
    \(\displaystyle 18 = 9 + \cosh(\omega t) \implies \cosh(\omega t) = 9\)，从而
    \(\displaystyle \omega t = \operatorname{arccosh}(9) = \ln(9 + \sqrt{81 - 1}) = \ln(9 + \sqrt{80}) = \ln(9 + 4\sqrt{5})\)。
    最后，将 $\omega = \frac{\sqrt{g}}{3}$ 移过去，得 \(\displaystyle \mathbf{t = \frac{3}{\sqrt{g}}\ln(9+4\sqrt{5})\ \text{(s)}}\)。
\end{enumerate}
\end{solution}

\begin{exercise}[平衡位置的改变及振动建模]
设有长度为 $l$ 的弹簧, 其上端固定, 用五个都为 $m$ 的重物同时挂于弹簧下端, 使弹簧伸长了 $5a$. 今突然取去其中的一个重物, 使弹簧由静止状态开始振动. 若不计弹簧本身重量, 求所挂重物的运动规律.
\end{exercise}
\begin{solution}
\begin{enumerate}
    \item \textbf{求出弹簧的力学参数：}\\
    根据题设条件，挂载 $5m$ 的总质量时，静力平衡使得拉伸了 $5a$。因此弹簧弹力抵消重力 $k(5a) = 5mg \implies k = \frac{mg}{a}$。

    \item \textbf{确定振动方程与新平衡点：}\\
    突然取走一个重物后，系统转变为总质量为 $4m$ 的弹簧振子模型。
    此时，新的系统如果达到绝对静止，其理想的伸长量应为 $x_{eq} = \frac{4mg}{k} = 4a$。
    为了消去重力引起的常数偏置项，我们在物理上通常直接以\textbf{这个新的平衡位置}（即向下伸长 $4a$ 的位置）为原点，建立坐标系 $x$。此时新产生的微分方程是不含常数项的纯齐次方程：
    \(\displaystyle 4m x'' = -k x \implies 4m x'' + \frac{mg}{a}x = 0 \implies x'' + \frac{g}{4a}x = 0\)。
    这是一个固有角频率为 $\omega = \sqrt{\frac{g}{4a}}$ 的无阻尼自由振动，其通解为
    \(\displaystyle x(t) = C_1\cos\left(\sqrt{\frac{g}{4a}} t\right) + C_2\sin\left(\sqrt{\frac{g}{4a}} t\right)\)。

    \item \textbf{代入初值确立运动规律：}\\
    在突然取下重物的瞬间（$t=0$），剩余的重物由于惯性仍然停留在“旧平衡位置”，即距离悬挂点伸长 $5a$ 处。
    因此相对于我们刚才设定的新原点（$4a$ 处），它的初始物理位移是 $x(0) = 5a - 4a = a$。
    因为“由静止状态开始振动”，意味着它的初速度为 $0$，即 $x'(0) = 0$。
    将这组初值条件代入通解：
    \(\displaystyle x(0) = C_1 = a,\qquad
    x'(t) = -\omega C_1\sin(\omega t) + \omega C_2\cos(\omega t) \implies x'(0) = \omega C_2 = 0 \implies C_2=0\)。
    所以，相对于新平衡点，重物的规律为
    \(\displaystyle \mathbf{x(t) = a\cos\left(\frac{1}{2}\sqrt{\frac{g}{a}} t\right)},\qquad
    Y(t) = 4a + a\cos\left(\frac{1}{2}\sqrt{\frac{g}{a}} t\right)\)。
\end{enumerate}
\end{solution}

\begin{exercise}[一阶速度模型的降阶求解]
一质点徐徐沉入液体, 当下沉时, 液体的阻力与下沉速度成正比, 求此质点的运动规律.
\end{exercise}
\begin{solution}
\begin{enumerate}
    \item \textbf{建立运动微分方程：}\\
    设质点的质量为 $m$，液体对它的恒定浮力为 $F_b$。取竖直向下为正方向，设下沉深度的坐标为 $x(t)$。
    运动中质点共受三个力：向下的重力 $mg$，向上的浮力 $F_b$，以及向上的粘性阻力 $F_{fric} = kv$（其中 $k>0$ 为阻力比例系数，$v = x'$ 为质点速度）。
    由牛顿第二定律列出二阶微分方程。为了使方程的物理含义更清楚，令常量 $G' = mg - F_b$ 为“有效重力”，并定义 $g' = \frac{G'}{m}$，则方程可改写为：
    \(\displaystyle m x'' = mg - F_b - kx' \implies x'' + \frac{k}{m} x' = g'\)。

    \item \textbf{方程降阶求解速度：}\\
    观察发现，方程中并未出现位移 $x$ 本身。我们可以降一阶，将其视为关于速度 $v(t)$ 的一阶线性微分方程。先求齐次方程 $v'+\frac{k}{m}v=0$ 得通解 $C_1\mathrm{e}^{-\frac{k}{m}t}$，再取常数特解 $v^* = \frac{mg'}{k}$，于是
    \(\displaystyle v' + \frac{k}{m}v = g',\qquad v(t) = \frac{mg'}{k} + C_1\mathrm{e}^{-\frac{k}{m}t}\)。
    因为题目描述为“徐徐沉入”，合理的物理抽象是在液面刚刚下沉瞬间 $t=0$ 时初速度极小（通常按 $v(0)=0$ 理想化处理计算）。
    \(\displaystyle 0 = \frac{mg'}{k} + C_1 \implies C_1 = -\frac{mg'}{k},\qquad
    v(t) = \frac{mg'}{k}\left(1 - \mathrm{e}^{-\frac{k}{m}t}\right)\)。

    \item \textbf{积分求解运动位置规律：}\\
    位移就是速度对时间的积分（设初始位移 $x(0)=0$）：
    \[
    x(t)=\int_0^t v(\tau)\mathrm{d}\tau
    =\int_0^t \left[ \frac{mg'}{k} - \frac{mg'}{k}\mathrm{e}^{-\frac{k}{m}\tau} \right] \mathrm{d}\tau,
    \]
    从而
    \[
    x(t)=\mathbf{\frac{mg'}{k}t - \frac{m^2g'}{k^2}\left(1 - \mathrm{e}^{-\frac{k}{m}t}\right)}.
    \]
\end{enumerate}
\end{solution}

\begin{exercise}[频率计算（保留重力加速度符号 $g$）]
弹簧的弹性力与其伸缩量成正比. 设长度增加 $1\ \mathrm{cm}$ 时, 弹性力等于 $1\ \mathrm{kg}$. 现把 $2\ \mathrm{kg}$ 的重物悬挂在弹簧上, 如果先稍微把重物往下拉, 然后放开它, 求重物由此所产生的振动周期.
\end{exercise}
\begin{solution}
\begin{enumerate}
    \item \textbf{明确单位制并提取参数：}\\
    题目中所描述的“力等于 $1\ \mathrm{kg}$”指的是千克力（$\mathrm{kgf}$），即质量为 $1\ \mathrm{kg}$ 物体的重力 $1\cdot g$。悬挂重物的质量为 $m = 2\ \mathrm{kg}$。
    根据条件“长度增加 $\Delta x = 1$ 时，弹力 $F = 1 \cdot g$”，可以确定弹簧的劲度系数 $k = \frac{F}{\Delta x} = \frac{g}{1} = g$。

    \item \textbf{建立简谐振动模型求周期：}\\
    将重物置于弹簧末端发生振动时，其围绕静平衡位置的微分运动方程为：
    \(\displaystyle m x'' + k x = 0 \implies 2x'' + gx = 0 \implies x'' + \frac{g}{2}x = 0\)。
    根据标准的简谐振动方程 $x'' + \omega^2 x = 0$，对照系数可得
    \(\displaystyle \omega^2 = \frac{g}{2} \implies \omega = \sqrt{\frac{g}{2}},\qquad
    \mathbf{T = \frac{2\pi}{\omega} = \frac{2\pi}{\sqrt{\frac{g}{2}}} = 2\pi\sqrt{\frac{2}{g}}}\)。
\end{enumerate}
\end{solution}

\begin{exercise}[共振与受迫振动建模（保留符号 $g$）]
设重量为 $4\ \mathrm{kg}$ 的物体悬挂在弹簧上, 使弹簧伸长 $1\ \mathrm{cm}$. 如果弹簧的上端作铅直的简谐振动, 即 $y_0=2\sin 30t$ ($\mathrm{cm}$), 且在初始时刻重物处于静止状态(介质阻力不计), 求重物的运动规律.
\end{exercise}
\begin{solution}
\begin{enumerate}
    \item \textbf{确立受迫振动微分方程：}\\
    重量为 $4\ \mathrm{kgf}$ 的物体能产生大小为 $4g$ 的重力，已知其挂载使得弹簧伸长 $1\ \mathrm{cm}$ 达到平衡，故劲度系数 $k = \frac{4g}{1} = 4g$。系统的质量为 $m = 4$。
    以初始静平衡点为原点向下建轴，重物的绝对位移记为 $y(t)$。
    由于弹簧悬挂端并非静止，而是受到外部激励运动 $y_0(t) = 2\sin 30t$，这导致弹簧的实际形变相对变化了 $y - y_0$。
    根据牛顿定律建立非齐次微分方程，除以 4 并移项后化为标准形式：
    \(\displaystyle m y'' = -k(y - y_0) \implies 4y'' = -4g(y - 2\sin 30t) \implies y'' + gy = 2g\sin 30t\)。

    \item \textbf{求解通解规律：}\\
    对应的齐次方程 $y'' + gy = 0$ 的特征根为 $\pm i\sqrt{g}$。于是齐次通解为：
    \(\displaystyle y_h(t) = C_1\cos(\sqrt{g}t) + C_2\sin(\sqrt{g}t)\)。
    对于右侧由于外驱力产生的非齐次项 $2g\sin 30t$，若激励频率 $\omega=30$ 未与固有频率吻合产生共振（即假设 $g \neq 900$），可直接设特解结构为 \(\displaystyle y^* = A\sin 30t\)。
    对其求两次导后代入原方程 $(y^*)'' + g y^* = 2g\sin 30t$：
    \(\displaystyle -900A\sin 30t + gA\sin 30t = 2g\sin 30t \implies A(g - 900) = 2g \implies A = \frac{2g}{g - 900}\)。
    从而方程的完整通解为
    \(\displaystyle y(t) = C_1\cos(\sqrt{g}t) + C_2\sin(\sqrt{g}t) + \frac{2g}{g - 900}\sin 30t\)。

    \item \textbf{代入初值求解确定方程：}\\
    题设条件“初始时刻重物处于静止状态”等价于初值 $y(0) = 0$ 且 $y'(0) = 0$。
    首先代入 $t=0$ 得 $y(0) = C_1 \cdot 1 + C_2 \cdot 0 + 0 = 0 \implies C_1 = 0$。
    因此 $y(t)$ 简化为 $C_2\sin(\sqrt{g}t) + \frac{2g}{g - 900}\sin 30t$。求导数：
    求导得 \(\displaystyle y'(t)=C_2\sqrt{g}\cos(\sqrt{g}t) + \frac{60g}{g - 900}\cos 30t\)。
    由 \( \displaystyle y'(0)=C_2\sqrt{g} + \frac{60g}{g - 900} = 0\) 得
    \(\displaystyle C_2\sqrt{g} = -\frac{60g}{g - 900} \implies C_2 = -\frac{60\sqrt{g}}{g - 900}\)，故
    \(\displaystyle \mathbf{y(t)}=\mathbf{\frac{1}{g - 900}\left(2g\sin 30t - 60\sqrt{g}\sin(\sqrt{g}t)\right)}\)。
\end{enumerate}
\end{solution}

\begin{exercise}[积分微分方程转化为常系数方程]
设 $f(0)=0$, $f'(x)=1+\int_0^x[6\sin^2 t-f(t)]\mathrm{d}t$, 其中, $f(x)$ 二次可微, 求 $f(x)$.
\end{exercise}
\begin{solution}
\begin{enumerate}
    \item \textbf{通过求导消除积分号：}\\
    面对积分方程，常用的方法是对等式两边连续求导。利用变上限积分的求导定理 $\frac{\mathrm{d}}{\mathrm{d}x} \int_0^x g(t)\mathrm{d}t = g(x)$：
    得 \(\displaystyle f''(x) = \frac{\mathrm{d}}{\mathrm{d}x}(1) + 6\sin^2 x - f(x) = 6\sin^2 x - f(x),\qquad f''(x) + f(x) = 6\sin^2 x = 3(1-\cos 2x) = 3 - 3\cos 2x\)。

    \item \textbf{求解非齐次微分方程：}\\
    方程特征根为 $\pm i$，固有齐次通解为 $f_h(x) = C_1\cos x + C_2\sin x$。
    根据叠加原理将非齐次项分为 $3$ 与 $-3\cos 2x$ 分别猜特解（均不属于共振）：
    对常数项 $3$，设特解 $y_1^* = A$，直接代入得 $A=3$。
    对余弦项 $-3\cos 2x$，设特解 $y_2^* = B\cos 2x$，二次求导为 $-4B\cos 2x$：
    代入得 $-4B\cos 2x + B\cos 2x = -3B\cos 2x = -3\cos 2x \implies B = 1$。
    把所有的解合并，得到方程的通解结构：
    \(\displaystyle f(x) = C_1\cos x + C_2\sin x + 3 + \cos 2x\)。

    \item \textbf{发掘隐含初值以定常数：}\\
    除了题干明确给出的 $f(0) = 0$ 之外，我们需要第二个初值条件来求解两个常数。
    回到最初未求导的恒等式：$f'(x)=1+\int_0^x[6\sin^2 t-f(t)]\mathrm{d}t$。直接在该式中令 $x=0$，此时积分上限等于下限，整个积分项直接归零：
    $f'(0) = 1 + \int_0^0 [\cdots]\mathrm{d}t = 1 + 0 = 1$。
    现在结合这组条件 $(f(0)=0, f'(0)=1)$ 对通解定系数。
    对通解求一次导 $f'(x) = -C_1\sin x + C_2\cos x - 2\sin 2x$，再代入 $x=0$，得
    \(\displaystyle f(0)=C_1 + 4 = 0 \implies C_1 = -4,\qquad
    f'(0)=C_2 = 1 \implies C_2 = 1\)，故
    \(\displaystyle \mathbf{f(x)}=\mathbf{-4\cos x + \sin x + 3 + \cos 2x}\)。
\end{enumerate}
\end{solution}

% ============================================================
% §5.8 微分方程组的基本理论与常系数线性组的消元法
% ============================================================

\subsection{第七节习题：微分方程组的基本理论}

\begin{exercise}[基础：概念辨析]
判断下列方程组是否为线性微分方程组。若是，写出其矩阵形式 $\mathbf{x}'=A\mathbf{x}+\mathbf{f}(t)$（其中 $\mathbf{x}=(x,y)^{\mathsf T}$）：\\
(1) $\begin{cases}x'=xy+t,\\ y'=x^{2}-y^{2}\end{cases}$~~(2) $\begin{cases}x'=x+\sin t,\\ y'=y+\cos t\end{cases}$~~(3)$\begin{cases}x'=3x+y,\\ y'=x-y+z'\end{cases}$
\end{exercise}
\begin{solution}
\begin{enumerate}
    \item[(1)] \textbf{非线性}。因为方程中存在未知函数的乘积项 $xy$ 以及平方项 $x^2, y^2$，不满足未知函数及其导数次数均为一次的要求。
    \item[(2)] \textbf{线性}。其未知函数 $x, y$ 均为一次方，独立项只含有自变量 $t$。
    可以写成矩阵形式：
    \[
    \begin{pmatrix} x' \\ y' \end{pmatrix} = \begin{pmatrix} 1 & 0 \\ 0 & 1 \end{pmatrix} \begin{pmatrix} x \\ y \end{pmatrix} + \begin{pmatrix} \sin t \\ \cos t \end{pmatrix}
    \]
    \item[(3)] \textbf{非二维标准线性系统}。虽然各项看起来是一次的，但第二式中引入了一个新的未知函数导数 $z'$，这就需要额外提供关于 $z$ 的微分方程才能构成封闭的微分方程组。对于仅含 $x, y$ 的系统而言是不完备的。
\end{enumerate}
\end{solution}

\begin{exercise}[基础：消元法练习]
用消元法求解下列常系数线性方程组：
\begin{shortexenum}
  \shortchoice{(1)}{$\begin{cases}x'=x+2y,\\y'=4x+3y\end{cases}$} &
  \shortchoice{(2)}{$\begin{cases}x'=-2x+y,\\ y'=-x-2y\end{cases},\;x(0)=1,\;y(0)=0$} \\
  \shortchoice{(3)}{$\begin{cases}x'=x+y+z,\\ y'=x+2y-z,\\ z'=2x+3y+z\end{cases}$} &
\end{shortexenum}
\end{exercise}
\begin{solution}
\begin{enumerate}
    \item[(1)] \textbf{思路：消去 $y$。}\\
    由第一式得 $2y = x' - x \implies y = \frac{1}{2}x' - \frac{1}{2}x$。
    求导得 $y' = \frac{1}{2}x'' - \frac{1}{2}x'$。
    代入第二式得
    \(\displaystyle \frac{1}{2}x'' - \frac{1}{2}x' = 4x + 3\left(\frac{1}{2}x' - \frac{1}{2}x\right) = \frac{5}{2}x + \frac{3}{2}x'\)。
    整理得 $x'' - 4x' - 5x = 0$。特征方程 $\lambda^2 - 4\lambda - 5 = 0$，解得 $\lambda_1 = 5, \lambda_2 = -1$。
    故 $x = C_1\mathrm{e}^{5t} + C_2\mathrm{e}^{-t}$。
    回代得
    \(\displaystyle y = \frac{1}{2}(5C_1\mathrm{e}^{5t} - C_2\mathrm{e}^{-t}) - \frac{1}{2}(C_1\mathrm{e}^{5t} + C_2\mathrm{e}^{-t}) = 2C_1\mathrm{e}^{5t} - C_2\mathrm{e}^{-t}\)。
    通解为 \(\displaystyle x(t)=C_1\mathrm{e}^{5t}+C_2\mathrm{e}^{-t}\)，且
    \(\displaystyle y(t)=2C_1\mathrm{e}^{5t}-C_2\mathrm{e}^{-t}\)。

    \item[(2)] \textbf{思路：消去 $y$。}\\
    由第一式得 $y = x' + 2x \implies y' = x'' + 2x'$。
    代入第二式得
    \(\displaystyle x'' + 2x' = -x - 2(x' + 2x) = -5x - 2x' \implies x'' + 4x' + 5x = 0\)。
    特征方程 $\lambda^2 + 4\lambda + 5 = 0 \implies \lambda = -2 \pm i$。
    故 $x = \mathrm{e}^{-2t}(C_1\cos t + C_2\sin t)$。
    代入初始条件 $x(0) = 1 \implies C_1 = 1$。所以 $x = \mathrm{e}^{-2t}(\cos t + C_2\sin t)$。
    求导得
    \(\displaystyle x' = -2\mathrm{e}^{-2t}(\cos t + C_2\sin t) + \mathrm{e}^{-2t}(-\sin t + C_2\cos t) = \mathrm{e}^{-2t}[(-2+C_2)\cos t - (2C_2+1)\sin t]\)。
    回代得
    \(\displaystyle y = x' + 2x = \mathrm{e}^{-2t}[(-2+C_2)\cos t - (2C_2+1)\sin t] + 2\mathrm{e}^{-2t}(\cos t + C_2\sin t) = \mathrm{e}^{-2t}[C_2\cos t - \sin t]\)。
    代入初始条件 $y(0) = 0 \implies C_2 = 0$。
    特解为 \(\displaystyle x(t) = \mathrm{e}^{-2t}\cos t,\ y(t) = -\mathrm{e}^{-2t}\sin t\)。

    \item[(3)] \textbf{思路：对三维系统进行连续消元。}\\
    从第一式解出 $z = x' - x - y$。对 $t$ 求导得 $z' = x'' - x' - y'$。
    将 $z$ 代入第二式，得
    \(\displaystyle y'=x+2y-(x'-x-y)=2x+3y-x' \implies y'-3y=2x-x' \quad \text{--- (i)}\)。
    将 $z$ 和 $z'$ 代入第三式，得
    \(\displaystyle x''-x'-y'=2x+3y+(x'-x-y)=x+2y+x' \implies y'+2y=x''-2x'-x \quad \text{--- (ii)}\)。
    由 (ii) 减 (i) 消去 $y'$，得
    \(\displaystyle 5y=x''-x'-3x \implies y=\frac{1}{5}(x''-x'-3x) \quad \text{--- (iii)}\)。
    求导得 $y' = \frac{1}{5}(x''' - x'' - 3x')$。
    将 $y$ 和 $y'$ 代回式 (i)，得
    \(\displaystyle \frac{1}{5}(x'''-x''-3x')-\frac{3}{5}(x''-x'-3x)=2x-x' \implies x'''-4x''+5x'-x=0\)。
    其对应的特征方程为 $\lambda^3 - 4\lambda^2 + 5\lambda - 1 = 0$。
    （注：此三次方程无有理根，可用卡尔丹公式求出三个实根 $\lambda_1, \lambda_2, \lambda_3$）。
    在此假设已解出根，则 $x(t) = \sum_{i=1}^3 C_i \mathrm{e}^{\lambda_i t}$。
    之后，利用式 (iii) 可求得 $y(t)$，再代入 $z = x' - x - y$ 即可求得 $z(t)$。
\end{enumerate}
\end{solution}

\begin{exercise}[提高：非齐次方程组]
求解非齐次线性方程组 \(\displaystyle x'-2x+y=\mathrm{e}^{2t},\qquad y'+x-2y=\mathrm{e}^{3t}\)。
\end{exercise}
\begin{solution}
\begin{enumerate}
    \item \textbf{通过消元法化为高阶非齐次方程：}\\
    由第一式解出 \( \displaystyle y = \mathrm{e}^{2t} - x' + 2x \)，从而
    \(\displaystyle y' = 2\mathrm{e}^{2t} - x'' + 2x'\)。
    将 $y$ 和 $y'$ 代入第二式，得
    \(\displaystyle (2\mathrm{e}^{2t}-x''+2x')+x-2(\mathrm{e}^{2t}-x'+2x)=\mathrm{e}^{3t}\implies x''-4x'+3x=-\mathrm{e}^{3t}\)。

    \item \textbf{求解二阶非齐次方程：}\\
    特征方程 $\lambda^2 - 4\lambda + 3 = 0 \implies \lambda = 1, 3$。
    齐次通解 $x_h = C_1\mathrm{e}^t + C_2\mathrm{e}^{3t}$。
    由于非齐次项 $-\mathrm{e}^{3t}$ 的指数 $3$ 是特征单根，发生共振。设特解为 $x^* = A t \mathrm{e}^{3t}$。
    求导代入方程得
    \(\displaystyle (D-1)(D-3)[At\mathrm{e}^{3t}] = -\mathrm{e}^{3t} \implies A = -\frac{1}{2}\)，故
    \(\displaystyle x(t) = C_1\mathrm{e}^t + C_2\mathrm{e}^{3t} - \frac{1}{2}t\mathrm{e}^{3t}\)。

    \item \textbf{回代求 $y$：}\\
    先求 \(\displaystyle x' = C_1\mathrm{e}^t + 3C_2\mathrm{e}^{3t} - \frac{1}{2}\mathrm{e}^{3t} - \frac{3}{2}t\mathrm{e}^{3t}\)。
    代入 $y = \mathrm{e}^{2t} - x' + 2x$，整理得
    \(\displaystyle y = C_1\mathrm{e}^t - C_2\mathrm{e}^{3t} + \mathrm{e}^{2t} + \frac{1}{2}\mathrm{e}^{3t} + \frac{1}{2}t\mathrm{e}^{3t}\)。
    通解为 \(\displaystyle x(t)=C_1\mathrm{e}^t+C_2\mathrm{e}^{3t}-\frac{1}{2}t\mathrm{e}^{3t}\)，且
    \(\displaystyle y(t)=C_1\mathrm{e}^t-C_2\mathrm{e}^{3t}+\mathrm{e}^{2t}+\frac{1}{2}\mathrm{e}^{3t}+\frac{1}{2}t\mathrm{e}^{3t}\)。
\end{enumerate}
\end{solution}

\begin{exercise}[解微分方程组综合]
求下列微分方程组的通解或特解:\\
  (1) $\begin{cases}\dfrac{\mathrm{d}x}{\mathrm{d}t} = y, \\ \dfrac{\mathrm{d}y}{\mathrm{d}t} = x; \end{cases}$ ~
  (2) $\begin{cases}\dfrac{\mathrm{d}x}{\mathrm{d}t} = y,\, x(0)=0, \\\dfrac{\mathrm{d}y}{\mathrm{d}t} = -x,\, y(0)=1;\end{cases}$ ~
  (3) $\begin{cases} \dfrac{\mathrm{d}x}{\mathrm{d}t} = x+2y, \\ \dfrac{\mathrm{d}y}{\mathrm{d}t} = 4x+3y; \end{cases}$ ~
  (4) $\begin{cases} \dfrac{\mathrm{d}x}{\mathrm{d}t} = x-5y, \\ \dfrac{\mathrm{d}y}{\mathrm{d}t} = 2x-y. \end{cases}$
\end{exercise}
\begin{solution}
\begin{enumerate}
    \item[(1)] \textbf{思路：相互代入法。}\\
    由第一式 $y = x'$，求导得 $y' = x''$。代入第二式得 $x'' = x \implies x'' - x = 0$。
    特征根为 $\lambda = \pm 1$，故 $x(t) = C_1\mathrm{e}^t + C_2\mathrm{e}^{-t}$。
    回代求 $y$：$y(t) = x' = C_1\mathrm{e}^t - C_2\mathrm{e}^{-t}$。
    通解为 \(\displaystyle x(t) = C_1\mathrm{e}^t + C_2\mathrm{e}^{-t},\ y(t) = C_1\mathrm{e}^t - C_2\mathrm{e}^{-t}\)。

    \item[(2)] \textbf{思路：典型谐振子方程组。}\\
    由第一式 $y = x'$，求导得 $y' = x''$。代入第二式得 $x'' = -x \implies x'' + x = 0$。
    特征根为 $\lambda = \pm i$，故 $x(t) = C_1\cos t + C_2\sin t$。
    代入初值 $x(0) = 0 \implies C_1 = 0$，故 $x(t) = C_2\sin t$。
    回代得 $y(t) = x' = C_2\cos t$。
    代入初值 $y(0) = 1 \implies C_2 = 1$。
    特解为 \(\displaystyle x(t) = \sin t,\ y(t) = \cos t\)。

    \item[(3)] \textbf{思路：此题与前面“基础：消元法练习”的(1)完全一致。}\\
    直接给出结果：\(\displaystyle x(t) = C_1\mathrm{e}^{5t} + C_2\mathrm{e}^{-t},\ y(t) = 2C_1\mathrm{e}^{5t} - C_2\mathrm{e}^{-t}\)。

    \item[(4)] \textbf{思路：消元法。}\\
    由第一式解出 $y$：$5y = x - x' \implies y = \frac{1}{5}(x - x')$。
    求导得 $y' = \frac{1}{5}(x' - x'')$。
    代入第二式得
    \(\displaystyle \frac{1}{5}(x' - x'') = 2x - \frac{1}{5}(x - x') \implies x' - x'' = 10x - x + x' \implies -x'' = 9x \implies x'' + 9x = 0\)。
    特征根为 $\lambda = \pm 3i$。故 $x(t) = C_1\cos 3t + C_2\sin 3t$。
    回代得
    \(\displaystyle y = \frac{1}{5}\left[ C_1\cos 3t + C_2\sin 3t - (-3C_1\sin 3t + 3C_2\cos 3t) \right]\)。
    整理得通解
    \(\displaystyle x(t)=C_1\cos 3t+C_2\sin 3t\)，
    \(\displaystyle y(t)=\frac{C_1-3C_2}{5}\cos 3t+\frac{3C_1+C_2}{5}\sin 3t\)。
\end{enumerate}
\end{solution}

\begin{exercise}[微分方程组物理建模：炮弹受阻运动]
炮弹以初速度 $v_0$ 且与水平线成 $\alpha$ 角从炮口射出. 设空气阻力与速度成正比, 求炮弹的运动方程.
\end{exercise}
\begin{solution}
\begin{enumerate}
    \item \textbf{力学建模与推导：}\\
    设炮弹质量为 $m$，空气阻力系数为 $k > 0$。在二维平面上建立直角坐标系，以原点为出发点。
    重力沿 $y$ 轴负向，阻力 $\vec{f} = -k\vec{v} = -k(x'\vec{i} + y'\vec{j})$。
    根据牛顿第二定律分别在 $x, y$ 轴投影，得
    \(\displaystyle m x''=-k x'\)，\(\displaystyle m y''=-mg-k y'\)。
    初始条件为：$x(0)=0, y(0)=0$，且初始速度分量为 $x'(0) = v_0\cos\alpha, y'(0) = v_0\sin\alpha$。
    这是一个相互解耦的微分方程组，可以分别求解。

    \item \textbf{求解 $x$ 轴方程：}\\
    $x'' + \frac{k}{m}x' = 0$，特征根为 $0, -\frac{k}{m}$。
    $x'(t) = C_1\mathrm{e}^{-\frac{k}{m}t}$，代入初速 $x'(0) = v_0\cos\alpha \implies C_1 = v_0\cos\alpha$。
    所以
    \(\displaystyle x(t) = \int v_0\cos\alpha\,\mathrm{e}^{-\frac{k}{m}t}\mathrm{d}t
    = -\frac{mv_0\cos\alpha}{k}\mathrm{e}^{-\frac{k}{m}t} + C_2
    = \frac{mv_0\cos\alpha}{k}\left(1 - \mathrm{e}^{-\frac{k}{m}t}\right)\)。

    \item \textbf{求解 $y$ 轴方程：}\\
    $y'' + \frac{k}{m}y' = -g$。
    齐次通解为 $y_h' = C_3\mathrm{e}^{-\frac{k}{m}t}$。特解为 $y_p' = -\frac{mg}{k}$。
    所以 $y'(t) = C_3\mathrm{e}^{-\frac{k}{m}t} - \frac{mg}{k}$。
    代入初速 $y'(0) = v_0\sin\alpha \implies C_3 - \frac{mg}{k} = v_0\sin\alpha \implies C_3 = v_0\sin\alpha + \frac{mg}{k}$。
    积分求位移：
    \(\displaystyle y(t) = -\frac{m}{k}\left(v_0\sin\alpha + \frac{mg}{k}\right)\mathrm{e}^{-\frac{k}{m}t} - \frac{mg}{k}t + C_4
    = \frac{m}{k}\left(v_0\sin\alpha + \frac{mg}{k}\right)\left(1 - \mathrm{e}^{-\frac{k}{m}t}\right) - \frac{mg}{k}t\)。

    \item \textbf{最终结论：}\\
    炮弹的运动方程参数形式为：
    \[
    \boxed{
    \begin{cases}
    x(t) = \frac{m v_0 \cos\alpha}{k}\left(1 - \mathrm{e}^{-\frac{k}{m}t}\right) \\[6pt]
    y(t) = \frac{m}{k}\left(v_0 \sin\alpha + \frac{mg}{k}\right)\left(1 - \mathrm{e}^{-\frac{k}{m}t}\right) - \frac{mg}{k}t
    \end{cases}
    }
    \]
\end{enumerate}
\end{solution}

\begin{exercise}[有心力场中的粒子轨迹]
设质点 $M$ 受力心 $O$ 的吸引力与距离成正比, 今有一质点从与力心相距 $a$ 的点 $A$ 出发, 以垂直于线段 $OA$ 的初速度 $v_0$ 开始运动, 求质点 $M$ 的轨迹.
\end{exercise}
\begin{solution}
\begin{enumerate}
    \item \textbf{建立坐标系与方程：}\\
    以力心 $O$ 为原点建立直角坐标系。设点 $A$ 在 $x$ 轴正半轴上，则初始位置坐标为 $x(0) = a, y(0) = 0$。
    由于初速度 $v_0$ 垂直于 $OA$（即垂直于 $x$ 轴），所以初始速度条件为 $x'(0) = 0, y'(0) = v_0$。
    受力心吸引且引力与距离成正比，即 $\vec{F} = -k\vec{r} = -k(x\vec{i} + y\vec{j})$，其中 $k>0$。
    由牛顿第二定律，运动微分方程为解耦的形式：\(\displaystyle m x''=-kx,\ m y''=-ky\Longrightarrow x''+\omega^2 x=0,\ y''+\omega^2 y=0\)。
    其中 $\omega = \sqrt{\frac{k}{m}}$ 为系统自然角频率。
    \item \textbf{求解运动方程：}\\
    两个方程均为简谐振动方程，通解为：
    即 \(\displaystyle x(t) = C_1\cos\omega t + C_2\sin\omega t,\qquad
    y(t) = C_3\cos\omega t + C_4\sin\omega t\)。
    代入 $x$ 轴的初始条件 $x(0)=a, x'(0)=0$：\\
    $C_1 = a$, $C_2\omega = 0 \implies C_2 = 0$。故 $x(t) = a\cos\omega t$。\\
    代入 $y$ 轴的初始条件 $y(0)=0, y'(0)=v_0$：\\
    $C_3 = 0$, $C_4\omega = v_0 \implies C_4 = \frac{v_0}{\omega}$。故 $y(t) = \frac{v_0}{\omega}\sin\omega t$。
    \item \textbf{消去参数 $t$ 求几何轨迹：}\\
    由运动方程有
    即 \(\displaystyle \frac{x}{a}=\cos\omega t,\qquad \frac{\omega y}{v_0}=\sin\omega t
    \implies \boxed{\frac{x^2}{a^2}+\frac{y^2}{\left(\frac{v_0}{\omega}\right)^2}=1}\)。
    由于这是标准方程的形式，因此质点的运动轨迹是一个\textbf{椭圆}，力心 $O$ 位于椭圆的中心。
\end{enumerate}
\end{solution}

\subsection{本章总习题}

\begin{exercise}[基本概念与方程求解]
填空题:
\begin{shortexenum}
    \shortchoice{(1)}{$(y''')^4+(y')^4+x^3+x+1=0$ 是 \_\_\_\_\_\_ 阶微分方程.} &
    \shortchoice{(2)}{曲线族 $y=\cos(x+C)$ ($C$ 为任意常数)所满足的一阶微分方程是 \_\_\_\_\_\_.} \\
    \shortchoice{(3)}{已知线性齐次方程的基础解系为 $y_1=e^x, y_2=xe^x$, 则该方程为 \_\_\_\_\_\_.} &
    \shortchoice{(4)}{设方程 $y''+a(x)y'+b(x)y=0$ 有非零特解 $y=u(x)$, 则与其线性无关的特解为 \_\_\_\_\_\_.} \\
    \shortchoice{(5)}{设 $y_1=3, y_2=3+x^2, y_3=3+x^2+e^x$ 都是方程 \(\displaystyle (x^2-2x)y''-(x^2-2)y'+(2x-2)y=6x-6\) 的解, 则方程的通解为 \_\_\_\_\_\_.} &
\end{shortexenum}
\end{exercise}
\begin{solution}
\begin{enumerate}
    \item[(1)] \textbf{解析}：微分方程的阶数由方程中未知函数的最高阶导数决定。观察方程 $(y''')^4+(y')^4+x^3+x+1=0$，其中出现的最高阶导数为 $y'''$（三阶导数），因此该方程是\textbf{3}阶微分方程。

    \item[(2)] \textbf{解析}：要求曲线族满足的微分方程，需要通过求导消去任意常数 $C$。
    \begin{enumerate}
        \item 对 $y=\cos(x+C)$ 求导得：$y'=-\sin(x+C)$。
        \item 利用三角恒等式 $\sin^2(x+C)+\cos^2(x+C)=1$，将 $y$ 和 $y'$ 的表达式代入得 \(\displaystyle (-y')^2 + y^2 = 1 \implies (y')^2 + y^2 = 1\)。
        所以该曲线族满足的一阶微分方程是 \textbf{$(y')^2 + y^2 = 1$}。
    \end{enumerate}

    \item[(3)] \textbf{解析}：基础解系为 $y_1=e^x, y_2=xe^x$ 表明特征方程存在二重实根 $\lambda_{1,2}=1$。
    \begin{enumerate}
        \item 构造对应的特征方程：\(\displaystyle (\lambda-1)^2 = 0 \implies \lambda^2 - 2\lambda + 1 = 0 \implies y'' - 2y' + y = 0\)。
        \item 将特征方程还原为二阶常系数线性齐次微分方程。
        故该方程为 \textbf{$y'' - 2y' + y = 0$}。
    \end{enumerate}

    \item[(4)] \textbf{解析}：已知二阶线性齐次方程的一个非零特解 $y=u(x)$，可以利用刘维尔公式（Liouville's formula）或降阶法求出另一个线性无关的特解。
    \begin{enumerate}
        \item 设另一解为 $y = u(x)v(x)$。
        \item 根据刘维尔公式，直接写出与其线性无关的特解公式：\(\displaystyle y_2 = u(x) \int \frac{1}{u^2(x)} e^{-\int a(x)\mathrm{d}x} \mathrm{d}x\)。
        填空答案为 \textbf{$u(x) \int \frac{1}{u^2(x)} e^{-\int a(x)\mathrm{d}x} \mathrm{d}x$}。
    \end{enumerate}

    \item[(5)] \textbf{解析}：利用线性微分方程解的结构定理。
    \begin{enumerate}
        \item 已知 $y_1, y_2, y_3$ 均为非齐次方程的解，则它们的差是对应齐次方程的解。
        \item 构造齐次方程的解：
        \(\displaystyle Y_1 = y_2 - y_1 = (3+x^2) - 3 = x^2,\qquad
        Y_2 = y_3 - y_2 = (3+x^2+e^x) - (3+x^2) = e^x\)。
        \item 显然 $x^2$ 与 $e^x$ 线性无关，构成了齐次方程的基础解系。
        \item 非齐次方程的通解等于齐次方程的通解加上非齐次方程的一个特解（可任取 $y_1=3$ 作为特解）。
        故通解为 \textbf{$y = C_1 x^2 + C_2 e^x + 3$} （$C_1, C_2$ 为任意常数）。
    \end{enumerate}
\end{enumerate}
\end{solution}

\begin{exercise}[微分方程概念辨析]
选择题(四个答案中只有一个是正确的):
\begin{enumerate}
    \item[(1)] 若 $y_1, y_2$ 是非齐次线性方程 $y''+a y'+b y=f(x)$ 的两个特解, 则下面结论中正确的是( ).
    \begin{shortexenum}
        \shortchoice{(A)}{$y_1+y_2$ 是非齐次线性方程的解} &
        \shortchoice{(B)}{$y_1-y_2$ 是非齐次线性方程的解} \\
        \shortchoice{(C)}{$y_1+y_2$ 是 $y''+a y'+b y=0$ 的解} &
        \shortchoice{(D)}{$y_1-y_2$ 是 $y''+a y'+b y=0$ 的解}
    \end{shortexenum}
    \item[(2)] 方程 $(y-\ln x)\mathrm{d}x+x\mathrm{d}y=0$ 是 ( ).
    \begin{shortexenum}
        \shortchoice{(A)}{可分离变量方程} &
        \shortchoice{(B)}{一阶线性非齐次方程} \\
        \shortchoice{(C)}{一阶线性齐次方程} &
        \shortchoice{(D)}{非线性方程}
    \end{shortexenum}
    \item[(3)] 下述函数中( )可能是二阶微分方程 $f(x,y',y'')=0$ 的通解.
    \begin{shortexenum}
        \shortchoice{(A)}{$C_1x+C_2y=0$} &
        \shortchoice{(B)}{$y=C_1e^x+C_2$（按原题等价整理）} \\
        \shortchoice{(C)}{$y=C_1e^x+C_2e^{-2x}$} &
        \shortchoice{(D)}{$y=C_1\ln x+C_2x+C_3$}
    \end{shortexenum}
    \item[(4)] 某种植物的年生长率随着其当前高度和成熟高度与当前高度之差的乘积而变化, 要给这一问题以恰当的数学描述, 应选取参量( ).
    \begin{shortexenum}
        \shortchoice{(A)}{年变化率, 当前高度, 成熟高度} &
        \shortchoice{(B)}{年变化率, 当前高度, 比例常数} \\
        \shortchoice{(C)}{当前高度, 成熟高度} &
        \shortchoice{(D)}{当前高度, 成熟高度, 比例常数}
    \end{shortexenum}
    \item[(5)] 上述问题(4)可描述为( ).
    \begin{shortexenum}
        \shortchoice{(A)}{$y'(t)=y(t)[H(t)-y(t)]$} &
        \shortchoice{(B)}{$y'(t)=y(t)[H(t)-y(t)]$} \\
        \shortchoice{(C)}{$y'(t)=ky(t)[H-y(t)]$} &
        \shortchoice{(D)}{$y'(t)=ky(t)[H(t)-y(t)]$}
    \end{shortexenum}
\end{enumerate}
\end{exercise}
\begin{solution}
\begin{enumerate}
    \item[(1)] \textbf{解析}：记微分算子 $L[y] = y''+a y'+b y$。已知 $L[y_1]=f(x)$ 且 $L[y_2]=f(x)$。
    \begin{enumerate}
        \item 根据线性算子的性质：\(\displaystyle L[y_1-y_2] = L[y_1] - L[y_2] = f(x) - f(x) = 0\)。
        \item 因此 $y_1-y_2$ 是对应齐次方程 $y''+a y'+b y=0$ 的解。故选 \textbf{(D)}。
    \end{enumerate}

    \item[(2)] \textbf{解析}：对原方程进行变形。
    \begin{enumerate}
        \item 将方程两边同除以 $x \mathrm{d}x$，得
        \(\displaystyle \frac{y-\ln x}{x} + \frac{\mathrm{d}y}{\mathrm{d}x} = 0 \implies \frac{\mathrm{d}y}{\mathrm{d}x} + \frac{1}{x}y = \frac{\ln x}{x}\)。
        \item 这是一个标准形式的 $y' + P(x)y = Q(x)$，其中 $P(x)=\frac{1}{x}$，$Q(x)=\frac{\ln x}{x} \neq 0$。
        \item 故其为一阶线性非齐次方程。选 \textbf{(B)}。
    \end{enumerate}

    \item[(3)] \textbf{解析}：分析缺 $y$ 的二阶微分方程的通解结构。
    \begin{enumerate}
        \item 方程 $f(x,y',y'')=0$ 缺变量 $y$。令 $p=y'$，方程降阶为一阶方程 $f(x,p,p')=0$。
        \item 设该一阶方程的通解为 $p = \varphi(x, C_1)$，即 $y' = \varphi(x, C_1)$。
        \item 积分一次得到 \(\displaystyle y = \int \varphi(x, C_1) \mathrm{d}x + C_2\)。
        \item 这说明通解中必须包含两个独立常数，且其中一个常数 $C_2$ 必定以单独相加的形式出现（即形如 $y = F(x, C_1) + C_2$）。
        \item 观察选项：(A) 只有 1 个独立常数；(C) 两个常数都作为函数的系数出现；(D) 含有 3 个常数；而 (B) $y=C_1e^x+C_2$ 恰好符合上述结构，因此选 \textbf{(B)}。
    \end{enumerate}

    \item[(4)] \textbf{解析}：分析题意提取建模所需变量。
    “生长率”是导数，“随着其当前高度和成熟高度与当前高度之差的乘积而变化”意味着导数正比于这个乘积。要将“随……而变化（成正比）”写成等式，必须引入一个比例常数。因此需要的参量有：当前高度 $y(t)$、固定的成熟高度 $H$、以及比例常数 $k$。选 \textbf{(D)}。

    \item[(5)] \textbf{解析}：根据第(4)题的文字描述直接写出方程。
    设当前高度为 $y(t)$，成熟高度为常数 $H$，比例常数为 $k$。生长率为 $y'(t)$。
    乘积为 $y(t)[H-y(t)]$。
    故方程为 $y'(t) = k y(t) [H - y(t)]$。选 \textbf{(C)}。
\end{enumerate}
\end{solution}

\begin{exercise}[人口增长模型]
已知意大利在 1980 年的人口总数为 5700 万. 假定人口的年增长率保持为 $2\%$, 试求 2000 年的人口总数.
\end{exercise}
\begin{solution}
\begin{enumerate}
    \item \textbf{建立微分方程模型}：设 $t$ 年的人口总数为 $P(t)$，以 1980 年为时间起点 $t=0$，此时 $P(0) = 5700$（万人）。
    年增长率为 $2\%$，即人口的变化率与当前人口成正比：\(\displaystyle \frac{\mathrm{d}P}{\mathrm{d}t} = 0.02 P\)。

    \item \textbf{求解微分方程}：分离变量并积分，
    \(\displaystyle \frac{\mathrm{d}P}{P} = 0.02 \mathrm{d}t \implies \ln P = 0.02 t + C \implies P(t) = C_1 e^{0.02 t}\)；
    由 $P(0)=5700$ 得 \(\displaystyle P(t) = 5700 e^{0.02 t}\)。

    \item \textbf{计算 2000 年的人口数}：2000 年对应 $t = 2000 - 1980 = 20$。
    代入方程得 \( \displaystyle P(20)=5700e^{0.02\times20}=5700e^{0.4}\)。
    查表或计算可知 \( \displaystyle e^{0.4}\approx1.4918\)，从而
    \(\displaystyle P(20)\approx5700\times1.4918\approx8503.26\)。
    所以 2000 年的人口总数约为 8503.26 万。
\end{enumerate}
\end{solution}

\begin{exercise}[一阶线性微分方程解的结构]
设 $y_1, y_2, y_3$ 是线性方程 $y'+P(x)y=Q(x)$ 的三个相异特解. 证明 $\frac{y_3-y_1}{y_2-y_1}$ 是常数.
\end{exercise}
\begin{solution}
\begin{enumerate}
    \item \textbf{运用线性方程解的性质}：根据线性微分方程的结构定理，非齐次方程的两个解之差必定是对应齐次方程的解。
    对应齐次方程为 $y' + P(x)y = 0$。

    \item \textbf{构造齐次方程的解}：
    令 $u(x) = y_2 - y_1,\ v(x)=y_3-y_1$，则它们都是 $y' + P(x)y = 0$ 的解，即 \(\displaystyle u' + P(x)u = 0 \implies \frac{u'}{u} = -P(x)\)，从而 \(\displaystyle u(x) = C_1 e^{-\int P(x)\mathrm{d}x},\quad v(x) = C_2 e^{-\int P(x)\mathrm{d}x}\)。
    这里 $C_1, C_2$ 是非零常数（因为解相异）。

    \item \textbf{作比求证}：
    考察两者之比：
    \[
    \frac{y_3 - y_1}{y_2 - y_1}
    = \frac{v(x)}{u(x)}
    = \frac{C_2 e^{-\int P(x)\mathrm{d}x}}{C_1 e^{-\int P(x)\mathrm{d}x}}
    = \frac{C_2}{C_1}.
    \]
    由于 $C_1$ 和 $C_2$ 均为常数，所以其比值也是常数。证明完毕。
\end{enumerate}
\end{solution}

\begin{exercise}[Logistic 模型应用]
在某一人群中推广新技术是通过其中已掌握新技术的人进行的. 设该人群的总人数为 $N$. 在 $t=0$ 时刻已掌握新技术的人数为 $x_0$, 在任意时刻 $t$ 已掌握新技术的人数为 $x(t)$ (将 $x(t)$ 视为连续可微变量), 其变化率与已掌握新技术人数和未掌握新技术人数之积成正比, 比例常数 $k>0$, 求 $x(t)$.
\end{exercise}
\begin{solution}
\begin{enumerate}
    \item \textbf{建立微分方程}：根据题意，已掌握新技术人数为 $x(t)$，未掌握人数为 $N - x(t)$。变化率 $\frac{\mathrm{d}x}{\mathrm{d}t}$ 与两者的乘积成正比：
    \(\displaystyle \frac{\mathrm{d}x}{\mathrm{d}t} = k x(t) [N - x(t)]\)。
    初始条件为 $x(0) = x_0$。

    \item \textbf{分离变量求解}：
    先分离变量并作部分分式分解：
    \(\displaystyle \frac{\mathrm{d}x}{x(N-x)}=k\mathrm{d}t,\quad
    \frac{1}{N}\left(\frac{1}{x}+\frac{1}{N-x}\right)\mathrm{d}x=k\mathrm{d}t\)。
    两边同时积分，得 \(\displaystyle \frac{1}{N} \left( \ln|x| - \ln|N-x| \right) = kt + C
    \implies \ln \left| \frac{x}{N-x} \right| = Nkt + NC\)。
    由于 $0 < x < N$，绝对值可以直接去掉：
    \(\displaystyle \frac{x}{N-x} = C_1 e^{Nkt} \quad (C_1 = e^{NC})\)。

    \item \textbf{代入初值求常数}：
    当 $t=0$ 时，$x(0) = x_0$，故 \(\displaystyle C_1 = \frac{x_0}{N-x_0}\)。

    \item \textbf{化简得出特解}：
    将 $C_1$ 代回并解出 $x(t)$：
    \begin{align*}
    x(t)&=(N-x(t)) \frac{x_0}{N-x_0} e^{Nkt},\\
    x(t)\left(1+\frac{x_0}{N-x_0}e^{Nkt}\right)
    &=N\frac{x_0}{N-x_0}e^{Nkt},\\
    x(t)&=\frac{N x_0 e^{Nkt}}{(N-x_0)+x_0 e^{Nkt}}
    =\frac{N x_0}{x_0+(N-x_0)e^{-Nkt}},
    \end{align*}
    此即典型的 Logistic 曲线方程。
\end{enumerate}
\end{solution}

\begin{exercise}[微分方程的几何应用]
设曲线 $L$ 的极坐标方程为 $r=r(\theta), M(r,\theta)$ 为 $L$ 上任一点, $M_0(2,0)$ 为 $L$ 上一定点. 若极径 $OM_0, OM$ 与曲线 $L$ 构成的曲边扇形面积等于 $L$ 上 $M_0, M$ 两点间弧长值的一半, 试求曲线 $L$ 的方程.
\end{exercise}
\begin{solution}
\begin{enumerate}
    \item \textbf{写出几何量的积分表达式}：
    极坐标下的曲边扇形面积公式为 \(\displaystyle S = \frac{1}{2} \int_0^\theta r^2(\phi) \mathrm{d}\phi\)，弧长为 \(\displaystyle s = \int_0^\theta \sqrt{r^2(\phi) + [r'(\phi)]^2} \mathrm{d}\phi\)。

    \item \textbf{建立并化简积分方程}：
    题设要求 $S = \frac{1}{2} s$，即
    \(\displaystyle \frac{1}{2} \int_0^\theta r^2(\phi) \mathrm{d}\phi = \frac{1}{2} \int_0^\theta \sqrt{r^2(\phi) + [r'(\phi)]^2} \mathrm{d}\phi\)。
    因而
    \(\displaystyle r^2 = \sqrt{r^2 + (r')^2}\)，进一步得
    \(\displaystyle (r')^2 = r^4 - r^2 = r^2(r^2 - 1)\)。

    \item \textbf{求解微分方程}：
    开平方得 $r' = \pm r\sqrt{r^2-1}$。由于 $r \ge 1$，分离变量：
    \(\displaystyle \frac{\mathrm{d}r}{r\sqrt{r^2-1}} = \pm \mathrm{d}\theta\)。
    利用三角换元积分，令 $r = \sec t$（$t \in [0, \pi/2)$），则 $\mathrm{d}r = \sec t \tan t \mathrm{d}t$，$\sqrt{r^2-1} = \tan t$。
    左边积分为 $\int \frac{\sec t \tan t}{\sec t \tan t} \mathrm{d}t = \int \mathrm{d}t = t = \arccos\left(\frac{1}{r}\right)$ （亦可记作 $\text{arcsec}(r)$）。
    所以有 \(\displaystyle \arccos\left(\frac{1}{r}\right) = \pm \theta + C\)。

    \item \textbf{代入已知点定常数}：
    曲线经过点 $M_0(2,0)$，即 $\theta=0$ 时 $r=2$。
    \(\displaystyle \arccos\left(\frac{1}{2}\right) = C \implies C = \frac{\pi}{3}\)。
    所以 $\arccos\left(\frac{1}{r}\right) = \pm \theta + \frac{\pi}{3}$。

    \item \textbf{写出最终方程}：
    取余弦得 \(\displaystyle \frac{1}{r} = \cos\left(\pm \theta + \frac{\pi}{3}\right) = \cos\left(\theta \pm \frac{\pi}{3}\right)\)，
    因而 \(\displaystyle r = \sec\left(\theta \pm \frac{\pi}{3}\right)\)。
\end{enumerate}
\end{solution}

\begin{exercise}[微分方程在物理化学中的应用]
有一种球状的降压血药, 直径为 $0.50 \text{ cm}$. 它在胃中溶解时, 其半径变小的速度与药丸的表面积成正比, 实验室观察表明, 当药丸吞下去两分钟后, 直径从 $0.50 \text{ cm}$ 减少到 $0.38 \text{ cm}$. 问要花多少分钟药丸的直径才会减小到 $0.02 \text{ cm}$ (即几乎完全溶解)? (单位球的表面积为 $4\pi$.)
\end{exercise}
\begin{solution}
\begin{enumerate}
    \item \textbf{提取物理量建立方程}：
    设药丸半径为 $r(t)$。表面积公式为 $S = 4\pi r^2$。
    “半径变小的速度”指的是 $-\frac{\mathrm{d}r}{\mathrm{d}t}$。其与表面积成正比，比例常数为 $\alpha > 0$：
    \(\displaystyle -\frac{\mathrm{d}r}{\mathrm{d}t} = \alpha S = \alpha (4\pi r^2) = K r^2 \quad (\text{令 } K = 4\pi\alpha > 0)\)。

    \item \textbf{分离变量求通解}：
    分离变量并积分，得 \(\displaystyle -\frac{\mathrm{d}r}{r^2} = K \mathrm{d}t \implies \frac{1}{r} = Kt + C\)。

    \item \textbf{代入实验数据求解常数}：
    已知初始时 $t=0$，直径 $D = 0.50$，所以半径 $r_0 = 0.25 = \frac{1}{4}$。
    代入方程得 \(\displaystyle \frac{1}{0.25} = K(0) + C \implies C = 4\)。
    已知 $t=2$ 分钟时，直径 $D = 0.38$，所以半径 $r(2) = 0.19 = \frac{19}{100}$。
    代入方程，得
    \(\displaystyle \frac{1}{0.19}=2K+4 \implies \frac{100}{19}=2K+4 \implies 2K=\frac{24}{19} \implies K=\frac{12}{19}\)，
    从而 \(\displaystyle \frac{1}{r}=\frac{12}{19} t+4\)。

    \item \textbf{计算目标时间}：
    求直径减小到 $D = 0.02 \text{ cm}$，即半径 $r = 0.01 = \frac{1}{100}$ 时的 $t$：
    \(\displaystyle \frac{1}{0.01}=100=\frac{12}{19} t+4 \implies \frac{12}{19} t=96 \implies t=96\times\frac{19}{12}=152\)。
    故大约需要 \textbf{152 分钟}药丸才会减小到 $0.02 \text{ cm}$。
\end{enumerate}
\end{solution}

\begin{exercise}[动力学与微分方程]
物体下落时, 空气阻力可以认为与速度的平方成正比. 设初速为零. 求运动规律.
\end{exercise}
\begin{solution}
\begin{enumerate}
    \item \textbf{建立动力学微分方程}：
    设物体质量为 $m$，下落速度为 $v(t)$，位移为 $x(t)$（向下为正）。
    重力向下，大小为 $mg$；阻力向上，与速度平方成正比，大小为 $kv^2$（$k>0$ 为比例常数）。
    根据牛顿第二定律，\(\displaystyle m \frac{\mathrm{d}v}{\mathrm{d}t} = mg - kv^2\)。
    初始条件为：$v(0) = 0, x(0) = 0$。

    \item \textbf{求速度规律 $v(t)$}：
    提取公因子化简方程，令 $a^2 = \frac{mg}{k}$，则
    \[
    \frac{\mathrm{d}v}{\mathrm{d}t}
    =g-\frac{k}{m}v^2
    =\frac{k}{m}\left(\frac{mg}{k}-v^2\right)
    =\frac{k}{m}(a^2-v^2).
    \]
    分离变量积分：
    \begin{align*}
    \frac{\mathrm{d}v}{a^2-v^2}&=\frac{k}{m}\,\mathrm{d}t,\\
    \frac{1}{2a}\int\left(\frac{1}{a+v}+\frac{1}{a-v}\right)\mathrm{d}v
    &=\int\frac{k}{m}\,\mathrm{d}t,\\
    \frac{1}{2a}\ln\left|\frac{a+v}{a-v}\right|
    &=\frac{k}{m}t+C_1.
    \end{align*}
    由于物体刚开始下落，速度小于极限速度 $a$，所以 $\frac{a+v}{a-v} > 0$。代入 $v(0)=0$，得 $C_1 = 0$。
    解出 $v(t)$：
    \begin{align*}
    \ln \frac{a+v}{a-v}&=\frac{2ak}{m}t,\\
    \frac{a+v}{a-v}&=e^{\frac{2ak}{m}t},\\
    v(t)&=a\frac{e^{\frac{2ak}{m}t}-1}{e^{\frac{2ak}{m}t}+1}
    =\sqrt{\frac{mg}{k}}\tanh\left(\sqrt{\frac{kg}{m}}\,t\right).
    \end{align*}

    \item \textbf{求位移规律 $x(t)$}：
    由于 $v(t) = \frac{\mathrm{d}x}{\mathrm{d}t}$，对其进行积分。利用公式 $\int \tanh(u) \mathrm{d}u = \ln(\cosh u)$，令 $u = \sqrt{\frac{kg}{m}} \tau$，得
    \[
    x(t)=\int_0^t \sqrt{\frac{mg}{k}}
    \tanh\left(\sqrt{\frac{kg}{m}}\,\tau\right)\mathrm{d}\tau
    =\frac{m}{k}\ln\left[\cosh\left(\sqrt{\frac{kg}{m}}\,t\right)\right].
    \]
    这就是物体下落的运动位移规律。
\end{enumerate}
\end{solution}

\begin{exercise}[常系数线性非齐次方程的特解与通解]
设二阶常系数线性微分方程 $y''+\alpha y'+\beta y=\gamma e^x$ 的一个特解为 $y=e^{2x}+(1+x)e^x$, 试确定常数 $\alpha, \beta, \gamma$ 并求该方程的通解.
\end{exercise}
\begin{solution}
\begin{enumerate}
    \item \textbf{理解解的结构}：
    给定的特解 $y^* = e^{2x} + e^x + x e^x$ 是一个多项式和指数函数的组合。根据非齐次线性微分方程解的叠加原理，它必由两部分组成：对应齐次方程的解和对应非齐次项 $\gamma e^x$ 的特解。
    由于等号右侧只有 $e^x$ 这一项，特解中出现了 $x e^x$，说明 $\lambda=1$ 必定是特征方程的根，此时发生了共振，产生的特解形式为 $A x e^x$。
    同时出现的 $e^{2x}$ 和独立的 $e^x$ 必然是齐次方程 $y''+\alpha y'+\beta y=0$ 的解，即特征根为 $\lambda_1=2, \lambda_2=1$。

    \item \textbf{确定 $\alpha$ 和 $\beta$}：
    已知特征根为 1 和 2，特征方程为
    \(\displaystyle (\lambda-1)(\lambda-2) = 0 \implies \lambda^2 - 3\lambda + 2 = 0,
    \qquad \alpha = -3,\ \beta = 2\)。

    \item \textbf{确定 $\gamma$}：
    现在方程的形式为 $y'' - 3y' + 2y = \gamma e^x$。
    把原题给的整个解 $y = e^{2x} + (1+x)e^x$ 代入该方程验证并求 $\gamma$。
    求导得
    \(\displaystyle y' = 2e^{2x} + e^x + (1+x)e^x = 2e^{2x} + (2+x)e^x\)，
    \(\displaystyle y'' = 4e^{2x} + e^x + (2+x)e^x = 4e^{2x} + (3+x)e^x\)。
    代入方程左边：
    \(\displaystyle y'' - 3y' + 2y = [4e^{2x} + (3+x)e^x] - 3[2e^{2x} + (2+x)e^x] + 2[e^{2x} + (1+x)e^x]\)。
    比较系数可得：$e^{2x}$ 项系数 $4 - 6 + 2 = 0$，$x e^x$ 项系数 $1 - 3 + 2 = 0$，$e^x$ 项系数 $3 - 6 + 2 = -1$。
    所以方程左边算出来等于 $-e^x$。对应右侧 $\gamma e^x$，得 \(\displaystyle \gamma = -1\)。

    \item \textbf{写出通解}：
    方程为 $y'' - 3y' + 2y = -e^x$。
    其齐次通解为 $C_1 e^{2x} + C_2 e^x$。通过观察特解部分，非齐次特解的共振升幂项为 $-x e^x$（把 $(1+x)e^x$ 拆开，$e^x$ 吸收进齐次解）。
    因此，通解为 \(\displaystyle y = C_1 e^{2x} + C_2 e^x - x e^x\)。
\end{enumerate}
\end{solution}

\begin{exercise}[二阶齐次方程的初值问题]
求方程 $y''-4y'+3y=0$ 的积分曲线, 使其在点 $(0,2)$ 处与直线 $2x-2y+9=0$ 相切.
\end{exercise}
\begin{solution}
\begin{enumerate}
    \item \textbf{将几何条件翻译为初始条件}：
    积分曲线经过点 $(0,2)$，说明当 $x=0$ 时有 \(\displaystyle y(0) = 2\)。
    在点 $(0,2)$ 处与直线 $2x-2y+9=0$ 相切。该直线方程可化为 $y = x + 4.5$，斜率为 $1$。相切意味着在该点处曲线的导数等于直线的斜率：
    \(\displaystyle y'(0) = 1\)。

    \item \textbf{求解齐次微分方程通解}：
    特征方程为 $\lambda^2 - 4\lambda + 3 = 0 \implies (\lambda-1)(\lambda-3) = 0$。
    特征根为 $\lambda_1 = 1, \lambda_2 = 3$。
    通解为 \(\displaystyle y(x) = C_1 e^x + C_2 e^{3x}\)。

    \item \textbf{代入初值求解待定常数}：
    求导得：$y'(x) = C_1 e^x + 3C_2 e^{3x}$。
    根据初始条件有 \(\displaystyle C_1 + C_2 = 2,\ C_1 + 3C_2 = 1\)。
    两式相减得 \(\displaystyle 2C_2 = -1 \implies C_2 = -\frac{1}{2}\)，再代回第一式得
    \(\displaystyle C_1 = 2 - \left(-\frac{1}{2}\right) = \frac{5}{2}\)。

    \item \textbf{得出积分曲线方程}：\(\displaystyle y = \frac{5}{2} e^x - \frac{1}{2} e^{3x}\)。
\end{enumerate}
\end{solution}

\begin{exercise}[自变量与因变量互换求解]
求方程 $y''+(4x+e^{2y})(y')^3=0$ 的通解.
\end{exercise}
\begin{solution}
\begin{enumerate}
    \item \textbf{观察方程结构进行变量代换}：
    方程中 $y$ 和 $y'$ 的关系非线性，但如果把 $y$ 视为自变量，$x$ 视为未知函数，将导数反转可能可以将其线性化。
    由复合函数反函数求导法则，
    \(\displaystyle y' = \frac{\mathrm{d}y}{\mathrm{d}x} = \frac{1}{\frac{\mathrm{d}x}{\mathrm{d}y}} = \frac{1}{x'}\)；
    再次对 $x$ 求导得到
    \(\displaystyle y'' = \frac{\mathrm{d}}{\mathrm{d}x}(y') = \frac{\mathrm{d}}{\mathrm{d}y}\left(\frac{1}{x'}\right) \cdot \frac{\mathrm{d}y}{\mathrm{d}x} = -\frac{x''}{(x')^2} \cdot \frac{1}{x'} = -\frac{x''}{(x')^3}\)。

    \item \textbf{代入原方程化简}：
    将 $y'$ 和 $y''$ 代入原方程 $y''+(4x+e^{2y})(y')^3=0$，得
    \(\displaystyle -\frac{x''}{(x')^3} + (4x+e^{2y}) \left(\frac{1}{x'}\right)^3 = 0\)。
    当 $y' \neq 0$ 时（排除了常数解 $y=C$），两边同乘 $(x')^3$：
    \(\displaystyle -x'' + 4x + e^{2y} = 0 \implies x'' - 4x = e^{2y}\)。
    现在这是一个以 $y$ 为自变量、$x$ 为因变量的二阶常系数线性非齐次微分方程。

    \item \textbf{求解新的常系数线性方程}：
    齐次方程 $x'' - 4x = 0$ 的特征方程为 $\lambda^2 - 4 = 0 \implies \lambda = \pm 2$。
    齐次通解为：$x_h = C_1 e^{2y} + C_2 e^{-2y}$。
    对于非齐次项 $e^{2y}$，其特征指数为 2，恰好与特征根重合，发生共振。设特解为 $x^* = A y e^{2y}$。
    求导得 \(\displaystyle (x^*)' = A(1+2y)e^{2y}, \quad (x^*)'' = A(4+4y)e^{2y}\)。
    代入 $x'' - 4x = e^{2y}$，得
    \(\displaystyle A(4+4y)e^{2y} - 4A y e^{2y} = 4A e^{2y} = e^{2y} \implies 4A = 1 \implies A = \frac{1}{4}\)。
    因此特解为 $x^* = \frac{1}{4} y e^{2y}$。

    \item \textbf{写出通解}：所以 $x(y)$ 的通解为
    \(\displaystyle x = C_1 e^{2y} + C_2 e^{-2y} + \frac{1}{4} y e^{2y}\)。
\end{enumerate}
\end{solution}

\begin{exercise}[一阶线性方程与积分不等式]
设 $f(x)$ 为连续函数,
\begin{shortsingleenum}
    \shortsinglechoice{(1)}{求初值问题 $y'+ay=f(x), y(0)=0$ 的解 $y(x)$, 其中 $a>0$ 为常数;}
    \shortsinglechoice{(2)}{若 $|f(x)|\leqslant  K$ ($K$ 为常数), 证明当 $x\ge 0$ 时, 有 $|y(x)|\leqslant  \frac{K}{a}(1-e^{-ax})$.}
\end{shortsingleenum}
\end{exercise}
\begin{solution}
\begin{enumerate}
    \item[(1)] \textbf{求解一阶线性方程}：
    原方程为标准的一阶线性常微分方程。利用积分因子法，积分因子为 $e^{\int a \mathrm{d}x} = e^{ax}$。
    方程两边同乘 $e^{ax}$，得
    \(\displaystyle y'e^{ax} + a y e^{ax} = f(x)e^{ax} \implies \frac{\mathrm{d}}{\mathrm{d}x} (y e^{ax}) = f(x)e^{ax}\)。
    从 $0$ 积分到 $x$，得
    \(\displaystyle y(x)e^{ax} - y(0)e^0 = \int_0^x f(t)e^{at} \mathrm{d}t\)。
    利用初始条件 $y(0) = 0$，得到解的积分表达式
    \(\displaystyle y(x) = e^{-ax} \int_0^x f(t)e^{at} \mathrm{d}t\)。

    \item[(2)] \textbf{放缩法证明不等式}：
    对(1)中求得的 $y(x)$ 取绝对值，并利用积分绝对值不等式
    \(\displaystyle \left| \int g(t) \,\mathrm{d}t \right| \leqslant \int |g(t)| \,\mathrm{d}t\)
    （因为 $x \ge 0$ 积分上限大于下限），得
    \(\displaystyle |y(x)|
    \leqslant e^{-ax} \int_0^x |f(t)| e^{at} \mathrm{d}t
    \leqslant e^{-ax} \int_0^x K e^{at} \mathrm{d}t
    = \frac{K}{a} (1 - e^{-ax})\)。
    证明完毕。
\end{enumerate}
\end{solution}

\begin{exercise}[微分不等式比较定理]
设函数 $u_1(x)$ 和 $u_2(x)$ 分别满足
\(\displaystyle u_1' = a(x)u_1+v(x), \quad u_1(0)=C\)
以及
\(\displaystyle u_2' \leqslant  a(x)u_2+v(x), \quad u_2(0)=C\)，
其中, $a(x),v(x)$ 在 $x\ge 0$ 上连续, $C$ 为常数. 证明不等式
\(\displaystyle u_2(x) \leqslant  u_1(x)\ (x\ge 0)\).
\end{exercise}
\begin{solution}
\begin{enumerate}
    \item \textbf{构造差函数}：
    令 $w(x) = u_1(x) - u_2(x)$。要证当 $x \ge 0$ 时 $u_2(x) \leqslant  u_1(x)$，即证明 $w(x) \ge 0$。
    \item \textbf{建立差函数满足的微分不等式}：
    对 $w(x)$ 求导，并利用已知条件得
    \(\displaystyle w'(x) = u_1'(x) - u_2'(x) \ge [a(x)u_1(x) + v(x)] - [a(x)u_2(x) + v(x)] = a(x)[u_1(x) - u_2(x)]\)，
    于是 $w(x)$ 满足一阶线性微分不等式 \(\displaystyle w'(x) - a(x)w(x) \ge 0\)。
    初始条件为：$w(0) = u_1(0) - u_2(0) = C - C = 0$。
    \item \textbf{利用积分因子法求解}：
    两边同乘正数积分因子 $\exp\left(-\int_0^x a(t)\mathrm{d}t\right)$，得
    \(\displaystyle [w'(x) - a(x)w(x)] e^{-\int_0^x a(t)\mathrm{d}t} \ge 0
    \implies \frac{\mathrm{d}}{\mathrm{d}x} \left[ w(x) e^{-\int_0^x a(t)\mathrm{d}t} \right] \ge 0\)。
    这说明函数 $H(x) = w(x) e^{-\int_0^x a(t)\mathrm{d}t}$ 在 $x \ge 0$ 上是单调不减的。
    \item \textbf{得出结论}：
    由于 $H(x)$ 是单调不减函数，所以对任意 $x \ge 0$，有
    \(\displaystyle H(x) \ge H(0) = w(0) e^0 = 0 \implies w(x) e^{-\int_0^x a(t)\mathrm{d}t} \ge 0\)。
    由于指数函数永远大于零，因此必有 $w(x) \ge 0$，即 $u_1(x) \ge u_2(x)$ 对于 $x \ge 0$ 恒成立。证明完毕。
\end{enumerate}
\end{solution}

\begin{exercise}[二阶微分方程解的单调性分析]
给定方程 \(\displaystyle y''-a(x)y = 0,\qquad y(0)=y_0,\quad y'(0)=y_0'\)（1），
其中, $a(x)\in C[0, +\infty)$, 且 $a(x)>0$. 设 $y(x)$ 是方程(1) 的解, 其中, $y_0>0, y_0'>0$. 证明: 当 $x\ge 0$ 时, 函数 $y(x)y'(x)$ 及 $y(x)$ 都是递增的正函数.
\end{exercise}
\begin{solution}
\begin{enumerate}
    \item \textbf{引入能量函数考察 $y(x)y'(x)$ 的单调性}：
    记待证函数 $E(x) = y(x)y'(x)$。我们要证明 $E(x)$ 递增，即证 $E'(x) > 0$。
    对其求导。由方程(1)知 \( \displaystyle y''(x)=a(x)y(x)\)，于是
    \(\displaystyle E'(x)=\frac{\mathrm{d}}{\mathrm{d}x}[y(x)y'(x)]=[y'(x)]^2+y(x)y''(x)=[y'(x)]^2+a(x)[y(x)]^2\)。
    已知 $a(x) > 0$。只要 $y$ 和 $y'$ 不全为零，上式必大于 0。

    \item \textbf{分析初始时刻}：
    当 $x=0$ 时，
    \(\displaystyle E(0) = y(0)y'(0) = y_0 y_0' > 0 \quad (\text{因 } y_0>0, y_0'>0)\)。
    此时 $E'(0) = (y_0')^2 + a(0)y_0^2 > 0$。由于 $E'(x)$ 严格大于 0，$E(x)$ 严格递增。
    因此对于所有 $x \ge 0$，有 $E(x) \ge E(0) > 0$，即
    \(\displaystyle y(x)y'(x) > 0 \quad (x \ge 0)\)。

    \item \textbf{证明 $y(x)$ 和 $y'(x)$ 的符号保持正}：
    已知 $y(0) > 0$，若存在某点 $y(x_1) \leqslant  0$，根据连续性，必然存在一个点 $x^*$ 使得 $y(x^*) = 0$。
    但在 $x^*$ 处，$E(x^*) = y(x^*)y'(x^*) = 0 \cdot y'(x^*) = 0$。
    这与 $E(x) > 0$ 恒成立相矛盾。
    因此，对于所有 $x \ge 0$，必然有 $y(x) > 0$。
    又因为 $y(x)y'(x) > 0$，所以 $y'(x) > 0$ 恒成立。

    \item \textbf{总结单调性结论}：
    由于 $y'(x) > 0$，说明函数 $y(x)$ 是严格递增的正函数。
    同时，由于 $E'(x) = [y'(x)]^2 + a(x)[y(x)]^2 > 0$ 恒成立，说明 $y(x)y'(x)$ 也是严格递增的正函数。证明完毕。
\end{enumerate}
\end{solution}

\begin{exercise}[由函数方程导出微分方程]
设连续函数 $f(x)$ 具有性质 $f(x+y)=\frac{f(x)+f(y)}{1-f(x)f(y)}$, 且 $f'(0)$ 存在. 试导出 $f(x)$ 所满足的微分方程, 并求出 $f(x)$.
\end{exercise}
\begin{solution}
\begin{enumerate}
    \item \textbf{求特殊点 $f(0)$ 的值}：
    在恒等式中令 $x = 0, y = 0$：
    \(\displaystyle f(0) = \frac{2f(0)}{1-f^2(0)} \implies f(0) - f^3(0) = 2f(0) \implies f^3(0) + f(0) = 0 \implies f(0)[f^2(0)+1] = 0\)。
    因为 $f(x)$ 为实值函数，解得 $f(0) = 0$。

    \item \textbf{利用偏导数导出微分方程}：
    将恒等式 $f(x+y) = \frac{f(x)+f(y)}{1-f(x)f(y)}$ 两边对变量 $y$ 求导（将 $x$ 视为常数）：
    \(\displaystyle f'(x+y) \cdot 1 = \frac{f'(y)[1-f(x)f(y)] - [f(x)+f(y)][-f(x)f'(y)]}{[1-f(x)f(y)]^2}\)。
    提取分子中公因子 $f'(y)$：
    \(\displaystyle f'(x+y) = f'(y) \frac{1 - f(x)f(y) + f^2(x) + f(x)f(y)}{[1-f(x)f(y)]^2} = f'(y) \frac{1 + f^2(x)}{[1-f(x)f(y)]^2}\)。
    现在，在上述等式中令 $y = 0$，利用我们求出的 $f(0) = 0$，得
    \(\displaystyle f'(x) = f'(0) \frac{1 + f^2(x)}{[1-0]^2} = f'(0) [1 + f^2(x)]\)。
    设已知常数 $f'(0) = k$，则导出的微分方程为 \(\displaystyle f'(x) = k [1 + f^2(x)]\)。

    \item \textbf{求解微分方程}：
    这是一个可分离变量的方程。分离变量并积分，再代入初始条件 $f(0) = 0$，得
    \(\displaystyle \frac{\mathrm{d}f}{1 + f^2} = k \mathrm{d}x \implies \arctan f(x) = kx + C\)；
    又由 \(\displaystyle \arctan(0) = 0 + C\) 得 $C = 0$。
    所以 $\arctan f(x) = kx$。两边取正切即得
    \(\displaystyle f(x) = \tan(kx)\ (\text{其中 } k = f'(0))\)。
\end{enumerate}
\end{solution}

\begin{exercise}[积微分方程求解]
设微分方程的初值问题
\(\displaystyle \frac{\mathrm{d}u(t)}{\mathrm{d}t} = u(t) + \int_0^t u(s)\mathrm{d}s,\quad u(0)=1\)
有唯一解, 求 $u(t)$.
\end{exercise}
\begin{solution}
\begin{enumerate}
    \item \textbf{将积微分方程转化为二阶常微分方程}：
    对原方程两边关于 $t$ 求导，利用变上限积分的求导法则得
    \(\displaystyle u''(t) = u'(t) + u(t) \implies u'' - u' - u = 0\)。
    这就转化成了一个二阶常系数线性齐次微分方程。

    \item \textbf{确定微分方程的初值条件}：
    题干已给出一个条件：$u(0) = 1$。
    我们需要找第二个条件 $u'(0)$。将 $t=0$ 代入原积微分方程，得
    \(\displaystyle u'(0) = u(0) + \int_0^0 u(s)\mathrm{d}s = 1 + 0 = 1\)。
    所以初值条件为：$u(0) = 1, u'(0) = 1$。

    \item \textbf{求解微分方程}：
    特征方程为 $r^2 - r - 1 = 0$。由求根公式得
    \(\displaystyle r = \frac{1 \pm \sqrt{(-1)^2 - 4(1)(-1)}}{2} = \frac{1 \pm \sqrt{5}}{2}\)。
    设 $r_1 = \frac{1+\sqrt{5}}{2}, r_2 = \frac{1-\sqrt{5}}{2}$，则通解为
    \(\displaystyle u(t) = C_1 e^{r_1 t} + C_2 e^{r_2 t}\)。

    \item \textbf{代入初值求解常数}：
    由初值条件得到 \(\displaystyle C_1 + C_2 = 1,\quad C_1 r_1 + C_2 r_2 = 1\)。
    将 $C_2 = 1 - C_1$ 代入第二式：
    即 \(\displaystyle C_1 r_1 + (1 - C_1) r_2 = 1 \implies C_1(r_1 - r_2) = 1 - r_2 \implies C_1 = \frac{1 - r_2}{r_1 - r_2}\)。
    因为 $r_1 - r_2 = \sqrt{5}$，且 $1 - r_2 = 1 - \frac{1-\sqrt{5}}{2} = \frac{1+\sqrt{5}}{2}$，所以
    \(\displaystyle C_1 = \frac{\frac{1+\sqrt{5}}{2}}{\sqrt{5}} = \frac{5+\sqrt{5}}{10}\)；
    同理
    \(\displaystyle C_2 = 1 - \frac{5+\sqrt{5}}{10} = \frac{5-\sqrt{5}}{10}\)。

    \item \textbf{得出最终解}：
    \(\displaystyle u(t) = \frac{5+\sqrt{5}}{10} e^{\frac{1+\sqrt{5}}{2}t} + \frac{5-\sqrt{5}}{10} e^{\frac{1-\sqrt{5}}{2}t}\)。
\end{enumerate}
\end{solution}

\begin{exercise}[变上限积分方程求解]
求满足条件 $\int_0^x tf(t)\mathrm{d}t = nf(x)-1$ ($n>0, n\neq 1$) 的连续函数 $f(x)$.
\end{exercise}
\begin{solution}
\begin{enumerate}
    \item \textbf{对方程两边求导构造微分方程}：
    原方程 $\int_0^x tf(t)\mathrm{d}t = nf(x)-1$ 两边关于 $x$ 求导，整理得
    \(\displaystyle x f(x) = n f'(x) \implies f'(x) - \frac{x}{n} f(x) = 0\)。

    \item \textbf{求解一阶微分方程}：
    分离变量并积分，得
    \(\displaystyle \frac{\mathrm{d}f}{f} = \frac{x}{n} \mathrm{d}x \implies \ln|f| = \frac{x^2}{2n} + C \implies f(x) = C_1 e^{\frac{x^2}{2n}}\)。

    \item \textbf{利用原积分方程提取初始条件}：
    将 $x=0$ 代入原方程，得 \(\displaystyle f(0) = \frac{1}{n}\)；
    再由 \(\displaystyle f(0) = C_1 e^0 = C_1\) 得 \(C_1 = \frac{1}{n}\)。

    \item \textbf{得出最终答案}：
    因此该连续函数为 \(\displaystyle f(x) = \frac{1}{n} e^{\frac{x^2}{2n}}\)。
\end{enumerate}
\end{solution}

\begin{exercise}[含卷积的积分方程求解]
设 $f(x)=\sin x - \int_0^x (x-t)f(t)\mathrm{d}t$, 其中 $f(x)$ 为连续函数. 求 $f(x)$.
\end{exercise}
\begin{solution}
\begin{enumerate}
    \item \textbf{拆分积分项以备求导}：
    变上限积分中含有外部变量 $x$，求导前必须先把它拆分开。
    即 \(\displaystyle f(x) = \sin x - \left( x\int_0^x f(t)\mathrm{d}t - \int_0^x tf(t)\mathrm{d}t \right)
    = \sin x - x\int_0^x f(t)\mathrm{d}t + \int_0^x tf(t)\mathrm{d}t\)。

    \item \textbf{进行第一轮求导}：
    对方程两边关于 $x$ 求导。利用乘积法则和变上限积分求导法则，并化简得
    即 \(\displaystyle f'(x) = \cos x - \left( 1 \cdot \int_0^x f(t)\mathrm{d}t + x f(x) \right) + xf(x) = \cos x - \int_0^x f(t)\mathrm{d}t\)。

    \item \textbf{进行第二轮求导}：
    对上述结果再次关于 $x$ 求导，得 \(\displaystyle f'' = -\sin x - f\)。
    因而 \(\displaystyle f'' + f = -\sin x\)。

    \item \textbf{确定初值条件}：
    需要两个条件 $f(0)$ 和 $f'(0)$。回到求导前原式：
    将 $x=0$ 代入 $f(x) = \sin x - \int_0^x (x-t)f(t)\mathrm{d}t$，积分上限为 0，得
    \(\displaystyle f(0) = \sin 0 - 0 = 0\)；
    再将 $x=0$ 代入一次导数式 $f'(x) = \cos x - \int_0^x f(t)\mathrm{d}t$，得
    \(\displaystyle f'(0) = \cos 0 - 0 = 1\)。

    \item \textbf{求解非齐次微分方程}：
    对应的齐次方程 $f'' + f = 0$ 的特征根为 $\pm i$，故齐次通解为 $C_1\cos x + C_2\sin x$。
    对于非齐次项 $-\sin x$，其频率与特征根吻合，发生共振。设特解为 $f^* = x(A\cos x + B\sin x)$。
    求二阶导数并代入方程 $f'' + f = -\sin x$，得
    \(\displaystyle (f^*)'' = -x(A\cos x + B\sin x) - 2A\sin x + 2B\cos x\)，
    从而 \(\displaystyle -2A\sin x + 2B\cos x = -\sin x \implies A = \frac{1}{2},\ B = 0\)。
    特解为 $f^* = \frac{1}{2}x\cos x$。
    通解为 \(\displaystyle f(x) = C_1\cos x + C_2\sin x + \frac{1}{2}x\cos x\)。

    \item \textbf{代入初值}：
    由 $f(0) = 0 \implies C_1 = 0$。此时 $f(x) = C_2\sin x + \frac{1}{2}x\cos x$。
    求导：$f'(x) = C_2\cos x + \frac{1}{2}\cos x - \frac{1}{2}x\sin x$。
    由 $f'(0) = 1 \implies C_2 + \frac{1}{2} = 1 \implies C_2 = \frac{1}{2}$。
    最终得出函数 $f(x)$：
    \(\displaystyle \boxed{f(x) = \frac{1}{2}\sin x + \frac{1}{2}x\cos x}\)。
\end{enumerate}
\end{solution}

\chapter{多元函数微分学}

\section{知识讲解}

本章知识讲解依次介绍对象识别、极限连续、一阶微分、复合求导、梯度几何、隐函数、泰勒与极值、偏微分方程。各专题围绕极限、连续、求导、几何量、隐函数、极值和变量替换等内容展开，并在定义、定理、证明与例题之间保持对应关系。

\subsection{\texorpdfstring{$\mathbb{R}^n$、点集与区域}{Rn、点集与区域}}

\begin{definition}[$n$ 维实坐标空间] \label{def:rn-set}
称
\[
\boxed{\;
\mathbb R^n:=\bigl\{(x_1,x_2,\dots,x_n)\mid x_i\in\mathbb R,\;i=1,2,\dots,n\bigr\}\;}
\]
为 \textbf{$n$ 维实坐标空间}（或简称 \textbf{$n$ 维空间}），其中每个元素
$x=(x_1,x_2,\dots,x_n)$ 称为一个\textbf{点}（Point），也称为一个\textbf{向量}（Vector）；
$x_i$ 称为 $x$ 的第 $i$ 个\textbf{分量}（Component）或\textbf{坐标}（Coordinate）。
特别地，$\boldsymbol 0:=(0,0,\dots,0)$ 称为\textbf{零向量}（Zero Vector）或\textbf{原点}（Origin）。
\end{definition}

\begin{definition}[内积 / Inner Product] \label{def:inner-product}
对任意 $x=(x_1,\dots,x_n),\;y=(y_1,\dots,y_n)\in\mathbb R^n$，定义
\(\displaystyle \boxed{\;\langle x,y\rangle:=\sum_{i=1}^{n}x_i y_i\;}\)
为 $x$ 与 $y$ 的\textbf{内积}（或称为\textbf{点积} Dot Product、\textbf{标量积} Scalar Product）。
\end{definition}


\begin{definition}[$n$ 维欧氏空间 / Euclidean Space] \label{def:euclidean-space}
赋定了内积 $\langle\cdot,\cdot\rangle$ 的线性空间 $\mathbb R^n$ 称为\textbf{$n$ 维欧几里得空间}（Euclidean Space），
记作 $(\mathbb R^n,\langle\cdot,\cdot\rangle)$ 或简记为 $\mathbb E^n$。
\end{definition}

\begin{definition}[范数 / Norm] \label{def:norm}
对 $x\in\mathbb R^n$，称
\(\displaystyle \boxed{\;\|x\|:=\sqrt{\langle x,x\rangle} = \sqrt{\sum_{i=1}^{n}x_i^2}\;}\)
为 $x$ 的\textbf{范数}（Norm）（或称为\textbf{模长} Length、\textbf{模} Magnitude）。
\end{definition}

当 $n=2$ 时，$\|(x,y)\|=\sqrt{x^2+y^2}$ 正是平面解析几何中熟知的 ``原点到点 $(x,y)$ 的距离''。

\begin{property}[范数的基本性质] \label{prop:norm-properties}
\begin{enumerate}
  \item\textbf{(非负性)}: $\|x\|\geqslant 0$，且 $\|x\|=0\Longleftrightarrow x=\boldsymbol 0$；
  \item\textbf{(正齐次性)}: $\forall\lambda\in\mathbb R,\;\|\lambda x\|=|\lambda|\cdot\|x\|$；
  \item\textbf{(三角不等式)}: $\|x+y\|\leqslant\|x\|+\|y\|$，对所有 $x,y\in\mathbb R^n$ 成立。
\end{enumerate}
\end{property}

\begin{theorem}[柯西--施瓦茨不等式 / Cauchy--Schwarz Inequality] \label{thm:cauchy-schwarz}
对任意 $x,y\in\mathbb R^n$，有 \(\displaystyle \boxed{\;|\langle x,y\rangle|\leqslant\|x\|\cdot\|y\|\;}\)。等号成立当且仅当 $x$ 与 $y$ \textbf{线性相关}（即其中一个为零向量，或者一个是另一个的非零常数倍）。
\end{theorem}

\begin{proof}
若 $x=\boldsymbol 0$，则两边均为 $0$，等号显然成立。

下设 $x\neq\boldsymbol 0$。考虑关于实变量 $t$ 的二次函数
\(\displaystyle f(t):=\langle tx+y,\,tx+y\rangle\geqslant 0\)，
其中非负性来自内积的恒正性。展开得
\(\displaystyle f(t)=t^2\langle x,x\rangle+2t\langle x,y\rangle+\langle y,y\rangle\)。
这是 $t$ 的二次多项式（首项系数 $\langle x,x\rangle=\|x\|^2>0$），且对所有 $t\in\mathbb R$ 取值非负。
故其判别式必非正：\(\displaystyle \Delta=4\langle x,y\rangle^2-4\langle x,x\rangle\langle y,y\rangle\leqslant  0\)，即 $\langle x,y\rangle^2\leqslant\|x\|^2\|y\|^2$。开方即得欲证不等式。

等号成立条件分析：$\Delta=0$ 当且仅当 $f(t)=0$ 有唯一实根 $t_0$。此时 $\langle t_0x+y,t_0x+y\rangle=0 \implies t_0x+y=\boldsymbol 0$，即 $x$ 与 $y$ 线性相关。
\end{proof}

\begin{theorem}[三角不等式 / Triangle Inequality] \label{thm:triangle-ineq}
对任意 $x,y\in\mathbb R^n$：$\|x+y\|\leqslant\|x\|+\|y\|$。
\end{theorem}

\begin{proof}
\begin{align*}
\|x+y\|^2 &= \langle x+y,x+y\rangle = \|x\|^2+2\langle x,y\rangle+\|y\|^2 \\
&\leqslant \|x\|^2+2|\langle x,y\rangle|+\|y\|^2 \leqslant \|x\|^2+2\|x\|\|y\|+\|y\|^2 \\
&= (\|x\|+\|y\|)^2 .
\end{align*}
两边开平方根即得证。
\end{proof}

\begin{exercise}[反三角不等式 / Reverse Triangle Inequality] \label{ex:norm-reverse-triangle}
证明：$\bigl|\|x\|-\|y\|\bigr|\leqslant\|x-y\|$ 对所有 $x,y\in\mathbb R^n$ 成立。
\end{exercise}
\begin{proof}
待证不等式等价于双边估计 $-\|x-y\|\leqslant\|x\|-\|y\|\leqslant\|x-y\|$。
\begin{itemize}
\item 右边不等式：$\|x\|=\|x-y+y\|\leqslant\|x-y\|+\|y\| \implies \|x\|-\|y\|\leqslant\|x-y\|$。
\item 左边不等式：将 $x$ 与 $y$ 互换，同理可得 $\|y\|-\|x\|\leqslant\|y-x\|=\|x-y\|$。
\end{itemize}
合并两边即得 $\bigl|\|x\|-\|y\|\bigr|\leqslant\|x-y\|$。
\end{proof}

\begin{definition}[向量夹角与正交] \label{def:vector-angle} \label{def:orthogonal}
设 $x,y\in\mathbb R^n$ 均为非零向量。存在唯一的 $\theta\in[0,\pi]$ 使得 \(\displaystyle \cos\theta=\frac{\langle x,y\rangle}{\|x\|\|y\|}\)。称 $\theta$ 为向量 $x$ 与 $y$ 的\textbf{夹角}（Angle）。
若 $\langle x,y\rangle=0$（等价于夹角 $\theta=\dfrac{\pi}{2}$），则称 $x$ 与 $y$ \textbf{正交}（Orthogonal）或\textbf{垂直}，记作 $x\perp y$。
\end{definition}


\begin{definition}[标准基 / Standard Basis] \label{def:standard-basis}
取 $\mathbb R^n$ 中如下 $n$ 个特殊向量：\(\displaystyle e_1=(1,0,0,\dots,0)\)，\(\displaystyle e_2=(0,1,0,\dots,0)\)，\(\dots\)，\(\displaystyle e_n=(0,0,\dots,0,1)\)。称 $\{e_1,e_2,\dots,e_n\}$ 为 $\mathbb R^n$ 的\textbf{标准基}。它们满足 $\langle e_i,e_j\rangle=\delta_{ij}$（克罗内克 delta 符号），即两两正交且模长为 1。
\end{definition}

至此，我们从纯集合出发，依次赋予了 $\mathbb R^n$ 三层结构：
\[
\underbrace{\mathbb R^n}_{\text{集合}} \;\xrightarrow{\;\text{加法}+\text{数乘}\;}\; \underbrace{\text{线性空间}}_{\text{代数}} \;\xrightarrow{\;\text{内积}\;}\; \underbrace{\text{欧氏空间}}_{\text{度量}}.
\]

\begin{definition}[开球与闭球 / Open Ball \& Closed Ball] \label{def:neighborhood}
设 $p_0\in\mathbb R^n$，$r>0$。
\textbf{开球}（Open Ball）：称 \(\displaystyle B_r(p_0):=\bigl\{p\in\mathbb R^n\mid\|p-p_0\|<r\bigr\}\) 为以 $p_0$ 为中心、$r$ 为半径的\textbf{$r$ 邻域}，也叫\textbf{开球}。
\textbf{闭球}（Closed Ball）：称 \(\displaystyle \overline{B}_r(p_0):=\bigl\{p\in\mathbb R^n\mid\|p-p_0\|\leqslant  r\bigr\}\) 为以 $p_0$ 为中心、$r$ 为半径的\textbf{闭球}。
\end{definition}

\begin{definition}[内点、外点与边界点] \label{def:interior-point}
设 $A\subset\mathbb R^n$，$p\in\mathbb R^n$。按照 $p$ 相对于 $A$ 的位置关系分为三类：
\begin{enumerate}
  \item\textbf{内点（Interior Point）}:若 $\exists\,r>0$ 使 $B_r(p)\subset A$，则称 $p$ 为 $A$ 的\textbf{内点}。全体内点集合记为 $A^\circ$。
  \item\textbf{外点（Exterior Point）}:若 $\exists\,r>0$ 使 $B_r(p)\subset A^c$，则称 $p$ 为 $A$ 的\textbf{外点}。
  \item\textbf{边界点（Boundary Point）}:若对 $\forall\,r>0$ 都有 $B_r(p)\cap A\neq\varnothing$ 且 $B_r(p)\cap A^c\neq\varnothing$，则称 $p$ 为 $A$ 的\textbf{边界点}。全体边界点集合记为 $\partial A$。
\end{enumerate}
\end{definition}

\begin{definition}[聚点与孤立点] \label{def:accumulation-point}
设 $A\subset\mathbb R^n$，$p\in\mathbb R^n$。
若对 $\forall\,r>0$ 都有 $(B_r(p)\setminus\{p\})\cap A\neq\varnothing$，则称 $p$ 为 $A$ 的\textbf{聚点}（Accumulation Point）。
若 $p\in A$ 但 $p$ 不是 $A$ 的聚点，则称 $p$ 为 $A$ 的\textbf{孤立点}（Isolated Point）。
\end{definition}


\begin{definition}[开集与闭集] \label{def:open-set}
设 $A\subset\mathbb R^n$。
\begin{itemize}
  \item 若 $A=A^\circ$（即 $A$ 的每个点都是内点），则称 $A$ 为\textbf{开集}（Open Set）；
  \item 若 $A^c$（$A$ 的补集）是开集，则称 $A$ 为\textbf{闭集}（Closed Set）。
\end{itemize}
\end{definition}

\begin{theorem}[闭集的聚点刻画] \label{thm:closed-equiv}
$A$ 为闭集 $\Longleftrightarrow$ $A$ 包含其所有的聚点（即 $A'\subset A$，其中 $A'$ 表示 $A$ 的导集）。
\end{theorem}

\begin{theorem}[闭包是闭集] \label{thm:closure-closed}
对任意 $A\subset\mathbb R^n$，其闭包 $\bar A=A\cup A'$ 必为闭集。
\end{theorem}


\begin{definition}[区域 / Region / Domain] \label{def:region}
设 $D\subset\mathbb R^n$ 是一个\textbf{开集}。若对 $D$ 中任意两点 $x,y\in D$，都存在一条完全位于 $D$ 内部的折线连接 $x$ 与 $y$，则称 $D$ 是\textbf{连通的}（Connected）。
\textbf{连通的开集}称为\textbf{开区域}，开区域的闭包称为\textbf{闭区域}。
\end{definition}

\begin{note}[术语提醒]
函数本身可以定义在任意点集上，并不要求定义域一开始就是区域。这里只是把后续教材里最常出现的一类定义域单独命名出来。等到讨论极限、方向导数、可微性、曲线曲面积分等问题时，我们往往需要在点的附近向各个方向做局部扰动，这时“开集”或“区域”的假设才会频繁出现。
\end{note}

\begin{exercise}[开球是区域] \label{ex:ball-connected}
证明：$\mathbb R^n$ 中的开球 $B_r(a)$ 是区域（即既是开集又是连通的）。
\end{exercise}
\begin{solution}
\textbf{开集性}：开球中每一点都是内点，故是开集。

\textbf{连通性}：任取 $p,q\in B_r(a)$，构造连接它们的直线段 $\Phi(t)=(1-t)p+tq,\; 0\leqslant  t\leqslant  1$。由三角不等式，对任意 $t\in[0,1]$，\(\displaystyle \|\Phi(t)-a\| = \|(1-t)(p-a)+t(q-a)\|\leqslant (1-t)\|p-a\|+t\|q-a\|< (1-t)r+tr = r\)。
故整条线段完全含于球内，因此 $B_r(a)$ 连通。
\end{solution}

\begin{remark}[凸性蕴含连通性]
这里先补充一个术语：若集合 $E\subset\mathbb R^n$ 满足对任意 $p,q\in E$ 与任意 $t\in[0,1]$，都有
\[
  (1-t)p+tq\in E,
\]
则称 $E$ 为\textbf{凸集}（Convex Set）。也就是说，集合中任意两点之间的整条线段都仍在集合内。上述证明实际上展示了更强的性质：\textbf{开球是凸集}。
凸集一定是连通的（因为线段是一种特殊的折线），但反之不成立——例如 $\mathbb R^2\setminus\{0\}$ 是连通的但不是凸的。
\end{remark}

\begin{exercise}[开但不连通的集合] \label{ex:not-region}
设 $D=(0,1)\times(1,2)\,\cup\,(1,2)\times(0,1)\subset\mathbb R^2$。则 $D$ 是开集但不是区域（不连通）。
\end{exercise}
\begin{solution}
$D$ 可以写成 \(\displaystyle D=U\cup V\)，其中
\(\displaystyle U=(0,1)\times(1,2),\ V=(1,2)\times(0,1)\)。
$U$ 与 $V$ 都是开矩形，因此都是开集；开集的并仍是开集，所以 $D$ 是开集。

另一方面，$U$ 与 $V$ 非空、互不相交，且都相对 $D$ 开，于是 $D$ 被分成了两个互不连通的部分。因此 $D$ 不是连通集，也就不是区域。
\end{solution}

\subsection{多元函数、极限与连续}

\begin{definition}[$n$ 元函数 / Function of $n$ Variables] \label{def:multivar-func}
设 $D\subset\mathbb{R}^n$ 是一个\textbf{区域}（Region，即连通的开集）。称映射
\(\displaystyle f:\;D\implies\mathbb{R},\ (x_1,x_2,\dots,x_n)\longmapsto f(x_1,x_2,\dots,x_n)\) 为一个\textbf{$n$ 元函数}（Function of $n$ Variables）。其中：
\begin{itemize}
  \item $D$ 称为 $f$ 的\textbf{定义域}（Domain）；
  \item $f(D)\subset\mathbb{R}$ 称为 $f$ 的\textbf{值域}（Range / Codomain）；
  \item $(x_1,x_2,\dots,x_n)$ 称为\textbf{自变量}（Independent Variables）；
  \item $f(x_1,x_2,\dots,x_n)$ 或简记为 $f(x)$（其中 $x=(x_1,\dots,x_n)$）称为\textbf{因变量}。
\end{itemize}
\end{definition}

\begin{note}[记号惯例]
特别值得注意的是关于值域连通性的一个事实：若 $D$ 是区域且 $f$ 连续（连续性的严格定义将在本节后文给出），则 $f(D)$ 作为连续映射下的像集也是 $\mathbb{R}$ 中的连通子集——即 $f(D)$ 必为一个区间。这一性质将在介值定理（定理 \ref{thm:cont-intermediate}）中得到体现。
\end{note}

\begin{exercise}[求二元函数的定义域] \label{ex:domain-finding}
求下列二元函数的自然定义域：(1)$z=\sqrt{xy}$；(2) $\displaystyle z=\frac{\sqrt{2x-x^2-y^2}}{\sqrt{x^2+y^2-1}}$。
\end{exercise}
\begin{solution}
\begin{enumerate}
  \item \textbf{求解题 (1)}：被开方量须非负，即要求 $xy\geqslant 0$。这意味着 $x$ 与 $y$ 必须同号（或至少有一个为零）。由此得到定义域 \(\displaystyle D_1=\{(x,y)\mid x\geqslant 0,\;y\geqslant 0\}\cup\{(x,y)\mid x\leqslant  0,\;y\leqslant  0\}\)。
  \textit{几何意义}：这是平面直角坐标系中的第一象限和第三象限（含坐标轴），呈两个对顶的 ``无限角形'' 区域。

  \item \textbf{求解题 (2)}：将函数分子和分母的条件分别进行分析：
  \begin{itemize}
    \item \textbf{分子部分}：由于处于根号内，要求 $2x-x^2-y^2\geqslant 0$。通过配方化简可得 \(\displaystyle x^2-2x+y^2\leqslant  0 \Longleftrightarrow (x-1)^2+y^2\leqslant  1\)。
    这表示以 $(1,0)$ 为圆心、半径为 $1$ 的闭圆盘。
    \item \textbf{分母部分}：不仅在根号内需要大于等于零，且作为分母不能为零，故要求 $x^2+y^2-1>0$。化简可得 \(\displaystyle x^2+y^2>1\)。
    这表示单位圆的外部（不含边界圆周）。
  \end{itemize}
  \textbf{综合条件}：取分子与分母条件的交集，得到最终定义域 \(\displaystyle D_2=\bigl\{(x,y)\mid (x-1)^2+y^2\leqslant  1\bigr\}\cap\bigl\{(x,y)\mid x^2+y^2>1\bigr\}\)。
  几何意义：这是两个相交圆的公共部分，即一个月牙形的闭区域去除了其内侧边界（不包含 $x^2+y^2=1$ 的圆弧部分）。
\end{enumerate}
\end{solution}

对于一元函数 $y=f(x)$，其图像是 $\mathbb{R}^2$ 平面上的一条曲线。对于二元函数 $z=f(x,y)$，情况变得更加丰富：当点 $(x,y,0)$ 在定义域 $D$ 上连续变动时，对应的空间点 $(x,y,f(x,y))$ 在三维空间中描绘出的轨迹构成一张\textbf{曲面}。这张曲面称为二元函数 $z=f(x,y)$ 的\textbf{图像}，记作 \(\displaystyle \Gamma_f:=\bigl\{(x,y,f(x,y))\mid(x,y)\in D\bigr\}\subset\mathbb{R}^3\)。

\begin{exercise}[双曲抛物面] \label{ex:hyperbolic-paraboloid}
函数 $z=x^2-y^2$ 的图像是一张著名的二次曲面——\textbf{双曲抛物面}（Hyperbolic Paraboloid），因其形似马鞍又俗称\textbf{马鞍面}（Saddle Surface）。
沿 $x$ 轴方向 ($y=0$)：$z=x^2$ ——开口向上的抛物线；沿 $y$ 轴方向 ($x=0$)：$z=-y^2$ ——开口向下的抛物线。``一向上一下向'' 的交叉形态正是马鞍面的特征。
\end{exercise}

\begin{definition}[等高线 / Contour Line] \label{def:contour-line}
设 $z=f(x,y)$ 是定义在 $D\subset\mathbb{R}^2$ 上的二元函数。对给定的常数 $c\in\mathbb{R}$，方程 \(\displaystyle f(x,y)=c\) 确定的平面点集（通常是 $\mathbb{R}^2$ 中的一条或若干条曲线）称为函数 $f$ 的\textbf{值为 $c$ 的等高线}（Contour Line / Level Curve），也称为\textbf{等位线}或\textbf{水平集}（Level Set）。

直观地说，等高线就是水平面 $z=c$ 与曲面 $z=f(x,y)$ 的交线在 $xy$ 平面上的投影。
\end{definition}

\begin{exercise}[旋转抛物面的等高线]
考虑函数 $f(x,y)=1-x^2-y^2$（开口向下的旋转抛物面，顶点在 $(0,0,1)$）。令 $z=c$ 得 $x^2+y^2=1-c$：
\begin{itemize}
  \item 当 $c<1$ 时，$1-c>0$，等高线是以原点为中心、半径为 $\sqrt{1-c}$ 的圆；
  \item 当 $c=1$ 时，退化为一点 $(0,0)$（对应抛物面的顶点）；
  \item 当 $c>1$ 时，无实解（水平面高于顶点，不相交）。
\end{itemize}
因此该函数的等高线族是一组同心圆，半径随 $c$ 减小而增大。特别地，取 $c=0.75$ 则等高线为 $x^2+y^2=0.25$，即半径为 $0.5$ 的圆。
\end{exercise}

将等高线的思想推广到三元函数：

\begin{definition}[等位面 / Level Surface] \label{def:level-surface}
设 $u=f(x,y,z)$ 是定义在 $\Omega\subset\mathbb{R}^3$ 上的三元函数。对常数 $c\in\mathbb{R}$，方程 \(\displaystyle f(x,y,z)=c\) 确定的 $\mathbb{R}^3$ 中的曲面称为 $f$ 的\textbf{值为 $c$ 的等位面}（Level Surface）。
\end{definition}

\begin{exercise}[球面等位面族]
函数 $f(x,y,z)=x^2+y^2+z^2$ 的等位面方程为 $x^2+y^2+z^2=c$：
\begin{itemize}
  \item 当 $c>0$ 时，是以原点 $(0,0,0)$ 为中心、$\sqrt{c}$ 为半径的球面；
  \item 当 $c=0$ 时，退化为原点这一个孤立点；
  \item 当 $c<0$ 时，空集（无实数解）。
\end{itemize}
当 $c$ 从 $+\infty$ 变化到 $0$ 时得到一族 ``由外向内收缩'' 的同心球面——这正是一个以原点为 ``源'' 的径向场的等位面分布。
\end{exercise}


\subsubsection{多元函数的极限}

有了邻域的概念之后，我们可以将一元函数极限的 $\varepsilon$-$\delta$ 语言自然地推广到多元情形。然而多元极限远比一元极限复杂—— ``趋近方式 '' 从 ``左右两侧 '' 爆炸式地变为 ``无穷多条路径 ''。

\begin{definition}[$n$ 元函数的极限] \label{def:multivar-limit}
设 $f:D\to\mathbb{R}$ 是定义在区域 $D\subset\mathbb{R}^n$ 上的 $n$ 元函数，$a=(a_1,a_2,\dots,a_n)\in D'$（即 $a$ 是 $D$ 的一个聚点），$A\in\mathbb{R}$ 是一个给定常数。

如果对于任意给定的 $\varepsilon>0$（无论多么小），总存在 $\delta>0$，使得对所有满足 $0<\|x-a\|<\delta$ 且 $x\in D$ 的 $x$，都有 \(\displaystyle |f(x)-A|<\varepsilon\)，则称 $A$ 为 $f(x)$ 当 $x\to a$ 时的\textbf{极限}（Limit），记作 \(\displaystyle \lim_{x\to a}f(x)=A\quad\text{或}\quad f(x)\to A\;\;(\|x-a\|\to 0)\)。
\end{definition}

当 $n=2$ 时，上述极限特称为\textbf{二重极限}（Double Limit），以区别于后面将要介绍的 ``累次极限''。在一元函数极限中，我们只需检查 ``左极限'' 和 ``右极限'' 是否相等——因为自变量只能沿着实数轴从一个点的左侧或右侧趋近该点。但在多元情形下，点 $x$ 向 $a$ 的趋近可以有\textbf{无穷多种方式}：
\begin{itemize}
  \item 沿着不同的直线方向（斜率不同）；
  \item 沿着各种曲线路径（抛物线 $y=kx^2$、螺线 $r=e^{-\theta}$ 等）；
  \item 甚至沿着 ``锯齿状'' 的振荡路径趋近。
\end{itemize}

下面先给出极限计算的正面例子，再展示几个不存在的典型反例。

\begin{exercise}[极坐标下的判零] \label{ex:limit-squeeze-basic}
设 \(p=(x,y)\neq(0,0)\)，\(\displaystyle f(p)=f(x,y)=\frac{xy}{\sqrt{x^2+y^2}}\)，求 $\displaystyle\lim_{p\to 0}f(p)$。
\end{exercise}
\begin{solution}
令 \(x=r\cos\theta,\ y=r\sin\theta\)，则 \(\sqrt{x^2+y^2}=r\)，且 \(p\to0\) 等价于 \(r\to0^+\)。于是
\[
  f(r\cos\theta,r\sin\theta)
  =\frac{r^2\cos\theta\sin\theta}{r}
  =r\cos\theta\sin\theta.
\]
由于 \(|\cos\theta\sin\theta|\leqslant \frac12\)，有
\[
  |f(r\cos\theta,r\sin\theta)|\leqslant \frac12r\to0.
\]
故 \(\displaystyle \boxed{\lim_{p\to 0}f(p)=0}\)。
\end{solution}

\begin{theorem}[海涅归结原理 / Heine's Principle for Multivariable Functions] \label{thm:heine-multivar}
设 $D\subset\mathbb{R}^n$，$f:D\to\mathbb{R}$，$a\in D'$ 为 $D$ 的聚点。则以下两个命题等价：
\[
\lim_{x\to a}f(x)=A
\quad\Longleftrightarrow\quad
\forall\,\{x_m\}_{m=1}^\infty\subset D,\;
x_m \neq a,\;
\lim_{m\to\infty}x_m=a:
\;\implies\;
\lim_{m\to\infty}f(x_m)=A.
\]
也就是说，$f(x)$ 在 $a$ 处有极限 $A$ 当且仅当对任何趋向于 $a$ 的（异于 $a$ 的）点列，对应的函数值序列都收敛于同一个值 $A$。
\end{theorem}

\begin{proof}
\textbf{``$\Rightarrow$'' 方向}：由极限定义直接可得。对任意 $\varepsilon>0$，存在 $\delta>0$ 使 $0<\|x-a\|<\delta$ 时 $|f(x)-A|<\varepsilon$。由于 $x_m\to a$，存在 $N$ 使得当 $m>N$ 时 $0<\|x_m-a\|<\delta$，从而 $m>N$ 时 $|f(x_m)-A|<\varepsilon$。这就证明了 $f(x_m) \to A$。

\textbf{``$\Leftarrow$'' 方向}：（反证法）假设 $\displaystyle\lim_{x\to a}f(x) \neq A$（或极限不存在）。根据极限定义的否定形式，存在某个 $\varepsilon_0>0$ 使得对任意 $m\in\mathbb{N}$（无论 $\delta = 1/m$ 多小），都能在去心邻域内找到 $x_m\in D$ 同时满足
\(\displaystyle 0<\|x_m-a\|<\frac1m\quad\text{但}\quad |f(x_m)-A|\geqslant\varepsilon_0\)。
这样构造的点列 $\{x_m\}$ 显然满足 $x_m\to a$（因为 $\|x_m-a\|<1/m\to 0$），但 $\{f(x_m)\}$ 不收敛于 $A$（因为每项离 $A$ 至少有 $\varepsilon_0$ 的距离），这与前提条件矛盾。
\end{proof}

\begin{note}
海涅原理提供了证明极限\textbf{不存在}的理论依据：只要能找到趋向于 $a$ 的不同方式，使函数值趋于不同结果，就可断言极限不存在。接下来两个例子用极坐标把这种差异写出来。
\end{note}

\begin{exercise}[角度不同导致极限不同] \label{ex:limit-not-exist-linear}
讨论函数 \(\displaystyle f(x,y)=\frac{xy}{x^2+y^2}\) 在原点处的极限是否存在，其中 $(x,y)\neq(0,0)$。
\end{exercise}
\begin{solution}
我们引入极坐标变换来讨论该重极限。令 $x = r\cos\theta$, $y = r\sin\theta$，当 $(x,y) \to (0,0)$ 时，等价于极径 $r \to 0^+$，而极角 $\theta$ 可以是任意方向。
\begin{enumerate}
  \item \textbf{代入极坐标化简：}
  将极坐标代入原函数中进行整理：\(\displaystyle f(r\cos\theta, r\sin\theta)= \frac{(r\cos\theta)(r\sin\theta)}{(r\cos\theta)^2+(r\sin\theta)^2}= \frac{r^2\cos\theta\sin\theta}{r^2(\cos^2\theta+\sin^2\theta)}= \frac{1}{2}\sin(2\theta)\)。

  \item \textbf{分析极限结果：}
  我们发现，化简后的表达式中 $r$ 已经被完全消去，函数值仅依赖于趋近原点时的极角 $\theta$。
  这意味着，当我们沿着不同角度的射线趋向原点时，会得到截然不同的极限值。例如：
  \begin{itemize}
      \item 当沿 $x$ 轴（即 $\theta = 0$）趋近时，$\lim_{r\to 0^+} f = \frac{1}{2}\sin 0 = 0$；
      \item 当沿直线 $y=x$（即 $\theta = \frac{\pi}{4}$）趋近时，$\lim_{r\to 0^+} f = \frac{1}{2}\sin\left(2 \cdot \frac{\pi}{4}\right) = \frac{1}{2}$。
  \end{itemize}

  \item \textbf{得出结论：}
  由于沿不同方向（不同 $\theta$）趋近原点时，极限值不唯一，根据重极限存在的充要条件（极限值必须与趋近路径无关），该极限不存在，即 \(\displaystyle \boxed{\lim_{(x,y)\to(0,0)}\frac{xy}{x^2+y^2}\text{ 不存在}}\)。
\end{enumerate}
\end{solution}

\begin{exercise}[角度随半径变化的情形] \label{ex:limit-path-curve}
讨论函数 \(\displaystyle f(p)=f(x,y)=\frac{2x^2y}{x^4+y^2},\qquad p=(x,y) \neq (0,0)\) 在原点处的极限是否存在。
\end{exercise}
\begin{solution}
我们引入极坐标变换 $x = r\cos\theta$, $y = r\sin\theta$，探讨当 $r \to 0^+$ 时的极限行为。
\begin{enumerate}
  \item \textbf{代入极坐标化简：}
  将极坐标代入原函数并提取公因式：
  \[
  f(r\cos\theta, r\sin\theta)
  =\frac{2(r\cos\theta)^2(r\sin\theta)}{(r\cos\theta)^4+(r\sin\theta)^2}
  =\frac{2r\cos^2\theta\sin\theta}{r^2\cos^4\theta+\sin^2\theta}.
  \]

  \item \textbf{让角度随 $r$ 变化：}
  观察极坐标分母 $r^2\cos^4\theta + \sin^2\theta$，为了让分子分母同阶，可令 $\sin\theta = k r\cos^2\theta$（其中 $k$ 为非零常数，这在直角坐标系中对应 $y=kx^2$）。
  将此关系代入化简后的表达式中，得
  \[
  \frac{2r\cos^2\theta \cdot (kr\cos^2\theta)}{r^2\cos^4\theta+(kr\cos^2\theta)^2}
  =\frac{2k}{1+k^2}.
  \]

  \item \textbf{得出结论：}
  当 $\sin\theta = k r\cos^2\theta$ 时，极限值是一个与参数 $k$ 相关的常数 $\frac{2k}{1+k^2}$。例如：
  \begin{itemize}
      \item 取 $k=1$ （对应 $y=x^2$），极限为 $\frac{2}{1+1} = 1$；
      \item 取 $k=2$ （对应 $y=2x^2$），极限为 $\frac{4}{1+4} = \frac{4}{5}$。
  \end{itemize}
  由于所得结果随 $k$ 改变，因此 \(\displaystyle \boxed{\lim_{(x,y)\to(0,0)}\frac{2x^2y}{x^4+y^2}\text{ 不存在}}\)。
\end{enumerate}
\begin{note}
本例提醒我们：固定角度时看不出的差异，可能会在 $\theta$ 随 $r$ 变化时出现。因此判断二重极限时，极坐标下的表达式还要留意分母中不同项的阶数平衡。
\end{note}
\end{solution}

\begin{definition}[累次极限] \label{def:iterated-limit}
设 $f(x,y)$ 在 $(x_0,y_0)$ 的某去心邻域内有定义。两种\textbf{累次极限}记为
\[
L_1=\lim_{y\to y_0}\!\lim_{x\to x_0}f(x,y),\qquad
L_2=\lim_{x\to x_0}\!\lim_{y\to y_0}f(x,y).
\]
其中第一个表示 ``先固定 $y$，对 $x$ 取极限（得到中间函数 $\varphi(y)$），再对 $y$ 取极限''；第二个同理但顺序相反。

这两种累次极限与前面定义的 $\displaystyle\lim_{(x,y)\to(x_0,y_0)}f(x,y)$（\textbf{二重极限}）是\textbf{完全不同的概念}。
\end{definition}

\begin{note}[两者的区别]
\par\noindent{\centering
\renewcommand{\arraystretch}{1.4}
\begin{tabular}{@{}p{4.8cm}p{7cm}@{}}
\textbf{二重极限} & \textbf{累次极限} \\
\makecell[l]{$(x,y)$ 同时以\textbf{任意方式}\\趋于 $(x_0,y_0)$} &
\makecell[l]{先让一个变量趋于极限，\\再让另一个变量趋于极限} \\
只涉及一个极限过程 & 涉及两个串行的极限过程 \\
结果如果存在则必唯一 & 两个顺序的结果可能不同 \\
不依赖坐标系选取 & 依赖变量顺序的选择 \\
\end{tabular}
\par}
\end{note}

\begin{exercise}[二重极限不存在但累次极限存在] \label{ex:iterated-exists-double-no}
回顾例 \ref{ex:limit-path-curve} 中的函数 $f(x,y)=\dfrac{2x^2y}{x^4+y^2}$（$(x,y) \neq (0,0)$）。

已知其二重极限在 $(0,0)$ 处\textbf{不存在}。但两个累次极限均存在且相等：先固定 $y$ 得 \(\displaystyle \lim_{y\to 0}\lim_{x\to 0}\frac{2x^2y}{x^4+y^2}=\lim_{y\to 0}0=0\)，先固定 $x$ 得 \(\displaystyle \lim_{x\to 0}\lim_{y\to 0}\frac{2x^2y}{x^4+y^2}=\lim_{x\to 0}0=0\)。
\end{exercise}

\begin{exercise}[二重极限存在但累次极限不全存在] \label{ex:double-exists-iterated-partial}
设 \(\displaystyle f(x,y)=y\sin\frac1x\quad(x \neq 0)\)。讨论其在 $(0,0)$ 处的二重极限与累次极限。
\end{exercise}
\begin{solution}
\begin{enumerate}
  \item \textbf{检验二重极限}：令 $x=r\cos\theta,\ y=r\sin\theta$。在定义域内有 $\cos\theta\neq0$，并且
  \[
    f(r\cos\theta,r\sin\theta)
    =r\sin\theta\sin\frac{1}{r\cos\theta}.
  \]
  因为 \(\left|\sin\dfrac{1}{r\cos\theta}\right|\leqslant1\)，所以
  \[
    |f(r\cos\theta,r\sin\theta)|\leqslant r|\sin\theta|\leqslant r\to0.
  \]
  因此 \(\displaystyle\lim_{(x,y)\to(0,0)}f(x,y)=0\)，二重极限存在且为 \(0\)。

  \item \textbf{检验累次极限}：
  \begin{itemize}
    \item 先对 $x$ 求极限再对 $y$ 求极限：
    \[
    \lim_{y\to 0}\lim_{x\to 0}y\sin\frac1x
    =\lim_{y\to 0}\left(y\cdot \lim_{x\to 0}\sin\frac1x\right).
    \]
    由于内层极限 $\lim_{x\to 0}\sin(1/x)$ 在 $[-1,1]$ 间无限振荡而不存在，因此这个累次极限整体\textbf{不存在}。
    \item 先对 $y$ 求极限再对 $x$ 求极限：\(\displaystyle \lim_{x\to 0}\lim_{y\to 0}y\sin\frac1x
    =\lim_{x\to 0}\left(0\cdot\sin\frac1x\right)
    =\lim_{x\to 0}0=0\)。
    此累次极限存在且等于 $0$。
  \end{itemize}
\end{enumerate}

\end{solution}



\begin{exercise}[含三角函数的极限] \label{ex:limit-squeeze}
求 $\displaystyle\lim_{(x,y)\to(0,0)}\frac{\sin(x^3+y^3)}{x^2+y^2}$。
\end{exercise}
\begin{solution}
令 \(x=r\cos\theta,\ y=r\sin\theta\)。则
\[
  \left|\frac{\sin(x^3+y^3)}{x^2+y^2}\right|
  =\left|\frac{\sin\bigl(r^3(\cos^3\theta+\sin^3\theta)\bigr)}{r^2}\right|.
\]
由 \(|\sin u|\leqslant |u|\)，得
\[
  \left|\frac{\sin\bigl(r^3(\cos^3\theta+\sin^3\theta)\bigr)}{r^2}\right|
  \leqslant r|\cos^3\theta+\sin^3\theta|
  \leqslant 2r\to0.
\]
因此
\[
  \boxed{\displaystyle\lim_{(x,y)\to(0,0)}\frac{\sin(x^3+y^3)}{x^2+y^2}=0}.
\]
\end{solution}

\begin{exercise}[含分段定义的极限] \label{ex:limit-piecewise}
设
\[
f(x,y)=
\begin{cases}
\dfrac{\sin(xy)}{x},&x \neq 0,\\[10pt]
y,&x=0.
\end{cases}
\]
求 $\displaystyle\lim_{(x,y)\to(0,0)}f(x,y)$。
\end{exercise}
\begin{solution}
令 \(x=r\cos\theta,\ y=r\sin\theta\)。若 \(x\neq0\)，则
\[
  |f(x,y)|=\left|\frac{\sin(xy)}{x}\right|
  \leqslant \frac{|xy|}{|x|}
  =|y|=r|\sin\theta|\leqslant r.
\]
若 \(x=0\)，则 \(f(0,y)=y\)，同样有 \(|f(0,y)|=|y|=r\)。所以在两种情形下均有 \(|f(x,y)|\leqslant r\to0\)，从而
\[
  \boxed{\displaystyle\lim_{(x,y)\to(0,0)}f(x,y)=0}.
\]
\end{solution}

对于形如 $f(x,y)^{g(x,y)}$ 的不定式，常用的策略是取对数将其化为指数形式 $e^{g \ln f}$。

\begin{exercise}[幂指型极限] \label{ex:limit-power-exp}
求 $\displaystyle\lim_{(x,y)\to(0,0)}(x^2+y^2)^{xy}$。
\end{exercise}
\begin{solution}
\begin{enumerate}
  \item \textbf{指数化变形}：将原表达式改写为以 $e$ 为底的指数形式：\(\displaystyle (x^2+y^2)^{xy}=e^{\,xy\ln(x^2+y^2)}\)。
  由于指数函数连续，我们只需计算其指数部分的极限 $\displaystyle\lim_{(x,y)\to(0,0)}xy\ln(x^2+y^2)$。

  \item \textbf{代入极坐标}：令 \(x=r\cos\theta,\ y=r\sin\theta\)，则
  \[
    xy\ln(x^2+y^2)
    =r^2\cos\theta\sin\theta\ln r^2.
  \]
  因为 \(|\cos\theta\sin\theta|\leqslant\frac12\)，所以
  \[
    \bigl|r^2\cos\theta\sin\theta\ln r^2\bigr|
    \leqslant \frac12 r^2|\ln r^2|\to0.
  \]

  \item \textbf{得出最终结果}：由于指数部分的极限趋于 $0$，因此原极限趋于 $e^0=1$：\(\displaystyle \boxed{\lim_{(x,y)\to(0,0)}(x^2+y^2)^{xy}=1}\)。
\end{enumerate}
\end{solution}

\begin{exercise}[含小角三角函数的极限] \label{ex:limit-product}
求 $\displaystyle\lim_{(x,y)\to(0,0)}\frac{\sin(x^2y)}{x^2+y^2}$。
\end{exercise}
\begin{solution}
令 \(x=r\cos\theta,\ y=r\sin\theta\)。由 \(|\sin u|\leqslant |u|\)，有
\[
  \left|\frac{\sin(x^2y)}{x^2+y^2}\right|
  =\left|\frac{\sin\bigl(r^3\cos^2\theta\sin\theta\bigr)}{r^2}\right|
  \leqslant r|\cos^2\theta\sin\theta|
  \leqslant r\to0.
\]
故
\[
  \boxed{\displaystyle\lim_{(x,y)\to(0,0)}\frac{\sin(x^2y)}{x^2+y^2}=0}.
\]
\end{solution}


\subsubsection{多元函数的连续性}

\begin{definition}[$n$ 元函数的连续性] \label{def:multivar-cont}
设 $f:D\to\mathbb{R}$ 定义在区域 $D\subset\mathbb{R}^n$ 上，$p_0=(x_1^{(0)},\dots,x_n^{(0)})\in D$ 为 $D$ 的聚点。若
\begin{align*}
\lim_{p\to p_0}f(p)&=f(p_0),\\
\text{即}\quad
\forall\varepsilon>0,\;\exists\delta>0:\quad
\|p-p_0\|<\delta&\implies |f(p)-f(p_0)|<\varepsilon,
\end{align*}
则称 $f$ 在 $p_0$ 处\textbf{连续}（Continuous），$p_0$ 为 $f$ 的一个\textbf{连续点}（Point of Continuity）。若 $f$ 在 $D$ 的每个点处都连续，则称 $f$ 在 $D$ 上\textbf{连续}，记作 $f\in C(D)$。
\end{definition}

\begin{note}[连续性的三个要点]
\begin{enumerate}
  \item[(1)] $f$ 在 $p_0$ 处连续首先要求 $f$ 在 $p_0$ 处\textbf{有定义}——这与极限不同（极限不考虑该点本身的值）。
  \item[(2)] 若 $p_0$ 是 $D$ 的\textbf{孤立点}（Isolated Point），则 $f$ 在 $p_0$ 处\textbf{自动连续}。因为孤立点的去心邻域不含 $D$ 中其他点，极限定义 ``平凡满足''。特别地，若 $D$ 是有限点集，则任何 $f:D\to\mathbb{R}$ 都是连续的。
  \item[(3)] 多元连续函数保持了一元连续函数的所有优良性质：四则运算封闭性、复合运算封闭性等均可平行推广。
\end{enumerate}
\end{note}

\begin{property}[边缘连续性与整体连续性]
若 $f(x,y)$ 在 $(x_0,y_0)$ 处连续，则它的两个 ``边缘函数''（Partial Functions）\(\displaystyle g(x):=f(x,y_0),\qquad h(y):=f(x_0,y)\) 分别在 $x_0$ 和 $y_0$ 处（作为一元函数）连续。

\textbf{反之不真！} 边缘函数的连续性\textbf{不能推出} 二元函数的整体连续性。反例见例 \ref{ex:cont-edge-not-total}。
\end{property}

\begin{exercise}[基本连续函数类] \label{ex:basic-continuous}
\begin{shortsingleenum}
  \shortsingleitem{1}{常值函数 $f(x)=c$ 在 $\mathbb{R}^n$ 上连续（显然）；}
  \shortsingleitem{2}{\textbf{投影函数}（Coordinate Projection）：对 $x=(x_1,x_2,\dots,x_n)\in\mathbb{R}^n$，令 $f(x)=x_i$（取出第 $i$ 个分量），则 $f$ 在 $\mathbb{R}^n$ 上处处连续。验证：对任意 $p\in\mathbb{R}^n$，$f(p)=p_i$，于是 \(\displaystyle |f(x)-f(p)|=|x_i-p_i|\leqslant\sqrt{\sum_{j=1}^{n}(x_j-p_j)^2}=\|x-p\|\)。给定 $\varepsilon>0$，取 $\delta=\varepsilon$ 即可满足连续性定义。}
\end{shortsingleenum}
\end{exercise}

\begin{exercise}[可去间断点] \label{ex:removable-disc}
讨论函数 \(\displaystyle f(p)=\frac{x^2y^2}{x^2+y^2},\quad p=(x,y) \neq (0,0)\) 的连续性。
\end{exercise}
\begin{solution}
\begin{enumerate}
  \item \textbf{原点之外的区域}：当 $(x,y) \neq (0,0)$ 时，$f$ 是两个多项式的商且分母非零。由代数运算的连续性可知，$f$ 在 $\mathbb{R}^2\setminus\{(0,0)\}$ 上是连续的。
  \item \textbf{考察原点处的极限}：令 $x=r\cos\theta,\ y=r\sin\theta$，则
  \[
    f(r\cos\theta,r\sin\theta)
    =\frac{r^4\cos^2\theta\sin^2\theta}{r^2}
    =r^2\cos^2\theta\sin^2\theta.
  \]
  因为 \(0\leqslant\cos^2\theta\sin^2\theta\leqslant1\)，故 \(|f(r\cos\theta,r\sin\theta)|\leqslant r^2\to0\)，所以 \(\displaystyle\lim_{p\to0}f(p)=0\)。
  \item \textbf{结论}：若我们补充定义 $f(0,0)=0$，则函数值等于极限值，此时 $f$ 在 $\mathbb{R}^2$ 上就变得\textbf{处处连续}了。因此，原点是该函数的一个\textbf{可去间断点}（Removable Discontinuity）。
\end{enumerate}
\end{solution}

\begin{exercise}[复合函数的连续性] \label{ex:composite-continuous}
证明 $f(x,y,z)=\sin(xy+z)$ 在 $\mathbb{R}^3$ 上连续。
\end{exercise}
\begin{solution}
\begin{enumerate}
  \item \textbf{分解复合结构}：我们将 $f$ 分解为复合函数的链式结构，即
  \(\displaystyle (x,y,z)\;\xmapsto{\;g\;}\; xy+z\;\xmapsto{\;\sin\;}\; \sin(xy+z)=f(x,y,z)\)。
  \item \textbf{分析子函数连续性}：
  \begin{itemize}
    \item 内部函数 $g(x,y,z)=xy+z$ 是关于变量 $x, y, z$ 的多项式函数，多项式在全空间 $\mathbb{R}^3$ 上是连续的。
    \item 外部函数 $\sin t$ 是一元基本初等函数，在 $\mathbb{R}$ 上是连续的。
  \end{itemize}
  \item \textbf{应用复合定理}：根据\textbf{多元复合函数的连续性定理}（一元连续函数与多元连续函数的复合仍连续），最终函数 $f=\sin\circ g$ 在 $\mathbb{R}^3$ 上处处连续。
\end{enumerate}
\end{solution}

\begin{exercise}[边缘连续但不整体连续] \label{ex:cont-edge-not-total}
设
\[
f(x,y)=
\begin{cases}
\dfrac{xy}{x^2+y^2},&(x,y) \neq (0,0),\\[8pt]
0,&(x,y)=(0,0).
\end{cases}
\]
讨论 $f$ 的连续性以及边缘函数 $f(\cdot,y)$ 和 $f(x,\cdot)$ 的连续性。
\end{exercise}
\begin{solution}
\begin{enumerate}
  \item \textbf{整体连续性检验}：由前文例 \ref{ex:limit-not-exist-linear} 已知，令 $x=r\cos\theta,\ y=r\sin\theta$ 后，
  \[
    \frac{xy}{x^2+y^2}=\cos\theta\sin\theta=\frac12\sin2\theta,
  \]
  结果依赖于角度 $\theta$，故原点处二重极限不存在。尽管 $f(0,0)$ 有定义，函数 $f$ 在 $(0,0)$ 处仍不连续。（除原点外，其余点处函数为分母非零的有理式，故连续。）

  \item \textbf{边缘函数 $g(x):=f(x,y_0)$ 的连续性检验}（固定 $y_0$ 看 $x$）：
  \begin{itemize}
    \item 当 $y_0 \neq 0$ 时：对任意 $x\in\mathbb{R}$，$g(x)=\dfrac{xy_0}{x^2+y_0^2}$ 是关于 $x$ 的有理函数，且分母 $x^2+y_0^2 \geqslant y_0^2 > 0$ 恒不为零，因此 $g(x)$ 在整个 $\mathbb{R}$ 上连续。
    \item 当 $y_0=0$ 时：代入定义得 $g(x)=f(x,0)=0$（对所有 $x$）。作为常值函数，显然也是连续的。特别地，$\displaystyle\lim_{x\to 0}f(x,0)=0=f(0,0)$。
  \end{itemize}
  故对于任何固定的 $y_0$，边缘函数 $g(x)$ 都是一元连续函数。

  \item \textbf{边缘函数 $h(y):=f(x_0,y)$ 的连续性检验}（固定 $x_0$ 看 $y$）：由于原函数 $f(x,y)$ 关于 $x$ 和 $y$ 是高度对称的，同理可证 $h(y)$ 在 $\mathbb{R}$ 上对每个固定的 $x_0$ 也是连续的一元函数。
\end{enumerate}
\begin{note}
$f$ 在原点处整体不连续，但沿着任何平等于坐标轴的直线切过去（即边缘函数），得到的一维函数却全都是连续的！这生动地印证了 ``边缘连续 $\not\Rightarrow$ 整体连续''。
\end{note}

\end{solution}

\subsubsection{有界闭区域上连续函数的性质}

一元函数在有界闭区间 $[a,b]$ 上的三大定理——有界性定理、最大最小值定理、介值定理（以及由此推出的零点定理和一致连续性）——都建立在同一个核心事实上：$\mathbb{R}^n$ 中的有界闭区域是紧集。

\begin{theorem}[有界性定理] \label{thm:cont-bounded}
设 $D\subset\mathbb{R}^n$ 是\textbf{有界闭区域}（Bounded Closed Region），$f:D\to\mathbb{R}$ 连续。则 $f$ 在 $D$ 上\textbf{有界}，即存在 $M>0$ 使得 $|f(p)|\leqslant  M$ 对所有 $p\in D$ 成立。
\end{theorem}

\begin{theorem}[最大值最小值定理 / Extreme Value Theorem] \label{thm:cont-extreme}
设 $D\subset\mathbb{R}^n$ 是有界闭区域，$f:D\to\mathbb{R}$ 连续。则 $f$ 在 $D$ 上\textbf{必定能取到}最大值和最小值：存在 $p_{\min},p_{\max}\in D$ 使得 \(\displaystyle f(p_{\min})\leqslant  f(p)\leqslant  f(p_{\max}),\quad\forall p\in D\)。
\end{theorem}

\begin{theorem}[介值定理 / Intermediate Value Theorem] \label{thm:cont-intermediate}
设 $D\subset\mathbb{R}^n$ 是有界闭区域（实际上只需 $D$ 是\textbf{道路连通}的），$f:D\to\mathbb{R}$ 连续。若 $a,b\in D$ 且实数 $r$ 满足 $f(a)<r<f(b)$，则存在 $c\in D$ 使得 $f(c)=r$.
\end{theorem}

\begin{theorem}[一致连续性 / Uniform Continuity] \label{thm:cont-uniform}
设 $D\subset\mathbb{R}^n$ 是有界闭区域，$f:D\to\mathbb{R}$ 连续。则 $f$ 在 $D$ 上\textbf{一致连续}（Uniformly Continuous），即
\[
\forall\varepsilon>0,\;\exists\delta>0:\;
\forall p,q\in D,\;\|p-q\|<\delta\;\implies\;|f(p)-f(q)|<\varepsilon.
\]
\end{theorem}
\begin{note}[四大定理之间的关系]
最值定理是最强的（蕴含其余三个）；介值定理和一致连续性彼此独立；有界性是最弱的推论。这四个定理的共同前提都是 ``$D$ 有界闭 + $f$ 连续''。
\end{note}

\begin{exercise}[无穷远处有极限的一致连续判定] \label{ex:inf-zero-extreme-unif}
设 $f\in C(\mathbb{R}^n)$，并且存在常数 $A\in\mathbb{R}$，使得
\[
  \lim_{\|\boldsymbol{p}\|\to\infty} f(\boldsymbol{p})=A.
\]
证明：$f$ 在 $\mathbb{R}^n$ 上一致连续。
\end{exercise}

\begin{proof}
任给 $\varepsilon>0$。由
\[
  \lim_{\|\boldsymbol{p}\|\to\infty}f(\boldsymbol{p})=A
\]
可知，存在 $R>0$，使得当 $\|\boldsymbol{p}\|>R$ 时，有
\[
  |f(\boldsymbol{p})-A|<\frac{\varepsilon}{2}.
\]
在有界闭集 $\overline{B}_{R+1}(\boldsymbol{0})$ 上，$f$ 连续，故由定理 \ref{thm:cont-uniform}，$f$ 在该闭球上一致连续。因此存在 $\delta\in(0,1)$，使得当
\[
  \boldsymbol{p},\boldsymbol{q}\in\overline{B}_{R+1}(\boldsymbol{0}),
  \qquad \|\boldsymbol{p}-\boldsymbol{q}\|<\delta
\]
时，有
\[
  |f(\boldsymbol{p})-f(\boldsymbol{q})|<\varepsilon.
\]

现在任取 $\boldsymbol{p},\boldsymbol{q}\in\mathbb{R}^n$，且 $\|\boldsymbol{p}-\boldsymbol{q}\|<\delta$。若 $\|\boldsymbol{p}\|>R$ 且 $\|\boldsymbol{q}\|>R$，则由三角不等式得
\[
  |f(\boldsymbol{p})-f(\boldsymbol{q})|
  \leqslant |f(\boldsymbol{p})-A|+|f(\boldsymbol{q})-A|
  <\varepsilon.
\]
若至少有一点满足范数不超过 $R$，不妨设 $\|\boldsymbol{p}\|\leqslant R$，则
\[
  \|\boldsymbol{q}\|\leqslant \|\boldsymbol{q}-\boldsymbol{p}\|+\|\boldsymbol{p}\|
  <\delta+R<R+1.
\]
于是 $\boldsymbol{p},\boldsymbol{q}\in\overline{B}_{R+1}(\boldsymbol{0})$，从而
\[
  |f(\boldsymbol{p})-f(\boldsymbol{q})|<\varepsilon.
\]
综上，对任意 $\varepsilon>0$，存在 $\delta>0$，使得 $\boldsymbol{p},\boldsymbol{q}\in\mathbb{R}^n$ 且 $\|\boldsymbol{p}-\boldsymbol{q}\|<\delta$ 时恒有 $|f(\boldsymbol{p})-f(\boldsymbol{q})|<\varepsilon$。故 $f$ 在 $\mathbb{R}^n$ 上一致连续。
\end{proof}

\begin{remark}
原题中 $\lim_{\|\boldsymbol{p}\|\to\infty}f(\boldsymbol{p})=0$ 只是上述判定定理的一个特殊情形。这里真正起作用的是：远处的函数值都靠近同一个常数，近处则落在某个有界闭集上，可以由一致连续性定理控制。
\end{remark}

\subsection{偏导数、方向导数、可微与全微分、nabla算子，梯度}

\subsubsection{方向导数}

\begin{definition}[方向导数 / Directional Derivative] \label{def:directional-derivative}
设 $D\subset\mathbb R^n$ 为开集，$f:D\to\mathbb R$ 为多元函数，$\vec l=(l_1,l_2,\dots,l_n)$ 是 $\mathbb R^n$ 中的一个\textbf{单位向量}（即 $\|\vec l\|=1$）。取定一点 $p_0\in D$。考虑一元辅助函数 \(\displaystyle u(t):=f(p_0+t\vec l),\qquad t\in(-\varepsilon,\varepsilon)\)
（其中 $\varepsilon>0$ 充分小以保证 $p_0+t\vec l\in D$；注意 $t$ 可正可负，对应两个相反的射线方向）。
若极限
\[\displaystyle u'(0)=\lim_{t\to 0}\frac{f(p_0+t\vec l)-f(p_0)}{t}\]
存在且有限，则称此极限为 $f$ 在点 $p_0$ 处\textbf{沿方向} $\vec l$ 的\textbf{方向导数}，记作
\(\displaystyle \frac{\partial f}{\partial \vec l}(p_0)\)。
\end{definition}

直观地说，$\dfrac{\partial f}{\partial \vec l}(p_0)$ 就是函数值 $f$ 当自变量从 $p_0$ 出发、以单位速度沿方向 $\vec l$ 移动时的\textbf{瞬时变化率}。同时要求 $D$ 是开集，是为了保证从 $p_0$ 沿 $\pm\vec l$ 走一小段时仍留在定义域内；否则差商甚至可能没有定义。

\begin{exercise}[方向导数的直接计算] \label{ex:dir-deriv-basic}
设 $f(x,y)=xy$，求 $f$ 在点 $(1,1)$ 处沿方向
$\displaystyle\vec l=\Bigl(\cos\frac{\pi}{4},\,\sin\frac{\pi}{4}\Bigr) =\Bigl(\frac{\sqrt2}{2},\,\frac{\sqrt2}{2}\Bigr)$ 的方向导数。
\end{exercise}
\begin{solution}
\begin{enumerate}
  \item \textbf{构造辅助函数并求导}：
  按定义，将 $p_0=(1,1)$ 和 $\vec l=\bigl(\frac{\sqrt2}{2},\frac{\sqrt2}{2}\bigr)$ 代入，构造一元函数
  \(\displaystyle u(t)=f\!\left(1+\frac{\sqrt2}{2}t,\;1+\frac{\sqrt2}{2}t\right)
  =\left(1+\frac{\sqrt2}{2}t\right)^{\!2}\)，
  从而
  \(\displaystyle u'(t)=2\left(1+\frac{\sqrt2}{2}t\right)\cdot\frac{\sqrt2}{2}
  \implies u'(0)=\sqrt2\)。
  因此 $\displaystyle\frac{\partial f}{\partial\vec l}(1,1)=\sqrt2$。
  \textbf{物理解释}：在点 $(1,1)$ 处，函数值 $f=1$；沿 $45^\circ$ 方向（东北方向）移动时，$f$ 以速率 $\sqrt2$ 增长——这比单独沿 $x$ 或 $y$ 轴方向的增速都要大（因为同时利用了两个变量的增长）
\end{enumerate}
\end{solution}

\subsubsection{偏导数}

在所有可能的无限多个方向中，有两个方向具有特殊的地位——那就是沿着坐标轴正向的单位向量 $\vec i=(1,0)$ 和 $\vec j=(0,1)$。

\begin{definition}[一阶偏导数 / First-order Partial Derivative] \label{def:partial-derivative}
设 $f:D\to\mathbb R$ 定义在开集 $D\subset\mathbb R^2$ 上，$p_0=(x_0,y_0)\in D$。

沿 $\vec i=(1,0)$ 方向的方向导数称为 $f$ 在 $p_0$ 处对 $x$ 的\textbf{一阶偏导数}（Partial Derivative with respect to $x$）；类似地，沿 $\vec j=(0,1)$ 方向的方向导数称为 $f$ 在 $p_0$ 处对 $y$ 的\textbf{一阶偏导数}。它们分别记作 \[\displaystyle \frac{\partial f}{\partial x}(p_0),\ f_x(p_0),\ f_1'(p_0),\ \frac{\partial f}{\partial x}\biggr|_{(x_0,y_0)}~~~~~~~~\displaystyle \frac{\partial f}{\partial y}(p_0),\ f_y(p_0),\ f_2'(p_0)\]
\end{definition}

把方向导数的定义放到 $x,y$ 方向，就是 \[\displaystyle \frac{\partial f}{\partial x}(x_0,y_0)=\lim_{t\to 0}\frac{f(x_0+t,\,y_0)-f(x_0,y_0)}{t},~~\displaystyle \frac{\partial f}{\partial y}(x_0,y_0)=\lim_{t\to 0}\frac{f(x_0,\,y_0+t)-f(x_0,y_0)}{t}\]
固定 $y=y_0$，令 $\varphi(x):=f(x,y_0)$（这就是我们讨论过的 ``边缘函数''），或者固定 $x=x_0$ 并令 $\psi(y):=f(x_0,y)$：
\[\displaystyle \frac{\partial f}{\partial x}(x_0,y_0) =\lim_{t\to 0}\frac{\varphi(x_0+t)-\varphi(x_0)}{t} =\varphi'(x_0),~~\dfrac{\partial f}{\partial y}(x_0,y_0)=\psi'(y_0)\]

\begin{exercise}[在某点处求偏导数值] \label{ex:pd-at-point}
设 $f(x,y)=\dfrac{y^2}{1+x}$，求 $f_x(3,2)$ 和 $f_y(3,2)$。
\end{exercise}
\begin{solution}
\begin{enumerate}
  \item \textbf{分别求两个偏导数}：
  固定 $y=2$，得到边缘函数 $\varphi(x)=f(x,2)=\dfrac{4}{1+x}$；固定 $x=3$，得到边缘函数 $\psi(y)=f(3,y)=\dfrac{y^2}{4}$。于是
  \(\displaystyle \varphi'(x)=\frac{-4}{(1+x)^2}\Longrightarrow f_x(3,2)=\varphi'(3)=\frac{-4}{16}=-\frac{1}{4}\)，且 \(\displaystyle \psi'(y)=\frac{2y}{4}=\frac{y}{2}\Longrightarrow f_y(3,2)=\psi'(2)=1\)。
\end{enumerate}
\end{solution}

\begin{exercise}[求偏导函数（解析表达式）] \label{ex:pd-analytic}
求 $z=\arctan\dfrac{y}{x}$ 的两个偏导函数。
\end{exercise}
\begin{solution}
\begin{enumerate}
  \item \textbf{对 $x$ 求偏导}：
  视 $y$ 为常数，应用复合函数链式法则：
  \(\displaystyle \frac{\partial z}{\partial x}
  =\frac{1}{1+(y/x)^2}\cdot\left(\frac{y}{x}\right)'_{\!x}
  =\frac{x^2}{x^2+y^2}\cdot\left(-\frac{y}{x^2}\right)
  =\boxed{-\frac{y}{x^2+y^2}}\)。
  \item \textbf{对 $y$ 求偏导}：
  视 $x$ 为常数，同理应用链式法则：\(\displaystyle \frac{\partial z}{\partial y} =\frac{1}{1+(y/x)^2}\cdot\frac{1}{x} =\frac{x^2}{x^2+y^2}\cdot\frac{1}{x} =\boxed{\frac{x}{x^2+y^2}}\)。
\end{enumerate}
\end{solution}

\begin{exercise}[复合函数的偏导数] \label{ex:pd-composite}
求下列函数的所有一阶偏导数：
\begin{shortexenum}
  \shortitem{1}{$f(x,y)=x^2\sin 3y$；} &
  \shortitem{2}{$f(x,y)=(3xy+2x)^5$。}
\end{shortexenum}
\end{exercise}
\begin{solution}
\begin{enumerate}
  \item[(1)] \textbf{处理 $f(x,y)=x^2\sin 3y$}：
  对 $x$ 求导时视 $\sin 3y$ 为常系数，由幂法则得
  \(\displaystyle \frac{\partial f}{\partial x}=2x\sin 3y\)。
  对 $y$ 求导时视 $x^2$ 为常系数，内层函数 $3y$ 对 $y$ 的导数为 $3$，由链式法则得
  \(\displaystyle \frac{\partial f}{\partial y}=x^2\cdot 3\cos 3y=3x^2\cos 3y\)。
  \item[(2)] \textbf{处理 $f(x,y)=(3xy+2x)^5$}：
  令中间变量 $u=3xy+2x$，则 $f=u^5$。对 $x$ 求导得 \(\displaystyle \frac{\partial f}{\partial x}=5u^4\cdot\frac{\partial u}{\partial x} =5(3xy+2x)^4\cdot(3y+2)\)。
  对 $y$ 求导得 $\dfrac{\partial f}{\partial y}=5u^4\cdot\dfrac{\partial u}{\partial y} =5(3xy+2x)^4\cdot 3x=15x(3xy+2x)^4$。
\end{enumerate}
\end{solution}

\begin{remark}[偏导数的几何意义] \label{remark:pd-geometry}
将二元函数写成 $z=f(x,y)\;((x,y)\in D)$，其图像是 $\mathbb R^3$ 中的一张曲面 $\Gamma_f=\{(x,y,f(x,y))\mid(x,y)\in D\}$。取定点 $p_0=(x_0,y_0)\in D$，考虑过此点的两条特殊平面截线：
\begin{itemize}
  \item 平面 $y=y_0$ 与曲面的交线 $\begin{cases}z=f(x,y),\\y=y_0,\end{cases}$ 是曲面上一条 ``纵向'' 曲线。偏导数 $\dfrac{\partial f}{\partial x}(x_0,y_0)$ 正是这条曲线在点 $(x_0,y_0,f(x_0,y_0))$ 处切线的\textbf{斜率}（相对于 $x$ 轴）。
  \item 平面 $x=x_0$ 与曲面的交线 $\begin{cases}z=f(x,y),\\x=x_0,\end{cases}$，类似地，$\dfrac{\partial f}{\partial y}(x_0,y_0)$ 是这条 ``横向'' 截线的切线斜率。
\end{itemize}

这两条切线（一条 ``东西走向''，一条 ``南北走向''）张成的唯一平面恰好就是曲面在该点的\textbf{切平面}（Tangent Plane）——这一结论将在本节后半部分的可微性讨论中得到严格的数学表述。
\end{remark}

\begin{exercise}[仅在某些方向存在方向导数] \label{ex:pdd-exists-limited}
讨论分段函数
\[
f(x,y)=
\begin{cases}
\dfrac{2xy}{x^2+y^2},&(x,y) \neq (0,0),\\[10pt]
1,&(x,y)=(0,0),
\end{cases}
\]
在原点处沿各个方向的方向导数是否存在。
\end{exercise}
\begin{solution}
\begin{enumerate}
  \item \textbf{参数化方向与辅助函数}：
  任意方向可表示为单位向量 $\vec l=(\cos\theta,\sin\theta)$，$\theta\in[0,2\pi)$。构造原点 $p_0=(0,0)$ 处的辅助函数：\(\displaystyle u(t)=f(t\cos\theta,t\sin\theta)=\frac{2t^2\cos\theta\sin\theta}{t^2(\cos^2\theta+\sin^2\theta)}=\sin 2\theta\ (t \neq 0),\quad u(0)=f(0,0)=1\)。

  \item \textbf{分析方向导数的定义式}：
  计算差商的极限：\(\displaystyle u'(0)=\lim_{t\to 0}\frac{u(t)-u(0)}{t}=\lim_{t\to 0}\frac{\sin 2\theta-1}{t}\)。
  我们分情况讨论该极限：
  \begin{itemize}
    \item \textbf{当 $\sin 2\theta=1$ 时}（即 $\theta=\pi/4$ 或 $5\pi/4$）：分子恒为零，因此 $u'(0)=0$ —— 方向导数存在且为零。这两个方向正是 $\vec l=\pm\bigl(\frac{\sqrt2}{2},\frac{\sqrt2}{2}\bigr)$。
    \item \textbf{当 $\sin 2\theta \neq 1$ 时}：分子是非零常数，而分母趋于零，极限不存在（发散到 $\pm\infty$）—— 此时方向导数不存在。
  \end{itemize}

  \item \textbf{得出结论}：
  此外注意到 $t \neq 0$ 时 $u(t)\equiv\sin 2\theta$ 是不随 $t$ 变化的常数，但 $u(0)=1 \neq \sin 2\theta$（除非 $\sin 2\theta=1$），故 $u(t)$ 在 $t=0$ 处甚至不连续，当然不可导。结论是：$f$ 在原点处仅在上述两个特定方向上存在方向导数
\end{enumerate}
\end{solution}

\begin{note}
这个例子说明即使一个函数在某点沿某些特殊方向有方向导数，也完全不保证它在其他方向上有良好的行为。更不用说 ``所有方向导数都存在'' 了。这也预示了 ``偏导数存在'' 和 ``可微'' 之间存在着深刻的差距——下面我们就来系统地探讨这些关系。
\end{note}

\subsubsection{梯度的定义}

有了上面的链式法则公式后，一个自然的问题出现了：\textbf{沿哪个方向的方向导数最大？最大值是多少？}将方向导数视为单位方向 $\vec l$ 的函数，考虑最大方向导数问题 \(\displaystyle \max_{\|\vec l\|=1}\frac{\partial f}{\partial\vec l}(p_0)\)。由 Cauchy--Schwarz 不等式：
\begin{align*}
\left(\sum_{i=1}^{n}\frac{\partial f}{\partial x_i}(p_0)\,l_i\right)^{\!2}
\leqslant
   \left(\sum_{i=1}^{n}\left(\frac{\partial f}{\partial x_i}(p_0)\right)^{\!2}\right)
   \left(\sum_{i=1}^{n}l_i^2\right)=
   \sum_{i=1}^{n}\left(\frac{\partial f}{\partial x_i}(p_0)\right)^{\!2}
   \qquad(\text{因为 }\|\vec l\|^2=1).
\end{align*}

等号成立当且仅当 $\vec l$ 与向量
$\displaystyle\left(
  \frac{\partial f}{\partial x_1}(p_0),\,
  \frac{\partial f}{\partial x_2}(p_0),\,
  \dots,\,
  \frac{\partial f}{\partial x_n}(p_0)
\right)$
同向平行。

\begin{definition}[梯度 / Gradient] \label{def:gradient}
设 $D\subset\mathbb R^n$ 为开区域，$f:D\to\mathbb R$ 在 $p_0\in D$ 处可微。
称向量
\[
\boxed{
\operatorname{grad}f(p_0)
:=\left(
     \frac{\partial f}{\partial x_1}(p_0),\,
     \frac{\partial f}{\partial x_2}(p_0),\,
     \dots,\,
     \frac{\partial f}{\partial x_n}(p_0)
   \right)}
\]
为 $f$ 在 $p_0$ 处的\textbf{梯度}（Gradient），也记作 $\nabla f(p_0)$。
\end{definition}

\begin{note}[名称来源]
“梯度”一词源自英文 Gradient，原意为“坡度、倾斜度”。直观地说，梯度指向函数值增长最陡峭的方向——就像山坡上最陡的上坡方向。
\end{note}

\begin{theorem}[梯度的最优性] \label{thm:gradient-max}
设 $f$ 在 $p_0$ 处可微，$\nabla f(p_0)\neq\boldsymbol 0$。则：
\begin{enumerate}
  \item 沿梯度同方向（即 $\vec l_{\max}=\dfrac{\nabla f(p_0)}{\| \nabla f(p_0)\|}$），
        方向导数取\textbf{最大值} $\|\nabla f(p_0)\|$；
  \item 沿梯度反方向（即 $\vec l_{\min}=-\dfrac{\nabla f(p_0)}{\| \nabla f(p_0)\|}$），
        方向导数取\textbf{最小值} $-\|\nabla f(p_0)\|$；
  \item 沿任何与梯度垂直的方向（即 $\nabla f(p_0)\cdot\vec l=0$），方向导数为 $0$。
\end{enumerate}
\end{theorem}

\begin{remark}[梯度的几何意义] \label{remark:gradient-geometry}
梯度 $\nabla f(p_0)$ 包含了丰富的几何信息：

\textbf{(1) 最速上升方向}：$\nabla f(p_0)$ 指向函数 $f$ 增长最快的方向，
其模长 $\|\nabla f(p_0)\|$ 就是该最大增长率。
反之，$-\nabla f(p_0)$ 指向函数下降最快的方向。

\textbf{(2) 等值面的法向量}：梯度 $\nabla f(p_0)$ 垂直于过 $p_0$ 点的等值面
$\{x\mid f(x)=f(p_0)\}$。
这是因为沿等值面的切方向移动时函数值不变（方向导数为零），
而梯度恰与所有切方向垂直。

\textbf{(3) 变化率的度量}：$\|\nabla f(p_0)\|$ 越大，说明函数在该点附近的“地形”越陡峭；
$\|\nabla f(p_0)\|=0$ 则说明该点是“平地”（驻点/临界点）。
\end{remark}


\begin{property}[梯度的代数运算法则] \label{property:gradient-algebra}
设 $f,g:D\to\mathbb R$ 均为可微函数，$c\in\mathbb R$ 为常数，$\phi:\mathbb R\to\mathbb R$
为一元可微函数。则：
\begin{enumerate}
  \item\textbf{(齐次性)} :$\nabla(cf)=c\,\nabla f$；
  \item\textbf{(加法法则)}: $\nabla(f\pm g)=\nabla f\pm \nabla g$；
  \item\textbf{(乘积法则)}: $\nabla(fg)=f\,\nabla g+g\,\nabla f$；
        （注意与一元乘积法则 $(fg)'=f'g+fg'$ 完全类比）
  \item\textbf{(复合函数法则)} :$\nabla(\phi\circ f)=(\phi'\circ f)\,\nabla f$。
        （也称为\textbf{链式法则}：外层求导乘以内层梯度）
\end{enumerate}
\end{property}


\begin{exercise}[范数函数的梯度] \label{ex:grad-norm}
设 $p=(x,y,z)$，$f(p)=\|p\|$（即点到原点的距离）。求 $\nabla f(p)$（当 $p\neq\boldsymbol 0$ 时）。
\end{exercise}
\begin{solution}
\begin{enumerate}
  \item \textbf{写出函数的显式表达式}：
  $f(p)=\sqrt{x^2+y^2+z^2}=(x^2+y^2+z^2)^{1/2}$。

  \item \textbf{逐个求偏导数}：
  注意 $p\neq\boldsymbol 0$ 保证分母非零。
  $\dfrac{\partial f}{\partial x} = \dfrac{1}{2}(x^2+y^2+z^2)^{-1/2}\cdot 2x = \dfrac{x}{\sqrt{x^2+y^2+z^2}} = \dfrac{x}{\|p\|}$，
  同理 $\dfrac{\partial f}{\partial y} = \dfrac{y}{\|p\|}$，$\dfrac{\partial f}{\partial z} = \dfrac{z}{\|p\|}$。

  \item \textbf{构造梯度向量}：
  组合成分量形式：\(\displaystyle \boxed{\nabla f(p)=\frac{1}{\|p\|}(x,y,z)=\frac{p}{\|p\|}=\nabla\|p\|}\)。

  \item \textbf{物理解释}：距离函数 $f(p)=\|p\|$ 的梯度指向远离原点的径向外侧方向，
  模长恒为 $1$。这意味着从空间任一点（原点除外）
  沿径向移动单位距离，距离值正好增加 $1$ —— 这符合“距离等于半径”的直觉。
\end{enumerate}
\end{solution}

\subsubsection{方向导数的链式法则表示}

\begin{property}[方向导数的链式法则公式] \label{prop:directional-via-chain}
设 $f:D\to\mathbb R$ 在 $p_0\in D$ 处可微，$\vec l=(l_1,l_2,\dots,l_n)$ 为单位向量。
则 $f$ 在 $p_0$ 处沿 $\vec l$ 的方向导数为
\[
\boxed{\;
\frac{\partial f}{\partial\vec l}(p_0)
=\sum_{i=1}^{n}\frac{\partial f}{\partial x_i}(p_0)\cdot l_i
=\nabla f(p_0)\cdot\vec l
\;},
\]
其中 $\nabla f(p_0)=\left(\dfrac{\partial f}{\partial x_1}(p_0),\dots,
\dfrac{\partial f}{\partial x_n}(p_0)\right)$ 为 $f$ 在 $p_0$ 处的\textbf{梯度行向量}。
\end{property}

特别地，对于 $n=3$ 的情形，
\(\displaystyle \dfrac{\partial f}{\partial\vec l}(p_0)
=\dfrac{\partial f}{\partial x}(p_0)\cos\alpha
+\dfrac{\partial f}{\partial y}(p_0)\cos\beta
+\dfrac{\partial f}{\partial z}(p_0)\cos\gamma\)。

\begin{exercise}[沿特定方向的方向导数] \label{ex:grad-exp}
设 $z=xe^{xy}$，$p_0=(1,1)$。求沿与向量 $(1,1)$ 同向的方向导数 $\dfrac{\partial z}{\partial\vec l}(1,1)$。
\end{exercise}
\begin{solution}
\begin{enumerate}
  \item \textbf{计算偏导数及梯度}：
  \(\displaystyle \frac{\partial z}{\partial x}=e^{xy}+xy\cdot e^{xy}=e^{xy}(1+xy),\qquad
  \frac{\partial z}{\partial y}=x\cdot e^{xy}\cdot x=x^2e^{xy}\)。
  在 $p_0(1,1)$ 处取值可得
  $\left.\dfrac{\partial z}{\partial x}\right|_{(1,1)}=e^1(1+1)=2e$，
  $\left.\dfrac{\partial z}{\partial y}\right|_{(1,1)}=1^2\cdot e=e$。
  故梯度为 $\nabla z(1,1)=(2e,\;e)$。

  \item \textbf{确定单位方向向量}：
  方向 $\vec l$ 应与 $(1,1)$ 同方向，而 $\|(1,1)\|=\sqrt2$，故
  $\vec l=\dfrac{1}{\sqrt2}(1,1)=\left(\dfrac{\sqrt2}{2},\,\dfrac{\sqrt2}{2}\right)$。

  \item \textbf{做内积求方向导数}：
  利用公式 $\dfrac{\partial z}{\partial\vec l} = \nabla z \cdot \vec l$，得
  $\dfrac{\partial z}{\partial\vec l}(1,1) = \nabla z(1,1)\cdot\vec l
  =(2e,\;e)\cdot\left(\dfrac{\sqrt2}{2},\,\dfrac{\sqrt2}{2}\right)
  =2e\cdot\dfrac{\sqrt2}{2}+e\cdot\dfrac{\sqrt2}{2}
  =\boxed{\dfrac{3\sqrt2}{2}\,e}$。
\end{enumerate}
\end{solution}

\begin{exercise}[给定角度] \label{ex:grad-log}
设 $u=\ln(x^2+y^2)$，$p_0=(1,1)$。方向 $\vec l$ 与 $x$ 轴正向夹角为 $60^\circ$。求 $\left.\dfrac{\partial u}{\partial\vec l}\right|_{p_0}$。
\end{exercise}
\begin{solution}
\begin{enumerate}
  \item \textbf{方法一：利用梯度公式法}
  先计算偏导数：$\dfrac{\partial u}{\partial x}=\dfrac{2x}{x^2+y^2}$，$\dfrac{\partial u}{\partial y}=\dfrac{2y}{x^2+y^2}$。
  在 $(1,1)$ 处与对应的单位方向向量分别为
  \(\displaystyle \left.\nabla u\right|_{(1,1)}=\left(\frac{2}{2},\,\frac{2}{2}\right)=(1,\,1),\qquad
  \vec l=(\cos 60^\circ,\;\sin 60^\circ) = \left(\frac12,\;\frac{\sqrt3}{2}\right)\)。
  于是方向导数为
  $\left.\dfrac{\partial u}{\partial\vec l}\right|_{(1,1)}
  =(1,1)\cdot\left(\dfrac12,\dfrac{\sqrt3}{2}\right)
  = \dfrac{1+\sqrt3}{2}$。

  \item \textbf{方法二：按定义验证（辅助函数法）}
  构造辅助函数 $\varphi(t)=u(p_0+t\vec l)=u\left(1+\dfrac{t}{2},\;1+\dfrac{\sqrt3\,t}{2}\right)$，
  得
  \(\displaystyle \varphi(t)= \ln\left[\left(1+\frac{t}{2}\right)^{\!2} +\left(1+\frac{\sqrt3\,t}{2}\right)^{\!2}\right] = \ln\left[t^2+(\sqrt3+1)t+2\right]\)。
  求导并在 $t=0$ 处取值，得
  $\varphi'(t)=\dfrac{2t+(\sqrt3+1)}{t^2+(\sqrt3+1)t+2}$，
  从而 $\varphi'(0)=\dfrac{1+\sqrt3}{2}$。
  两种方法结果完全一致，最终答案为 $\boxed{\dfrac{1+\sqrt3}{2}}$。
\end{enumerate}
\end{solution}

\begin{remark}[两种方法的比较]
定义法虽然在计算效率上不占优势，但它提供了“回归定义”的机会——当对梯度公式的适用条件有疑虑时，用定义直接计算是最可靠的检验手段。
\end{remark}

\begin{exercise}[梯度为零的方向导数] \label{ex:grad-zero-dir}
设 $f(x,y)=xe^y$。求 $(1,1)$ 处的梯度，并求沿 $\vec i-\vec j$ 方向的方向导数。
\end{exercise}
\begin{solution}
\begin{enumerate}
  \item \textbf{计算梯度}：
  $\dfrac{\partial f}{\partial x}=e^y$，$\dfrac{\partial f}{\partial y}=xe^y$。
  在 $(1,1)$ 处：
  $\nabla f(1,1)=(e^1,\;1\cdot e^1)=(e,\;e)$。

  \item \textbf{计算方向导数}：
  方向为 $\vec i-\vec j=(1,-1)$，归一化得单位向量
  $\vec l=\dfrac{(1,-1)}{\sqrt{1^2+(-1)^2}} =\left(\dfrac{\sqrt2}{2},\,-\dfrac{\sqrt2}{2}\right)$。
  做内积：
  $\dfrac{\partial f}{\partial\vec l}(1,1) = (e,e)\cdot\left(\dfrac{\sqrt2}{2},-\dfrac{\sqrt2}{2}\right)= \dfrac{\sqrt2}{2}e-\dfrac{\sqrt2}{2}e = \boxed{0}$。

  \item \textbf{几何解释}：
  方向 $(1,-1)$ 恰好与梯度 $(e,e)$ 垂直（因为内积为 0）。
  因此沿此方向移动时，函数值的一阶变化率为零——这正是等值线（此处为 $xe^y=e$）的切线方向！
  事实上，$f(x,y)=xe^y$ 过点 $(1,1)$ 的等值线在其上的切线斜率可以通过隐函数求导验证为 $-1$，与方向 $(1,-1)$ 完全一致。
\end{enumerate}
\end{solution}

\subsubsection{可微性与全微分}

\begin{definition}[二元函数的可微性] \label{def:differentiability}
设 $f:D\to\mathbb R$ 定义在开集 $D\subset\mathbb R^2$ 上，$p_0=(x_0,y_0)\in D$。要验证该函数在该点可微，需要验证：
  \[
  \lim_{(h_1,h_2)\to(0,0)} \frac{f(x_0+h_1,y_0+h_2) - f(x_0,y_0) - (f_x(x_0,y_0)h_1 + f_y(x_0,y_0)h_2)}{\sqrt{h_1^2+h_2^2}} = 0.
  \]
其中$f_x(x_0,y_0),f_y(x_0,y_0)$有时可以正常求，但有时可能需要借助定义去求。此时称线性映射 $\mathrm{d}f(p_0):\;\boldsymbol h\longmapsto A\boldsymbol h=a_1h_1+a_2h_2$ 为 $f$ 在 $p_0$ 处的\textbf{微分}。
\end{definition}

\begin{exercise}[证明某函数在原点可微] \label{ex:diff-prove}
证明函数
\[
f(x,y)=
\begin{cases}
x+y+(x^2+y^2)\sin\dfrac{1}{x^2+y^2},&(x,y) \neq (0,0),\\[8pt]
0,&(x,y)=(0,0)
\end{cases}
\]
在原点处可微，并求其微分。
\end{exercise}
\begin{solution}
\begin{enumerate}
  \item 首先，通过定义计算函数在原点 $(0,0)$ 处的偏导数：
  \begin{align*}
  f_x(0,0) &= \lim_{x\to 0} \frac{f(x,0)-f(0,0)}{x} = \lim_{x\to 0} \frac{x + x^2\sin\frac{1}{x^2}}{x} = \lim_{x\to 0} \left(1 + x\sin\frac{1}{x^2}\right) = 1,\\
  f_y(0,0) &= \lim_{y\to 0} \frac{f(0,y)-f(0,0)}{y} = \lim_{y\to 0} \left(1 + y\sin\frac{1}{y^2}\right) = 1.
  \end{align*}
  其中用到了 $\left|x\sin\frac{1}{x^2}\right| \leqslant |x|$ 以及同样的 $y$ 向估计。
  若函数可微，其在原点的线性增量应为 $A h_1 + B h_2 = 1\cdot h_1 + 1\cdot h_2 = h_1 + h_2$。
  \item   根据二元函数可微的定义，我们需要考察增量误差除以距离的极限是否为 $0$。令 $\rho = \sqrt{h_1^2+h_2^2}$，需要验证：
  \[
  \lim_{(h_1,h_2)\to(0,0)} \frac{f(h_1,h_2) - f(0,0) - (f_x(0,0)h_1 + f_y(0,0)h_2)}{\sqrt{h_1^2+h_2^2}} = 0.
  \]
  将函数表达式及偏导数值代入该极限的表达式中：
  \begin{align*}
  \text{原式} &= \lim_{(h_1,h_2)\to(0,0)} \frac{\left( h_1+h_2+(h_1^2+h_2^2)\sin\frac{1}{h_1^2+h_2^2} \right) - 0 - (h_1 + h_2)}{\sqrt{h_1^2+h_2^2}} \\
  &= \lim_{(h_1,h_2)\to(0,0)} \frac{(h_1^2+h_2^2)\sin\frac{1}{h_1^2+h_2^2}}{\sqrt{h_1^2+h_2^2}}.
  \end{align*}
  利用 $\rho = \sqrt{h_1^2+h_2^2}$ 进行替换，当 $(h_1,h_2)\to(0,0)$ 时，$\rho \to 0^+$。极限转化为 \(\displaystyle \lim_{\rho\to 0^+} \frac{\rho^2 \sin\frac{1}{\rho^2}}{\rho} = \lim_{\rho\to 0^+} \rho \sin\frac{1}{\rho^2}\)。
  \item  由于正弦函数在实数域上有界，即 $\left|\sin\frac{1}{\rho^2}\right| \leqslant 1$，因此可以放缩：\(\displaystyle 0 \leqslant \left| \rho \sin\frac{1}{\rho^2} \right| \leqslant \rho\)。当 $\rho \to 0^+$ 时，$\rho \to 0$。由夹逼定理可知 \(\displaystyle \lim_{\rho\to 0^+} \rho \sin\frac{1}{\rho^2} = 0\)。
  因为该极限等于 $0$，说明函数的全增量与线性部分之差是关于 $\rho$ 的高阶无穷小。因此，函数 $f(x,y)$ 在原点处可微，且其微分为 \(\displaystyle \boxed{\mathrm{d}f(0,0) = \mathrm{d}x + \mathrm{d}y}\)（或写为增量形式：\(\displaystyle \mathrm{d}f(0,0)=h_1+h_2\)）。
\end{enumerate}
\end{solution}

\begin{theorem}[可微可导连续的关系] \label{thm:diff-cont}
\begin{itemize}
\item 可微 $\Rightarrow$ 可导 且 可微 $\Rightarrow$ 连续（可微是强条件）。
\item 可导 $\nRightarrow$ 连续，连续 $\nRightarrow$ 可导（两者无因果关系）。
\item 可导 且 连续 $\nRightarrow$ 可微：即使函数在某点是连续的，且两个偏导数都存在，它依然可能不可微。偏导数只限制了平行于坐标轴（$x$轴和 $y$轴）两个方向的变化行为。即使这两个方向有切线且函数连续，也无法保证其他斜向（如 $y=x$ 方向）的变化能与这两个方向协调，从而无法形成一个统一的、平滑的“切平面”（即不满足可微的定义）。
\item 偏导数连续 $\Rightarrow$ 可微（充分不必要条件）：如果偏导数不仅存在，而且在某点邻域内是连续变化的（没有跳跃突变），那么函数在该点一定可微。
\end{itemize}
\end{theorem}

\begin{theorem}[可微的充分条件] \label{thm:diff-sufficient}
设 $f$ 在 $p_0=(x_0,y_0)$ 的某个邻域内存在一阶偏导数 $f_x$, $f_y$，且 $f_x$, $f_y$ 在 $p_0$ 处\textbf{连续}，则 $f$ 在 $p_0$ 处可微。
\end{theorem}

\begin{proof}
考虑函数增量，将其拆分为两步（先变 $x$ 再变 $y$）：
  \[
  f(x_0+h_1,\,y_0+h_2)-f(x_0,y_0)
  =\underbrace{f(x_0+h_1,\,y_0+h_2)-f(x_0,\,y_0+h_2)}_{\Delta_x}
  +\underbrace{f(x_0,\,y_0+h_2)-f(x_0,y_0)}_{\Delta_y}.
  \]
对 $\Delta_x$ 应用一元拉格朗日中值定理（固定第二个变量 $y_0+h_2$，对第一个变量用中值定理），得 \(\displaystyle \Delta_x = f_x(x_0+\theta_1 h_1,\,y_0+h_2)\cdot h_1, \qquad\theta_1\in(0,1)\)。对 $\Delta_y$ 同理，\(\displaystyle \Delta_y = f_y(x_0,\,y_0+\theta_2 h_2)\cdot h_2, \qquad\theta_2\in(0,1)\)。
合并得：
  \(\displaystyle f(p_0+\boldsymbol h)-f(p_0) = f_x(p_0)h_1+f_y(p_0)h_2 +\alpha h_1+\beta h_2\)。
其中
  \(\displaystyle \alpha:=f_x(x_0+\theta_1 h_1,\,y_0+h_2)-f_x(x_0,y_0),\qquad
  \beta:=f_y(x_0,\,y_0+\theta_2 h_2)-f_y(x_0,y_0)\)。
由 $f_x$, $f_y$ 在 $p_0$ 处的连续性，当 $h_1,h_2\to 0$ 时 $\alpha\to 0$, $\beta\to 0$。
又
  \(\displaystyle \frac{|\alpha h_1+\beta h_2|}{\|\boldsymbol h\|} \leqslant  |\alpha|+|\beta|\xrightarrow[h\to 0]{}0\)。
故 $\alpha h_1+\beta h_2=o(\|\boldsymbol h\|)$，即 $f$ 在 $p_0$ 处可微
\end{proof}

\begin{exercise}[偏导存在但不可微] \label{ex:pd-exist-not-diff}
考察函数
\[
f(x,y)=
\begin{cases}
\dfrac{xy}{x^2+y^2},&(x,y) \neq (0,0),\\[10pt]
0,&(x,y)=(0,0).
\end{cases}
\]
\end{exercise}
\begin{solution}
\begin{enumerate}
  \item \textbf{计算偏导数以确定线性主部候选}：
  首先，我们在原点处通过定义计算偏导数：
  \[
  f_x(0,0)=\lim_{x\to 0}\frac{f(x,0)-f(0,0)}{x}=0,\qquad
  f_y(0,0)=\lim_{y\to 0}\frac{f(0,y)-f(0,0)}{y}=0.
  \]
  若函数在原点可微，其在原点的线性增量必定为 \(\displaystyle \mathrm{d}f = f_x(0,0)h_1 + f_y(0,0)h_2 = 0\)。
  \item \textbf{利用极限定义验证可微性}：
  根据二元函数可微的定义，需要考察全增量减去线性主部后，除以距离 $\rho = \sqrt{h_1^2+h_2^2}$ 的极限是否为 $0$。将函数表达式及偏导数值代入可微性检验式，得到
  \begin{align*}
  \lim_{(h_1,h_2)\to(0,0)}
  \frac{f(h_1,h_2)-f(0,0)-(f_x(0,0)h_1+f_y(0,0)h_2)}{\sqrt{h_1^2+h_2^2}}
  = \lim_{(h_1,h_2)\to(0,0)} \frac{h_1 h_2}{(h_1^2+h_2^2)^{\frac{3}{2}}}.
  \end{align*}
  \item \textbf{计算并分析该极限}：
  我们让点 $(h_1,h_2)$ 沿直线 $h_2 = h_1$（且 $h_1 > 0$）趋近于 $(0,0)$，代入上式计算：
  \[
  \lim_{h_1\to 0^+} \frac{h_1 \cdot h_1}{(h_1^2+h_1^2)^{\frac{3}{2}}}
  = \lim_{h_1\to 0^+} \frac{h_1^2}{2\sqrt{2}h_1^3}
  = \lim_{h_1\to 0^+} \frac{1}{2\sqrt{2}h_1}
  = +\infty.
  \]
  由于沿该路径趋近时，极限趋于无穷大（显然不等于 $0$，甚至极限不存在），说明函数的全增量与线性部分之差并不是关于距离 $\rho$ 的高阶无穷小。
  \item \textbf{得出结论}：
  由可微的定义可知，由于误差项除以距离的极限不为 $0$，因此函数 $f(x,y)$ 在 $(0,0)$ 处虽然偏导数存在，但\textbf{不可微}。
\end{enumerate}
\begin{note}[重要教训]
\textbf{偏导数存在 $\nRightarrow$ 函数可微！}甚至不能保证函数连续！这是因为偏导数只捕捉了“沿坐标轴方向”的变化信息，完全忽略了“其他方向”（如对角方向）的行为。
\end{note}
\end{solution}

\begin{exercise}[可微但偏导数不连续] \label{ex:diff-but-pd-disc}
设
\[
f(x,y)=
\begin{cases}
(x^2+y^2)\sin\dfrac{1}{x^2+y^2},&(x,y) \neq (0,0),\\[8pt]
0,&(x,y)=(0,0).
\end{cases}
\]
证明：(1) $f$ 在原点处偏导数存在；(2) $f_x$, $f_y$ 在原点处不连续；(3) 但 $f$ 在原点处可微。
\end{exercise}
\begin{solution}
\begin{enumerate}
  \item \textbf{验证原点处的偏导数存在}：
  利用定义计算：\(\displaystyle f_x(0,0)=\lim_{t\to 0}\frac{f(t,0)-0}{t}
  =\lim_{t\to 0}t\sin\frac{1}{t^2}=0\)，且 \(\displaystyle f_y(0,0)=\lim_{t\to 0}\frac{f(0,t)-0}{t}=0\)。
  偏导存在且都为 $0$。
  \item \textbf{验证偏导数在原点处不连续}：
  当 $(x,y) \neq (0,0)$ 时，计算一般的偏导数表达式：
  \begin{align*}
  f_x(x,y)&=2x\sin\frac{1}{x^2+y^2} +(x^2+y^2)\cos\frac{1}{x^2+y^2}\cdot\!\left(-\frac{2x}{(x^2+y^2)^2}\right)\\
  &=2x\sin\frac{1}{x^2+y^2}-\frac{2x}{x^2+y^2}\cos\frac{1}{x^2+y^2},\\[8pt]
  f_y(x,y)&=2y\sin\frac{1}{x^2+y^2}-\frac{2y}{x^2+y^2}\cos\frac{1}{x^2+y^2}.
  \end{align*}
  引入极坐标变换 $x = r\cos\theta, y = r\sin\theta$。当 $(x,y) \to (0,0)$ 时，等价于 $r \to 0^+$。代入 $f_x$ 的表达式中：
  \(\displaystyle f_x(r\cos\theta, r\sin\theta)
  = 2r\cos\theta\sin\frac{1}{r^2} - \frac{2r\cos\theta}{r^2}\cos\frac{1}{r^2}
  = 2r\cos\theta\sin\frac{1}{r^2} - \frac{2\cos\theta}{r}\cos\frac{1}{r^2}\)。
  当 $r \to 0^+$ 时，第一项由于有界量乘无穷小量，因此 $2r\cos\theta\sin\frac{1}{r^2} \to 0$。
  但是对于第二项，若我们固定角度 $\theta$ 使得 $\cos\theta \neq 0$（例如取 $\theta=0$ 沿 $x$ 轴趋近），该项变为 $-\frac{2}{r}\cos\frac{1}{r^2}$。随着 $r \to 0^+$，由于 $\frac{1}{r} \to +\infty$，而 $\cos\frac{1}{r^2}$ 在 $-1$ 到 $1$ 之间剧烈振荡，导致整个表达式不仅无界且发散。
  因此，极限 $\lim\limits_{(x,y)\to(0,0)} f_x(x,y)$ 不存在，从而 $f_x$ 在原点处不连续。同理可证 $f_y$ 在原点处也不连续。
  \item \textbf{验证函数在原点处可微}：
  需检查余项是否为 $o(\|\boldsymbol h\|)$。
  \[
  \frac{f(h_1,h_2)-f(0,0)-0\cdot h_1-0\cdot h_2}{\|\boldsymbol h\|} =\frac{(h_1^2+h_2^2)\sin\dfrac{1}{h_1^2+h_2^2}}{\sqrt{h_1^2+h_2^2}} =\sqrt{h_1^2+h_2^2}\sin\frac{1}{h_1^2+h_2^2} \xrightarrow[\|\boldsymbol h\|\to 0]{}0.
  \]
  由于极限为 $0$，所以 $f$ 在原点处依然是可微的。
\end{enumerate}
\end{solution}

\begin{definition}[全微分 / Total Differential] \label{property:total-diff-formula}
若 $z=f(x,y)$ 在点 $(x,y)$ 处可微，则其\textbf{全微分}为 \(\displaystyle \boxed{\; \mathrm{d}z=\mathrm{d}f=\frac{\partial z}{\partial x}\,\mathrm{d}x+\frac{\partial z}{\partial y}\,\mathrm{d}y\;}\)。
\end{definition}

这里 $\mathrm{d}x=h_1$ 和 $\mathrm{d}y=h_2$ 分别是自变量 $x$ 和 $y$ 的\textbf{微分}（即增量的线性主部）。特别地，取 $f(x,y)=x$ 则 $\mathrm{d}z=\mathrm{d}x$；取 $f(x,y)=y$ 则 $\mathrm{d}z=\mathrm{d}y$。这与一元情形下的符号约定保持一致。

对三元函数 $u=f(x,y,z)$，若可微则全微分为 $\mathrm{d}u=\frac{\partial u}{\partial x}\mathrm{d}x+\frac{\partial u}{\partial y}\mathrm{d}y+\frac{\partial u}{\partial z}\mathrm{d}z$。

\begin{exercise}[多项式的全微分] \label{ex:total-diff-poly}
求 $z=x^2+4xy^2+y^4$ 的全微分。
\end{exercise}
\begin{solution}
\begin{enumerate}
  \item \textbf{求各个偏导数}：\(\displaystyle \frac{\partial z}{\partial x}=2x+4y^2,\qquad \frac{\partial z}{\partial y}=8xy+4y^3\)。
  \item \textbf{套用全微分公式}：\(\displaystyle \boxed{\mathrm{d}z=(2x+4y^2)\mathrm{d}x+(8xy+4y^3)\mathrm{d}y}\)。
\end{enumerate}
\end{solution}

\begin{exercise}[混合型函数的全微分] \label{ex:total-diff-mixed}
求 $u=2x+\cos y+e^{yz}$ 的全微分。
\end{exercise}
\begin{solution}
\begin{enumerate}
  \item \textbf{求对三个变量的偏导数}：$\frac{\partial u}{\partial x}=2,\ \frac{\partial u}{\partial y}=-\sin y+z e^{yz},\ \frac{\partial u}{\partial z}=ye^{yz}$。
  \item \textbf{构造全微分}：\(\displaystyle \boxed{\mathrm{d}u=2\,\mathrm{d}x+\bigl(-\sin y+z e^{yz}\bigr)\mathrm{d}y+ye^{yz}\mathrm{d}z}\)。
\end{enumerate}
\end{solution}

\begin{exercise}[利用全微分做近似计算] \label{ex:diff-approx}
利用微分法求 $(1.03)^{1.98}$ 的近似值。
\end{exercise}
\begin{solution}
\begin{enumerate}
  \item \textbf{构造多元函数与基准点}：
  令 $f(x,y)=x^y$。观察到目标值附近的整数基准点非常容易计算，故选定 $(x_0,y_0)=(1,2)$。增量为：
  $1.03=1+0.03 \implies h_1=0.03$；$1.98=2-0.02 \implies h_2=-0.02$。

  \item \textbf{写出偏导数与局部线性化公式}：
  偏导数为 $f_x=y x^{y-1}$，$f_y=x^y\ln x$。
  代入基准点 $(x_0,y_0)=(1,2)$ 得到：
  $f(1,2)=1$，$f_x(1,2)=2\cdot 1^1=2$，$f_y(1,2)=1^2\ln 1=0$。

  \item \textbf{执行近似计算}：
  由局部线性化公式 $f(x_0+h_1,\,y_0+h_2)\approx f(x_0,y_0)+f_x(x_0,y_0)\,h_1+f_y(x_0,y_0)\,h_2$，有 \(\displaystyle (1.03)^{1.98}\approx 1+2\times 0.03+0\times(-0.02)=1+0.06=\boxed{1.06}\)。
  （对比：精确值约 $1.0597\dots$，误差极小。）
\end{enumerate}
\end{solution}

\begin{remark}[近似计算的误差估计]
全微分近似 $f(p_0+\boldsymbol h)\approx f(p_0)+ \nabla f(p_0)\cdot\boldsymbol h$ 的误差量级为 $o(\|\boldsymbol h\|)$——即误差是 $\|\boldsymbol h\|$ 的高阶无穷小。这意味着当 $\|\boldsymbol h\|$ 很小时，相对误差会很小。实际应用中常需配合二阶泰勒展开或 Lagrange 余项来做更精细的误差界估计。
\end{remark}

\subsection{向量值函数，雅可比矩阵}

\begin{definition}[向量值函数 / Vector-valued Function] \label{def:vec-func}
设 $D\subset\mathbb R^n$ 为一个点集。称映射 $f:\;D\to\mathbb R^m\ (m\geqslant 2)$ 为定义在 $D$ 上、在 $\mathbb R^m$ 中取值的\textbf{向量值函数}（或称\textbf{向量函数}、\textbf{映射}）。

将输入和输出均写成列向量形式：
\(\displaystyle \boldsymbol x=(x_1,x_2,\dots,x_n)^{\mathsf T}\in\mathbb R^n,\quad
\boldsymbol y=(y_1,y_2,\dots,y_m)^{\mathsf T}\in\mathbb R^m\)，
则 $\boldsymbol y=f(\boldsymbol x)$ 可以具体展开为 $m$ 个分量方程：
\[
\begin{cases}
y_1=f_1(x_1,x_2,\dots,x_n),\\
y_2=f_2(x_1,x_2,\dots,x_n),\\
\quad\vdots\\
y_m=f_m(x_1,x_2,\dots,x_n).
\end{cases}
\tag{$\star$}
\]
其中每个 $f_i:D\to\mathbb R$ 都是一个 $n$ 元\textbf{实值}函数，称为 $f$ 的\textbf{第 $i$ 个分量函数}。
\end{definition}

\begin{note}[关于“点集”与“开区域”的区分]
这里先把定义域写成一般点集，是因为“向量值函数”这个概念本身并不要求定义域一定开或连通。等到后面讨论极限、连续时，只需要聚点；讨论可微时，则必须把点放在开区域中，才能在各个方向上作局部扰动。这个区分与前面实值多元函数的情形完全一致。
\end{note}

\begin{note}[三种函数类型的比较]
\par\noindent{\centering
\renewcommand{\arraystretch}{1.4}
\begin{tabular}{@{}l|c|c|c@{}}
\textbf{类型} & \textbf{记号} & \textbf{自变量} & \textbf{因变量} \\
一元实值函数 & $f:\mathbb R\to\mathbb R$ & 标量 $x$ & 标量 $y$ \\
多元实值函数 & $f:\mathbb R^n\to\mathbb R$ & 向量 $\boldsymbol x$ & 标量 $y$ \\
向量值函数 & $f:\mathbb R^n\to\mathbb R^m$ & 向量 $\boldsymbol x$ & 向量 $\boldsymbol y$
\end{tabular}
\par}
向量值函数是前两种情形的自然推广：当 $m=1$ 时退化为多元实值函数；
当 $n=m=1$ 时退为一元函数。因此我们期望前述理论的大部分结论都能
“分量地”推广过来。
\end{note}

\begin{exercise}[空间螺旋线] \label{ex:helix}
向量值函数 \(\displaystyle f(t)=(\cos t,\;\sin t,\;t)^{\mathsf T}\quad (t\in\mathbb R)\) 表示三维空间中的一条\textbf{螺旋线}。若令 $\vec i=(1,0,0)^{\mathsf T}$，$\vec j=(0,1,0)^{\mathsf T}$，
$\vec k=(0,0,1)^{\mathsf T}$（标准基向量），则该函数也可表示为 \(\displaystyle f(t)=\vec i\cos t+\vec j\sin t+\vec k\,t\)。
\end{exercise}

\subsubsection{向量值函数的极限和连续性}

对任何函数，极限的 $\varepsilon$-$\delta$ 定义在本质结构上是一致的——
“当自变量趋近目标点时，函数值可以任意接近极限值”。
唯一的区别在于“接近程度”的度量方式：
对于向量值函数，我们需要使用相应空间中的\textbf{欧几里得范数}（即距离）来衡量。

\begin{definition}[向量值函数的极限] \label{def:vec-limit}
设 $D\subset\mathbb R^n$ 为区域，$f:D\to\mathbb R^m$ 为 $n$ 元 $m$ 维向量值函数。
设 $\boldsymbol a=(a_1,a_2,\dots,a_n)^{\mathsf T}\in D'$ 为 $D$ 的聚点，
$\boldsymbol A=(A_1,A_2,\dots,A_m)^{\mathsf T}\in\mathbb R^m$ 为给定向量。

若 $\forall\varepsilon>0$，$\exists\delta>0$，使得对满足 $0<\|\boldsymbol x-\boldsymbol a\|_n<\delta$ 的一切 $\boldsymbol x\in D$，均有 \(\displaystyle \|f(\boldsymbol x)-\boldsymbol A\|_m<\varepsilon\)，则称 $\boldsymbol A$ 为 $f(\boldsymbol x)$ 当 $\boldsymbol x\to\boldsymbol a$ 时的\textbf{极限}，记作 \(\displaystyle \lim_{\boldsymbol x\to\boldsymbol a}f(\boldsymbol x)=\boldsymbol A\)。

其中 $\|\cdot\|_n$ 和 $\|\cdot\|_m$ 分别是 $\mathbb R^n$ 和 $\mathbb R^m$ 中的欧几里得范数：
\[
\|\boldsymbol x-\boldsymbol a\|_n = \sqrt{\sum_{k=1}^{n}(x_k-a_k)^2},\qquad
\|f(\boldsymbol x)-\boldsymbol A\|_m = \sqrt{\sum_{k=1}^{m}\bigl(f_k(\boldsymbol x)-A_k\bigr)^2}.
\]
\end{definition}

\begin{property}[极限的分量等价性] \label{prop:vec-limit-component}
向量值函数的极限存在当且仅当每个分量函数的极限同时存在：
\[
\lim_{\boldsymbol x\to\boldsymbol a}f(\boldsymbol x)=\boldsymbol A
\quad\Longleftrightarrow\quad
\lim_{\boldsymbol x\to\boldsymbol a}f_i(\boldsymbol x)=A_i,\quad i=1,2,\dots,m.
\]
\end{property}

\begin{proof}
“$\Rightarrow$”（必要性）：由欧几里得距离的性质，对任意 $i=1,\dots,m$，恒有 $|f_i(\boldsymbol x)-A_i| \leqslant  \|f(\boldsymbol x)-\boldsymbol A\|_m$。
已知 $\|f(\boldsymbol x)-\boldsymbol A\|_m<\varepsilon$，故 $|f_i(\boldsymbol x)-A_i|<\varepsilon$ 成立，即各分量极限存在且为 $A_i$。

“$\Leftarrow$”（充分性）：设每个 $\lim\limits_{\boldsymbol x\to\boldsymbol a}f_i(\boldsymbol x)=A_i$。
对任意给定的 $\varepsilon>0$，对每个分量 $i$，都存在 $\delta_i>0$ 使得当 $0<\|\boldsymbol x-\boldsymbol a\|_n<\delta_i$ 时有 $|f_i(\boldsymbol x)-A_i|<\dfrac\varepsilon{\sqrt{m}}$。
取 $\delta=\min\{\delta_1,\dots,\delta_m\}>0$，则当 $0<\|\boldsymbol x-\boldsymbol a\|_n<\delta$ 时，所有分量的不等式同时成立，此时：
\[
\|f(\boldsymbol x)-\boldsymbol A\|_m
=\sqrt{\sum_{i=1}^{m}|f_i(\boldsymbol x)-A_i|^2}
<\sqrt{\sum_{i=1}^{m}\left(\frac{\varepsilon}{\sqrt{m}}\right)^2}
=\sqrt{m\cdot\frac{\varepsilon^2}{m}}=\varepsilon.
\]
充分性得证。
\end{proof}

\begin{definition}[向量值函数的连续性] \label{def:vec-continuity}
设 $D\subset\mathbb R^n$ 为区域，$f:D\to\mathbb R^m$ 为向量值函数，
$\boldsymbol a\in D$ 为 $D$ 的聚点。

若 $\lim\limits_{\boldsymbol x\to\boldsymbol a}f(\boldsymbol x)=f(\boldsymbol a)$，
即 $\forall\varepsilon>0$，$\exists\delta>0$，使得 $\|\boldsymbol x-\boldsymbol a\|_n<\delta$ 时
$\|f(\boldsymbol x)-f(\boldsymbol a)\|_m<\varepsilon$，
则称 $f$ 在 $\boldsymbol a$ 处\textbf{连续}。

若 $f$ 在 $D$ 的每一点都连续，则称 $f$ 在 $D$ 上连续，记作 $f\in C(D;\mathbb R^m)$。
\end{definition}

\begin{property}[连续性的分量等价性] \label{prop:vec-cont-component}
$f$ 在 $\boldsymbol a$ 处连续 $\Longleftrightarrow$每个分量函数 $f_i$ 在 $\boldsymbol a$ 处连续（$i=1,2,\dots,m$）。
\end{property}

\begin{proof}
直接由连续性的定义和极限的性质 \ref{prop:vec-limit-component} 得到：
$\lim\limits_{\boldsymbol x\to\boldsymbol a}f(\boldsymbol x)=f(\boldsymbol a)$
当且仅当对每个 $i$ 都有
$\lim\limits_{\boldsymbol x\to\boldsymbol a}f_i(\boldsymbol x)=f_i(\boldsymbol a)$。
\end{proof}

\subsubsection{向量值函数的可微性与 Jacobian 矩阵}

向量值函数的“导数”是什么？回顾前述两种情形中“可微”的统一表述。可微性本质上是用一个线性映射去逼近函数在某一点的增量：
\[
f(\boldsymbol a+\boldsymbol h)-f(\boldsymbol a)
=L(\boldsymbol h)+o(\|\boldsymbol h\|),\qquad\|\boldsymbol h\|\to 0,
\]
其中 $L$ 是一个\textbf{线性映射}。

\begin{itemize}
  \item 一元实值函数 $f:\mathbb R\to\mathbb R$：线性映射表达为 $L(h)=f'(a)\cdot h$
        （退化为“乘以标量 $f'(a)$”）；
  \item 多元实值函数 $f:\mathbb R^n\to\mathbb R$：线性映射表达为
        $L(\boldsymbol h)= \nabla f(\boldsymbol a)\cdot\boldsymbol h$
        （退化为“与梯度行向量的内积”）。
\end{itemize}

对于向量值函数 $f:\mathbb R^n\to\mathbb R^m$，线性映射 $L:\mathbb R^n\to\mathbb R^m$
在标准基下应当用一个 $m\times n$ 矩阵来表示——这个核心矩阵就是 \textbf{Jacobian 矩阵}。

\begin{definition}[向量值函数的可微性] \label{def:vec-diff}
设 $D\subset\mathbb R^n$ 为开区域$f:D\to\mathbb R^m$ 为向量值函数，
$\boldsymbol a\in D$，向量增量记作 $\boldsymbol h=(h_1,h_2,\dots,h_n)^{\mathsf T}$。
若存在一个只与 $\boldsymbol a$ 有关的 $m\times n$ 常数矩阵 $A=(a_{ij})_{m\times n}$ 使得
\[
f(\boldsymbol a+\boldsymbol h)-f(\boldsymbol a)=A\boldsymbol h+\boldsymbol r(\boldsymbol h),\qquad
\lim_{\|\boldsymbol h\|_n\to 0}\frac{\|\boldsymbol r(\boldsymbol h)\|_m}{\|\boldsymbol h\|_n}=0,
\]
展开就是\[\begin{pmatrix}
f_1(\boldsymbol a+\boldsymbol h)-f_1(\boldsymbol a)\\
f_2(\boldsymbol a+\boldsymbol h)-f_2(\boldsymbol a)\\
\vdots\\
f_m(\boldsymbol a+\boldsymbol h)-f_m(\boldsymbol a)
\end{pmatrix}=\begin{pmatrix}
a_{11}h_1+a_{12}h_2+\cdots+a_{1n}h_n+r_1(\boldsymbol h)\\
a_{21}h_1+a_{22}h_2+\cdots+a_{2n}h_n+r_2(\boldsymbol h)\\
\vdots\\
a_{m1}h_1+a_{m2}h_2+\cdots+a_{mn}h_n+r_m(\boldsymbol h)
\end{pmatrix}=A\boldsymbol h+\boldsymbol r(\boldsymbol h)\]将两端逐行比较可以得到，对每个 $i=1,2,\dots,m$，
\[
f_i(\boldsymbol a+\boldsymbol h)-f_i(\boldsymbol a)
=\sum_{j=1}^{n}a_{ij}h_j+r_i(\boldsymbol h)
=(a_{i1},a_{i2},\dots,a_{in})\boldsymbol h+r_i(\boldsymbol h). 
\]
同时，由于余项向量的极限为零，各分量必然也满足 \(\displaystyle 0 \leqslant  \frac{|r_i(\boldsymbol h)|}{\|\boldsymbol h\|_n}
\leqslant\frac{\|\boldsymbol r(\boldsymbol h)\|_m}{\|\boldsymbol h\|_n}\to 0\)。恰好说明了：每个分量函数 $f_i$ 都在 $\boldsymbol a$ 处可微，且其线性主部系数向量就是矩阵 $A$ 的第 $i$ 行 $(a_{i1},a_{i2},\dots,a_{in})$。这组系数正是 $f_i$ 在 $\boldsymbol a$ 处的梯度：
\[
(a_{i1},a_{i2},\dots,a_{in})
=\left(\frac{\partial f_i}{\partial x_1}(\boldsymbol a),\;
       \frac{\partial f_i}{\partial x_2}(\boldsymbol a),\;\dots,\;
       \frac{\partial f_i}{\partial x_n}(\boldsymbol a)\right).
\]
则称 $f$ 在 $\boldsymbol a$ 处\textbf{可微}（Differentiable），
并称 $\mathrm{d}f(\boldsymbol a)=A\boldsymbol h$ 为 $f$ 在 $\boldsymbol a$ 处的\textbf{微分}（Differential）。
\end{definition}

\begin{theorem}[向量值函数可微的充分条件] \label{thm:vec-diff-sufficient}
设 $D\subset\mathbb R^n$ 为开区域，$f:D\to\mathbb R^m$ 为向量值函数。
若 $f$ 的所有偏导数 $\dfrac{\partial f_i}{\partial x_j}$（$i=1,\dots,m$；$j=1,\dots,n$）
在 $\boldsymbol a$ 的某邻域内存在，且在 $\boldsymbol a$ 处\textbf{连续}，
则 $f$ 在 $\boldsymbol a$ 处可微。
\end{theorem}

\begin{theorem}[Jacobian 矩阵公式] \label{thm:jacobian-formula}
若向量值函数 $f:D\to\mathbb R^m$（$D\subset\mathbb R^n$ 为开区域）
在 $\boldsymbol a$ 处可微，则其导数矩阵为
\[
\boxed{
Jf(\boldsymbol a)= \dfrac{\partial(f_1,f_2,\dots,f_m)}{\partial(x_1,x_2,\dots,x_n)}\biggr|_{\boldsymbol a}=\begin{pmatrix}
\nabla f_1(\boldsymbol a) \\[4pt]
\nabla f_2(\boldsymbol a) \\[4pt]
\vdots \\[4pt]
\nabla f_m(\boldsymbol a)
\end{pmatrix}
=\begin{pmatrix}
\dfrac{\partial f_1}{\partial x_1}(\boldsymbol a) &
\dfrac{\partial f_1}{\partial x_2}(\boldsymbol a) &
\cdots &
\dfrac{\partial f_1}{\partial x_n}(\boldsymbol a) \\[10pt]
\dfrac{\partial f_2}{\partial x_1}(\boldsymbol a) &
\dfrac{\partial f_2}{\partial x_2}(\boldsymbol a) &
\cdots &
\dfrac{\partial f_2}{\partial x_n}(\boldsymbol a) \\[10pt]
\vdots & \vdots & \dots & \vdots \\[8pt]
\dfrac{\partial f_m}{\partial x_1}(\boldsymbol a) &
\dfrac{\partial f_m}{\partial x_2}(\boldsymbol a) &
\cdots &
\dfrac{\partial f_m}{\partial x_n}(\boldsymbol a)
\end{pmatrix}_{m\times n}
}
\]
\end{theorem}
特别地，Jacobian 矩阵统一了以下几种情况：
\begin{itemize}
  \item 当 $n=1$, $m=1$ 时：$Jf(a)=f'(a)$（退化为普通标量导数）；
  \item 当 $n\geqslant 2$, $m=1$ 时：$Jf(\boldsymbol a)= \nabla f(\boldsymbol a)$（退化为 $1 \times n$ 梯度行向量）；
  \item 当 $n=1$, $m\geqslant 2$ 时：\(\displaystyle Jf(a)=(f_1'(a),f_2'(a),\dots,f_m'(a))^{\mathsf T}\)（退化为 $m \times 1$ 列向量，表示空间曲线的“速度向量”）。
\end{itemize}

\paragraph*{矩阵大小的度量：Frobenius 范数}

在向量值函数的可微性中，Jacobian 矩阵可能是一般的 $m\times n$ 矩阵。为了估计线性主部和余项的大小，需要给矩阵也配上一个方便计算的范数。

\begin{definition}[矩阵的 Frobenius 范数] \label{def:matrix-norm}
设 $A=(a_{ij})_{m\times n}$ 为 $m\times n$ 实矩阵。定义
\[
  \boxed{\;
  \|A\|_{\mathrm F}:=
  \left(\sum_{i=1}^{m}\sum_{j=1}^{n}a_{ij}^{2}\right)^{\!1/2}
  \;}
\]
为矩阵 $A$ 的\textbf{Frobenius 范数}（Frobenius Norm），也称\textbf{Hilbert--Schmidt 范数}。若上下文不会混淆，也常简记为 $\|A\|$。
\end{definition}

\begin{note}
Frobenius 范数的本质是把 $m\times n$ 矩阵“展平”成一个 $mn$ 维向量后取欧几里得范数。
它是 $\mathbb R^{m\times n}$ 这个 $mn$ 维线性空间上很自然的度量。
虽然还有其他类型的矩阵范数（如算子范数），但 Frobenius 范数因其计算简便而在微分学中最为常用。
\end{note}

\begin{property}[Frobenius 范数的基本性质] \label{prop:norm-basic}
设 $A,B\in\mathbb R^{m\times n}$，$k\in\mathbb R$。则
\begin{enumerate}
  \item \textbf{(非负性)}: $\|A\|_{\mathrm F}\geqslant 0$，且 $\|A\|_{\mathrm F}=0\Longleftrightarrow A=O$（零矩阵）；
  \item \textbf{(正齐次性)}: $\|kA\|_{\mathrm F}=|k|\cdot\|A\|_{\mathrm F}$；
  \item \textbf{(三角不等式)}: $\|A+B\|_{\mathrm F}\leqslant\|A\|_{\mathrm F}+\|B\|_{\mathrm F}$。
\end{enumerate}
\end{property}

这三条性质的验证与向量范数完全类似，直接由定义可得。

\begin{property}[Frobenius 范数的次可乘性] \label{prop:norm-submultiplicative}
设 $A\in\mathbb R^{m\times n}$，$B\in\mathbb R^{n\times p}$，则
\[
  \boxed{\;\|AB\|_{\mathrm F}\leqslant\|A\|_{\mathrm F}\cdot\|B\|_{\mathrm F}\;} .
\]
\end{property}

\begin{proof}
\begin{enumerate}
  \item \textbf{写出矩阵乘积的元素：}
  设 $AB=C=(c_{ij})_{m\times p}$，其中
  \[
    c_{ij}=\sum_{k=1}^{n}a_{ik}b_{kj},
    \qquad i=1,\dots,m,\quad j=1,\dots,p.
  \]

  \item \textbf{应用柯西-施瓦茨不等式：}
  对每个 $c_{ij}^2$，利用向量内积的 Cauchy--Schwarz 不等式进行放缩：
  \[
    c_{ij}^2
    =\left(\sum_{k=1}^{n}a_{ik}b_{kj}\right)^{\!2}
    \leqslant
    \left(\sum_{k=1}^{n}a_{ik}^2\right)
    \left(\sum_{k=1}^{n}b_{kj}^2\right).
  \]

  \item \textbf{对所有元素求和：}
  将上述不等式对所有的 $i$ 和 $j$ 求和，得到 Frobenius 范数的平方：
  \begin{align*}
  \|AB\|_{\mathrm F}^2
  &=\sum_{i=1}^{m}\sum_{j=1}^{p}c_{ij}^2
  \leqslant\sum_{i=1}^{m}\sum_{j=1}^{p}
    \left(\sum_{k=1}^{n}a_{ik}^2\right)\!\left(\sum_{k=1}^{n}b_{kj}^2\right)\\
  &=\left(\sum_{i=1}^{m}\sum_{k=1}^{n}a_{ik}^2\right)\!
    \left(\sum_{k=1}^{n}\sum_{j=1}^{p}b_{kj}^2\right)
  =\|A\|_{\mathrm F}^2\cdot\|B\|_{\mathrm F}^2.
  \end{align*}

  \item \textbf{得出结论：}
  两边同时开平方即得 $\|AB\|_{\mathrm F}\leqslant\|A\|_{\mathrm F}\|B\|_{\mathrm F}$。
\end{enumerate}
\end{proof}

\begin{remark}[次可乘性的意义]
这一不等式在链式法则的证明中扮演关键角色：
当我们需要估计矩阵乘法或线性映射作用后的大小时，
次可乘性给出统一控制：
\[
  \|A\boldsymbol h\|
  \leqslant \|A\|_{\mathrm F}\,\|\boldsymbol h\|.
\]
这里可把列向量 $\boldsymbol h$ 看作一个 $n\times 1$ 矩阵。直观地说，线性映射作用后的增量大小，可以由矩阵本身的大小和输入增量的大小共同控制。
\end{remark}

\begin{exercise}[一元向量函数的导数——曲线速度] \label{ex:curve-deriv}
设 $D\subset\mathbb R$，$f:D\to\mathbb R^3$ 为一元三维向量值函数。
对 $a\in D$ 及充分小的 $h\in\mathbb R$，若 $f$ 在 $a$ 处可微，则定义给出 $f(a+h)-f(a)=Jf(a)\cdot h+\boldsymbol r(h)$。

将其写成分量形式：
\[
\begin{pmatrix}
f_1(a+h)\\f_2(a+h)\\f_3(a+h)
\end{pmatrix}
-\begin{pmatrix}
f_1(a)\\f_2(a)\\f_3(a)
\end{pmatrix}
=\begin{pmatrix}
f_1'(a)\\f_2'(a)\\f_3'(a)
\end{pmatrix}h
+\begin{pmatrix}
r_1(h)\\r_2(h)\\r_3(h)
\end{pmatrix}.
\]

由于自变量只有一个，对 $x$ 的偏导数就是普通导数。因此 Jacobian 矩阵为一个 $3\times 1$ 列向量，可写为 \(\displaystyle Jf(a)=\bigl(f_1'(a),\;f_2'(a),\;f_3'(a)\bigr)^{\mathsf T}\)。
\end{exercise}

让我们回到开篇提到的螺旋线的例子，进行具体计算验证：

\begin{exercise}[螺旋线的速度向量] \label{ex:helix-velocity}
对螺旋线 $f(t)=(\cos t,\sin t,t)^{\mathsf T}$，计算其导数并分析其几何特征。
\end{exercise}
\begin{solution}
\begin{enumerate}
  \item \textbf{计算导数（速度向量）：}
  使用标准基表示为 $f(t)=\vec i\cos t+\vec j\sin t+\vec k\,t$。对各个分量分别关于 $t$ 求导，得 \(\displaystyle f'(t)=(-\sin t,\;\cos t,\;1)^{\mathsf T}\)，也即 \(\displaystyle f'(t)=-\vec i\sin t+\vec j\cos t+\vec k\)。

  \item \textbf{几何特征分析：}
  \begin{itemize}
    \item 速度在 $xy$ 平面上的投影为水平分量 $(-\sin t,\cos t)^{\mathsf T}$。其长度恒为 $\sqrt{(-\sin t)^2+(\cos t)^2}=1$，且方向时刻与圆周相切，指向旋转的切线方向。
    \item 速度在 $z$ 轴上的投影为竖直分量 $1$。由于它是正的常数，表示质点在 $z$ 轴方向上做匀速上升运动。
    \item 质点运动的总速率为 \(\displaystyle \|f'(t)\|=\sqrt{(-\sin t)^2+(\cos t)^2+1^2}=\sqrt{1+1}=\sqrt{2}\)。
  \end{itemize}
\end{enumerate}
\end{solution}

\begin{exercise}[Jacobian 矩阵的完整计算] \label{ex:jacobian-compute}
求向量值函数
\[
f(x,y,z)=\begin{pmatrix}
3x+e^y z\\
x^3+y^2\sin z
\end{pmatrix}
\]
在点 $\left(\dfrac12,1,\pi\right)$ 处的 Jacobian 矩阵。
\end{exercise}
\begin{solution}
此函数是从 $\mathbb R^3$ 到 $\mathbb R^2$ 的映射，包含两个分量函数 $f_1(x,y,z)=3x+e^yz,\; f_2(x,y,z)=x^3+y^2\sin z$。

\begin{enumerate}
  \item \textbf{求各分量的偏导函数：}
  直接求得 \(\displaystyle \frac{\partial f_1}{\partial x}=3,\ \frac{\partial f_1}{\partial y}=e^y z,\ \frac{\partial f_1}{\partial z}=e^y\)，且 \(\displaystyle \frac{\partial f_2}{\partial x}=3x^2,\ \frac{\partial f_2}{\partial y}=2y\sin z,\ \frac{\partial f_2}{\partial z}=y^2\cos z\)。

  \item \textbf{在指定点 $\bigl(\frac12,1,\pi\bigr)$ 处求值：}
  代入 $x=\frac12, y=1, z=\pi$ 到上述偏导函数中，得第一行偏导值为 \(3,e\pi,e\)，第二行偏导值为 \(\displaystyle \frac34,0,-1\)。

  \item \textbf{组装 Jacobian 矩阵：}
  将 $f_1$ 的偏导数作为第一行，$f_2$ 的偏导数作为第二行，得到 $2 \times 3$ 的矩阵：
  \[
  \boxed{Jf\!\left(\frac12,1,\pi\right)
  =\begin{pmatrix}
  3 & e\pi & e\\[6pt]
  \dfrac34 & 0 & -1
  \end{pmatrix}_{2\times 3}}.
  \]

  \item \textbf{验证（线性近似的体现）：}
  在这个点附近产生微小增量 $\boldsymbol h = (h_1, h_2, h_3)^{\mathsf T}$ 时，函数的线性近似为
  \(\displaystyle f\!\left(\tfrac12+h_1,\,1+h_2,\,\pi+h_3\right)
  \approx f\!\left(\tfrac12,1,\pi\right)
  +Jf\!\left(\tfrac12,1,\pi\right)\boldsymbol h\)。
  具体写出就是：
  \[
  \begin{pmatrix}
  \frac32+e\pi\\
  \frac18+0
  \end{pmatrix}
  +\begin{pmatrix}
  3 & e\pi & e\\[6pt]
  \frac34 & 0 & -1
  \end{pmatrix}
  \begin{pmatrix}
  h_1\\h_2\\h_3
  \end{pmatrix}.
  \]
\end{enumerate}
\end{solution}




\subsection{链式法则、高阶偏导与二阶全微分}

\subsubsection{复合求导的链式法则}

\begin{theorem}[链式法则 / Chain Rule] \label{thm:chain-rule-matrix}
设
\[
z=g(y_1,\dots,y_m),\qquad
y_i=y_i(x_1,\dots,x_k),\quad i=1,\dots,m,
\]
也就是
\[
(x_1,\dots,x_k)^{\mathsf T}\longrightarrow (y_1,\dots,y_m)^{\mathsf T}\longrightarrow z.
\]
若内层映射
\[
f=(y_1,\dots,y_m)^{\mathsf T}:D\subset\mathbb R^k\to\mathbb R^m
\]
在点 $\boldsymbol x_0\in D$ 可微，外层函数 $g$ 在对应点 $f(\boldsymbol x_0)$ 可微，
则复合函数 $z=z(x_1,\dots,x_k)$ 在 $\boldsymbol x_0$ 处可微。

对任意一个自变量 $x_j$，从 $x_j$ 到 $z$ 有 $m$ 条路径：
\[
  x_j\longrightarrow y_1\longrightarrow z,\qquad
  x_j\longrightarrow y_2\longrightarrow z,\qquad
  \dots,\qquad
  x_j\longrightarrow y_m\longrightarrow z.
\]
因此
\[
\boxed{\;
\frac{\partial z}{\partial x_j}
=
\frac{\partial z}{\partial y_1}
\frac{\partial y_1}{\partial x_j}
+
\frac{\partial z}{\partial y_2}
\frac{\partial y_2}{\partial x_j}
+
\cdots
+
\frac{\partial z}{\partial y_m}
\frac{\partial y_m}{\partial x_j},
\qquad j=1,2,\dots,k.
\;}
\]
写开就是
\[
\begin{aligned}
\frac{\partial z}{\partial x_1}
&=
\frac{\partial z}{\partial y_1}\frac{\partial y_1}{\partial x_1}
+
\frac{\partial z}{\partial y_2}\frac{\partial y_2}{\partial x_1}
+
\cdots
+
\frac{\partial z}{\partial y_m}\frac{\partial y_m}{\partial x_1},\\[4pt]
\frac{\partial z}{\partial x_2}
&=
\frac{\partial z}{\partial y_1}\frac{\partial y_1}{\partial x_2}
+
\frac{\partial z}{\partial y_2}\frac{\partial y_2}{\partial x_2}
+
\cdots
+
\frac{\partial z}{\partial y_m}\frac{\partial y_m}{\partial x_2},\\[-2pt]
&\hspace{5em}\vdots\\[-2pt]
\frac{\partial z}{\partial x_k}
&=
\frac{\partial z}{\partial y_1}\frac{\partial y_1}{\partial x_k}
+
\frac{\partial z}{\partial y_2}\frac{\partial y_2}{\partial x_k}
+
\cdots
+
\frac{\partial z}{\partial y_m}\frac{\partial y_m}{\partial x_k}.
\end{aligned}
\]
把这些式子排成矩阵，就是
\[
\boxed{
\begin{pmatrix}
  \dfrac{\partial z}{\partial x_1} &
  \dfrac{\partial z}{\partial x_2} &
  \cdots &
  \dfrac{\partial z}{\partial x_k}
\end{pmatrix}
=
\begin{pmatrix}
  \dfrac{\partial z}{\partial y_1} &
  \dfrac{\partial z}{\partial y_2} &
  \cdots &
  \dfrac{\partial z}{\partial y_m}
\end{pmatrix}
\begin{pmatrix}
  \dfrac{\partial y_1}{\partial x_1} &
  \dfrac{\partial y_1}{\partial x_2} &
  \cdots &
  \dfrac{\partial y_1}{\partial x_k}\\[8pt]
  \dfrac{\partial y_2}{\partial x_1} &
  \dfrac{\partial y_2}{\partial x_2} &
  \cdots &
  \dfrac{\partial y_2}{\partial x_k}\\[-2pt]
  \vdots & \vdots & \ddots & \vdots\\[2pt]
  \dfrac{\partial y_m}{\partial x_1} &
  \dfrac{\partial y_m}{\partial x_2} &
  \cdots &
  \dfrac{\partial y_m}{\partial x_k}
\end{pmatrix}
}.
\]
\end{theorem}

\begin{note}[注意：可微性条件不可省略！]
使用链式法则必须保证中间映射 $f$ \textbf{可微}（而不仅仅是偏导数存在）。
例如函数
\[
f(x,y)=
\begin{cases}
\dfrac{x^2y}{x^2+y^2},&(x,y)\neq(0,0),\\[10pt]
0,&(x,y)=(0,0),
\end{cases}
\]
在 $(0,0)$ 处偏导数存在但不可微（甚至不连续），若对其直接套用链式法则将得到错误结论。
\end{note}

\begin{exercise}[一元复合：中间变量含 $t$] \label{ex:chain-1var}
设 $z=y\sin x$，$x=t^3$，$y=5t+2$，求 $\dfrac{\mathrm{d}z}{\mathrm{d}t}$。
\end{exercise}
\begin{solution}
把 $x,y$ 看作中间变量，直接套矩阵形式：
\[
\begin{pmatrix}
\dfrac{\mathrm{d}z}{\mathrm{d}t}
\end{pmatrix}
=
\begin{pmatrix}
\dfrac{\partial z}{\partial x} & \dfrac{\partial z}{\partial y}
\end{pmatrix}
\begin{pmatrix}
\dfrac{\mathrm{d}x}{\mathrm{d}t}\\[6pt]
\dfrac{\mathrm{d}y}{\mathrm{d}t}
\end{pmatrix}
=
\begin{pmatrix}
y\cos x & \sin x
\end{pmatrix}
\begin{pmatrix}
3t^2\\[4pt]
5
\end{pmatrix}.
\]
因此
\[
\frac{\mathrm{d}z}{\mathrm{d}t}
=3t^2y\cos x+5\sin x
=\boxed{3t^2(5t+2)\cos(t^3)+5\sin(t^3)}.
\]
\end{solution}

\begin{exercise}[二元复合：中间变量含 $u,v$] \label{ex:chain-2var}
设 $w=x^2e^y$，$x=4u$，$y=3u^2-2v$，求 $\dfrac{\partial w}{\partial u}$ 和 $\dfrac{\partial w}{\partial v}$。
\end{exercise}
\begin{solution}
把 $x,y$ 看作中间变量，直接套矩阵形式：
\[
\begin{pmatrix}
\dfrac{\partial w}{\partial u} & \dfrac{\partial w}{\partial v}
\end{pmatrix}
=
\begin{pmatrix}
\dfrac{\partial w}{\partial x} & \dfrac{\partial w}{\partial y}
\end{pmatrix}
\begin{pmatrix}
\dfrac{\partial x}{\partial u} & \dfrac{\partial x}{\partial v}\\[8pt]
\dfrac{\partial y}{\partial u} & \dfrac{\partial y}{\partial v}
\end{pmatrix}.
\]
各项偏导为
\[
\begin{pmatrix}
\dfrac{\partial w}{\partial x} & \dfrac{\partial w}{\partial y}
\end{pmatrix}
=
\begin{pmatrix}
2xe^y & x^2e^y
\end{pmatrix},
\qquad
\begin{pmatrix}
\dfrac{\partial x}{\partial u} & \dfrac{\partial x}{\partial v}\\[8pt]
\dfrac{\partial y}{\partial u} & \dfrac{\partial y}{\partial v}
\end{pmatrix}
=
\begin{pmatrix}
4 & 0\\
6u & -2
\end{pmatrix}.
\]
所以
\[
\begin{pmatrix}
w_u & w_v
\end{pmatrix}
=
\begin{pmatrix}
2xe^y & x^2e^y
\end{pmatrix}
\begin{pmatrix}
4 & 0\\
6u & -2
\end{pmatrix}
=
\begin{pmatrix}
8xe^y+6ux^2e^y & -2x^2e^y
\end{pmatrix}.
\]
代入 $x=4u,\ y=3u^2-2v$，得
\[
\boxed{
\frac{\partial w}{\partial u}=32u(1+3u^2)e^{3u^2-2v},
\qquad
\frac{\partial w}{\partial v}=-32u^2e^{3u^2-2v}
}.
\]
\end{solution}

\begin{exercise}[抽象复合函数的偏导数] \label{ex:chain-abstract}
设 $z=f\!\left(xy,\;\dfrac{x}{y}\right)$，其中 $f$ 具有连续偏导数。
求 $\dfrac{\partial z}{\partial x}$ 和 $\dfrac{\partial z}{\partial y}$。
\end{exercise}
\begin{solution}
令
\[
u=xy,\qquad v=\frac{x}{y},\qquad z=f(u,v).
\]
直接套矩阵形式：
\[
\begin{pmatrix}
\dfrac{\partial z}{\partial x} & \dfrac{\partial z}{\partial y}
\end{pmatrix}
=
\begin{pmatrix}
\dfrac{\partial z}{\partial u} & \dfrac{\partial z}{\partial v}
\end{pmatrix}
\begin{pmatrix}
\dfrac{\partial u}{\partial x} & \dfrac{\partial u}{\partial y}\\[8pt]
\dfrac{\partial v}{\partial x} & \dfrac{\partial v}{\partial y}
\end{pmatrix}.
\]
这里
\[
\begin{pmatrix}
\dfrac{\partial z}{\partial u} & \dfrac{\partial z}{\partial v}
\end{pmatrix}
=
\begin{pmatrix}
f_u & f_v
\end{pmatrix},
\qquad
\begin{pmatrix}
\dfrac{\partial u}{\partial x} & \dfrac{\partial u}{\partial y}\\[8pt]
\dfrac{\partial v}{\partial x} & \dfrac{\partial v}{\partial y}
\end{pmatrix}
=
\begin{pmatrix}
y & x\\[4pt]
\dfrac1y & -\dfrac{x}{y^2}
\end{pmatrix}.
\]
所以
\[
\begin{pmatrix}
z_x & z_y
\end{pmatrix}
=
\begin{pmatrix}
f_u & f_v
\end{pmatrix}
\begin{pmatrix}
y & x\\[4pt]
\dfrac1y & -\dfrac{x}{y^2}
\end{pmatrix}
=
\begin{pmatrix}
yf_u+\dfrac1y f_v &
xf_u-\dfrac{x}{y^2}f_v
\end{pmatrix}.
\]
故
\[
\boxed{
\frac{\partial z}{\partial x}=yf_u+\frac1y f_v,\qquad
\frac{\partial z}{\partial y}=xf_u-\frac{x}{y^2}f_v
},
\]
其中 $f_u,f_v$ 均在 $(u,v)=\left(xy,\dfrac{x}{y}\right)$ 处取值。
\end{solution}


\subsubsection{一阶全微分形式的不变性}

链式法则的一个重要推论是\textbf{全微分形式的不变性}——
无论 $x,y$ 是自变量还是中间变量，全微分 $\mathrm{d}z$ 的表达式在形式上完全相同。

\begin{theorem}[一阶全微分形式的不变性] \label{thm:diff-invariance}
设 $z=f(x,y)$，无论 $x,y$ 是自变量还是中间变量，都有 \(\displaystyle \boxed{\;\mathrm{d}z=\frac{\partial z}{\partial x}\,\mathrm{d}x+\frac{\partial z}{\partial y}\,\mathrm{d}y\;}\)。
\end{theorem}

\begin{exercise}[利用全微分形式不变性求偏导] \label{ex:diff-invariance-1}
设 $z=\sin(2x+y)$，利用全微分形式不变性求 $\dfrac{\partial z}{\partial x}$ 和 $\dfrac{\partial z}{\partial y}$。
\end{exercise}
\begin{solution}
\begin{enumerate}
  \item \textbf{计算全微分 $\mathrm{d}z$}：
  令中间变量 $u=2x+y$，则 $z=\sin u$。
  利用一元函数的微分公式和全微分的线性性质：\(\displaystyle \mathrm{d}z=\mathrm{d}(\sin u)=\cos u\,\mathrm{d}u=\cos(2x+y)(2\,\mathrm{d}x+\mathrm{d}y)=2\cos(2x+y)\,\mathrm{d}x+\cos(2x+y)\,\mathrm{d}y\)。

  \item \textbf{比较系数提取偏导数}：
  由全微分定义比较系数可得
  \(\displaystyle \frac{\partial z}{\partial x}=2\cos(2x+y),\qquad
  \frac{\partial z}{\partial y}=\cos(2x+y)\)。
\end{enumerate}
\end{solution}

\begin{exercise}[三元函数的全微分求偏导] \label{ex:diff-invariance-2}
设 $u=\ln\sqrt{x^2+y^2+z^2}$，求 $u_x$, $u_y$, $u_z$。
\end{exercise}
\begin{solution}
  利用对数性质将根号提出：$u=\dfrac12\ln(x^2+y^2+z^2)$。将括号内整体看作一个变量，应用 $\mathrm{d}(\ln v) = \dfrac{\mathrm{d}v}{v}$，得
\begin{align*}
\mathrm{d}u=\frac{x\,\mathrm{d}x+y\,\mathrm{d}y+z\,\mathrm{d}z}{x^2+y^2+z^2},~~
  u_x=\frac{x}{x^2+y^2+z^2},~~
  u_y=\frac{y}{x^2+y^2+z^2},~~
  u_z=\frac{z}{x^2+y^2+z^2}.
\end{align*}
  上式已经写成 $\mathrm{d}u = u_x\,\mathrm{d}x + u_y\,\mathrm{d}y + u_z\,\mathrm{d}z$ 的形式，直接读系数即可。
\end{solution}

\subsubsection{高阶偏导数与 Hesse 矩阵}

\begin{definition}[二阶偏导数与 Hesse 矩阵] \label{def:hessian}
设 $f:D\to\mathbb R$（$D\subset\mathbb R^2$ 为开区域）具有一阶偏导数 $f_x$, $f_y$。
若 $f_x$, $f_y$ 本身也具有偏导数，则称其为 $f$ 的\textbf{二阶偏导数}：
\begin{align*}
\frac{\partial^2 f}{\partial x^2}&=f_{xx}
  :=\frac{\partial}{\partial x}\!\left(\frac{\partial f}{\partial x}\right),\qquad
\frac{\partial^2 f}{\partial y\,\partial x}=f_{xy}
  :=\frac{\partial}{\partial y}\!\left(\frac{\partial f}{\partial x}\right),\\[8pt]
\frac{\partial^2 f}{\partial x\,\partial y}&=f_{yx}
  :=\frac{\partial}{\partial x}\!\left(\frac{\partial f}{\partial y}\right),\qquad
\frac{\partial^2 f}{\partial y^2}=f_{yy}
  :=\frac{\partial}{\partial y}\!\left(\frac{\partial f}{\partial y}\right).
\end{align*}

其中 $f_{xy}$ 和 $f_{yx}$ 称为\textbf{混合偏导数}（Mixed Partial Derivatives），
注意求导顺序：$f_{xy}$ 表示“先对 $x$ 求偏导，再对 $y$ 求偏导”。

矩阵
\[
H_f(\boldsymbol x_0)
=\begin{pmatrix}
f_{xx}(\boldsymbol x_0) & f_{xy}(\boldsymbol x_0)\\[6pt]
f_{yx}(\boldsymbol x_0) & f_{yy}(\boldsymbol x_0)
\end{pmatrix}
\]
称为 $f$ 在 $\boldsymbol x_0$ 处的\textbf{Hesse 矩阵}（Hessian Matrix）。
\end{definition}

\begin{theorem}[复合函数的二阶偏导数]
设
\[
z=f(u_1,\dots,u_m),\qquad
u_i=u_i(x_1,\dots,x_k),\quad i=1,\dots,m,
\]
且 $f,u_1,\dots,u_m$ 均具有连续二阶偏导数。下面所有关于 $f$ 的偏导数都在点
\((u_1(x),\dots,u_m(x))\) 处取值。

记
\[
J=\left(\frac{\partial u_i}{\partial x_j}\right)_{m\times k},\qquad
H_{\boldsymbol u}(f)=\left(f_{u_i u_j}\right)_{m\times m},\qquad
H_{\boldsymbol x}(u_i)=\left((u_i)_{x_p x_q}\right)_{k\times k}.
\]
则复合函数 $z=f(u_1(x),\dots,u_m(x))$ 关于 $\boldsymbol x=(x_1,\dots,x_k)^{\mathsf T}$ 的 Hesse 矩阵为
\[
\boxed{\;
H_{\boldsymbol x}(z)
=J^{\mathsf T}H_{\boldsymbol u}(f)J
+f_{u_1}H_{\boldsymbol x}(u_1)+\cdots+f_{u_m}H_{\boldsymbol x}(u_m).
\;}
\]

把这个矩阵公式写开，就是
\[
\resizebox{\linewidth}{!}{%
$\displaystyle
\boxed{
\begin{aligned}
&
\begin{pmatrix}
z_{x_1x_1} & z_{x_1x_2} & \cdots & z_{x_1x_k}\\[4pt]
z_{x_2x_1} & z_{x_2x_2} & \cdots & z_{x_2x_k}\\
\vdots & \vdots & \ddots & \vdots\\
z_{x_kx_1} & z_{x_kx_2} & \cdots & z_{x_kx_k}
\end{pmatrix} \\[6pt]
&=
\begin{pmatrix}
(u_1)_{x_1} & (u_2)_{x_1} & \cdots & (u_m)_{x_1}\\[4pt]
(u_1)_{x_2} & (u_2)_{x_2} & \cdots & (u_m)_{x_2}\\
\vdots & \vdots & \ddots & \vdots\\
(u_1)_{x_k} & (u_2)_{x_k} & \cdots & (u_m)_{x_k}
\end{pmatrix}
\begin{pmatrix}
f_{u_1u_1} & f_{u_1u_2} & \cdots & f_{u_1u_m}\\[4pt]
f_{u_2u_1} & f_{u_2u_2} & \cdots & f_{u_2u_m}\\
\vdots & \vdots & \ddots & \vdots\\
f_{u_mu_1} & f_{u_mu_2} & \cdots & f_{u_mu_m}
\end{pmatrix}
\begin{pmatrix}
(u_1)_{x_1} & (u_1)_{x_2} & \cdots & (u_1)_{x_k}\\[4pt]
(u_2)_{x_1} & (u_2)_{x_2} & \cdots & (u_2)_{x_k}\\
\vdots & \vdots & \ddots & \vdots\\
(u_m)_{x_1} & (u_m)_{x_2} & \cdots & (u_m)_{x_k}
\end{pmatrix} \\[6pt]
&\quad
+f_{u_1}
\begin{pmatrix}
(u_1)_{x_1x_1} & (u_1)_{x_1x_2} & \cdots & (u_1)_{x_1x_k}\\
(u_1)_{x_2x_1} & (u_1)_{x_2x_2} & \cdots & (u_1)_{x_2x_k}\\
\vdots & \vdots & \ddots & \vdots\\
(u_1)_{x_kx_1} & (u_1)_{x_kx_2} & \cdots & (u_1)_{x_kx_k}
\end{pmatrix}
+\cdots
+f_{u_m}
\begin{pmatrix}
(u_m)_{x_1x_1} & (u_m)_{x_1x_2} & \cdots & (u_m)_{x_1x_k}\\
(u_m)_{x_2x_1} & (u_m)_{x_2x_2} & \cdots & (u_m)_{x_2x_k}\\
\vdots & \vdots & \ddots & \vdots\\
(u_m)_{x_kx_1} & (u_m)_{x_kx_2} & \cdots & (u_m)_{x_kx_k}
\end{pmatrix}.
\end{aligned}
}
$%
}
\]

特别地，当 $z=f(u,v)$，$u=u(x,y)$，$v=v(x,y)$ 时，上式化为
\[
H_{(x,y)}(z)
=
\begin{pmatrix}
u_x & v_x\\
u_y & v_y
\end{pmatrix}
\begin{pmatrix}
f_{uu} & f_{uv}\\
f_{vu} & f_{vv}
\end{pmatrix}
\begin{pmatrix}
u_x & u_y\\
v_x & v_y
\end{pmatrix}
+f_u
\begin{pmatrix}
u_{xx} & u_{xy}\\
u_{yx} & u_{yy}
\end{pmatrix}
+f_v
\begin{pmatrix}
v_{xx} & v_{xy}\\
v_{yx} & v_{yy}
\end{pmatrix}.
\]
若 $f$ 的二阶混合偏导连续，则 $f_{uv}=f_{vu}$，于是可读出
\begin{align*}
z_{xx}&=f_{uu}u_x^2+2f_{uv}u_xv_x+f_{vv}v_x^2+f_u u_{xx}+f_v v_{xx},\\
z_{xy}&=f_{uu}u_xu_y+f_{uv}(u_xv_y+u_yv_x)+f_{vv}v_xv_y+f_u u_{xy}+f_v v_{xy},\\
z_{yy}&=f_{uu}u_y^2+2f_{uv}u_yv_y+f_{vv}v_y^2+f_u u_{yy}+f_v v_{yy}.
\end{align*}
\end{theorem}

\begin{proof}
由一阶链式法则，对任意 $p=1,\dots,k$，
\[
z_{x_p}=\sum_{i=1}^{m}f_{u_i}(u_i)_{x_p}.
\]
再对 $x_q$ 求偏导。由于 $f_{u_i}$ 仍是 $u_1,\dots,u_m$ 的复合函数，有
\[
\frac{\partial f_{u_i}}{\partial x_q}
=\sum_{j=1}^{m}f_{u_i u_j}(u_j)_{x_q}.
\]
因此
\begin{align*}
z_{x_p x_q}=\sum_{i=1}^{m}
\left[
\left(\frac{\partial f_{u_i}}{\partial x_q}\right)(u_i)_{x_p}
+f_{u_i}(u_i)_{x_p x_q}
\right]=\sum_{i=1}^{m}\sum_{j=1}^{m}
f_{u_i u_j}(u_i)_{x_p}(u_j)_{x_q}
+\sum_{i=1}^{m}f_{u_i}(u_i)_{x_p x_q}.
\end{align*}
第一项正是矩阵 $J^{\mathsf T}H_{\boldsymbol u}(f)J$ 的第 $(p,q)$ 个元素，
第二项正是矩阵 $\sum\limits_{i=1}^{m}f_{u_i}H_{\boldsymbol x}(u_i)$ 的第 $(p,q)$ 个元素。
由于 $p,q$ 任意，矩阵公式成立。
\end{proof}


\begin{theorem}[混合偏导数交换定理 / Clairaut's Theorem] \label{thm:mixed-pd-exchange}
若 $f_{xy}$ 和 $f_{yx}$ 在点 $(x_0,y_0)$ 的某邻域内存在且在 $(x_0,y_0)$ 处\textbf{连续}，
则 \(\displaystyle \boxed{\;f_{xy}(x_0,y_0)=f_{yx}(x_0,y_0)\;}\)。
\end{theorem}

\begin{note}[证明思路]
证明的核心工具是\textbf{二次差商}：
考虑 $f(x_0+h,y_0+k)-f(x_0+h,y_0)-f(x_0,y_0+k)+f(x_0,y_0)$。
对这个量分别用“先 $x$ 后 $y$”和“先 $y$ 后 $x$”两种方式应用拉格朗日中值定理，
再令 $h,k\to 0$ 取极限并利用二阶偏导数的连续性即可得证。
\end{note}

交换定理要求 $f_{xy}$ 和 $f_{yx}$ \textbf{连续}。
若连续性不满足，混合偏导数可能不相等。

\begin{exercise}[混合偏导数不相等] \label{ex:mixed-pd-unequal}
设
\[
f(x,y)=
\begin{cases}
xy\,\dfrac{x^2-y^2}{x^2+y^2},&(x,y)\neq(0,0),\\[10pt]
0,&(x,y)=(0,0).
\end{cases}
\]
验证 $f_{xy}(0,0)\neq f_{yx}(0,0)$。
\end{exercise}
\begin{solution}
\begin{enumerate}
  \item \textbf{在坐标轴上求一阶偏导数 $f_x(0,y)$}：
  针对非原点的纵轴上的点（此时 $y\neq 0$）：
  应用导数的定义直接计算极限：
  \[
  f_x(0,y)=\lim_{h\to 0}\frac{f(h,y)-f(0,y)}{h}
  =\lim_{h\to 0}\frac{hy\,\dfrac{h^2-y^2}{h^2+y^2}-0}{h}
  =\lim_{h\to 0} y\cdot\frac{h^2-y^2}{h^2+y^2} = y\cdot\frac{-y^2}{y^2}=-y.
  \]

  \item \textbf{在坐标轴上求一阶偏导数 $f_y(x,0)$}：
  针对非原点的横轴上的点（此时 $x\neq 0$）：
  同理应用导数定义：
  \[
  f_y(x,0)=\lim_{k\to 0}\frac{f(x,k)-f(x,0)}{k}
  =\lim_{k\to 0}\frac{xk\,\dfrac{x^2-k^2}{x^2+k^2}-0}{k}
  =\lim_{k\to 0} x\cdot\frac{x^2-k^2}{x^2+k^2} = x\cdot\frac{x^2}{x^2}=x.
  \]

  \item \textbf{计算原点处的混合偏导数 $f_{xy}(0,0)$}：
  先对 $x$ 求偏导，再对 $y$ 差商：
  \[
  f_{xy}(0,0)=\lim_{k\to 0}\frac{f_x(0,k)-f_x(0,0)}{k}
  =\lim_{k\to 0}\frac{-k-0}{k}=-1.
  \]

  \item \textbf{计算原点处的混合偏导数 $f_{yx}(0,0)$}：
  先对 $y$ 求偏导，再对 $x$ 差商：
  \[
  f_{yx}(0,0)=\lim_{h\to 0}\frac{f_y(h,0)-f_y(0,0)}{h}
  =\lim_{h\to 0}\frac{h-0}{h}=1.
  \]

  \item \textbf{比对结论}：
  发现结果不相等：$\boxed{f_{xy}(0,0)=-1\neq 1=f_{yx}(0,0)}$。
  这说明 $f_{xy}$ 和 $f_{yx}$ 在原点处必有一方或双方不连续。
\end{enumerate}
\end{solution}

\begin{remark}
这个经典反例揭示了：
“偏导数存在”和“偏导数连续”之间存在本质差距。
在实际应用中，当我们需要交换混合偏导数的顺序时，
\textbf{务必先验证连续性条件}，否则可能得出错误的结论。
\end{remark}

\begin{exercise}[二阶偏导数的完整计算] \label{ex:hessian-basic}
设 $f(x,y)=y\cos x+3x^2e^y$，求所有二阶偏导数。
\end{exercise}
\begin{solution}
把它看成复合的平凡情形：
\[
z=F(u,v),\qquad u=x,\qquad v=y,\qquad
F(u,v)=v\cos u+3u^2e^v.
\]
此时公式中需要的矩阵为
\[
\begin{pmatrix}
u_x&v_x\\
u_y&v_y
\end{pmatrix}
=
\begin{pmatrix}
u_x&u_y\\
v_x&v_y
\end{pmatrix}
=
\begin{pmatrix}
1&0\\
0&1
\end{pmatrix},\qquad
H_{(x,y)}(u)=0,\qquad H_{(x,y)}(v)=0.
\]
外层函数的 Hesse 矩阵为
\[
H_{(u,v)}(F)=
\begin{pmatrix}
F_{uu}&F_{uv}\\
F_{vu}&F_{vv}
\end{pmatrix}
=
\begin{pmatrix}
-v\cos u+6e^v&-\sin u+6ue^v\\
-\sin u+6ue^v&3u^2e^v
\end{pmatrix}.
\]
代入
\[
H_{(x,y)}(z)
=
\begin{pmatrix}
1&0\\
0&1
\end{pmatrix}
H_{(u,v)}(F)
\begin{pmatrix}
1&0\\
0&1
\end{pmatrix}
+F_uH_{(x,y)}(u)+F_vH_{(x,y)}(v),
\]
得
\[
H_{(x,y)}(f)=
\begin{pmatrix}
-y\cos x+6e^y&-\sin x+6xe^y\\
-\sin x+6xe^y&3x^2e^y
\end{pmatrix}.
\]
因此
\[
\boxed{
f_{xx}=-y\cos x+6e^y,\quad
f_{xy}=f_{yx}=-\sin x+6xe^y,\quad
f_{yy}=3x^2e^y
}.
\]
\end{solution}

\begin{exercise}[arctan 的二阶偏导数] \label{ex:hessian-arctan}
设 $f(x,y)=\arctan\dfrac{y}{x}$，求所有二阶偏导数。
\end{exercise}
\begin{solution}
设
\[
z=F(u,v),\qquad u=x,\qquad v=y,\qquad F(u,v)=\arctan\frac{v}{u}.
\]
同样有
\[
\begin{pmatrix}
u_x&v_x\\
u_y&v_y
\end{pmatrix}
=
\begin{pmatrix}
u_x&u_y\\
v_x&v_y
\end{pmatrix}
=
\begin{pmatrix}
1&0\\
0&1
\end{pmatrix},\qquad
H_{(x,y)}(u)=0,\qquad H_{(x,y)}(v)=0.
\]
先列外层的 Hesse 矩阵：
\[
H_{(u,v)}(F)=
\begin{pmatrix}
\dfrac{2uv}{(u^2+v^2)^2}&\dfrac{v^2-u^2}{(u^2+v^2)^2}\\[10pt]
\dfrac{v^2-u^2}{(u^2+v^2)^2}&-\dfrac{2uv}{(u^2+v^2)^2}
\end{pmatrix}.
\]
由
\[
H_{(x,y)}(z)
=
\begin{pmatrix}
1&0\\
0&1
\end{pmatrix}
H_{(u,v)}(F)
\begin{pmatrix}
1&0\\
0&1
\end{pmatrix}
+F_uH_{(x,y)}(u)+F_vH_{(x,y)}(v)
\]
可得
\[
H_{(x,y)}(f)=
\begin{pmatrix}
\dfrac{2xy}{(x^2+y^2)^2}&\dfrac{y^2-x^2}{(x^2+y^2)^2}\\[10pt]
\dfrac{y^2-x^2}{(x^2+y^2)^2}&-\dfrac{2xy}{(x^2+y^2)^2}
\end{pmatrix}.
\]
因此
\[
\boxed{
f_{xx}=\frac{2xy}{(x^2+y^2)^2},\quad
f_{xy}=f_{yx}=\frac{y^2-x^2}{(x^2+y^2)^2},\quad
f_{yy}=-\frac{2xy}{(x^2+y^2)^2}
}.
\]
\end{solution}

\begin{exercise}[抽象函数的二阶偏导数] \label{ex:chain-2nd-order}
设 $z=f\!\left(xy,\;\dfrac{x}{y}\right)$，$f$ 具有二阶连续偏导数。
求 $\dfrac{\partial^2 z}{\partial x^2}$，$\dfrac{\partial^2 z}{\partial x\partial y}$，$\dfrac{\partial^2 z}{\partial y^2}$。
\end{exercise}
\begin{solution}
令
\[
u=xy,\qquad v=\frac{x}{y},\qquad z=f(u,v).
\]
先列矩阵：
\[
\begin{pmatrix}
u_x&v_x\\
u_y&v_y
\end{pmatrix}
=
\begin{pmatrix}
y&\dfrac1y\\[4pt]
x&-\dfrac{x}{y^2}
\end{pmatrix},\qquad
\begin{pmatrix}
u_x&u_y\\
v_x&v_y
\end{pmatrix}
=
\begin{pmatrix}
y&x\\[4pt]
\dfrac1y&-\dfrac{x}{y^2}
\end{pmatrix},
\]
\[
H_{(x,y)}(u)=
\begin{pmatrix}
0&1\\
1&0
\end{pmatrix},\qquad
H_{(x,y)}(v)=
\begin{pmatrix}
0&-\dfrac1{y^2}\\[8pt]
-\dfrac1{y^2}&\dfrac{2x}{y^3}
\end{pmatrix}.
\]
由公式
\[
\begin{aligned}
H_{(x,y)}(z)
={}
\begin{pmatrix}
y&\dfrac1y\\[4pt]
x&-\dfrac{x}{y^2}
\end{pmatrix}
\begin{pmatrix}
f_{uu}&f_{uv}\\
f_{vu}&f_{vv}
\end{pmatrix}
\begin{pmatrix}
y&x\\[4pt]
\dfrac1y&-\dfrac{x}{y^2}
\end{pmatrix}
+f_u
\begin{pmatrix}
0&1\\
1&0
\end{pmatrix}
+f_v
\begin{pmatrix}
0&-\dfrac1{y^2}\\[8pt]
-\dfrac1{y^2}&\dfrac{2x}{y^3}
\end{pmatrix}.
\end{aligned}
\]
矩阵相乘后得到
\[
H_{(x,y)}(z)=
\begin{pmatrix}
y^2f_{uu}+2f_{uv}+\dfrac1{y^2}f_{vv}
&
xyf_{uu}-\dfrac{x}{y^3}f_{vv}+f_u-\dfrac1{y^2}f_v\\[10pt]
xyf_{uu}-\dfrac{x}{y^3}f_{vv}+f_u-\dfrac1{y^2}f_v
&
x^2f_{uu}-\dfrac{2x^2}{y^2}f_{uv}+\dfrac{x^2}{y^4}f_{vv}+\dfrac{2x}{y^3}f_v
\end{pmatrix}.
\]
因此
\[
\boxed{
\begin{aligned}
z_{xx}&=y^2f_{uu}+2f_{uv}+\frac1{y^2}f_{vv},\\
z_{xy}&=xyf_{uu}-\frac{x}{y^3}f_{vv}+f_u-\frac1{y^2}f_v,\\
z_{yy}&=x^2f_{uu}-\frac{2x^2}{y^2}f_{uv}+\frac{x^2}{y^4}f_{vv}+\frac{2x}{y^3}f_v.
\end{aligned}}
\]
其中所有关于 $f$ 的偏导数均在 \(\left(xy,\dfrac{x}{y}\right)\) 处取值。
\end{solution}

\begin{exercise}[另一个抽象复合函数] \label{ex:chain-2nd-order-2}
设 $z=f(2x-y,\;y\sin x)$，$f$ 具有二阶连续偏导数。
求 $\dfrac{\partial z}{\partial x}$ 和 $\dfrac{\partial^2 z}{\partial x\partial y}$。
\end{exercise}
\begin{solution}
令
\[
u=2x-y,\qquad v=y\sin x,\qquad z=f(u,v).
\]
一阶导数先用矩阵公式：
\[
\begin{pmatrix}
z_x&z_y
\end{pmatrix}
=
\begin{pmatrix}
f_u&f_v
\end{pmatrix}
\begin{pmatrix}
u_x&u_y\\
v_x&v_y
\end{pmatrix}
=
\begin{pmatrix}
f_u&f_v
\end{pmatrix}
\begin{pmatrix}
2&-1\\
y\cos x&\sin x
\end{pmatrix}.
\]
因此
\[
\boxed{z_x=2f_u+y\cos x\,f_v}.
\]
再列二阶矩阵：
\[
\begin{pmatrix}
u_x&v_x\\
u_y&v_y
\end{pmatrix}
=
\begin{pmatrix}
2&y\cos x\\
-1&\sin x
\end{pmatrix},\qquad
\begin{pmatrix}
u_x&u_y\\
v_x&v_y
\end{pmatrix}
=
\begin{pmatrix}
2&-1\\
y\cos x&\sin x
\end{pmatrix},
\]
\[
H_{(x,y)}(u)=
\begin{pmatrix}
0&0\\
0&0
\end{pmatrix},
\]
\[
H_{(x,y)}(v)=
\begin{pmatrix}
-y\sin x&\cos x\\
\cos x&0
\end{pmatrix}.
\]
代入
\[
\begin{aligned}
H_{(x,y)}(z)
={}&
\begin{pmatrix}
2&y\cos x\\
-1&\sin x
\end{pmatrix}
\begin{pmatrix}
f_{uu}&f_{uv}\\
f_{vu}&f_{vv}
\end{pmatrix}
\begin{pmatrix}
2&-1\\
y\cos x&\sin x
\end{pmatrix}\\
&+f_u
\begin{pmatrix}
0&0\\
0&0
\end{pmatrix}
+f_v
\begin{pmatrix}
-y\sin x&\cos x\\
\cos x&0
\end{pmatrix}.
\end{aligned}
\]
取其中 $(1,2)$ 位置，得
\[
\boxed{
z_{xy}
=-2f_{uu}+(2\sin x-y\cos x)f_{uv}
+y\sin x\cos x\,f_{vv}
+\cos x\,f_v
}.
\]
其中所有关于 $f$ 的偏导数均在 \((2x-y,\;y\sin x)\) 处取值。
\end{solution}

\subsubsection{二阶全微分}

\begin{definition}[二阶全微分] \label{def:second-total-diff}
设 $z=f(x,y)$ 具有二阶连续偏导数。
对一阶全微分 $\mathrm{d}z=f_x\,\mathrm{d}x+f_y\,\mathrm{d}y$ 再求一次微分（视自变量微元 $\mathrm{d}x,\mathrm{d}y$ 为独立常数），
得到\textbf{二阶全微分}：\(\displaystyle \mathrm{d}^2z=\mathrm{d}(\mathrm{d}z)=f_{xx}\,\mathrm{d}x^2+2f_{xy}\,\mathrm{d}x\,\mathrm{d}y+f_{yy}\,\mathrm{d}y^2\)。
其中 $\mathrm{d}x^2=(\mathrm{d}x)^2$，$\mathrm{d}y^2=(\mathrm{d}y)^2$。
\end{definition}

\begin{exercise}[二阶全微分的计算] \label{ex:2nd-diff}
设 $z=x\sin y$，求 $\mathrm{d}z$ 和 $\mathrm{d}^2z$。
\end{exercise}
\begin{solution}
\begin{enumerate}
  \item \textbf{计算一阶偏导数及一阶全微分}：
  $\dfrac{\partial z}{\partial x}=\sin y$，$\;\dfrac{\partial z}{\partial y}=x\cos y$。
  由此得 $\mathrm{d}z=\sin y\,\mathrm{d}x+x\cos y\,\mathrm{d}y$。

  \item \textbf{计算二阶偏导数}：
  继续求二阶导数：
  $f_{xx}=0$，$\;f_{xy}=\cos y$，$\;f_{yy}=-x\sin y$。

  \item \textbf{代入二阶全微分公式}：
  将求得的二阶导数代入定义式即得 \(\displaystyle \mathrm{d}^2z=\boxed{2\cos y\,\mathrm{d}x\,\mathrm{d}y-x\sin y\,\mathrm{d}y^2}\)。
\end{enumerate}
\end{solution}


\subsection{切线、法平面、切平面与法线}

\begin{itemize}
  \item 给定一条空间曲线，如何求出它在某点处的 \textbf{切线} 和 \textbf{法平面}？
  \item 给定一个曲面，如何求出它在某点处的 \textbf{切平面} 和 \textbf{法线}？
\end{itemize}

\subsubsection{空间曲线的切线与法平面}

\begin{definition}[空间参数曲线] \label{def:param-curve}
称映射 $\vec{r}:[\alpha,\beta]\to\mathbb R^3$：
\(\displaystyle \vec{r}(t)=\bigl(\phi(t),\;\psi(t),\;\omega(t)\bigr),\ \alpha\leqslant t\leqslant\beta\)，为一条 \textbf{参数曲线}（Parametric Curve），
其中 $\phi$, $\psi$, $\omega$ 称为 \textbf{坐标函数}（Component Functions）。
若三者均在 $(\alpha,\beta)$ 内连续可微且 $\vec{r}'(t) \neq \boldsymbol{0}$ 对所有 $t$ 成立，则称 $L$ 为 \textbf{光滑曲线}（Smooth Curve）。
\end{definition}

\begin{property}[切线方程（对称式 / 点向式）] \label{property:tangent-line-param}
设光滑参数曲线 $\vec{r}(t)=(\phi(t),\psi(t),\omega(t))$ 在
$p_0=(x_0,y_0,z_0)=\vec{r}(t_0)$ 处有切向量
$\vec{r}'(t_0)=(\phi'(t_0),\psi'(t_0),\omega'(t_0)) \neq \boldsymbol{0}$。
则曲线在 $p_0$ 处的\textbf{切线}（Tangent Line）方程为 \[\displaystyle \boxed{\frac{x-x_0}{\phi'(t_0)}=\frac{y-y_0}{\psi'(t_0)}=\frac{z-z_0}{\omega'(t_0)}}\]
\textbf{约定}：若某个分母为零（例如 $\phi'(t_0)=0$），则对应的分子也为零
（即该坐标保持不变），理解为相应的比例式退化为约束条件 $x-x_0=0$。等价的参数形式写法为
$\boldsymbol{p}-p_0=k\cdot\vec{r}'(t_0)$，$k\in(-\infty,+\infty)$。
\end{property}

\begin{property}[法平面方程] \label{property:normal-plane}
它是以切向量 $\vec{r}'(t_0)$ 为\textbf{法向量}、过点 $p_0$ 的平面。其方程为 \[\displaystyle \boxed{\phi'(t_0)(x-x_0)+\psi'(t_0)(y-y_0)+\omega'(t_0)(z-z_0)=0}\]
\end{property}

上述公式要求 $\vec{r}'(t_0) \neq \boldsymbol{0}$。
那么当 $\vec{r}'(t_0)=\boldsymbol{0}$ 时会发生什么？

\begin{remark}[奇点]
若 $\vec{r}'(t_0)=(\phi'(t_0),\psi'(t_0),\omega'(t_0))=\boldsymbol{0}$，
则我们 \textbf{无法直接判断} 曲线在该点的切线是否存在。
使得 $\vec{r}'(t)=\boldsymbol{0}$ 的点 $t$ 称为曲线的 \textbf{奇点}（Singular Point）。

注意：$\vec{r}'(t_0)=\boldsymbol{0}$ 只说明“此参数化在该点的瞬时速度为零”，
并不等同于“曲线本身在该点没有切线”。
判断奇点处是否有切线需要更精细的分析（例如考察高阶导数或换一种参数化方式）。
这就是为什么“光滑曲线”的定义要求 $\vec{r}'(t) \neq \boldsymbol{0}$
对 \textbf{所有} $t$ 成立——这保证了参数化的正则性与曲线本身的正则性相一致。
\end{remark}

\begin{exercise}[奇点处切线行为] \label{ex:singular-points}
考察以下两条曲线在原点处的行为：

\textbf{曲线 $L_1$}：
\[
\vec{r}_1(t)=
\begin{cases}
(t^3,\;t^3,\;t^3),      & t\geqslant 0,\\
(t^3,\;2t^3,\;3t^3),    & t<0.
\end{cases}
\]

\textbf{曲线 $L_2$}：$\vec{r}_2(t)=(t^3,t^3,t^3)$，$t\in(-\infty,+\infty)$。
\end{exercise}
\begin{solution}
\begin{enumerate}
  \item \textbf{计算两曲线在原点处的一阶导数：}
  对于 $L_1$：当 $t>0$ 时 $\vec{r}_1'(t)=(3t^2,3t^2,3t^2)\to(0,0,0)$；当 $t<0$ 时 $\vec{r}_1'(t)=(3t^2,6t^2,9t^2)\to(0,0,0)$。故 $\vec{r}_1'(0)=\boldsymbol{0}$。
  对于 $L_2$：$\vec{r}_2'(t)=(3t^2,3t^2,3t^2)$，同样 $\vec{r}_2'(0)=\boldsymbol{0}$。
  两者在原点处都是奇点。

  \item \textbf{分析 $L_2$ 的几何特征：}
  $L_2$ 的参数方程本质上刻画的是直线 $x=y=z$（只是以 $t^3$ 作为非均匀的参数化变量）。
  尽管参数化的“速度”在原点处停顿了（一阶导数为零），但曲线本身是一条笔直的直线 —— 因此在原点处 \textbf{切线确实存在}（就是这条直线自身，方向向量为 $(1,1,1)$ 或其任意非零倍数）。

  \item \textbf{分析 $L_1$ 的几何特征：}
  分段来看 $L_1$ 的轨迹：
  当 $t\geqslant 0$ 时，令 $s=t^3\geqslant 0$，轨迹为射线 $(s,s,s)$；
  当 $t<0$ 时，令 $s=t^3<0$，轨迹为射线 $(s,2s,3s)$。
  $L_1$ 是两条从原点出发但方向不同的射线的并集 —— 具体地，一条沿 $(1,1,1)$ 方向，另一条沿 $(1,2,3)$ 方向。在原点处形成一个“尖角”，因此\textbf{切线不存在}。
\end{enumerate}
\end{solution}

\begin{exercise}[参数曲线的基本计算] \label{ex:param-basic}
求曲线 $x=t^2+t,\ y=t^2-t,\ z=t^2$ 在点 $p_0=(6,2,4)$ 处的切线方程以及法平面方程。
\end{exercise}
\begin{solution}
\begin{enumerate}
  \item \textbf{确定对应的参数值 $t_0$：}
  将 $p_0=(6,2,4)$ 代入参数方程。从第三个方程得到 $z=t^2=4 \implies t=\pm 2$。
  检验 $t=2$：$x=4+2=6$ (符合)，$y=4-2=2$ (符合)，$z=4$ (符合)。
  检验 $t=-2$：$x=4-2=2 \neq 6$ (不符合)。
  故 $t_0=2$。

  \item \textbf{计算切向量：}
  分别对 $x, y, z$ 求导数得 $x'(t)=2t+1,\ y'(t)=2t-1,\ z'(t)=2t$。
  代入 $t_0=2$ 得 $\vec{r}'(2)=(5,3,4) \neq \boldsymbol{0}$。

  \item \textbf{写出切线和法平面方程：}
  切线方程（对称式）为 \(\displaystyle \boxed{\frac{x-6}{5}=\frac{y-2}{3}=\frac{z-4}{4}}\)。
  法平面方程为 $5(x-6)+3(y-2)+4(z-4)=0$，
  展开整理得 $\boxed{5x+3y+4z-52=0}$。
\end{enumerate}
\end{solution}


\begin{definition}[显式曲线的参数化] \label{def:explicit-curve}
有时曲线不是直接给出参数方程，而是以两个显式函数的形式给出：
\[
\begin{cases}
y=y(x),\\
z=z(x),
\end{cases}
\qquad x\in[a,b].\tag{$\ast$}
\]
这种情况下只需将其视为以自变量 $x$ 为参数的特殊参数曲线即可。
对 ($\ast$) 式定义的曲线，令 $\vec{r}(x)=\bigl(x,\;y(x),\;z(x)\bigr)$，其中 $x\in[a,b]$。
则切向量为 \(\displaystyle \vec{r}'(x)=\bigl(1,y'(x),z'(x)\bigr)\)。
\end{definition}

\begin{exercise}[显式曲线的切线与法平面] \label{ex:explicit-curve-ex}
求曲线 $y=x$, $z=x^2$ 在点 $M=(1,1,1)$ 处的切线方程和法平面方程。
\end{exercise}
\begin{solution}
\begin{enumerate}
  \item \textbf{引入参数化形式：}
  取 $x$ 为参数，得到曲线的参数方程形式为 $\vec{r}(x)=(x,x,x^2)$。

  \item \textbf{确定参数与计算切向量：}
  在点 $M=(1,1,1)$ 处，对应参数 $x_0=1$。
  对各分量求导得出切向量 $\vec{r}'(x) = (1,\;1,\;2x)$。
  代入 $x=1$ 得 $\vec{r}'(1)=(1,1,2)$。

  \item \textbf{写出几何方程：}
  \textbf{切线方程}（对称式）：\(\displaystyle \boxed{\frac{x-1}{1}=\frac{y-1}{1}=\frac{z-1}{2}}\)。
  \textbf{法平面方程}为 $(x-1)+(y-1)+2(z-1)=0$，
  展开整理得 $\boxed{x+y+2z-4=0}$。
\end{enumerate}
\end{solution}

在实际问题中，空间曲线常常以 \textbf{两个曲面的交集} 形式给出，此时如何求切线？关键在于利用隐函数定理将“隐式”转化为已知的“参数/显式”形式。

\begin{proposition}[交线切向量的 Jacobi 行列式公式] \label{prop:intersection-curve-tangent}
设曲线 $L$ 由方程组 $F_1(x,y,z)=0,\ F_2(x,y,z)=0$ 定义，$p_0=(x_0,y_0,z_0)\in L$，
且 $F_1$, $F_2$ 在 $p_0$ 的邻域内具有连续偏导数。

若 Jacobi 行列式 $\displaystyle J:=\left.\frac{\partial(F_1,F_2)}{\partial(y,z)}\right|_{p_0} \neq 0$，
则在 $p_0$ 的某邻域内$\begin{cases}F_1(x,y,z)=0,\\F_2(x,y,z)=0.\end{cases}$可确定一组隐函数 $y=y(x)$, $z=z(x)$，且 $L$ 在 $p_0$ 处的切向量可表示为
\[
\boxed{
\left.\left(
\frac{\partial(F_1,F_2)}{\partial(y,z)},\;
\frac{\partial(F_1,F_2)}{\partial(z,x)},\;
\frac{\partial(F_1,F_2)}{\partial(x,y)}
\right)\right|_{p_0}
:= (A,B,C)}.
\]
由此可得：切线方程为
\(\displaystyle \frac{x-x_0}{A}=\frac{y-y_0}{B}=\frac{z-z_0}{C}\)，
法平面方程为
\(\displaystyle A(x-x_0)+B(y-y_0)+C(z-z_0)=0\)。
\end{proposition}

\begin{exercise}[两柱面交线的切线与法平面] \label{ex:cylinder-intersection}
求两柱面 $x^2+y^2=a^2$ 与 $x^2+z^2=a^2$ 的交线
在点 $p_0=\left(\dfrac{a}{\sqrt{2}},\,\dfrac{a}{\sqrt{2}},\,\dfrac{a}{\sqrt{2}}\right)$
处的切线方程和法平面方程。
\end{exercise}
\begin{solution}
\begin{enumerate}
  \item \textbf{识别曲面并定义辅助函数：}
  令 $F_1(x,y,z)=x^2+y^2-a^2=0$（母线平行于 $z$ 轴的圆柱面），
  $F_2(x,y,z)=x^2+z^2-a^2=0$（母线平行于 $y$ 轴的圆柱面）。

  \item \textbf{计算三个 Jacobi 行列式（切向量的分量）：}
  \begin{align*}
  A &= \left.\frac{\partial(F_1,F_2)}{\partial(y,z)}\right|_{p_0}
   = \begin{vmatrix}
      \dfrac{\partial F_1}{\partial y}&\dfrac{\partial F_1}{\partial z}\\[8pt]
      \dfrac{\partial F_2}{\partial y}&\dfrac{\partial F_2}{\partial z}
    \end{vmatrix}_{p_0}
   = \begin{vmatrix} 2y & 0 \\ 0 & 2z \end{vmatrix}_{p_0}
   = 4yz\Big|_{p_0} = 4\left(\frac{a}{\sqrt{2}}\right)\left(\frac{a}{\sqrt{2}}\right) = 2a^2,\\[10pt]
  B &= \left.\frac{\partial(F_1,F_2)}{\partial(z,x)}\right|_{p_0}
   = \begin{vmatrix}
      \dfrac{\partial F_1}{\partial z}&\dfrac{\partial F_1}{\partial x}\\[8pt]
      \dfrac{\partial F_2}{\partial z}&\dfrac{\partial F_2}{\partial x}
    \end{vmatrix}_{p_0}
   = \begin{vmatrix} 0 & 2x \\ 2z & 2x \end{vmatrix}_{p_0}
   = -4xz\Big|_{p_0} = -4\left(\frac{a}{\sqrt{2}}\right)\left(\frac{a}{\sqrt{2}}\right) = -2a^2,\\[10pt]
  C &= \left.\frac{\partial(F_1,F_2)}{\partial(x,y)}\right|_{p_0}
   = \begin{vmatrix}
      \dfrac{\partial F_1}{\partial x}&\dfrac{\partial F_1}{\partial y}\\[8pt]
      \dfrac{\partial F_2}{\partial x}&\dfrac{\partial F_2}{\partial y}
    \end{vmatrix}_{p_0}
   = \begin{vmatrix} 2x & 2y \\ 2x & 0 \end{vmatrix}_{p_0}
   = -4xy\Big|_{p_0} = -4\left(\frac{a}{\sqrt{2}}\right)\left(\frac{a}{\sqrt{2}}\right) = -2a^2.
  \end{align*}
  切向量为 $(2a^2, -2a^2, -2a^2)$，提取公约数 $2a^2$ 后等价切向量为 $(1, -1, -1)$。

  \item \textbf{写出切线方程和法平面方程：}
  \textbf{切线方程}：\(\displaystyle \boxed{\frac{x-a/\sqrt2}{1}=\frac{y-a/\sqrt2}{-1}=\frac{z-a/\sqrt2}{-1}}\)。
  \textbf{法平面方程}为 $1\cdot\left(x-\dfrac{a}{\sqrt2}\right) - 1\cdot\left(y-\dfrac{a}{\sqrt2}\right) - 1\cdot\left(z-\dfrac{a}{\sqrt2}\right)=0$，
  即 $\boxed{x-y-z+\frac{\sqrt2}{2}a=0}$。
\end{enumerate}
\end{solution}

\begin{proposition}[向量叉乘法求切向量] \label{prop:cross-product-tangent}
设空间曲线 $L$ 由两个隐式曲面 $S_1: F_1(x,y,z)=0$ 与 $S_2: F_2(x,y,z)=0$ 交成。若 $p_0 \in L$ 且两个曲面在 $p_0$ 处的梯度向量 $\nabla F_1(p_0)$ 与 $\nabla F_2(p_0)$ 不平行，则曲线 $L$ 在 $p_0$ 处的切向量 $\vec{T}$ 可以表示为：
\(\displaystyle \boxed{\vec{T} = \nabla F_1(p_0) \times \nabla F_2(p_0)}\)。
其中，$\nabla F_1 = (F_{1x}, F_{1y}, F_{1z})$，$\nabla F_2 = (F_{2x}, F_{2y}, F_{2z})$。
\end{proposition}

\subsubsection{曲面的切平面与法线}

\begin{theorem}[显式曲面的法线与切平面方程] \label{thm:tangent-plane-explicit}
设 $z=f(x,y)$ 在 $(x_0,y_0)$ 处 \textbf{可微}，记
\[
z_0=f(x_0,y_0),\qquad p_0=(x_0,y_0,z_0).
\]
先把显式曲面改写成等值面
\[
F(x,y,z)=f(x,y)-z=0.
\]
在 $p_0$ 处，
\[
\nabla F(p_0)=\bigl(f_x(x_0,y_0),\,f_y(x_0,y_0),\,-1\bigr).
\]
为什么梯度是法线方向？任取一条位于曲面上且经过 $p_0$ 的光滑曲线
\[
\boldsymbol r(t)=(x(t),y(t),z(t)),\qquad \boldsymbol r(t_0)=p_0.
\]
由于曲线始终在曲面上，有 $F(\boldsymbol r(t))\equiv 0$。对 $t$ 求导得
\[
\nabla F(p_0)\cdot \boldsymbol r'(t_0)=0.
\]
而 $\boldsymbol r'(t_0)$ 正是曲面在 $p_0$ 处的切方向。也就是说，$\nabla F(p_0)$ 垂直于曲面上的一切切方向，因此它就是法线方向。
曲面 $z=f(x,y)$ 在 $p_0=(x_0,y_0,z_0)$ 处的\textbf{法线}（Normal Line）
是以法向量 $\vec{n}(p_0)=(-f_x,-f_y,1)\big|_{p_0}$ 为方向向量的直线，
其方程为 \(\displaystyle \boxed{\frac{x-x_0}{-f_x(x_0,y_0)}=\frac{y-y_0}{-f_y(x_0,y_0)}=\frac{z-z_0}{1}}\)。
切平面就是过 $p_0$ 且垂直于该法向量的平面，因此
\[
\vec n\cdot\bigl((x,y,z)-p_0\bigr)=0\implies 
\boxed{z-z_0=f_x(x_0,y_0)(x-x_0)+f_y(x_0,y_0)(y-y_0)},
\]
或写成行列式形式
\[
\begin{vmatrix}
x-x_0 & y-y_0 & z-z_0\\
1 & 0 & f_x(x_0,y_0)\\
0 & 1 & f_y(x_0,y_0)
\end{vmatrix}=0.
\]
\end{theorem}

\begin{definition}[方向余弦公式] \label{property:direction-cosines}
设法向量 $\vec{n}=(-f_x,-f_y,1)$ 与 $x$, $y$, $z$ 轴正向的夹角分别为
$\alpha$, $\beta$, $\gamma$（即\textbf{方向角}）。
则\textbf{方向余弦}（Direction Cosines）为
\begin{align*}
\cos\alpha=\frac{-f_x}{\pm\sqrt{1+f_x^2+f_y^2}},~~
\cos\beta=\frac{-f_y}{\pm\sqrt{1+f_x^2+f_y^2}},~~
\cos\gamma=\frac{1}{\pm\sqrt{1+f_x^2+f_y^2}},
\end{align*}
其中根号前的符号必须同时取正号或同时取负号：
\begin{itemize}
  \item 若法向量指 \textbf{向上}（与 $z$ 轴正向成锐角，$\gamma\in(0,\pi/2)$），
        则 $\cos\gamma>0$，三者 \textbf{均取正号}；
  \item 反之（法向量向下），三者 \textbf{均取负号}。
\end{itemize}
\end{definition}

\begin{proposition}[隐式曲面的切平面与法线] \label{prop:implicit-surface}
设 $F(x,y,z)=0$ 定义了隐式曲面 $S$，$p_0=(x_0,y_0,z_0)\in S$。
若 $F$ 在 $p_0$ 处可微且梯度 $\nabla F(p_0)=(F_x(p_0),F_y(p_0),F_z(p_0)) \neq \boldsymbol{0}$，则切平面方程与法线方程分别为
\begin{align*}
&\boxed{F_x(p_0)(x-x_0)+F_y(p_0)(y-y_0)+F_z(p_0)(z-z_0)=0},\\
&\boxed{\frac{x-x_0}{F_x(p_0)}
=\frac{y-y_0}{F_y(p_0)}
=\frac{z-z_0}{F_z(p_0)}}.
\end{align*}
\end{proposition}

\begin{remark}[奇点问题]
与参数曲线的情形类似，若 $\nabla F(p_0)=\boldsymbol{0}$（即
$F_x(p_0)=F_y(p_0)=F_z(p_0)=0$），
则不能判定曲面 $S$ 在 $p_0$ 处是否存在切平面。
使 $\nabla F=\boldsymbol{0}$ 的点称为隐式曲面 $S$ 的\textbf{奇点}（Singular Point）。
例如圆锥面 $x^2+y^2-z^2=0$ 在顶点 $(0,0,0)$ 处就是一个奇点。
\end{remark}

\begin{exercise}[显式曲面：反正切曲面] \label{ex:arctan-surface}
求曲面 $z=\arctan\dfrac{y}{x}$ 在点 $p_0=\left(1,1,\dfrac{\pi}{4}\right)$ 处
的切平面方程和法线方程。
\end{exercise}
\begin{solution}
\begin{enumerate}
  \item \textbf{计算一阶偏导数：}
  应用复合函数求导法则，\(\displaystyle z_x=\frac{\partial}{\partial x}\arctan\frac{y}{x}=-\frac{y}{x^2+y^2}\)，
  \(\displaystyle z_y=\frac{\partial}{\partial y}\arctan\frac{y}{x}=\frac{x}{x^2+y^2}\)。

  \item \textbf{在 $p_0$ 处取值：}
  代入 $x=1, y=1$，得 $z_x(1,1)=-\dfrac{1}{1+1}=-\dfrac{1}{2}, \quad z_y(1,1)=\dfrac{1}{1+1}=\dfrac{1}{2}$。

  \item \textbf{确定法向量和写出方程：}
  法向量（显式曲面形式）为 $\vec{n}=(-z_x,-z_y,1)\big|_{p_0}=\left(\dfrac{1}{2},-\dfrac{1}{2},1\right)$。
  \textbf{切平面方程}为 $\dfrac{1}{2}(x-1)-\dfrac{1}{2}(y-1)+\left(z-\dfrac{\pi}{4}\right)=0$，
  整理得 $\boxed{x-y+2z-\dfrac{\pi}{2}-1=0}$，或写成 $\boxed{x-y+2z=1+\dfrac{\pi}{2}}$。
  \textbf{法线方程}（消去分母后）为 $\boxed{\dfrac{x-1}{1}=\dfrac{y-1}{-1}=\dfrac{z-\pi/4}{2}}$。
\end{enumerate}
\end{solution}

\begin{exercise}[隐式曲面：椭球面] \label{ex:ellipsoid-tangent}
求椭球面 $\dfrac{x^2}{a^2}+\dfrac{y^2}{b^2}+\dfrac{z^2}{c^2}=1$
在点 $p_0=(x_0,y_0,z_0)$（椭球面上任一点）处的切平面方程。
\end{exercise}
\begin{solution}
\begin{enumerate}
  \item \textbf{构造辅助隐式函数：}
  令 $F(x,y,z)=\dfrac{x^2}{a^2}+\dfrac{y^2}{b^2}+\dfrac{z^2}{c^2}-1=0$。

  \item \textbf{计算梯度：}
  计算偏导数得 $F_x=\dfrac{2x}{a^2},\ F_y=\dfrac{2y}{b^2},\ F_z=\dfrac{2z}{c^2}$。
  在 $p_0$ 处，法向量为 $\nabla F(p_0)=\left(\dfrac{2x_0}{a^2},\,\dfrac{2y_0}{b^2},\,\dfrac{2z_0}{c^2}\right)$。由于 $p_0$ 在椭球面上，至少有一个坐标分量非零，故 $\nabla F(p_0) \neq \boldsymbol{0}$ 成立。

  \item \textbf{代入切平面公式并化简：}
  切平面方程为 $\dfrac{2x_0}{a^2}(x-x_0)+\dfrac{2y_0}{b^2}(y-y_0)+\dfrac{2z_0}{c^2}(z-z_0)=0$。
  利用 $p_0$ 满足椭球面方程 $\dfrac{x_0^2}{a^2}+\dfrac{y_0^2}{b^2}+\dfrac{z_0^2}{c^2}=1$，上式可整理为 $\boxed{\dfrac{xx_0}{a^2}+\dfrac{yy_0}{b^2}+\dfrac{zz_0}{c^2}=1}$。
\end{enumerate}
\end{solution}

\begin{remark}
这个结果具有明显的对称性：左端三项的结构与椭球面方程本身相同，只是每个变量的平方被替换成“当前坐标与切点坐标的乘积”。这种形式在二次曲面的切平面理论中经常出现。例如当 $a=b=c=R$ 时为球面，切平面方程为 $xx_0+yy_0+zz_0=R^2$。
\end{remark}

\subsection{隐函数专题：单方程、方程组与雅可比行列式}

\subsubsection{单方程隐函数存在定理}

\begin{theorem}[隐函数存在定理 / Implicit Function Theorem] \label{thm:implicit-function}
设二元函数 $F(x,y)$ 在点 $(x_0,y_0)$ 的邻域内满足以下三条件：
\begin{enumerate}
  \item[(A)] 在矩形区域 $D=\{(x,y):|x-x_0|<a,\;|y-y_0|<b\}$ 内，$F$ 关于 $x$ 和 $y$ 具有\textbf{连续偏导数}；
  \item[(B)] $F(x_0,y_0)=0$；
  \item[(C)] $F_y(x_0,y_0) \neq 0$。
\end{enumerate}
则以下结论成立：
\begin{enumerate}
  \item[\textbf{(I)}] \textbf{存在性}：存在 $\eta>0$ 和定义在 $B_\eta(x_0)=\{x:|x-x_0|<\eta\}$ 上的唯一函数 $y=f(x)$，使得 $F\bigl(x,f(x)\bigr)\equiv 0$ 对所有 $x\in B_\eta(x_0)$ 成立，且 $y_0=f(x_0)$；
  \item[\textbf{(II)}] \textbf{连续性}：$f$ 在 $B_\eta(x_0)$ 内\textbf{连续}；
  \item[\textbf{(III)}] \textbf{可微性与求导公式}：$f$ 在 $B_\eta(x_0)$ 内具有\textbf{连续导数}，且 \(\displaystyle \boxed{\;f'(x)=-\frac{F_x(x,y)}{F_y(x,y)}\Big|_{y=f(x)}\;}\)。
\end{enumerate}
\end{theorem}

\begin{property}[隐函数求导公式 / Implicit Derivative Formula] \label{property:implicit-deriv-formula}
设方程 $F(x,y)=0$ 满足隐函数存在定理的条件，确定了可微隐函数 $y=f(x)$。
则在 $f$ 有定义的范围内：\(\displaystyle \boxed{\;
\frac{\mathrm{d}y}{\mathrm{d}x}
=-\frac{F_x(x,y)}{F_y(x,y)}
\Biggr|_{y=f(x)}
\;}\)。
\end{property}

\begin{note}
既然方程 $F(x,y)=0$ 确定了隐函数 $y=f(x)$，这就意味着如果我们把 $f(x)$ 代回原方程，必定会得到一个关于 $x$ 的恒等式 $F(x, f(x)) \equiv 0$。
现在，我们将等号左边看作是一个关于一元变量 $x$ 的复合函数。我们对等式两边同时关于 $x$ 求导。根据多元函数链式法则，自变量 $x$ 通过两条路径影响着 $F$ 的值（一条是直接作为第一变量 $x$，另一条是间接通过第二变量 $y$）：
\[
\frac{\mathrm{d}}{\mathrm{d}x} F(x, f(x)) = \frac{\partial F}{\partial x} \cdot \frac{\mathrm{d}x}{\mathrm{d}x} + \frac{\partial F}{\partial y} \cdot \frac{\mathrm{d}y}{\mathrm{d}x} = 0
\]
因为 $\dfrac{\mathrm{d}x}{\mathrm{d}x} = 1$，且为了书写简便，将偏导数记为 $F_x$ 和 $F_y$，上式即变为大家熟悉的 $F_x + F_y \cdot f'(x) = 0$。只要 $F_y \neq 0$，移项即可解出隐函数的导数：$f'(x) = \dfrac{\mathrm{d}y}{\mathrm{d}x} = -\dfrac{F_x}{F_y}$。
\end{note}

上述定理可以直接推广到 $n$ 元函数的情形。

\begin{theorem}[多元隐函数存在定理] \label{thm:implicit-function-nvar}
设开集 $D\subset\mathbb R^n$，$F:D\to\mathbb R$ 满足：
\begin{enumerate}
  \item 在 $U=\{(x_1,\dots,x_n):|x_i-x_i^{(0)}|<a_i,\;i=1,\dots,n\}\subset D$ 内，$F$ 对各变量具有\textbf{连续偏导数}；
  \item $F\bigl(x_1^{(0)},\dots,x_n^{(0)}\bigr)=0$；
  \item $F_{x_n}\bigl(x_1^{(0)},\dots,x_n^{(0)}\bigr) \neq 0$。
\end{enumerate}
则在点 $\boldsymbol{x}^{(0)}=(x_1^{(0)},\dots,x_n^{(0)})$ 的某邻域内，方程 $F(x_1,\dots,x_n)=0$ 唯一确定一个可微函数 $x_n=f(x_1,\dots,x_{n-1})$，满足 $x_n^{(0)}=f(x_1^{(0)},\dots,x_{n-1}^{(0)})$，且 $f$ 具有连续偏导数：
\[
\boxed{
\frac{\partial f}{\partial x_i}
=-\frac{F_{x_i}(x_1,\dots,x_{n-1},\,f(\cdots))}
       {F_{x_n}(x_1,\dots,x_{n-1},\,f(\cdots))},
\qquad i=1,\dots,n-1.
}
\]
\end{theorem}

特别地，对于三元函数 $F(x,y,z)=0$ 确定隐函数 $z=z(x,y)$ 的情形：
若 $F_z \neq 0$，则 $\dfrac{\partial z}{\partial x}=-\dfrac{F_x}{F_z}$，$\dfrac{\partial z}{\partial y}=-\dfrac{F_y}{F_z}$。
进而全微分为
$\mathrm{d}z=\dfrac{\partial z}{\partial x}\,\mathrm{d}x + \dfrac{\partial z}{\partial y}\,\mathrm{d}y
    =-\dfrac{F_x}{F_z}\,\mathrm{d}x-\dfrac{F_y}{F_z}\,\mathrm{d}y$。

\begin{remark}[关于隐函数定理的若干要点]
\begin{enumerate}
  \item \textbf{条件的充分性}：定理的三条条件仅是\textbf{充分的}而非必要的。例如方程 $F(x,y)=y^3-x^3=0$ 确定唯一的隐函数 $y=x$（处处可微），但在原点处 $F_y(0,0)=3y^2|_{(0,0)}=0$，不满足条件 (C)。反之，若条件 (C) 不满足，结论确实可能失效（见下一点）。
  \item \textbf{条件不满足时可能失效}：考虑 $F(x,y)=(x^2+y^2)^2-x^2+y^2=0$（双纽线的一部分），则 $F_y(0,0)=0$。在 $(0,0)$ 的无论多小的邻域内，同一个 $x \neq 0$ 总是对应多个 $y$ 值，故不存在以 $x$ 为自变量的隐函数 $y=f(x)$。
  \item \textbf{全微分法的优势}：在实际计算中，往往不需要先确认定理的所有条件（通常题目已隐含“隐函数存在且可微”的前提），直接对方程两边求全微分，再整理出所求偏导数即可。这种方法的优点是不需要事先区分谁是自变量谁是因变量，统一的微分运算自动处理了所有的依赖关系。
\end{enumerate}
\end{remark}

\begin{exercise}[隐函数的全微分] \label{ex:implicit-dz}
设 $z=z(x,y)$ 是由方程 $z-y-x+x\,\mathrm{e}^{\,z-y-x}=0$（记为 $(\dagger)$）确定的隐函数。求 $\mathrm{d}z$。
\end{exercise}
\begin{solution}
设方程为 $z-y-x+x\,\mathrm{e}^{z-y-x}=0 $。
\begin{enumerate}
    \item[\textbf{方法一：}] \textbf{隐函数求导公式法}
    设 $F(x,y,z) = z - y - x + x\mathrm{e}^{z-y-x}$，分别求三个变量的偏导数：
    \begin{itemize}
        \item $F_x = -1 + \mathrm{e}^{z-y-x} + x\mathrm{e}^{z-y-x}(-1) = -1 + (1-x)\mathrm{e}^{z-y-x}$；
        \item $F_y = -1 + x\mathrm{e}^{z-y-x}(-1) = -(1 + x\mathrm{e}^{z-y-x})$；
        \item $F_z = 1 + x\mathrm{e}^{z-y-x} \cdot 1 = 1 + x\mathrm{e}^{z-y-x}$。
    \end{itemize}
    根据隐函数求导公式 $z_x = -\frac{F_x}{F_z}$ 和 $z_y = -\frac{F_y}{F_z}$：
    \begin{align*}
        z_x &= -\frac{-1 + (1-x)\mathrm{e}^{z-y-x}}{1 + x\mathrm{e}^{z-y-x}} = \frac{1 + (x-1)\mathrm{e}^{z-y-x}}{1 + x\mathrm{e}^{z-y-x}}, \\
        z_y &= -\frac{-(1 + x\mathrm{e}^{z-y-x})}{1 + x\mathrm{e}^{z-y-x}} = 1, \\
        \mathrm{d}z &= z_x \mathrm{d}x + z_y \mathrm{d}y = \frac{1 + (x-1)\mathrm{e}^{z-y-x}}{1 + x\mathrm{e}^{z-y-x}}\mathrm{d}x + \mathrm{d}y.
    \end{align*}
    \item[\textbf{方法二：}] \textbf{全微分法（直接微分法）}
    直接对原方程两边取全微分：
    \begin{align*}
    \mathrm{d}z-\mathrm{d}y-\mathrm{d}x
    +\mathrm{e}^{z-y-x}\mathrm{d}x
    +x\mathrm{e}^{z-y-x}(\mathrm{d}z-\mathrm{d}y-\mathrm{d}x)=0.
    \end{align*}
    整理得
    \begin{align*}
    (1+x\mathrm{e}^{z-y-x})\mathrm{d}z
    -(1+x\mathrm{e}^{z-y-x})\mathrm{d}y
    +\bigl[-1+(1-x)\mathrm{e}^{z-y-x}\bigr]\mathrm{d}x=0.
    \end{align*}
    因此
    \begin{align*}
    \mathrm{d}z
    &=\mathrm{d}y+
    \frac{1+(x-1)\mathrm{e}^{z-y-x}}{1+x\mathrm{e}^{z-y-x}}\,\mathrm{d}x.
    \end{align*}
\end{enumerate}
\end{solution}

\subsubsection{方程组确定的隐函数} \label{def:implicit-system}

上面讨论的是一个方程确定一个隐函数的情形。在实际问题中更常出现的是\textbf{方程组}（System of Equations）——多个方程联合确定多个隐函数的情况。
考虑一般的方程组。设有 $m$ 个方程涉及 $m+n$ 个变量：
\[
\begin{cases}
F_1(x_1,\dots,x_n,\,u_1,\dots,u_m)=0,\\
F_2(x_1,\dots,x_n,\,u_1,\dots,u_m)=0,\\
\qquad\vdots\\
F_m(x_1,\dots,x_n,\,u_1,\dots,u_m)=0.
\end{cases}
\tag{$\ast\ast$}
\]
直觉上，$m$ 个方程可以对 $m+n$ 个变量施加 $m$ 个独立的约束，剩下 $n$ 个“自由度”——这 $n$ 个变量可以作为\textbf{自变量}，而被约束的 $m$ 个变量则成为这些自变量的\textbf{（隐）函数}。具体而言：
\begin{itemize}
  \item 若方程组 ($\ast\ast$) 含 $3$ 个变量 $(x,y,z)$ 和 $2$ 个方程 $\Rightarrow$ 可确定 $2$ 个一元函数 $y=y(x)$，$z=z(x)$；
  \item 若含 $4$ 个变量 $(x,y,u,v)$ 和 $2$ 个方程 $\Rightarrow$ 可确定 $2$ 个二元函数 $u=u(x,y)$，$v=v(x,y)$。
\end{itemize}

\begin{definition}[Jacobi 行列式 / Jacobian Determinant] \label{def:jacobi-det}
设 $F,G:\mathbb R^{n+m}\to\mathbb R$ 为两个可微函数。称行列式
\[
\frac{\partial(F,G)}{\partial(u,v)}
:=\begin{vmatrix}
\dfrac{\partial F}{\partial u} & \dfrac{\partial F}{\partial v}\\[8pt]
\dfrac{\partial G}{\partial u} & \dfrac{\partial G}{\partial v}
\end{vmatrix}
\]
为 $F,G$ 关于变量组 $(u,v)$ 的\textbf{Jacobi 行列式}（或称\textbf{函数行列式}）。
\end{definition}

Jacobi 行列式在方程组隐函数理论中的作用，类似于单个方程情形下的 $F_y \neq 0$ 条件：它是“能否将 $(u,v)$ 解出为其他变量函数”的判别标准。
现在讨论最常见的情形：四个变量 $(x,y,u,v)$ 通过两个方程
\[
\begin{cases}
F(x,y,u,v)=0,\\[1ex]
G(x,y,u,v)=0
\end{cases}
\tag{$\ast\ast\ast$}
\]
确定了两个二元函数 $u=u(x,y)$，$v=v(x,y)$。
\begin{note}[Jacobian 比值公式的来源]
一旦确认 $(u,v)$ 是 $(x,y)$ 的隐函数，对方程组关于 $x$ 求导就得到
\[
\begin{cases}
F_u u_x + F_v v_x = -F_x,\\
G_u u_x + G_v v_x = -G_x.
\end{cases}
\]
于是分母为什么总是 $\dfrac{\partial(F,G)}{\partial(u,v)}$ 就很清楚了：它本来就是这个二元一次方程组关于未知量 $u_x,v_x$ 的系数行列式。
若要求的是 $u_x$，按 Cramer 法则就把第一列换成右端常数列 $(-F_x,-G_x)^{\mathrm T}$，所以分子变成
\[
\begin{vmatrix}
-F_x & F_v\\
-G_x & G_v
\end{vmatrix}
=-
\begin{vmatrix}
F_x & F_v\\
G_x & G_v
\end{vmatrix}
=-
\frac{\partial(F,G)}{\partial(x,v)}.
\]
因此
$u_x=
-\dfrac{\dfrac{\partial(F,G)}{\partial(x,v)}}
       {\dfrac{\partial(F,G)}{\partial(u,v)}}$。
这里前面的负号也不是额外规定出来的，而是因为关于 $x$ 的项在求导后被移到了等号右边。
\end{note}

\begin{theorem}[方程组形式的隐函数存在定理] \label{thm:implicit-system-exist}
设 $P_0=(x_0,y_0,u_0,v_0)\in\mathbb R^4$。若函数 $F(x,y,u,v)$ 和 $G(x,y,u,v)$ 满足：
\begin{enumerate}
  \item[(A')] $F,G$ 在 $P_0$ 的某邻域内对各变量具有\textbf{连续偏导数}；
  \item[(B')] $F(P_0)=0$ 且 $G(P_0)=0$；
  \item[(C')] 在 $P_0$ 处的 Jacobi 行列式
        $\displaystyle\left.\frac{\partial(F,G)}{\partial(u,v)}\right|_{P_0}
        =
        \begin{vmatrix}
          F_u(P_0)&F_v(P_0)\\
          G_u(P_0)&G_v(P_0)
        \end{vmatrix}
        \neq 0$。
\end{enumerate}
则存在 $P_0$ 的一个邻域，在此邻域内方程组 $\begin{cases}F(x,y,u,v)=0,\\G(x,y,u,v)=0\end{cases}$ 唯一确定一组可微函数 $u=u(x,y)$，$v=v(x,y)$，满足 $u_0=u(x_0,y_0)$，$v_0=v(x_0,y_0)$，且 $u,v$ 具有关于 $x,y$ 的\textbf{连续偏导数}。
\end{theorem}

\begin{exercise}[线性方程组的隐函数] \label{ex:implicit-system1}
设方程组
\[
\begin{cases}
xu-yv=0,\\
yu+xv=1.
\end{cases}
\]
求 $\dfrac{\partial u}{\partial x}$，$\dfrac{\partial v}{\partial x}$，$\dfrac{\partial u}{\partial y}$，$\dfrac{\partial v}{\partial y}$。
\end{exercise}
\begin{solution}
\textbf{识别变量}：这里 $(x,y)$ 为自变量，$(u,v)$ 为因变量（隐函数）。
方程组可写为 $F(x,y,u,v)=xu-yv=0$，$G(x,y,u,v)=yu+xv-1=0$。

\begin{enumerate}
  \item[\textbf{方法一：}] \textbf{全微分法}

  对两方程分别求全微分，并将含 $\mathrm{d}u$, $\mathrm{d}v$ 的项留在左边：
  \begin{align*}
  \begin{cases}
  x\,\mathrm{d}u+u\,\mathrm{d}x-y\,\mathrm{d}v-v\,\mathrm{d}y=0,\\[6pt]
  y\,\mathrm{d}u+u\,\mathrm{d}y+x\,\mathrm{d}v+v\,\mathrm{d}x=0.
  \end{cases}
  &\Longrightarrow
  \begin{cases}
  x\,\mathrm{d}u-y\,\mathrm{d}v=-u\,\mathrm{d}x+v\,\mathrm{d}y,\\[6pt]
  y\,\mathrm{d}u+x\,\mathrm{d}v=-v\,\mathrm{d}x-u\,\mathrm{d}y.
  \end{cases}
  \end{align*}

  系数矩阵（Jacobi 矩阵）为 $A=\begin{pmatrix} x & -y \\ y & x \end{pmatrix}$，其行列式 $J=\det A=x^2+y^2$（当 $(x,y) \neq (0,0)$ 时非零）。

  由 Cramer 法则解得：
  \begin{align*}
  \mathrm{d}u &= \frac{1}{J}\det\begin{pmatrix} -u\,\mathrm{d}x+v\,\mathrm{d}y & -y \\ -v\,\mathrm{d}x-u\,\mathrm{d}y & x \end{pmatrix}
    = \frac{x(-u\,\mathrm{d}x+v\,\mathrm{d}y) - (-y)(-v\,\mathrm{d}x-u\,\mathrm{d}y)}{x^2+y^2} \\
    &= \frac{-(xu+yv)\,\mathrm{d}x + (xv-yu)\,\mathrm{d}y}{x^2+y^2}. \\[8pt]
  \mathrm{d}v &= \frac{1}{J}\det\begin{pmatrix} x & -u\,\mathrm{d}x+v\,\mathrm{d}y \\ y & -v\,\mathrm{d}x-u\,\mathrm{d}y \end{pmatrix}
    = \frac{x(-v\,\mathrm{d}x-u\,\mathrm{d}y) - y(-u\,\mathrm{d}x+v\,\mathrm{d}y)}{x^2+y^2} \\
    &= \frac{(yu-xv)\,\mathrm{d}x - (xu+yv)\,\mathrm{d}y}{x^2+y^2}.
  \end{align*}

  分别读取 $\mathrm{d}x$ 和 $\mathrm{d}y$ 的系数即可得到各偏导数：
  \[
  \begin{aligned}
  \frac{\partial u}{\partial x} &= -\frac{xu+yv}{x^2+y^2},&
  \frac{\partial u}{\partial y} &= \frac{xv-yu}{x^2+y^2},\\
  \frac{\partial v}{\partial x} &= \frac{yu-xv}{x^2+y^2},&
  \frac{\partial v}{\partial y} &= -\frac{xu+yv}{x^2+y^2}.
  \end{aligned}
  \]

  \item[\textbf{方法二：}] \textbf{逐项偏导验证}

  对原方程组关于 $x$ 求偏导：
  \begin{align*}
  \begin{cases}
  u+xu_x-yv_x=0,\\
  yu_x+v+xv_x=0,
  \end{cases}
  &\Longrightarrow
  \begin{cases}
  xu_x-yv_x=-u,\\
  yu_x+xv_x=-v,
  \end{cases}
  \end{align*}
  利用常规二元一次方程组的解法可得
  $u_x=-\dfrac{xu+yv}{x^2+y^2}$，$v_x=\dfrac{yu-xv}{x^2+y^2}$。
  若关于 $y$ 求偏导，则必须重新把“与 $y$ 有关的显式项”一并求出，而不是简单口头略过。对原方程组关于 $y$ 求偏导得
  \begin{align*}
  \begin{cases}
  xu_y-v-yv_y=0,\\
  u+yu_y+xv_y=0,
  \end{cases}
  &\Longrightarrow
  \begin{cases}
  xu_y-yv_y=v,\\
  yu_y+xv_y=-u.
  \end{cases}
  \end{align*}
  再解这个线性方程组，得到
  $u_y=\dfrac{xv-yu}{x^2+y^2}$，$v_y=-\dfrac{xu+yv}{x^2+y^2}$。
  这与方法一完全一致。
\end{enumerate}
\end{solution}

\begin{exercise}[非线性方程组的隐函数] \label{ex:implicit-system2}
设
\[
\begin{cases}
x=-u^2+v+z,\\
y=u+vz.
\end{cases}
\]
求 $\dfrac{\partial u}{\partial x}$，$\dfrac{\partial v}{\partial x}$，$\dfrac{\partial u}{\partial z}$。
\end{exercise}
\begin{solution}
识别变量：此处 $(x,y,z)$ 为自变量，$(u,v)$ 为隐函数。

\begin{enumerate}
  \item \textbf{求全微分}
  分别对两式求全微分：
  \[
  \begin{cases}
  \mathrm{d}x = -2u\,\mathrm{d}u+\mathrm{d}v+\mathrm{d}z,\\[6pt]
  \mathrm{d}y = \mathrm{d}u+z\,\mathrm{d}v+v\,\mathrm{d}z.
  \end{cases}
  \]

  \item \textbf{整理为标准形式}
  将含 $\mathrm{d}u$, $\mathrm{d}v$ 的项移到左边以求解因变量的微分：
  \[
  \begin{cases}
  2u\,\mathrm{d}u-\mathrm{d}v = -\mathrm{d}x+\mathrm{d}z,\\[6pt]
  -\mathrm{d}u+z\,\mathrm{d}v = \mathrm{d}y-v\,\mathrm{d}z.
  \end{cases}
  \]
  此时左侧未知数的系数矩阵为 $A=\begin{pmatrix} 2u & -1 \\ -1 & z \end{pmatrix}$，Jacobi 行列式 $J=\det A=2uz-1$（假设 $2uz \neq 1$）。

  \item \textbf{用 Cramer 法则求解并提取系数}
  \begin{align*}
  \mathrm{d}u &= \frac{1}{J}\det\begin{pmatrix}
    -\mathrm{d}x+\mathrm{d}z & -1\\
    \mathrm{d}y-v\,\mathrm{d}z & z
  \end{pmatrix}
   = \frac{z(-\mathrm{d}x+\mathrm{d}z)-(-1)(\mathrm{d}y-v\,\mathrm{d}z)}{2uz-1}\\[8pt]
  &= \frac{-z\,\mathrm{d}x+\mathrm{d}y+(z-v)\,\mathrm{d}z}{2uz-1}. \\[12pt]
  \mathrm{d}v &= \frac{1}{J}\det\begin{pmatrix}
    2u & -\mathrm{d}x+\mathrm{d}z\\
    -1 & \mathrm{d}y-v\,\mathrm{d}z
  \end{pmatrix}
   = \frac{2u(\mathrm{d}y-v\,\mathrm{d}z)-(-1)(-\mathrm{d}x+\mathrm{d}z)}{2uz-1}\\[8pt]
  &= \frac{-\mathrm{d}x+2u\,\mathrm{d}y+(1-2uv)\,\mathrm{d}z}{2uz-1}.
  \end{align*}

  比较 $\mathrm{d}u$, $\mathrm{d}v$ 表达式中各微分的系数（即各项偏导数）：
  \[
  \boxed{\frac{\partial u}{\partial x} = -\frac{z}{2uz-1},\qquad
  \frac{\partial v}{\partial x} = -\frac{1}{2uz-1},\qquad
  \frac{\partial u}{\partial z} = \frac{z-v}{2uz-1}.}
  \]
\end{enumerate}
\end{solution}

\subsection{泰勒专题：展开、余项与海森矩阵}

先通过一个具体例子感受一下线性近似的精度极限。

\begin{exercise}[线性近似计算] \label{ex:taylor-linear}
求 $Q:=\dfrac{1.03^2}{\sqrt{0.98}\,\sqrt[3]{1.06}}$ 的近似值。
\end{exercise}
\begin{solution}
\begin{enumerate}
  \item \textbf{构造多元函数与确定基准点：}
  构造三元函数 $f(x,y,z)=\dfrac{x^2}{\sqrt{y}\,\sqrt[3]{z}}$（定义域为 $y>0$, $z>0$），则问题归结为求 $f(1.03,\;0.98,\;1.06)$ 的值。
  取展开中心 $p_0=(1,1,1)^{\mathsf T}$，对应的增量向量为 $\boldsymbol h=(0.03,\,-0.02,\,0.06)^{\mathsf T}$。此时 $\|\boldsymbol h\|=\sqrt{0.03^2+(-0.02)^2+0.06^2}=\sqrt{0.0049}\approx 0.07$，确实较小。

  \item \textbf{计算基准点处的函数值与梯度：}
  由于函数表达式为幂次乘积，采用对数求导法更为简便：$\ln f = 2\ln x - \frac{1}{2}\ln y - \frac{1}{3}\ln z$。
  对两边同时求各个变量的偏导数可得 $\frac{1}{f}f_x = \frac{2}{x}, \ \frac{1}{f}f_y = -\frac{1}{2y}, \ \frac{1}{f}f_z = -\frac{1}{3z}$。
  因此 $\nabla f(x,y,z) = f(x,y,z)\left( \frac{2}{x},\; -\frac{1}{2y},\; -\frac{1}{3z} \right)$。
  代入 $p_0=(1,1,1)$ 且 $f(1,1,1)=1$，得 $\nabla f(1,1,1) = 1 \cdot \left(2,\,-\frac{1}{2},\,-\frac{1}{3}\right) = \left(2,\,-\frac{1}{2},\,-\frac{1}{3}\right)$。

  \item \textbf{运用一阶线性近似：}
  \begin{align*}
  Q &\approx f(p_0) + \nabla f(p_0)\cdot\boldsymbol h = 1 + \left(2,\,-\frac{1}{2},\,-\frac{1}{3}\right) \cdot (0.03,\,-0.02,\,0.06) \\
  &= 1 + 2(0.03) - \frac{1}{2}(-0.02) - \frac{1}{3}(0.06) = 1 + 0.06 + 0.01 - 0.02 = 1.05.
  \end{align*}
  因此近似值为 $\boxed{1.05}$。

  \item \textbf{精度检验：} 实际精确值约为 $1.0489\dots$，绝对误差约 $0.0011$，相对误差约 $0.11\%$ ——对于大多数工程目的已经足够好。
\end{enumerate}
\end{solution}

\subsubsection{二元函数的 Taylor 定理}

\begin{theorem}[Taylor 定理 / Taylor's Theorem] \label{thm:taylor}
设 $z=f(x,y)$ 在点 $p_0=(x_0,y_0)^{\mathsf T}$ 的某个开邻域 $D$ 内具有直到 $n+1$ 阶的\textbf{连续偏导数}。记
\[
\Delta x=x-x_0,\qquad \Delta y=y-y_0,
\qquad \boldsymbol h=(\Delta x,\Delta y)^{\mathsf T}.
\]
则存在 $\rho>0$，使得当 $\|\boldsymbol h\|<\rho$ 时，可以按下面的规律作 Taylor 展开：
\begin{align*}
f(x_0+\Delta x,\ y_0+\Delta y)
&=f(x_0,y_0)\\
&\quad+\bigl(f_x\Delta x+f_y\Delta y\bigr)\\
&\quad+\frac1{2!}\bigl(f_{xx}(\Delta x)^2+2f_{xy}\Delta x\Delta y+f_{yy}(\Delta y)^2\bigr)\\
&\quad+\frac1{3!}\bigl(f_{xxx}(\Delta x)^3+3f_{xxy}(\Delta x)^2\Delta y
      +3f_{xyy}\Delta x(\Delta y)^2+f_{yyy}(\Delta y)^3\bigr)\\
&\quad+\cdots\\
&\quad+\mathcal R_n.
\end{align*}
上式中，除最后的余项外，所有偏导数都在 $(x_0,y_0)$ 处取值；若只展开到一阶或二阶，就在对应阶数处截断即可。一般地，第 $k$ 阶项为
\[
\frac1{k!}\sum_{j=0}^{k}\binom{k}{j}
\frac{\partial^k f}{\partial x^{\,k-j}\partial y^{\,j}}(x_0,y_0)
(\Delta x)^{k-j}(\Delta y)^j,\qquad k=1,2,\dots,n.
\]

其中 Lagrange 余项为
\[
\begin{aligned}
\mathcal R_n
&=\frac1{(n+1)!}\Biggl[
\frac{\partial^{n+1}f}{\partial x^{\,n+1}}(p_\theta)(\Delta x)^{n+1}\\
&\quad+\binom{n+1}{1}\frac{\partial^{n+1}f}{\partial x^{\,n}\partial y}(p_\theta)(\Delta x)^n\Delta y+\cdots\\
&\quad+\binom{n+1}{j}\frac{\partial^{n+1}f}{\partial x^{\,n+1-j}\partial y^{\,j}}(p_\theta)
(\Delta x)^{n+1-j}(\Delta y)^j+\cdots\\
&\quad+\frac{\partial^{n+1}f}{\partial y^{\,n+1}}(p_\theta)(\Delta y)^{n+1}
\Biggr],
\end{aligned}
\]
这里
\[
p_\theta=(x_0+\theta\Delta x,\ y_0+\theta\Delta y),\qquad 0<\theta<1.
\]
\end{theorem}

\begin{remark}[微分算子与前面展开式的对应关系] \label{prop:diff-op-power}
前面 Taylor 公式已经把一阶、二阶、三阶项写开了。这里的“微分算子”不是新的公式，只是把这些展开项收拢成一个短记号。

记
\[
D:=\Delta x\,\frac{\partial}{\partial x}
+\Delta y\,\frac{\partial}{\partial y}.
\]
这里要特别注意：\(\Delta x,\Delta y\) 是从展开中心出发的固定增量。用 \(D\) 去作用时，真正被求偏导的是 \(f\) 以及 \(f\) 的各阶偏导数；\(\Delta x,\Delta y\) 只当作常数系数保留下来。

先作用一次：
\[
Df=f_x\Delta x+f_y\Delta y,
\]
这正是 Taylor 公式的一次项。

再算二次项。所谓 \(D^2f\)，就是把 \(D\) 作用到 \(Df\) 上：
\begin{align*}
D^2f
&=D(Df)\\
&=\left(\Delta x\frac{\partial}{\partial x}
+\Delta y\frac{\partial}{\partial y}\right)
\bigl(f_x\Delta x+f_y\Delta y\bigr)\\
&=\Delta x\frac{\partial}{\partial x}
\bigl(f_x\Delta x+f_y\Delta y\bigr)
+\Delta y\frac{\partial}{\partial y}
\bigl(f_x\Delta x+f_y\Delta y\bigr)\\
&=\Delta x\bigl(f_{xx}\Delta x+f_{yx}\Delta y\bigr)
+\Delta y\bigl(f_{xy}\Delta x+f_{yy}\Delta y\bigr)\\
&=f_{xx}(\Delta x)^2+(f_{xy}+f_{yx})\Delta x\Delta y+f_{yy}(\Delta y)^2\\
&=f_{xx}(\Delta x)^2+2f_{xy}\Delta x\Delta y+f_{yy}(\Delta y)^2.
\end{align*}
最后一步用了混合偏导相等 \(f_{xy}=f_{yx}\)。这就解释了二次项中间为什么有系数 \(2\)。

三次项同理，不要跳步地写就是
\begin{align*}
D^3f
&=D(D^2f)\\
&=D\bigl(f_{xx}(\Delta x)^2+2f_{xy}\Delta x\Delta y
+f_{yy}(\Delta y)^2\bigr)\\
&=\Delta x\frac{\partial}{\partial x}
\bigl(f_{xx}(\Delta x)^2+2f_{xy}\Delta x\Delta y+f_{yy}(\Delta y)^2\bigr)\\
&\quad+\Delta y\frac{\partial}{\partial y}
\bigl(f_{xx}(\Delta x)^2+2f_{xy}\Delta x\Delta y+f_{yy}(\Delta y)^2\bigr)\\
&=\Delta x\bigl(f_{xxx}(\Delta x)^2+2f_{xxy}\Delta x\Delta y
+f_{xyy}(\Delta y)^2\bigr)\\
&\quad+\Delta y\bigl(f_{xxy}(\Delta x)^2+2f_{xyy}\Delta x\Delta y
+f_{yyy}(\Delta y)^2\bigr)\\
&=f_{xxx}(\Delta x)^3+3f_{xxy}(\Delta x)^2\Delta y
+3f_{xyy}\Delta x(\Delta y)^2+f_{yyy}(\Delta y)^3.
\end{align*}
这里同样用了三阶混合偏导可以交换次序，例如 \(f_{xxy}=f_{xyx}=f_{yxx}\)，\(f_{xyy}=f_{yxy}=f_{yyx}\)。这就解释了三次项中的系数 \(1,3,3,1\)。

所以 Taylor 公式可以简写成
\[
f(x_0+\Delta x,y_0+\Delta y)
=f(x_0,y_0)+Df+\frac1{2!}D^2f+\frac1{3!}D^3f+\cdots+\mathcal R_n,
\]
其中这些 $Df,D^2f,D^3f,\dots$ 都是在 $(x_0,y_0)$ 处取值。

一般地，
\[
D^kf
=\sum_{j=0}^{k}\binom{k}{j}
\frac{\partial^k f}{\partial x^{\,k-j}\partial y^{\,j}}
(\Delta x)^{k-j}(\Delta y)^j.
\]
因此，前面 Taylor 公式里的第 $k$ 阶项就是
\[
\frac1{k!}D^kf.
\]
\end{remark}

\subsubsection{余项估计}

\begin{property}[余项的量级估计] \label{prop:remainder-bound}
在定理 \ref{thm:taylor} 的条件下，Lagrange 余项 $\mathcal R_n$ 满足
\[
\boxed{\;\mathcal R_n=O\bigl(\|\boldsymbol h\|^{n+1}\bigr),\qquad \|\boldsymbol h\|\to 0\;}
\]
更精确地，存在常数 $C>0$ 使得当 $\|\boldsymbol h\|$ 充分小时
\(\displaystyle |\mathcal R_n|\leqslant C\,\|\boldsymbol h\|^{n+1}\)。
\end{property}

余项估计告诉我们：当 $\|\boldsymbol h\|$ 缩小一个数量级时，Taylor 近似的误差会缩小大约 $\|\boldsymbol h\|$ 的 $(n+1)$ 次方量级。例如，若使用一阶近似（$n=1$），误差是 $O(\|\boldsymbol h\|^2)$ ——当步长缩小 10 倍时，误差大致缩小 $100$ 倍；若使用二阶近似（$n=2$），误差是 $O(\|\boldsymbol h\|^3)$ ——同样的 10 倍步长缩减会使误差缩小约 $1000$ 倍。

\subsection{无约束极值、闭区域最值与最小二乘}

\subsubsection{极值的概念与必要条件}

\begin{definition}[极值 / Extreme Value] \label{def:extremum}
设二元函数 $f(x,y)$ 在 $p_0=(x_0,y_0)$ 的某邻域内有定义。若存在 $\delta>0$，使得对任意 $p=(x,y)\in B_\delta(p_0)\setminus\{p_0\}$：
\begin{itemize}
  \item 若恒有 $f(p)<f(p_0)$，则称 $f$ 在 $p_0$ 处取 \textbf{严格极大值}
        （Strict Local Maximum），$p_0$ 为\textbf{严格极大值点}；
  \item 若恒有 $f(p)>f(p_0)$，则称 $f$ 在 $p_0$ 处取 \textbf{严格极小值}
        （Strict Local Minimum），$p_0$ 为\textbf{严格极小值点}；
  \item 若将上述“$<$”/“$>$”替换为“$\leqslant$”/“$\geqslant$”，
        则分别称为（非严格的）\textbf{极大值}和\textbf{极小值}，
        对应的 $p_0$ 称为\textbf{极大值点}/\textbf{极小值点}。
\end{itemize}
\end{definition}
\begin{definition}[临界点与驻点] \label{def:critical-stationary-saddle}
对二元函数 $f(x,y)$：
\begin{itemize}
  \item 使 $\nabla f \neq \boldsymbol 0$ 不成立（即梯度为零或梯度不存在）的点，
        统称为 $f$ 的 \textbf{临界点}（Critical Point）；

  \item 其中满足 $\nabla f=\boldsymbol 0$ 的点特称为 \textbf{驻点}
        （Stationary Point / Critical Point in narrow sense）；

  \item 若 $p_0$ 是驻点但不是极值点，则称其为 \textbf{鞍点}
        （Saddle Point） —— 因为其附近的曲面形状像马鞍。
\end{itemize}
\end{definition}
\begin{theorem}[极值的必要条件] \label{thm:necessary-condition-extreme}
设二元函数 $f(x,y)$ 在 $p_0=(x_0,y_0)$ 处取极值（无论严格与否），
且 Jacobian 矩阵 $Jf(p_0)$ 存在（即两个一阶偏导数均存在）。
则
\[
\boxed{\;Jf(p_0)=
\bigl(f_x(p_0),\,f_y(p_0)\bigr)=\boldsymbol 0\;}
\quad\Longleftrightarrow\quad
\nabla f(p_0)=\boldsymbol 0.
\]换句话说，可微函数的极值点 \textbf{一定} 是梯度为零的点。但逆命题不成立！梯度为零只是“候选条件”，并非充分条件。偏导数 \textbf{不存在} 的点也可能是极值点！例如 $f(x,y)=\sqrt{x^2+y^2}$（圆锥面）在 $(0,0)$ 处两个偏导数均不存在（沿各方向的方向导数为无穷），但 $(0,0)$ 显然是该函数的全局最小值点。这提醒我们搜索极值时不能遗漏偏导不存在的点。
\end{theorem}

\begin{exercise}[求驻点——二次函数] \label{ex:extreme-basic}
求函数
$f(x,y)=x^2-2x+y^2-4y+5$
的所有驻点。
\end{exercise}
\begin{solution}
\begin{enumerate}
  \item \textbf{分析可导性：}
  该函数为多项式，在全平面 $\mathbb R^2$ 上无穷次可微。

  \item \textbf{计算一阶偏导数并令其为零：}
  \begin{itemize}
    \item $f_x=2x-2=0 \implies x=1$
    \item $f_y=2y-4=0 \implies y=2$
  \end{itemize}

  \item \textbf{得出结论：}
  故 $f$ 在 $\mathbb R^2$ 上有唯一的驻点 $(1,2)$。
  （配方可知 $f(x,y)=(x-1)^2+(y-2)^2$，这是以 $(1,2)$ 为顶点开口向上的旋转抛物面，故 $(1,2)$ 确实是严格全局极小值点）。
\end{enumerate}
\end{solution}

\begin{exercise}[鞍点的典型例子] \label{ex:saddle-point-basic}
讨论函数 $f(x,y)=y^2-x^2$ 的临界点性质。
\end{exercise}
\begin{solution}
\begin{enumerate}
  \item \textbf{计算驻点：}
  计算偏导数 $f_x=-2x$, $f_y=2y$。
  解方程组 $f_x=0, f_y=0$，得到唯一驻点 $(0,0)$。

  \item \textbf{考察驻点邻域内的函数值变化：}
  \begin{itemize}
    \item 沿 $x$ 轴方向（$y=0$）：$f(x,0)=-x^2<0$（当 $x \neq 0$ 时）；
    \item 沿 $y$ 轴方向（$x=0$）：$f(0,y)=y^2>0$（当 $y \neq 0$ 时）。
  \end{itemize}

  \item \textbf{得出结论：}
  在原点的任意小的邻域内，既存在使 $f<0$ 的点（如 $(\varepsilon,0)$），也存在使 $f>0$ 的点（如 $(0,\varepsilon)$）。
  因此 $f(0,0)=0$ 既不是极大值也不是极小值。$(0,0)$ 是一个典型的 \textbf{鞍点}（双曲抛物面）。
\end{enumerate}
\end{solution}


\subsubsection{极值的充分条件：Hesse 矩阵判别法}

上一小节给出了“极值点 $\implies$ 驻点”的必要条件。
现在要回答反向的问题：\textbf{已知一个点是驻点，如何判断它是否为极值点？}
答案是利用 Taylor 展开的二阶项 —— 即 Hesse 矩阵的定性分析。

\begin{theorem}[正定负定判别法] \label{thm:sylvester-criterion}
设 $A$ 是 $n$ 阶实对称矩阵，记其前 $k$ 阶顺序主子式为
\[
\Delta_k=
\det
\begin{pmatrix}
a_{11} & \cdots & a_{1k}\\
\vdots & \ddots & \vdots\\
a_{k1} & \cdots & a_{kk}
\end{pmatrix},
\qquad k=1,2,\dots,n.
\]
则：
\begin{compactenum}
  \item $A$ 正定当且仅当 $\Delta_k>0\ (k=1,2,\dots,n)$；
  \item $A$ 负定当且仅当 $(-1)^k\Delta_k>0\ (k=1,2,\dots,n)$。
\end{compactenum}
对于二元函数的 Hesse 矩阵，上述判据可化为后文所用的 $f_{xx}(p_0)$ 与 $\Delta=f_{xx}(p_0)f_{yy}(p_0)-f_{xy}(p_0)^2$ 的分类准则。
\end{theorem}

\begin{theorem}[极值的充分条件：Hesse 矩阵判别法] \label{thm:sufficient-condition-hessian}
设 $f$ 在 $p_0$ 的某邻域内具有连续的二阶偏导数，且 $\nabla f(p_0)=\boldsymbol 0$（即 $p_0$ 为驻点）。
若 $H_f(p_0)$ \textbf{正定}，则 $f$ 在 $p_0$ 处取\textbf{严格极小值}；若 $H_f(p_0)$ \textbf{负定}，则 $f$ 在 $p_0$ 处取\textbf{严格极大值}；若 $H_f(p_0)$ \textbf{不定}，则 $f$ 在 $p_0$ 处\textbf{没有极值}（$p_0$ 为鞍点）。
\end{theorem}

对于二元函数的实际计算，我们通常不需要去算特征值，而是使用基于顺序主子式的等价判据：

\begin{property}[Hesse 矩阵的实用判别法] \label{property:discriminant-four-cases}
记 $a:=f_{xx}(p_0)$, $b:=f_{xy}(p_0)=f_{yx}(p_0)$, $c:=f_{yy}(p_0)$，
以及判别式 $\Delta:=ac-b^2$。则：

\renewcommand{\arraystretch}{1.6}
\par\noindent{\centering
\begin{tabular}{@{}clll@{}}
\hline
情形 & 条件 & $H_f(p_0)$ 性质 & 结论 \\
\hline
(i)   & $a>0$ 且 $\Delta>0$ & 正定 & \textbf{严格极小值点} \\
(ii)  & $a<0$ 且 $\Delta>0$ & 负定 & \textbf{严格极大值点} \\
(iii) & $\Delta<0$           & 不定 & \textbf{鞍点（非极值）} \\
(iv)  & $\Delta=0$           & 退化 & \textbf{无法判断} \\
\hline
\end{tabular}
\par}
\end{property}

\begin{note}[关于情形 (iv) 的说明]
当 $\Delta=ac-b^2=0$ 时，二阶信息不足以确定极值性 ——
需要借助更高阶的 Taylor 项或其他方法。
此时三种结果均有可能出现（见下文例~\ref{ex:degenerate-delta-zero}），
这也是为什么这个判别法被称为“充分的”而非“充分必要的”。
\end{note}

\begin{exercise}[退化的 Hesse 判别：$\Delta=0$ 的三种可能性] \label{ex:degenerate-delta-zero}
讨论函数 \(\displaystyle f_1(x,y)=x^2y^2\)、\(\displaystyle f_2(x,y)=-x^2y^2\)、\(\displaystyle f_3(x,y)=x^2y^3\) 在原点处的极值情况。
\end{exercise}
\begin{solution}
三个函数在 $(0,0)$ 处均有 $f_x=f_y=0$（直接验证即可），且其二阶偏导数均为零：$a=b=c=0$，从而 $\Delta=0$ —— Hesse 判别法 \textbf{失效}。需要逐一进行具体分析。

\begin{enumerate}
  \item \textbf{分析函数 $f_1(x,y)=x^2y^2$}：
  显然 $f_1(x,y)=(xy)^2\geqslant 0$ 对所有 $(x,y)$ 成立，且 $f_1(x,y)=0$ 当且仅当 $x=0$ 或 $y=0$。
  因此在 $(0,0)$ 的任何邻域内，$f_1(x,y)\geqslant 0=f_1(0,0)$。
  $(0,0)$ 是 $f_1$ 的 \textbf{（非严格）极小值点}（实际上也是全局最小值点）。

  \item \textbf{分析函数 $f_2(x,y)=-x^2y^2$}：
  与 $f_1$ 相反，$f_2(x,y)=-(xy)^2\leqslant 0$ 对所有 $(x,y)$ 成立，且 $f_2(0,0)=0$ 为最大值。
  $(0,0)$ 是 $f_2$ 的 \textbf{（非严格）极大值点}。

  \item \textbf{分析函数 $f_3(x,y)=x^2y^3$}：
  此函数在原点附近变号！
  \begin{itemize}
    \item 当 $y>0$ 且 $x \neq 0$ 时：$f_3=x^2y^3>0>f_3(0,0)=0$；
    \item 当 $y<0$ 且 $x \neq 0$ 时：$f_3=x^2y^3<0<f_3(0,0)=0$。
  \end{itemize}
  因此在原点的任何邻域内既有大于 $f(0,0)$ 的点也有小于 $f(0,0)$ 的点。
  $(0,0)$ \textbf{不是极值点}（属于广义的鞍点）。
\end{enumerate}
总结：三者均为 $\Delta=0$ 的退化情形，却分别对应极小值、极大值和无极值三种截然不同的结果。这充分说明了高阶分析的必要性。
\end{solution}

\begin{exercise}[含三次项的多项式函数] \label{ex:cubic-polynomial-extreme}
求函数
$f(x,y)=\dfrac{1}{2} x^2+3y^3+9y^2-3xy+9y-9x$
的所有驻点，并判断其类型。
\end{exercise}
\begin{solution}
\begin{enumerate}
  \item \textbf{求驻点：}
  计算一阶偏导数得 $f_x = x-3y-9,\ f_y = 9y^2+18y-3x+9$。令 $f_x=f_y=0$，即 $x-3y-9=0,\ 9y^2+18y-3x+9=0$。由第一式 $x=3y+9$ 代入第二式，得 $9y^2+9y-18 = 0$，即 $(y+2)(y-1) = 0$。因此 $y=-2$ 或 $y=1$，对应驻点分别为 $p_1=(3,-2)$ 与 $p_2=(12,1)$。

  \item \textbf{计算 Hesse 矩阵元素：}
  二阶偏导数为 $f_{xx}=1,\ f_{yy}=18y+18,\ f_{xy}=-3$。

  \item \textbf{逐个判别：}
  在 $p_1=(3,-2)$ 处，$a=1,\ b=-3,\ c=-18$，故 $\Delta=ac-b^2=-27<0$，所以 $p_1$ 为 \textbf{鞍点}；在 $p_2=(12,1)$ 处，$a=1,\ b=-3,\ c=36$，故 $\Delta=27>0$ 且 $a>0$，所以 $p_2$ 为 \textbf{严格极小值点}。
\end{enumerate}
\end{solution}

\begin{exercise}[超越函数的无穷多极值点] \label{ex:trig-exp-extreme}
求函数 $z=(1+\mathrm{e}^{y})\cos x-y\mathrm{e}^{y}$ 的全部极值点。
\end{exercise}
\begin{solution}
这是一个周期性超越函数与指数函数的复合，预期可能有无穷多个极值点。

\begin{enumerate}
  \item \textbf{求驻点方程组：}
  \(\displaystyle z_x=-(1+\mathrm{e}^{y})\sin x = 0,\qquad
  z_y=\mathrm{e}^{y}\cos x - \mathrm{e}^{y} - y\mathrm{e}^{y} = \mathrm{e}^{y}(\cos x - 1 - y) = 0\)。
  由 $z_x=0$ 且 $1+\mathrm{e}^{y}>0$（恒成立），得 $\sin x=0$，故 $x=k\pi$（$k\in\mathbb Z$）。
  由 $z_y=0$ 且 $\mathrm{e}^{y}>0$，得 $\cos x-1-y=0$，即 $y=\cos x-1$。
  若 $k$ 为偶数，即 $x=2k\pi$，则 $\cos x = 1$，故 $y=0$，驻点为 $p_k^{(+)}=(2k\pi,\;0)$；若 $k$ 为奇数，即 $x=(2k+1)\pi$，则 $\cos x = -1$，故 $y=-2$，驻点为 $q_k^{(-)}=((2k+1)\pi,\;-2)$。

  \item \textbf{Hesse 判别：}
  先写出通用的二阶偏导数：
  \[
  z_{xx}=-(1+\mathrm{e}^{y})\cos x,\qquad
  z_{xy}=-\mathrm{e}^{y}\sin x,\qquad
  z_{yy}=\mathrm{e}^{y}(\cos x-2-y).
  \]
  对 $p_k^{(+)}=(2k\pi,0)$，有 $a=-2<0,\ b=0,\ c=-1$，故 $\Delta=2>0$，所以每个 $p_k^{(+)}$ 都是 \textbf{严格极大值点}，且极大值为 $z(2k\pi,0)=2$；对 $q_k^{(-)}=((2k+1)\pi,-2)$，有 \(\displaystyle a=1+\mathrm{e}^{-2}>0,\quad b=0,\quad c=-\mathrm{e}^{-2}<0\)，故 $\Delta=-\mathrm{e}^{-2}(1+\mathrm{e}^{-2})<0$，所以每个 $q_k^{(-)}$ 都是 \textbf{鞍点}。

  \item \textbf{结论：}
  函数有无穷多个严格极大值点 $(2k\pi,0)$（极大值恒为 $2$），无穷多个鞍点 $((2k+1)\pi,-2)$，但 \textbf{没有任何极小值点}。
\end{enumerate}
\end{solution}

\begin{exercise}["猴鞍" 点（Monkey Saddle）] \label{ex:monkey-saddle}
求函数 $z=(y-x^2)(y-2x^2)$ 的极值点。
\end{exercise}
\begin{solution}
展开原式以便于求导：$z=y^2-3x^2y+2x^4$。

\begin{enumerate}
  \item \textbf{求驻点：}
  $z_x = -6xy+8x^3=2x(4x^2-3y),\ z_y = 2y-3x^2$。令 $z_x=z_y=0$，由 $z_y=0$ 得 $y=\dfrac{3}{2}x^2$，代入 $z_x$ 中得 $2x\left(4x^2-3\cdot\dfrac{3}{2} x^2\right) = -x^3=0$，故 $x=0$，从而 $y=0$。
  唯一驻点为 $(0,0)$。

  \item \textbf{Hesse 判别：}
  二阶偏导数为 $z_{xx}=-6y+24x^2,\ z_{xy}=-6x,\ z_{yy}=2$。
  在 $(0,0)$ 处：$a=0$, $b=0$, $c=2$，故 $\Delta=0\cdot 2-0=0$。\textbf{判别法退化，无法判断！}

  \item \textbf{直接分析法：}
  将 $z$ 重新因式分解：$z=(y-x^2)(y-2x^2)$。观察原点邻域内 $z$ 的符号分布：
  在区域 I（两抛物线之间）$x^2<y<2x^2$（$x \neq 0$）时，两因子一负一正，故 $z<0$；在区域 II（上方）$y>2x^2$ 时，两因子皆正，故 $z>0$；在区域 III（下方）$y<x^2$ 时，两因子皆负，故 $z>0$。
  因此在原点的任意邻域内同时存在 $z>0$ 和 $z<0$ 的点。$(0,0)$ \textbf{不是极值点}。
\end{enumerate}
\begin{note}
  \textbf{几何意义}：这个曲面在原点处形成所谓的 \textbf{猴鞍}（Monkey Saddle）——比普通鞍点“多一条谷”：普通鞍点（如 $y^2-x^2$）有两条“下降通道”（沿 $x$ 轴方向），而猴鞍有 \textbf{三条}（沿三个间隔 $120^\circ$ 的方向）。
\end{note}
\end{solution}

\begin{exercise}[提高：Hesse 矩阵与极值判定]
设 $f(x,y,z)=x^3+y^3+z^3-3xyz-6$。
\begin{shortexenum}
  \shortitem{1}{求 $f$ 在 $(1,1,1)$ 处的梯度和 Hesse 矩阵；} &
  \shortitem{2}{写出 $f$ 在 $(1,1,1)$ 处的二阶 Taylor 展开式；} \\
  \shortitem{3}{判断 $f$ 在 $(1,1,1)$ 处是否有极值？若有，是极大还是极小？} &
\end{shortexenum}
\end{exercise}
\begin{solution}
\begin{enumerate}
  \item[(1)] \textbf{求梯度与 Hesse 矩阵：}
  一阶偏导数（并代入$x=y=z=1$）：
  $f_x = 3x^2 - 3yz,\ f_y = 3y^2 - 3xz,\ f_z = 3z^2 - 3xy$，故 $f_i(1,1,1)=0\ (i=x,y,z)$。
  即梯度为零向量 $\boldsymbol{0}$（此点为驻点）。
  二阶偏导数在该点满足 $f_{xx}=f_{yy}=f_{zz}=6,\ f_{xy}=f_{xz}=f_{yz}=-3$。
  Hesse 矩阵为：
  \[
  H_f(1,1,1) = \begin{pmatrix} 6 & -3 & -3 \\ -3 & 6 & -3 \\ -3 & -3 & 6 \end{pmatrix}.
  \]

  \item[(2)] \textbf{写出二阶 Taylor 展开式：}
  在 $(1,1,1)$ 处，$f(1,1,1) = 1+1+1-3-6 = -6$。
  由多变量 Taylor 公式：
  \begin{align*}
  f(1+h_1, 1+h_2, 1+h_3) &= -6 + \frac{1}{2} \begin{pmatrix} h_1 & h_2 & h_3 \end{pmatrix} \begin{pmatrix} 6 & -3 & -3 \\ -3 & 6 & -3 \\ -3 & -3 & 6 \end{pmatrix} \begin{pmatrix} h_1 \\ h_2 \\ h_3 \end{pmatrix} + \mathcal{R}_2 \\
  &= -6 + 3(h_1^2+h_2^2+h_3^2) - 3(h_1h_2+h_2h_3+h_3h_1) + \mathcal{R}_2.
  \end{align*}

  \item[(3)] \textbf{极值判定：}
  考察二阶项的二次型：
  \[
  3(h_1^2+h_2^2+h_3^2)-3(h_1h_2+h_2h_3+h_3h_1)
  =\frac{3}{2}\left[(h_1-h_2)^2+(h_2-h_3)^2+(h_3-h_1)^2\right]\geqslant 0.
  \]
  这说明该 Hesse 矩阵是\textbf{半正定}的（特征值为 0, 9, 9），所以无法通过常规的严格正定性直接判别为严格极小值点。
  但是，我们可以通过代数恒等式：
  \[
  f(x,y,z)-(-6)=x^3+y^3+z^3-3xyz
  =\frac{1}{2}(x+y+z)\left[(x-y)^2+(y-z)^2+(z-x)^2\right].
  \]
  在 $(1,1,1)$ 的充分小邻域内，$x+y+z \approx 3 > 0$，故对于邻域内任意点总有 $f(x,y,z) \geqslant -6$ 恒成立。
  因此，$f$ 在 $(1,1,1)$ 处\textbf{有局部极小值}，但不具备严格局部极小值（因为沿 $x=y=z$ 路径函数值常为 -6）。
\end{enumerate}
\end{solution}

\subsubsection{最值的求法}
在实际应用中，我们往往更关心函数在某个区域上的\textbf{全局最大值}和 \textbf{全局最小值}。

\begin{theorem}[闭区域上连续函数最值的求法] \label{theorem:global-extrema-theorem}
设二元函数 $z=f(x,y)$ 在 \textbf{有界闭区域} $D$ 上连续。
由多元连续函数的最值定理（Weierstrass 定理的推广），$f$ 在 $D$ 上必定能取到最大值和最小值。
最值只可能出现在 $D$ 的内部临界点或边界上。因而求全局最值时，应先找出内部驻点和偏导数不存在的临界点，并计算这些点处的函数值；再把边界 $\partial D$ 作为曲线段处理，通常借助一元函数方法或参数化求其最大值和最小值；最后把内部和边界得到的所有候选值统一比较，其中最大者就是 $f$ 在 $D$ 上的最大值，最小者就是最小值。
\end{theorem}

\begin{exercise}[多边形域上线性函数的最值] \label{ex:linear-polygon-minimum}
求函数 $u(x,y)=x+y$ 在区域 \(\displaystyle D=\{(x,y):x\geqslant 0,\;y\geqslant 0,\;2x+y\geqslant 3,\;x+3y\geqslant 4\}\) 上的最小值。
\end{exercise}
\begin{solution}
\begin{enumerate}
  \item \textbf{分析区域 $D$：}
  $D$ 是第一象限内由四条直线围成的无界凸多边形（凸集的交集）。四条边界分别为：
  $L_1:x=0$, $L_2:y=0$, $L_3:2x+y=3$, $L_4:x+3y=4$。
  \item \textbf{求内部临界点：}
  $u_x=1 \neq 0$, $u_y=1 \neq 0$ —— $u$ 在全平面上无驻点。因此最值只能在 \textbf{边界}上取得。
  \item \textbf{分析边界上的最值：}
  逐一考察四条边界的贡献：
  \begin{itemize}
    \item \textbf{边界 $L_1$}（$x=0,\,y\geqslant 3$）：$u(0,y)=y\geqslant 3$，最小值为 $3$（在 $(0,3)$ 处取到）。
    \item \textbf{边界 $L_2$}（$y=0,\,x\geqslant 4$）：$u(x,0)=x\geqslant 4$，最小值为 $4$（在 $(4,0)$ 处取到）。
    \item \textbf{边界 $L_3$}（$2x+y=3,\,0\leqslant x\leqslant 1$）：将约束代入目标函数：$u=x+(3-2x)=3-x$。当 $x\in[0,1]$ 时，$u$ 从 $3$ 递减到 $2$。最小值为 $2$（在端点 $(1,1)$ 处取到）。
    \item \textbf{边界 $L_4$}（$x+3y=4,\,0\leqslant y\leqslant 1$）：代入目标函数：$u=(4-3y)+y=4-2y$。当 $y\in[0,1]$ 时，$u$ 从 $4$ 递减到 $2$。最小值为 $2$（在端点 $(1,1)$ 处取到）。
  \end{itemize}
  \item \textbf{得出结论：}
  所有候选极小值：$3$, $4$, $2$, $2$。
  最小者为 $\boxed{2}$，在点 $(1,1)$ 处取到（它是 $L_3$ 与 $L_4$ 的交点）。
\end{enumerate}
\end{solution}

\begin{exercise}[思考题：三角函数在多边形域上的最值]
在闭区域 $D=\{(x,y):x\geqslant 0,\;y\geqslant 0,\;x+y\leqslant\pi\}$ 上求函数 $f(x,y)=\sin x\,\sin y\,\sin(x+y)$ 的最大值和最小值。
\end{exercise}
\begin{solution}
\begin{enumerate}
  \item \textbf{分析边界点：}
  闭区域 $D$ 是以 $(0,0), (\pi,0), (0,\pi)$ 为顶点的直角三角形。
  在边界 $x=0$ 上有 $f(0,y)=0$，在边界 $y=0$ 上有 $f(x,0)=0$，在边界 $x+y=\pi$ 上有 $f(x,\pi-x)=0$。
  因此，$f$ 在所有边界上的值均为 0。
  \item \textbf{分析区域内部：}
  在区域内部，有 $0 < x < \pi$，$0 < y < \pi$ 且 $0 < x+y < \pi$。
  此时 $\sin x > 0$，$\sin y > 0$，$\sin(x+y) > 0$，因此 $f(x,y) > 0$。
  由此可知，$f$ 的最小值为边界上的 \textbf{0}，而最大值必定在区域内部的驻点取得。
  \item \textbf{求内部驻点：}
  计算偏导数并令其为 0（利用乘积法则）：
  \begin{align*}
  f_x &= \cos x \sin y \sin(x+y) + \sin x \sin y \cos(x+y) = \sin y \left[ \cos x \sin(x+y) + \sin x \cos(x+y) \right] \\
      &= \sin y \sin(2x+y) = 0 \\
  f_y &= \sin x \cos y \sin(x+y) + \sin x \sin y \cos(x+y) = \sin x \left[ \cos y \sin(x+y) + \sin y \cos(x+y) \right] \\
      &= \sin x \sin(x+2y) = 0
  \end{align*}
  因为在区域内部 $\sin x \neq 0$ 且 $\sin y \neq 0$，故必须有 $\sin(2x+y) = 0,\ \sin(x+2y) = 0$。
  由于 $x,y > 0$ 且 $x+y < \pi$，可知 $0 < 2x+y < 2\pi$ 且 $0 < x+2y < 2\pi$。
  在 $(0,2\pi)$ 内正弦值为 0 的角仅有 $\pi$。
  所以 $2x+y = \pi,\ x+2y = \pi$。
  解此方程组得到唯一内部驻点：$x = \dfrac{\pi}{3}$，$y = \dfrac{\pi}{3}$。
  \item \textbf{计算最大值：}
  将驻点坐标代入目标函数：
  $f\left(\dfrac{\pi}{3},\dfrac{\pi}{3}\right) = \sin\left(\dfrac{\pi}{3}\right)\sin\left(\dfrac{\pi}{3}\right)\sin\left(\dfrac{2\pi}{3}\right) = \dfrac{3\sqrt{3}}{8}$。
  故最大值为 $\boxed{\dfrac{3\sqrt{3}}{8}}$，最小值为 $\boxed{0}$。
\end{enumerate}
\end{solution}

\subsection{约束极值与拉格朗日乘数法}

先从一个具体例子感受一下“约束”的存在方式以及消元法的局限性。

\begin{exercise}[长方体的表面积最小化] \label{ex:box-surface-minimum}
要做一个体积为 $2$ 的长方体盒子，问做成怎样的尺寸能使表面积最小？
\end{exercise}
\begin{solution}
\textbf{先说选法：}这道题其实既可以用消元法，也可以用 Lagrange 乘数法。由于约束 $xyz=2$ 很容易直接解出一个变量，本题先用消元法更直观；这样也便于和后面真正需要乘数法的情形做对照。
\begin{enumerate}
  \item \textbf{建立数学模型：} 设盒子的长、宽、高分别为 $x>0$, $y>0$, $z>0$，则目标函数（需最小化）为表面积 $S(x,y,z)=2(xy+yz+zx)$，约束条件为 $xyz=2$。
   \item \textbf{使用消元法：} 由约束 $xyz=2$ 解出 $z=\dfrac{2}{xy}$，代回得 \(\displaystyle S=2\left(xy+\frac{2}{x}+\frac{2}{y}\right)\)。这样就把三元问题化成了定义域 $x>0,\ y>0$ 上的二元无约束极值。
   \item \textbf{求无约束驻点：} 令偏导为零，得到 \(\displaystyle y=\dfrac{2}{x^2}\) 与 \(\displaystyle x=\dfrac{2}{y^2}\)。把前式代入后式，得 \(\displaystyle x^3=2\)，所以 \(\displaystyle x=y=\sqrt[3]{2}\)，进而 \(\displaystyle z=\sqrt[3]{2}\)。
   \item \textbf{得出结论并回到原问题复核：} 候选点只有 \(\bigl(\sqrt[3]{2},\sqrt[3]{2},\sqrt[3]{2}\bigr)\)，代回表面积公式得 \(\displaystyle S_{\min}=6\cdot 2^{2/3}\)。因此体积固定时，表面积最小的长方体是立方体。
\end{enumerate}
这个例子也提示：若约束能轻松解出一个变量，消元法通常是最短路径；但若约束复杂、多重约束同时存在，或题目具有明显的对称结构，就更适合直接切换到 Lagrange 乘数法。
\end{solution}

\begin{definition}[Lagrangian / 辅助函数] \label{def:lagrangian}
考虑如下 \textbf{条件极值问题}：在约束条件 $g(x,y)=0$ 下求目标函数 $f(x,y)$ 的极值。
构造辅助函数 \(\displaystyle \boxed{\;L(x,y;\lambda):=f(x,y)+\lambda\,g(x,y)\;}\)，其中 $\lambda\in\mathbb R$ 为待定参数，称为 \textbf{Lagrange 乘子}（Lagrange Multiplier）；$L$ 称为 \textbf{Lagrangian 函数}（简称 \textbf{Lagrange 函数}）。
注意我们将 $\lambda$ 写在分号后面以示区别：$(x,y)$ 是原变量，而 $\lambda$ 是新引入的辅助变量。
\end{definition}

\begin{theorem}[Lagrange 乘数法（单约束 / 二元情形）] \label{thm:lagrange-single}
设目标函数 $f:\mathbb R^2\to\mathbb R$ 和约束函数 $g:\mathbb R^2\to\mathbb R$ 均在 $p_0=(x_0,y_0)$ 的某邻域内具有一阶连续偏导数。若 $p_0$ 是 $f$ 在约束 $g(x,y)=0$ 下的 \textbf{条件极值点}，且 $\nabla g(p_0) \neq \boldsymbol{0}$（正则性条件），

则存在唯一的实数 $\lambda_0\in\mathbb R$，使得 $(x_0,y_0)$ 是 Lagrange 函数 $L(x,y;\lambda_0)$ 的一个 \textbf{驻点}。即：
\[
\boxed{
\begin{cases}
\dfrac{\partial L}{\partial x}(p_0) =f_x(p_0)+\lambda_0\,g_x(p_0)=0,\\[10pt]
\dfrac{\partial L}{\partial y}(p_0) =f_y(p_0)+\lambda_0\,g_y(p_0)=0,\\[10pt]
g(x_0,y_0)=0.\quad\text{（原始约束）}
\end{cases}}
\]
\end{theorem}

\begin{exercise}[圆周上的乘积极值] \label{ex:circle-product-extreme}
求函数 $f(x,y)=xy$ 在圆周 $x^2+y^2=4$ 上的最大值和最小值。
\end{exercise}
\begin{solution}
\textbf{思路：}圆周约束 $x^2+y^2=4$ 天然适合两种做法。若要体现一般性的约束极值方法，优先用 Lagrange 乘数法；若只求这一题的结果，参数化成三角函数会更快。

\textbf{方法一：Lagrange 乘数法}
\begin{enumerate}
  \item \textbf{构造 Lagrange 函数：}
  目标函数为 $f(x,y)=xy$，约束函数为 $g(x,y)=x^2+y^2-4=0$，故 $L(x,y;\lambda)=xy+\lambda(x^2+y^2-4)$。

  \item \textbf{求驻点方程组：}
  \[
  \begin{cases}
  L_x = y + 2\lambda x = 0, &\quad\text{(i)}\\
  L_y = x + 2\lambda y = 0, &\quad\text{(ii)}\\
  \quad g: \;x^2+y^2 = 4. &\quad\text{(iii)}
  \end{cases}
  \]

  \item \textbf{求解方程组寻候选点：}
  由 (i) 得 $y=-2\lambda x$。代入 (ii) 得
  \(\displaystyle x+2\lambda(-2\lambda x)=x(1-4\lambda^2)=0\)。
  若 $x=0$，由 (iii) 得 $y=\pm 2$，但此时 (ii) 给出 $\lambda=0$，回代 (i) 时 $\pm 2 = 0$ 矛盾，故 $x \neq 0$。于是 \(\displaystyle 1-4\lambda^2=0,\qquad \lambda=\pm\dfrac{1}{2}\)。再将 $y=-2\lambda x$ 代入 (iii) 得 $x^2+4\lambda^2x^2=2x^2=4$，故 $x=\pm\sqrt{2}$，从而得到四组有效点：$(\sqrt{2},-\sqrt{2})$, $(-\sqrt{2},\sqrt{2})$, $(\sqrt{2},\sqrt{2})$, $(-\sqrt{2},-\sqrt{2})$。

  \item \textbf{比较得出结论：}
  代入目标函数得 $f(\sqrt{2},\sqrt{2}) = f(-\sqrt{2},-\sqrt{2}) = 2$，$f(\sqrt{2},-\sqrt{2}) = f(-\sqrt{2},\sqrt{2}) = -2$，故最大值为 $2$，最小值为 $-2$。
\end{enumerate}

\textbf{方法二：代数参数化法}
\begin{enumerate}
  \item \textbf{参数化转化：}约束本身就是半径为 $2$ 的圆周，因此最自然的参数化是 \(\displaystyle x=2\cos t,\ y=2\sin t\)。
  它的优点是：约束条件被自动满足，问题直接降为一元最值。
  \item \textbf{化为一元函数：}\(\displaystyle f=xy=4\sin t\cos t=2\sin 2t\)。
  \item \textbf{求解一元最值：}由于正弦函数的值域是 $[-1,1]$，故 \(\displaystyle -2\le 2\sin 2t\le 2,\ f_{\max}=2,\ f_{\min}=-2\)。
  当 $\sin 2t=1$ 时取到最大值，对应点为 $(\sqrt2,\sqrt2)$、$(-\sqrt2,-\sqrt2)$；当 $\sin 2t=-1$ 时取到最小值，对应点为 $(\sqrt2,-\sqrt2)$、$(-\sqrt2,\sqrt2)$。
\end{enumerate}
这一题说明，Lagrange 法更适合推广到一般曲线、曲面约束；参数化法则适合“约束本身就有天然参数”的情形。约束能否被自然参数化，常常决定哪一种解法更简洁。
\end{solution}

\subsubsection{Lagrange 乘数法的推广}

\begin{proposition}[三元单约束的 Lagrange 方法] \label{prop:lagrange-3d-one-constraint}
设 $f,g:\mathbb R^3\to\mathbb R$ 具有一阶连续偏导数，$\nabla g(p_0) \neq \boldsymbol{0}$。
若 $p_0=(x_0,y_0,z_0)$ 是 $f$ 在约束 $g(x,y,z)=0$ 下的条件极值点，则存在 $\lambda\in\mathbb R$ 使得 $p_0$ 是以下 Lagrange 函数的驻点：
\(\displaystyle L(x,y,z;\lambda)=f(x,y,z)+\lambda\,g(x,y,z)\)，即满足
\[
\begin{cases}
f_x(p_0)+\lambda g_x(p_0)=0,\\
f_y(p_0)+\lambda g_y(p_0)=0,\\
f_z(p_0)+\lambda g_z(p_0)=0,\\
g(x_0,y_0,z_0)=0.
\end{cases}
\]
\end{proposition}

\begin{proposition}[三元双约束的 Lagrange 方法] \label{prop:lagrange-3d-two-constraints}
设 $f,g_1,g_2:\mathbb R^3\to\mathbb R$ 具有一阶连续偏导数，$\nabla g_1(p_0)$ 与 $\nabla g_2(p_0)$ \textbf{线性无关}（即两者不平行）。
若 $p_0=(x_0,y_0,z_0)$ 是 $f$ 在联合约束
$\begin{cases}g_1(x,y,z)=0\\g_2(x,y,z)=0\end{cases}$
下的条件极值点，则存在 $\lambda,\mu\in\mathbb R$ 使得
\(\displaystyle L(x,y,z;\lambda,\mu):=f(x,y,z)+\lambda g_1(x,y,z)+\mu g_2(x,y,z)\) 以 $(x_0,y_0,z_0)$ 为驻点，即：
\[
\begin{cases}
f_x(p_0)+\lambda g_{1x}(p_0)+\mu g_{2x}(p_0)=0,\\
f_y(p_0)+\lambda g_{1y}(p_0)+\mu g_{2y}(p_0)=0,\\
f_z(p_0)+\lambda g_{1z}(p_0)+\mu g_{2z}(p_0)=0,\\
g_1(p_0)=0,\;g_2(p_0)=0.
\end{cases}
\]
\end{proposition}

\begin{exercise}[长方体表面积的 Lagrange 解法] \label{ex:box-lagrangian}
在 $x,y,z>0$ 且 $xyz=2$ 的条件下，求 $S(x,y,z)=2(xy+yz+zx)$ 的最小值点。
\end{exercise}
\begin{solution}
\begin{enumerate}
  \item \textbf{构造 Lagrange 函数：} 目标是最小化 $S(x,y,z)=2(xy+yz+zx)$，约束是 $g(x,y,z)=xyz-2=0$，取
  \(\displaystyle L(x,y,z;\lambda)=2(xy+yz+zx)+\lambda(xyz-2)\)。
  \item \textbf{求驻点方程组：} 求偏导并令为零，得 $L_x=2(y+z)+\lambda yz=0$，$L_y=2(x+z)+\lambda xz=0$，$L_z=2(x+y)+\lambda xy=0$。
  \item \textbf{解方程组：}
  由 (i) 得 $\lambda=-\dfrac{2(y+z)}{yz}$，由 (ii) 得 $\lambda=-\dfrac{2(x+z)}{xz}$，由 (iii) 得 $\lambda=-\dfrac{2(x+y)}{xy}$。
  三式右边由于同等对称，建立等式有 \(\displaystyle \frac{y+z}{yz}=\frac{x+z}{xz}=\frac{x+y}{xy}\)。
  取第一等号化简：$\dfrac{1}{z}+\dfrac{1}{y}=\dfrac{1}{z}+\dfrac{1}{x} \implies x=y$。
  同理第二等号给出 $y=z$。因此 $x=y=z$。
  代入约束 $xyz=2$：$x^3=2$，得 $x=y=z=\sqrt[3]{2}$。
  \item \textbf{得出结论：} 唯一候选点为 $(\sqrt[3]{2},\sqrt[3]{2},\sqrt[3]{2})$（立方体）。由问题的物理意义（边长趋于零导致面积在开边界趋于无穷大）可知此必然为全局最小值点。
\end{enumerate}
\end{solution}

\begin{remark}[与算术--几何平均不等式的对照]
本题其实可以直接用 \textbf{AM--GM 不等式}解决：
 \(\displaystyle \frac{S}{2}=xy+yz+zx\geqslant 3\sqrt[3]{(xy)(yz)(zx)} =3\sqrt[3]{(xyz)^2}=3\sqrt[3]{4}\)，
等号成立当且仅当 $xy=yz=zx$，即 $x=y=z=\sqrt[3]{2}$。
因此 $S_{\min}=6\sqrt[3]{4}$，与 Lagrange 法结果完全一致。
这揭示了一个深刻的联系：\textbf{许多经典不等式本质上都是某个优化问题的最优性条件}。Lagrange 乘数法提供了一种“从零推导”这些不等式的统一框架。
\end{remark}

\begin{exercise}[算术--几何平均不等式的最优化证明] \label{ex:am-gm-lagrangian}
将正数 $a>0$ 分成 $n$ 个正数之和，使这 $n$ 个数的乘积最大。
\end{exercise}
\begin{solution}
\begin{enumerate}
  \item \textbf{数学建模：}
  决策变量为 $x_1,x_2,\dots,x_n>0$，约束条件为 $x_1+x_2+\cdots+x_n=a$，目标函数为 $P(x_1,\dots,x_n)=x_1x_2\cdots x_n$。
  \item \textbf{构造 Lagrange 函数：}
  \(\displaystyle L(x_1,\dots,x_n;\lambda)=x_1x_2\cdots x_n+\lambda(x_1+x_2+\cdots+x_n-a)\)。
  \item \textbf{求驻点方程组：}
  对第 $k$ 个变量 $x_k$ 求偏导并令其为 0：
  \(\displaystyle \frac{\partial L}{\partial x_k}=\frac{x_1x_2\cdots x_n}{x_k}+\lambda=0\implies \frac{P}{x_k}=-\lambda\)。
    注意到右端 $-\lambda$ 与指标 $k$ 无关，这意味着对所有 $k$，$\dfrac{P}{x_k}$ 必须取相同的值。
  由于约束要求正数乘积 $P>0$，两边同除以 $P$ 得 $x_1=x_2=\cdots=x_n$。
  \item \textbf{得出结论：}
  结合约束 $\sum x_i=a$，得到唯一的内部驻点：$x_1=x_2=\cdots=x_n=\dfrac{a}{n}$。
  此时的乘积值为 $P_{\max}=\left(\dfrac{a}{n}\right)^n$。
  当且仅当 $a$ 被均匀分割时，乘积达到最大值。这正是著名的算术--几何平均不等式的实质内涵。
\end{enumerate}
\end{solution}

\begin{note}[不等式与优化的深刻联系]
这个例子展示了优化理论的一个“反向应用”：我们不是用已知的不等式来解题，而是通过求解优化问题来 \textbf{证明}不等式。事实上，大量重要的不等式（Cauchy--Schwarz、Hölder、Minkowski 等）都可以从相应的凸优化问题的 KKT 条件推导出来。
\end{note}

\subsection{偏微分方程专题：验证、变量替换与坐标变换}

\subsubsection{验证函数满足给定的偏微分方程}

验证 ``某函数是否满足某 PDE'' 的思路非常直接：将函数代入方程的左端，通过计算偏导数验证等式是否成立。虽然原理简单，但计算过程中需要仔细处理复合函数的链式法则、对称性简化以及隐函数求导等技术细节。

\begin{definition}[偏微分方程的基本概念] \label{def:pde-intro}
设 $u=u(x_1,\dots,x_n)$ 是一个未知的多元函数。含有 $u$ 及其若干偏导数 $\dfrac{\partial^{\alpha}u}{\partial x_1^{a_1}\cdots\partial x_n^{a_n}}$（其中 $|\alpha|=a_1+\cdots+a_n\geqslant 1$）的方程称为 \textbf{偏微分方程}（PDE）。

特别地，\(\displaystyle \Delta u:=\sum_{i=1}^{n}\frac{\partial^2 u}{\partial x_i^2}=0\) 称为 $n$ 维空间中的 \textbf{Laplace 方程}（或位势方程/调和方程），其中 $\displaystyle\Delta=\sum_{i=1}^n\frac{\partial^2}{\partial x_i^2}$ 称为 \textbf{Laplace 算子}（或记作 $\nabla^2$）。满足 $\Delta u=0$ 的函数称为 \textbf{调和函数}（Harmonic Function）。

对比地，若方程中未知函数仅依赖 \textbf{一个} 自变量且只出现常导数，则称为 \textbf{常微分方程}（ODE）。
\end{definition}

\begin{exercise}[Laplace 方程的基本解] \label{ex:laplace-fundamental}
证明函数 \(\displaystyle u(x,y,z)=\frac{1}{\sqrt{x^2+y^2+z^2}}\) 在其定义域 $\mathbb R^3\setminus\{(0,0,0)\}$ 上满足三维 Laplace 方程 $\Delta u=0$。
\end{exercise}

\begin{proof}
引入球坐标中的径向距离符号以简化书写：\(\displaystyle r:=\sqrt{x^2+y^2+z^2},\ \text{于是 }u=\frac{1}{r}\)。

\begin{enumerate}
  \item \textbf{计算一阶偏导数。}

  由链式法则，$\dfrac{\partial r}{\partial x}=\dfrac{x}{r}$，故
  \[
  u_x=\frac{\partial}{\partial x}\!\left(r^{-1}\right)
  =(-1)r^{-2}\cdot\frac{\partial r}{\partial x}
  =-\frac{1}{r^2}\cdot\frac{x}{r}
  =-\frac{x}{r^3}.
  \]
  利用 $(x,y,z)$ 的完全对称性，同理可得 $u_y=-\dfrac{y}{r^3}$，$u_z=-\dfrac{z}{r^3}$。

  \item \textbf{计算二阶偏导数。}

  对 $u_x$ 再次关于 $x$ 求导（使用商法则）：
  \begin{align*}
  u_{xx}&=\frac{\partial}{\partial x}\!\left(-\frac{x}{r^3}\right)
  =-\frac{1\cdot r^3-x\cdot\dfrac{\partial}{\partial x}(r^3)}{r^6}\\[6pt]
  &=-\frac{1}{r^3}+\frac{x\cdot 3r^2\cdot\dfrac{x}{r}}{r^6}
  =-\frac{1}{r^3}+\frac{3x^2}{r^5}
  =\frac{3x^2-r^2}{r^5}.
  \end{align*}
  同样由对称性可得 $u_{yy}=\dfrac{3y^2-r^2}{r^5}$，$u_{zz}=\dfrac{3z^2-r^2}{r^5}$。

  \item \textbf{验证 Laplace 方程。}
  \begin{align*}
  \Delta u&=u_{xx}+u_{yy}+u_{zz}=\frac{(3x^2-r^2)+(3y^2-r^2)+(3z^2-r^2)}{r^5}\\
  &=\frac{3(x^2+y^2+z^2)-3r^2}{r^5}=\frac{3r^2-3r^2}{r^5}=0.
  \end{align*}
  恒等式成立，证毕。
\end{enumerate}
\end{proof}

\begin{exercise}[一维波动方程的通解验证] \label{ex:wave-dalembert}
设 $\phi$, $\psi:\mathbb R\to\mathbb R$ 为任意二阶可微函数，$a>0$ 为常数。证明函数 \(u(x,t)=\phi(x-at)+\psi(x+at)\) 满足一维\textbf{弦振动方程} \(\displaystyle a^2u_{xx}=u_{tt}\)。
\end{exercise}

\begin{proof}
这是一个典型的 ``行波解'' 形式：$\phi(x-at)$ 代表以速度 $a$ 向右传播的波（波形不变，只是位置随时间推移），而 $\psi(x+at)$ 代表以速度 $a$ 向左传播的波。直接计算各阶偏导数：

\begin{enumerate}
  \item \textbf{计算关于 $x$ 的导数}：
  \(\displaystyle u_x=\phi'(x-at)+\psi'(x+at),\quad u_{xx}=\phi''(x-at)+\psi''(x+at)\)。

  \item \textbf{计算关于 $t$ 的导数}（注意内层函数对 $t$ 的导数带因子 $\mp a$）：
  \(\displaystyle u_t=-a\,\phi'(x-at)+a\,\psi'(x+at)\)，再求一次导得 \(\displaystyle u_{tt}=a^2\bigl[\phi''(x-at)+\psi''(x+at)\bigr]\)。

  \item \textbf{比较两端验证方程}：
  \(\displaystyle u_{tt}=a^2\bigl[\phi''(x-at)+\psi''(x+at)\bigr]=a^2u_{xx}\)。
  等式恒成立，证毕。
\end{enumerate}
\end{proof}

\begin{note}[d'Alembert 公式的意义]
这个结果实际上是 \textbf{d'Alembert 公式}（1746 年）的一维情形特例。对于初值问题
\[
\begin{cases}
u_{tt}-a^2u_{xx}=0,&(x,t)\in\mathbb R\times(0,+\infty),\\
u(x,0)=f(x),\;u_t(x,0)=g(x),
\end{cases}
\]
其解可以写成
\[
u(x,t)=\frac{f(x+at)+f(x-at)}{2}+\frac{1}{2a}\int_{x-at}^{x+at}g(s)\,\mathrm{d}s,
\]
\end{note}


有时待验证的表达式不是显式的 $z=z(x,y)$，而是隐式地由某个方程确定。

\begin{exercise}[隐函数形式的 PDE 恒等式] \label{ex:implicit-pde-identity}
设 $z=z(x,y)$ 由方程 $F\!\left(x+\dfrac{z}{y},\;y+\dfrac{z}{x}\right)=0$ 唯一确定（假设隐函数条件满足）。证明恒等式：
\(\displaystyle x\,\frac{\partial z}{\partial x}+y\,\frac{\partial z}{\partial y}=z-xy\)。
\end{exercise}

\begin{proof}
\textbf{方法一：按复合关系直接求导。}

\textbf{先说选法：}原方程中的 $F$ 并不是直接作用在 $(x,y,z)$ 上，而是作用在两个复合量 \(\displaystyle x+\frac{z}{y},\ y+\frac{z}{x}\) 上。若不先把这两个整体单独记号化，后续对 $x$、$y$ 求偏导时链式法则会显得非常杂乱。因此我们先压缩记号，把“真正进入 $F$ 的自变量”单独抽出来。

为了便于书写，令中间变量 \(\displaystyle \xi=x+\frac{z}{y},\ \eta=y+\frac{z}{x}\)，则约束方程为 $F(\xi,\eta)=0$。

\begin{enumerate}
  \item \textbf{对 $x$ 求全导数：}

  将 $z$ 视为 $x,y$ 的函数，对方程 $F(\xi,\eta)=0$ 两边关于 $x$ 求偏导：\(\displaystyle F_\xi \cdot \frac{\partial \xi}{\partial x} + F_\eta \cdot \frac{\partial \eta}{\partial x} = 0\)。
  其中各中间变量对 $x$ 的偏导数为 $\dfrac{\partial \xi}{\partial x} = 1 + \dfrac{z_x}{y}$，$\dfrac{\partial \eta}{\partial x} = -\dfrac{z}{x^2} + \dfrac{z_x}{x}$。
  代入后整理提取 $z_x$：
代入后整理，得 \(\displaystyle z_x\left(\frac{F_\xi}{y}+\frac{F_\eta}{x}\right)=-F_\xi+\frac{z}{x^2}F_\eta\)。

  \item \textbf{对 $y$ 求全导数：}

  同理，对方程 $F(\xi,\eta)=0$ 两边关于 $y$ 求偏导：\(\displaystyle F_\xi \cdot \frac{\partial \xi}{\partial y} + F_\eta \cdot \frac{\partial \eta}{\partial y} = 0\)。
  其中 $\dfrac{\partial \xi}{\partial y} = -\dfrac{z}{y^2} + \dfrac{z_y}{y}$，$\dfrac{\partial \eta}{\partial y} = 1 + \dfrac{z_y}{x}$。
  代入后整理提取 $z_y$：
代入后整理，得 \(\displaystyle z_y\left(\frac{F_\xi}{y}+\frac{F_\eta}{x}\right)=\frac{z}{y^2}F_\xi-F_\eta\)。

  \item \textbf{构造目标式的左端并化简：}

  令公共系数 $K = \dfrac{F_\xi}{y} + \dfrac{F_\eta}{x} = \dfrac{x F_\xi + y F_\eta}{xy}$。

  我们分别将 $z_x$ 乘以 $x$，将 $z_y$ 乘以 $y$ 并相加：\(\displaystyle x z_x \cdot K = -x F_\xi + \frac{z}{x}F_\eta\)，\(\displaystyle y z_y \cdot K = \frac{z}{y}F_\xi - y F_\eta\)。
  将上述两式相加得：
  \[
  K \left(x z_x + y z_y\right)
  = -(x F_\xi + y F_\eta) + z\left(\frac{F_\xi}{y} + \frac{F_\eta}{x}\right)
  = -(x F_\xi + y F_\eta) + z K.
  \]

  \item \textbf{验证目标恒等式：}

  将上式两边同除以 $K$：\(\displaystyle x z_x + y z_y = -\frac{x F_\xi + y F_\eta}{K} + z\)。
  又因为 $K = \dfrac{x F_\xi + y F_\eta}{xy}$，所以 $\dfrac{x F_\xi + y F_\eta}{K} = xy$。

  代入即得 \(\displaystyle x z_x + y z_y = -xy + z\)。
  证毕。
\end{enumerate}

\medskip
\noindent\textbf{方法二：先直接求出 \(z_x,z_y\)，再代入比对。}

这个做法完全沿用前面单方程隐函数求导公式。先把原方程整体记为
\[
H(x,y,z)=F\!\left(x+\frac{z}{y},\;y+\frac{z}{x}\right)=0,
\]
其中
\[
\xi=x+\frac{z}{y},\qquad \eta=y+\frac{z}{x}.
\]
计算 \(H_x,H_y,H_z\) 时，先把 \(x,y,z\) 视为相互独立的变量；等到套隐函数公式时，再把 \(z\) 看成 \(z(x,y)\)。由链式法则，
\[
\begin{aligned}
H_x&=F_\xi\cdot 1+F_\eta\left(-\frac{z}{x^2}\right)
=F_\xi-\frac{z}{x^2}F_\eta,\\
H_y&=F_\xi\left(-\frac{z}{y^2}\right)+F_\eta\cdot 1
=-\frac{z}{y^2}F_\xi+F_\eta,\\
H_z&=F_\xi\cdot\frac1y+F_\eta\cdot\frac1x
=\frac{F_\xi}{y}+\frac{F_\eta}{x}.
\end{aligned}
\]
这里 \(F_\xi,F_\eta\) 都是在 \((\xi,\eta)\) 处取值。由单方程隐函数求导公式
\[
z_x=-\frac{H_x}{H_z},\qquad z_y=-\frac{H_y}{H_z},
\]
于是
\[
\boxed{
z_x=
\frac{-F_\xi+\dfrac{z}{x^2}F_\eta}
{\dfrac{F_\xi}{y}+\dfrac{F_\eta}{x}},
\qquad
z_y=
\frac{\dfrac{z}{y^2}F_\xi-F_\eta}
{\dfrac{F_\xi}{y}+\dfrac{F_\eta}{x}}
}.
\]
令
\[
K=\frac{F_\xi}{y}+\frac{F_\eta}{x}
=\frac{xF_\xi+yF_\eta}{xy}.
\]
把刚才求出的 \(z_x,z_y\) 直接代入目标式左端：
\[
\begin{aligned}
xz_x+yz_y
&=\frac{x\left(-F_\xi+\dfrac{z}{x^2}F_\eta\right)
y\left(\dfrac{z}{y^2}F_\xi-F_\eta\right)}{K}\\
&=\frac{-xF_\xi-yF_\eta
z\left(\dfrac{F_\xi}{y}+\dfrac{F_\eta}{x}\right)}{K}\\
&=\frac{zK-(xF_\xi+yF_\eta)}{K}
=z-\frac{xF_\xi+yF_\eta}{K}.
\end{aligned}
\]
又因为 \(\displaystyle K=\frac{xF_\xi+yF_\eta}{xy}\)，所以
\[
\frac{xF_\xi+yF_\eta}{K}=xy.
\]
因此
\[
xz_x+yz_y=z-xy.
\]
\end{proof}

\subsubsection{变量替换（坐标变换）}

\begin{proposition}[极坐标下的梯度模长与 Laplace 算子] \label{prop:polar-laplacian}
设 $u=f(x,y)$ 具有二阶连续偏导数。在极坐标变换 $x=r\cos\theta$, $y=r\sin\theta$ 下（此时 $u=f(r\cos\theta,r\sin\theta)=:F(r,\theta)$），以下两个重要表达式成立：

\textbf{(1)} 梯度的模长平方（保持不变）：
\[
\left(\frac{\partial u}{\partial x}\right)^{\!2}
+\left(\frac{\partial u}{\partial y}\right)^{\!2}
=\left(\frac{\partial u}{\partial r}\right)^{\!2}
+\frac{1}{r^2}\left(\frac{\partial u}{\partial\theta}\right)^{\!2}.
\]

\textbf{(2)} Laplace 算子的极坐标形式：
\[
\frac{\partial^2 u}{\partial x^2}+\frac{\partial^2 u}{\partial y^2}
=\frac{\partial^2 u}{\partial r^2}
+\frac{1}{r}\frac{\partial u}{\partial r}
+\frac{1}{r^2}\frac{\partial^2 u}{\partial\theta^2}
=\frac{1}{r^2}\left[r\frac{\partial}{\partial r}\!\left(r\frac{\partial u}{\partial r}\right)
+\frac{\partial^2 u}{\partial\theta^2}\right].
\]
\end{proposition}

\begin{proof}
\begin{enumerate}
  \item \textbf{预备工作：链式法则建立 $u_x$, $u_y$ 与 $u_r$, $u_\theta$ 的关系。}

  首先需要计算 $\dfrac{\partial r}{\partial x}$, $\dfrac{\partial r}{\partial y}$, $\dfrac{\partial\theta}{\partial x}$, $\dfrac{\partial\theta}{\partial y}$。

  由 $r^2=x^2+y^2$ 两边对 $x$ 求偏导：$2r\,r_x=2x \implies r_x=\dfrac{x}{r}=\cos\theta$。
  同理 $r_y=\dfrac{y}{r}=\sin\theta$。

  由 $\tan\theta=\dfrac{y}{x}$（在 $x>0$ 区域），两边对 $x$ 求导：$\sec^2\theta\cdot\theta_x=-\dfrac{y}{x^2}$，
  故 $\theta_x=\dfrac{-y/x^2}{\sec^2\theta} =-\dfrac{y}{x^2}\cdot\dfrac{x^2}{r^2}=-\dfrac{y}{r^2} =-\dfrac{\sin\theta}{r}$。
  类似地，$\theta_y=\dfrac{1/x}{\sec^2\theta} = \dfrac{x}{r^2} = \dfrac{\cos\theta}{r}$。

  总结四个变换系数：
  \[
    r_x=\cos\theta,\qquad r_y=\sin\theta,\qquad
    \theta_x=-\dfrac{\sin\theta}{r},\qquad \theta_y=\dfrac{\cos\theta}{r}.
  \]

  现在计算 $u_x$ 和 $u_y$（注意 $u$ 通过 $r,\theta$ 依赖于 $x,y$）：
  \begin{align}
  u_x&=u_r\,r_x+u_\theta\,\theta_x
  =u_r\cos\theta-u_\theta\frac{\sin\theta}{r}, \tag{A}\\[6pt]
  u_y&=u_r\,r_y+u_\theta\,\theta_y
  =u_r\sin\theta+u_\theta\frac{\cos\theta}{r}. \tag{B}
  \end{align}

  \item \textbf{证明 (1)：梯度模长平方的不变性。}
  \begin{align*}
  u_x^2+u_y^2
  &=\left(u_r\cos\theta-u_\theta\frac{\sin\theta}{r}\right)^2
   +\left(u_r\sin\theta+u_\theta\frac{\cos\theta}{r}\right)^2\\[6pt]
  &=u_r^2(\cos^2\theta+\sin^2\theta)
  +\frac{u_\theta^2}{r^2}(\sin^2\theta+\cos^2\theta)\\[4pt]
  &\quad-\frac{2u_ru_\theta\cos\theta\sin\theta}{r}
  +\frac{2u_ru_\theta\sin\theta\cos\theta}{r}\\[6pt]
  &=u_r^2+\frac{u_\theta^2}{r^2}.
  \end{align*}
  交叉项恰好抵消 —— 这正是 ``正交变换保长度'' 的体现。极坐标变换在每一点处都是局部正交的（$r$ 线与 $\theta$ 线垂直）。

  \item \textbf{证明 (2)：Laplace 算子的推导。}

  先对 $u_x$（式 (A)）再次关于 $x$ 求导：
  \begin{align*}
  u_{xx}&=\frac{\partial}{\partial x}(u_r\cos\theta)-\frac{\partial}{\partial x}\!\left(u_\theta\frac{\sin\theta}{r}\right)\\[4pt]
  &=\cos\theta\,(u_{rr}r_x+u_{r\theta}\theta_x) -u_r\sin\theta\cdot\theta_x\\[4pt]
  &\quad-\frac{\sin\theta}{r}(u_{\theta r}r_x+u_{\theta\theta}\theta_x) -u_\theta\cdot\frac{\partial}{\partial x}\!\left(\frac{\sin\theta}{r}\right). \tag{*}
  \end{align*}

  计算 $\dfrac{\partial}{\partial x}\!\left(\dfrac{\sin\theta}{r}\right)$：
  \begin{align*}
  \frac{\partial}{\partial x}\!\left(\frac{\sin\theta}{r}\right)
  &=\frac{r\cos\theta\cdot\theta_x-\sin\theta\cdot r_x}{r^2}
  =\frac{r\cos\theta\left(-\frac{\sin\theta}{r}\right)-\sin\theta\cos\theta}{r^2}\\[6pt]
  &=-\frac{2\sin\theta\cos\theta}{r^2}.
  \end{align*}

  代入所有系数，将 (*) 式展开：
  \begin{align*}
  u_{xx}&=u_{rr}\cos^2\theta-\frac{2\sin\theta\cos\theta}{r}u_{r\theta}
  +\frac{u_r\sin^2\theta}{r}
  +\frac{u_{\theta\theta}\sin^2\theta}{r^2}
  +\frac{2u_\theta\sin\theta\cos\theta}{r^2}. \tag{C}
  \end{align*}

  利用完全相同的推导模式（只需将 $\cos\leftrightarrow\sin$ 并注意符号变化），对 $u_y$（式 (B)）关于 $y$ 求导得：
  \begin{align*}
  u_{yy}&=u_{rr}\sin^2\theta+\frac{2\sin\theta\cos\theta}{r}u_{r\theta}
  +\frac{u_r\cos^2\theta}{r}
  +\frac{u_{\theta\theta}\cos^2\theta}{r^2}
  -\frac{2u_\theta\sin\theta\cos\theta}{r^2}. \tag{D}
  \end{align*}

  \item \textbf{合并结果：}

  将 (C) 和 (D) 相加，交叉导数项 $u_{r\theta}$ 以及一阶导数交叉项 $u_\theta$ 相互抵消，剩余各项按类型分组：
  \begin{align*}
  \Delta u&=u_{xx}+u_{yy}\\[4pt]
  &=u_{rr}(\cos^2\theta+\sin^2\theta)
  +u_r\left(\frac{\sin^2\theta}{r}+\frac{\cos^2\theta}{r}\right)
  +u_{\theta\theta}\left(\frac{\sin^2\theta}{r^2}+\frac{\cos^2\theta}{r^2}\right)\\[6pt]
  &=u_{rr}+\frac{1}{r}u_r+\frac{1}{r^2}u_{\theta\theta}.
  \end{align*}
  证毕。
\end{enumerate}
\end{proof}

\begin{note}[极坐标 Laplace 算子的标准形式]
结果 $u_{rr}+\dfrac{1}{r}u_r+\dfrac{1}{r^2}u_{\theta\theta}=0$ 通常写作更紧凑的形式：
\[
\frac{1}{r}\frac{\partial}{\partial r}\!\left(r\frac{\partial u}{\partial r}\right)
+\frac{1}{r^2}\frac{\partial^2 u}{\partial\theta^2}=0.
\]
这是因为 $\dfrac{1}{r}(ru_r)_r=\dfrac{1}{r}(u_r+ru_{rr})=\dfrac{u_r}{r}+u_{rr}$，正好对应前两项之和。这种写法在分离变量法求解时更为方便。
\end{note}

\begin{note}[极坐标 Laplace 算子为什么长成这样]
可以把 \(\displaystyle u_{rr}+\frac{1}{r}u_r+\frac{1}{r^2}u_{\theta\theta}\) 分成 ``径向变化'' 和 ``角向变化'' 两部分来看：$u_{rr}$ 是沿半径方向的二阶变化；$\dfrac1r u_r$ 来自半径变大时，同样的径向流量要分摊到更长的圆周上，所以会多出一个几何修正因子 $\dfrac1r$；而 $\dfrac1{r^2}u_{\theta\theta}$ 则是因为角变量 $\theta$ 不是实际长度，圆周上的实际弧长是 $r\,\mathrm{d}\theta$，所以换回真实长度后要再除一个 $r^2$。

因此这个式子不是凭空多出两个 $r$，而是 ``直角坐标里的均匀变化'' 改写到极坐标后，必须把圆周随半径扩张的几何效应一起算进去。
\end{note}


上面的例子展示了最常用的坐标变换（直角 $\to$ 极坐标）。在实际问题中，还可能出现更一般的变量替换 —— 新自变量和新因变量都可能以复杂的非线性方式依赖于旧变量。

\begin{exercise}[一般变量替换下的 PDE 变换] \label{ex:variable-change-pde}
作变量替换
作变换 \(\displaystyle u=\frac{y}{x},\quad v=y\)（新自变量），\(\displaystyle w=yz-x\)（新因变量），
将方程 \(\displaystyle xz_{xx}+2z_x=\frac{2}{y}\) 变换为新坐标系 $(u,v)$ 下关于 $w$ 的方程。
\end{exercise}
\begin{solution}
\begin{enumerate}
  \item \textbf{理清变量依赖关系。}

  旧体系的自变量是 $(x,y)$，因变量是 $z=z(x,y)$；新体系的自变量是 $(u,v)$，因变量是 $w=w(u,v)$，并且
  \(\displaystyle u=\frac{y}{x},\quad v=y,\quad w=yz-x\)。

  \item \textbf{将 $z$ 用新变量表示。}

  由 $w=yz-x$ 和 $v=y$ 得 $w=zv-x$，即 \(\displaystyle z=\frac{x+w}{v}=\frac{x}{v}+\frac{w}{v}\)。

  \item \textbf{用链式法则计算 $z_x$。}

  这里的关键是 $w$ 通过 $(u,v)$ 间接依赖于 $(x,y)$，而 $u,v$ 又依赖于 $(x,y)$：
  \begin{align*}
  z_x&=\frac{\partial}{\partial x}\!\left(\frac{x}{v}+\frac{w}{v}\right)
  =\frac{1}{v}+\frac{1}{v}(w_u u_x+w_v v_x)-w\cdot\frac{v_x}{v^2}\\[6pt]
  &=\frac{1}{v}+\frac{w_u u_x+w_v v_x}{v}-\frac{wv_x}{v^2}.
  \end{align*}
  需要先算出 $u_x$ 与 $v_x$：$u_x=\dfrac{\partial}{\partial x}\!\left(\dfrac{y}{x}\right)=-\dfrac{y}{x^2}=-\dfrac{v}{x^2}$，$v_x=0$。
  代入得（注意 $v_x=0$ 大大简化了计算）：
  \[
  z_x=\frac{1}{v}+\frac{1}{v}\!\left(w_u\!\left(-\frac{v}{x^2}\right)+w_v\cdot 0\right)-0
  =\frac{1}{v}-\frac{vw_u}{x^2v}=\frac{1}{v}-\frac{w_u}{x^2}.
  \]

  \item \textbf{计算 $z_{xx}$。}
  \begin{align*}
  z_{xx}&=\frac{\partial}{\partial x}\!\left(\frac{1}{v}-\frac{w_u}{x^2}\right)
  =0-\frac{\partial}{\partial x}\!\left(w_u\cdot x^{-2}\right)\\[6pt]
  &=-\Bigl[(w_{uu}u_x+w_{uv}v_x)\cdot x^{-2}+w_u\cdot(-2x^{-3})\Bigr]\\[6pt]
  &=-w_{uu}\!\left(-\frac{v}{x^2}\right)\!\cdot\frac{1}{x^2} +\frac{2w_u}{x^3}
  =\frac{vw_{uu}}{x^4}+\frac{2w_u}{x^3}.
  \end{align*}

  \item \textbf{代回原方程。}
  \begin{align*}
  xz_{xx}+2z_x
  &=x\!\left(\frac{vw_{uu}}{x^4}+\frac{2w_u}{x^3}\right)
  +2\!\left(\frac{1}{v}-\frac{w_u}{x^2}\right)\\[6pt]
  &=\frac{vw_{uu}}{x^3}+\frac{2w_u}{x^2}+\frac{2}{v}-\frac{2w_u}{x^2}\\[6pt]
  &=\frac{vw_{uu}}{x^3}+\frac{2}{v}.
  \end{align*}
  原方程要求上式等于 $\dfrac{2}{y}=\dfrac{2}{v}$，即 \(\displaystyle \frac{vw_{uu}}{x^3}+\frac{2}{v}=\frac{2}{v}\implies \frac{vw_{uu}}{x^3}=0\)。
  由于 $v=y \neq 0$ 且 $x \neq 0$（在定义域内），最终得到变换后的方程 \(w_{uu}=0\)。
\end{enumerate}
\end{solution}

\begin{note}[结果的解读]
变换后的方程 $w_{uu}=0$ 是一个二阶常微分方程，其通解为 $w(u,v)=A(v)\cdot u+B(v)$，其中 $A,B$ 是关于 $v$ 的任意函数。

这意味着原 PDE 的解 $z(x,y)$ 必须满足某种特定的结构 —— 具体来说，当以 $u=y/x$ 和 $v=y$ 作为 ``自然变量'' 来观察时，$w=yz-x$ 对 $u$ 的二阶变化率为零。这个例子展示了变量替换的效果：一个看似复杂的变系数 PDE，经过恰当的变量替换后变成了可以直接积分的形式。
\end{note}

\section{习题}

\subsection{\texorpdfstring{第一节习题：$n$ 维欧氏空间的基本概念}{第一节习题：n 维欧氏空间的基本概念}}

\subsubsection{A组}

\begin{exercise}[基本概念问答]
回答下列问题：
\begin{shortexenum}
  \shortitem{1}{$n$ 维欧氏空间 $\mathbb{R}^n$ 如何定义？} &
  \shortitem{2}{$\mathbb{R}^2$ 中点的邻域如何定义？} \\
  \shortitem{3}{什么叫集合 $E$ 的内点和外点、边界点？} &
  \shortitem{4}{集合 $E$ 的聚点怎样定义？聚点和边界点有何不同？} \\
  \shortitem{5}{$\mathbb{R}^2$ 中的开集和闭集怎样定义？} &
  \shortitem{6}{闭包是指什么？} \\
  \shortitem{7}{什么叫开区域和闭区域？} &
\end{shortexenum}
\end{exercise}
\begin{solution}
\begin{enumerate}
  \item[(1)] \textbf{思维过程}：回顾笛卡尔积和欧氏空间的定义。首先定义 $n$ 维实坐标空间为 $\mathbb R$ 的 $n$ 次笛卡尔积，再赋予其内积结构。\\
  \textbf{解答}：$n$ 维欧氏空间 $\mathbb{R}^n$ 首先是集合 $\mathbb{R} \times \mathbb{R} \times \dots \times \mathbb{R}$ 的全体 $n$ 元有序实数组 $(x_1, x_2, \dots, x_n)$ 构成的线性空间，并且在该空间上定义了内积运算 $\langle x,y \rangle = \sum_{i=1}^n x_i y_i$ 及其诱导的欧氏距离。

  \item[(2)] \textbf{思维过程}：邻域通常指开球。在二维平面中即为不含边界的圆盘。\\
  \textbf{解答}：在 $\mathbb{R}^2$ 中，点 $P_0(x_0,y_0)$ 的 $r$ 邻域（$r>0$）定义为到 $P_0$ 的距离严格小于 $r$ 的所有点的集合，即开球 $B_r(P_0) = \{P(x,y) \in \mathbb{R}^2 \mid \sqrt{(x-x_0)^2+(y-y_0)^2} < r\}$。

  \item[(3)] \textbf{思维过程}：根据点的一个开邻域与该集合的位置关系来刻画。\\
  \textbf{解答}：对于集合 $E$ 和点 $p$：
  \begin{itemize}
    \item \textbf{内点}：若存在 $r>0$ 使得 $B_r(p) \subset E$，则 $p$ 是 $E$ 的内点。
    \item \textbf{外点}：若存在 $r>0$ 使得 $B_r(p) \cap E = \varnothing$（即 $B_r(p) \subset E^c$），则 $p$ 是 $E$ 的外点。
    \item \textbf{边界点}：若对于任意 $r>0$，$B_r(p)$ 中既含有 $E$ 中的点，又含有不属于 $E$ 的点，则 $p$ 是 $E$ 的边界点。
  \end{itemize}

  \item[(4)] \textbf{思维过程}：聚点强调“无限逼近”，需要用去心邻域。边界点强调“内外相交”。这两者可能有重叠但也存在差异。\\
  \textbf{解答}：
  \textbf{聚点定义}：若点 $p$ 的任意去心邻域 $(B_r(p) \setminus \{p\})$ 内都含有集合 $E$ 中的点，则称 $p$ 为 $E$ 的聚点。\\
  \textbf{不同之处}：边界点反映的是点相对于集合内外的拓扑“位置”，而聚点反映的是集合中点列的“极限”行为。例如有限集（如 $\{a\}$）有边界点 $a$，但没有任何聚点；反之，开集内的点是聚点但不是边界点。只有既不在集合内部、又被集合内的点无限逼近的点，才既是边界点又是聚点。

  \item[(5)] \textbf{思维过程}：开集是没有边界的集合，闭集是包含所有边界的集合（或补集是开集）。\\
  \textbf{解答}：
  \begin{itemize}
    \item \textbf{开集}：若集合 $E$ 中的每一个点都是它的内点，则称 $E$ 为开集。
    \item \textbf{闭集}：若集合 $E$ 的补集 $E^c$（即 $\mathbb{R}^2 \setminus E$）是开集，则称 $E$ 为闭集。等价地，若 $E$ 包含其所有的聚点，则 $E$ 为闭集。
  \end{itemize}

  \item[(6)] \textbf{思维过程}：闭包等于原集合并上其导集（全体聚点），或者并上其边界。\\
  \textbf{解答}：集合 $E$ 的闭包，记为 $\bar{E}$，定义为集合 $E$ 与其所有聚点构成的集合（导集 $E'$）的并集，即 $\bar{E} = E \cup E'$。它也可以理解为包含 $E$ 的最小闭集。

  \item[(7)] \textbf{思维过程}：连通开集为开区域，加上边界就是闭区域。\\
  \textbf{解答}：
  \begin{itemize}
    \item \textbf{开区域}：如果一个开集是连通的（即其中任意两点都可以用完全位于该集合内的折线相连），则称该开集为开区域（或简称区域）。
    \item \textbf{闭区域}：开区域加上它所有的边界点（即开区域的闭包），称为闭区域。
  \end{itemize}
\end{enumerate}
\end{solution}

\begin{exercise}[范数不等式的证明与几何意义]
证明不等式：
\begin{shortexenum}
  \shortitem{1}{$\|x-z\| \leqslant  \|x-y\| + \|y-z\|$；} &
  \shortitem{2}{$|\|x\| - \|y\|| \leqslant  \|x-y\|$。}
\end{shortexenum}
其中，$x,y,z \in \mathbb{R}^n$，并在 $\mathbb{R}^2$ 中给出不等式(1)的几何解释。
\end{exercise}
\begin{solution}
\begin{enumerate}
  \item[(1)] 令 \(u=x-y\), \(v=y-z\)，则 \(u+v=x-z\)。由基本三角不等式 \(\displaystyle \|x-z\|=\|u+v\|\leqslant \|u\|+\|v\|=\|x-y\|+\|y-z\|\)。在 \(\mathbb{R}^2\) 中，这正是三角形两边之和不小于第三边。

  \item[(2)] 先写
  \(\displaystyle \|x\|=\|(x-y)+y\|\leqslant \|x-y\|+\|y\|\)，且
  \(\displaystyle \|y\|=\|(y-x)+x\|\leqslant \|x-y\|+\|x\|\)。
  两式分别移项，就得到
  \(\displaystyle -\|x-y\|\leqslant \|x\|-\|y\|\leqslant \|x-y\|\)，
  也就是
  \(\displaystyle \bigl|\|x\|-\|y\|\bigr|\leqslant \|x-y\|\)。
\end{enumerate}
\end{solution}

\begin{exercise}[闭集的判定]
集合 $B=\{(x,y) \mid x^2+y^2=1\}$ 是否为 $\mathbb{R}^2$ 中的闭集？为什么？
\end{exercise}
\begin{note}
判断一个集合是否为闭集，最直接的方法是考察其补集是否为开集，或者验证该集合是否包含它所有的边界点 / 聚点。集合 $B$ 是平面上的单位圆周。直观上它的补集是圆内和圆外的两块开区域组成的，因此它是闭集。
\end{note}
\begin{solution}
集合 $B$ \textbf{是} $\mathbb{R}^2$ 中的闭集。
\begin{itemize}
\item \textbf{证明方法一（利用补集是开集）}：
考察集合 $B$ 的补集
\[
B^c=\{(x,y) \mid x^2+y^2 \neq 1\}
=\{(x,y) \mid x^2+y^2 < 1\}\cup \{(x,y) \mid x^2+y^2 > 1\}.
\]
设 $D_1 = \{(x,y) \mid x^2+y^2 < 1\}$，$D_2 = \{(x,y) \mid x^2+y^2 > 1\}$。
对于 $D_1$ 中的任意一点 $p(x_0, y_0)$，令 $d = \sqrt{x_0^2+y_0^2} < 1$。取邻域半径 $r = 1 - d > 0$，则 $B_r(p) \subset D_1$。因此 $D_1$ 中每一点都是内点，即 $D_1$ 是开集。
同理，对于 $D_2$ 中的任意一点 $q(x_1, y_1)$，令 $d' = \sqrt{x_1^2+y_1^2} > 1$。取半径 $r' = d' - 1 > 0$，则 $B_{r'}(q) \subset D_2$。因此 $D_2$ 也是开集。
由于两个开集的并集 $B^c = D_1 \cup D_2$ 仍然是开集，根据闭集的定义，其补集 $B$ 必定是闭集。
\item \textbf{证明方法二（利用聚点）}：
设 $(x_0, y_0)$ 为集合 $B$ 的任意一个聚点。根据聚点的性质，存在 $B$ 中的点列 $\{(x_n, y_n)\}$ 使得当 $n \to \infty$ 时，$(x_n, y_n) \to (x_0, y_0)$。
因为对于所有的 $n$，都有 $x_n^2 + y_n^2 = 1$。
对等式两边取极限，利用二元函数 $f(x,y) = x^2+y^2$ 的连续性，可得 \(\displaystyle \lim_{n \to \infty} (x_n^2 + y_n^2) = x_0^2 + y_0^2 = 1\)。
这说明聚点 $(x_0, y_0)$ 也满足集合 $B$ 的定义方程，故 $(x_0, y_0) \in B$。
由于集合 $B$ 包含了它的所有聚点，因此 $B$ 是闭集。
\end{itemize}
\end{solution}

\begin{exercise}[点集拓扑基本概念的计算]
求下列点集的内点、外点和边界点：
\begin{shortexenum}
  \shortitem{1}{$A=\{(x,y) \mid y < x^2\}$；} &
  \shortitem{2}{$B=\left\{(x,y) \mid 1 \leqslant  \frac{x^2}{3} + \frac{y^2}{4} < 5\right\}$；} \\
  \shortitem{3}{$E=\{(x,y) \mid 0 < x^2+y^2 < 1\}$；} &
  \shortitem{4}{$F=\{(x,y) \mid x \text{ 与 } y \text{ 都是有理数}\}$。}
\end{shortexenum}
\end{exercise}
\begin{solution}
\begin{enumerate}
  \item[(1)] \textbf{思维过程}：集合 $A$ 表示抛物线 $y = x^2$ 下方的无边界区域。不等式严格（$<$）意味着它本身就是开集。
  \textbf{解答}：
  \begin{itemize}
    \item \textbf{内点集}：抛物线严格下方的所有点，即 $A^\circ = \{(x,y) \mid y < x^2\}$。
    \item \textbf{边界点集}：抛物线上的所有点构成的曲线，即 $\partial A = \{(x,y) \mid y = x^2\}$。
    \item \textbf{外点集}：抛物线严格上方的所有点，即 $(A^c)^\circ = \{(x,y) \mid y > x^2\}$。
  \end{itemize}

  \item[(2)] \textbf{思维过程}：集合 $B$ 是由两个同心椭圆界定的环形区域（内含边界，外不含边界）。内部去掉了所有的等号即可得到内点集。边界包括内外两个椭圆的轮廓。
  \textbf{解答}：
  \begin{itemize}
    \item \textbf{内点集}：去掉等号的严格不等式部分，即 \(\displaystyle B^\circ = \left\{(x,y) \mid 1 < \frac{x^2}{3} + \frac{y^2}{4} < 5\right\}\)。
    \item \textbf{边界点集}：等号成立的两个椭圆周，即
    \[
    \partial B =
    \left\{(x,y) \mid \frac{x^2}{3} + \frac{y^2}{4} = 1\right\}
    \cup
    \left\{(x,y) \mid \frac{x^2}{3} + \frac{y^2}{4} = 5\right\}.
    \]
    \item \textbf{外点集}：椭圆环外部和中心空洞内部的开区域，即
    \[
    (B^c)^\circ=
    \left\{(x,y) \mid \frac{x^2}{3} + \frac{y^2}{4} < 1\right\}
    \cup
    \left\{(x,y) \mid \frac{x^2}{3} + \frac{y^2}{4} > 5\right\}.
    \]
  \end{itemize}

  \item[(3)] \textbf{思维过程}：集合 $E$ 是一个去掉原点的开单位圆盘（去心开圆）。
  \textbf{解答}：
  \begin{itemize}
    \item \textbf{内点集}：由于 $E$ 中的不等式都是严格的，它本身就是开集，因此内点集为其自身：$E^\circ = \{(x,y) \mid 0 < x^2+y^2 < 1\}$。
    \item \textbf{边界点集}：包括外圈的圆周和被挖去的中心点原点，即 $\partial E = \{(x,y) \mid x^2+y^2 = 1\} \cup \{(0,0)\}$。
    \item \textbf{外点集}：圆周外部严格区域，即 $(E^c)^\circ = \{(x,y) \mid x^2+y^2 > 1\}$。
  \end{itemize}

  \item[(4)] \textbf{思维过程}：$F$ 为平面上的有理数点集。根据实数理论，任何一个点的任意小的邻域内都同时存在有理点和无理点。这意味着没有任何邻域能被完全包在 $F$ 中，也没有任何邻域能完全避开 $F$。
  \textbf{解答}：
  \begin{itemize}
    \item \textbf{内点集}：$\varnothing$（无内点，因为任何开球内都有无理点，不能完全属于 $F$）。
    \item \textbf{外点集}：$\varnothing$（无外点，因为任何开球内都有有理点，不能完全属于 $F^c$）。
    \item \textbf{边界点集}：整个平面 $\mathbb{R}^2$。平面上任意一点的任意邻域中都既包含 $F$ 中的点（有理点）又包含 $F^c$ 中的点（无理点）。
  \end{itemize}
\end{enumerate}
\end{solution}

\subsubsection{B组}

\begin{exercise}[开集性质的证明]
证明 $\mathbb{R}^2$ 中每一个开球都是开集。
\end{exercise}
\begin{note}
要证明开球 $B_r(p_0)$ 是开集，必须依据开集的定义：证明开球内的任意一点 $q$ 都是内点。
这就需要找到一个以 $q$ 为中心、半径为 $\delta$ 的更小的新开球 $B_\delta(q)$，使得这个新开球完全包含在原开球 $B_r(p_0)$ 内部。
利用几何直观（或三角不等式），因为 $q$ 距离中心 $p_0$ 的距离 $d = \|q - p_0\| < r$，只要我们取 $\delta = r - d$，那么从 $q$ 往任何方向走不超过 $\delta$ 的距离，总距离 $p_0$ 就不会超过 $d + \delta = r$。
\end{note}
\begin{solution}
设 $B_r(p_0) = \{x \in \mathbb{R}^2 \mid \|x - p_0\| < r\}$ 为 $\mathbb{R}^2$ 中的任意一个开球（$r > 0$）。
任取 $q \in B_r(p_0)$，依据开球定义可知 \(\displaystyle d = \|q - p_0\| < r\)。我们要证明 $q$ 是 $B_r(p_0)$ 的内点。
令 $\delta = r - d$。因为 $d < r$，所以 $\delta > 0$。
考虑以 $q$ 为中心，$\delta$ 为半径的开球 $B_\delta(q)$。任取 $z \in B_\delta(q)$，满足 \(\displaystyle \|z - q\| < \delta\)。
我们来计算 $z$ 到原开球中心 $p_0$ 的距离。利用范数的三角不等式，有 \(\displaystyle
\|z - p_0\| = \|(z - q) + (q - p_0)\| \leqslant \|z - q\| + \|q - p_0\| < \delta + d = (r - d) + d = r\)。
因为 $\|z - p_0\| < r$，所以点 $z$ 必定属于原开球 $B_r(p_0)$。
由于 $z$ 是 $B_\delta(q)$ 中的任意一点，故有 \(\displaystyle B_\delta(q) \subset B_r(p_0)\)。
这说明点 $q$ 具有一个完全包含在 $B_r(p_0)$ 中的邻域，即 $q$ 是 $B_r(p_0)$ 的内点。
由于 $q$ 是 $B_r(p_0)$ 中的任意一点，因此 $B_r(p_0)$ 中的每一点都是内点。
根据开集的定义，开球 $B_r(p_0)$ 是开集。证明完毕。
\end{solution}

\begin{exercise}[邻域概念的高维推广]
试仿照 $\mathbb{R}^2$ 中邻域的定义，写出 $\mathbb{R}^n$ 中点的邻域的定义。
\end{exercise}
\begin{note}
在 $\mathbb{R}^2$ 中，点 $P_0(x_0, y_0)$ 的 $r$ 邻域定义为满足 $\sqrt{(x-x_0)^2+(y-y_0)^2} < r$ 的全体点 $P(x,y)$ 构成的集合，实质上就是用二维范数定义的距离小于 $r$ 的点的集合。推广到 $\mathbb{R}^n$ 中，只需将点从二元组扩充为 $n$ 元组向量，并将二维距离（二维范数）替换为 $n$ 维的欧氏范数。
\end{note}
\begin{solution}
设 $P_0(a_1, a_2, \dots, a_n)$ 是 $n$ 维欧氏空间 $\mathbb{R}^n$ 中的一个定点，$r$ 为一个正实数（$r>0$）。
在 $\mathbb{R}^n$ 中，以 $P_0$ 为中心、$r$ 为半径的 \textbf{$r$ 邻域}（或称为\textbf{开球}）定义为 \(\displaystyle
B_r(P_0) = \left\{ P(x_1, x_2, \dots, x_n) \in \mathbb{R}^n \;\middle|\; \sqrt{\sum_{i=1}^n (x_i - a_i)^2} < r \right\}\)。
若使用范数符号记作 $P=(x_1, \dots, x_n)$ 和 $P_0=(a_1, \dots, a_n)$，则定义可以更加简洁地写作 \(\displaystyle B_r(P_0) = \bigl\{ P \in \mathbb{R}^n \mid \|P - P_0\| < r \bigr\}\)。
这表示 $\mathbb{R}^n$ 空间中所有到点 $P_0$ 的欧氏距离严格小于 $r$ 的点构成的集合。
\end{solution}

% ============================================================
% §7.2 多元函数的基本概念、极限与连续性
% ============================================================

\subsection{第二节习题：多元函数的基本概念}

\begin{exercise}[基础：定义域与图像]
\begin{shortexenum}
  \shortitem{1}{求下列函数的自然定义域并画草图：
        (a) $z=\ln(x+y)$；(b) $z=\sqrt{1-x^2}+\sqrt{y^2-1}$；
        (c) $z=\arcsin(x^2+y^2)$。} &
  \shortitem{2}{描述 $z=xy$ 和 $z=e^{-(x^2+y^2)}$ 的等高线的形状。}
\end{shortexenum}
\end{exercise}
\begin{solution}
\begin{enumerate}
  \item[(1)] \textbf{求自然定义域：}
  \begin{enumerate}
    \item[(a)] 对于对数函数 $z=\ln(x+y)$，真数必须大于 $0$。故定义域为 $x+y>0$。
    图形上：这是平面内直线 $y=-x$ 右上方的半平面（不含边界直线 $y=-x$）。
    \item[(b)] 对于偶次根号函数，被开方数必须非负。故要求 $1-x^2 \geqslant 0$ 且 $y^2-1 \geqslant 0$。
    解得：$|x| \leqslant  1$ 且 $|y| \geqslant 1$。
    图形上：这是夹在两条平行垂直直线 $x=-1$ 和 $x=1$ 之间的竖直带状区域，但要剔除掉位于水平直线 $y=-1$ 和 $y=1$ 之间的部分（即只包含 $y\geqslant 1$ 或 $y\leqslant  -1$ 的上下两段带状闭区域）。
    \item[(c)] 对于反正弦函数 $z=\arcsin(x^2+y^2)$，其变量必须落在 $[-1, 1]$ 之间。故要求 $-1 \leqslant  x^2+y^2 \leqslant  1$。
    由于 $x^2+y^2 \geqslant 0$ 自动成立，定义域等价于 $x^2+y^2 \leqslant  1$。
    图形上：这是以原点为圆心，半径为 $1$ 的闭圆盘。
  \end{enumerate}
  \item[(2)] \textbf{描述等高线形状：}
  \begin{itemize}
    \item 对于 $z=xy$：等高线方程为 $xy=c$。
    当 $c \neq 0$ 时，这是一族反比例函数的图像，即中心在原点、渐近线为坐标轴的\textbf{等轴双曲线}（$c>0$ 位于一三象限，$c<0$ 位于二四象限）。
    当 $c=0$ 时，退化为两坐标轴 $x=0$ 或 $y=0$。
    \item 对于 $z=e^{-(x^2+y^2)}$：由于 $x^2+y^2 \geqslant 0$，函数值域为 $(0, 1]$。等高线方程为 $e^{-(x^2+y^2)}=c$。
    取对数得 $-(x^2+y^2)=\ln c \implies x^2+y^2 = -\ln c$。
    当 $0<c<1$ 时，$-\ln c > 0$，这是一族以原点为圆心，半径为 $\sqrt{-\ln c}$ 的\textbf{同心圆}。
    当 $c=1$ 时，退化为原点 $(0,0)$。
    当 $c\leqslant  0$ 或 $c>1$ 时，等高线为空集。
  \end{itemize}
\end{enumerate}
\end{solution}

\begin{exercise}[基础：极坐标下的极限计算]
求下列极限：\\
(1)$\displaystyle\lim_{(x,y)\to(0,0)}\frac{x^3+y^3}{x^2+y^2}$;(2)$\displaystyle\lim_{(x,y)\to(0,0)}\frac{x^2\sin^2y}{x^2+2y^2}$;(3)$\displaystyle\lim_{(x,y)\to(0,0)}(x^2+y^2)\sin\frac1{x^2+y^2}$
\end{exercise}
\begin{solution}
统一令 $x=r\cos\theta,\ y=r\sin\theta$，则 $(x,y)\to(0,0)$ 等价于 $r\to0^+$。
\begin{enumerate}
  \item[(1)] 代入得
  \[
    \frac{x^3+y^3}{x^2+y^2}
    =r(\cos^3\theta+\sin^3\theta).
  \]
  由于 \(|\cos^3\theta+\sin^3\theta|\leqslant2\)，该式随 \(r\to0^+\) 趋于 \(0\)。

  \item[(2)] 代入并用 \(|\sin u|\leqslant |u|\)，有
  \[
    \left|\frac{x^2\sin^2y}{x^2+2y^2}\right|
    \leqslant
    \frac{r^2\cos^2\theta\cdot r^2\sin^2\theta}
         {r^2(\cos^2\theta+2\sin^2\theta)}
    \leqslant r^2\to0.
  \]
  故原极限为 \(0\)。

  \item[(3)] 代入后为
  \[
    r^2\sin\frac1{r^2},
  \]
  其绝对值不超过 \(r^2\)，故极限为 \(0\)。
\end{enumerate}
\end{solution}

\begin{exercise}[基础：极限存在性判定]
判断下列函数在 $(0,0)$ 处的极限是否存在。若不存在，给出证明：
\begin{shortexenum}
  \shortitem{1}{$f(x,y)=\dfrac{x^2-y^2}{x^2+y^2}$} &
  \shortitem{2}{$f(x,y)=\dfrac{x^3y}{x^6+y^2}$} \\
  \shortitem{3}{$f(x,y)=\dfrac{xy}{|x|+|y|}$} &
  \shortitem{4}{$f(x,y)=(x^2+y^2)\cos\dfrac{1}{xy}$（$xy \neq 0$）}
\end{shortexenum}
\end{exercise}
\begin{solution}
引入极坐标变换 $x = r\cos\theta, y = r\sin\theta$。当 $(x,y) \to (0,0)$ 时，等价于 $r \to 0^+$。
\begin{enumerate}
  \item[(1)] \textbf{不存在}。将极坐标代入得
  \[
  f(r\cos\theta,r\sin\theta)
  =\frac{r^2\cos^2\theta-r^2\sin^2\theta}{r^2\cos^2\theta+r^2\sin^2\theta}
  =\cos^2\theta-\sin^2\theta=\cos2\theta.
  \]
  当 $r \to 0^+$ 时，极限结果 $\cos 2\theta$ 仍然依赖于趋近原点时的角度 $\theta$。例如，当 $\theta = 0$ 时极限为 $1$；当 $\theta = \frac{\pi}{4}$ 时极限为 $0$。由于沿不同方向趋近原点得到的极限不同，故重极限不存在。
  \item[(2)] \textbf{不存在}。将极坐标代入得
  \[
  f(r\cos\theta,r\sin\theta)
  =\frac{(r\cos\theta)^3(r\sin\theta)}{(r\cos\theta)^6+(r\sin\theta)^2}
  =\frac{r^2\cos^3\theta\sin\theta}{r^4\cos^6\theta+\sin^2\theta}.
  \]
  若沿直线射线 $\theta = \text{const}$（且 $\sin\theta \neq 0$）趋近原点，$r \to 0^+$ 时极限为 $0$。
  但是，也可以让角度随 $r$ 变化。令 $\sin\theta = r^2\cos^3\theta$（对应直角坐标系下的 $y=x^3$），当 $r \to 0^+$ 时 $\theta \to 0$ 或 $\pi$，仍满足趋于原点的条件。代入上式得
  \[
  \lim_{r\to0^+}
  \frac{r^2\cos^3\theta(r^2\cos^3\theta)}
       {r^4\cos^6\theta+(r^2\cos^3\theta)^2}
  =\lim_{r\to0^+}\frac{r^4\cos^6\theta}{2r^4\cos^6\theta}
  =\frac12.
  \]
  由于这种角度变化方式给出的极限值为 $1/2$，而固定 \(\sin\theta\neq0\) 的角度时极限值为 $0$，故重极限不存在。
  \item[(3)] \textbf{存在，且为 $\boldsymbol{0}$}。将极坐标代入，并考察其绝对值。由于
  \(\displaystyle (|\cos\theta|+|\sin\theta|)^2
  =1+|\sin2\theta|\geqslant 1\)，所以 $|\cos\theta|+|\sin\theta|\geqslant 1$。因此
  \(\displaystyle |f(r\cos\theta,r\sin\theta)|=r\frac{|\cos\theta\sin\theta|}{|\cos\theta|+|\sin\theta|}\leqslant \frac12 r\)。当 $r \to 0^+$ 时，$\frac{1}{2}r \to 0$。由夹逼定理可知，原极限存在且等于 $0$。
  \item[(4)] \textbf{存在，且为 $\boldsymbol{0}$}。将极坐标代入，考察其绝对值：
  \(\displaystyle |f(r\cos\theta,r\sin\theta)|=\left|r^2\cos\left(\frac{1}{r^2\cos\theta\sin\theta}\right)\right|\leqslant r^2\)。当 $r \to 0^+$ 时，$r^2 \to 0$。由夹逼定理可知，原极限存在且等于 $0$。
\end{enumerate}
\end{solution}

\begin{exercise}[提高：累次极限与二重极限]
设 $f(x,y)=\dfrac{x^2-y^2}{x^2+y^2}$（$(x,y) \neq (0,0)$）。
\begin{shortexenum}
  \shortchoice{(1)}{求两个累次极限 $\lim_{y\to 0}\lim_{x\to 0}f(x,y)$ 和 $\lim_{x\to 0}\lim_{y\to 0}f(x,y)$；} &
  \shortchoice{(2)}{二重极限是否存在？说明理由；} \\
  \shortchoice{(3)}{若定义 $f(0,0)=0$，讨论 $f$ 的连续性。} &
\end{shortexenum}
\end{exercise}
\begin{solution}
\begin{enumerate}
  \item[(1)] \textbf{计算第一个累次极限}：先对 $x$ 求极限，此时将 $y$ 视为常数：\(\displaystyle
  \lim_{x\to 0} \frac{x^2-y^2}{x^2+y^2} = \frac{0-y^2}{0+y^2} = -1 \quad (y \neq 0)\)。
  再对 $y$ 求极限：$\displaystyle\lim_{y\to 0}(-1) = -1$。
  \textbf{计算第二个累次极限}：先对 $y$ 求极限，此时将 $x$ 视为常数：
  \(\displaystyle \lim_{y\to 0} \frac{x^2-y^2}{x^2+y^2} = \frac{x^2-0}{x^2+0} = 1 \quad (x \neq 0)\)。
  再对 $x$ 求极限：$\displaystyle\lim_{x\to 0}(1) = 1$。

  \item[(2)] \textbf{二重极限不存在}。一方面，根据定理，若二重极限存在，且两个累次极限也都存在，则三者必须相等；而此处两个累次极限分别为 $-1$ 与 $1$。另一方面，令 $x=r\cos\theta,\ y=r\sin\theta$，则
  \[
    f(r\cos\theta,r\sin\theta)=\cos2\theta,
  \]
  当 $r\to0^+$ 时结果仍随 $\theta$ 改变，也说明二重极限不存在。

  \item[(3)] \textbf{不连续}。因为二重极限 $\displaystyle\lim_{(x,y)\to(0,0)}f(x,y)$ 不存在，函数在原点处连拥有极限的前提都不满足，所以无论如何补充定义 $f(0,0)$ 的值，该函数在原点 $(0,0)$ 处都是不连续的。但在原点之外的区域 $\mathbb{R}^2 \setminus \{(0,0)\}$ 处，由于是有理函数且分母非零，它是连续的。
\end{enumerate}
\end{solution}

\begin{exercise}[提高：连续性综合]
\begin{shortsingleenum}
  \shortsingleitem{1}{设 $f(x,y)=\begin{cases}\dfrac{x^3+y^3}{x^2+y^2},&(x,y) \neq (0,0)\\0,&(x,y)=(0,0)\end{cases}$。证明 $f$ 在 $\mathbb{R}^2$ 上连续。}
  \shortsingleitem{2}{设 $f(x,y)=\begin{cases}\dfrac{xy}{\sqrt{x^2+y^2}},&(x,y) \neq (0,0)\\0,&(x,y)=(0,0)\end{cases}$。证明 $f$ 在 $(0,0)$ 处不可微（可微性定义将在下章给出，此处只需验证偏导数的存在性并思考为何 ``不可微''）。}
  \shortsingleitem{3}{设 $f:\mathbb{R}^n\to\mathbb{R}$ 满足 Lipschitz 条件：$|f(x)-f(y)|\leqslant  L\|x-y\|$ 对某个 $L>0$ 和所有 $x,y$ 成立。证明 $f$ 在 $\mathbb{R}^n$ 上一致连续。}
\end{shortsingleenum}
\end{exercise}
\begin{solution}
\begin{enumerate}
  \item[(1)] \textbf{证明在 $\mathbb{R}^2$ 上连续}：
  首先，在 $\mathbb{R}^2 \setminus \{(0,0)\}$ 上，函数是有理多项式的商，且分母不为零，由初等函数的连续性知其连续。
  接着，考察原点 $(0,0)$ 处的连续性。令 $x=r\cos\theta,\ y=r\sin\theta$，则
  \[
    f(r\cos\theta,r\sin\theta)=r(\cos^3\theta+\sin^3\theta),
  \]
  且其绝对值不超过 \(2r\)，故 \(\displaystyle\lim_{(x,y)\to(0,0)}f(x,y)=0\)。
  由于极限值 $0$ 刚好等于给定该点的函数值 $f(0,0)=0$，因此 $f$ 在 $(0,0)$ 处也是连续的。综上，$f$ 在 $\mathbb{R}^2$ 上处处连续。

  \item[(2)] \textbf{探讨偏导与不可微性}：
  对第(2)问，先考察函数 \(\displaystyle f(x,y)=\frac{xy}{\sqrt{x^2+y^2}}\) 在原点处的偏导数是否可计算。依据偏导数定义，\(\displaystyle
  f_x(0,0)=\lim_{h\to 0}\frac{f(h,0)-f(0,0)}{h}
  =\lim_{h\to 0}\frac{0-0}{h}=0\)。
  同理可得 $f_y(0,0)=0$，即偏导数都存在。
  如果要函数可微，则全增量与线性逼近之差必须是距离的高阶无穷小，即要求 \(\displaystyle
  \lim_{(x,y)\to(0,0)} \frac{f(x,y) - f(0,0) - (f_x(0,0)x + f_y(0,0)y)}{\sqrt{x^2+y^2}} = 0\)。
  代入已知值，上述极限式变为判断下式是否趋于 0：
  \(\displaystyle
  \lim_{(x,y)\to(0,0)} \frac{\frac{xy}{\sqrt{x^2+y^2}} - 0 - 0}{\sqrt{x^2+y^2}} = \lim_{(x,y)\to(0,0)} \frac{xy}{x^2+y^2}\)。
  对这个极限令 \(x=r\cos\theta,\ y=r\sin\theta\)，可得 \(\displaystyle \frac{xy}{x^2+y^2}=\frac12\sin2\theta\)，结果随角度改变，故二重极限不存在。
  因此误差项无法构成高阶无穷小，函数在 $(0,0)$ 处不可微。

  \item[(3)] \textbf{证明一致连续}：
  要证明 $f$ 在 $\mathbb{R}^n$ 上一致连续，我们需要证明对任意给定的 $\varepsilon>0$，存在一个不依赖于具体点位的 $\delta>0$，使得对所有的 $x,y \in \mathbb{R}^n$，只要 $\|x-y\|<\delta$，就有 $|f(x)-f(y)|<\varepsilon$。
  已知 $f$ 满足 Lipschitz 条件：$|f(x)-f(y)|\leqslant  L\|x-y\|$。
  取 $\delta = \dfrac{\varepsilon}{L}$（因为 $L>0$ 为常数，此 $\delta$ 只与 $\varepsilon$ 有关，对空间内所有点统一有效）。
  则当任意两点满足 $\|x-y\|<\delta$ 时，我们有 \(\displaystyle |f(x)-f(y)| \leqslant  L\|x-y\| < L \cdot \frac{\varepsilon}{L} = \varepsilon\)。
  根据一致连续性的定义，结论得证。
\end{enumerate}
\end{solution}

\subsubsection{A组}
\begin{exercise}[基本概念问答]
回答下列问题：
\begin{shortexenum}
  \shortitem{1}{二元函数的定义是什么？它的图像是怎样的一个点集？} &
  \shortitem{2}{什么是函数 $f(x,y)$ 的等高线？} \\
  \shortitem{3}{什么是函数 $f(x,y,z)$ 的等位面？} &
  \shortitem{4}{二元函数 $f(x,y)$ 在点 $(x_0,y_0)$ 处的极限怎样定义？} \\
  \shortitem{5}{二元函数 $f(x,y)$ 在点 $(x_0,y_0)$ 处连续怎样定义？} &
\end{shortexenum}
\end{exercise}
\begin{solution}
\begin{enumerate}
  \item[(1)] \textbf{二元函数的定义}：设 $D$ 为 $\mathbb{R}^2$ 中的一个非空平面点集。若按照某种对应法则 $f$，对于 $D$ 中的每一个点 $(x,y)$，都有唯一确定的实数值 $z$ 与之对应，则称 $f$ 为定义在 $D$ 上的二元函数，记为 $z=f(x,y)$。
  \textbf{图像点集}：其图像是三维空间 $\mathbb{R}^3$ 中的一个点集 $\Gamma = \{(x,y,z) \mid z=f(x,y), (x,y) \in D\}$，通常表示三维空间中的一张曲面。
  \item[(2)] \textbf{等高线}：设二元函数 $z=f(x,y)$，对于常数 $c$，使得 $f(x,y)=c$ 成立的平面点集 $\{(x,y) \in D \mid f(x,y)=c\}$ 称为函数 $f$ 的值为 $c$ 的等高线。它在几何上反映了使得函数值为常数的所有输入自变量的集合。
  \item[(3)] \textbf{等位面}：对于三元函数 $u=f(x,y,z)$，方程 $f(x,y,z)=c$（$c$ 为常数）在三维空间 $\mathbb{R}^3$ 中确定的曲面集合 $\{(x,y,z) \in \Omega \mid f(x,y,z)=c\}$ 称为该函数的等位面（或等值面）。
  \item[(4)] \textbf{极限的定义}：设 $f(x,y)$ 在点 $p_0=(x_0,y_0)$ 的某去心邻域内有定义。若存在常数 $A$，对任意 $\varepsilon>0$，总存在 $\delta>0$，使得当 $0 < \sqrt{(x-x_0)^2+(y-y_0)^2} < \delta$ 时，均有 $|f(x,y)-A| < \varepsilon$，则称常数 $A$ 为函数当 $(x,y) \to (x_0,y_0)$ 时的二重极限。
  \item[(5)] \textbf{连续的定义}：若点 $(x_0,y_0)$ 是函数的定义域内的一点，并且极限 $\displaystyle\lim_{(x,y)\to(x_0,y_0)}f(x,y)$ 存在，且该极限值正好等于函数在该点的取值 $f(x_0,y_0)$，即 \(\displaystyle \lim_{(x,y)\to(x_0,y_0)}f(x,y) = f(x_0,y_0)\)，则称函数 $f(x,y)$ 在点 $(x_0,y_0)$ 处连续。
\end{enumerate}
\end{solution}

\begin{exercise}[求定义域]
确定并画出下列二元函数的定义域：\\
\begin{tabular}{lll}
  (1) $z = x + \sqrt{y}$ &
  (2) $z = \sqrt{1-x^2} + \sqrt{y^2-1}$ &
  (3) $z = \sqrt{1-x^2-y^2}$ \\
  (4) $z = \dfrac{1}{\sqrt{x^2+y^2-1}}$ &
  (5) $z = \ln(-x-y)$ &
  (6) $z = \arcsin\dfrac{y}{x}$
\end{tabular}
\end{exercise}
\begin{solution}
\begin{enumerate}
  \item[(1)] 根号下要求非负：$y \geqslant 0$。而 $x$ 可以取任意实数。
  \textbf{图形}：这是 $x$ 轴以及 $x$ 轴上方的全体平面区域（上半平面）。
  \item[(2)] 分别要求根号下非负：$1-x^2 \geqslant 0$ 且 $y^2-1 \geqslant 0$。解得 $|x| \leqslant  1$ 且 $|y| \geqslant 1$。
  \textbf{图形}：由直线 $x=1, x=-1$ 夹住的无穷带，去掉位于直线 $y=1$ 和 $y=-1$ 之间的中心水平部分。呈现为两条互不连通的垂直宽带状半闭区域。
  \item[(3)] 根号下非负：$1-x^2-y^2 \geqslant 0 \implies x^2+y^2 \leqslant  1$。
  \textbf{图形}：以原点为圆心，半径为 1 的单位闭圆盘区域（包含圆周边界）。
  \item[(4)] 分母中不能为零且根号内非负：$x^2+y^2-1 > 0 \implies x^2+y^2 > 1$。
  \textbf{图形}：以原点为圆心，半径为 1 的单位圆外部全体区域（不含圆周边界）。
  \item[(5)] 对数函数的真数必须为正：$-x-y > 0 \implies x+y < 0$。
  \textbf{图形}：经过原点的直线 $y = -x$ 的左下方完全开放半平面区域（不含直线 $y=-x$）。
  \item[(6)] 反正弦函数的自变量在 $[-1, 1]$ 之间，同时分母不能为零：
  $-1 \leqslant  \dfrac{y}{x} \leqslant  1$ 且 $x \neq 0$。
  若 $x>0$，解得 $-x \leqslant  y \leqslant  x$；若 $x<0$，解得 $x \leqslant  y \leqslant  -x$。综合可写为 $|y| \leqslant  |x|,\; x \neq 0$。
  \textbf{图形}：位于直线 $y=x$ 和 $y=-x$ 之间，并且包含 $x$ 轴但不包含 $y$ 轴中心的两个无限扇形区域（开口在左右方向），去掉 $y$ 轴这一整条直线。
\end{enumerate}
\end{solution}

\begin{exercise}[齐次函数性质]
设函数 $f(x,y) = x^2 + y^2 - xy\arctan\dfrac{x}{y}$，求 $f(tx,ty)$。
\end{exercise}
\begin{solution}
\begin{enumerate}
  \item \textbf{代入新变量}：将 $tx$ 和 $ty$ 代替原函数中的 $x$ 和 $y$：\(\displaystyle f(tx,ty) = (tx)^2 + (ty)^2 - (tx)(ty)\arctan\frac{tx}{ty}\)。
  \item \textbf{化简代数项与超越项}：
  \begin{itemize}
    \item 多项式项化简：$(tx)^2 = t^2 x^2$，$(ty)^2 = t^2 y^2$，乘积项 $(tx)(ty) = t^2 xy$。
    \item 反三角项化简：由于 $t$ 从分子分母上相约（设 $t \neq 0$），有 $\arctan\frac{tx}{ty} = \arctan\frac{x}{y}$。
  \end{itemize}
  \item \textbf{提取公因式}：将上述化简结果统一提出 $t^2$，得 \(\displaystyle f(tx,ty)=t^2\left( x^2+y^2-xy\arctan\frac{x}{y}\right)\)。
  \item \textbf{得出结果}：括号内的正是原函数本身，因此 \(\displaystyle f(tx,ty) = t^2 f(x,y)\)。
  这说明该函数是一个 $2$ 次齐次函数。
\end{enumerate}
\end{solution}

\begin{exercise}[二元函数的辨析]
问下列表达式是否是 $a,b$ 的二元函数：\\
(1) $I = \int_0^1 (a+bx)^2 \mathrm{d}x$~~~(2) $I = \int_a^b (1+x)^2 \mathrm{d}x$~~~(3)$I = \begin{cases} 1, & a>b, \\ 0, & a=b, \\ -1, & a<b. \end{cases}$
\end{exercise}
\begin{solution}
\begin{enumerate}
  \item[(1)] \textbf{是}。这里的积分变量是 $x$，而 $a,b$ 只是参数。
  积分做完以后，$x$ 被消去，剩下的是只含 $a,b$ 的数值，例如 $I=\int_0^1(a+bx)^2\,\mathrm{d}x$。对每一组确定的 $(a,b)$ 都给出唯一结果，所以它是二元函数。
  \item[(2)] \textbf{是}。虽然被积函数只显含 $x$，但上下限本身由 $a,b$ 决定。
  按牛顿-莱布尼茨公式，$I=F(b)-F(a)$，其中 $F(x)=\frac{(1+x)^3}{3}$，所以最终结果仍是关于 $(a,b)$ 的唯一对应法则。
  \item[(3)] \textbf{是}。这个表达式实际上按 $a-b$ 的符号分段定义。
  对任意给定的 $(a,b)$，三种情形 $a>b,\ a=b,\ a<b$ 中必恰有一种发生，因此 $I$ 的值也唯一确定。
  从映射角度看，它就是二维情形下的 $\operatorname{sgn}(a-b)$。
\end{enumerate}
\end{solution}

\begin{exercise}[构造函数表达式]
根据下列对应关系，写出函数表达式：
\begin{shortexenum}
  \shortchoice{(1)}{$\forall (x,y)\in\mathbb{R}^2$，用该点到原点的距离与之对应；} &
  \shortchoice{(2)}{$\forall (x,y)\in\mathbb{R}^2$，用该点的坐标之和与之对应。}
\end{shortexenum}
\end{exercise}
\begin{solution}
\begin{enumerate}
  \item[(1)] 点 $(x,y)$ 到原点 $(0,0)$ 的距离为 $\sqrt{(x-0)^2+(y-0)^2}$。因此函数表达式为 $f(x,y) = \sqrt{x^2+y^2}$。
  \item[(2)] “坐标之和”即横坐标与纵坐标的代数加和运算。因此直接写出表达式为 $f(x,y) = x+y$。
\end{enumerate}
\end{solution}

\begin{exercise}[多元函数极限计算]
求下列函数极限：\\
(1) $\lim\limits_{\substack{x\to 0 \\ y\to 1}} \dfrac{1-xy}{x^2+y^2}$~(2)$\lim\limits_{\substack{x\to 2 \\ y\to 0}} \dfrac{1}{x^2+y^2}$~(3)$\lim\limits_{(x,y)\to(0,0)} \dfrac{\sin(xy)}{x}$~(4)$\lim\limits_{\substack{x\to +\infty \\ y\to +\infty}} (x^2+y^2)e^{-(x+y)}$
\end{exercise}
\begin{solution}
\begin{enumerate}
  \item[(1)] \textbf{解析}：以目标点 $(0,1)$ 为中心写 \(x=r\cos\theta,\ y=1+r\sin\theta\)。当 \(r\to0^+\) 时，
  \[
    \frac{1-xy}{x^2+y^2}
    =\frac{1-r\cos\theta(1+r\sin\theta)}
           {r^2\cos^2\theta+(1+r\sin\theta)^2}
    \to1.
  \]

  \item[(2)] \textbf{解析}：以目标点 $(2,0)$ 为中心写 \(x=2+r\cos\theta,\ y=r\sin\theta\)。于是
  \[
    \frac{1}{x^2+y^2}
    =\frac{1}{(2+r\cos\theta)^2+r^2\sin^2\theta}
    \to\frac14.
  \]

  \item[(3)] \textbf{解析}：令 \(x=r\cos\theta,\ y=r\sin\theta\)。在定义域 $x\neq0$ 内，由 \(|\sin u|\leqslant|u|\) 得
  \[
    \left|\frac{\sin(xy)}{x}\right|
    \leqslant\frac{|xy|}{|x|}
    =|y|
    =r|\sin\theta|\leqslant r\to0.
  \]
  故该极限为 \(0\)。

  \item[(4)] \textbf{解析}：令 \(x=r\cos\theta,\ y=r\sin\theta\)，其中 \(\theta\in(0,\frac{\pi}{2})\)，且 \(r\to+\infty\)。此时
  \[
    (x^2+y^2)e^{-(x+y)}
    =r^2e^{-r(\cos\theta+\sin\theta)}.
  \]
  在第一象限内 \(\cos\theta+\sin\theta\geqslant1\)，故
  \[
    0\leqslant r^2e^{-r(\cos\theta+\sin\theta)}
    \leqslant r^2e^{-r}\to0.
  \]
  原极限为 \(0\)。
\end{enumerate}
\end{solution}

\begin{exercise}[连续性判定]
证明函数 $f(x,y) = \begin{cases} \dfrac{x^2y}{x^4+y^2}, & x^2+y^2\neq 0, \\ 0, & x^2+y^2=0 \end{cases}$ 在原点 $(0,0)$ 处不连续。
\end{exercise}
\begin{solution}
\begin{enumerate}
  \item \textbf{回顾连续性条件}：
  要使 $f(x,y)$ 在 $(0,0)$ 处连续，必须满足 \(\displaystyle \lim_{(x,y)\to(0,0)} f(x,y) = f(0,0) = 0\)。
  \item \textbf{极坐标下考察}：令 $x=r\cos\theta,\ y=r\sin\theta$，则
  \[
    f(r\cos\theta,r\sin\theta)
    =\frac{r\cos^2\theta\sin\theta}{r^2\cos^4\theta+\sin^2\theta}.
  \]
  为了观察分子分母同阶的情形，可取随 $r$ 变化的角度，使 \(\sin\theta=r\cos^2\theta\)（对应直角坐标下 \(y=x^2\)）。代入上式得
  \[
    f(r\cos\theta,r\sin\theta)=\frac12.
  \]
  \item \textbf{对比结论得出判定}：沿上述趋近方式函数值趋于 \(\frac12\)，不等于 \(f(0,0)=0\)。因此函数在原点 $(0,0)$ 处不连续。
\end{enumerate}
\end{solution}

\begin{exercise}[求间断点集]
求下列函数的间断点集. 设 $f(0,0)=0$。\\
(1)$f(x,y) = \dfrac{1}{\sqrt{x^2+y^2}}$~~~~(2)$f(x,y) = \dfrac{xy}{x^2+y^2}$~~~~(3)$f(x,y) = \dfrac{x+y}{x^3+y^3}$
\end{exercise}
\begin{solution}
\begin{enumerate}
  \item[(1)] \textbf{考察分式结构}：该函数在定义域内为有理式和无理式的复合，对于分母非零处均连续。
  分母在 $(0,0)$ 处值为 0。令 $x=r\cos\theta,\ y=r\sin\theta$，则 \(f=\frac1r\)，当 \(r\to0^+\) 时无界，故原点处不连续。
  故在 $(0,0)$ 处不连续，间断点集仅为此一个孤立点：$\{(0,0)\}$。
  \item[(2)] \textbf{考察分母结构与原点极限}：函数在分母非零的区域处处连续。分母 $x^2+y^2=0$ 的唯一解即原点 $(0,0)$。
  极坐标下有 \(\displaystyle \frac{xy}{x^2+y^2}=\frac12\sin2\theta\)，当 \(r\to0^+\) 时结果仍随 \(\theta\) 改变，说明极限本身不存在。
  因为重极限不存在，尽管人为规定了 $f(0,0)=0$，这仍然是一个间断点。间断点集为：$\{(0,0)\}$。
  \item[(3)] \textbf{考察分母为零的所有位置}：该函数在分母 $x^3+y^3 \neq 0$ 处连续。
  令分母为 0 寻找可能的所有奇点：\(\displaystyle x^3+y^3 = 0 \implies (x+y)(x^2-xy+y^2) = 0\)。
  这在实数域意味着 $x+y=0$，即直线 $y=-x$。
  虽然在此直线上，似乎有 $x+y$ 构成了 $\frac{0}{0}$ 不定式（分子也是 $x+y$），但经过约分（当 $x+y \neq 0$ 时）有 $f(x,y) = \frac{1}{x^2-xy+y^2}$。
  当点 $(x,y)$ 趋于直线 $y=-x$ 上的某一点 $(x_0, -x_0)$（其中 $x_0 \neq 0$）时，约分后的分母为 $3x_0^2 \neq 0$。函数虽然存在极限 $\frac{1}{3x_0^2}$ 但这在原式中根本无定义，更别提这本身也是原式的不可去代数奇点位置的瑕点，所以是不连续点。
  此外考察原点：约分后在极坐标下为
  \[
    \frac{1}{x^2-xy+y^2}
    =\frac{1}{r^2(\cos^2\theta-\cos\theta\sin\theta+\sin^2\theta)}.
  \]
  其中角度因子恒为正且不可能为 \(0\)，故 \(r\to0^+\) 时函数无界，原点依然是不连续点（且已经包含在直线方程 $y=-x$ 内）。
  因此间断点集为空间中一整条直线上的所有点：$\{(x,y) \mid y = -x\}$。
\end{enumerate}
\end{solution}

\begin{exercise}[实际问题中的极限]
当 $x$ 摩尔硫酸和 $y$ 摩尔水混合时，产生的热量为 \(\displaystyle Q(x,y)=\dfrac{17.86xy}{1.798x+y}\)，其中 $x>0,\ y>0$。
试问 $Q$ 在 $(0,0)$ 点的极限是否存在？若存在，试求出极限值。
\end{exercise}
\begin{solution}
\begin{enumerate}
  \item \textbf{分析问题条件及函数结构}：限定条件为第一象限 $x>0, y>0$。令 \(x=r\cos\theta,\ y=r\sin\theta\)，其中 \(0<\theta<\frac{\pi}{2}\)，且 \(r\to0^+\)。
  \item \textbf{代入极坐标}：
  \[
    Q(r\cos\theta,r\sin\theta)
    =\frac{17.86r^2\cos\theta\sin\theta}
           {r(1.798\cos\theta+\sin\theta)}
    =r\frac{17.86\cos\theta\sin\theta}
           {1.798\cos\theta+\sin\theta}.
  \]
  在第一象限内 \(1.798\cos\theta+\sin\theta\geqslant1\)，且 \(\cos\theta\sin\theta\leqslant\frac12\)，所以 \(0<Q(x,y)\leqslant 8.93r\to0\)。
  \item \textbf{得出结论}：极限存在，并且产生热量在无限接近混合物均无的情况时趋于零：
  \[
    \boxed{\displaystyle\lim_{\substack{x\to 0^+ \\ y\to 0^+}} Q(x,y) = 0}.
  \]
\end{enumerate}
\end{solution}

\subsubsection{B组}

\begin{exercise}[绘制等高线]
画出下列函数的等高线，并写出等高线方程：\\
(1) $z = x+y$~~(2)$z = |x|+y$~~(3)$z = x^2+y^2$~~(4)$z = \dfrac{1}{x^2+2y^2}$
\end{exercise}
\begin{solution}
\begin{enumerate}
  \item[(1)] \textbf{等高线方程}：令 $z=c$，得到 $x+y=c \implies y = -x+c$。
  \textbf{图形描绘}：这是一族斜率为 $-1$ 的平行直线。不同 $c$ 对应不同的截距。函数图像为空间的一个倾斜平面。
  \item[(2)] \textbf{等高线方程}：令 $z=c$，得到 $|x|+y=c \implies y = -|x|+c$。
  \textbf{图形描绘}：对于每个特定的 $c$，这是一条由两段射线拼成的顶点在 $y$ 轴 $(0,c)$ 上、开口向下成直角的折线（类似于倒过来的字母“V”）。
  \item[(3)] \textbf{等高线方程}：令 $z=c$，得到 $x^2+y^2=c$。
  \textbf{图形描绘}：只有当 $c\geqslant 0$ 时有意义。当 $c>0$ 时，这是一族以原点为圆心、半径为 $\sqrt{c}$ 的同心圆；当 $c=0$ 时退化为一个原点 $(0,0)$。
  \item[(4)] \textbf{等高线方程}：令 $z=c$，考虑到分母形式必然有 $c>0$，得到 $\dfrac{1}{x^2+2y^2} = c \implies x^2+2y^2 = \dfrac{1}{c}$。
  整理为标准形式：$\dfrac{x^2}{(1/\sqrt{c})^2} + \dfrac{y^2}{(1/\sqrt{2c})^2} = 1$。
  \textbf{图形描绘}：这是一族以原点为中心的同心椭圆族，其长半轴沿 $x$ 轴方向（半径为 $1/\sqrt{c}$），短半轴沿 $y$ 轴方向（半径为 $1/\sqrt{2c}$）。常数 $c$ 越大，表示的高程越高，椭圆等高线越向内缩小。
\end{enumerate}
\end{solution}

\begin{exercise}[求解等位面]
求下列函数的等位面并写出其方程：
\begin{shortexenum}
  \shortitem{1}{$u = x+y+z$} &
  \shortitem{2}{$f(x,y,z) = x^2+y^2+z^2$} \\
  \shortitem{3}{$g(x,y,z) = y-z$} &
  \shortitem{4}{$u = \operatorname{sgn}\sin(x^2+y^2+z^2)$}
\end{shortexenum}
\end{exercise}
\begin{solution}
\begin{enumerate}
  \item[(1)] \textbf{等位面方程}：$x+y+z=c$。
  \textbf{几何描述}：这是一族三维空间中互相平行的平面，其共同的法向量为 $(1,1,1)$。
  \item[(2)] \textbf{等位面方程}：$x^2+y^2+z^2=c$。
  \textbf{几何描述}：当 $c>0$ 时，这是以原点为球心、半径为 $\sqrt{c}$ 的同心球面族。$c=0$ 时退化为一点，$c<0$ 时为空集。
  \item[(3)] \textbf{等位面方程}：$y-z=c$。
  \textbf{几何描述}：这是一族平行平面。由于方程中缺少 $x$，说明该平面在 $x$ 方向是无限延伸无变化的，该族平面平行于 $x$ 轴。法向量为 $(0,1,-1)$。
  \item[(4)] \textbf{等位面分析}：符号函数 $\operatorname{sgn}(\cdot)$ 只能取三个离散值：$1, 0, -1$。因此只存在三种等位面。令 $r^2 = x^2+y^2+z^2$。
  \begin{itemize}
    \item \textbf{当 $c=1$ 时}：要求 $\sin(r^2) > 0$。即 $2k\pi < r^2 < (2k+1)\pi$（$k=0,1,2,\dots$）。这对应空间中无穷多个嵌套的、有厚度的球壳形空心区域。
    \item \textbf{当 $c=0$ 时}：要求 $\sin(r^2) = 0$。即 $r^2 = k\pi$（$k=0,1,2,\dots$）。这对应一系列离散的无穷多层同心球面。
    \item \textbf{当 $c=-1$ 时}：要求 $\sin(r^2) < 0$。即 $(2k-1)\pi < r^2 < 2k\pi$（$k=1,2,\dots$）。这对应夹在 $c=1$ 区域之间的其它球壳形空心区域。
  \end{itemize}
\end{enumerate}
\end{solution}

\begin{exercise}[三元函数概念推广]
试给出三元函数 $u=f(x,y,z)$ 在点 $(x_0,y_0,z_0)$ 的极限定义及连续性定义。
\end{exercise}
\begin{solution}
\begin{enumerate}
  \item \textbf{三元函数极限的定义（$\boldsymbol{\varepsilon-\delta}$ 语言）}：
  设三元函数 $f(x,y,z)$ 在点 $P_0(x_0,y_0,z_0)$ 的某去心邻域内有定义。如果存在一个实常数 $A$，使得对于任意给定的正数 $\varepsilon > 0$，总存在相应的正数 $\delta > 0$，使得当空间点 $P(x,y,z)$ 在函数定义域内且满足
  \(\displaystyle 0 < \rho(P,P_0) = \sqrt{(x-x_0)^2+(y-y_0)^2+(z-z_0)^2} < \delta\)
  时，都有 \(\displaystyle |f(x,y,z) - A| < \varepsilon\)，则称常数 $A$ 为函数 $f(x,y,z)$ 当 $(x,y,z) \to (x_0,y_0,z_0)$ 时的三重极限，记作
  \(\displaystyle \lim_{(x,y,z)\to(x_0,y_0,z_0)}f(x,y,z) = A\)。
  \item \textbf{三元函数连续性的定义}：
  如果三元函数 $f(x,y,z)$ 在点 $P_0(x_0,y_0,z_0)$ 处有定义，并且其在此处的极限存在且等于该点处的函数值，即满足 $\displaystyle\lim_{(x,y,z)\to(x_0,y_0,z_0)}f(x,y,z) = f(x_0,y_0,z_0)$，
  则称函数 $f(x,y,z)$ 在点 $(x_0,y_0,z_0)$ 处是连续的。
\end{enumerate}
\end{solution}

\begin{exercise}[证明极限不存在]
证明下列极限不存在：(1) $\lim\limits_{\substack{x\to 0 \\ y\to 0}} \dfrac{y}{x}$~~~(2) $\lim\limits_{\substack{x\to 0 \\ y\to 0}} \dfrac{x^2-y^2}{x^2+y^2}$
\end{exercise}
\begin{solution}
引入极坐标变换 $x = r\cos\theta, y = r\sin\theta$。当 $(x,y) \to (0,0)$ 时，等价于 $r \to 0^+$。
\begin{enumerate}
  \item[(1)] \textbf{证明极限 (1) 不存在}：
  将极坐标代入函数中得
  \(\displaystyle f(r\cos\theta, r\sin\theta) = \frac{r\sin\theta}{r\cos\theta} = \tan\theta \quad (\cos\theta \neq 0)\)。
  当 $r \to 0^+$ 时，极限结果为 $\tan\theta$。由于该极限值依赖于趋近原点时的角度 $\theta$，并非一个确定的常数。例如：
  \begin{itemize}
    \item 当沿射线 $\theta = \frac{\pi}{4}$ 趋近原点时，极限为 $\tan\frac{\pi}{4} = 1$。
    \item 当沿射线 $\theta = 0$ 趋近原点时，极限为 $\tan 0 = 0$。
  \end{itemize}
  因为沿不同角度趋近原点得到的极限值不一致，所以该重极限不存在。
  \item[(2)] \textbf{证明极限 (2) 不存在}：
  将极坐标代入函数中得 \(\displaystyle
  f(r\cos\theta, r\sin\theta)
  = \frac{r^2\cos^2\theta - r^2\sin^2\theta}{r^2\cos^2\theta + r^2\sin^2\theta}
  = \cos 2\theta\)。
  当 $r \to 0^+$ 时，极限结果为 $\cos 2\theta$。该结果同样依赖于趋近原点时的角度 $\theta$。例如：
  \begin{itemize}
    \item 当沿射线 $\theta = 0$ 趋近原点时（沿 $x$ 轴），极限为 $\cos 0 = 1$。
    \item 当沿射线 $\theta = \frac{\pi}{2}$ 趋近原点时（沿 $y$ 轴），极限为 $\cos\pi = -1$。
  \end{itemize}
  由于沿不同方向趋近原点得到的极限值不相等，因此该二重极限不存在。
\end{enumerate}
\end{solution}

% ============================================================
% §7.3 偏导数与全微分
% ============================================================

\subsection{第三节习题：偏导数与全微分}

\begin{exercise}[方向导数与偏导数的计算]
\begin{shortexenum}
  \shortchoice{(1)}{设 $f(x,y)=x^2e^{-y}+y\ln x$，求 $\dfrac{\partial f}{\partial x}(1,0)$ 和 $\dfrac{\partial f}{\partial y}(1,0)$；} &
  \shortchoice{(2)}{设 $f(x,y,z)=\dfrac{x}{y}+\dfrac{y}{z}+\dfrac{z}{x}$，求所有一阶偏导函数；} \\
  \shortchoice{(3)}{设 $f(x,y)=\arctan(xy)$，求 $f$ 在 $(1,1)$ 处沿方向 $\vec l=\bigl(\frac35,\frac45\bigr)$ 的方向导数。} &
\end{shortexenum}
\end{exercise}
\begin{solution}
\begin{enumerate}
  \item[(1)] \textbf{解析}：先求出一般的偏导函数，然后代入点 $(1,0)$。
  对 $x$ 求偏导（视 $y$ 为常数）：
  \(\displaystyle \frac{\partial f}{\partial x} = 2xe^{-y} + \frac{y}{x}\)。
  代入 $(1,0)$ 得 $\dfrac{\partial f}{\partial x}(1,0) = 2(1)e^0 + 0 = 2$。
  对 $y$ 求偏导（视 $x$ 为常数）：
  \(\displaystyle \frac{\partial f}{\partial y} = -x^2e^{-y} + \ln x\)。
  代入 $(1,0)$ 得 $\dfrac{\partial f}{\partial y}(1,0) = -(1)e^0 + \ln 1 = -1$。

  \item[(2)] \textbf{解析}：分别对 $x, y, z$ 求导，其余变量视作常数。
  \(\displaystyle f_x(x,y,z)= \frac{1}{y} - \frac{z}{x^2},\qquad
  f_y(x,y,z)= -\frac{x}{y^2} + \frac{1}{z},\qquad
  f_z(x,y,z)= -\frac{y}{z^2} + \frac{1}{x}\)。

  \item[(3)] \textbf{解析}：先求梯度，再与方向向量做点积。因 $\vec l$ 已是单位向量（$\sqrt{(3/5)^2+(4/5)^2}=1$），有
  \(\displaystyle f_x = \frac{y}{1+(xy)^2}, \qquad
  f_y = \frac{x}{1+(xy)^2}, \qquad
  \nabla f(1,1) = \left(\frac{1}{2}, \frac{1}{2}\right)\)，从而
  \(\displaystyle \frac{\partial f}{\partial \vec l}(1,1)
  = \nabla f(1,1) \cdot \vec l
  = \frac{1}{2} \cdot \frac{3}{5} + \frac{1}{2} \cdot \frac{4}{5}
  = \frac{7}{10}\)。
\end{enumerate}
\end{solution}

\begin{exercise}[可微性判定]
判断下列函数在 $(0,0)$ 处是否可微。若是，写出全微分：
\begin{shortsingleenum}
  \shortsinglechoice{(1)}{$f(x,y)=\sqrt{|xy|}$；}
  \shortsinglechoice{(2)}{$f(x,y)=\begin{cases} \dfrac{x^3y}{x^6+y^2},&(x,y) \neq (0,0),\\[8pt] 0,&(x,y)=(0,0); \end{cases}$}
  \shortsinglechoice{(3)}{$f(x,y)=(x^2+y^2)\cos\dfrac{1}{\sqrt{x^2+y^2}}$（补充 $f(0,0)=0$）。}
\end{shortsingleenum}
\end{exercise}
\begin{solution}
\begin{enumerate}
  \item[(1)] \textbf{解析}：
  先考察 $(0,0)$ 处的偏导数：$f_x(0,0)=\lim\limits_{x \to 0} \dfrac{f(x,0)-0}{x}=0$，且 $f_y(0,0)=\lim\limits_{y \to 0} \dfrac{f(0,y)-0}{y}=0$。
  若可微，则增量公式中的余项极限需为零。验证余项：\(\displaystyle \lim_{(x,y) \to (0,0)} \frac{f(x,y) - f(0,0) - (0 \cdot x + 0 \cdot y)}{\sqrt{x^2+y^2}} = \lim_{(x,y) \to (0,0)} \frac{\sqrt{|xy|}}{\sqrt{x^2+y^2}}\)。
  令 \(x=r\cos\theta,\ y=r\sin\theta\)，该比值化为 \(\sqrt{|\cos\theta\sin\theta|}\)。取固定角度 \(\theta=\frac{\pi}{4}\) 时得到 \(\frac{1}{\sqrt2}\neq0\)，故余项极限不为 \(0\)。因此 $f$ 在 $(0,0)$ \textbf{不可微}。

  \item[(2)] \textbf{解析}：
  同样先求 $(0,0)$ 处偏导数：$f_x(0,0)=\lim\limits_{x \to 0} \dfrac{f(x,0)-0}{x}=0$，且 $f_y(0,0)=\lim\limits_{y \to 0} \dfrac{f(0,y)-0}{y}=0$。
  验证余项极限：\(\displaystyle \lim_{(x,y) \to (0,0)} \frac{x^3y}{(x^6+y^2)\sqrt{x^2+y^2}}\)。
  令 \(x=r\cos\theta,\ y=r\sin\theta\)，该比值为
  \[
    \frac{r\cos^3\theta\sin\theta}{r^4\cos^6\theta+\sin^2\theta}.
  \]
  令 \(\sin\theta=r^2\cos^3\theta\)（对应 \(y=x^3\)），则上式等于 \(\dfrac{1}{2r}\)，当 \(r\to0^+\) 时无界。故余项极限不存在且不为零，$f$ 在 $(0,0)$ \textbf{不可微}。

  \item[(3)] \textbf{解析}：
  求 $(0,0)$ 处偏导数：$f_x(0,0) = \lim_{x \to 0} x \cos(1/|x|) = 0$（有界函数乘无穷小），同理 $f_y(0,0) = 0$。
  验证可微性的余项极限，令 $\rho = \sqrt{x^2+y^2} \to 0$，则 \(\displaystyle \frac{f(x,y)-0-0}{\rho}=\rho\cos(1/\rho)\to0\)。
  完全符合可微的定义要求，故 $f$ 在 $(0,0)$ \textbf{可微}，其全微分为 $\mathrm{d}f(0,0) = 0\,\mathrm{d}x + 0\,\mathrm{d}y = \boxed{0}$。
\end{enumerate}
\end{solution}

\begin{exercise}[偏导数连续性的分析]
设 $f(x,y)=\begin{cases} \dfrac{x^3}{x^2+y^2},&(x,y) \neq (0,0),\\[8pt] 0,&(x,y)=(0,0). \end{cases}$
\begin{shortexenum}
  \shortitem{1}{求 $f(x,y)$ 当 $(x,y) \neq (0,0)$ 时的偏导函数 $f_x(x,y)$ 和 $f_y(x,y)$；} &
  \shortitem{2}{讨论 $f_x$ 和 $f_y$ 在 $(0,0)$ 处的存在性和连续性；} \\
  \shortitem{3}{$f$ 在 $(0,0)$ 处是否可微？} &
\end{shortexenum}
\end{exercise}
\begin{solution}
\begin{enumerate}
  \item[(1)] \textbf{解析}：
  利用除法求导法则（$(x,y) \neq (0,0)$）：
  \begin{align*}
  f_x(x,y) &= \frac{3x^2(x^2+y^2) - x^3(2x)}{(x^2+y^2)^2} = \frac{x^4 + 3x^2y^2}{(x^2+y^2)^2},\\
  f_y(x,y) &= x^3 \cdot \left(-\frac{2y}{(x^2+y^2)^2}\right) = -\frac{2x^3y}{(x^2+y^2)^2}.
  \end{align*}
  \item[(2)] \textbf{解析}：
  先考察在 $(0,0)$ 处的存在性（用定义）：$f_x(0,0) = \lim_{x \to 0} \dfrac{f(x,0)-0}{x} = 1,\; f_y(0,0) = \lim_{y \to 0} \dfrac{f(0,y)-0}{y} = 0$。
  再考察连续性。令 $x = r\cos\theta, y = r\sin\theta$，则
  \(\displaystyle f_x(r\cos\theta, r\sin\theta)=\cos^4\theta + 3\cos^2\theta\sin^2\theta,\qquad
  f_y(r\cos\theta, r\sin\theta)=-2\cos^3\theta\sin\theta\)。
  当 $r \to 0^+$ 时，两式的极限都依赖于角度 $\theta$：前者在 $\theta=0$ 时为 $1$、在 $\theta=\frac{\pi}{2}$ 时为 $0$；后者在 $\theta=0$ 时为 $0$、在 $\theta=\frac{\pi}{4}$ 时为 $-\frac{1}{2}$。因此 $\lim\limits_{(x,y)\to(0,0)} f_x(x,y)$ 与 $\lim\limits_{(x,y)\to(0,0)} f_y(x,y)$ 都不存在，所以 $f_x,f_y$ 在 $(0,0)$ 处均不连续。
  \item[(3)] \textbf{解析}：
  即便偏导数不连续，我们仍需通过定义验证可微性。考察余项极限：
  \[
  \lim_{(x,y) \to (0,0)} \frac{f(x,y) - f(0,0) - (1 \cdot x + 0 \cdot y)}{\sqrt{x^2+y^2}}
  = \lim_{(x,y) \to (0,0)} \frac{\frac{x^3}{x^2+y^2} - x}{\sqrt{x^2+y^2}} = \lim_{(x,y) \to (0,0)} \frac{-xy^2}{(x^2+y^2)^{3/2}}.
  \]
  令 \(x=r\cos\theta,\ y=r\sin\theta\)，上式化为 \(-\cos\theta\sin^2\theta\)。取固定角度 \(\theta=\frac{\pi}{4}\)，得到 \(-\frac{1}{2\sqrt2}\neq0\)。因此 $f$ 在 $(0,0)$ \textbf{不可微}。
\end{enumerate}
\end{solution}

\begin{exercise}[方向导数与梯度的关系预告]
虽然本节尚未正式引入梯度算子 $\nabla$，但你已经具备了足够的知识来猜测以下结论：

\textbf{猜想}：若 $f$ 在 $p_0$ 处可微，则沿任意单位方向 $\vec l=(l_1,l_2)$ 的方向导数为
\(\displaystyle \frac{\partial f}{\partial\vec l}(p_0)=f_x(p_0)\,l_1+f_y(p_0)\,l_2 = \nabla f(p_0)\cdot\vec l\)。
请利用可微的定义式 ($\star$) 直接证明这一公式。
进一步思考：若 $\nabla f(p_0) \neq \boldsymbol 0$，沿哪个方向的方向导数最大？最大值是多少？这揭示了什么几何事实？
\end{exercise}
\begin{solution}
\begin{enumerate}
  \item \textbf{证明猜想公式}：
  由 $f$ 在 $p_0$ 处可微的定义，对增量 $\boldsymbol h = t\vec l = (tl_1, tl_2)$ 有
  \[
  f(p_0 + t\vec l) - f(p_0)
  = f_x(p_0)(tl_1) + f_y(p_0)(tl_2) + o(\|t\vec l\|)
  = t\bigl(f_x(p_0)l_1 + f_y(p_0)l_2\bigr) + o(|t|).
  \]
  其中 $\|t\vec l\| = |t|$。两边除以 $t$ 并令 $t \to 0$，便得
  \(\displaystyle \frac{\partial f}{\partial\vec l}(p_0) = f_x(p_0)l_1 + f_y(p_0)l_2 = \nabla f(p_0)\cdot\vec l\)。

  \item \textbf{进一步思考}：
  方向导数可以写为两个向量的点积：$D_{\vec l} f = \nabla f \cdot \vec l = \|\nabla f\| \|\vec l\| \cos\theta = \|\nabla f\| \cos\theta$，其中 $\theta$ 是梯度与方向 $\vec l$ 之间的夹角。
  要使该值最大，只需让 $\cos\theta = 1$，即 $\theta = 0$。因此当 $\vec l$ 与梯度 $\nabla f(p_0)$ \textbf{同向}时，方向导数取得最大值，且最大值为 $\|\nabla f(p_0)\|$。
  \textbf{几何事实}：梯度向量指向了函数值增长最快的方向，其模长正是这个最大增长率。
\end{enumerate}
\end{solution}

\subsubsection{A组}

\begin{exercise}[基本概念问答]
回答下列问题：
\begin{shortexenum}
  \shortitem{1}{二元函数 $f(x,y)$ 的偏导数怎样定义？} &
  \shortitem{2}{三元函数 $f(x,y,z)$ 的偏导数怎样定义？} \\
  \shortitem{3}{$f(x,y)$ 的偏导数的几何意义是什么？} &
  \shortitem{4}{函数 $z=f(x,y)$ 在点 $(a,b)$ 可微怎样定义？全微分是指什么？} \\
  \shortitem{5}{函数 $f(x,y)$ 在某点附近局部线性化是什么意思？} &
  \shortitem{6}{函数 $f(x,y)$ 在一点的可微性与连续性有什么关系？} \\
  \shortitem{7}{函数 $f(x,y)$ 可微的必要条件是什么？充分条件是什么？} &
\end{shortexenum}
\end{exercise}
\begin{solution}
\begin{enumerate}
  \item[(1)] \textbf{解析}：见定义 \ref{def:partial-derivative}，将其中一个变量固定为常量，对剩余变量求一元导数的极限过程即为偏导数。
  \item[(2)] \textbf{解析}：与二元函数类似，固定其中两个变量，对第三个变量求导数，即通过构造极限
  \[
  \lim_{\Delta x\to0}\frac{f(x+\Delta x,y,z)-f(x,y,z)}{\Delta x}
  \]
  来定义。
  \item[(3)] \textbf{解析}：见备注 \ref{remark:pd-geometry}。几何上代表曲面 $z=f(x,y)$ 被平行于坐标平面的平面（如 $y=y_0$ 或 $x=x_0$）截得的截线在特定点处的切线斜率。
  \item[(4)] \textbf{解析}：见定义 \ref{def:differentiability}。若增量可写成 $f(a+\Delta x, b+\Delta y)-f(a,b) = A\Delta x + B\Delta y + o(\rho)$ 形式（$\rho=\sqrt{\Delta x^2+\Delta y^2}$），则在此点可微。其线性部分 $A\Delta x + B\Delta y$ 即全微分。
  \item[(5)] \textbf{解析}：指在一个足够小的邻域内，复杂的曲面（函数）可以由该点处的切平面（线性映射）加上高阶无穷小误差来近似描述。
  \item[(6)] \textbf{解析}：可微必连续，但不连续必不可微。连续是可微的必要非充分条件。
  \item[(7)] \textbf{解析}：可微的必要条件是偏导数存在（且微分系数恰好等于对应的偏导数值）。可微的充分条件是该点附近偏导数存在，且偏导数在该点连续。
\end{enumerate}
\end{solution}

\begin{exercise}[偏导数的计算]
求下列函数的偏导数：
\par\noindent
\begin{tabular}{@{}p{0.3\textwidth}p{0.3\textwidth}p{0.3\textwidth}@{}}
(1) $z = x^2y - y^2x$ & (2) $y = \sin(2t - 5x)$ & (3) $f(x,y) = xy + \dfrac{x}{y}$ \\
(4) $z = \dfrac{x}{y^2}$ & (5) $z = \dfrac{\cos x^2}{y}$ & (6) $z = \tan\dfrac{x^2}{y}$ \\
(7) $z = x^y$ & (8) $g(x,y) = \ln\sqrt{x^2+y^2}$ & (9) $u = \dfrac{1}{x^2+y^2+z^2}$ \\
(10) $u = z^{xy}$ & (11) $u = (xy)^z$ & (12) $u = x^{\frac{y}{z}}$
\end{tabular}
\end{exercise}
\begin{solution}
\begin{enumerate}
  \item[(1)] \textbf{解析}：视 $y$ 或 $x$ 分别为常数。
  $\dfrac{\partial z}{\partial x} = 2xy - y^2, \quad \dfrac{\partial z}{\partial y} = x^2 - 2xy$.
  \item[(2)] \textbf{解析}：这里的变量显然是 $t$ 和 $x$。
  $\dfrac{\partial y}{\partial t} = 2\cos(2t - 5x), \quad \dfrac{\partial y}{\partial x} = -5\cos(2t - 5x)$.
  \item[(3)] \textbf{解析}：
  $f_x = y + \dfrac{1}{y}, \quad f_y = x - \dfrac{x}{y^2}$.
  \item[(4)] \textbf{解析}：
  $\dfrac{\partial z}{\partial x} = \dfrac{1}{y^2}, \quad \dfrac{\partial z}{\partial y} = -\dfrac{2x}{y^3}$.
  \item[(5)] \textbf{解析}：对 $x$ 求导用到复合函数法则。
  $\dfrac{\partial z}{\partial x} = \dfrac{-2x\sin x^2}{y}, \quad \dfrac{\partial z}{\partial y} = -\dfrac{\cos x^2}{y^2}$.
  \item[(6)] \textbf{解析}：内层为 $\dfrac{x^2}{y}$。
  $\dfrac{\partial z}{\partial x} = \sec^2\left(\frac{x^2}{y}\right) \cdot \frac{2x}{y}, \quad \dfrac{\partial z}{\partial y} = \sec^2\left(\frac{x^2}{y}\right) \cdot \left(-\frac{x^2}{y^2}\right)$.
  \item[(7)] \textbf{解析}：对 $x$ 是幂函数，对 $y$ 是指数函数。
  $\dfrac{\partial z}{\partial x} = yx^{y-1}, \quad \dfrac{\partial z}{\partial y} = x^y\ln x$.
  \item[(8)] \textbf{解析}：化简 $g = \frac{1}{2}\ln(x^2+y^2)$ 后求导更为简便。
  $g_x = \dfrac{x}{x^2+y^2}, \quad g_y = \dfrac{y}{x^2+y^2}$.
  \item[(9)] \textbf{解析}：可改写为 $u = (x^2+y^2+z^2)^{-1}$。因此
  \[
  u_x = \dfrac{-2x}{(x^2+y^2+z^2)^2},\quad
  u_y = \dfrac{-2y}{(x^2+y^2+z^2)^2},\quad
  u_z = \dfrac{-2z}{(x^2+y^2+z^2)^2}.
  \]
  \item[(10)] \textbf{解析}：对于三个变量 $x, y, z$ 的函数，对 $x$ 和 $y$ 是指数形式求导，对 $z$ 是幂形式。
  \[
  u_x = z^{xy}\ln z \cdot y = yz^{xy}\ln z,\quad
  u_y = z^{xy}\ln z \cdot x = xz^{xy}\ln z,\quad
  u_z = xy z^{xy-1}.
  \]
  \item[(11)] \textbf{解析}：对 $x, y$ 是幂形式，对 $z$ 是指数形式。
  \[
  u_x = z(xy)^{z-1} \cdot y = yz(xy)^{z-1},\quad
  u_y = z(xy)^{z-1} \cdot x = xz(xy)^{z-1},\quad
  u_z = (xy)^z \ln(xy).
  \]
  \item[(12)] \textbf{解析}：对 $x$ 为幂函数，对 $y, z$ 整体为指数函数。
  \[
  u_x = \dfrac{y}{z} x^{\frac{y}{z}-1},\quad
  u_y = x^{\frac{y}{z}}\ln x \cdot \dfrac{1}{z},\quad
  u_z = x^{\frac{y}{z}}\ln x \cdot \left(-\dfrac{y}{z^2}\right).
  \]
\end{enumerate}
\end{solution}

\begin{exercise}[指定点的偏导数计算]
求下列函数在指定点的各个偏导数：
\par\noindent
\begin{tabular}{@{}p{0.45\textwidth}p{0.45\textwidth}@{}}
(1) $z = \dfrac{x}{\sqrt{x^2+y^2}}$，在 $(1,0),(0,1)$ 处 & (2) $z = e^{-x}\sin(x+2y)$，在 $(0,\dfrac{\pi}{4})$ 处 \\
(3) $f(x,y,z) = \sqrt[x]{\dfrac{x}{y}}$，在 $(1,1,1)$ 处 & (4) $z = \sin x\ln(y+1) + \cos y\ln(1-x)$，在 $(0,0)$ 处
\end{tabular}
\end{exercise}
\begin{solution}
\begin{enumerate}
  \item[(1)] \textbf{解析}：
  先求出偏导数解析式：\(\displaystyle z_x=\frac{y^2}{(x^2+y^2)^{3/2}},\quad z_y=-\frac{xy}{(x^2+y^2)^{3/2}}\)。
  代入点 $(1,0)$：$z_x(1,0) = 0, z_y(1,0) = 0$。
  代入点 $(0,1)$：$z_x(0,1) = 1, z_y(0,1) = 0$。

  \item[(2)] \textbf{解析}：
  先求出偏导数：\(\displaystyle z_x = -e^{-x}\sin(x+2y) + e^{-x}\cos(x+2y),\qquad z_y = e^{-x} \cdot 2\cos(x+2y) = 2e^{-x}\cos(x+2y)\)。
  代入 $(0, \pi/4)$，注意 $x+2y = \pi/2$。
  $z_x(0, \pi/4) = -e^0\sin(\pi/2) + e^0\cos(\pi/2) = -1 + 0 = -1$；
  $z_y(0, \pi/4) = 2e^0\cos(\pi/2) = 0$。

  \item[(3)] \textbf{解析}：
  函数可写作 $f = (x/y)^{1/x} = x^{1/x}y^{-1/x}$。注意题目给定的点是 $(1,1,1)$，这里并没有 $z$ 变量，说明只需对 $x$ 和 $y$ 求导即可。
  $f_y = x^{1/x} \left(-\frac{1}{x}\right) y^{-1/x-1}$，代入 $(1,1)$ 得 $f_y(1,1) = 1 \cdot (-1) \cdot 1 = -1$。
  对 $f_x$ 可以用对数求导法：$\ln f = \frac{1}{x}(\ln x - \ln y)$，则 $\frac{f_x}{f} = -\frac{1}{x^2}(\ln x - \ln y) + \frac{1}{x}\cdot\frac{1}{x} = \frac{1-\ln x+\ln y}{x^2}$。
  代入 $(1,1)$ 得 $\frac{f_x(1,1)}{f(1,1)} = \frac{1-0+0}{1} = 1$。由于 $f(1,1)=1$，故 $f_x(1,1) = 1$。

  \item[(4)] \textbf{解析}：
  求导：\(\displaystyle z_x = \cos x\ln(y+1) - \dfrac{\cos y}{1-x},\qquad z_y = \dfrac{\sin x}{y+1} - \sin y\ln(1-x)\)。
  代入 $(0,0)$：
  $z_x(0,0) = 1\cdot 0 - \frac{1}{1} = -1$；
  $z_y(0,0) = \frac{0}{1} - 0 \cdot 0 = 0$。
\end{enumerate}
\end{solution}

\begin{exercise}[分段函数的偏导数]
求函数 $f(x,y) = \begin{cases} \dfrac{xy}{\sqrt{x^2+y^2}}, & x^2+y^2 \neq 0, \\ 0, & x^2+y^2 = 0 \end{cases}$ 的偏导数。
\end{exercise}
\begin{solution}
\begin{enumerate}
  \item \textbf{当 $(x,y) \neq (0,0)$ 时}：应用普通的除法求导法则。
  \begin{align*}
  f_x(x,y) &= \frac{y\sqrt{x^2+y^2} - xy\frac{x}{\sqrt{x^2+y^2}}}{x^2+y^2} = \frac{y(x^2+y^2) - x^2y}{(x^2+y^2)^{3/2}} = \frac{y^3}{(x^2+y^2)^{3/2}}. \\[6pt]
  f_y(x,y) &= \frac{x\sqrt{x^2+y^2} - xy\frac{y}{\sqrt{x^2+y^2}}}{x^2+y^2} = \frac{x^3}{(x^2+y^2)^{3/2}}.
  \end{align*}
  \item \textbf{当 $(x,y) = (0,0)$ 时}：只能用偏导数定义，且 \(\displaystyle f_x(0,0) = \lim_{x \to 0} \frac{f(x,0) - f(0,0)}{x} = \lim_{x \to 0} \frac{0 - 0}{x} = 0,\qquad f_y(0,0) = \lim_{y \to 0} \frac{f(0,y) - f(0,0)}{y} = 0\)。
  \item \textbf{得出完整结论}：
  综合以上两点，偏导数为分段函数形式：
  \begin{align*}
  f_x(x,y) &= \begin{cases} \dfrac{y^3}{(x^2+y^2)^{3/2}}, & (x,y) \neq (0,0), \\ 0, & (x,y) = (0,0), \end{cases}\\
  f_y(x,y) &= \begin{cases} \dfrac{x^3}{(x^2+y^2)^{3/2}}, & (x,y) \neq (0,0), \\ 0, & (x,y) = (0,0). \end{cases}
  \end{align*}
\end{enumerate}
\end{solution}

\begin{exercise}[多元复合函数的偏导数]
求下列函数的偏导数：\\
（1） $u = \ln(x_1+x_2+\cdots+x_n)$；（2）$u = \arcsin(x_1^2+x_2^2+\cdots+x_n^2)$。
\end{exercise}
\begin{solution}
\begin{enumerate}
  \item[(1)] \textbf{解析}：对任意 $k \in \{1,2,\dots,n\}$，视其余自变量为常数，由链式法则直接可得 \(\displaystyle \frac{\partial u}{\partial x_k}=\frac{1}{\sum_{i=1}^n x_i}\)。
  \item[(2)] \textbf{解析}：同样取定 $x_k$，外层函数是 $\arcsin(\cdot)$，内层函数中仅 $x_k^2$ 求导非零，故 \(\displaystyle \frac{\partial u}{\partial x_k}=\frac{2x_k}{\sqrt{1-\left(\sum_{i=1}^n x_i^2\right)^2}}\)。
\end{enumerate}
\end{solution}

\begin{exercise}[全微分的计算]
求下列函数的全微分：
\par\noindent
\begin{tabular}{@{}p{0.3\textwidth}p{0.3\textwidth}p{0.3\textwidth}@{}}
(1) $z = e^{-x}\cos y$ & (2) $f(x,y) = \sin(xy)$ & (3) $g(u,v) = u^2+uv$ \\
(4) $u = x^{yz}$ & (5) $z = \dfrac{y}{\sqrt{x^2+y^2}}$ & (6) $z = x^2y + \dfrac{x}{y}$
\end{tabular}
\end{exercise}
\begin{solution}
\begin{enumerate}
  \item[(1)] \textbf{解析}：$z_x = -e^{-x}\cos y$，$z_y = -e^{-x}\sin y$。
  $\mathrm{d}z = -e^{-x}\cos y \,\mathrm{d}x - e^{-x}\sin y \,\mathrm{d}y$.

  \item[(2)] \textbf{解析}：$f_x = y\cos(xy)$，$f_y = x\cos(xy)$。
  $\mathrm{d}f = y\cos(xy)\,\mathrm{d}x + x\cos(xy)\,\mathrm{d}y$.

  \item[(3)] \textbf{解析}：这里自变量是 $u$ 和 $v$，故求导并在公式中用 $\mathrm{d}u, \mathrm{d}v$。$g_u = 2u+v$，$g_v = u$。
  $\mathrm{d}g = (2u+v)\,\mathrm{d}u + u\,\mathrm{d}v$.

  \item[(4)] \textbf{解析}：$u_x = yz x^{yz-1}$，$u_y = x^{yz}\ln x \cdot z$，$u_z = x^{yz}\ln x \cdot y$。
  $\mathrm{d}u = yzx^{yz-1}\,\mathrm{d}x + zx^{yz}\ln x\,\mathrm{d}y + yx^{yz}\ln x\,\mathrm{d}z$.

  \item[(5)] \textbf{解析}：$z = y(x^2+y^2)^{-1/2}$，因此 \(\displaystyle z_x = -\frac{xy}{(x^2+y^2)^{3/2}}\)，且 \(\displaystyle z_y = \frac{x^2}{(x^2+y^2)^{3/2}}\)。
  $\mathrm{d}z = -\dfrac{xy}{(x^2+y^2)^{3/2}}\,\mathrm{d}x + \dfrac{x^2}{(x^2+y^2)^{3/2}}\,\mathrm{d}y$.
  \item[(6)] \textbf{解析}：$z_x = 2xy + \frac{1}{y}$，$z_y = x^2 - \frac{x}{y^2}$。
  $\mathrm{d}z = \left(2xy + \frac{1}{y}\right)\mathrm{d}x + \left(x^2 - \frac{x}{y^2}\right)\mathrm{d}y$.
\end{enumerate}
\end{solution}

\begin{exercise}[指定点的全微分]
求下列函数在给定点的全微分：
\par\noindent
\begin{tabular}{@{}p{0.45\textwidth}p{0.45\textwidth}@{}}
(1) $f(x,y) = xe^{-y}$，在 $(1,0)$ 点 & (2) $g(x,t) = x^2\sin2t$，在 $(2,\dfrac{\pi}{4})$ 点 \\
(3) $F(m,r) = \dfrac{Gm}{r^2}$，在 $(100,10)$ 点 & (4) $f(x,y) = \ln(1+x^2+y^2)$，在 $(1,2)$ 点
\end{tabular}
\end{exercise}
\begin{solution}
\begin{enumerate}
  \item[(1)] \textbf{解析}：$f_x = e^{-y}$, $f_y = -xe^{-y}$。
  在 $(1,0)$ 处：$f_x(1,0) = 1$, $f_y(1,0) = -1$。
  $\mathrm{d}f|_{(1,0)} = \mathrm{d}x - \mathrm{d}y$.

  \item[(2)] \textbf{解析}：$g_x = 2x\sin2t$, $g_t = 2x^2\cos2t$。
  在 $(2,\pi/4)$ 处：$g_x = 4\sin(\pi/2) = 4$, $g_t = 8\cos(\pi/2) = 0$。
  $\mathrm{d}g|_{(2,\pi/4)} = 4\,\mathrm{d}x + 0\,\mathrm{d}t = 4\,\mathrm{d}x$.

  \item[(3)] \textbf{解析}：将 $G$ 视为常量。$F_m = G/r^2$, $F_r = -2Gm/r^3$。
  在 $(100,10)$ 处：$F_m = G/100$, $F_r = -2G(100)/1000 = -G/5$。
  $\mathrm{d}F|_{(100,10)} = \frac{G}{100}\,\mathrm{d}m - \frac{G}{5}\,\mathrm{d}r$.

  \item[(4)] \textbf{解析}：\(\displaystyle f_x = \frac{2x}{1+x^2+y^2},\qquad f_y = \frac{2y}{1+x^2+y^2}\)。
  在 $(1,2)$ 处，分母为 $1+1+4=6$，所以 \(\displaystyle f_x(1,2)=\frac13,\qquad f_y(1,2)=\frac23\)。
  因而 \(\displaystyle \mathrm{d}f|_{(1,2)} = \frac{1}{3}\,\mathrm{d}x + \frac{2}{3}\,\mathrm{d}y\)。
\end{enumerate}
\end{solution}

\begin{exercise}[可微性判定与偏导数有界性]
证明：函数
\[
f(x,y) = \begin{cases}
\dfrac{xy}{\sqrt{x^2+y^2}}, & x^2+y^2 \neq 0, \\
0, & x^2+y^2 = 0
\end{cases}
\]
在点 $(0,0)$ 的邻域内连续，且有有界的偏导数 $f_x(x,y)$ 与 $f_y(x,y)$，但此函数在点 $(0,0)$ 不可微。
\end{exercise}
\begin{solution}
\begin{enumerate}
  \item \textbf{证明在 $(0,0)$ 处连续}：
  令 \(x=r\cos\theta,\ y=r\sin\theta\)，则
  \[
    f(r\cos\theta,r\sin\theta)=r\cos\theta\sin\theta.
  \]
  其绝对值不超过 \(\frac12r\)，故 \(\displaystyle\lim_{(x,y)\to (0,0)} f(x,y) = 0 = f(0,0)\)。又因其在去心邻域是由初等函数构成显然连续，故函数在全平面连续。

  \item \textbf{证明偏导数存在且有界}：
  由前面的习题已计算得当 $(x,y) \neq (0,0)$ 时，$f_x(x,y) = \dfrac{y^3}{(x^2+y^2)^{3/2}}$，$f_y(x,y) = \dfrac{x^3}{(x^2+y^2)^{3/2}}$；而 $f_x(0,0)=f_y(0,0)=0$。
  考察偏导数的绝对值：$|y^3| = |y|^3 \leqslant  (\sqrt{x^2+y^2})^3 = (x^2+y^2)^{3/2}$。
  因此对任意 $(x,y)\neq(0,0)$，有 $|f_x(x,y)| \leqslant  1$。同理 $|f_y(x,y)| \leqslant  1$。并且在原点均为 0，所以 $f_x$ 和 $f_y$ 在全平面内均有界（界为 1）。

  \item \textbf{证明在 $(0,0)$ 处不可微}：
  利用可微定义进行验证。此时 $\mathrm{d}f(0,0) = f_x(0,0)\Delta x + f_y(0,0)\Delta y = 0$。
  如果可微，则极限必须为 $0$：
  \[
    \lim_{(x,y)\to(0,0)} \frac{f(x,y) - f(0,0)}{\sqrt{x^2+y^2}}
    = \lim_{(x,y)\to(0,0)} \frac{xy}{x^2+y^2}.
  \]
  极坐标下右端为 \(\frac12\sin2\theta\)，随角度改变而没有二重极限，更不可能等于 \(0\)。这说明函数的增量无法由切平面做高阶近似。因此 $f$ 在 $(0,0)$ \textbf{不可微}。
\end{enumerate}
\end{solution}

\subsubsection{B组}

\begin{exercise}[偏导数的计算与求导法则]
计算下列各题：
\par\noindent
\begin{tabular}{@{}p{0.35\textwidth}p{0.3\textwidth}p{0.25\textwidth}@{}}
(1) $\dfrac{\partial}{\partial m}\left(\dfrac{1}{2}mv^2\right)$ & (2) $\dfrac{\partial}{\partial T}\left(\dfrac{2\pi r}{T}\right)$ & (3) $\dfrac{\partial}{\partial x}(a\sqrt{x}),\ x>0$ \\
(4) $\dfrac{\partial}{\partial T}\left(\ln\dfrac{T+3}{V}\right)$ & (5) $\dfrac{\partial}{\partial y}\left(\dfrac{3x^2y^7-y^2}{15xy-8}\right)$ & (6) $\dfrac{\partial}{\partial \lambda}\left(\dfrac{x^2y\lambda-3\lambda^5}{\sqrt{\lambda^2-2\lambda+5}}\right)$ \\
(7) $\dfrac{\partial}{\partial \theta}\left[3\cos(5\theta-1)-\tan(7t\theta^2)\right]$ & (8) $\dfrac{\partial}{\partial w}\left(\sqrt{2\pi xyw-13x^7y^3v}\right)$ & \\
\end{tabular}
\end{exercise}
\begin{solution}
\begin{enumerate}
  \item[(1)] \textbf{解析}：视 $v$ 为常数，只对 $m$ 求导。结果为 $\dfrac{1}{2}v^2$。
  \item[(2)] \textbf{解析}：视 $r$ 为常数，对 $T$ 求导。$\dfrac{\partial}{\partial T}\left(2\pi r T^{-1}\right) = -2\pi r T^{-2} = -\dfrac{2\pi r}{T^2}$。
  \item[(3)] \textbf{解析}：视 $a$ 为常数。$\dfrac{\partial}{\partial x}(ax^{1/2}) = \dfrac{a}{2\sqrt{x}}$。
  \item[(4)] \textbf{解析}：利用对数性质拆分 $\ln(T+3) - \ln V$ 极易求解。对 $T$ 求导仅剩前一部分，即 $\dfrac{1}{T+3}$。
  \item[(5)] \textbf{解析}：视 $x$ 为常数，应用除法法则对 $y$ 求导。分子部分为
  \[
  \frac{(21x^2y^6-2y)(15xy-8) - (3x^2y^7-y^2)(15x)}{(15xy-8)^2},
  \]
  将其化简，分子的乘积展开略微复杂，但按步骤进行即可。
  \item[(6)] \textbf{解析}：考察函数 $\dfrac{P(\lambda)}{Q(\lambda)}$，其中 $x, y$ 只是常系数。直接用一元函数除法法则得
  \begin{align*}
  \frac{\partial}{\partial\lambda}\frac{x^2y\lambda-3\lambda^5}{\sqrt{\lambda^2-2\lambda+5}}
  &=
  \frac{(x^2y - 15\lambda^4)\sqrt{\lambda^2-2\lambda+5}}{\lambda^2-2\lambda+5}\\
  &\quad-\frac{(x^2y\lambda-3\lambda^5)(2\lambda-2)}{2(\lambda^2-2\lambda+5)^{3/2}},
  \end{align*}
  化简繁琐保留即可。
  \item[(7)] \textbf{解析}：两项分别对 $\theta$ 求导：
  前项：$-3\sin(5\theta-1)\cdot 5 = -15\sin(5\theta-1)$。
  后项：$-\sec^2(7t\theta^2) \cdot (14t\theta)$。
  相加即可。
  \item[(8)] \textbf{解析}：外层为根号函数，内层针对 $w$ 求导，因此 \(\displaystyle \frac{\partial}{\partial w}\left(\sqrt{2\pi xyw-13x^7y^3v}\right)=\frac{\pi xy}{\sqrt{2\pi xyw-13x^7y^3v}}\)。
\end{enumerate}
\end{solution}

\begin{exercise}[分段函数在指定点的偏导数]
求下列函数在指定点的偏导数：
\begin{shortsingleenum}
  \shortsinglechoice{(1)}{$f(x,y) = \begin{cases} \dfrac{x^3+y^3}{x^2+y^2}, & x^2+y^2 \neq 0, \\ 0, & x^2+y^2 = 0 \end{cases}$，在 $(0,0)$ 点；}
  \shortsinglechoice{(2)}{$f(x,y) = \begin{cases} \dfrac{x^2y^3}{x^2+4y^3}, & x^2+y^2 \neq 0, \\ 0, & x^2+y^2 = 0 \end{cases}$，在 $(0,0)$ 点。}
\end{shortsingleenum}
\end{exercise}
\begin{solution}
\begin{enumerate}
  \item[(1)] \textbf{解析}：在分界点处严格使用偏导数定义。由 \(\displaystyle f_x(0,0) = \lim_{x \to 0} \frac{f(x,0) - f(0,0)}{x} = \lim_{x \to 0} \frac{x^3/x^2 - 0}{x} = 1\)
  可知 $f_x(0,0)=1$；由于关于 $x$ 和 $y$ 对称，同理 $f_y(0,0)=1$。

  \item[(2)] \textbf{解析}：同样使用偏导数定义，并有 \(\displaystyle f_x(0,0) = \lim_{x \to 0} \frac{f(x,0) - 0}{x} = 0, \qquad
  f_y(0,0) = \lim_{y \to 0} \frac{f(0,y) - 0}{y} = 0\)。
\end{enumerate}
\end{solution}

\begin{exercise}[全微分近似值计算]
求下列各式的近似值：(1)~$(0.97)^{1.05}$；~~(2)~$\sin29^\circ\tan46^\circ$。
\end{exercise}
\begin{solution}
\begin{enumerate}
  \item[(1)] \textbf{解析}：
  令 $f(x,y) = x^y$，基准点取易于计算的整数 $(x_0,y_0)=(1,1)$。
  则增量为 $h_1 = -0.03$，$h_2 = 0.05$。
  先计算偏导数及基准点的值：$f(1,1)=1$；$f_x = yx^{y-1} \implies f_x(1,1)=1$；$f_y = x^y\ln x \implies f_y(1,1)=0$。
  利用全微分公式：$f(1-0.03, 1+0.05) \approx f(1,1) + f_x(1,1)h_1 + f_y(1,1)h_2 = 1 + 1(-0.03) + 0(0.05) = \boxed{0.97}$。

  \item[(2)] \textbf{解析}：
  令 $f(x,y) = \sin x\tan y$，基准点取特殊角化为弧度 $(x_0,y_0)=(\pi/6, \pi/4)$，即对应 $30^\circ$ 和 $45^\circ$。
  对应增量需转弧度：$h_1 = -1^\circ = -\pi/180$，$h_2 = 1^\circ = \pi/180$。
  基准点的值：$f(\pi/6, \pi/4) = \sin(30^\circ)\tan(45^\circ) = \frac{1}{2} \cdot 1 = 0.5$。
  偏导数：$f_x = \cos x\tan y \implies f_x(\pi/6, \pi/4) = \frac{\sqrt{3}}{2} \cdot 1 = \frac{\sqrt{3}}{2}$。
  $f_y = \sin x\sec^2 y \implies f_y(\pi/6, \pi/4) = \frac{1}{2} \cdot 2 = 1$。
  近似计算：\(\displaystyle f(29^\circ, 46^\circ) \approx 0.5 + \frac{\sqrt{3}}{2}\left(-\frac{\pi}{180}\right) + 1\left(\frac{\pi}{180}\right) = 0.5 + \frac{(2-\sqrt{3})\pi}{360}\)。
  代入近似常数计算得具体数值即可。
\end{enumerate}
\end{solution}

\begin{exercise}[偏导数不连续但可微的证明]
设函数 $f(x,y) = \begin{cases} xy\sin\dfrac{1}{\sqrt{x^2+y^2}}, & x^2+y^2 \neq 0, \\ 0, & x^2+y^2 = 0 \end{cases}$，证明：
\begin{shortexenum}
  \shortitem{1}{$f_x(0,0)$ 与 $f_y(0,0)$ 存在；} &
  \shortitem{2}{$f_x(x,y)$ 与 $f_y(x,y)$ 在点 $(0,0)$ 不连续；} \\
  \shortitem{3}{$f(x,y)$ 在点 $(0,0)$ 可微。} &
\end{shortexenum}
\end{exercise}
\begin{solution}
\begin{enumerate}
  \item[(1)] \textbf{验证偏导数存在}：
  利用定义可得 $f_x(0,0) = \lim\limits_{x \to 0} \dfrac{x\cdot 0 \cdot \sin(1/|x|) - 0}{x} = 0$；同理，由对称性 $f_y(0,0) = 0$。两者均存在。
  \item[(2)] \textbf{验证偏导数不连续}：
  求 $(x,y) \neq (0,0)$ 时的偏导数表达式。对 $x$ 求导：
  \begin{align*}
  f_x(x,y)
  &= y\sin\frac{1}{\sqrt{x^2+y^2}} + xy\cos\frac{1}{\sqrt{x^2+y^2}} \cdot \left(-\frac{1}{2}(x^2+y^2)^{-3/2} \cdot 2x\right)\\
  &= y\sin\frac{1}{\sqrt{x^2+y^2}} - \frac{x^2y}{(x^2+y^2)^{3/2}}\cos\frac{1}{\sqrt{x^2+y^2}}.
  \end{align*}
  引入极坐标变换 $x = r\cos\theta, y = r\sin\theta$。当 $(x,y) \to (0,0)$ 时，等价于 $r \to 0^+$。将其代入偏导数表达式得 \(\displaystyle f_x(r\cos\theta, r\sin\theta)=r\sin\theta\sin\frac{1}{r}-\cos^2\theta\sin\theta\cos\frac{1}{r}\)。
  当 $r \to 0^+$ 时，第一项由于是有界量乘无穷小量，故 $r\sin\theta\sin\frac{1}{r} \to 0$。
  对于第二项，若我们固定角度 $\theta$ 使得 $\cos^2\theta\sin\theta \neq 0$（例如取 $\theta=\frac{\pi}{4}$，该系数为 $\frac{\sqrt{2}}{4}$），随着 $r \to 0^+$，$\cos\frac{1}{r}$ 在 $-1$ 到 $1$ 之间不断振荡，导致整个表达式无确定的极限。
  因此，极限 $\lim\limits_{(x,y)\to(0,0)} f_x(x,y)$ 不存在，从而 $f_x$ 在 $(0,0)$ \textbf{不连续}。同理可证 $f_y$ 也不连续。
  \item[(3)] \textbf{验证可微性}：
  可微定义要求余项趋于 $0$。因 $f_x(0,0)=0, f_y(0,0)=0$，线性部分为 0。
  考察比值极限 \(\displaystyle \lim_{(x,y)\to(0,0)} \frac{xy\sin\dfrac{1}{\sqrt{x^2+y^2}}}{\sqrt{x^2+y^2}}\)。
  我们有 $|\sin(\cdot)| \leqslant 1$，且 $\dfrac{|xy|}{\sqrt{x^2+y^2}} \leqslant \dfrac{x^2+y^2}{2\sqrt{x^2+y^2}} = \dfrac{1}{2}\sqrt{x^2+y^2}$。
  当 $(x,y)\to(0,0)$ 时，上式右端趋于 $0$。根据夹逼定理，整体余项极限为 $0$。
  这表明增量公式成立，因此 $f(x,y)$ 在 $(0,0)$ \textbf{可微}。
\end{enumerate}
\end{solution}

\begin{exercise}[绝对值与函数可微性]
函数 $f(x,y) = |x|+\sin(xy)$ 在点 $(0,0)$ 是否可微？为什么？
\end{exercise}
\begin{solution}
\begin{enumerate}
  \item \textbf{考察偏导数的存在性}：
  按定义考察在 $(0,0)$ 处对 $x$ 的偏导数：\(\displaystyle f_x(0,0) = \lim_{\Delta x \to 0} \frac{f(\Delta x, 0) - f(0,0)}{\Delta x} = \lim_{\Delta x \to 0} \frac{|\Delta x| + \sin(0) - 0}{\Delta x} = \lim_{\Delta x \to 0} \frac{|\Delta x|}{\Delta x}\)。
  由于左右极限不相等（右侧趋于 $1$，左侧趋于 $-1$），所以 $f_x(0,0)$ \textbf{不存在}。
  \item \textbf{得出结论}：
  由定理 \ref{thm:diff-cont}（可微必有偏导数）可知，因为函数在原点对 $x$ 的偏导数不存在，故该函数在 $(0,0)$ 处必然\textbf{不可微}。
\end{enumerate}
\end{solution}

% ============================================================
% §7.8 向量值函数的导数
% ============================================================

\subsection{第四节习题：向量值函数及其微分}

\begin{exercise}[基础：向量值函数的极限与连续]
\begin{shortexenum}
  \shortchoice{(1)}{设 $f:\mathbb R\to\mathbb R^2$，$f(t)=(t^2,1-t^3)^{\mathsf T}$。求 $\displaystyle\lim_{t\to 1}f(t)$。} &
  \shortchoice{(2)}{设 $f:\mathbb R^2\to\mathbb R^2$，$f(x,y)=(e^{x+y},\sin(xy))^{\mathsf T}$。证明 $f$ 在 $\mathbb R^2$ 上连续。}
\end{shortexenum}
\end{exercise}
\begin{solution}
\begin{enumerate}
  \item[(1)] \textbf{求向量值函数的极限。}
  根据极限的分量等价性定理，向量值函数的极限等于各个分量函数极限组成的向量。
  计算各个分量在 $t \to 1$ 时的极限：
  \begin{itemize}
    \item 第一分量：$\lim\limits_{t\to 1} t^2 = 1^2 = 1$。
    \item 第二分量：$\lim\limits_{t\to 1} (1-t^3) = 1 - 1^3 = 0$。
  \end{itemize}
  因此，组合得到 $\lim_{t\to 1}f(t) = \begin{pmatrix} 1 \\ 0 \end{pmatrix}$。

  \item[(2)] \textbf{证明向量值函数的连续性。}
  根据连续性的分量等价性定理，$f$ 在 $\mathbb R^2$ 上连续当且仅当它的每个分量函数在 $\mathbb R^2$ 上连续。
  考察其两个分量函数：
  \begin{itemize}
    \item $f_1(x,y) = e^{x+y}$：这是多元多项式 $x+y$ 与指数函数 $e^u$ 的复合。由于多项式处处连续，指数函数处处连续，所以复合函数 $f_1$ 在整个 $\mathbb R^2$ 上连续。
    \item $f_2(x,y) = \sin(xy)$：这是多元多项式 $xy$ 与正弦函数 $\sin(u)$ 的复合。同理，它在整个 $\mathbb R^2$ 上也连续。
  \end{itemize}
  既然两个分量函数处处连续，那么向量值函数 $f(x,y)$ 在 $\mathbb R^2$ 上必然连续。
\end{enumerate}
\end{solution}

\begin{exercise}[基础：Jacobian 矩阵计算]
求下列向量值函数的 Jacobian 矩阵：
\begin{shortexenum}
  \shortchoice{(1)}{$f:\mathbb R^2\to\mathbb R^2$，$f(x,y)=(x^2-y^2,\;2xy)^{\mathsf T}$（复平面上 $z\mapsto z^2$ 的实化表示）；} &
  \shortchoice{(2)}{$f:\mathbb R\to\mathbb R^3$，$f(t)=(t,\,t^2,\,t^3)^{\mathsf T}$（三次挠曲线 Twisted Cubic）；} \\
  \shortchoice{(3)}{$f:\mathbb R^2\to\mathbb R^3$，$f(u,v)=(u\cos v,\;u\sin v,\;u)^{\mathsf T}$（圆锥面的参数化）。} &
\end{shortexenum}
\end{exercise}
\begin{solution}
\begin{enumerate}
  \item[(1)] \textbf{求 $f(x,y)=(x^2-y^2,\;2xy)^{\mathsf T}$ 的 Jacobian 矩阵。}
  分量函数为 $f_1 = x^2-y^2$ 和 $f_2 = 2xy$。分别对 $x, y$ 求偏导并按行排布，可得
  \[
  Jf(x,y) = \begin{pmatrix} \dfrac{\partial f_1}{\partial x} & \dfrac{\partial f_1}{\partial y} \\[6pt] \dfrac{\partial f_2}{\partial x} & \dfrac{\partial f_2}{\partial y} \end{pmatrix}
  = \begin{pmatrix} 2x & -2y \\ 2y & 2x \end{pmatrix}.
  \]

  \item[(2)] \textbf{求 $f(t)=(t,\,t^2,\,t^3)^{\mathsf T}$ 的 Jacobian 矩阵。}
  自变量只有一个 $t$，因此直接对每个分量求一阶导数，结果为一个 $3 \times 1$ 的列向量：\(\displaystyle Jf(t)=\bigl(1,2t,3t^2\bigr)^{\mathsf T}\)。

  \item[(3)] \textbf{求 $f(u,v)=(u\cos v,\;u\sin v,\;u)^{\mathsf T}$ 的 Jacobian 矩阵。}
  分量函数为 $f_1 = u\cos v$, $f_2 = u\sin v$, $f_3 = u$。
  分别对 $u, v$ 求偏导数：
  \begin{itemize}
    \item 第一行（$f_1$）：$\dfrac{\partial f_1}{\partial u} = \cos v$, \quad $\dfrac{\partial f_1}{\partial v} = -u\sin v$
    \item 第二行（$f_2$）：$\dfrac{\partial f_2}{\partial u} = \sin v$, \quad $\dfrac{\partial f_2}{\partial v} = u\cos v$
    \item 第三行（$f_3$）：$\dfrac{\partial f_3}{\partial u} = 1$, \quad $\dfrac{\partial f_3}{\partial v} = 0$
  \end{itemize}
  按行排布，得到 $3 \times 2$ 矩阵：
  \[
  Jf(u,v) = \begin{pmatrix}
  \cos v & -u\sin v \\
  \sin v & u\cos v \\
  1 & 0
  \end{pmatrix}.
  \]
\end{enumerate}
\end{solution}

\begin{exercise}[提高：极坐标变换的 Jacobian]
\begin{shortsingleenum}
  \shortsinglechoice{(1)}{设 $f:\mathbb R^2\to\mathbb R^2$ 为极坐标变换 $f(r,\theta)=(r\cos\theta,\;r\sin\theta)^{\mathsf T}$，计算 $Jf(r,\theta)$ 并求 $\det Jf(r,\theta)$；}
  \shortsinglechoice{(2)}{球坐标变换 $f(\rho,\varphi,\theta)=(\rho\sin\varphi\cos\theta,\;\rho\sin\varphi\sin\theta,\;\rho\cos\varphi)^{\mathsf T}$，计算 $\det Jf(\rho,\varphi,\theta)$。}
\end{shortsingleenum}
（提示：这两个行列式在多重积分的变量代换公式中扮演核心角色。）
\end{exercise}
\begin{solution}
\begin{enumerate}
  \item[(1)] \textbf{极坐标变换的 Jacobian 及其行列式。}
  分量函数为 $x(r,\theta) = r\cos\theta$ 和 $y(r,\theta) = r\sin\theta$。
  分别求偏导并组装矩阵：
  \begin{align*}
  Jf(r,\theta)
  &= \begin{pmatrix}
  \dfrac{\partial x}{\partial r} & \dfrac{\partial x}{\partial \theta} \\
  \dfrac{\partial y}{\partial r} & \dfrac{\partial y}{\partial \theta}
  \end{pmatrix}
  = \begin{pmatrix} \cos\theta & -r\sin\theta \\ \sin\theta & r\cos\theta \end{pmatrix},\\
  \det Jf(r,\theta)
  &= (\cos\theta)(r\cos\theta) - (-r\sin\theta)(\sin\theta)
  = r\cos^2\theta + r\sin^2\theta
  = r.
  \end{align*}

  \item[(2)] \textbf{球坐标变换的 Jacobian 行列式。}
  分量函数为 $x = \rho\sin\varphi\cos\theta$，$y = \rho\sin\varphi\sin\theta$，$z = \rho\cos\varphi$。
  对其分别关于 $\rho, \varphi, \theta$ 求偏导，得到 $3 \times 3$ 的 Jacobian 矩阵：
  \[
  Jf(\rho,\varphi,\theta) = \begin{pmatrix}
  \sin\varphi\cos\theta & \rho\cos\varphi\cos\theta & -\rho\sin\varphi\sin\theta \\
  \sin\varphi\sin\theta & \rho\cos\varphi\sin\theta & \rho\sin\varphi\cos\theta \\
  \cos\varphi & -\rho\sin\varphi & 0
  \end{pmatrix}.
  \]
  按第三行展开计算行列式：
  \begin{align*}
  \det Jf &= \cos\varphi \begin{vmatrix} \rho\cos\varphi\cos\theta & -\rho\sin\varphi\sin\theta \\ \rho\cos\varphi\sin\theta & \rho\sin\varphi\cos\theta \end{vmatrix}
  - (-\rho\sin\varphi) \begin{vmatrix} \sin\varphi\cos\theta & -\rho\sin\varphi\sin\theta \\ \sin\varphi\sin\theta & \rho\sin\varphi\cos\theta \end{vmatrix} + 0 \\
  &= \cos\varphi \left[ (\rho\cos\varphi\cos\theta)(\rho\sin\varphi\cos\theta) - (-\rho\sin\varphi\sin\theta)(\rho\cos\varphi\sin\theta) \right] \\
  &\quad + \rho\sin\varphi \left[ (\sin\varphi\cos\theta)(\rho\sin\varphi\cos\theta) - (-\rho\sin\varphi\sin\theta)(\sin\varphi\sin\theta) \right] \\
  &= \cos\varphi \left[ \rho^2\sin\varphi\cos\varphi(\cos^2\theta + \sin^2\theta) \right] + \rho\sin\varphi \left[ \rho\sin^2\varphi(\cos^2\theta + \sin^2\theta) \right] \\
  &= \cos\varphi (\rho^2\sin\varphi\cos\varphi) + \rho\sin\varphi (\rho\sin^2\varphi) \\
  &= \rho^2\sin\varphi\cos^2\varphi + \rho^2\sin^3\varphi \\
  &= \rho^2\sin\varphi(\cos^2\varphi + \sin^2\varphi) = \rho^2\sin\varphi.
  \end{align*}
  即球坐标变换的 Jacobian 行列式为 $\rho^2\sin\varphi$。
\end{enumerate}
\end{solution}

\begin{exercise}[Jacobian 行列式与坐标变换]
当 $m=n$ 时，Jacobian 矩阵是方阵，其行列式 $\det Jf(\boldsymbol a)$ 称为
\textbf{Jacobian 行列式}（或\textbf{函数行列式}），记作
$\dfrac{\partial(f_1,f_2,\dots,f_n)}{\partial(x_1,x_2,\dots,x_n)}$。

证明以下关于 Jacobian 行列式的重要性质：
\begin{shortsingleenum}
  \shortsinglechoice{(1)}{若 $f$ 和 $g$ 都是从 $\mathbb R^n$ 到 $\mathbb R^n$ 的可微映射，则 $\det J(g\circ f)=\det(Jg)\cdot\det(Jf)$（链式法则取行列式——矩阵乘法的行列式等于行列式之积）；}
  \shortsinglechoice{(2)}{若 $f$ 有可微的逆映射 $f^{-1}$，则 $\det J(f^{-1})=\dfrac{1}{\det Jf}$；}
  \shortsinglechoice{(3)}{在极坐标变换 $x=r\cos\theta$, $y=r\sin\theta$ 下，$\det Jf=r$。解释为什么 “面积元素 $\mathrm{d}x\,\mathrm{d}y=r\,\mathrm{d}r\,\mathrm{d}\theta$”。}
\end{shortsingleenum}
\end{exercise}
\begin{solution}
\begin{enumerate}
  \item[(1)] \textbf{证明复合映射的 Jacobian 行列式公式。}
  由上题证明的多元链式法则可知，复合映射 $g \circ f$ 的 Jacobian 矩阵为其各自 Jacobian 矩阵的乘积：\(\displaystyle J(g \circ f)(\boldsymbol x)=Jg(f(\boldsymbol x))\cdot Jf(\boldsymbol x)\)，从而 \(\displaystyle \det \bigl[J(g \circ f)(\boldsymbol x)\bigr]=\det \bigl[Jg(f(\boldsymbol x))\bigr]\cdot \det \bigl[Jf(\boldsymbol x)\bigr]\)。
  命题得证。

  \item[(2)] \textbf{证明逆映射的 Jacobian 行列式。}
  设恒等映射为 $I(\boldsymbol x) = \boldsymbol x$。显然恒等映射的 Jacobian 矩阵是单位矩阵 $I_n$，其行列式为 $1$。
  若 $f$ 具有可微的逆映射 $f^{-1}$，则依据定义有 $f^{-1} \circ f = I$。对该等式两边同时求 Jacobian 矩阵并应用链式法则，得 \(\displaystyle J(f^{-1})(f(\boldsymbol x)) \cdot Jf(\boldsymbol x)=J(I)(\boldsymbol x)=I_n\)，从而 \(\displaystyle \det\bigl[J(f^{-1})\bigr]\cdot\det\bigl[Jf\bigr]=\det(I_n)=1\)。
  当 $\det Jf \neq 0$ 时，即可得出 $\det J(f^{-1}) = \dfrac{1}{\det Jf}$。

  \item[(3)] \textbf{极坐标面积元素的几何解释。}
  在多元积分的变量代换中，Jacobian 行列式的绝对值 $|\det Jf|$ 刻画了由坐标变换所引起的\textbf{局部面积（或体积）的放缩比例}。
  对于极坐标变换，我们在网格平面 $(r, \theta)$ 上取一个长为 $\Delta r$，宽为 $\Delta \theta$ 的微小矩形区域，其面积为 $\Delta r \Delta \theta$。
  经过变换 $f$ 映射到 $(x, y)$ 平面后，这个区域变成了一个圆环扇形，其半径介于 $r$ 和 $r+\Delta r$ 之间，圆心角为 $\Delta \theta$。
  该扇形区域在平面上的真实几何面积约为 弧长 $\times$ 宽度，即 $(r\Delta\theta) \times \Delta r = r\,\Delta r\Delta\theta$。
  这说明映射将局部面积放大/缩小了 $r$ 倍。而计算得出的 $\det Jf = r$ 恰好反映了这个\textbf{面积伸缩因数}。
  因此，在微分意义下，积分的微面积元素满足等式 $\mathrm{d}x\,\mathrm{d}y = |\det Jf|\,\mathrm{d}r\,\mathrm{d}\theta = r\,\mathrm{d}r\,\mathrm{d}\theta$。
\end{enumerate}
\end{solution}

\begin{note}[Jacobian 行列式在第四章的核心应用]
上面第 (3) 小题揭示的"$|\det Jf|$ = 局部面积放缩倍数"正是第四章\textbf{重积分换元公式}（定理~\ref{thm:double-integral-substitution}）的理论基础：
\[
\iint_D f(x,y)\,\mathrm{d}x\,\mathrm{d}y
= \iint_{D'} f\bigl(x(u,v),\,y(u,v)\bigr)\,
\left|\frac{\partial(x,y)}{\partial(u,v)}\right|\mathrm{d}u\,\mathrm{d}v.
\]
换言之，第三章学的 Jacobian 矩阵不只是"线性近似的矩阵"，它的行列式直接告诉你换元后微元面积（体积）要乘多少。到第四章时只需把这个结论拿来用即可。
\end{note}

\begin{exercise}[概念与基本性质问答]
回答下列问题：
\begin{shortexenum}
  \shortitem{1}{什么叫做向量值函数？} &
  \shortitem{2}{向量值函数的极限与连续如何定义？} \\
  \shortitem{3}{雅可比矩阵是什么？向量值函数 $f:D\subset\mathbb{R}^n\to\mathbb{R}^m$ 在点 $\boldsymbol x^\circ$ 的导数怎样定义？} &
\end{shortexenum}
\end{exercise}
\begin{solution}
\begin{enumerate}
  \item[(1)] \textbf{向量值函数的定义：}
  若一个映射 $f: D \to \mathbb{R}^m$ 将区域 $D \subset \mathbb{R}^n$ 中的任意一点 $\boldsymbol x = (x_1, \dots, x_n)^{\mathsf T}$ 对应到一个 $m$ 维向量 $\boldsymbol y = (y_1, \dots, y_m)^{\mathsf T} \in \mathbb{R}^m$，则称 $f$ 为定义在 $D$ 上的向量值函数。它本质上是由 $m$ 个共同拥有自变量 $\boldsymbol x$ 的实值分量函数组装而成的。

  \item[(2)] \textbf{极限与连续的定义：}
  \begin{itemize}
    \item \textbf{极限：} 设 $\boldsymbol a$ 为 $D$ 的聚点。若对任意 $\varepsilon > 0$，都存在 $\delta > 0$，使得当 $0 < \|\boldsymbol x - \boldsymbol a\|_n < \delta$ 且 $\boldsymbol x \in D$ 时，恒有函数值与向量 $\boldsymbol A$ 的欧氏距离 $\|f(\boldsymbol x) - \boldsymbol A\|_m < \varepsilon$，则称当 $\boldsymbol x \to \boldsymbol a$ 时，$f(\boldsymbol x)$ 的极限为 $\boldsymbol A$。
    \item \textbf{连续：} 如果极限值恰好等于该点处的函数值，即 $\lim_{\boldsymbol x\to\boldsymbol a}f(\boldsymbol x)=f(\boldsymbol a)$，则称 $f$ 在点 $\boldsymbol a$ 处连续。这两者都可以等价转化为对其所有分量函数逐一要求极限和连续性。
  \end{itemize}

  \item[(3)] \textbf{雅可比矩阵与导数：}
  \begin{itemize}
    \item \textbf{导数（可微性）的定义：} 对于点 $\boldsymbol x^\circ$，若能找到一个 $m \times n$ 的常数矩阵 $A$，使得 $f(\boldsymbol x^\circ + \boldsymbol h) - f(\boldsymbol x^\circ) = A\boldsymbol h + \boldsymbol r(\boldsymbol h)$，且余项满足 $\lim_{\|\boldsymbol h\| \to 0} \frac{\|\boldsymbol r(\boldsymbol h)\|}{\|\boldsymbol h\|} = 0$，则称 $f$ 在 $\boldsymbol x^\circ$ 可微，这个矩阵 $A$ 就是 $f$ 在该点的导数。
    \item \textbf{雅可比矩阵（Jacobian 矩阵）：} 这个发挥导数作用的矩阵 $A$ 就叫做雅可比矩阵，记作 $Jf(\boldsymbol x^\circ)$。它的第 $i$ 行、第 $j$ 列的元素是由第 $i$ 个分量函数对第 $j$ 个自变量求偏导得到的值，即 $\dfrac{\partial f_i}{\partial x_j}(\boldsymbol x^\circ)$。
  \end{itemize}
\end{enumerate}
\end{solution}

\begin{exercise}[Jacobian 矩阵计算（二维）]
求向量值函数 $f(x,y) = \begin{bmatrix} x^2+y^2 \\ 3xy \end{bmatrix}$ 的导数 $f'(x,y)$。
\end{exercise}
\begin{solution}
\begin{enumerate}
  \item \textbf{识别分量并计算偏导数：}
  该函数有两个分量 $f_1(x,y) = x^2+y^2,\; f_2(x,y) = 3xy$，故 \(\displaystyle \frac{\partial f_1}{\partial x}=2x,\ \frac{\partial f_1}{\partial y}=2y,\ \frac{\partial f_2}{\partial x}=3y,\ \frac{\partial f_2}{\partial y}=3x\)。

  \item \textbf{组装导数矩阵（Jacobian 矩阵）：}
  将上述偏导数按所在分量的顺序排成矩阵：
  \[
  f'(x,y) = Jf(x,y) = \begin{bmatrix} 2x & 2y \\ 3y & 3x \end{bmatrix}.
  \]
\end{enumerate}
\end{solution}

\begin{exercise}[点处 Jacobian 矩阵计算]
设 $f:\mathbb{R}^3\to\mathbb{R}^2$ 定义为 $f(x,y,z) = \begin{bmatrix} e^x\cos y + e^y z^2 \\ 2x\sin y - 3y z^3 \end{bmatrix}$，求 $f'\left(0,\frac{\pi}{2},1\right)$。
\end{exercise}
\begin{solution}
\begin{enumerate}
  \item \textbf{列出分量函数：}
  $f_1(x,y,z) = e^x\cos y + e^y z^2$，$f_2(x,y,z) = 2x\sin y - 3y z^3$。

  \item \textbf{计算 $f_1$ 的偏导数并在点 $\left(0,\frac{\pi}{2},1\right)$ 求值：}
  \begin{itemize}
    \item 对 $x$ 偏导：$\dfrac{\partial f_1}{\partial x} = e^x\cos y$
          代入值 $\Rightarrow e^0\cos(\pi/2) = 1 \cdot 0 = 0$
    \item 对 $y$ 偏导：$\dfrac{\partial f_1}{\partial y} = -e^x\sin y + e^y z^2$
          代入值 $\Rightarrow -e^0\sin(\pi/2) + e^{\pi/2}(1^2) = -1 + e^{\pi/2}$
    \item 对 $z$ 偏导：$\dfrac{\partial f_1}{\partial z} = 2z e^y$
          代入值 $\Rightarrow 2(1)e^{\pi/2} = 2e^{\pi/2}$
  \end{itemize}

  \item \textbf{计算 $f_2$ 的偏导数并在点 $\left(0,\frac{\pi}{2},1\right)$ 求值：}
  \begin{itemize}
    \item 对 $x$ 偏导：$\dfrac{\partial f_2}{\partial x} = 2\sin y$
          代入值 $\Rightarrow 2\sin(\pi/2) = 2(1) = 2$
    \item 对 $y$ 偏导：$\dfrac{\partial f_2}{\partial y} = 2x\cos y - 3z^3$
          代入值 $\Rightarrow 2(0)\cos(\pi/2) - 3(1^3) = 0 - 3 = -3$
    \item 对 $z$ 偏导：$\dfrac{\partial f_2}{\partial z} = -9y z^2$
          代入值 $\Rightarrow -9(\pi/2)(1^2) = -\frac{9\pi}{2}$
  \end{itemize}

  \item \textbf{组装最终的导数矩阵：}
  将算出的值排成 $2 \times 3$ 的矩阵：
  \[
  f'\left(0,\frac{\pi}{2},1\right) = \begin{bmatrix}
  0 & e^{\pi/2}-1 & 2e^{\pi/2} \\[6pt]
  2 & -3 & -\dfrac{9\pi}{2}
  \end{bmatrix}.
  \]
\end{enumerate}
\end{solution}

\begin{exercise}[复杂函数的点处 Jacobian 矩阵计算]
设 $f:\mathbb{R}^3\to\mathbb{R}^3$ 定义为
\[
f(x,y,z) = \begin{bmatrix} \sin(x^2-y^2) \\ \ln(x^2+z^2) \\ \dfrac{1}{\sqrt{y^2+z^2}} \end{bmatrix}
\]
求 $f'(1,1,1)$。
\end{exercise}
\begin{solution}
\begin{enumerate}
  \item \textbf{分别写出分量并求关于各变量的偏导数。}
  代入的点为 $(x, y, z) = (1, 1, 1)$。

  \item \textbf{第一分量 $f_1(x,y,z) = \sin(x^2-y^2)$：}
  \begin{itemize}
    \item $\dfrac{\partial f_1}{\partial x} = \cos(x^2-y^2) \cdot 2x$ \quad $\Rightarrow$ \quad 代入 $(1,1,1)$ 得 $2(1)\cos(1-1) = 2\cos(0) = 2$
    \item $\dfrac{\partial f_1}{\partial y} = \cos(x^2-y^2) \cdot (-2y)$ \quad $\Rightarrow$ \quad 代入得 $-2(1)\cos(0) = -2$
    \item $\dfrac{\partial f_1}{\partial z} = 0$ \quad $\Rightarrow$ \quad 显然值为 $0$
  \end{itemize}

  \item \textbf{第二分量 $f_2(x,y,z) = \ln(x^2+z^2)$：}
  \begin{itemize}
    \item $\dfrac{\partial f_2}{\partial x} = \dfrac{1}{x^2+z^2} \cdot 2x$ \quad $\Rightarrow$ \quad 代入得 $\dfrac{2(1)}{1^2+1^2} = \dfrac{2}{2} = 1$
    \item $\dfrac{\partial f_2}{\partial y} = 0$ \quad $\Rightarrow$ \quad 显然值为 $0$
    \item $\dfrac{\partial f_2}{\partial z} = \dfrac{1}{x^2+z^2} \cdot 2z$ \quad $\Rightarrow$ \quad 代入得 $\dfrac{2(1)}{1^2+1^2} = 1$
  \end{itemize}

  \item \textbf{第三分量 $f_3(x,y,z) = (y^2+z^2)^{-1/2}$：}
  这里运用链式法则：
  \begin{itemize}
    \item $\dfrac{\partial f_3}{\partial x} = 0$ \quad $\Rightarrow$ \quad 显然值为 $0$
    \item $\dfrac{\partial f_3}{\partial y} = -\dfrac{1}{2}(y^2+z^2)^{-3/2} \cdot 2y = -y(y^2+z^2)^{-3/2}$ \quad $\Rightarrow$ \quad 代入得 $-1(1^2+1^2)^{-3/2}  = -\dfrac{\sqrt{2}}{4}$
    \item $\dfrac{\partial f_3}{\partial z} = -\dfrac{1}{2}(y^2+z^2)^{-3/2} \cdot 2z = -z(y^2+z^2)^{-3/2}$ \quad $\Rightarrow$ \quad 代入得 $-1(2)^{-3/2} = -\dfrac{\sqrt{2}}{4}$
  \end{itemize}

  \item \textbf{组装导数矩阵：}
  将三行偏导数值排成 $3 \times 3$ 的矩阵：
  \[
  f'(1,1,1) = \begin{bmatrix}
  2 & -2 & 0 \\[6pt]
  1 & 0 & 1 \\[6pt]
  0 & -\dfrac{\sqrt{2}}{4} & -\dfrac{\sqrt{2}}{4}
  \end{bmatrix}.
  \]
\end{enumerate}
\end{solution}

% ============================================================
% §7.4 复合函数的求导法则
% ============================================================

\subsection{第五节习题：复合函数的求导法则}

\subsubsection{A组}

\begin{exercise}[基本概念问答]
回答下列问题：
\begin{shortexenum}
  \shortitem{1}{你知道哪些情形下的求复合函数偏导数的链式法则？} &
  \shortitem{2}{一阶全微分形式的不变性指的是什么？} \\
  \shortitem{3}{高阶偏导数怎样定义？} &
  \shortitem{4}{在什么条件下混合偏导数相等？} \\
  \shortitem{5}{高阶全微分怎样定义？} &
\end{shortexenum}
\end{exercise}
\begin{solution}
\begin{enumerate}
  \item[(1)] \textbf{链式法则的情形：}
  \begin{itemize}
    \item \textbf{一元复合情形}：当最终自变量只有一个时，偏导数变为全导数（即对多条支路上的普通导数求和）。
    \item \textbf{多元复合情形}：当变量间构成多对多的嵌套层级时，每求一个自变量的偏导，需要从因变量出发遍历所有通向该自变量的路径并对各路径中间环节偏导数相乘后求和。
    \item \textbf{一般形式}：可用 Jacobian 矩阵相乘的形式来高度概括上述所有复杂映射链。
  \end{itemize}

  \item[(2)] \textbf{一阶全微分形式的不变性：}
  设函数 $z = f(x, y)$，全微分恒可表达为 $\mathrm{d}z = \dfrac{\partial z}{\partial x}\,\mathrm{d}x + \dfrac{\partial z}{\partial y}\,\mathrm{d}y$。其不变性是指无论 $x, y$ 本身是最底层的独立自变量，还是由其它更深层变量复合而成的中间函数，这个等式的形式始终有效，不会发生改变。

  \item[(3)] \textbf{高阶偏导数的定义：}
  多元函数对某变量求出的一阶偏导函数，如果对原各变量依然偏导数存在，则可继续求导。连续进行两次偏导运算即可得到二阶偏导数（包括两次对同变量的纯偏导，和对不同变量的混合偏导），以此类推可定义任意高阶的偏导数。

  \item[(4)] \textbf{混合偏导数相等的条件：}
  如果函数 $f(x,y)$ 的两个混合偏导数 $f_{xy}$ 和 $f_{yx}$ 在考察点 $(x_0, y_0)$ 及其某邻域内存在，并且在考察点处都是\textbf{连续}的，那么它们在此处的值一定相等。这就是 Clairaut 交换定理。

  \item[(5)] \textbf{高阶全微分的定义：}
  对函数的一阶全微分式 $\mathrm{d}z = f_x\,\mathrm{d}x + f_y\,\mathrm{d}y$ 视自变量增量（如 $\mathrm{d}x, \mathrm{d}y$）为独立常数再次执行微分操作得到的式子称为二阶全微分，即 $\mathrm{d}^2z = \mathrm{d}(\mathrm{d}z)$。依此类推递归定义更高阶的全微分。
\end{enumerate}
\end{solution}

\begin{exercise}[链式法则求一阶导数]
求下列函数的导数 $\dfrac{\mathrm{d}z}{\mathrm{d}t}$：
\par\noindent
\begin{tabular}{@{}p{0.45\textwidth}p{0.45\textwidth}@{}}
(1) $z = xy^2,\ x = e^{-t},\ y = \sin t$ & (2) $z = x\sin y + y\sin x,\ x = t^2,\ y = \ln t$ \\
(3) $z = \ln(x^2+y^2),\ x = \dfrac{1}{t},\ y = \sqrt{t}$ & (4) $z = \sin\dfrac{x}{y},\ x = 2t,\ y = 1-t^2$ \\
(5) $z = (x+y)e^y,\ x = 2t,\ y = 1-t^2$ & (6) $z = \arctan\dfrac{x}{y},\ x = 2t,\ y = 1-t^2$
\end{tabular}
\end{exercise}
\begin{solution}
\begin{enumerate}
  \item[(1)] \textbf{分析：} $z=xy^2$。先求偏导：$z_x = y^2$，$z_y = 2xy$。中间变量对 $t$ 导数：$x' = -e^{-t}$，$y' = \cos t$。
  \textbf{代入链式公式：}
  \begin{align*}
  \frac{\mathrm{d}z}{\mathrm{d}t} &= z_x\frac{\mathrm{d}x}{\mathrm{d}t} + z_y\frac{\mathrm{d}y}{\mathrm{d}t} = y^2(-e^{-t}) + 2xy(\cos t) \\
  &= (\sin^2 t)(-e^{-t}) + 2(e^{-t})(\sin t)(\cos t) \\
  &= e^{-t}\sin t(2\cos t - \sin t).
  \end{align*}

  \item[(2)] \textbf{分析：} $z=x\sin y + y\sin x$。偏导数：$z_x = \sin y + y\cos x$，$z_y = x\cos y + \sin x$。内层导数：$x' = 2t$，$y' = \dfrac{1}{t}$。
  \textbf{代入式中：}
  \begin{align*}
  \frac{\mathrm{d}z}{\mathrm{d}t}
  &= (\sin y + y\cos x)(2t) + (x\cos y + \sin x)\left(\frac{1}{t}\right) \\
  &= 2t(\sin(\ln t) + \ln t\cdot \cos(t^2)) + \frac{1}{t}(t^2\cos(\ln t) + \sin(t^2)).
  \end{align*}

  \item[(3)] \textbf{分析：} $z=\ln(x^2+y^2)$。此题代入合并为一元函数再求导会更简单：
  因为 $x = \dfrac{1}{t}$，$y = \sqrt{t}$，所以 $x^2+y^2 = \dfrac{1}{t^2} + t = \dfrac{t^3+1}{t^2}$。
  $z = \ln\left(\dfrac{t^3+1}{t^2}\right) = \ln(t^3+1) - 2\ln t$。
  \textbf{直接求导：}
  \(\displaystyle \frac{\mathrm{d}z}{\mathrm{d}t} = \frac{3t^2}{t^3+1} - \frac{2}{t} = \frac{3t^3 - 2(t^3+1)}{t(t^3+1)} = \frac{t^3-2}{t(t^3+1)}\)。

  \item[(4)] \textbf{分析：} $z=\sin\dfrac{x}{y}$。偏导数：$z_x = \dfrac{1}{y}\cos\dfrac{x}{y}$，$z_y = -\dfrac{x}{y^2}\cos\dfrac{x}{y}$。内层：$x'=2$，$y'=-2t$。
  \textbf{组装链式公式：}
  \begin{align*}
  \frac{\mathrm{d}z}{\mathrm{d}t} &= \left(\frac{1}{y}\cos\dfrac{x}{y}\right)(2) + \left(-\frac{x}{y^2}\cos\dfrac{x}{y}\right)(-2t)
  = \left(\frac{2}{y} + \frac{2xt}{y^2}\right)\cos\dfrac{x}{y} \\
  &= \left(\frac{2}{1-t^2} + \frac{4t^2}{(1-t^2)^2}\right)\cos\left(\frac{2t}{1-t^2}\right) \\
  &= \frac{2-2t^2+4t^2}{(1-t^2)^2}\cos\left(\frac{2t}{1-t^2}\right) = \frac{2(1+t^2)}{(1-t^2)^2}\cos\left(\frac{2t}{1-t^2}\right).
  \end{align*}

  \item[(5)] \textbf{分析：} $z=(x+y)e^y$。$z_x = e^y$，$z_y = e^y + (x+y)e^y = (x+y+1)e^y$。内层：$x'=2$，$y'=-2t$。
  \textbf{组装：}\(\displaystyle
  \frac{\mathrm{d}z}{\mathrm{d}t}
  =e^y(2)+(x+y+1)e^y(-2t)=2e^y [1-t(x+y+1)]
  =2e^{1-t^2}(1-2t-2t^2+t^3)\)。

  \item[(6)] \textbf{分析：} $z=\arctan\dfrac{x}{y}$。偏导：$z_x = \dfrac{1/y}{1+(x/y)^2} = \dfrac{y}{x^2+y^2}$，$z_y = \dfrac{-x/y^2}{1+(x/y)^2} = \dfrac{-x}{x^2+y^2}$。
  $x'=2$，$y'=-2t$。
  \textbf{代入：}
  \begin{align*}
  \frac{\mathrm{d}z}{\mathrm{d}t} &= \left(\frac{y}{x^2+y^2}\right)(2) + \left(\frac{-x}{x^2+y^2}\right)(-2t) = \frac{2y + 2xt}{x^2+y^2} \\
  &= \frac{2(1-t^2) + 2(2t)t}{(2t)^2 + (1-t^2)^2} = \frac{2 - 2t^2 + 4t^2}{4t^2 + 1 - 2t^2 + t^4} \\
  &= \frac{2(1+t^2)}{(1+t^2)^2} = \frac{2}{1+t^2}.
  \end{align*}
\end{enumerate}
\end{solution}

\begin{exercise}[链式法则求偏导数]
在以下各题中求 $\dfrac{\partial z}{\partial u}$ 和 $\dfrac{\partial z}{\partial v}$：
\begin{shortexenum}
\shortitem{1}{$z = xe^{-y} + ye^{-x},\ x = u\sin v,\ y = v\cos u$} &
\shortitem{2}{$z = xe^y,\ x = \ln u,\ y = v$} \\
\shortitem{3}{$z = \dfrac{\cos x}{y},\ x = \dfrac{v}{u},\ y = u^2-v^2$} &
\shortitem{4}{$z = \arctan(x+y),\ x = 2u-v^2,\ y = u^2v$} \\
\shortitem{5}{$z = e^{x\sin y},\ x = uv,\ y = \ln(u+v)$} &
\shortitem{6}{$z = \arctan\dfrac{x}{y},\ x = u^2+v^2,\ y = u^2-v^2$}
\end{shortexenum}
\end{exercise}
\begin{solution}
由于把自变量完全代入后表达式会较长，因此这类题通常允许答案中同时保留 $x,y$ 与 $u,v$（或在最后说明再代回 $x,y$ 的表达式）。
\begin{enumerate}
  \item[(1)] \textbf{分别求一阶偏导：} $z_x = e^{-y} - ye^{-x}$，$z_y = -xe^{-y} + e^{-x}$。
  内层偏导：$x_u = \sin v,\ x_v = u\cos v,\ y_u = -v\sin u,\ y_v = \cos u$。
  \textbf{代入法则：}
  \(\displaystyle \frac{\partial z}{\partial u}=z_xx_u+z_yy_u=(e^{-y}-ye^{-x})\sin v+(e^{-x}-xe^{-y})(-v\sin u)\)，
  \(\displaystyle \frac{\partial z}{\partial v}=z_xx_v+z_yy_v=(e^{-y}-ye^{-x})(u\cos v)+(e^{-x}-xe^{-y})\cos u\)。

  \item[(2)] \textbf{求一阶偏导：} $z_x = e^y$，$z_y = xe^y$。
  内层偏导：$x_u = \dfrac{1}{u},\ x_v = 0,\ y_u = 0,\ y_v = 1$。
  \textbf{代入：}
  $\dfrac{\partial z}{\partial u} = e^y\left(\dfrac{1}{u}\right) + xe^y(0) = \dfrac{e^v}{u}$，
  $\dfrac{\partial z}{\partial v} = e^y(0) + xe^y(1) = e^v\ln u$。

  \item[(3)] \textbf{求一阶偏导：} $z_x = -\dfrac{\sin x}{y}$，$z_y = -\dfrac{\cos x}{y^2}$。
  内层：$x_u = -\dfrac{v}{u^2},\ x_v = \dfrac{1}{u},\ y_u = 2u,\ y_v = -2v$。
  \textbf{代入：}
  \(\displaystyle \frac{\partial z}{\partial u}=\left(-\frac{\sin x}{y}\right)\left(-\frac{v}{u^2}\right)+\left(-\frac{\cos x}{y^2}\right)(2u)=\frac{v\sin x}{u^2y}-\frac{2u\cos x}{y^2}\)，
  \(\displaystyle \frac{\partial z}{\partial v}=\left(-\frac{\sin x}{y}\right)\left(\frac{1}{u}\right)+\left(-\frac{\cos x}{y^2}\right)(-2v)=-\frac{\sin x}{uy}+\frac{2v\cos x}{y^2}\)。

  \item[(4)] \textbf{求一阶偏导：} $z_x = \dfrac{1}{1+(x+y)^2} = z_y$。
  内层：$x_u = 2,\ x_v = -2v,\ y_u = 2uv,\ y_v = u^2$。
  \textbf{代入提取公因子：}
  \(\displaystyle \frac{\partial z}{\partial u}=\frac{2}{1+(x+y)^2}+\frac{2uv}{1+(x+y)^2}=\frac{2(1+uv)}{1+(x+y)^2}\)，
  \(\displaystyle \frac{\partial z}{\partial v}=\frac{-2v}{1+(x+y)^2}+\frac{u^2}{1+(x+y)^2}=\frac{u^2-2v}{1+(x+y)^2}\)。

  \item[(5)] \textbf{求一阶偏导：} $z_x = (\sin y)e^{x\sin y}$，$z_y = (x\cos y)e^{x\sin y}$。
  内层：$x_u = v,\ x_v = u,\ y_u = \dfrac{1}{u+v} = y_v$。
  \textbf{代入：}
  \(\displaystyle \frac{\partial z}{\partial u}= e^{x\sin y}\left[ v\sin y + \frac{x\cos y}{u+v} \right],\qquad
  \frac{\partial z}{\partial v}= e^{x\sin y}\left[ u\sin y + \frac{x\cos y}{u+v} \right]\)。

  \item[(6)] \textbf{求一阶偏导：} $z_x = \dfrac{y}{x^2+y^2}$，$z_y = \dfrac{-x}{x^2+y^2}$。
  内层：$x_u = 2u,\ x_v = 2v,\ y_u = 2u,\ y_v = -2v$。
  \textbf{代入：}
  \(\displaystyle \frac{\partial z}{\partial u}=\frac{y(2u)-x(2u)}{x^2+y^2}=\frac{2u(y-x)}{x^2+y^2}=-\frac{2uv^2}{u^4+v^4}\)，
  \(\displaystyle \frac{\partial z}{\partial v}=\frac{y(2v)-x(-2v)}{x^2+y^2}=\frac{2v(y+x)}{x^2+y^2}=\frac{2u^2v}{u^4+v^4}\)。
\end{enumerate}
\end{solution}

\begin{exercise}[抽象函数链式法则推导]
设 $u = f(x,y,z),\ x = x(s,t),\ y = y(s,t),\ z = z(s,t)$。写出求复合函数 $u = f(x(s,t),y(s,t),z(s,t))$ 的偏导数 $\dfrac{\partial u}{\partial s}$ 和 $\dfrac{\partial u}{\partial t}$ 的链式法则。
\end{exercise}
\begin{solution}
\begin{enumerate}
  \item \textbf{梳理变量的拓扑路径：}
  从目标因变量 $u$ 到底层自变量 $s, t$，一共需要跨越三条并行路径，即中间变量有三个：$x, y, z$。

  \item \textbf{对 $s,t$ 应用法则展开求和}：
  顺着三个中间变量分别对 $s,t$ 求导累加：
  \begin{align*}
  \frac{\partial u}{\partial s} &= \frac{\partial f}{\partial x}\frac{\partial x}{\partial s} + \frac{\partial f}{\partial y}\frac{\partial y}{\partial s} + \frac{\partial f}{\partial z}\frac{\partial z}{\partial s},\\
  \frac{\partial u}{\partial t} &= \frac{\partial f}{\partial x}\frac{\partial x}{\partial t} + \frac{\partial f}{\partial y}\frac{\partial y}{\partial t} + \frac{\partial f}{\partial z}\frac{\partial z}{\partial t}.
  \end{align*}
\end{enumerate}
\end{solution}

\begin{exercise}[点处偏导数与全导数计算]
设 $u = f(x,y,z) = 3xy+yz,\ x = \ln s+\cos t,\ y = 1+s\sin t,\ z = st$。
\begin{shortsingleenum}
  \shortsinglechoice{(1)}{求 $\dfrac{\partial u}{\partial s}$ 和 $\dfrac{\partial u}{\partial t}$ 在 $(s,t) = (1,\pi)$ 点的值；}
  \shortsinglechoice{(2)}{假定 $s,t$ 还是 $w$ 的函数，即 $s = 1+\sin(\pi w),\ t = \pi w^2$，试求 $\left.\dfrac{\mathrm{d}u}{\mathrm{d}w}\right|_{w=1}$。}
\end{shortsingleenum}
\end{exercise}
\begin{solution}
\begin{enumerate}
  \item[(1)] \textbf{步骤 1：定位状态点坐标。}
  当 $s=1, t=\pi$ 时，代入内层表达式计算中间变量：
  $x = \ln 1 + \cos\pi = -1$；$y = 1 + 1\cdot\sin\pi = 1$；$z = 1\cdot\pi = \pi$。

  \textbf{步骤 2：在目标点处计算一阶外层偏导数。}
  $u_x = 3y \implies u_x(-1,1,\pi) = 3$。
  $u_y = 3x + z \implies u_y(-1,1,\pi) = 3(-1) + \pi = \pi - 3$。
  $u_z = y \implies u_z(-1,1,\pi) = 1$。

  \textbf{步骤 3：在目标点处计算内层偏导数。}
  $x_s = \dfrac{1}{s} \implies x_s(1,\pi) = 1$； $x_t = -\sin t \implies x_t(1,\pi) = 0$。
  $y_s = \sin t \implies y_s(1,\pi) = 0$； $y_t = s\cos t \implies y_t(1,\pi) = -1$。
  $z_s = t \implies z_s(1,\pi) = \pi$； $z_t = s \implies z_t(1,\pi) = 1$。

  \textbf{步骤 4：组装求值。}
  \begin{align*}
  \left.\frac{\partial u}{\partial s}\right|_{(1,\pi)} &= u_x x_s + u_y y_s + u_z z_s = 3(1) + (\pi-3)(0) + 1(\pi) = \pi + 3. \\
  \left.\frac{\partial u}{\partial t}\right|_{(1,\pi)} &= u_x x_t + u_y y_t + u_z z_t = 3(0) + (\pi-3)(-1) + 1(1) = 4 - \pi.
  \end{align*}

  \item[(2)] \textbf{步骤 1：计算新点映射。}
  当 $w=1$ 时：$s = 1+\sin\pi = 1$；$t = \pi(1)^2 = \pi$。此点恰好为前一小问计算完毕的点！

  \textbf{步骤 2：计算对 $w$ 的导数。}
  $s_w = \pi\cos(\pi w) \implies s_w(1) = \pi(-1) = -\pi$。
  $t_w = 2\pi w \implies t_w(1) = 2\pi$。

  \textbf{步骤 3：运用全导数链式法则。}
  利用我们在第 (1) 问已求得的 $\dfrac{\partial u}{\partial s}$ 和 $\dfrac{\partial u}{\partial t}$ 在此点的值，代入：
  \[
  \left.\frac{\mathrm{d}u}{\mathrm{d}w}\right|_{w=1}
  =\frac{\partial u}{\partial s}\frac{\mathrm{d}s}{\mathrm{d}w}
  +\frac{\partial u}{\partial t}\frac{\mathrm{d}t}{\mathrm{d}w}
  =(\pi+3)(-\pi)+(4-\pi)(2\pi)=5\pi-3\pi^2.
  \]
\end{enumerate}
\end{solution}

\begin{exercise}[抽象复合函数求导]
求下列复合函数的偏导数：\\
(1) $z = f(\sqrt{x^2+y^2})$ ~~ (2) $u = f(x^2+y^2+z^2)$ ~~(3) $z = f\left(x,\dfrac{x}{y}\right)$ ~~ (4) $u = f(x,xy,xyz)$
\end{exercise}
\begin{solution}
约定用 $f'$ 或 $f_i'$ 记作对第 $i$ 个位置变量求导的记号。
\begin{enumerate}
  \item[(1)] \textbf{分析与计算：} 中间变量只有一个 $v = \sqrt{x^2+y^2}$。有 \(\displaystyle v_x = \dfrac{x}{\sqrt{x^2+y^2}},\quad v_y = \dfrac{y}{\sqrt{x^2+y^2}}\)。因此偏导数为 \(\displaystyle z_x = f'(v) \dfrac{x}{\sqrt{x^2+y^2}},\quad z_y = f'(v) \dfrac{y}{\sqrt{x^2+y^2}}\)。

  \item[(2)] \textbf{分析与计算：} 中间变量只有一个 $v = x^2+y^2+z^2$。对三个独立变量偏导完全对称：
  $u_x = 2x f'(v)$，$u_y = 2y f'(v)$，$u_z = 2z f'(v)$。

  \item[(3)] \textbf{分析与计算：} 设第一位置变量 $u_1=x$，第二位置变量 $u_2=\dfrac{x}{y}$。于是 \(\displaystyle z_x = f_1'\cdot(1) + f_2'\cdot\left(\frac{1}{y}\right) = f_1' + \frac{1}{y}f_2'\)，且 \(\displaystyle z_y = f_1'\cdot(0) + f_2'\cdot\left(-\frac{x}{y^2}\right) = -\frac{x}{y^2}f_2'\)。

  \item[(4)] \textbf{分析与计算：} $u_1=x, u_2=xy, u_3=xyz$，所以 \(\displaystyle u_x=f_1'+yf_2'+yzf_3'\)，\(\displaystyle u_y=xf_2'+xzf_3'\)，\(\displaystyle u_z=xyf_3'\)。
\end{enumerate}
\end{solution}

\begin{exercise}[求具体函数的二阶偏导数]
求下列函数的二阶偏导数：
\par\noindent
\begin{tabular}{@{}p{0.2\textwidth}p{0.2\textwidth}p{0.2\textwidth}p{0.2\textwidth}@{}}
(1) $z = \sin\left(\dfrac{x}{y}\right)$ & (2) $z = xe^y$ & (3) $z = \sin(x^2+y^2)$ & (4) $z = \sqrt{x^2+y^2}$ \\
(5) $z = \dfrac{xe^y}{x+y}$ & (6) $z = \arctan(x+y)$ & (7) $f(x,y) = \ln(xy)$ & (8) $f(x,y) = x^y$
\end{tabular}
\end{exercise}
\begin{solution}
\begin{enumerate}
  \item[(1)] $z_x = \dfrac{1}{y}\cos\dfrac{x}{y}$, \quad $z_y = -\dfrac{x}{y^2}\cos\dfrac{x}{y}$。
  \begin{align*}
  z_{xx} &= -\frac{1}{y^2}\sin\dfrac{x}{y},\\
  z_{xy} &= -\frac{1}{y^2}\cos\dfrac{x}{y} + \frac{x}{y^3}\sin\dfrac{x}{y},\\
  z_{yy} &= \frac{2x}{y^3}\cos\dfrac{x}{y} - \frac{x^2}{y^4}\sin\dfrac{x}{y}.
  \end{align*}

  \item[(2)] $z_x = e^y$, \quad $z_y = xe^y$。
  $z_{xx} = 0$，$z_{xy} = e^y$，$z_{yy} = xe^y$。

  \item[(3)] $z_x = 2x\cos(x^2+y^2)$, \quad $z_y = 2y\cos(x^2+y^2)$。
  \(\displaystyle z_{xx}=2\cos(x^2+y^2)-4x^2\sin(x^2+y^2)\)，
  \(\displaystyle z_{xy}=-4xy\sin(x^2+y^2)\)，
  \(\displaystyle z_{yy}=2\cos(x^2+y^2)-4y^2\sin(x^2+y^2)\)。

  \item[(4)] $z_x = \dfrac{x}{\sqrt{x^2+y^2}}$, \quad $z_y = \dfrac{y}{\sqrt{x^2+y^2}}$。记 $r = \sqrt{x^2+y^2}$。
  \begin{align*}
  z_{xx} &= \frac{1\cdot r - x\cdot(x/r)}{r^2} = \frac{y^2}{r^3} = \frac{y^2}{(x^2+y^2)^{3/2}},\\
  z_{xy} &= \frac{0\cdot r - x\cdot(y/r)}{r^2} = -\frac{xy}{(x^2+y^2)^{3/2}},\\
  z_{yy} &= \frac{x^2}{(x^2+y^2)^{3/2}}.
  \end{align*}

  \item[(5)] 一阶偏导为 \(\displaystyle z_x = \frac{y e^y}{(x+y)^2},\qquad z_y = \frac{x(x+y-1)e^y}{(x+y)^2}\)。二阶偏导为
  \[
  z_{xx}=\frac{-2y e^y}{(x+y)^3},\qquad
  z_{xy}=\frac{e^y(x-y)}{(x+y)^3},\qquad
  z_{yy}=\frac{x(x^2+xy-2)e^y}{(x+y)^3}.
  \]

  \item[(6)] $z_x = \dfrac{1}{1+(x+y)^2} = z_y$。
  $z_{xx} = z_{xy} = z_{yy} = -\dfrac{2(x+y)}{(1+(x+y)^2)^2}$。

  \item[(7)] $f = \ln(xy) = \ln x + \ln y$（限定 $x,y$ 同号）。故 \(\displaystyle f_x=\frac1x\)、\(\displaystyle f_y=\frac1y\)、\(\displaystyle f_{xx}=-\frac1{x^2}\)、\(\displaystyle f_{xy}=0\)、\(\displaystyle f_{yy}=-\frac1{y^2}\)。

  \item[(8)] $f = x^y = e^{y\ln x}$。因此 \(\displaystyle f_x=yx^{y-1}\)、\(\displaystyle f_y=x^y\ln x\)，并且 \(\displaystyle f_{xx}=y(y-1)x^{y-2}\)、\(\displaystyle f_{xy}=x^{y-1}(1+y\ln x)\)、\(\displaystyle f_{yy}=x^y(\ln x)^2\)。
\end{enumerate}
\end{solution}

\begin{exercise}[二阶偏导数在原点的值]
设 $f(x,y) = a+bx+cy+\mathrm{d}x^2+exy+ky^2$，其中 $a,b,c,d,e,k$ 均为常数。验证 $a = f(0,0),\ b = \dfrac{\partial f}{\partial x}(0,0),\ c = \dfrac{\partial f}{\partial y}(0,0),\ d = \dfrac{1}{2}\dfrac{\partial^2 f}{\partial x^2}(0,0),\ e = \dfrac{\partial^2 f}{\partial x\partial y}(0,0),\ k = \dfrac{1}{2}\dfrac{\partial^2 f}{\partial y^2}(0,0)$。
\end{exercise}
\begin{solution}
\begin{enumerate}
  \item \textbf{代入原点检验函数值：}
  直接令 $x=0, y=0$ 得到：
  $f(0,0) = a + 0 + 0 + 0 + 0 + 0 = a$。得证。

  \item \textbf{求解并检验一阶偏导数在原点的值：}
  $f_x = b + 2\mathrm{d}x + ey \implies f_x(0,0) = b$。得证。
  $f_y = c + ex + 2ky \implies f_y(0,0) = c$。得证。

  \item \textbf{求解并检验二阶偏导数在原点的值：}
  $f_{xx} = \dfrac{\partial}{\partial x}(b + 2\mathrm{d}x + ey) = 2d \implies f_{xx}(0,0) = 2d \implies d = \dfrac{1}{2}f_{xx}(0,0)$。得证。
  $f_{xy} = \dfrac{\partial}{\partial y}(b + 2\mathrm{d}x + ey) = e \implies f_{xy}(0,0) = e$。得证。
  $f_{yy} = \dfrac{\partial}{\partial y}(c + ex + 2ky) = 2k \implies f_{yy}(0,0) = 2k \implies k = \dfrac{1}{2}f_{yy}(0,0)$。得证。
\end{enumerate}
这正是多元函数在原点进行二阶泰勒展开的一个体现。
\end{solution}

\begin{exercise}[偏微分方程待定常数]
设 $z = x^k e^{-\frac{y}{x}}$ 满足关系式 $\dfrac{\partial z}{\partial x} = y\dfrac{\partial^2 z}{\partial y^2} + \dfrac{\partial z}{\partial y}$，试求常数 $k$。
\end{exercise}
\begin{solution}
\begin{enumerate}
  \item \textbf{求解各项所需的偏导数：}
  \begin{align*}
  \frac{\partial z}{\partial x} &= kx^{k-1}e^{-\frac{y}{x}} + x^k e^{-\frac{y}{x}}\left(\frac{y}{x^2}\right) = x^{k-1}e^{-\frac{y}{x}}\left(k + \frac{y}{x}\right) = x^{k-2}e^{-\frac{y}{x}}(kx+y). \\
  \frac{\partial z}{\partial y} &= x^k e^{-\frac{y}{x}}\left(-\frac{1}{x}\right) = -x^{k-1}e^{-\frac{y}{x}}. \\
  \frac{\partial^2 z}{\partial y^2} &= \frac{\partial}{\partial y}\left(-x^{k-1}e^{-\frac{y}{x}}\right) = -x^{k-1}e^{-\frac{y}{x}}\left(-\frac{1}{x}\right) = x^{k-2}e^{-\frac{y}{x}}.
  \end{align*}

  \item \textbf{代入偏微分方程：}
  方程右端化简为 \(\displaystyle y\frac{\partial^2 z}{\partial y^2}+\frac{\partial z}{\partial y}=y\left(x^{k-2}e^{-\frac{y}{x}}\right)-x^{k-1}e^{-\frac{y}{x}}=x^{k-2}e^{-\frac{y}{x}}(y-x)\)。

  \item \textbf{比对系数求 $k$：}
  使方程两边相等，即要求 \(\displaystyle x^{k-2}e^{-\frac{y}{x}}(kx+y) \equiv x^{k-2}e^{-\frac{y}{x}}(y-x)\)。
  由于指数项绝不为零，消去公因式得出恒等关系：$kx+y \equiv y-x$。由此可知必定有 $k = -1$。
\end{enumerate}
\end{solution}

\begin{exercise}[验证偏微分方程]
验证函数 $u = x\varphi\left(\dfrac{y}{x}\right) + \psi\left(\dfrac{y}{x}\right)$ 满足方程 $x^2u_{xx} + 2xyu_{xy} + y^2u_{yy} = 0$。
\end{exercise}
\begin{solution}
\begin{enumerate}
  \item \textbf{改写成二元复合函数。}
  令 $s=x$，$t=\dfrac{y}{x}$，$F(s,t)=s\varphi(t)+\psi(t)$。
  则原函数就是 $u=F(s,t)$。外层函数的偏导为 \(\displaystyle F_s=\varphi(t),\ F_t=s\varphi'(t)+\psi'(t),\ F_{ss}=0,\ F_{st}=\varphi'(t),\ F_{tt}=s\varphi''(t)+\psi''(t)\)。
  内层函数的偏导为
  \begin{align*}
  s_x&=1,\quad s_y=0,\quad s_{xx}=s_{xy}=s_{yy}=0,\\
  t_x&=-\frac{y}{x^2}=-\frac{t}{x},\quad
  t_y=\frac1x,\quad
  t_{xx}=\frac{2y}{x^3}=\frac{2t}{x^2},\quad
  t_{xy}=-\frac1{x^2},\quad
  t_{yy}=0.
  \end{align*}

  \item \textbf{套用二阶复合函数结论。}
  对 $u=F(s,t)$ 使用 \(\displaystyle u_{xx}=F_{ss}s_x^2+2F_{st}s_xt_x+F_{tt}t_x^2+F_s s_{xx}+F_t t_{xx}\)，
  以及对应的 $u_{xy},u_{yy}$ 公式，得到
  \begin{align*}
  u_{xx}
  &=2\varphi'(t)\left(-\frac{t}{x}\right)
  +\bigl(x\varphi''(t)+\psi''(t)\bigr)\frac{t^2}{x^2}
  +\bigl(x\varphi'(t)+\psi'(t)\bigr)\frac{2t}{x^2}\\
  &=\frac{t^2}{x}\varphi''(t)+\frac{2t}{x^2}\psi'(t)+\frac{t^2}{x^2}\psi''(t),
  \end{align*}
  \begin{align*}
  u_{xy}
  &=\varphi'(t)\frac1x
  +\bigl(x\varphi''(t)+\psi''(t)\bigr)\left(-\frac{t}{x^2}\right)
  +\bigl(x\varphi'(t)+\psi'(t)\bigr)\left(-\frac1{x^2}\right)\\
  &=-\frac{t}{x}\varphi''(t)-\frac1{x^2}\psi'(t)-\frac{t}{x^2}\psi''(t),\\
  u_{yy}
  &=\bigl(x\varphi''(t)+\psi''(t)\bigr)\frac1{x^2}
   =\frac1x\varphi''(t)+\frac1{x^2}\psi''(t).
  \end{align*}

  \item \textbf{代入目标组合。}
  因为 $y=tx$，所以
  \begin{align*}
  x^2u_{xx}
  &=xt^2\varphi''(t)+2t\psi'(t)+t^2\psi''(t),\\
  2xyu_{xy}
  &=2x^2t\left(-\frac{t}{x}\varphi''(t)-\frac1{x^2}\psi'(t)-\frac{t}{x^2}\psi''(t)\right)\\
  &=-2xt^2\varphi''(t)-2t\psi'(t)-2t^2\psi''(t),\\
  y^2u_{yy}
  &=x^2t^2\left(\frac1x\varphi''(t)+\frac1{x^2}\psi''(t)\right)
  =xt^2\varphi''(t)+t^2\psi''(t).
  \end{align*}
  三式相加，$\varphi''(t)$、$\psi'(t)$、$\psi''(t)$ 的系数分别为
  $xt^2-2xt^2+xt^2$，$2t-2t$，$t^2-2t^2+t^2$，全部为 $0$。故
  $x^2u_{xx}+2xyu_{xy}+y^2u_{yy}=0$。
\end{enumerate}
\end{solution}

\subsubsection{B组}

\begin{note}
对于 $z=f(u,v)$，其二阶偏导的标准结构为：
\begin{align*}
z_{xx} &= f_{11}'' u_x^2 + 2f_{12}'' u_x v_x + f_{22}'' v_x^2 + f_1' u_{xx} + f_2' v_{xx}, \\
z_{xy} &= f_{11}'' u_x u_y + f_{12}'' (u_x v_y + u_y v_x) + f_{22}'' v_x v_y + f_1' u_{xy} + f_2' v_{xy}, \\
z_{yy} &= f_{11}'' u_y^2 + 2f_{12}'' u_y v_y + f_{22}'' v_y^2 + f_1' u_{yy} + f_2' v_{yy}.
\end{align*}
对于 $z=f(u,v,w)$，它的二阶偏导数依然遵循“外层二次型 + 内层曲率”的结构。对应公式为
\begin{align*}
z_{xx}
&= f_{11}'' u_x^2 + f_{22}'' v_x^2 + f_{33}'' w_x^2
   + 2f_{12}'' u_x v_x + 2f_{13}'' u_x w_x + 2f_{23}'' v_x w_x \\
&\quad + f_1' u_{xx} + f_2' v_{xx} + f_3' w_{xx}, \\
z_{yy}
&= f_{11}'' u_y^2 + f_{22}'' v_y^2 + f_{33}'' w_y^2
   + 2f_{12}'' u_y v_y + 2f_{13}'' u_y w_y + 2f_{23}'' v_y w_y \\
&\quad + f_1' u_{yy} + f_2' v_{yy} + f_3' w_{yy}, \\
z_{xy}
&= f_{11}'' u_x u_y + f_{22}'' v_x v_y + f_{33}'' w_x w_y
   + f_{12}''(u_x v_y + u_y v_x) \\
&\quad + f_{13}''(u_x w_y + u_y w_x)
   + f_{23}''(v_x w_y + v_y w_x) \\
&\quad + f_1' u_{xy} + f_2' v_{xy} + f_3' w_{xy}.
\end{align*}
\end{note}
\begin{exercise}[抽象函数的二阶偏导数]
求下列函数的二阶偏导数：
\par\noindent
\begin{tabular}{@{}p{0.3\textwidth}p{0.3\textwidth}p{0.3\textwidth}@{}}
(1) $z = f(ax,by)$ & (2) $u = f(x^2+y^2)$ & (3) $z = f(x+y,xy)$ \\
(4) $z = f(x+y,x-y)$ & (5) $z = f(xy^2,x^2y)$ & (6) $z = f(\sin x,\cos y,e^{x+y})$
\end{tabular}
\par
\end{exercise}
\begin{solution}
\begin{enumerate}
  \item[(1)] 令 $u=ax, v=by$。
  内层导数：$u_x=a, u_y=0$；$v_x=0, v_y=b$。所有内层的二阶导数均为 $0$。
  一阶：$z_x = af_1'$，$z_y = bf_2'$。
  二阶（代入结构公式，内层曲率项为0）：
于是 \(\displaystyle z_{xx}=a^2f_{11}''\)，\(\displaystyle z_{xy}=abf_{12}''\)，\(\displaystyle z_{yy}=b^2f_{22}''\)。

  \item[(2)] 令 $v=x^2+y^2$，$u=f(v)$。此为一元抽象函数复合。
  内层导数：$v_x = 2x, v_y = 2y$。内层二阶导数：$v_{xx} = 2, v_{xy} = 0, v_{yy} = 2$。
  一阶：$u_x = 2xf'(v)$，$u_y = 2yf'(v)$。
  二阶（根据一元链式法则 $u_{xx} = f''(v)v_x^2 + f'(v)v_{xx}$）：\(\displaystyle
  u_{xx}=4x^2 f''+2f',\qquad u_{xy}=4xy f'',\qquad u_{yy}=4y^2 f''+2f'\)。

  \item[(3)] 令 $u=x+y, v=xy$。
  内层一阶导数：$u_x=1, u_y=1$；$v_x=y, v_y=x$。
  内层二阶导数：$u_{xx}=u_{xy}=u_{yy}=0$；$v_{xx}=0, v_{xy}=1, v_{yy}=0$。
   一阶：$z_x = f_1' + yf_2'$，$z_y = f_1' + xf_2'$。二阶直接代入结构公式，得 \(\displaystyle z_{xx}=f_{11}''+2yf_{12}''+y^2f_{22}''\)；\(\displaystyle z_{xy}=f_{11}''+(x+y)f_{12}''+xyf_{22}''+f_2'\)；\(\displaystyle z_{yy}=f_{11}''+2xf_{12}''+x^2f_{22}''\)。

  \item[(4)] 令 $u=x+y, v=x-y$。
  内层导数：$u_x=1, u_y=1$；$v_x=1, v_y=-1$。所有内层二阶导均为 $0$。
  一阶：$z_x = f_1' + f_2'$，$z_y = f_1' - f_2'$。
  二阶：
  \begin{align*}
  z_{xx} &= f_{11}''(1)^2 + 2f_{12}''(1)(1) + f_{22}''(1)^2 = f_{11}'' + 2f_{12}'' + f_{22}'' \\
  z_{xy} &= f_{11}''(1)(1) + f_{12}''(1(-1) + 1(1)) + f_{22}''(1)(-1) = f_{11}'' - f_{22}'' \\
  z_{yy} &= f_{11}''(1)^2 + 2f_{12}''(1)(-1) + f_{22}''(-1)^2 = f_{11}'' - 2f_{12}'' + f_{22}''
  \end{align*}

  \item[(5)] 令 $u=xy^2, v=x^2y$。
  内层一阶：$u_x=y^2, u_y=2xy$；$v_x=2xy, v_y=x^2$。
  内层二阶：$u_{xx}=0, u_{xy}=2y, u_{yy}=2x$；$v_{xx}=2y, v_{xy}=2x, v_{yy}=0$。
  一阶：$z_x = y^2f_1' + 2xyf_2'$，$z_y = 2xyf_1' + x^2f_2'$。
  二阶：
  \begin{align*}
  z_{xx} &= f_{11}''(y^2)^2 + 2f_{12}''(y^2)(2xy) + f_{22}''(2xy)^2 + f_1'(0) + f_2'(2y) \\
         &= y^4 f_{11}'' + 4xy^3 f_{12}'' + 4x^2y^2 f_{22}'' + 2y f_2' \\
  z_{xy} &= f_{11}''(y^2)(2xy) + f_{12}''(y^2 x^2 + 2xy\cdot 2xy) + f_{22}''(2xy)(x^2) + f_1'(2y) + f_2'(2x) \\
         &= 2xy^3 f_{11}'' + 5x^2y^2 f_{12}'' + 2x^3y f_{22}'' + 2yf_1' + 2xf_2' \\
  z_{yy} &= f_{11}''(2xy)^2 + 2f_{12}''(2xy)(x^2) + f_{22}''(x^2)^2 + f_1'(2x) + f_2'(0) \\
         &= 4x^2y^2 f_{11}'' + 4x^3y f_{12}'' + x^4 f_{22}'' + 2xf_1'
  \end{align*}

  \item[(6)] 三元抽象函数 $z = f(u,v,w)$，令 $u=\sin x, v=\cos y, w=e^{x+y}$。
  结构公式扩展为三元完全平方展开加内层二阶导。
  内层一阶：$u_x=\cos x, u_y=0$；$v_x=0, v_y=-\sin y$；$w_x=e^{x+y}, w_y=e^{x+y}$。
  内层二阶：$u_{xx}=-\sin x, v_{yy}=-\cos y, w_{xx}=w_{xy}=w_{yy}=e^{x+y}$，其余交叉项为 $0$。
  一阶：$z_x = (\cos x)f_1' + e^{x+y}f_3'$，$z_y = (-\sin y)f_2' + e^{x+y}f_3'$。
  二阶：
  \begin{align*}
  z_{xx} &= f_{11}''(\cos x)^2 + f_{33}''(e^{x+y})^2 + 2f_{13}''(\cos x)(e^{x+y}) + f_1'(-\sin x) + f_3'(e^{x+y}) \\
         &= \cos^2 x f_{11}'' + e^{2(x+y)} f_{33}'' + 2e^{x+y}\cos x f_{13}'' - \sin x f_1' + e^{x+y}f_3' \\
  z_{xy} &= f_{12}''(\cos x)(-\sin y) + f_{13}''(\cos x)(e^{x+y}) + f_{23}''(-\sin y)(e^{x+y}) + f_{33}''(e^{x+y})^2 + f_3'(e^{x+y}) \\
         &= -\cos x\sin y f_{12}'' + e^{x+y}\cos x f_{13}'' - e^{x+y}\sin y f_{23}'' + e^{2(x+y)} f_{33}'' + e^{x+y}f_3' \\
  z_{yy} &= f_{22}''(-\sin y)^2 + f_{33}''(e^{x+y})^2 + 2f_{23}''(-\sin y)(e^{x+y}) + f_2'(-\cos y) + f_3'(e^{x+y}) \\
         &= \sin^2 y f_{22}'' + e^{2(x+y)} f_{33}'' - 2e^{x+y}\sin y f_{23}'' - \cos y f_2' + e^{x+y}f_3'
  \end{align*}
\end{enumerate}
\end{solution}
\begin{exercise}[混合偏导数不相等]
设 $f(x,y) = \begin{cases} \dfrac{x^3y-xy^3}{x^2+y^2}, & (x,y)\neq(0,0), \\ 0, & (x,y)=(0,0) \end{cases}$，试证：$f_{xy}(0,0)\neq f_{yx}(0,0)$。
\end{exercise}
\begin{solution}
\begin{enumerate}
  \item \textbf{在轴上利用差商求一阶偏导：}
  由于原点为分段点边界，我们需要利用极限定义进行差商：
  对于 $x=0$，当 $y\neq 0$ 时，直接代入 $x=0$ 可见 $f(0,y) = 0$。同理 $f(x,0) = 0$。
  现在求在轴非零点上交叉方向的导数：
  \begin{align*}
  f_x(0,y) &= \lim_{h\to 0} \frac{f(h,y)-f(0,y)}{h} = \lim_{h\to 0} \frac{h^3y - hy^3}{h(h^2+y^2)} = \lim_{h\to 0} \frac{h^2y - y^3}{h^2+y^2} = -y. \\
  f_y(x,0) &= \lim_{k\to 0} \frac{f(x,k)-f(x,0)}{k} = \lim_{k\to 0} \frac{x^3k - xk^3}{k(x^2+k^2)} = \lim_{k\to 0} \frac{x^3 - xk^2}{x^2+k^2} = x.
  \end{align*}

  \item \textbf{在原点利用差商求二阶混合偏导：}
  \(\displaystyle f_{xy}(0,0)=\lim_{k\to 0} \frac{f_x(0,k)-f_x(0,0)}{k}=-1,\qquad
  f_{yx}(0,0)=\lim_{h\to 0} \frac{f_y(h,0)-f_y(0,0)}{h}=1\)。

  \item \textbf{结论：}
  由于 $-1\neq 1$，故 $f_{xy}(0,0)\neq f_{yx}(0,0)$。
  这说明若混合偏导数缺乏适当的连续性条件，则二者未必相等。
\end{enumerate}
\end{solution}

\begin{exercise}[欧拉齐次函数定理]
若函数 $f(x,y)$ 对任意正实数 $t$ 满足关系 $f(tx,ty) = t^n f(x,y)$，则称 $f(x,y)$ 为 $n$ 次齐次函数。设 $f(x,y)$ 可微，试证明 $f(x,y)$ 为 $n$ 次齐次函数的充要条件是
$x\dfrac{\partial f}{\partial x} + y\dfrac{\partial f}{\partial y} = nf(x,y)$。
\end{exercise}
\begin{solution}
\begin{enumerate}
  \item \textbf{证明必要性（$\Rightarrow$）：}
  已知恒等式 $f(tx, ty) = t^n f(x,y)$ 对所有 $t>0$ 成立。我们将 $t$ 视作自变量对其两边求导。
  左端求导运用复合函数链式法则，右端求导则是普通一元函数导数：
  \begin{align*}
  \frac{\mathrm d}{\mathrm{d}t} f(tx,ty)
  &= \frac{\partial f}{\partial (tx)} \cdot \frac{\partial (tx)}{\partial t}
   + \frac{\partial f}{\partial (ty)} \cdot \frac{\partial (ty)}{\partial t}
   = x f_1'(tx,ty) + y f_2'(tx,ty),\\
  \frac{\mathrm d}{\mathrm{d}t} [t^n f(x,y)]
  &= n t^{n-1} f(x,y).
  \end{align*}
  两端相连得到：$x f_1'(tx,ty) + y f_2'(tx,ty) = n t^{n-1} f(x,y)$。
  特别地，由于恒等式对任意 $t$ 成立，令 $t=1$，则 $tx=x, ty=y$，代入即得 \(\displaystyle x \frac{\partial f}{\partial x}(x,y) + y \frac{\partial f}{\partial y}(x,y) = n f(x,y)\)。

  \item \textbf{证明充分性（$\Leftarrow$）：}
  已知对于空间任意点 $(X, Y)$ 都满足偏微分方程 \(\displaystyle X f_X(X,Y) + Y f_Y(X,Y) = n f(X,Y)\)。为了探究 $f(tx,ty)$ 与 $t^n f(x,y)$ 的关系，构造辅助函数 \(\displaystyle g(t)=\frac{f(tx,ty)}{t^n}=t^{-n} f(tx,ty)\)（此处将 $x, y$ 视作固定的参数）。
  对 $g(t)$ 求导以判断其是否为常数，得 \(\displaystyle
  g'(t)=-n t^{-n-1} f(tx,ty)+t^{-n}\left[x f_1'(tx,ty)+y f_2'(tx,ty)\right]\)。
  在已知等式中令 $X=tx, Y=ty$，则有 $tx f_1'(tx,ty) + ty f_2'(tx,ty) = n f(tx,ty)$。提取 $t$ 即 $x f_1' + y f_2' = \dfrac{n}{t} f(tx,ty)$。将此关系代入 $g'(t)$：\(\displaystyle g'(t) = -n t^{-n-1} f(tx,ty) + t^{-n} \left[ \frac{n}{t} f(tx,ty) \right] = 0\)。
  这表明 $g(t)$ 是与 $t$ 无关的常数，即 $g(t) \equiv g(1)$。
  由于 $g(1) = 1^{-n} f(x,y) = f(x,y)$，所以对于一切 $t>0$，有 $t^{-n}f(tx,ty) = f(x,y)$，移项即为 $f(tx,ty) = t^n f(x,y)$。这证明了它是 $n$ 次齐次函数。
\end{enumerate}
\end{solution}

\begin{exercise}[抽象函数验证偏微分方程]
验证函数 $u(x,y) = x^n f\left(\dfrac{y}{x^2}\right)$ 满足方程 $x\dfrac{\partial u}{\partial x} + 2y\dfrac{\partial u}{\partial y} = nu$。
\end{exercise}
\begin{solution}
\begin{enumerate}
  \item \textbf{利用乘积求导法则计算偏导数：}
  令中间变量 $t = \dfrac{y}{x^2}$。注意 $x$ 即在外部多项式里又在复合函数里。
  \begin{align*}
  \frac{\partial u}{\partial x} &= n x^{n-1} f(t) + x^n f'(t) \cdot \left(-\frac{2y}{x^3}\right) = n x^{n-1} f(t) - 2y x^{n-3} f'(t), \\
  \frac{\partial u}{\partial y} &= x^n f'(t) \cdot \left(\frac{1}{x^2}\right) = x^{n-2} f'(t).
  \end{align*}

  \item \textbf{将求导结果代入偏微分方程：}
  将各项乘上指定系数后加和：
  \begin{align*}
  x\frac{\partial u}{\partial x} + 2y\frac{\partial u}{\partial y} &= x\left[n x^{n-1} f(t) - 2y x^{n-3} f'(t)\right] + 2y\left[x^{n-2} f'(t)\right] \\
  &= n x^n f(t) - 2y x^{n-2} f'(t) + 2y x^{n-2} f'(t) \\
  &= n x^n f(t)
   = n \left[x^n f\left(\frac{y}{x^2}\right)\right] = nu.
  \end{align*}

  \item \textbf{化简后得出结论：}
  由于后两项互为相反数，正好相消，方程得证。
  这也表明此式具备某种特殊的广义齐次性。
\end{enumerate}
\end{solution}

% ============================================================
% §7.5 梯度与方向导数
% ============================================================

\subsection{第六节习题：梯度与方向导数}

\subsubsection{A组}

\begin{exercise}[基本概念问答]
回答下列问题：
\begin{shortexenum}
  \shortitem{1}{函数 $z=f(x,y)$ 在点 $p$ 沿方向 $l=(\cos\alpha,\sin\alpha)$ 的方向导数如何定义？} &
  \shortitem{2}{什么叫梯度？叙述它的准确定义.} \\
  \shortitem{3}{函数 $f(x,y,z)$ 在一点沿什么方向的变化率最大？最小？为零？} &
\end{shortexenum}
\end{exercise}
\begin{solution}
\begin{enumerate}
  \item[(1)] \textbf{方向导数的定义：}
  设函数 $z=f(x,y)$ 在点 $p=(x_0,y_0)$ 的某一邻域内有定义。自点 $p$ 引一条射线 $l$，其单位方向向量为 $l=(\cos\alpha, \sin\alpha)$。若极限 \(\displaystyle \lim_{\rho\to0^+}\frac{f(x_0+\rho\cos\alpha,y_0+\rho\sin\alpha)-f(x_0,y_0)}{\rho}\) 存在，则称此极限为函数 $f(x,y)$ 在点 $p$ 沿方向 $l$ 的方向导数，记作 $\dfrac{\partial f}{\partial l}\bigg|_p$。

  \item[(2)] \textbf{梯度的定义：}
  设函数 $f(x,y,z)$ 在点 $P_0(x_0,y_0,z_0)$ 处可微，则由该点处的偏导数作为分量所构成的向量
  $\left( f_x(x_0,y_0,z_0), f_y(x_0,y_0,z_0), f_z(x_0,y_0,z_0) \right)$
  称为函数 $f(x,y,z)$ 在点 $P_0$ 处的梯度（Gradient），记作 $\operatorname{grad}f(P_0)$ 或 $\nabla f(P_0)$。

  \item[(3)] \textbf{变化率的极值方向：}
  根据方向导数与梯度的关系公式 $\dfrac{\partial f}{\partial l} = \nabla f \cdot l = \|\nabla f\| \cos\theta$：
  \begin{itemize}
    \item 沿\textbf{梯度方向}（即与梯度同向，$\theta=0$），函数的变化率（增加率）最大；
    \item 沿\textbf{负梯度方向}（即与梯度反向，$\theta=\pi$），函数的变化率最小（即减小得最快）；
    \item 沿\textbf{垂直于梯度的方向}（即等值面的切方向，$\theta=\pi/2$ 或 $3\pi/2$），函数的变化率为零。
  \end{itemize}
\end{enumerate}
\end{solution}

\begin{exercise}[利用方向角计算方向导数]
在下列各题中，$\alpha$ 表示 $l$ 与 $x$ 轴正向的夹角. 试求函数 $z=x^2y$ 在点 $(1,1)$ 沿方向 $l=(\cos\alpha,\sin\alpha)$ 的方向导数：
\par\noindent
\begin{tabular}{@{}p{0.3\textwidth}p{0.3\textwidth}p{0.3\textwidth}@{}}
(1) $\alpha=0$ & (2) $\alpha=\pi$ & (3) $\alpha=\frac{\pi}{4}$ \\
(4) $\alpha=\frac{\pi}{2}$ & (5) $\alpha=\frac{3\pi}{2}$ & (6) $\alpha=\frac{5\pi}{4}$
\end{tabular}
\end{exercise}
\begin{solution}
\begin{enumerate}
  \item \textbf{先求函数在指定点处的梯度：}
  计算一阶偏导数：$z_x = 2xy$，$z_y = x^2$。
  在点 $(1,1)$ 处取值，得到梯度向量 $\nabla z(1,1) = (2(1)(1), 1^2) = (2, 1)$。

  \item \textbf{写出方向导数的一般表达式：}
  由于函数可微，任一方向 $l=(\cos\alpha, \sin\alpha)$ 的方向导数为
  $\dfrac{\partial z}{\partial l} = \nabla z \cdot l = 2\cos\alpha + \sin\alpha$。

  \item \textbf{分别代入不同的 $\alpha$ 计算具体值：}
  \begin{itemize}
    \item[(1)] $\alpha=0$：$\dfrac{\partial z}{\partial l} = 2\cos 0 + \sin 0 = 2$.
    \item[(2)] $\alpha=\pi$：$\dfrac{\partial z}{\partial l} = 2\cos\pi + \sin\pi = -2$.
    \item[(3)] $\alpha=\dfrac{\pi}{4}$：$\dfrac{\partial z}{\partial l} = 2\cos\dfrac{\pi}{4} + \sin\dfrac{\pi}{4} = 2\left(\dfrac{\sqrt{2}}{2}\right) + \dfrac{\sqrt{2}}{2} = \dfrac{3\sqrt{2}}{2}$.
    \item[(4)] $\alpha=\dfrac{\pi}{2}$：$\dfrac{\partial z}{\partial l} = 2\cos\dfrac{\pi}{2} + \sin\dfrac{\pi}{2} = 0 + 1 = 1$.
    \item[(5)] $\alpha=\dfrac{3\pi}{2}$：$\dfrac{\partial z}{\partial l} = 2\cos\dfrac{3\pi}{2} + \sin\dfrac{3\pi}{2} = 0 - 1 = -1$.
    \item[(6)] $\alpha=\dfrac{5\pi}{4}$：$\dfrac{\partial z}{\partial l} = 2\cos\dfrac{5\pi}{4} + \sin\dfrac{5\pi}{4} = 2\left(-\dfrac{\sqrt{2}}{2}\right) - \dfrac{\sqrt{2}}{2} = -\dfrac{3\sqrt{2}}{2}$.
  \end{itemize}
\end{enumerate}
\end{solution}

\begin{exercise}[计算具体点的梯度]
计算下列函数在指定点的 $\nabla f$：
\begin{shortexenum}
  \shortitem{1}{$f(x,y)=x^2y$，在 $(2,5)$ 和 $(3,1)$；} &
  \shortitem{2}{$f(x,y)=\dfrac{1}{\sqrt{x^2+y^2}}$，在 $(1,2)$ 和 $(3,0)$。}
\end{shortexenum}
\end{exercise}
\begin{solution}
\begin{enumerate}
  \item[(1)] \textbf{对于 $f(x,y)=x^2y$：}
  偏导数为 $f_x = 2xy$，$f_y = x^2$。即 $\nabla f(x,y) = (2xy, x^2)$。
  \begin{itemize}
    \item 在点 $(2,5)$ 处，$\nabla f(2,5) = (2\cdot 2\cdot 5, 2^2) = (20, 4)$。
    \item 在点 $(3,1)$ 处，$\nabla f(3,1) = (2\cdot 3\cdot 1, 3^2) = (6, 9)$。
  \end{itemize}

  \item[(2)] \textbf{对于 $f(x,y)=(x^2+y^2)^{-1/2}$：}
  应用复合函数求导法则：\(\displaystyle f_x = -\frac{1}{2}(x^2+y^2)^{-3/2}(2x) = -\frac{x}{(x^2+y^2)^{3/2}}\)，\(\displaystyle f_y = -\frac{1}{2}(x^2+y^2)^{-3/2}(2y) = -\frac{y}{(x^2+y^2)^{3/2}}\)。
  即 $\nabla f(x,y) = -\dfrac{1}{(x^2+y^2)^{3/2}}(x,y)$。
  \begin{itemize}
    \item 在点 $(1,2)$ 处，分母为 $(1^2+2^2)^{3/2} = 5^{3/2} = 5\sqrt{5}$。
          $\nabla f(1,2) = -\dfrac{1}{5\sqrt{5}}(1, 2) = \left(-\dfrac{\sqrt{5}}{25}, -\dfrac{2\sqrt{5}}{25}\right)$。
    \item 在点 $(3,0)$ 处，分母为 $(3^2+0^2)^{3/2} = 9^{3/2} = 27$。
          $\nabla f(3,0) = -\dfrac{1}{27}(3, 0) = \left(-\frac{1}{9}, 0\right)$。
  \end{itemize}
\end{enumerate}
\end{solution}

\begin{exercise}[指定向量方向上的方向导数]
求函数 $z=e^{x+2y}$ 在点 $(0,0)$ 沿方向 $A=2\boldsymbol{i}+3\boldsymbol{j}$ 的方向导数.
\end{exercise}
\begin{solution}
\begin{enumerate}
  \item \textbf{求给定点的梯度：}
  计算偏导数：
  $z_x = e^{x+2y}$，$z_y = 2e^{x+2y}$。
  代入点 $(0,0)$：
  $\nabla z(0,0) = (e^0, 2e^0) = (1, 2)$。

  \item \textbf{将方向向量单位化：}
  给定方向为 $A = (2, 3)$。其模长为 $\|A\| = \sqrt{2^2+3^2} = \sqrt{13}$。
  单位方向向量为 $l = \left(\dfrac{2}{\sqrt{13}}, \dfrac{3}{\sqrt{13}}\right)$。

  \item \textbf{计算方向导数：}
  应用点积公式计算，得
  $\dfrac{\partial z}{\partial l} = \nabla z(0,0) \cdot l
  = (1, 2) \cdot \left(\dfrac{2}{\sqrt{13}}, \dfrac{3}{\sqrt{13}}\right)
  = \dfrac{2}{\sqrt{13}} + \dfrac{6}{\sqrt{13}} = \dfrac{8}{\sqrt{13}} = \dfrac{8\sqrt{13}}{13}$。
\end{enumerate}
\end{solution}

\begin{exercise}[三元函数的梯度]
设 $f(x,y,z)=x^2+y^2+z^2$。计算在点 $(2,0,0),(0,2,0)$ 和 $(0,0,2)$ 处的梯度 $\operatorname{grad}f$。
\end{exercise}
\begin{solution}
\begin{enumerate}
  \item \textbf{求一般的梯度表达式：}
  对各个变量求偏导数：
  $f_x = 2x$, $f_y = 2y$, $f_z = 2z$。
  因此，梯度向量为 $\operatorname{grad}f(x,y,z) = (2x, 2y, 2z)$。

  \item \textbf{分别代入指定的点：}
  \begin{itemize}
    \item 在点 $(2,0,0)$ 处：$\operatorname{grad}f(2,0,0) = (2(2), 2(0), 2(0)) = (4, 0, 0)$。
    \item 在点 $(0,2,0)$ 处：$\operatorname{grad}f(0,2,0) = (2(0), 2(2), 2(0)) = (0, 4, 0)$。
    \item 在点 $(0,0,2)$ 处：$\operatorname{grad}f(0,0,2) = (2(0), 2(0), 2(2)) = (0, 0, 4)$。
  \end{itemize}
\end{enumerate}
\end{solution}

\begin{exercise}[三元函数在多个方向的方向导数]
求函数 $f(x,y,z)=x^3y^2z$ 在点 $(1,1,1)$ 沿下列方向的方向导数：
\begin{itemize}
  \item (1) $\boldsymbol{i}$，(2) $\boldsymbol{j}$，(3) $\boldsymbol{k}$。
  \item (4) $-\boldsymbol{i}$，(5) $\boldsymbol{i}+\boldsymbol{j}+\boldsymbol{k}$。
\end{itemize}
\end{exercise}
\begin{solution}
\begin{enumerate}
  \item \textbf{求指定点处的梯度：}
  偏导数为 $f_x = 3x^2y^2z$，$f_y = 2x^3yz$，$f_z = x^3y^2$。
  在点 $(1,1,1)$ 处，梯度为 $\nabla f(1,1,1) = (3, 2, 1)$。

  \item \textbf{与各给定单位方向做内积：}
  \begin{itemize}
    \item[(1)] 沿 $\boldsymbol{i} = (1,0,0)$：方向导数即为 $f_x$ 的值，为 $3$。
    \item[(2)] 沿 $\boldsymbol{j} = (0,1,0)$：方向导数即为 $f_y$ 的值，为 $2$。
    \item[(3)] 沿 $\boldsymbol{k} = (0,0,1)$：方向导数即为 $f_z$ 的值，为 $1$。
    \item[(4)] 沿 $-\boldsymbol{i} = (-1,0,0)$：方向导数 $= (3,2,1)\cdot(-1,0,0) = -3$。
    \item[(5)] 沿 $\boldsymbol{i}+\boldsymbol{j}+\boldsymbol{k} = (1,1,1)$：
          首先单位化，单位向量 $l = \dfrac{1}{\sqrt{3}}(1,1,1)$。
          方向导数 $= (3,2,1)\cdot\left(\dfrac{1}{\sqrt{3}}, \dfrac{1}{\sqrt{3}}, \dfrac{1}{\sqrt{3}}\right) = \dfrac{3+2+1}{\sqrt{3}} = \dfrac{6}{\sqrt{3}} = 2\sqrt{3}$。
  \end{itemize}
\end{enumerate}
\end{solution}

\begin{exercise}[方向导数为零的几何探究]
(1) 设 $f_x(a,b,c)=2, f_y(a,b,c)=3, f_z(a,b,c)=1$。求三个不同的单位向量 $\boldsymbol{l}$，使得 $\dfrac{\partial f(a,b,c)}{\partial l}$ 为零.\\
(2) 有多少个单位向量 $\boldsymbol{l}$ 使 $\dfrac{\partial f}{\partial l}$ 在点 $(a,b,c)$ 的值为零？
\end{exercise}
\begin{solution}
\begin{enumerate}
  \item \textbf{分析方向导数为零的条件：}
  已知该点的梯度向量为 $\nabla f = (2, 3, 1)$。
  设所求的单位向量为 $\boldsymbol{l} = (u, v, w)$，满足 $u^2+v^2+w^2=1$。
  方向导数为零等价于 $\nabla f \cdot \boldsymbol{l} = 0$，即需要寻找使得 $2u + 3v + w = 0$ 成立的单位向量。

  \item \textbf{构造三个具体的单位向量：}
  我们可以通过依次令某一个分量为 0，求解另两个分量的比例，然后归一化来找解。
  \begin{itemize}
    \item 令 $w=0$，则 $2u+3v=0$。可取 $(u,v) = (-3, 2)$。归一化得到 $\boldsymbol{l_1} = \left(-\dfrac{3}{\sqrt{13}}, \dfrac{2}{\sqrt{13}}, 0\right)$。
    \item 令 $u=0$，则 $3v+w=0$。可取 $(v,w) = (-1, 3)$。归一化得到 $\boldsymbol{l_2} = \left(0, -\dfrac{1}{\sqrt{10}}, \dfrac{3}{\sqrt{10}}\right)$。
    \item 令 $v=0$，则 $2u+w=0$。可取 $(u,w) = (1, -2)$。归一化得到 $\boldsymbol{l_3} = \left(\dfrac{1}{\sqrt{5}}, 0, -\dfrac{2}{\sqrt{5}}\right)$。
  \end{itemize}

  \item \textbf{回答第 (2) 问：}
  所有满足 $2u + 3v + w = 0$ 且 $u^2+v^2+w^2=1$ 的向量 $\boldsymbol{l} = (u,v,w)$，在几何上表示过原点且法向量为 $(2,3,1)$ 的平面与单位球面的交线。该交线是一个大圆（great circle），大圆上有\textbf{无数个}点。
  因此，有\textbf{无数个}单位向量 $\boldsymbol{l}$ 可以使方向导数为零。
\end{enumerate}
\end{solution}

\begin{exercise}[梯度的几何与物理综合应用]
在平面上任一点 $(x,y)$ 处的温度由函数 \(\displaystyle T(x,y)=\frac{100}{x^2+y^2+1}\) 表示。
\begin{shortexenum}
  \shortitem{1}{平面上何处最热？在那里温度是多少？} &
  \shortitem{2}{求温度在点 $(3,2)$ 处增加最大的方向，这一方向上的最大变化率是多少？} \\
  \shortitem{3}{求温度在点 $(3,2)$ 处减少最快的方向。} &
  \shortitem{4}{在 (3) 中所得的方向指向原点吗？} \\
  \shortitem{5}{在点 $(3,2)$ 求一个方向，使得在这个方向上，温度不增不减。} &
  \shortitem{6}{$T$ 的等高线是什么样子的？}
\end{shortexenum}
\end{exercise}
\begin{solution}
\begin{enumerate}
  \item[(1)] \textbf{寻找最热点：}
  由于分子 $100$ 为常数，要使得 $T(x,y)$ 取得最大值，分母 $x^2+y^2+1$ 必须取最小值。
  显然 $x^2+y^2 \geqslant 0$，当且仅当 $x=0, y=0$ 时分母最小为 1。
  因此平面上\textbf{原点 $(0,0)$ 处最热}，那里的温度是 $T(0,0) = 100$。

  \item[(2)] \textbf{求增加最大的方向及变化率：}
  增加最快的方向即梯度的方向。先求梯度向量 $\nabla T(x,y)$：
  \(\displaystyle T_x = -100(x^2+y^2+1)^{-2}(2x) = -\frac{200x}{(x^2+y^2+1)^2}\)，
  \(\displaystyle T_y = -100(x^2+y^2+1)^{-2}(2y) = -\frac{200y}{(x^2+y^2+1)^2}\)。
  代入点 $(3,2)$，其中 $x^2+y^2+1 = 3^2+2^2+1 = 14$，得
  $\nabla T(3,2) = \left( -\dfrac{200(3)}{14^2}, -\dfrac{200(2)}{14^2} \right)
  = \left(-\dfrac{600}{196}, -\dfrac{400}{196}\right) = \left(-\dfrac{150}{49}, -\dfrac{100}{49}\right)$。
  该向量与向量 $(-3, -2)$ 同向。因此，增加最大的方向为向量 \textbf{$(-3,-2)$ 的方向}。
  增加最大的变化率为梯度的模长：
  $\|\nabla T(3,2)\| = \sqrt{\left(-\dfrac{150}{49}\right)^2 + \left(-\dfrac{100}{49}\right)^2}
  = \dfrac{50}{49}\sqrt{(-3)^2+(-2)^2} = \dfrac{50\sqrt{13}}{49}$。

  \item[(3)] \textbf{求减少最快的方向：}
  减少最快的方向即\textbf{负梯度方向}。由 (2) 可知 $\nabla T$ 与 $(-3,-2)$ 同向，故 $-\nabla T$ 与 \textbf{$(3,2)$ 的方向} 同向。

  \item[(4)] \textbf{判断负梯度方向是否指向原点：}
  在第 (3) 问中，减少最快的方向是向着点 $(3,2)$ 在外部延伸的径向方向移动，也就是背离原点的方向。
  若要指向原点，位移向量应为 $(0,0) - (3,2) = (-3,-2)$。
  可见，\textbf{减少最快的方向（指向外侧）并没有指向原点}；反而，增加最快的方向 $(-3,-2)$ 是指向原点的。

  \item[(5)] \textbf{求温度不增不减的方向：}
  不增不减意味着方向导数为零，即需寻找与梯度垂直的方向。
  由于梯度方向为 $(-3,-2)$，在二维平面上与其点积为零的方向向量可以互换坐标并改变其中一个符号得出，即 \textbf{$(2,-3)$ 或 $(-2,3)$} 的方向。

  \item[(6)] \textbf{分析等高线的形状：}
  等高线满足方程 $T(x,y) = C$（$C$ 为正数且 $C \leqslant 100$）。
  即 $\dfrac{100}{x^2+y^2+1} = C \implies x^2+y^2 = \dfrac{100}{C} - 1$。
  因为右侧为常数，所以等高线是\textbf{以原点为中心的一族同心圆}。
\end{enumerate}
\end{solution}

\begin{exercise}[最速上升方向求解]
求下列函数在指定点处函数值增加最快的方向：
\begin{shortexenum}
  \shortitem{1}{$f(x,y)=e^x(\cos y+\sin y)$，点 $(0,0)$} &
  \shortitem{2}{$f(x,y,z)=\ln(x^2+y^2+z^2)$，点 $(2,0,1)$} \\
  \shortitem{3}{$f(x,y,z)=\cos(xyz)$，点 $\left(\frac{1}{3},\frac{1}{2},\pi\right)$} &
  \shortitem{4}{$f(x,y,z)=3x^2+4y^2$，点 $(-1,1)$}
\end{shortexenum}
\end{exercise}
\begin{solution}
函数值增加最快的方向即为其梯度 $\nabla f$ 的方向，我们只需计算该点处的梯度向量并（根据需要）单位化或提取其方向基准。
\begin{enumerate}
  \item[(1)] \textbf{求 $f(x,y)=e^x(\cos y+\sin y)$ 在 $(0,0)$ 处的梯度方向：}
  \(\displaystyle f_x=e^x(\cos y+\sin y) \implies f_x(0,0) = e^0(1+0) = 1,\qquad
  f_y=e^x(-\sin y+\cos y) \implies f_y(0,0) = e^0(0+1) = 1\)。
  梯度为 $(1,1)$。函数值增加最快的方向即为 \textbf{$(1,1)$} 的方向。

  \item[(2)] \textbf{求 $f(x,y,z)=\ln(x^2+y^2+z^2)$ 在 $(2,0,1)$ 处的梯度方向：}
  \(\displaystyle f_x=\frac{2x}{x^2+y^2+z^2} \implies f_x(2,0,1) = \frac{4}{5},\qquad
  f_y=\frac{2y}{x^2+y^2+z^2} \implies f_y(2,0,1) = 0,\qquad
  f_z=\frac{2z}{x^2+y^2+z^2} \implies f_z(2,0,1) = \frac{2}{5}\)。
  梯度为 $\left(\dfrac{4}{5}, 0, \dfrac{2}{5}\right)$，与 \textbf{$(2,0,1)$} 同向。

  \item[(3)] \textbf{求 $f(x,y,z)=\cos(xyz)$ 在 $\left(\frac{1}{3},\frac{1}{2},\pi\right)$ 处的梯度方向：}
  记该点处 $xyz = \dfrac{1}{3}\cdot\dfrac{1}{2}\cdot\pi = \dfrac{\pi}{6}$。其正弦值为 $\sin\left(\dfrac{\pi}{6}\right) = \dfrac{1}{2}$。
  \(\displaystyle f_x=-yz\sin(xyz) \implies f_x = -\frac{\pi}{4},\qquad
  f_y=-xz\sin(xyz) \implies f_y = -\frac{\pi}{6},\qquad
  f_z=-xy\sin(xyz) \implies f_z = -\frac{1}{12}\)。
  梯度为 $\left(-\dfrac{\pi}{4}, -\dfrac{\pi}{6}, -\dfrac{1}{12}\right)$。提取公因式 $-1/12$，该方向与向量 \textbf{$(-3\pi, -2\pi, -1)$} 的方向相同。

  \item[(4)] \textbf{求 $f(x,y)=3x^2+4y^2$ 在 $(-1,1)$ 处的梯度方向：}
  （注：题目原写的是三元函数 $f(x,y,z)$ 但表达式中无 $z$，因此作为二维/柱面函数处理皆可，这里按其二维形式处理）
  $f_x = 6x \implies f_x(-1,1) = -6$，$f_y = 8y \implies f_y(-1,1) = 8$。
  梯度为 $(-6, 8)$，其指向与向量 \textbf{$(-3, 4)$} 的方向相同。
\end{enumerate}
\end{solution}

\begin{exercise}[最速下降方向求解]
求下列函数在指定点处函数值减小最快的方向：
\begin{shortexenum}
\shortitem{1}{$f(x,y)=\sin(\pi xy)$，点 $\left(\frac{1}{2},\frac{2}{3}\right)$} &
\shortitem{2}{$f(x,y,z)=\dfrac{x-y}{y+z}$，点 $(-1,1,3)$}
\end{shortexenum}
\end{exercise}
\begin{solution}
函数值减小最快的方向即为其\textbf{负梯度 $-\nabla f$} 的方向。
\begin{enumerate}
  \item[(1)] \textbf{求 $\sin(\pi xy)$ 的负梯度：}
  在点 $\left(\dfrac{1}{2},\dfrac{2}{3}\right)$ 处，$\pi xy = \pi\left(\dfrac{1}{2}\right)\left(\dfrac{2}{3}\right) = \dfrac{\pi}{3}$，$\cos\left(\dfrac{\pi}{3}\right) = \dfrac{1}{2}$。
  此时 \(\displaystyle f_x=\pi y\cos(\pi xy)=\pi\left(\frac{2}{3}\right)\left(\frac{1}{2}\right)=\frac{\pi}{3}\)。同理，\(\displaystyle f_y=\pi x\cos(\pi xy)=\pi\left(\frac{1}{2}\right)\left(\frac{1}{2}\right)=\frac{\pi}{4}\)。
  梯度为 $\left(\dfrac{\pi}{3}, \dfrac{\pi}{4}\right)$。因此，减小最快的方向为\textbf{负梯度方向 $\left(-\dfrac{\pi}{3}, -\dfrac{\pi}{4}\right)$}，即与 $(-4, -3)$ 同向的方向。

  \item[(2)] \textbf{求 $\dfrac{x-y}{y+z}$ 的负梯度：}
  代入点 $(-1, 1, 3)$，此时分母 $y+z = 1+3 = 4$。
  \begin{align*}
  f_x &= \frac{1}{y+z} \implies f_x = \frac{1}{4} \\
  f_y &= \frac{-1\cdot(y+z) - (x-y)\cdot 1}{(y+z)^2} = \frac{-x-z}{(y+z)^2} \implies f_y = \frac{-(-1)-3}{4^2} = \frac{-2}{16} = -\frac{1}{8} \\
  f_z &= (x-y)\cdot(-1)(y+z)^{-2} = \frac{y-x}{(y+z)^2} \implies f_z = \frac{1-(-1)}{16} = \frac{2}{16} = \frac{1}{8}
  \end{align*}
  梯度为 $\left(\dfrac{1}{4}, -\dfrac{1}{8}, \dfrac{1}{8}\right) = \dfrac{1}{8}(2, -1, 1)$。
  减小最快的负梯度方向即为 \textbf{$(-2, 1, -1)$} 的方向。
\end{enumerate}
\end{solution}

\subsubsection{B组}

\begin{exercise}[梯度与等高线正交性]
设 $f(x,y)=x^2+y^2$。证明：如果 $(a,b)$ 是等高线 $x^2+y^2=9$ 上任一点，则在点 $(a,b)$ 的梯度 $\nabla f$ 必垂直于等高线在此点的切线.
\end{exercise}
\begin{solution}
\begin{enumerate}
  \item \textbf{求解梯度向量：}
  计算函数 $f(x,y) = x^2+y^2$ 在点 $(a,b)$ 的梯度：
  $\nabla f(a,b) = (2x, 2y)\big|_{(a,b)} = (2a, 2b)$。

  \item \textbf{求等高线的切线向量：}
  等高线方程为 $x^2+y^2=9$。对其两边对 $x$ 隐函数求导，得
  $2x + 2y \dfrac{\mathrm{d}y}{\mathrm{d}x} = 0 \implies \dfrac{\mathrm{d}y}{\mathrm{d}x} = -\dfrac{x}{y}$。
  在点 $(a,b)$ 处（假设 $b \neq 0$），切线的斜率为 $k = -\dfrac{a}{b}$。
  因此，该切线方向上的一个切向向量可以取为 $\vec{\tau} = (b, -a)$。

  \item \textbf{验证垂直关系：}
  将梯度向量与切向向量做内积运算：
  $\nabla f(a,b) \cdot \vec{\tau} = (2a, 2b) \cdot (b, -a) = 2ab - 2ab = 0$。
  内积为零，这就从解析几何的层面上严格证明了：梯度必然垂直于该点处的等高线的切线。
  （注：若 $b=0$，则点为 $(\pm3, 0)$，等高线在此处存在垂直切线即向量 $(0, 1)$，而梯度为 $(\pm6, 0)$，同样正交。）
\end{enumerate}
\end{solution}

\begin{exercise}[向径表示法计算梯度]
设 $f(x,y,z)=1/\sqrt{x^2+y^2+z^2}$，令 $\boldsymbol{r}=x\boldsymbol{i}+y\boldsymbol{j}+z\boldsymbol{k}$，试用 $\boldsymbol{r}$ 来表示 $\nabla f$。
\end{exercise}
\begin{solution}
\begin{enumerate}
  \item \textbf{引入向量及其模长的记号：}
  令 $\boldsymbol{r} = (x,y,z)$，其模长 $r = \|\boldsymbol{r}\| = \sqrt{x^2+y^2+z^2}$。
  于是我们要考察的函数可以非常紧凑地表示为关于模长的一元函数：$f = \dfrac{1}{r} = r^{-1}$。

  \item \textbf{计算距离函数 $r$ 的梯度：}
  由之前的例子（例~\ref{ex:grad-norm}）已知，距离函数 $r = \sqrt{x^2+y^2+z^2}$ 的梯度为
  $\nabla r = \left( \dfrac{x}{r}, \dfrac{y}{r}, \dfrac{z}{r} \right) = \dfrac{\boldsymbol{r}}{r}$。

  \item \textbf{利用梯度链式法则：}
  应用性质 $\nabla(\phi(r)) = \phi'(r)\nabla r$。在这里 $\phi(r) = r^{-1}$，故 $\phi'(r) = -r^{-2}$。
  因而 $\nabla f = \nabla (r^{-1}) = -r^{-2} \nabla r = -r^{-2} \left(\dfrac{\boldsymbol{r}}{r}\right) = -\dfrac{\boldsymbol{r}}{r^3}$。
\end{enumerate}
\end{solution}

\begin{exercise}[温度场的方向寻优实际应用]
顶点位于 $(1,1),(5,1),(1,3)$ 和 $(5,3)$ 的金属盘被来自于原点的火焰加热；盘上各点的温度反比于此点到原点的距离. 试问点 $(3,2)$ 处的蚂蚁应朝哪个方向爬行，才能最快到达凉处？
\end{exercise}
\begin{solution}
\begin{enumerate}
  \item \textbf{建立物理量的数学模型：}
  根据题意，温度 $T$ 反比于到原点的距离 $r = \sqrt{x^2+y^2}$。故可设
  $T(x,y) = \dfrac{k}{r} = k(x^2+y^2)^{-1/2}$
  其中比例系数 $k > 0$（原点处为火焰，所以温度随距离增加而下降，符合热源向外辐射的一般常识）。

  \item \textbf{确定“最快到达凉处”的数学含义：}
  蚂蚁需要尽快降低温度，意味着它应当沿着温度\textbf{减小最快}的方向爬行，即沿着负梯度方向 $-\nabla T$ 爬行。

  \item \textbf{计算负梯度方向：}
  沿用上一题推导的复合求导结果，由于此处为二维坐标：
  $\nabla T = \nabla(k r^{-1}) = -k r^{-2} \nabla r = -k r^{-2} \left(\dfrac{\boldsymbol{r}}{r}\right) = -k\dfrac{\boldsymbol{r}}{r^3}$，
  因此，负梯度向量为 $-\nabla T = k\dfrac{\boldsymbol{r}}{r^3}$。

  \item \textbf{代入已知点解释几何意义：}
  上述负梯度向量始终和向径 $\boldsymbol{r} = (x,y)$ 也就是从原点指向该点的向量同向。
  在点 $(3,2)$ 处，向径即为 $(3,2)$。
  这在直观上极易理解：温度源位于原点，为了最快降温，蚂蚁当然应该\textbf{沿着远离原点的射线方向}向外逃离。
  所以蚂蚁应朝 \textbf{$(3,2)$} 的方向爬行。
\end{enumerate}
\end{solution}

\begin{exercise}[曲面高程函数的最速下降]
设某山峰可由曲面 $z=5-x^2-2y^2$ 表示. 位于点 $\left(\frac{1}{2},-\frac{1}{2},\frac{17}{4}\right)$ 处的登山者发现其供氧面具漏气，他应沿哪个方向才能最快到达山底？
\end{exercise}
\begin{solution}
\begin{enumerate}
  \item \textbf{理解几何意义：}
  “到达山底”也就是要让高度函数 $z(x,y)$ 下降得最快。
  对于高程函数 $z = f(x,y)$ 来说，“下山最快的方向”就是该函数在所处位置在 $xy$ 投影平面上的\textbf{负梯度方向}。

  \item \textbf{计算高度函数的梯度：}
  $z_x = -2x$，$z_y = -4y$。
  在当前点对应的水平坐标 $(x,y) = \left(\dfrac{1}{2}, -\dfrac{1}{2}\right)$ 处，
  $z_x\left(\dfrac{1}{2}, -\dfrac{1}{2}\right) = -2\left(\dfrac{1}{2}\right) = -1$，
  $z_y\left(\dfrac{1}{2}, -\dfrac{1}{2}\right) = -4\left(-\dfrac{1}{2}\right) = 2$。
  因此高程函数的梯度为 $\nabla z = (-1, 2)$。

  \item \textbf{确定最速下降方向：}
  登山者应沿负梯度方向投影行动，即 $-\nabla z = (1, -2)$。
  所以他应当在地图投影上朝着 \textbf{$(1,-2)$} 的方向（或者说是东南侧某特定方位角）下山。
\end{enumerate}
\end{solution}

\subsection{第七节习题：隐函数微分法}

\begin{exercise}[隐函数的基本概念]
回答下列问题：
\begin{shortexenum}
  \shortitem{1}{由方程确定的隐函数，这个“隐”字是什么含义？} &
  \shortitem{2}{$n$ 个方程的方程组可确定多少个隐函数？自变量的个数如何确定？}
\end{shortexenum}
\end{exercise}
\begin{solution}
\begin{enumerate}
  \item[(1)] \textbf{“隐”字的含义：}
  在显函数关系（如 $y=f(x)$）中，因变量可以直接从式子里读出。
  而在隐函数关系中，变量间的对应规律被“藏”在方程 $F(x,y)=0$ 里，往往不能直接把 $y$ 单独解出来。
  因此“隐”强调的是：函数关系由方程约束隐含给出，而不是以显式解析式直接写出。

  \item[(2)] \textbf{隐函数及自变量个数的确定：}
  若一个方程组包含 $n$ 个相互独立的方程，并共同约束 $m+n$ 个变量，
  那么在雅可比行列式不为零等条件下，可把其中 $n$ 个变量看成另外 $m$ 个变量的函数。
  换言之，方程组可确定 $n$ 个隐函数，而剩余的 $m$ 个变量就是自变量，因为总变量数减去约束方程数恰好等于自由度数：$(m+n)-n=m$。
\end{enumerate}
\end{solution}

\begin{exercise}[单方程隐函数求导]
求下列方程确定的隐函数的导数 $\dfrac{\mathrm{d}y}{\mathrm{d}x}$ 和 $\dfrac{\mathrm{d}^2y}{\mathrm{d}x^2}$：
\par\noindent
\begin{tabular}{@{}p{0.45\textwidth}p{0.45\textwidth}@{}}
(1) $x^2 + 2xy - y^2 = a^2$ & (2) $\ln\sqrt{x^2+y^2} = \arctan\dfrac{y}{x}$ \\
(3) $y = 2x\arctan\dfrac{y}{x}$ & (4) $x^y = y^x\ (x \neq y)$
\end{tabular}
\end{exercise}
\begin{solution}
\begin{enumerate}
  \item[(1)] \textbf{求解 $x^2 + 2xy - y^2 = a^2$：}
  对方程两端关于 $x$ 求导，得
  $2x + 2y + 2xy' - 2yy' = 0 \implies x + y + (x-y)y' = 0 \implies y' = \dfrac{x+y}{y-x}$。
  继续对上式关于 $x$ 隐式求导以求二阶导（将上式写为 $1 + y' + (1-y')y' + (x-y)y'' = 0$），得
  $1 + 2y' - (y')^2 + (x-y)y'' = 0 \implies y'' = \dfrac{(y')^2 - 2y' - 1}{x-y}$。
  将 $y' = \dfrac{x+y}{y-x}$ 代入上式：
  \begin{align*}
  (y')^2 - 2y' - 1 &= \left(\frac{x+y}{y-x}\right)^2 + 2\frac{x+y}{x-y} - 1 = \frac{(x+y)^2 + 2(x+y)(x-y) - (x-y)^2}{(y-x)^2} \\
  &= \frac{x^2+2xy+y^2 + 2x^2-2y^2 - (x^2-2xy+y^2)}{(y-x)^2} = \frac{2x^2+4xy-2y^2}{(y-x)^2} \\
  &= \frac{2(x^2+2xy-y^2)}{(y-x)^2}.
  \end{align*}
  注意到原方程有 $x^2+2xy-y^2=a^2$，将其代入得分子为 $2a^2$，故 $y'' = \dfrac{2a^2}{(y-x)^3}$。

  \item[(2)] \textbf{求解 $\ln\sqrt{x^2+y^2} = \arctan\dfrac{y}{x}$：}
  原式可化为 $\dfrac{1}{2}\ln(x^2+y^2) = \arctan\dfrac{y}{x}$。两端对 $x$ 求导：\(\displaystyle \frac{x+yy'}{x^2+y^2}=\frac{1}{1+(y/x)^2}\cdot\frac{y'x-y}{x^2}=\frac{xy'-y}{x^2+y^2}\)。
  约去分母得 $x+yy' = xy'-y \implies y'(x-y) = x+y \implies y' = \dfrac{x+y}{x-y}$.
  再对 $y'$ 用商的求导法则求 $y''$：
  \begin{align*}
  y'' &= \frac{(1+y')(x-y) - (x+y)(1-y')}{(x-y)^2} = \frac{x-y+xy'-yy'-x+xy'-y+yy'}{(x-y)^2} \\
  &= \frac{2xy'-2y}{(x-y)^2} = \frac{2\left(x\frac{x+y}{x-y} - y\right)}{(x-y)^2} = \frac{2(x^2+xy - xy+y^2)}{(x-y)^3} = \frac{2(x^2+y^2)}{(x-y)^3}.
  \end{align*}

  \item[(3)] \textbf{求解 $y = 2x\arctan\dfrac{y}{x}$：}
  此方程具特殊性。若令 $u = \dfrac{y}{x}$，则有 $xu = 2x\arctan u \implies u = 2\arctan u$。
  这说明比值 $u$ 是一个只满足固定超越方程的常数 $c$。
  故 $y = cx$，这代表一条过原点的直线。对其求导直接得 $y' = c = \dfrac{y}{x}$，$y'' = 0$。

  \item[(4)] \textbf{求解 $x^y = y^x$：}
  两端取对数得 $y\ln x = x\ln y$。对 $x$ 求导得 \(\displaystyle
  y'\ln x+\frac{y}{x}=\ln y+x\frac{y'}{y}
  \implies y'\left(\ln x-\frac{x}{y}\right)=\ln y-\frac{y}{x}\)。
  因此得 $y' = \dfrac{\ln y - y/x}{\ln x - x/y} = \dfrac{y(x\ln y - y)}{x(y\ln x - x)}$。
  继续对方程 $y'\left(\ln x - \frac{x}{y}\right) = \ln y - \frac{y}{x}$ 两端对 $x$ 求导，并展开移项合并含 $y''$ 的项：
  \begin{align*}
  y''\left(\ln x - \frac{x}{y}\right) + y'\left(\frac{1}{x} - \frac{y - xy'}{y^2}\right)
  &= \frac{y'}{y} - \frac{y'x - y}{x^2},\\
  y''\left(\ln x - \frac{x}{y}\right) &= \frac{y'}{y} - \frac{y'}{x} + \frac{y}{x^2} - \frac{y'}{x} + \frac{y'}{y} - x\frac{(y')^2}{y^2} \\
  &= \frac{2y'}{y} - \frac{2y'}{x} + \frac{y}{x^2} - x\left(\frac{y'}{y}\right)^2.
  \end{align*}
  解得：$y'' = \dfrac{\dfrac{2y'}{y} - \dfrac{2y'}{x} + \dfrac{y}{x^2} - x\left(\dfrac{y'}{y}\right)^2}{\ln x - \dfrac{x}{y}}$ （其中 $y'$ 代入上式结果）。
\end{enumerate}
\end{solution}

\begin{exercise}[三元单方程隐函数的二阶偏导数]
设 $x + y + z = \mathrm{e}^z$，求 $\dfrac{\partial^2 z}{\partial x^2}$，$\dfrac{\partial^2 z}{\partial x\partial y}$ 和 $\dfrac{\partial^2 z}{\partial y^2}$。
\end{exercise}
\begin{solution}
\begin{enumerate}
  \item \textbf{求一阶偏导数：}
  设 $F(x,y,z) = x+y+z-\mathrm{e}^z = 0$。根据公式或直接偏导：
  对方程两端关于 $x$ 隐式求导（视 $y$ 为常数）：\(\displaystyle 1 + z_x = \mathrm{e}^z z_x \implies z_x(\mathrm{e}^z - 1) = 1 \implies z_x = \frac{1}{\mathrm{e}^z - 1}\)。
  同理，因为原方程对 $x, y$ 完全对称，$z_y = \dfrac{1}{\mathrm{e}^z - 1}$。

  \item \textbf{求二阶纯偏导数：}
  对 $z_x$ 再次关于 $x$ 求导，利用复合函数求导法：
  \[ z_{xx} = \frac{\partial}{\partial x}\left( (\mathrm{e}^z - 1)^{-1} \right) = -(\mathrm{e}^z - 1)^{-2} \cdot \mathrm{e}^z \cdot z_x = \frac{-\mathrm{e}^z}{(\mathrm{e}^z - 1)^2} \cdot \frac{1}{\mathrm{e}^z - 1} = -\frac{\mathrm{e}^z}{(\mathrm{e}^z - 1)^3}. \]
  同理对 $y$ 的情况，$z_{yy} = -\dfrac{\mathrm{e}^z}{(\mathrm{e}^z - 1)^3}$。

  \item \textbf{求二阶混合偏导数：}
  对 $z_x$ 关于 $y$ 求导：
  \[ z_{xy} = \frac{\partial}{\partial y}\left( (\mathrm{e}^z - 1)^{-1} \right) = -(\mathrm{e}^z - 1)^{-2} \cdot \mathrm{e}^z \cdot z_y = \frac{-\mathrm{e}^z}{(\mathrm{e}^z - 1)^2} \cdot \frac{1}{\mathrm{e}^z - 1} = -\frac{\mathrm{e}^z}{(\mathrm{e}^z - 1)^3}. \]
  （可见三个二阶偏导数均相等）。
\end{enumerate}
\end{solution}

\begin{exercise}[方程组隐函数求导]
设 $x + y + z = 0$，$x^2 + y^2 + z^2 = 1$，求 $\dfrac{\mathrm{d}x}{\mathrm{d}z}$ 和 $\dfrac{\mathrm{d}y}{\mathrm{d}z}$。
\end{exercise}
\begin{solution}
\begin{enumerate}
  \item \textbf{确定自变量和因变量：}
  根据要求，自变量为 $z$，因变量为 $x$ 和 $y$。

  \item \textbf{对整个方程组关于 $z$ 求导：}
  \[
  \begin{cases}
  x' + y' + 1 = 0 \\
  2xx' + 2yy' + 2z = 0 \implies xx' + yy' + z = 0
  \end{cases}
  \]

  \item \textbf{解线性代数方程组：}
  由第一式得到 $y' = -1 - x'$。将其代入第二式：\(\displaystyle xx'+y(-1-x')+z=0\implies x'(x-y)-y+z=0\implies x'=\frac{y-z}{x-y}\)。
  再将 $x'$ 回代求解 $y'$：
  \(\displaystyle y' = -1 - \frac{y-z}{x-y} = \frac{-(x-y) - (y-z)}{x-y} = \frac{-x+y-y+z}{x-y} = \frac{z-x}{x-y}\)。

  \item \textbf{给出结果：}
  即 $\dfrac{\mathrm{d}x}{\mathrm{d}z} = \dfrac{y-z}{x-y}$，$\dfrac{\mathrm{d}y}{\mathrm{d}z} = \dfrac{z-x}{x-y}$。
\end{enumerate}
\end{solution}

\begin{exercise}[抽象隐函数求导]
设 $x^2 + y^2 + z^2 = yf\left(\dfrac{z}{y}\right)$，求 $z_x$ 和 $z_y$。
\end{exercise}
\begin{solution}
\begin{enumerate}
  \item \textbf{关于 $x$ 隐式求导：}
  两端关于 $x$ 求偏导（注意 $y$ 为常数参量）：
  \(\displaystyle 2x + 2z z_x = y f'\left(\frac{z}{y}\right) \cdot \frac{z_x}{y} = f'\left(\frac{z}{y}\right) z_x\)。
  移项合并含 $z_x$ 的项：
  \(\displaystyle z_x\left(2z - f'\left(\frac{z}{y}\right)\right) = -2x \implies z_x = \frac{2x}{f'\left(\dfrac{z}{y}\right) - 2z}\)。

  \item \textbf{关于 $y$ 隐式求导：}
  两端关于 $y$ 求偏导（这里需使用乘积和商的复合求导法则）：
  \(\displaystyle 2y + 2z z_y = f\left(\frac{z}{y}\right) + y f'\left(\frac{z}{y}\right) \cdot \frac{z_y y - z}{y^2} = f\left(\frac{z}{y}\right) + f'\left(\frac{z}{y}\right) \frac{y z_y - z}{y}\)。
  展开整理：
  \(\displaystyle 2y + 2z z_y = f\left(\frac{z}{y}\right) + f'\left(\frac{z}{y}\right) z_y - f'\left(\frac{z}{y}\right)\frac{z}{y}\)。
  移项合并含 $z_y$ 的项：
  \(\displaystyle z_y\left(2z - f'\left(\frac{z}{y}\right)\right) = f\left(\frac{z}{y}\right) - f'\left(\frac{z}{y}\right)\frac{z}{y} - 2y\)。
  解得：
  \(\displaystyle z_y = \frac{f\left(\dfrac{z}{y}\right) - f'\left(\dfrac{z}{y}\right)\dfrac{z}{y} - 2y}{2z - f'\left(\dfrac{z}{y}\right)} = \frac{2y + f'\left(\dfrac{z}{y}\right)\dfrac{z}{y} - f\left(\dfrac{z}{y}\right)}{f'\left(\dfrac{z}{y}\right) - 2z}\)。
\end{enumerate}
\end{solution}

\begin{exercise}[抽象隐函数求导二]
设 $F(x-y, y-z, z-x) = 0$，求 $\dfrac{\partial z}{\partial x}$ 和 $\dfrac{\partial z}{\partial y}$。
\end{exercise}
\begin{solution}
\begin{enumerate}
  \item \textbf{引入中间变量并求偏导：}
  令中间变量 $u = x-y,\ v = y-z,\ w = z-x$。根据复合函数求导与隐函数求导的结合，对方程两边分别关于 $x$ 和 $y$ 求偏导（将 $z$ 视为 $x,y$ 的函数）：

  \item \textbf{关于 $x$ 求导：}
  \begin{align*}
  F_1\frac{\partial u}{\partial x}+F_2\frac{\partial v}{\partial x}+F_3\frac{\partial w}{\partial x}&=0,\\
  F_1(1)+F_2(-z_x)+F_3(z_x-1)&=0,\\
  z_x(F_3-F_2)+F_1-F_3&=0
  \implies z_x=\frac{F_3-F_1}{F_3-F_2}.
  \end{align*}

  \item \textbf{关于 $y$ 求导：}
  \begin{align*}
  F_1\frac{\partial u}{\partial y}+F_2\frac{\partial v}{\partial y}+F_3\frac{\partial w}{\partial y}&=0,\\
  F_1(-1)+F_2(1-z_y)+F_3z_y&=0,\\
  z_y(F_3-F_2)-F_1+F_2&=0
  \implies z_y=\frac{F_1-F_2}{F_3-F_2}.
  \end{align*}
\end{enumerate}
\end{solution}

\begin{exercise}[隐函数偏导数的轮换对称性]
设 $x = x(y,z)$，$y = y(x,z)$，$z = z(x,y)$ 都是由方程 $F(x,y,z) = 0$ 所确定的具有连续偏导数的函数，证明 $\dfrac{\partial x}{\partial y} \cdot \dfrac{\partial y}{\partial z} \cdot \dfrac{\partial z}{\partial x} = -1$。
\end{exercise}
\begin{solution}
\begin{enumerate}
  \item \textbf{利用单方程隐函数定理写出各偏导：}
  根据隐函数存在的条件（假设 $F_x, F_y, F_z$ 均不为零），可以将对应的变量分别视为隐函数。其偏导数公式如下：
  当 $x = x(y,z)$ 时，有 $\dfrac{\partial x}{\partial y} = -\dfrac{F_y}{F_x}$。
  当 $y = y(x,z)$ 时，有 $\dfrac{\partial y}{\partial z} = -\dfrac{F_z}{F_y}$。
  当 $z = z(x,y)$ 时，有 $\dfrac{\partial z}{\partial x} = -\dfrac{F_x}{F_z}$。

  \item \textbf{连乘证明结果：}
  将这三个式子相乘：
  $\dfrac{\partial x}{\partial y} \cdot \dfrac{\partial y}{\partial z} \cdot \dfrac{\partial z}{\partial x}
  = \left(-\dfrac{F_y}{F_x}\right) \cdot \left(-\dfrac{F_z}{F_y}\right) \cdot \left(-\dfrac{F_x}{F_z}\right)
  = (-1)^3 \dfrac{F_y F_z F_x}{F_x F_y F_z} = -1$。
  结论得证。这在热力学中被称为三重积规则（Triple product rule）。
\end{enumerate}
\end{solution}

\begin{exercise}[方程组隐函数的高阶导数]
设 $x^2 + y^2 = \dfrac{1}{2}z^2$，$x + y + z = 2$，求 $\dfrac{\mathrm{d}x}{\mathrm{d}z}$ 和 $\dfrac{\mathrm{d}^2x}{\mathrm{d}z^2}$。
\end{exercise}
\begin{solution}
\begin{enumerate}
  \item \textbf{求一阶导数：}
  将方程组整体关于自变量 $z$ 求导，得
  \(\displaystyle 2xx'+2yy'=z,\quad x'+y'+1=0\implies y'=-1-x'\)。
  将 $y'$ 表达式代入上式消去 $y'$，得
  $2x x' + 2y(-1 - x') = z \implies 2x'(x-y) - 2y = z \implies x' =\dfrac{\mathrm{d}x}{\mathrm{d}z} = \dfrac{z+2y}{2(x-y)}$。

  \item \textbf{求二阶导数准备工作：}
  对方程组的导数形式再关于 $z$ 求一次导，得
  \(\displaystyle x''+y''=0\implies y''=-x'',\quad 2(x')^2+2xx''+2(y')^2+2yy''=1\)。
  将 $y''$ 代入第二式，得
  $2(x')^2 + 2x x'' + 2(y')^2 - 2y x'' = 1 \implies 2(x-y)x'' = 1 - 2(x')^2 - 2(y')^2$。

  \item \textbf{代入并化简一阶导平方和：}
  由于 $y' = -1 - x'$，将其代入化简右侧项：\(\displaystyle 1 - 2(x')^2 - 2(-1-x')^2=1 - 2(x')^2 - 2(1 + 2x' + (x')^2)=-1 - 4x' - 4(x')^2=-(1 + 2x')^2\)。
  因此有 $x'' = -\dfrac{(1+2x')^2}{2(x-y)}$。

  \item \textbf{最后化简：}
  注意到 $1+2x' = 1 + \dfrac{z+2y}{x-y} = \dfrac{x-y+z+2y}{x-y} = \dfrac{x+y+z}{x-y}$。
  又原方程给出 $x+y+z=2$，故 $1+2x' = \dfrac{2}{x-y}$。
  代入上式最终得到
  $x'' = -\dfrac{(2/(x-y))^2}{2(x-y)} = -\dfrac{4/(x-y)^2}{2(x-y)} = -\dfrac{2}{(x-y)^3}$。
\end{enumerate}
\end{solution}

\begin{exercise}[抽象方程组隐函数求导]
设
\[
\begin{cases}
u = f(xu, v+y), \\
v = g(u-x, v^2y),
\end{cases}
\]
其中 $f,g$ 具有一阶连续偏导数。求 $\dfrac{\partial u}{\partial x}$ 和 $\dfrac{\partial v}{\partial x}$。
\end{exercise}
\begin{solution}
\begin{enumerate}
  \item \textbf{对整个方程组关于 $x$ 隐式求导：}
  将 $u,v$ 均视作 $x,y$ 的函数。设 $f$ 和 $g$ 对其第一、第二位置变量的偏导数分别为 $f_1', f_2'$ 和 $g_1', g_2'$。
  对第一个方程关于 $x$ 求偏导：
  \(\displaystyle u_x = f_1' \cdot (1\cdot u + x u_x) + f_2' \cdot (v_x)\)。
  整理将未知数 $u_x, v_x$ 放到左边，得 \(\displaystyle (1 - x f_1')u_x - f_2' v_x = u f_1'\)，记为 (1)。
  对第二个方程关于 $x$ 求偏导：
  \(\displaystyle v_x = g_1' \cdot (u_x - 1) + g_2' \cdot (2v v_x y)\)。
  整理未知数项，得 \(\displaystyle -g_1' u_x + (1 - 2vy g_2')v_x = -g_1'\)，记为 (2)。
  \item \textbf{解线性方程组求偏导：}
  将 (1) 和 (2) 组成关于未知数 $u_x, v_x$ 的二元一次方程组，其系数行列式 $\Delta$ 为：
  \[
  \Delta = \begin{vmatrix} 1-x f_1' & -f_2' \\ -g_1' & 1-2vy g_2' \end{vmatrix}
  = (1-x f_1')(1-2vy g_2') - f_2' g_1'.
  \]
  利用克莱姆法则求解 $u_x$：
  \[
  u_x = \frac{1}{\Delta} \begin{vmatrix} u f_1' & -f_2' \\ -g_1' & 1-2vy g_2' \end{vmatrix}
  = \frac{u f_1'(1-2vy g_2') - f_2' g_1'}{(1-x f_1')(1-2vy g_2') - f_2' g_1'}.
  \]
  求解 $v_x$：
  \[
  v_x = \frac{1}{\Delta} \begin{vmatrix} 1-x f_1' & u f_1' \\ -g_1' & -g_1' \end{vmatrix}
  = \frac{-g_1'(1-x f_1') - u f_1'(-g_1')}{\Delta}
  = \frac{g_1'(u f_1' + x f_1' - 1)}{(1-x f_1')(1-2vy g_2') - f_2' g_1'}.
  \]
\end{enumerate}
\end{solution}

\subsection{第八节习题：Taylor 公式}

\begin{exercise}[基础：Taylor 展开计算]
\begin{shortsingleenum}
  \shortsingleitem{1}{将 $f(x,y)=\mathrm{e}^{x}\cos(y)$ 在 $(0,0)$ 处展出到二阶（含 Lagrange 余项）；}
  \shortsingleitem{2}{利用二阶 Taylor 展开求 $\sqrt{(1.02)^3+(1.97)^3}$ 的近似值，并给出误差的大致量级估计；}
  \shortsingleitem{3}{设 $f(x,y)=\ln(1+x+2y)$，在 $(0,0)$ 处求二阶 Taylor 多项式 $P_2(x,y)$，使得 $|f(x,y)-P_2(x,y)|=o(x^2+y^2)$ 当 $(x,y)\to(0,0)$。}
\end{shortsingleenum}
\end{exercise}
\begin{solution}
\begin{enumerate}
  \item[(1)] \textbf{求 $\mathrm{e}^{x}\cos(y)$ 的二阶展开：}
  在 $(0,0)$ 处有 $f(0,0) = 1$。
  一阶偏导数为 $f_x = \mathrm{e}^x\cos y,\allowbreak\ f_y = -\mathrm{e}^x\sin y$，
  因此 $f_x(0,0)=1,\allowbreak\ f_y(0,0)=0$。
  二阶偏导数为 $f_{xx} = \mathrm{e}^x\cos y,\allowbreak\ f_{xy} = -\mathrm{e}^x\sin y,\allowbreak\ f_{yy} = -\mathrm{e}^x\cos y$，
  因此 $f_{xx}(0,0)=1,\allowbreak\ f_{xy}(0,0)=0,\allowbreak\ f_{yy}(0,0)=-1$。
  Taylor 多项式为 $P_2(x,y) = 1 + x + \dfrac{1}{2}(x^2 - y^2)$。
  Lagrange 余项为（其中 $\theta \in (0,1)$）：
  \begin{align*}
  \mathcal{R}_2 &= \frac{1}{3!} \left( x\frac{\partial}{\partial x} + y\frac{\partial}{\partial y} \right)^3 f(\theta x, \theta y) \\
  &= \frac{1}{6} \left( x^3 \mathrm{e}^{\theta x}\cos(\theta y) - 3x^2y \mathrm{e}^{\theta x}\sin(\theta y) - 3xy^2 \mathrm{e}^{\theta x}\cos(\theta y) + y^3 \mathrm{e}^{\theta x}\sin(\theta y) \right).
  \end{align*}

  \item[(2)] \textbf{求 $\sqrt{(1.02)^3+(1.97)^3}$ 的近似值：}
  设 $f(x,y) = \sqrt{x^3+y^3}$，在 $p_0 = (1,2)$ 处展开，$\boldsymbol{h} = (0.02, -0.03)$。
  $f(1,2) = \sqrt{1+8} = 3$。
  计算一阶与二阶偏导数：
  \begin{align*}
  f_x &= \dfrac{3x^2}{2\sqrt{x^3+y^3}}, & f_x(1,2) &= \dfrac{3}{2(3)} = 0.5,\\
  f_y &= \dfrac{3y^2}{2\sqrt{x^3+y^3}}, & f_y(1,2) &= \dfrac{12}{2(3)} = 2,\\
  f_{xx} &= \dfrac{3x}{\sqrt{x^3+y^3}} - \dfrac{9x^4}{4(x^3+y^3)^{3/2}},\\
  f_{xx}(1,2) &= 1 - \dfrac{9}{108} = \dfrac{11}{12},\\
  f_{xy} &= -\dfrac{9x^2y^2}{4(x^3+y^3)^{3/2}}, & f_{xy}(1,2) &= -\dfrac{36}{108} = -\dfrac{1}{3},\\
  f_{yy} &= \dfrac{3y}{\sqrt{x^3+y^3}} - \dfrac{9y^4}{4(x^3+y^3)^{3/2}},\\
  f_{yy}(1,2) &= 2 - \dfrac{144}{108} = \dfrac{2}{3}.
  \end{align*}
  代入二阶 Taylor 公式：
  \begin{align*}
  f &\approx 3 + [0.5(0.02) + 2(-0.03)] + \frac{1}{2}\left[ \frac{11}{12}(0.02)^2 + 2\left(-\frac{1}{3}\right)(0.02)(-0.03) + \frac{2}{3}(-0.03)^2 \right] \\
  &= 3 - 0.05 + \frac{1}{2}\left[ \frac{11}{12}(0.0004) + 0.0004 + 0.0006 \right] \\
  &= 2.95 + \frac{1}{2}\left[ 0.0003666\dots + 0.0010 \right] \\
  &= 2.95 + 0.00068333\dots \approx 2.95068.
  \end{align*}
  误差量级估计：Lagrange 余项 $\mathcal{R}_2 = O(\|\boldsymbol{h}\|^3)$。因为 $\|\boldsymbol{h}\| = \sqrt{0.02^2+(-0.03)^2} \approx 0.036$，所以误差大约在 $O(10^{-5})$ 的量级。

  \item[(3)] \textbf{利用一元级数间接求 Taylor 多项式：}
  由于求偏导较为繁琐，我们可以利用已知的一元 Maclaurin 展开式 $\ln(1+u) = u - \dfrac{u^2}{2} + o(u^2)$。
  令 $u = x+2y$，当 $(x,y) \to (0,0)$ 时 $u \to 0$。
  将其代入展开式，得 $f(x,y) = (x+2y) - \dfrac{1}{2}(x+2y)^2 + o((x+2y)^2) = x + 2y - \dfrac{1}{2}x^2 - 2xy - 2y^2 + o(x^2+y^2)$。
  （这里因 $(x+2y)^2 \leqslant 2(x^2+4y^2) \leqslant 8(x^2+y^2)$，故 $o((x+2y)^2) = o(x^2+y^2)$）
  因此，二阶 Taylor 多项式为 $P_2(x,y) = x + 2y - \dfrac{1}{2}x^2 - 2xy - 2y^2$。
\end{enumerate}
\end{solution}

\begin{exercise}[提高：三元函数的高阶展开]
设 $u(x,y,z)=\sin(x+y+z)$。
\begin{shortexenum}
  \shortitem{1}{计算 $u$ 在 $(0,0,0)$ 处的所有三阶偏导数；} &
  \shortitem{2}{写出三阶 Taylor 展开式并给出余项形式；} \\
  \shortitem{3}{利用此展开式估算 $\sin(0.1+0.2-0.05)$ 的值，并与精确值比较。} &
\end{shortexenum}
\end{exercise}
\begin{solution}
\begin{enumerate}
  \item[(1)] \textbf{计算三阶偏导数：}
  令 $w = x+y+z$。则 $u = \sin w$。
  无论对 $x,y,z$ 进行何种组合的偏导求导，由链式法则可知其本质等价于对 $w$ 连续求导。
  因此，所有的三阶偏导数都满足 $\dfrac{\partial^3 u}{\partial x^i \partial y^j \partial z^k} = \dfrac{\mathrm{d}^3}{\mathrm{d}w^3}(\sin w) = -\cos w$。
  在 $(0,0,0)$ 处，$w=0$，故所有的三阶偏导数值均为 $-\cos(0) = -1$。

  \item[(2)] \textbf{写出三阶 Taylor 展开式：}
  由于所有对应高阶导数均等于一元函数 $\sin w$ 的导数展开，我们可以直接写出
  $u(x,y,z) = (x+y+z) - \dfrac{1}{3!}(x+y+z)^3 + \mathcal{R}_3$。
  其中 $\mathcal{R}_3 = \dfrac{\sin(\theta(x+y+z))}{4!}(x+y+z)^4$（$\theta \in (0,1)$）为 Lagrange 余项。

  \item[(3)] \textbf{估算并比较误差：}
  估算 $\sin(0.1+0.2-0.05) = \sin(0.25)$。
  代入三阶泰勒展开式得到 $u \approx 0.25 - \dfrac{(0.25)^3}{6} = 0.25 - \dfrac{0.015625}{6} = 0.25 - 0.00260416\dots \approx 0.247396$。
  精确值为 $\sin(0.25) \approx 0.247404$。
  绝对误差约为 $0.000008$，与余项 $\mathcal{R}_3$ 预估的量级 $O(0.25^4) \approx O(10^{-3})$ 高度吻合且近似效果极佳。
\end{enumerate}
\end{solution}

\begin{exercise}[Taylor 级数的收敛性与解析延拓]
考虑复变函数 $f(z)=\dfrac{1}{1-z}$（$z\in\mathbb C,\;|z|<1$）。
\begin{shortsingleenum}
  \shortsingleitem{1}{求 $f$ 在 $z=0$ 处的 Taylor 级数 $\displaystyle f(z)=\sum_{n=0}^{\infty}z^n$（提示：几何级数求和）；}
  \shortsingleitem{2}{该级数的收敛半径是多少？与 $f(z)$ 的奇点有何关系？}
  \shortsingleitem{3}{若截断到 $N$ 项得到部分和 $S_N(z)$，证明误差 $|f(z)-S_N(z)|\leqslant\dfrac{|z|^{N+1}}{1-|z|}$。这说明了什么？（提示：这与实 Taylor 多项式的余项有何异同？）}
\end{shortsingleenum}
\end{exercise}
\begin{solution}
\begin{enumerate}
  \item[(1)] \textbf{求 Taylor 级数：}
  根据等比数列求和公式，当 $|z|<1$ 时，公比为 $z$ 的无穷等比数列和为 $\dfrac{1}{1-z}$。
  故 $f(z)$ 在 $z=0$ 处的 Taylor 级数即为 $f(z) = \sum_{n=0}^{\infty} z^n = 1 + z + z^2 + \cdots + z^n + \cdots$。

  \item[(2)] \textbf{收敛半径与奇点的关系：}
  该级数由等比数列得来，显然收敛半径为 $R=1$。
  考察函数 $f(z) = \dfrac{1}{1-z}$，其在复平面上唯一的奇点位于 $z=1$。
  从展开中心 $z=0$ 到最近奇点 $z=1$ 的距离恰好为 $|1-0|=1$。这说明复变函数 Taylor 级数的收敛圆恰好是以展开点为圆心，延伸至触碰到最近的奇点为止。

  \item[(3)] \textbf{截断误差分析：}
  部分和为 $S_N(z) = \sum_{n=0}^N z^n$。根据等比数列求和公式，\(\displaystyle f(z) - S_N(z) = \sum_{n=0}^{\infty} z^n - \sum_{n=0}^N z^n = \sum_{n=N+1}^{\infty} z^n = \frac{z^{N+1}}{1-z}\)。
  对其取绝对值并放缩，得 $|f(z) - S_N(z)| = \dfrac{|z|^{N+1}}{|1-z|}$。
  由三角不等式 $|1-z| \ge 1 - |z|$，故 $\dfrac{1}{|1-z|} \leqslant \dfrac{1}{1-|z|}$，于是得到 $|f(z) - S_N(z)| \leqslant \dfrac{|z|^{N+1}}{1-|z|}$。
  \textbf{结论说明：} 当 $|z|<1$ 时，$\lim_{N \to \infty} |z|^{N+1} = 0$，由此严格证明了级数在圆内处处收敛。
  复变解析函数的 Taylor 级数误差严格趋于 0，这意味着函数就是其 Taylor 级数（即具有解析性）；而实变函数中，无限可微但非解析的函数（如 $\mathrm{e}^{-1/x^2}$）的 Taylor 余项在某点领域可能并不收敛到 0。这突显了复变解析函数的极高刚性。
\end{enumerate}
\end{solution}

\subsubsection{A组}

\begin{exercise}[基本 Taylor 展开一]
求 $z = \sin2x + \cos y$ 在点 $(0,0)$ 处的二阶泰勒多项式.
\end{exercise}
\begin{solution}
\begin{enumerate}
  \item \textbf{求展开点处的函数值与偏导数：}
  在 $(0,0)$ 处，函数值 $z(0,0) = 0 + 1 = 1$。
  一阶偏导数：
  $z_x = 2\cos 2x \implies z_x(0,0) = 2$。
  $z_y = -\sin y \implies z_y(0,0) = 0$。

  \item \textbf{求展开点处的二阶偏导数：}
  二阶纯偏导数：
  $z_{xx} = -4\sin 2x \implies z_{xx}(0,0) = 0$。
  $z_{yy} = -\cos y \implies z_{yy}(0,0) = -1$。
  二阶混合偏导数：
  $z_{xy} = 0 \implies z_{xy}(0,0) = 0$。

  \item \textbf{组装二阶泰勒多项式：}
  根据泰勒公式 $P_2(x,y) = f(0,0) + f_x x + f_y y + \dfrac{1}{2}(f_{xx}x^2 + 2f_{xy}xy + f_{yy}y^2)$，得
  $P_2(x,y) = 1 + 2x + 0\cdot y + \dfrac{1}{2}\left( 0\cdot x^2 + 0\cdot xy + (-1)y^2 \right) = 1 + 2x - \dfrac{1}{2}y^2$。
\end{enumerate}
\end{solution}

\begin{exercise}[基本 Taylor 展开二]
求 $z = \dfrac{1}{xy}$ 在点 $(1,2)$ 处的二阶泰勒多项式.
\end{exercise}
\begin{solution}
\begin{enumerate}
  \item \textbf{求展开点处的函数值与偏导数：}
  在 $(1,2)$ 处，函数值 $z(1,2) = \dfrac{1}{2}$。
  一阶偏导数：
  $z_x = -\dfrac{1}{x^2y} \implies z_x(1,2) = -\dfrac{1}{2}$。
  $z_y = -\dfrac{1}{xy^2} \implies z_y(1,2) = -\dfrac{1}{4}$。

  \item \textbf{求展开点处的二阶偏导数：}
  二阶纯偏导数：
  $z_{xx} = \dfrac{2}{x^3y} \implies z_{xx}(1,2) = \dfrac{2}{2} = 1$。
  $z_{yy} = \dfrac{2}{xy^3} \implies z_{yy}(1,2) = \dfrac{2}{8} = \dfrac{1}{4}$。
  二阶混合偏导数：
  $z_{xy} = \dfrac{1}{x^2y^2} \implies z_{xy}(1,2) = \dfrac{1}{4}$。

  \item \textbf{组装二阶泰勒多项式：}
  代入点 $(1,2)$ 处的展开式结构：
  \begin{align*}
  P_2(x,y) &= z(1,2) + z_x(x-1) + z_y(y-2) \\
  &\quad + \frac{1}{2}\left[ z_{xx}(x-1)^2 + 2z_{xy}(x-1)(y-2) + z_{yy}(y-2)^2 \right] \\
  &= \frac{1}{2} - \frac{1}{2}(x-1) - \frac{1}{4}(y-2) + \frac{1}{2}\left[ 1(x-1)^2 + 2\left(\frac{1}{4}\right)(x-1)(y-2) + \frac{1}{4}(y-2)^2 \right] \\
  &= \frac{1}{2} - \frac{1}{2}(x-1) - \frac{1}{4}(y-2) + \frac{1}{2}(x-1)^2 + \frac{1}{4}(x-1)(y-2) + \frac{1}{8}(y-2)^2.
  \end{align*}
\end{enumerate}
\end{solution}

\subsubsection{B组}

\begin{exercise}[一阶 Taylor 近似应用一]
试利用一阶泰勒多项式求 $\sqrt{(3.012)^2 + (3.997)^2}$ 的近似值.
\end{exercise}
\begin{solution}
\begin{enumerate}
  \item \textbf{构造函数与寻找基准点：}
  设二元函数 $f(x,y) = \sqrt{x^2+y^2}$。
  要计算的是点 $(3.012, 3.997)$ 处的函数值。
  我们选取最近的整数基准点 $(x_0, y_0) = (3, 4)$，此处的函数值易于计算：$f(3,4) = \sqrt{3^2+4^2} = 5$。
  增量为 $\Delta x = 3.012 - 3 = 0.012$，$\Delta y = 3.997 - 4 = -0.003$。

  \item \textbf{计算偏导数：}
  求偏导得 $f_x = \dfrac{x}{\sqrt{x^2+y^2}},\allowbreak\ f_y = \dfrac{y}{\sqrt{x^2+y^2}}$。
  代入基准点得 $f_x(3,4) = \dfrac{3}{5} = 0.6,\allowbreak\ f_y(3,4) = \dfrac{4}{5} = 0.8$。

  \item \textbf{运用一阶泰勒多项式求近似：}
  根据一阶全微分近似公式 $f(x,y) \approx f(x_0,y_0) + f_x \Delta x + f_y \Delta y$，得
  $f(3.012, 3.997) \approx 5 + 0.6(0.012) + 0.8(-0.003) = 5 + 0.0072 - 0.0024 = 5.0048$。
\end{enumerate}
\end{solution}

\begin{exercise}[一阶 Taylor 近似应用二]
电阻 $R_1$ 和 $R_2$ 并联后的电阻为 $R = \dfrac{R_1R_2}{R_1+R_2}$，利用一阶泰勒多项式计算当 $R_1 = 2.013\ \Omega$，$R_2 = 5.972\ \Omega$ 时，$R$ 的近似值.
\end{exercise}
\begin{solution}
\begin{enumerate}
  \item \textbf{构造函数与基准点：}
  并联电阻函数为 $R(R_1,\allowbreak R_2) = \dfrac{R_1R_2}{R_1+R_2}$。
  取整数基准点 $(R_{10},R_{20})=(2,6)$，此时 $R(2,6)=1.5\ \Omega$。
  自变量增量分别为 $\Delta R_1 = 0.013\ \Omega$ 与 $\Delta R_2 = -0.028\ \Omega$。

  \item \textbf{计算偏导数：}
  对 $R_1$ 求偏导得 \(\displaystyle \frac{\partial R}{\partial R_1}=\frac{R_2^2}{(R_1+R_2)^2}\)。
  在基准点处：$\left. \dfrac{\partial R}{\partial R_1} \right|_{(2,6)} = \dfrac{6^2}{(2+6)^2} = \dfrac{36}{64} = \dfrac{9}{16}$。
  同理，对 $R_2$ 求偏导：\(\displaystyle \frac{\partial R}{\partial R_2} = \frac{R_1^2}{(R_1+R_2)^2}\)。
  在基准点处：$\left. \dfrac{\partial R}{\partial R_2} \right|_{(2,6)} = \dfrac{2^2}{(2+6)^2} = \dfrac{4}{64} = \dfrac{1}{16}$。

  \item \textbf{利用一阶泰勒公式近似：}
  根据一阶泰勒展开式 $R \approx R(2,6) + \dfrac{\partial R}{\partial R_1}\Delta R_1 + \dfrac{\partial R}{\partial R_2}\Delta R_2$ 进行代入计算：\(\displaystyle R \approx 1.5 + \frac{9}{16}(0.013) + \frac{1}{16}(-0.028)=1.5+\frac{0.117-0.028}{16}=1.5055625\ \Omega\)。
  即所求并联电阻的近似值为 $1.5055625\ \Omega$。
\end{enumerate}
\end{solution}

\begin{exercise}[抽象函数的一阶 Taylor 多项式]
试写出函数 $u = f(x,y,z)$ 在点 $(a,b,c)$ 的一阶泰勒多项式.
\end{exercise}
\begin{solution}
\begin{enumerate}
  \item \textbf{明确一阶多项式的代数结构：}
  多元函数的一阶 Taylor 多项式正是函数在该点处的函数值叠加其全微分（线性主部）。
  \item \textbf{写出一般公式：}
  根据多元微分学定义，三元函数在点 $(a,b,c)$ 的一阶泰勒多项式为
  $P_1(x,y,z) = f(a,b,c) + f_x(a,b,c)(x-a) + f_y(a,b,c)(y-b) + f_z(a,b,c)(z-c)$。
\end{enumerate}
\end{solution}

\begin{exercise}[纸箱容积近似计算]
一长方体纸箱外侧尺寸分别为 $14\ \mathrm{cm},14\ \mathrm{cm},28\ \mathrm{cm}$，厚度为 $0.125\ \mathrm{cm}$，试利用习题3的结果计算纸箱容积的近似值.
\end{exercise}
\begin{solution}
\begin{enumerate}
  \item \textbf{分析题意与建立模型：}
  纸箱“容积”指的是内部空腔的体积。设纸箱的外侧长、宽、高分别为 $x, y, z$，内部容积为 $V$。
  由于纸箱是有厚度的闭合长方体，内部长、宽、高分别为 $x - 2\Delta x, y - 2\Delta y, z - 2\Delta z$，其中单侧厚度即缩减增量绝对值 $\Delta x = \Delta y = \Delta z = 0.125\ \mathrm{cm}$。
  建立容积函数 $V(x,y,z) = xyz$。

  \item \textbf{计算基准点的值与偏导数：}
  取外侧尺寸作为展开的基准点 $(x_0, y_0, z_0) = (14, 14, 28)$。
  在此处，基准容积 $V_0 = 14 \times 14 \times 28 = 5488\ \mathrm{cm}^3$。
  计算偏导数：\(\displaystyle V_x = yz \implies V_x(14,14,28) = 14 \times 28 = 392\)，
  \(\displaystyle V_y = xz \implies V_y(14,14,28) = 14 \times 28 = 392\)，
  \(\displaystyle V_z = xy \implies V_z(14,14,28) = 14 \times 14 = 196\)。

  \item \textbf{利用一阶泰勒近似求容积增量：}
  各维度的微小总增量（向内）为 $\mathrm{d}x = \mathrm{d}y = \mathrm{d}z = -2 \times 0.125 = -0.25\ \mathrm{cm}$。
  由前述习题 3 的一阶泰勒多项式（全微分近似），有 \(\displaystyle \Delta V \approx V_x \mathrm{d}x + V_y \mathrm{d}y + V_z \mathrm{d}z=392(-0.25)+392(-0.25)+196(-0.25)=-245\ \mathrm{cm}^3\)。

  \item \textbf{得出最终近似值：}
  内部近似容积为 $V \approx V_0 + \Delta V = 5488 - 245 = 5243\ \mathrm{cm}^3$。
\end{enumerate}
\end{solution}

% ============================================================
% §7.9 偏导数在几何上的应用
% ============================================================

\clearpage

\subsection{第九节习题：偏导数在几何上的应用}

\subsubsection{A组}

\begin{exercise}[曲线切向量的基本计算]
对下列曲线写出已知点处的切向量：
\begin{shortsingleenum}
    \shortsinglechoice{(1)}{$y=x,\,z=x^2,\,M_0(1,1,1)$；}
    \shortsinglechoice{(2)}{$x=a\cos\alpha\cos t,\,y=a\sin\alpha\cos t,\,z=a\sin t,\,t=t_0$；}
    \shortsinglechoice{(3)}{$x^2+z^2=10,\,y^2+z^2=10,\,M_0(1,1,3)$；}
    \shortsinglechoice{(4)}{$z=f(x,y),\,\dfrac{x-x_0}{\cos\alpha}=\dfrac{y-y_0}{\sin\alpha},\,f$ 可微，在曲线上的点 $M_0(x_0,y_0,z_0)$ 处。}
\end{shortsingleenum}
\end{exercise}
\begin{solution}
\begin{enumerate}
    \item[(1)] 取 $x$ 为参数，曲线参数化为 $\vec{r}(x) = (x, x, x^2)$。
    切向量 $\vec{r}'(x) = (1, 1, 2x)$。在 $M_0(1,1,1)$ 处 $x=1$，代入得切向量 $\vec{T} = (1, 1, 2)$。

    \item[(2)] 直接对参数求导，得 $x'(t)=-a\cos\alpha\sin t,\ y'(t)=-a\sin\alpha\sin t,\ z'(t)=a\cos t$，故在 $t=t_0$ 时切向量为 $\vec{T} = (-a\cos\alpha\sin t_0,\, -a\sin\alpha\sin t_0,\, a\cos t_0)$。

    \item[(3)] 该曲线是两个隐式圆柱面的交线。令 $F_1 = x^2+z^2-10 = 0$ 和 $F_2 = y^2+z^2-10 = 0$，则 $\nabla F_1 = (2x, 0, 2z)$，在 $(1,1,3)$ 处为 $(2, 0, 6) \sim (1, 0, 3)$；$\nabla F_2 = (0, 2y, 2z)$，在 $(1,1,3)$ 处为 $(0, 2, 6) \sim (0, 1, 3)$。交线切向量与两法向量均垂直，故可取 $\vec{T} = (1, 0, 3) \times (0, 1, 3) = (-3, -3, 1)$。

    \item[(4)] 引入参数 $t$，直线方程可参数化为 $x = x_0 + t\cos\alpha$，$y = y_0 + t\sin\alpha$。
    代入曲面方程得曲线 $z(t) = f(x_0+t\cos\alpha, y_0+t\sin\alpha)$。点 $M_0$ 对应 $t=0$。
    对 $t$ 求导得切向量 $\vec{r}'(t) = (\cos\alpha,\, \sin\alpha,\, f_x\cos\alpha + f_y\sin\alpha)$。
    在 $M_0$ 处，切向量为 $\vec{T} = (\cos\alpha,\, \sin\alpha,\, f_x(x_0,y_0)\cos\alpha + f_y(x_0,y_0)\sin\alpha)$。
\end{enumerate}
\end{solution}

\begin{exercise}[切线与法平面的几何计算]
求下列曲线在指定点的切线和法平面方程：
\begin{shortexenum}
    \shortitem{1}{$x=t,\,y=t^2,\,z=t^3,\,M(1,1,1)$；} &
    \shortitem{2}{$x=t-\sin t,\,y=1-\cos t,\,z=4\sin\dfrac{t}{2},\,M\left(\dfrac{\pi}{2}-1,1,2\sqrt{2}\right)$；} \\
    \shortitem{3}{$x=a\sin^2 t,\,y=b\sin t\cos t,\,z=c\cos^2 t,\,t=\dfrac{\pi}{4}$；} &
    \shortitem{4}{$x=a\cos t,\,y=a\sin t,\,z=ct,\,M(a,0,0)$。}
\end{shortexenum}
\end{exercise}
\begin{solution}
\begin{enumerate}
    \item[(1)] \textbf{确定切向量：} 参数 $t=1$。求导得 $\vec{r}'(t) = (1, 2t, 3t^2)$，代入 $t=1$ 得切向量 $\vec{T} = (1, 2, 3)$，故切线方程为 $\dfrac{x-1}{1} = \dfrac{y-1}{2} = \dfrac{z-1}{3}$，法平面方程为 $1(x-1) + 2(y-1) + 3(z-1) = 0 \implies x+2y+3z-6=0$。

    \item[(2)] \textbf{确定参数与切向量：} 点 $M$ 对应 $t = \pi/2$。对参数求导得 \(\displaystyle x'(t) = 1-\cos t,\quad y'(t) = \sin t,\quad z'(t) = 2\cos(t/2)\)。故 $x'(\pi/2)=1,\ y'(\pi/2)=1,\ z'(\pi/2)=2\left(\dfrac{\sqrt{2}}{2}\right)=\sqrt{2}$，从而切向量 $\vec{T} = (1, 1, \sqrt{2})$。于是切线方程为 \(\displaystyle \dfrac{x-(\pi/2-1)}{1} = \dfrac{y-1}{1} = \dfrac{z-2\sqrt{2}}{\sqrt{2}}\)，
    法平面方程为
    \[
    1(x-\pi/2+1) + 1(y-1) + \sqrt{2}(z-2\sqrt{2}) = 0
    \implies x+y+\sqrt{2}z - \dfrac{\pi}{2} - 4 = 0.
    \]

    \item[(3)] \textbf{计算坐标和切向量：} $t=\pi/4$。此时 $x(\pi/4) = a/2,\ y(\pi/4) = b/2,\ z(\pi/4) = c/2$。又有 \(\displaystyle x'(t)=a\sin 2t,\qquad y'(t)=b\cos 2t,\qquad z'(t)=-c\sin 2t\)。
    故 $x'(\pi/4) = a,\ y'(\pi/4) = 0,\ z'(\pi/4) = -c$，于是切向量 $\vec{T} = (a, 0, -c)$。从而切线方程为 \(\displaystyle \dfrac{x-a/2}{a} = \dfrac{y-b/2}{0} = \dfrac{z-c/2}{-c}\)，
    法平面方程为 \(\displaystyle a(x-a/2) - c(z-c/2) = 0\implies ax - cz - \dfrac{a^2-c^2}{2} = 0\)。

    \item[(4)] \textbf{确定参数与切向量：} 点 $M(a,0,0)$ 对应 $t=0$。求导得 $\vec{r}'(t) = (-a\sin t,\, a\cos t,\, c)$，代入 $t=0$ 得 $\vec{T} = (0, a, c)$，故切线方程为 $\dfrac{x-a}{0} = \dfrac{y}{a} = \dfrac{z}{c}$，法平面方程为 $0(x-a) + a(y-0) + c(z-0) = 0 \implies ay+cz=0$。
\end{enumerate}
\end{solution}

\begin{exercise}[寻找满足特定切线条件的点]
在曲线 $y=x^2,\,z=x^3$ 上求一点，使该点处的切线平行于平面 $x+2y+z=4$。
\end{exercise}
\begin{solution}
\begin{enumerate}
    \item \textbf{参数化曲线并求切向量：}
    取 $x$ 为参数，得到曲线向量 $\vec{r}(x) = (x, x^2, x^3)$。
    对其求导得到切向量表达式 $\vec{T}(x) = (1, 2x, 3x^2)$。

    \item \textbf{运用平行条件建立方程：}
    给定平面 $x+2y+z=4$ 的法向量为 $\vec{n} = (1, 2, 1)$。要使切线平行于该平面，只需曲线切向量垂直于该法向量，即 $\vec{T} \cdot \vec{n} = 0$，从而 $1(1) + 2x(2) + 3x^2(1) = 0$，即 $3x^2 + 4x + 1 = 0$。

    \item \textbf{解方程得坐标点：}
    解二次方程得 $(3x+1)(x+1) = 0 \implies x = -1$ 或 $x = -1/3$。当 $x = -1$ 时，$y=1,\ z=-1$，点为 $(-1, 1, -1)$；当 $x = -1/3$ 时，$y=1/9,\ z=-1/27$，点为 $\left(-\dfrac{1}{3}, \dfrac{1}{9}, -\dfrac{1}{27}\right)$。此两点均满足条件。
\end{enumerate}
\end{solution}

\begin{exercise}[隐式交线切线一]
求曲线 $\begin{cases} x^2+y^2+z^2-3x=0, \\ 2x-3y+5z-4=0 \end{cases}$ 在点 $(1,1,1)$ 处的切线和法平面方程。
\end{exercise}
\begin{solution}
\begin{enumerate}
    \item \textbf{计算隐式曲面的梯度：}
    令 $F_1 = x^2+y^2+z^2-3x = 0$ 和 $F_2 = 2x-3y+5z-4 = 0$。则 $\nabla F_1 = (2x-3, 2y, 2z)$，在点 $(1,1,1)$ 处为 $(-1, 2, 2)$；$\nabla F_2 = (2, -3, 5)$，由于它是平面，法向量处处相同。

    \item \textbf{通过叉积求切向量：}
    曲线切向量与两个法向量垂直，因此可取 $\vec{T} = \nabla F_1 \times \nabla F_2 = (16, 9, -1)$。

    \item \textbf{写出结果方程：}
    \textbf{切线方程}为 $\dfrac{x-1}{16} = \dfrac{y-1}{9} = \dfrac{z-1}{-1}$，\textbf{法平面方程}为 $16(x-1) + 9(y-1) - 1(z-1) = 0 \implies 16x+9y-z-24=0$。
\end{enumerate}
\end{solution}

\begin{exercise}[隐式交线切线二]
求曲线 $\begin{cases} x^2+y^2+z^2=6, \\ x+y+z=0 \end{cases}$ 在点 $(1,-2,1)$ 处的切线和法平面方程。
\end{exercise}
\begin{solution}
\begin{enumerate}
    \item \textbf{计算隐式曲面的梯度：}
    令 $F_1 = x^2+y^2+z^2-6 = 0$，$F_2 = x+y+z = 0$。则 $\nabla F_1 = (2x, 2y, 2z) \sim (x,y,z)$，在 $(1,-2,1)$ 处可取法向量 $(1, -2, 1)$；$\nabla F_2 = (1, 1, 1)$。

    \item \textbf{通过叉积求切向量：}
    $\vec{T} = (1, -2, 1) \times (1, 1, 1) = (-3, 0, 3)$，可提取公因式 $-3$，取等效方向向量 $\vec{T} = (-1, 0, 1)$。

    \item \textbf{写出结果方程：}
    \textbf{切线方程}为 $\dfrac{x-1}{-1} = \dfrac{y+2}{0} = \dfrac{z-1}{1}$，\textbf{法平面方程}为 $-(x-1) + 0(y+2) + 1(z-1) = 0 \implies x-z=0$。
\end{enumerate}
\end{solution}

\begin{exercise}[曲面法向量的求解]
对下列曲面写出在已知点的法向量：
\begin{shortexenum}
    \shortitem{1}{$z=x^2+y^2,\,M_0(1,2,5)$；} &
    \shortitem{2}{$z=\arctan\dfrac{y}{x},\,M_0\left(1,1,\dfrac{\pi}{4}\right)$；} \\
    \shortitem{3}{$2^{\frac{y}{x}}+2^{\frac{y}{z}}=8,\,M_0(2,2,1)$；} &
    \shortitem{4}{$F\left(\dfrac{x-a}{z-c},\dfrac{y-b}{z-c}\right)=0,\,M_0(x_0,y_0,z_0)$。}
\end{shortexenum}
\end{exercise}
\begin{solution}
\begin{enumerate}
    \item[(1)] 转化为隐式 $F = x^2+y^2-z$。梯度 $\nabla F = (2x, 2y, -1)$。代入 $(1,2,5)$，得到法向量 $\vec{n} = (2, 4, -1)$。

    \item[(2)] 转化为隐式 $F = \arctan(y/x) - z$，则 $F_x = \dfrac{-y}{x^2+y^2},\ F_y = \dfrac{x}{x^2+y^2},\ F_z = -1$。代入点得法向量为 $\vec{n} = \left(-\dfrac{1}{2}, \dfrac{1}{2}, -1\right) \sim (1, -1, 2)$。

    \item[(3)] 设 $F = 2^{y/x}+2^{y/z}-8 = 0$。由链式法则
    \begin{align*}
        F_x &= 2^{y/x}\ln 2 \left(-\dfrac{y}{x^2}\right),\\
        F_y &= 2^{y/x}\ln 2 \left(\dfrac{1}{x}\right) + 2^{y/z}\ln 2 \left(\dfrac{1}{z}\right),\\
        F_z &= 2^{y/z}\ln 2 \left(-\dfrac{y}{z^2}\right).
    \end{align*}
    在 $(2,2,1)$ 处分别得 $F_x = -\ln 2,\ F_y = 5\ln 2,\ F_z = -8\ln 2$，故法向量为 $(-\ln 2, 5\ln 2, -8\ln 2) \sim (1, -5, 8)$。

    \item[(4)] 令中间变量 $u = \dfrac{x-a}{z-c},\ v = \dfrac{y-b}{z-c}$。利用链式法则得 $n_x = F_1 \cdot \dfrac{1}{z-c},\ n_y = F_2 \cdot \dfrac{1}{z-c},\ n_z = F_1 \left(-\dfrac{x-a}{(z-c)^2}\right) + F_2 \left(-\dfrac{y-b}{(z-c)^2}\right)$。将法向量各项同时乘以 $(z-c)$，可得到更简单的法向量方向比例式 \(\displaystyle \vec n\sim\left(F_1,F_2,-F_1\dfrac{x-a}{z-c}-F_2\dfrac{y-b}{z-c}\right)\)。
\end{enumerate}
\end{solution}

\begin{exercise}[综合求切平面和法线方程]
求下列曲面在指定点的切平面和法线方程：
\par\noindent
\begin{tabular}{@{}p{0.4\textwidth}p{0.4\textwidth}@{}}
(1) $z=x^2+y^2-1$ 于 $P_0(2,1,4)$ 处 & (2) $z=\mathrm{e}^y+x+x^2+6$ 于 $P_0(1,0,9)$ 处 \\
(3) $z=y\mathrm{e}^{x/y}$ 于 $P_0(1,1,\mathrm{e})$ 处 & (4) $z=\dfrac{1}{2}(x^2+4y^2)$ 于 $P_0(2,1,4)$ 处 \\
(5) $x^2+y^2+z^2=169$ 于 $P_0(3,4,12)$ 处 & (6) $3x^2+2y^2=2z+1$ 于 $P_0(1,1,2)$ 处 \\
(7) $z=y+\ln\dfrac{x}{z}$ 于 $P_0(1,1,1)$ 处 & (8) $\mathrm{e}^z-z+xy=3$ 于 $P_0(2,1,0)$ 处
\end{tabular}
\end{exercise}
\begin{solution}
\begin{enumerate}
    \item[(1)] 构造隐函数 $F = x^2+y^2-1-z$。$\nabla F = (2x, 2y, -1) = (4, 2, -1)$，故切平面方程为 $4(x-2)+2(y-1)-(z-4)=0 \implies 4x+2y-z-6=0$，法线方程为 $\dfrac{x-2}{4}=\dfrac{y-1}{2}=\dfrac{z-4}{-1}$。

    \item[(2)] 构造隐函数 $F = \mathrm{e}^y+x+x^2+6-z$。$\nabla F = (1+2x, \mathrm{e}^y, -1) = (3, 1, -1)$，故切平面方程为 $3(x-1)+1(y-0)-(z-9)=0 \implies 3x+y-z+6=0$，法线方程为 $\dfrac{x-1}{3}=\dfrac{y}{1}=\dfrac{z-9}{-1}$。

    \item[(3)] 构造隐函数 $F = y\mathrm{e}^{x/y}-z$，则 $F_x = \mathrm{e}^{x/y},\ F_y = \mathrm{e}^{x/y} + y\mathrm{e}^{x/y}(-x/y^2) = \mathrm{e}^{x/y}(1-x/y),\ F_z = -1$，在 $P_0(1,1,\mathrm{e})$ 处得 $\nabla F = (\mathrm{e}, 0, -1)$，故切平面方程为 $\mathrm{e}(x-1)+0(y-1)-1(z-\mathrm{e})=0 \implies \mathrm{e}x-z=0$，法线方程为 $\dfrac{x-1}{\mathrm{e}}=\dfrac{y-1}{0}=\dfrac{z-\mathrm{e}}{-1}$。

    \item[(4)] 构造隐函数 $F = x^2+4y^2-2z$。$\nabla F = (2x, 8y, -2) \sim (x, 4y, -1) = (2, 4, -1)$，故切平面方程为 $2(x-2)+4(y-1)-1(z-4)=0 \implies 2x+4y-z-4=0$，法线方程为 $\dfrac{x-2}{2}=\dfrac{y-1}{4}=\dfrac{z-4}{-1}$。

    \item[(5)] 构造隐函数 $F = x^2+y^2+z^2-169$。$\nabla F = (2x, 2y, 2z) \sim (x,y,z) = (3, 4, 12)$，故切平面方程为 $3(x-3)+4(y-4)+12(z-12)=0 \implies 3x+4y+12z-169=0$，法线方程为 $\dfrac{x-3}{3}=\dfrac{y-4}{4}=\dfrac{z-12}{12}$。

    \item[(6)] 构造隐函数 $F = 3x^2+2y^2-2z-1$。$\nabla F = (6x, 4y, -2) \sim (3x, 2y, -1) = (3, 2, -1)$，故切平面方程为 $3(x-1)+2(y-1)-1(z-2)=0 \implies 3x+2y-z-3=0$，法线方程为 $\dfrac{x-1}{3}=\dfrac{y-1}{2}=\dfrac{z-2}{-1}$。

    \item[(7)] 构造隐函数 $F = y+\ln x - \ln z - z$，则 $F_x = 1/x = 1,\ F_y = 1,\ F_z = -1/z - 1 = -2$，从而 $\nabla F = (1, 1, -2)$，故切平面方程为 $1(x-1)+1(y-1)-2(z-1)=0 \implies x+y-2z=0$，法线方程为 $\dfrac{x-1}{1}=\dfrac{y-1}{1}=\dfrac{z-1}{-2}$。

    \item[(8)] 构造隐函数 $F = \mathrm{e}^z-z+xy-3$，则 $F_x = y = 1,\ F_y = x = 2,\ F_z = \mathrm{e}^z-1 = 0$，从而 $\nabla F = (1, 2, 0)$，故切平面方程为 $1(x-2)+2(y-1)+0(z-0)=0 \implies x+2y-4=0$，法线方程为 $\dfrac{x-2}{1}=\dfrac{y-1}{2}=\dfrac{z}{0}$。
\end{enumerate}
\end{solution}

\begin{exercise}[法向量平行应用]
求曲面 $2x^2+3y^2+z^2=9$ 的切平面，使之平行于平面 $2x-3y+2z=1$。
\end{exercise}
\begin{solution}
\begin{enumerate}
    \item \textbf{提取相关法向量：}
    隐式曲面 $F = 2x^2+3y^2+z^2-9=0$。其任意点处法向量为 $\nabla F = (4x, 6y, 2z)$。
    已知平面的法向量为 $\vec{n} = (2, -3, 2)$。

    \item \textbf{建立平行条件方程：}
    由于切平面平行于已知平面，两法向量必须共线，即 $(4x, 6y, 2z) = k(2, -3, 2)$，从而 $x=\dfrac{k}{2},\ y=-\dfrac{k}{2},\ z=k$。

    \item \textbf{代回曲面方程求坐标点：}
    将 $x, y, z$ 的表达式代入曲面方程，得 $2\left(\dfrac{k^2}{4}\right) + 3\left(\dfrac{k^2}{4}\right) + k^2 = 9$，即 $\dfrac{9}{4}k^2 = 9$，故 $k=\pm 2$。于是当 $k=2$ 时切点为 $(1, -1, 2)$，当 $k=-2$ 时切点为 $(-1, 1, -2)$。

    \item \textbf{写出切平面方程：}
    点 $(1, -1, 2)$ 处的切平面方程为 $2(x-1)-3(y+1)+2(z-2)=0 \implies 2x-3y+2z-9=0$；点 $(-1, 1, -2)$ 处的切平面方程为 $2(x+1)-3(y-1)+2(z+2)=0 \implies 2x-3y+2z+9=0$。
\end{enumerate}
\end{solution}

\subsubsection{B组}

\begin{exercise}[曲线与曲面交角的证明]
证明曲线 $x=a\mathrm{e}^t\cos t,\,y=a\mathrm{e}^t\sin t,\,z=a\mathrm{e}^t$（其中 $a$ 为常数）与锥面 $z^2=x^2+y^2$ 的母线相交的角的大小是相同的。
\end{exercise}
\begin{solution}
\begin{enumerate}
    \item \textbf{验证曲线位于锥面上：}
    将曲线坐标代入得 $x^2+y^2 = a^2\mathrm{e}^{2t}(\cos^2 t + \sin^2 t) = a^2\mathrm{e}^{2t}$，而 $z^2 = (a\mathrm{e}^t)^2 = a^2\mathrm{e}^{2t}$，两者相等，说明该曲线确实位于锥面上。

    \item \textbf{写出曲线的切向量与母线的方向向量：}
    对曲线的参数方程求导得到切向量
    \[
    \vec{T} = \vec{r}'(t) = a\mathrm{e}^t(\cos t-\sin t,\, \sin t+\cos t,\, 1).
    \]
    过曲线上一点 $(a\mathrm{e}^t\cos t,\, a\mathrm{e}^t\sin t,\, a\mathrm{e}^t)$ 的锥面母线，是连接原点与该点的直线，其方向向量 $\vec{g}$ 可取为位置向量本身，即 \(\displaystyle \vec{g} = a\mathrm{e}^t(\cos t,\, \sin t,\, 1)\)。

    \item \textbf{计算两个向量的夹角余弦值：}
    分别计算得
    \begin{align*}
    |\vec{T}|&=a\mathrm{e}^t\sqrt{(\cos t-\sin t)^2+(\sin t+\cos t)^2+1}=a\mathrm{e}^t\sqrt{3},\\
    |\vec{g}|&=a\mathrm{e}^t\sqrt{\cos^2 t+\sin^2 t+1}=a\mathrm{e}^t\sqrt{2}.
    \end{align*}
    且
    \[
    \vec{T} \cdot \vec{g}
    = a^2\mathrm{e}^{2t} \left[ \cos t(\cos t-\sin t) + \sin t(\sin t+\cos t) + 1\cdot 1 \right]
    = 2a^2\mathrm{e}^{2t}.
    \]
    设相交角为 $\alpha$，则
    \[
    \cos\alpha
    = \frac{\vec{T} \cdot \vec{g}}{|\vec{T}| |\vec{g}|}
    = \frac{2a^2\mathrm{e}^{2t}}{\sqrt{3}a\mathrm{e}^t \cdot \sqrt{2}a\mathrm{e}^t}
    = \frac{2}{\sqrt{6}}
    = \frac{\sqrt{6}}{3}.
    \]
    因为余弦值是一个与 $t$ 无关的常数，所以曲线与该锥面各处母线的相交角大小均相同。
\end{enumerate}
\end{solution}

\begin{exercise}[求两曲面交角]
求曲面 $3x^2+y^2+z^2=16$ 与 $x^2+y^2+z^2=14$ 在点 $(-1,2,3)$ 的交角 $\theta$。
\end{exercise}
\begin{solution}
\begin{enumerate}
    \item \textbf{定义交角及计算各曲面的法向量：}
    曲面的交角即为在该相交点处二者法向量（即切平面法线）之间的夹角。设 $F_1 = 3x^2+y^2+z^2-16 = 0$，$F_2 = x^2+y^2+z^2-14 = 0$，则 $\nabla F_1 = (6x, 2y, 2z)$，在 $(-1,2,3)$ 处可取 $\vec{n}_1 = (-6, 4, 6) \sim (-3, 2, 3)$；$\nabla F_2 = (2x, 2y, 2z)$，在 $(-1,2,3)$ 处可取 $\vec{n}_2 = (-2, 4, 6) \sim (-1, 2, 3)$。

    \item \textbf{计算夹角的余弦值：}
    \begin{align*}
    \cos\theta &= \frac{|\vec{n}_1 \cdot \vec{n}_2|}{|\vec{n}_1| |\vec{n}_2|} \\
    &= \frac{|(-3)(-1) + (2)(2) + (3)(3)|}{\sqrt{(-3)^2+2^2+3^2} \sqrt{(-1)^2+2^2+3^2}} \\
    &= \frac{3+4+9}{\sqrt{22} \sqrt{14}} = \frac{16}{\sqrt{308}} = \frac{16}{2\sqrt{77}} = \frac{8\sqrt{77}}{77}.
    \end{align*}
    交角为 $\theta = \arccos\left(\dfrac{8\sqrt{77}}{77}\right)$。
\end{enumerate}
\end{solution}

\begin{exercise}[曲面切平面与两平面的垂直关系]
求曲面 $x^2+y^2+z^2=x$ 的切平面，使它垂直于平面 $x-y-z=2$ 和 $x-y-\dfrac{1}{2}z=2$。
\end{exercise}
\begin{solution}
\begin{enumerate}
    \item \textbf{由两平面的法向量求所需切平面的法向量：}
    已知两个平面的法向量分别为 $\vec{n}_1 = (1, -1, -1)$ 和 $\vec{n}_2 = \left(1, -1, -\dfrac{1}{2}\right) \sim (2, -2, -1)$。切平面垂直于这两个平面，这意味着切平面的法向量 $\vec{n}$ 必须同时垂直于 $\vec{n}_1$ 和 $\vec{n}_2$，故可取 $\vec{n} = \vec{n}_1 \times \vec{n}_2 = (-1, -1, 0) \sim (1, 1, 0)$。

    \item \textbf{根据法向量求出曲面上的切点：}
    曲面隐式方程 $F = x^2-x+y^2+z^2 = 0$。其法向量为 $\nabla F = (2x-1, 2y, 2z)$。
    该法向量应与 $\vec{n}$ 共线，即 $(2x-1, 2y, 2z) = k(1, 1, 0)$，从而 $2x-1 = k,\ 2y = k,\ 2z = 0$，即 $z=0,\ y=\dfrac{k}{2},\ x=\dfrac{k+1}{2}$。代入曲面方程得 $\left(\dfrac{k+1}{2}\right)^2 - \dfrac{k+1}{2} + \left(\dfrac{k}{2}\right)^2 = 0$，即 $2k^2 - 1 = 0$，故 $k = \pm\dfrac{\sqrt{2}}{2}$。

    \item \textbf{写出切平面方程：}
    点坐标为 $\left( \dfrac{2\pm\sqrt{2}}{4},\, \pm\dfrac{\sqrt{2}}{4},\, 0 \right)$，法向量为 $(1, 1, 0)$，故
    \(\displaystyle
        1\cdot\left(x - \dfrac{2\pm\sqrt{2}}{4}\right) + 1\cdot\left(y \mp \dfrac{\sqrt{2}}{4}\right) = 0
    \)，即
    \(\displaystyle
        x+y - \dfrac{2\pm\sqrt{2}\pm\sqrt{2}}{4} = 0
    \)。
    因为 $x$ 取正时 $y$ 亦取正，所以两处的 $\pm$ 同号，最终得到两个平行的切平面：$x+y - \dfrac{1+\sqrt{2}}{2} = 0$ 与 $x+y - \dfrac{1-\sqrt{2}}{2} = 0$。
\end{enumerate}
\end{solution}

\begin{exercise}[抽象曲面切平面的几何特性推导]
设 $F(u,v)$ 是可微函数，求证：曲面 $F(ax-bz,ay-cz)=0\,(abc\neq 0)$ 的切平面平行于某定直线。
\end{exercise}
\begin{solution}
\begin{enumerate}
    \item \textbf{计算曲面法向量：}
    设曲面方程为 $G(x,y,z) = F(ax-bz, ay-cz) = 0$。令 $u = ax-bz,\ v = ay-cz$。由复合函数求导法得
    \begin{align*}
    G_x &= F_u \cdot \dfrac{\partial u}{\partial x}
    + F_v \cdot \dfrac{\partial v}{\partial x} = a F_u,\\
    G_y &= F_u \cdot \dfrac{\partial u}{\partial y}
    + F_v \cdot \dfrac{\partial v}{\partial y} = a F_v,\\
    G_z &= F_u \cdot \dfrac{\partial u}{\partial z}
    + F_v \cdot \dfrac{\partial v}{\partial z} = -b F_u - c F_v.
    \end{align*}
    所以该曲面上任意一点处的法向量为 \(\displaystyle \vec{n} = (a F_u,\, a F_v,\, -b F_u - c F_v)\)。

    \item \textbf{寻找定直线：}
    若切平面始终平行于某定直线，等价于法向量始终垂直于该定直线的方向向量 $\vec{l} = (l_1, l_2, l_3)$，即 $\vec{n} \cdot \vec{l} \equiv 0$。代入并整理得
    \begin{align*}
    \vec{n} \cdot \vec{l}
    &= a F_u l_1 + a F_v l_2 - b F_u l_3 - c F_v l_3\\
    &= F_u(a l_1 - b l_3) + F_v(a l_2 - c l_3)=0.
    \end{align*}
    要使该式对于任意可能的偏导数值 $F_u, F_v$ 均恒成立，必须有 $a l_1 - b l_3 = 0$ 和 $a l_2 - c l_3 = 0$，即 \(\displaystyle l_1 = \dfrac{b}{a} l_3,\qquad l_2 = \dfrac{c}{a} l_3\)。

    \item \textbf{得出结论：}
    不妨取 $l_3 = a$（由于 $abc \neq 0$，则 $a \neq 0$），那么 $l_1 = b,\, l_2 = c$。
    故总是可以找到一个固定的方向向量 $\vec{l} = (b, c, a)$，使得该曲面无论处于哪一点，其切平面的法向量都与其垂直，即切平面平行于方向为 $(b, c, a)$ 的定直线。证毕。
\end{enumerate}
\end{solution}

% ============================================================
% §7.10 无约束的最优化问题
% ============================================================

\subsection{第十节习题：无约束的最优化问题}

\begin{exercise}[基础：求驻点与极值判定]
\begin{shortexenum}
  \shortitem{1}{求 $f(x,y)=x^3+y^3-3xy$ 的所有驻点并判定类型；} &
  \shortitem{2}{求 $f(x,y)=x^3-3xy+y^3$ 的所有驻点并判定类型；该函数有几个鞍点？} \\
  \shortitem{3}{讨论 $f(x,y)=xy\mathrm{e}^{-(x^2+y^2)}$ 在 $(0,0)$ 处是否取极值？} &
\end{shortexenum}
\end{exercise}
\begin{solution}
\begin{enumerate}
  \item[(1)] 一阶偏导数：$f_x = 3x^2-3y = 0 \implies y=x^2$。$f_y = 3y^2-3x = 0 \implies x=y^2$。
  代入得 $x = x^4 \implies x(x^3-1)=0 \implies x=0$ 或 $x=1$。对应得到驻点 $(0,0)$ 和 $(1,1)$。
  二阶偏导数：$f_{xx} = 6x$，$f_{xy} = -3$，$f_{yy} = 6y$。
  在 $(0,0)$ 处，$A=0,\ B=-3,\ C=0$，故 $\Delta = AC - B^2 = -9 < 0$，此点为\textbf{鞍点}；在 $(1,1)$ 处，$A=6,\ B=-3,\ C=6$，故 $\Delta = 27 > 0$ 且 $A > 0$，此点为\textbf{极小值点}。

  \item[(2)] 此函数与 (1) 完全相同，只是交换了中间项的顺序。
  因此，驻点及其类型完全一致。该函数有 \textbf{1 个鞍点}（位于原点）。

  \item[(3)]  一阶偏导数：
  $f_x = y\mathrm{e}^{-(x^2+y^2)} + xy\mathrm{e}^{-(x^2+y^2)}(-2x) = y\mathrm{e}^{-(x^2+y^2)}(1-2x^2)$，$f_y = x\mathrm{e}^{-(x^2+y^2)} + xy\mathrm{e}^{-(x^2+y^2)}(-2y) = x\mathrm{e}^{-(x^2+y^2)}(1-2y^2)$。
  在 $(0,0)$ 处，$f_x(0,0)=0$ 且 $f_y(0,0)=0$，故 $(0,0)$ 是驻点。
  计算二阶偏导数在 $(0,0)$ 处的值：
  由于 $f_{xx}$ 和 $f_{yy}$ 的各项中均包含因子 $y$ 和 $x$，代入 $(0,0)$ 后显然有 $A = f_{xx}(0,0) = 0$，$C = f_{yy}(0,0) = 0$。
  而 $f_{xy} = \mathrm{e}^{-(x^2+y^2)}(1-2x^2) + y\mathrm{e}^{-(x^2+y^2)}(1-2x^2)(-2y) = \mathrm{e}^{-(x^2+y^2)}(1-2x^2)(1-2y^2)$。
  代入 $(0,0)$，有 $B = f_{xy}(0,0) = 1$。
  $\Delta = AC - B^2 = -1 < 0$。
  所以 $(0,0)$ 为\textbf{鞍点}，\textbf{不取极值}。
\end{enumerate}
\end{solution}

\begin{exercise}[基础：Hesse 判别法熟练运用]
\begin{shortexenum}
  \shortitem{1}{求 $f(x,y)=x^3+y^3-3x^2-3y^2+9x$ 的极值；} &
  \shortitem{2}{证明函数 $f(x,y)=(x^2+y^2)\mathrm{e}^{x^2-y^2}$ 在 $(0,0)$ 处取极小值，并求出该极小值；} \\
  \shortitem{3}{设 $f(x,y)=x^4+y^4-2x^2-2\sqrt2\,xy+2y^2$，求 $f$ 的所有驻点并分类。} &
\end{shortexenum}
\end{exercise}
\begin{solution}
\begin{enumerate}
  \item[(1)]
  计算一阶偏导数：
  $f_x = 3x^2-6x+9 = 3(x^2-2x+3) = 3((x-1)^2+2) > 0$ 对于所有实数 $x$ 恒成立。
  由于 $f_x$ 永不为 0，该函数在定义域内\textbf{没有驻点}，因此也\textbf{没有极值}。

  \item[(2)]
  由于对于任何 $(x,y) \neq (0,0)$，有 $x^2+y^2 > 0$ 且 $\mathrm{e}^{x^2-y^2} > 0$，所以 $f(x,y) > 0$。
  而在 $(0,0)$ 处，$f(0,0) = 0$。显然 $f(x,y) \geqslant f(0,0)$ 恒成立，故 $(0,0)$ 是全局严格极小值点，极小值为 \textbf{0}。
  利用 Hesse 矩阵验证：
  $f_x = 2x\mathrm{e}^{x^2-y^2} + (x^2+y^2)\mathrm{e}^{x^2-y^2}(2x) = 2x\mathrm{e}^{x^2-y^2}(1+x^2+y^2)$，$f_y = 2y\mathrm{e}^{x^2-y^2} + (x^2+y^2)\mathrm{e}^{x^2-y^2}(-2y) = 2y\mathrm{e}^{x^2-y^2}(1-x^2-y^2)$。
  显然在 $(0,0)$ 处 $f_x=f_y=0$。求二阶偏导数：
  $f_{xx}(0,0) = \left. \left[ 2\mathrm{e}^{x^2-y^2}(1+x^2+y^2) + \dots \right] \right|_{(0,0)} = 2$，$f_{yy}(0,0) = \left. \left[ 2\mathrm{e}^{x^2-y^2}(1-x^2-y^2) + \dots \right] \right|_{(0,0)} = 2$，且 $f_{xy}(0,0) = 0$。
  $\Delta = 4 > 0$ 且 $A=2 > 0$。验证了 $(0,0)$ 为\textbf{极小值点}。

  \item[(3)]
  一阶偏导数：
  $f_x = 4x^3 - 4x - 2\sqrt{2}y = 0 \implies 2x^3 - 2x - \sqrt{2}y = 0$，$f_y = 4y^3 + 4y - 2\sqrt{2}x = 0 \implies 2y^3 + 2y - \sqrt{2}x = 0$。
  将第一式乘 $x$，第二式乘 $y$ 后相减：$2x^4 - 2y^4 - 2x^2 - 2y^2 = 0 \implies (x^2+y^2)(x^2-y^2-1) = 0$。
  由于 $x^2+y^2 \geqslant 0$，解得 $(x,y)=(0,0)$ 或 $x^2-y^2=1$。
  将 $x^2 = y^2+1$ 代入 $f_y=0$ 中有 $2y(y^2+1) = \sqrt{2}x \implies 2yx^2 = \sqrt{2}x$。
  由于 $x^2 = y^2+1 \geqslant 1$，知 $x \neq 0$，故 $xy = \dfrac{\sqrt{2}}{2}$。
  联立 $x^2-y^2=1$ 与 $x^2y^2=\dfrac{1}{2}$，将 $y^2 = x^2-1$ 代入：$x^2(x^2-1) = \dfrac{1}{2} \implies 2x^4 - 2x^2 - 1 = 0 \implies x^2 = \dfrac{1+\sqrt{3}}{2}$。
  则 $y^2 = \dfrac{\sqrt{3}-1}{2}$。由于 $xy > 0$，$x,y$ 同号。
  得驻点 $P_0(0,0)$，$P_1\left(\sqrt{\dfrac{\sqrt{3}+1}{2}}, \sqrt{\dfrac{\sqrt{3}-1}{2}}\right)$，$P_2\left(-\sqrt{\dfrac{\sqrt{3}+1}{2}}, -\sqrt{\dfrac{\sqrt{3}-1}{2}}\right)$。
  二阶偏导数：$f_{xx} = 12x^2-4$，$f_{yy} = 12y^2+4$，$f_{xy} = -2\sqrt{2}$。
  在 $P_0(0,0)$ 处，$A=-4,\ C=4,\ B=-2\sqrt{2}$，故 $\Delta = -24 < 0$，所以是\textbf{鞍点}；在 $P_1, P_2$ 处，$A = 6\sqrt{3}+2,\ C = 6\sqrt{3}-2,\ B = -2\sqrt{2}$，故 $\Delta = (6\sqrt{3}+2)(6\sqrt{3}-2) - 8 = 96 > 0$，且 $A > 0$，所以这两点均为\textbf{极小值点}。
\end{enumerate}
\end{solution}

\begin{exercise}[提高：闭域上的最值]
\begin{shortexenum}
  \shortitem{1}{求 $f(x,y)=x^2+xy+y^2+x-y+1$ 在闭区域 $D=\{(x,y):|x|\leqslant 1,\;|y|\leqslant 1\}$ 上的最大值和最小值；} &
  \shortitem{2}{求 $f(x,y)=\sin x+\sin y+\sin(x+y)$ 在闭区域 $D=[0,\pi]\times[0,\pi]$ 上的最值；} \\
  \shortitem{3}{求表面积一定时体积最大的长方体的长宽高之比。} &
\end{shortexenum}
\end{exercise}
\begin{solution}
\begin{enumerate}
  \item[(1)] \textbf{求矩形区域的最值：}
  \begin{itemize}
    \item \textbf{内部驻点：} $f_x = 2x+y+1 = 0$，$f_y = x+2y-1 = 0$。解得驻点 $(-1,1)$，函数值为 $f(-1,1) = 1-1+1-1-1+1 = 0$。
    \item \textbf{边界分析：}
          边界 $x=1, y \in [-1,1]$：$g(y) = 1+y+y^2+1-y+1 = y^2+3$。最小值为 3（在 $y=0$），最大值为 4（在 $y=\pm 1$）。
          边界 $x=-1, y \in [-1,1]$：$g(y) = 1-y+y^2-1-y+1 = y^2-2y+1 = (y-1)^2$。最小值为 0（在 $y=1$），最大值为 4（在 $y=-1$）。
          边界 $y=1, x \in [-1,1]$：$g(x) = x^2+x+1+x-1+1 = x^2+2x+1 = (x+1)^2$。最小值为 0（在 $x=-1$），最大值为 4（在 $x=1$）。
          边界 $y=-1, x \in [-1,1]$：$g(x) = x^2-x+1+x+1+1 = x^2+3$。最小值为 3（在 $x=0$），最大值为 4（在 $x=\pm 1$）。
    \item \textbf{结论：} 比较所有候选值，\textbf{最大值为 4}（在 $(1,\pm 1)$ 和 $(-1,-1)$），\textbf{最小值为 0}（在 $(-1,1)$）。
  \end{itemize}

  \item[(2)] \textbf{求正方形区域的最值：}
  \begin{itemize}
    \item \textbf{内部驻点：} $f_x = \cos x+\cos(x+y) = 0$，$f_y = \cos y+\cos(x+y) = 0$。推得 $\cos x = \cos y \implies x=y$（在开区间内）。
          代入得 $\cos x + \cos 2x = 0 \implies 2\cos^2x + \cos x - 1 = 0 \implies (2\cos x-1)(\cos x+1) = 0$。
          内部驻点为 $x=y=\pi/3$。此时 $f(\pi/3, \pi/3) = \dfrac{\sqrt{3}}{2} + \dfrac{\sqrt{3}}{2} + \dfrac{\sqrt{3}}{2} = \dfrac{3\sqrt{3}}{2} \approx 2.598$。
    \item \textbf{边界分析：}
          $x=0$ 或 $y=0$ 时：$f = 2\sin y$ 或 $2\sin x$，最大值为 2，最小值为 0。
          $x=\pi$ 或 $y=\pi$ 时：由于 $\sin \pi = 0$ 以及 $\sin(\pi+t) = -\sin t$，$f$ 恒为 0。
    \item \textbf{结论：} 比较极值点，\textbf{最大值为 $\dfrac{3\sqrt{3}}{2}$}，\textbf{最小值为 0}。
  \end{itemize}

  \item[(3)] \textbf{条件极值问题转化为无条件极值：}
  设长、宽、高分别为 $x,y,z$，体积 $V = xyz$，约束表面积 $S = 2(xy+yz+zx)$。
  由约束提取 $z = \dfrac{S-2xy}{2(x+y)}$。
  体积函数 $V(x,y) = \dfrac{xy(S-2xy)}{2(x+y)}$。
  因为体积公式对于 $x, y, z$ 的轮换对称性，在物理和几何意义上当这三个变量取等值时达到对称最值。
  当 $x=y=z = \sqrt{S/6}$ 时体积最大。故表面积一定时体积最大的长方体必定是\textbf{正方体}，长宽高之比为 \textbf{1:1:1}。
\end{enumerate}
\end{solution}

\begin{exercise}[提高：最小二乘法的推广]
\begin{shortsingleenum}
  \shortsingleitem{1}{将线性最小二乘法推广到 \textbf{多项式拟合}：给定数据 $(x_i,y_i)$，$i=1,\dots,n$，求二次多项式 $y=ax^2+bx+c$ 使得残差平方和最小。导出相应的正规方程组（三元一次方程组）；}
  \shortsingleitem{2}{对数据点 $(0,1)$, $(1,0)$, $(2,1)$, $(3,4)$，用 (1) 的方法求最佳二次拟合曲线；}
  \shortsingleitem{3}{证明：若数据点 $(x_i,y_i)$ 恰好共线（存在 $a^*,b^*$ 使 $y_i=a^*x_i+b^*$ 对所有 $i$ 成立），则最小二乘解就是 $(a^*,b^*)$ 且残差平方和为零。}
\end{shortsingleenum}
\end{exercise}
\begin{solution}
\begin{enumerate}
  \item[(1)] \textbf{推导二次拟合的正规方程组：}
  构造残差平方和函数 $\Phi(a,b,c) = \sum_{i=1}^n (ax_i^2+bx_i+c-y_i)^2$。分别对 $a, b, c$ 求偏导并令其为 0：
  \begin{itemize}
    \item $\Phi_a = 2 \sum (ax_i^2+bx_i+c-y_i)x_i^2 = 0 \implies a\sum x_i^4 + b\sum x_i^3 + c\sum x_i^2 = \sum x_i^2 y_i$
    \item $\Phi_b = 2 \sum (ax_i^2+bx_i+c-y_i)x_i = 0 \implies a\sum x_i^3 + b\sum x_i^2 + c\sum x_i = \sum x_i y_i$
    \item $\Phi_c = 2 \sum (ax_i^2+bx_i+c-y_i) = 0 \implies a\sum x_i^2 + b\sum x_i + cn = \sum y_i$
  \end{itemize}
  这就是二次多项式拟合的三元正规方程组。

  \item[(2)] \textbf{具体数据点的拟合计算：}
  对给定的点集 $n=4$，计算各统计和：
  $\sum x_i = 6$，$\sum x_i^2 = 14$，$\sum x_i^3 = 36$，$\sum x_i^4 = 98$。
  $\sum y_i = 6$，$\sum x_iy_i = 14$，$\sum x_i^2y_i = 40$。
  代入正规方程组：
  \[
  \begin{cases}
  98a + 36b + 14c = 40 \implies 49a + 18b + 7c = 20 \\
  36a + 14b + 6c = 14 \implies 18a + 7b + 3c = 7 \\
  14a + 6b + 4c = 6 \implies 7a + 3b + 2c = 3
  \end{cases}
  \]
  解该三元一次方程组，得 $a=1$, $b=-2$, $c=1$。
  所以最佳拟合曲线为 $\boxed{y = x^2-2x+1 = (x-1)^2}$。
    \item[(3)] \textbf{严格共线条件下的证明：}
  由于所有数据点完全落在直线 $y=a^*x+b^*$ 上，残差平方和函数的理论最小值极值下界为 $0$。
  而将 $(a^*, b^*)$ 代入目标函数 $\Phi(a,b) = \sum(ax_i+b-y_i)^2$，得到 $\Phi(a^*,b^*) = \sum(0)^2 = 0$。
  由于平方和函数是一个严格下凸的正定二次型，其最小值必须且仅在该唯一的全导数为 0 处（即驻点）取到。因此，该最小二乘法的最优解必然且唯一地是 $(a^*, b^*)$，残差平方和为零。
\end{enumerate}
\end{solution}

\subsubsection*{A组}

\begin{exercise}[概念问答]
回答下列问题：
\begin{shortexenum}
    \shortitem{1}{$f(x,y)$ 的极值怎样定义？} &
    \shortitem{2}{$f(x,y)$ 在定义域上的最大(小)值怎样定义？} \\
    \shortitem{3}{$f(x,y)$ 取极值的必要条件是什么？} &
    \shortitem{4}{怎样判断函数的临界点是否为极值点？} \\
    \shortitem{5}{求 $f(x,y)$ 在闭区域 $D$ 上的最大(小)值的一般步骤是什么？} &
\end{shortexenum}
\end{exercise}
\begin{solution}
\begin{enumerate}
  \item[(1)] \textbf{极值的定义：}
  极值是局部概念。
  若存在点 $P_0$ 的某个邻域，使该邻域内一切异于 $P_0$ 的点 $P$ 都满足 $f(P)<f(P_0)$（或 $f(P)>f(P_0)$），则称 $f(P_0)$ 为极大值（或极小值）。
  \item[(2)] \textbf{最值的定义：}
  最值是全局概念。
  若对定义域中所有点 $P$ 都有 $f(P)\le f(P_0)$（或 $f(P)\ge f(P_0)$），则称 $f(P_0)$ 为最大值（或最小值）。
  \item[(3)] \textbf{极值的必要条件：}
  对于可导函数，若在内部点取得极值，则梯度必须为零，即 \(\displaystyle f_x(P_0)=0,\ f_y(P_0)=0\)。
  但这只是必要条件，不是充分条件；另外偏导不存在的点也可能成为极值点。
  \item[(4)] \textbf{临界点的极值判断：}
  先算 $\Delta=f_{xx}f_{yy}-(f_{xy})^2$。
  当 $\Delta>0$ 且 $f_{xx}>0$ 时为极小点；$\Delta>0$ 且 $f_{xx}<0$ 时为极大点；$\Delta<0$ 时为鞍点；$\Delta=0$ 时二阶判别法失效，需另作分析。
  \item[(5)] \textbf{闭区域求最值的依据：}
  闭区域上的最大值、最小值只能出现在内部临界点或边界上，因此需要把区域内部的驻点、不可导临界点以及边界候选点的函数值统一比较。
\end{enumerate}
\end{solution}
\begin{exercise}[临界点的分类]
求函数 $f(x,y)=8xy-\dfrac{1}{4}(x+y)^4$ 的所有临界点，并判定它们是极大点、极小点还是鞍点。
\end{exercise}
\begin{solution}
\begin{enumerate}
    \item \textbf{求解一阶偏导数并建立驻点方程组：}
    \begin{itemize}
        \item 对 $x$ 求偏导：$f_x(x,y) = 8y - \frac{1}{4} \cdot 4(x+y)^3 = 8y - (x+y)^3$。
        \item 对 $y$ 求偏导：$f_y(x,y) = 8x - \frac{1}{4} \cdot 4(x+y)^3 = 8x - (x+y)^3$。
        \item 令一阶偏导数同时为零，得到方程组：
        \[
        \begin{cases}
        8y - (x+y)^3 = 0 \quad (1)\\
        8x - (x+y)^3 = 0 \quad (2)
        \end{cases}
        \]
    \end{itemize}

    \item \textbf{求解方程组寻找临界点：}
    \begin{itemize}
        \item 将方程 $(1)$ 减去方程 $(2)$，消去 $(x+y)^3$ 项，得到 \(\displaystyle 8y - 8x = 0 \implies 8(y-x) = 0 \implies y = x\)。
        \item 将 $y = x$ 代回方程 $(2)$ 中，得 \(\displaystyle 8x - (x+x)^3 = 0 \implies 8x - (2x)^3 = 0 \implies 8x - 8x^3 = 0\)。
        \item 提取公因式：\(\displaystyle 8x(1-x^2) = 0 \implies 8x(1-x)(1+x) = 0\)。
        \item 解得三个实数根：$x_1 = 0, x_2 = 1, x_3 = -1$。
        \item 由于 $y = x$，对应得到三个临界点：$P_1(0,0), P_2(1,1), P_3(-1,-1)$。
    \end{itemize}

    \item \textbf{计算二阶偏导数以构造 Hesse 矩阵：}
    \begin{itemize}
        \item $A = f_{xx}(x,y) = \frac{\partial}{\partial x}[8y - (x+y)^3] = -3(x+y)^2 \cdot 1 = -3(x+y)^2$。
        \item $C = f_{yy}(x,y) = \frac{\partial}{\partial y}[8x - (x+y)^3] = -3(x+y)^2 \cdot 1 = -3(x+y)^2$。
        \item $B = f_{xy}(x,y) = \frac{\partial}{\partial y}[8y - (x+y)^3] = 8 - 3(x+y)^2$。
    \end{itemize}

    \item \textbf{利用 Hesse 判别法 $\Delta = AC - B^2$ 对各临界点进行分类：}
    \begin{itemize}
        \item \textbf{判定点 $P_1(0,0)$：}
        代入 $x=0, y=0$：
        $A = -3(0)^2 = 0$；
        $B = 8 - 3(0)^2 = 8$；
        $C = -3(0)^2 = 0$。
        计算判别式：$\Delta = AC - B^2 = 0 \cdot 0 - 8^2 = -64$。
        因为 $\Delta < 0$，所以 $(0,0)$ 是\textbf{鞍点}。

        \item \textbf{判定点 $P_2(1,1)$：}
        代入 $x=1, y=1$，此时 $x+y = 2$：
        $A = -3(2)^2 = -12$；
        $B = 8 - 3(2)^2 = 8 - 12 = -4$；
        $C = -3(2)^2 = -12$。
        计算判别式：$\Delta = AC - B^2 = (-12) \cdot (-12) - (-4)^2 = 144 - 16 = 128$。
        因为 $\Delta > 0$ 且 $A = -12 < 0$，所以 $(1,1)$ 是\textbf{极大值点}。

        \item \textbf{判定点 $P_3(-1,-1)$：}
        代入 $x=-1, y=-1$，此时 $x+y = -2$：
        $A = -3(-2)^2 = -12$；
        $B = 8 - 3(-2)^2 = 8 - 12 = -4$；
        $C = -3(-2)^2 = -12$。
        计算判别式：$\Delta = AC - B^2 = (-12) \cdot (-12) - (-4)^2 = 144 - 16 = 128$。
        因为 $\Delta > 0$ 且 $A = -12 < 0$，所以 $(-1,-1)$ 也是\textbf{极大值点}。
    \end{itemize}
\end{enumerate}
\end{solution}

\begin{exercise}[各类函数的极值判定]
求下列函数的极值：
\par\noindent
\begin{tabular}{@{}p{0.45\textwidth}p{0.45\textwidth}@{}}
(1) $f(x,y)=(x+y)(xy+1)$ & (2) $f(x,y)=\mathrm{e}^{2x}(x+y^2+2y)$ \\
(3) $f(x,y)=x^3+y^3-3(x^2+y^2)$ & (4) $f(x,y)=\sin x+\sin y+\sin(x+y)$ \\
\end{tabular}\\
注：第 (4) 小题中限制变量范围为 $0 < x < \pi, 0 < y < \pi$。
\end{exercise}
\begin{solution}
\begin{enumerate}
    \item[\textbf{解答(1)}] \textbf{对于函数 $f(x,y)=(x+y)(xy+1) = x^2y + x + xy^2 + y$：}
    \begin{enumerate}
        \item \textbf{求驻点：}
        $f_x = 2xy + 1 + y^2 = (x+y)^2 - x^2 + 1 = y^2 + 2xy + 1 = 0$；
        $f_y = x^2 + 2xy + 1 = 0$。
        两式相减得：$y^2 - x^2 = 0 \implies (y-x)(y+x) = 0 \implies y = x$ 或 $y = -x$。
        \begin{itemize}
            \item 若 $y = x$，代入 $f_x=0$ 得 $x^2 + 2x^2 + 1 = 0 \implies 3x^2 = -1$，在实数范围内无解。
            \item 若 $y = -x$，代入 $f_x=0$ 得 $(-x)^2 + 2x(-x) + 1 = 0 \implies x^2 - 2x^2 + 1 = 0 \implies x^2 = 1 \implies x = \pm 1$。
        \end{itemize}
        当 $x=1$ 时，$y=-1$；当 $x=-1$ 时，$y=1$。故驻点为 $P_1(1,-1)$ 和 $P_2(-1,1)$。

        \item \textbf{计算二阶偏导数：}
        $A = f_{xx} = 2y$；
        $B = f_{xy} = 2y + 2x = 2(x+y)$；
        $C = f_{yy} = 2x$。

        \item \textbf{判别：}
        在 $P_1(1,-1)$ 处，$A=-2,\ B=0,\ C=2$，故 $\Delta = AC - B^2 = -4 < 0$，此点为\textbf{鞍点}；在 $P_2(-1,1)$ 处，$A=2,\ B=0,\ C=-2$，故 $\Delta = AC - B^2 = -4 < 0$，此点也为\textbf{鞍点}。因此该函数\textbf{无极值}。
    \end{enumerate}

    \item[\textbf{解答(2)}] \textbf{对于函数 $f(x,y)=\mathrm{e}^{2x}(x+y^2+2y)$：}
    \begin{enumerate}
        \item \textbf{求驻点：}
        利用乘积求导法则：
        $f_x = 2\mathrm{e}^{2x}(x+y^2+2y) + \mathrm{e}^{2x}(1) = \mathrm{e}^{2x}(2x + 2y^2 + 4y + 1) = 0$；
        $f_y = \mathrm{e}^{2x}(2y+2) = 0$。
        因为 $\mathrm{e}^{2x} \neq 0$，由 $f_y=0$ 必须有 $2y+2 = 0 \implies y = -1$。
        将 $y = -1$ 代入 $f_x=0$ 的括号部分得：$2x + 2(-1)^2 + 4(-1) + 1 = 0 \implies 2x + 2 - 4 + 1 = 0 \implies 2x - 1 = 0 \implies x = 1/2$。
        故唯一驻点为 $P(1/2, -1)$。

        \item \textbf{计算二阶偏导数：}
        $A = f_{xx} = 2\mathrm{e}^{2x}(2x + 2y^2 + 4y + 1) + \mathrm{e}^{2x}(2) = 2\mathrm{e}^{2x}(2x + 2y^2 + 4y + 2) = 4\mathrm{e}^{2x}(x + y^2 + 2y + 1)$；
        $B = f_{xy} = \frac{\partial}{\partial y}[\mathrm{e}^{2x}(2x + 2y^2 + 4y + 1)] = \mathrm{e}^{2x}(4y + 4) = 4\mathrm{e}^{2x}(y+1)$；
        $C = f_{yy} = \frac{\partial}{\partial y}[\mathrm{e}^{2x}(2y+2)] = 2\mathrm{e}^{2x}$。

        \item \textbf{判别：}
        在 $P(1/2, -1)$ 处，$2x = 1$，故 $\mathrm{e}^{2x} = \mathrm{e}$。
        $A = 4\mathrm{e}(1/2 + 1 - 2 + 1) = 4\mathrm{e}(1/2) = 2\mathrm{e}$；
        $B = 4\mathrm{e}(-1+1) = 0$；
        $C = 2\mathrm{e}$。
        $\Delta = AC - B^2 = (2\mathrm{e})(2\mathrm{e}) - 0 = 4\mathrm{e}^2 > 0$。
        因为 $\Delta > 0$ 且 $A = 2\mathrm{e} > 0$，故该点为\textbf{极小值点}。
        \textbf{结论：} 极小值为 $f(1/2, -1) = \mathrm{e}^1(1/2 + 1 - 2) = \mathrm{e}(-1/2) = \mathbf{-\dfrac{\mathrm{e}}{2}}$。
    \end{enumerate}

    \item[\textbf{解答(3)}] \textbf{对于函数 $f(x,y)=x^3+y^3-3(x^2+y^2)$：}
    \begin{enumerate}
        \item \textbf{求驻点：}
        $f_x = 3x^2 - 6x = 3x(x-2) = 0 \implies x=0$ 或 $x=2$；
        $f_y = 3y^2 - 6y = 3y(y-2) = 0 \implies y=0$ 或 $y=2$。
        交叉组合得到四个驻点：$P_1(0,0), P_2(0,2), P_3(2,0), P_4(2,2)$。

        \item \textbf{计算二阶偏导数：}
        $A = f_{xx} = 6x - 6 = 6(x-1)$；
        $B = f_{xy} = 0$；
        $C = f_{yy} = 6y - 6 = 6(y-1)$。
        由于 $B=0$，判别式 $\Delta = AC = 36(x-1)(y-1)$。

        \item \textbf{逐一判别：}
        在 $P_1(0,0)$ 处，$A=-6<0,\ C=-6,\ \Delta=36>0$，故 $(0,0)$ 为\textbf{极大值点}，极大值为 $\mathbf{0}$；在 $P_2(0,2)$ 处，$A=-6,\ C=6,\ \Delta=-36<0$，为\textbf{鞍点}；在 $P_3(2,0)$ 处，$A=6,\ C=-6,\ \Delta=-36<0$，为\textbf{鞍点}；在 $P_4(2,2)$ 处，$A=6>0,\ C=6,\ \Delta=36>0$，故 $(2,2)$ 为\textbf{极小值点}，极小值为 $8+8-3(4+4)=\mathbf{-8}$。
    \end{enumerate}

    \item[\textbf{解答(4)}] \textbf{对于函数 $f(x,y)=\sin x+\sin y+\sin(x+y)$：}
    \begin{enumerate}
        \item \textbf{求驻点：}
        $f_x = \cos x + \cos(x+y) = 0 \implies \cos x = -\cos(x+y) \implies \cos x = \cos(\pi - (x+y))$；
        $f_y = \cos y + \cos(x+y) = 0 \implies \cos y = -\cos(x+y)$。
        由这两式可知 $\cos x = \cos y$。因为 $0 < x, y < \pi$，在此区间内余弦函数是单调递减的一一映射，所以必有 $x = y$。
        将 $y = x$ 代回 $f_x=0$ 得到方程：$\cos x + \cos 2x = 0$。
        利用倍角公式 $\cos 2x = 2\cos^2 x - 1$，代入得：
        $2\cos^2 x + \cos x - 1 = 0 \implies (2\cos x - 1)(\cos x + 1) = 0$。
        因为 $0 < x < \pi$，$\cos x \neq -1$（即 $x \neq \pi$），所以必须满足 $2\cos x - 1 = 0 \implies \cos x = 1/2 \implies x = \pi/3$。
        故唯一驻点为 $P(\pi/3, \pi/3)$。

        \item \textbf{计算二阶偏导数：}
        $A = f_{xx} = -\sin x - \sin(x+y)$；
        $B = f_{xy} = -\sin(x+y)$；
        $C = f_{yy} = -\sin y - \sin(x+y)$。

        \item \textbf{判别：}
        在 $P(\pi/3, \pi/3)$ 处，$x+y = 2\pi/3$。
        $A = -\sin(\pi/3) - \sin(2\pi/3) = -\frac{\sqrt{3}}{2} - \frac{\sqrt{3}}{2} = -\sqrt{3}$；
        $B = -\sin(2\pi/3) = -\frac{\sqrt{3}}{2}$；
        $C = -\sin(\pi/3) - \sin(2\pi/3) = -\sqrt{3}$。
        计算判别式：$\Delta = AC - B^2 = (-\sqrt{3})(-\sqrt{3}) - (-\frac{\sqrt{3}}{2})^2 = 3 - \frac{3}{4} = \frac{9}{4} > 0$。
        因为 $\Delta > 0$ 且 $A = -\sqrt{3} < 0$，该点为\textbf{极大值点}。
        \textbf{结论：} 极大值为 $f(\pi/3, \pi/3) = \sin(\pi/3) + \sin(\pi/3) + \sin(2\pi/3) = \frac{\sqrt{3}}{2} + \frac{\sqrt{3}}{2} + \frac{\sqrt{3}}{2} = \mathbf{\dfrac{3\sqrt{3}}{2}}$。
    \end{enumerate}
\end{enumerate}
\end{solution}

\begin{exercise}[极值点与鞍点的求解]
求下列函数的极值点和鞍点。
\begin{shortexenum}
  \shortitem{1}{$f(x,y)=x^2+y^3-3xy$；} &
  \shortitem{2}{$f(x,y)=xy+\ln x+y^2-10\ (x>0)$。}
\end{shortexenum}
\end{exercise}
\begin{solution}
\begin{enumerate}
    \item[\textbf{解答(1)}] \textbf{对于函数 $f(x,y)=x^2+y^3-3xy$：}
    \begin{enumerate}
        \item \textbf{求驻点：}
        $f_x = 2x - 3y = 0 \implies x = \frac{3}{2}y$；
        $f_y = 3y^2 - 3x = 0 \implies x = y^2$。
        将 $x=\frac{3}{2}y$ 代入 $x=y^2$ 中，得 $y^2=\frac{3}{2}y \implies y\left(y-\frac{3}{2}\right)=0$，故 $y=0$ 或 $y=\frac{3}{2}$；对应地 $x=0$ 或 $x=\frac{9}{4}$，所以驻点为 $P_1(0,0)$ 和 $P_2\left(\frac{9}{4}, \frac{3}{2}\right)$。

        \item \textbf{计算二阶偏导数：}
        $A = f_{xx} = 2$；
        $B = f_{xy} = -3$；
        $C = f_{yy} = 6y$。

        \item \textbf{判别分类：}
        在 $P_1(0,0)$ 处，$A=2,\ B=-3,\ C=0$，故 $\Delta = AC - B^2 = -9 < 0$，因此 $(0,0)$ 是\textbf{鞍点}；在 $P_2\left(\frac{9}{4}, \frac{3}{2}\right)$ 处，$A=2,\ B=-3,\ C=9$，故 $\Delta = AC - B^2 = 18 - 9 = 9 > 0$，且 $A=2>0$，因此 $\left(\frac{9}{4}, \frac{3}{2}\right)$ 是\textbf{极小值点}。
    \end{enumerate}

    \item[\textbf{解答(2)}] \textbf{对于函数 $f(x,y)=xy+\ln x+y^2-10 \quad (x>0)$：}
    \begin{enumerate}
        \item \textbf{求驻点：}
        $f_x = y + \frac{1}{x} = 0 \implies y = -\frac{1}{x}$；
        $f_y = x + 2y = 0 \implies x = -2y$。
        将 $y = -\frac{1}{x}$ 代入第二个方程，得：
        $x = -2\left(-\frac{1}{x}\right) = \frac{2}{x} \implies x^2 = 2$。
        由于定义域限制 $x > 0$，所以取正根 $x = \sqrt{2}$。
        将其代回求 $y$：$y = -\frac{1}{\sqrt{2}} = -\frac{\sqrt{2}}{2}$。
        故唯一驻点为 $P(\sqrt{2}, -\sqrt{2}/2)$。

        \item \textbf{计算二阶偏导数：}
        $A = f_{xx} = -\frac{1}{x^2}$；
        $B = f_{xy} = 1$；
        $C = f_{yy} = 2$。

        \item \textbf{判别分类：}
        在 $P(\sqrt{2}, -\sqrt{2}/2)$ 处：
        $A = -\frac{1}{(\sqrt{2})^2} = -\frac{1}{2}$；
        $B = 1$；
        $C = 2$。
        计算判别式：$\Delta = AC - B^2 = \left(-\frac{1}{2}\right)(2) - 1^2 = -1 - 1 = -2 < 0$。
        因为 $\Delta < 0$，该点为\textbf{鞍点}。该函数不存在极值点。
    \end{enumerate}
\end{enumerate}
\end{solution}

\begin{exercise}[超越函数的极值特性证明]
证明函数 $z=(1+\mathrm{e}^y)\cos x-y\mathrm{e}^y$ 有无穷多个极大值点，但没有任何极小值点。
\end{exercise}
\begin{solution}
\begin{enumerate}
    \item 求解一阶偏导数并确定驻点族：
    \[
    z_x = -(1+\mathrm{e}^y)\sin x,\qquad
    z_y = \mathrm{e}^y\cos x - (\mathrm{e}^y + y\mathrm{e}^y)
    = \mathrm{e}^y(\cos x - 1 - y).
    \]
    令 $z_x = 0$，由于 $\mathrm{e}^y > 0$ 且 $1+\mathrm{e}^y \neq 0$，只能有 $\sin x = 0$，即 $x = k\pi$（$k \in \mathbb{Z}$）；令 $z_y = 0$，则 $\cos x - 1 - y = 0$，即 $y = \cos x - 1$。将 $x = k\pi$ 代入后，若 $k=2m$，则 $\cos x = 1,\ y=0$，得到第一类驻点族 $P_m(2m\pi, 0)$；若 $k=2m+1$，则 $\cos x = -1,\ y=-2$，得到第二类驻点族 $Q_m((2m+1)\pi, -2)$。

    \item \textbf{计算二阶偏导数建立 Hesse 行列式参数：}
    \begin{align*}
    A &= z_{xx} = -(1+\mathrm{e}^y)\cos x,\qquad
    B = z_{xy} = -\mathrm{e}^y\sin x,\\
    C &= z_{yy} = \mathrm{e}^y(\cos x - 1 - y) - \mathrm{e}^y
    = \mathrm{e}^y(\cos x - 2 - y).
    \end{align*}

    \item \textbf{对两类驻点族分别进行极值判定：}
    对第一类驻点族 $P_m(2m\pi, 0)$，由于 $\cos(2m\pi) = 1,\ \sin(2m\pi) = 0,\ y = 0$，代入得 $A = -2,\ B = 0,\ C = -1$，故 $\Delta = AC - B^2 = 2 > 0$，且 $A=-2<0$，所以所有点 $P_m(2m\pi, 0)$ 都是\textbf{严格的极大值点}；又因整数 $m$ 有无穷多个，故原函数有无穷多个极大值。
    对第二类驻点族 $Q_m((2m+1)\pi, -2)$，有 $\cos((2m+1)\pi) = -1,\ \sin((2m+1)\pi) = 0,\ y = -2$，故 $A = 1+\mathrm{e}^{-2},\ B = 0,\ C = -\mathrm{e}^{-2}$，从而 $\Delta = AC - B^2 = -\mathrm{e}^{-2} - \mathrm{e}^{-4} < 0$，所以所有点 $Q_m((2m+1)\pi, -2)$ 都是\textbf{鞍点}。
    \textbf{结论：} 穷举了所有可能取得极值的驻点，发现它们要么是极大值点，要么是鞍点，没有任何一个点满足极小值的条件（即 $\Delta > 0$ 且 $A > 0$）。证毕。
\end{enumerate}
\end{solution}

\begin{exercise}[有界闭区域求最值]
求函数 $f(x,y)=x^3-4x^2+2xy-y^2$ 在由不等式 $-5\leqslant x\leqslant5, -1\leqslant y\leqslant1$ 所确定的矩形闭区域 $D$ 上的最大值与最小值。
\end{exercise}
\begin{solution}
\begin{enumerate}
    \item 寻找区域 $D$ 内部的可能极值点（驻点）
    由 $f_x = 3x^2 - 8x + 2y = 0,\ f_y = 2x - 2y = 0$ 得 $y = x$；代入前式得 $3x^2 - 6x = 0$，即 $x = 0$ 或 $x = 2$，对应驻点为 $P_1(0,0)$ 和 $P_2(2,2)$。其中 $(0,0)$ 位于区域内部，而 $(2,2)$ 因 $y = 2 > 1$ 位于区域外部，故只保留 $(0,0)$，且 $f(0,0)=\mathbf{0}$。

    \item 区域 $D$ 是矩形，有四条边界线段。分别化为一元函数考察：对右边界 $L_1$（$x = 5,\ y \in [-1, 1]$），有 $g_1(y) = f(5,y) = -y^2 + 10y + 25$，它在区间 $[-1,1]$ 上单调递增，因此 $g_1(-1)=\mathbf{14},\ g_1(1)=\mathbf{34}$。对左边界 $L_2$（$x = -5,\ y \in [-1, 1]$），有 $g_2(y) = f(-5,y) = -y^2 - 10y - 225$，它在区间 $[-1,1]$ 上单调递减，因此 $g_2(-1)=\mathbf{-216},\ g_2(1)=\mathbf{-236}$。

    对上边界 $L_3$（$y = 1,\ x \in [-5, 5]$），有 $g_3(x) = f(x,1) = x^3 - 4x^2 + 2x - 1$。由 $g_3'(x) = 3x^2 - 8x + 2 = 0$ 得驻点 $x_{1,2} = \frac{4 \pm \sqrt{10}}{3}$，其中极大值点为 $x_1 = \frac{4 - \sqrt{10}}{3}$，极小值点为 $x_2 = \frac{4 + \sqrt{10}}{3}$。又
    \[
      x_{right} = x_2 + \frac{x_2 - x_1}{2} = \frac{4 + 2\sqrt{10}}{3}<5,\qquad
      x_{left} = x_1 - \frac{x_2 - x_1}{2} = \frac{4 - 2\sqrt{10}}{3}>-5,
    \]
    故区间端点 $x=\pm 5$ 均落在 $[x_{left}, x_{right}]$ 外部，所以最值在端点取到，且 $g_3(5)=\mathbf{34},\ g_3(-5)=\mathbf{-236}$。

    对下边界 $L_4$（$y = -1,\ x \in [-5, 5]$），有 $g_4(x) = f(x,-1) = x^3 - 4x^2 - 2x - 1$。由 $g_4'(x) = 3x^2 - 8x - 2 = 0$ 得驻点 $x_{1,2} = \frac{4 \pm \sqrt{22}}{3}$，其中极大值点为 $x_1 = \frac{4 - \sqrt{22}}{3}$，极小值点为 $x_2 = \frac{4 + \sqrt{22}}{3}$。同理有
    \[
      x_{right} = x_2 + \frac{x_2 - x_1}{2} = \frac{4 + 2\sqrt{22}}{3}<5,\qquad
      x_{left} = x_1 - \frac{x_2 - x_1}{2} = \frac{4 - 2\sqrt{22}}{3}>-5,
    \]
    因而区间端点 $x=\pm 5$ 也都在上述阈值范围之外，所以最值仍只可能在端点取得，且 $g_4(5)=\mathbf{14},\ g_4(-5)=\mathbf{-216}$。
    \item 内部极值：$f(0,0) = 0$。边界极值集中在矩形的四个顶点：$f(5, 1) = 34$；$f(5, -1) = 14$；$f(-5, 1) = -236$；$f(-5, -1) = -216$。比较得出最大值为 \textbf{34}，在点 $(5, 1)$ 处取得。最小值为 \textbf{-236}，在点 $(-5, 1)$ 处取得。
\end{enumerate}
\end{solution}

\begin{exercise}[利用最小二乘法进行线性拟合]
设水深与流速间的测量数据如下：
\par\noindent{\centering
\begin{tabular}{|c|c|c|c|c|c|}
\hline
水深 $x$/m & 0 & 0.1 & 0.2 & 0.3 & 0.4 \\
\hline
流速 $y$/(m/s) & 3.195 & 3.229 9 & 3.253 2 & 3.261 1 & 3.251 6 \\
\hline
\end{tabular}
\par\vspace{0.5em}
\begin{tabular}{|c|c|c|c|c|c|}
\hline
水深 $x$/m & 0.5 & 0.6 & 0.7 & 0.8 & 0.9 \\
\hline
流速 $y$/(m/s) & 3.228 2 & 3.180 7 & 3.126 6 & 3.059 4 & 2.975 7 \\
\hline
\end{tabular}
\par}
试用最小二乘法求最佳拟合直线 $y=ax+b$。
\end{exercise}
\begin{solution}
\begin{enumerate}
  \item \textbf{计算各统计求和量}：
  这里数据点总数 $n=10$，并有 \(\displaystyle \sum_{i=1}^{10} x_i = 4.5,\ \sum_{i=1}^{10} x_i^2 = 2.85,\ \sum_{i=1}^{10} y_i = 31.7614,\ \sum_{i=1}^{10} x_i y_i = 14.08939\)。

  \item \textbf{构建正规方程组：}
  将算得的统计量代入最小二乘的系数公式体系（正规方程）：
  \[
  \begin{cases}
  10b + 4.5a = 31.7614 \\
  4.5b + 2.85a = 14.08939
  \end{cases}
  \]

  \item \textbf{解线性方程组求出参数：}
  对第一式两边同乘以 $0.45$，并将其从第二式中减去，经代数运算得 $a \approx -0.2464,\ b \approx 3.2870$。

  \item \textbf{写出拟合直线方程：}
  故最佳拟合直线近似为 $\boxed{y = -0.2464x + 3.2870}$。
\end{enumerate}
\end{solution}

\begin{exercise}[实际问题中的条件最值]
试设计一个宽 $w$，长 $l$，高 $h$ 的长方形纸板盒，使其具有 $512\ \mathrm{cm}^3$ 的容积。这个盒子的侧边的成本是 $0.1$ 元/$\mathrm{cm}^2$，而上下底的成本是 $0.2$ 元/$\mathrm{cm}^2$。试求这个盒子的尺寸，使得所用材料的成本最低。
\end{exercise}
\begin{solution}
\begin{enumerate}
  \item \textbf{建立约束条件与成本目标函数：}
  容积恒定约束：$V = lwh = 512 \implies h = \dfrac{512}{lw}$。
  表面材料成本函数为：$C = 0.1(2lh + 2wh) + 0.2(2lw) = 0.2(lh+wh) + 0.4lw$。

  \item \textbf{消去高 $h$ 转为二元函数求极值：}
  将 $h$ 代入成本函数，得关于 $l$ 和 $w$ 的二元函数 \(\displaystyle C(l,w)=\frac{102.4}{w}+\frac{102.4}{l}+0.4lw\)。

  \item \textbf{求偏导并寻找驻点：}
  对 $l, w$ 分别求导，得 $C_l = -\dfrac{102.4}{l^2} + 0.4w = 0 \implies wl^2 = 256$，$C_w = -\dfrac{102.4}{w^2} + 0.4l = 0 \implies lw^2 = 256$。

  \item \textbf{解方程得出最优尺寸：}
  两式相比（由于边长必大于 0），有 $\dfrac{l}{w} = 1 \implies l=w$。
  代入得 $l^3 = 256 \implies l = 2^{8/3} = 4\sqrt[3]{4} \approx 6.35\ \mathrm{cm}$。则 $w = 4\sqrt[3]{4}$。
  高度 $h = \dfrac{512}{l^2} = \dfrac{512}{256^{2/3}} = \dfrac{2^9}{2^{16/3}} = 2^{11/3} = 8\sqrt[3]{4}\ \mathrm{cm}$。
  根据物理意义，所求出的驻点即为全局成本最小值的尺寸。
  因此，尺寸设计应为：\textbf{长与宽均为 $4\sqrt[3]{4}\ \mathrm{cm}$，高为 $8\sqrt[3]{4}\ \mathrm{cm}$}。
\end{enumerate}
\end{solution}

\subsubsection*{B组}

\begin{exercise}[含指数与三角的极值判定]
求函数 $f(x,y)=\mathrm{e}^x(1-\cos y)$ 的临界点，并判定它们是极大(小)值点还是鞍点。
\end{exercise}
\begin{solution}
\begin{enumerate}
  \item \textbf{求临界点方程：}
  由 $f_x = \mathrm{e}^x(1-\cos y) = 0$ 得 $\cos y = 1$，即 $y = 2k\pi$（$k \in \mathbb Z$）；由 $f_y = \mathrm{e}^x\sin y = 0$ 得 $\sin y = 0$，即 $y = n\pi$。两式同时成立时必须有 $y = 2k\pi$，而 $x$ 可取任意实数，所以临界点是一系列水平直线 $(x, 2k\pi)$。

  \item \textbf{分析临界点的性质：}
  由于这并非孤立驻点，Hesse 判别法 $\Delta=0$ 会失效（因 $f_{yy}=\mathrm{e}^x\cos(2k\pi)=\mathrm{e}^x$，而 $f_{xy}=0$，$f_{xx}=0$，得 $\Delta=0$）。
  只能使用直接判断法：函数在此直线上值为 $f(x, 2k\pi) = 0$。
  而当 $y \neq 2k\pi$ 时，无论 $x$ 取何值，$1-\cos y > 0$ 且 $\mathrm{e}^x > 0$，故整体上 $f(x,y) > 0$ 恒成立。
  因此，在此点集的任意邻域内，函数值总是不小于 0 的。结论为：整条线上的点均为\textbf{非严格极小值点}。
\end{enumerate}
\end{solution}

\begin{exercise}[二阶曲面的隐函数极值一]
求由方程 $2x^2+2y^2+z^2+8xz-z+8=0$ 所确定的隐函数 $z=z(x,y)$ 的极值。
\end{exercise}
\begin{solution}
\begin{enumerate}
  \item \textbf{利用隐函数一阶求导法则寻驻点：}
  设 $F = 2x^2+2y^2+z^2+8xz-z+8 = 0$。隐函数的极值点必须满足偏导为 0；由 $z_x = -\dfrac{F_x}{F_z} = -\dfrac{4x+8z}{2z+8x-1} = 0$ 得 $x = -2z$，由 $z_y = -\dfrac{F_y}{F_z} = -\dfrac{4y}{2z+8x-1} = 0$ 得 $y = 0$。

  \item \textbf{代回方程求可能发生极值的极值点坐标：}
  将 $x=-2z, y=0$ 代回原方程 $F=0$，得 \(\displaystyle 2(-2z)^2+z^2+8(-2z)z-z+8=0\implies -7z^2-z+8=0\)。
  即 $7z^2 + z - 8 = 0 \implies (7z+8)(z-1) = 0$。
  解得 $z=1$ 或 $z=-8/7$。
  对于 $z=1$，对应 $P_1(-2,0)$；对于 $z=-8/7$，对应 $P_2(16/7, 0)$。

  \item \textbf{求二阶隐偏导数并代入判别：}
  在驻点处，由于 $z_x = z_y = 0$，二阶偏导只需考虑显式二阶梯度的比值，且 \(\displaystyle F_{xx} = 4,\qquad F_{xy} = 0,\qquad F_{yy} = 4\)。
  在 $P_1(-2,0)$ 处（$z=1$），$F_z = 2(1)+8(-2)-1 = -15$，故
  \[
  z_{xx} = -F_{xx}/F_z = 4/15 > 0,\qquad
  z_{yy} = 4/15,\qquad z_{xy} = 0.
  \]
  于是 $\Delta > 0,\ z_{xx} > 0$，此点处 $z$ 有\textbf{极小值 $1$}。在 $P_2(16/7, 0)$ 处（$z=-8/7$），$F_z = 2(-8/7)+8(16/7)-1 = 15$，故
  \(\displaystyle z_{xx} = -4/15 < 0,\quad z_{yy} = -4/15,\quad z_{xy} = 0\)。于是 $\Delta > 0,\ z_{xx} < 0$，此点处 $z$ 有\textbf{极大值 $-\dfrac{8}{7}$}。
\end{enumerate}
\end{solution}

\begin{exercise}[三角形面积最大值]
在半径为 $R$ 的圆内求一内接三角形，使其面积最大。
\end{exercise}
\begin{solution}
\begin{enumerate}
  \item 设该内接三角形的三个顶点为 $A, B, C$，圆心为 $O$。要使三角形面积最大，圆心 $O$ 必然在三角形内部（若圆心在边界或外部，三角形面积不到圆面积的一半，显然不是最大）。
  设三角形的三条边 $AB, BC, CA$ 所对的圆心角分别为 $\alpha, \beta, \gamma$。
  由于圆心在三角形内部，这三个圆心角满足 \(\displaystyle \alpha>0,\ \beta>0,\ \gamma>0,\ \alpha+\beta+\gamma=2\pi\)。
  连接 $OA, OB, OC$，整个三角形 $ABC$ 的面积 $S$ 可以拆分为三个等腰三角形面积之和：
  \(\displaystyle S=S_{\triangle OAB}+S_{\triangle OBC}+S_{\triangle OCA}=\frac{1}{2}R^2(\sin\alpha+\sin\beta+\sin\gamma)\)。
  \item 利用约束条件 $\gamma = 2\pi - (\alpha + \beta)$ 代入面积公式，且由诱导公式知 $\sin(2\pi - (\alpha+\beta)) = -\sin(\alpha+\beta)$。
  目标函数转化为关于 $\alpha, \beta$ 的二元函数 \(\displaystyle S(\alpha,\beta)=\frac{1}{2}R^2\left[\sin\alpha+\sin\beta-\sin(\alpha+\beta)\right]\)。
  其定义域为开区域 $D = \{ (\alpha, \beta) \mid \alpha > 0, \beta > 0, \alpha+\beta < 2\pi \}$。
  \item 对 $\alpha, \beta$ 分别求偏导数，并令其等于零：
  \[
  S_\alpha = \dfrac{1}{2}R^2 \left[ \cos\alpha - \cos(\alpha+\beta) \right] = 0,\qquad
  S_\beta = \dfrac{1}{2}R^2 \left[ \cos\beta - \cos(\alpha+\beta) \right] = 0.
  \]
  因此 $\cos\alpha = \cos(\alpha+\beta)$，$\cos\beta = \cos(\alpha+\beta)$，从而 $\cos\alpha = \cos\beta$。因为 $\alpha, \beta \in (0, 2\pi)$，所以 $\alpha = \beta$ 或 $\alpha + \beta = 2\pi$；后者会导致 $\gamma = 0$，不能构成三角形，故只能取 $\alpha = \beta$。将其代回第一个方程，得：
  将其代回第一个方程，得 \(\displaystyle \cos\alpha = \cos(2\alpha)\)。利用二倍角公式 $\cos(2\alpha) = 2\cos^2\alpha - 1$，得到关于 $\cos\alpha$ 的一元二次方程：
  \(\displaystyle 2\cos^2\alpha - \cos\alpha - 1 = 0 \implies (2\cos\alpha + 1)(\cos\alpha - 1) = 0\)。
  解得 $\cos\alpha = 1$ 或 $\cos\alpha = -\dfrac{1}{2}$。
  因为 $\alpha > 0$，故 $\cos\alpha \neq 1$。所以只能是 $\cos\alpha = -\dfrac{1}{2}$，解得唯一驻点 $\alpha = \dfrac{2\pi}{3}$。
  从而 $\beta = \dfrac{2\pi}{3}$。
  \item 计算二阶偏导数：
  在驻点处，\(\displaystyle A=S_{\alpha\alpha}=-\dfrac{\sqrt{3}}{2}R^2\)，
  \(\displaystyle B=S_{\alpha\beta}=-\dfrac{\sqrt{3}}{4}R^2\)，
  \(\displaystyle C=S_{\beta\beta}=-\dfrac{\sqrt{3}}{2}R^2\)。
  计算判别式：\(\displaystyle \Delta=AC-B^2=\frac{3}{4}R^4-\frac{3}{16}R^4=\frac{9}{16}R^4\)。
  因为 $\Delta > 0$ 且 $A < 0$，该驻点确实为\textbf{局部极大值点}。
  \item 由于在开区域 $D$ 内仅存在这一个极值点，此时 $\alpha = \beta = \dfrac{2\pi}{3}$，故 $\gamma = 2\pi - \dfrac{4\pi}{3} = \dfrac{2\pi}{3}$。
  三角形三条边所对的圆心角全部相等（均为 $120^\circ$），这意味着该内接三角形为\textbf{正三角形}。
\end{enumerate}
\end{solution}

\begin{exercise}[积分求极值]
求 $a,b$，使积分 $I=\int_0^1 (a+bx-x^2)^2\mathrm{d}x$ 的值最小。
\end{exercise}
\begin{solution}
  目标函数 $I(a,b)$ 是含有积分的变元函数，根据含参积分求导法则（直接在积分号内求导），有
  \(\displaystyle I_a=2 \int_0^1 (a+bx-x^2)\,\mathrm{d}x
  =2\left( a + \frac{b}{2} - \frac{1}{3} \right)=0\)，
  \(\displaystyle I_b=2 \int_0^1 (a+bx-x^2)x\,\mathrm{d}x
  =2\left( \frac{a}{2} + \frac{b}{3} - \frac{1}{4} \right)=0\)。
  将上述求导令为零的结果转化为关于 $a,b$ 的方程组
  \(\displaystyle 6a+3b=2,\quad 6a+4b=3\)。
  第二式减去第一式得 $b = 1$。
  将 $b=1$ 代入第一式得 $6a + 3 = 2 \implies 6a = -1 \implies a = -\dfrac{1}{6}$。
\end{solution}

% ============================================================
% §7.11 约束最优化问题
% ============================================================

\clearpage

\subsection{第十一节习题：约束最优化问题}

\begin{exercise}[基础：单约束 Lagrange 乘数法]
求下列函数在指定约束下的极值与最值：
\begin{shortexenum}
  \shortitem{1}{求 $f(x,y)=x+y$ 在约束 $x^2+y^2=1$ 下的最大值和最小值；} &
  \shortitem{2}{求 $f(x,y)=\mathrm{e}^{xy}$ 在约束 $x^2+2y^2=1$ 下的极值点；} \\
  \shortitem{3}{用 Lagrange 乘数法求原点到直线 $\dfrac{x}{a}+\dfrac{y}{b}=1$ 的距离。} &
\end{shortexenum}
\end{exercise}
\begin{solution}
\begin{enumerate}
  \item[(1)] \textbf{求 $f(x,y)=x+y$ 在圆上的极值：}
  \begin{enumerate}
    \item \textbf{构造辅助函数：}
    \(\displaystyle L(x,y;\lambda) = x+y+\lambda(x^2+y^2-1)\)。
    \item \textbf{求偏导并令为零：}
    由 $L_x = 1+2\lambda x = 0$ 与 $L_y = 1+2\lambda y = 0$ 得 $x = y = -\dfrac{1}{2\lambda}$。代入约束方程 $x^2+y^2=1$，得 $2x^2=1$，故 $x=y=\pm \dfrac{\sqrt{2}}{2}$。
    \item \textbf{比较函数值：}
    \(\displaystyle f\left(\frac{\sqrt{2}}{2}, \frac{\sqrt{2}}{2}\right)=\sqrt{2},\qquad f\left(-\frac{\sqrt{2}}{2}, -\frac{\sqrt{2}}{2}\right)=-\sqrt{2}\)。
    因此在约束 $x^2+y^2=1$ 下，\(\displaystyle \boxed{f_{\max}=\sqrt{2},\ f_{\min}=-\sqrt{2}}\)。
  \end{enumerate}

  \item[(2)] \textbf{求 $f(x,y)=\mathrm{e}^{xy}$ 的极值：}
  \begin{enumerate}
    \item \textbf{构造辅助函数：}
    \(\displaystyle L(x,y;\lambda)=\mathrm{e}^{xy}+\lambda(x^2+2y^2-1)\)。
    \item \textbf{求偏导并令为零：}
    由 $L_x = y\mathrm{e}^{xy} + 2\lambda x = 0$ 与 $L_y = x\mathrm{e}^{xy} + 4\lambda y = 0$，若 $x=0$ 则必有 $y=0$（否则矛盾），但 $(0,0)$ 不满足约束，故 $x \neq 0,\ y \neq 0$。
    \item \textbf{求解极值点：}
两式移项后相除，消去 $\lambda$ 与 $\mathrm{e}^{xy}$，得 \(\displaystyle \dfrac{y}{x} = \dfrac{-2\lambda x}{-4\lambda y} = \dfrac{x}{2y}\)，
    即 $x^2 = 2y^2$。代入约束 $x^2+2y^2=1$，得 $4y^2=1$，故 $y = \pm \dfrac{1}{2}$，对应 $x = \pm \dfrac{\sqrt{2}}{2}$，交叉配对共 $4$ 个驻点。
    \item \textbf{比较函数值：}
    由于指数函数单调递增，只需比较 $xy$ 的大小。当 $x,y$ 同号时，$xy=\dfrac{\sqrt{2}}{4}$，故 \(\displaystyle f_{\max}=\mathrm{e}^{\sqrt{2}/4}\)，
    在 $\left(\dfrac{\sqrt{2}}{2}, \dfrac{1}{2}\right)$ 与 $\left(-\dfrac{\sqrt{2}}{2}, -\dfrac{1}{2}\right)$ 处取得；当 $x,y$ 异号时，$xy=-\dfrac{\sqrt{2}}{4}$，故
    \(\displaystyle f_{\min}=\mathrm{e}^{-\sqrt{2}/4}\)，
    在 $\left(-\dfrac{\sqrt{2}}{2}, \dfrac{1}{2}\right)$ 与 $\left(\dfrac{\sqrt{2}}{2}, -\dfrac{1}{2}\right)$ 处取得。
  \end{enumerate}

  \item[(3)] \textbf{求原点到直线的距离：}
  \begin{enumerate}
    \item \textbf{建立数学模型：}
    将几何目标转为最小化距离平方 $D^2=x^2+y^2$，构造辅助函数
    \(\displaystyle L(x,y;\lambda)=x^2+y^2+\lambda\left(\frac{x}{a}+\frac{y}{b}-1\right)\)。
    由 $L_x = 2x + \dfrac{\lambda}{a} = 0$ 与 $L_y = 2y + \dfrac{\lambda}{b} = 0$ 得 $x = -\dfrac{\lambda}{2a}$，$y = -\dfrac{\lambda}{2b}$；再代入直线约束方程，得 \(\displaystyle -\frac{\lambda}{2a^2}-\frac{\lambda}{2b^2}=1 \implies \lambda=\frac{-2}{\dfrac{1}{a^2}+\dfrac{1}{b^2}}\)。
    \item \textbf{计算最短距离：}
    \(\displaystyle D^2=x^2+y^2=\frac{\lambda^2}{4a^2}+\frac{\lambda^2}{4b^2}=\frac{\lambda^2}{4}\left(\frac{1}{a^2}+\frac{1}{b^2}\right)\)。
    代入 $\lambda$ 得 \(\displaystyle D^2=\frac{4}{\left(\dfrac{1}{a^2}+\dfrac{1}{b^2}\right)^2}\cdot\frac{\dfrac{1}{a^2}+\dfrac{1}{b^2}}{4}\)，即 \(\displaystyle D^2=\frac{1}{\dfrac{1}{a^2}+\dfrac{1}{b^2}}=\frac{a^2b^2}{a^2+b^2}\)，
    开根号得距离 $D = \dfrac{|ab|}{\sqrt{a^2+b^2}}$，与解析几何距离公式完全一致。
  \end{enumerate}
\end{enumerate}
\end{solution}

\begin{exercise}[基础：三维单约束问题]
\begin{shortexenum}
  \shortitem{1}{求椭球面 $\dfrac{x^2}{a^2}+\dfrac{y^2}{b^2}+\dfrac{z^2}{c^2}=1$ 内接长方体的最大体积；} &
  \shortitem{2}{求原点到平面 $Ax+By+Cz+D=0$ 的最短距离；} \\
  \shortitem{3}{在抛物面 $z=x^2+y^2$ 上求距平面 $x+y+z+1=0$ 最近的点。} &
\end{shortexenum}
\end{exercise}
\begin{solution}
\begin{enumerate}
  \item[(1)] \textbf{椭球面内接长方体体积：}
  \begin{enumerate}
    \item \textbf{建立模型：}
    设长方体位于第一卦限的顶点为 $(x,y,z)$ 且 $x,y,z>0$，则总体积为 $V = 8xyz$。
    \item \textbf{利用均值不等式法求解：}
    由约束条件应用 AM-GM 不等式：\(\displaystyle 1=\frac{x^2}{a^2}+\frac{y^2}{b^2}+\frac{z^2}{c^2}\ge 3\sqrt[3]{\frac{x^2y^2z^2}{a^2b^2c^2}}\)，因此 \(\displaystyle xyz\leqslant \frac{abc}{3\sqrt3}\)。
    \item \textbf{整理结论：}
    因此 \(\displaystyle V_{\max}=8\left(\frac{abc}{3\sqrt3}\right)=\frac{8abc}{3\sqrt3}\)。
    当且仅当 \(\displaystyle \frac{x^2}{a^2}=\frac{y^2}{b^2}=\frac{z^2}{c^2}=\frac13\) 时取等号，即 \(\displaystyle x=\frac{a}{\sqrt3},\ y=\frac{b}{\sqrt3},\ z=\frac{c}{\sqrt3}\)。
  \end{enumerate}

  \item[(2)] \textbf{点到平面距离：}
  \begin{enumerate}
    \item \textbf{构造辅助函数：}
    最小化 $f(x,y,z) = x^2+y^2+z^2$，约束为 $g = Ax+By+Cz+D=0$。构造 \(\displaystyle L(x,y,z;\lambda) = x^2+y^2+z^2 + \lambda(Ax+By+Cz+D)\)。
    \item \textbf{求偏导并代入约束：}
    \begin{itemize}
      \item $L_x = 2x+\lambda A=0 \implies x=-\dfrac{\lambda A}{2}$
      \item $L_y = 2y+\lambda B=0 \implies y=-\dfrac{\lambda B}{2}$
      \item $L_z = 2z+\lambda C=0 \implies z=-\dfrac{\lambda C}{2}$
    \end{itemize}
    代入平面方程：\(\displaystyle -\frac{\lambda}{2}(A^2+B^2+C^2) + D = 0 \implies \lambda = \frac{2D}{A^2+B^2+C^2}\)。
    \item \textbf{计算最短距离：}
    最小距离的平方为 \(\displaystyle x^2+y^2+z^2 = \frac{\lambda^2}{4}(A^2+B^2+C^2) = \frac{D^2}{A^2+B^2+C^2}\)。
    开根号即得最短距离为 $d = \dfrac{|D|}{\sqrt{A^2+B^2+C^2}}$。
  \end{enumerate}

  \item[(3)] \textbf{曲面到平面的距离最近点：}
  \begin{enumerate}
    \item \textbf{利用法向量平行性：}
    当距离极小化时，曲面在最近点处的切平面必平行于目标平面。
    \begin{itemize}
      \item 目标平面法向量为 $\vec{n} = (1, 1, 1)$。
      \item 抛物面隐式方程为 $F(x,y,z) = x^2+y^2-z=0$，其法向量为 $\nabla F = (2x, 2y, -1)$。
    \end{itemize}
    \item \textbf{建立比例关系：}
    两者必须共线，故存在非零常数 $k$ 使 \(\displaystyle 2x=k,\ 2y=k,\ -1=k\)。由此得 $k=-1$，所以 \(\displaystyle x=y=-\frac{1}{2}\)。
    \item \textbf{求出对应的 $z$ 坐标：}
    代回曲面方程求得 \(\displaystyle z=\left(-\frac{1}{2}\right)^2+\left(-\frac{1}{2}\right)^2=\frac{1}{2}\)，所以最近的点坐标为 $\left(-\dfrac{1}{2}, -\dfrac{1}{2}, \dfrac{1}{2}\right)$。
  \end{enumerate}
\end{enumerate}
\end{solution}

\begin{exercise}[提高：双约束 Lagrange 乘数法]
\begin{shortsingleenum}
  \shortsingleitem{1}{求 $f(x,y,z)=x+y+z$ 在约束 $\begin{cases}x^2+y^2=1\\x+z=1\end{cases}$ 下的极值；}
  \shortsingleitem{2}{求函数 $u=xyz$ 在约束 $\begin{cases}x+y+z=5\\xy+yz+zx=8\end{cases}$ 下的最大值和最小值；}
  \shortsingleitem{3}{证明椭圆 $\dfrac{x^2}{a^2}+\dfrac{y^2}{b^2}=1$ 的内接三角形中，面积最大的三角形是单位圆内接正三角形在变换 $x=ax',\ y=by'$ 下的像。}
\end{shortsingleenum}
\end{exercise}
\begin{solution}
\begin{enumerate}
  \item[(1)] \textbf{双约束的线性函数极值：}
  \begin{enumerate}
    \item \textbf{构造辅助函数：}\(\displaystyle L = x+y+z + \lambda(x^2+y^2-1) + \mu(x+z-1)\)。
    \item \textbf{求偏导并令为零：}
    由 $L_x = 1+2\lambda x + \mu = 0$，$L_y = 1+2\lambda y = 0 \implies y = -\dfrac{1}{2\lambda}$，$L_z = 1+\mu = 0 \implies \mu = -1$。
    \item \textbf{求解极值点：}
    将 $\mu=-1$ 代回第一式，得 \(\displaystyle 1+2\lambda x - 1 = 0 \implies 2\lambda x = 0\)。
    因为由 $L_y=0$ 知 $\lambda \neq 0$，故必定有 $x=0$。
    代入第一约束 $x^2+y^2=1 \implies y^2=1 \implies y=\pm 1$。
    代入第二约束 $x+z=1 \implies z=1$。
    \item \textbf{比较函数值：}
    两个候选点分别为 $(0,1,1)$ 与 $(0,-1,1)$，对应 \(\displaystyle f(0,1,1)=2,\ f(0,-1,1)=0\)。
    因此最大值为 $2$，最小值为 $0$。
  \end{enumerate}

  \item[(2)] \textbf{对称约束的极值问题：}
  \begin{enumerate}
    \item \textbf{构造辅助函数：}\(\displaystyle L = xyz + \lambda(x+y+z-5) + \mu(xy+yz+zx-8)\)。
    \item \textbf{求偏导并令为 0：}
    由 $L_x = yz + \lambda + \mu(y+z) = 0$，$L_y = xz + \lambda + \mu(x+z) = 0$，$L_z = xy + \lambda + \mu(x+y) = 0$。
    \item \textbf{利用对称性解方程：}
    第一式与第二式相减得 \(\displaystyle z(y-x) + \mu(y-x) = 0 \implies (z+\mu)(y-x) = 0\)。
    根据问题的彻底轮换对称性，不失一般性假定必有两个自变量相等。假设 $x=y$。
    代入第一约束得 $2x+z=5 \implies z=5-2x$。
    代入第二约束得 \(\displaystyle x^2 + 2x(5-2x) = 8 \implies x^2 + 10x - 4x^2 = 8 \implies 3x^2 - 10x + 8 = 0\)。
    因式分解 $(3x-4)(x-2)=0$，解得 $x=2$ 或 $x=4/3$。
    \item \textbf{计算极值：}
    当 $x=y=2$ 时，$z=1$，乘积 $u = xyz = 4$；当 $x=y=\dfrac{4}{3}$ 时，$z=5-\dfrac{8}{3}=\dfrac{7}{3}$，乘积 $u = xyz = \left(\dfrac{4}{3}\right)\left(\dfrac{4}{3}\right)\left(\dfrac{7}{3}\right) = \dfrac{112}{27} \approx 4.148$。比较得\textbf{最大值为 $\dfrac{112}{27}$，最小值为 $4$}。
  \end{enumerate}

  \item[(3)] \textbf{内接三角形面积最值分析：}
  \begin{enumerate}
    \item \textbf{仿射变换与结论：}
    通过坐标缩放变换 $x'=\dfrac{x}{a},\ y'=\dfrac{y}{b}$，椭圆变为单位圆 $x'^2+y'^2=1$。任意内接三角形面积的放缩比例是定值（因为雅可比行列式恒为 $ab$），故面积大小关系保持不变。又已知单位圆内面积最大者为内接正三角形，所以椭圆内面积最大的内接三角形，正是单位圆内接正三角形在变换 $(x',y')\mapsto (ax',by')$ 下的像。
  \end{enumerate}
\end{enumerate}
\end{solution}

\begin{exercise}[KKT 条件与不等式约束预告]
到目前为止，我们的 Lagrange 乘数法只处理了 \textbf{等式约束} $g=0$。但在实际问题中，更多出现的是 \textbf{不等式约束}（如 $x\geqslant 0$, $g(x)\leqslant 0$ 等）。这类问题由 \textbf{KKT 条件}（Karush--Kuhn--Tucker 条件）处理。
\begin{enumerate}
  \item[(1)] 考虑不等式约束优化问题：$\max f(x,y)$ s.t. $g(x,y)\leqslant 0$。引入“松弛变量”$s\geqslant 0$ 将 $g\leqslant 0$ 化为等式 $g+s^2=0$，再对此应用 Lagrange 乘数法，证明在最优点处满足：
        \[
        \begin{cases}
          \nabla f+\lambda\nabla g=\boldsymbol 0,\\
          \lambda\leqslant 0,\quad g\leqslant 0,\quad \lambda g=0.
        \end{cases}
        \]
        最后一个条件 $\lambda g=0$ 称为 \textbf{互补松弛性}；它的含义是：要么约束“绷紧”（$g=0$ 且 $\lambda$ 可非零），要么约束“松弛”（$g<0$ 且 $\lambda=0$，该约束不起作用）。
  \item[(2)] 求 $\min(x^2+y^2)$ s.t. $x+y\geqslant 1$。
  \item[(3)] 思考：为什么 KKT 条件被称为“Lagrange 乘数法的非线性推广”？它与微积分中的“区间端点检查”有何类比？
\end{enumerate}
\end{exercise}
\begin{solution}
\begin{enumerate}
  \item[(1)] \textbf{互补松弛性的严格推导：}
  \begin{enumerate}
    \item 将不等式改为等式约束，构造辅助函数 \(\displaystyle L(x,y,s;\lambda) = f(x,y) + \lambda(g(x,y) + s^2)\)，
    并由 $L_x = f_x + \lambda g_x = 0$，$L_y = f_y + \lambda g_y = 0$，$L_s = 2\lambda s = 0$ 建立驻点条件。
    \item 从 $L_s=0$ 可以得出，要么 $\lambda=0$，要么 $s=0$。由于 $s^2 = -g$，当 $s=0$ 时必定 $g=0$。所以无论哪种情况总有 $\lambda g = 0$ 成立，这证明了互补松弛性。
    \item 至于 $\lambda \leqslant  0$：由于我们求最大化，在边界处目标函数梯度应当指引方向向外（趋于禁止区域），而约束梯度的正方向亦朝外，两者反向抵抗，故需采用特定符号约束（非正）。
  \end{enumerate}

  \item[(2)] \textbf{通过 KKT 求解简单二次规划：}
  \begin{enumerate}
    \item 转化为标准形式：$\min x^2+y^2$ s.t. $1-x-y \leqslant  0$，构造拉格朗日函数 \(\displaystyle L = x^2+y^2 + \lambda(1-x-y)\)。
    由 $L_x = 2x - \lambda = 0$ 与 $L_y = 2y - \lambda = 0$ 得 $x=y=\dfrac{\lambda}{2}$。若约束松弛，则 $\lambda=0$，得到 $(0,0)$，但它不满足 $x+y\ge 1$；故约束必须起作用，即 $g=0 \implies x+y=1$，从而 $\lambda=1$，解得 $x=\dfrac{1}{2},\ y=\dfrac{1}{2}$。
    \item 目标函数的极小值为 $\left(\dfrac{1}{2}\right)^2+\left(\dfrac{1}{2}\right)^2 = \dfrac{1}{2}$。
  \end{enumerate}

  \item[(3)] \textbf{与端点检查的类比启示：}
  KKT 条件包含了一种智能的穷举选择逻辑。它实际上在问：极值点到底是在可行域的“内部”（此时完全没有碰到边界约束，这等价于传统求极值的驻点法则，乘子失效为 0），还是正好撞在“边界”上（此时转化为传统等式约束求极值，乘子发挥效力）。这与高中一元函数闭区间求最值必须比较内部导数为 0 处及两侧端点处值的做法有着深刻的逻辑同构性。
\end{enumerate}
\end{solution}

\subsubsection{A组}

\begin{exercise}[概念问答]
回答下列问题：
\begin{shortexenum}
    \shortitem{1}{目标函数 $z=f(x,y)$ 在条件 $g(x,y)=0$ 下取得极值的必要条件是什么？} &
    \shortitem{2}{什么是拉格朗日乘数法？}
\end{shortexenum}
\end{exercise}
\begin{solution}
\begin{enumerate}
    \item[(1)] \textbf{条件极值的必要条件：}
    目标函数 $z=f(x,y)$ 在条件 $g(x,y)=0$ 下取得极值时，若约束曲线在该点光滑，即 $\nabla g(x_0,y_0)\neq \boldsymbol 0$，则存在常数 $\lambda$ 使 \(\displaystyle \nabla f(x_0,y_0)+\lambda\nabla g(x_0,y_0)=\boldsymbol 0\)。
    几何上，这表示极值点处目标函数的等值线与约束曲线相切，因此它们的法向量（梯度）共线。

    \item[(2)] \textbf{拉格朗日乘数法定义：}
    拉格朗日乘数法是一种把“带约束的极值问题”改造成“无约束驻点问题”的方法。
    它通过引入乘子 $\lambda$，构造辅助函数 \(\displaystyle L(x,y;\lambda)=f(x,y)+\lambda g(x,y),\qquad L_x=0,\qquad L_y=0,\qquad g(x,y)=0\) 来求候选点。
\end{enumerate}
\end{solution}

\begin{exercise}[拉格朗日乘数法求条件极值]
求下列函数在给定的约束条件下的最大值、最小值：
\begin{shortexenum}
    \shortitem{1}{$f(x,y)=x^2+y,\ x^2-y^2=1$} &
    \shortitem{2}{$f(x,y)=xy,\ x+y=1$} \\
    \shortitem{3}{$f(x,y,z)=x+3y+5z,\ x^2+y^2+z^2=1$} &
    \shortitem{4}{$f(x,y,z)=2x+y+4z,\ x^2+y^2+z^2=16$}
\end{shortexenum}
\end{exercise}
\begin{solution}
\begin{enumerate}
    \item[(1)] \textbf{求解 $f(x,y)=x^2+y,\ x^2-y^2=1$：}
    \begin{enumerate}
        \item \textbf{构造拉格朗日函数：}
        \(\displaystyle L(x,y;\lambda) = x^2+y + \lambda(x^2-y^2-1)\)。
        \item \textbf{求偏导：}
        \begin{itemize}
            \item $L_x = 2x + 2\lambda x = 0 \implies x(1+\lambda)=0 \implies x=0$ 或 $\lambda=-1$。
            \item 若 $x=0$，代入约束得 $-y^2=1$，无实解。
            \item 若 $\lambda=-1$，$L_y = 1 - 2\lambda y = 1 + 2y = 0 \implies y=-\dfrac{1}{2}$。
        \end{itemize}
        \item \textbf{求解坐标：}
        代入约束得 $x^2 - \dfrac{1}{4} = 1 \implies x^2 = \dfrac{5}{4} \implies x = \pm \dfrac{\sqrt{5}}{2}$。
        \item \textbf{判断最值：}
        在点 $\left(\pm \dfrac{\sqrt{5}}{2}, -\dfrac{1}{2}\right)$ 处，函数值为 $f = \dfrac{5}{4} - \dfrac{1}{2} = \dfrac{3}{4}$。
        由于约束是一条无界的双曲线，且 $f(x,y) = 1+y^2+y = \left(y+\dfrac{1}{2}\right)^2 + \dfrac{3}{4} \ge \dfrac{3}{4}$，当 $y \to \infty$ 时 $f \to \infty$。
        因此，\textbf{最小值为 $\dfrac{3}{4}$}，无最大值。
    \end{enumerate}

    \item[(2)] \textbf{求解 $f(x,y)=xy,\ x+y=1$：}
    \begin{enumerate}
        \item \textbf{消元法：}
        约束可得 $y = 1-x$。
        \item \textbf{构造一元函数：}
        目标函数转化为 \(\displaystyle f(x) = x(1-x) = -x^2+x = -\left(x-\frac{1}{2}\right)^2 + \frac{1}{4}\)。
        \item \textbf{结论：}
        该函数是开口向下的抛物线，因此在约束 $x+y=1$ 下，\(\displaystyle \boxed{f_{\max}=\frac14}\)
        在 $x=y=\dfrac12$ 处取得；同时当 $|x|\to\infty$ 时，$f(x,1-x)\to-\infty$，故无最小值。
    \end{enumerate}

    \item[(3)] \textbf{求解 $f(x,y,z)=x+3y+5z,\ x^2+y^2+z^2=1$：}
    \begin{enumerate}
        \item \textbf{构造拉格朗日函数：}\(\displaystyle L = x+3y+5z + \lambda(x^2+y^2+z^2-1)\)。
        \item \textbf{求偏导：}
        \begin{itemize}
            \item $L_x = 1+2\lambda x = 0 \implies x = -\dfrac{1}{2\lambda}$
            \item $L_y = 3+2\lambda y = 0 \implies y = -\dfrac{3}{2\lambda}$
            \item $L_z = 5+2\lambda z = 0 \implies z = -\dfrac{5}{2\lambda}$
        \end{itemize}
        \item \textbf{求解乘子与坐标：}
        代入约束平方和公式：\(\displaystyle \left(-\frac{1}{2\lambda}\right)^2 (1+9+25) = 1 \implies \frac{35}{4\lambda^2} = 1 \implies \lambda = \pm \frac{\sqrt{35}}{2}\)。
        解出的驻点群为 $\pm\left(\dfrac{1}{\sqrt{35}}, \dfrac{3}{\sqrt{35}}, \dfrac{5}{\sqrt{35}}\right)$。
        \item \textbf{比较函数值：}
        代入目标函数可得 \(\displaystyle f=\pm\sqrt{35}\)。
        因此最大值为 $\sqrt{35}$，最小值为 $-\sqrt{35}$。
    \end{enumerate}

    \item[(4)] \textbf{求解 $f(x,y,z)=2x+y+4z,\ x^2+y^2+z^2=16$：}
    \begin{enumerate}
        \item \textbf{应用拉格朗日法解题：}\(\displaystyle L = 2x+y+4z + \lambda(x^2+y^2+z^2-16)\)。
        同理求出偏导对应关系：
        \begin{itemize}
            \item $x = -\dfrac{1}{\lambda}$
            \item $y = -\dfrac{1}{2\lambda}$
            \item $z = -\dfrac{2}{\lambda}$
        \end{itemize}
        \item \textbf{代入约束求解：}
        代入约束方程：\(\displaystyle \frac{1}{\lambda^2} + \frac{1}{4\lambda^2} + \frac{4}{\lambda^2} = 16 \implies \frac{21}{4\lambda^2} = 16 \implies \lambda^2 = \frac{21}{64} \implies \lambda = \pm \frac{\sqrt{21}}{8}\)。
        \item \textbf{比较函数值：}
        将驻点代入目标函数，得到 \(\displaystyle
        f=\pm\sqrt{2^2+1^2+4^2}\cdot\sqrt{16}
        =\pm4\sqrt{21}\)。
        因此最大值为 $4\sqrt{21}$，最小值为 $-4\sqrt{21}$。
    \end{enumerate}
\end{enumerate}
\end{solution}

\begin{exercise}[几何约束最优化一]
从斜边之长为 $l$ 的一切直角三角形中，求有最大周界的直角三角形。
\end{exercise}
\begin{solution}
\begin{enumerate}
    \item 设直角三角形的两直角边分别为 $x, y$。目标是最大化周长 $C(x,y) = x+y+l$，约束条件为 $x^2+y^2=l^2$ 且 $x>0, y>0$。
    \item \textbf{构造拉格朗日函数：}\(\displaystyle L(x,y;\lambda) = x+y+l + \lambda(x^2+y^2-l^2)\)。
    \item \textbf{求驻点：}
    分别求偏导
    \begin{itemize}
        \item $L_x = 1+2\lambda x = 0 \implies x = -\dfrac{1}{2\lambda}$
        \item $L_y = 1+2\lambda y = 0 \implies y = -\dfrac{1}{2\lambda}$
    \end{itemize}
    通过比对可知必须要求 $x=y$。代入约束方程得 $2x^2=l^2 \implies x = y = \dfrac{l}{\sqrt{2}}$。

    \item \textbf{结论：}
    当 $x=y=\dfrac{l}{\sqrt2}$ 时，直角三角形为等腰直角三角形，此时周长达到最大值 \(\displaystyle C_{\max}=x+y+l=\sqrt2\,l+l=(\sqrt2+1)l\)。
\end{enumerate}
\end{solution}

\begin{exercise}[几何约束最优化二]
求内接于半径为 $a$ 的球且有最大体积的长方体。
\end{exercise}
\begin{solution}
\begin{enumerate}
    \item 设长方体的长、宽、高分别为 $2x, 2y, 2z$ ($x,y,z>0$)。体积 $V = 8xyz$。长方体对角线等于球直径，故在坐标原点位于几何中心的布局下有约束 \(\displaystyle 4x^2+4y^2+4z^2=4a^2\)，即 \(\displaystyle x^2+y^2+z^2=a^2\)。

    \item 根据 AM-GM 不等式，\(\displaystyle x^2+y^2+z^2\ge 3\sqrt[3]{x^2y^2z^2}\)。代入约束得 \(\displaystyle a^2\ge 3(xyz)^{2/3}\)，因此 \(\displaystyle xyz\leqslant \left(\frac{a^2}{3}\right)^{3/2}=\frac{a^3}{3\sqrt3}\)。
    等号成立当且仅当 $x^2=y^2=z^2$，即 \(\displaystyle x=y=z=\frac{a}{\sqrt3}\)。
    因而体积最大的内接长方体是正方体，其边长为 $\dfrac{2a}{\sqrt3}$，最大体积为 \(\displaystyle \boxed{V_{\max}=\frac{8a^3}{3\sqrt3}}\)。
\end{enumerate}
\end{solution}
\begin{exercise}[距离最小化]
在椭圆 $x^2+4y^2=4$ 上求一点，使其到直线 $2x+3y-6=0$ 的距离最短。
\end{exercise}
\begin{solution}
\begin{enumerate}
    \item 设椭圆上任意一点的坐标为 $(x,y)$，其到直线 $2x+3y-6=0$ 的距离是 \(\displaystyle \mathrm{d}(x,y)=\frac{|2x+3y-6|}{\sqrt{13}}\)。
    因而最小距离问题等价于：在约束 \(\displaystyle \varphi(x,y)=x^2+4y^2-4=0\) 下，使 \(\displaystyle f(x,y)=2x+3y\) 尽量接近 $6$。

    \item 引入拉格朗日乘数 $\lambda$，构造辅助函数 \(\displaystyle L(x,y,\lambda)=f(x,y)+\lambda\varphi(x,y)=2x+3y+\lambda(x^2+4y^2-4)\)。

    \item 分别对 $x, y, \lambda$ 求偏导数并令其等于零，得到拉格朗日方程组：
    \[
    \begin{cases}
    L_x = 2 + 2\lambda x = 0 \quad &(1)\\
    L_y = 3 + 8\lambda y = 0 \quad &(2)\\
    L_\lambda = x^2 + 4y^2 - 4 = 0 \quad &(3)
    \end{cases}
    \]

    \item \textbf{求解非线性方程组：}
    \begin{itemize}
        \item 显然 $\lambda \neq 0$（否则方程 $(1)$ 变为 $2=0$，矛盾）。由方程 $(1)$ 和 $(2)$ 解出 $x$ 和 $y$ 关于 $\lambda$ 的表达式：
        \(\displaystyle x = -\frac{1}{\lambda}, \quad y = -\frac{3}{8\lambda}\)
        \item 将这两个表达式代入约束方程 $(3)$ 中：
        \(\displaystyle \left(-\frac{1}{\lambda}\right)^2 + 4\left(-\frac{3}{8\lambda}\right)^2 - 4 = 0\)，即
        \(\displaystyle \frac{1}{\lambda^2} + 4 \cdot \frac{9}{64\lambda^2} = 4 \implies \frac{1}{\lambda^2} + \frac{9}{16\lambda^2} = 4\)。
        \item 提取公因式并通分：
        \(\displaystyle \frac{16 + 9}{16\lambda^2} = 4 \implies \frac{25}{16\lambda^2} = 4 \implies 64\lambda^2 = 25 \implies \lambda^2 = \frac{25}{64} \implies \lambda = \pm \frac{5}{8}\)。
        \item \textbf{情况一：} 当 $\lambda = \dfrac{5}{8}$ 时，代入得 $x = -\dfrac{8}{5}, y = -\dfrac{3}{5}$。得到候选点 $P_1\left(-\dfrac{8}{5}, -\dfrac{3}{5}\right)$。
        \item \textbf{情况二：} 当 $\lambda = -\dfrac{5}{8}$ 时，代入得 $x = \dfrac{8}{5}, y = \dfrac{3}{5}$。得到候选点 $P_2\left(\dfrac{8}{5}, \dfrac{3}{5}\right)$。
    \end{itemize}

    \item \textbf{计算距离并比较判定最值：}
    \begin{itemize}
        \item 在点 $P_1\left(-\dfrac{8}{5}, -\dfrac{3}{5}\right)$ 处，距离为：
        \(\displaystyle d_1 = \frac{\left| 2\left(-\dfrac{8}{5}\right) + 3\left(-\dfrac{3}{5}\right) - 6 \right|}{\sqrt{13}} = \frac{\left| -\dfrac{16}{5} - \dfrac{9}{5} - \dfrac{30}{5} \right|}{\sqrt{13}} = \frac{11}{\sqrt{13}}\)。
        \item 在点 $P_2\left(\dfrac{8}{5}, \dfrac{3}{5}\right)$ 处，距离为：
        \(\displaystyle d_2 = \frac{\left| 2\left(\dfrac{8}{5}\right) + 3\left(\dfrac{3}{5}\right) - 6 \right|}{\sqrt{13}} = \frac{\left| \dfrac{16}{5} + \dfrac{9}{5} - \dfrac{30}{5} \right|}{\sqrt{13}} = \frac{1}{\sqrt{13}}\)。
    \end{itemize}

    \item \textbf{结论：}
     由于 $d_2<d_1$，故所求点为 \(\displaystyle \boxed{\left(\frac85,\frac35\right)},\qquad \text{最短距离为 }\boxed{\frac{1}{\sqrt{13}}}\)。
\end{enumerate}
\end{solution}

\begin{exercise}[成本极值]
欲造一无盖的长方体容器，已知底部造价为 3 元/$\mathrm{m}^2$，侧面造价均为 1 元/$\mathrm{m}^2$。现想用 36 元造一个容积最大的容器，求它的尺寸。
\end{exercise}
\begin{solution}
\begin{enumerate}
    \item \textbf{建立目标函数与约束条件：}
    设该长方体容器的底面长为 $x$，宽为 $y$，高为 $z$（其中 $x>0, y>0, z>0$）。
    容器的容积作为我们需要最大化的目标函数：\(\displaystyle V(x,y,z) = xyz\)。
    成本方面：底面面积为 $xy$，单价 3 元；前后面面积共 $2xz$，单价 1 元；左右面面积共 $2yz$，单价 1 元。
    总造价的等式约束条件为：\(\displaystyle \varphi(x,y,z) = 3xy + 2xz + 2yz - 36 = 0\)。

    \item \textbf{构造拉格朗日函数：}
    引入拉格朗日乘数 $\lambda$，构造辅助函数
    \(\displaystyle L(x,y,z,\lambda)=xyz+\lambda(3xy+2xz+2yz-36)\)。

    \item \textbf{求解偏导数并建立方程组：}
    分别对 $x, y, z$ 求偏导数并令其等于零，得到：
    \[
    \begin{cases}
    L_x = yz + \lambda(3y + 2z) = 0 \quad &(1)\\
    L_y = xz + \lambda(3x + 2z) = 0 \quad &(2)\\
    L_z = xy + \lambda(2x + 2y) = 0 \quad &(3)\\
    L_\lambda = 3xy + 2xz + 2yz - 36 = 0 \quad &(4)
    \end{cases}
    \]

    \item \textbf{使用对称化方法求解方程组：}
    因为容器有体积，必然有 $x \neq 0, y \neq 0, z \neq 0$。我们将方程 $(1), (2), (3)$ 分别同乘 $x, y, z$，进行移项：
    \begin{itemize}
        \item 方程 $(1) \times x \implies xyz + \lambda(3xy + 2xz) = 0 \implies -xyz = \lambda(3xy + 2xz)$
        \item 方程 $(2) \times y \implies xyz + \lambda(3xy + 2yz) = 0 \implies -xyz = \lambda(3xy + 2yz)$
        \item 方程 $(3) \times z \implies xyz + \lambda(2xz + 2yz) = 0 \implies -xyz = \lambda(2xz + 2yz)$
    \end{itemize}
    若 $\lambda = 0$，则 $xyz = 0$，容积为零，不符合最大化需求，故 $\lambda \neq 0$。
    等式左边完全相同，从而我们可以令右边括号内的各项相等：
    \begin{itemize}
        \item 由 $\lambda(3xy + 2xz) = \lambda(3xy + 2yz)$，消去 $\lambda$ 和公共项 $3xy$，得：
        \(\displaystyle 2xz = 2yz \implies x = y \quad (\text{因 } z \neq 0)\)。
        这说明为实现造价最优化，底面必须是正方形。

        \item 由 $\lambda(3xy + 2yz) = \lambda(2xz + 2yz)$，消去 $\lambda$ 和公共项 $2yz$，得：
        \(\displaystyle 3xy = 2xz \implies 3y = 2z \implies z = \frac{3}{2}y \quad (\text{因 } x \neq 0)\)。
    \end{itemize}

    \item \textbf{代回约束方程求绝对尺寸：}
    将 $x = y$ 和 $z = \dfrac{3}{2}y$ 统一代回预算约束方程 $(4)$，得
    \(\displaystyle 3y^2+2y\left(\frac{3}{2}y\right)+2y\left(\frac{3}{2}y\right)=36
    \Longrightarrow 9y^2=36\Longrightarrow y^2=4\)。
    由于边长必须为正数，解得 $y = 2$。
    顺势反推出其他变量：
    \(\displaystyle x = y = 2,\quad z = \frac{3}{2}(2) = 3\)。

    \item \textbf{结论：}
    由求得的唯一正驻点可知 \(\displaystyle \boxed{x=2,\ y=2,\ z=3}\)，且 \(\displaystyle \boxed{V_{\max}=xyz=12\ \mathrm{m}^3}\)。
\end{enumerate}
\end{solution}

\subsubsection{B组}

\begin{exercise}[闭区域最值]
求二元函数 $u=f(x,y)=x^2y(4-x-y)$ 在由直线 $x+y=6$、$x$ 轴和 $y$ 轴所围成的闭区域 $D$ 上的极值、最大值和最小值。
\end{exercise}
\begin{solution}
\begin{enumerate}
    \item \textbf{分析区域及边界特征值遍历：}
    闭区域 $D$ 是由点 $(0,0), (6,0), (0,6)$ 连线围成的封闭实心三角形及其内部面域。
    \begin{itemize}
        \item 在 $x$ 轴（$y=0$）上，函数恒为 $u=0$。
        \item 在 $y$ 轴（$x=0$）上，同样有 $u=0$。
        \item 在 $x+y=6$ 直线上，有置换关系 $4-x-y = -2$，代入得 $u = -2x^2y = -2x^2(6-x) = 2x^3 - 12x^2$。
              对于 $x \in [0,6]$ 内求其最值，执行一元导数测试：$6x^2-24x = 0 \implies x=0$ 或 $x=4$。
              在 $x=4$ 处关联求得 $y=2$，代入算得边界局部极小波谷为 $-64$。各个交角端点处回落均为 0。
    \end{itemize}

    \item \textbf{求区域内部的驻点：}
    展开形态为 $u(x,y) = 4x^2y - x^3y - x^2y^2$。
    分别求偏导：
    \begin{itemize}
        \item $u_x = 8xy - 3x^2y - 2xy^2 = xy(8 - 3x - 2y) = 0$
        \item $u_y = 4x^2 - x^3 - 2x^2y = x^2(4 - x - 2y) = 0$
    \end{itemize}
    在区域内部有 $x>0,\ y>0$，故只能令括号内的因子为零：\(\displaystyle 3x+2y=8,\ x+2y=4\)。解此线性方程组得 $x=2,\ y=1$。因而区域内部唯一的驻点为 $(2,1)$，且 \(\displaystyle u(2,1)=2^2\cdot1\cdot(4-2-1)=4\)。

    \item \textbf{比较各候选点的函数值：}
    区域内部驻点给出函数值 $4$；两条坐标轴边界上函数值恒为 $0$；在边界 $x+y=6$ 上最小值为 $-64$。
    因此在闭区域 $D$ 上，\(\displaystyle \boxed{u_{\max}=4}\)（在 $(2,1)$ 处取得），\(\displaystyle \boxed{u_{\min}=-64}\)（在 $(4,2)$ 处取得）。
\end{enumerate}
\end{solution}

\begin{exercise}[几何极值应用]
证明：周长一定的三角形中以等边三角形的面积最大。
\end{exercise}
\begin{solution}
设该三角形的三边长度分别为 $x,y,z$，有 $2p = x+y+z$（视作常数不变量）。
由 Heron 公式，\(\displaystyle S^2 = p(p-x)(p-y)(p-z)\)。
任意三角形其较短的两边之和大于最长第三边，$p-x>0, p-y>0, p-z>0$ 必定满足。由均值不等式得 \(\displaystyle (p-x)(p-y)(p-z) \leqslant  \left( \frac{(p-x)+(p-y)+(p-z)}{3} \right)^{\!3} = \left( \frac{3p - (x+y+z)}{3} \right)^{\!3} =  \frac{p^3}{27}\)。
    等号条件为 \(\displaystyle p-x = p-y = p-z \implies x = y = z = \frac{2p}{3}\)。
因此在周长约束下，当且仅当构建三角形的三段实体边长相等（即等边三角形）时，其面积能够达最大值。
\end{solution}

\begin{exercise}[切平面构成的体积最值]
在第一卦限内作椭球面 $\dfrac{x^2}{a^2}+\dfrac{y^2}{b^2}+\dfrac{z^2}{c^2}=1$ 的切平面，使切平面与三个坐标面所围成的四面体的体积最小，求切点的坐标。
\end{exercise}
\begin{solution}
假设切点坐标为 $P_0(x_0, y_0, z_0)$，$x_0, y_0, z_0 > 0$。
    该驻点对应的切平面特征方程直接写明可为 \(\displaystyle \frac{xx_0}{a^2} + \frac{yy_0}{b^2} + \frac{zz_0}{c^2} = 1\)。
    截距分别为：
    \begin{itemize}
        \item 横轴截距 $X = \dfrac{a^2}{x_0}$
        \item 纵轴截距 $Y = \dfrac{b^2}{y_0}$
        \item 垂轴截距 $Z = \dfrac{c^2}{z_0}$
    \end{itemize}
    四面体体积公式为 \(\displaystyle V = \frac{1}{6} XYZ = \frac{a^2 b^2 c^2}{6 x_0 y_0 z_0}\)。
    想要使体积 $V$ 最小，就必然要分母 $x_0 y_0 z_0$ 在约束 $\dfrac{x_0^2}{a^2}+\dfrac{y_0^2}{b^2}+\dfrac{z_0^2}{c^2}=1$ 的前提下寻找极值最大的点。
    调用 AM-GM 不等式：
    \begin{align*}
    1&=\frac{x_0^2}{a^2}+\frac{y_0^2}{b^2}+\frac{z_0^2}{c^2}
    \ge 3\sqrt[3]{\frac{x_0^2y_0^2z_0^2}{a^2b^2c^2}}
    \implies x_0^2y_0^2z_0^2\leqslant\frac{a^2b^2c^2}{27},\\
    \frac{x_0^2}{a^2}&=\frac{y_0^2}{b^2}=\frac{z_0^2}{c^2}=\frac{1}{3}
    \implies x_0=\frac{a}{\sqrt3},\quad y_0=\frac{b}{\sqrt3},\quad z_0=\frac{c}{\sqrt3}.
    \end{align*}
所要求点坐标为：$\left( \dfrac{a}{\sqrt{3}}, \dfrac{b}{\sqrt{3}}, \dfrac{c}{\sqrt{3}} \right)$
\end{solution}

% ============================================================
% §7.12 偏导数在偏微分方程中的应用
% ============================================================

\subsection{第十二节习题：偏导数在偏微分方程中的应用}

\begin{exercise}[PDE 验证]
\begin{shortsingleenum}
  \shortsingleitem{1}{证明 $u=\ln(x^2+y^2)$ 满足二维 Laplace 方程 $u_{xx}+u_{yy}=0$（在 $(x,y) \neq (0,0)$ 处）；}
  \shortsingleitem{2}{设 $u=\mathrm{e}^{ax}\sin(by)$（$a,b$ 为常数）且满足热方程 $u_t=k(u_{xx}+u_{yy})$。讨论使该式成立的条件。更一般地，求使 $u=\mathrm{e}^{\alpha x+\beta t}\sin(\gamma x+\delta t)$ 成为 $u_t=u_{xx}$ 解的条件；}
  \shortsingleitem{3}{设 $u=f(x+at)+g(x-at)$，其中 $f,g$ 四次可微。证明 $u$ 满足 $a^4u_{xxxx}-u_{tttt}=0$（四阶波动型方程）。}
\end{shortsingleenum}
\end{exercise}
\begin{solution}
\begin{enumerate}
  \item[(1)] \textbf{解析：}
  \begin{enumerate}
      \item 求一阶偏导数：
      \begin{itemize}
          \item $u_x = \dfrac{2x}{x^2+y^2}$
          \item $u_y = \dfrac{2y}{x^2+y^2}$
      \end{itemize}
      \item 求二阶偏导数，应用除法求导法则：
      \begin{itemize}
          \item $u_{xx} = \dfrac{2(x^2+y^2) - 2x(2x)}{(x^2+y^2)^2} = \dfrac{2y^2-2x^2}{(x^2+y^2)^2}$
          \item $u_{yy} = \dfrac{2(x^2+y^2) - 2y(2y)}{(x^2+y^2)^2} = \dfrac{2x^2-2y^2}{(x^2+y^2)^2}$
      \end{itemize}
      \item 验证方程：\(\displaystyle u_{xx}+u_{yy}=\frac{2y^2-2x^2+2x^2-2y^2}{(x^2+y^2)^2}=0\)。
      证毕。
  \end{enumerate}

  \item[(2)] \textbf{解析：}
  \begin{enumerate}
      \item 对于第一部分，$u=\mathrm{e}^{ax}\sin(by)$ 中不显含自变量 $t$，故 $u_t = 0$。
      分别对 $x,y$ 求二阶导：
      \(\displaystyle u_x=a\mathrm{e}^{ax}\sin(by),\quad u_{xx}=a^2\mathrm{e}^{ax}\sin(by)=a^2u\)，
      \(\displaystyle u_y=b\mathrm{e}^{ax}\cos(by),\quad u_{yy}=-b^2\mathrm{e}^{ax}\sin(by)=-b^2u\)。
      代入方程 $u_t=k(u_{xx}+u_{yy})$，得 \(\displaystyle 0 = k(a^2u - b^2u) = k(a^2-b^2)u\)。
      要使方程对任意 $(x,y)$ 恒成立且 $u \not\equiv 0$、$k \neq 0$，必须满足 $a^2 = b^2$，即 $a = \pm b$。

      \item 对于一般情况 $u=\mathrm{e}^{\alpha x+\beta t}\sin(\gamma x+\delta t)$ 满足 $u_t=u_{xx}$：
      求对 $t$ 的偏导数：
      \(\displaystyle u_t = \beta \mathrm{e}^{\alpha x+\beta t}\sin(\gamma x+\delta t) + \delta \mathrm{e}^{\alpha x+\beta t}\cos(\gamma x+\delta t)\)。
      求对 $x$ 的一阶与二阶偏导数：
      \(\displaystyle u_x = \alpha \mathrm{e}^{\alpha x+\beta t}\sin(\gamma x+\delta t) + \gamma \mathrm{e}^{\alpha x+\beta t}\cos(\gamma x+\delta t)\)。
      \begin{align*}
      u_{xx} &= \alpha \left[ \alpha \mathrm{e}^{\dots}\sin(\dots) + \gamma \mathrm{e}^{\dots}\cos(\dots) \right] + \gamma \left[ \alpha \mathrm{e}^{\dots}\cos(\dots) - \gamma \mathrm{e}^{\dots}\sin(\dots) \right] \\[6pt]
      &= (\alpha^2-\gamma^2)\mathrm{e}^{\alpha x+\beta t}\sin(\gamma x+\delta t) + 2\alpha\gamma \mathrm{e}^{\alpha x+\beta t}\cos(\gamma x+\delta t).
      \end{align*}
      比较 $u_t$ 和 $u_{xx}$ 在正弦和余弦项前的系数（由于 $\sin$ 和 $\cos$ 线性无关），得
      \(\displaystyle \beta = \alpha^2 - \gamma^2,\qquad \delta = 2\alpha\gamma\)。
      这就是所需条件。
  \end{enumerate}

  \item[(3)] \textbf{解析：}
  \begin{enumerate}
      \item 多次应用链式法则可得
      \begin{align*}
      u_x &= f'(x+at) + g'(x-at), & u_{xx} &= f''(x+at) + g''(x-at),\\
      u_{xxxx} &= f^{(4)}(x+at) + g^{(4)}(x-at), & u_t &= a f'(x+at) - a g'(x-at),\\
      u_{tt} &= a^2 f''(x+at) + a^2 g''(x-at), & u_{ttt} &= a^3 f'''(x+at) - a^3 g'''(x-at),\\
      u_{tttt} &= a^4 f^{(4)}(x+at) + a^4 g^{(4)}(x-at).
      \end{align*}
      \item 验证四阶方程：\(\displaystyle a^4 u_{xxxx} = a^4 \left[ f^{(4)}(x+at) + g^{(4)}(x-at) \right] = u_{tttt}\)。
      故 $a^4u_{xxxx}-u_{tttt}=0$ 成立，证毕。
  \end{enumerate}
\end{enumerate}
\end{solution}

\begin{exercise}[坐标变换]
\begin{shortsingleenum}
  \shortsingleitem{1}{利用前面命题的结论，将二维 Laplace 方程 $u_{xx}+u_{yy}=0$ 化为极坐标形式，并验证 $u=r^n\cos(n\theta)$ 和 $u=r^n\sin(n\theta)$（$n$ 为整数）都是其解；}
  \shortsingleitem{2}{设 $u=f(x,y)$ 具有三阶连续偏导数。计算 $\dfrac{\partial^3 u}{\partial x^3}+\dfrac{\partial^3 u}{\partial x\partial y^2} +\dfrac{\partial^3 u}{\partial y\partial x^2}+\dfrac{\partial^3 u}{\partial y^3}$ 在极坐标下的表达式（提示：利用已知的 $u_{xx}+u_{yy}$ 结果并对其继续求导）；}
  \shortsingleitem{3}{作抛物线坐标变换 $x=\dfrac{u^2-v^2}{2}$，$y=uv$。计算在此变换下 Laplace 算子 $\Delta=\partial_x^2+\partial_y^2$ 的表达式。}
\end{shortsingleenum}
\end{exercise}
\begin{solution}
\begin{enumerate}
  \item[(1)] \textbf{解析：}
  \begin{enumerate}
      \item 由命题~\ref{prop:polar-laplacian}，极坐标下的 Laplace 方程为 \(\displaystyle u_{rr} + \frac{1}{r}u_r + \frac{1}{r^2}u_{\theta\theta} = 0\)。
      \item 验证 $u=r^n\cos(n\theta)$ 时，\(\displaystyle u_r=n r^{n-1}\cos(n\theta)\)、\(\displaystyle u_{rr}=n(n-1)r^{n-2}\cos(n\theta)\)、\(\displaystyle u_\theta=-n r^n\sin(n\theta)\)、\(\displaystyle u_{\theta\theta}=-n^2 r^n\cos(n\theta)\)。代入得 \(\displaystyle \Delta u=\bigl[n(n-1)+n-n^2\bigr]r^{n-2}\cos(n\theta)=0\)。
      \item 同理，\(u=r^n\sin(n\theta)\) 的验证只是在 $\theta$ 方向再出现一个 $-n^2$ 因子，完全同样可得零。
  \end{enumerate}

  \item[(2)] \textbf{解析：}
  \begin{enumerate}
      \item 观察待求式的结构，可将其提取出公共算子：
      \begin{align*}
      \frac{\partial^3 u}{\partial x^3} + \frac{\partial^3 u}{\partial x\partial y^2} + \frac{\partial^3 u}{\partial y\partial x^2} + \frac{\partial^3 u}{\partial y^3}
      &= \frac{\partial}{\partial x}\left(\frac{\partial^2 u}{\partial x^2} + \frac{\partial^2 u}{\partial y^2}\right) + \frac{\partial}{\partial y}\left(\frac{\partial^2 u}{\partial x^2} + \frac{\partial^2 u}{\partial y^2}\right) \\[6pt]
      &= \left(\frac{\partial}{\partial x} + \frac{\partial}{\partial y}\right)(\Delta u).
      \end{align*}
      \item 已知 $\Delta u = u_{rr} + \dfrac{1}{r}u_r + \dfrac{1}{r^2}u_{\theta\theta}$。令 $W = \Delta u$，则我们需要计算 $W_x + W_y$。
      \item 由链式法则，\(\displaystyle W_x = W_r r_x + W_\theta \theta_x = W_r \cos\theta - W_\theta \frac{\sin\theta}{r}\)，\(\displaystyle W_y = W_r r_y + W_\theta \theta_y = W_r \sin\theta + W_\theta \frac{\cos\theta}{r}\)。从而 \(\displaystyle W_x + W_y = W_r(\cos\theta + \sin\theta) + \frac{W_\theta}{r}(\cos\theta - \sin\theta)\)。
      将 $W = u_{rr} + \dfrac{1}{r}u_r + \dfrac{1}{r^2}u_{\theta\theta}$ 代入即可完成展开。
  \end{enumerate}

  \item[(3)] \textbf{解析：}
  此为正交曲线坐标系的算子变换，通过直接应用链式法则推导：
  \begin{enumerate}
      \item 利用复合函数求导公式，分别表达对新变量 $u,v$ 的一阶偏导算子（注意这里自变量为 $u,v$ 函数为 $f$）：
      \begin{itemize}
          \item $\dfrac{\partial}{\partial u} = \dfrac{\partial x}{\partial u}\dfrac{\partial}{\partial x} + \dfrac{\partial y}{\partial u}\dfrac{\partial}{\partial y} = u\dfrac{\partial}{\partial x} + v\dfrac{\partial}{\partial y}$
          \item $\dfrac{\partial}{\partial v} = \dfrac{\partial x}{\partial v}\dfrac{\partial}{\partial x} + \dfrac{\partial y}{\partial v}\dfrac{\partial}{\partial y} = -v\dfrac{\partial}{\partial x} + u\dfrac{\partial}{\partial y}$
      \end{itemize}

      \item 计算对 $u$ 的二阶导数 $\dfrac{\partial^2 f}{\partial u^2}$：
      \begin{align*}
      \frac{\partial^2 f}{\partial u^2} &= \frac{\partial}{\partial u}\left(u f_x + v f_y\right) = f_x + u \frac{\partial f_x}{\partial u} + v \frac{\partial f_y}{\partial u} \\[6pt]
      &= f_x + u(u f_{xx} + v f_{xy}) + v(u f_{yx} + v f_{yy}) \\[6pt]
      &= f_x + u^2 f_{xx} + 2uv f_{xy} + v^2 f_{yy}.
      \end{align*}

      \item 计算对 $v$ 的二阶导数 $\dfrac{\partial^2 f}{\partial v^2}$：
      \begin{align*}
      \frac{\partial^2 f}{\partial v^2} &= \frac{\partial}{\partial v}\left(-v f_x + u f_y\right) = -f_x - v \frac{\partial f_x}{\partial v} + u \frac{\partial f_y}{\partial v} \\[6pt]
      &= -f_x - v(-v f_{xx} + u f_{xy}) + u(-v f_{yx} + u f_{yy}) \\[6pt]
      &= -f_x + v^2 f_{xx} - 2uv f_{xy} + u^2 f_{yy}.
      \end{align*}

      \item 结合上述结果求 Laplace 算子：
      将两式相加，一阶导数和混合导数项刚好抵消，并解出抛物线坐标下的 Laplace 算子：
      \begin{align*}
      \frac{\partial^2 f}{\partial u^2}+\frac{\partial^2 f}{\partial v^2}
      &=(u^2+v^2)f_{xx}+(u^2+v^2)f_{yy}
      =(u^2+v^2)\left(\frac{\partial^2 f}{\partial x^2}+\frac{\partial^2 f}{\partial y^2}\right),\\
      \Delta f
      &=\frac{\partial^2 f}{\partial x^2}+\frac{\partial^2 f}{\partial y^2}
      =\frac{1}{u^2+v^2}\left(\frac{\partial^2 f}{\partial u^2}+\frac{\partial^2 f}{\partial v^2}\right).
      \end{align*}
  \end{enumerate}
\end{enumerate}
\end{solution}

\begin{exercise}[隐函数 PDE 与变量替换的综合运用]
\begin{shortsingleenum}
  \shortsinglechoice{(1)}{设 $z=z(x,y)$ 由方程 $F\!\left(\dfrac{z}{x},\dfrac{z}{y}\right)=0$ 确定。证明 $x\dfrac{\partial z}{\partial x}+y\dfrac{\partial z}{\partial y}=z$；}
  \shortsinglechoice{(2)}{作变量替换 $\xi=x-y$, $\eta=xy$，$w=z^2$。将方程 $z_{xx}+z_{yy}=0$ 变换为关于 $w(\xi,\eta)$ 的方程；}
  \shortsinglechoice{(3)}{设 $u(x,t)=f(x+t)+g(x-t)+xt$。求 $u_{tt}-u_{xx}$ 的表达式。}
\end{shortsingleenum}
\end{exercise}
\begin{solution}
\begin{enumerate}
  \item[(1)] \textbf{解析：}
  \begin{enumerate}
      \item 引入中间变量 $\alpha = \dfrac{z}{x}, \beta = \dfrac{z}{y}$，则方程化为 $F(\alpha, \beta) = 0$。
      \item 分别对 $x,y$ 求全导数，得
      \begin{align*}
      F_\alpha \cdot \frac{z_x x - z}{x^2} + F_\beta \cdot \frac{z_x}{y} &= 0 \implies z_x\left(\frac{F_\alpha}{x} + \frac{F_\beta}{y}\right) = \frac{z}{x^2}F_\alpha,\\
      F_\alpha \cdot \frac{z_y}{x} + F_\beta \cdot \frac{z_y y - z}{y^2} &= 0 \implies z_y\left(\frac{F_\alpha}{x} + \frac{F_\beta}{y}\right) = \frac{z}{y^2}F_\beta.
      \end{align*}
      \item 令 $K = \dfrac{F_\alpha}{x} + \dfrac{F_\beta}{y}$，计算目标式左端：
      \begin{align*}
      x z_x + y z_y &= x \cdot \frac{\frac{z}{x^2}F_\alpha}{K} + y \cdot \frac{\frac{z}{y^2}F_\beta}{K} \\[6pt]
      &= \frac{\frac{z}{x}F_\alpha + \frac{z}{y}F_\beta}{K} = \frac{z\left(\frac{F_\alpha}{x} + \frac{F_\beta}{y}\right)}{K} \\[6pt]
      &= \frac{z \cdot K}{K} = z.
      \end{align*}
      原命题得证。
  \end{enumerate}

  \item[(2)] \textbf{解析：}
  \begin{enumerate}
      \item 确立各变量间的关系：$z = \sqrt{w}$。
      计算 $\xi_x = 1, \xi_y = -1$ 以及 $\eta_x = y, \eta_y = x$。

      \item 计算一阶导数：\(\displaystyle z_x=\frac{1}{2z}(w_\xi+yw_\eta),\qquad z_y=\frac{1}{2z}(-w_\xi+xw_\eta)\)。

      \item 计算二阶导数：
      利用乘积法则 $\partial_x \left( \dfrac{A}{2z} \right) = \dfrac{\partial_x A}{2z} - \dfrac{A z_x}{2z^2}$。
      对于 $z_{xx}$，令 $A = w_\xi + y w_\eta$：
      \begin{align*}
      \partial_x A
      &= w_{\xi\xi} \xi_x + w_{\xi\eta} \eta_x + y(w_{\eta\xi} \xi_x + w_{\eta\eta} \eta_x)
      = w_{\xi\xi} + 2y w_{\xi\eta} + y^2 w_{\eta\eta},\\
      z_{xx}
      &= -\frac{1}{4z^3}(w_\xi + y w_\eta)^2 + \frac{1}{2z}(w_{\xi\xi} + 2y w_{\xi\eta} + y^2 w_{\eta\eta}).
      \end{align*}
      对于 $z_{yy}$，令 $B = -w_\xi + x w_\eta$：
      \begin{align*}
      \partial_y B
      &= -(w_{\xi\xi}(-1) + w_{\xi\eta} x) + x(w_{\eta\xi}(-1) + w_{\eta\eta} x)
      = w_{\xi\xi} - 2x w_{\xi\eta} + x^2 w_{\eta\eta},\\
      z_{yy}
      &= -\frac{1}{4z^3}(-w_\xi + x w_\eta)^2 + \frac{1}{2z}(w_{\xi\xi} - 2x w_{\xi\eta} + x^2 w_{\eta\eta}).
      \end{align*}

      \item 合并并利用原方程：
      $z_{xx}+z_{yy}=0$，等式两端同乘 $4z^3 = 4w\sqrt{w}$（消去分母）：
      \begin{align*}
      0 = &\, 2z^2 \left[ 2w_{\xi\xi} + 2(y-x)w_{\xi\eta} + (x^2+y^2)w_{\eta\eta} \right] \\[4pt] - \left[ (w_\xi + y w_\eta)^2 + (-w_\xi + x w_\eta)^2 \right].
      \end{align*}

      \item 利用新变量 $\xi, \eta$ 替换原变量 $x, y$：
      已知 $y-x = -\xi$，且 $x^2+y^2 = (x-y)^2 + 2xy = \xi^2 + 2\eta$。
      整理平方项：
      \[
      (w_\xi + y w_\eta)^2 + (-w_\xi + x w_\eta)^2 = 2w_\xi^2 + 2(y-x)w_\xi w_\eta + (x^2+y^2)w_\eta^2 = 2w_\xi^2 - 2\xi w_\xi w_\eta + (\xi^2+2\eta)w_\eta^2.
      \]
      代入最终方程：
      \[
      2w \left[ 2w_{\xi\xi} - 2\xi w_{\xi\eta} + (\xi^2+2\eta)w_{\eta\eta} \right] - \left[ 2w_\xi^2 - 2\xi w_\xi w_\eta + (\xi^2+2\eta)w_\eta^2 \right] = 0.
      \]
  \end{enumerate}

  \item[(3)] \textbf{解析：}
  \begin{enumerate}
      \item 计算得
      \begin{align*}
      u_t &= f'(x+t) - g'(x-t) + x,\qquad
      u_{tt} = f''(x+t) + g''(x-t),\\
      u_x &= f'(x+t) + g'(x-t) + t,\qquad
      u_{xx} = f''(x+t) + g''(x-t),\\
      u_{tt} - u_{xx}
      &= \left[ f''(x+t) + g''(x-t) \right] - \left[ f''(x+t) + g''(x-t) \right] = 0.
      \end{align*}
  \end{enumerate}
\end{enumerate}
\end{solution}

\begin{exercise}[柱坐标系下的 Laplace 算子]
直角坐标系与柱坐标系之间的变换式为 $x=r\cos\theta,\,y=r\sin\theta,\,z=z$。试将拉普拉斯方程 \(\displaystyle \frac{\partial^2 u}{\partial x^2} + \frac{\partial^2 u}{\partial y^2} + \frac{\partial^2 u}{\partial z^2} = 0\) 变换到柱坐标系下表示。
\end{exercise}
\begin{solution}
柱坐标系就是 $xy$ 平面内的极坐标变换外加 $z$ 轴恒等变换，因此由命题~\ref{prop:polar-laplacian}，再把未变的 $\dfrac{\partial^2 u}{\partial z^2}$ 直接加上，得
\begin{align*}
\frac{\partial^2 u}{\partial x^2}+\frac{\partial^2 u}{\partial y^2}
&=\frac{\partial^2 u}{\partial r^2}+\frac{1}{r}\frac{\partial u}{\partial r}+\frac{1}{r^2}\frac{\partial^2 u}{\partial\theta^2},\\
\frac{\partial^2 u}{\partial r^2}+\frac{1}{r}\frac{\partial u}{\partial r}
&+\frac{1}{r^2}\frac{\partial^2 u}{\partial\theta^2}+\frac{\partial^2 u}{\partial z^2}=0.
\end{align*}
\end{solution}

\begin{note}[柱坐标 Laplace 算子的结构]
柱坐标只是在 $xy$ 平面把直角坐标换成极坐标，$z$ 方向完全没有变。因此
\[
\Delta
=\left(\frac{\partial^2}{\partial r^2}+\frac{1}{r}\frac{\partial}{\partial r}+\frac{1}{r^2}\frac{\partial^2}{\partial\theta^2}\right)
+\frac{\partial^2}{\partial z^2}.
\]
也就是说，前面三项负责横截面里的扩散，最后的 $u_{zz}$ 负责轴向变化。
\end{note}

\begin{exercise}[特征线变换下的波动方程]
作自变量变换 $\xi=x+t,\,\eta=x-t$，变换弦振动方程
\(\displaystyle \frac{\partial^2 u}{\partial x^2} = \frac{\partial^2 u}{\partial t^2}\)。
\end{exercise}
\begin{solution}
根据变换 $\xi=x+t,\ \eta=x-t$，有 $\xi_x=\eta_x=1,\ \xi_t=1,\ \eta_t=-1$，从而 \(\displaystyle u_x = u_\xi + u_\eta,\ u_{xx} = u_{\xi\xi} + 2u_{\xi\eta} + u_{\eta\eta}\)，且 \(\displaystyle u_t = u_\xi - u_\eta,\ u_{tt} = u_{\xi\xi} - 2u_{\xi\eta} + u_{\eta\eta}\)。
代入 $u_{xx}=u_{tt}$，得
\[
u_{\xi\xi} + 2u_{\xi\eta} + u_{\eta\eta} = u_{\xi\xi} - 2u_{\xi\eta} + u_{\eta\eta}
\implies 4u_{\xi\eta}=0
\implies \frac{\partial^2 u}{\partial \xi \partial \eta}=0.
\]
这正是波动方程在特征坐标系下的最简标准型。
\end{solution}

% ============================================================
% 第七章 总习题
% ============================================================

\subsection{本章总习题}

% ============================================================
% 一、填空题（一）
% ============================================================

\begin{exercise}[复合函数、极限与间断点]
\begin{shortsingleenum}
  \shortsingleitem{1}{已知 $f(x+y,x-y)=xy+y^2$，则 $f(x,y)=\underline{\hspace{1cm}}$；}
  \shortsingleitem{2}{极限 $\displaystyle\lim_{\substack{x\to\infty \\ y\to 0}} x\sin y=\underline{\hspace{1cm}}$；}
  \shortsingleitem{3}{函数 $f(x,y)=\dfrac{1}{x^3y-xy^2}$ 的间断点集为\underline{\hspace{1cm}}；}
  \shortsingleitem{4}{函数 $z=2x^3+y^3$ 在点 $(1,1)$ 处沿梯度方向的方向导数为$\underline{\hspace{1cm}}$；}
  \shortsingleitem{5}{设 $f(x,y,z)=xy^2z^3$，其中 $z=z(x,y)$ 是由方程 $x^2+y^2+z^2=3xyz$ 所确定的函数，则 $f_x(1,1,1)=\underline{\hspace{1cm}}$；}
  \shortsingleitem{6}{设 $y=y(x)$ 为由方程 $f(x^2+y^2,xy)=0$ 确定的隐函数，其中 $f(u,v)$ 为二元可微函数，则 $\dfrac{\mathrm{d}y}{\mathrm{d}x}=\underline{\hspace{1cm}}$；}
  \shortsingleitem{7}{函数 $z=\arctan\dfrac{y}{x}+\ln\sqrt{x^2+y^2}$ 在点 $(1,0)$ 处的梯度 $\left.\mathrm{grad}\,z\right|_{(1,0)}=\underline{\hspace{1cm}}$；}
  \shortsingleitem{8}{由方程 $xyz+\sqrt{x^2+y^2+z^2}=\sqrt{2}$ 所确定的函数 $z=z(x,y)$ 在点 $(1,0,-1)$ 处的全微分 $\mathrm{d}z=\underline{\hspace{1cm}}$；}
  \shortsingleitem{9}{设 $\boldsymbol a_0$ 是常向量，$\lambda(t)$ 是数值函数，则 $\dfrac{\mathrm{d}}{\mathrm{d}t}\bigl(\lambda(t)\boldsymbol a_0\bigr)=\underline{\hspace{1cm}}$，$\dfrac{\mathrm{d}}{\mathrm{d}t}\bigl(\boldsymbol a_0\cdot\boldsymbol r(t)\bigr)=\underline{\hspace{1cm}}$。}
\end{shortsingleenum}
\end{exercise}
\begin{solution}
\begin{enumerate}
  \item \textbf{解析}：令 $u=x+y, v=x-y$，则 $x=\dfrac{u+v}{2}, y=\dfrac{u-v}{2}$。代入原式得 \(\displaystyle f(u,v) = \frac{u+v}{2} \cdot \frac{u-v}{2} + \left(\frac{u-v}{2}\right)^2 = \frac{u^2-v^2}{4} + \frac{u^2-2uv+v^2}{4} = \frac{u^2-uv}{2}\)。
  故填 $\mathbf{\dfrac{x^2-xy}{2}}$。

  \item \textbf{解析}：考察不同路径。若沿 $y = \dfrac{1}{x}$ 趋近，原式变为 $x\sin\dfrac{1}{x} \to 1$；若沿 $y = \dfrac{1}{x^2}$ 趋近，原式变为 $x\sin\dfrac{1}{x^2} \to 0$。极限不唯一，故填 \textbf{不存在}。

  \item \textbf{解析}：间断点即分母为零的点。$x^3y-xy^2 = xy(x^2-y) = 0 \implies x=0$ 或 $y=0$ 或 $y=x^2$。
  故填 $\mathbf{\{(x,y) \mid xy(x^2-y)=0\}}$。

  \item \textbf{解析}：梯度为 $\nabla z = (6x^2, 3y^2)$，故在 $(1,1)$ 处有 $\nabla z=(6,3)$。沿梯度方向的方向导数等于梯度的模长，
  \(\displaystyle \|\nabla z(1,1)\| = \sqrt{6^2+3^2} = \sqrt{45} = 3\sqrt{5}\)。
  故填 $\mathbf{3\sqrt{5}}$。

  \item \textbf{解析}：对约束方程两边求导：$2x + 2z \cdot z_x = 3yz + 3xy \cdot z_x \implies z_x = \dfrac{3yz-2x}{2z-3xy}$。
  在 $(1,1,1)$ 处，$z_x = \dfrac{3-2}{2-3} = -1$。
  $f_x = y^2z^3 + 3xy^2z^2 z_x$，代入 $(1,1,1)$ 得 $f_x = 1 + 3(-1) = -2$。故填 $\mathbf{-2}$。

  \item \textbf{解析}：方程两边对 $x$ 求全导数：$f_1 \cdot (2x+2y y') + f_2 \cdot (y+x y') = 0$。解得 $y' = -\dfrac{2x f_1 + y f_2}{2y f_1 + x f_2}$。故填 $\mathbf{-\dfrac{2x f_1 + y f_2}{2y f_1 + x f_2}}$。

  \item \textbf{解析}：分别求偏导得 \(\displaystyle z_x = \frac{-y/x^2}{1+(y/x)^2} + \frac{x}{x^2+y^2} = \frac{x-y}{x^2+y^2},\qquad z_y = \frac{1/x}{1+(y/x)^2} + \frac{y}{x^2+y^2} = \frac{x+y}{x^2+y^2}\)。
  代入点 $(1,0)$，得 $z_x = 1,\ z_y = 1$。故填 $\mathbf{\{1, 1\}}$ 或 $\mathbf{\boldsymbol i + \boldsymbol j}$。

  \item \textbf{解析}：令 $F = xyz+\sqrt{x^2+y^2+z^2}-\sqrt{2} = 0$，在点 $(1,0,-1)$ 处 $\sqrt{x^2+y^2+z^2} = \sqrt{2}$，于是 \(\displaystyle F_x = \frac{\sqrt{2}}{2},\qquad F_y = -1,\qquad F_z = -\frac{\sqrt{2}}{2}\)。
  从而 $z_x = -\dfrac{F_x}{F_z} = 1$，$z_y = -\dfrac{F_y}{F_z} = -\sqrt{2}$，故全微分 $\mathrm{d}z = \mathbf{\mathrm{d}x - \sqrt{2}\mathrm{d}y}$。

  \item \textbf{解析}：向量求导直接应用法则，故填 $\mathbf{\lambda'(t)\boldsymbol a_0}$ 与 $\mathbf{\boldsymbol a_0\cdot\boldsymbol r'(t)}$。
\end{enumerate}
\end{solution}

\begin{exercise}[极值、条件极值与曲面几何]
\begin{shortexenum}
  \shortitem{1}{函数 $z=x^3+y^3-3xy$ 的极小值为\underline{\hspace{1cm}}；} &
  \shortitem{2}{点 $(0,0)$ 是函数 $z=x^2-4xy+5y^2-1$ 的极\underline{\hspace{1cm}}点；} \\
  \shortitem{3}{函数 $z=xy$ 在条件 $x+y=1$ 下的极大值为\underline{\hspace{1cm}}；} &
  \shortitem{4}{在椭圆抛物面 $\dfrac{z}{c}=\dfrac{x^2}{a^2}+\dfrac{y^2}{b^2}$（$z\leqslant c$）的一段中，嵌入有最大体积的直角平行六面体，则该六面体的尺寸为长$=\underline{\hspace{2cm}}$，宽$=\underline{\hspace{2cm}}$，高$=\underline{\hspace{2cm}}$；} \\
  \shortitem{5}{函数 $z=3x^2+3y^2-x^3$ 在闭区域 $D:x^2+y^2\leqslant 16$ 上的最小值是\underline{\hspace{1cm}}；} &
  \shortitem{6}{当 $a>0$ 时，曲线 $\begin{cases}x^2+y^2+z^2=4a^2,\\ x^2+y^2=2ax\end{cases}$ 在点 $M(a,a,\sqrt{2}\,a)$ 处的切线方程为\underline{\hspace{1cm}}，法平面方程为\underline{\hspace{1cm}}；} \\
  \shortitem{7}{设 $F$ 为可微函数，$a,b,c$ 为非零常数，则由方程 $F(cx-az,\,cy-bz)=0$ 给出的曲面 $S$ 上任一点处的法向量为 $\boldsymbol n=\underline{\hspace{1cm}}$；} &
  \shortitem{8}{椭球面 $2x^2+3y^2+z^2=6$ 在点 $P(1,1,1)$ 处的切平面方程为\underline{\hspace{1cm}}，法线方程为$\underline{\hspace{1cm}}$。}
\end{shortexenum}
\end{exercise}
\begin{solution}
\begin{enumerate}
  \item \textbf{解析}：求驻点，令 $z_x=3x^2-3y=0, z_y=3y^2-3x=0$，得驻点 $(0,0)$ 和 $(1,1)$。在 $(1,1)$ 处，$A=6, B=-3, C=6, AC-B^2 = 27 > 0$ 且 $A>0$，取极小值 $z(1,1)=-1$。故填 $\mathbf{-1}$。
  \item \textbf{解析}：$z_x = 2x-4y, z_y = -4x+10y$，在 $(0,0)$ 处均取 $0$。$A=2, B=-4, C=10$，$AC-B^2 = 20-16=4>0, A>0$。故填 \textbf{小值}。
  \item \textbf{解析}：代入得 $z=x(1-x)=-x^2+x$。当 $x=\dfrac{1}{2}$ 时取极大值 $\dfrac{1}{4}$。故填 $\mathbf{\dfrac{1}{4}}$。
  \item \textbf{解析}：设第一象限顶点为 $(x,y,z)$，体积 $V=8xy(c-z)=8cxy(1-x^2/a^2-y^2/b^2)$。利用均值不等式，当 $\dfrac{x^2}{a^2} = \dfrac{y^2}{b^2} = 1 - \dfrac{x^2}{a^2} - \dfrac{y^2}{b^2}$ 即 $x=\dfrac{a}{2}, y=\dfrac{b}{2}, z=\dfrac{c}{2}$ 时体积最大。故长宽高等于 $2x, 2y, c-z$。故填 $\mathbf{a}$、$\mathbf{b}$、$\mathbf{c/2}$。
  \item \textbf{解析}：求区域内部驻点：$z_x=6x-3x^2=0, z_y=6y=0 \implies (0,0)$ 与 $(2,0)$。值分别为 $0$ 和 $4$。考察边界 $y^2=16-x^2$：$z=3x^2+3(16-x^2)-x^3 = 48-x^3 \ (x \in [-4,4])$。当 $x=4$ 时取最小值 $-16$。故填 $\mathbf{-16}$。
  \item \textbf{解析}：两曲面的法向量在 $M$ 处分别平行于 $(1,1,\sqrt{2})$ 与 $(0,1,0)$，故切向 $\boldsymbol T = (-\sqrt{2}, 0, 1)$。于是切线可写为 $\mathbf{(x,y,z)=(a,a,\sqrt{2}a)+t(-\sqrt{2},0,1)}$，法平面方程为 $\mathbf{\sqrt{2}x - z = 0}$。
  \item \textbf{解析}：令 $u=cx-az, v=cy-bz$，则法向量 $\boldsymbol n = (F_x, F_y, F_z) = (c F_u, c F_v, -a F_u - b F_v)$。故填 $\mathbf{(c F_1, c F_2, -a F_1 - b F_2)}$。
  \item \textbf{解析}：法向量 $\boldsymbol n = (4x, 6y, 2z)|_P = (4, 6, 2) \parallel (2, 3, 1)$，故切平面方程为 $2(x-1)+3(y-1)+1(z-1)=0 \implies \mathbf{2x+3y+z-6=0}$，法线方程为 $\mathbf{\dfrac{x-1}{2} = \dfrac{y-1}{3} = \dfrac{z-1}{1}}$。
\end{enumerate}
\end{solution}

% ============================================================
% 三、选择题（基础概念）
% ============================================================

\begin{exercise}[函数性质、连续性与可微性]
以下每题只有一个正确答案。
\begin{enumerate}
  \item 设 $z_1=x+y$, $z_2=(\sqrt{x+y})^2$, $z_3=\sqrt{(x+y)^2}$，则（\qquad）。
    \begin{shortexenum}
      \shortchoice{(A)}{$z_1$ 与 $z_2$ 是相同的函数} &
      \shortchoice{(B)}{$z_1$ 与 $z_3$ 是相同的函数} \\
      \shortchoice{(C)}{$z_2$ 与 $z_3$ 是相同的函数} &
      \shortchoice{(D)}{其中任何两个都不是相同的函数}
    \end{shortexenum}

  \item 极限 $\displaystyle\lim_{\substack{x\to x_0 \\ y\to y_0}}f(x,y)$ 存在是 $f(x,y)$ 在点 $(x_0,y_0)$ 连续的（\qquad）。
    \begin{shortexenum}
      \shortchoice{(A)}{必要条件} &
      \shortchoice{(B)}{充分条件} \\
      \shortchoice{(C)}{充分必要条件} &
      \shortchoice{(D)}{既非必要条件又非充分条件}
    \end{shortexenum}

  \item 函数 $z=\arcsin\dfrac{1}{x^2+y^2}+\ln(1-x^2-y^2)$ 的定义域为（\qquad）。
    \begin{shortexenum}
      \shortchoice{(A)}{空集} &
      \shortchoice{(B)}{$\{(x,y)\mid x^2+y^2\leqslant 1\}$} \\
      \shortchoice{(C)}{$\{(0,0)\}$} &
      \shortchoice{(D)}{$\{(x,y)\mid x^2+y^2=1\}$}
    \end{shortexenum}

  \item 二元函数 $f(x,y)=\begin{cases}\dfrac{xy}{x^2+y^2}, & (x,y) \neq (0,0),\\[4pt] 0, & (x,y)=(0,0)\end{cases}$
        在点 $(0,0)$ 处（\qquad）。
    \begin{shortexenum}
      \shortchoice{(A)}{连续，偏导数存在} &
      \shortchoice{(B)}{连续，偏导数不存在} \\
      \shortchoice{(C)}{不连续，偏导数存在} &
      \shortchoice{(D)}{不连续，偏导数不存在}
    \end{shortexenum}

  \item 设 $f(x,y)$ 在点 $(x_0,y_0)$ 处的两个偏导数 $f_x(x_0,y_0)$ 及 $f_y(x_0,y_0)$ 存在，
        则（\qquad）。
    \begin{shortexenum}
      \shortchoice{(A)}{$f(x,y)$ 在点 $(x_0,y_0)$ 必连续} &
      \shortchoice{(B)}{$f(x,y)$ 在点 $(x_0,y_0)$ 必可微} \\
      \shortchoice{(C)}{$\lim\limits_{x\to x_0}f(x,y_0)$ 与 $\lim\limits_{y\to y_0}f(x_0,y)$ 都存在} &
      \shortchoice{(D)}{$\lim\limits_{\substack{x\to x_0 \\ y\to y_0}}f(x,y)$ 存在}
    \end{shortexenum}

  \item 二元函数 $f(x,y)$ 在点 $(x_0,y_0)$ 处两个偏导数 $f_x(x_0,y_0)$ 及 $f_y(x_0,y_0)$ 存在
        是 $f(x,y)$ 在该点连续的（\qquad）。
    \begin{shortexenum}
      \shortchoice{(A)}{充分而非必要条件} &
      \shortchoice{(B)}{必要而非充分条件} \\
      \shortchoice{(C)}{充分必要条件} &
      \shortchoice{(D)}{既非充分又非必要条件}
    \end{shortexenum}

  \item 设 $z=f(x,y)$ 可微，且 $f(x,x^2)=\dfrac{1}{2}$，$f_x(x,y)\big|_{y=x^2}=\dfrac{1}{x}$，则 $f_y(x,y)\big|_{y=x^2}=$（\qquad）。
    \begin{shortexenum}
      \shortchoice{(A)}{$-\dfrac{1}{x^2}$} &
      \shortchoice{(B)}{$\dfrac{1}{x^2}$} \\
      \shortchoice{(C)}{$\dfrac{1}{2x^2}$} &
      \shortchoice{(D)}{$-\dfrac{1}{2x^2}$}
    \end{shortexenum}

  \item 设 $f(x,y)$ 是可微函数，且 $f(0,0)=0$, $f_x(0,0)=a$, $f_y(0,0)=b$。
        令 $\varphi(t)=f(t,f(t,t))$，则 $\varphi'(0)=$（\qquad）。
    \begin{shortexenum}
      \shortchoice{(A)}{$a$} &
      \shortchoice{(B)}{$a+b(a+b)$} \\
      \shortchoice{(C)}{$a+1$} &
      \shortchoice{(D)}{$\dfrac{a}{1-b}$}
    \end{shortexenum}

  \item 设 $E=\{(x,y)\mid xy>0\}$，则 $E$ 是 $\mathbb R^2$ 中的（\qquad）。
    \begin{shortexenum}
      \shortchoice{(A)}{开集} &
      \shortchoice{(B)}{闭集} \\
      \shortchoice{(C)}{区域（连通开集）} &
      \shortchoice{(D)}{有界集}
    \end{shortexenum}
\end{enumerate}
\end{exercise}
\begin{solution}
\begin{enumerate}
  \item \textbf{解析}：$z_1$ 定义域为 $\mathbb R^2$；$z_2=x+y$ 定义域为 $x+y \geqslant 0$；$z_3 = |x+y|$ 定义域为 $\mathbb R^2$ 但解析式与 $z_1$ 不同。三者均不同，选 \textbf{D}。
  \item \textbf{解析}：连续要求极限存在且等于该点函数值。因此极限存在只是连续的必要条件，选 \textbf{A}。
  \item \textbf{解析}：定义域须满足 $\left|\dfrac{1}{x^2+y^2}\right| \leqslant 1 \implies x^2+y^2 \geqslant 1$，同时要求 $1-x^2-y^2 > 0 \implies x^2+y^2 < 1$。两者矛盾，选 \textbf{A}。
  \item \textbf{解析}：极坐标下 \(\displaystyle \frac{xy}{x^2+y^2}=\frac12\sin2\theta\)，当 \(r\to0^+\) 时结果仍依赖角度 \(\theta\)，故极限不存在、不连续。但 $f_x(0,0) = \lim_{x\to 0}\frac{0}{x} = 0, f_y(0,0) = 0$，偏导存在。选 \textbf{C}。
  \item \textbf{解析}：偏导数存在仅保证沿平行于坐标轴方向可导（因而连续），故 $\lim_{x\to x_0}f(x,y_0)$ 等存在，但不足以推断整体连续或可微。选 \textbf{C}。
  \item \textbf{解析}：如上题，偏导数存在不蕴含连续；而连续也不蕴含偏导数存在（如锥面顶点）。两者无必然因果，选 \textbf{D}。
  \item \textbf{解析}：利用隐函数链式法则，对 $f(x,x^2)=\dfrac{1}{2}$ 两边关于 $x$ 求导，得
  \(\displaystyle \frac{\mathrm{d}}{\mathrm{d}x}f(x,x^2) = f_x(x,x^2) + 2x\,f_y(x,x^2) = 0\)。
  代入已知条件 $f_x = \dfrac{1}{x}$，得 $\dfrac{1}{x} + 2x f_y = 0 \implies f_y = -\dfrac{1}{2x^2}$。选 \textbf{D}。（注：原题干可能有排版笔误，已调整为常见形式以契合选项）
  \item \textbf{解析}：由链式求导法则，\(\displaystyle \varphi'(t) = f_x \cdot 1 + f_y \cdot [f_x \cdot 1 + f_y \cdot 1]\)。
  代入 $t=0$ 及各偏导数值，得 $\varphi'(0) = a + b(a+b)$。选 \textbf{B}。
  \item \textbf{解析}：$E$ 代表第一和第三象限的内部区域，不包含坐标轴边界，因而是开集，但不是连通区域，选 \textbf{A}。
\end{enumerate}
\end{solution}

% ============================================================
% 四、选择题（计算与应用）
% ============================================================

\begin{exercise}[极值判定、优化问题与几何应用]
以下每题只有一个正确答案。
\begin{enumerate}
  \item 设 $z=f(x,y)$ 在点 $(x_0,y_0)$ 的某邻域内具有直到二阶的连续偏导数，则 $f(x_0,y_0)$ 为函数的极大值的充分条件是（\qquad）。
    \begin{shortexenum}
      \shortchoice{(A)}{$f_x(x_0,y_0)=0$, $f_y(x_0,y_0)=0$} &
      \shortchoice{(B)}{$[f_{xy}(x_0,y_0)]^2-f_{xx}(x_0,y_0)f_{yy}(x_0,y_0)<0$} \\
      \shortchoice{(C)}{$f_x(x_0,y_0)=f_y(x_0,y_0)=0$, $[f_{xy}(x_0,y_0)]^2-f_{xx}(x_0,y_0)f_{yy}(x_0,y_0)<0$, $f_{xx}(x_0,y_0)>0$} &
      \shortchoice{(D)}{$f_x(x_0,y_0)=f_y(x_0,y_0)=0$, $[f_{xy}(x_0,y_0)]^2-f_{xx}(x_0,y_0)f_{yy}(x_0,y_0)<0$, $f_{xx}(x_0,y_0)<0$}
    \end{shortexenum}

  \item 已知矩形的周长是 $2p$，将它绕其一边旋转而形成一个旋转体，当此旋转体的体积为最大时，矩形两边的长分别为（\qquad）。
    \begin{shortexenum}
      \shortchoice{(A)}{$\dfrac{p}{3},\;\dfrac{2p}{3}$} &
      \shortchoice{(B)}{$\dfrac{p}{2},\;\dfrac{p}{2}$} \\
      \shortchoice{(C)}{$\dfrac{p}{4},\;\dfrac{3p}{4}$} &
      \shortchoice{(D)}{$\dfrac{2p}{5},\;\dfrac{3p}{5}$}
    \end{shortexenum}

  \item 设 $z=f(x,y)=x^4+y^4-x^2-2xy-y^2$，由 $f_x(x,y)=0$ 及 $f_y(x,y)=0$ 求得临界点 $M_0(0,0), M_1(1,1)$ 及 $M_2(-1,-1)$，则（\qquad）。
    \begin{shortexenum}
      \shortchoice{(A)}{$f(M_0)$ 是极大值} &
      \shortchoice{(B)}{$f(M_1)$ 与 $f(M_2)$ 都是极大值} \\
      \shortchoice{(C)}{$f(M_0)$ 是极小值} &
      \shortchoice{(D)}{$f(M_1)$ 与 $f(M_2)$ 都是极小值}
    \end{shortexenum}

  \item 平面 $\dfrac{x}{3}+\dfrac{y}{4}+\dfrac{z}{5}=1$ 和柱面 $x^2+y^2=1$ 的交线上与 $Oxy$ 平面距离最短的点的坐标为（\qquad）。
    \begin{shortexenum}
      \shortchoice{(A)}{$\left(\dfrac{3}{5},\;\dfrac{4}{5},\;\dfrac{35}{12}\right)$} &
      \shortchoice{(B)}{$\left(\dfrac{4}{5},\;\dfrac{3}{5},\;\dfrac{35}{12}\right)$} \\
      \shortchoice{(C)}{$\left(-\dfrac{3}{5},\;-\dfrac{4}{5},\;\dfrac{85}{12}\right)$} &
      \shortchoice{(D)}{$\left(-\dfrac{4}{5},\;-\dfrac{3}{5},\;\dfrac{85}{12}\right)$}
    \end{shortexenum}

  \item 抛物线 $y^2=4x$ 上与直线 $x-y+4=0$ 相距最近的点是（\qquad）。
    \begin{shortexenum}
      \shortchoice{(A)}{$(0,0)$} &
      \shortchoice{(B)}{$(1,1)$} \\
      \shortchoice{(C)}{$(1,2)$} &
      \shortchoice{(D)}{$(2,2\sqrt{2})$}
    \end{shortexenum}

  \item 曲线 $\begin{cases}x^2+y^2+z^2=6,\\ x+y+z=0\end{cases}$ 在点 $(1,-2,1)$ 的切线一定平行于（\qquad）。
    \begin{shortexenum}
      \shortchoice{(A)}{$Oxy$ 平面} &
      \shortchoice{(B)}{$Oyz$ 平面} \\
      \shortchoice{(C)}{$Ozx$ 平面} &
      \shortchoice{(D)}{平面 $x+y+z=0$}
    \end{shortexenum}

  \item 在曲线 $C:\begin{cases}x=t,\\ y=-t^2,\\ z=t^3\end{cases}$ 的所有切线中，与平面 $x+2y+z=4$ 平行的切线（\qquad）。
    \begin{shortexenum}
      \shortchoice{(A)}{只有 $1$ 条} &
      \shortchoice{(B)}{只有 $2$ 条} \\
      \shortchoice{(C)}{至少有 $3$ 条} &
      \shortchoice{(D)}{不存在}
    \end{shortexenum}

  \item 已知曲面 $z=4-x^2-y^2$ 上点 $P$ 处的切平面平行于平面 $2x+2y+z-1=0$，则点 $P$ 的坐标是（\qquad）。
    \begin{shortexenum}
      \shortchoice{(A)}{$(1,-1,2)$} &
      \shortchoice{(B)}{$(-1,1,2)$} \\
      \shortchoice{(C)}{$(1,1,2)$} &
      \shortchoice{(D)}{$(-1,-1,2)$}
    \end{shortexenum}

  \item 曲面 $3x^2+y^2+z^2=12$ 上点 $M(-1,0,3)$ 处的切平面与平面 $z=0$ 的夹角是（\qquad）。
    \begin{shortexenum}
      \shortchoice{(A)}{$\dfrac{\pi}{6}$} &
      \shortchoice{(B)}{$\dfrac{\pi}{4}$} \\
      \shortchoice{(C)}{$\dfrac{\pi}{3}$} &
      \shortchoice{(D)}{$\dfrac{\pi}{2}$}
    \end{shortexenum}

  \item 曲线 $C:x=a\mathrm{e}^t\sin t,\;y=a\mathrm{e}^t\cos t,\;z=a\mathrm{e}^t$ 上任意一点处的切线与（\qquad）形成定角。
    \begin{shortexenum}
      \shortchoice{(A)}{$z$ 轴} &
      \shortchoice{(B)}{$x$ 轴} \\
      \shortchoice{(C)}{$y$ 轴} &
      \shortchoice{(D)}{锥面 $x^2+y^2=z^2$ 的各母线}
    \end{shortexenum}
\end{enumerate}
\end{exercise}
\begin{solution}
\begin{enumerate}
  \item \textbf{解析}：极大值充分条件为驻点处二阶偏导构成的判别式 $AC-B^2 > 0$（即 $B^2-AC < 0$）且 $A = f_{xx} < 0$。选 \textbf{D}。
  \item \textbf{解析}：设边长为 $x, y$，则 $x+y=p$。绕 $x$ 边旋转，体积 $V = \pi y^2 x$。需最大化 $y^2(p-y)$，求导得 $2py - 3y^2 = 0 \implies y = \dfrac{2p}{3}, x = \dfrac{p}{3}$。选 \textbf{A}。
  \item \textbf{解析}：计算二阶偏导数 $A = 12x^2-2, B = -2, C = 12y^2-2$。
  在 $(1,1)$ 和 $(-1,-1)$ 处，$AC-B^2 = 10 \times 10 - 4 = 96 > 0$ 且 $A=10>0$，故均为极小值点。在 $(0,0)$ 处虽然 $AC-B^2=0$，但为鞍点。选 \textbf{D}。
  \item \textbf{解析}：目标是最小化 $z = 5\left(1 - \dfrac{x}{3} - \dfrac{y}{4}\right)$，即最大化 $\dfrac{x}{3} + \dfrac{y}{4}$ 在 $x^2+y^2=1$ 上的值。应用柯西不等式或参数方程易得极大值点为 $\left(\dfrac{4}{5}, \dfrac{3}{5}\right)$，对应 $z = \dfrac{35}{12}$。选 \textbf{B}。
  \item \textbf{解析}：转化为求函数 $f(y) = \dfrac{y^2}{4} - y + 4$ 的极值点，求导得 $\dfrac{y}{2} - 1 = 0 \implies y=2, x=1$。选 \textbf{C}。
  \item \textbf{解析}：切向量等于两曲面法向量的叉乘，$\boldsymbol T = (2, -4, 2) \times (1, 1, 1) = (-6, 0, 6) \parallel (-1, 0, 1)$。由于切向量 $y$ 分量为 $0$，故平行于 $Ozx$ 平面。选 \textbf{C}。
  \item \textbf{解析}：切向量 $\boldsymbol T = (1, -2t, 3t^2)$。与平面平行即与法向量垂直，$\boldsymbol T \cdot (1,2,1) = 1 - 4t + 3t^2 = 0$。解得 $t=1$ 和 $t=\dfrac{1}{3}$，对应 $2$ 条切线。选 \textbf{B}。
  \item \textbf{解析}：曲面法向量 $\boldsymbol n = (2x, 2y, 1)$，平行于平面法向量 $(2, 2, 1)$，故 $x=1, y=1$。代入曲面方程得 $z = 4-1-1 = 2$。选 \textbf{C}。
  \item \textbf{解析}：曲面法向量 $\boldsymbol n_1 = (6x, 2y, 2z)|_M = (-6, 0, 6)$。平面 $z=0$ 法向量 $\boldsymbol n_2 = (0, 0, 1)$。夹角余弦 $\cos\theta = \dfrac{|-6 \times 0 + 0 \times 0 + 6 \times 1|}{\sqrt{36+36} \times 1} = \dfrac{6}{6\sqrt{2}} = \dfrac{\sqrt{2}}{2}$，故夹角为 $\dfrac{\pi}{4}$。选 \textbf{B}。
  \item \textbf{解析}：求切向量 $\boldsymbol r' = (a\mathrm{e}^t(\sin t+\cos t), a\mathrm{e}^t(\cos t-\sin t), a\mathrm{e}^t)$，其模长 $|\boldsymbol r'| = \sqrt{3}a\mathrm{e}^t$。它与 $z$ 轴正向夹角余弦为 $\dfrac{a\mathrm{e}^t}{\sqrt{3}a\mathrm{e}^t} = \dfrac{1}{\sqrt{3}}$（常数）。选 \textbf{A}。
\end{enumerate}
\end{solution}

% ============================================================
% 五、极限计算
% ============================================================

\begin{exercise}[多元函数极限的计算]
求下列函数极限：
\begin{shortexenum}
  \shortitem{1}{$\displaystyle\lim_{\substack{x\to 0 \\ y\to 0}} \frac{3-\sqrt{9+xy}}{xy}$；} &
  \shortitem{2}{$\displaystyle\lim_{\substack{x\to 0 \\ y\to 0}} \frac{x^3+y^3}{x^2+y^2}$；} \\
  \shortitem{3}{$\displaystyle\lim_{\substack{x\to\infty \\ y\to\infty}} \frac{x+y}{x^2-xy+y^2}$。} &
\end{shortexenum}
\end{exercise}
\begin{solution}
\begin{enumerate}
  \item \textbf{解析}：令 \(x=r\cos\theta,\ y=r\sin\theta\)，并对分子有理化：
  \[
  \frac{3-\sqrt{9+xy}}{xy}
  =\frac{-1}{3+\sqrt{9+xy}}
  =\frac{-1}{3+\sqrt{9+r^2\cos\theta\sin\theta}}
  \to-\frac16.
  \]

  \item \textbf{解析}：使用极坐标代换 $x = r\cos\theta, y = r\sin\theta$，则 \(\displaystyle 0\leqslant \left|\frac{r^3(\cos^3\theta+\sin^3\theta)}{r^2}\right|=r|\cos^3\theta+\sin^3\theta|\leqslant2r\)。
  当 $r \to 0$ 时，$2r \to 0$，由夹逼定理知原极限为 $0$。

  \item \textbf{解析}：当 \(x,y\to+\infty\) 时，令 \(x=r\cos\theta,\ y=r\sin\theta\)，其中 \(0<\theta<\frac{\pi}{2}\)，且 \(r\to+\infty\)。则
  \[
    \frac{x+y}{x^2-xy+y^2}
    =\frac{r(\cos\theta+\sin\theta)}
           {r^2(\cos^2\theta-\cos\theta\sin\theta+\sin^2\theta)}.
  \]
  在第一象限内 \(\cos^2\theta-\cos\theta\sin\theta+\sin^2\theta
  =1-\cos\theta\sin\theta\geqslant\frac12\)，且 \(\cos\theta+\sin\theta\leqslant\sqrt2\)。故其绝对值不超过 \(\dfrac{2\sqrt2}{r}\to0\)，原极限为 \(0\)。
\end{enumerate}
\end{solution}

% ============================================================
% 六、连续性的研究
% ============================================================

\begin{exercise}[分段函数的连续性分析]
研究下列函数的连续性：
\begin{shortsingleenum}
  \shortsingleitem{1}{$f(x,y)=\begin{cases}\dfrac{x^2y}{x^2+y^2}, & x^2+y^2 \neq 0,\\[4pt] 0, & x^2+y^2=0.\end{cases}$}
  \shortsingleitem{2}{$f(x,y)=\begin{cases}\dfrac{xy}{\sqrt{x^2+y^2}}, & x^2+y^2 \neq 0,\\[4pt] 0, & x^2+y^2=0.\end{cases}$}
  \shortsingleitem{3}{$f(x,y)=\begin{cases}x\sin\dfrac{1}{y}, & y \neq 0,\\[4pt] 0, & y=0.\end{cases}$}
\end{shortsingleenum}
\end{exercise}
\begin{solution}
\begin{enumerate}
  \item \textbf{解析}：当 $(x,y) \neq (0,0)$ 时由初等函数构成必连续。考察 $(0,0)$ 处，
  令 \(x=r\cos\theta,\ y=r\sin\theta\)，则
  \[
    \frac{x^2y}{x^2+y^2}=r\cos^2\theta\sin\theta,
  \]
  其绝对值不超过 \(r\)，故原点处极限为 \(0=f(0,0)\)。
  故函数在全平面连续。

  \item \textbf{解析}：当 $(x,y) \neq (0,0)$ 时由初等函数构成必连续。考察 $(0,0)$ 处，
  极坐标下
  \[
    \frac{xy}{\sqrt{x^2+y^2}}=r\cos\theta\sin\theta,
  \]
  其绝对值不超过 \(\frac12r\)，故原点处极限为 \(0=f(0,0)\)。
  故函数在全平面连续。

  \item \textbf{解析}：当 $y \neq 0$ 时由初等函数构成必连续。考察原点，令 \(x=r\cos\theta,\ y=r\sin\theta\)（其中 \(\sin\theta\neq0\) 时按原式），则
  \[
    \left|x\sin\frac1y\right|
    \leqslant r|\cos\theta|\leqslant r\to0,
  \]
  故 $(0,0)$ 处连续。而在点 $(x_0,0)$（$x_0 \neq 0$）处，当 $y \to 0$ 时 $\sin\dfrac{1}{y}$ 振荡发散，极限不存在。故函数在除 $x$ 轴上非原点外的所有点均连续，在 $\{(x,0) \mid x \neq 0\}$ 处不连续。
\end{enumerate}
\end{solution}

% ============================================================
% 七、微分学的核心概念辨析
% ============================================================

\begin{exercise}[可微性、偏导数与方向导数的关系]
设 $z=f(x,y)$ 定义于点 $P_0(a,b)$ 的某邻域内。
试总结下列诸性质之间的蕴涵关系，并对其中不可逆蕴涵者举反例说明：
\begin{shortexenum}
  \shortitem{1}{$f$ 在 $P_0$ 连续；} &
  \shortitem{2}{$f$ 在 $P_0$ 偏导数存在；} \\
  \shortitem{3}{$f$ 在 $P_0$ 可微；} &
  \shortitem{4}{$f$ 的偏导数在 $P_0$ 连续。}
\end{shortexenum}
\end{exercise}
\begin{solution}
\textbf{解析}：(4) 偏导连续 $\implies$ (3) 可微 $\implies$ (1) 连续，且 (3) 也蕴含 (2) 偏导存在。不可逆之处可用下列反例说明：
\begin{itemize}
\item (1) $\nRightarrow$ (2) 且 (1) $\nRightarrow$ (3)：取 \(\displaystyle f(x,y)=\sqrt{x^2+y^2}\)。它在 $(0,0)$ 连续，但偏导数不存在，当然也不可微。
\item (2) $\nRightarrow$ (1) 且 (2) $\nRightarrow$ (3)：取十字交叉函数 \(\displaystyle f(x,y)=\dfrac{xy}{x^2+y^2}\quad\text{（原点外）},\qquad f(0,0)=0\)。则原点偏导存在，但不连续也不可微。
\item (3) $\nRightarrow$ (4)：取 \(\displaystyle f(x,y)=(x^2+y^2)\sin\dfrac{1}{\sqrt{x^2+y^2}},\qquad f(0,0)=0\)。则原点可微（微分为 $0$），但偏导数在该点附近振荡，在原点不连续。
\end{itemize}
\end{solution}

% ============================================================
% 八、方向导数的存在性
% ============================================================

\begin{exercise}[方向导数与偏导数存在性的关系]
由下列每一条件能否断定 $f(x,y)$ 在点 $(a,b)$ 有偏导数 $f_x(a,b)$？
\begin{shortexenum}
  \shortitem{1}{$f(x,y)$ 在点 $(a,b)$ 沿任何方向的方向导数存在且相等；} &
  \shortitem{2}{$f(x,y)$ 在点 $(a,b)$ 沿 $x$ 轴正向和负向的方向导数存在且相等。}
\end{shortexenum}
\end{exercise}
\begin{solution}
\begin{enumerate}
  \item \textbf{解析}：能断定，因为偏导数 $f_x(a,b)$ 就是沿 $x$ 轴正方向单位向量 $(1,0)$ 的方向导数，即 $f_x(a,b)=D_{(1,0)}f(a,b)$。题设说沿任何方向的方向导数都存在且相等，于是沿 $(1,0)$ 的方向导数当然也存在，从而 $f_x(a,b)$ 存在。
  \item \textbf{解析}：不能断定。若沿 $x$ 轴正方向、负方向的方向导数都等于 $A$，则 \(\displaystyle \lim_{t\to0^+}\frac{f(a+t,b)-f(a,b)}{t}=A,\quad \lim_{t\to0^+}\frac{f(a-t,b)-f(a,b)}{t}=A\)，但偏导数左极限应写成 \(\displaystyle \lim_{t\to0^+}\frac{f(a-t,b)-f(a,b)}{-t}=-A\)，
  因而只有当 $A=0$ 时左右极限才可能相等。例如 $f(x,y)=|x-a|$ 在点 $(a,b)$ 两侧方向导数都等于 $1$，但偏导数并不存在。
\end{enumerate}
\end{solution}

% ============================================================
% 九、隐函数求导恒等式
% ============================================================

\begin{exercise}[隐函数求导公式的验证]
设 $z=f(x,y)$ 满足方程 $x-az=\varphi(y-bz)$，其中 $\varphi$ 可微，$a,b$ 为常数。求证：$a\dfrac{\partial z}{\partial x}+b\dfrac{\partial z}{\partial y}=1$。
\end{exercise}
\begin{solution}
令 $F(x,y,z) = x - az - \varphi(y-bz) = 0$。由隐函数求导公式，
\begin{align*}
F_x&=1,\qquad F_y=-\varphi',\qquad F_z=-a+b\varphi',\\
\frac{\partial z}{\partial x} &=-\frac{F_x}{F_z} = \frac{1}{a-b\varphi'},\qquad
\frac{\partial z}{\partial y} = -\frac{F_y}{F_z} = -\frac{\varphi'}{a-b\varphi'}.
\end{align*}
因此 \(\displaystyle a\frac{\partial z}{\partial x} + b\frac{\partial z}{\partial y}
= a\left(\frac{1}{a-b\varphi'}\right) + b\left(-\frac{\varphi'}{a-b\varphi'}\right)
= \frac{a-b\varphi'}{a-b\varphi'} = 1\)。
恒等式成立。
\end{solution}

% ============================================================
% 十、齐次函数与 Euler 定理
% ============================================================

\begin{exercise}[Euler 齐次函数定理]
证明函数 $u=x^k F\!\left(\dfrac{z}{x},\,\dfrac{y}{x}\right)$ 满足方程 $x\dfrac{\partial u}{\partial x}+y\dfrac{\partial u}{\partial y}+z\dfrac{\partial u}{\partial z}=ku$，其中 $k$ 为常数。
\end{exercise}
\begin{solution}
令 $s = \dfrac{z}{x},\ t = \dfrac{y}{x}$，则 $u = x^k F(s,t)$。由链式法则，\(\displaystyle u_x = kx^{k-1}F - x^{k-2}z F_s - x^{k-2}y F_t\)，\(\displaystyle u_y = x^{k-1} F_t\)，\(\displaystyle u_z = x^{k-1} F_s\)。于是
\[
x u_x + y u_y + z u_z
= x(kx^{k-1}F - x^{k-2}z F_s - x^{k-2}y F_t) + y(x^{k-1} F_t) + z(x^{k-1} F_s)
= kx^kF = ku.
\]
恒等式成立。
\end{solution}

% ============================================================
% 十一、高阶偏导数的链式法则
% ============================================================

\begin{exercise}[混合偏导数的链式法则]
设 $u=yf\!\left(\dfrac{y}{x}\right)+xg\!\left(\dfrac{y}{x}\right)$，其中 $f,g$ 具有一阶连续偏导数。求 $x\dfrac{\partial^2 u}{\partial x^2}+y\dfrac{\partial^2 u}{\partial x\,\partial y}$。
\end{exercise}
\begin{solution}
对任意非零常数 $t$，有
\[
u(tx,ty)=tyf(y/x)+txg(y/x)
=t\left[yf(y/x)+xg(y/x)\right]=tu(x,y),
\]
故 $u$ 是一次齐次函数。由 Euler 定理，\(\displaystyle xu_x+yu_y=u\)。
再对上式关于 $x$ 求偏导，得 \(\displaystyle \frac{\partial}{\partial x}(xu_x + yu_y)=u_x\)，即
\[
u_x+xu_{xx}+yu_{yx}=u_x,\qquad
x \frac{\partial^2 u}{\partial x^2} + y \frac{\partial^2 u}{\partial x \partial y}=0,
\]
故结果为 $\mathbf{0}$。
\end{solution}

% ============================================================
% 十二、全导数（多层复合函数）
% ============================================================

\begin{exercise}[全导数的树状链式法则]
设 $w=F(x,y,z)$, $z=f(x,y)$, $y=\varphi(x)$。求 $\dfrac{\mathrm{d}w}{\mathrm{d}x}$。
\end{exercise}
\begin{solution}
这里 $w=F(x,y,z)$，其中 $z=f(x,y)$ 且 $y=\varphi(x)$，所以所有变量最终都依赖于 $x$。由链式法则逐层求导，
\begin{align*}
\frac{\mathrm{d}w}{\mathrm{d}x} &= \frac{\partial F}{\partial x} \frac{\mathrm{d}x}{\mathrm{d}x} + \frac{\partial F}{\partial y} \frac{\mathrm{d}y}{\mathrm{d}x} + \frac{\partial F}{\partial z} \frac{\mathrm{d}z}{\mathrm{d}x} \\[4pt]
&= F_x + F_y \varphi'(x) + F_z \left[ \frac{\partial f}{\partial x} + \frac{\partial f}{\partial y}\frac{\mathrm{d}y}{\mathrm{d}x} \right] \\[4pt]
&= F_x + F_y \varphi'(x) + F_z \left[ f_x + f_y \varphi'(x) \right].
\end{align*}
\end{solution}

% ============================================================
% 十三、变系数 PDE
% ============================================================

\begin{exercise}[通过 PDE 确定函数形式]
设函数 $f(u)$ 具有二阶连续导数，而 $z=f(\mathrm{e}^x\sin y)$ 满足方程 $\dfrac{\partial^2 z}{\partial x^2}+\dfrac{\partial^2 z}{\partial y^2}=\mathrm{e}^{2x}z$。求 $f(u)$。
\end{exercise}
\begin{solution}
\begin{enumerate}
  \item \textbf{令中间变量求导}：令 $u = \mathrm{e}^x\sin y$，则 $z = f(u)$。
  \begin{itemize}
    \item $z_x = f'(u) \cdot u_x = f'(u) \mathrm{e}^x\sin y = u f'(u)$
    \item $z_y = f'(u) \cdot u_y = f'(u) \mathrm{e}^x\cos y$
  \end{itemize}

  \item \textbf{计算二阶偏导数}：
  \begin{itemize}
    \item $z_{xx} = \dfrac{\partial}{\partial x}(u f'(u)) = f'(u) u_x + u f''(u) u_x = u f'(u) + u^2 f''(u)$
    \item $z_{yy} = \dfrac{\partial}{\partial y}(f'(u) \mathrm{e}^x\cos y) = f''(u) (\mathrm{e}^x\cos y)^2 + f'(u) (-\mathrm{e}^x\sin y) = f''(u) \mathrm{e}^{2x}\cos^2 y - u f'(u)$
  \end{itemize}

  \item \textbf{代入 PDE 方程化简}：
  将求得的二阶偏导相加：
  \[
  z_{xx} + z_{yy}=u^2 f''(u) + f''(u) \mathrm{e}^{2x}\cos^2 y
  = f''(u)(\mathrm{e}^{2x}\sin^2 y + \mathrm{e}^{2x}\cos^2 y)
  = f''(u)\mathrm{e}^{2x}.
  \]
  根据题设条件，上式等于 $\mathrm{e}^{2x}z = \mathrm{e}^{2x}f(u)$。

  \item \textbf{解常微分方程}：
  消去正系数 $\mathrm{e}^{2x}$，得到 $f''(u) = f(u)$，即常系数线性齐次微分方程 $f''(u) - f(u) = 0$。
  其特征方程 $\lambda^2 - 1 = 0 \implies \lambda = \pm 1$。
  故通解为 $f(u) = \mathbf{C_1 \mathrm{e}^u + C_2 \mathrm{e}^{-u}}$（$C_1, C_2$ 为任意常数）。
\end{enumerate}
\end{solution}

% ============================================================
% 十四、隐函数组求全导数
% ============================================================

\begin{exercise}[隐函数组的全导数]
设函数 $u=f(x,y,z)$ 有连续偏导数，$y=y(x)$ 和 $z=z(x)$ 分别由方程 $\mathrm{e}^{xy}=y$ 和 $\mathrm{e}^z=xz$ 确定。求 $\dfrac{\mathrm{d}u}{\mathrm{d}x}$。
\end{exercise}
\begin{solution}
\begin{enumerate}
  \item \textbf{通过隐函数求导求 $y'$ 和 $z'$}：
  \begin{itemize}
    \item 对 $\mathrm{e}^{xy}=y$ 两边关于 $x$ 求导：$\mathrm{e}^{xy}(y + x y') = y' \implies y' = \dfrac{y\mathrm{e}^{xy}}{1-x\mathrm{e}^{xy}} = \dfrac{y^2}{1-xy}$（利用了 $\mathrm{e}^{xy}=y$ 替换）。
    \item 对 $\mathrm{e}^z=xz$ 两边关于 $x$ 求导：$\mathrm{e}^z z' = z + x z' \implies z' = \dfrac{z}{\mathrm{e}^z-x} = \dfrac{z}{xz-x} = \dfrac{z}{x(z-1)}$（利用了 $\mathrm{e}^z=xz$ 替换）。
  \end{itemize}

  \item \textbf{应用全导数公式}：
  利用多元函数复合求导法则展开 \(\displaystyle \frac{\mathrm{d}u}{\mathrm{d}x} = f_x + f_y y' + f_z z' = f_x + f_y \left( \frac{y^2}{1-xy} \right) + f_z \left( \frac{z}{x(z-1)} \right)\)。
  这就是所求的导数表达式。
\end{enumerate}
\end{solution}

% ============================================================
% 十五、参数方程的二阶偏导数
% ============================================================

\begin{exercise}[柱坐标变换下的二阶偏导数]
求由方程组
\[
\begin{cases}
x=u\cos v,\\[4pt]
y=u\sin v,\\[4pt]
z=v
\end{cases}
\]
所确定的函数 $z=z(x,y)$ 的所有二阶偏导数。
\end{exercise}
\begin{solution}
\begin{enumerate}
  \item \textbf{消去参数得到显函数}：
  由 $x=u\cos v,\ y=u\sin v$ 可推得 $\tan v = \dfrac{y}{x}$，于是 $v = \arctan\dfrac{y}{x}$，故 \(z = \arctan\dfrac{y}{x}\)。
  \item \textbf{计算一阶偏导数}：
  \(\displaystyle z_x = -\frac{y}{x^2+y^2},\qquad z_y = \frac{x}{x^2+y^2}\)。
  \item \textbf{计算二阶偏导数}：再由除法求导法则可得
  \(\displaystyle z_{xx} = \frac{2xy}{(x^2+y^2)^2},\qquad
  z_{yy} = -\frac{2xy}{(x^2+y^2)^2},\qquad
  z_{xy} = \frac{y^2-x^2}{(x^2+y^2)^2}\)。
\end{enumerate}
\end{solution}

% ============================================================
% 十六、方向导数的具体计算
% ============================================================

\begin{exercise}[方向导数的基本计算]
设 $u=\left(\dfrac{y}{x}\right)^{\!z}$，向量 $\boldsymbol l$ 与向量 $\{2,1,-1\}$ 同向。求 $\left.\dfrac{\partial u}{\partial \boldsymbol l}\right|_{(1,1,1)}$。
\end{exercise}
\begin{solution}
\begin{enumerate}
  \item \textbf{求该点处的梯度}：计算偏导并代入 $(1,1,1)$，得
  \(\displaystyle u_x = -1,\qquad u_y = 1,\qquad u_z = 0\)。
  故在该点处梯度为 $\nabla u = (-1, 1, 0)$。
  \item \textbf{标准化方向向量}：
  与 $\{2,1,-1\}$ 同向的单位向量为
  \[
  \boldsymbol l^\circ = \dfrac{1}{\sqrt{4+1+1}}(2, 1, -1)
  = \left(\dfrac{2}{\sqrt{6}}, \dfrac{1}{\sqrt{6}}, -\dfrac{1}{\sqrt{6}}\right).
  \]
  \item \textbf{计算方向导数}：利用公式 $D_{\boldsymbol l} u = \nabla u \cdot \boldsymbol l^\circ$，
  \[
  \left.\frac{\partial u}{\partial \boldsymbol l}\right|_{(1,1,1)}
  = -1\left(\frac{2}{\sqrt{6}}\right)
  + 1\left(\frac{1}{\sqrt{6}}\right)
  + 0\left(-\frac{1}{\sqrt{6}}\right)
  =-\frac{\sqrt{6}}{6}.
  \]
\end{enumerate}
\end{solution}

% ============================================================
% 十七、梯度与方向导数
% ============================================================

\begin{exercise}[梯度与方向导数的联合计算]
求 $u=xy+yz^3$ 在点 $M(2,-1,1)$ 处的梯度和在向量 $\boldsymbol l=\{2,2,-1\}$ 方向的方向导数。
\end{exercise}
\begin{solution}
\begin{enumerate}
  \item \textbf{求梯度}：分别对 $x,y,z$ 求导并代入点 $M(2,-1,1)$，得 $u_x=-1,\ u_y=3,\ u_z=-3$，故梯度为 $\mathbf{\mathrm{grad}\,u = (-1, 3, -3)}$。
  \item \textbf{标准化方向向量}：单位方向向量 $\boldsymbol l^\circ = \dfrac{(2,2,-1)}{\sqrt{4+4+1}} = \left(\dfrac{2}{3}, \dfrac{2}{3}, -\dfrac{1}{3}\right)$。
  \item \textbf{计算方向导数}：作点积：
  \[
  \frac{\partial u}{\partial \boldsymbol l} = \nabla u \cdot \boldsymbol l^\circ = (-1)\left(\frac{2}{3}\right) + 3\left(\frac{2}{3}\right) + (-3)\left(-\frac{1}{3}\right) = -\frac{2}{3} + 2 + 1 = \frac{7}{3}.
  \]
\end{enumerate}
\end{solution}

% ============================================================
% 十八、Laplace 方程的不变性
% ============================================================

\begin{exercise}[Laplace 方程的坐标不变性]
设 $f(\xi,\eta)$ 具有连续的二阶偏导数，且满足 Laplace 方程 $\dfrac{\partial^2 f}{\partial\xi^2}+\dfrac{\partial^2 f}{\partial\eta^2}=0$。证明：函数 $z=f(x^2-y^2,\;2xy)$ 也满足 Laplace 方程 $\dfrac{\partial^2 z}{\partial x^2}+\dfrac{\partial^2 z}{\partial y^2}=0$。
\end{exercise}
\begin{solution}
\begin{enumerate}
  \item \textbf{引入中间变量}：令 $\xi = x^2-y^2, \eta = 2xy$。
  \item \textbf{求一阶与二阶偏导数}：由链式法则，
  \begin{align*}
  z_x &= 2x f_\xi + 2y f_\eta, & z_{xx} &= 2f_\xi + 4x^2 f_{\xi\xi} + 8xy f_{\xi\eta} + 4y^2 f_{\eta\eta},\\
  z_y &= -2y f_\xi + 2x f_\eta, & z_{yy} &= -2f_\xi + 4y^2 f_{\xi\xi} - 8xy f_{\xi\eta} + 4x^2 f_{\eta\eta}.
  \end{align*}
  \item \textbf{证明等式成立}：将两式相加，交叉项 $8xy f_{\xi\eta}$ 和一阶项 $2f_\xi$ 刚好抵消，得 \(\displaystyle z_{xx}+z_{yy}=4(x^2+y^2)f_{\xi\xi}+4(x^2+y^2)f_{\eta\eta}=4(x^2+y^2)(f_{\xi\xi}+f_{\eta\eta})\)。
  因原题设 $f_{\xi\xi} + f_{\eta\eta} = 0$，故上式为 $0$。命题得证。
\end{enumerate}
\end{solution}

% ============================================================
% 十九、无约束极值问题
% ============================================================

\begin{exercise}[多元函数的极值计算]
求下列函数的极值：
\begin{shortexenum}
  \shortitem{1}{$z=3axy-x^3-y^3\;(a>0)$；} &
  \shortitem{2}{$z=\sin x+\cos y+\cos(x-y)$，其中 $0\leqslant x,y\leqslant\dfrac{\pi}{2}$。}
\end{shortexenum}
\end{exercise}
\begin{solution}
\begin{enumerate}
  \item \textbf{求解题一}：令 $z_x = 3ay - 3x^2 = 0,\ z_y = 3ax - 3y^2 = 0$。由 $y = x^2/a$ 代入第二式得 $x^4/a^2 = ax \implies x(x^3-a^3)=0$，故驻点为 $(0,0)$ 和 $(a,a)$。又 $A = -6x,\ B = 3a,\ C = -6y$，于是 $(0,0)$ 处有 $AC-B^2 = -9a^2 < 0$，为鞍点；$(a,a)$ 处有 $AC-B^2 = 27a^2 > 0$ 且 $A = -6a < 0$，故在此处取得极大值，且 $z(a,a) = a^3$。
  \item \textbf{求解题二}：驻点满足
  \(\displaystyle z_x = \cos x - \sin(x-y) = 0,\qquad z_y = -\sin y + \sin(x-y) = 0\)。
  在 $0\leqslant x,y\leqslant \dfrac{\pi}{2}$ 内，由 $\sin y = \sin(x-y)$ 得 $x=2y$，再由 $\cos(2y)=\sin y$ 得 $y=\dfrac{\pi}{6},\ x=\dfrac{\pi}{3}$。此时 \(\displaystyle A = -\sin x - \cos(x-y) = -\sqrt{3},\quad B = \cos(x-y) = \frac{\sqrt{3}}{2},\quad C = -\cos y - \cos(x-y) = -\sqrt{3}\)。
  因而 $AC-B^2 = \dfrac{9}{4} > 0$ 且 $A < 0$，故取得极大值 $\dfrac{3\sqrt{3}}{2}$。
\end{enumerate}
\end{solution}

% ============================================================
% 二十、隐函数极值问题
% ============================================================

\begin{exercise}[隐函数确定的极值]
求由下列方程所确定的隐函数 $z=z(x,y)$ 的极值：
\begin{shortexenum}
  \shortitem{1}{$x^2+xyz-x^2y-xy^2-9=0$；} &
  \shortitem{2}{$x^2+y^2+z^2-2x+4y-6z-11=0$。}
\end{shortexenum}
\end{exercise}
\begin{solution}
\begin{enumerate}
  \item \textbf{求解题一}：
  \begin{enumerate}
    \item 设 $F = x^2+xyz-x^2y-xy^2-9=0$。极值点满足 $z_x = 0, z_y = 0$，即 $F_x = 0, F_y = 0$。
    \item 求偏导方程：
    \begin{itemize}
      \item $F_x = 2x + yz - 2xy - y^2 = 0$
      \item $F_y = xz - x^2 - 2xy = 0 \implies x(z-x-2y) = 0$
    \end{itemize}
    \item 解方程组：由原方程知 $x \neq 0$（否则 $-9=0$ 矛盾）。故 $z = x+2y$。
    代入 $F_x=0$ 得 $2x+y(x+2y)-2xy-y^2 = 2x-xy+y^2=0 \implies x = \dfrac{y^2}{y-2}$。
    代入原方程 $F=0$：$x^2 + xy(x+2y) - x^2y - xy^2 - 9 = 0 \implies x^2+xy^2-9=0$。
    将其联立化简得 $x^2(y-1)=9$。将 $x$ 替换，经过代数运算可解得极值点坐标 $(x,y)$ 及其对应的 $z$ 值。
  \end{enumerate}

  \item \textbf{求解题二}：
  \begin{enumerate}
    \item 配方得：$(x-1)^2 + (y+2)^2 + (z-3)^2 = 25$。这是几何中心在 $(1,-2,3)$、半径为 $5$ 的球面。
    \item 利用几何意义：球面上的点到 $z=0$ 平面的最高点和最低点显然就是其 $z$ 坐标达到极限的地方。
    \item 得出结论：当 $x=1, y=-2$ 时，取最大值 $z = 3+5 = 8$ 和最小值 $z = 3-5 = -2$。分别对应极大值和极小值。
  \end{enumerate}
\end{enumerate}
\end{solution}

% ============================================================
% 二十一、实际最优化问题
% ============================================================

\begin{exercise}[植物生长率的优化 —— 二元二次模型]
生物学家发现某种植物的生长率 $g$ 是温度 $x$ 和湿度 $y$ 的函数。根据大量实验数据，生物学家认为温室内的生长率由 \(\displaystyle g(x,y)=-2x^2+196x+3y^2+242y+2xy-16360\) 给出，其中约束区域为 $75\leqslant x\leqslant 85$, $50\leqslant y\leqslant 75$。

试求使植物保持最大生长率的温度和湿度，并求此最大生长率。
\end{exercise}
\begin{solution}
\begin{enumerate}
  \item \textbf{求内部驻点}：偏导数满足 $g_x = -4x + 2y + 196 = 0,\ g_y = 2x + 6y + 242 = 0$。联立解得驻点位于考察区域外部，故极值只能在边界上取得。
  \item \textbf{利用单调性}：对所有 $x>0,\ y>0$，都有 $g_y = 6y + 2x + 242 > 0$，故 $g(x,y)$ 随 $y$ 单调递增，因此最大值必取在上边界 $y=75$。
  \item \textbf{转化为一元极值问题}：令 $y=75$，得 $h(x) = g(x,75) = -2x^2 + 346x + \text{常数}$，于是 $h'(x) = -4x + 346 = 0 \implies x = 86.5$。由于抛物线开口向下且 $x \leqslant 85$，故在区间 $[75,85]$ 上最大值取于 $x=85$。
  \item \textbf{代入计算最大值}：将 $x=85,\ y=75$ 代入原方程，可算出最大值为 $\mathbf{33625}$。
\end{enumerate}
\end{solution}

% ============================================================
% 二十二、点到曲线的距离优化
% ============================================================

\begin{exercise}[点到空间曲线的最短距离]
求曲线
\[
\begin{cases}
z=x^2+y^2,\\[4pt]
y=\dfrac{1}{x},
\end{cases}
\]
上到 $Oxy$ 平面距离最近的点。
\end{exercise}
\begin{solution}
\begin{enumerate}
  \item \textbf{构建目标函数}：点到 $Oxy$ 平面的距离是 $|z|$。因 $z = x^2+y^2 \geqslant 0$，故只需最小化
  \(\displaystyle z = x^2 + \frac{1}{x^2} \geqslant 2\sqrt{x^2 \cdot \frac{1}{x^2}} = 2\)。
  等号成立当且仅当 $x^2 = \dfrac{1}{x^2}$，即 $x^4 = 1 \implies x = \pm 1$。于是当 $x=1$ 时，$y=1,\ z=2$；当 $x=-1$ 时，$y=-1,\ z=2$，故距离最近的点为 $\mathbf{(1,1,2)}$ 和 $\mathbf{(-1,-1,2)}$。
\end{enumerate}
\end{solution}

% ============================================================
% 二十三、椭圆面积优化
% ============================================================

\begin{exercise}[椭圆面积的最小化]
求 $a,b$ 的值，使得包含圆 $(x-1)^2+y^2=1$ 在其内部的椭圆 \(\displaystyle \frac{x^2}{a^2}+\frac{y^2}{b^2}=1\ (a>0,\;b>0,\;a \neq b)\) 有最小的面积。
\end{exercise}
\begin{solution}
\begin{enumerate}
  \item \textbf{建立相切模型}：
  椭圆内含该圆，面积要最小，必然使得椭圆与圆相内切。在第一象限的切点处，椭圆与圆有相同的斜率和坐标。
  \item \textbf{利用切线方程和交点条件}：
  对椭圆和圆隐函数求导，椭圆斜率 $y' = -\dfrac{b^2 x}{a^2 y}$，圆斜率 $y' = -\dfrac{x-1}{y}$。
  相切使得两者相等：$\dfrac{b^2 x}{a^2 y} = \dfrac{x-1}{y} \implies b^2 x = a^2(x-1)$。
  交点须满足圆系方程：$y^2 = 2x-x^2$。代入椭圆消元后可得相切时的关系式：$a^2 = \dfrac{b^4}{b^2-1}$。
  \item \textbf{最优化面积目标}：
  椭圆面积为 $S = \pi ab$。最优化 $S^2$ 等价于最优化 $f(b) = a^2b^2 = \dfrac{b^6}{b^2-1}$。
  对 $b^2$ 求导并令其为 $0$，解得 $b^2 = \dfrac{3}{2}$，从而求出 $b = \dfrac{\sqrt{6}}{2}$，进而代回求得 $a = \dfrac{3\sqrt{2}}{2}$。
\end{enumerate}
\end{solution}

% ============================================================
% 二十四、AM--GM 不等式的优化证明
% ============================================================

\begin{exercise}[AM--GM 不等式的 Lagrange 乘子法证明]
求函数 $z=\dfrac{1}{2}(x^n+y^n)$ 在条件 $x+y=l\;(l>0,\;n\geqslant 1)$ 下的极值，
并由此证明：当 $a\geqslant 0$, $b\geqslant 0$, $n\geqslant 1$ 时，
\(\displaystyle \left(\frac{a+b}{2}\right)^{\!n}\!\leqslant\!\frac{a^n+b^n}{2}\)。
\end{exercise}
\begin{solution}
\begin{enumerate}
  \item \textbf{构造拉格朗日函数}：令 $L(x,y,\lambda) = \dfrac{1}{2}(x^n+y^n) - \lambda(x+y-l)$。
  \item \textbf{求解临界点}：由 \(\displaystyle L_x = \frac{n}{2}x^{n-1} - \lambda = 0,\qquad L_y = \frac{n}{2}y^{n-1} - \lambda = 0,\qquad L_\lambda = -(x+y-l) = 0\)
  得 $x^{n-1}=y^{n-1}$，故 $x=y$，再代入约束得极值点 $x=y=\dfrac{l}{2}$。
  \item \textbf{证明不等式}：由 $x^n$ 的凸性知此处取得极小值，且最小值为 \(\displaystyle \frac{1}{2}\left[\left(\frac{l}{2}\right)^n+\left(\frac{l}{2}\right)^n\right]=\left(\frac{l}{2}\right)^n\)。
  将任意满足 $a+b=l$ 的非负实数代入，即得 $\dfrac{a^n+b^n}{2} \geqslant \left(\dfrac{a+b}{2}\right)^n$。
\end{enumerate}
\end{solution}

% ============================================================
% 二十五、不等式证明
% ============================================================

\begin{exercise}[含绝对值项的不等式证明]
设 $x,y,z$ 为实数，且 $\mathrm{e}^x+y^2+|z|=3$。求证：$\mathrm{e}^x\,y^2\,|z|\leqslant 1$。
\end{exercise}
\begin{solution}
\begin{enumerate}
  \item \textbf{变量替换简化模型}：令 $u=\mathrm{e}^x,\ v=y^2,\ w=|z|$，则 $u+v+w=3$ 且 $u>0,\ v\geqslant 0,\ w\geqslant 0$，目标化为证明 $uvw \leqslant 1$。
  \item \textbf{应用均值不等式（AM-GM）}：由算术-几何平均值不等式，\(\displaystyle uvw \leqslant \left(\frac{u+v+w}{3}\right)^3 = 1\)。
  \item \textbf{得出结论}：故 $\mathrm{e}^x\,y^2\,|z| \leqslant 1$ 恒成立；当且仅当 $u=v=w=1$，即 $x=0,\ y=\pm 1,\ z=\pm 1$ 时等号成立。
\end{enumerate}
\end{solution}

% ============================================================
% 二十六、球面上的对数最值
% ============================================================

\begin{exercise}[对数函数在球面上的极值与 AM--GM 推广]
当 $x>0$, $y>0$, $z>0$ 时，求函数 $t(x,y,z)=\ln x+\ln y+3\ln z$ 在球面 $x^2+y^2+z^2=5R^2$ 上的最大值，并由此证明：当 $a>0$, $b>0$, $c>0$ 时，\(\displaystyle abc^3\leqslant 27\left(\frac{a+b+c}{5}\right)^{\!5}\)。
\end{exercise}
\begin{solution}
\begin{enumerate}
  \item \textbf{使用 Lagrange 乘子法}：令 $L = \ln x + \ln y + 3\ln z - \lambda(x^2+y^2+z^2-5R^2)$。由 \(\displaystyle \frac{1}{x}-2\lambda x=0,\quad \frac{1}{y}-2\lambda y=0,\quad \frac{3}{z}-2\lambda z=0\)
  得 $x^2 = \dfrac{1}{2\lambda},\ y^2 = \dfrac{1}{2\lambda},\ z^2 = \dfrac{3}{2\lambda}$。
  \item \textbf{求解极大值坐标和对应的最大值}：由 $x^2+y^2+z^2 = \dfrac{5}{2\lambda} = 5R^2$ 得 $\lambda = \dfrac{1}{2R^2}$，再结合 $x,y,z>0$ 可知极值点为 $x=R,\ y=R,\ z=\sqrt{3}R$，最大值为 \(\displaystyle \ln(R \cdot R \cdot (\sqrt{3}R)^3) = \ln(3\sqrt{3}R^5)\)。
  因而球面上恒有 $\ln(xyz^3) \leqslant \ln(3\sqrt{3}R^5)$，即 $x y z^3 \leqslant 3\sqrt{3}R^5$，进而 $x^2 y^2 z^6 \leqslant 27R^{10}$。
  \item \textbf{推广至不等式}：对任意正数 $a,b,c$，令 $x^2=a,\ y^2=b,\ z^2=c$，并把 $5R^2$ 换成 $a+b+c$，即可得到所求不等式。
\end{enumerate}
\end{solution}

% ============================================================
% 二十七、圆柱螺旋线的几何性质
% ============================================================

\begin{exercise}[圆柱螺旋线的切线几何性质]
证明：圆柱螺旋线 $x=a\cos t,\quad y=a\sin t,\quad z=bt$ 的切线与 $z$ 轴的夹角为常数。
\end{exercise}
\begin{solution}
对位置向量 $\boldsymbol r(t) = (a\cos t, a\sin t, bt)$ 求导得切向量
\(\displaystyle \boldsymbol r'(t) = (-a\sin t, a\cos t, b)\)，从而
\(\displaystyle |\boldsymbol r'(t)| = \sqrt{(-a\sin t)^2 + (a\cos t)^2 + b^2} = \sqrt{a^2+b^2}\)。
设 $z$ 轴单位方向向量为 $\boldsymbol k=(0,0,1)$，则切线与 $z$ 轴夹角 $\theta$ 满足
\(\displaystyle \cos\theta = \frac{\boldsymbol r'(t) \cdot \boldsymbol k}{|\boldsymbol r'(t)| |\boldsymbol k|} = \frac{b}{\sqrt{a^2+b^2}}\)，
右端为常数，故切线与 $z$ 轴的夹角为常数。
\end{solution}

% ============================================================
% 二十八、球面上具有指定法向量的曲线
% ============================================================

\begin{exercise}[球面上的定向曲线构造]
在球面 $x^2+y^2+z^2=R^2$ 上求一条曲线，使其上每一点的法线与平面 $2x+3y+6z-1=0$ 的夹角为 $30^\circ$。
\end{exercise}
\begin{solution}
\begin{enumerate}
  \item \textbf{提取法向量并设定夹角条件}：球面上任一点 $(x,y,z)$ 的法向量是径向向量 $\boldsymbol n_1=(x,y,z)$，给定平面的法向量为 $\boldsymbol n_2=(2,3,6)$，且两者夹角满足 $\cos\theta = \cos 30^\circ = \dfrac{\sqrt{3}}{2}$。
  \item \textbf{通过点积公式列出方程}：利用向量点积公式，
  \(\displaystyle \frac{|\boldsymbol n_1 \cdot \boldsymbol n_2|}{|\boldsymbol n_1| |\boldsymbol n_2|} = \frac{|2x+3y+6z|}{R \cdot \sqrt{4+9+36}} = \frac{|2x+3y+6z|}{7R} = \frac{\sqrt{3}}{2}\)。
  \item \textbf{得到交线所满足的方程组}：因此所求点同时满足球面方程和
  \[
  \begin{cases}
    x^2+y^2+z^2=R^2 \\[4pt]
    2x+3y+6z=\pm \dfrac{7\sqrt{3}R}{2}
  \end{cases}
  \]
  所求曲线即为该球面被这两个平行平面所截得的两条圆周。
\end{enumerate}
\end{solution}

% ============================================================
% 二十九、等位面的法向量性质（复习巩固）
% ============================================================

\begin{exercise}[梯度作为等位面的法向量 —— 几何视角的严格证明]
证明：函数 $F(x,y,z)$ 在点 $P_0(x_0,y_0,z_0)$ 的梯度向量是函数 $F(x,y,z)$ 在点 $P_0$ 的等位面的法向量。
\end{exercise}
\begin{solution}
设过点 $P_0$ 的等位面 $F(x,y,z)=C$ 上任意一条光滑曲线参数方程为 $x=x(t),\ y=y(t),\ z=z(t)$，且在 $t=t_0$ 时对应点 $P_0$。沿该曲线恒有 $F(x(t),y(t),z(t))=C$，故对 $t$ 求全导数得
\(\displaystyle F_x \cdot x'(t) + F_y \cdot y'(t) + F_z \cdot z'(t) = 0\)。
在 $t=t_0$ 时，这可写成
\(\displaystyle \nabla F(P_0) \cdot \boldsymbol r'(t_0) = 0\)。
其中 $\boldsymbol r'(t_0)$ 是任意过 $P_0$ 的切向量，因此梯度向量与等位面上一切切向量都垂直，故 $\nabla F(P_0)$ 就是该等位面的法向量。
\end{solution}

% ============================================================
% 三十、平行切平面的确定
% ============================================================

\begin{exercise}[椭球面的平行切平面族]
求曲面 $x^2+2y^2+3z^2=21$ 的所有平行于平面 $x+4y+6z=0$ 的切平面。
\end{exercise}
\begin{solution}
\begin{enumerate}
  \item \textbf{利用平行条件确定切点关系}：曲面的法向量 $\boldsymbol N = (2x, 4y, 6z)$ 必须与平面法向量 $\boldsymbol n = (1, 4, 6)$ 共线。设 $(2x,4y,6z)=k(1,4,6)$，则
  \(\displaystyle x = \frac{k}{2},\qquad y = k,\qquad z = k\)。
  \item \textbf{代回曲面方程求解 $k$}：代入 $x^2+2y^2+3z^2=21$ 得
  \(\displaystyle \left(\frac{k}{2}\right)^2 + 2k^2 + 3k^2 = \frac{21}{4}k^2 = 21 \implies k^2 = 4 \implies k = \pm 2\)，
  故切点为 $(1,2,2)$ 和 $(-1,-2,-2)$。
  \item \textbf{写出切平面方程}：对应切平面分别为 $\mathbf{x+4y+6z-21=0}$ 和 $\mathbf{x+4y+6z+21=0}$。
\end{enumerate}
\end{solution}

\begin{exercise}[过直线的切平面]
求椭球面 $x^2+2y^2+3z^2=21$ 上某处的切平面 $\Pi$，使平面 $\Pi$ 过已知直线 \(\displaystyle L:\frac{x-6}{2}=\frac{y-3}{1}=\frac{2z-2}{1}\)。
\end{exercise}
\begin{solution}
\begin{enumerate}
  \item \textbf{建立切平面的一般方程}：
  设切点为 $P_0(x_0, y_0, z_0)$，则切平面 $\Pi$ 的方程为 \(\displaystyle x_0 x + 2y_0 y + 3z_0 z = 21\)。
  \item \textbf{利用直线代入条件}：
  直线 $L$ 过点 $A(6,3,1)$（对应 $t=0$）且过另一点，比如取参数得另一个特例点（此处将导致代数矛盾，因为该直线实际穿过或处于椭球包络范围使得纯粹代数可能无解）。不妨利用特例点 $B(0,0,-0.5)$ 代入：
  \begin{itemize}
    \item 点 $A$ 代入：$6x_0 + 6y_0 + 3z_0 = 21 \implies 2x_0+2y_0+z_0=7$。
    \item 点 $B$ 代入：$3z_0(-0.5) = 21 \implies z_0 = -14$。
  \end{itemize}
  \item \textbf{发现代数矛盾}：
  若 $z_0 = -14$，此切点须在椭球面上，满足 $x_0^2 + 2y_0^2 + 3(-14)^2 = x_0^2 + 2y_0^2 + 588 = 21$。
  这显然无实数解，表明给定的这条直线在椭球面的“视野”内与曲面直接相交，无法找到从外部贴合包覆此直线的实切平面。故此题\textbf{无实数解切平面}。
\end{enumerate}
\end{solution}

% ============================================================
% 三十二、方向导数的几何应用
% ============================================================

\begin{exercise}[外法向量的方向导数]
给定曲面 $2x^2+3y^2+z^2=6$ 在点 $P(1,1,1)$ 处指向外侧的法向量 $\boldsymbol n$。求函数 $u=\dfrac{\sqrt{6x^2+8y^2}}{z}$ 在点 $P$ 处沿方向 $\boldsymbol n$ 的方向导数。
\end{exercise}
\begin{solution}
\begin{enumerate}
  \item \textbf{求方向向量及其标准化}：
  设 $F = 2x^2+3y^2+z^2-6$，其外法向量 $\boldsymbol n = (4x, 6y, 2z)|_P = (4, 6, 2)$。单位化得 \(\displaystyle \hat{\boldsymbol n} = \frac{(4, 6, 2)}{\sqrt{16+36+4}} = \left(\frac{2}{\sqrt{14}}, \frac{3}{\sqrt{14}}, \frac{1}{\sqrt{14}}\right)\)。
  \item \textbf{求目标函数梯度}：在 $P(1,1,1)$ 处，\(\displaystyle u_x = \frac{6}{\sqrt{14}},\qquad u_y = \frac{8}{\sqrt{14}},\qquad u_z = -\sqrt{14}\)。
  \item \textbf{计算点积求方向导数}：
  \[
  D_{\hat{\boldsymbol n}} u = \left(\frac{6}{\sqrt{14}}, \frac{8}{\sqrt{14}}, -\sqrt{14}\right) \cdot \left(\frac{2}{\sqrt{14}}, \frac{3}{\sqrt{14}}, \frac{1}{\sqrt{14}}\right) = \frac{12}{14} + \frac{24}{14} - \frac{14}{14} = \frac{22}{14} = \mathbf{\frac{11}{7}}.
  \]
\end{enumerate}
\end{solution}

\begin{exercise}[平移变换下切平面过定点]
给定曲面 $F\!\left(\dfrac{x-a}{z-c},\;\dfrac{y-b}{z-c}\right)=0$（$a,b,c$ 为常数），或由它确定的曲面 $z=z(x,y)$。

证明：曲面的切平面通过一个定点。
\end{exercise}
\begin{solution}
\begin{enumerate}
  \item \textbf{令中间变量并求偏导}：令 $u = \dfrac{x-a}{z-c},\ v = \dfrac{y-b}{z-c}$，则 \(\displaystyle F_x = \frac{F_u}{z-c}\)，\(\displaystyle F_y = \frac{F_v}{z-c}\)，\(\displaystyle F_z = -F_u \frac{x-a}{(z-c)^2} - F_v \frac{y-b}{(z-c)^2}\)。
  \item \textbf{代入切平面一般方程}：在任意切点 $(x_0,y_0,z_0)$ 处，切平面为 $F_x(x-x_0) + F_y(y-y_0) + F_z(z-z_0) = 0$。代入上述偏导并两边同乘 $(z_0-c)^2$，可改写为
  \[
  F_u \cdot \left[ (z_0-c)(x-a) - (x_0-a)(z-c) \right] + F_v \cdot \left[ (z_0-c)(y-b) - (y_0-b)(z-c) \right] = 0.
  \]
  \item \textbf{寻找定点}：将 $(x,y,z) = (a,b,c)$ 代入，上式中两个方括号都为 $0$，故等式恒成立。因此无论切点在何处，其切平面必过定点 $\mathbf{(a,b,c)}$。
\end{enumerate}
\end{solution}

\chapter{重积分}

\section{知识讲解}

\subsection{二重积分基础专题：概念、性质与区域语言}

\subsubsection{二重积分的严格定义和重要性质}

\begin{definition}[二重积分] \label{def:double-integral}
设 $D$ 是 $\mathbb{R}^2$ 中的有界闭区域，函数 $f(x, y)$ 在 $D$ 上有定义。将 $D$ 任意分割成 $n$ 个小区域 $\Delta\sigma_i$（其面积也记为 $\Delta\sigma_i$），设 $\lambda$ 为各小区域直径的最大值。在每个 $\Delta\sigma_i$ 内任取一点 $(\xi_i, \eta_i)$，作黎曼和 \(\displaystyle \sum_{i=1}^n f(\xi_i, \eta_i) \Delta\sigma_i\)。
如果当 $\lambda \to 0$ 时，该黎曼和的极限存在，且\textbf{该极限值与闭区域 $D$ 的分割方式无关，也与点 $(\xi_i, \eta_i)$ 的取法无关}，则称此极限值为函数 $f(x, y)$ 在闭区域 $D$ 上的\textbf{二重积分}，记作
\(\displaystyle \iint\limits_D f(x, y) \, \mathrm{d}\sigma\)。
此时称函数 $f(x, y)$ 在 $D$ 上\textbf{可积}。其中：
\begin{itemize}
    \item $\iint$ 称为双重积分号；$D$ 称为积分区域；
    \item $f(x, y)$ 称为被积函数；$x, y$ 称为积分变量；
    \item $\mathrm{d}\sigma$ 称为面积元素（在直角坐标系下，通常写作 $\mathrm{d}x\,\mathrm{d}y$，即 $\iint\limits_D f(x, y) \, \mathrm{d}x\,\mathrm{d}y$）。
\end{itemize}
\end{definition}

\begin{remark}[关于可积性]
并不是所有函数都是可积的。但在高等数学中，我们主要研究以下两类可积函数：
\begin{itemize}
    \item 若函数 $f(x, y)$ 在有界闭区域 $D$ 上\textbf{连续}，则 $f(x, y)$ 在 $D$ 上必然可积。
    \item 若函数 $f(x, y)$ 在有界闭区域 $D$ 上\textbf{有界}，且其不连续点仅仅分布在有限个点或有限条光滑曲线上，则 $f(x, y)$ 在 $D$ 上也是可积的。
\end{itemize}
\end{remark}

设 $f(x, y)$ 和 $g(x, y)$ 是在闭区域 $D$ 上的可积函数，二重积分满足以下核心性质：

\begin{property}[基本代数性质与区域可加性] \label{prop:double-integral-basic}
\begin{itemize}
    \item \textbf{线性性质}：常数因子可以提出，和差的积分等于积分的和差，即
    \[
    \iint\limits_D [a f(x, y) \pm b g(x, y)] \, \mathrm{d}\sigma
    = a \iint\limits_D f(x, y) \, \mathrm{d}\sigma
    \pm b \iint\limits_D g(x, y) \, \mathrm{d}\sigma.
    \]
    \item \textbf{区域可加性}：若闭区域 $D$ 被划分为两个无公共内部的闭区域 $D_1$ 和 $D_2$（即 $D = D_1 \cup D_2$），则积分可以在区域上累加：
    \[
    \iint\limits_D f(x, y) \, \mathrm{d}\sigma
    = \iint\limits_{D_1} f(x, y) \, \mathrm{d}\sigma
    + \iint\limits_{D_2} f(x, y) \, \mathrm{d}\sigma.
    \]
\end{itemize}
\end{property}

\begin{property}[比较定理与估值定理] \label{prop:double-integral-estimate}
\begin{itemize}
    \item \textbf{面积积分}：若被积函数恒为 $1$，则二重积分的值恰好等于积分区域 $D$ 的面积 $\sigma$，即 \(\displaystyle \iint\limits_D 1 \, \mathrm{d}\sigma = \sigma\)。
    \item \textbf{保号性（比较定理）}：如果在区域 $D$ 上恒有 $f(x, y) \leqslant g(x, y)$，那么它们的积分也保持相同的不等号关系：\(\displaystyle \iint\limits_D f(x, y) \, \mathrm{d}\sigma
    \leqslant \iint\limits_D g(x, y) \, \mathrm{d}\sigma\)。
    \item \textbf{绝对值不等式}：函数绝对值的积分，大于等于函数积分的绝对值：
    \[
    \left| \iint\limits_D f(x, y) \, \mathrm{d}\sigma \right|
    \leqslant
    \iint\limits_D |f(x, y)| \, \mathrm{d}\sigma.
    \]
    \item \textbf{估值定理}：设 $m$ 和 $M$ 分别是连续函数 $f(x, y)$ 在闭区域 $D$ 上的最小值和最大值，$\sigma$ 为 $D$ 的面积，则有 \(\displaystyle m\sigma \leqslant \iint\limits_D f(x, y) \, \mathrm{d}\sigma \leqslant M\sigma\)。
\end{itemize}
\end{property}

\begin{theorem}[二重积分中值定理] \label{thm:double-integral-mean}
设 $f(x, y)$ 是闭区域 $D$ 上的连续函数，$\sigma$ 是区域 $D$ 的面积，则在 $D$ 内至少存在一点 $(\xi, \eta)$，使得 \(\displaystyle \iint\limits_D f(x, y) \, \mathrm{d}\sigma = f(\xi, \eta) \cdot \sigma\)。
（该定理表明，连续函数的二重积分等于该区域面积与区域内某一点函数值的乘积。此函数值 $f(\xi, \eta)$ 也常被称为函数在区域上的“平均高度”或“平均值”）。
\end{theorem}

\subsubsection{可积性、几何意义与估值}

\begin{exercise}[不可积函数的反例]
设区域为单位正方形 $D = [0, 1] \times [0, 1]$。已知二维的狄利克雷（Dirichlet）函数定义为：
\[
f(x, y) =
\begin{cases}
1, & (x, y) \text{ 两个坐标均为有理数时} \\
0, & \text{其它情况}
\end{cases}
\]
请严格依据定义证明：该函数在区域 $D$ 上是不可积的。
\end{exercise}
\begin{solution}
\begin{compactenum}
    \item \textbf{回顾可积定义}：根据二重积分的定义，函数可积的前提是：不论如何分割区域，也不论在小区域内如何选取采样点 $(\xi_i, \eta_i)$，黎曼和的极限都必须存在且等于同一个常数。若能找出两种不同的取点方式导致极限值不同，则说明函数不可积。
    \item \textbf{执行任意分割}：将正方形区域 $D$ 任意分割成 $n$ 个小区域 $\Delta\sigma_i$ $(i = 1, 2, \dots, n)$。
    \item \textbf{第一种取点策略（全取有理点）}：根据实数系中有理数的稠密性，任何一个小区域 $\Delta\sigma_i$ 内都必定包含无穷多个坐标全为有理数的点。我们强制要求在每一个小区域内选取的点 $(\xi_i, \eta_i)$ 都是有理点。此时 \(\displaystyle f(\xi_i, \eta_i)=1\qquad(\text{对所有 }i)\)，且 \(\displaystyle \sum_{i=1}^n f(\xi_i, \eta_i)\Delta\sigma_i=\sum_{i=1}^n\Delta\sigma_i=1\)。
    当 $\lambda \to 0$ 时，其极限值显然为 $1$。
    \item \textbf{第二种取点策略（全取无理点）}：同样根据无理数的稠密性，我们在每一个小区域内强制选取包含至少一个无理坐标的点作为 $(\xi_i, \eta_i)$。此时 \(\displaystyle f(\xi_i,\eta_i)=0\qquad(\text{对所有 }i)\)，且 \(\displaystyle \sum_{i=1}^n f(\xi_i,\eta_i)\Delta\sigma_i=0\)。
    当 $\lambda \to 0$ 时，其极限值显然为 $0$。
    \item \textbf{得出结论}：由于针对同一种分割，采用了不同的取点方式，导致黎曼和的极限值不同（$1 \neq 0$）。这违反了积分定义中“极限值与点的取法无关”的要求，因此函数 $f(x, y)$ 在区域 $D$ 上不可积。
\end{compactenum}
\end{solution}

\begin{exercise}[利用几何意义免计算求积分]
已知 $D$ 为平面上的圆域 $x^2 + y^2 \leqslant R^2$，试不通过微积分计算，直接求出二重积分 \(\displaystyle \iint\limits_D \sqrt{R^2 - x^2 - y^2} \, \mathrm{d}\sigma\) 的值。
\end{exercise}
\begin{solution}
\begin{compactenum}
    \item \textbf{分析被积函数的几何形态}：观察给定的二重积分，令三维空间中的竖坐标 $z = \sqrt{R^2 - x^2 - y^2}$。
    \item \textbf{代数变形与曲面识别}：将上述等式两边平方，可得 $z^2 = R^2 - x^2 - y^2$，即 $x^2 + y^2 + z^2 = R^2$。这是一个以原点 $(0,0,0)$ 为球心、$R$ 为半径的球面方程。由于原函数前有正号的算术平方根，所以 $z \geqslant 0$，这意味着曲面是位于 $Oxy$ 平面上方的\textbf{上半球面}。
    \item \textbf{联立积分区域}：积分底面区域 $D$ 是 $x^2 + y^2 \leqslant R^2$，这恰好是上述球体在 $Oxy$ 平面上的赤道截面。因此，该二重积分在几何上表示的是：以原点为中心、半径为 $R$ 的\textbf{上半球体的体积}。
    \item \textbf{利用已知体积公式求解}：半径为 $R$ 的完整球体体积为 $\dfrac{4}{3}\pi R^3$，故上半球体积 \(\displaystyle V = \frac{1}{2} \times \frac{4}{3}\pi R^3 = \frac{2}{3}\pi R^3\)。
    因而积分结果为 $\dfrac{2}{3}\pi R^3$。
\end{compactenum}
\end{solution}

\begin{exercise}[利用估值定理估计积分范围]
利用二重积分的性质估计积分 $I = \iint\limits_D \mathrm{e}^{x^2+y^2} \, \mathrm{d}\sigma$ 的取值范围，其中积分区域 $D$ 是由椭圆 $\dfrac{x^2}{a^2} + \dfrac{y^2}{b^2} \leqslant 1$ 围成的闭区域（假定 $0 < b < a$）。
\end{exercise}
\begin{solution}
\begin{compactenum}
    \item \textbf{明确估值定理的要求}：要使用估值定理 $m\sigma \leqslant I \leqslant M\sigma$，我们需要先求出被积函数 $f(x, y) = \mathrm{e}^{x^2+y^2}$ 在区域 $D$ 上的最大值 $M$、最小值 $m$，以及区域 $D$ 的面积 $\sigma$。
    \item \textbf{分析被积函数的最值}：由于自然指数函数 $u \mapsto \mathrm{e}^u$ 是严格单调递增的，所以求 $\mathrm{e}^{x^2+y^2}$ 的最值，等价于求指数部分 $u = x^2+y^2$ 的最值。在几何上，$x^2+y^2$ 表示区域内点 $(x, y)$ 到原点 $(0, 0)$ 的距离的平方。
    \item \textbf{结合椭圆区域求距离极值}：区域 $D$ 包含原点 $(0,0)$，故 $u_{\min}=0$，对应最小值 $m=\mathrm{e}^0=1$；又由于 $a>b$，椭圆上离原点最远的点是 $(\pm a,0)$，故 $u_{\max}=a^2$，对应最大值 $M=\mathrm{e}^{a^2}$。
    \item \textbf{计算区域的面积}：长半轴为 $a$、短半轴为 $b$ 的椭圆所围成的区域面积为标准公式 $\sigma = \pi a b$。
    \item \textbf{代入估值定理公式}：于是 \(\displaystyle \pi a b \leqslant
    \iint\limits_D \mathrm{e}^{x^2+y^2} \, \mathrm{d}\sigma
    \leqslant \pi a b \,\mathrm{e}^{a^2}\)。
\end{compactenum}
\end{solution}

\begin{exercise}[中值定理与极限的综合应用]
设函数 $f(x, y)$ 是在有界闭区域 $D: x^2 + y^2 \leqslant \rho_0^2$ 上连续的函数，求极限
\[
\lim_{\rho \to 0^+} \frac{1}{\pi\rho^2}
\iint\limits_{x^2+y^2\leqslant\rho^2} f(x, y) \, \mathrm{d}\sigma
\qquad (0 < \rho \leqslant \rho_0).
\]
\end{exercise}
\begin{solution}
\begin{compactenum}
    \item \textbf{识别表达式的结构}：观察极限号后的表达式。分母 $\pi\rho^2$ 恰好是积分区域 $D_\rho: x^2 + y^2 \leqslant \rho^2$ 的面积 $\sigma$。因此，整体表达式 $\dfrac{1}{\sigma} \iint_{D_\rho} f(x, y) \, \mathrm{d}\sigma$ 代表的是连续函数 $f(x, y)$ 在圆域 $D_\rho$ 上的“积分平均值”。
    \item \textbf{应用二重积分中值定理}：因为 $f(x,y)$ 在圆域 $D_\rho$ 上连续，故存在 $(\xi_\rho, \eta_\rho)\in D_\rho$，使得
    \[
    \iint\limits_{x^2+y^2\leqslant\rho^2} f(x, y) \, \mathrm{d}\sigma
    =f(\xi_\rho, \eta_\rho)\cdot \pi\rho^2.
    \]
    \item \textbf{化简极限表达式}：代回原式即得
    \[
    \lim_{\rho \to 0^+} \frac{1}{\pi\rho^2}
    \iint\limits_{x^2+y^2\leqslant\rho^2} f(x, y) \, \mathrm{d}\sigma
    =\lim_{\rho \to 0^+} f(\xi_\rho, \eta_\rho).
    \]
    \item \textbf{分析点的位置变化并利用连续性求解}：因为点 $(\xi_\rho, \eta_\rho)$ 始终位于半径为 $\rho$ 的圆内（即满足 $\xi_\rho^2 + \eta_\rho^2 \leqslant \rho^2$），当极限过程 $\rho \to 0^+$ 发生时，这个圆域会向着原点 $(0, 0)$ 不断收缩，导致内部的中值点 $(\xi_\rho, \eta_\rho)$ 必然被迫无限趋近于原点 $(0, 0)$。
    \item \textbf{利用函数的连续性}：由连续性，
    \[
    \lim_{\rho \to 0^+} f(\xi_\rho, \eta_\rho)
    =f\left( \lim_{\rho \to 0^+} \xi_\rho, \lim_{\rho \to 0^+} \eta_\rho \right)
    =f(0, 0),
    \]
    故原极限结果为 $f(0, 0)$。
\end{compactenum}
\end{solution}


\subsection{二重积分计算专题：累次积分、换序与极坐标}

\subsubsection{矩形区域上的二重积分}

最简单的二维积分区域是边平行于坐标轴的矩形。我们首先研究在这类区域上如何将二重积分转化为两次单变量定积分。

\begin{theorem}[矩形区域上的二重积分 (Fubini 定理)] \label{thm:fubini-rect}
设闭矩形区域 $D = \{ (x,y) : a \leqslant x \leqslant b, c \leqslant y \leqslant d \} = [a,b] \times [c,d]$，且二元函数 $f(x, y)$ 在 $D$ 上可积，则：
\begin{compactenum}
    \item 若对任意给定的 $y \in [c, d]$，函数 $f(x, y)$ 作为 $x$ 的一元函数在 $[a, b]$ 上可积。令 $\phi(y) = \int_{a}^{b} f(x, y) \,\mathrm{d}x$，则 $\phi(y)$ 在 $[c, d]$ 上也可积，并且：
    \[
    \int_{c}^{d} \phi(y) \,\mathrm{d}y = \int_{c}^{d} \left[ \int_{a}^{b} f(x, y) \,\mathrm{d}x \right] \mathrm{d}y = \iint\limits_{D} f(x, y) \,\mathrm{d}\sigma
    \]
    \item 同理，若对任意给定的 $x \in [a, b]$，函数 $f(x, y)$ 作为 $y$ 的一元函数在 $[c, d]$ 上可积。令 $\psi(x) = \int_{c}^{d} f(x, y) \,\mathrm{d}y$，则 $\psi(x)$ 在 $[a, b]$ 上也可积，并且：
    \[
    \int_{a}^{b} \psi(x) \,\mathrm{d}x = \int_{a}^{b} \left[ \int_{c}^{d} f(x, y) \,\mathrm{d}y \right] \mathrm{d}x = \iint\limits_{D} f(x, y) \,\mathrm{d}\sigma
    \]
\end{compactenum}
\end{theorem}

定理的核心是把二重积分化成两个嵌套的一元定积分。内层积分先固定一个 $y$，沿着水平线段 $x\in[a,b]$ 把这一条切片上的函数值累加起来，因此输出只剩关于 $y$ 的函数；外层积分再把所有切片总量继续累加，最终得到整个二维区域上的总和。为强调积分变量的归属，常写作
\[
\int_{c}^{d} \left[ \int_{a}^{b} f(x, y) \,\mathrm{d}x \right] \mathrm{d}y \equiv \int_{c}^{d} \mathrm{d}y \int_{a}^{b} f(x, y) \,\mathrm{d}x.
\]
当 $f(x, y) \geqslant 0$ 时，二重积分还可看成曲顶柱体体积：用平面 $x = x_0$ 截该柱体，截面面积是 $\psi(x_0) = \int_{c}^{d} f(x_0, y) \,\mathrm{d}y$，再沿 $x$ 轴累加这些截面面积，就得到总体积。

\subsubsection{一般区域上的二重积分：X型与Y型区域}

并不是所有的积分区域都是矩形。对于边界为曲线的一般区域，我们常常通过作平行于坐标轴的直线，将其划分为便于积分的“标准区域”。

\begin{definition}[积分区域的分类] \label{def:double-integral-region-types}
\textbf{$x$-型区域（上下穿透型）}：若区域可表示为 $D = \{ (x, y) : \phi_{1}(x) \leqslant y \leqslant \phi_{2}(x), a \leqslant x \leqslant b \}$。即区域在 $x$ 轴上的投影为区间 $[a,b]$，且垂直于 $x$ 轴的直线穿过区域内部时，与边界最多只有两个交点。

\textbf{$y$-型区域（左右穿透型）}：若区域可表示为 $D = \{ (x, y) : \psi_{1}(y) \leqslant x \leqslant \psi_{2}(y), c \leqslant y \leqslant d \}$。即区域在 $y$ 轴上的投影为区间 $[c,d]$，且水平线穿过内部时最多只有左右两个交点。
\end{definition}

\begin{theorem}[一般区域上的二重积分计算] \label{thm:double-integral-general-region}
\begin{compactenum}
    \item 若 $f(x, y)$ 在 $x$-型区域 $D$ 上可积，则化为先对 $y$ 后对 $x$ 的累次积分：\[\displaystyle \iint\limits_{D} f(x, y) \,\mathrm{d}\sigma = \int_{a}^{b} \mathrm{d}x \int_{\phi_{1}(x)}^{\phi_{2}(x)} f(x, y) \,\mathrm{d}y\]
    \item 若 $f(x, y)$ 在 $y$-型区域 $D$ 上可积，则化为先对 $x$ 后对 $y$ 的累次积分：\[\displaystyle \iint\limits_{D} f(x, y) \,\mathrm{d}\sigma = \int_{c}^{d} \mathrm{d}y \int_{\psi_{1}(y)}^{\psi_{2}(y)} f(x, y) \,\mathrm{d}x\]
\end{compactenum}
\end{theorem}

\begin{note}
微分号后面是$x$，积分上下限和$y$有关；微分号后面是$y$，积分上下限和$x$有关。
\end{note}

\subsubsection{直角坐标系下的计算范例}

\begin{exercise}[矩形区域下的直接计算]
求 $\iint\limits_{D} (x+y^{2}) \,\mathrm{d}\sigma$，其中积分区域为矩形 $D = [0,1] \times [0,3]$。
\end{exercise}
\begin{solution}
\begin{compactenum}
    \item \textbf{确定积分次序}：区域 $D$ 为标准矩形，被积函数多项式无奇点，选择先对 $x$ 后对 $y$ 或者相反均可。此处选择先对 $x$ 积分。
    \item \textbf{转换为累次积分}：\(\displaystyle \iint\limits_{D} (x+y^{2}) \,\mathrm{d}\sigma = \int_{0}^{3} \mathrm{d}y \int_{0}^{1} (x+y^{2}) \,\mathrm{d}x\)。
    \item \textbf{计算内层关于 $x$ 的积分}（视 $y$ 为常数）：\(\displaystyle \int_{0}^{1} (x+y^{2}) \,\mathrm{d}x = \left[ \frac{1}{2}x^2 + xy^2 \right]_{x=0}^{x=1} = \frac{1}{2} + y^2\)。
    \item \textbf{计算外层关于 $y$ 的积分}：\(\displaystyle \int_{0}^{3} \left( \frac{1}{2} + y^2 \right) \mathrm{d}y = \left[ \frac{1}{2}y + \frac{1}{3}y^3 \right]_0^3 = \frac{3}{2} + 9 = \frac{21}{2}\)。
\end{compactenum}
\end{solution}

\begin{exercise}[复杂区域的类型划分]
设 $D = \{ (x, y) : 2y \leqslant x^{2}+y^{2} \leqslant 4y, x \geqslant 0 \}$，试将 $D$ 分别表示为 $x$-型区域和 $y$-型区域的并集。
\end{exercise}
\begin{solution}
\begin{compactenum}
    \item \textbf{分析边界方程的几何特征}：
    \begin{itemize}
        \item 内部边界 $x^2+y^2 = 2y \implies x^2+(y-1)^2 = 1$，是以 $(0,1)$ 为圆心，半径为 $1$ 的圆。
        \item 外部边界 $x^2+y^2 = 4y \implies x^2+(y-2)^2 = 4$，是以 $(0,2)$ 为圆心，半径为 $2$ 的圆。
        \item 约束条件 $x \geqslant 0$ 表示区域仅位于 $y$ 轴的右侧。
        \item 因此，$D$ 是位于第一象限的两个相切圆之间的月牙形区域。
    \end{itemize}
    \item \textbf{划分并表示为 $x$-型区域}：
    \begin{itemize}
        \item 整体的 $x$ 跨度为大圆的最右端，即 $x \in [0, 2]$。但在 $x \in [0, 1]$ 之间存在小圆的阻挡。
        \item 从大圆方程解出上下半支：$y_{\mathrm{outer,up}} = 2+\sqrt{4-x^2}$，$y_{\mathrm{outer,low}} = 2-\sqrt{4-x^2}$。
        \item 从小圆方程解出上下半支：$y_{\mathrm{inner,up}} = 1+\sqrt{1-x^2}$，$y_{\mathrm{inner,low}} = 1-\sqrt{1-x^2}$。
        \item 当 $0 \leqslant x \leqslant 1$ 时，穿透区域的竖线被打断成两截。分为上下两块：\\
        $D_1 = \{ (x,y) : 0 \leqslant x \leqslant 1, 2-\sqrt{4-x^2} \leqslant y \leqslant 1-\sqrt{1-x^2} \}$ \\
        $D_2 = \{ (x,y) : 0 \leqslant x \leqslant 1, 1+\sqrt{1-x^2} \leqslant y \leqslant 2+\sqrt{4-x^2} \}$
        \item 当 $1 < x \leqslant 2$ 时，小圆已经不再与竖直切片相交，所以竖线只会与大圆边界相交两次，即从大圆下缘进入、从大圆上缘穿出：\\
        $D_3 = \{ (x,y) : 1 \leqslant x \leqslant 2, 2-\sqrt{4-x^2} \leqslant y \leqslant 2+\sqrt{4-x^2} \}$
        \item $D = D_1 \cup D_2 \cup D_3$。
    \end{itemize}
    \item \textbf{划分并表示为 $y$-型区域}：
    \begin{itemize}
        \item 整体的 $y$ 跨度为大圆的极值，即 $y \in [0, 4]$。解出右半支方程（$x \geqslant 0$）：大圆 $x = \sqrt{4y-y^2}$，小圆 $x = \sqrt{2y-y^2}$。
        \item 当 $0 \leqslant y \leqslant 2$ 时，水平线从小圆右侧进入，从大圆右侧穿出：\\
        $D_4 = \{ (x,y) : 0 \leqslant y \leqslant 2, \sqrt{2y-y^2} \leqslant x \leqslant \sqrt{4y-y^2} \}$
        \item 当 $2 < y \leqslant 4$ 时，小圆已经结束，水平线从 $y$ 轴（$x=0$）进入，从大圆右侧穿出：\\
        $D_5 = \{ (x,y) : 2 \leqslant y \leqslant 4, 0 \leqslant x \leqslant \sqrt{4y-y^2} \}$
        \item $D = D_4 \cup D_5$。相比之下，$y$-型划分更加简洁。
    \end{itemize}
\end{compactenum}
\end{solution}

\begin{exercise}
设 $D$ 是由直线 $y=x$ 与抛物线 $y=x^{2}$ 围成的闭区域，求 \(\displaystyle I = \iint\limits_{D} \frac{1}{2}(2-x-y) \,\mathrm{d}\sigma\)。
\end{exercise}
\begin{solution}
\begin{compactenum}
    \item \textbf{明确边界与交点}：联立方程 $y=x$ 与 $y=x^2$，解得交点为 $(0,0)$ 和 $(1,1)$。在 $(0,1)$ 区间内，直线 $y=x$ 恒在抛物线 $y=x^2$ 的上方。
    \item \textbf{视角一：将 $D$ 视为 $x$-型区域进行积分}：
    \begin{itemize}
        \item 投影范围：$x \in [0, 1]$。下边界 $y = x^2$，上边界 $y = x$。
        \item 建立累次积分并计算内层：
        \[
        I = \frac{1}{2} \int_{0}^{1} \mathrm{d}x \int_{x^{2}}^{x} (2-x-y) \,\mathrm{d}y
        = \frac{1}{2} \int_{0}^{1} \left[ (2-x)y - \frac{1}{2}y^2 \right]_{y=x^2}^{y=x} \mathrm{d}x
        \]
        \item 代入上下限整理被积函数：\(\displaystyle \left( (2-x)x - \frac{1}{2}x^2 \right) - \left( (2-x)x^2 - \frac{1}{2}x^4 \right) = 2x - \frac{7}{2}x^2 + x^3 + \frac{1}{2}x^4\)。
        \item 计算外层积分：\(\displaystyle I=\frac12\int_0^1\left(2x-\frac72x^2+x^3+\frac12x^4\right)\mathrm{d}x=\frac12\left(1-\frac76+\frac14+\frac1{10}\right)=\frac{11}{120}\)。
    \end{itemize}
    \item \textbf{视角二：将 $D$ 视为 $y$-型区域进行积分验证}：
    \begin{itemize}
        \item 投影范围：$y \in [0, 1]$。左边界 $x = y$（由直线反解），右边界 $x = \sqrt{y}$（由抛物线在第一象限反解）。
        \item 建立累次积分并计算内层：
        \[
        I = \frac{1}{2} \int_{0}^{1} \mathrm{d}y \int_{y}^{\sqrt{y}} (2-x-y) \,\mathrm{d}x
        = \frac{1}{2} \int_{0}^{1} \left[ (2-y)x - \frac{1}{2}x^2 \right]_{x=y}^{x=\sqrt{y}} \mathrm{d}y
        \]
        \item 代入计算后其结果依然为 $\frac{11}{120}$。这验证了 Fubini 定理下积分次序交换的正确性。
    \end{itemize}
\end{compactenum}
\end{solution}

\begin{exercise}[不可积函数的积分次序选择]
设三角形平面薄片区域 $D$ 的三个顶点分别为 $(0,0)$、$(2,1)$、$(0,1)$。
其面密度分布为 $\rho(x, y) = \mathrm{e}^{y^{2}}$。求总质量 \(\displaystyle M = \iint_{D} \rho(x, y) \,\mathrm{d}\sigma\)。
\end{exercise}
\begin{solution}
    \begin{itemize}
        \item 若先对 $x$ 积分：投影区间 $y \in [0, 1]$。对于任意固定的高度 $y$，左边界为 $x = 0$，右边界由斜线反解得 $x = 2y$。
    \item 建立新的累次积分：\(\displaystyle M = \int_{0}^{1} \mathrm{d}y \int_{0}^{2y} \mathrm{e}^{y^{2}} \,\mathrm{d}x\)。
        \item 内层关于 $x$ 的积分极易完成（此时 $\mathrm{e}^{y^{2}}$ 视为常数）：
        \[
        \int_{0}^{2y} \mathrm{e}^{y^{2}} \,\mathrm{d}x
        = \left[ x \mathrm{e}^{y^{2}} \right]_{x=0}^{x=2y}
        = 2y\mathrm{e}^{y^{2}}.
        \]
        \item 惊喜出现了：内层积分提供的额外变量 $2y$，恰好构成了指数部分 $y^2$ 的导数 $\mathrm{d}(y^2) = 2y\,\mathrm{d}y$。
        \item 顺理成章地完成外层定积分：令 $s=y^2$，则 \(\displaystyle M=\int_0^1 \mathrm{e}^{s}\,\mathrm{d}s=\mathrm{e}-1\)。
    \end{itemize}
本题真正的核心是\textbf{能否识别谁不该放在内层}。一旦看到 $\mathrm{e}^{y^2}$ 没有初等原函数，就应让它留在外层，把另一变量的积分先做掉。
\end{solution}

\begin{exercise}[累次积分的次序交换]
计算给定的累次积分：
\(\displaystyle A = \int_{0}^{1} \mathrm{d}y \int_{y}^{\sqrt{y}} \frac{\sin x}{x} \,\mathrm{d}x\)。
\end{exercise}
\begin{solution}
\begin{compactenum}
    \item \textbf{识别困难与策略制定}：原式先对 $x$ 积分，但 $\frac{\sin x}{x}$ 无法积出初等原函数。必须交换积分次序。
    \item \textbf{由积分限还原二维积分区域}：
    \begin{itemize}
        \item 外层提供 $y$ 的绝对范围：$y \in [0, 1]$。
        \item 内层提供 $x$ 的相对范围：$x$ 在 $y$ 与 $\sqrt{y}$ 之间，即 $y \leqslant x \leqslant \sqrt{y}$。
        \item 这个约束条件实际上意味着区域 $D$ 的边界为抛物线 $x = \sqrt{y}$（即 $y = x^2$ 的第一象限分支）和直线 $x = y$（即 $y = x$）。
    \end{itemize}
    \item \textbf{改写为另一型的积分限}：
    \begin{itemize}
        \item 区域 $D$ 投影到 $x$ 轴上的范围是 $x \in [0, 1]$。
        \item 对于任一给定的 $x$，在 $D$ 内部作竖直射线，从下边界 $y = x^2$ 穿入，从上边界 $y = x$ 穿出。
        \item 所以区域重写为：\(\displaystyle D=\{(x,y):0\leqslant x\leqslant 1,\ x^2\leqslant y\leqslant x\}\)。
    \end{itemize}
    \item \textbf{建立新次序并求解}：
    \begin{itemize}
    \item 交换次序后的二重积分为 \(\displaystyle A = \int_{0}^{1} \mathrm{d}x \int_{x^{2}}^{x} \frac{\sin x}{x} \,\mathrm{d}y\)。
        \item 先对 $y$ 积分（把含有 $x$ 的部分当作系数提出去）：
        \[
        \int_{x^{2}}^{x} \frac{\sin x}{x} \,\mathrm{d}y
        = \frac{\sin x}{x} \left( x - x^{2} \right)
        = (1 - x)\sin x.
        \]
        \item 进行外层关于 $x$ 的定积分计算（使用分部积分法）：\(\displaystyle A=\int_0^1(1-x)\sin x\,\mathrm{d}x=\int_0^1(1-x)\,\mathrm \mathrm{d}(-\cos x)=1-\sin1\)。
    \end{itemize}
\end{compactenum}
\end{solution}

\begin{exercise}[利用对称性简化积分计算]
设区域 $D$ 是由圆周 $x^{2}+y^{2}=4$ 以及抛物线 $y=-x^{2}+1$ 和 $y=x^{2}-1$（取 $|x| \leqslant 1$ 的部分）共同围成的复合闭区域，求
\(\displaystyle I = \iint\limits_{D} (x^{2}+y) \,\mathrm{d}\sigma\)。
\end{exercise}
\begin{solution}
\begin{compactenum}
    \item \textbf{分析区域对称性与被积函数性质}：
    \begin{itemize}
        \item 观察圆方程 $x^2+y^2=4$ 关于 $x$ 轴和 $y$ 轴均对称。
        \item 两段抛物线边界 $y=-x^2+1$ 和 $y=x^2-1$ 互为相反数，表明它们共同构成的“纺锤形”内凹部分也完全关于 $x$ 轴上下对称，同时抛物线本身关于 $y$ 轴左右对称。
        \item 整体积分区域 $D$ 是一个\textbf{既关于 $x$ 轴对称，又关于 $y$ 轴对称}的图形。
        \item 被积函数 $f(x,y) = x^2 + y$ 可拆分为 $f_1(x,y) = y$ 和 $f_2(x,y) = x^2$ 两部分独立计算。
    \end{itemize}
    \item \textbf{利用奇偶性简化第一部分}：
    \begin{itemize}
        \item 对于 $f_1(x,y) = y$，有 $f_1(x, -y) = -y = -f_1(x,y)$，说明它是关于 $y$ 的\textbf{奇函数}。
        \item 由于区域 $D$ 关于 $x$ 轴对称，区域中每个点 $(x,y)$ 都能和 $(x,-y)$ 配对，而这两个点对应的函数值互为相反数，因此上下两部分的贡献完全抵消：\(\displaystyle \iint\limits_D y\,\mathrm{d}\sigma=0\)。
    \end{itemize}
    \item \textbf{利用偶次对称简化第二部分}：
    \begin{itemize}
        \item 对于 $f_2(x,y) = x^2$，它既是关于 $x$ 的偶函数也是关于 $y$ 的偶函数。
        \item 这意味着四个象限中对应小块的函数值完全相同，所以可以只计算第一象限部分 $D_1 = D \cap \{x \geqslant 0, y \geqslant 0\}$，再乘以 $4$。
        \item 现在关键不是立刻列积分，而是看清第一象限里的区域结构：当 $0 \leqslant x \leqslant 1$ 时，圆域下方有一块被抛物线 $y=1-x^2$ 挖去；当 $1 \leqslant x \leqslant 2$ 时，抛物线已经落到 $x$ 轴下方，不再影响第一象限中的区域，所以竖直切片会从 $y=0$ 一直走到圆周 $y=\sqrt{4-x^2}$。
    \end{itemize}
    \item \textbf{切分 $D_1$ 并建立积分}：
    \[
    I = 4 \iint\limits_{D_1} x^2 \,\mathrm{d}\sigma = 4 \left( \int_{0}^{1} \mathrm{d}x \int_{-x^{2}+1}^{\sqrt{4-x^{2}}} x^{2} \,\mathrm{d}y + \int_{1}^{2} \mathrm{d}x \int_{0}^{\sqrt{4-x^{2}}} x^{2} \,\mathrm{d}y \right)
    \]
    上式是按“分段切片”直接写出的。如果从运算量考虑，用“第一象限整块圆域减去被挖去的小块”更简洁：
    \[
    I = 4 \left( \int_{0}^{2} \mathrm{d}x \int_{0}^{\sqrt{4-x^{2}}} x^2 \,\mathrm{d}y - \int_{0}^{1} \mathrm{d}x \int_{0}^{-x^{2}+1} x^2 \,\mathrm{d}y \right)
    \]
    \item \textbf{逐步计算定积分}：
    \begin{itemize}
        \item 先完成内层积分：
        \[
        I=4\int_{0}^{2}x^2\sqrt{4-x^2}\,\mathrm{d}x
        -4\int_{0}^{1}x^2(1-x^2)\,\mathrm{d}x.
        \]
        \item 计算前一部分（运用三角代换，令 $x = 2\sin t$, $\mathrm{d}x = 2\cos t \,\mathrm{d}t$）：
        \[
        4 \int_{0}^{\frac{\pi}{2}} (4\sin^2 t)(2\cos t)(2\cos t) \,\mathrm{d}t
        = 64 \int_{0}^{\frac{\pi}{2}} \sin^2 t \cos^2 t \,\mathrm{d}t.
        \]
        利用半角公式 $\sin^2 t \cos^2 t = \frac{1}{4}\sin^2 2t = \frac{1}{8}(1-\cos 4t)$，积分为
        \[
        64\times\frac18\left[t-\frac{\sin4t}{4}\right]_0^{\frac{\pi}{2}}
        =8\times\frac{\pi}{2}=4\pi.
        \]
        \item 计算后一部分（基础多项式）：
        \[
        4\int_{0}^{1}(x^2-x^4)\,\mathrm{d}x
        =4\left[\frac{x^3}{3}-\frac{x^5}{5}\right]_0^1
        =4\times\frac{2}{15}=\frac{8}{15}.
        \]
    \end{itemize}
    \item \textbf{得出最终结果}：将两部分相减，得到 $I = 4\pi - \frac{8}{15}$。

    本题的计算还说明两点：第一，先把被积函数拆成“奇函数部分 + 偶函数部分”，用对称性先消掉一部分；第二，复合区域在某个象限内常可用“整块减去挖空部分”的补集写法，把分段积分压缩成更整齐的表达式。
\end{compactenum}
\end{solution}

\subsubsection{二重积分的一般换元法}

\begin{definition}[曲线坐标与雅可比行列式] \label{def:curvilinear-coordinates-jacobian}
设有两个连续可导的多元函数 $x = x(u, v)$ 和 $y = y(u, v)$，并且该映射将 $uv$-平面上的区域 $D'$ 与 $xy$-平面上的区域 $D$ 建立了一一对应的关系。我们将网格在映射下导致的面积畸变率（微小平行四边形面积的放大或缩小倍数）定义为雅可比 (Jacobian) 行列式的绝对值，雅可比式定义为：
\[
\frac{\partial(x, y)}{\partial(u, v)} = \det \begin{pmatrix}
\frac{\partial x}{\partial u} & \frac{\partial y}{\partial u} \\
\frac{\partial x}{\partial v} & \frac{\partial y}{\partial v}
\end{pmatrix} = \frac{\partial x}{\partial u}\frac{\partial y}{\partial v} - \frac{\partial x}{\partial v}\frac{\partial y}{\partial u}
\]
在曲线坐标下，新的面积微元必须附加这个变形系数：$ \mathrm{d}\sigma = \left| \frac{\partial(x, y)}{\partial(u, v)} \right| \mathrm{d}u\,\mathrm{d}v $。
\end{definition}

\begin{theorem}[二重积分的一般换元公式] \label{thm:double-integral-substitution}
设变换 $x=x(u, v)$, $y=y(u, v)$ 将区域 $D'$ 映射到区域 $D$ 且满足上述条件，则对于 $D$ 上可积的函数 $f(x, y)$，有：
\[
\iint\limits_{D} f(x, y) \,\mathrm{d}x\,\mathrm{d}y = \iint\limits_{D'} f\big(x(u, v), y(u, v)\big) \left| \frac{\partial(x, y)}{\partial(u, v)} \right| \mathrm{d}u\,\mathrm{d}v
\]
\end{theorem}

\begin{exercise}[极坐标变换的面积微元推导]
请推导标准极坐标变换 $x = r \cos \phi$, $y = r \sin \phi$ 下的面积元素，并写出换元公式。
\end{exercise}
\begin{solution}
\begin{compactenum}
    \item \textbf{明确变换关系}：以 $r$ 和 $\phi$ 分别扮演定理中的 $u$ 和 $v$ 角色。
    \item \textbf{计算偏导数}：
    \begin{itemize}
        \item 对 $r$ 求偏导：$\frac{\partial x}{\partial r} = \cos \phi$，$\frac{\partial y}{\partial r} = \sin \phi$。
        \item 对 $\phi$ 求偏导：$\frac{\partial x}{\partial \phi} = -r\sin \phi$，$\frac{\partial y}{\partial \phi} = r\cos \phi$。
    \end{itemize}
    \item \textbf{构造雅可比行列式}：
    \[
    \frac{\partial(x, y)}{\partial(r,\phi)} = \det \begin{pmatrix}
    \cos \phi & -r \sin \phi \\
    \sin \phi & r \cos \phi
    \end{pmatrix} = r(\cos^2\phi + \sin^2\phi) = r
    \]
    \item \textbf{得出面积微元与公式}：由于极径 $r \geqslant 0$，其绝对值 $|r| = r$。所以极坐标系下的面积元素为 $r \,\mathrm{d}r\,\mathrm{d}\phi$。一般换元公式变为著名的极坐标积分形式：\(\displaystyle \iint\limits_{D} f(x, y) \,\mathrm{d}\sigma = \iint\limits_{D'} f(r \cos \phi, r \sin \phi) \,r \,\mathrm{d}r\,\mathrm{d}\phi\)。
\end{compactenum}
\end{solution}

\begin{exercise}[极坐标与高斯积分计算]
设区域 $D = \{ (x, y) : x^{2}+y^{2} \leqslant a^{2} \}$，
\begin{shortexenum}
    \shortitem{1}{利用极坐标变换求积分 $I_a = \iint\limits_{D} \mathrm{e}^{-(x^{2}+y^{2})} \,\mathrm{d}\sigma$；} &
    \shortitem{2}{证明：$\int_{0}^{+\infty} \mathrm{e}^{-x^{2}} \,\mathrm{d}x = \frac{\sqrt{\pi}}{2}$。}
\end{shortexenum}
\end{exercise}
\begin{solution}
\begin{compactenum}
    \item \textbf{解答第一问：直接在极坐标下计算}：
    \begin{itemize}
        \item 区域 $D$ 是半径为 $a$ 的中心圆域，在极坐标下描述极为简单：$0 \leqslant \phi \leqslant 2\pi$，且 $0 \leqslant r \leqslant a$。
        \item 被积函数中包含特征项 $x^2+y^2 = r^2$。代入极坐标面积微元：\(\displaystyle I_a = \int_{0}^{2\pi} \mathrm{d}\phi \int_{0}^{a} \mathrm{e}^{-r^{2}} r \,\mathrm{d}r\)。
        \item 外层角度积分与内层变量无关，可直接提出：$\int_0^{2\pi} \mathrm{d}\phi = 2\pi$。
        \item 内层径向积分利用凑微分：
        \[
        \int_{0}^{a} \mathrm{e}^{-r^{2}} r \,\mathrm{d}r = -\frac{1}{2} \int_{0}^{a} \mathrm{e}^{-r^{2}} \,\mathrm{d}(-r^2) = -\frac{1}{2} \left[ \mathrm{e}^{-r^{2}} \right]_0^a = -\frac{1}{2} (\mathrm{e}^{-a^{2}} - 1) = \frac{1}{2}(1 - \mathrm{e}^{-a^{2}})
        \]
        \item 二者相乘得到圆域积分结果：$I_a = \pi(1 - \mathrm{e}^{-a^{2}})$。
    \end{itemize}
    \item \textbf{解答第二问：利用夹逼定理推导无穷限一元积分}：
    \begin{itemize}
        \item 构造三个嵌套的闭区域：
        \begin{itemize}
            \item 内切圆域 $D$：$x^2+y^2 \leqslant a^2$
            \item 正方形区域 $D_1$：$[-a, a] \times [-a, a]$
            \item 外接圆域 $D_2$：$x^2+y^2 \leqslant 2a^2$
        \end{itemize}
        显然有空间包含关系：$D \subset D_1 \subset D_2$。
        \item 由于被积函数恒正 $\mathrm{e}^{-(x^2+y^2)} > 0$，区域面积越大积分值越大（保号性衍生）：
        \[
        \iint\limits_{D} \mathrm{e}^{-(x^{2}+y^{2})} \,\mathrm{d}\sigma
        \leqslant \iint\limits_{D_1} \mathrm{e}^{-(x^{2}+y^{2})} \,\mathrm{d}\sigma
        \leqslant \iint\limits_{D_2} \mathrm{e}^{-(x^{2}+y^{2})} \,\mathrm{d}\sigma.
        \]
        \item 我们已经算出圆域上的积分公式，代入两端得：\(\displaystyle \pi(1 - \mathrm{e}^{-a^{2}}) \leqslant \iint\limits_{D_1} \mathrm{e}^{-(x^{2}+y^{2})} \,\mathrm{d}\sigma \leqslant \pi(1 - \mathrm{e}^{-2a^{2}})\)。
        \item 对于中间正方形区域，运用 Fubini 定理将其分离为两个一元积分的乘积：
        \[
        \iint\limits_{D_1} \mathrm{e}^{-x^{2}}\mathrm{e}^{-y^{2}} \,\mathrm{d}x\,\mathrm{d}y = \left( \int_{-a}^{a} \mathrm{e}^{-x^{2}} \,\mathrm{d}x \right) \left( \int_{-a}^{a} \mathrm{e}^{-y^{2}} \,\mathrm{d}y \right) = \left( \int_{-a}^{a} \mathrm{e}^{-x^{2}} \,\mathrm{d}x \right)^2
        \]
        \item 于是建立起了不等式连缀关系：\(\displaystyle \pi(1 - \mathrm{e}^{-a^{2}}) \leqslant \left( \int_{-a}^{a} \mathrm{e}^{-x^{2}} \,\mathrm{d}x \right)^2 \leqslant \pi(1 - \mathrm{e}^{-2a^{2}})\)。
        \item 对此不等式两端令 $a \to +\infty$，根据夹逼定理，左右极限均为 $\pi$，即 \(\displaystyle \lim_{a \to +\infty} \left( \int_{-a}^{a} \mathrm{e}^{-x^{2}} \,\mathrm{d}x \right)^2 = \pi\)，因此 \(\displaystyle \int_{-\infty}^{+\infty} \mathrm{e}^{-x^{2}} \,\mathrm{d}x = \sqrt{\pi}\)。
        \item 由偶性，半轴积分等于全轴积分的一半，故 \(\displaystyle \int_{0}^{+\infty}\mathrm{e}^{-x^2}\,\mathrm{d}x=\frac{\sqrt{\pi}}{2}\)。
    \end{itemize}
\end{compactenum}
\end{solution}

\begin{exercise}[偏心圆域的极坐标积分]
求二重积分 \(\displaystyle I=\iint_D\sqrt{x^2+y^2}\,\mathrm{d}x\,\mathrm{d}y\)。其中 $D$ 由上半圆 $(x-a)^2+y^2=a^2$、上半圆 $(x-2a)^2+y^2=4a^2$ 与直线 $y=x$ 围成，且 $a>0$。
\end{exercise}
\begin{solution}
\begin{compactenum}
    \item \textbf{将直角边界方程转化为极坐标方程}：
    \begin{itemize}
        \item 内侧小圆：$(x-a)^2+y^2=a^2 \implies x^2+y^2-2ax=0 \implies r^2-2ar\cos\phi = 0 \implies r = 2a\cos\phi$。
        \item 外侧大圆：$(x-2a)^2+y^2=4a^2 \implies x^2+y^2-4ax=0 \implies r^2-4ar\cos\phi = 0 \implies r = 4a\cos\phi$。
        \item 放射状边界线：$y=x$ 对应于极角 $\phi = \frac{\pi}{4}$；而“上半圆”约束意味着图形最高止于 $y$ 轴，即极角上限为 $\phi = \frac{\pi}{2}$。
    \end{itemize}
    \item \textbf{构建极坐标下的累次积分域（$\phi$-型区域）}：
    \begin{itemize}
        \item 角度的辐射范围为固定的夹角：$\frac{\pi}{4} \leqslant \phi \leqslant \frac{\pi}{2}$。
        \item 对于此区间内的任意射线 $\phi$，它首先穿过内圆 $r=2a\cos\phi$，然后到达外圆 $r=4a\cos\phi$ 停止。
        \item 极径范围为：$2a\cos\phi \leqslant r \leqslant 4a\cos\phi$。
    \end{itemize}
    \item \textbf{设定积分并计算内层}：
    \begin{itemize}
        \item 注意被积函数为 $\sqrt{r^2}=r$，面积微元为 $r\mathrm{d}r\,\mathrm{d}\phi$，总被积体为 $r^2$：
        \(\displaystyle I = \int_{\frac{\pi}{4}}^{\frac{\pi}{2}} \mathrm{d}\phi \int_{2a\cos\phi}^{4a\cos\phi} r^2 \,\mathrm{d}r\)。
        \item 计算关于 $r$ 的内层积分：
        \[ \int_{2a\cos\phi}^{4a\cos\phi} r^2 \,\mathrm{d}r = \left[ \frac{r^3}{3} \right]_{2a\cos\phi}^{4a\cos\phi} = \frac{(4a\cos\phi)^3 - (2a\cos\phi)^3}{3} = \frac{64a^3 - 8a^3}{3}\cos^3\phi = \frac{56a^3}{3}\cos^3\phi \]
    \end{itemize}
    \item \textbf{完成外层三角函数定积分}：
    \begin{itemize}
        \item 将常数提至外侧：$I = \frac{56a^3}{3} \int_{\frac{\pi}{4}}^{\frac{\pi}{2}} \cos^3\phi \,\mathrm{d}\phi$。
        \item 降次凑微分法处理奇次幂：$\cos^3\phi = \cos^2\phi \cdot \cos\phi = (1-\sin^2\phi)\cos\phi$：
        \(\displaystyle \int \cos^3\phi \,\mathrm{d}\phi = \int (1-\sin^2\phi) \,\mathrm{d}(\sin\phi) = \sin\phi - \frac{1}{3}\sin^3\phi\)。
        \item 代入积分限 $[\frac{\pi}{4}, \frac{\pi}{2}]$：
        \[
        \left( 1 - \frac{1}{3}(1)^3 \right) - \left( \frac{\sqrt{2}}{2} - \frac{1}{3}\left(\frac{\sqrt{2}}{2}\right)^3 \right) = \frac{2}{3} - \left( \frac{\sqrt{2}}{2} - \frac{2\sqrt{2}}{24} \right) = \frac{2}{3} - \frac{5\sqrt{2}}{12}
        \]
    \end{itemize}
    \item \textbf{最终合并结果}：
    \[
    I = \frac{56a^3}{3} \left( \frac{8 - 5\sqrt{2}}{12} \right)
    = \frac{14a^3(8 - 5\sqrt{2})}{9}
    = \frac{112 - 70\sqrt{2}}{9}a^3.
    \]
\end{compactenum}
\end{solution}


\subsection{广义重积分专题：无界区域与反常性判断}

\begin{definition}[无界区域上的广义二重积分] \label{def:improper-double-integral}
设 $D \subset \mathbb{R}^{2}$ 是一个无界区域，函数 $f(x, y)$ 在 $D$ 上有定义且在任何有界子区域上可积。
我们在 $D$ 内任意取定一个边界为光滑闭曲线的有界闭区域 $\sigma$。
\begin{compactenum}
    \item \textbf{构造截断积分}：在有界闭区域 $\sigma$ 上计算常规的二重积分 $\iint\limits_{\sigma} f(x, y) \,\mathrm{d}\sigma$。
    \item \textbf{向无界区域扩张}：让有界区域 $\sigma$ 的边界曲线连续变动，使得 $\sigma$ 无限向外扩展并最终趋近于整个无界区域 $D$。
    \item \textbf{取极限判定收敛性}：无论 $\sigma$ 的形状如何，也不论其扩展趋于 $D$ 的方式和路径怎样，如果极限 \(\displaystyle \lim_{\sigma \to D}\iint\limits_{\sigma}f(x,y)\,\mathrm{d}\sigma\) 恒存在且等于同一个常数 $I$，则称该极限值 $I$ 为函数 $f(x, y)$ 在无界区域 $D$ 上的\textbf{广义二重积分}，记作 \(\displaystyle I=\iint\limits_D f(x,y)\,\mathrm{d}\sigma\)。此时我们称该广义二重积分\textbf{收敛}；如果极限不存在（或趋于无穷），则称该积分\textbf{发散}。
\end{compactenum}
\end{definition}

为了简化计算，对于非负函数，我们可以直接利用类似于 Fubini 定理的极限形式进行累次积分。以下是两个判定与计算非负函数广义二重积分的核心定理。

\begin{exercise}[半无界扇形区域积分]
计算广义二重积分 \(\displaystyle \iint\limits_{\substack{0 \leqslant x \leqslant y}} \mathrm{e}^{-(x+y)} \,\mathrm{d}\sigma\)。
\end{exercise}
\begin{solution}
\begin{compactenum}
    \item \textbf{分析并写成累次积分}：区域 $D=\{(x,y):0\le x\le y<+\infty\}$ 是第一象限中位于直线 $y=x$ 上方的无界楔形区域，且被积函数非负，故可直接用广义 Fubini 定理写成
    \(\displaystyle I = \int_{0}^{+\infty} \mathrm{d}x\int_{x}^{+\infty} \mathrm{e}^{-(x+y)} \,\mathrm{d}y\)。
    \item \textbf{计算}：先算内层并代回，\(\displaystyle \int_{x}^{+\infty} \mathrm{e}^{-(x+y)} \,\mathrm{d}y=\mathrm{e}^{-x}\int_{x}^{+\infty}\mathrm{e}^{-y}\,\mathrm{d}y=\mathrm{e}^{-2x}\)。因此 \(\displaystyle I=\int_{0}^{+\infty}\mathrm{e}^{-2x}\,\mathrm{d}x=\left[-\frac{1}{2}\mathrm{e}^{-2x}\right]_{0}^{+\infty}=\frac{1}{2}\)。
\end{compactenum}
\end{solution}

\begin{exercise}[全平面高斯积分与概率积分推导]
计算广义二重积分 \(\displaystyle I = \int_{-\infty}^{+\infty} \int_{-\infty}^{+\infty} \mathrm{e}^{-(x^{2}+y^{2})} \,\mathrm{d}x\,\mathrm{d}y\)，
并由此证明概率积分公式 $\displaystyle \frac{1}{\sqrt{\pi}} \int_{-\infty}^{+\infty} \mathrm{e}^{-x^{2}} \,\mathrm{d}x = 1$。
\end{exercise}
\begin{solution}
\begin{compactenum}
    \item \textbf{引入极坐标变换计算二重反常积分}：在极坐标下，整个平面写成 $\phi \in [0,2\pi],\ r \in [0,+\infty)$，故
    \[
    I = \int_{0}^{2\pi} \mathrm{d}\phi \int_{0}^{+\infty} \mathrm{e}^{-r^{2}} r \,\mathrm{d}r.
    \]
    内层积分给出 \(\displaystyle \int_{0}^{+\infty} \mathrm{e}^{-r^2} r \,\mathrm{d}r
    = \left[-\frac{1}{2}\mathrm{e}^{-r^2}\right]_{0}^{+\infty}
    = \frac{1}{2}\)，
    从而 $I=\pi$。
    \item \textbf{导出概率积分结果}：回到直角坐标，因 $\mathrm{e}^{-(x^2+y^2)}=\mathrm{e}^{-x^2}\mathrm{e}^{-y^2}$，故
    \[
    I=\left( \int_{-\infty}^{+\infty} \mathrm{e}^{-x^2}\,\mathrm{d}x \right)
      \left( \int_{-\infty}^{+\infty} \mathrm{e}^{-y^2}\,\mathrm{d}y \right).
    \]
    设这两个相等的积分都为 $J$，则 $J^2=I=\pi$。又因被积函数恒正，所以 $J=\sqrt{\pi}$，即
    \[
    \int_{-\infty}^{+\infty} \mathrm{e}^{-x^2}\,\mathrm{d}x = \sqrt{\pi},\qquad
    \frac{1}{\sqrt{\pi}} \int_{-\infty}^{+\infty} \mathrm{e}^{-x^2}\,\mathrm{d}x = 1.
    \]
\end{compactenum}
\end{solution}

\begin{exercise}[含 $\min$ 函数的全平面反常积分（直角坐标法）]
计算广义二重积分
\[
I = \int_{-\infty}^{+\infty} \int_{-\infty}^{+\infty}
\min(x, y) \mathrm{e}^{-(x^{2}+y^{2})} \,\mathrm{d}x\,\mathrm{d}y .
\]
\end{exercise}
\begin{solution}
\begin{compactenum}
    \item \textbf{依据 $\min$ 函数定义拆分无界区域}：
    \begin{itemize}
        \item 积分区域为全平面 $\mathbb{R}^2$。直线 $y=x$ 是判断 $\min(x,y)$ 取值的自然分界线。
        \item 令 $D_1 = \{ (x, y) : y \geqslant x \}$，在此区域内 $x \leqslant y$，故 $\min(x, y) = x$。
        \item 令 $D_2 = \{ (x, y) : y \leqslant x \}$，在此区域内 $y \leqslant x$，故 $\min(x, y) = y$。
        \item 积分拆分为 \(\displaystyle I = \iint\limits_{D_1} x \mathrm{e}^{-(x^2+y^2)} \,\mathrm{d}\sigma+\iint\limits_{D_2} y \mathrm{e}^{-(x^2+y^2)} \,\mathrm{d}\sigma = I_1 + I_2\)。
    \end{itemize}
    \item \textbf{在直角坐标系下设定 $I_1$ 的积分限}：
    \begin{itemize}
        \item 将 $D_1$ 视为 $y$-型推广区域：对整个 $y$ 轴投影，$y \in (-\infty, +\infty)$；固定 $y$ 时，$x$ 从 $-\infty$ 到边界 $x = y$。
        \item 于是 \(\displaystyle I_1 = \int_{-\infty}^{+\infty} \mathrm{d}y \int_{-\infty}^{y} x \mathrm{e}^{-x^2-y^2} \,\mathrm{d}x=\int_{-\infty}^{+\infty} \mathrm{e}^{-y^2} \mathrm{d}y \int_{-\infty}^{y} x \mathrm{e}^{-x^2} \,\mathrm{d}x\)。
    \end{itemize}
    \item \textbf{计算 $I_1$ 的内层与外层反常积分}：
    \begin{itemize}
        \item 内层凑微分：
        \[
        \int_{-\infty}^{y} x \mathrm{e}^{-x^2} \,\mathrm{d}x
        =\left[ -\frac{1}{2}\mathrm{e}^{-x^2} \right]_{-\infty}^y
        =-\frac{1}{2}\mathrm{e}^{-y^2}.
        \]
        \item 代入外层：
        \[
        I_1 = \int_{-\infty}^{+\infty} \mathrm{e}^{-y^2}
        \left(-\frac{1}{2}\mathrm{e}^{-y^2}\right) \mathrm{d}y
        =-\frac{1}{2} \int_{-\infty}^{+\infty} \mathrm{e}^{-2y^2} \,\mathrm{d}y.
        \]
        \item 进行变量代换求外层：令 $\sqrt{2}y = u$，则 $\mathrm{d}y = \frac{1}{\sqrt{2}}\mathrm{d}u$。积分区间不变，得
        \[
        I_1 = -\frac{1}{2\sqrt{2}} \int_{-\infty}^{+\infty} \mathrm{e}^{-u^2} \,\mathrm{d}u
        = -\frac{\sqrt{2\pi}}{4}.
        \]
    \end{itemize}
    \item \textbf{利用对称性求 $I_2$ 并汇总}：
    \begin{itemize}
        \item 对于 $I_2 = \iint\limits_{D_2} y \mathrm{e}^{-(x^2+y^2)} \,\mathrm{d}\sigma$，若作变量互换 $(x,y)\mapsto(y,x)$，那么区域条件 $y\leqslant x$ 会变成 $x\leqslant y$，也就是恰好变为 $D_1$。
        \item 同时，被积函数中的 $y\mathrm{e}^{-(x^2+y^2)}$ 也随变量互换变成 $x\mathrm{e}^{-(x^2+y^2)}$，因此整个积分完全转化为 $I_1$。
        \item 所以 $I_2 = I_1 = -\frac{\sqrt{2\pi}}{4}$。
        \item 总积分 $I = I_1 + I_2 = -\frac{\sqrt{2\pi}}{4} - \frac{\sqrt{2\pi}}{4} = -\frac{\sqrt{2\pi}}{2}$。
    \end{itemize}

    遇到 $\min(x,y)$、$\max(x,y)$ 这类分段函数时，先找出“分界直线或分界曲线”；若分界后两个子区域能通过交换坐标彼此对应，就往往可以只算一半，再用对称或变量互换补回另一半。
\end{compactenum}
\end{solution}


\subsection{三重积分专题：直角坐标、柱坐标与球坐标}

\begin{definition}[三重积分的定义] \label{def:triple-integral}
设 $\Omega \subset \mathbb{R}^{3}$ 为一个有界闭区域，函数 $f(x,y,z)$ 在 $\Omega$ 上有定义。
把 $\Omega$ 任意分成 $n$ 个小块 $\Delta V_i$，在每个 $\Delta V_i$ 内任取一点
$(\xi_i,\eta_i,\zeta_i)$。记 $\lambda$ 为所有小块直径的最大者。若极限 \(\displaystyle \lim_{\lambda \to 0}\sum_{i=1}^{n} f(\xi_i,\eta_i,\zeta_i)\Delta V_i\) 存在，且极限值与分割方式及取点无关，则称该极限值为函数 $f(x,y,z)$ 在闭区域 $\Omega$ 上的\textbf{三重积分}，记作
\[
\iiint\limits_{\Omega} f(x,y,z)\,\mathrm{d}V
\quad\text{或}\quad
\iiint\limits_{\Omega} f(x,y,z)\,\mathrm{d}x\,\mathrm{d}y\,\mathrm{d}z.
\]
其中 $f(x,y,z)$ 称为被积函数，$\Omega$ 称为积分区域，
\(\mathrm{d}V\)（在直角坐标系下即 \(\mathrm{d}x\,\mathrm{d}y\,\mathrm{d}z\)）称为体积元素。
\end{definition}

\begin{theorem}[三重积分的存在性定理] \label{thm:triple-integral-existence}
若函数 $f(x, y, z)$ 在有界闭区域 $\Omega$ 上连续，则 $f(x, y, z)$ 在 $\Omega$ 上的三重积分必定存在。
\end{theorem}

\begin{remark}[三重积分的几何意义退化]
当被积函数恒为 $1$，即 $f(x, y, z) \equiv 1$ 时，三重积分 $\iiint\limits_{\Omega} 1 \,\mathrm{d}V$ 在数值上恰好等于空间闭区域 $\Omega$ 的\textbf{体积}。
\end{remark}

\subsubsection{利用直角坐标系计算三重积分}

我们完全可以模仿二重积分的思路，将三重积分的理论建立起来。计算的核心思想是将多重积分转化为多次单变量积分（累次积分）。
与二重积分（只有 $2$ 种积分顺序）不同的是，三重积分在直角坐标系下共有 $3! = 6$ 种不同的积分顺序（如先 $z$ 后 $y$ 再 $x$，先 $x$ 后 $z$ 再 $y$ 等）。

写积分限真正要掌握的不是背几个名字，而是会把区域按某个次序一层一层切开。二重积分是“外层投影 + 内层线段”，三重积分就是“外层投影 + 中层截面 + 内层线段”。

\begin{conclusion}[积分上下限的写法]
先定积分次序，再从外到内写限。最靠近被积函数的微分对应最先积分的变量；写在最左边的积分限对应最外层变量。

若决定按
\[
\mathrm{d}z\ \mathrm{d}y\ \mathrm{d}x
\]
的顺序积分，也就是先对 $z$ 积分，再对 $y$ 积分，最后对 $x$ 积分，那么标准写法应是
\[
\iiint_{\Omega}f(x,y,z)\,\mathrm{d}V
=
\int_{a}^{b}\mathrm{d}x
\int_{\varphi_1(x)}^{\varphi_2(x)}\mathrm{d}y
\int_{\psi_1(x,y)}^{\psi_2(x,y)}
f(x,y,z)\,\mathrm{d}z.
\]
它翻译成区域语言就是
\[
\Omega=
\left\{
(x,y,z):
a\le x\le b,\quad
\varphi_1(x)\le y\le \varphi_2(x),\quad
\psi_1(x,y)\le z\le \psi_2(x,y)
\right\}.
\]
这里有一个很实用的检查规则：$x$ 的限只能是常数；$y$ 的限可以含 $x$，但不能含 $y,z$；$z$ 的限可以含 $x,y$，但不能含 $z$。换成别的积分次序时，也按同样原则检查：某个变量的上下限只能含比它更外层的变量。
\end{conclusion}

\begin{exercise}[用先一再一再一法计算四面体积分]
设积分区域 $\Omega$ 由坐标面 $x=0, y=0, z=0$ 以及平面 $x+y+z=1$ 围成。分别求 \(\displaystyle I=\iiint\limits_{\Omega} z^{2}\,\mathrm{d}V,\ J=\iiint\limits_{\Omega} x^{2}\,\mathrm{d}V\)。
\end{exercise}
\begin{solution}
\begin{compactenum}
    \item \textbf{先把区域逐层写出来}：平面方程为 \(x+y+z=1\)，且四面体位于第一卦限，所以
    \[
    0\le x\le 1,\qquad
    0\le y\le 1-x,\qquad
    0\le z\le 1-x-y.
    \]
    这就是先定一个变量的总范围，再定第二个变量的范围，最后定第三个变量的范围。按这个次序写成
    \[
    \iiint_{\Omega}f(x,y,z)\,\mathrm{d}V
    =
    \int_0^1\mathrm{d}x
    \int_0^{1-x}\mathrm{d}y
    \int_0^{1-x-y}f(x,y,z)\,\mathrm{d}z.
    \]
    \item \textbf{计算 $I$}：把 \(f(x,y,z)=z^2\) 代入上式，得
    \[
    \begin{aligned}
    I
    &=\int_0^1\mathrm{d}x
      \int_0^{1-x}\mathrm{d}y
      \int_0^{1-x-y}z^2\,\mathrm{d}z\\
    &=\frac13\int_0^1\mathrm{d}x
      \int_0^{1-x}(1-x-y)^3\,\mathrm{d}y\\
    &=\frac1{12}\int_0^1(1-x)^4\,\mathrm{d}x
    =\frac{1}{60}.
    \end{aligned}
    \]
    \item \textbf{计算 $J$}：继续用同一组积分限，不再换方法：
    \[
    \begin{aligned}
    J
    &=\int_0^1\mathrm{d}x
      \int_0^{1-x}\mathrm{d}y
      \int_0^{1-x-y}x^2\,\mathrm{d}z\\
    &=\int_0^1\mathrm{d}x
      \int_0^{1-x}x^2(1-x-y)\,\mathrm{d}y\\
    &=\frac12\int_0^1 x^2(1-x)^2\,\mathrm{d}x\\
    &=\frac12\left(\frac13-\frac12+\frac15\right)
    =\frac{1}{60}.
    \end{aligned}
    \]
    当然，由于区域 \(x+y+z\le1,\ x,y,z\ge0\) 对 \(x,y,z\) 完全对称，也可在算出 \(I\) 后直接得到 \(J=I\)。但初学列限时，最好先按同一套限完整写一遍，手感会更稳。
\end{compactenum}
\end{solution}

\begin{exercise}[利用对称性简化椭球体积分]
计算三重积分 \(\displaystyle I=\iiint\limits_{\Omega}(x+y+z)^{2}\,\mathrm{d}x\,\mathrm{d}y\,\mathrm{d}z\)，其中积分区域 $\Omega$ 为实心椭球体：
\(\displaystyle \Omega = \left\{ (x, y, z): \frac{x^{2}}{a^{2}} + \frac{y^{2}}{b^{2}} + \frac{z^{2}}{c^{2}} \leqslant 1 \right\}\)。
\end{exercise}
\begin{solution}
\begin{compactenum}
    \item \textbf{先用对称性消去交叉项}：展开得 \(\displaystyle (x+y+z)^2=x^2+y^2+z^2+2xy+2yz+2xz\)。
    椭球 $\Omega$ 关于三个坐标平面都对称，故奇函数项积分全为 $0$，即 \(\displaystyle \iiint_{\Omega}2xy\,\mathrm{d}V=\iiint_{\Omega}2yz\,\mathrm{d}V=\iiint_{\Omega}2xz\,\mathrm{d}V=0\)，从而
    \(\displaystyle I=\iiint\limits_{\Omega}x^{2}\,\mathrm{d}V+\iiint\limits_{\Omega}y^{2}\,\mathrm{d}V+\iiint\limits_{\Omega}z^{2}\,\mathrm{d}V=:I_{1}+I_{2}+I_{3}\)。
    \item \textbf{计算 $I_{3}$}：对固定 $z\in[-c,c]$，截面 \(\displaystyle \frac{x^2}{a^2}+\frac{y^2}{b^2}\leqslant 1-\frac{z^2}{c^2}\) 是椭圆，其面积为 $S(z)=\pi ab\left(1-\frac{z^{2}}{c^{2}}\right)$，故
    \(\displaystyle I_{3}=\int_{-c}^{c}z^{2}S(z)\,\mathrm{d}z
    =2\pi ab\int_{0}^{c}\left(z^{2}-\frac{z^{4}}{c^{2}}\right)\mathrm{d}z
    =\frac{4}{15}\pi abc^{3}\)。
    \item \textbf{轮换坐标得到 $I_{1},I_{2}$}：完全类似地，\(\displaystyle I_{1}=\frac{4}{15}\pi a^{3}bc,\ I_{2}=\frac{4}{15}\pi ab^{3}c\)。
    \item \textbf{合并结果}：\(\displaystyle I=I_{1}+I_{2}+I_{3}=\frac{4}{15}\pi abc(a^{2}+b^{2}+c^{2})\)。
\end{compactenum}
\end{solution}

\subsubsection{三重积分的换元法}

类似于二重积分的极坐标变换，在三维空间中，当积分区域是圆柱体、圆锥体、球体或它们的组合时，使用\textbf{柱坐标系}或\textbf{球坐标系}能够极大地简化积分限和被积函数。

\begin{theorem}[三重积分的一般换元公式] \label{thm:triple-integral-substitution}
令 $x=x(u, v, w), y=y(u, v, w), z=z(u, v, w)$，且该变换在区域内部是一一对应的，具有连续的一阶偏导数。若雅可比行列式：
\[ J = \frac{\partial(x, y, z)}{\partial(u, v, w)} = \det \begin{pmatrix} x_{u} & x_{v} & x_{w} \\ y_{u} & y_{v} & y_{w} \\ z_{u} & z_{v} & z_{w} \end{pmatrix} \neq 0 \]
则三重积分的换元公式为：
\[ \iiint\limits_{\Omega} f(x,y,z) \,\mathrm{d}x\,\mathrm{d}y\,\mathrm{d}z = \iiint\limits_{\Omega'} f\big(x(u,v,w), \dots\big) |J| \,\mathrm{d}u\,\mathrm{d}v\,\mathrm{d}w \]
\end{theorem}

柱坐标系 $(r, \theta, z)$ 实际上就是：在 $Oxy$ 平面上采用极坐标 $(r, \theta)$，而竖直高度 $z$ 保持不变。
\begin{itemize}
    \item \textbf{变换方程}：$x = r \cos \theta, \quad y = r \sin \theta, \quad z = z$
    \item \textbf{参数范围}：$0 \leqslant r < +\infty, \quad 0 \leqslant \theta \leqslant 2\pi, \quad -\infty < z < +\infty$
    \item \textbf{雅可比行列式绝对值}：$|J| = r$
    \item \textbf{体积元素}：$\mathrm{d}V = r \,\mathrm{d}r\,\mathrm{d}\theta\,\mathrm{d}z$
\end{itemize}

\begin{exercise}[柱坐标下抛物面与球面的结合]
计算三重积分 \(\displaystyle I = \iiint\limits_{\Omega} z \,\mathrm{d}x\,\mathrm{d}y\,\mathrm{d}z\)，
其中区域 $\Omega$ 由球面 $x^{2}+y^{2}+z^{2}=4$ 和旋转抛物面 $x^{2}+y^{2}=3z$ 共同围成（假设 $z \geqslant 0$ 的部分）。
\end{exercise}
\begin{solution}
\begin{compactenum}
    \item \textbf{求交线并确定投影}：联立
    \[
    \begin{cases}
    x^{2}+y^{2}+z^{2}=4,\\
    x^{2}+y^{2}=3z
    \end{cases}
    \]
    得 $z^{2}+3z-4=0$，故 $z=1$ 或 $z=-4$。又因抛物面满足 $z\geqslant 0$，故只取 $z=1$，于是交线为 $x^{2}+y^{2}=3$，投影区域为 $D_{xy}=\{x^{2}+y^{2}\leqslant 3\}$。
    \item \textbf{改写为柱坐标}：令 $x=r\cos\theta,\ y=r\sin\theta$，则 \(\displaystyle 0\leqslant \theta \leqslant 2\pi,\qquad 0\leqslant r\leqslant \sqrt{3},\qquad \frac{r^{2}}{3}\leqslant z\leqslant \sqrt{4-r^{2}}\)。
    \item \textbf{建立并计算积分}：
    \begin{align*}
    I&=\int_{0}^{2\pi}\mathrm{d}\theta\int_{0}^{\sqrt{3}}\mathrm{d}r\int_{r^{2}/3}^{\sqrt{4-r^{2}}}zr\,\mathrm{d}z
    =2\pi\int_{0}^{\sqrt{3}}\frac{r}{2}\left(4-r^{2}-\frac{r^{4}}{9}\right)\mathrm{d}r\\
    &=2\pi\int_{0}^{\sqrt{3}}\left(2r-\frac{1}{2}r^{3}-\frac{1}{18}r^{5}\right)\mathrm{d}r
    =2\pi\left[r^{2}-\frac{1}{8}r^{4}-\frac{1}{108}r^{6}\right]_{0}^{\sqrt{3}}
    =\frac{13\pi}{4}.
    \end{align*}
\end{compactenum}
\end{solution}


球坐标系 $(r,\varphi,\theta)$ 中，$r$ 是点到原点的距离，$\varphi$ 是向径与 $z$ 轴正向的夹角（天顶角），$\theta$ 是向径在 $Oxy$ 平面上的投影与 $x$ 轴正向的夹角（方位角）。其变换公式与取值范围为
\[
\begin{array}{c}
x=r\sin\varphi\cos\theta,\qquad
y=r\sin\varphi\sin\theta,\qquad
z=r\cos\varphi,\\[2pt]
0\leqslant r<+\infty,\qquad
0\leqslant \varphi\leqslant \pi,\qquad
0\leqslant \theta\leqslant 2\pi,\qquad
|J|=r^{2}\sin\varphi.
\end{array}
\]
故体积元素为 \(\displaystyle \mathrm{d}V=r^{2}\sin\varphi\,\mathrm{d}r\,\mathrm{d}\varphi\,\mathrm{d}\theta\)。

\begin{exercise}[球坐标下计算锥面与球面间的积分]
计算三重积分 \(\displaystyle I = \iiint\limits_{\Omega} (x^{2}+y^{2}+z^{2}) \,\mathrm{d}V\)，
其中区域 $\Omega = \{ (x, y, z): x^{2}+y^{2}+z^{2} \leqslant 2z \}$。
\end{exercise}
\begin{solution}
\begin{compactenum}
    \item \textbf{确定球坐标积分域}：由
    \(\displaystyle x^{2}+y^{2}+z^{2}\leqslant 2z
    \iff x^{2}+y^{2}+(z-1)^{2}\leqslant 1\)
    可知 $\Omega$ 是球心在 $(0,0,1)$、半径为 $1$ 的实心球。改写成球坐标后，边界为 $r^{2}=2r\cos\varphi$，即 $r=2\cos\varphi$，因此 \(\displaystyle 0\leqslant \theta\leqslant 2\pi,\ 0\leqslant \varphi\leqslant \frac{\pi}{2},\ 0\leqslant r\leqslant 2\cos\varphi\)。
    \item \textbf{建立并计算积分}：被积函数化为 $r^{2}$，于是
    \begin{align*}
    I&=\int_{0}^{2\pi}\mathrm{d}\theta\int_{0}^{\pi/2}\mathrm{d}\varphi\int_{0}^{2\cos\varphi}r^{4}\sin\varphi\,\mathrm{d}r
    =2\pi\int_{0}^{\pi/2}\frac{32}{5}\cos^{5}\varphi\sin\varphi\,\mathrm{d}\varphi\\
    &=-\frac{64\pi}{5}\int_{0}^{\pi/2}\cos^{5}\varphi\,\mathrm{d}(\cos\varphi)
    =\frac{32\pi}{15}.
    \end{align*}
\end{compactenum}
\end{solution}

\begin{exercise}[体积计算 / 两圆柱体相交的牟合方盖问题]
求由两个底圆半径均为 $1$ 的直交圆柱面 $x^2+y^2=1$ 与 $x^2+z^2=1$ 所共同围成的立体 $\Omega$ 的体积。
\end{exercise}
\begin{solution}
\begin{compactenum}
    \item \textbf{分析立体的对称性}：
    这两个圆柱面在空间中正交相交，所形成的立体（中国古代数学家刘徽称其为“牟合方盖”）具有极高的对称性。它关于三个坐标平面 $x=0, y=0, z=0$ 均完全对称。
    这意味着立体被三个坐标平面切成了全等的 $8$ 块，因此只计算第一卦限（$x \geqslant 0, y \geqslant 0, z \geqslant 0$）内的体积，再乘以 $8$，不会遗漏也不会重复。
    \item \textbf{确定积分区域与上下边界}：

    第一卦限内的底面是四分之一圆域 \(\displaystyle D=\{(x,y):x^2+y^2\leqslant 1,\ x\geqslant 0,\ y\geqslant 0\}\)。对任意 $(x,y)\in D$，下边界为 $z=0$，而由圆柱面 $x^2+z^2=1$ 可得上边界 $z=\sqrt{1-x^2}$。
    \item \textbf{建立二重积分并化为累次积分}：
    根据体积公式，第一卦限的体积为 $\iint_D \sqrt{1-x^2} \,\mathrm{d}x\,\mathrm{d}y$。总体积 $V$ 为：\(\displaystyle V = 8 \iint\limits_{D} \sqrt{1-x^2} \,\mathrm{d}x\,\mathrm{d}y = 8 \int_{0}^{1} \mathrm{d}x \int_{0}^{\sqrt{1-x^2}} \sqrt{1-x^2} \,\mathrm{d}y\)。
    \textit{（这里特意先对 $y$ 积分，是因为被积函数只含 $x$，这样内层积分只是“常数乘以切片长度”；若反过来先对 $x$，反而会把根式留到更麻烦的位置。）}
    \item \textbf{计算积分求得体积}：

    内层关于 $y$ 的积分可直接算出：
    \(\displaystyle \int_{0}^{\sqrt{1-x^2}} \sqrt{1-x^2} \,\mathrm{d}y=1-x^2\)。
    进而 \(\displaystyle V=8\int_{0}^{1}(1-x^2)\,\mathrm{d}x
    =8\left[x-\frac{1}{3}x^3\right]_{0}^{1}=\frac{16}{3}\)。

    对称立体求体积时，可先看能否缩到一个卦限或一个扇区；缩区以后，再比较哪一个变量会让切片长度最简单。这类题很多时候不是靠高级换元，而是靠对称性加上正确的切片方向。
\end{compactenum}
\end{solution}

\begin{exercise}[体积计算 / 柱坐标系下的复杂曲面夹层]
求由向下开口的旋转抛物面 $z = 6-x^2-y^2$ 和向上开口的圆锥面 $z = \sqrt{x^2+y^2}$ 所共同围成立体的体积。
\end{exercise}
\begin{solution}
\begin{compactenum}
    \item \textbf{联立方程寻找交线及投影区域}：
    \begin{itemize}
        \item 先判断上下位置：在交线内部取 $r=0$，此时抛物面给出 $z=6$，圆锥面给出 $z=0$，所以抛物面在上方构成“盖”，圆锥面在下方构成“底”。
        \item 为找到两曲面的交线，联立它们的方程。令极径 $r = \sqrt{x^2+y^2}$（且 $r \geqslant 0$），则两方程变为：上边界 $z = 6-r^2$，下边界 $z = r$。
        \item 联立得：$6-r^2 = r \implies r^2+r-6 = 0 \implies (r+3)(r-2) = 0$。
        \item 因为 $r \geqslant 0$，解得交线投影圆的半径为 $r = 2$。立体在 $Oxy$ 平面上的总投影 $D$ 是圆域 $x^2+y^2 \leqslant 4$。
    \end{itemize}
    \item \textbf{引入柱坐标系构建三重积分}：
    \begin{itemize}
        \item 利用柱坐标系 $(r, \theta, z)$ 来描述该立体 $\Omega'$ 非常自然，此时 \(\displaystyle 0 \leqslant \theta \leqslant 2\pi,\ 0 \leqslant r \leqslant 2,\ r \leqslant z \leqslant 6-r^2\)。
        \item 体积 $V = \iiint_{\Omega'} r \,\mathrm{d}r\,\mathrm{d}\theta\,\mathrm{d}z$。
    \end{itemize}
    \item \textbf{转化为累次积分并计算}：
    \begin{itemize}
        \item 根据区域的不等式边界，建立三次积分：\(\displaystyle V = \int_{0}^{2\pi} \mathrm{d}\theta \int_{0}^{2} \mathrm{d}r \int_{r}^{6-r^2} 1 \cdot r \,\mathrm{d}z\)。
        \item 外层对 $\theta$ 积分独立，直接提出 $2\pi$；内层对 $z$ 积分为 \(\displaystyle \int_{r}^{6-r^2} r \,\mathrm{d}z = r [z]_r^{6-r^2} = 6r-r^3-r^2\)。
        \item 代入中间层关于 $r$ 的积分：\(\displaystyle V = 2\pi \int_{0}^{2} (6r - r^2 - r^3) \,\mathrm{d}r = 2\pi \left[ 3r^2 - \frac{1}{3}r^3 - \frac{1}{4}r^4 \right]_0^2\)。
        \item 代入上限 $r=2$ 计算：\(\displaystyle V = 2\pi \left( 12 - \frac{8}{3} - 4 \right)=\frac{32\pi}{3}\)。
    \end{itemize}
\end{compactenum}
\end{solution}

\begin{note}[空间物体的质心坐标公式]
设有一连续分布的物体，占据空间区域 $V$，其密度函数为 $\rho(x,y,z)$。则该物体的总质量为 \(\displaystyle M = \iiint\limits_{V} \rho(x,y,z) \,\mathrm{d}V\)。其质心坐标 $(\bar{x}, \bar{y}, \bar{z})$ 分别为物体对三个坐标平面的静力矩除以总质量：
\[ \bar{x} = \frac{1}{M}\iiint\limits_{V} x\rho(x,y,z)\,\mathrm{d}V, \quad \bar{y} = \frac{1}{M}\iiint\limits_{V} y\rho(x,y,z)\,\mathrm{d}V, \quad \bar{z} = \frac{1}{M}\iiint\limits_{V} z\rho(x,y,z)\,\mathrm{d}V \]
\textit{（若物体为几何均匀分布，则 $\rho$ 为常数，可在分子分母中消去，此时质心即为几何中心（形心）。）}
\end{note}

\begin{exercise}[质心计算 / 利用对称性与柱坐标系]
设均匀物体 $\Omega$ 由旋转抛物面 $z=x^2+y^2$ 与顶平面 $z=1$ 围成。求该物体的质心坐标。
\end{exercise}
\begin{solution}
\begin{compactenum}
    \item \textbf{分析几何对称性并确定部分坐标}：
    因为物体是均匀的（设恒定密度为 $\rho$），且区域 $\Omega$ 由 $z=x^2+y^2$ 与 $z=1$ 构成，图形关于 $xOz$ 平面与 $yOz$ 平面均完全对称（即关于 $z$ 轴旋转对称）。
    更具体地说，若把区域中任一点 $(x,y,z)$ 关于平面 $yOz$ 反射，就得到另一点 $(-x,y,z)$，而这两个点质量相同、对 $\bar{x}$ 的贡献互为相反数，所以 $x$ 方向的平均坐标必为零；同理可得 $\bar{y}=0$。因此质心只能落在 $z$ 轴上，我们只需求 $\bar{z}$。
    \item \textbf{建立柱坐标系描述区域}：
    立体在 $Oxy$ 面的投影是单位圆域；在柱坐标系下，\(\displaystyle 0\leqslant \theta\leqslant 2\pi,\qquad 0\leqslant r\leqslant 1,\qquad r^2\leqslant z\leqslant 1\)，
    即 $\Omega'=\{(r,\theta,z):0\leqslant \theta\leqslant 2\pi,\ 0\leqslant r\leqslant 1,\ r^2\leqslant z\leqslant 1\}$。
    \item \textbf{计算物体的总质量 $M$}：
    \[
    M = \iiint\limits_{\Omega} \rho \,\mathrm{d}V = \rho \int_{0}^{2\pi} \mathrm{d}\theta \int_{0}^{1} \mathrm{d}r \int_{r^2}^{1} r \,\mathrm{d}z
    = 2\pi\rho \int_{0}^{1} r(1-r^2) \,\mathrm{d}r = \frac{\pi\rho}{2}.
    \]
    \item \textbf{计算物体对 $xOy$ 面的静力矩 $M_{xy}$}：
    \[
    M_{xy} = \iiint\limits_{\Omega} z\rho \,\mathrm{d}V = \rho \int_{0}^{2\pi} \mathrm{d}\theta \int_{0}^{1} \mathrm{d}r \int_{r^2}^{1} z \cdot r \,\mathrm{d}z
    = 2\pi\rho \int_{0}^{1} r \left[ \frac{1}{2}z^2 \right]_{r^2}^1 \mathrm{d}r
    = \pi\rho \int_{0}^{1} (r-r^5)\,\mathrm{d}r = \frac{\pi\rho}{3}.
    \]
    \item \textbf{求取 $\bar{z}$ 坐标得出最终结果}：
    \(\displaystyle \bar{z} = \frac{M_{xy}}{M} = \frac{\pi\rho/3}{\pi\rho/2} = \frac{2}{3}\)
    综上所述，该均匀物体的质心坐标为 $\left( 0, 0, \frac{2}{3} \right)$。

    求质心时遇到旋转体或关于坐标平面对称的立体，都应先用对称性消掉一部分坐标分量，再把主要精力放在剩余分量的积分上。这样既能减少计算量，也能避免一开始就列三个几乎相同的积分。
\end{compactenum}
\end{solution}

\begin{note}[空间与平面物体的转动惯量]
\textbf{空间物体（三重积分）}对坐标轴的转动惯量为
\[
I_x = \iiint\limits_{V} (y^2+z^2)\rho(x,y,z)\,\mathrm{d}V, \quad I_y = \iiint\limits_{V} (x^2+z^2)\rho(x,y,z)\,\mathrm{d}V, \quad I_z = \iiint\limits_{V} (x^2+y^2)\rho(x,y,z)\,\mathrm{d}V,
\]
其中被积函数中的多项式正是空间点到对应坐标轴距离的平方。\textbf{平面薄片（二重积分）}分布在 $Oxy$ 平面时，上式退化为 \(\displaystyle I_x = \iint\limits_{D} y^2\rho(x,y)\,\mathrm{d}\sigma,\quad I_y = \iint\limits_{D} x^2\rho(x,y)\,\mathrm{d}\sigma\)。
\end{note}

\begin{exercise}[平面区域转动惯量 / 半圆薄片问题]
求半径为 $a$ 的均匀半圆薄片对其直边直径所在轴的转动惯量（假设密度 $\rho \equiv 1$）。
\end{exercise}
\begin{solution}
\begin{compactenum}
    \item \textbf{建立坐标系与区域描述}：
    为了计算简便，我们以半圆的直径所在的直线作为 $x$ 轴，半圆的圆心作为原点。
    由于薄片是半圆形且位于 $x$ 轴一侧（设位于上半平面），积分区域为
    \(D = \{ (x, y) : x^2+y^2 \leqslant a^2, y \geqslant 0 \}\)。
    \item \textbf{转化为极坐标系表示}：
    在极坐标系下，上半圆的极角范围为 $\theta \in [0, \pi]$，极径范围为 $r \in [0, a]$。
    要求对 $x$ 轴的转动惯量 $I_x$，根据公式，被积函数为 $y^2 = (r\sin\theta)^2 = r^2\sin^2\theta$。面积微元为 $r\mathrm{d}r\,\mathrm{d}\theta$。
    \item \textbf{构建并计算二重积分}：
    \(\displaystyle I_x = \iint\limits_{D} y^2 \,\mathrm{d}\sigma = \int_{0}^{\pi} \mathrm{d}\theta \int_{0}^{a} (r^2\sin^2\theta) \cdot r \,\mathrm{d}r = \int_{0}^{\pi} \sin^2\theta \,\mathrm{d}\theta \int_{0}^{a} r^3 \,\mathrm{d}r\)。
    \item \textbf{分别计算独立积分}：

    两个独立积分分别为 \(\displaystyle \int_{0}^{a} r^3 \,\mathrm{d}r = \left[ \frac{1}{4}r^4 \right]_0^a = \frac{1}{4}a^4,\quad
    \int_{0}^{\pi} \sin^2\theta \,\mathrm{d}\theta = \frac{1}{2} \int_{0}^{\pi} (1 - \cos 2\theta) \,\mathrm{d}\theta = \frac{\pi}{2}\)。
    \item \textbf{得出最终结果}：\(\displaystyle I_x = \frac{\pi}{2} \times \frac{a^4}{4} = \frac{\pi a^4}{8}\)。
\end{compactenum}
\end{solution}

\begin{note}[空间物体的引力计算公式]
设空间中占据区域 $V$ 的物体具有连续密度分布 $\rho(x,y,z)$。依据牛顿万有引力定律，体积微元 $\mathrm{d}V$ 内的微质量对位于点 $P_0(x_0,y_0,z_0)$ 处单位质量质点产生的引力微元为
\(\displaystyle \mathrm{d}\vec{F} = K \frac{\rho(x,y,z) \mathrm{d}V}{\|\vec{l}\|^3} \vec{l}\)。
其中 $K$ 为万有引力常数，$\vec{l} = (x-x_0, y-y_0, z-z_0)$ 是从目标点 $P_0$ 指向引力源微元点 $P(x,y,z)$ 的向量，$\|\vec{l}\| = \sqrt{(x-x_0)^2+(y-y_0)^2+(z-z_0)^2}$ 为两点间的距离。
将向量拆分为三个坐标分量，通过三重积分对全区域累加即可得到总引力向量：
\[ F_x = \iiint\limits_{V} K\rho(x,y,z)\frac{x-x_0}{\|\vec{l}\|^3}\,\mathrm{d}V, \quad F_y = \iiint\limits_{V} K\rho(x,y,z)\frac{y-y_0}{\|\vec{l}\|^3}\,\mathrm{d}V, \quad F_z = \iiint\limits_{V} K\rho(x,y,z)\frac{z-z_0}{\|\vec{l}\|^3}\,\mathrm{d}V \]
\end{note}

\begin{exercise}[引力计算 / 球顶圆锥对顶点的吸引力]
设有一均匀的空间立体 $\Omega$（密度设为 $\rho \equiv 1$），其形状为一个球顶圆锥体（如同冰淇淋的顶部）：其底部锥面的顶点位于原点 $O(0,0,0)$，对称轴为 $z$ 轴正半轴，锥面母线与 $z$ 轴的夹角为常数 $\alpha$；其顶部被以原点为球心、半径为 $R$ 的球面封顶。
求该立体对位于原点（圆锥顶点）处单位质量质点的引力向量 $\vec{F}$。
\end{exercise}
\begin{solution}
\begin{compactenum}
    \item \textbf{分析几何对称性并简化分量}：

    被吸引点在原点 $P_0(0,0,0)$，此时 $\vec{l}=(x,y,z)$，$\|\vec{l}\|=\sqrt{x^2+y^2+z^2}$。又因 $\Omega$ 关于 $xOz$ 平面和 $yOz$ 平面对称，故
    \(\displaystyle F_x = \iiint_{\Omega} K\frac{x}{(x^2+y^2+z^2)^{3/2}} \,\mathrm{d}V = 0,\ F_y=0\)，
    从而只需求 \(\displaystyle F_z = \iiint\limits_{\Omega} K \frac{z}{(x^2+y^2+z^2)^{3/2}} \,\mathrm{d}V\)。
    \item \textbf{建立球坐标系界定区域}：

    在球坐标系下，圆锥面母线与 $z$ 轴夹角为 $\alpha$，故 $0\leqslant \varphi\leqslant \alpha$；顶面是半径为 $R$ 的球面，故 $0\leqslant r\leqslant R$；再绕 $z$ 轴转一周，故 $0\leqslant \theta\leqslant 2\pi$。
    \item \textbf{将被积函数与微元转为球坐标}：

    由 $z=r\cos\varphi$、$(x^2+y^2+z^2)^{3/2}=r^3$ 以及 $\mathrm{d}V=r^2\sin\varphi\,\mathrm{d}r\,\mathrm{d}\varphi\,\mathrm{d}\theta$，得
    \[
    F_z = \iiint\limits_{\Omega'} K \frac{r\cos\varphi}{r^3} (r^2\sin\varphi) \,\mathrm{d}r\,\mathrm{d}\varphi\,\mathrm{d}\theta
    = \iiint\limits_{\Omega'} K \sin\varphi\cos\varphi \,\mathrm{d}r\,\mathrm{d}\varphi\,\mathrm{d}\theta,
    \]
    即径向因子完全抵消。
    \item \textbf{分离变量计算积分}：
    因而
    \[
    F_z = K \left( \int_{0}^{2\pi} 1 \,\mathrm{d}\theta \right) \left( \int_{0}^{\alpha} \sin\varphi\cos\varphi \,\mathrm{d}\varphi \right) \left( \int_{0}^{R} 1 \,\mathrm{d}r \right)
    = K\cdot 2\pi \cdot \frac{1}{2}\sin^2\alpha \cdot R
    = \pi K R \sin^2\alpha.
    \]
    \item \textbf{得出最终引力向量}：
    该立体对顶点的总引力方向垂直向上，引力向量为 $\vec{F} = (0, 0, \pi K R \sin^2\alpha)$。

    空间向量积分题里，不必急着把三个分量都算一遍。只要区域具有旋转对称或关于坐标平面对称，通常都能先用“奇函数在对称区域上积分为零”的思想消掉横向分量，只留下真正需要计算的主轴方向。
\end{compactenum}
\end{solution}

\begin{exercise}[利用雅可比行列式的一般换元]
计算二重积分：
\(\displaystyle I = \iint\limits_{D} (x+y) \,\mathrm{d}x\,\mathrm{d}y\)
其中闭区域 $D$ 由开口向右的抛物线 $y^{2}=2x$ 以及两条互相平行的直线 $x+y=4$ 和 $x+y=2$ 所围成。
\end{exercise}
\begin{solution}
\begin{compactenum}
    \item \textbf{寻求解题思路并设定变换参数}：
    由于边界中出现两条平行直线 $x+y=4,\ x+y=2$，自然取
    \(\displaystyle u=x+y,\ v=y\)，于是 \(\displaystyle x=u-v,\ y=v\)。
    \item \textbf{计算雅可比变形系数}：
    \begin{align*}
    J
    &= \frac{\partial(x, y)}{\partial(u, v)}
    = \det \begin{pmatrix} \frac{\partial x}{\partial u} & \frac{\partial x}{\partial v} \\ \frac{\partial y}{\partial u} & \frac{\partial y}{\partial v} \end{pmatrix} \\
    &= \det \begin{pmatrix} 1 & -1 \\ 0 & 1 \end{pmatrix}
    = 1\times 1 - 0\times(-1) = 1.
    \end{align*}
    由于 $|J| = 1$，说明这个非正交变换并未发生面积放缩变形，$\mathrm{d}x\,\mathrm{d}y = \mathrm{d}u\,\mathrm{d}v$。
    \item \textbf{化边界曲线至 $u-v$ 平面以确定新区域 $D'$}：
    两条直线直接化为 $u=2,\ u=4$；而抛物线 $y^2=2x$ 变为
    \(\displaystyle v^2=2(u-v)\iff v^2+2v-2u=0\iff (v+1)^2=2u+1\)，因而
    \(\displaystyle v=-1\pm \sqrt{2u+1}\)。
    \item \textbf{建立 $u-v$ 平面上的累次积分}：
    被积函数 $x+y$ 变成 $u$，故
    \(\displaystyle I = \int_{2}^{4} \mathrm{d}u \int_{-1-\sqrt{2u+1}}^{-1+\sqrt{2u+1}} u \,\mathrm{d}v\)。
    \item \textbf{计算内层与外层积分}：
    内层积分只需取上下限之差：
    \(\displaystyle I=\int_{2}^{4}u\cdot 2\sqrt{2u+1}\,\mathrm{d}u\)。
    再令 $t=\sqrt{2u+1}$，则 $u=\dfrac{t^2-1}{2}$，$\mathrm{d}u=t\,\mathrm{d}t$，且 $u=2,4$ 分别对应 $t=\sqrt{5},3$，于是
    \(\displaystyle I = \int_{\sqrt{5}}^{3} (t^2-1)t^2 \,\mathrm{d}t = \int_{\sqrt{5}}^{3} (t^4 - t^2) \,\mathrm{d}t\)。
    \item \textbf{代入上下限得出确切数值}：
    \(\displaystyle
    I = \left[ \frac{1}{5}t^5 - \frac{1}{3}t^3 \right]_{\sqrt{5}}^3 = \left( \frac{243}{5} - 9 \right) - \left( \frac{25\sqrt{5}}{5} - \frac{5\sqrt{5}}{3} \right)
    = \frac{198}{5} - \frac{10\sqrt{5}}{3}\)。
\end{compactenum}
\end{solution}

\begin{exercise}[利用对称性与奇偶性计算积分]
计算以下二重积分：
\(\displaystyle I = \iint\limits_{D} x \left[ 1 + y f(x^{2}+y^{2}) \right] \mathrm{d}x\,\mathrm{d}y\)
其中积分区域 $D$ 由立方抛物线 $y=x^{3}$、水平直线 $y=1$ 以及铅直直线 $x=-1$ 所围成（假定函数 $f$ 是连续函数）。
\end{exercise}
\begin{solution}
\begin{compactenum}
    \item \textbf{将原积分进行代数拆分}：
    由于被积函数存在两个加项，我们利用积分线性性质将其一分为二：
    \(\displaystyle I = \iint\limits_{D} x \,\mathrm{d}x\,\mathrm{d}y + \iint\limits_{D} xy f(x^{2}+y^{2}) \,\mathrm{d}x\,\mathrm{d}y = I_1 + I_2\)。
    \item \textbf{计算常规多项式积分 $I_1$}：
    把区域视为 $x$-型区域，则 $x\in[-1,1]$，且固定 $x$ 时 $y$ 从抛物线 $y=x^3$ 到直线 $y=1$，所以
    \(\displaystyle I_1 = \int_{-1}^{1} \mathrm{d}x \int_{x^3}^{1} x \,\mathrm{d}y
    = \int_{-1}^{1} (x - x^4) \,\mathrm{d}x\)。
    其中奇项 $x$ 在对称区间上积分为 $0$，故
    \(\displaystyle I_1 = -2 \int_{0}^{1} x^4 \,\mathrm{d}x = -2 \left[ \frac{1}{5}x^5 \right]_0^1 = -\frac{2}{5}\)。
    \item \textbf{运用微积分基本定理分析 $I_2$}：
    设 $F'(t)=f(t)$，则对内层变量 $y$ 有
    \[
    \int_{x^3}^{1} y f(x^2+y^2) \,\mathrm{d}y
    = \frac{1}{2} \int_{y=x^3}^{y=1} f(x^2+y^2) \,\mathrm{d}(x^2+y^2)
    = \frac{1}{2} \left( F(x^2+1) - F(x^2+x^6) \right).
    \]
    因而
    \[
    I_2 = \frac{1}{2} \int_{-1}^{1}
    \underbrace{x \left[ F(x^2+1) - F(x^2+x^6) \right]}_{g(x)} \mathrm{d}x.
    \]
    又
    \[
    g(-x)=(-x)\left[ F((-x)^2+1)-F((-x)^2+(-x)^6) \right]
    =-x\left[ F(x^2+1)-F(x^2+x^6) \right]=-g(x),
    \]
    所以 $g(x)$ 为奇函数，在对称区间 $[-1,1]$ 上积分为零，即 $I_2=0$。
    \item \textbf{得出最终结论}：\(\displaystyle I=I_1+I_2=-\frac{2}{5}+0=-\frac{2}{5}\)。
\end{compactenum}
\end{solution}

\begin{exercise}[极坐标下圆外心脏线内区域的求积]
计算二重积分 \(\displaystyle I = \iint\limits_{D} \sqrt{x^{2}+y^{2}} \,\mathrm{d}\sigma\)，
其中区域 $D$ 是由心脏线 $r=1+\cos\theta$ 和圆 $r=1$ 所围成的位于圆外的闭区域。
\end{exercise}
\begin{solution}
\begin{compactenum}
    \item \textbf{寻找交点与极角范围}：
    联立 $1+\cos\theta=1$ 得 $\cos\theta=0$，故交点对应的极角为 $\theta=\pm \dfrac{\pi}{2}$；并且在区间 $\left[-\dfrac{\pi}{2},\dfrac{\pi}{2}\right]$ 上有 $1+\cos\theta\geqslant 1$，正好对应“位于圆外而在心脏线内”的那部分区域。
    \item \textbf{确定区域的极坐标不等式描述}：
    因此区域可写为 \(\displaystyle -\frac{\pi}{2}\leqslant \theta \leqslant \frac{\pi}{2},\ 1 \leqslant r \leqslant 1+\cos\theta\)。
    \item \textbf{建立极坐标积分并计算}：
    被积函数 $\sqrt{x^2+y^2}=r$，乘上极面元后总被积项为 $r^2$，于是展开并利用偶性可得
    \begin{align*}
    I
    &= \int_{-\frac{\pi}{2}}^{\frac{\pi}{2}} \mathrm{d}\theta \int_{1}^{1+\cos\theta} r^2 \,\mathrm{d}r
    = \frac{1}{3} \int_{-\frac{\pi}{2}}^{\frac{\pi}{2}} \left( (1+\cos\theta)^3 - 1 \right) \mathrm{d}\theta\\
    &= \frac{2}{3} \int_{0}^{\frac{\pi}{2}} \left( 3\cos\theta + 3\cos^2\theta + \cos^3\theta \right) \mathrm{d}\theta.
    \end{align*}
    其中 \(\displaystyle 3 \int_0^{\pi/2} \cos\theta \,\mathrm{d}\theta = 3,\quad
    3 \int_0^{\pi/2} \cos^2\theta \,\mathrm{d}\theta = \frac{3\pi}{4},\quad
    \int_0^{\pi/2} \cos^3\theta \,\mathrm{d}\theta = \frac{2}{3}\)。
    \item \textbf{合并并化简结果}：\(\displaystyle I = \frac{2}{3} \left( 3 + \frac{3\pi}{4} + \frac{2}{3} \right) = \frac{2}{3} \left( \frac{11}{3} + \frac{3\pi}{4} \right) = \frac{22}{9} + \frac{\pi}{2}\)。
\end{compactenum}
\end{solution}

\begin{exercise}[利用柱坐标计算带旋转面的体积]
设 $\Omega$ 由抛物柱面绕 $z$ 轴旋转而成的曲面（即 $x^{2}+y^{2}=2z$）与平面 $z=4$ 围成，求：
\(\displaystyle \iiint\limits_{\Omega} (x^{2}+y^{2}+z) \,\mathrm{d}V\)。
\end{exercise}
\begin{solution}
\begin{compactenum}
    \item \textbf{确定区域的边界与坐标系选择}：
    曲面 $x^2+y^2=2z$ 是开口向上的旋转抛物面，顶部由平面 $z=4$ 封顶。联立得交线 $x^2+y^2=8$，故立体在 $Oxy$ 平面上的投影是半径 $2\sqrt{2}$ 的圆域。由于区域旋转对称且被积函数含有 $x^2+y^2$，取柱坐标最方便。
    \item \textbf{在柱坐标系下界定积分区域}：
    因而 \(\displaystyle 0\leqslant \theta\leqslant 2\pi,\ 0\leqslant r\leqslant 2\sqrt{2},\ \frac{r^2}{2}\leqslant z\leqslant 4\)。
    \item \textbf{建立并计算三重积分}：
    被积函数化为 $r^2+z$，体积微元为 $r\,\mathrm{d}r\,\mathrm{d}\theta\,\mathrm{d}z$，故
    \begin{align*}
    I
    &= \int_{0}^{2\pi}\mathrm{d}\theta
       \int_{0}^{2\sqrt{2}}\mathrm{d}r
       \int_{\frac{r^2}{2}}^{4} (r^2+z)r \,\mathrm{d}z \\
    &= 2\pi \int_{0}^{2\sqrt{2}} \left( 8r + 4r^3 - \frac{5}{8}r^5 \right)\mathrm{d}r \\
    &= 2\pi \left[ 4r^2 + r^4 - \frac{5}{48}r^6 \right]_0^{2\sqrt{2}}
     = 2\pi \left( 4(8) + 64 - \frac{5}{48}(512) \right)
     = \frac{256\pi}{3}.
    \end{align*}
\end{compactenum}
\end{solution}

\begin{exercise}[含参抽象函数的极限与球坐标积分]
设 $f(x)$ 在 $x=0$ 处可导，且 $f(0)=0$。求极限
\[
\lim_{t \to 0} \frac{1}{t^{4}} \iiint\limits_{\Omega_t}
f\left( \sqrt{x^{2}+y^{2}+z^{2}} \right) \mathrm{d}x\,\mathrm{d}y\,\mathrm{d}z ,
\]
其中积分区域为以原点为球心、$t$ 为半径的球体 $\Omega_t: x^{2}+y^{2}+z^{2} \leqslant t^{2}$。
\end{exercise}
\begin{solution}
\begin{compactenum}
    \item \textbf{利用球坐标系化简三重积分}：
    由于 $\Omega_t$ 是半径为 $t$ 的原心球，取球坐标，有
    \[
      0\leqslant \theta\leqslant 2\pi,\qquad
      0\leqslant \varphi\leqslant \pi,\qquad
      0\leqslant r\leqslant t.
    \]
    又 \(f\!\left(\sqrt{x^{2}+y^{2}+z^{2}}\right)=f(r)\)，故
    \[
      I(t)
      =\int_{0}^{2\pi} \mathrm{d}\theta
       \int_{0}^{\pi} \sin\varphi\,\mathrm{d}\varphi
       \int_{0}^{t} f(r)r^2 \,\mathrm{d}r
      =4\pi \int_{0}^{t} r^2 f(r) \,\mathrm{d}r.
    \]
    \item \textbf{构造极限并应用洛必达法则}：
    因而原极限化为
    \(\displaystyle L = \lim_{t \to 0} \frac{4\pi \int_{0}^{t} r^2 f(r) \,\mathrm{d}r}{t^4}\)。
    这是 $\frac{0}{0}$ 型，利用变上限积分求导法则与洛必达法则，得
    \(\displaystyle L = \lim_{t \to 0} \frac{4\pi \cdot t^2 f(t)}{4t^3} = \pi \lim_{t \to 0} \frac{f(t)}{t}\)。
    \item \textbf{利用导数定义求出最终结果}：
    由 $f(0)=0$，可把上式写成
    \(\displaystyle L = \pi \lim_{t \to 0} \frac{f(t) - f(0)}{t - 0}\)。
    因为 $f$ 在 $0$ 处可导，所以 $L=\pi f'(0)$。
\end{compactenum}
\end{solution}

\begin{exercise}[利用切片法与柱坐标系求积分]
计算三重积分 \(\displaystyle I=\iiint\limits_{\Omega}\frac{\mathrm{e}^{z}}{\sqrt{x^{2}+y^{2}}}\,\mathrm{d}x\,\mathrm{d}y\,\mathrm{d}z\)，其中区域 $\Omega$ 由旋转锥面 $z=\sqrt{x^{2}+y^{2}}$ 与两平行平面 $z=1$ 和 $z=2$ 所围成。
\end{exercise}
\begin{solution}
\begin{compactenum}
    \item \textbf{解析积分区域结构并定限}：
    区域位于平面 $z=1$ 与 $z=2$ 之间，侧面由锥面 $z=\sqrt{x^2+y^2}$ 围成，因此最自然的是按 $z$ 切片：$1\leqslant z\leqslant 2$。固定 $z$ 时，截面
    \(D_z=\{(x,y):x^2+y^2\leqslant z^2\}\) 是半径为 $z$ 的圆域。
    \item \textbf{建立圆盘截面上的极坐标二重积分}：
    在截面 $D_z$ 上再用极坐标，取 $0\leqslant \theta\leqslant 2\pi,\ 0\leqslant r\leqslant z$。此时
    \(\displaystyle \frac{\mathrm{e}^z}{\sqrt{x^2+y^2}}=\frac{\mathrm{e}^z}{r}\)，
    乘上面积微元 $r\,\mathrm{d}r\,\mathrm{d}\theta$ 后，得到
    \(\displaystyle J(z) = \iint\limits_{D_z} \frac{\mathrm{e}^z}{r} \,\mathrm{d}x\,\mathrm{d}y
    = \int_{0}^{2\pi} \mathrm{d}\theta \int_{0}^{z} \left(\frac{\mathrm{e}^z}{r}\right) r \,\mathrm{d}r
    = 2\pi z \mathrm{e}^z\)。
    \item \textbf{计算最外层的一元定积分}：
    最外层积分为 \(\displaystyle I = \int_{1}^{2} J(z) \,\mathrm{d}z = 2\pi \int_{1}^{2} z \mathrm{e}^z \,\mathrm{d}z\)。
    又 \(\displaystyle \int z \mathrm{e}^z \,\mathrm{d}z = \int z \,\mathrm{d}(\mathrm{e}^z) = z\mathrm{e}^z - \int \mathrm{e}^z \,\mathrm{d}z = \mathrm{e}^z(z-1)\)，
    所以 \(\displaystyle I = 2\pi \Big[ \mathrm{e}^z(z-1) \Big]_1^2 = 2\pi \mathrm{e}^2\)。
\end{compactenum}
\end{solution}


\section{习题}

\subsection{第一节习题：二重积分的概念与性质}

\subsubsection{A组}

\begin{exercise}[基本概念简答]
回答下列问题：
\begin{shorttrienum}
    \shortitem{1}{什么叫曲顶柱体？它的体积怎样求？} &
    \shortitem{2}{二重积分的概念是怎样的？} &
    \shortitem{3}{二重积分有哪些重要性质？}
\end{shorttrienum}
\end{exercise}
\begin{solution}
\begin{compactenum}
    \item \textbf{曲顶柱体及其体积求法}：以 $Oxy$ 平面上的有界闭区域 $D$ 为底面，侧面由沿 $D$ 边界且平行于 $z$ 轴的母线构成，顶部是定义在 $D$ 上的非负连续曲面 $z=f(x,y)$，三者围成的立体称为曲顶柱体。按微元法分割底面并累加小柱体体积，取极限后得 \(\displaystyle V = \iint_D f(x, y) \, \mathrm{d}\sigma\)。
    \item \textbf{二重积分的概念}：二重积分本质上是黎曼和的极限，即把定义在区域 $D$ 上的函数值与小面积元 $\Delta\sigma_i$ 相乘后求和，当分割无限细化且极限与分割和取点方式无关时，这个极限就定义为函数在 $D$ 上的二重积分。它在物理上可表示质量，在几何上可表示体积。
    \item \textbf{二重积分的重要性质}：包括线性性质、区域可加性、比较性质（保号性）、估值定理与中值定理。
\end{compactenum}
\end{solution}

\begin{exercise}[物理应用：薄片质量的积分建模]
设有一平面薄片，它占有 $Oxy$ 平面上的闭区域 $D$，在 $D$ 上的每点 $(x, y)$ 处的面密度为 $\rho(x, y)$，其中 $\rho(x, y)$ 是 $D$ 上的连续正值函数。试利用本节所介绍的求曲顶柱体体积的思想及方法，详细推导并求出这块平面薄片的质量。
\end{exercise}
\begin{solution}
\begin{compactenum}
    \item \textbf{进行网格化分割}：由于薄片非均匀，不能直接用密度乘面积，故先把区域 $D$ 任意分割为 $n$ 个互不重叠的小区域块 $\Delta\sigma_1,\dots,\Delta\sigma_n$。
    \item \textbf{执行局部均匀化近似}：在每个小块 $\Delta\sigma_i$ 内任选一点 $(\xi_i,\eta_i)$。当小块足够小时，可近似认为该块密度恒为 $\rho(\xi_i,\eta_i)$，于是 \(\displaystyle \Delta M_i \approx \rho(\xi_i, \eta_i) \cdot \Delta\sigma_i,\ M \approx \sum_{i=1}^n \rho(\xi_i, \eta_i) \Delta\sigma_i\)。
    \item \textbf{取极限获取精确值}：令最大分割直径 $\lambda \to 0$，便得 \(\displaystyle M = \lim_{\lambda \to 0} \sum_{i=1}^n \rho(\xi_i, \eta_i) \Delta\sigma_i = \iint\limits_D \rho(x, y) \, \mathrm{d}\sigma\)。
\end{compactenum}
\end{solution}

\begin{exercise}[积分表达式的实际意义识别]
试对下列二重积分给出几何上或物理上的一种解释：设 $D$ 是平面闭区域，
\begin{shortexenum}
    \shortitem{1}{$\displaystyle \iint\limits_D 1 \, \mathrm{d}\sigma$} &
    \shortitem{2}{$\displaystyle \iint\limits_D \rho(x, y) \, \mathrm{d}\sigma$，其中 $\rho(x, y)$ 为密度} \\
    \shortitem{3}{$\displaystyle \iint\limits_D f(x, y) \, \mathrm{d}\sigma$，其中 $f(x, y)$ 为曲面上点的竖坐标且 $f \geqslant 0$} &
\end{shortexenum}
\end{exercise}
\begin{solution}
\begin{compactenum}
    \item \textbf{对于表达式 (1)}：当被积函数恒为 $1$ 时，微元 $1\cdot \mathrm{d}\sigma=\mathrm{d}\sigma$ 就是一个小面积元，因此整个积分表示区域 $D$ 的总面积。
    \item \textbf{对于表达式 (2)}：若 $\rho(x,y)$ 是面密度，则微元 $\rho(x,y)\,\mathrm{d}\sigma$ 表示该小面元上的微元质量，积分即薄片总质量。
    \item \textbf{对于表达式 (3)}：当 $f(x,y)\ge0$ 表示曲面高度时，微元 $f(x,y)\,\mathrm{d}\sigma$ 可看成底面积为 $\mathrm{d}\sigma$、高为 $f(x,y)$ 的小柱体体积，因此积分给出曲顶柱体体积。
\end{compactenum}
\end{solution}

\begin{exercise}[利用对称性判定积分符号]
根据二重积分的几何意义与对称性质，说明下列积分值是大于零、小于零还是等于零：
\par\noindent
\begin{tabular}{@{}p{0.48\linewidth}p{0.48\linewidth}@{}}
(1) $\displaystyle \iint\limits_{\substack{|x|\leqslant 1 \\ |y|\leqslant 1}} (1+x)\mathrm{d}\sigma$ &
(2) $\displaystyle \iint\limits_{\substack{|x|\leqslant 1 \\ |y|\leqslant 1}} (x-1)\mathrm{d}\sigma$ \\
(3) $\displaystyle \iint\limits_{x^2+y^2\leqslant 1} x\mathrm{d}\sigma$ &
(4) $\displaystyle \iint\limits_{x^2+y^2\leqslant 1} x^2\mathrm{d}\sigma$ \\
(5) $\displaystyle \iint\limits_{x^2+y^2\leqslant 1} xy\mathrm{d}\sigma$ &
(6) $\displaystyle \iint\limits_{\substack{0\leqslant x\leqslant 1 \\ 0\leqslant y\leqslant 1}} y\mathrm{d}\sigma$ \\
\end{tabular}
\par
\end{exercise}
\begin{solution}
\begin{compactenum}
    \item \textbf{解第(1)题}：积分区域 $D$ 是边长为 $2$ 的正方形 $[-1, 1] \times [-1, 1]$。在这个区域内，由于 $x \geqslant -1$，故被积函数 $(1+x) \geqslant 0$。且除了边界线 $x = -1$ 外，函数值均严格大于 $0$。根据二重积分的保号性质，积分值\textbf{大于零}（正体积）。
    \item \textbf{解第(2)题}：积分区域同上。在该区域内，由于 $x \leqslant 1$，故被积函数 $(x-1) \leqslant 0$。除了边界线 $x = 1$ 外，函数值严格小于 $0$。根据保号性质，积分累加的都是负值，故积分值\textbf{小于零}。
    \item \textbf{解第(3)题}：积分区域是单位圆 $x^2 + y^2 \leqslant 1$。该区域几何形状关于 $y$ 轴左右完全对称。而被积函数 $f(x, y) = x$ 满足 $f(-x, y) = -x = -f(x, y)$，即关于变量 $x$ 是\textbf{奇函数}。根据奇函数在对称区间上的积分定理（左半圆的负体积恰好抵消右半圆的正体积），积分值\textbf{等于零}。
    \item \textbf{解第(4)题}：积分区域是单位圆。被积函数 $f(x, y) = x^2$。在整个闭区域内恒有 $x^2 \geqslant 0$，且除了 $y$ 轴上的点（$x=0$）之外，均严格大于 $0$。恒正函数的体积不为零，故积分值\textbf{大于零}。
    \item \textbf{解第(5)题}：积分区域是单位圆，关于 $x$ 轴和 $y$ 轴均对称。被积函数 $f(x, y) = xy$。当把点 $(x,y)$ 关于 $y$ 轴对称到 $(-x,y)$ 时，函数值变为 $f(-x,y)=-xy=-f(x,y)$，因此右半圆与左半圆的贡献互相抵消；等价地，也可以说它关于 $x$ 或 $y$ 都是奇的。一三象限的正值与二四象限的负值完全抵消，故积分值\textbf{等于零}。
    \item \textbf{解第(6)题}：积分区域是第一象限内的正方形 $[0, 1] \times [0, 1]$。在该区域内，坐标 $y$ 满足 $0 \leqslant y \leqslant 1$，因此被积函数 $f(x, y) = y$ 恒非负，且非恒为零。累加非负体积元，故积分值\textbf{大于零}。
\end{compactenum}
\end{solution}

\subsubsection{B组}

\begin{exercise}[积分绝对值不等式的严谨证明]
设函数 $f(x,y)$ 在闭区域 $D$ 上连续，请严格证明绝对值不等式性质 \(\displaystyle \left|\iint\limits_D f(x,y)\mathrm{d}\sigma\right| \leqslant \iint\limits_D |f(x,y)|\mathrm{d}\sigma\)。
\end{exercise}
\begin{solution}
\begin{compactenum}
    \item \textbf{构造基础函数不等式}：根据实数绝对值的基本性质，对于区域 $D$ 内的任意一点 $(x, y)$，其函数值 $f(x, y)$ 必然介于其绝对值的相反数与其绝对值之间。即：
    \(\displaystyle -|f(x, y)| \leqslant f(x, y) \leqslant |f(x, y)|\)
    \item \textbf{应用积分保号性}：由于 $f(x, y)$ 在 $D$ 上连续，那么 $|f(x, y)|$ 和 $-|f(x, y)|$ 在 $D$ 上也必然连续，因而在 $D$ 上都是可积的。根据二重积分的保号性质（即不等关系在积分后不变），对上述不等式的三边同时在区域 $D$ 上取二重积分，方向不变：\(\displaystyle \iint\limits_D \left(-|f(x, y)|\right) \, \mathrm{d}\sigma \leqslant \iint\limits_D f(x, y) \, \mathrm{d}\sigma \leqslant \iint\limits_D |f(x, y)| \, \mathrm{d}\sigma\)。
    \item \textbf{提取常数进行化简}：根据二重积分的线性性质，可以把最左边积分式内部的负号（即常数 $-1$）提取到积分号外面：\(\displaystyle -\iint\limits_D |f(x, y)| \, \mathrm{d}\sigma \leqslant \iint\limits_D f(x, y) \, \mathrm{d}\sigma \leqslant \iint\limits_D |f(x, y)| \, \mathrm{d}\sigma\)。
    \item \textbf{转化为绝对值形式}：观察上式，中间的实数值 $\iint_D f(x, y) \, \mathrm{d}\sigma$ 被夹在正数 $\iint_D |f(x, y)| \, \mathrm{d}\sigma$ 与其相反数之间。根据实数不等式的等价转化法则（若 $-A \leqslant B \leqslant A$，则 $|B| \leqslant A$），即可直接得出结论：
    \(\displaystyle \left| \iint\limits_D f(x, y) \, \mathrm{d}\sigma \right| \leqslant \iint\limits_D |f(x, y)| \, \mathrm{d}\sigma\)
    证明完毕。
\end{compactenum}
\end{solution}

\begin{exercise}[利用最值估计比较积分大小]
根据二重积分的性质，通过严谨的逻辑推理比较下列各组积分的大小：
\begin{shortsingleenum}
    \shortsingleitem{1}{比较 \(\displaystyle I_1=\iint\limits_D (x+y)^2\mathrm{d}\sigma,\qquad I_2=\iint\limits_D (x+y)^3\mathrm{d}\sigma\)，其中 $D$ 是由圆周 $(x-2)^2+(y-1)^2=2$ 所围成的闭区域；}
    \shortsingleitem{2}{比较 \(\displaystyle J_1=\iint\limits_D \ln(x+y)\mathrm{d}\sigma,\qquad J_2=\iint\limits_D [\ln(x+y)]^2\mathrm{d}\sigma\)，其中 $D$ 是矩形闭区域：$3\leqslant x\leqslant 5,\ 0\leqslant y\leqslant 1$。}
\end{shortsingleenum}
\end{exercise}
\begin{solution}
\begin{compactenum}
    \item \textbf{求解第(1)题，比较圆域上的积分大小}：
    \begin{itemize}
        \item \textbf{几何分析区域}：区域 $D$ 是一个圆，圆心位于 $C(2, 1)$，半径为 $R = \sqrt{2}$。
        \item \textbf{确定被积函数底数的范围}：我们需要评估在区域 $D$ 内，函数 $u(x, y) = x+y$ 的取值范围。在几何上，令直线系方程为 $x+y = k$，我们求该直线与圆有交点时 $k$ 的范围。
        \item \textbf{利用点到直线距离公式}：圆心 $(2, 1)$ 到直线 $x+y-k=0$ 的距离必须小于等于半径 $\sqrt{2}$，即 \(\displaystyle d=\frac{|2+1-k|}{\sqrt{1^2+1^2}}=\frac{|3-k|}{\sqrt{2}}\leqslant\sqrt{2}\)。
        解这个绝对值不等式：$|3 - k| \leqslant 2$，推得 $-2 \leqslant 3 - k \leqslant 2$，进一步得到 $1 \leqslant k \leqslant 5$。
        \item \textbf{得出函数范围并比较积分}：这就意味着，在整个闭区域 $D$ 上，恒有 $x+y \geqslant 1$。对于任意实数 $u \geqslant 1$，必有 $u^3 \geqslant u^2$。因此在 $D$ 上恒成立 $(x+y)^3 \geqslant (x+y)^2$。根据积分的保号性，得出结论 \(\displaystyle \iint\limits_D (x+y)^3 \, \mathrm{d}\sigma > \iint\limits_D (x+y)^2 \, \mathrm{d}\sigma\)，即 $I_2>I_1$。
    \end{itemize}
    \item \textbf{求解第(2)题，比较矩形区域上的积分大小}：
    \begin{itemize}
        \item \textbf{确定独立变量范围}：已知区域 $D$ 定义为 $x \in [3, 5]$ 且 $y \in [0, 1]$。
        \item \textbf{计算和式变量的范围}：根据不等式加法规则，两个变量相加后的范围是 $(x+y) \in [3+0, 5+1]$，即 $x+y \in [3, 6]$。
        \item \textbf{代入对数函数进行评估}：我们需要比较的是以 $\ln(x+y)$ 为底数的函数大小。因为 $x+y \geqslant 3$，而自然对数的底数 $\mathrm{e} \approx 2.718 < 3$。对数函数是单调递增的，所以有 \(\displaystyle \ln(x+y) \geqslant \ln(3) > \ln(\mathrm{e}) = 1\)。
        \item \textbf{得出结论}：我们得到在区域 $D$ 上，被积函数的底数 $\ln(x+y)$ 严格大于 $1$。对于任何大于 $1$ 的实数 $u$，必然有 $u^2 > u$。因此在 $D$ 上恒有 $[\ln(x+y)]^2 > \ln(x+y)$。根据保号性，得出结论 \(\displaystyle \iint\limits_D [\ln(x+y)]^2 \, \mathrm{d}\sigma > \iint\limits_D \ln(x+y) \, \mathrm{d}\sigma\)，即 $J_2>J_1$。
    \end{itemize}
\end{compactenum}
\end{solution}


\subsection{第二节习题：直角坐标系下二重积分的计算}

\subsubsection{A组}

\begin{exercise}[累次积分的转换与表达]
将二重积分 $\iint\limits_{D} f(x,y) \,\mathrm{d}\sigma$ 转化为累次积分的形式，其中积分区域 $D$ 分别为：
\begin{shortexenum}
    \shortitem{1}{顶点位于 $(1,1), (1,2), (4,1), (4,2)$ 的矩形区域；} &
    \shortitem{2}{顶点位于 $(1,0), (1,2), (4,1), (4,2)$ 的梯形区域。}
\end{shortexenum}
\end{exercise}
\begin{solution}
\begin{compactenum}
    \item \textbf{解答小题 (1)：标准矩形区域}
    \begin{itemize}
        \item 观察顶点坐标，区域在 $x$ 方向的跨度是从 $x=1$ 到 $x=4$；在 $y$ 方向的跨度是从 $y=1$ 到 $y=2$。并且边界互相垂直平行，恰好构成闭矩形区域 $D = [1,4] \times [1,2]$。
    \item 按照 Fubini 定理，直接写出两种等价的累次积分形式：
    \begin{align*}
    &\int_{1}^{4} \mathrm{d}x \int_{1}^{2} f(x,y) \,\mathrm{d}y,\\
    &\int_{1}^{2} \mathrm{d}y \int_{1}^{4} f(x,y) \,\mathrm{d}x .
    \end{align*}
    \end{itemize}
    \item \textbf{解答小题 (2)：梯形区域的划分}
    \begin{itemize}
        \item 将四个顶点描点绘图，发现梯形的左右边界是铅直的（左边直线 $x=1$，右边直线 $x=4$），上边界是水平的（从 $(1,2)$ 连向 $(4,2)$ 的直线 $y=2$）。而下边界是斜线，连接点 $(1,0)$ 到 $(4,1)$。
        \item 我们需要求出这根倾斜的下边界的解析式。由两点式直线方程可得：$\frac{y - 0}{1 - 0} = \frac{x - 1}{4 - 1} \implies y = \frac{1}{3}(x-1)$。
        \item \textbf{情况一：作 $x$-型区域视角}。投影到 $x$ 轴范围为 $[1, 4]$。对于给定横坐标 $x$，对应的垂线段从下方的倾斜直线出发，一直延伸到上方的水平横线终止。故 $y$ 的范围是 $[\frac{1}{3}(x-1), 2]$。化为累次积分为 \(\displaystyle \int_{1}^{4} \mathrm{d}x \int_{\frac{1}{3}(x-1)}^{2} f(x,y) \,\mathrm{d}y\)。
        \item \textbf{情况二：作 $y$-型区域视角}。该梯形需要沿水平方向切割，这会产生分段的情况（左侧边界是一段竖直和一段倾斜的拼接）。由于主要目的是了解如何转换，我们这里仅提供最直观、无需分割的 $x$-型区域结果即可。
    \end{itemize}
\end{compactenum}
\end{solution}

\begin{exercise}[多重区域的积分限安置]
在下列各题中，对二重积分 $I = \iint\limits_{D} f(x,y)\,\mathrm{d}x\,\mathrm{d}y$ 按所给的区域 $D$，请依两个不同的顺序（先对 $x$ 再对 $y$，和先对 $y$ 再对 $x$）分别写出积分的上下限表达：
\begin{shortexenum}
  \shortitem{1}{$D$：以 $(0,0),(1,0),(1,1)$ 为顶点的三角形；} &
  \shortitem{2}{$D$：圆 $x^2+y^2\leqslant a^2$；}\\
  \shortitem{3}{$D$：以 $(0,0),(2,1),(-2,1)$ 为顶点的三角形；} &
  \shortitem{4}{$D$：圆 $x^2+y^2\leqslant y$；}\\
  \shortitem{5}{$D$：由曲线 $y=x^2$ 与直线 $y=1$ 所围成。} &
\end{shortexenum}
\end{exercise}
\begin{solution}
\begin{compactenum}
    \item \textbf{小题 (1) 的解析}：
    \begin{itemize}
        \item 三角形三条边界边分别为：底边 $y=0$，右直角边 $x=1$，以及连接原点和右上顶点的斜边 $y=x$。
        \item \textbf{先对 $y$ 后对 $x$ ($x$-型)}：投影区间为 $x \in [0,1]$。在固定 $x$ 下，$y$ 沿着竖线从下界 $y=0$ 走到上边界 $y=x$。表达式为：$\int_{0}^{1} \mathrm{d}x \int_{0}^{x} f(x,y) \,\mathrm{d}y$。
        \item \textbf{先对 $x$ 后对 $y$ ($y$-型)}：投影区间为 $y \in [0,1]$。在固定 $y$ 下，$x$ 沿着横线从左边界 $x=y$ 走到右直角边 $x=1$。表达式为：$\int_{0}^{1} \mathrm{d}y \int_{y}^{1} f(x,y) \,\mathrm{d}x$。
    \end{itemize}
    \item \textbf{小题 (2) 的解析}：
    \begin{itemize}
        \item 半径为 $a$ 的中心圆（假设 $a>0$），边界方程 $y = \pm\sqrt{a^2-x^2}$，反解为 $x = \pm\sqrt{a^2-y^2}$。
        \item \textbf{先对 $y$ 后对 $x$}：$\int_{-a}^{a} \mathrm{d}x \int_{-\sqrt{a^2-x^2}}^{\sqrt{a^2-x^2}} f(x,y) \,\mathrm{d}y$。
        \item \textbf{先对 $x$ 后对 $y$}：$\int_{-a}^{a} \mathrm{d}y \int_{-\sqrt{a^2-y^2}}^{\sqrt{a^2-y^2}} f(x,y) \,\mathrm{d}x$。
    \end{itemize}
    \item \textbf{小题 (3) 的解析}：
    \begin{itemize}
        \item 顶点为 $(0,0), (-2,1), (2,1)$，这是底边水平（$y=1$），尖点朝下顶着原点的等腰三角形。右斜边方程为 $x=2y$；左斜边为 $x=-2y$。
        \item \textbf{先对 $x$ 后对 $y$ ($y$-型)}：直接整体投影到 $y$ 轴，范围 $y \in [0,1]$。横向扫描从左斜边进入，右斜边出：$\int_{0}^{1} \mathrm{d}y \int_{-2y}^{2y} f(x,y) \,\mathrm{d}x$。
        \item \textbf{先对 $y$ 后对 $x$ ($x$-型)}：
        投影到 $x$ 轴的范围为 $x \in [-2, 2]$。下边界在原点处发生转折，
        因此需以 $y$ 轴为界把区域拆成左右两个三角形：\(\displaystyle
        \int_{-2}^{0} \mathrm{d}x \int_{-x/2}^{1} f(x,y) \,\mathrm{d}y
        + \int_{0}^{2} \mathrm{d}x \int_{x/2}^{1} f(x,y) \,\mathrm{d}y\)。
    \end{itemize}
    \item \textbf{小题 (4) 的解析}：
    \begin{itemize}
        \item 边界圆方程配方：$x^2+y^2-y=0 \implies x^2+(y-\frac{1}{2})^2 = \frac{1}{4}$。这是一个圆心在 $(0, 0.5)$、半径为 $0.5$、位于第一二象限与原点相切的偏心圆。
        \item 解出显式函数表达：上下半支为 $y = \frac{1}{2} \pm \sqrt{\frac{1}{4}-x^2}$；左右半支为 $x = \pm\sqrt{y-y^2}$。
        \item \textbf{先对 $y$ 后对 $x$}：圆的横向最宽处于极值点，范围是 \(x \in [-\frac12, \frac12]\)。因此
        \[
        \int_{-\frac12}^{\frac12}\mathrm{d}x
        \int_{\frac12-\sqrt{\frac14-x^2}}^{\frac12+\sqrt{\frac14-x^2}} f(x,y)\,\mathrm{d}y .
        \]
        \item \textbf{先对 $x$ 后对 $y$}：纵向极端范围为原点到最高点，即 \(y\in[0,1]\)。表达式为 \(\displaystyle \int_0^1\mathrm{d}y\int_{-\sqrt{y-y^2}}^{\sqrt{y-y^2}} f\,\mathrm{d}x\)。
    \end{itemize}
    \item \textbf{小题 (5) 的解析}：
    \begin{itemize}
        \item 被一条开口朝上的抛物线 $y=x^2$ 和一根天花板横线 $y=1$ 在第一、二象限封出了一个“碗形”区域。交点为 $(-1,1)$ 和 $(1,1)$。
        \item \textbf{先对 $y$ 后对 $x$}：区域横跨 \(x \in [-1, 1]\)，从抛物线底向上走到天花板：\(\displaystyle \int_{-1}^{1} \mathrm{d}x \int_{x^2}^{1} f(x,y) \,\mathrm{d}y\)。
        \item \textbf{先对 $x$ 后对 $y$}：整体被压倒在 $y$ 轴上的投影区间是 \(y \in [0, 1]\)。对任意高度 $y$，抛物线从左支 \(x=-\sqrt{y}\) 走到右支 \(x=\sqrt{y}\)：\(\displaystyle \int_{0}^{1} \mathrm{d}y \int_{-\sqrt{y}}^{\sqrt{y}} f(x,y) \,\mathrm{d}x\)。
    \end{itemize}
\end{compactenum}
\end{solution}

\begin{exercise}[直接计算二重积分]
计算下列二重积分：
\begin{shortsingleenum}
    \shortsingleitem{1}{$\displaystyle \iint\limits_D \sqrt{x+y} \,\mathrm{d}\sigma$，其中 $D$ 是矩形区域：$0 \leqslant x \leqslant 1, 0 \leqslant y \leqslant 2$；}
    \shortsingleitem{2}{$\displaystyle \iint\limits_D (5x^2+1)\sin 3y \,\mathrm{d}\sigma$，其中 $D: -1 \leqslant x \leqslant 1, 0 \leqslant y \leqslant \dfrac{\pi}{3}$；}
    \shortsingleitem{3}{$\displaystyle \iint\limits_D (2x+3y)^2 \,\mathrm{d}\sigma$，其中 $D$ 是以 $(-1,0), (0,1)$ 和 $(1,0)$ 为顶点的三角形区域。}
\end{shortsingleenum}
\end{exercise}
\begin{solution}
\begin{compactenum}
    \item \textbf{解答小题 (1)}：
    \begin{itemize}
        \item 积分区域为矩形，被积函数无奇点，可直接化为累次积分：
        \(\displaystyle I_1 = \int_{0}^{1} \mathrm{d}x \int_{0}^{2} \sqrt{x+y} \,\mathrm{d}y\)。
        \item 先对 $y$ 积分，将 $x$ 视作常数，利用幂函数积分公式 $\int u^{1/2}\mathrm{d}u = \frac{2}{3}u^{3/2}$：
        \(\displaystyle \int_{0}^{2} (x+y)^{1/2} \,\mathrm{d}y = \left[ \frac{2}{3}(x+y)^{3/2} \right]_{y=0}^{y=2} = \frac{2}{3} \left( (x+2)^{3/2} - x^{3/2} \right)\)。
        \item 再对 $x$ 积分：
        \(\displaystyle I_1 = \frac{2}{3} \int_{0}^{1} \left( (x+2)^{3/2} - x^{3/2} \right) \mathrm{d}x = \frac{2}{3} \left[ \frac{2}{5}(x+2)^{5/2} - \frac{2}{5}x^{5/2} \right]_0^1\)。
        \item 代入上下限计算最终结果：
        \(\displaystyle I_1 = \frac{4}{15} \left( (3^{5/2} - 1^{5/2}) - (2^{5/2} - 0) \right) = \frac{4}{15} \left( 9\sqrt{3} - 4\sqrt{2} - 1 \right)\)。
    \end{itemize}

    \item \textbf{解答小题 (2)}：
    \begin{itemize}
        \item 区域为矩形 $D = [-1, 1] \times [0, \frac{\pi}{3}]$，且被积函数可以分离变量 $f(x,y) = g(x)h(y)$。
        \item 利用分离变量积分法，将二重积分拆分为两个一元定积分的乘积：
        \(\displaystyle I_2 = \left( \int_{-1}^{1} (5x^2+1) \,\mathrm{d}x \right) \cdot \left( \int_{0}^{\frac{\pi}{3}} \sin 3y \,\mathrm{d}y \right)\)。
        \item 计算关于 $x$ 的积分：\(\displaystyle \int_{-1}^{1} (5x^2+1) \,\mathrm{d}x = \left[ \frac{5}{3}x^3 + x \right]_{-1}^{1} = \frac{16}{3}\)。
        \item 计算关于 $y$ 的积分：\(\displaystyle \int_{0}^{\frac{\pi}{3}} \sin 3y \,\mathrm{d}y = \left[ -\frac{1}{3}\cos 3y \right]_0^{\frac{\pi}{3}} = \frac{2}{3}\)。
        \item 两者相乘得最终结果：$I_2 = \frac{16}{3} \cdot \frac{2}{3} = \frac{32}{9}$。
    \end{itemize}

    \item \textbf{解答小题 (3)}：
    \begin{itemize}
        \item 首先确定三角形区域 $D$ 的边界方程。底边位于 $x$ 轴上区间 $[-1, 1]$，两斜边方程分别为 $y = x+1$（左斜边）和 $y = -x+1$（右斜边）。
        \item 将其看作 $y$-型区域更为简便。投影到 $y$ 轴的范围是 $y \in [0, 1]$。对于固定的 $y$，$x$ 从左边界 $x = y-1$ 变化到右边界 $x = 1-y$。
        \item 建立累次积分：\(\displaystyle I_3=\int_0^1\mathrm{d}y\int_{y-1}^{1-y}(2x+3y)^2\,\mathrm{d}x\)，展开后为 \(\displaystyle \int_0^1\mathrm{d}y\int_{y-1}^{1-y}(4x^2+12xy+9y^2)\,\mathrm{d}x\)。
        \item 分析内层关于 $x$ 的积分的对称性：积分区间 $[y-1, 1-y]$ 关于原点对称。被积函数中的奇次项 $12xy$ 在对称区间的积分为 $0$。而偶次项 $4x^2 + 9y^2$ 可化为两倍的半区间积分：\(\displaystyle \int_{y-1}^{1-y} (4x^2 + 9y^2) \,\mathrm{d}x = 2 \int_{0}^{1-y} (4x^2 + 9y^2) \,\mathrm{d}x\)。
        \item 计算简化后的内层积分：\(\displaystyle 2 \left[ \frac{4}{3}x^3 + 9y^2 x \right]_{x=0}^{x=1-y} = \frac{8}{3}(1-y)^3 + 18(y^2-y^3)\)。
        \item 计算外层关于 $y$ 的积分：
        \[
        I_3=\int_0^1\left[\frac83(1-y)^3+18y^2-18y^3\right]\mathrm{d}y
        =\left[-\frac23(1-y)^4+6y^3-\frac92y^4\right]_0^1 .
        \]
        \item 代入上下限：\(\displaystyle I_3=\left(0+6-\frac{9}{2}\right)-\left(-\frac{2}{3}+0-0\right)=\frac{3}{2}+\frac{2}{3}=\frac{13}{6}\)。
    \end{itemize}
\end{compactenum}
\end{solution}

\begin{exercise}[累次积分计算]
计算下列重积分：
\par\noindent
\begin{tabular}{@{}p{0.48\linewidth}p{0.48\linewidth}@{}}
(1) $\displaystyle \int_{0}^{2}\mathrm{d}x\int_{x}^{2}\mathrm{e}^{y^2}\mathrm{d}y$ &
(2) $\displaystyle \int_{1}^{4}\mathrm{d}y\int_{\sqrt{y}}^{y}x^2y^3\mathrm{d}x$ \\
(3) $\displaystyle \int_{1}^{2}\mathrm{d}y\int_{y}^{y^2}\mathrm{e}^{x/y}\mathrm{d}x$
{\footnotesize (注：原书排版 $e^{x+y}$ 无法积分，依惯例修正为 $e^{x/y}$)} &
(4) $\displaystyle \int_{0}^{1}\mathrm{d}x\int_{0}^{x}y\sin(\pi x)\mathrm{d}y$ \\
\end{tabular}
\par
\end{exercise}
\begin{solution}
\begin{compactenum}
    \item \textbf{解答小题 (1)}：
    \begin{itemize}
        \item \textbf{分析}：内层积分 $\int \mathrm{e}^{y^2} \mathrm{d}y$ 原函数非初等函数，直接计算不可行，必须交换积分次序。
        \item \textbf{确定区域}：由原积分限知，$x \in [0, 2]$，且 $x \leqslant y \leqslant 2$。这表明区域是由 $y=x$, $y=2$ 和 $x=0$ 围成的三角形。
        \item \textbf{交换次序}：改为先对 $x$ 积分。投影到 $y$ 轴，范围 $y \in [0, 2]$。对给定的 $y$，$x$ 从左边界 $0$ 到右边界 $y$：
        \(\displaystyle I_1 = \int_{0}^{2} \mathrm{d}y \int_{0}^{y} \mathrm{e}^{y^2} \,\mathrm{d}x\)。
        \item \textbf{计算}：
        \[ I_1 = \int_{0}^{2} \left[ x\mathrm{e}^{y^2} \right]_{x=0}^{x=y} \mathrm{d}y = \int_{0}^{2} y\mathrm{e}^{y^2} \,\mathrm{d}y = \frac{1}{2}\int_{0}^{2} \mathrm{e}^{y^2} \,\mathrm{d}(y^2) = \frac{1}{2} \left[ \mathrm{e}^{y^2} \right]_0^2 = \frac{1}{2}(\mathrm{e}^4 - 1) \]
    \end{itemize}

    \item \textbf{解答小题 (2)}：
    \begin{itemize}
        \item \textbf{分析}：被积函数为简单的多项式，直接按原有次序积分即可。
        \item \textbf{内层积分}（对 $x$ 积分）：
        \(\displaystyle \int_{\sqrt{y}}^{y} x^2y^3 \,\mathrm{d}x = y^3 \left[ \frac{1}{3}x^3 \right]_{x=\sqrt{y}}^{x=y} = \frac{1}{3}y^3 \left( y^3 - y^{3/2} \right) = \frac{1}{3}(y^6 - y^{9/2})\)。
        \item \textbf{外层积分}（对 $y$ 积分）：\(\displaystyle I_2=\frac13\int_1^4(y^6-y^{9/2})\,\mathrm{d}y=\frac13\left(\frac{16383}{7}-\frac{2}{11}(2048-1)\right)\)。
        从而 \(\displaystyle I_2=\frac{1}{3} \left( \frac{16383}{7} - \frac{4094}{11} \right) = \frac{151555}{231}\)。
    \end{itemize}

    \item \textbf{解答小题 (3)}：
    \begin{itemize}
        \item \textbf{分析}：被积函数 $\mathrm{e}^{x/y}$ 关于 $x$ 的积分可以直接计算，因此按原次序处理即可。
        \item \textbf{内层积分}：
        \(\displaystyle \int_{y}^{y^2} \mathrm{e}^{x/y} \,\mathrm{d}x = \left[ y\mathrm{e}^{x/y} \right]_{x=y}^{x=y^2} = y\left( \mathrm{e}^{y^2/y} - \mathrm{e}^{y/y} \right) = y(\mathrm{e}^y - \mathrm{e})\)。
        \item \textbf{外层积分}：利用分部积分法：
        \(\displaystyle I_3 = \int_{1}^{2} y(\mathrm{e}^y - \mathrm{e}) \,\mathrm{d}y = \int_{1}^{2} y\mathrm{e}^y \,\mathrm{d}y - \int_{1}^{2} \mathrm{e}y \,\mathrm{d}y\)。
        其中 $\int_{1}^{2} y\mathrm{e}^y \,\mathrm{d}y = [y\mathrm{e}^y]_1^2 - \int_{1}^{2} \mathrm{e}^y \,\mathrm{d}y = (2\mathrm{e}^2 - \mathrm{e}) - (\mathrm{e}^2 - \mathrm{e}) = \mathrm{e}^2$。
        \item \textbf{组合结果}：
        \(\displaystyle I_3 = \mathrm{e}^2 - \left[ \frac{\mathrm{e}}{2}y^2 \right]_1^2 = \mathrm{e}^2 - \left( 2\mathrm{e} - \frac{\mathrm{e}}{2} \right) = \mathrm{e}^2 - \frac{3\mathrm{e}}{2}\)。
    \end{itemize}

    \item \textbf{解答小题 (4)}：
    \begin{itemize}
        \item \textbf{分析}：内层积分对 $y$ 进行，被积函数中含有 $\sin(\pi x)$，可将其视为常数。
        \item \textbf{内层积分}：
        \(\displaystyle \int_{0}^{x} y\sin(\pi x) \,\mathrm{d}y = \sin(\pi x) \left[ \frac{1}{2}y^2 \right]_0^x = \frac{1}{2}x^2\sin(\pi x)\)
        \item \textbf{外层积分}：需用到两次分部积分法：
        \begin{align*}
        I_4
        &= \frac{1}{2} \int_{0}^{1} x^2\sin(\pi x) \,\mathrm{d}x
        = \frac{1}{2} \int_{0}^{1} x^2 \,\mathrm{d}\left(-\frac{\cos(\pi x)}{\pi}\right) \\
        &= \frac{1}{2} \left[ -\frac{x^2\cos(\pi x)}{\pi} \right]_0^1 - \frac{1}{2} \int_{0}^{1} \left(-\frac{\cos(\pi x)}{\pi}\right) \cdot 2x \,\mathrm{d}x .
        \end{align*}
        代入第一部分 $[ \dots ]_0^1 = -\frac{1\cdot(-1)}{\pi} - 0 = \frac{1}{\pi}$。第二部分继续分部：
        \[
        \frac{1}{\pi} \int_{0}^{1} x\cos(\pi x) \,\mathrm{d}x
        = \frac{1}{\pi} \left[ \frac{x\sin(\pi x)}{\pi} \right]_0^1 - \frac{1}{\pi^2} \int_{0}^{1} \sin(\pi x) \,\mathrm{d}x.
        \]
        第一项在 $1$ 和 $0$ 处正弦皆为 $0$。只剩下 $-\frac{1}{\pi^2} \left[ -\frac{\cos(\pi x)}{\pi} \right]_0^1 = \frac{1}{\pi^3}(-1 - 1) = -\frac{2}{\pi^3}$。
        \item \textbf{最终合并}：
        \(\displaystyle I_4 = \frac{1}{2\pi} - \frac{2}{\pi^3}\)
    \end{itemize}
\end{compactenum}
\end{solution}

\begin{exercise}[交换积分次序并计算]
在下列积分中先交换积分的次序，然后算出积分值：
\par\noindent
\begin{tabular}{@{}p{0.48\linewidth}p{0.48\linewidth}@{}}
(1) $\displaystyle \int_{0}^{1}\mathrm{d}y\int_{y}^{1}\mathrm{e}^{x^2}\mathrm{d}x$ &
(2) $\displaystyle \int_{0}^{3}\mathrm{d}y\int_{y^2}^{9}y\sin(x^2)\mathrm{d}x$ \\
(3) $\displaystyle \int_{0}^{1}\mathrm{d}y\int_{\sqrt{y}}^{1}\sqrt{2+x^3}\mathrm{d}x$ &
(4) \shortstack[l]{$\displaystyle \int_{-4}^{0}\mathrm{d}x\int_{0}^{2x+8}(8-y)^{1/3}\mathrm{d}y$\\
$\displaystyle + \int_{0}^{4}\mathrm{d}x\int_{0}^{-2x+8}(8-y)^{1/3}\mathrm{d}y$} \\
\end{tabular}
\par
\end{exercise}
\begin{solution}
\begin{compactenum}
    \item \textbf{解答小题 (1)}：
    \begin{itemize}
        \item 原区域 $D$ 满足 $0 \leqslant y \leqslant 1,\ y \leqslant x \leqslant 1$。
        \item 交换次序后，先取 $x\in[0,1]$，再对固定的 $x$ 取 $y\in[0,x]$。
        \item 交换次序计算：\(\displaystyle
        I_1=\int_{0}^{1} \mathrm{d}x \int_{0}^{x} \mathrm{e}^{x^2} \,\mathrm{d}y
        =\int_{0}^{1} x\mathrm{e}^{x^2} \,\mathrm{d}x
        =\frac{1}{2}(\mathrm{e}-1)\)。
    \end{itemize}

    \item \textbf{解答小题 (2)}：
    \begin{itemize}
        \item 原区域 $D$: $0 \leqslant y \leqslant 3$, $y^2 \leqslant x \leqslant 9$。这是一条抛物线 $x = y^2$ 和直线 $x=9, y=0$ 围成的区域。
        \item 改为先对 $y$ 积分：$x$ 的取值范围是 $[0, 9]$，对于给定的 $x$，$y$ 从 $0$ 上升到 $\sqrt{x}$。
        \item 交换次序计算：\(\displaystyle
        I_2=\int_{0}^{9} \mathrm{d}x \int_{0}^{\sqrt{x}} y\sin(x^2) \,\mathrm{d}y
        =\frac{1}{2} \int_{0}^{9} x\sin(x^2) \,\mathrm{d}x
        =\frac{1}{4}(1-\cos 81)\)。
    \end{itemize}

    \item \textbf{解答小题 (3)}：
    \begin{itemize}
        \item 原区域 $D$: $0 \leqslant y \leqslant 1$, $\sqrt{y} \leqslant x \leqslant 1$。边界是 $x = \sqrt{y}$（即 $y = x^2$）和 $x=1$。
        \item 改为先对 $y$ 积分：$x$ 范围为 $[0, 1]$，对于固定的 $x$，$y$ 从 $0$ 到 $x^2$。
        \item 交换次序计算：\(\displaystyle I_3=\int_0^1\mathrm{d}x\int_0^{x^2}\sqrt{2+x^3}\,\mathrm{d}y=\int_0^1x^2\sqrt{2+x^3}\,\mathrm{d}x=\frac29(3\sqrt3-2\sqrt2)\)。
    \end{itemize}

    \item \textbf{解答小题 (4)}：
    \begin{itemize}
        \item 分析区域：原式是两个积分之和。
        \item 第一块对应
        $x \in [-4, 0]$、$y \in [0, 2x+8]$，
        顶点为 $(-4,0)$、$(0,0)$、$(0,8)$。
        \item 第二块对应
        $x \in [0, 4]$、$y \in [0, -2x+8]$，
        顶点为 $(0,0)$、$(4,0)$、$(0,8)$。
        \item 合并后，
        整个区域 $D$ 是以 $(-4,0)$、$(4,0)$、$(0,8)$ 为顶点的大三角形。
        \item 交换次序：从 $y$-轴方向看，整体 $y \in [0, 8]$。对于任一高度 $y$，水平线从左侧边界穿到右侧边界。
        左边界 $y = 2x+8 \implies x = \frac{y-8}{2}$。
        右边界 $y = -2x+8 \implies x = \frac{8-y}{2}$。
        \item 建立新积分并计算内层：\(\displaystyle I_4 = \int_{0}^{8} \mathrm{d}y \int_{\frac{y-8}{2}}^{\frac{8-y}{2}} (8-y)^{1/3} \,\mathrm{d}x = \int_{0}^{8} (8-y)^{1/3} \left[ x \right]_{\frac{y-8}{2}}^{\frac{8-y}{2}} \mathrm{d}y\)。
        上限减下限：$\frac{8-y}{2} - \frac{y-8}{2} = 8-y$。
        \item 计算外层定积分：\(\displaystyle
        I_4 = \int_{0}^{8} (8-y)^{4/3} \,\mathrm{d}y
        = -\left[ \frac{3}{7}(8-y)^{7/3} \right]_0^8
        = \frac{384}{7}\)。
    \end{itemize}
\end{compactenum}
\end{solution}

\begin{exercise}[含有绝对值的二重积分]
求 $I = \displaystyle \iint\limits_D |y-x^2| \,\mathrm{d}\sigma$，其中 $D: -1 \leqslant x \leqslant 1, 0 \leqslant y \leqslant 1$。
\end{exercise}
\begin{solution}
\begin{compactenum}
    \item \textbf{消除绝对值符号}：由于绝对值的特性，需根据 $y-x^2$ 的正负号拆分积分区域。
    令抛物线 $y = x^2$ 成为分界线。在矩形 $D$ 中：
    \begin{itemize}
        \item 区域 $D_1$: $y > x^2$ 的部分（抛物线上方），此时 $|y-x^2| = y-x^2$。
        \item 区域 $D_2$: $y < x^2$ 的部分（抛物线下方），此时 $|y-x^2| = x^2-y$。
    \end{itemize}
    \item \textbf{分别设立积分限}：以 $x$-型区域视角进行划分。
    \begin{itemize}
        \item 对 $D_1$：$x \in [-1, 1]$，$y$ 从 $x^2$ 到 $1$。
        \item 对 $D_2$：$x \in [-1, 1]$，$y$ 从 $0$ 到 $x^2$。
    \end{itemize}
    \item \textbf{建立积分式并利用对称性}：由于积分区域关于 $y$ 轴对称，且被积函数均为偶函数，可仅计算 $[0, 1]$ 区间再乘 2：\(\displaystyle I = 2 \left( \int_{0}^{1} \mathrm{d}x \int_{x^2}^{1} (y-x^2) \,\mathrm{d}y + \int_{0}^{1} \mathrm{d}x \int_{0}^{x^2} (x^2-y) \,\mathrm{d}y \right)\)。
    \item \textbf{计算两个区域并相加}：
    \begin{align*}
    \int_{x^2}^{1}(y-x^2)\,\mathrm{d}y
    &=\left[\frac{1}{2}y^2-x^2y\right]_{x^2}^{1}
    =\frac{1}{2}-x^2+\frac{1}{2}x^4,\\
    \int_{0}^{x^2}(x^2-y)\,\mathrm{d}y
    &=\left[x^2y-\frac{1}{2}y^2\right]_0^{x^2}=\frac{1}{2}x^4,\\
    I&=2\int_{0}^{1}\left(\frac{1}{2}-x^2+x^4\right)\mathrm{d}x
    =2\left(\frac{1}{2}-\frac{1}{3}+\frac{1}{5}\right)=\frac{11}{15}.
    \end{align*}
\end{compactenum}
\end{solution}

\begin{exercise}[极坐标与绝对值的综合应用]
求 $I = \displaystyle \iint\limits_D |xy| \,\mathrm{d}\sigma$，其中 $D: x^2+y^2 \leqslant a^2$。
\end{exercise}
\begin{solution}
\begin{compactenum}
    \item \textbf{先用对称性去绝对值}：
    区域 $D$ 关于 $x$ 轴、$y$ 轴都对称，而 $|xy|$ 在四个象限的值相同，所以只算第一象限 $D_1$ 就够了：
    \(\displaystyle I=4\iint\limits_{D_1}xy\,\mathrm{d}\sigma\)。
    \item \textbf{改写为极坐标}：
    在 $D_1$ 内有 \(0\leqslant \theta\leqslant \frac{\pi}{2},\ 0\leqslant r\leqslant a\)，且 \(\mathrm{d}\sigma=r\,\mathrm{d}r\,\mathrm{d}\theta\)。
    \item \textbf{直接计算}：\(\displaystyle
    I=4\int_{0}^{\frac{\pi}{2}}\int_{0}^{a}r^3\sin\theta\cos\theta\,\mathrm{d}r\,\mathrm{d}\theta
    =4\left(\int_{0}^{\frac{\pi}{2}}\sin\theta\cos\theta\,\mathrm{d}\theta\right)\left(\int_{0}^{a}r^3\,\mathrm{d}r\right)
    =\frac{a^4}{2}\)。
\end{compactenum}
\end{solution}

\begin{exercise}[利用区域结构简化凑微分]
求 $I = \displaystyle \iint\limits_D \frac{\sin y}{y} \,\mathrm{d}x\,\mathrm{d}y$，$D$ 由 $y^2=x$ 及 $y=x$ 围成。
\end{exercise}
\begin{solution}
\begin{compactenum}
    \item \textbf{分析积分区域}：
    联立方程 $y=x$ 与 $x=y^2$。消去 $x$ 得 $y^2 = y \implies y(y-1) = 0$。
    交点为 $(0,0)$ 和 $(1,1)$。
    由于在区间 $y \in [0, 1]$ 之间有 $y^2 \leqslant y$，这表示直线在抛物线的右方，抛物线构成左边界。
    \item \textbf{选择最佳积分次序}：
    如果先对 $y$ 积分，原函数 $\frac{\sin y}{y}$ 是无法写出初等表达式的超越函数。
    因此，必须先对 $x$ 积分，即采用 $y$-型区域视角。
    投影区间为 \(\displaystyle y \in [0, 1],\ x \in [y^2, y]\)。
    \item \textbf{建立积分并求解内层}：\(\displaystyle I = \int_{0}^{1} \mathrm{d}y \int_{y^2}^{y} \frac{\sin y}{y} \,\mathrm{d}x = \int_{0}^{1} \frac{\sin y}{y} (y - y^2) \,\mathrm{d}y\)。
    \item \textbf{计算化简后的外层积分}：
    消去分母上的 $y$，得 \(\displaystyle I = \int_{0}^{1} (1 - y)\sin y \,\mathrm{d}y\)。
    应用分部积分法：
    \begin{align*}
    I&=\int_{0}^{1} (1-y) \,\mathrm{d}(-\cos y)\\
    &= \left[ -(1-y)\cos y \right]_0^1
    - \int_{0}^{1} (-\cos y)(-1) \,\mathrm{d}y\\
    &=1-\sin 1.
    \end{align*}
\end{compactenum}
\end{solution}

\begin{exercise}[极坐标常数限条件分析]
在怎样的情况下，当变换为极坐标之后，二重积分的上下限全是常数？
\end{exercise}
\begin{solution}
\begin{compactenum}
    \item \textbf{回顾极坐标积分的基本形式}：
    化为极坐标后的累次积分通常写成
    \[
    \int_{\alpha}^{\beta}\mathrm{d}\theta
    \int_{r_1(\theta)}^{r_2(\theta)}
    f(r\cos\theta, r\sin\theta)r \,\mathrm{d}r.
    \]
    \item \textbf{分析常数限的几何含义}：
    要使得积分限全部是常数，必须满足：
    \begin{itemize}
        \item 极角 $\theta$ 的范围独立，即区域夹在两条由原点出发的固定射线 $\theta = \alpha$ 与 $\theta = \beta$ 之间。
        \item 极径 $r$ 的范围独立于 $\theta$，即下边界为固定半径圆弧 $r = r_1$（常数），上边界为固定半径圆弧 $r = r_2$（常数）。
    \end{itemize}
    \item \textbf{得出最终结论}：
    当且仅当积分区域 $D$ 是由以原点为圆心的两个同心圆（或一个圆和原点）以及从原点引出的两条射线所围成的扇形、圆环或者圆环扇形区域时，变换为极坐标后的积分限才全都是常数。
\end{compactenum}
\end{solution}

\begin{exercise}[极坐标积分基础练习]
利用极坐标计算下列二重积分：
\par\noindent
\begin{tabular}{@{}p{\linewidth}@{}}
(1) $\displaystyle \iint\limits_D \mathrm{e}^{x^2+y^2}\mathrm{d}\sigma,\quad D:x^2+y^2\leqslant 4$ \\
(2) $\displaystyle \int_{0}^{1}\mathrm{d}x\int_{0}^{\sqrt{1-x^2}}x^3\mathrm{d}y$ \\
(3) $\displaystyle \iint\limits_D \ln(1+x^2+y^2)\mathrm{d}\sigma$，$D$ 为圆周 $x^2+y^2=1$ 及坐标轴在第一象限围成区域 \\
(4) $\displaystyle \iint\limits_D \arctan\dfrac{y}{x}\mathrm{d}\sigma$，$D$ 是圆环 $1\leqslant x^2+y^2\leqslant 4$ 与射线 $y=0,y=x$ 所围成的第一象限闭区域 \\
\end{tabular}
\par
\end{exercise}
\begin{solution}
\begin{compactenum}
    \item \textbf{解答小题 (1)}：积分区域是以原点为中心、半径为 $2$ 的圆域，故极坐标下 $\theta \in [0,2\pi],\ r \in [0,2]$，于是 \(\displaystyle I_1=\int_{0}^{2\pi}\mathrm{d}\theta\int_{0}^{2}\mathrm{e}^{r^2}r\,\mathrm{d}r=2\pi\int_{0}^{2}\mathrm{e}^{r^2}r\,\mathrm{d}r=\pi(\mathrm{e}^4-1)\)。
    \item \textbf{解答小题 (2)}：区域是第一象限单位四分之一圆域，故 $\theta \in \left[0,\frac{\pi}{2}\right],\ r \in [0,1]$，且 $x^3=r^3\cos^3\theta$。因此 \(\displaystyle I_2=\left(\int_{0}^{\frac{\pi}{2}}\cos^3\theta\,\mathrm{d}\theta\right)\left(\int_{0}^{1}r^4\,\mathrm{d}r\right)=\frac{2}{15}\)。
    \item \textbf{解答小题 (3)}：同样在四分之一圆域上，\(\displaystyle I_3=\frac{\pi}{2}\int_{0}^{1}\ln(1+r^2)r\,\mathrm{d}r=\frac{\pi}{4}(2\ln 2-1)\)。
    其中令 $u=1+r^2$，即 $r\,\mathrm{d}r=\dfrac{1}{2}\mathrm{d}u$。
    \item \textbf{解答小题 (4)}：区域是极角 $\theta \in \left[0,\frac{\pi}{4}\right]$、极径 $r\in[1,2]$ 的圆环扇形，且在第一象限有 $\arctan\!\left(\dfrac{y}{x}\right)=\theta$。于是 \(\displaystyle I_4=\int_{0}^{\frac{\pi}{4}}\theta\,\mathrm{d}\theta\int_{1}^{2}r\,\mathrm{d}r=\frac{3\pi^2}{64}\)。
\end{compactenum}
\end{solution}

\begin{exercise}[利用对称性判定等式]
下列等式是否成立？请说明理由，其中积分区域 $D: x^2+y^2 \leqslant 1; D_1: x^2+y^2 \leqslant 1, x \geqslant 0, y \geqslant 0$。
\par\noindent
\begin{tabular}{@{}p{0.48\linewidth}p{0.48\linewidth}@{}}
(1) $\displaystyle \iint\limits_D x\ln(x^2+y^2)\mathrm{d}\sigma=0$ &
(2) \makecell[l]{$\displaystyle \iint\limits_D \sqrt{1-x^2-y^2}\mathrm{d}\sigma$\\$\displaystyle = 4\iint\limits_{D_1}\sqrt{1-x^2-y^2}\mathrm{d}\sigma$} \\
(3) $\displaystyle \iint\limits_D xy\mathrm{d}\sigma=4\iint\limits_{D_1}xy\mathrm{d}x\,\mathrm{d}y$ &
(4) $\displaystyle \iint\limits_D |xy|\mathrm{d}\sigma=4\iint\limits_{D_1}xy\mathrm{d}\sigma$ \\
\end{tabular}
\par
\end{exercise}
\begin{solution}
\begin{compactenum}
    \item \textbf{解答小题 (1) —— 成立}。
    理由：全圆域 $D$ 关于 $y$ 轴对称。被积函数 $f(x,y) = x\ln(x^2+y^2)$ 满足 $f(-x,y) = -x\ln((-x)^2+y^2) = -f(x,y)$，说明其是关于 $x$ 的奇函数。利用奇函数在对称区域的积分定理，积分为 $0$。等式成立。

    \item \textbf{解答小题 (2) —— 成立}。
    理由：全圆域 $D$ 关于 $x$ 轴和 $y$ 轴均对称。第一象限 $D_1$ 是全区域的四分之一。被积函数 $f(x,y) = \sqrt{1-x^2-y^2}$ 是关于 $x$ 的偶函数，也是关于 $y$ 的偶函数。所以一四象限体积相等，一二象限体积相等。总积分为第一象限的 $4$ 倍。等式成立。

    \item \textbf{解答小题 (3) —— 不成立}。
    理由：全圆域 $D$ 关于 $y$ 轴对称，而被积函数 $f(x,y) = xy$ 满足 $f(-x, y) = -xy = -f(x,y)$，是奇函数，故等式左侧 $\iint_D xy \,\mathrm{d}\sigma = 0$。然而在第一象限内 $x>0, y>0$，等式右侧 $4\iint_{D_1} xy \,\mathrm{d}x\,\mathrm{d}y > 0$。左侧不等于右侧，等式不成立。

    \item \textbf{解答小题 (4) —— 成立}。
    理由：与上题不同，加入绝对值后，被积函数变为 $|xy|$，它关于坐标轴对称且非负。利用这种对称性，全圆积分可化为第一象限积分的 4 倍。由于在 $D_1$ 内 $x,y \geqslant 0$，绝对值可直接去掉，故等式成立。
\end{compactenum}
\end{solution}

\subsubsection{B组}

\begin{exercise}[超越区域上的常规运算]
求 $I = \displaystyle \iint\limits_D (x^2-y^2) \,\mathrm{d}\sigma$，其中 $D: 0 \leqslant y \leqslant \sin x, 0 \leqslant x \leqslant \pi$。
\end{exercise}
\begin{solution}
\begin{compactenum}
    \item \textbf{化为累次积分}：区域 $D$ 是典型且标准的 $x$-型区域，其上界由超越函数正弦曲线构成，即 \(\displaystyle I = \int_{0}^{\pi} \mathrm{d}x \int_{0}^{\sin x} (x^2-y^2) \,\mathrm{d}y\)。
    \item \textbf{对 $y$ 进行内层积分}：\(\displaystyle \int_{0}^{\sin x} (x^2-y^2) \,\mathrm{d}y = \left[ x^2y - \frac{1}{3}y^3 \right]_0^{\sin x} = x^2\sin x - \frac{1}{3}\sin^3 x\)。
    \item \textbf{将外层积分分为两部分独立计算}：\(\displaystyle I = \int_{0}^{\pi} x^2\sin x \,\mathrm{d}x - \frac{1}{3} \int_{0}^{\pi} \sin^3 x \,\mathrm{d}x\)。
    \item \textbf{利用分部积分计算第一部分}：
    \begin{align*}
    \int_{0}^{\pi} x^2\sin x \,\mathrm{d}x
    &= \int_{0}^{\pi} x^2 \,\mathrm{d}(-\cos x)\\
    &= \left[ -x^2\cos x \right]_0^\pi
    - \int_{0}^{\pi} (-\cos x) \cdot 2x \,\mathrm{d}x .
    \end{align*}
    代入得 $(\pi^2 \cdot 1 - 0) = \pi^2$。对后续项再分部积分：\(\displaystyle 2\int_{0}^{\pi} x\cos x \,\mathrm{d}x =2 \left( \left[ x\sin x \right]_0^\pi - \int_{0}^{\pi} \sin x \,\mathrm{d}x \right)=-4\)。
    第一部分总计为 $\pi^2 - 4$。
    \item \textbf{计算第二部分 $\sin^3 x$ 的积分}：\(\displaystyle \int_{0}^{\pi} \sin^3 x \,\mathrm{d}x =\int_{0}^{\pi}(1-\cos^2x)\sin x\,\mathrm{d}x =-\left[\cos x-\frac{1}{3}\cos^3x\right]_0^\pi=\frac{4}{3}\)。
    所以第二项的值为 $- \frac{1}{3} \times \frac{4}{3} = -\frac{4}{9}$。
    \item \textbf{总计结果}：\(\displaystyle I = \pi^2 - 4 - \frac{4}{9} = \pi^2 - \frac{40}{9}\)。
\end{compactenum}
\end{solution}

\begin{exercise}[转换直角坐标区域至极坐标系]
画出积分区域，把积分 $\iint\limits_D f(x,y)\,\mathrm{d}x\,\mathrm{d}y$ 表示为极坐标形式的二次积分，其中积分区域 $D$ 是：
\begin{shortexenum}
    \shortitem{1}{$x^2+y^2 \leqslant a^2 \ (a>0)$；} &
    \shortitem{2}{$x^2+y^2 \leqslant 2x$；} \\
    \shortitem{3}{$a^2 \leqslant x^2+y^2 \leqslant b^2 \ (0<a<b)$；} &
    \shortitem{4}{$0 \leqslant y \leqslant 1-x, 0 \leqslant x \leqslant 1$。}
\end{shortexenum}
\end{exercise}
\begin{solution}
\begin{compactenum}
    \item \textbf{解答小题 (1) —— 中心圆}：原点为圆心、半径为 $a$ 的圆域在极坐标下直接写成 $\theta \in [0,2\pi],\ r \in [0,a]$，故 \(\displaystyle \int_{0}^{2\pi} \mathrm{d}\theta \int_{0}^{a} f(r\cos\theta, r\sin\theta) \,r\,\mathrm{d}r\)。
    \item \textbf{解答小题 (2) —— 偏心圆}：由 $x^2+y^2 \leqslant 2x$ 得 $(x-1)^2+y^2 \leqslant 1$，这是与 $y$ 轴相切于原点的偏心圆。极角范围为 $\theta \in \left[-\dfrac{\pi}{2},\dfrac{\pi}{2}\right]$，边界方程为 $r \leqslant 2\cos\theta$，故 \(\displaystyle \int_{-\frac{\pi}{2}}^{\frac{\pi}{2}} \mathrm{d}\theta \int_{0}^{2\cos\theta} f(r\cos\theta, r\sin\theta) \,r\,\mathrm{d}r\)。
    \item \textbf{解答小题 (3) —— 同心圆环}：同心圆环直接给出 \(\theta\in[0,2\pi]\)、\(r\in[a,b]\)。故极坐标积分为
    \[
    \int_{0}^{2\pi} \mathrm{d}\theta \int_{a}^{b} f(r\cos\theta, r\sin\theta) \,r\,\mathrm{d}r .
    \]
    \item \textbf{解答小题 (4) —— 直角三角形}：区域是顶点 $(0,0),(1,0),(0,1)$ 的等腰直角三角形，故 $\theta \in \left[0,\dfrac{\pi}{2}\right]$，斜边 $x+y=1$ 化为 \(r=\dfrac{1}{\cos\theta+\sin\theta}\)。于是积分为 \(\displaystyle \int_{0}^{\frac{\pi}{2}} \mathrm{d}\theta \int_{0}^{\frac{1}{\cos\theta+\sin\theta}} f(r\cos\theta, r\sin\theta) \,r\,\mathrm{d}r\)。
\end{compactenum}
\end{solution}

\begin{exercise}[累次积分的坐标系转换]
将下列直角坐标下的累次积分变换为极坐标下的二次积分：
\par\noindent
\begin{tabular}{@{}p{0.06\linewidth}p{0.88\linewidth}@{}}
(1) & \(\displaystyle \int_{0}^{1}\mathrm{d}x\int_{0}^{1}f(x,y)\mathrm{d}y\)；\\
(2) & \(\displaystyle \int_{0}^{1}\mathrm{d}x\int_{1-x}^{\sqrt{1-x^2}}f(x,y)\mathrm{d}y\)。
\end{tabular}
\par
\end{exercise}
\begin{solution}
\begin{compactenum}
    \item \textbf{解答小题 (1) —— 正方形转换}：区域是正方形 $[0,1]\times[0,1]$。以对角线 $y=x$（即 $\theta=\frac{\pi}{4}$）为界，前半段射线先撞到右边界 $x=1$，故 $r=\sec\theta$；后半段先撞到上边界 $y=1$，故 $r=\csc\theta$。因此积分等于
    \[
    \int_{0}^{\frac{\pi}{4}} \mathrm{d}\theta \int_{0}^{\sec\theta}
    f(r\cos\theta, r\sin\theta) \,r\,\mathrm{d}r
    +\int_{\frac{\pi}{4}}^{\frac{\pi}{2}} \mathrm{d}\theta \int_{0}^{\csc\theta}
    f(r\cos\theta, r\sin\theta) \,r\,\mathrm{d}r .
    \]
    \item \textbf{解答小题 (2) —— 弓形区域转换}：区域由直线 \(x+y=1\) 与单位圆弧 \(x^2+y^2=1\) 在第一象限围成，故 \(\theta \in \left[0,\dfrac{\pi}{2}\right]\)，且 \(\displaystyle r=\frac{1}{\cos\theta+\sin\theta}\) 到 \(r=1\)。所以极坐标形式为
    \[
    \int_{0}^{\frac{\pi}{2}} \mathrm{d}\theta
    \int_{\frac{1}{\cos\theta+\sin\theta}}^{1}
    f(r\cos\theta, r\sin\theta) \,r\,\mathrm{d}r .
    \]
\end{compactenum}
\end{solution}

\begin{exercise}[利用极坐标计算累次积分]
利用极坐标计算二重积分：\(\displaystyle \int_{0}^{1}\mathrm{d}x\int_{0}^{x}\sqrt{x^2+y^2}\mathrm{d}y\)。
\end{exercise}
\begin{solution}
\begin{compactenum}
    \item \textbf{根据积分限绘制区域}：原区域由 $y=0,x=1,y=x$ 围成，是第一象限内的直角三角形；又因被积函数 $\sqrt{x^2+y^2}$ 在极坐标下直接化为 $r$，故改用极坐标最自然。
    \item \textbf{求解极坐标边界}：极角位于 $x$ 轴与直线 $y=x$ 之间，故 \(\theta \in \left[0,\dfrac{\pi}{4}\right]\)。极径由边界 \(x=1\) 给出，即 \(r\cos\theta=1\)，所以 \(r=\sec\theta\)。
    \item \textbf{建立积分并计算}：有 \(\displaystyle I = \int_{0}^{\frac{\pi}{4}} \mathrm{d}\theta \int_{0}^{\sec\theta} r^2 \,\mathrm{d}r = \frac{1}{3}\int_{0}^{\frac{\pi}{4}} \sec^3\theta \,\mathrm{d}\theta\)。
    设 $J=\int \sec^3\theta\,\mathrm{d}\theta$，分部积分得 \(\displaystyle J=\int \sec\theta\,\mathrm{d}(\tan\theta) = \sec\theta\tan\theta - \int \sec\theta(\sec^2\theta-1)\,\mathrm{d}\theta\)，
    故 $2J=\sec\theta\tan\theta+\ln|\sec\theta+\tan\theta|$。于是 \(\displaystyle I = \frac{1}{6}\left[\sec\theta\tan\theta+\ln|\sec\theta+\tan\theta|\right]_{0}^{\frac{\pi}{4}} = \frac{1}{6}\left(\sqrt{2}+\ln(\sqrt{2}+1)\right)\)。
\end{compactenum}
\end{solution}

\begin{exercise}[证明积分恒等式(狄利克雷公式特例)]
设 $f(x) \in C[a,b]$，请证明交换积分次序后的等式 \(\displaystyle \int_{a}^{b}\mathrm{d}x\int_{a}^{x}f(y)\mathrm{d}y= \int_{a}^{b}f(y)(b-y)\mathrm{d}y\)。
\end{exercise}
\begin{solution}
\begin{compactenum}
    \item \textbf{解读区域并交换次序}：左侧累次积分对应区域
    \(\displaystyle D=\{(x,y):a\le y\le x\le b\}\)，
    即由直线 $x=b,\ y=a,\ y=x$ 围成的三角形。改写为 $y$-型后可表示为 $a\le y\le b,\ y\le x\le b$，故 \(\displaystyle \int_{a}^{b}\mathrm{d}x\int_{a}^{x}f(y)\,\mathrm{d}y
    = \int_{a}^{b}\mathrm{d}y\int_{y}^{b}f(y)\,\mathrm{d}x\)。
    \item \textbf{计算内层积分}：由于 $f(y)$ 与 $x$ 无关，\(\displaystyle \int_{a}^{b}\mathrm{d}y\int_{y}^{b}f(y)\,\mathrm{d}x
    = \int_{a}^{b} f(y)\left(\int_{y}^{b}1\,\mathrm{d}x\right)\mathrm{d}y
    = \int_{a}^{b} f(y)(b-y)\,\mathrm{d}y\)。
    两端恒等，证毕。
\end{compactenum}
\end{solution}

\begin{exercise}[利用雅可比行列式的一般换元法]
求由抛物线 $y^2=px, y^2=qx\ (0<p<q)$ 以及双曲线 $xy=a, xy=b\ (0<a<b)$ 共同所围成的闭区域的面积。
\end{exercise}
\begin{solution}
\begin{compactenum}
    \item \textbf{分析并选取新变量}：四条边界可整理为 \(\displaystyle \frac{y^2}{x}=p,\qquad \frac{y^2}{x}=q,\qquad xy=a,\qquad xy=b\)，因此令 \(\displaystyle u=\frac{y^2}{x},\qquad v=xy\)。这样原曲边四边形在 $uv$ 平面上就变成标准长方形 \(\displaystyle D'=\{(u,v):p\le u\le q,\ a\le v\le b\}\)。
    \item \textbf{计算雅可比}：由 $uv=y^3$ 得 $y=(uv)^{1/3}$，再得 $x=u^{-1/3}v^{2/3}$。更方便的是先算逆雅可比：
    \[
    \frac{\partial(u,v)}{\partial(x,y)}
    = \det\begin{pmatrix} -y^2/x^2 & 2y/x \\ y & x \end{pmatrix}
    = -\frac{y^2}{x} - \frac{2y^2}{x}
    = -3u,
    \]
    故 \(\displaystyle \left|\frac{\partial(x,y)}{\partial(u,v)}\right|=\frac{1}{3u}\)。
    \item \textbf{代入换元公式求面积}：于是
    \[
    A = \iint_D 1\,\mathrm{d}\sigma
    = \iint_{D'} \frac{1}{3u}\,\mathrm{d}u\,\mathrm{d}v
    = \left(\int_p^q \frac{1}{3u}\,\mathrm{d}u\right)\left(\int_a^b 1\,\mathrm{d}v\right)
    = \frac{1}{3}(b-a)\ln\left(\frac{q}{p}\right).
    \]
\end{compactenum}
\end{solution}


\subsection{第三节习题：无界区域上的广义二重积分}

\begin{exercise}[无界圆外区域的广义积分]
求广义二重积分 \(\displaystyle I = \iint\limits_{x^2+y^2 \geqslant 1} \mathrm{e}^{-(x^2+y^2)} \,\mathrm{d}x\,\mathrm{d}y\)。
\end{exercise}
\begin{solution}
\begin{compactenum}
    \item \textbf{选取坐标系并改写积分}：区域 \(D=\{(x,y):x^2+y^2\geqslant 1\}\) 是单位圆外的无界区域，边界与被积函数都只含 \(x^2+y^2\)，故宜用极坐标 \(x=r\cos\theta,\ y=r\sin\theta\)，其中 \(\mathrm{d}x\,\mathrm{d}y=r\,\mathrm{d}r\,\mathrm{d}\theta\)，且 \(\theta\in[0,2\pi],\ r\in[1,+\infty)\)。于是 \(\displaystyle I=\int_{0}^{2\pi}\mathrm{d}\theta\int_{1}^{+\infty}\mathrm{e}^{-r^2}r\,\mathrm{d}r\)。
    \item \textbf{计算积分}：由于变量可分离，
    \[
    I=\left(\int_{0}^{2\pi}\mathrm{d}\theta\right)
    \left(\int_{1}^{+\infty}\mathrm{e}^{-r^2}r\,\mathrm{d}r\right)
    =2\pi\cdot \frac{1}{2}\lim_{A\to+\infty}\left[-\mathrm{e}^{-r^2}\right]_{1}^{A}
    =\frac{\pi}{\mathrm{e}} .
    \]
\end{compactenum}
\end{solution}

\begin{exercise}[受抛物线约束的无界区域积分]
求广义二重积分 \(\displaystyle I=\iint\limits_D x\mathrm{e}^{-y^2}\,\mathrm{d}x\,\mathrm{d}y\)，其中 $D$ 是由曲线 $y=4x^2$ 和 $y=9x^2$ 在第一象限所围成的无界区域。
\end{exercise}
\begin{solution}
\begin{compactenum}
    \item \textbf{确定积分次序与积分限}：区域 \(D\) 位于第一象限且夹在抛物线 \(y=4x^2,\ y=9x^2\) 之间。若先对 \(y\) 积分，则 \(\mathrm{e}^{-y^2}\) 无初等原函数，故取 \(y\) 为外层变量、\(x\) 为内层变量。对任意 \(y\geqslant 0\)，有 \(\displaystyle \frac{1}{3}\sqrt{y}\leqslant x\leqslant \frac{1}{2}\sqrt{y}\)，从而
    \[
    I=\int_{0}^{+\infty}\mathrm{d}y
    \int_{\frac{1}{3}\sqrt{y}}^{\frac{1}{2}\sqrt{y}}x\mathrm{e}^{-y^2}\,\mathrm{d}x .
    \]
    \item \textbf{计算积分}：先对 \(x\) 积分得
    \[
    \int_{\frac{1}{3}\sqrt{y}}^{\frac{1}{2}\sqrt{y}}x\mathrm{e}^{-y^2}\,\mathrm{d}x
    =\frac{5}{72}y\mathrm{e}^{-y^2}.
    \]
    因此
    \[
    I=\frac{5}{72}\int_{0}^{+\infty}y\mathrm{e}^{-y^2}\,\mathrm{d}y
    =-\frac{5}{144}\lim_{A\to+\infty}\left[\mathrm{e}^{-y^2}\right]_{0}^{A}
    =\frac{5}{144}.
    \]
\end{compactenum}
\end{solution}

\begin{exercise}[含 $\min$ 函数的全平面反常积分（极坐标替代解法）]
用极坐标替代解法计算广义二重积分 \(\displaystyle I=\int_{-\infty}^{+\infty}\int_{-\infty}^{+\infty}\min\{x,y\}\mathrm{e}^{-(x^2+y^2)}\,\mathrm{d}x\,\mathrm{d}y\)。
\end{exercise}
\begin{solution}
\begin{compactenum}
    \item \textbf{按极角分区}：函数 $\min\{x,y\}$ 的分界线是 $y=x$。在极坐标 $x=r\cos\theta,\ y=r\sin\theta$ 下，这对应于 $\theta=\frac{\pi}{4},\frac{5\pi}{4}$，故
    \begin{align*}
    D_1&:\theta\in\left[\frac{\pi}{4},\frac{5\pi}{4}\right],
    &\min\{x,y\}&=x=r\cos\theta,\\
    D_2&:\theta\in\left[-\frac{3\pi}{4},\frac{\pi}{4}\right],
    &\min\{x,y\}&=y=r\sin\theta.
    \end{align*}
    \item \textbf{建立极坐标积分}：记 \(\displaystyle J=\int_{0}^{+\infty}r^{2}\mathrm{e}^{-r^{2}}\,\mathrm{d}r\)，
    则 \(\displaystyle I=\left(\int_{\pi/4}^{5\pi/4}\cos\theta\,\mathrm{d}\theta+\int_{-3\pi/4}^{\pi/4}\sin\theta\,\mathrm{d}\theta\right)J\)。
    \item \textbf{分别计算角向与径向部分}：由分部积分并利用 $\int_{0}^{+\infty}\mathrm{e}^{-r^{2}}\,\mathrm{d}r=\frac{\sqrt{\pi}}{2}$，得
    \[
    J=\lim_{A\to+\infty}\left(\left[-\frac{1}{2}r\mathrm{e}^{-r^{2}}\right]_{0}^{A}
    +\frac{1}{2}\int_{0}^{A}\mathrm{e}^{-r^{2}}\,\mathrm{d}r\right)=\frac{\sqrt{\pi}}{4},
    \]
    且 \(\displaystyle \int_{\pi/4}^{5\pi/4}\cos\theta\,\mathrm{d}\theta=-\sqrt{2}\)，\(\displaystyle \int_{-3\pi/4}^{\pi/4}\sin\theta\,\mathrm{d}\theta=-\sqrt{2}\)。
    \item \textbf{合并结果}：\(\displaystyle I=(-\sqrt{2}-\sqrt{2})\cdot \frac{\sqrt{\pi}}{4}=-\frac{\sqrt{2\pi}}{2}\)。
\end{compactenum}
\end{solution}


\subsection{第四节习题：三重积分}

\subsubsection{A组}

\begin{exercise}[在直角坐标系下化为累次积分]
在直角坐标系下，将三重积分 $I = \iiint\limits_{\Omega} f(x,y,z)\,\mathrm{d}V$ 转化为三次累次积分，其中积分区域 $\Omega$ 分别为：
\begin{shortexenum}
    \shortitem{1}{$\Omega$ 由三个坐标面及平面 $x+y+z=1$ 所围成的四面体；} &
    \shortitem{2}{$\Omega$ 由平面 $z=R \ (R>0)$ 与锥面 $z=\sqrt{x^2+y^2}$ 围成；}\\
    \shortitem{3}{$\Omega$ 由平面 $z=2,z=4$ 及旋转抛物面 $z=x^2+y^2$ 围成；} &
    \shortitem{4}{$\Omega$ 是两个球面 $x^2+y^2+z^2 \leqslant R^2$ 与 $x^2+y^2+(z-R)^2 \leqslant R^2$ 的公共部分，其中 $R>0$。}
\end{shortexenum}
\end{exercise}
\begin{solution}
\begin{compactenum}
    \item \textbf{四面体区域}

    由三个坐标面及平面 $x+y+z=1$ 围成，区域位于第一卦限。投影到 $Oxy$ 平面得
    \(\displaystyle D_{xy}=\{(x,y)\mid x\geqslant 0,\ y\geqslant 0,\ x+y\leqslant 1\}\)，
    且对任意 $(x,y)\in D_{xy}$，有 $0\leqslant z\leqslant 1-x-y$，故
    \(\displaystyle I=\int_{0}^{1}\mathrm{d}x\int_{0}^{1-x}\mathrm{d}y\int_{0}^{1-x-y}f(x,y,z)\,\mathrm{d}z\)。

    \item \textbf{锥面与顶平面的封闭区域}

    下边界为圆锥面 $z=\sqrt{x^2+y^2}$，上边界为平面 $z=R$。两者交线在 $Oxy$ 平面上的投影是圆 $x^2+y^2=R^2$，因此
    \[
    -R\leqslant x\leqslant R,\qquad
    -\sqrt{R^2-x^2}\leqslant y\leqslant \sqrt{R^2-x^2},\qquad
    \sqrt{x^2+y^2}\leqslant z\leqslant R.
    \]
    从而
    \[
    I=\int_{-R}^{R}\mathrm{d}x
    \int_{-\sqrt{R^2-x^2}}^{\sqrt{R^2-x^2}}\mathrm{d}y
    \int_{\sqrt{x^2+y^2}}^{R}f(x,y,z)\,\mathrm{d}z.
    \]

    \item \textbf{抛物面与平面的封闭区域（纠错与详细分析）}

    原解答若只看 $z=x^2+y^2$ 与 $z=2$，会漏掉上平面 $z=4$。此题更适合按 $z$ 切片：立体位于 $2\leqslant z\leqslant 4$ 之间，而固定 $z$ 时，截面由 $x^2+y^2\leqslant z$ 给出，故
    \[
    -\sqrt{z}\leqslant x\leqslant \sqrt{z},\qquad
    -\sqrt{z-x^2}\leqslant y\leqslant \sqrt{z-x^2}.
    \]
    因此
    \[
    I=\int_{2}^{4}\mathrm{d}z
    \int_{-\sqrt{z}}^{\sqrt{z}}\mathrm{d}x
    \int_{-\sqrt{z-x^2}}^{\sqrt{z-x^2}}f(x,y,z)\,\mathrm{d}y.
    \]
    这种顺序避免了向 $Oxy$ 平面投影时的分块讨论。
    \item \textbf{两个球体的相交公共部分}

    公共部分由球面 \(\displaystyle x^2+y^2+z^2\leqslant R^2,\qquad x^2+y^2+(z-R)^2\leqslant R^2\) 决定。联立边界得交线高度 $z=\frac{R}{2}$，故其在 $Oxy$ 平面上的投影满足
    \(\displaystyle x^2+y^2\leqslant R^2-\left(\frac{R}{2}\right)^2=\frac{3}{4}R^2\)。
    对投影圆盘内任意 $(x,y)$，下边界来自上部球体的下半球面 $z=R-\sqrt{R^2-x^2-y^2}$，上边界来自下部球体的上半球面 $z=\sqrt{R^2-x^2-y^2}$，故
    \[
    I=\int_{-\frac{\sqrt{3}}{2}R}^{\frac{\sqrt{3}}{2}R}\mathrm{d}x
    \int_{-\sqrt{\frac{3}{4}R^2-x^2}}^{\sqrt{\frac{3}{4}R^2-x^2}}\mathrm{d}y
    \int_{R-\sqrt{R^2-x^2-y^2}}^{\sqrt{R^2-x^2-y^2}} f(x,y,z)\,\mathrm{d}z.
    \]
\end{compactenum}
\end{solution}

\begin{exercise}[在柱坐标系下化为累次积分]
在柱坐标系 $(r, \theta, z)$ 下，将三重积分 $I = \iiint\limits_{\Omega} f(x,y,z)\,\mathrm{d}V$ 转化为三次累次积分，其中区域 $\Omega$ 分别为：
\begin{shortexenum}
    \shortitem{1}{$\Omega$ 由曲面 $x^2+y^2=2z$ 与平面 $z=2$ 围成；} &
    \shortitem{2}{$\Omega$ 由曲面 $z=\sqrt{2-x^2-y^2}$ 与旋转抛物面 $z=x^2+y^2$ 围成；} \\
    \shortitem{3}{$\Omega$ 是由圆柱面 $x^2+y^2=16$ 与两水平平面 $z=2, z=8$ 所围成的立体；} &
    \shortitem{4}{$\Omega$ 由锥面 $z=\sqrt{x^2+y^2}$ 与上半球面 $z=\sqrt{1-x^2-y^2}$ 围成。}
\end{shortexenum}
\end{exercise}
\begin{solution}
\begin{compactenum}
    \item \textbf{旋转抛物面与平面的封闭体}：

    柱坐标下，下界为 $z=\frac{1}{2}r^2$，上界为 $z=2$，交线满足 $r^2=4$，故
    \[
    0\leqslant \theta\leqslant 2\pi,\qquad
    0\leqslant r\leqslant 2,\qquad
    \frac{1}{2}r^2\leqslant z\leqslant 2,
    \]
    并有
    \[
    I=\int_{0}^{2\pi}\mathrm{d}\theta\int_{0}^{2}\mathrm{d}r
    \int_{\frac{1}{2}r^2}^{2}f(r\cos\theta,r\sin\theta,z)\,r\,\mathrm{d}z.
    \]
    \item \textbf{球面与抛物面的封闭体（冰淇淋筒状）}：

    此时上界为 $z=\sqrt{2-r^2}$，下界为 $z=r^2$。由 $\sqrt{2-r^2}=r^2$ 得 $r=1$，故
    \[
    I=\int_{0}^{2\pi}\mathrm{d}\theta\int_{0}^{1}\mathrm{d}r
    \int_{r^2}^{\sqrt{2-r^2}}f(r\cos\theta,r\sin\theta,z)\,r\,\mathrm{d}z.
    \]
    \item \textbf{被截断的标准圆柱体}：

    这是标准圆柱体，故 \(\displaystyle 0\leqslant \theta\leqslant 2\pi,\quad 0\leqslant r\leqslant 4,\quad 2\leqslant z\leqslant 8\)，并且
    \[
    I=\int_{0}^{2\pi}\mathrm{d}\theta\int_{0}^{4}\mathrm{d}r
    \int_{2}^{8}f(r\cos\theta,r\sin\theta,z)\,r\,\mathrm{d}z.
    \]
    \item \textbf{锥面与球面的结合体}：

    下界为 $z=r$，上界为 $z=\sqrt{1-r^2}$，由 $r=\sqrt{1-r^2}$ 得 $r=\frac{\sqrt{2}}{2}$，于是
    \[
    I=\int_{0}^{2\pi}\mathrm{d}\theta\int_{0}^{\frac{\sqrt{2}}{2}}\mathrm{d}r
    \int_{r}^{\sqrt{1-r^2}}f(r\cos\theta,r\sin\theta,z)\,r\,\mathrm{d}z.
    \]
\end{compactenum}
\end{solution}

\begin{exercise}[在球坐标系下化为累次积分]
在球坐标系 $(r, \varphi, \theta)$ 下，将三重积分 $I = \iiint\limits_{\Omega} f(x,y,z)\,\mathrm{d}V$ 转化为三次累次积分，其中区域 $\Omega$ 分别为：
\begin{shortexenum}
    \shortitem{1}{$\Omega$ 为偏心实心球体 $x^2+y^2+z^2 \leqslant 2Rz \ (R>0)$；} &
    \shortitem{2}{$\Omega$ 为偏心球体 $x^2+y^2+(z-a)^2 \leqslant a^2$ 与圆锥体 $x^2+y^2 \leqslant z^2$ 内部的公共部分；}\\
    \shortitem{3}{$\Omega$ 由锥面 $z=\sqrt{x^2+y^2}$ 与平面 $z=1$ 围成；} &
    \shortitem{4}{$\Omega$ 位于两同心球面间 $a^2 \leqslant x^2+y^2+z^2 \leqslant 4a^2 \ (a>0)$，且约束在锥面内部 $\sqrt{x^2+y^2} \leqslant z$。}
\end{shortexenum}
\end{exercise}
\begin{solution}
\begin{compactenum}
    \item \textbf{偏心球体的球坐标界定}：
    \begin{itemize}
        \item 球心在 $(0,0,R)$，半径为 $R$，其底部与原点相切。整体位于上半空间，故 $\theta\in[0,2\pi]$，$\varphi\in[0,\frac{\pi}{2}]$。
        \item 将不等式 $x^2+y^2+z^2 \leqslant 2Rz$ 代入球坐标关系，可得 $r^2\leqslant 2Rr\cos\varphi$，即 $0\leqslant r\leqslant 2R\cos\varphi$。
        \item 累次积分式为：
        \[ I = \int_{0}^{2\pi}\mathrm{d}\theta\int_{0}^{\frac{\pi}{2}}\mathrm{d}\varphi\int_{0}^{2R\cos\varphi} f(r\sin\varphi\cos\theta, \dots) r^2\sin\varphi\,\mathrm{d}r \]
    \end{itemize}
    \item \textbf{偏心球体与锥体的相交区域}：
    \begin{itemize}
        \item 球体方程同上，边界为 $r = 2a\cos\varphi$。
        \item 圆锥体内部约束为 $x^2+y^2 \leqslant z^2$，化为球坐标即 $(r\sin\varphi)^2 \leqslant (r\cos\varphi)^2 \implies \tan^2\varphi \leqslant 1$。因为整体在上半空间，故天顶角 $\varphi \leqslant \frac{\pi}{4}$。
        \item 因此积分域受到圆锥体的张角限制，极角为 $\varphi \in [0, \frac{\pi}{4}]$，而外表面由球面提供，极径为 $r \in [0, 2a\cos\varphi]$。
        \item 累次积分式为：
        \[ I = \int_{0}^{2\pi}\mathrm{d}\theta\int_{0}^{\frac{\pi}{4}}\mathrm{d}\varphi\int_{0}^{2a\cos\varphi} f(r\sin\varphi\cos\theta, \dots) r^2\sin\varphi\,\mathrm{d}r \]
        \textit{（注：若被积区域为并集等其他情形则需拆分，此处为公共部分交集，由张角与向径直接决定，无须分段。）}
    \end{itemize}
    \item \textbf{平顶圆锥体}：
    \begin{itemize}
        \item 下界锥面 $z=\sqrt{x^2+y^2}$，对应于天顶角 $\varphi = \frac{\pi}{4}$。锥体内部即 $\varphi \in [0, \frac{\pi}{4}]$。
        \item 顶部由水平平面 $z=1$ 封顶。代入球坐标：$r\cos\varphi = 1 \implies r = \frac{1}{\cos\varphi} = \sec\varphi$。
        \item 累次积分式为：\(\displaystyle I = \int_{0}^{2\pi}\mathrm{d}\theta\int_{0}^{\frac{\pi}{4}}\mathrm{d}\varphi\int_{0}^{\sec\varphi} f(r\sin\varphi\cos\theta, \dots) r^2\sin\varphi\,\mathrm{d}r\)。
    \end{itemize}
    \item \textbf{锥内同心球壳}：
    \begin{itemize}
        \item 同心球壳提供直接的常数极径界限：内径 $r=a$，外径 $r=2a$，即 $r \in [a, 2a]$。
        \item 锥面约束 $\sqrt{x^2+y^2} \leqslant z$ 即 $\tan\varphi \leqslant 1$，由于 $z \geqslant 0$，这划定了天顶角的范围为 $\varphi \in [0, \frac{\pi}{4}]$。
        \item 累次积分式为：\(\displaystyle I = \int_{0}^{2\pi}\mathrm{d}\theta\int_{0}^{\frac{\pi}{4}}\mathrm{d}\varphi\int_{a}^{2a} f(r\sin\varphi\cos\theta, \dots) r^2\sin\varphi\,\mathrm{d}r\)。
    \end{itemize}
\end{compactenum}
\end{solution}

\begin{exercise}[三重积分的常规计算]
计算下列三重积分：
\par\noindent
\begin{tabular}{@{}p{0.48\linewidth}p{0.48\linewidth}@{}}
(1) $\displaystyle \iiint\limits_{\Omega} x\,\mathrm{d}x\,\mathrm{d}y\,\mathrm{d}z$ &
(2) $\displaystyle \iiint\limits_{\Omega} \frac{\mathrm{d}V}{(1+x+y+z)^3}$ \\
{\footnotesize $\Omega$ 为三个坐标面及平面 $x+2y+z=1$ 围成} &
{\footnotesize $\Omega$ 是由三个坐标面及 $x+y+z=1$ 所围成的四面体} \\
(3) $\displaystyle \iiint\limits_{\Omega} (x+y+z)\,\mathrm{d}V$ &
(4) $\displaystyle \iiint\limits_{\Omega} z\,\mathrm{d}V$ \\
{\footnotesize $\Omega: x \geqslant 0, y \geqslant 0, z \geqslant 0, x^2+y^2+z^2 \leqslant a^2$} &
{\footnotesize $\Omega$ 由抛物面 $z=x^2+y^2$ 及平面 $z=1, z=2$ 围成} \\
\end{tabular}
\par
\end{exercise}
\begin{solution}
\begin{compactenum}
    \item \textbf{解答小题 (1)}：
    \begin{itemize}
        \item 积分区域 $\Omega$ 位于第一卦限，由截距式知其在 $Oxy$ 平面的投影为 $D_{xy}: x \geqslant 0, y \geqslant 0, x+2y \leqslant 1$。高度范围 $z \in [0, 1-x-2y]$。
        \item 逐层写出积分限：\(\displaystyle I_1 = \int_{0}^{1}\mathrm{d}x\int_{0}^{\frac{1-x}{2}}\mathrm{d}y\int_{0}^{1-x-2y} x\,\mathrm{d}z\)。
        \item 计算内层关于 $z$ 的积分：\(\displaystyle \int_{0}^{1-x-2y}x\,\mathrm{d}z=x(1-x-2y)=x-x^2-2xy\)。
        \item 计算中层关于 $y$ 的积分：\(\displaystyle \int_{0}^{\frac{1-x}{2}} (x-x^2-2xy) \,\mathrm{d}y = \left[ (x-x^2)y - xy^2 \right]_{y=0}^{y=\frac{1-x}{2}} = \frac{x(1-x)^2}{4}\)。
        \item 计算外层关于 $x$ 的积分：\(\displaystyle I_1 = \frac{1}{4} \int_{0}^{1} x(1-x)^2 \,\mathrm{d}x = \frac{1}{4} \int_{0}^{1} (x - 2x^2 + x^3) \,\mathrm{d}x = \frac{1}{48}\)。
    \end{itemize}

    \item \textbf{解答小题 (2)}：
    \begin{itemize}
        \item 积分域为标准单纯形，投影到 $Oxy$ 面为 $D_{xy}: x \geqslant 0, y \geqslant 0, x+y \leqslant 1$。
        \item 化为累次积分：\(\displaystyle I_2 = \int_{0}^{1}\mathrm{d}x\int_{0}^{1-x}\mathrm{d}y\int_{0}^{1-x-y} (1+x+y+z)^{-3} \,\mathrm{d}z\)。
        \item 内层关于 $z$ 的积分凑微分：\(\displaystyle \left[ -\frac{1}{2}(1+x+y+z)^{-2} \right]_{z=0}^{z=1-x-y} = \frac{1}{2(1+x+y)^2} - \frac{1}{8}\)。
        \item 中层关于 $y$ 的积分：\(\displaystyle \int_{0}^{1-x} \left( \frac{1}{2(1+x+y)^2} - \frac{1}{8} \right) \mathrm{d}y = \left[ -\frac{1}{2(1+x+y)} - \frac{y}{8} \right]_{y=0}^{y=1-x}\)。
        代入得 $\left( -\frac{1}{4} - \frac{1-x}{8} \right) - \left( -\frac{1}{2(1+x)} - 0 \right) = \frac{1}{2(1+x)} - \frac{3-x}{8}$。
        \item 外层关于 $x$ 的积分：
        \[
        I_2 = \int_{0}^{1} \left( \frac{1}{2(1+x)} - \frac{3}{8} + \frac{x}{8} \right)\mathrm{d}x
        = \left[ \frac{1}{2}\ln(1+x) - \frac{3}{8}x + \frac{1}{16}x^2 \right]_0^1
        = \frac{1}{2}\ln 2 - \frac{5}{16}
        \]
    \end{itemize}

    \item \textbf{解答小题 (3)}：
    \begin{itemize}
        \item 区域 $\Omega$ 是第一卦限内半径为 $a$ 的八分之一球体。它对于 $x, y, z$ 的地位完全对称。
        \item 由对称性知：$\iiint_{\Omega} x\,\mathrm{d}V = \iiint_{\Omega} y\,\mathrm{d}V = \iiint_{\Omega} z\,\mathrm{d}V$。原积分等于 $3 \iiint_{\Omega} z\,\mathrm{d}V$。
        \item 引入球坐标系：$\theta \in [0, \frac{\pi}{2}]$, $\varphi \in [0, \frac{\pi}{2}]$, $r \in [0, a]$。被积函数 $z = r\cos\varphi$。
        \begin{align*}
        I_3
        &= 3\int_{0}^{\frac{\pi}{2}}\mathrm{d}\theta
           \int_{0}^{\frac{\pi}{2}}\mathrm{d}\varphi
           \int_{0}^{a} r^3\sin\varphi\cos\varphi\,\mathrm{d}r \\
        &= 3\left( \int_{0}^{\frac{\pi}{2}}1\,\mathrm{d}\theta \right)
           \left( \int_{0}^{\frac{\pi}{2}} \sin\varphi\cos\varphi\,\mathrm{d}\varphi \right)
           \left( \int_{0}^{a} r^3\,\mathrm{d}r \right)
        \end{align*}
        \item 分别计算三个简单积分：$\frac{\pi}{2}$、$[\frac{1}{2}\sin^2\varphi]_0^{\pi/2} = \frac{1}{2}$ 和 $\frac{a^4}{4}$。
        \item 相乘即得结果：$I_3 = 3 \cdot \frac{\pi}{2} \cdot \frac{1}{2} \cdot \frac{a^4}{4} = \frac{3\pi}{16}a^4$。
    \end{itemize}

    \item \textbf{解答小题 (4)}：
    \begin{itemize}
        \item 积分区域被限制在平行截面 $z=1$ 和 $z=2$ 之间，先固定 $z$ 看截面最清楚。
        \item $z$ 在区间 $[1, 2]$。对于特定的 $z$，截面 $D_z$ 是由 $x^2+y^2 \leqslant z$ 给出的圆盘，其面积 $S(z) = \pi (\sqrt{z})^2 = \pi z$。
        \item 构建一元定积分：\(\displaystyle I_4 = \int_{1}^{2} \mathrm{d}z \iint\limits_{D_z} z\,\mathrm{d}x\,\mathrm{d}y = \int_{1}^{2} z \cdot S(z) \,\mathrm{d}z = \int_{1}^{2} \pi z^2 \,\mathrm{d}z = \frac{7\pi}{3}\)。
    \end{itemize}
\end{compactenum}
\end{solution}

\begin{exercise}[柱坐标系下的三重积分计算]
利用柱坐标系计算下列三重积分：
\par\noindent
\begin{tabular}{@{}p{0.48\linewidth}p{0.48\linewidth}@{}}
(1) $\displaystyle \iiint\limits_{\Omega} (x^2+y^2+z)\,\mathrm{d}V$ &
(2) $\displaystyle \iiint\limits_{\Omega} (x^2+y^2)\,\mathrm{d}V$ \\
{\footnotesize $\Omega$ 由曲面 $x^2+y^2=2z$ 与平面 $z=4$ 围成} &
{\footnotesize $\Omega$ 由曲面 $z=x^2+y^2$ 及平面 $z=4, z=16$ 围成} \\
(3) $\displaystyle \iiint\limits_{\Omega} z\,\mathrm{d}V$ &
(4) $\displaystyle \iiint\limits_{\Omega} (\sqrt{x^2+y^2})^3\,\mathrm{d}V$ \\
{\footnotesize $\Omega$ 由半球 $z=\sqrt{2-x^2-y^2}$ 及抛物面 $z=x^2+y^2$ 围成} &
{\footnotesize $\Omega$ 由圆柱 $x^2+y^2=9, 16$ 及 $z^2=x^2+y^2, z \geqslant 0$ 围成} \\
\end{tabular}
\par
\end{exercise}
\begin{solution}
\begin{compactenum}
    \item \textbf{解答小题 (1)}：
    \begin{itemize}
        \item 区域上界平面 $z=4$，下界抛物面 $z=\frac{1}{2}r^2$。交线 $r^2=8 \implies r=2\sqrt{2}$。
        \item 积分限：$\theta \in [0, 2\pi]$, $r \in [0, 2\sqrt{2}]$, $z \in [\frac{r^2}{2}, 4]$。被积函数变为 $(r^2+z)$。
        \item 建立柱坐标积分并计算：
        \begin{align*}
        I_1
        &=\int_{0}^{2\pi}\mathrm{d}\theta\int_{0}^{2\sqrt{2}} r\mathrm{d}r
        \int_{\frac{r^2}{2}}^{4}(r^2+z)\,\mathrm{d}z\\
        &=2\pi\int_{0}^{2\sqrt{2}}r\left[r^2z+\frac{1}{2}z^2\right]_{z=r^2/2}^{z=4}\mathrm{d}r\\
        &=2\pi\int_{0}^{2\sqrt{2}}\left(8r+4r^3-\frac{5}{8}r^5\right)\mathrm{d}r\\
        &=2\pi\left[4r^2+r^4-\frac{5}{48}r^6\right]_0^{2\sqrt{2}}
        =\frac{256\pi}{3}.
        \end{align*}
    \end{itemize}

    \item \textbf{解答小题 (2)}：
    \begin{itemize}
        \item 区域 $\Omega$ 夹在两平行平面 $z=4$ 与 $z=16$ 之间，且由抛物面 $z=x^2+y^2$ 包围。利用切片法与极坐标结合更为简捷。
        \item 截面 $D_z$ 为 $x^2+y^2 \leqslant z$，即极径 $r \in [0, \sqrt{z}]$。
        \item 先求截面上的二重积分：\(\displaystyle
        \iint\limits_{D_z} (x^2+y^2) \,\mathrm{d}x\,\mathrm{d}y
        =\int_{0}^{2\pi}\mathrm{d}\theta\int_{0}^{\sqrt{z}} r^3\,\mathrm{d}r
        =\frac{\pi}{2}z^2\)。
        \item 再对 $z$ 积分：
        \(\displaystyle I_2=\int_{4}^{16}\frac{\pi}{2}z^2\,\mathrm{d}z
        =\frac{\pi}{6}(4096-64)=672\pi\)。
        \textit{（注：原书参考答案给出的 $5440\pi$ 是错将题干中的被积函数看作了 $(x^2+y^2)^2$ 导致的。若被积函数为平方，积分则为 $\int_{4}^{16} \frac{\pi}{3}z^3\mathrm{d}z = \frac{\pi}{12}(16^4-4^4) = 5440\pi$。本解析严格遵循题干表达式推导，正确结果确为 $672\pi$。）}
    \end{itemize}

    \item \textbf{解答小题 (3)}：
    \begin{itemize}
        \item 联立方程求交线：$\sqrt{2-r^2} = r^2 \implies r=1$。故投影是单位圆。
        \item 积分域为：$\theta \in [0, 2\pi], r \in [0, 1], z \in [r^2, \sqrt{2-r^2}]$。
        \item 计算积分：
        \begin{align*}
        I_3
        &= \int_{0}^{2\pi}\mathrm{d}\theta\int_{0}^{1} r\mathrm{d}r\int_{r^2}^{\sqrt{2-r^2}} z\,\mathrm{d}z \\
        &= 2\pi \int_{0}^{1} r \left[ \frac{1}{2}z^2 \right]_{r^2}^{\sqrt{2-r^2}} \mathrm{d}r \\
        &= \pi \int_{0}^{1} (2r-r^3-r^5)\,\mathrm{d}r \\
        &= \pi \left[ r^2 - \frac{1}{4}r^4 - \frac{1}{6}r^6 \right]_0^1 \\
        &= \frac{7\pi}{12}.
        \end{align*}
    \end{itemize}

    \item \textbf{解答小题 (4)}：
    \begin{itemize}
        \item 区域被内圆柱 $r=3$ 和外圆柱 $r=4$ 夹住。下界为 $z=0$ 面，上界由锥面 $z=r$ 封顶。
        \item 积分限为：$\theta \in [0, 2\pi], r \in [3, 4], z \in [0, r]$。被积函数化简为 $(r^2)^{3/2} = r^3$。
        \item 构建并计算积分：\(\displaystyle I_4=\int_{0}^{2\pi}\mathrm{d}\theta\int_{3}^{4} r\mathrm{d}r\int_{0}^{r} r^3\,\mathrm{d}z\)。内层给出 $r^4$，故
        \(\displaystyle I_4=2\pi \int_{3}^{4} r^5 \,\mathrm{d}r=\frac{3367\pi}{3}\)。
    \end{itemize}
\end{compactenum}
\end{solution}

\begin{exercise}[球坐标系下的三重积分计算]
利用球坐标系计算下列三重积分：
\begin{shortexenum}
    \shortitem{1}{$\displaystyle \iiint\limits_{\Omega} (x^2+y^2+z^2)\,\mathrm{d}V$，$\Omega: x^2+y^2+z^2 \leqslant 1$。} &
    \shortitem{2}{$\displaystyle \iiint\limits_{\Omega} (x+y+z)\,\mathrm{d}V$，$\Omega: x^2+y^2+z^2 \leqslant R^2$。}\\
    \shortitem{3}{$\displaystyle \iiint\limits_{\Omega} \sqrt{x^2+y^2+z^2}\,\mathrm{d}V$，$\Omega$ 由锥面 $z=\sqrt{x^2+y^2}$ 及平面 $z=1$ 围成。} &
\end{shortexenum}
\end{exercise}
\begin{solution}
\begin{compactenum}
    \item \textbf{解答小题 (1)}：

    单位球在球坐标下满足 $0\leqslant \theta\leqslant 2\pi,\ 0\leqslant \varphi\leqslant \pi,\ 0\leqslant r\leqslant 1$，被积函数化为 $r^2$，故 \(\displaystyle
    I_1=\left(\int_{0}^{2\pi}1\,\mathrm{d}\theta\right)\left(\int_{0}^{\pi}\sin\varphi\,\mathrm{d}\varphi\right)\left(\int_{0}^{1}r^4\,\mathrm{d}r\right)
    =\frac{4\pi}{5}\)。

    \item \textbf{解答小题 (2)}：

    区域 $\Omega$ 是关于三个坐标平面都对称的完整球体，而 $x,y,z$ 分别都是相应变量的奇函数，所以
    \(\displaystyle \iiint_{\Omega}x\,\mathrm{d}V=\iiint_{\Omega}y\,\mathrm{d}V=\iiint_{\Omega}z\,\mathrm{d}V=0\)，
    因而 $I_2=0$。

    \item \textbf{解答小题 (3)}：

    锥面 $z=\sqrt{x^2+y^2}$ 对应天顶角 $\varphi=\frac{\pi}{4}$，其内部满足 $0\leqslant \varphi\leqslant \frac{\pi}{4}$；平面 $z=1$ 化为 $r\cos\varphi=1$，即 $0\leqslant r\leqslant \sec\varphi$。于是 \(\displaystyle
    I_3=\int_{0}^{2\pi}\mathrm{d}\theta\int_{0}^{\frac{\pi}{4}}\mathrm{d}\varphi\int_{0}^{\sec\varphi}r^3\sin\varphi\,\mathrm{d}r
    =\frac{\pi}{2}\int_{0}^{\frac{\pi}{4}}\frac{\sin\varphi}{\cos^{4}\varphi}\,\mathrm{d}\varphi\)。
    令 $u=\cos\varphi$，则 $\mathrm{d}u=-\sin\varphi\,\mathrm{d}\varphi$，故
    \(\displaystyle I_3=\frac{\pi}{2}\int_{\frac{\sqrt{2}}{2}}^{1}u^{-4}\,\mathrm{d}u
    =\frac{\pi}{6}(2\sqrt{2}-1)\)。
\end{compactenum}
\end{solution}

\subsubsection{B组}

\begin{exercise}[利用球坐标解决偏心及复杂曲面区域]
计算下列三重积分：
\begin{shortsingleenum}
    \shortsingleitem{1}{$\displaystyle \iiint\limits_{\Omega} z\,\mathrm{d}V$，$\Omega: x^2+y^2+(z-a)^2 \leqslant a^2, x^2+y^2 \leqslant z^2$。}
    \shortsingleitem{2}{$\displaystyle \iiint\limits_{\Omega} |\sqrt{x^2+y^2+z^2}-1|\,\mathrm{d}V$，$\Omega$ 由锥面 $z=\sqrt{x^2+y^2}$ 及平面 $z=1$ 围成。}
    \shortsingleitem{3}{$\displaystyle \iiint\limits_{\Omega} (x+y+z)\,\mathrm{d}V$，$\Omega: \frac{x^2}{a^2}+\frac{y^2}{b^2}+\frac{z^2}{c^2} \leqslant 1$。}
\end{shortsingleenum}
\end{exercise}
\begin{solution}
\begin{compactenum}
    \item \textbf{解答小题 (1)}：

    采用球坐标 $x=r\sin\varphi\cos\theta,\ y=r\sin\varphi\sin\theta,\ z=r\cos\varphi$，则偏心球面给出 $r\leqslant 2a\cos\varphi$，圆锥内部给出 $0\leqslant \varphi\leqslant \frac{\pi}{4}$，并且 $0\leqslant \theta\leqslant 2\pi$。由于被积函数为 $z=r\cos\varphi$，故
    \[
    I_1=\int_{0}^{2\pi}\mathrm{d}\theta\int_{0}^{\frac{\pi}{4}}\mathrm{d}\varphi\int_{0}^{2a\cos\varphi}(r\cos\varphi)\,r^2\sin\varphi\,\mathrm{d}r
    =8\pi a^4\int_{0}^{\frac{\pi}{4}}\cos^5\varphi\,\mathrm{d}(-\cos\varphi)
    =\frac{7}{6}\pi a^4.
    \]

    \item \textbf{解答小题 (2)}：

    锥面对应 $0\leqslant \varphi\leqslant \frac{\pi}{4}$，平面 $z=1$ 对应 $0\leqslant r\leqslant \sec\varphi$，且 $0\leqslant \theta\leqslant 2\pi$。由于 $\sec\varphi\geqslant 1$，处理被积函数 $|r-1|$ 时需在 $r=1$ 处分段。记
    \(\displaystyle I_2=\int_{0}^{2\pi}\mathrm{d}\theta\int_{0}^{\frac{\pi}{4}}\sin\varphi\,\mathrm{d}\varphi\int_{0}^{\sec\varphi}|r-1|r^2\,\mathrm{d}r\)，
    并设内层积分为 $K(\varphi)$，则
    \begin{align*}
    K(\varphi)
    &= \int_{0}^{1} (1-r)r^2\,\mathrm{d}r+\int_{1}^{\sec\varphi}(r-1)r^2\,\mathrm{d}r \\
    &= \left[\frac{1}{3}r^3-\frac{1}{4}r^4\right]_0^1+\left[\frac{1}{4}r^4-\frac{1}{3}r^3\right]_1^{\sec\varphi} \\
    &= \frac{1}{6}+\frac{1}{4}\sec^4\varphi-\frac{1}{3}\sec^3\varphi.
    \end{align*}
    于是 \(\displaystyle \int_0^{\pi/4}K(\varphi)\sin\varphi\,\mathrm{d}\varphi=\frac{\sqrt2-1}{12}\)，
    故 \(\displaystyle I_2=2\pi\cdot\frac{\sqrt2-1}{12}=\frac{\pi}{6}(\sqrt2-1)\)。

    \item \textbf{解答小题 (3)}：

    椭球体 $\Omega:\dfrac{x^2}{a^2}+\dfrac{y^2}{b^2}+\dfrac{z^2}{c^2}\leqslant 1$ 关于三个坐标平面都对称，而
    \(\displaystyle I_3=\iiint\limits_{\Omega}x\,\mathrm{d}V+\iiint\limits_{\Omega}y\,\mathrm{d}V+\iiint\limits_{\Omega}z\,\mathrm{d}V\)，
    中三项分别都是奇函数在对应对称区域上的积分，故都为 $0$，从而 $I_3=0$。
\end{compactenum}
\end{solution}

\begin{exercise}[积分恒等式的对称性证明]
设 $\Omega$ 为标准单位球体 $x^2+y^2+z^2 \leqslant 1$。已知 $f(x)$ 满足可积条件，求证
\(\displaystyle \iiint\limits_{\Omega} f(x)\,\mathrm{d}x\,\mathrm{d}y\,\mathrm{d}z = \pi \int_{-1}^{1} f(z)(1-z^2)\,\mathrm{d}z\)。
\end{exercise}
\begin{solution}
\begin{compactenum}
    \item \textbf{利用球体对称性替换变量}：

    单位球 $\Omega$ 关于原点及各坐标轴都对称，因此交换坐标轴不会改变积分区域，从而
    \(\displaystyle \iiint\limits_{\Omega} f(x) \,\mathrm{d}V = \iiint\limits_{\Omega} f(y) \,\mathrm{d}V = \iiint\limits_{\Omega} f(z) \,\mathrm{d}V\)。
    \item \textbf{固定 $z$ 写截面并计算等效积分}：

    既然三者相等，只需计算最便于切片的 $\iiint_{\Omega} f(z)\,\mathrm{d}V$。沿 $z$ 轴切片时，$z\in[-1,1]$，截面 \(\displaystyle D_z=\{(x,y):x^2+y^2\leqslant 1-z^2\}\) 是半径为 $\sqrt{1-z^2}$ 的圆，所以 \(\displaystyle S(z)=\iint\limits_{D_z}1\,\mathrm{d}x\,\mathrm{d}y=\pi(1-z^2)\)。
    \item \textbf{化三重积分为一元积分并完成证明}：

    由切片法，
    \[
    \iiint\limits_{\Omega} f(z) \,\mathrm{d}V
    =\int_{-1}^{1}\mathrm{d}z\iint\limits_{D_z}f(z)\,\mathrm{d}x\,\mathrm{d}y
    =\int_{-1}^{1}f(z)\left(\iint\limits_{D_z}1\,\mathrm{d}x\,\mathrm{d}y\right)\mathrm{d}z
    =\pi\int_{-1}^{1}f(z)(1-z^2)\,\mathrm{d}z.
    \]
    再结合第 1 步中的变量对调，即得
    \(\displaystyle \iiint\limits_{\Omega} f(x)\,\mathrm{d}x\,\mathrm{d}y\,\mathrm{d}z = \pi \int_{-1}^{1} f(z)(1-z^2)\,\mathrm{d}z\)。
\end{compactenum}
\end{solution}


\subsection{第五节习题：重积分的应用}

\subsubsection{A组}

\begin{exercise}[曲面围成立体的体积计算]
求由下列曲面所围立体的体积：
\par\noindent
\begin{tabular}{@{}p{0.05\linewidth}p{0.9\linewidth}@{}}
(1) & 由 $x=0, y=0, x=1, y=1$ 围成的柱体被 $z=0$ 及 $2x+3y+z=6$ 截得的立体； \\
(2) & $z=x^2+y^2, y=x^2, y=1, z=0$ 所围成立体； \\
(3) & 锥面 $z=\sqrt{x^2+y^2}$ 与抛物面 $z=x^2+y^2$ 所围成立体； \\
(4) & 球面 $x^2+y^2+z^2=2az \ (a>0)$ 及锥面 $x^2+y^2=z^2$（含有 $z$ 轴的部分）所围成立体； \\
(5) & 球面 $x^2+y^2+z^2=1$, $x^2+y^2+z^2=2z$ 以及锥面 $z=\sqrt{x^2+y^2}$ 围成且位于锥面内部的部分。 \\
\end{tabular}
\par
\end{exercise}
\begin{solution}
\begin{compactenum}
    \item \textbf{解答小题 (1) —— 截断方柱体体积}：

    投影域是单位正方形 $D=[0,1]\times[0,1]$，下底为 $z=0$，上顶为 $z=6-2x-3y$，故
    \(\displaystyle V_1 = \int_{0}^{1}\mathrm{d}x\int_{0}^{1} (6-2x-3y) \,\mathrm{d}y
    = \int_{0}^{1} \left( \frac{9}{2} - 2x \right) \mathrm{d}x
    = \frac{7}{2}\)。

    \item \textbf{解答小题 (2) —— 抛物面与平底的封闭体}：

    顶面是抛物面 $z=x^2+y^2$，底面是 $z=0$；投影域由 $y=x^2$ 与 $y=1$ 围成，其交点为 $x=\pm 1$。取 $x$-型积分，得
    \(\displaystyle V_2 = \int_{-1}^{1}\mathrm{d}x\int_{x^2}^{1} (x^2+y^2) \,\mathrm{d}y
    = \int_{-1}^{1} \left( x^2 + \frac{1}{3} - x^4 - \frac{1}{3}x^6 \right) \mathrm{d}x\)。
    被积函数是偶函数，故 \(\displaystyle V_2 = 2 \int_{0}^{1} \left( \frac{1}{3} + x^2 - x^4 - \frac{1}{3}x^6 \right) \mathrm{d}x
    = \frac{88}{105}\)。

    \item \textbf{解答小题 (3) —— 锥面与抛物面交界体}：

    联立 $\sqrt{x^2+y^2}=x^2+y^2$，在极坐标下得 $r=r^2$，故交线对应 $r=1$，投影是单位圆。又在 $0\leqslant r\leqslant 1$ 上有 $r\geqslant r^2$，因此
    \(\displaystyle V_3 = \int_{0}^{2\pi}\mathrm{d}\theta\int_{0}^{1} (r - r^2) r \,\mathrm{d}r
    = 2\pi \int_{0}^{1} (r^2 - r^3) \,\mathrm{d}r
    = \frac{\pi}{6}\)。

    \item \textbf{解答小题 (4) —— 球面与锥面结合体（冰淇淋）}：

    “含有 $z$ 轴的部分”即锥体内部。球坐标下，锥面 $x^2+y^2=z^2$ 对应 $\varphi=\frac{\pi}{4}$，故 $0\leqslant \varphi\leqslant \frac{\pi}{4}$；球面 $x^2+y^2+z^2=2az$ 对应 $0\leqslant r\leqslant 2a\cos\varphi$。因此
    \begin{align*}
    V_4 &= \int_{0}^{2\pi}\mathrm{d}\theta\int_{0}^{\pi/4}\mathrm{d}\varphi\int_{0}^{2a\cos\varphi} r^2\sin\varphi \,\mathrm{d}r \\
        &= 2\pi \int_{0}^{\pi/4} \sin\varphi \left[ \frac{1}{3}r^3 \right]_0^{2a\cos\varphi}\mathrm{d}\varphi \\
        &= \frac{16\pi a^3}{3} \int_{0}^{\pi/4} \cos^3\varphi\sin\varphi \,\mathrm{d}\varphi
         = \frac{16\pi a^3}{3} \left[ -\frac{1}{4}\cos^4\varphi \right]_0^{\pi/4}
         = \pi a^3.
    \end{align*}

    \item \textbf{解答小题 (5) —— 双球面与锥面构成的复杂体}：

    球坐标下，内球面为 $r=1$，外偏心球面为 $r=2\cos\varphi$，锥体内部满足 $0\leqslant \varphi\leqslant \frac{\pi}{4}$。由于在该区间内 $2\cos\varphi\geqslant \sqrt{2}>1$，故内外边界始终分别为 $r=1$ 与 $r=2\cos\varphi$。于是
    \[
    V_5=\int_0^{2\pi}\mathrm{d}\theta
    \int_0^{\pi/4}\mathrm{d}\varphi
    \int_1^{2\cos\varphi}r^2\sin\varphi\,\mathrm{d}r
    =\frac{\pi(1+\sqrt2)}{3}.
    \]
\end{compactenum}
\end{solution}

\begin{exercise}[平面图形的质心计算]
求下列平面均匀薄片（设密度 $\rho \equiv 1$）的质心：
\begin{shortexenum}
    \shortitem{1}{由曲线 $\sqrt{x}+\sqrt{y}=\sqrt{a}$ 与两坐标轴所围成的部分；} &
    \shortitem{2}{椭圆 $\dfrac{x^2}{a^2}+\dfrac{y^2}{b^2} \leqslant 1$ 在第一象限的部分。}
\end{shortexenum}
\end{exercise}
\begin{solution}
\begin{compactenum}
    \item \textbf{解答小题 (1) —— 根式曲线图形}：

    曲线与坐标轴交于 $(a,0)$、$(0,a)$，并且关于 $x,y$ 对称，所以 $\bar{x}=\bar{y}$。由边界方程
    \(\displaystyle y=(\sqrt{a}-\sqrt{x})^2,\ 0\le x\le a\)，得
    \(\displaystyle A=\int_{0}^{a}(\sqrt{a}-\sqrt{x})^2\,\mathrm{d}x=\frac{a^2}{6}\)，
    \(\displaystyle M_y=\iint\limits_D x\,\mathrm{d}x\,\mathrm{d}y=\int_{0}^{a}x(\sqrt{a}-\sqrt{x})^2\,\mathrm{d}x=\frac{a^3}{30}\)。
    因而 \(\displaystyle \bar{x}=\frac{M_y}{A}=\frac{a}{5},\quad \bar{y}=\frac{a}{5}\)，
    即质心为 $\left( \frac{a}{5}, \frac{a}{5} \right)$。

    \item \textbf{解答小题 (2) —— 四分之一椭圆}：

    第一象限部分面积为 \(\displaystyle A=\frac{\pi ab}{4}\)。令 $x = ar\cos\theta,\ y = br\sin\theta$，则 $|J|=abr$，且区域化为 $0\leqslant \theta\leqslant \frac{\pi}{2},\ 0\leqslant r\leqslant 1$。于是
    \(\displaystyle M_y = \iint\limits_{D} x \,\mathrm{d}x\,\mathrm{d}y
    = \int_{0}^{\pi/2}\mathrm{d}\theta \int_{0}^{1} (ar\cos\theta)\cdot abr \,\mathrm{d}r
    = \frac{a^2b}{3}\)，
    同理 $M_x=\dfrac{ab^2}{3}$。因此 \(\displaystyle \bar{x}=\frac{M_y}{A}=\frac{4a}{3\pi},\quad
    \bar{y}=\frac{M_x}{A}=\frac{4b}{3\pi}\)，
    即质心为 $\left( \frac{4a}{3\pi}, \frac{4b}{3\pi} \right)$。
\end{compactenum}
\end{solution}

\begin{exercise}[半球壳区域的质心计算]
求区域 $\Omega: a^2 \leqslant x^2+y^2+z^2 \leqslant b^2$ 与 $x \geqslant 0$ 的质心（密度 $\rho \equiv 1$）。
\end{exercise}
\begin{solution}
\begin{compactenum}
    \item \textbf{分析几何对称性}：

    该区域是球壳被平面 $x=0$ 截成的“右半球壳”，关于 $xOy$ 平面和 $xOz$ 平面都对称，所以 $\bar{y}=0,\ \bar{z}=0$，只需求 $\bar{x}$。
    \item \textbf{确定区域体积 $V$}：

    体积为半个外球减半个内球，即
    \(\displaystyle V=\frac{1}{2}\cdot\frac{4}{3}\pi b^3-\frac{1}{2}\cdot\frac{4}{3}\pi a^3=\frac{2}{3}\pi(b^3-a^3)\)。
    \item \textbf{在球坐标系下计算 $x$ 的力矩}：

    由 $x\geqslant 0$ 得 $\theta\in\left[-\frac{\pi}{2},\frac{\pi}{2}\right]$，并有 $r\in[a,b],\ \varphi\in[0,\pi]$。于是
    \[
    M_{yz} = \iiint\limits_{\Omega} x\,\mathrm{d}V
    = \int_{-\pi/2}^{\pi/2}\mathrm{d}\theta \int_{0}^{\pi}\mathrm{d}\varphi \int_{a}^{b} (r\sin\varphi\cos\theta) r^2\sin\varphi\,\mathrm{d}r.
    \]
    分离变量后得到
    \[
    M_{yz}
    = \left( \int_{-\pi/2}^{\pi/2} \cos\theta \,\mathrm{d}\theta \right)
      \left( \int_{0}^{\pi} \sin^2\varphi \,\mathrm{d}\varphi \right)
      \left( \int_{a}^{b} r^3 \,\mathrm{d}r \right)
    = 2\cdot \frac{\pi}{2}\cdot \frac{b^4-a^4}{4}
    = \frac{\pi}{4}(b^4-a^4).
    \]
    \item \textbf{求取质心坐标}：
    \(\displaystyle \bar{x}=\frac{M_{yz}}{V}
    =\frac{\frac{\pi}{4}(b^4-a^4)}{\frac{2}{3}\pi(b^3-a^3)}
    =\frac{3}{8}\frac{b^4-a^4}{b^3-a^3}\)。
    综上，质心坐标为 $\left( \frac{3(b^4-a^4)}{8(b^3-a^3)}, 0, 0 \right)$。
\end{compactenum}
\end{solution}

\begin{exercise}[变密度球体的质心计算]
求偏心球体 $x^2+y^2+z^2 \leqslant 2az \ (a>0)$ 的质心，假设球体内各点处的体密度等于该点到坐标原点的距离的平方。
\end{exercise}
\begin{solution}
\begin{compactenum}
    \item \textbf{分析区域对称性及密度函数}：

    区域 $\Omega$ 是球心在 $(0,0,a)$、半径为 $a$ 的球，关于 $x=0$ 和 $y=0$ 对称；密度函数 \(\displaystyle \rho(x,y,z)=x^2+y^2+z^2=r^2\) 也是关于这两个平面的偶函数，故 $\bar{x}=0,\ \bar{y}=0$，只需求 $\bar{z}$。
    \item \textbf{建立球坐标系并求总质量 $M$}：

    球面边界在球坐标下为 $r\leqslant 2a\cos\varphi$，且 $0\leqslant \varphi\leqslant \frac{\pi}{2}$，因此
    \[
    M = \iiint\limits_{\Omega} \rho \,\mathrm{d}V
    = \int_{0}^{2\pi}\mathrm{d}\theta \int_{0}^{\pi/2}\mathrm{d}\varphi \int_{0}^{2a\cos\varphi} r^4\sin\varphi\,\mathrm{d}r
    = \frac{64\pi a^5}{5} \int_{0}^{\pi/2} \cos^5\varphi\sin\varphi\,\mathrm{d}\varphi
    = \frac{32\pi a^5}{15}.
    \]
    \item \textbf{求对 $Oxy$ 平面的静力矩 $M_{xy}$}：

    再乘上坐标 $z=r\cos\varphi$，得
    \begin{align*}
    M_{xy}
    &=\iiint\limits_{\Omega} z\rho\,\mathrm{d}V
    =\int_{0}^{2\pi}\mathrm{d}\theta\int_{0}^{\pi/2}\mathrm{d}\varphi
    \int_{0}^{2a\cos\varphi}(r\cos\varphi)r^4\sin\varphi\,\mathrm{d}r\\
    &=\frac{64\pi a^6}{3}\int_{0}^{\pi/2}\cos^7\varphi\sin\varphi\,\mathrm{d}\varphi
    =\frac{8\pi a^6}{3}.
    \end{align*}
    \item \textbf{得出质心 $z$ 坐标}：
    \(\displaystyle \bar{z}=\frac{M_{xy}}{M}
    =\frac{8\pi a^6/3}{32\pi a^5/15}
    =\frac{5}{4}a\)。
    即质心坐标为 $\left( 0, 0, \frac{5}{4}a \right)$。
\end{compactenum}
\end{solution}

\begin{exercise}[平面区域的转动惯量计算]
设均匀平面薄片所占区域 $D$ 由抛物线 $y^2=\dfrac{9}{2}x$ 与直线 $x=2$ 围成。求其对坐标轴的转动惯量 $I_x$ 和 $I_y$（设密度 $\rho \equiv 1$）。
\end{exercise}
\begin{solution}
\begin{compactenum}
    \item 将抛物线写成 $x=\frac{2}{9}y^2$ 后，区域 $D$ 可表示为 \(\displaystyle -3\leqslant y\leqslant 3,\qquad \frac{2}{9}y^2\leqslant x\leqslant 2\)，
    这样直接采用 $y$-型区域最方便。
    \item 于是对 $x$ 轴的转动惯量为
    \[
    I_x=\iint_D y^2\,\mathrm{d}x\,\mathrm{d}y
    =\int_{-3}^{3}\mathrm{d}y\int_{2y^2/9}^{2}y^2\,\mathrm{d}x
    =2\int_{0}^{3}\left(2y^2-\frac{2}{9}y^4\right)\mathrm{d}y
    =\frac{72}{5}.
    \]
    \item 对 $y$ 轴同理有
    \[
    I_y=\iint_D x^2\,\mathrm{d}x\,\mathrm{d}y
    =\int_{-3}^{3}\mathrm{d}y\int_{2y^2/9}^{2}x^2\,\mathrm{d}x
    =\frac{2}{3}\int_{0}^{3}\left(8-\frac{8}{729}y^6\right)\mathrm{d}y
    =\frac{96}{7}.
    \]
\end{compactenum}
\end{solution}

\begin{exercise}[非均匀球体的转动惯量]
设球体在动点 $P(x,y,z)$ 的密度与该点到球心距离成正比。已知该球体 $x^2+y^2+z^2 \leqslant R^2$ 的总质量为 $m$，求其对于其任一直径的转动惯量。
\end{exercise}
\begin{solution}
\begin{compactenum}
    \item 设 $r=\sqrt{x^2+y^2+z^2}$，则密度为 $\rho(r)=kr$。由总质量
    \[
    m=\iiint\limits_{\Omega}kr\,\mathrm{d}V
    =\int_{0}^{2\pi}\mathrm{d}\theta\int_{0}^{\pi}\sin\varphi\,\mathrm{d}\varphi\int_{0}^{R}kr^3\,\mathrm{d}r
    =\pi kR^4
    \]
    得 $k=\dfrac{m}{\pi R^4}$。
    \item 由球体的对称性，只需求其对 $z$ 轴的转动惯量：
    \[
    I_z=\iiint_{\Omega}(x^2+y^2)\rho(r)\,\mathrm{d}V
    =\iiint_{\Omega}r^2\sin^2\varphi\cdot kr\cdot r^2\sin\varphi\,\mathrm{d}r\,\mathrm{d}\varphi\,\mathrm{d}\theta.
    \]
    \item 因而
    \[
    I_z
    =k\left(\int_{0}^{2\pi}\mathrm{d}\theta\right)\left(\int_{0}^{\pi}\sin^3\varphi\,\mathrm{d}\varphi\right)\left(\int_{0}^{R}r^5\,\mathrm{d}r\right)
    =k\cdot 2\pi\cdot \frac{4}{3}\cdot \frac{R^6}{6}
    =\frac{4\pi kR^6}{9}
    =\frac{4}{9}mR^2.
    \]
\end{compactenum}
\end{solution}

\subsubsection{B组}

\begin{exercise}[重积分计算复杂立体的体积]
求下列各曲面所围立体的体积：
\begin{shortexenum}
    \shortitem{1}{$z=xy, x+y+z=1$ 及坐标平面 $y=0, x=0$；} &
    \shortitem{2}{$z=x^2+2y^2$ 与 $z=6-2x^2-y^2$。}
\end{shortexenum}
\end{exercise}
\begin{solution}
\begin{compactenum}
    \item \textbf{解答小题 (1) —— 鞍面与截面的组合体积}：立体位于第一卦限，下底为 $z=xy$，上顶为 $z=1-x-y$。两曲面交线在 $Oxy$ 面上的投影满足 \(\displaystyle xy=1-x-y \iff (x+1)(y+1)=2 \iff y=\frac{1-x}{1+x}\)，
    故投影域 $D$ 由 $x$ 轴、$y$ 轴与该双曲线围成，且 $0\leqslant x\leqslant 1$。于是令 $t=x+1$，则 $t\in[1,2]$，并且
    \begin{align*}
    V_1=\int_{0}^{1}\mathrm{d}x\int_{0}^{\frac{1-x}{1+x}}(1-x-y-xy)\,\mathrm{d}y
    &=\int_{0}^{1}\frac{(1-x)^2(1+3x)}{2(1+x)^2}\,\mathrm{d}x,\\
    \frac{(1-x)^2(1+3x)}{2(1+x)^2}
    &=\frac{(2-t)^2(3t-2)}{2t^2}
    =\frac{3}{2}t-7+\frac{10}{t}-\frac{4}{t^2},
    \end{align*}
    从而 \(\displaystyle V_1=\left[\frac{3}{4}t^2-7t+10\ln t+\frac{4}{t}\right]_{1}^{2}
    =10\ln 2-\frac{27}{4}\)。

    \item \textbf{解答小题 (2) —— 椭圆抛物面间的封闭域}：下界为 $z_1=x^2+2y^2$，上界为 $z_2=6-2x^2-y^2$，交线满足 \(\displaystyle x^2+2y^2=6-2x^2-y^2 \iff x^2+y^2=2\)，
    故投影域 $D$ 是半径为 $\sqrt{2}$ 的圆域。于是
    \(\displaystyle V_2=\iint_{D}(z_2-z_1)\,\mathrm{d}x\,\mathrm{d}y
    =\iint_{D}(6-3x^2-3y^2)\,\mathrm{d}x\,\mathrm{d}y
    =3\int_{0}^{2\pi}\mathrm{d}\theta\int_{0}^{\sqrt{2}}(2-r^2)r\,\mathrm{d}r
    =6\pi\)。
\end{compactenum}
\end{solution}

\begin{exercise}[利用极坐标与对称性求质心]
求位于两圆 $r=2\sin\theta$ 和 $r=4\sin\theta$ 之间的均匀薄片（密度 $\rho\equiv 1$）的质心。
\end{exercise}
\begin{solution}
\begin{compactenum}
    \item 曲线 $r=2\sin\theta$、$r=4\sin\theta$ 分别表示圆心在 $(0,1)$、$(0,2)$，半径为 $1$、$2$ 的两个圆，薄片夹在两圆之间，并关于 $y$ 轴对称，故 $\bar{x}=0$。
    \item 面积与对 $x$ 轴的静力矩可直接由“外圆减内圆”得到。半径为 $R$、圆心为 $(0,R)$ 的均匀圆面积为 $\pi R^2$，其质心在 $(0,R)$，故 \(\displaystyle A=\pi(2^2)-\pi(1^2)=3\pi,\ M_x=\pi\cdot 2^3-\pi\cdot 1^3=7\pi\)。
    因而 \(\displaystyle \bar{y}=\frac{M_x}{A}=\frac{7\pi}{3\pi}=\frac{7}{3}\)，
    所以质心坐标为 $\left(0,\frac{7}{3}\right)$。
\end{compactenum}
\end{solution}

\begin{exercise}[转动惯量的定积分公式推导]
证明：由 $x=a, x=b, y=f(x)$ 及 $x$ 轴所围的平面图形绕 $x$ 轴旋转一周所成的旋转体（设密度 $\rho\equiv 1$），对 $x$ 轴的转动惯量公式为 $I_x = \frac{\pi}{2}\int_a^b [f(x)]^4 \,\mathrm{d}x$。
\end{exercise}
\begin{solution}
\begin{compactenum}
    \item 设旋转体为 $\Omega$。对固定的 $x\in[a,b]$，其截面 $D_x$ 是位于与 $Oyz$ 平面平行的平面内、以 $x$ 轴为中心、半径为 $f(x)$ 的圆盘，因此
    \(\displaystyle I_x=\iiint_{\Omega}(y^2+z^2)\,\mathrm{d}V
    =\int_{a}^{b}\mathrm{d}x\iint_{D_x}(y^2+z^2)\,\mathrm{d}y\,\mathrm{d}z\)。
    \item 在圆盘 $D_x$ 上取极坐标 $y=r\cos\theta,\ z=r\sin\theta$，则 $y^2+z^2=r^2$，面积元为 $r\,\mathrm{d}r\,\mathrm{d}\theta$，故
    \(\displaystyle \iint\limits_{D_x}(y^2+z^2)\,\mathrm{d}y\,\mathrm{d}z
    =\int_{0}^{2\pi}\mathrm{d}\theta\int_{0}^{f(x)}r^3\,\mathrm{d}r
    =\frac{\pi}{2}[f(x)]^4\)。
    \item 代回即得 \(\displaystyle I_x=\int_{a}^{b}\frac{\pi}{2}[f(x)]^4\,\mathrm{d}x =\frac{\pi}{2}\int_{a}^{b}[f(x)]^4\,\mathrm{d}x\)。
\end{compactenum}
\end{solution}

\begin{exercise}[平面引力模型的极坐标化简]
求面密度为常量（设为 $\rho$），半径为 $R$ 的均匀圆形薄片 $x^2+y^2 \leqslant R^2, z=0$，对位于 $z$ 轴上点 $M(0,0,a) \ (a>0)$ 处单位质量质点的引力向量。
\end{exercise}
\begin{solution}
\begin{compactenum}
    \item 由对称性知水平分力抵消，引力只沿 $z$ 轴并指向薄片。若圆盘上微元的坐标为 $(x,y,0)$，则 \(\displaystyle F_z=\iint\limits_{D}K\rho\frac{-a}{(x^2+y^2+a^2)^{3/2}}\,\mathrm{d}\sigma=-K\rho a\iint\limits_{D}(x^2+y^2+a^2)^{-3/2}\,\mathrm{d}\sigma\)。
    \item 在极坐标下 $D:\ 0\leqslant r\leqslant R,\ 0\leqslant \theta\leqslant 2\pi$，于是
    \begin{align*}
    F_z=-K\rho a\int_{0}^{2\pi}\mathrm{d}\theta\int_{0}^{R}(r^2+a^2)^{-3/2}r\,\mathrm{d}r
    &=-2\pi K\rho a\int_{0}^{R}r(r^2+a^2)^{-3/2}\,\mathrm{d}r,\\
    \int r(r^2+a^2)^{-3/2}\,\mathrm{d}r
    &=\frac{1}{2}\int u^{-3/2}\,\mathrm{d}u
    =-\frac{1}{\sqrt{r^2+a^2}},\\
    F_z=-2\pi K\rho a\left[-\frac{1}{\sqrt{r^2+a^2}}\right]_{0}^{R}
    &=-2\pi K\rho\left(1-\frac{a}{\sqrt{R^2+a^2}}\right).
    \end{align*}
    其中第二行令 $u=r^2+a^2$。
    \item 因而引力向量为
    \(\displaystyle \vec{F}=\left(0,0,-2\pi K\rho\left(1-\frac{a}{\sqrt{R^2+a^2}}\right)\right)\)。
\end{compactenum}
\end{solution}

\begin{exercise}[立体引力模型计算]
设有一顶角为 $90^\circ$、高为 $1$ 的均匀正圆锥体（设密度为 $1$），求此正圆锥体与位于其顶点处的质量为 $1$ 的质点之间的引力。
\end{exercise}
\begin{solution}
\begin{compactenum}
    \item 取锥体顶点为原点，对称轴沿 $z$ 轴正向。顶角为 $90^\circ$，故半顶角为 $\alpha=\frac{\pi}{4}$；高为 $1$，故顶面为 $z=1$。由对称性只需计算 \(\displaystyle F_z=\iiint_{\Omega}K\frac{z}{(x^2+y^2+z^2)^{3/2}}\,\mathrm{d}V\)。
    \item 在球坐标下有 $0\leqslant \theta\leqslant 2\pi$，$0\leqslant \varphi\leqslant \frac{\pi}{4}$，而平面 $z=1$ 给出上界 $r\cos\varphi=1$，即 $r=\sec\varphi$。于是
    \[
    F_z=\int_{0}^{2\pi}\mathrm{d}\theta\int_{0}^{\pi/4}\mathrm{d}\varphi\int_{0}^{\sec\varphi}
    K\frac{r\cos\varphi}{r^3}r^2\sin\varphi\,\mathrm{d}r
    =2\pi K\int_{0}^{\pi/4}\cos\varphi\sin\varphi\left[\int_{0}^{\sec\varphi}1\,\mathrm{d}r\right]\mathrm{d}\varphi.
    \]
    \item 因为内层积分等于 $\sec\varphi$，故
    \(\displaystyle F_z=2\pi K\int_{0}^{\pi/4}\sin\varphi\,\mathrm{d}\varphi
    =2\pi K\left[-\cos\varphi\right]_{0}^{\pi/4}
    =\pi K(2-\sqrt{2})\)。
    因而锥体对其顶点处质点的引力沿中心轴指向内部，大小为 $\pi K(2-\sqrt{2})$。
\end{compactenum}
\end{solution}


\subsection{本章总习题}

\begin{exercise}[填空题综合考查]
填空题：
\par\noindent
\begin{tabular}{@{}p{0.48\linewidth}p{0.48\linewidth}@{}}
(1) 积分 $\displaystyle \int_{0}^{2}\mathrm{d}x\int_{x}^{2}\mathrm{e}^{-y^{2}}\mathrm{d}y$ 的值等于 \underline{\qquad}. &
(2) \makecell[l]{交换积分的顺序：\\$\displaystyle \int_{1}^{\mathrm{e}}\mathrm{d}x\int_{0}^{\ln x}f(x,y)\mathrm{d}y=$ \underline{\qquad}.} \\
(3) \makecell[l]{交换积分的顺序：\\$\displaystyle \int_{0}^{1}\mathrm{d}y$\\$\displaystyle \int_{-y}^{\sqrt{2y-y^{2}}}f(x,y)\mathrm{d}x = \underline{\qquad}$}. &
(4) 由 $Oxy$ 平面及曲面 $z=x^{2}+y^{2}, x^{2}+y^{2}=4$ 围成立体的体积 $V=$ \underline{\qquad}. \\
\end{tabular}
\par
\end{exercise}
\begin{solution}
\begin{compactenum}
    \item \textbf{解答小题 (1) —— 交换积分次序}：原积分中 $\mathrm{e}^{-y^2}$ 没有初等原函数，故先换序。区域 $D$ 由 $0\leqslant x\leqslant 2,\ x\leqslant y\leqslant 2$ 给出，是顶点为 $(0,0),(0,2),(2,2)$ 的三角形；改写成 $y$-型区域即 $0\leqslant y\leqslant 2,\ 0\leqslant x\leqslant y$，于是
        \[
        I_1=\int_{0}^{2}\mathrm{d}y\int_{0}^{y}\mathrm{e}^{-y^2}\,\mathrm{d}x
        =\int_{0}^{2}y\mathrm{e}^{-y^2}\,\mathrm{d}y
        =-\frac{1}{2}\int_{0}^{2}\mathrm{e}^{-y^2}\,\mathrm{d}(-y^2)
        =\frac{1}{2}(1-\mathrm{e}^{-4}).
        \]

    \item \textbf{解答小题 (2) —— 积分限的反解}：原区域由 $1\leqslant x\leqslant \mathrm{e},\ 0\leqslant y\leqslant \ln x$ 给出。换序后先看 $y$ 的范围：$0\leqslant y\leqslant 1$；再由 $y=\ln x$ 反解得左边界 $x=\mathrm{e}^y$，右边界仍为 $x=\mathrm{e}$，所以 \(\displaystyle \int_{1}^{\mathrm{e}}\mathrm{d}x\int_{0}^{\ln x}f(x,y)\,\mathrm{d}y=\int_{0}^{1}\mathrm{d}y\int_{\mathrm{e}^y}^{\mathrm{e}}f(x,y)\,\mathrm{d}x\)。

    \item \textbf{解答小题 (3) —— 复杂边界的分段换序}：原区域满足 $0\leqslant y\leqslant 1,\ -y\leqslant x\leqslant \sqrt{2y-y^2}$。右边界曲线可化为 \(\displaystyle x^2=2y-y^2 \iff x^2+(y-1)^2=1\)，即圆心 $(0,1)$、半径 $1$ 的圆右半部分。整体投影为 $-1\leqslant x\leqslant 1$，且须在 $x=0$ 处分段：当 $-1\leqslant x\leqslant 0$ 时，$-x\leqslant y\leqslant 1$；当 $0\leqslant x\leqslant 1$ 时，$1-\sqrt{1-x^2}\leqslant y\leqslant 1$。故
        \[
        \int_{0}^{1}\mathrm{d}y\int_{-y}^{\sqrt{2y-y^2}}f(x,y)\,\mathrm{d}x
        =\int_{-1}^{0}\mathrm{d}x\int_{-x}^{1}f(x,y)\,\mathrm{d}y
        +\int_{0}^{1}\mathrm{d}x\int_{1-\sqrt{1-x^2}}^{1}f(x,y)\,\mathrm{d}y.
        \]

    \item \textbf{解答小题 (4) —— 二重积分求体积}：底面是圆域 $x^2+y^2\leqslant 4$，顶面为 $z=x^2+y^2$，故 \(\displaystyle V=\iint\limits_{D}(x^2+y^2)\,\mathrm{d}\sigma=\int_{0}^{2\pi}\mathrm{d}\theta\int_{0}^{2}r^3\,\mathrm{d}r=8\pi\)。
\end{compactenum}
\end{solution}

\begin{exercise}[选择题综合概念与性质判定]
选择题(只有一个答案是正确的)：
\begin{compactenum}
    \item 设 $D$ 是 $Oxy$ 平面上以 $(1,1),(-1,1)$ 和 $(-1,-1)$ 为顶点的三角形区域， $D_1$ 为第一象限的部分，则 $\displaystyle \iint\limits_D (xy+\cos x\sin y)\,\mathrm{d}x\,\mathrm{d}y$ 等于 (\quad) 。\\
\begin{tabular}{ll}
    (A) $\displaystyle 2\iint\limits_{D_1}\cos x\sin y\,\mathrm{d}x\,\mathrm{d}y$ & (B) $\displaystyle 2\iint\limits_{D_1}xy\,\mathrm{d}x\,\mathrm{d}y$ \\[10pt]
    (C) $\displaystyle 4\iint\limits_{D_1}(xy+\cos x\sin y)\,\mathrm{d}x\,\mathrm{d}y$ & (D) $0$
\end{tabular}

    \item 设有空间区域 $\Omega_1: x^2+y^2+z^2 \leqslant R^2, z \geqslant 0$； $\Omega_2: x^2+y^2+z^2 \leqslant R^2, x \geqslant 0, y \geqslant 0, z \geqslant 0$。则 (\quad) 。\\
\begin{tabular}{ll}
  (A) $\displaystyle \iiint\limits_{\Omega_1}x\,\mathrm{d}V = 4\iiint\limits_{\Omega_2}x\,\mathrm{d}V$
  & (B) $\displaystyle \iiint\limits_{\Omega_1}y\,\mathrm{d}V = 4\iiint\limits_{\Omega_2}y\,\mathrm{d}V$ \\[12pt]
  (C) $\displaystyle \iiint\limits_{\Omega_1}z\,\mathrm{d}V = 4\iiint\limits_{\Omega_2}z\,\mathrm{d}V$
  & (D) $\displaystyle \iiint\limits_{\Omega_1}xyz\,\mathrm{d}V = 4\iiint\limits_{\Omega_2}xyz\,\mathrm{d}V$
\end{tabular}

    \item 设 $f(x,y)$ 是有界闭区域 $D: x^2+y^2 \leqslant a^2$ 上的连续函数，则当 $a \to 0$ 时， $\displaystyle \frac{1}{\pi a^2}\iint\limits_D f(x,y)\,\mathrm{d}x\,\mathrm{d}y$ 的极限 (\quad) 。\\
    (A) 不存在 \hfill (B) 等于 $f(0,0)$ \hfill (C) 等于 $f(1,1)$ \hfill (D) 等于 $f(1,0)$

    \item 设区域 $D: 0 \leqslant y \leqslant x, 0 \leqslant x \leqslant \pi$，则二重积分 $\displaystyle \iint\limits_D \sqrt{1-\sin^2 x}\,\mathrm{d}x\,\mathrm{d}y =$ (\quad) 。\\
    (A) $0$ \hfill (B) $\frac{\pi}{2}$ \hfill (C) $2$ \hfill (D) $\pi$
\end{compactenum}
\end{exercise}
\begin{solution}
\begin{compactenum}
    \item \textbf{解答小题 (1)}：正确答案选 \textbf{(A)}。
    将积分拆成 \(\displaystyle \iint_Dxy\,\mathrm{d}\sigma+\iint_D\cos x\sin y\,\mathrm{d}\sigma\)，
    其中 $D=\{(x,y):-1\leqslant x\leqslant 1,\ x\leqslant y\leqslant 1\}$。第一项化为 \(\displaystyle \int_{-1}^{1}x\,\mathrm{d}x\int_{x}^{1}y\,\mathrm{d}y
    =\int_{-1}^{1}\frac{x-x^3}{2}\,\mathrm{d}x=0\)，因为被积函数是奇函数。第二项化为
    \(\displaystyle \int_{-1}^{1}\cos x\,\mathrm{d}x\int_{x}^{1}\sin y\,\mathrm{d}y
    =\int_{-1}^{1}\cos x(\cos x-\cos 1)\,\mathrm{d}x\)。
    这里 $g(x)=\cos x(\cos x-\cos 1)$ 为偶函数，所以该积分等于
    \(\displaystyle 2\int_{0}^{1}\cos x(\cos x-\cos 1)\,\mathrm{d}x
    =2\iint\limits_{D_1}\cos x\sin y\,\mathrm{d}x\,\mathrm{d}y\)。

    \item \textbf{解答小题 (2)}：正确答案选 \textbf{(C)}。
    $\Omega_1$ 是上半球，$\Omega_2$ 是第一卦限中的八分之一球。对 $x,y,xyz$ 而言，它们在 $\Omega_1$ 内都存在正负对消，所以
    \[
    \iiint\limits_{\Omega_1}x\,\mathrm{d}V
    =\iiint\limits_{\Omega_1}y\,\mathrm{d}V
    =\iiint\limits_{\Omega_1}xyz\,\mathrm{d}V=0,
    \]
    因而 (A)、(B)、(D) 都不成立；只有 $z\geqslant 0$ 在四个卦限中的分布完全相同，故
    \(\displaystyle \iiint\limits_{\Omega_1}z\,\mathrm{d}V=4\iiint\limits_{\Omega_2}z\,\mathrm{d}V\)。

    \item \textbf{解答小题 (3)}：正确答案选 \textbf{(B)}。
    这里 $\pi a^2$ 正是圆域 $D$ 的面积 $\sigma$，故
    \[
    \frac{1}{\pi a^2}\iint\limits_D f(x,y)\,\mathrm{d}x\,\mathrm{d}y
    =\frac{1}{\sigma}\iint\limits_D f(x,y)\,\mathrm{d}\sigma
    \]
    是 $f$ 在 $D$ 上的平均值。由积分中值定理，存在 $(\xi,\eta)\in D$ 使其等于 $f(\xi,\eta)$；当 $a\to 0$ 时，$D$ 收缩到原点，于是 $(\xi,\eta)\to(0,0)$，再由连续性得极限为 $f(0,0)$。

    \item \textbf{解答小题 (4)}：正确答案选 \textbf{(D)}。
    因为 $\sqrt{1-\sin^2x}=|\cos x|$，故
    \[
    \iint\limits_D\sqrt{1-\sin^2x}\,\mathrm{d}x\,\mathrm{d}y
    =\int_{0}^{\pi}\mathrm{d}x\int_{0}^{x}|\cos x|\,\mathrm{d}y
    =\int_{0}^{\pi}x|\cos x|\,\mathrm{d}x.
    \]
    在 $\frac{\pi}{2}$ 处分段，并用分部积分公式 $\int x\cos x\,\mathrm{d}x=x\sin x+\cos x$，于是
    \begin{align*}
    \int_{0}^{\pi}x|\cos x|\,\mathrm{d}x
    &=\int_{0}^{\pi/2}x\cos x\,\mathrm{d}x-\int_{\pi/2}^{\pi}x\cos x\,\mathrm{d}x\\
    &=[x\sin x+\cos x]_{0}^{\pi/2}-[x\sin x+\cos x]_{\pi/2}^{\pi}
    =\left(\frac{\pi}{2}-1\right)+\left(1+\frac{\pi}{2}\right)=\pi.
    \end{align*}
\end{compactenum}
\end{solution}

\begin{exercise}[多重积分的对称性分析]
设 $\Omega:x^2+y^2+z^2 \leqslant a^2$（全球）；$\Omega_1:x^2+y^2+z^2 \leqslant a^2, z \geqslant 0$（上半球）；$\Omega_2:x^2+y^2+z^2 \leqslant a^2, x \geqslant 0, y \geqslant 0, z \geqslant 0$（第一卦限部分球）。\\
问下列等式是否成立，请详细说明理由。
\par\noindent
\begin{tabular}{@{}p{0.06\linewidth}p{0.88\linewidth}@{}}
(1) & \(\displaystyle \iiint\limits_{\Omega}x\,\mathrm{d}V=0, \quad \iiint\limits_{\Omega}z\,\mathrm{d}V=0\)；\\
(2) & \(\displaystyle \iiint\limits_{\Omega_1}x\,\mathrm{d}V=4\iiint\limits_{\Omega_2}x\,\mathrm{d}V, \quad \iiint\limits_{\Omega_1}z\,\mathrm{d}V=4\iiint\limits_{\Omega_2}z\,\mathrm{d}V\)；\\
(3) & \(\displaystyle \iiint\limits_{\Omega_1}xy\,\mathrm{d}V=\iiint\limits_{\Omega_1}yz\,\mathrm{d}V=\iiint\limits_{\Omega_1}zx\,\mathrm{d}V=0\)。
\end{tabular}
\par
\end{exercise}
\begin{solution}
\begin{compactenum}
    \item \textbf{解答小题 (1) —— 成立}。全球区域 $\Omega$ 关于 $yOz$ 平面对称，而被积函数 $x$ 满足 $f(-x,y,z)=-f(x,y,z)$，故
    \(\displaystyle \iiint\limits_{\Omega}x\,\mathrm{d}V=0\)。同理，$\Omega$ 关于 $xOy$ 平面对称，函数 $z$ 关于 $z$ 是奇的，因此
    \(\displaystyle \iiint\limits_{\Omega}z\,\mathrm{d}V=0\)。

    \item \textbf{解答小题 (2) —— 前错后对}。上半球 $\Omega_1$ 关于 $yOz$ 平面仍对称，所以
    \(\displaystyle \iiint\limits_{\Omega_1}x\,\mathrm{d}V=0\)；
    但在第一卦限的 $\Omega_2$ 内 $x\geqslant 0$ 且不恒为零，所以 $\iiint_{\Omega_2}x\,\mathrm{d}V>0$，前一个等式不成立。至于 $z$，它在 $\Omega_1$ 的四个卦限子区域中分布完全相同，而 $\Omega_2$ 恰是其中之一，因此 \(\displaystyle \iiint\limits_{\Omega_1}z\,\mathrm{d}V=4\iiint\limits_{\Omega_2}z\,\mathrm{d}V\) 成立。

    \item \textbf{解答小题 (3) —— 成立}。虽然 $\Omega_1$ 只限制了 $z\geqslant 0$，但它关于 $x=0$ 和 $y=0$ 仍然对称，因此 $xy$ 与 $zx$ 关于 $x$ 为奇函数，$yz$ 关于 $y$ 为奇函数，于是 \(\displaystyle \iiint\limits_{\Omega_1}xy\,\mathrm{d}V=\iiint\limits_{\Omega_1}yz\,\mathrm{d}V=\iiint\limits_{\Omega_1}zx\,\mathrm{d}V=0\)。
\end{compactenum}
\end{solution}

\begin{exercise}[极坐标求二重积分]
求二重积分 $\displaystyle I=\iint\limits_D \frac{x}{x^2+y^2}\,\mathrm{d}x\,\mathrm{d}y$，其中 $D$ 由曲线 $y=\frac{x^2}{2}$ 与直线 $y=x$ 围成。
\end{exercise}
\begin{solution}
\begin{compactenum}
    \item 联立 $y=x$ 与 $y=\frac{x^2}{2}$，得交点为 $(0,0)$ 和 $(2,2)$；在此区间内抛物线位于直线下方，区域处于第一象限，适合改用极坐标。
    \item 直线 $y=x$ 对应 $\theta=\frac{\pi}{4}$，抛物线化为 \(\displaystyle r\sin\theta=\frac{1}{2}r^2\cos^2\theta \iff r=\frac{2\sin\theta}{\cos^2\theta}\)，
    因而边界为 $0\leqslant \theta\leqslant \frac{\pi}{4}$，$0\leqslant r\leqslant \frac{2\sin\theta}{\cos^2\theta}$。
    \item 又 \(\displaystyle \frac{x}{x^2+y^2}=\frac{r\cos\theta}{r^2}=\frac{\cos\theta}{r}\)，
    与面积元 $r\,\mathrm{d}r\,\mathrm{d}\theta$ 相乘后恰好约去 $r$，也就是 \(\displaystyle \frac{x}{x^2+y^2}\,r\,\mathrm{d}r\,\mathrm{d}\theta=\cos\theta\,\mathrm{d}r\,\mathrm{d}\theta\)，故 \(\displaystyle I=\int_0^{\pi/4}\mathrm{d}\theta\int_0^{\frac{2\sin\theta}{\cos^2\theta}}\cos\theta\,\mathrm{d}r=\int_0^{\pi/4}\frac{2\sin\theta}{\cos\theta}\,\mathrm{d}\theta=\ln2\)。
\end{compactenum}
\end{solution}

\begin{exercise}[换序拯救无法计算的反常积分]
计算：
\[ I = \int_{1}^{2}\mathrm{d}x\int_{\sqrt{x}}^{x}\sin\frac{\pi x}{2y}\,\mathrm{d}y + \int_{2}^{4}\mathrm{d}x\int_{\sqrt{x}}^{2}\sin\frac{\pi x}{2y}\,\mathrm{d}y \]
\end{exercise}
\begin{solution}
\begin{compactenum}
    \item 两段积分在 $xOy$ 平面上无缝拼成区域 $D$，其上下边界分别为 $y=x$ 与 $y=\sqrt{x}$。改写成 $y$-型区域后，$1\leqslant y\leqslant 2$，且对每个 $y$ 有 $y\leqslant x\leqslant y^2$，于是
    \(\displaystyle I=\int_{1}^{2}\mathrm{d}y\int_{y}^{y^2}\sin\left(\frac{\pi x}{2y}\right)\mathrm{d}x\)。
    \item 把 $y$ 看作常数，令 $k=\frac{\pi}{2y}$，则
    \[
    \int_{y}^{y^2}\sin(kx)\,\mathrm{d}x
    =\left[-\frac{1}{k}\cos(kx)\right]_{y}^{y^2}
    =-\frac{2y}{\pi}\cos\left(\frac{\pi y}{2}\right),
    \]
    因而 \(\displaystyle I=-\frac{2}{\pi}\int_{1}^{2}y\cos\left(\frac{\pi y}{2}\right)\mathrm{d}y\)。
    \item 再用分部积分，得 \(\displaystyle I=\frac{8}{\pi^3}+\frac{4}{\pi^2}=\frac{4(2+\pi)}{\pi^3}\)。
    \textit{（重要纠误指出：原书参考答案在此处给出的为 $\frac{4}{\pi^2}(2+\pi)$，实则在推导过程中遗漏了分母的 $\pi$ 因子，发生阶次错误，正确结果确系分母带有 $\pi^3$ 的式子。）}
\end{compactenum}
\end{solution}

\begin{exercise}[含参极坐标的代数换元]
求 $\displaystyle I = \iint\limits_D \frac{1-x^2-y^2}{1+x^2+y^2}\,\mathrm{d}x\,\mathrm{d}y$，其中 $D$ 是区域 $x^2+y^2 \leqslant 1, x \geqslant 0, y \geqslant 0$。
\end{exercise}
\begin{solution}
\begin{compactenum}
    \item 区域 $D$ 是第一象限的单位四分之一圆域，因此 $0\leqslant \theta\leqslant \frac{\pi}{2}$，$0\leqslant r\leqslant 1$，积分化为
    \[
    I=\int_{0}^{\pi/2}\mathrm{d}\theta\int_{0}^{1}\frac{1-r^2}{1+r^2}r\,\mathrm{d}r
    =\frac{\pi}{2}\int_{0}^{1}\frac{1-r^2}{1+r^2}r\,\mathrm{d}r.
    \]
    \item 令 $u=r^2$，则 $r\,\mathrm{d}r=\frac{1}{2}\mathrm{d}u$，于是
    \[
    I=\frac{\pi}{4}\int_{0}^{1}\frac{1-u}{1+u}\,\mathrm{d}u
    =\frac{\pi}{4}\int_{0}^{1}\left(\frac{2}{1+u}-1\right)\mathrm{d}u
    =\frac{\pi}{4}\left[2\ln(1+u)-u\right]_{0}^{1}
    =\frac{\pi}{2}\ln 2-\frac{\pi}{4}.
    \]
\end{compactenum}
\end{solution}

\begin{exercise}[利用雅可比行列式的指数代换]
设 $D$ 是曲线 $\sqrt{\frac{x}{a}}+\sqrt{\frac{y}{b}}=1$ 与 $x$ 轴、$y$ 轴所围成的区域（$a>0,b>0$）。计算二重积分 \(\displaystyle I = \iint\limits_D y \,\mathrm{d}x\,\mathrm{d}y\)。
\end{exercise}
\begin{solution}
\begin{compactenum}
    \item 令
    \(\displaystyle u=\sqrt{\frac{x}{a}},\qquad v=\sqrt{\frac{y}{b}}\)，
    则 $x=au^2,\ y=bv^2$，并且原边界 $\sqrt{x/a}+\sqrt{y/b}=1$ 恰化为直线 $u+v=1$，所以新区域为
    \(\displaystyle D'=\{(u,v):u\geqslant 0,\ v\geqslant 0,\ u+v\leqslant 1\}\)。
    \item 雅可比行列式为
    \[
    J=\frac{\partial(x,y)}{\partial(u,v)}
    =\det\begin{pmatrix}2au&0\\0&2bv\end{pmatrix}
    =4abuv,
    \]
    从而 \(\displaystyle I=\iint\limits_{D'}bv^2\cdot 4abuv\,\mathrm{d}u\,\mathrm{d}v=4ab^2\iint\limits_{D'}uv^3\,\mathrm{d}u\,\mathrm{d}v\)。
    \item 化为累次积分后，\(\displaystyle I=4ab^2\int_{0}^{1}\mathrm{d}u\int_{0}^{1-u}uv^3\,\mathrm{d}v=ab^2\int_{0}^{1}u(1-u)^4\,\mathrm{d}u\)。再令 $t=1-u$，便得 \(\displaystyle \int_{0}^{1}u(1-u)^4\,\mathrm{d}u=\int_0^1(t^4-t^5)\,\mathrm{d}t=\frac{1}{30}\)，
    故 \(\displaystyle I=\frac{ab^2}{30}\)。
\end{compactenum}
\end{solution}

\begin{exercise}[利用区域质心对称性免算积分]
求 $\displaystyle I = \iint\limits_D (x+y) \,\mathrm{d}x\,\mathrm{d}y$，其中 $D: x^2+y^2 \leqslant x+y+1$。
\end{exercise}
\begin{solution}
\begin{compactenum}
    \item 配方得 \(\displaystyle x^2-x+y^2-y\leqslant1\iff \left(x-\frac{1}{2}\right)^2+\left(y-\frac{1}{2}\right)^2\leqslant\frac{3}{2}\)。
    所以 $D$ 是圆心 $\left(\frac{1}{2},\frac{1}{2}\right)$、半径 $\sqrt{\frac{3}{2}}$ 的圆域。
    \item 因而面积 \(\displaystyle A=\pi\cdot \frac{3}{2}=\frac{3\pi}{2}\)，且圆域的形心就是其圆心，所以 $\bar{x}=\bar{y}=\frac{1}{2}$。
    \item 由质心公式，\(\displaystyle I=\iint\limits_Dx\,\mathrm{d}\sigma+\iint\limits_Dy\,\mathrm{d}\sigma=(\bar{x}+\bar{y})A=\frac{3\pi}{2}\)。
\end{compactenum}
\end{solution}

\begin{exercise}[灵活规划三重积分的三维切片路径]
设 $\Omega$ 由 $y=\sqrt{x}, y=0, z=0, x+z=\frac{\pi}{2}$ 围成，计算三重积分
\(\displaystyle I = \iiint\limits_{\Omega} \frac{y\sin x}{x} \,\mathrm{d}x\,\mathrm{d}y\,\mathrm{d}z\)。
\end{exercise}
\begin{solution}
\begin{compactenum}
    \item 由边界 $z=0$ 与 $x+z=\frac{\pi}{2}$ 可知最内层宜取 $z$，且 $0\leqslant z\leqslant \frac{\pi}{2}-x$；投影到 $Oxy$ 面后，再由 $y=\sqrt{x}$ 与 $y=0$ 得 $0\leqslant x\leqslant \frac{\pi}{2}$，$0\leqslant y\leqslant \sqrt{x}$。因此
    \(\displaystyle I=\int_{0}^{\pi/2}\mathrm{d}x\int_{0}^{\sqrt{x}}\mathrm{d}y\int_{0}^{\frac{\pi}{2}-x}\frac{y\sin x}{x}\,\mathrm{d}z\)。
    \item 先对 $z$ 积分得到因子 \(\displaystyle \frac{\pi}{2}-x\)，再对 $y$ 积分用 \(\displaystyle \int_0^{\sqrt{x}}y\,\mathrm{d}y=\frac{x}{2}\)，故
    \(\displaystyle I=\frac{1}{2}\int_{0}^{\pi/2}\sin x\left(\frac{\pi}{2}-x\right)\mathrm{d}x\)。
    \item 即 \(\displaystyle I=\frac{1}{2}\int_{0}^{\pi/2}\left(\frac{\pi}{2}\sin x-x\sin x\right)\mathrm{d}x\)。其中
    \(\displaystyle \frac{\pi}{2}\int_{0}^{\pi/2}\sin x\,\mathrm{d}x=\frac{\pi}{2},\quad
    \int_{0}^{\pi/2}x\sin x\,\mathrm{d}x=1\)，
    所以 \(\displaystyle I=\frac{1}{2}\left(\frac{\pi}{2}-1\right)=\frac{\pi}{4}-\frac{1}{2}\)。
\end{compactenum}
\end{solution}

\begin{exercise}[柱与球坐标系统下的繁杂边界重塑]
将三次累次积分：
\[ I = \int_{0}^{1}\mathrm{d}y\int_{-\sqrt{y-y^2}}^{\sqrt{y-y^2}}\mathrm{d}x\int_{0}^{\sqrt{3(x^2+y^2)}} f\left(\sqrt{x^2+y^2+z^2}\right) \mathrm{d}z \]
变换为完整的柱坐标及球坐标形式的累次积分。
\end{exercise}
\begin{solution}
\begin{compactenum}
    \item 由外中两层积分限知，投影域满足 \(\displaystyle y\in[0,1],\ x^2\leqslant y-y^2\iff x^2+\left(y-\frac{1}{2}\right)^2\leqslant \frac{1}{4}\)。
    这是与原点相切、位于第一二象限的圆；在极坐标下即 $r\leqslant \sin\theta$，并且 $0\leqslant \theta\leqslant \pi$。内层积分又给出 $0\leqslant z\leqslant \sqrt{3}r$。
    \item 所以柱坐标形式为
    \[
    I_{\text{柱}}
    =\int_{0}^{\pi}\mathrm{d}\theta\int_{0}^{\sin\theta}r\,\mathrm{d}r\int_{0}^{\sqrt{3}r}f\left(\sqrt{r^2+z^2}\right)\mathrm{d}z.
    \]
    \item 若改用球坐标，则平面 $z=0$ 对应 $\varphi=\frac{\pi}{2}$，锥面 $z=\sqrt{3}r$ 对应 \(\displaystyle \rho\cos\varphi=\sqrt{3}\rho\sin\varphi\iff \tan\varphi=\frac{1}{\sqrt{3}}\iff \varphi=\frac{\pi}{6}\)，
    故 $\frac{\pi}{6}\leqslant \varphi\leqslant \frac{\pi}{2}$；方位角仍为 $0\leqslant \theta\leqslant \pi$，而边界 $r=\sin\theta$ 变为 $\rho\sin\varphi=\sin\theta$，即 $\rho\leqslant \frac{\sin\theta}{\sin\varphi}$。于是
    \[
    I_{\text{球}}
    =\int_{0}^{\pi}\mathrm{d}\theta\int_{\pi/6}^{\pi/2}\sin\varphi\,\mathrm{d}\varphi
    \int_{0}^{\frac{\sin\theta}{\sin\varphi}}f(\rho)\rho^2\,\mathrm{d}\rho.
    \]
    \textit{（勘误注：部分教材的参考答案在此处将天顶角下限错标为了 $\pi/3$。若为 $\pi/3$ 则对应的下锥面方程为 $z = \frac{1}{\sqrt{3}}r$，与题干矛盾。在此依照严密的几何对应法则更正为正确的界限 $\pi/6$。）}
\end{compactenum}
\end{solution}

\begin{exercise}[利用球坐标系处理截断旋转体体积]
计算积分 \(\displaystyle I = \iiint\limits_{\Omega} (x^2+y^2) \,\mathrm{d}V\)，其中 $\Omega$ 是由平面曲线 $y^2=2z, x=0$ 绕 $z$ 轴旋转一周形成的曲面与平面 $z=8$ 所共同围成的区域。
\end{exercise}
\begin{solution}
\begin{compactenum}
    \item 曲线 $y^2=2z$ 绕 $z$ 轴旋转后变为抛物面 \(\displaystyle x^2+y^2=2z\)，
    再与平面 $z=8$ 围成立体；其在 $Oxy$ 面上的投影满足 $x^2+y^2\leqslant 16$，故最适合用柱坐标。
    \item 在柱坐标下，$0\leqslant \theta\leqslant 2\pi$，$0\leqslant r\leqslant 4$，且对固定 $r$ 有 $\frac{r^2}{2}\leqslant z\leqslant 8$。被积函数 $x^2+y^2$ 化为 $r^2$，乘上体积元中的 $r$，得 \(\displaystyle I=2\pi\int_{0}^{4}r^3\left(8-\frac{r^2}{2}\right)\mathrm{d}r\)。继续计算，\(\displaystyle I=2\pi\int_{0}^{4}\left(8r^3-\frac{1}{2}r^5\right)\mathrm{d}r=2\pi\left[2r^4-\frac{1}{12}r^6\right]_{0}^{4}=\frac{1024\pi}{3}\)。
\end{compactenum}
\end{solution}

\begin{exercise}[分离变量化简三重球壳域积分]
求积分 $\displaystyle I = \iiint\limits_{\Omega} (x+z)\mathrm{e}^{-(x^2+y^2+z^2)}\,\mathrm{d}V$，其中 $\Omega: 1 \leqslant x^2+y^2+z^2 \leqslant 4, x \geqslant 0, y \geqslant 0, z \geqslant 0$。
\end{exercise}
\begin{solution}
\begin{compactenum}
    \item $\Omega$ 是第一卦限内的同心球壳，对 $x$ 与 $z$ 完全对称。因此
    \[
    \iiint\limits_{\Omega}x\mathrm{e}^{-(x^2+y^2+z^2)}\,\mathrm{d}V
    =
    \iiint\limits_{\Omega}z\mathrm{e}^{-(x^2+y^2+z^2)}\,\mathrm{d}V,
    \]
故 \(\displaystyle I=2\iiint\limits_{\Omega}z\mathrm{e}^{-(x^2+y^2+z^2)}\,\mathrm{d}V\)。
    \item 取球坐标，则 $0\leqslant \theta\leqslant \frac{\pi}{2}$，$0\leqslant \varphi\leqslant \frac{\pi}{2}$，$1\leqslant \rho\leqslant 2$，并有 $z=\rho\cos\varphi$，故
    \[
    I=2\int_{0}^{\pi/2}\mathrm{d}\theta\int_{0}^{\pi/2}\mathrm{d}\varphi\int_{1}^{2}
    \rho\cos\varphi\,\mathrm{e}^{-\rho^2}\rho^2\sin\varphi\,\mathrm{d}\rho.
    \]
    \item 三部分可分离，且 \(\displaystyle \int_{0}^{\pi/2}1\,\mathrm{d}\theta=\frac{\pi}{2},\qquad \int_{0}^{\pi/2}\sin\varphi\cos\varphi\,\mathrm{d}\varphi=\frac{1}{2}\)。
    对径向积分令 $u=\rho^2$，则
    \[
    \int_{1}^{2}\rho^3\mathrm{e}^{-\rho^2}\,\mathrm{d}\rho
    =\frac{1}{2}\int_{1}^{4}u\mathrm{e}^{-u}\,\mathrm{d}u
    =\frac{1}{2}\left(2\mathrm{e}^{-1}-5\mathrm{e}^{-4}\right).
    \]
    因而
    \[
    I=2\cdot \frac{\pi}{2}\cdot \frac{1}{2}\cdot \frac{1}{2}
    \left(2\mathrm{e}^{-1}-5\mathrm{e}^{-4}\right)
    =\frac{\pi(2\mathrm{e}^3-5)}{4\mathrm{e}^4}.
    \]
\end{compactenum}
\end{solution}

\chapter{曲线积分与曲面积分}

\section{知识讲解}

本章知识讲解围绕“先判对象，再判结构，最后计算”展开。曲线或曲面题中的积分元包括 $\mathrm{d}s,\mathrm{d}S$ 以及 $\mathrm{d}x,\mathrm{d}y,\mathrm{d}z$、$\mathrm{d}y\,\mathrm{d}z,\mathrm{d}z\,\mathrm{d}x,\mathrm{d}x\,\mathrm{d}y$；它们先告诉我们题目属于哪一类积分。方向、闭合性、边界、补边、封口、换面和奇点再决定是否使用保守场、Green、Stokes 或 Gauss。直接计算负责把微元落到普通积分，结构公式负责把对象换成更合适的内部积分或端点差，二者不是同一层级的“并列招式”。

\subsection{总路线：先识别对象，再看结构，再选工具}

做曲线积分和曲面积分时，第一反应不应是背某个公式，而应先把题目翻译成四层问题：
\[
\boxed{
\text{定义层语言}
\longrightarrow
\text{结构主工具}
\longrightarrow
\text{修正手段}
\longrightarrow
\text{最终计算}
}.
\]
定义层语言回答“我到底在积什么微元”；结构主工具回答“题目是否能由端点差、Green、Stokes 或 Gauss 接管”；补边、封口、换面、挖洞、分块等修正手段负责让结构公式合法可用；最后才把问题落到一元、二重或三重积分中。直接计算是定义层语言，也是某些题的主路，但不是万能保底；只有题面给法明确、参数化或投影后积分可控时，它才适合作为计算主路。

\subsubsection{识别对象：四类积分与对应微元}

题面里的微分元已经把积分对象说得很清楚。先看它是沿曲线累加，还是穿过曲面累加；是标量密度，还是带方向的向量场积分。

\begin{property}[四类积分的入口]
\begin{itemize}
    \item \(\displaystyle \int_\Gamma f\,\mathrm{d}s\) 是第一类曲线积分，只沿曲线累加标量密度，方向不改号。它通常以直接参数化和弧长元为主。
    \item \(\displaystyle \iint_\Sigma f\,\mathrm{d}S\) 是第一类曲面积分，只在曲面上累加标量密度，与取向无关。它通常以曲面参数化、投影面积元和对称性为主。
    \item \(\displaystyle \int_\Gamma P\,\mathrm{d}x+Q\,\mathrm{d}y+R\,\mathrm{d}z\) 是第二类曲线积分，也就是 \(\displaystyle \int_\Gamma \vec F\cdot\mathrm{d}\vec r\)。它沿曲线看切向作用，方向一反结果变号；开曲线先查保守场，闭曲线再看 Green 或 Stokes。
    \item \(\displaystyle \iint_\Sigma P\,\mathrm{d}y\,\mathrm{d}z+Q\,\mathrm{d}z\,\mathrm{d}x+R\,\mathrm{d}x\,\mathrm{d}y\) 是第二类曲面积分，也就是 \(\displaystyle \iint_\Sigma \vec F\cdot\mathrm{d}\vec S\)。它计算通量，必须先定法向；闭曲面先想 Gauss，开曲面先想投影公式或封口。
\end{itemize}
\end{property}

这里要分清两件事：第一类积分多数时候以直接计算为主，因为 \(\mathrm{d}s,\mathrm{d}S\) 本身就是标量累加；第二类积分不能只停在参数化层面。若题目出现开曲线端点、平面闭曲线、空间闭曲线、闭曲面、可封口曲面等信号，就已经在提示我们检查保守场、Green、Stokes 或 Gauss。直接计算仍然必须会，因为它解释微元、方向和符号；但在这些结构信号明确时，结构公式通常就是主路，而不是可有可无的花样。

\subsubsection{结构：方向、闭合、边界与奇点}

第二步看题目是否带有能“换对象”的结构。闭合性决定 Green、Stokes、Gauss 是否进入候选；方向决定结果是否变号；奇点决定这些公式能不能合法使用。

\begin{conclusion}[环量、通量、旋度、散度的对应关系]
\begin{compactenum}
\item \textbf{环量：沿着边走。}
若 \(\vec F=(P,Q,R)\)，则
\[
\oint_C \vec F\cdot\mathrm{d}\vec r
=\oint_C P\,\mathrm{d}x+Q\,\mathrm{d}y+R\,\mathrm{d}z.
\]
它看的是向量场在曲线切线方向上的累计作用。曲线反向，\(\mathrm{d}\vec r\) 变号，环量也变号。

\item \textbf{通量：穿过面。}
若曲面 \(\Sigma\) 已选定正法向 \(\vec n\)，则
\[
\iint_\Sigma \vec F\cdot\vec n\,\mathrm{d}S
=\iint_\Sigma
P\,\mathrm{d}y\,\mathrm{d}z
+Q\,\mathrm{d}z\,\mathrm{d}x
+R\,\mathrm{d}x\,\mathrm{d}y.
\]
它看的是向量场在法向上的累计穿过量。法向一反，通量就变号。

\item \textbf{旋度管环量。}
\[
\operatorname{rot}\vec F=\nabla\times\vec F.
\]
旋度描述局部打转趋势。二维场 \(\vec F=(P,Q)\) 的旋度密度就是 \(Q_x-P_y\)，所以 Green 公式和 Stokes 公式本质上都在处理环量。

\item \textbf{散度管通量。}
\[
\operatorname{div}\vec F=\nabla\cdot\vec F=P_x+Q_y+R_z.
\]
散度描述局部净流出趋势，所以 Gauss 公式本质上在处理通量。
\end{compactenum}
\end{conclusion}

\begin{conclusion}[\texorpdfstring{\(\nabla\) 总表：哪里用，就在哪里展开}{nabla 总表：哪里用，就在哪里展开}]
\(\nabla\) 只是把梯度、旋度、散度放进同一套速查语言中。此处只记公式入口，具体计算放回对应场景。

\par\noindent{\centering
\renewcommand{\arraystretch}{1.35}
\begin{tabular}{|c|c|c|}
\hline
\textbf{算子} & \textbf{读法} & \textbf{对应主线} \\
\hline
\(\nabla u=\operatorname{grad}u\) & 标量场 \(\to\) 向量场 & 保守场端点差：\(\int_A^B\nabla u\cdot\mathrm{d}\vec r=u(B)-u(A)\) \\
\hline
\(\nabla\times\vec F=\operatorname{rot}\vec F\) & 向量场 \(\to\) 向量场 & Green/Stokes 环量：闭曲线先算旋度 \\
\hline
\(\nabla\cdot\vec F=\operatorname{div}\vec F\) & 向量场 \(\to\) 标量场 & Gauss 通量：闭曲面先算散度 \\
\hline
\end{tabular}
\par}
\end{conclusion}

判断结构时记四件事：闭合、方向、补边/封口、奇点。Green、Stokes、Gauss 都要求相应区域内有足够好的连续可微性，不能把奇点直接吞进公式里。

\subsubsection{选工具：结构公式先于计算路线}

选工具不是按公式名字排队，更不是把直接计算和 Green、Stokes、Gauss 摆成平级选项，而是按题面结构触发。
\[
\boxed{\text{Green：平面闭曲线的环量 = 平面区域内的旋度总和}},
\]
\[
\boxed{\text{Stokes：空间闭曲线的环量 = 张成曲面上的旋度通量}},
\]
\[
\boxed{\text{Gauss：闭曲面的通量 = 所围体内的散度总和}}.
\]

\begin{property}[选法的落笔顺序]
\begin{enumerate}
    \item 先写清楚“这是哪一类积分”：是 \(\mathrm{d}s\)、\(\mathrm{d}S\)，还是 \(P\,\mathrm{d}x+Q\,\mathrm{d}y+R\,\mathrm{d}z\)，还是 \(P\,\mathrm{d}y\,\mathrm{d}z+Q\,\mathrm{d}z\,\mathrm{d}x+R\,\mathrm{d}x\,\mathrm{d}y\)。
    \item 再写清楚方向：曲线正向、曲面法向、外侧或内侧；方向不明时，先翻译题目的“从某轴看去”。
    \item 然后看触发条件：开曲线端点题查保守场，平面闭曲线查 Green，空间闭曲线查 Stokes，闭曲面查 Gauss；不闭时先判断补边、封口、换面或分块能否把题目送进这些公式。
    \item 最后比较合法性与计算量：结构公式合法且能降维或简化时优先使用；若公式暂时不合法，先考虑挖洞、补边、封口、换面、分块等修正手段。只有曲线或曲面给法明确，且参数化或投影积分确实可控时，才采用直接计算。
\end{enumerate}
\end{property}

\begin{conclusion}[什么时候必须认真检查结构公式]
\begin{itemize}
    \item \(P\,\mathrm{d}x+Q\,\mathrm{d}y\) 沿平面闭曲线积分时，必须先查 Green 公式；如果边界分段多、参数化麻烦，而 \(Q_x-P_y\) 很简单，直接逐段积分就是绕远。
    \item \(P\,\mathrm{d}x+Q\,\mathrm{d}y+R\,\mathrm{d}z\) 沿空间闭曲线积分时，必须先查 Stokes 公式；如果同边界能换成平面圆盘、三角形或其他简单曲面，通常应换面。
    \item \(P\,\mathrm{d}y\,\mathrm{d}z+Q\,\mathrm{d}z\,\mathrm{d}x+R\,\mathrm{d}x\,\mathrm{d}y\) 在闭曲面上积分时，必须先查 Gauss 公式；如果曲面未闭但只差一个简单补面，也应先考虑封口。
    \item 开曲线第二类积分若只给端点，或题目暗示“路径无关”，必须先查保守场和势函数；若能找到势函数，就不应再沿某条具体路径硬算。
\end{itemize}
\end{conclusion}

普通第二类曲面积分首先应读成通量，不要只凭“曲面有边界”就跳到 Stokes。只有当被积表达式本身是某个旋度通量，或者题目要求在线积分和曲面积分之间互换时，才自然进入 Stokes。

若题目只给端点、任意路径、复杂闭边界、空间交线、闭曲面通量或含奇点区域，这些信号往往要求先去检查保守场、Green、Stokes、Gauss，或先做补边、封口、换面、挖洞、分块等处理。

\subsection{第一类曲线积分与弧长元}

\begin{property}[第一类曲线积分的三个直接结论]
\begin{itemize}
    \item \(\displaystyle \int_{\Gamma}\mathrm{d}s\) 就是曲线 \(\Gamma\) 的弧长，所以第一类曲线积分与方向无关。
    \item 若 \(\Gamma:\vec r(t)\,(a\le t\le b)\)，则 \(\displaystyle \int_{\Gamma}f\,\mathrm{d}s=\int_a^b f(\vec r(t))\|\vec r\,'(t)\|\,\mathrm{d}t\)。
    \item 看到对称曲线时，先想奇偶性；奇函数常直接判零，偶函数常可化为半段积分的两倍。
\end{itemize}
\end{property}

\begin{definition}[第一型曲线积分]\label{def:curve-integral-type1}
设曲线 $\Gamma$ 的参数方程为 $\vec{r}=\vec{r}(t)=(x(t),y(t),z(t))$，$t \in [\alpha, \beta]$ ($\alpha<\beta$)。由弧长公式可得 $\mathrm{d}s=\|\vec{r}^{\,\prime}(t)\|\,\mathrm{d}t$，从而
\[
s(t)=\int_{\alpha}^{t}\|\vec{r}^{\,\prime}(\tau)\|\,\mathrm{d}\tau,
\qquad
\int_{\Gamma}f(p)\,\mathrm{d}s
=\int_{\alpha}^{\beta}f(\vec{r}(t))\|\vec{r}^{\,\prime}(t)\|\,\mathrm{d}t.
\]也就是\[
\frac{\mathrm{d}s}{\mathrm{d}t}
=\sqrt{x^{\prime2}(t)+y^{\prime2}(t)+z^{\prime2}(t)},\qquad
\mathrm{d}s=\sqrt{x^{\prime2}(t)+y^{\prime2}(t)+z^{\prime2}(t)}\,\mathrm{d}t.
\]
\end{definition}

\begin{example}
求曲线 $\Gamma$: $\vec{r}(t)=(a\cos t, a\sin t, bt)$，$0\le t\le2\pi$，$a>0; b>0$ 的质量，其中线密度 $\rho(x,y,z)=x^{2}+y^{2}+z^{2}$。
\end{example}
\begin{solution}
题目给了线密度和 $\mathrm{d}s$，这是第一类曲线积分，直接套参数公式 \(\displaystyle I=\int_0^{2\pi}\rho(\vec r(t))\|\vec r\,'(t)\|\,\mathrm{d}t\)。由 \(\displaystyle \vec{r}(t)=(a\cos t,a\sin t,bt)\) 得
\(\displaystyle \vec{r}^{\,\prime}(t)=(-a\sin t,a\cos t,b),\quad
\|\vec{r}^{\,\prime}(t)\|=\sqrt{a^2+b^2}\)。
又 \(\displaystyle \rho(\vec r(t))=a^2\cos^2 t+a^2\sin^2 t+b^2t^2=a^2+b^2t^2\)，
于是
\[
I=\int_{0}^{2\pi}(a^2+b^2t^2)\sqrt{a^2+b^2}\,\mathrm{d}t
=\sqrt{a^2+b^2}\left(2\pi a^2+\frac{8}{3}\pi^3b^2\right).
\]
\end{solution}

\begin{example}
求 \(\displaystyle \int_{\Gamma}(x^{2}+y^{2})\,\mathrm{d}s\)，其中 $\Gamma$ 为曲线 $(x(t), y(t)) = (a(\cos t+t\sin t), a(\sin t-t\cos t))$，$t\in[0,2\pi]$，$a>0$。
\end{example}
\begin{solution}
先求曲线的导数：\(\displaystyle x'(t)=at\cos t,\qquad y'(t)=at\sin t\)。因此 \(\displaystyle \mathrm{d}s=\sqrt{x'^2(t)+y'^2(t)}\,\mathrm{d}t=at\,\mathrm{d}t\)，同时 \(\displaystyle x^2+y^2=a^2(\cos t+t\sin t)^2+a^2(\sin t-t\cos t)^2=a^2(1+t^2)\)。代入第一类曲线积分公式得
\begin{align*}
\int_{\Gamma}(x^{2}+y^{2})\,\mathrm{d}s&=\int_{0}^{2\pi}\left(a^{2}(\cos t+t\sin t)^{2}+a^{2}(\sin t-t\cos t)^{2}\right)\sqrt{a^{2}t^{2}\cos^{2}t+a^{2}t^{2}\sin^{2}t}\,\mathrm{d}t\\
&=\int_{0}^{2\pi}a^{3}t(1+t^{2})\,\mathrm{d}t=2\pi^{2}a^{3}(1+2\pi^{2}).
\end{align*}
\end{solution}

\begin{example}
求一个均匀半圆周对位于圆心的单位质量的引力，这里 $\rho=1$。
\end{example}
\begin{solution}
设曲线 $\Gamma$ 的方程为 $x=a\cos t$，$y=a\sin t$，$0\le t\le\pi$，$a>0$。如图3，由微元法可知 \(\displaystyle \vec{F}=(F_{x},F_{y})=\left(K\int_{\Gamma}\frac{x}{a^{3}}\,\mathrm{d}s,\,K\int_{\Gamma}\frac{y}{a^{3}}\,\mathrm{d}s\right)\)。
由于半圆关于 $y$ 轴对称，取关于 $y$ 轴对称的两点时，它们的水平分力大小相等、方向相反，所以 $F_x=0$；而竖直分力同向相加，因此总力只剩竖直分量需要计算。下面计算
\begin{align*}
F_{y}&=K\int_{\Gamma}\frac{y}{a^{3}}\,\mathrm{d}s=\frac{K}{a^{3}}\int_{0}^{\pi}a\sin t\sqrt{a^{2}\sin^{2}t+a^{2}\cos^{2}t}\,\mathrm{d}t=\frac{K}{a}\int_{0}^{\pi}\sin t\,\mathrm{d}t=\frac{2K}{a}.
\end{align*}
\end{solution}

\begin{example}
计算 $\int_{\Gamma}x^{2}\,\mathrm{d}s$，其中 $\Gamma$ 为球面 $x^{2}+y^{2}+z^{2}=a^{2}$ 与平面 $x+y+z=0$ 交成的圆周。
\end{example}
\begin{solution}
我们先给出曲线 $\Gamma$ 的参数方程。记平面 $x+y+z=0$ 为 $\Pi$，它的法向量为
\(\vec{n}=(1,1,1)\)。由于 $\Gamma$ 是球面与过原点的平面 $\Pi$ 的交线，故 $\Gamma$ 是位于
$\Pi$ 内、圆心在原点、半径为 $a$ 的圆。为了写出这个圆的参数方程，需要先在平面 $\Pi$ 内找一组
正交单位基。先取与法向量垂直的一个单位向量：
\[
\vec{e}_{1}=\frac{\vec{u}_{1}}{|\vec{u}_{1}|}=\frac{1}{\sqrt{6}}(1,-2,1).
\]
第二个基底方向可由叉乘直接得到。令
\[
\vec{u}_{2}=\vec{n}\times\vec{u}_{1}
=\begin{vmatrix}
\vec{i}&\vec{j}&\vec{k}\\
1&1&1\\
1&-2&1
\end{vmatrix}
=(3,0,-3)=3(1,0,-1).
\]
叉乘所得的 $\vec{u}_{2}$ 同时垂直于 $\vec{n}$ 和 $\vec{u}_{1}$，因此 $\vec{u}_{2}$ 也在平面 $\Pi$ 内，且与
$\vec{u}_{1}$ 垂直。把它单位化，得
\[
\vec{e}_{2}=\frac{\vec{u}_{2}}{|\vec{u}_{2}|}=\frac{1}{\sqrt{2}}(1,0,-1).
\]
起点在原点终点在 $\Gamma$ 上的任何一个向量 $\vec{r}=(x,y,z)$ 均位于 $\vec{e}_{1}$，$\vec{e}_{2}$ 张成的平面内，因此向量 $\vec{r}$ 可以被向量 $\vec{e}_{1}$，$\vec{e}_{2}$ 线性组合出来，设$\Gamma$ 的向量方程是 \[\displaystyle \vec{r}=a(\vec{e}_{1}\cos\theta+\vec{e}_{2}\sin\theta)=\left(a\left(\frac{1}{\sqrt{6}}\cos\theta+\frac{1}{\sqrt{2}}\sin\theta\right),a\left(-\frac{2}{\sqrt{6}}\cos\theta\right),a\left(\frac{1}{\sqrt{6}}\cos\theta-\frac{1}{\sqrt{2}}\sin\theta\right)\right)\]。此时，弧长微元是 $\mathrm{d}s=a\,\mathrm{d}\theta$，所以
\begin{align*}
\int_{\Gamma}x^{2}\,\mathrm{d}s&=a\int_{0}^{2\pi}a^{2}\left(\frac{1}{\sqrt{6}}\cos\theta+\frac{1}{\sqrt{2}}\sin\theta\right)^{2}\mathrm{d}\theta\\
&=a^{3}\int_{0}^{2\pi}\left[\frac{1}{6}\cos^{2}\theta+\frac{1}{\sqrt{3}}\sin\theta\cos\theta+\frac{1}{2}\sin^{2}\theta\right]\mathrm{d}\theta=\frac{2a^{3}\pi}{3}.
\end{align*}
此外，由 $\Gamma$ 的参数表示很容易求出 \(\displaystyle \int_{\Gamma}y^{2}\,\mathrm{d}s=\frac{2a^{3}\pi}{3},\ \int_{\Gamma}z^{2}\,\mathrm{d}s=\frac{2a^{3}\pi}{3}\)。
\end{solution}

\begin{example}
求 $\oint_{\Gamma}(x+y^{3})\,\mathrm{d}s$，其中 $\Gamma$ 是圆周 $x^{2}+y^{2}=R^{2}$，$R>0$。
\end{example}
\begin{solution}
取参数方程 \(\displaystyle x=R\cos t,\qquad y=R\sin t,\qquad 0\le t\le 2\pi,\qquad \mathrm{d}s=R\,\mathrm{d}t\)。
于是
\[
\oint_{\Gamma}x\,\mathrm{d}s=R^2\int_0^{2\pi}\cos t\,\mathrm{d}t=0,\qquad
\oint_{\Gamma}y^3\,\mathrm{d}s=R^4\int_0^{2\pi}\sin^3 t\,\mathrm{d}t=0.
\]
故 \(\displaystyle \oint_{\Gamma}(x+y^3)\,\mathrm{d}s=0\)。
\end{solution}

\begin{example}
计算 \(\displaystyle \int_{\Gamma}|y|\,\mathrm{d}s\)，其中 $\Gamma$ 是右半圆周 $x^{2}+y^{2}=R^{2}$，$x\ge0$ (图5)。
\end{example}
\begin{solution}
把右半圆参数化为 \(\displaystyle x=R\cos t,\ y=R\sin t,\ -\frac{\pi}{2}\le t\le\frac{\pi}{2}\)，于是 \(\displaystyle \mathrm{d}s=\sqrt{(-R\sin t)^2+(R\cos t)^2}\,\mathrm{d}t=R\,\mathrm{d}t,\ |y|=R|\sin t|\)。
所以 \(\displaystyle \int_{\Gamma}|y|\,\mathrm{d}s
=R^2\int_{-\pi/2}^{\pi/2}|\sin t|\,\mathrm{d}t
=2R^2\int_0^{\pi/2}\sin t\,\mathrm{d}t = 2R^2\)。
\end{solution}


\subsection{第一类曲面积分与面积元}

\begin{definition}[正则曲面]\label{def:regular-surface}
设曲面 $\Sigma$ 有参数方程 $\vec{r}=\vec{r}(u,v)=(x(u,v),y(u,v),z(u,v))$，其中 $(u,v)\in D$。若 $x(u,v)$，$y(u,v)$，$z(u,v)$ 在 $D$ 上具有连续的一阶偏导数，并且
\[
\vec{r}_{u}(u,v)\times\vec{r}_{v}(u,v)\ne\vec{0},\qquad (u,v)\in D,
\]
则称 $\Sigma$ 为正则曲面。这里 $\vec{r}_{u}$，$\vec{r}_{v}$ 张成曲面的切平面，条件 $\vec{r}_{u}\times\vec{r}_{v}\ne\vec{0}$ 表明两个切向量不共线，从而面积元不会退化为零。
\end{definition}

\begin{definition}[曲面面积]\label{def:surface-area}
设正则曲面 $\Sigma$ 有参数方程 $\vec{r}=\vec{r}(u,v), (u,v)\in D$。我们称 \(\displaystyle S(\Sigma)=\iint_{D}\|\vec{r}_{u}\times\vec{r}_{v}\|\mathrm{d}u\,\mathrm{d}v\) 为曲面 $\Sigma$ 的面积，称 \(\displaystyle \mathrm{d}S=\|\vec{r}_{u}\times\vec{r}_{v}\|\mathrm{d}u\,\mathrm{d}v\) 为曲面 $\Sigma$ 的面积元。令
\[
\vec{r}_{u}=(x_{u},y_{u},z_{u}),\qquad \vec{r}_{v}=(x_{v},y_{v},z_{v}).
\]
由叉乘的坐标公式，
\[
\vec{r}_{u}\times\vec{r}_{v}=(y_{u}z_{v}-z_{u}y_{v},\ z_{u}x_{v}-x_{u}z_{v},\ x_{u}y_{v}-y_{u}x_{v}).
\]
因此
\begin{align*}
\|\vec{r}_{u}\times\vec{r}_{v}\|^{2}
={}&(y_{u}z_{v}-z_{u}y_{v})^{2}+(z_{u}x_{v}-x_{u}z_{v})^{2}+(x_{u}y_{v}-y_{u}x_{v})^{2}\\
={}&y_{u}^{2}z_{v}^{2}+z_{u}^{2}y_{v}^{2}-2y_{u}z_{u}y_{v}z_{v}
+z_{u}^{2}x_{v}^{2}+x_{u}^{2}z_{v}^{2}-2z_{u}x_{u}x_{v}z_{v}+x_{u}^{2}y_{v}^{2}+y_{u}^{2}x_{v}^{2}-2x_{u}y_{u}x_{v}y_{v}\\
={}&(x_{u}^{2}+y_{u}^{2}+z_{u}^{2})(x_{v}^{2}+y_{v}^{2}+z_{v}^{2})-(x_{u}x_{v}+y_{u}y_{v}+z_{u}z_{v})^{2}\\
={}&\|\vec{r}_{u}\|^{2}\|\vec{r}_{v}\|^{2}-(\vec{r}_{u}\cdot\vec{r}_{v})^{2}.
\end{align*}
令 \(\displaystyle E=\|\vec{r}_{u}\|^{2}=x_{u}^{2}+y_{u}^{2}+z_{u}^{2}\)，
\(\displaystyle G=\|\vec{r}_{v}\|^{2}=x_{v}^{2}+y_{v}^{2}+z_{v}^{2}\)，
\(\displaystyle F=\vec{r}_{u}\cdot\vec{r}_{v}=x_{u}x_{v}+y_{u}y_{v}+z_{u}z_{v}\)。
所以 $\mathrm{d}S=\sqrt{EG-F^{2}}\,\mathrm{d}u\,\mathrm{d}v$ 为面积元素。若 \(\Sigma:z=f(x,y)\)，则 \[\mathrm{d}S=\sqrt{1+f_x^2+f_y^2}\,\mathrm{d}x\,\mathrm{d}y\]特别的，平面图形也可以看成一张曲面。如果 $D$ 是 $xy$ 平面上的点集，有参数方程 $\vec{r}(x,y)=(x,y,0)$，其中 $(x,y)\in D$，则
\[\displaystyle \mathrm{d}S=\|\vec{r}_{x}\times\vec{r}_{y}\|\mathrm{d}x\,\mathrm{d}y=\mathrm{d}x\,\mathrm{d}y\]
\end{definition}


\begin{property}[第一类曲面积分的直接性质]
(1) 第一类曲面积分具有线性性以及对积分曲面的可加性。

(2) 若 $\Sigma$ 是一个封闭曲面，则 $\Sigma$ 上的第一类曲面积分记为 $\displaystyle\oint_{\Sigma}f(x,y,z)\,\mathrm{d}S$。
\end{property}


\begin{example}
求球面 $x^{2}+y^{2}+z^{2}=a^{2}$ 的面积。
\end{example}
\begin{solution}
求曲面面积就是第一类曲面积分 \(\displaystyle S=\iint_\Sigma \mathrm{d}S\)。
球面最稳的写法是球坐标参数化。
取
\[
\vec r(\phi,\theta)=(a\sin\phi\cos\theta,\ a\sin\phi\sin\theta,\ a\cos\phi),
\quad
0\le \phi\le \pi,\ 0\le \theta\le 2\pi.
\]
则
有 \(\displaystyle E=x_\phi^2+y_\phi^2+z_\phi^2=a^2\)，
\(\displaystyle G=x_\theta^2+y_\theta^2+z_\theta^2=a^2\sin^2\phi\)，
\(\displaystyle F=x_\phi x_\theta+y_\phi y_\theta+z_\phi z_\theta=0\)。
故面积元为
\begin{align*}
\mathrm{d}S&=\sqrt{EG-F^2}\,\mathrm{d}\phi\,\mathrm{d}\theta
=a^2\sin\phi\,\mathrm{d}\phi\,\mathrm{d}\theta,\\
S&=\int_0^\pi\int_0^{2\pi}a^2\sin\phi\,\mathrm{d}\theta\,\mathrm{d}\phi
=2\pi a^2\int_0^\pi \sin\phi\,\mathrm{d}\phi
=4\pi a^2.
\end{align*}

$\phi\in[0,\pi]$、$\theta\in[0,2\pi]$ 恰好把球面覆盖一遍，所以面积为 \(\displaystyle S=4\pi a^2\)。
\end{solution}

\begin{example}
计算曲面积分 \(\displaystyle I=\iint_{\Sigma}(x^{2}+y^{2}+z^{2})\,\mathrm{d}S\)，
其中 $\Sigma$ 是球面 $x^{2}+y^{2}+z^{2}=a^{2}$。
\end{example}
\begin{solution}
令 $x=a\sin\phi\cos\theta, y=a\sin\phi\sin\theta, z=a\cos\phi$，其中 $\phi\in[0,\pi],\theta\in[0,2\pi]$，则
\[
I = \iint_{\Sigma}(x^{2}+y^{2}+z^{2})\,\mathrm{d}S
= \int_{0}^{\pi}\mathrm{d}\phi\int_{0}^{2\pi}a^{2}\cdot a^{2}\sin\phi\,\mathrm{d}\theta
= 4\pi a^{4}.
\]
\end{solution}

\begin{example}
计算曲面积分 \(\displaystyle I=\iint_{\Sigma}\frac{2x^{2}+2y^{2}+z}{\sqrt{4x^{2}+4y^{2}+1}}\,\mathrm{d}S\)，
其中 $\Sigma$ 是曲面 $z=9-x^{2}-y^{2}$，$z\ge0$。
\end{example}
\begin{solution}
如图3，曲面的参数方程可以表示为 \(\displaystyle \vec{r}(x,y)=(x,y,9-x^{2}-y^{2}),\ (x,y)\in D_{xy}=\{(x,y):x^{2}+y^{2}\le9\}\)，且 \(\displaystyle \mathrm{d}S=\sqrt{1+4x^{2}+4y^{2}}\,\mathrm{d}x\,\mathrm{d}y\)。
因此 \(\displaystyle I=\iint_{D_{xy}}(9+x^{2}+y^{2})\,\mathrm{d}x\,\mathrm{d}y
=\int_{0}^{3}\mathrm{d}r\int_{0}^{2\pi}(9+r^{2})r\,\mathrm{d}\theta
=\frac{243\pi}{2}\)。
\end{solution}

\begin{example}
求密度 $\rho\equiv1$ 的均匀球壳 $x^{2}+y^{2}+z^{2}=a^{2}$，$z\ge0$ 对于 $oz$ 轴的转动惯量。
\end{example}
\begin{solution}
令 $x=a\sin\phi\cos\theta, y=a\sin\phi\sin\theta, z=a\cos\phi$，其中 $\phi\in\left[0,\frac{\pi}{2}\right],\theta\in[0,2\pi]$，则
\(\displaystyle I_{z}=\iint_{\Sigma}(x^{2}+y^{2})\,\mathrm{d}S
=\int_{0}^{\frac{\pi}{2}}\mathrm{d}\phi\int_{0}^{2\pi}a^{4}\sin^{3}\phi\,\mathrm{d}\theta
=\frac{4}{3}\pi a^{4}\)。
\end{solution}

\begin{example}
如图4，计算积分 $I=\iint_{\Sigma}x\,\mathrm{d}S$，其中 $\Sigma$ 是柱面 $\Sigma_{1}:x^{2}+y^{2}=1$，平面 $\Sigma_{2}:z=1+x$ 以及 $\Sigma_{3}:z=0$ 所围成的空间立体的表面。
\end{example}
\begin{solution}
第一类曲面积分对曲面具有可加性，所以把边界曲面拆成三块后，积分也随之拆成三项：
\[
I=\iint_{\Sigma}x\,\mathrm{d}S
=\iint_{\Sigma_{1}}x\,\mathrm{d}S+\iint_{\Sigma_{2}}x\,\mathrm{d}S+\iint_{\Sigma_{3}}x\,\mathrm{d}S
:=I_{1}+I_{2}+I_{3}.
\]
因此下面分别计算 $I_{1}, I_{2}, I_{3}$。

(1) 求 $I_{1}$：取
\begin{align*}
\vec{r}_{1}(\theta,z)&=(\cos\theta,\sin\theta,z),\quad 0\le\theta\le2\pi,\quad 0\le z\le1+\cos\theta,\\
\frac{\partial\vec{r}_{1}}{\partial\theta}&=(-\sin\theta,\cos\theta,0),\quad \frac{\partial\vec{r}_{1}}{\partial z}=(0,0,1),\quad E=G=1,\ F=0,\ \mathrm{d}S=\mathrm{d}\theta\,\mathrm{d}z.
\end{align*}
故 \(\displaystyle I_{1}=\iint_{\Sigma_{1}}x\,\mathrm{d}S
=\int_{0}^{2\pi}\int_{0}^{1+\cos\theta}\cos\theta\,\mathrm{d}z\,\mathrm{d}\theta
=\int_{0}^{2\pi}\cos\theta(1+\cos\theta)\,\mathrm{d}\theta=\pi\)。

(2) 求 $I_{2}$：取
\begin{align*}
\vec{r}_{2}(x,y)&=(x,y,1+x),\quad (x,y)\in D_{xy}=\{(x,y):x^{2}+y^{2}\le1\},\\
\frac{\partial\vec{r}_{2}}{\partial x}&=(1,0,1),\quad \frac{\partial\vec{r}_{2}}{\partial y}=(0,1,0),\quad E=2,\ G=1,\ F=0,\ \mathrm{d}S=\sqrt{2}\,\mathrm{d}x\,\mathrm{d}y.
\end{align*}
故 \(\displaystyle I_{2}=\iint_{\Sigma_{2}}x\,\mathrm{d}S=\iint_{D_{xy}}\sqrt{2}x\,\mathrm{d}x\,\mathrm{d}y=\sqrt{2}\int_{0}^{1}r^{2}\mathrm{d}r\int_{0}^{2\pi}\cos\theta\,\mathrm{d}\theta=0\)。

(3) 求 $I_{3}$：取
\begin{align*}
\vec{r}_{3}(x,y)&=(x,y,0),\quad (x,y)\in D_{xy}=\{(x,y):x^{2}+y^{2}\le1\},\\
\frac{\partial\vec{r}_{3}}{\partial x}&=(1,0,0),\quad \frac{\partial\vec{r}_{3}}{\partial y}=(0,1,0),\quad E=G=1,\ F=0,\ \mathrm{d}S=\mathrm{d}x\,\mathrm{d}y.
\end{align*}
故 \(\displaystyle I_{3}=\iint_{\Sigma_{3}}x\,\mathrm{d}S=\iint_{D_{xy}}x\,\mathrm{d}x\,\mathrm{d}y=\int_{0}^{1}r^{2}\mathrm{d}r\int_{0}^{2\pi}\cos\theta\,\mathrm{d}\theta=0\)。
所以 $I=I_{1}+I_{2}+I_{3}=\pi$。
\end{solution}

\begin{example}
计算积分 \(\displaystyle A=\int_{\Sigma}(x+y+z)\,\mathrm{d}S\)，
其中 $\Sigma$ 是上半球面 $x^{2}+y^{2}+z^{2}=a^{2}$，$z\ge0$。
\end{example}
\begin{solution}
取参数方程 \(\displaystyle \vec{r}(\phi,\theta) = (x(\phi,\theta),y(\phi,\theta),z(\phi,\theta)),\ \phi\in\left[0,\frac{\pi}{2}\right],\ \theta\in[0,2\pi]\)。由线性性 \(\displaystyle A=\int_{\Sigma}x\,\mathrm{d}S+\int_{\Sigma}y\,\mathrm{d}S+\int_{\Sigma}z\,\mathrm{d}S\)。
又因对称性
\begin{align*}
\int_{\Sigma}x\,\mathrm{d}S&=\int_{0}^{2\pi}\mathrm{d}\theta\int_{0}^{\frac{\pi}{2}}a^{3}\sin^{2}\phi\cos\theta\,\mathrm{d}\phi=0,~~
\int_{\Sigma}y\,\mathrm{d}S=\int_{0}^{2\pi}\mathrm{d}\theta\int_{0}^{\frac{\pi}{2}}a^{3}\sin^{2}\phi\sin\theta\,\mathrm{d}\phi=0,
\end{align*}
从而 \(\displaystyle A=\int_{\Sigma}z\,\mathrm{d}S=\int_{0}^{2\pi}\mathrm{d}\theta\int_{0}^{\frac{\pi}{2}}a^{3}\sin\phi\cos\phi\,\mathrm{d}\phi=\pi a^{3}\)。另解。仍由对称性知前两项为零，因此
\begin{align*}
A=\int_{\Sigma}z\,\mathrm{d}S=\iint_{x^{2}+y^{2}\le a^{2}}z\sqrt{1+\left(\frac{x}{z}\right)^{2}+\left(\frac{y}{z}\right)^{2}}\mathrm{d}x\,\mathrm{d}y=\iint_{x^{2}+y^{2}\le a^{2}}z\cdot\frac{a}{z}\mathrm{d}x\,\mathrm{d}y=\pi a^{3}.
\end{align*}
\end{solution}

\begin{example}
设 $\Sigma$ 为单位球面 $x^{2}+y^{2}+z^{2}=1$，证明 Poisson 公式
\[
\int_{\Sigma}f(ax+by+cz)\,\mathrm{d}S=2\pi\int_{-1}^{1}f(\sqrt{a^{2}+b^{2}+c^{2}}t)\,\mathrm{d}t.
\]
\end{example}
\begin{solution}
记 \(R=\sqrt{a^{2}+b^{2}+c^{2}}\)。若 \(R=0\)，则等式两边都等于 \(4\pi f(0)\)，结论显然成立。下面设 \(R>0\)。取新的直角坐标系 \(uvw\)，使 \(w\) 轴正方向就是向量 \((a,b,c)\) 的方向，且 \(u,v,w\) 构成右手正交系。此时单位球面仍为$u^{2}+v^{2}+w^{2}=1$，而球面上任一点 \((x,y,z)\) 在 \(w\) 轴上的坐标可以用投影得到。因为新的 \(w\) 轴单位方向向量为 \(\dfrac{(a,b,c)}{R}\)，所以位置向量 \((x,y,z)\) 在这个方向上的有向投影长度就是它与该单位方向向量的点积，即
\[
w=(x,y,z)\cdot\frac{(a,b,c)}{R}=\frac{ax+by+cz}{R},~~~
ax+by+cz=Rw.
\]

于是被积函数只与新坐标 \(w\) 有关。固定 \(w\in[-1,1]\)，平面 \(w=\text{常数}\) 截单位球面得到半径为 \(\sqrt{1-w^{2}}\) 的圆。引入角度 \(\theta\)，将球面参数化为
\[
\boldsymbol r(w,\theta)=\bigl(\sqrt{1-w^{2}}\cos\theta,\sqrt{1-w^{2}}\sin\theta,w\bigr),
\qquad -1\le w\le1,\quad 0\le\theta\le2\pi.
\]
直接计算得
\[
\boldsymbol r_w=\left(-\frac{w}{\sqrt{1-w^{2}}}\cos\theta,-\frac{w}{\sqrt{1-w^{2}}}\sin\theta,1\right),
\quad
\boldsymbol r_\theta=\left(-\sqrt{1-w^{2}}\sin\theta,\sqrt{1-w^{2}}\cos\theta,0\right),
\]
从而 \(\|\boldsymbol r_w\times\boldsymbol r_\theta\|=1\)，即 \(\mathrm{d}S=\mathrm{d}w\,\mathrm{d}\theta\)。所以
\begin{align*}
\int_{\Sigma}f(ax+by+cz)\,\mathrm{d}S=\int_{-1}^{1}\int_{0}^{2\pi}f(Rw)\,\mathrm{d}\theta\,\mathrm{d}w=2\pi\int_{-1}^{1}f(\sqrt{a^{2}+b^{2}+c^{2}}\,t)\,\mathrm{d}t.
\end{align*}
\end{solution}

\begin{example}
设 $\Sigma(t)$ 为平面 $x+y+z=t$ 被单位球面 $x^{2}+y^{2}+z^{2}=1$ 截下的部分。设 \(\displaystyle F(x,y,z)=1-(x^{2}+y^{2}+z^{2})\)。证明：对于 $|t|\le\sqrt{3}$ 时，有
\(\displaystyle \int_{\Sigma(t)}F(x,y,z)\,\mathrm{d}S=\frac{\pi}{18}(3-t^{2})^{2}\)。
\end{example}
\begin{solution}
取新的坐标系 $uvw$，使得平面 $\pi:x+y+z=0$ 为 $w=0$，再在 $\pi$ 中取互相垂直的两个单位向量 $u,v$，使得 $uvw$ 构成新的右手正交系。$w$ 上的单位向量是 $\dfrac{1}{\sqrt{3}}(1,1,1)$，所以
\(\displaystyle w=(x,y,z)\cdot\frac{(1,1,1)}{\sqrt{3}}=\frac{x+y+z}{\sqrt{3}}=\frac{t}{\sqrt{3}}\)。
故平面 $x+y+z=t$ 在新的坐标系中为 $w=\dfrac{t}{\sqrt{3}}$，故 $\Sigma(t)$ 变为 $w=\dfrac{t}{\sqrt{3}}$ 被球面 $u^{2}+v^{2}+w^{2}=1$ 截下的部分，它在 $w=0$ 平面上的投影为 $u^{2}+v^{2}=1-w^{2}=1-\dfrac{t^{2}}{3}$。这样所讨论的积分就是
\begin{align*}
\int_{\Sigma(t)}F(x,y,z)\,\mathrm{d}S
&=\int_{0}^{\sqrt{1-\frac{t^{2}}{3}}}\int_{0}^{2\pi}\left(1-\frac{t^{2}}{3}-r^{2}\right)r\,\mathrm{d}r\,\mathrm{d}\phi
=\frac{\pi}{18}(3-t^{2})^{2}.
\end{align*}
\end{solution}

\subsection{第二类曲线积分与环量}
\begin{note}[记号先分清]
记号 \(\displaystyle \mathrm{d}s\)、\(\displaystyle \mathrm{d}\vec r\)、\(\displaystyle \vec e_T\,\mathrm{d}s\) 以及
\(\displaystyle P\,\mathrm{d}x+Q\,\mathrm{d}y+R\,\mathrm{d}z\) 看起来都像“沿曲线的小东西”，但含义并不相同：
\begin{itemize}
\item $\mathrm{d}s$ 只记录长度，没有方向。
\item $\mathrm{d}\vec r=(\mathrm{d}x,\mathrm{d}y,\mathrm{d}z)$ 记录沿曲线正向的位移，有方向。
\item $\vec e_T\,\mathrm{d}s$ 是“单位切向量 $\times$ 长度”，本质上就是位移微元 $\mathrm{d}\vec r$。
\item \(P\,\mathrm{d}x+Q\,\mathrm{d}y+R\,\mathrm{d}z\) 不是三项彼此无关的拼接，而是同一个点积 $\vec F\cdot \mathrm{d}\vec r$ 在坐标上的展开。
\end{itemize}
凡是题目问“作功”或“环量”，积的一定是位移而不是长度；如果把 $\mathrm{d}s$ 和 $\mathrm{d}\vec r$ 混掉，方向与符号就会全乱。
\end{note}

\begin{theorem}[第二型曲线积分的计算]\label{thm:curve-integral-type2-calc}
设区域 $D\subset\mathbb{R}^{3}$，连续向量场 $\vec{F}=(P,Q,R):D\rightarrow\mathbb{R}^{3}$，又设 $\Gamma\subset D$ 是一条有向光滑曲线，参数方程为
\[
\vec{r}(t)=(x(t),y(t),z(t)),\qquad \alpha\le t\le\beta,
\]
并且 $t$ 增加时 $\vec{r}(t)$ 的运动方向就是 $\Gamma$ 的正方向。由于
\[
\mathrm{d}p=\vec{r}^{\,\prime}(t)\,\mathrm{d}t
=(x^{\prime}(t),y^{\prime}(t),z^{\prime}(t))\,\mathrm{d}t
=(\mathrm{d}x,\mathrm{d}y,\mathrm{d}z),
\]
所以
\[
\begin{aligned}
\int_{\Gamma}\vec{F}(p)\cdot\mathrm{d}p=\int_{\alpha}^{\beta}\vec{F}(\vec{r}(t))\cdot\vec{r}^{\,\prime}(t)\,\mathrm{d}t=\int_{\alpha}^{\beta}\bigl[P(\vec{r}(t))x^{\prime}(t)+Q(\vec{r}(t))y^{\prime}(t)+R(\vec{r}(t))z^{\prime}(t)\bigr]\,\mathrm{d}t.
\end{aligned}
\]
写成直角坐标形式，就是
\[
\int_{\Gamma}\vec{F}\cdot\mathrm{d}p
=\int_{\Gamma}(P,Q,R)\cdot(\mathrm{d}x,\mathrm{d}y,\mathrm{d}z)
=\int_{\Gamma}P\,\mathrm{d}x+Q\,\mathrm{d}y+R\,\mathrm{d}z.
\]
这里的 $P\,\mathrm{d}x+Q\,\mathrm{d}y+R\,\mathrm{d}z$ 是同一个点积的坐标展开，不是三个互不相关的一元积分；沿曲线前进时，$\mathrm{d}x,\mathrm{d}y,\mathrm{d}z$ 都由同一个参数和同一个定向决定。

另一方面，若 $\vec e_T$ 表示沿 $\Gamma$ 正方向的单位切向量，则
\[
\vec e_T=\frac{\vec r^{\,\prime}(t)}{\|\vec r^{\,\prime}(t)\|},\qquad
\mathrm{d}s=\|\vec r^{\,\prime}(t)\|\,\mathrm{d}t,
\]
于是 $\vec e_T\,\mathrm{d}s=\vec r^{\,\prime}(t)\,\mathrm{d}t=\mathrm{d}p$。因此第二型曲线积分也可以理解为“向量场沿切向的分量，再乘弧长累加”：
\[
\int_{\Gamma}\vec{F}(p)\cdot\mathrm{d}p
=\int_{\Gamma}\vec{F}(p)\cdot\vec e_T\,\mathrm{d}s.
\]
若进一步写 $\vec e_T=(\cos\alpha,\cos\beta,\cos\gamma)$，则
\[
\int_{\Gamma}\vec{F}(p)\cdot\mathrm{d}p
=\int_{\Gamma}(P\cos\alpha+Q\cos\beta+R\cos\gamma)\,\mathrm{d}s.
\]
曲线反向时，$\mathrm{d}p$ 与 $\vec e_T$ 同时变号，所以第二型曲线积分也随之变号。
\end{theorem}

\begin{example}
设 $\Gamma:x^{2}+y^{2}=a^{2}$，$y\ge0$，$a>0$，起点为 $(a,0)$，终点为 $(-a,0)$ (图2)。求 \(\displaystyle I=\int_{\Gamma}x\,\mathrm{d}x+y\,\mathrm{d}y\)。
\end{example}
\begin{solution}
这是第二类曲线积分，而且被积式正好长成 $x\,\mathrm{d}x+y\,\mathrm{d}y$，先查全微分。


\(\displaystyle x\,\mathrm{d}x+y\,\mathrm{d}y=\frac12\,\mathrm{d}(x^2+y^2)\)，所以 \(\displaystyle I=\frac12\int_\Gamma \mathrm{d}(x^2+y^2)=\frac12\bigl[(x^2+y^2)\bigr]_{(a,0)}^{(-a,0)}\)。因此直接走端点差，不做参数化。又因为曲线 $\Gamma$ 在半圆 $x^2+y^2=a^2$ 上，起点终点都满足 $x^2+y^2=a^2$，故
\(\displaystyle I=\frac12(a^2-a^2)=0\)。

这题与路径无关，只看端点；两个端点落在同一条圆周上，所以结果为 $0$。
\end{solution}

\begin{conclusion}[全微分识别的物理解释]
由图2可知力 $\vec{F}=(x,y)$ 总是与质点的运动方向相垂直，故力作的功为零。
\end{conclusion}

\begin{example}
如图3，设曲线 $\Gamma$ 为 (1) 折线 $\overrightarrow{OAB}$；(2) 直线 $\overrightarrow{OB}$。求 \(\displaystyle I=\int_{\Gamma}(x^{2}+y^{2})\mathrm{d}x+(x^{2}-y^{2})\mathrm{d}y\)。
线段 $\overrightarrow{OA}$：$x=x, y=x, x\in[0,1]$；线段 $\overrightarrow{AB}$：$x=x, y=-x+2, x\in[1,2]$；线段 $\overrightarrow{OB}$：$x=x, y=0, x\in[0,2]$。
\end{example}
\begin{solution}
(1) 折线 $\Gamma=\overrightarrow{OAB}$ 由两段组成。沿 $\overrightarrow{OA}$ 时，$y=x,\ \mathrm{d}y=\mathrm{d}x$，故
\[
\int_{\overrightarrow{OA}}(x^{2}+y^{2})\mathrm{d}x+(x^{2}-y^{2})\mathrm{d}y
=\int_0^1 \bigl[(x^2+x^2)+(x^2-x^2)\bigr]\mathrm{d}x
=\int_0^1 2x^2\,\mathrm{d}x
=\frac{2}{3}.
\]
沿 $\overrightarrow{AB}$ 时，$y=2-x,\ \mathrm{d}y=-\mathrm{d}x$，故
\[
\int_{\overrightarrow{AB}}(x^{2}+y^{2})\mathrm{d}x+(x^{2}-y^{2})\mathrm{d}y
=\int_1^2 \bigl[(x^2+y^2)-(x^2-y^2)\bigr]\mathrm{d}x
=\int_1^2 2(2-x)^2\,\mathrm{d}x
=\frac{2}{3}.
\]
故 $I=\frac{2}{3}+\frac{2}{3}=\frac{4}{3}$。

(2) 沿 $\Gamma=\overrightarrow{OB}$ 时有 $y=0,\ \mathrm{d}y=0$，故
$I=\int_{0}^{2}x^{2}\mathrm{d}x=\frac{8}{3}$。
\end{solution}

\begin{example}
计算 \(\displaystyle I=\int_{\Gamma}\frac{x\mathrm{d}x+y\mathrm{d}y}{\sqrt{x^{2}+y^{2}}}\)，
其中 $\Gamma$ 是上半平面的光滑曲线 $\vec{r}(t)=(x(t),y(t))$，起点 $A=(1,0)$，终点 $B=(6,8)$。
\end{example}
\begin{solution}
先看分子：
\(\displaystyle x\,\mathrm{d}x+y\,\mathrm{d}y=\frac{1}{2}\mathrm{d}(x^2+y^2)\)，且 \(\displaystyle \frac{x\,\mathrm{d}x+y\,\mathrm{d}y}{\sqrt{x^{2}+y^{2}}}=\mathrm{d}\sqrt{x^2+y^2}\)，故 \(\displaystyle I=\sqrt{x^2+y^2}\Big|_{A}^{B}=10-1=9\)。
\end{solution}

\begin{example}
计算 \(\displaystyle I=\int_{\Gamma}x^{3}\mathrm{d}x+y^{3}\mathrm{d}y+z^{3}\mathrm{d}z\)，
其中 $\Gamma$ 是从点 $A=(3,2,1)$ 到点 $B=(0,0,0)$ 的直线段 $AB$。
\end{example}
\begin{solution}
直线段 $AB$ 可参数化为 \(\displaystyle \vec{r}(t)=(3t,2t,t),\ t\in[1,0]\)。这里 $t=1$ 对应点 $A=(3,2,1)$，$t=0$ 对应点 $B=(0,0,0)$，方向与题意一致，又有
\[
\mathrm{d}x=3\,\mathrm{d}t,\qquad \mathrm{d}y=2\,\mathrm{d}t,\qquad \mathrm{d}z=\mathrm{d}t,
\qquad
I=\int_{1}^{0}(3^{4}t^{3}+2^{4}t^{3}+t^{3})\mathrm{d}t=\int_{1}^{0}98t^{3}\mathrm{d}t=-\frac{49}{2}.
\]
\end{solution}

\begin{conclusion}[闭曲线的定向约定]
对平面上的封闭曲线，计算曲线积分时，可取闭曲线上任意一点为起点，曲线积分的积分值与起点无关。一般的，规定逆时针方向为曲线的正向。
\end{conclusion}

\begin{example}\label{ex:unit-circle}
计算 $\displaystyle\oint_C \frac{x\,\mathrm{d}y - y\,\mathrm{d}x}{x^2+y^2}$，其中 $C$ 为单位圆 $x^2+y^2=1$，取正向。
\end{example}
\begin{solution}
这是闭曲线上的第二类曲线积分，曲线本身是标准单位圆，直接按正向参数化即可。

取单位圆正向参数方程
\[
x=\cos t,\qquad y=\sin t,\qquad 0\le t\le 2\pi,
\]
则
\[
\mathrm{d}x=-\sin t\,\mathrm{d}t,\qquad
\mathrm{d}y=\cos t\,\mathrm{d}t,\qquad
x^2+y^2=1.
\]
因此
\[
\oint_C \frac{x\,\mathrm{d}y-y\,\mathrm{d}x}{x^2+y^2}
=\int_0^{2\pi}\bigl(\cos^2 t+\sin^2 t\bigr)\,\mathrm{d}t
=\int_0^{2\pi}\mathrm{d}t=2\pi.
\]

$t$ 的增加方向对应逆时针正向，且被积式在圆周上有定义，所以结果为
\[
\oint_C \frac{x\,\mathrm{d}y-y\,\mathrm{d}x}{x^2+y^2}=2\pi.
\]
\end{solution}

\begin{example}
求 \(\displaystyle I=\oint_{\Gamma}(x+y)\mathrm{d}x+(x-y)\mathrm{d}y\)，其中 $\Gamma$ 为依逆时针方向通过的椭圆 \(\displaystyle \frac{x^{2}}{a^{2}}+\frac{y^{2}}{b^{2}}=1\)。
\end{example}
\begin{solution}
取椭圆的标准参数方程 \(\displaystyle x=a\cos\theta,\ y=b\sin\theta,\ \theta\in[0,2\pi],\ \mathrm{d}x=-a\sin\theta\,\mathrm{d}\theta,\ \mathrm{d}y=b\cos\theta\,\mathrm{d}\theta\)。
代入第二类曲线积分得 \(\displaystyle I=\int_0^{2\pi}[(a\cos\theta+b\sin\theta)(-a\sin\theta)+(a\cos\theta-b\sin\theta)(b\cos\theta)]\,\mathrm{d}\theta\)，
即 \(\displaystyle I=\int_0^{2\pi}\left(ab\cos2\theta-\frac{a^2+b^2}{2}\sin2\theta\right)\mathrm{d}\theta=0\)。
最后一步为零，是因为 $\cos2\theta$ 与 $\sin2\theta$ 在一个完整周期 $[0,2\pi]$ 上积分都为零。
\end{solution}

\begin{example}
设一个质量为 $m$ 的质点受重力作用沿空间某光滑曲线 $\Gamma$ 移动，求重力做的功。
\end{example}
\begin{solution}
如图4建立坐标系，则重力场为 $\vec{F}=(0,0,-mg)$。又令 $A=(x_{1},y_{1},z_{1}), B=(x_{2},y_{2},z_{2})$，且 $\Gamma$ 的参数方程为 $\vec{r}(t)=(x(t),y(t),z(t))$，$t\in[\alpha,\beta]$。则
\(\displaystyle W=\int_{\Gamma}\vec{F}\cdot\mathrm{d}p=\int_{\Gamma}(0,0,-mg)\cdot(\mathrm{d}x,\mathrm{d}y,\mathrm{d}z)=-\int_{\alpha}^{\beta}mgz^{\prime}(t)\mathrm{d}t=-mg(z_{2}-z_{1})\)。
这说明重力做功只与起点、终点的高度差有关，与质点沿哪条光滑曲线移动无关。
若终点比起点低，即 $z_2<z_1$，则 $W>0$；反之把物体抬高时，重力做负功。
\end{solution}

\begin{example}
设 $\Gamma$ 是由球面 $x^{2}+y^{2}+z^{2}=a^{2}$ 和平面 $x+y+z=0$ 交成的圆周，从第一卦限内看，它的方向是按反时针方向进行的。计算第二型曲线积分：(1) $\int_{\Gamma}p\cdot\mathrm{d}p$；(2) $\int_{\Gamma}z\mathrm{d}x+x\mathrm{d}y+y\mathrm{d}z$。
\end{example}
\begin{solution}
我们先补出曲线 $\Gamma$ 的参数方程。记平面 $x+y+z=0$ 为 $\Pi$，它的法向量为
\(\vec{n}=(1,1,1)\)。由于 $\Gamma$ 是球面与过原点的平面 $\Pi$ 的交线，故 $\Gamma$ 是位于
$\Pi$ 内、圆心在原点、半径为 $a$ 的圆。为了写出这个圆的参数方程，需要先在平面 $\Pi$ 内找一组
正交单位基。先取
\[
\vec{u}_{1}=(1,-2,1),
\]
则 $\vec{u}_{1}\cdot\vec{n}=1-2+1=0$，所以 $\vec{u}_{1}$ 在平面 $\Pi$ 内。把它单位化，得
\[
\vec{e}_{1}=\frac{\vec{u}_{1}}{|\vec{u}_{1}|}=\frac{1}{\sqrt{6}}(1,-2,1).
\]
第二个基底方向可由叉乘直接得到。令
\[
\vec{u}_{2}=\vec{n}\times\vec{u}_{1}
=\begin{vmatrix}
\vec{i}&\vec{j}&\vec{k}\\
1&1&1\\
1&-2&1
\end{vmatrix}
=(3,0,-3)=3(1,0,-1).
\]
叉乘所得的 $\vec{u}_{2}$ 同时垂直于 $\vec{n}$ 和 $\vec{u}_{1}$，因此 $\vec{u}_{2}$ 也在平面 $\Pi$ 内，且与
$\vec{u}_{1}$ 垂直。把它单位化，得
\[
\vec{e}_{2}=\frac{\vec{u}_{2}}{|\vec{u}_{2}|}=\frac{1}{\sqrt{2}}(1,0,-1).
\]
起点在原点终点在 $\Gamma$ 上的任何一个向量 $\vec{r}=(x,y,z)$ 均位于 $\vec{e}_{1}$，$\vec{e}_{2}$ 张成的平面内，因此向量 $\vec{r}$ 可以被向量 $\vec{e}_{1}$，$\vec{e}_{2}$ 线性组合出来。取
\[
\vec{r}=a(\vec{e}_{1}\cos\theta+\vec{e}_{2}\sin\theta),\qquad 0\le\theta\le2\pi,
\]
即
\begin{align*}
x=a\left(\frac{1}{\sqrt{6}}\cos\theta+\frac{1}{\sqrt{2}}\sin\theta\right),~y=a\left(-\frac{2}{\sqrt{6}}\cos\theta\right),~z=a\left(\frac{1}{\sqrt{6}}\cos\theta-\frac{1}{\sqrt{2}}\sin\theta\right)
\end{align*}
并且 $\vec{e}_{1}\times\vec{e}_{2}=\dfrac{1}{\sqrt{3}}(1,1,1)$ 指向第一卦限，所以 $\theta$ 增加时，从第一卦限内看正是反时针方向，与题设定向一致。

(1) 因为 $\|p\|^2=a^2$ 在 $\Gamma$ 上恒定，所以
\(\displaystyle \mathrm{d}\|p\|^2=2p\cdot\mathrm{d}p=0,\ \int_{\Gamma}p\cdot\mathrm{d}p=0\)。

(2) 由对称性
\begin{align*}
\int_{\Gamma}y\mathrm{d}z&=\frac{2}{\sqrt{6}}a^{2}\int_{0}^{2\pi}\left(\frac{1}{\sqrt{6}}\sin\theta+\frac{1}{\sqrt{2}}\cos\theta\right)\cos\theta\,\mathrm{d}\theta=\frac{2}{\sqrt{6}}a^{2}\frac{1}{\sqrt{2}}\int_{0}^{2\pi}\cos^{2}\theta\,\mathrm{d}\theta=\frac{\sqrt{3}}{3}\pi a^{2},
\end{align*}
按 $x,y,z$ 轮换可得
\(\displaystyle \int_{\Gamma}z\mathrm{d}x+x\mathrm{d}y+y\mathrm{d}z=\sqrt{3}\pi a^{2}\)。
\end{solution}

\subsection{第二类曲线积分：保守场、Green、Stokes}

\subsubsection{开曲线：保守场、势函数与路径无关}

开曲线题的核心是端点差。若存在数量函数 \(u\)，使
\[
\vec F=\operatorname{grad}u=\nabla u=(u_x,u_y,u_z),
\]
则称 \(\vec F\) 为梯度场，也就是保守场；\(u\) 称为势函数。此时沿任意从 \(A\) 到 \(B\) 的路径都有
\[
\int_A^B \vec F\cdot\mathrm{d}\vec r=\int_A^B \nabla u\cdot\mathrm{d}\vec r=u(B)-u(A).
\]
所以开曲线第二类积分如果只给端点，第一反应应是找势函数，而不是硬选一条路径。

梯度场必无旋：
\[
\nabla\times(\nabla u)=\vec0.
\]
在平面情形中，这就退化为 \(P_y=Q_x\)；在空间情形中，则退化为三组交叉偏导相等。它解释了为什么保守场判别总是先检查交叉偏导是否相等。

\begin{definition}[单连通区域]\label{def:simply-connected}
设 $\Omega$ 是连通区域。若 $\Omega$ 内任意一条闭曲线都可以始终留在 $\Omega$ 内连续缩成一个点，则称 $\Omega$ 为\textbf{单连通区域}。在平面区域中，可以等价理解为：$\Omega$ 内任意一条简单闭曲线所围成的内部区域仍全部属于 $\Omega$。直观地说，单连通区域就是没有洞的区域；圆盘、矩形、整个平面和整个空间都是单连通区域，而 $\mathbb R^2\setminus\{(0,0)\}$ 不是单连通区域。不是单连通的区域称为\textbf{多连通区域}（或复连通区域）。
\end{definition}

\begin{theorem}[保守场判别法]\label{thm:conservative-field-test}
设向量场各分量在单连通区域 $\Omega$ 内具有连续一阶偏导数。

若 $\Omega\subset\mathbb R^2$，$\vec F=(P,Q)$，则
\[
\vec F\text{ 为保守场}
\quad\Longleftrightarrow\quad
\frac{\partial P}{\partial y}=\frac{\partial Q}{\partial x}\quad(\text{在 }\Omega\text{ 内处处成立}).
\]

若 $\Omega\subset\mathbb R^3$，$\vec F=(P,Q,R)$，则
\[
\vec F\text{ 为保守场}
\quad\Longleftrightarrow\quad
\operatorname{rot}\vec F=\nabla\times\vec F=\vec0,
\]
即
\[
\frac{\partial P}{\partial y}=\frac{\partial Q}{\partial x},\qquad
\frac{\partial P}{\partial z}=\frac{\partial R}{\partial x},\qquad
\frac{\partial Q}{\partial z}=\frac{\partial R}{\partial y}.
\]
因此做题时先检查定义域是否单连通，再检查交叉偏导是否相等；若定义域不是单连通，上述等式仍是保守场的必要条件，但不能单凭它断定保守场。
\end{theorem}

\begin{example}
设 $\vec{F} = (yz, zx, xy)$，验证其为保守场，并求势函数。
\end{example}
\begin{solution}
按定理~\ref{thm:conservative-field-test}，先判定义域，再验交叉偏导。这里定义域是整个 $\mathbb R^3$，是单连通区域。记 $P = yz$，$Q = zx$，$R = xy$。
由于 \(\displaystyle \frac{\partial P}{\partial y} = z = \frac{\partial Q}{\partial x},\quad
\frac{\partial Q}{\partial z} = x = \frac{\partial R}{\partial y},\quad
\frac{\partial R}{\partial x} = y = \frac{\partial P}{\partial z}\)，故 $\vec{F}$ 为保守场。

由 $\frac{\partial u}{\partial x} = yz$，对 $x$ 积分得 $u = xyz + g(y,z)$。

由 $\frac{\partial u}{\partial y} = xz + g_y = xz$，得 $g_y = 0$，$g = h(z)$。

由 $\frac{\partial u}{\partial z} = xy + h'(z) = xy$，得 $h'(z) = 0$，$h = C$。

故势函数为 $u = xyz + C$。
\end{solution}

\begin{corollary}[保守场的两个直接结论]
若 \(\vec F=\nabla u\)，则对任意从 \(A\) 到 \(B\) 的分段光滑曲线 \(\Gamma\)，都有 \(\displaystyle \int_\Gamma \vec F\cdot \mathrm{d}\vec r=u(B)-u(A),\quad \oint_\Gamma \vec F\cdot \mathrm{d}\vec r=0\)。
\end{corollary}

\begin{property}[曲线积分版微积分基本定理]
(1) 使用保守场的端点差公式时，第二类曲线积分只需 $\Gamma$ 分段光滑即可，不要求曲线有高阶光滑性。

(2) 定理证明中的关键公式 $\int_{A}^{B} \mathrm{d}u = u\big|_{A}^{B} = u(B) - u(A)$ 正是微积分基本定理在曲线积分情形的推广。
\end{property}

\begin{example}
计算 $\int_{\Gamma} x\,\mathrm{d}x + y\,\mathrm{d}y$，其中 $\Gamma$ 为从 $A(1,1)$ 到 $B(4,3)$ 的任意分段光滑曲线。
\end{example}
\begin{solution}
$\vec{F} = (x, y)$，即 $P = x,\ Q = y$。先验算保守场条件：$\dfrac{\partial P}{\partial y}=0,\ \dfrac{\partial Q}{\partial x}=0$。两者相等，且定义域是整个平面，因此 $\vec{F}$ 为保守场，曲线积分只与端点有关。

设势函数为 $u(x,y)$，则 $u_x=x$。对 $x$ 积分得 $u=\dfrac{x^2}{2}+g(y)$；再由 $u_y=g'(y)=y$ 得 $g(y)=\dfrac{y^2}{2}$，故可取 $u(x,y)=\dfrac{x^2+y^2}{2}$。

于是 \(\displaystyle \int_{\Gamma} x\,\mathrm{d}x + y\,\mathrm{d}y = u(B) - u(A) = \frac{16 + 9}{2} - \frac{1 + 1}{2} = \frac{23}{2}\)。
\end{solution}


\begin{example}
验证 $\vec{F} = (y\cos x, \sin x)$ 为保守场。
\end{example}
\begin{solution}
题目要求“验证是否为保守场”，先做判别，再决定是否求势函数。

二维场在单连通区域上，按“交叉偏导相等 + 定义域单连通”判保守性。

记 \(\displaystyle P=y\cos x,\ Q=\sin x\)。则
\(\displaystyle \frac{\partial P}{\partial y}=\cos x=\frac{\partial Q}{\partial x}\)。
定义域是整个平面，属于单连通区域，所以 $\vec F$ 为保守场。

若进一步求势函数，令 $u_x=P$，则
\(\displaystyle u=\int y\cos x\,\mathrm{d}x=y\sin x+\varphi(y)\)。
再由 $u_y=Q$ 得 \(\displaystyle \sin x+\varphi'(y)=\sin x,\ \varphi'(y)=0\)，
故 \(\displaystyle \varphi(y)=C\)，可取势函数 \(\displaystyle u=y\sin x+C\)。

既核对了交叉偏导，又核对了单连通条件，所以“保守场”结论成立。
\end{solution}

\begin{example}[力场功与路径无关]
设 $\vec{F} = (x + y^2, 2xy - 8)$，证明力场做功与路径无关，并计算从 $A(0,0)$ 到 $B(1,2)$ 的功。
\end{example}
\begin{solution}
题目问“做功与路径无关”，这是保守场信号；一旦判成保守场，就把积分降成端点差。

先判保守性，再用待定函数法求势函数。

记 \(\displaystyle P=x+y^2,\ Q=2xy-8\)。则
\(\displaystyle \frac{\partial P}{\partial y}=2y=\frac{\partial Q}{\partial x}\)。
定义域为全平面，故 $\vec F$ 是保守场，做功与路径无关。

由 $u_x=P$ 得
\(\displaystyle u=\int (x+y^2)\,\mathrm{d}x=\frac{x^2}{2}+xy^2+\varphi(y)\)。
再由 $u_y=Q$ 得 \(\displaystyle 2xy+\varphi'(y)=2xy-8,\ \varphi'(y)=-8,\ \varphi(y)=-8y+C\)。
故可取势函数 \(\displaystyle u=\frac{x^2}{2}+xy^2-8y+C\)。于是 \(\displaystyle W=u(B)-u(A)=u(1,2)-u(0,0)=\left(\frac12+4-16\right)-0=-\frac{23}{2}\)。

这题最后只看端点 $A,B$，与所走路径无关，所以答案为 \(\displaystyle W=-\frac{23}{2}\)。
\end{solution}


\subsubsection{Green 公式、二维Green第一第二公式}

\begin{theorem}[Green 公式]\label{thm:green}
设平面闭区域 $D$ 的边界 $C=\partial D$ 由有限条分段光滑闭曲线组成，并取正向。若函数 $u(x,y),v(x,y)$ 在包含 $D$ 的开区域内具有连续二阶偏导数，则
\begin{equation}\label{eq:green}
\oint_C P\,\mathrm{d}x + Q\,\mathrm{d}y = \iint_D \left(\frac{\partial Q}{\partial x}-\frac{\partial P}{\partial y}\right)\mathrm{d}x\,\mathrm{d}y,
\end{equation}
其中 $C$ 按区域 $D$ 的正向边界方向积分。在 Green 公式中取 $P=0,\;Q=uv$，则
\begin{align*}
\frac{\partial Q}{\partial x}-\frac{\partial P}{\partial y}
&=\frac{\partial(uv)}{\partial x}
=u\frac{\partial v}{\partial x}+v\frac{\partial u}{\partial x},~~
\oint_C uv\,\mathrm{d}y= \iint_D \left(u\frac{\partial v}{\partial x}+v\frac{\partial u}{\partial x}\right)\mathrm{d}x\,\mathrm{d}y.
\end{align*}
在 Green 公式中取 $P=-y,\ Q=x$，则$Q_x-P_y=1-(-1)=2$，所以对正向边界 $C$，
\begin{equation}\label{eq:area}
S(D)=\frac12\oint_C (x\,\mathrm{d}y-y\,\mathrm{d}x).
\end{equation}
同理，取 $P=0,\ Q=x$ 可得 \(\displaystyle S(D)=\oint_C x\,\mathrm{d}y\)；取 $P=-y,\ Q=0$ 可得 \(\displaystyle S(D)=-\oint_C y\,\mathrm{d}x\)。这些公式都要求 $C$ 取正向，若取负向则结果变号。
\end{theorem}

\begin{conclusion}[Green 公式的做题流程]\label{rem:green}
看到平面闭曲线积分时，按下面顺序落笔。

\textbf{(1) 认对象：}必须是平面上的 \(\displaystyle \oint_C P\,\mathrm{d}x+Q\,\mathrm{d}y\)。若曲线不闭合，先考虑补边；补成闭曲线后，最后要减去补边上的积分。

\textbf{(2) 判方向：}Green 公式默认用正向边界。简单闭曲线的正向是逆时针；若题目给的是顺时针，最后结果要变号。多连通区域按
\[
\partial D=C_0^+-C_1^+-\cdots-C_m^+
\]
记账，即外边界逆时针，内边界顺时针。

\textbf{(3) 查奇点：}只有当 $P,Q$ 在区域 $D$ 内以及边界上都连续可偏导时，才能直接套 Green。若区域内部有奇点，不能把奇点直接吞进 $D$，应先挖去奇点附近的小圆或小区域，再在剩下的区域上使用 Green。

\textbf{(4) 算核心：}真正进入计算时只算
\[
Q_x-P_y.
\]
若边界复杂而 $Q_x-P_y$ 简单，Green 通常最省；若边界本身很标准而区域积分更麻烦，也可以直接参数化边界。
\end{conclusion}

下面的 Green 第一、第二公式是平面区域上的二维版本，由普通 Green 公式推出；后面 Gauss 公式后的 Green 第一、第二公式是空间闭区域上的三维版本。二者结构相同，但积分对象和微元不同。

\begin{theorem}[平面 Green 第一公式（二维）]\label{ex:greens-formulas}
设 $D$ 是由分段光滑简单闭曲线 $C$ 围成的有界闭区域。$u(x,y),v(x,y)$ 在 $D$ 上具有二阶连续偏导数，$\dfrac{\partial u}{\partial n}$ 表示 $u$ 沿 $C$ 的外法线方向的方向导数。证明：
\[
 \oint_C v\frac{\partial u}{\partial n}\,\mathrm{d}s =\iint_D v\Delta u\,\mathrm{d}x\,\mathrm{d}y+ \iint_D \left(\frac{\partial u}{\partial x}\frac{\partial v}{\partial x}+\frac{\partial u}{\partial y}\frac{\partial v}{\partial y}\right)\mathrm{d}x\,\mathrm{d}y.
\]从左往右推：借助方向导数和梯度的关系；从右往左推，凑微分。
\end{theorem}

\begin{solution}
由于 $C$ 取正向时区域在左手边，所以外法向在右手边。若切向量和外法向量为
\[
\vec T=\left(\frac{\mathrm{d}x}{\mathrm{d}s},\frac{\mathrm{d}y}{\mathrm{d}s}\right),~~\vec n=\left(\frac{\mathrm{d}y}{\mathrm{d}s},-\frac{\mathrm{d}x}{\mathrm{d}s}\right).
\]
于是由方向导数$\frac{\partial u}{\partial n}=\nabla u\cdot\vec n$：
\[
\frac{\partial u}{\partial n}\,\mathrm{d}s
=\nabla u\cdot\vec n\,\mathrm{d}s
=\left(u_x\frac{\mathrm{d}y}{\mathrm{d}s}-u_y\frac{\mathrm{d}x}{\mathrm{d}s}\right)\mathrm{d}s
=u_x\,\mathrm{d}y-u_y\,\mathrm{d}x\implies
\oint_C v\frac{\partial u}{\partial n}\,\mathrm{d}s
=\oint_C v u_x\,\mathrm{d}y-vu_y\,\mathrm{d}x.
\]
代入 Green 公式：
\begin{align*}
\oint_C v\frac{\partial u}{\partial n}\,\mathrm{d}s
&=\oint_C (-vu_y)\,\mathrm{d}x+(vu_x)\,\mathrm{d}y=\iint_D \left[(vu_x)_x-(-vu_y)_y\right]\mathrm{d}x\,\mathrm{d}y\\
&=\iint_D \left[(vu_x)_x+(vu_y)_y\right]\mathrm{d}x\,\mathrm{d}y=\iint_D \left(v_xu_x+vu_{xx}+v_yu_y+vu_{yy}\right)\mathrm{d}x\,\mathrm{d}y\\
&=\iint_D \left(v\Delta u+u_xv_x+u_yv_y\right)\mathrm{d}x\,\mathrm{d}y.
\end{align*}
\end{solution}

\begin{solution}
\textbf{从右往左推。}
从公式右边出发，先把二重积分中的被积函数凑成全微分形式：
\begin{align*}
&\iint_D v\Delta u\,\mathrm{d}x\,\mathrm{d}y
+\iint_D (u_xv_x+u_yv_y)\,\mathrm{d}x\,\mathrm{d}y\\
&=\iint_D \left[v(u_{xx}+u_{yy})+u_xv_x+u_yv_y\right]\mathrm{d}x\,\mathrm{d}y\\
&=\iint_D \left[(vu_x)_x+(vu_y)_y\right]\mathrm{d}x\,\mathrm{d}y.
\end{align*}
现在它已经是 Green 公式右端的形状。令
\[
P=-vu_y,\qquad Q=vu_x,
\]
则
\[
Q_x-P_y=(vu_x)_x-(-vu_y)_y=(vu_x)_x+(vu_y)_y.
\]
由 Green 公式，
\begin{align*}
\iint_D \left[(vu_x)_x+(vu_y)_y\right]\mathrm{d}x\,\mathrm{d}y=\oint_C (-vu_y)\,\mathrm{d}x+(vu_x)\,\mathrm{d}y=\oint_C v(u_x\,\mathrm{d}y-u_y\,\mathrm{d}x).
\end{align*}
又因为 $C$ 取正向时外单位法向量为
\[
\vec n=\left(\frac{\mathrm{d}y}{\mathrm{d}s},-\frac{\mathrm{d}x}{\mathrm{d}s}\right),
\]
所以
\[
\frac{\partial u}{\partial n}\,\mathrm{d}s
=\nabla u\cdot\vec n\,\mathrm{d}s
=u_x\,\mathrm{d}y-u_y\,\mathrm{d}x.
\]
因此
\[
\iint_D v\Delta u\,\mathrm{d}x\,\mathrm{d}y
+\iint_D (u_xv_x+u_yv_y)\,\mathrm{d}x\,\mathrm{d}y
=\oint_C v\frac{\partial u}{\partial n}\,\mathrm{d}s.
\]
\end{solution}


\begin{theorem}[平面 Green 第二公式（二维）]
设平面闭区域 $D$ 的边界 $C=\partial D$ 由有限条分段光滑闭曲线组成，并取正向。若函数 $u(x,y),v(x,y)$ 在包含 $D$ 的开区域内具有连续二阶偏导数，则
\begin{equation}\label{eq:green-second}
\iint_D (v\Delta u - u\Delta v)\,\mathrm{d}x\,\mathrm{d}y = \oint_C \left(v\frac{\partial u}{\partial n}-u\frac{\partial v}{\partial n}\right)\mathrm{d}s,
\end{equation}
其中 $\Delta u = \dfrac{\partial^2 u}{\partial x^2}+\dfrac{\partial^2 u}{\partial y^2}$ 为 Laplace 算子。
\end{theorem}
\begin{solution}
\textbf{从左往右推。}
从左边的二重积分出发，先把被积函数展开：
\begin{align*}
v\Delta u-u\Delta v=v(u_{xx}+u_{yy})-u(v_{xx}+v_{yy})=(vu_x-uv_x)_x+(vu_y-uv_y)_y.
\end{align*}
上式最后一步是因为
\[
(vu_x-uv_x)_x=v_xu_x+vu_{xx}-u_xv_x-uv_{xx}=vu_{xx}-uv_{xx},
\]
\[
(vu_y-uv_y)_y=v_yu_y+vu_{yy}-u_yv_y-uv_{yy}=vu_{yy}-uv_{yy}.
\]
于是
\[
\iint_D (v\Delta u-u\Delta v)\,\mathrm{d}x\,\mathrm{d}y
=\iint_D\left[(vu_x-uv_x)_x+(vu_y-uv_y)_y\right]\mathrm{d}x\,\mathrm{d}y.
\]
为了套 Green 公式，令
\[
P=uv_y-vu_y,\qquad Q=vu_x-uv_x.
\]
则
\[
Q_x-P_y=(vu_x-uv_x)_x+(vu_y-uv_y)_y.
\]
所以
\begin{align*}
\iint_D (v\Delta u-u\Delta v)\,\mathrm{d}x\,\mathrm{d}y
&=\oint_C (uv_y-vu_y)\,\mathrm{d}x+(vu_x-uv_x)\,\mathrm{d}y\\
&=\oint_C\left[v(u_x\,\mathrm{d}y-u_y\,\mathrm{d}x)-u(v_x\,\mathrm{d}y-v_y\,\mathrm{d}x)\right].
\end{align*}
又由
\[
\frac{\partial u}{\partial n}\,\mathrm{d}s=u_x\,\mathrm{d}y-u_y\,\mathrm{d}x,\qquad
\frac{\partial v}{\partial n}\,\mathrm{d}s=v_x\,\mathrm{d}y-v_y\,\mathrm{d}x,
\]
得到
\[
\iint_D (v\Delta u-u\Delta v)\,\mathrm{d}x\,\mathrm{d}y
=\oint_C\left(v\frac{\partial u}{\partial n}-u\frac{\partial v}{\partial n}\right)\mathrm{d}s.
\]
\end{solution}

\begin{solution}
\textbf{从右往左推。}
从右边的边界积分出发，先把法向导数改写成 $\mathrm{d}x,\mathrm{d}y$：
\[
\frac{\partial u}{\partial n}\,\mathrm{d}s=u_x\,\mathrm{d}y-u_y\,\mathrm{d}x,\qquad
\frac{\partial v}{\partial n}\,\mathrm{d}s=v_x\,\mathrm{d}y-v_y\,\mathrm{d}x.
\]
因此
\begin{align*}
&\oint_C\left(v\frac{\partial u}{\partial n}-u\frac{\partial v}{\partial n}\right)\mathrm{d}s=\oint_C\left[v(u_x\,\mathrm{d}y-u_y\,\mathrm{d}x)-u(v_x\,\mathrm{d}y-v_y\,\mathrm{d}x)\right]\\
&=\oint_C (uv_y-vu_y)\,\mathrm{d}x+(vu_x-uv_x)\,\mathrm{d}y.
\end{align*}
令
\[
P=uv_y-vu_y,\qquad Q=vu_x-uv_x.
\]
由 Green 公式，
\begin{align*}
&\oint_C (uv_y-vu_y)\,\mathrm{d}x+(vu_x-uv_x)\,\mathrm{d}y=\iint_D\left[(vu_x-uv_x)_x-(uv_y-vu_y)_y\right]\mathrm{d}x\,\mathrm{d}y\\
&=\iint_D\left[(v_xu_x+vu_{xx}-u_xv_x-uv_{xx})-(u_yv_y+uv_{yy}-v_yu_y-vu_{yy})\right]\mathrm{d}x\,\mathrm{d}y\\
&=\iint_D\left[vu_{xx}-uv_{xx}-uv_{yy}+vu_{yy}\right]\mathrm{d}x\,\mathrm{d}y=\iint_D (v\Delta u-u\Delta v)\,\mathrm{d}x\,\mathrm{d}y.
\end{align*}
这就从右边推回到了左边。
\end{solution}

\begin{exercise}[调和函数与 Green 第二公式]
设 $\Gamma$ 是光滑的封闭曲线，$D$ 是由 $\Gamma$ 所围成的区域。若 $u$ 在 $D$ 内调和，即 $\Delta u=0$ 在 $D$ 内成立，并且所需偏导在边界附近连续，证明：
\begin{compactenum}
    \item 当 $(x,y)\ne (x_0,y_0)$ 时，\(\displaystyle \Delta \ln r=0\)，其中
    \[
    r=\sqrt{(x-x_0)^2+(y-y_0)^2}
    \]
    是定点 $(x_0,y_0)$ 与 $D$ 上动点 $(x,y)$ 的距离；
    \item 对任意 $(x_0,y_0)\in D$，有
    \[
    u(x_0,y_0)=\frac{1}{2\pi}\int_{\Gamma}\left(u\frac{\partial \ln r}{\partial n}-\ln r\frac{\partial u}{\partial n}\right)\,\mathrm{d}s,
    \]
    其中 $\vec n$ 为 $\Gamma$ 的单位外法向。
\end{compactenum}
\end{exercise}
\begin{solution}
\begin{compactenum}
    \item 记
    \[
    r^2=(x-x_0)^2+(y-y_0)^2,\qquad
    \ln r=\frac12\ln\bigl[(x-x_0)^2+(y-y_0)^2\bigr].
    \]
    当 $(x,y)\ne(x_0,y_0)$ 时，$r\ne0$。先求一阶偏导：
    \[
    \frac{\partial \ln r}{\partial x}=\frac{x-x_0}{r^2},\qquad
    \frac{\partial \ln r}{\partial y}=\frac{y-y_0}{r^2}.
    \]
    再求二阶偏导：
    \[
    \frac{\partial^2 \ln r}{\partial x^2}
    =\frac{r^2-2(x-x_0)^2}{r^4},\qquad
    \frac{\partial^2 \ln r}{\partial y^2}
    =\frac{r^2-2(y-y_0)^2}{r^4}.
    \]
    因此
    \[
    \Delta\ln r
    =\frac{r^2-2(x-x_0)^2}{r^4}
    +\frac{r^2-2(y-y_0)^2}{r^4}
    =\frac{2r^2-2\bigl[(x-x_0)^2+(y-y_0)^2\bigr]}{r^4}=0.
    \]

    \item 设 $P_0=(x_0,y_0)$，并记 $w=\ln r$。由于 $w$ 在 $P_0$ 处没有定义，不能把 $P_0$ 直接放进积分区域。取充分小的 $\varepsilon>0$，使闭圆盘 $\overline{B_\varepsilon(P_0)}\subset D$，并记
    \[
    C_\varepsilon:\ r=\varepsilon,\qquad
    D_\varepsilon=D\setminus \overline{B_\varepsilon(P_0)}.
    \]
    在挖去小圆后的区域 $D_\varepsilon$ 上，$u$ 与 $w$ 都二阶连续可微，且
    \[
    \Delta u=0,\qquad \Delta w=0.
    \]

    现在不直接套 Green 第二公式，而是模仿它的证明，从边界上的方向导数拆起。对 $D_\varepsilon$ 的正向边界 $\partial D_\varepsilon$，有
    \[
    \frac{\partial w}{\partial n}\,\mathrm{d}s=w_x\,\mathrm{d}y-w_y\,\mathrm{d}x,\qquad
    \frac{\partial u}{\partial n}\,\mathrm{d}s=u_x\,\mathrm{d}y-u_y\,\mathrm{d}x.
    \]
    因此
    \begin{align*}
    &\int_{\partial D_\varepsilon}\left(u\frac{\partial w}{\partial n}-w\frac{\partial u}{\partial n}\right)\mathrm{d}s\\
    &=\oint_{\partial D_\varepsilon}\left[u(w_x\,\mathrm{d}y-w_y\,\mathrm{d}x)-w(u_x\,\mathrm{d}y-u_y\,\mathrm{d}x)\right]\\
    &=\oint_{\partial D_\varepsilon}(wu_y-uw_y)\,\mathrm{d}x+(uw_x-wu_x)\,\mathrm{d}y.
    \end{align*}
    令
    \[
    P=wu_y-uw_y,\qquad Q=uw_x-wu_x.
    \]
    用普通 Green 公式：
    \begin{align*}
    &\oint_{\partial D_\varepsilon}(wu_y-uw_y)\,\mathrm{d}x+(uw_x-wu_x)\,\mathrm{d}y\\
    &=\iint_{D_\varepsilon}\left[(uw_x-wu_x)_x-(wu_y-uw_y)_y\right]\mathrm{d}x\,\mathrm{d}y\\
    &=\iint_{D_\varepsilon}\left[(u_xw_x+uw_{xx}-w_xu_x-wu_{xx})\right.\\
    &\qquad\qquad\left.-(w_yu_y+wu_{yy}-u_yw_y-uw_{yy})\right]\mathrm{d}x\,\mathrm{d}y\\
    &=\iint_{D_\varepsilon}\left[u(w_{xx}+w_{yy})-w(u_{xx}+u_{yy})\right]\mathrm{d}x\,\mathrm{d}y\\
    &=\iint_{D_\varepsilon}(u\Delta w-w\Delta u)\,\mathrm{d}x\,\mathrm{d}y=0.
    \end{align*}
    所以
    \[
    0=\int_{\partial D_\varepsilon}\left(u\frac{\partial w}{\partial n}-w\frac{\partial u}{\partial n}\right)\mathrm{d}s.
    \]
    把边界拆成外边界 $\Gamma$ 与内边界 $C_\varepsilon$，即
    \[
    0=\int_{\Gamma}\left(u\frac{\partial\ln r}{\partial n}-\ln r\frac{\partial u}{\partial n}\right)\mathrm{d}s
    +\int_{C_\varepsilon}\left(u\frac{\partial\ln r}{\partial n}-\ln r\frac{\partial u}{\partial n}\right)\mathrm{d}s.
    \]

    下面计算小圆 $C_\varepsilon$ 上的积分。对挖洞后的区域 $D_\varepsilon$ 来说，$C_\varepsilon$ 上的外法向指向圆心，所以
    \[
    \frac{\partial}{\partial n}=-\frac{\partial}{\partial r},\qquad
    \frac{\partial\ln r}{\partial n}=-\frac1r=-\frac1\varepsilon.
    \]
    于是
    \begin{align*}
    \int_{C_\varepsilon}\left(u\frac{\partial\ln r}{\partial n}-\ln r\frac{\partial u}{\partial n}\right)\,\mathrm{d}s
    &=\int_{C_\varepsilon}\left(-\frac{u}{\varepsilon}+\ln\varepsilon\,\frac{\partial u}{\partial r}\right)\,\mathrm{d}s\\
    &=-\int_0^{2\pi}u(x_0+\varepsilon\cos\theta,y_0+\varepsilon\sin\theta)\,\mathrm{d}\theta\\
    &\quad+\varepsilon\ln\varepsilon\int_0^{2\pi}\frac{\partial u}{\partial r}(x_0+\varepsilon\cos\theta,y_0+\varepsilon\sin\theta)\,\mathrm{d}\theta.
    \end{align*}
    由 $u$ 的连续性，
    \[
    \int_0^{2\pi}u(x_0+\varepsilon\cos\theta,y_0+\varepsilon\sin\theta)\,\mathrm{d}\theta
    \to 2\pi u(x_0,y_0).
    \]
    又因为 $\dfrac{\partial u}{\partial r}$ 在 $P_0$ 附近有界，且 $\varepsilon\ln\varepsilon\to0$，所以
    \[
    \varepsilon\ln\varepsilon\int_0^{2\pi}\frac{\partial u}{\partial r}(x_0+\varepsilon\cos\theta,y_0+\varepsilon\sin\theta)\,\mathrm{d}\theta\to0.
    \]
    因而
    \[
    \int_{C_\varepsilon}\left(u\frac{\partial\ln r}{\partial n}-\ln r\frac{\partial u}{\partial n}\right)\,\mathrm{d}s
    \to -2\pi u(x_0,y_0).
    \]
    令 $\varepsilon\to0$，由
    \[
    0=\int_{\Gamma}\left(u\frac{\partial\ln r}{\partial n}-\ln r\frac{\partial u}{\partial n}\right)\,\mathrm{d}s
    +\int_{C_\varepsilon}\left(u\frac{\partial\ln r}{\partial n}-\ln r\frac{\partial u}{\partial n}\right)\,\mathrm{d}s
    \]
    得
    \[
    \int_{\Gamma}\left(u\frac{\partial\ln r}{\partial n}-\ln r\frac{\partial u}{\partial n}\right)\,\mathrm{d}s
    =2\pi u(x_0,y_0).
    \]
    即
    \[
    u(x_0,y_0)=\frac{1}{2\pi}\int_{\Gamma}\left(u\frac{\partial \ln r}{\partial n}-\ln r\frac{\partial u}{\partial n}\right)\,\mathrm{d}s.
    \]
\end{compactenum}
\end{solution}

\begin{example}[稳定平面流体的流量]\label{ex:fluid-flow}
设稳定平面流体的速度场为 $\vec{v}=P(x,y)\vec{i}+Q(x,y)\vec{j}$，密度为 $\rho=1$。求单位时间内从区域 $D$ 的正向边界 $C$ 流出的质量。
\end{example}
\begin{solution}
沿正向边界行进时，区域在左手边，外法向量在右手边。设弧长参数为 $s$，切向量和外单位法向量为
\[
\vec T=\left(\frac{\mathrm{d}x}{\mathrm{d}s},\frac{\mathrm{d}y}{\mathrm{d}s}\right),~~
\vec n=\left(\frac{\mathrm{d}y}{\mathrm{d}s},-\frac{\mathrm{d}x}{\mathrm{d}s}\right).
\]
故边界流出量为
\[
M=\oint_C \vec v\cdot\vec n\,\mathrm{d}s
=\oint_C\left(P\frac{\mathrm{d}y}{\mathrm{d}s}-Q\frac{\mathrm{d}x}{\mathrm{d}s}\right)\mathrm{d}s
=\oint_C P\,\mathrm{d}y-Q\,\mathrm{d}x.
\]
再由 Green 公式的通量形式，
\[
M=\iint_D (P_x+Q_y)\,\mathrm{d}x\,\mathrm{d}y.
\]
所以平面流体从边界流出的总量，等于区域内部散度 $P_x+Q_y$ 的总和。
\end{solution}

\begin{example}[椭圆的面积]\label{ex:ellipse-area}
求椭圆 $\dfrac{x^2}{a^2}+\dfrac{y^2}{b^2}=1$ 所围区域的面积。
\end{example}
\begin{solution}
题目要求闭曲线围成的面积，直接调用 Green 面积公式
\(\displaystyle S=\frac12\oint_\Gamma (x\,\mathrm{d}y-y\,\mathrm{d}x)\)。

边界是椭圆，最稳的是参数化边界，而不是回头做二重积分。

取 \(\displaystyle x=a\cos t,\qquad y=b\sin t,\qquad 0\le t\le 2\pi\)，则 \(\displaystyle \mathrm{d}x=-a\sin t\,\mathrm{d}t,\qquad \mathrm{d}y=b\cos t\,\mathrm{d}t\)。
代入面积公式：
于是 \(\displaystyle S=\frac12\int_0^{2\pi}ab(\cos^2 t+\sin^2 t)\,\mathrm{d}t=\frac12\int_0^{2\pi}ab\,\mathrm{d}t=\pi ab\)。

参数方向是逆时针正向，所以面积公式直接可用，结果为 \(\displaystyle S=\pi ab\)。
\end{solution}

\begin{example}\label{ex:green-zero}
证明 $\displaystyle\oint_C (2xy+\cos x)\,\mathrm{d}x+(x^2+\sin y)\,\mathrm{d}y = 0$，其中 $C$ 为任意分段光滑的简单闭曲线。
\end{example}
\begin{solution}
题目给任意简单闭曲线上的第二类曲线积分，而且 $P,Q$ 在全平面光滑，优先查 Green。

只要算出 $Q_x-P_y$，若为 $0$，闭路积分立刻为 $0$。

令 \(\displaystyle P=2xy+\cos x,\qquad Q=x^2+\sin y\)。则
\[
\frac{\partial Q}{\partial x}=2x,\qquad
\frac{\partial P}{\partial y}=2x,\qquad
\frac{\partial Q}{\partial x}-\frac{\partial P}{\partial y}=0.
\]
由 Green 公式，
\[
\oint_C (2xy+\cos x)\,\mathrm{d}x+(x^2+\sin y)\,\mathrm{d}y
=\iint_D (Q_x-P_y)\,\mathrm{d}x\,\mathrm{d}y=0.
\]

区域内无奇点，$C$ 又是闭曲线，所以 Green 可直接使用。
\end{solution}

\begin{example}\label{ex:any-curve-origin}
计算 $\displaystyle\oint_\Gamma \frac{x\,\mathrm{d}y - y\,\mathrm{d}x}{x^2+y^2}$，其中 $\Gamma$ 为任意包含原点的分段光滑简单闭曲线，取正向。
\end{example}
\begin{solution}
闭曲线积分，奇点原点在曲线内部，边界又是任意曲线，不能直接参数化。

用挖洞法，把原点周围挖去一个小圆，再在环域上套 Green。

以原点为圆心取充分小的圆 \(\displaystyle C_\varepsilon:\ x^2+y^2=\varepsilon^2\)，并记 $D'$ 为 $\Gamma$ 与 $C_\varepsilon$ 之间的环域。整体边界按“外逆内顺”记为 \(\displaystyle \partial D'=\Gamma^+-C_\varepsilon^+\)。在 $D'$ 上 \(\displaystyle P=\frac{-y}{x^2+y^2},\ Q=\frac{x}{x^2+y^2},\ Q_x-P_y=0\)。
由 Green 公式得
\[
\oint_{\Gamma}\frac{x\,\mathrm{d}y-y\,\mathrm{d}x}{x^2+y^2}
-\oint_{C_\varepsilon^+}\frac{x\,\mathrm{d}y-y\,\mathrm{d}x}{x^2+y^2}
=0.
\]
所以 \(\displaystyle \oint_{\Gamma}\frac{x\,\mathrm{d}y-y\,\mathrm{d}x}{x^2+y^2}=\oint_{C_\varepsilon^+}\frac{x\,\mathrm{d}y-y\,\mathrm{d}x}{x^2+y^2}\)。而右端正是例 \ref{ex:unit-circle} 的同型积分，值为 $2\pi$。故 \(\displaystyle \oint_{\Gamma}\frac{x\,\mathrm{d}y-y\,\mathrm{d}x}{x^2+y^2}=2\pi\)。

关键只查两件事：原点已被小圆挖掉；内边界写成顺时针，即公式里体现为 $-\oint_{C_\varepsilon^+}$。
\end{solution}

\begin{example}\label{ex:area-green}
利用 Green 公式证明以下结果：

\textbf{(1)} 区域 $D$ 的面积可表示为 \(\displaystyle S(D) = \pm\int_C x\,\mathrm{d}y\)，其中 $C$ 为 $D$ 的边界曲线，取适当的定向。

\textbf{(2)} 设 $\xi=\xi(x,y),\;\eta=\eta(x,y)$ 为 $C^1$ 类变换，且 Jacobi 行列式 $\dfrac{\partial(\xi,\eta)}{\partial(x,y)}\neq 0$，证明 \(\displaystyle \mathrm{d}x\,\mathrm{d}y = \left|\frac{\partial(x,y)}{\partial(\xi,\eta)}\right|\mathrm{d}\xi\,\mathrm{d}\eta\)。
\end{example}

\begin{proof}

\textbf{(1)} 在 Green 公式 \eqref{eq:green} 中取 $P=0,\;Q=x$，则 \(\displaystyle \oint_C x\,\mathrm{d}y = \iint_D \frac{\partial x}{\partial x}\,\mathrm{d}x\,\mathrm{d}y = \iint_D 1\,\mathrm{d}x\,\mathrm{d}y = S(D)\)。
因此 $S(D)=\displaystyle\int_{C^+} x\,\mathrm{d}y$，其中 $C^+$ 取正向（逆时针）。若取负向则加负号。

\textbf{(2)} 作变量替换 $x=x(\xi,\eta),\;y=y(\xi,\eta)$。设 $D_{\xi\eta}$ 为 $(\xi,\eta)$ 平面中对应的区域，其边界为 $C_{\xi\eta}$。由 (1) 的结论，\(\displaystyle S(D) = \int_{C^+} x\,\mathrm{d}y\)。

对右端作变量替换。由于 $y=y(\xi,\eta)$，有
\[
\mathrm{d}y = \frac{\partial y}{\partial \xi}\,\mathrm{d}\xi + \frac{\partial y}{\partial \eta}\,\mathrm{d}\eta,\qquad
x\,\mathrm{d}y = x\frac{\partial y}{\partial \xi}\,\mathrm{d}\xi + x\frac{\partial y}{\partial \eta}\,\mathrm{d}\eta.
\]
在 $(\xi,\eta)$ 平面上对 $C_{\xi\eta}$ 应用 Green 公式，取 $P=x\dfrac{\partial y}{\partial \xi},\;Q=x\dfrac{\partial y}{\partial \eta}$：
\[
\int_{C_{\xi\eta}^+} x\frac{\partial y}{\partial \xi}\,\mathrm{d}\xi + x\frac{\partial y}{\partial \eta}\,\mathrm{d}\eta
= \iint_{D_{\xi\eta}} \left[\frac{\partial}{\partial \xi}\left(x\frac{\partial y}{\partial \eta}\right)-\frac{\partial}{\partial \eta}\left(x\frac{\partial y}{\partial \xi}\right)\right]\mathrm{d}\xi\,\mathrm{d}\eta.
\]
展开被积函数：
\begin{align*}
\frac{\partial}{\partial \xi}\left(x\frac{\partial y}{\partial \eta}\right)
-\frac{\partial}{\partial \eta}\left(x\frac{\partial y}{\partial \xi}\right)
&= \frac{\partial x}{\partial \xi}\frac{\partial y}{\partial \eta}
+x\frac{\partial^2 y}{\partial \xi\partial \eta}  -\frac{\partial x}{\partial \eta}\frac{\partial y}{\partial \xi}
-x\frac{\partial^2 y}{\partial \eta\partial \xi} \\
&= \frac{\partial x}{\partial \xi}\frac{\partial y}{\partial \eta}-\frac{\partial x}{\partial \eta}\frac{\partial y}{\partial \xi}
= \frac{\partial(x,y)}{\partial(\xi,\eta)}.
\end{align*}

因此 \(\displaystyle S(D) = \iint_{D_{\xi\eta}} \frac{\partial(x,y)}{\partial(\xi,\eta)}\,\mathrm{d}\xi\,\mathrm{d}\eta\)。

另一方面，由二重积分的变量替换公式，$S(D)=\displaystyle\iint_D \mathrm{d}x\,\mathrm{d}y = \iint_{D_{\xi\eta}}\left|\dfrac{\partial(x,y)}{\partial(\xi,\eta)}\right|\mathrm{d}\xi\,\mathrm{d}\eta$。

比较两式，得 $\mathrm{d}x\,\mathrm{d}y = \left|\dfrac{\partial(x,y)}{\partial(\xi,\eta)}\right|\mathrm{d}\xi\,\mathrm{d}\eta$。
\end{proof}

\subsubsection{空间闭曲线：Stokes 公式与换面}

\begin{theorem}[Stokes公式]\label{thm:stokes}
设 $\Sigma$ 是有限块二阶连续可微光滑曲面拼接而成的定向曲面，边界 $\partial\Sigma$ 按右手法则取正向；函数 $P,Q,R$ 在包含 $\Sigma$ 的某个开区域内具有连续一阶偏导数，则
\begin{align*}
\int_{\partial\Sigma} P\,\mathrm{d}x+Q\,\mathrm{d}y+R\,\mathrm{d}z
={}& \iint_\Sigma \left(\frac{\partial R}{\partial y}-\frac{\partial Q}{\partial z}\right)\mathrm{d}y\,\mathrm{d}z+ \left(\frac{\partial P}{\partial z}-\frac{\partial R}{\partial x}\right)\mathrm{d}z\,\mathrm{d}x+ \left(\frac{\partial Q}{\partial x}-\frac{\partial P}{\partial y}\right)\mathrm{d}x\,\mathrm{d}y.
\end{align*}
Stokes公式可以用行列式记法简洁地表示为
\begin{align*}
\int_{\partial\Sigma}P\,\mathrm{d}x+Q\,\mathrm{d}y+R\,\mathrm{d}z
&= \iint_\Sigma
\begin{vmatrix}
\mathrm{d}y\,\mathrm{d}z & \mathrm{d}z\,\mathrm{d}x & \mathrm{d}x\,\mathrm{d}y \\
\dfrac{\partial}{\partial x} & \dfrac{\partial}{\partial y} & \dfrac{\partial}{\partial z} \\
P & Q & R
\end{vmatrix},
\end{align*}
或者用方向余弦表示为
\begin{align*}
\int_{\partial\Sigma}P\,\mathrm{d}x+Q\,\mathrm{d}y+R\,\mathrm{d}z
&= \iint_\Sigma
\begin{vmatrix}
\cos\alpha & \cos\beta & \cos\gamma \\
\dfrac{\partial}{\partial x} & \dfrac{\partial}{\partial y} & \dfrac{\partial}{\partial z} \\
P & Q & R
\end{vmatrix}\mathrm{d}S,
\end{align*}
其中$(\cos\alpha,\cos\beta,\cos\gamma)$为$\Sigma$的单位法向量.
\end{theorem}

\begin{example}
设$\Sigma$为平面$x+y+z=1$在第一卦限的部分，取上侧，$\vec{F}=(y^2,z^2,x^2)$。求力$\vec{F}$沿$\partial\Sigma$所做的功.
\end{example}
\begin{solution}
由Stokes公式，所求功为
\begin{align*}
W &= \int_{\partial\Sigma}y^2\,\mathrm{d}x+z^2\,\mathrm{d}y+x^2\,\mathrm{d}z\\
&= \iint_\Sigma (0-2z)\,\mathrm{d}y\,\mathrm{d}z+(0-2x)\,\mathrm{d}z\,\mathrm{d}x+(0-2y)\,\mathrm{d}x\,\mathrm{d}y \\
&= -2\iint_\Sigma z\,\mathrm{d}y\,\mathrm{d}z+x\,\mathrm{d}z\,\mathrm{d}x+y\,\mathrm{d}x\,\mathrm{d}y.
\end{align*}
这里三个积分项相等，记这个公共值为 $J$，则 \(\displaystyle W=-2(J+J+J)=-6J\)。
取最方便的一项 $\iint_\Sigma y\,\mathrm{d}x\,\mathrm{d}y$ 来计算。将$\Sigma$投影到$xOy$平面得
\(\displaystyle D_{xy}=\{(x,y)\mid x\geq 0,y\geq 0,x+y\leq 1\}\)。
且法向量向上，所以
\[
\iint_\Sigma y\,\mathrm{d}x\,\mathrm{d}y
=\iint_{D_{xy}}y\,\mathrm{d}x\,\mathrm{d}y
=\int_0^1\mathrm{d}x\int_0^{1-x}y\,\mathrm{d}y.
\]
计算得 \(\displaystyle \int_0^1\frac{(1-x)^2}{2}\mathrm{d}x=\frac{1}{6}\)。于是 $J=\dfrac{1}{6}$，从而 \(\displaystyle W=-6\cdot \frac{1}{6}=-1\)。
\end{solution}

\begin{example}
计算$\displaystyle\oint_L y\,\mathrm{d}x+z\,\mathrm{d}y+x\,\mathrm{d}z$，其中$L$为球面$x^2+y^2+z^2=a^2$与平面$x+y+z=0$的交线，从$x$轴正向看去为逆时针方向.
\end{example}
\begin{solution}
题目说“从 $x$ 轴正向看去为逆时针方向”，先把它翻译成法向判断：沿着法向看过去时，边界应是逆时针。由右手法则可知，此时法向量必须取 $x$ 分量为正的那一侧；平面 $x+y+z=0$ 的法向量可取 $\vec n=\dfrac{1}{\sqrt{3}}(1,1,1)$，它的三个方向余弦都为正，所以三个有向投影面积元都同号。于是取该法向，并记$\Sigma$为$L$所围的平面$x+y+z=0$上的圆盘。
\begin{align*}
\oint_L y\,\mathrm{d}x+z\,\mathrm{d}y+x\,\mathrm{d}z
&= \iint_\Sigma(0-1)\,\mathrm{d}y\,\mathrm{d}z+(0-1)\,\mathrm{d}z\,\mathrm{d}x+(0-1)\,\mathrm{d}x\,\mathrm{d}y \\
&= -\iint_\Sigma\mathrm{d}y\,\mathrm{d}z+\mathrm{d}z\,\mathrm{d}x+\mathrm{d}x\,\mathrm{d}y.
\end{align*}
平面$x+y+z=0$的单位法向量就是$\vec{n}=\dfrac{1}{\sqrt{3}}(1,1,1)$，故
\[
\mathrm{d}y\,\mathrm{d}z=\dfrac{1}{\sqrt{3}}\mathrm{d}S,\qquad
\mathrm{d}z\,\mathrm{d}x=\dfrac{1}{\sqrt{3}}\mathrm{d}S,\qquad
\mathrm{d}x\,\mathrm{d}y=\dfrac{1}{\sqrt{3}}\mathrm{d}S.
\]
因此 \(\displaystyle -\iint_\Sigma\frac{3}{\sqrt{3}}\mathrm{d}S=-\sqrt{3}\,\mathrm{Area}(\Sigma)\)。
又因为所取积分曲面就是球面与该平面截出的圆盘，而平面经过原点，所以这条交线是球的大圆，圆盘半径为 $a$。故$\mathrm{Area}(\Sigma)=\pi a^2$，所求积分为$-\sqrt{3}\pi a^2$.
\end{solution}

\begin{definition}[方向旋量]\label{def:directional-curl}
设 $\vec F$ 为空间向量场，$p_0$ 为场中一点，$\vec n$ 为单位向量。取过 $p_0$ 且法向为 $\vec n$ 的一小块有向曲面片 $D_\varepsilon$，其边界 $\partial D_\varepsilon$ 按右手法则取正向，记 $\vec e_T$ 为边界的单位切向量。若极限存在，则称
\[
\lim_{S(D_\varepsilon)\to 0}
\frac{1}{S(D_\varepsilon)}
\int_{\partial D_\varepsilon}\vec F\cdot\vec e_T\,\mathrm{d}s
\]
为 $\vec F$ 在点 $p_0$ 处沿方向 $\vec n$ 的\textbf{方向旋量}，其中 $S(D_\varepsilon)$ 表示曲面片 $D_\varepsilon$ 的面积。
\end{definition}

\begin{definition}[旋度]\label{def:curl}
设$\vec{F}=(P,Q,R)$，定义$\vec{F}$的\textbf{旋度}为
\[
\mathrm{rot}\,\vec{F} = \nabla\times\vec{F} = \left(\frac{\partial R}{\partial y}-\frac{\partial Q}{\partial z},\;\frac{\partial P}{\partial z}-\frac{\partial R}{\partial x},\;\frac{\partial Q}{\partial x}-\frac{\partial P}{\partial y}\right)=\begin{vmatrix} \vec{i} & \vec{j} & \vec{k} \\[6pt] \dfrac{\partial}{\partial x} & \dfrac{\partial}{\partial y} & \dfrac{\partial}{\partial z} \\[6pt] P & Q & R \end{vmatrix},
\]
由 Stokes 公式可以推导出方向旋量与旋度的关系：设 $p_0$ 为场中一点，$\vec n$ 为单位向量，则在 $p_0$ 处沿 $\vec n$ 方向的方向旋量为
\[
\lim_{S(D_\varepsilon)\to 0}
\frac{1}{S(D_\varepsilon)}
\int_{\partial D_\varepsilon}\vec{F}\cdot\vec{e}_T\,\mathrm{d}s
=(\mathrm{rot}\,\vec{F}(p_0))\cdot\vec{n}.
\]
\end{definition}

\begin{note}[为什么方向旋量要写成 $(\operatorname{rot}\vec F)\cdot \vec n$]
旋度记录的是各个方向上的局部旋转趋势；当只关心法向为 $\vec n$ 的那一小片时，就只取它在 $\vec n$ 方向上的分量，因此自然写成 $(\operatorname{rot}\vec F)\cdot \vec n$。这和方向导数写成 $\nabla f\cdot\vec l$ 的道理完全一样。
\end{note}

\begin{proposition}[旋度的几何意义]
旋度的物理意义：当小曲面片的法向 $\vec n$ 与 $\mathrm{rot}\,\vec F$ 同向时，方向旋量达到最大，最大值为 \(\|\mathrm{rot}\,\vec F\|\)；当 $\vec n$ 与 $\mathrm{rot}\,\vec F$ 垂直时，方向旋量为 $0$。由 \(\displaystyle \text{沿 }\vec n\text{ 的方向旋量}=\|\mathrm{rot}\,\vec F\|\cos\theta\) 可知，$\theta=0$ 时取到最大值，也就是 $\vec n$ 与旋度同向。因此旋度不仅给出三个分量，还直接指出了最强局部旋转轴的方向，而其模长就是最大单位面积环量。
\end{proposition}

\begin{example}[旋度的性质]
设$C$为常数，$\vec{F},\vec{F}_1,\vec{F}_2$为向量场，$\varphi$为数量场，则旋度具有以下性质：
\begin{compactenum}
\item $\nabla\times(C\vec{F})=C(\nabla\times\vec{F})$；
\item $\nabla\times(\vec{F}_1+\vec{F}_2)=\nabla\times\vec{F}_1+\nabla\times\vec{F}_2$；
\item $\nabla\times(\varphi\vec{F})=\nabla\varphi\times\vec{F}+\varphi(\nabla\times\vec{F})$；
\item $\nabla\cdot(\vec{F}_1\times\vec{F}_2)=(\nabla\times\vec{F}_1)\cdot\vec{F}_2-(\nabla\times\vec{F}_2)\cdot\vec{F}_1$.
\end{compactenum}
\end{example}
\begin{proof}
\textbf{性质(3)的证明.} 设$\vec{F}=(P,Q,R)$，则
\begin{align*}
\nabla\times(\varphi\vec{F}) &= \begin{vmatrix} \vec{i} & \vec{j} & \vec{k} \\[6pt] \dfrac{\partial}{\partial x} & \dfrac{\partial}{\partial y} & \dfrac{\partial}{\partial z} \\[6pt] \varphi P & \varphi Q & \varphi R \end{vmatrix} = \left(\frac{\partial(\varphi R)}{\partial y}-\frac{\partial(\varphi Q)}{\partial z},\;\frac{\partial(\varphi P)}{\partial z}-\frac{\partial(\varphi R)}{\partial x},\;\frac{\partial(\varphi Q)}{\partial x}-\frac{\partial(\varphi P)}{\partial y}\right).
\end{align*}
展开各分量，例如第一分量：
\begin{align*}
\frac{\partial(\varphi R)}{\partial y}-\frac{\partial(\varphi Q)}{\partial z}
&= \frac{\partial\varphi}{\partial y}R+\varphi\frac{\partial R}{\partial y}
-\frac{\partial\varphi}{\partial z}Q-\varphi\frac{\partial Q}{\partial z} = \varphi\left(\frac{\partial R}{\partial y}-\frac{\partial Q}{\partial z}\right)
+\left(\frac{\partial\varphi}{\partial y}R-\frac{\partial\varphi}{\partial z}Q\right).
\end{align*}
注意到$\nabla\varphi\times\vec{F}$的第一分量为$\dfrac{\partial\varphi}{\partial y}R-\dfrac{\partial\varphi}{\partial z}Q$，其余分量类似，故 \(\displaystyle \nabla\times(\varphi\vec{F})=\nabla\varphi\times\vec{F}+\varphi(\nabla\times\vec{F})\)。

\textbf{性质(4)的证明.} 设$\vec{F}_1=(P_1,Q_1,R_1)$，$\vec{F}_2=(P_2,Q_2,R_2)$，则
\(\displaystyle \vec{F}_1\times\vec{F}_2=(Q_1R_2-R_1Q_2,\;R_1P_2-P_1R_2,\;P_1Q_2-Q_1P_2)\)。
故
\begin{align*}
\nabla\cdot(\vec{F}_1\times\vec{F}_2) &= \frac{\partial}{\partial x}(Q_1R_2-R_1Q_2)+\frac{\partial}{\partial y}(R_1P_2-P_1R_2)+\frac{\partial}{\partial z}(P_1Q_2-Q_1P_2)\\
&= \left(\frac{\partial Q_1}{\partial x}R_2+Q_1\frac{\partial R_2}{\partial x}-\frac{\partial R_1}{\partial x}Q_2-R_1\frac{\partial Q_2}{\partial x}\right)\\
&\quad+\left(\frac{\partial R_1}{\partial y}P_2+R_1\frac{\partial P_2}{\partial y}-\frac{\partial P_1}{\partial y}R_2-P_1\frac{\partial R_2}{\partial y}\right)\\
&\quad+\left(\frac{\partial P_1}{\partial z}Q_2+P_1\frac{\partial Q_2}{\partial z}-\frac{\partial Q_1}{\partial z}P_2-Q_1\frac{\partial P_2}{\partial z}\right).
\end{align*}
将含$\vec{F}_1$偏导数的项和含$\vec{F}_2$偏导数的项分别整理：
\begin{align*}
&= \left(\frac{\partial R_1}{\partial y}-\frac{\partial Q_1}{\partial z}\right)P_2+\left(\frac{\partial P_1}{\partial z}-\frac{\partial R_1}{\partial x}\right)Q_2+\left(\frac{\partial Q_1}{\partial x}-\frac{\partial P_1}{\partial y}\right)R_2\\
&\quad-\left[\left(\frac{\partial R_2}{\partial y}-\frac{\partial Q_2}{\partial z}\right)P_1+\left(\frac{\partial P_2}{\partial z}-\frac{\partial R_2}{\partial x}\right)Q_1+\left(\frac{\partial Q_2}{\partial x}-\frac{\partial P_2}{\partial y}\right)R_1\right]\\
&= (\nabla\times\vec{F}_1)\cdot\vec{F}_2-(\nabla\times\vec{F}_2)\cdot\vec{F}_1.
\end{align*}
\end{proof}

\begin{example}[有心场的旋度]
设有心场$\vec{F}(p)=f(p)\vec{p}$，其中$p=\|\vec{p}\|>0$，$f$为可微函数. 证明$\mathrm{rot}\,\vec{F}=\vec{0}$.
\end{example}

\begin{proof}
设$\vec{p}=(x,y,z)$，则$p=\sqrt{x^2+y^2+z^2}$，$\vec{F}=(f(p)x,\,f(p)y,\,f(p)z)$.

计算旋度的第一分量：
\[
\frac{\partial(f(p)z)}{\partial y}-\frac{\partial(f(p)y)}{\partial z}
=f'(p)\frac{\partial p}{\partial y}z+f(p)\cdot 0
-f'(p)\frac{\partial p}{\partial z}y-f(p)\cdot 0.
\]
由$\dfrac{\partial p}{\partial y}=\dfrac{y}{p}$，$\dfrac{\partial p}{\partial z}=\dfrac{z}{p}$，得 \(\displaystyle f'(p)\frac{y}{p}z-f'(p)\frac{z}{p}y=0\)。由对称性，第二、三分量也为零，故$\mathrm{rot}\,\vec{F}=\vec{0}$.
\end{proof}

\begin{conclusion}[无旋场的直接结论]
\begin{compactenum}
\item 若$\mathrm{rot}\,\vec{F}=\vec{0}$，则称$\vec{F}$为\textbf{无旋场}.
\item 对任意可微径向函数 \(f(p)\)，有 \(\displaystyle \nabla\times\bigl(f(p)\vec p\bigr)=\vec0\)。特别地，\(\displaystyle \nabla\times\left(\frac{\vec p}{p^3}\right)=\vec0\ (p>0)\)，所以静电场$\vec{E}=\dfrac{q}{p^3}\vec{p}$（$p>0$）为无旋场.
\item 若各分量有连续二阶偏导，则 \(\displaystyle \nabla\cdot(\nabla\times\vec F)=0\)。这说明旋度场本身无源，可作为 Stokes 计算后的验算信号。
\end{compactenum}
\end{conclusion}

\subsection{第二类曲面积分与通量}

\begin{definition}[可定向曲面]\label{def:orientable-surface}
设正则曲面 \(\displaystyle \Sigma:\vec r=\vec r(u,v),(u,v)\in D\)。若能在整张曲面上连续选取一个单位法向量场，则称 $\Sigma$ 为\textbf{可定向曲面}或\textbf{双侧曲面}。

选定其中一个连续法向方向，就称为给曲面\textbf{定向}。被选中的一侧称为\textbf{正侧}，另一侧称为\textbf{负侧}。第二类曲面积分只能在可定向曲面上定义；若把曲面的定向整体反过来，积分值也会整体变号。
\end{definition}

\begin{property}[法向方向的判定]
(1) $\vec{r}_{u}\times\vec{r}_{v}$ 只给出了\textbf{某一个} 连续法向方向，它未必天然就是外法向或上法向。对封闭曲面，是否指向外侧需要借助几何位置或测试点来判断；对显表示曲面 $z=f(x,y)$，在参数化 $\vec r(x,y)=(x,y,f(x,y))$ 下，$\vec r_x\times\vec r_y=(-f_x,-f_y,1)$ 的第三分量为正，因此它确实指向上侧。

(2) 对椭球面这类以原点为中心的闭曲面，常用的判别方式是与位置向量作点积：若 $(\vec{r}_{\phi}\times\vec{r}_{\theta})\cdot\vec{r}(\phi,\theta)>0$，则该法向与从原点指向曲面点的方向同向，可判定它指向外侧。对于拼接曲面，若两块曲面共享一条边界，那么它们在这条公共边界上诱导出的环行方向必须相反，这样拼接后内部边界的通量才会相互抵消。常用的判据是右手法则：若右手四指沿边界正向弯曲，则拇指所指方向就是曲面的正法向。
\end{property}

\begin{example}
试确定 $\vec{r}=\vec{r}(\phi,\theta)=(a\sin\phi\cos\theta,b\sin\phi\sin\theta,c\cos\phi), \quad \phi\in[0,\pi], \theta\in[0,2\pi]$ 的外侧单位法向量。
\end{example}
\begin{solution}
由于 $\vec{n}=\epsilon\frac{\vec{r}_{\phi}\times\vec{r}_{\theta}}{\|\vec{r}_{\phi}\times\vec{r}_{\theta}\|}$。这里真正需要判断的是：$\vec r_\phi\times \vec r_\theta$ 给出的法向，到底已经指向外侧还是内侧。因为外法向在上半部分应当具有正的 $z$ 分量，所以当 $\phi\in(0,\frac{\pi}{2})$ 时，应满足 $\vec{n}\cdot(0,0,1)>0$。又
\begin{align*}
\vec{r}_{\phi}\times\vec{r}_{\theta}&=(a\cos\phi\cos\theta,b\cos\phi\sin\theta,-c\sin\phi)\times(-a\sin\phi\sin\theta,b\sin\phi\cos\theta,0)\\
&=\begin{vmatrix}\vec{i}&\vec{j}&\vec{k}\\ a\cos\phi\cos\theta&b\cos\phi\sin\theta&-c\sin\phi\\ -a\sin\phi\sin\theta&b\sin\phi\cos\theta&0\end{vmatrix}\\
&=(bc\sin^{2}\phi\cos\theta,ac\sin^{2}\phi\sin\theta,ab\cos\phi\sin\phi),\\
(\vec{r}_{\phi}\times\vec{r}_{\theta})\cdot(0,0,1)&=\frac{1}{2}ab\sin 2\phi>0, \quad \phi\in\left(0,\frac{\pi}{2}\right).
\end{align*}
即 $\epsilon=1$。
\end{solution}

\begin{theorem}[第二类曲面积分的计算]\label{thm:surface-integral-type2-calc}
设 $\vec n=(\cos\alpha,\cos\beta,\cos\gamma)$ 是指向曲面正侧的单位法向量。曲面面积元 $\mathrm{d}S$ 在三个坐标平面上的有向投影面积元定义为
\[
\mathrm{d}y\,\mathrm{d}z=\cos\alpha\,\mathrm{d}S,\qquad
\mathrm{d}z\,\mathrm{d}x=\cos\beta\,\mathrm{d}S,\qquad
\mathrm{d}x\,\mathrm{d}y=\cos\gamma\,\mathrm{d}S.
\]若 $\vec F=(P,Q,R)$，则
\begin{align*}
\iint_{\Sigma}(\vec F\cdot\vec n)\,\mathrm{d}S=\iint_{\Sigma}(P\cos\alpha+Q\cos\beta+R\cos\gamma)\,\mathrm{d}S=\iint_{\Sigma}P\,\mathrm{d}y\,\mathrm{d}z+Q\,\mathrm{d}z\,\mathrm{d}x+R\,\mathrm{d}x\,\mathrm{d}y.
\end{align*}
因此 $\mathrm{d}x\,\mathrm{d}y$ 不是普通面积大小，而是投影到 $xOy$ 平面后的有向面积；$\mathrm{d}y\,\mathrm{d}z,\mathrm{d}z\,\mathrm{d}x$ 同理。以 $\mathrm{d}x\,\mathrm{d}y$ 为例，它的符号由 $\vec n$ 与 $z$ 轴正向的夹角决定：$\gamma<\frac{\pi}{2}$ 时为正，$\gamma>\frac{\pi}{2}$ 时为负。法向一翻转，这三个有向投影面积元会同时变号。

对于参数曲面
\[
\vec r=\vec r(u,v)=(x(u,v),y(u,v),z(u,v)),
\]
有
\[
\begin{aligned}
\vec r_u\times\vec r_v=\begin{vmatrix}\vec i&\vec j&\vec k\\ x_u&y_u&z_u\\ x_v&y_v&z_v\end{vmatrix}=\left(
\frac{\partial(y,z)}{\partial(u,v)},
\frac{\partial(z,x)}{\partial(u,v)},
\frac{\partial(x,y)}{\partial(u,v)}
\right).
\end{aligned}
\]
若所取正侧与 $\vec r_u\times\vec r_v=\left(\dfrac{\partial(y,z)}{\partial(u,v)},\dfrac{\partial(z,x)}{\partial(u,v)},\dfrac{\partial(x,y)}{\partial(u,v)}\right)$ 同向，则
\[
\mathrm{d}y\,\mathrm{d}z=\frac{\partial(y,z)}{\partial(u,v)}\,\mathrm{d}u\,\mathrm{d}v,\qquad
\mathrm{d}z\,\mathrm{d}x=\frac{\partial(z,x)}{\partial(u,v)}\,\mathrm{d}u\,\mathrm{d}v,\qquad
\mathrm{d}x\,\mathrm{d}y=\frac{\partial(x,y)}{\partial(u,v)}\,\mathrm{d}u\,\mathrm{d}v.
\]
若所取正侧与 $\vec r_u\times\vec r_v=\left(\dfrac{\partial(y,z)}{\partial(u,v)},\dfrac{\partial(z,x)}{\partial(u,v)},\dfrac{\partial(x,y)}{\partial(u,v)}\right)$ 反向，则三式同时乘以 $-1$。代入题目给出的积分
\begin{align*}
I&=P\,\mathrm{d}x\,\mathrm{d}y+Q\,\mathrm{d}y\,\mathrm{d}z+R\,\mathrm{d}z\,\mathrm{d}x\\
&=\epsilon\iint_D\Bigl[
P(x,y,z)(x_u y_v-y_u x_v)
+Q(x,y,z)(y_u z_v-z_u y_v)
+R(x,y,z)(z_u x_v-x_u z_v)
\Bigr]\,\mathrm{d}u\,\mathrm{d}v\\
&=\epsilon\iint_D\Bigl[
P(x,y,z)\frac{\partial(x,y)}{\partial(u,v)}
+Q(x,y,z)\frac{\partial(y,z)}{\partial(u,v)}
+R(x,y,z)\frac{\partial(z,x)}{\partial(u,v)}
\Bigr]\,\mathrm{d}u\,\mathrm{d}v
\end{align*}
若 $\Sigma:z=f(x,y)$，$(x,y)\in D$，换元
\[
\frac{\partial(x,y)}{\partial(x,y)}=1,~\frac{\partial(y,z)}{\partial(x,y)}=\begin{vmatrix}y_x&z_x\\y_y&z_y\end{vmatrix}=y_xz_y-z_xy_y=-f_x,~\frac{\partial(z,x)}{\partial(x,y)}=\begin{vmatrix}z_x&x_x\\z_y&x_y\end{vmatrix}=z_xx_y-x_xz_y=-f_y
\]
取 $\epsilon=1$ 表示上侧，取 $\epsilon=-1$ 表示下侧。于是
\[
\mathrm{d}y\,\mathrm{d}z=-\epsilon f_x\,\mathrm{d}x\,\mathrm{d}y,~~
\mathrm{d}z\,\mathrm{d}x=-\epsilon f_y\,\mathrm{d}x\,\mathrm{d}y,~~
\mathrm{d}x\,\mathrm{d}y=\epsilon\,\mathrm{d}x\,\mathrm{d}y.
\]
其中右端的 $\mathrm{d}x\,\mathrm{d}y$ 是投影区域 $D$ 上的普通面积元。因此
\[
\begin{aligned}
&\iint_{\Sigma}P\,\mathrm{d}x\,\mathrm{d}y
+Q\,\mathrm{d}y\,\mathrm{d}z
+R\,\mathrm{d}z\,\mathrm{d}x\\
&\qquad
=\epsilon\iint_D\bigl[
P(x,y,f(x,y))-Q(x,y,f(x,y))f_x-R(x,y,f(x,y))f_y
\bigr]\,\mathrm{d}x\,\mathrm{d}y.
\end{aligned}
\]
若曲面以隐式方程 $\Phi(x,y,z)=0$ 给出，且 $\nabla\Phi\ne\vec0$，则先选一个偏导不为零的变量解出来，本质上仍是上面的雅可比换元。比如 $\Phi_z\ne0$ 时，可局部写成
\[
z=f(x,y).
\]
由恒等式 $\Phi(x,y,f(x,y))=0$ 对 $x,y$ 求偏导，得
\[
\Phi_x+\Phi_z f_x=0,\qquad \Phi_y+\Phi_z f_y=0,
\]
即
\[
f_x=-\frac{\Phi_x}{\Phi_z},\qquad f_y=-\frac{\Phi_y}{\Phi_z}.
\]
代入图形曲面的公式，便得到投影到 $xOy$ 平面时
\[
\mathrm{d}y\,\mathrm{d}z=\epsilon\frac{\Phi_x}{\Phi_z}\,\mathrm{d}x\,\mathrm{d}y,\qquad
\mathrm{d}z\,\mathrm{d}x=\epsilon\frac{\Phi_y}{\Phi_z}\,\mathrm{d}x\,\mathrm{d}y,\qquad
\mathrm{d}x\,\mathrm{d}y=\epsilon\,\mathrm{d}x\,\mathrm{d}y.
\]
这里右端的 $\Phi_x,\Phi_y,\Phi_z$ 都要在曲面点 $(x,y,f(x,y))$ 处取值。符号 $\epsilon$ 仍由题目给定的方向决定：若参数化
\[
\vec r(x,y)=(x,y,f(x,y))
\]
的法向量 $\vec r_x\times\vec r_y=(-f_x,-f_y,1)$ 指向题目所取一侧，则 $\epsilon=1$；反向则 $\epsilon=-1$。由于
\[
(-f_x,-f_y,1)=\left(\frac{\Phi_x}{\Phi_z},\frac{\Phi_y}{\Phi_z},1\right)=\frac{1}{\Phi_z}\nabla\Phi,
\]
所以当题目说取 $\nabla\Phi$ 所指一侧时，若 $\Phi_z>0$ 则 $\epsilon=1$，若 $\Phi_z<0$ 则 $\epsilon=-1$。
\end{theorem}

\begin{example}
计算曲面积分 \(\displaystyle I=\iint_{S}x^{2}\mathrm{d}y\,\mathrm{d}z+y^{2}\mathrm{d}z\,\mathrm{d}x+z^{2}\mathrm{d}x\,\mathrm{d}y\)，其中 $S$ 为平面 $x+y+z=1$ 在第一卦限部分，方向向上(图6)。
\end{example}
\begin{solution}
\[
\begin{aligned}
\frac{\partial(y,z)}{\partial(x,y)}=\begin{vmatrix}y_x&z_x\\y_y&z_y\end{vmatrix}=1,
\frac{\partial(z,x)}{\partial(x,y)}=\begin{vmatrix}z_x&x_x\\z_y&x_y\end{vmatrix}=1,
\frac{\partial(x,y)}{\partial(x,y)}=\begin{vmatrix}x_x&y_x\\x_y&y_y\end{vmatrix}=1.
\end{aligned}
\]
于是$\vec r_u\times\vec r_v=\left(\dfrac{\partial(y,z)}{\partial(u,v)},\dfrac{\partial(z,x)}{\partial(u,v)},\dfrac{\partial(x,y)}{\partial(u,v)}\right)=(1,1,1)$指向上侧，题目取上侧，故 $\epsilon=1$。于是把$\mathrm{d}y\,\mathrm{d}z=\mathrm{d}z\,\mathrm{d}x=\mathrm{d}x\,\mathrm{d}y$直接代入，且 $z=1-x-y$，得
\[
x^{2}\mathrm{d}y\,\mathrm{d}z+y^{2}\mathrm{d}z\,\mathrm{d}x+z^{2}\mathrm{d}x\,\mathrm{d}y
=\bigl[x^2+y^2+(1-x-y)^2\bigr]\mathrm{d}x\,\mathrm{d}y.
\]
所以
\[
\begin{aligned}
I=\iint_D\bigl[x^2+y^2+(1-x-y)^2\bigr]\,\mathrm{d}x\,\mathrm{d}y=\int_0^1\int_0^{1-x}\bigl[x^2+y^2+(1-x-y)^2\bigr]\,\mathrm{d}y\,\mathrm{d}x
=\frac14.
\end{aligned}
\]
\end{solution}

\begin{example}
计算 \(\displaystyle I=\iint_{S}x\mathrm{d}y\,\mathrm{d}z+y\mathrm{d}z\,\mathrm{d}x+z\mathrm{d}x\,\mathrm{d}y\)，其中 $S$ 为平面 $x+y+z=1$ 在第一卦限部分，方向向上
\end{example}
\begin{solution}
\[
\begin{aligned}
\frac{\partial(y,z)}{\partial(x,y)}=\begin{vmatrix}y_x&z_x\\y_y&z_y\end{vmatrix}=1,
\frac{\partial(z,x)}{\partial(x,y)}=\begin{vmatrix}z_x&x_x\\z_y&x_y\end{vmatrix}=1,
\frac{\partial(x,y)}{\partial(x,y)}=\begin{vmatrix}x_x&y_x\\x_y&y_y\end{vmatrix}=1.
\end{aligned}
\]
于是$\vec r_u\times\vec r_v=\left(\dfrac{\partial(y,z)}{\partial(u,v)},\dfrac{\partial(z,x)}{\partial(u,v)},\dfrac{\partial(x,y)}{\partial(u,v)}\right)=(1,1,1)$指向上侧，题目取上侧，故 $\epsilon=1$。把$\mathrm{d}y\,\mathrm{d}z=\mathrm{d}z\,\mathrm{d}x=\mathrm{d}x\,\mathrm{d}y$直接代入
\[
x\mathrm{d}y\,\mathrm{d}z+y\mathrm{d}z\,\mathrm{d}x+z\mathrm{d}x\,\mathrm{d}y
=\bigl[x+y+(1-x-y)\bigr]\mathrm{d}x\,\mathrm{d}y
=\mathrm{d}x\,\mathrm{d}y.
\]
于是
\[
I=\iint_D 1\,\mathrm{d}x\,\mathrm{d}y
=\int_0^1\int_0^{1-x}1\,\mathrm{d}y\,\mathrm{d}x
=\frac12.
\]
\end{solution}

\begin{example}
计算曲面积分 \(\displaystyle I=\iint_{S}(z^{2}+x)\mathrm{d}y\,\mathrm{d}z-z\mathrm{d}x\,\mathrm{d}y\)，其中 $S$ 为 $\left\{(x,y,z):z=\frac{1}{2}(x^{2}+y^{2}),0\le z\le2\right\}$，正向指向 $S$ 的下侧(图7)。
\end{example}
\begin{solution}
按定理~\ref{thm:surface-integral-type2-calc}，把曲面看成参数曲面
\[
\vec r(x,y)=\left(x,y,\frac12(x^2+y^2)\right),\qquad
D=\{(x,y)\mid x^2+y^2\le 4\}.
\]
先列偏导：
\[
x_x=1,\quad x_y=0,\quad y_x=0,\quad y_y=1,\quad z_x=x,\quad z_y=y.
\]
于是三个雅可比行列式为
\[
\begin{aligned}
\frac{\partial(y,z)}{\partial(x,y)}
&=\begin{vmatrix}y_x&z_x\\y_y&z_y\end{vmatrix}
=\begin{vmatrix}0&x\\1&y\end{vmatrix}
=0\cdot y-x\cdot1=-x,\\
\frac{\partial(z,x)}{\partial(x,y)}
&=\begin{vmatrix}z_x&x_x\\z_y&x_y\end{vmatrix}
=\begin{vmatrix}x&1\\y&0\end{vmatrix}
=x\cdot0-1\cdot y=-y,\\
\frac{\partial(x,y)}{\partial(x,y)}
&=\begin{vmatrix}x_x&y_x\\x_y&y_y\end{vmatrix}
=\begin{vmatrix}1&0\\0&1\end{vmatrix}=1.
\end{aligned}
\]
这组参数给出的法向对应上侧，而题目要求下侧，所以 $\epsilon=-1$。于是直接代入
\[
\mathrm{d}y\,\mathrm{d}z=x\,\mathrm{d}x\,\mathrm{d}y,\qquad
\mathrm{d}z\,\mathrm{d}x=y\,\mathrm{d}x\,\mathrm{d}y,\qquad
\mathrm{d}x\,\mathrm{d}y=-\,\mathrm{d}x\,\mathrm{d}y.
\]
题目中没有 $\mathrm{d}z\,\mathrm{d}x$ 项，所以只需代入 $\mathrm{d}y\,\mathrm{d}z$ 与 $\mathrm{d}x\,\mathrm{d}y$。将 $z=\dfrac12(x^2+y^2)$ 代入，有
\[
\begin{aligned}
(z^2+x)\mathrm{d}y\,\mathrm{d}z-z\mathrm{d}x\,\mathrm{d}y
&=\bigl[(z^2+x)x-z(-1)\bigr]\mathrm{d}x\,\mathrm{d}y\\
&=\bigl(xz^2+x^2+z\bigr)\mathrm{d}x\,\mathrm{d}y\\
&=\left[\frac14x(x^2+y^2)^2+x^2+\frac12(x^2+y^2)\right]\mathrm{d}x\,\mathrm{d}y.
\end{aligned}
\]
因此
\[
\begin{aligned}
I&=\iint_D\left[\frac14x(x^2+y^2)^2+x^2+\frac12(x^2+y^2)\right]\mathrm{d}x\,\mathrm{d}y\\
&=\int_0^{2\pi}\int_0^2
\left(\frac14r^5\cos\theta+r^2\cos^2\theta+\frac12r^2\right)r\,\mathrm{d}r\,\mathrm{d}\theta\\
&=0+4\pi+4\pi=8\pi.
\end{aligned}
\]
\end{solution}

\begin{example}
计算曲面积分 \(\displaystyle I=\iint_{S}x\mathrm{d}y\,\mathrm{d}z+y\mathrm{d}z\,\mathrm{d}x+z\mathrm{d}x\,\mathrm{d}y\)，其中 $S$ 为 $\{(x,y,z):x^{2}+y^{2}=1,x\ge0,y\ge0,0\le z\le3\}$，正向指向 $S$ 的外侧(图8)。
\end{example}
\begin{solution}
按定理~\ref{thm:surface-integral-type2-calc}，取圆柱侧面的参数化
\[
\vec r(\theta,t)=(\cos\theta,\sin\theta,t),\qquad
0\le \theta\le \frac{\pi}{2},\quad 0\le t\le 3.
\]
先列偏导：
\[
x_\theta=-\sin\theta,\quad x_t=0,\quad
y_\theta=\cos\theta,\quad y_t=0,\quad
z_\theta=0,\quad z_t=1.
\]
三个雅可比行列式为
\[
\begin{aligned}
\frac{\partial(y,z)}{\partial(\theta,t)}
&=\begin{vmatrix}y_\theta&z_\theta\\y_t&z_t\end{vmatrix}
=\begin{vmatrix}\cos\theta&0\\0&1\end{vmatrix}
=\cos\theta,\\
\frac{\partial(z,x)}{\partial(\theta,t)}
&=\begin{vmatrix}z_\theta&x_\theta\\z_t&x_t\end{vmatrix}
=\begin{vmatrix}0&-\sin\theta\\1&0\end{vmatrix}
=0\cdot0-(-\sin\theta)\cdot1=\sin\theta,\\
\frac{\partial(x,y)}{\partial(\theta,t)}
&=\begin{vmatrix}x_\theta&y_\theta\\x_t&y_t\end{vmatrix}
=\begin{vmatrix}-\sin\theta&\cos\theta\\0&0\end{vmatrix}=0.
\end{aligned}
\]
由这三个雅可比行列式组成的法向量为 $(\cos\theta,\sin\theta,0)$，正好指向圆柱外侧，故 $\epsilon=1$。因此
\[
\mathrm{d}y\,\mathrm{d}z=\cos\theta\,\mathrm{d}\theta\,\mathrm{d}t,\qquad
\mathrm{d}z\,\mathrm{d}x=\sin\theta\,\mathrm{d}\theta\,\mathrm{d}t,\qquad
\mathrm{d}x\,\mathrm{d}y=0.
\]
代入 $x=\cos\theta,\ y=\sin\theta$ 后
\[
x\mathrm{d}y\,\mathrm{d}z+y\mathrm{d}z\,\mathrm{d}x+z\mathrm{d}x\,\mathrm{d}y
=\bigl(\cos^2\theta+\sin^2\theta\bigr)\mathrm{d}\theta\,\mathrm{d}t
=\mathrm{d}\theta\,\mathrm{d}t.
\]
于是
\[
I=\int_0^3\int_0^{\pi/2}1\,\mathrm{d}\theta\,\mathrm{d}t
=\frac{3\pi}{2}.
\]
\end{solution}

\begin{example}
计算曲面积分 $I=\iint_{S}x\mathrm{d}y\,\mathrm{d}z+y\mathrm{d}z\,\mathrm{d}x+z\mathrm{d}x\,\mathrm{d}y$，其中 $S$ 为 $x^{2}+y^{2}+z^{2}=a^{2}$，正向指向 $S$ 的外侧(图9)。
\end{example}
\begin{solution}
按定理~\ref{thm:surface-integral-type2-calc}，取球面参数化
\[
\vec r(\varphi,\theta)
=(a\sin\varphi\cos\theta,\ a\sin\varphi\sin\theta,\ a\cos\varphi),
\quad 0\le \varphi\le \pi,\quad 0\le \theta\le 2\pi.
\]
先列偏导：
\[
\begin{aligned}
x_\varphi&=a\cos\varphi\cos\theta,& x_\theta&=-a\sin\varphi\sin\theta,\\
y_\varphi&=a\cos\varphi\sin\theta,& y_\theta&=a\sin\varphi\cos\theta,\\
z_\varphi&=-a\sin\varphi,& z_\theta&=0.
\end{aligned}
\]
然后逐个计算雅可比行列式：
\[
\begin{aligned}
\frac{\partial(y,z)}{\partial(\varphi,\theta)}&=\begin{vmatrix}y_\varphi&z_\varphi\\y_\theta&z_\theta\end{vmatrix}
=\begin{vmatrix}a\cos\varphi\sin\theta&-a\sin\varphi\\a\sin\varphi\cos\theta&0\end{vmatrix}=a^2\sin^2\varphi\cos\theta,\\
\frac{\partial(z,x)}{\partial(\varphi,\theta)}
&=\begin{vmatrix}z_\varphi&x_\varphi\\z_\theta&x_\theta\end{vmatrix}=\begin{vmatrix}-a\sin\varphi&a\cos\varphi\cos\theta\\
0&-a\sin\varphi\sin\theta\end{vmatrix}=a^2\sin^2\varphi\sin\theta,\\
\frac{\partial(x,y)}{\partial(\varphi,\theta)}
&=\begin{vmatrix}x_\varphi&y_\varphi\\x_\theta&y_\theta\end{vmatrix}
=\begin{vmatrix}
a\cos\varphi\cos\theta&a\cos\varphi\sin\theta\\
-a\sin\varphi\sin\theta&a\sin\varphi\cos\theta
\end{vmatrix}=a^2\sin\varphi\cos\varphi.
\end{aligned}
\]
组成的法向量$\left(a^2\sin^2\varphi\cos\theta,a^2\sin^2\varphi\sin\theta,a^2\sin\varphi\cos\varphi\right)$与位置向量同向，故指向球面外侧，题目也取外侧，所以 $\epsilon=1$。因此
\[
\mathrm{d}y\,\mathrm{d}z=a^2\sin^2\varphi\cos\theta\,\mathrm{d}\varphi\,\mathrm{d}\theta,~~
\mathrm{d}z\,\mathrm{d}x=a^2\sin^2\varphi\sin\theta\,\mathrm{d}\varphi\,\mathrm{d}\theta,~~
\mathrm{d}x\,\mathrm{d}y=a^2\sin\varphi\cos\varphi\,\mathrm{d}\varphi\,\mathrm{d}\theta.
\]
代入$x=a\sin\varphi\cos\theta,~y=a\sin\varphi\sin\theta,~z=a\cos\varphi$得
\[
x\mathrm{d}y\,\mathrm{d}z+y\mathrm{d}z\,\mathrm{d}x+z\mathrm{d}x\,\mathrm{d}y=a^3\bigl(\sin^3\varphi\cos^2\theta+\sin^3\varphi\sin^2\theta+\sin\varphi\cos^2\varphi\bigr)\mathrm{d}\varphi\,\mathrm{d}\theta=a^3\sin\varphi\,\mathrm{d}\varphi\,\mathrm{d}\theta.
\]
于是
\[
I=\int_0^{2\pi}\int_0^\pi a^3\sin\varphi\,\mathrm{d}\varphi\,\mathrm{d}\theta
=4\pi a^3.
\]
\end{solution}

\subsection{第二类曲面积分：投影、封口、Gauss}

\subsubsection{Gauss 公式：封口法的依据}

Gauss 公式把闭曲面的通量变成体区域中的散度积分。本小节围绕三个问题展开：散度是什么，为什么闭曲面通量可以换成体积分，以及怎样从散度运算自然推出 Laplace 算子和 Green 第一、第二公式。

设 $\vec{F}=(P,Q,R)$ 为向量场。通量的概念自然引出散度的定义。

\begin{proposition}[Gauss 公式的使用前提]
\begin{itemize}
    \item 只对封闭曲面直接成立，而且法向必须取外侧。
    \item 开曲面要先补成封闭曲面，再把补面的通量减回去。
    \item 向量场必须在所围体区域及其边界上具有连续一阶偏导数；若体内有奇点，不能直接把奇点包进去套 Gauss。
    \item 内部公共界面会因为法向相反而彼此抵消。
\end{itemize}
\end{proposition}

\begin{definition}[散度] \label{def:divergence}
设向量场 $\vec{F}=(P,Q,R)$ 在区域 $\Omega$ 上具有连续一阶偏导数, 定义 $\vec{F}$ 的\textbf{散度}为 \(\displaystyle \operatorname{div}\vec{F} = \frac{\partial P}{\partial x} + \frac{\partial Q}{\partial y} + \frac{\partial R}{\partial z}\)。
散度是一个标量场, 表示向量场在某点处单位体积的通量.
\end{definition}

\begin{note}
散度具有明确的物理意义：\(\operatorname{div}\vec F\) 在某点的值等于向量场在该点处单位体积的净流出量。
设 \(\vec F\) 在点 \(p_0\) 附近具有连续偏导数。取以 \(p_0\) 为心、\(\varepsilon\) 为半径的小球 \(B_\varepsilon(p_0)\)。由 Gauss 公式，
\[
\oiint_{\partial B_\varepsilon(p_0)} \vec F\cdot\vec n\,\mathrm{d}S
=\iiint_{B_\varepsilon(p_0)}\operatorname{div}\vec F\,\mathrm{d}V.
\]
两边除以小球体积 \(V(B_\varepsilon)=\frac43\pi\varepsilon^3\)，再令 \(\varepsilon\to0\)，得到
\[
\operatorname{div}\vec F(p_0)
=\lim_{\varepsilon\to0}
\frac{1}{V(B_\varepsilon(p_0))}
\oiint_{\partial B_\varepsilon(p_0)}\vec F\cdot\vec n\,\mathrm{d}S.
\]
这就是“单位体积通量”的严格含义。
\end{note}

\begin{theorem}[散度运算法则]
设 \(C\) 为常数，\(\varphi\) 为数量场，\(\vec F,\vec F_1,\vec F_2\) 为向量场，则
\[
\operatorname{div}(C\vec F)=C\operatorname{div}\vec F,\qquad
\operatorname{div}(\vec F_1+\vec F_2)=\operatorname{div}\vec F_1+\operatorname{div}\vec F_2,
\]
\[
\operatorname{div}(\varphi\vec F)
=\varphi\operatorname{div}\vec F+\vec F\cdot\nabla\varphi.
\]
\end{theorem}
\begin{proof}
前两式由偏导线性性直接得到。第三式设 \(\vec F=(P,Q,R)\)，则
\begin{align*}
\operatorname{div}(\varphi\vec F)
&=(\varphi P)_x+(\varphi Q)_y+(\varphi R)_z\\
&=\varphi(P_x+Q_y+R_z)+P\varphi_x+Q\varphi_y+R\varphi_z\\
&=\varphi\operatorname{div}\vec F+\vec F\cdot\nabla\varphi.
\end{align*}
\end{proof}

\begin{definition}[Laplace算子] \label{def:laplacian}
\textbf{Laplace 算子}定义为
\[
\Delta u=\nabla\cdot(\nabla u)=u_{xx}+u_{yy}+u_{zz}.
\]
若函数 \(u\) 满足 \(\Delta u=0\)，则称 \(u\) 为\textbf{调和函数}。
\end{definition}

\begin{example}[径向场散度计算]
设 \(\vec p=(x,y,z)\)，\(p=\|\vec p\|=\sqrt{x^2+y^2+z^2}\)，\(n\) 为实数。求 \(\operatorname{div}(p^n\vec p)\)。
\end{example}
\begin{solution}
由散度乘积法则，
\[
\operatorname{div}(p^n\vec p)
=p^n\operatorname{div}\vec p+\vec p\cdot\nabla(p^n).
\]
其中 \(\operatorname{div}\vec p=3\)。又
\[
\nabla(p^n)=np^{n-2}\vec p,
\qquad
\vec p\cdot\nabla(p^n)=np^n.
\]
所以
\[
\operatorname{div}(p^n\vec p)=(n+3)p^n.
\]
\end{solution}

\begin{theorem}[Gauss公式] \label{thm:gauss}
设空间有界闭区域 $\Omega$ 由分片光滑的闭曲面 $\partial\Omega$ 围成, 函数 $P(x,y,z),Q(x,y,z),R(x,y,z)$ 在 $\Omega$ 上具有连续一阶偏导数, 则
\begin{align*}
\iiint_\Omega \left(\frac{\partial P}{\partial x} + \frac{\partial Q}{\partial y} + \frac{\partial R}{\partial z}\right) \mathrm{d}x\,\mathrm{d}y\,\mathrm{d}z= \oiint_{\partial\Omega} P\mathrm{d}y\mathrm{d}z+Q\mathrm{d}z\mathrm{d}x+R\mathrm{d}x\mathrm{d}y
\end{align*}
其中 $\partial\Omega$ 的法向量 $\vec{n}=(\cos\alpha,\cos\beta,\cos\gamma)$ 指向外侧，第二式为向量形式。
\end{theorem}

\begin{example}
计算曲面积分 $\oiint_\Sigma x^3\mathrm{d}y\,\mathrm{d}z + y^3\mathrm{d}z\,\mathrm{d}x + z^3\mathrm{d}x\,\mathrm{d}y$, 其中 $\Sigma$ 为球面 $x^2+y^2+z^2 = a^2$ 的外侧.
\end{example}
\begin{solution}
由 Gauss 公式，设 $\Omega$ 为球体 $x^2+y^2+z^2 \le a^2$, 则
\begin{align*}
\oiint_\Sigma x^3\mathrm{d}y\,\mathrm{d}z + y^3\mathrm{d}z\,\mathrm{d}x + z^3\mathrm{d}x\,\mathrm{d}y
&= \iiint_\Omega (3x^2 + 3y^2 + 3z^2)\,\mathrm{d}x\,\mathrm{d}y\,\mathrm{d}z= 3\iiint_\Omega (x^2+y^2+z^2)\,\mathrm{d}x\,\mathrm{d}y\,\mathrm{d}z
\end{align*}
用球坐标 $x = r\sin\varphi\cos\theta,\,y = r\sin\varphi\sin\theta,\,z = r\cos\varphi$:
\begin{align*}
3\iiint_\Omega r^2 \cdot r^2\sin\varphi\,\mathrm{d}r\,\mathrm{d}\varphi\,\mathrm{d}\theta
&= 3\int_0^{2\pi}\mathrm{d}\theta\int_0^\pi\sin\varphi\,\mathrm{d}\varphi\int_0^a r^4\,\mathrm{d}r = 3\cdot 2\pi\cdot 2\cdot\frac{a^5}{5} = \frac{12\pi a^5}{5}
\end{align*}
\end{solution}

\begin{conclusion}[Gauss 公式]
\begin{itemize}
    \item 先看曲面是否封闭，再决定能不能直接套。
    \item 开曲面通常靠补面转化。
    \item 体积、通量和散度是一组最常用的互化关系。
\end{itemize}
\end{conclusion}

\subsubsection{三维Green第一第二公式}

\begin{theorem}[三维Green第一公式]\label{ex:green-first-3d}
设 \(\Omega\) 为光滑闭曲面 \(\partial\Omega\) 所围成的有界区域，函数 \(u,v\) 在 \(\Omega\) 上具有连续二阶偏导数，\(\dfrac{\partial}{\partial n}\) 表示沿外法向的方向导数，则
\[
\oiint_{\partial\Omega}v\frac{\partial u}{\partial n}\,\mathrm{d}S
=\iiint_\Omega\left(\nabla v\cdot\nabla u+v\Delta u\right)\mathrm{d}V,
\]
并且
\[
\oiint_{\partial\Omega}\left(v\frac{\partial u}{\partial n}
-u\frac{\partial v}{\partial n}\right)\mathrm{d}S
=\iiint_\Omega\left(v\Delta u-u\Delta v\right)\mathrm{d}V.
\]
\end{theorem}



\begin{theorem}[三维Green第二公式]
设 \(\Omega\) 为光滑闭曲面 \(\partial\Omega\) 所围成的有界区域，函数 \(u,v\) 在 \(\Omega\) 上具有连续二阶偏导数，\(\dfrac{\partial}{\partial n}\) 表示沿外法向的方向导数，则
\[
\oiint_{\partial\Omega}\left(v\frac{\partial u}{\partial n}
-u\frac{\partial v}{\partial n}\right)\mathrm{d}S
=\iiint_\Omega\left(v\Delta u-u\Delta v\right)\mathrm{d}V.
\]
\end{theorem}



\begin{example}[调和函数与奇点]
在 \(\mathbb R^3\) 中，设 \(p=\sqrt{x^2+y^2+z^2}\)。证明 \(\dfrac1p\) 在 \(\mathbb R^3\setminus\{0\}\) 上为调和函数。
\end{example}
\begin{proof}
直接计算：
\[
\frac{\partial}{\partial x}\left(\frac1p\right)=-\frac{x}{p^3},\qquad
\frac{\partial^2}{\partial x^2}\left(\frac1p\right)=-\frac1{p^3}+\frac{3x^2}{p^5}.
\]
类似地，
\[
\frac{\partial^2}{\partial y^2}\left(\frac1p\right)=-\frac1{p^3}+\frac{3y^2}{p^5},\qquad
\frac{\partial^2}{\partial z^2}\left(\frac1p\right)=-\frac1{p^3}+\frac{3z^2}{p^5}.
\]
三式相加得
\[
\Delta\left(\frac1p\right)
=-\frac3{p^3}+\frac{3(x^2+y^2+z^2)}{p^5}
=-\frac3{p^3}+\frac{3p^2}{p^5}=0.
\]
注意 \(p=0\) 是奇点；若体区域包含原点，不能把 \(\frac1p\) 当作处处光滑函数直接套 Gauss 或 Green 公式，应先挖去小球再取极限。
\end{proof}

\subsubsection{开曲面的封口与补面}

\begin{example}
计算曲面积分 $\oiint_\Sigma x\mathrm{d}y\,\mathrm{d}z + y\mathrm{d}z\,\mathrm{d}x + z\mathrm{d}x\,\mathrm{d}y$, 其中 $\Sigma$ 为平面 $x+y+z=1$ 在第一卦限部分的上侧, 法向量为 $(1,1,1)$ 方向.
\end{example}
\begin{solution}
曲面 $\Sigma$ 不是封闭曲面，不能直接使用 Gauss 公式，因此先补三个坐标面：$\Sigma_1$ 为 $x=0$ 上对应部分（外法向量朝 $-x$ 方向），$\Sigma_2$ 为 $y=0$ 上对应部分，$\Sigma_3$ 为 $z=0$ 上对应部分。令
$\Sigma' = \Sigma \cup \Sigma_1 \cup \Sigma_2 \cup \Sigma_3$，则 $\Sigma'$ 构成封闭曲面，围成的正是第一卦限内的四面体 $\Omega$。由 Gauss 公式，闭曲面总通量为 \(\displaystyle \oiint_{\Sigma'} x\mathrm{d}y\,\mathrm{d}z + y\mathrm{d}z\,\mathrm{d}x + z\mathrm{d}x\,\mathrm{d}y=\iiint_\Omega 3\,\mathrm{d}V=3V(\Omega)=3\cdot\frac{1}{6}=\frac{1}{2}\)。
再检查补上的三个面：
\begin{itemize}
    \item 在 $\Sigma_1$（$x=0$）上，$x=0$，故被积式中的 $x\,\mathrm{d}y\,\mathrm{d}z$ 这一项恒为 $0$；
    \item 在 $\Sigma_2$（$y=0$）上，$y=0$，故 $y\,\mathrm{d}z\,\mathrm{d}x$ 恒为 $0$；
    \item 在 $\Sigma_3$（$z=0$）上，$z=0$，故 $z\,\mathrm{d}x\,\mathrm{d}y$ 恒为 $0$。
\end{itemize}
其余两项之所以不出现，是因为相应坐标平面上的有向投影面积本来就为 $0$。因此三个补面上的通量都为 $0$。

因此 \(\displaystyle \oiint_\Sigma x\mathrm{d}y\,\mathrm{d}z + y\mathrm{d}z\,\mathrm{d}x + z\mathrm{d}x\,\mathrm{d}y = \frac{1}{2}\)。
\end{solution}

\begin{conclusion}[Gauss 公式中的封闭性与补面关系]
\begin{itemize}
    \item 外法向、封闭性、补面这三件事必须先核对。
    \item 若能把散度积成简单函数，优先走体积分。
    \item 若补面后新增部分特别简单，Gauss 公式通常比直接算通量更快。
\end{itemize}
\end{conclusion}

\subsection{算子计算小结：按使用场景回看}

本节不再单独堆放算子公式，只保留索引。读到具体题型时，应回到对应位置使用：

\par\noindent{\centering
\renewcommand{\arraystretch}{1.35}
\begin{tabular}{|c|c|}
\hline
\textbf{要查的内容} & \textbf{所在位置} \\
\hline
\(\nabla\) 总表、梯度/旋度/散度对应关系 & 5.1 总路线 \\
\hline
\(\operatorname{grad}u=\nabla u\)、保守场、\(\nabla\times(\nabla u)=\vec0\) & 开曲线：保守场、势函数与路径无关 \\
\hline
旋度定义、方向旋量、径向场无旋、\(\nabla\cdot(\nabla\times\vec F)=0\) & 空间闭曲线：Stokes 公式与换面 \\
\hline
散度定义、散度物理意义、散度运算法则 & Gauss 公式：封口法的依据 \\
\hline
\(\Delta u=\operatorname{div}(\nabla u)\)、Green 第一/第二公式、调和函数与奇点 & Gauss 公式的推论部分 \\
\hline
\end{tabular}
\par}

这样处理后，算子不是额外的一节新知识，而是跟着题面入口出现：开曲线看梯度，空间闭曲线看旋度，闭曲面和封口题看散度，Green 第一、第二公式看作 Gauss 公式加散度乘积法则的推论。

\subsection{综合应用：一题多解与方法优先级}

\begin{note}[为什么要单独写这一节]
Green、Gauss、Stokes 三条公式本来就是同一条思想在不同维度上的呈现：\textbf{边界上的积分 = 内部的"导数"积分}。本节把一题多解、公式优先级和奇点检查放在一起，不是为了说“随便哪种都行”，而是为了说明做题顺序：先看题面触发了什么结构，再看公式是否合法，最后才比较哪条计算更短。
\end{note}

\begin{conclusion}[综合题的选公式检查表]
\begin{itemize}
    \item 先看微元：\(\mathrm{d}s,\mathrm{d}S\) 多数按第一类积分直接落地；\(\mathrm{d}x,\mathrm{d}y,\mathrm{d}z\) 进入第二类曲线积分；\(\mathrm{d}y\,\mathrm{d}z,\mathrm{d}z\,\mathrm{d}x,\mathrm{d}x\,\mathrm{d}y\) 进入第二类曲面积分。
    \item 再看结构：开曲线端点题查保守场，平面闭曲线查 Green，空间闭曲线查 Stokes，闭曲面或可封口曲面查 Gauss。
    \item 再看合法性：方向、闭合性、光滑性、单连通性、奇点位置都要核对；不满足时先补边、封口、换面、挖洞或分块。
    \item 最后比较计算量：直接参数化主要用于第一类积分、显式给定且参数积分简单的第二类积分，或作为检查方向和符号的定义层方法；不能把它简单等同于结构公式失败后的保底路线。
\end{itemize}
\end{conclusion}

\subsubsection{三大公式的同一结构：边界换内部}

三条公式背后的共同结构是
\[
\text{边界上的积分}=\text{内部导数的积分}.
\]
若用较统一的语言，可以写成 \(\displaystyle \int_{\partial \Omega} \omega = \int_\Omega \mathrm{d}\omega\)，其中 $\Omega$ 是某个"区域"，$\partial\Omega$ 是它的边界，$\mathrm{d}\omega$ 是某种"微分"。初学时不必先记这个抽象式，只要看清下面这张表对应的题面形状即可：

\par\noindent{\centering
\renewcommand{\arraystretch}{1.4}
\begin{tabular}{|c|c|c|c|}
\hline
\textbf{定理} & \textbf{内部 $\Omega$} & \textbf{边界 $\partial\Omega$} & \textbf{内部"导数"} \\
\hline
Green & 平面区域 $D$（2维） & 闭曲线 $C$（1维） & $Q_x - P_y$ \\
\hline
Stokes & 曲面 $\Sigma$（2维） & 边界曲线 $\partial\Sigma$（1维） & $\nabla\times\vec F\cdot\vec n$ \\
\hline
Gauss & 体区域 $V$（3维） & 闭曲面 $\partial V$（2维） & $\nabla\cdot\vec F$ \\
\hline
\end{tabular}
\par}

把表格翻译回做题语言，就是三句话：平面闭曲线上的 \(P\,\mathrm{d}x+Q\,\mathrm{d}y\) 先找 Green；空间闭曲线上的 \(P\,\mathrm{d}x+Q\,\mathrm{d}y+R\,\mathrm{d}z\) 先找 Stokes；闭曲面或可封口曲面上的 \(P\,\mathrm{d}y\,\mathrm{d}z+Q\,\mathrm{d}z\,\mathrm{d}x+R\,\mathrm{d}x\,\mathrm{d}y\) 先找 Gauss。若曲面积分本身不是旋度通量，不要只因为曲面有边界就把它交给 Stokes。

\subsubsection{一题多解：直算、公式法与对称性比较}

一题多解在本章很常见，并不说明题目“没有固定方法”。同一个积分有时可以直接参数化，也可以用 Green、Stokes 或 Gauss；但考试落笔时要先看触发条件，而不是凭感觉任选。边界难而内部简单，就换到内部；空间边界曲线难而同边界曲面简单，就考虑 Stokes 换面；开曲面容易封口，就考虑封口后用 Gauss。若没有闭合、补边、封口、路径无关等结构信号，且曲线或曲面给法明确、参数积分可控，直接计算才是自然主路。

可以把全章主线压成一句话：\textbf{直接计算负责把微元落地，结构公式负责把对象换掉，补边、封口、换面、挖洞负责让结构公式合法可用；做题时先判对象和结构，再谈计算。}

\begin{conclusion}[公式不是题型标签]
\begin{itemize}
    \item Green、Stokes、Gauss 对应的是结构，不是固定题型名称。
    \item 看到闭合、边界、封口、奇点时，先问公式条件是否满足；条件暂时不满足时，先尝试补边、封口、换面、挖洞或分块，而不是立刻硬算。
    \item 直接计算、对称性和结构公式可以同时作为候选；考试书写时要选触发条件最明确、合法性最清楚、计算量最小的一条主路，并说明选择理由。
\end{itemize}
\end{conclusion}

\subsubsection{优先级：对象识别先于公式记忆}

\begin{example}[封口法 + Gauss 公式]\label{ex:gauss-cap}
计算曲面积分
\[
I=\iint_\Sigma x^2\,\mathrm{d}y\,\mathrm{d}z
+y^2\,\mathrm{d}z\,\mathrm{d}x
+z^2\,\mathrm{d}x\,\mathrm{d}y,
\]
其中 $\Sigma$ 为半球面 $z = \sqrt{a^2 - x^2 - y^2}$，取外侧（即法向量朝上）。
\end{example}
\begin{solution}
这是第二类曲面积分，但积分面只是上半球面，不是封闭曲面。

先补底面圆盘封口，再用 Gauss；最后减去补面贡献。

设 \(\displaystyle \vec F=(x^2,y^2,z^2)\)。补上底面 \(\displaystyle \Sigma_0:\ z=0,\ x^2+y^2\le a^2\)，并取下侧法向，使 $\Sigma\cup\Sigma_0$ 成为整体外向的闭曲面。记上半球体为 $V$，则 \(\displaystyle I+\iint_{\Sigma_0}\vec F\cdot\mathrm{d}\vec S=\iiint_V \nabla\cdot\vec F\,\mathrm{d}V\)。
而 \(\displaystyle \nabla\cdot\vec F=2x+2y+2z\)。
由对称性，\(\displaystyle \iiint_V x\,\mathrm{d}V=\iiint_V y\,\mathrm{d}V=0\)，所以 \(\displaystyle \iiint_V (2x+2y+2z)\,\mathrm{d}V=2\iiint_V z\,\mathrm{d}V\)。
改用柱坐标：
\[
2\int_0^{2\pi}\mathrm{d}\theta\int_0^a r\,\mathrm{d}r\int_0^{\sqrt{a^2-r^2}} z\,\mathrm{d}z
=2\pi\int_0^a (a^2r-r^3)\,\mathrm{d}r
=\frac{\pi a^4}{2}.
\]
再看补面 $\Sigma_0$：其上 $z=0$，且外法向朝下，所以 \(\displaystyle \iint_{\Sigma_0}\vec F\cdot\mathrm{d}\vec S=\iint_{\Sigma_0} z^2(-1)\,\mathrm{d}x\,\mathrm{d}y=0\)。
故 \(I=\dfrac{\pi a^4}{2}\)。

底面一定要取下侧法向，才能与半球面的外法向拼成整体外向。
\end{solution}

\begin{example}[Stokes 换面法]\label{ex:stokes-change-surface}
计算 \(\displaystyle I=\oint_L (y-z)\,\mathrm{d}x+(z-x)\,\mathrm{d}y+(x-y)\,\mathrm{d}z\)，其中 $L$ 为球面 $x^2+y^2+z^2=1$ 与平面 $x+y+z=0$ 的交线，从 $z$ 轴正向看去为逆时针方向。
\end{example}
\begin{solution}
这是空间闭曲线上的第二类曲线积分，直接参数化交线很麻烦。

换成同边界的最简单曲面，用 Stokes 把曲线积分改成旋度通量。

设 $\vec F=(y-z,\ z-x,\ x-y)$，则 \(\displaystyle \nabla\times\vec F=\bigl((-1)-1,\ (-1)-1,\ (-1)-1\bigr)=(-2,-2,-2)\)。取边界 $L$ 所围的平面圆盘 \(\displaystyle \Sigma:\ x+y+z=0,\qquad x^2+y^2+z^2\le 1\)。
从 $z$ 轴正向看去边界为逆时针，由右手法则应取 $z$ 分量为正的法向量 \(\displaystyle \vec n=\frac{1}{\sqrt3}(1,1,1)\)。
于是 \(\displaystyle (\nabla\times\vec F)\cdot\vec n=(-2,-2,-2)\cdot\frac{1}{\sqrt3}(1,1,1)=-2\sqrt3\)。
又该平面过球心，交线是大圆，故圆盘半径为 $1$，面积为 $\pi$。所以 \(\displaystyle I=\iint_\Sigma (\nabla\times\vec F)\cdot\vec n\,\mathrm{d}S=(-2\sqrt3)\pi\)。

这题最容易错的是法向方向；只要记住“从 $z$ 轴正向看逆时针”，就必须选 $z$ 分量为正的法向。
\end{solution}

\begin{example}[Green 公式 + 路径无关性判断]\label{ex:green-path-independence}
设 $\vec F = \left(\dfrac{-y}{x^2+y^2},\;\dfrac{x}{x^2+y^2}\right)$。
\begin{compactenum}
    \item 验证在去掉原点的平面上 $Q_x - P_y = 0$；
    \item 计算 $\oint_C \vec F\cdot\mathrm{d}\vec r$，其中 $C$ 为任意包含原点的正向简单闭曲线；
    \item 说明为什么 $Q_x = P_y$ 不能推出路径无关。
\end{compactenum}
\end{example}
\begin{solution}
这题同时考三个点：先算 $Q_x-P_y$，再处理含奇点的闭曲线积分，最后判断“是否路径无关”。

第 (1) 问直接求偏导；第 (2) 问用挖洞法；第 (3) 问回到单连通条件做结论。


\begin{compactenum}
    \item 记 \(\displaystyle P=\frac{-y}{x^2+y^2},\ Q=\frac{x}{x^2+y^2}\)。
    则在 $(x,y)\neq(0,0)$ 处，
    \(\displaystyle P_y=\frac{-(x^2+y^2)+2y^2}{(x^2+y^2)^2}
    =\frac{y^2-x^2}{(x^2+y^2)^2}\)，
    \(\displaystyle Q_x=\frac{(x^2+y^2)-2x^2}{(x^2+y^2)^2}
    =\frac{y^2-x^2}{(x^2+y^2)^2}\)。
所以 \(\displaystyle Q_x-P_y=0\)。

    \item 取小圆 $C_\varepsilon:\ x^2+y^2=\varepsilon^2$，在环域上由 Green 公式得
    \(\displaystyle \oint_C \vec F\cdot\mathrm{d}\vec r=\oint_{C_\varepsilon}\vec F\cdot\mathrm{d}\vec r\)。
    参数化 \(\displaystyle x=\varepsilon\cos t,\ y=\varepsilon\sin t,\ 0\le t\le 2\pi\)，
    则
    \(\displaystyle \oint_{C_\varepsilon}\frac{-y\,\mathrm{d}x+x\,\mathrm{d}y}{x^2+y^2}
    =\int_0^{2\pi}\mathrm{d}t
    =2\pi\)。
    因而 \(\displaystyle \oint_C \vec F\cdot\mathrm{d}\vec r=2\pi\)。

    \item 这里 $Q_x=P_y$ 只能说明局部上像保守场；但去掉原点后的平面有洞，不是单连通区域，所以“闭路积分为零”“路径无关”之间的全局推理断掉了。上面第 (2) 问已经给出反例：存在闭曲线积分等于 $2\pi\neq 0$。
\end{compactenum}

这题最后一定要把“奇点 + 非单连通”说全，否则就会误判成路径无关。
\end{solution}


\section{习题}

本章习题仍按知识点分节排列，但做题入口要按题面形状来选：第一、三节主要训练 \(\mathrm{d}s,\mathrm{d}S\) 型直接计算；第二节训练 \(P\,\mathrm{d}x+Q\,\mathrm{d}y(+R\,\mathrm{d}z)\) 的直接参数化；第五、六节分别对应平面闭路的 Green 公式与开路端点题的保守场方法；第七、八节分别对应通量题的 Gauss 公式与空间闭曲线的 Stokes 公式。遇到综合题时，先回到“认微元、看方向、看闭合、查奇点”这四步，再决定使用哪条公式链。

\subsection{第一节习题：第一类曲线积分}

\subsubsection{A组}

\begin{exercise}[基本概念]
回答下列问题：
\begin{shortexenum}
    \shortitem{1}{如何运用微元法求曲线的质量？} &
    \shortitem{2}{第一类曲线积分怎样定义？}
\end{shortexenum}
\end{exercise}
\begin{solution}
\begin{compactenum}
    \item 设曲线 $L$ 上的线密度为 $\rho(x,y,z)$。在曲线上取弧长微元 $\mathrm{d}s$，则对应的质量微元为 $\mathrm{d}m=\rho\,\mathrm{d}s$。把整条曲线上的质量微元累加起来，就得到总质量 \(\displaystyle m=\int_{L}\rho(x,y,z)\,\mathrm{d}s\)。
    这就是用微元法求曲线质量的基本思想。
    \item 设 $L$ 是一条可求长曲线，函数 $f$ 定义在 $L$ 上。把 $L$ 分成许多小弧段 $\Delta s_i$，在每段上任取一点 $\xi_i$，作和式 \(\displaystyle \sum_{i=1}^{n} f(\xi_i)\Delta s_i\)。
    当分割无限加细时，若上述和式的极限存在，且与分割方式和取点方式无关，则称该极限为函数 $f$ 在曲线 $L$ 上的第一类曲线积分，记作 $\int_L f\,\mathrm{d}s$。
\end{compactenum}
\end{solution}

\begin{exercise}[第一类曲线积分的直接计算]
计算下列曲线积分：
\begin{shortexenum}
    \shortitem{1}{$\displaystyle \int_{L}xy\,\mathrm{d}s$，其中 $L$ 是曲线 $x=t,\ y=\dfrac{t^2}{2},\ 0\le t\le 1$；} &
    \shortitem{2}{$\displaystyle \int_{L}(x+y)\,\mathrm{d}s$，其中 $L$ 是以点 $(0,0)$，$(1,0)$，$(1,1)$ 为顶点的三角形的边界；} \\
    \shortitem{3}{$\displaystyle \int_{C}(x+y)\,\mathrm{d}s$，其中 $C$ 是以 $a$ 为半径、圆心在原点的右半圆周；} &
    \shortitem{4}{$\displaystyle \int_{C}\sqrt{x^2+y^2}\,\mathrm{d}s$，其中 $C$ 为圆周 $x^2+y^2=ax$。}
\end{shortexenum}
\end{exercise}
\begin{solution}
\begin{compactenum}
    \item 由参数方程 $x=t,\ y=\dfrac{t^2}{2},\ 0\le t\le 1$ 得
    \begin{align*}
    \mathrm{d}s&=\sqrt{1+t^2}\,\mathrm{d}t,\qquad xy=\frac{t^3}{2},\\
    \int_L xy\,\mathrm{d}s
    &=\frac{1}{2}\int_0^1 t^3\sqrt{1+t^2}\,\mathrm{d}t
    =\frac{1}{4}\int_1^2 \left(u^{3/2}-u^{1/2}\right)\,\mathrm{d}u
    \quad (u=1+t^2) \\
    &=\left[\frac{1}{10}u^{5/2}-\frac{1}{6}u^{3/2}\right]_1^2
    =\frac{1+\sqrt{2}}{15}.
    \end{align*}
    \item 设三角形边界依次为三条边 $L_1,L_2,L_3$。由于第一类曲线积分与方向无关，且边界可拆成三段，只需分别计算后相加。
    在三段边界上分别有 \(\displaystyle \int_{L_1}(x+y)\,\mathrm{d}s=\int_0^1 x\,\mathrm{d}x=\frac{1}{2},\qquad \int_{L_2}(x+y)\,\mathrm{d}s=\int_0^1 (1+y)\,\mathrm{d}y=\frac{3}{2}\)。
又 $\mathrm{d}s=\sqrt{2}\,\mathrm{d}x$，故 \(\displaystyle \int_{L_3}(x+y)\,\mathrm{d}s=\int_0^1 2x\sqrt{2}\,\mathrm{d}x=\sqrt{2}\)。
    因而 $\int_L (x+y)\,\mathrm{d}s=2+\sqrt{2}$。
    \item 取参数方程 \(\displaystyle x=a\cos t,\ y=a\sin t,\ -\frac{\pi}{2}\le t\le \frac{\pi}{2},\ \mathrm{d}s=a\,\mathrm{d}t\)。
    所以 \(\displaystyle \int_C (x+y)\,\mathrm{d}s
    =a^2\int_{-\pi/2}^{\pi/2}(\cos t+\sin t)\,\mathrm{d}t
    =a^2\left[\sin t-\cos t\right]_{-\pi/2}^{\pi/2}
    =2a^2\)。
    \item 将圆周写成极坐标方程 \(\displaystyle r=a\cos\theta,\ -\frac{\pi}{2}\le \theta\le \frac{\pi}{2}\)。
    因为 \(\displaystyle \mathrm{d}s=\sqrt{r^2+\left(\frac{\mathrm{d}r}{\mathrm{d}\theta}\right)^2}\,\mathrm{d}\theta=a\,\mathrm{d}\theta\)，
    又 $\sqrt{x^2+y^2}=r=a\cos\theta$，故 \(\displaystyle \int_C \sqrt{x^2+y^2}\,\mathrm{d}s
    = a^2\int_{-\pi/2}^{\pi/2}\cos\theta\,\mathrm{d}\theta
    = 2a^2\)。
\end{compactenum}
\end{solution}

\begin{exercise}[右半圆上的第一类曲线积分]
计算曲线积分 $\displaystyle \int_C y\,\mathrm{d}s$ 与 $\displaystyle \int_C |y|\,\mathrm{d}s$，其中 $C$ 是右半圆周 $x^2+y^2=a^2,\ x\ge 0,\ a>0$。
\end{exercise}
\begin{solution}
取参数方程 \(\displaystyle x=a\cos t,\quad y=a\sin t,\quad -\frac{\pi}{2}\le t\le \frac{\pi}{2},\quad \mathrm{d}s=a\,\mathrm{d}t\)。于是 \(\displaystyle \int_C y\,\mathrm{d}s=a^2\int_{-\pi/2}^{\pi/2}\sin t\,\mathrm{d}t=0,\qquad \int_C |y|\,\mathrm{d}s=a^2\int_{-\pi/2}^{\pi/2}|\sin t|\,\mathrm{d}t=2a^2\)。
故 $\int_C y\,\mathrm{d}s=0,\ \int_C |y|\,\mathrm{d}s=2a^2$。
\end{solution}

\begin{exercise}[空间圆周上的平方函数积分]
计算曲线积分 $\displaystyle \int_L x^2\,\mathrm{d}s$，其中 $L$ 是球面 $x^2+y^2+z^2=a^2$ 与平面 $x+y+z=0$ 的交线。
\end{exercise}
\begin{solution}
由对称性，交圆上三个坐标的地位完全相同，因此 \(\displaystyle 3\int_L x^2\,\mathrm{d}s=\int_L (x^2+y^2+z^2)\,\mathrm{d}s=\int_L a^2\,\mathrm{d}s\)。
交线 $L$ 是半径为 $a$ 的大圆，其周长为 $2\pi a$，故
\(\displaystyle 3\int_L x^2\,\mathrm{d}s=a^2\cdot 2\pi a=2\pi a^3\)，
从而 $\boxed{\int_L x^2\,\mathrm{d}s=\frac{2\pi a^3}{3}}$。
\end{solution}

\subsubsection{B组}

\begin{exercise}[扇形边界上的第一类曲线积分]
求 \(\displaystyle \int_L \mathrm{e}^{\sqrt{x^2+y^2}}\,\mathrm{d}s\)，其中 $L$ 为圆周 $x^2+y^2=a^2$、直线 $y=x$ 及 $x$ 轴在第一象限内所围成的扇形的整个边界。
\end{exercise}
\begin{solution}
扇形边界 $L$ 由三部分组成：沿 $x$ 轴的线段 $L_1$，圆弧 $L_2:r=a,\ 0\le \theta\le \frac{\pi}{4}$，以及沿直线 $y=x$ 的线段 $L_3$。
在 $L_1$ 上，$\sqrt{x^2+y^2}=x$，且 $\mathrm{d}s=\mathrm{d}x$，故 \(\displaystyle \int_{L_1}\mathrm{e}^{\sqrt{x^2+y^2}}\,\mathrm{d}s=\int_0^a \mathrm{e}^x\,\mathrm{d}x=\mathrm{e}^a-1\)。
在 $L_2$ 上，$\sqrt{x^2+y^2}=a$，且 $\mathrm{d}s=a\,\mathrm{d}\theta$，故 \(\displaystyle \int_{L_2}\mathrm{e}^{\sqrt{x^2+y^2}}\,\mathrm{d}s=\int_0^{\pi/4}\mathrm{e}^a a\,\mathrm{d}\theta=\frac{\pi a}{4}\mathrm{e}^a\)。
在 $L_3$ 上可取极坐标参数 $r$，因为 $\theta=\dfrac{\pi}{4}$ 为常数，所以 $\mathrm{d}s=\mathrm{d}r$，从而
\[
\int_{L_3}\mathrm{e}^{\sqrt{x^2+y^2}}\,\mathrm{d}s=\int_0^a \mathrm{e}^r\,\mathrm{d}r=\mathrm{e}^a-1.
\qquad
\boxed{\int_L \mathrm{e}^{\sqrt{x^2+y^2}}\,\mathrm{d}s=2(\mathrm{e}^a-1)+\frac{\pi a}{4}\mathrm{e}^a}.
\]
\end{solution}

\begin{exercise}[空间曲线的质量]
求曲线 \(\displaystyle x=at,\ y=\frac{a}{2}t^2,\ z=\frac{a}{3}t^3\ (0\le t\le 1)\) 的质量，其中密度函数 $\rho=\sqrt{\dfrac{2y}{a}}$。
\end{exercise}
\begin{solution}
由 $y=\dfrac{a}{2}t^2$ 可知 $\rho=\sqrt{\dfrac{2y}{a}}=\sqrt{t^2}=t\ (0\le t\le 1)$。又
\[
\vec{r}(t)=\left(at,\frac{a}{2}t^2,\frac{a}{3}t^3\right),\quad
\vec{r}^{\,\prime}(t)=\left(a,at,at^2\right),\quad
\|\vec{r}^{\,\prime}(t)\|=a\sqrt{1+t^2+t^4}.
\]
于是曲线的质量为 \(\displaystyle m=\int_0^1 \rho\,\|\vec{r}^{\,\prime}(t)\|\,\mathrm{d}t=a\int_0^1 t\sqrt{1+t^2+t^4}\,\mathrm{d}t\)。
令 $u=t^2+\dfrac{1}{2}$，则 \(\displaystyle 1+t^2+t^4=u^2+\frac{3}{4},\ 2t\,\mathrm{d}t=\mathrm{d}u\)。
从而
\[
m=\frac{a}{2}\int_{1/2}^{3/2}\sqrt{u^2+\frac{3}{4}}\,\mathrm{d}u
=\frac{a}{4}\left[u\sqrt{u^2+\frac{3}{4}}+\frac{3}{4}\ln\left(u+\sqrt{u^2+\frac{3}{4}}\right)\right]_{1/2}^{3/2}.
\]
整理得 \(\displaystyle \boxed{
m=\frac{a}{16}\left(6\sqrt{3}-2+3\ln\frac{3+2\sqrt{3}}{3}\right)
}\)。
\end{solution}

\begin{exercise}[线质量分布的转动惯量与重心]
设在 $Oxy$ 平面内有一分布着质量的曲线 $L$，在点 $(x,y)$ 处它的线密度为 $\rho(x,y)$。试用第一类曲线积分分别表达：
\begin{shortexenum}
    \shortitem{1}{曲线 $L$ 对 $x$ 轴、对 $y$ 轴的转动惯量 $I_x,\ I_y$；} &
    \shortitem{2}{曲线 $L$ 的重心坐标 $\bar{x},\ \bar{y}$。}
\end{shortexenum}
\end{exercise}
\begin{solution}
设曲线总质量为 \(\displaystyle M=\int_L \rho(x,y)\,\mathrm{d}s\)，则
\begin{compactenum}
    \item 曲线 $L$ 对坐标轴的转动惯量为 \(\displaystyle I_x=\int_L y^2\rho(x,y)\,\mathrm{d}s,\quad I_y=\int_L x^2\rho(x,y)\,\mathrm{d}s\)。
    这是因为到 $x$ 轴的距离为 $|y|$，到 $y$ 轴的距离为 $|x|$。
    \item 曲线 $L$ 的重心坐标为 \(\displaystyle \bar{x}=\frac{1}{M}\int_L x\rho(x,y)\,\mathrm{d}s,\quad \bar{y}=\frac{1}{M}\int_L y\rho(x,y)\,\mathrm{d}s\)。
\end{compactenum}
\end{solution}

\begin{exercise}[圆弧的质心]
求半径为 $a$、中心角为 $2\varphi$ 的均匀圆弧（密度 $\rho=1$）的质心。
\end{exercise}
\begin{solution}
设圆心在原点，圆弧关于 $x$ 轴对称，其方程可写为 \(\displaystyle x=a\cos\theta,\ y=a\sin\theta,\ -\varphi\le \theta\le \varphi\)。因为圆弧均匀，线密度为 $1$，故总质量等于弧长：\(\displaystyle M=\int_L \mathrm{d}s=\int_{-\varphi}^{\varphi} a\,\mathrm{d}\theta=2a\varphi\)。
由对称性可知 $\bar{y}=0$。再计算 $\bar{x}$：\(\displaystyle \bar{x}
=\frac{1}{M}\int_L x\,\mathrm{d}s
=\frac{1}{2a\varphi}\int_{-\varphi}^{\varphi} a\cos\theta\cdot a\,\mathrm{d}\theta
=\frac{a}{2\varphi}\int_{-\varphi}^{\varphi}\cos\theta\,\mathrm{d}\theta
=\frac{a\sin\varphi}{\varphi}\)。
因此质心位于圆弧的对称轴上，到圆心的距离为 $\dfrac{a\sin\varphi}{\varphi}$。若取对称轴为 $x$ 轴，则 $\boxed{\left(\bar{x},\bar{y}\right)=\left(\frac{a\sin\varphi}{\varphi},\,0\right)}$。
\end{solution}


\subsection{第二节习题：第二类曲线积分}

\subsubsection{A组}

\begin{exercise}[基本概念]
回答下列问题：
\begin{shortexenum}
    \shortitem{1}{质点在变力 $\vec{F}$ 作用下沿曲线 $L$ 移动，变力 $\vec{F}$ 所作的功如何表示？} &
    \shortitem{2}{第二类曲线积分的定义是怎样的？} \\
    \shortitem{3}{第二类曲线积分有几个等价形式？} &
    \shortitem{4}{第二类曲线积分的计算公式是怎样的？}
\end{shortexenum}
\end{exercise}
\begin{solution}
\begin{compactenum}
    \item 若 $\vec{F}=(P,Q,R)$，则质点沿有向曲线 $L$ 运动时，变力所作的功表示为 \(\displaystyle W=\int_L \vec{F}\cdot \mathrm{d}\vec{r}=\int_L P\,\mathrm{d}x+Q\,\mathrm{d}y+R\,\mathrm{d}z\)。
    在平面情形中，就是 $W=\int_L P\,\mathrm{d}x+Q\,\mathrm{d}y$。
    \item 把有向曲线 $L$ 分成若干小弧段，记相应位移向量为 $\Delta \vec{r}_i$，在每段上任取一点 $\xi_i$，作和式 \(\displaystyle \sum_{i=1}^{n}\vec{F}(\xi_i)\cdot \Delta \vec{r}_i\)。当分割无限加细时，若该和式极限存在，则称该极限为向量场 $\vec{F}$ 沿有向曲线 $L$ 的第二类曲线积分，记作 $\displaystyle \int_L \vec{F}\cdot \mathrm{d}\vec{r}$。
    \item 常用的等价形式有四种：第一种是 \(\displaystyle \int_L \vec{F}\cdot \mathrm{d}\vec{r}\)，第二种是 \(\displaystyle \int_L P\,\mathrm{d}x+Q\,\mathrm{d}y+R\,\mathrm{d}z\)，第三种是 \(\displaystyle \int_L (P\cos\alpha+Q\cos\beta+R\cos\gamma)\,\mathrm{d}s\)，第四种是 \(\displaystyle \int_{\alpha}^{\beta}(Px'(t)+Qy'(t)+Rz'(t))\,\mathrm{d}t\)。
    \item 若曲线参数方程为 \(\displaystyle x=x(t),\ y=y(t),\ z=z(t),\ \alpha\le t\le \beta\)，则
    \begin{align*}
    \int_L P\,\mathrm{d}x+Q\,\mathrm{d}y+R\,\mathrm{d}z
    &= \int_{\alpha}^{\beta}\Bigl[
    P(x(t),y(t),z(t))x'(t) \\
    &\qquad +Q(x(t),y(t),z(t))y'(t)+R(x(t),y(t),z(t))z'(t)
    \Bigr]\mathrm{d}t.
    \end{align*}
\end{compactenum}
\end{solution}

\begin{exercise}[第二类曲线积分的计算]
计算下列第二类曲线积分：
\begin{compactenum}
    \item $\displaystyle \int_{OA}x\,\mathrm{d}y+y\,\mathrm{d}x$，其中 $O$ 为坐标原点，$A(1,2)$，并设：
    \begin{shortexenum}
        \shortitem{1}{$OA$ 为直线段；} &
        \shortitem{2}{$OA$ 为抛物线 $y=2x^2$；} \\
        \shortitem{3}{$OA$ 为由 $Ox$ 轴上的线段 $OB$ 与平行于 $Oy$ 轴的线段 $BA$ 所组成的折线。} &
    \end{shortexenum}
    \item $\displaystyle \int_L (x^2-2xy)\,\mathrm{d}x+(y^2-2xy)\,\mathrm{d}y$，其中 $L$ 是抛物线 $y^2=x$ 上从点 $(1,-1)$ 到点 $(1,1)$ 的一段弧；
    \item $\displaystyle \int_{AB} y\,\mathrm{d}x-x\,\mathrm{d}y$，其中 $A(1,1)$，$B(2,4)$，方向从 $A$ 到 $B$，并设：
    \begin{shortexenum}
        \shortitem{1}{$\widehat{AB}$ 为直线段 $AB$；} &
        \shortitem{2}{$\widehat{AB}$ 为抛物线段 $y=x^2\ (1\le x\le 2)$；} \\
        \shortitem{3}{$\widehat{AB}$ 为折线段 $\widehat{AC}+\widehat{CB}$，其中 $C(2,1)$。} &
    \end{shortexenum}
    \item $\displaystyle \oint_C \frac{(x+y)\,\mathrm{d}x-(x-y)\,\mathrm{d}y}{x^2+y^2}$，其中 $C$ 是依逆时针方向通过的圆周 $x^2+y^2=a^2$。
\end{compactenum}
\end{exercise}
\begin{solution}
\begin{compactenum}
    \item 注意到 \(\displaystyle \mathrm{d}(xy)=x\,\mathrm{d}y+y\,\mathrm{d}x,\ \int_{OA}x\,\mathrm{d}y+y\,\mathrm{d}x=(xy)\big|_O^A=2\)，
    因而三种路径下结果都为 $\boxed{2}$。
    \item 由抛物线方程 $x=y^2$，取参数 \(\displaystyle x=t^2,\ y=t,\ -1\le t\le 1,\ \mathrm{d}x=2t\,\mathrm{d}t,\ \mathrm{d}y=\mathrm{d}t\)。
    所以 \(\displaystyle I=\int_{-1}^{1}\left[(t^4-2t^3)2t+(t^2-2t^3)\right]\mathrm{d}t=\int_{-1}^{1}\left(2t^5-4t^4-2t^3+t^2\right)\mathrm{d}t\)。
    其中奇函数项积分为零，因此 \(\displaystyle I=\int_{-1}^{1}(t^2-4t^4)\,\mathrm{d}t=-\frac{14}{15}\)。
    \item 分别计算三条路径上的积分。
    \textbf{(a) 直线段 $AB$：} 由 $y=3x-2$，$1\le x\le 2$，得 \(\displaystyle \mathrm{d}y=3\,\mathrm{d}x,\ \int_{AB} y\,\mathrm{d}x-x\,\mathrm{d}y=\int_1^2[(3x-2)-3x]\,\mathrm{d}x=-2\)。
    \textbf{(b) 抛物线段 $y=x^2$：}\(\displaystyle \int_{AB} y\,\mathrm{d}x-x\,\mathrm{d}y=\int_1^2 (x^2-2x^2)\,\mathrm{d}x=-\frac{7}{3}\)。
    \textbf{(c) 折线段 $AC+CB$：}\(\displaystyle \int_{AC} y\,\mathrm{d}x-x\,\mathrm{d}y=\int_1^2 \mathrm{d}x=1\)，
    \(\displaystyle \int_{CB} y\,\mathrm{d}x-x\,\mathrm{d}y=-2\int_1^4 \mathrm{d}y=-6\)，
    因此 \(\displaystyle \int_{AC+CB} y\,\mathrm{d}x-x\,\mathrm{d}y=-5\)。
    综上，$\boxed{(a)\ -2,\qquad (b)\ -\frac{7}{3},\qquad (c)\ -5}$。
    \item 取参数方程 \(\displaystyle x=a\cos t,\ y=a\sin t,\ 0\le t\le 2\pi,\ \mathrm{d}x=-a\sin t\,\mathrm{d}t,\ \mathrm{d}y=a\cos t\,\mathrm{d}t\)。
    因而
    \begin{align*}
    \oint_C \frac{(x+y)\,\mathrm{d}x-(x-y)\,\mathrm{d}y}{x^2+y^2}
    &=\int_0^{2\pi}\frac{a(\cos t+\sin t)(-a\sin t)-a(\cos t-\sin t)(a\cos t)}{a^2}\,\mathrm{d}t \\
    &=\int_0^{2\pi}(-1)\,\mathrm{d}t
    =-2\pi.
    \end{align*}
\end{compactenum}
\end{solution}

\begin{exercise}[力场作功]
求力场 $\vec{F}$ 对运动的单位质点所作的功，此质点沿曲线 $C$ 从 $A$ 点运动到 $B$ 点：
\begin{shortsingleenum}
    \shortsingleitem{1}{$\vec{F}=(x-2xy^2)\vec{i}+(y-2x^2y)\vec{j}$，$C$ 为平面曲线 $y=x^2$，$A(0,0)$，$B(1,1)$；}
    \shortsingleitem{2}{$\vec{F}=(x+y)\vec{i}+xy\vec{j}$，$C$ 为平面曲线 $y=1-|1-x|$，$A(0,0)$，$B(2,0)$；}
    \shortsingleitem{3}{$\vec{F}=(x-y)\vec{i}+(y-z)\vec{j}+(z-x)\vec{k}$，$C$ 为空间曲线 $\vec{r}(t)=t\vec{i}+t^2\vec{j}+t^3\vec{k}$，$A(0,0,0)$，$B(1,1,1)$；}
    \shortsingleitem{4}{$\vec{F}=y^2\vec{i}+z^2\vec{j}+x^2\vec{k}$，$C$ 为曲线 $\vec{r}(t)=a\cos t\,\vec{i}+b\sin t\,\vec{j}+ct\,\vec{k}$，$a,b,c$ 为正常数，$A(a,0,0)$，$B(a,0,2\pi c)$。}
\end{shortsingleenum}
\end{exercise}
\begin{solution}
\begin{compactenum}
    \item 记 \(\displaystyle P=x-2xy^2,\ Q=y-2x^2y,\ \frac{\partial P}{\partial y}=-4xy,\ \frac{\partial Q}{\partial x}=-4xy\)，故该场为保守场。设势函数为 $u$，则 \(\displaystyle u_x=x-2xy^2\Longrightarrow
    u=\frac{x^2}{2}-x^2y^2+g(y),\quad
    u(x,y)=\frac{x^2+y^2}{2}-x^2y^2\)。
    所作的功为 $W=u(1,1)-u(0,0)=0$。
    \item 曲线 $C$ 分成两段：$C_1:\ y=x,\ 0\le x\le 1;\ C_2:\ y=2-x,\ 1\le x\le 2$。
    在 $C_1$ 上，\(\displaystyle \int_{C_1}(x+y)\,\mathrm{d}x+xy\,\mathrm{d}y=\int_0^1 (2x+x^2)\,\mathrm{d}x=\frac{4}{3}\)；
    在 $C_2$ 上，\(\displaystyle \int_{C_2}(x+y)\,\mathrm{d}x+xy\,\mathrm{d}y=\int_1^2 \left[2-x(2-x)\right]\mathrm{d}x=\frac{4}{3}\)。
    故 $W=\frac{8}{3}$。
    \item 取参数 $0\le t\le 1$，则 \(\displaystyle x=t,\ y=t^2,\ z=t^3,\ \mathrm{d}x=\mathrm{d}t,\ \mathrm{d}y=2t\,\mathrm{d}t,\ \mathrm{d}z=3t^2\,\mathrm{d}t\)。
    所以 \(\displaystyle W=\int_0^1 \Big[(t-t^2)+(t^2-t^3)2t+(t^3-t)3t^2\Big]\,\mathrm{d}t\)。化简被积函数后，\(\displaystyle W=\int_0^1 \left(t-t^2-t^3-2t^4+3t^5\right)\mathrm{d}t=\frac{1}{60}\)。
    \item 由参数方程 \(\displaystyle x=a\cos t,\quad y=b\sin t,\quad z=ct,\qquad 0\le t\le 2\pi\)，得 \(\displaystyle \mathrm{d}x=-a\sin t\,\mathrm{d}t,\quad
    \mathrm{d}y=b\cos t\,\mathrm{d}t,\quad
    \mathrm{d}z=c\,\mathrm{d}t\)。
    因而
    \begin{align*}
    W
    &=\int_0^{2\pi}\left(y^2\,\mathrm{d}x+z^2\,\mathrm{d}y+x^2\,\mathrm{d}z\right)
    =\int_0^{2\pi}\left[-ab^2\sin^3 t+b c^2 t^2\cos t+a^2 c\cos^2 t\right]\mathrm{d}t,\\
    \int_0^{2\pi}\sin^3 t\,\mathrm{d}t&=0,\qquad
    \int_0^{2\pi} t^2\cos t\,\mathrm{d}t=4\pi,\\
    \int_0^{2\pi}\cos^2 t\,\mathrm{d}t&=\pi,\qquad
    \therefore\quad W=\pi a^2 c+4\pi b c^2.
    \end{align*}
\end{compactenum}
\end{solution}

\begin{exercise}[含绝对值的闭曲线积分]
求 \(\displaystyle I=\oint_C |y|\,\mathrm{d}x+|x|\,\mathrm{d}y\)，
其中 $C$ 为以点 $A(1,0)$，$B(0,1)$ 及 $C(-1,0)$ 为顶点的三角形的边界曲线，取逆时针方向。
\end{exercise}
\begin{solution}
按逆时针方向将边界分成三段：$AB$，$BC$，$CA$。

在 $AB$ 上，$y=1-x$，$x$ 从 $1$ 变到 $0$，且 $x,y\ge 0$，故 \(\displaystyle |x|=x,\ |y|=y,\ \mathrm{d}y=-\mathrm{d}x\)，并且 \(\displaystyle \int_{AB}|y|\,\mathrm{d}x+|x|\,\mathrm{d}y
=\int_1^0 \big[(1-x)-x\big]\,\mathrm{d}x=0\)。

在 $BC$ 上，$y=x+1$，$x$ 从 $0$ 变到 $-1$，且 $x\le 0,\ y\ge 0$，故 \(\displaystyle |x|=-x,\ |y|=x+1,\ \mathrm{d}y=\mathrm{d}x\)，并且 \(\displaystyle \int_{BC}|y|\,\mathrm{d}x+|x|\,\mathrm{d}y
=\int_0^{-1}\big[(x+1)-x\big]\,\mathrm{d}x=-1\)。

在 $CA$ 上，$y=0$，故该段积分为 $0$。于是 $\boxed{I=-1}$。
\end{solution}

\subsubsection{B组}

\begin{exercise}[外法向量与法向导数的曲线积分表示]
设 $C$ 是一条简单的光滑闭曲线，取逆时针方向为正向，其参数方程为 $x=x(s),\ y=y(s)$，其中 $s$ 为弧长。
\begin{shortsingleenum}
    \shortsingleitem{1}{若记 $C$ 上的单位切向量为 \(\displaystyle \vec{e}_\tau=\frac{\mathrm{d}x}{\mathrm{d}s}\vec{i}+\frac{\mathrm{d}y}{\mathrm{d}s}\vec{j}\)，证明 $C$ 上的单位（外）法向量为 \(\displaystyle \vec{e}_n=\frac{\mathrm{d}y}{\mathrm{d}s}\vec{i}-\frac{\mathrm{d}x}{\mathrm{d}s}\vec{j}\)。}
    \shortsingleitem{2}{若 $u(x,y)$ 是可微函数，试将 \(\displaystyle \oint_C \frac{\partial u}{\partial n}\,\mathrm{d}s\) 化为对坐标的曲线积分。}
\end{shortsingleenum}
\end{exercise}
\begin{solution}
\begin{compactenum}
    \item 因为 $s$ 是弧长参数，所以 \(\displaystyle \left(\frac{\mathrm{d}x}{\mathrm{d}s}\right)^2+\left(\frac{\mathrm{d}y}{\mathrm{d}s}\right)^2=1\)。
    将单位切向量顺时针旋转 $90^\circ$（即将 $(a,b)$ 变为 $(b,-a)$），便得到单位外法向量 \(\displaystyle \vec{e}_n=\frac{\mathrm{d}y}{\mathrm{d}s}\vec{i}-\frac{\mathrm{d}x}{\mathrm{d}s}\vec{j}\)。
    它与 $\vec{e}_\tau$ 垂直且模长为 $1$。
    \item 由方向导数定义并同乘 $\mathrm{d}s$，得 \(\displaystyle \frac{\partial u}{\partial n}=\nabla u\cdot \vec{e}_n=u_x\frac{\mathrm{d}y}{\mathrm{d}s}-u_y\frac{\mathrm{d}x}{\mathrm{d}s}\)，
    从而 \(\displaystyle \frac{\partial u}{\partial n}\,\mathrm{d}s=u_x\,\mathrm{d}y-u_y\,\mathrm{d}x\)。
    因此 \(\displaystyle \boxed{\oint_C \frac{\partial u}{\partial n}\,\mathrm{d}s=\oint_C \left(u_x\,\mathrm{d}y-u_y\,\mathrm{d}x\right)}\)。
\end{compactenum}
\end{solution}

\begin{exercise}[平面流体的流量公式]
设有稳定的流体运动（即流速不随时间改变），流体层充分薄，可看成一个平面问题。每点处的流速可表示为向量 \(\displaystyle \vec{v}(x,y)=P(x,y)\vec{i}+Q(x,y)\vec{j}\)。平面上给定曲线 $C$，并给定了单位法向量的指向。
\begin{compactenum}
    \item 用微元法证明：单位时间内流出曲线 $C$ 的流量微元为
    \[
    \mathrm{d}q(x,y)=
    \left[P(x,y)\cos(\vec{e}_n,x)+Q(x,y)\cos(\vec{e}_n,y)\right]\mathrm{d}s;
    \]
    \item 用微元法证明：单位时间内从区域 $D$（$D$ 为 $C$ 所围区域）内渗出来或漏下去的流量微元为
    \[
    \mathrm{d}q'(x,y)=
    \left(\frac{\partial P}{\partial x}+\frac{\partial Q}{\partial y}\right)\mathrm{d}x\,\mathrm{d}y;
    \]
    \item 证明：流体通过 $C$ 的流量为
    \[
    \oint_C \left[P\cos(\vec{e}_n,x)+Q\cos(\vec{e}_n,y)\right]\mathrm{d}s
    =\iint_D\left(\frac{\partial P}{\partial x}+\frac{\partial Q}{\partial y}\right)\mathrm{d}x\,\mathrm{d}y.
    \]
\end{compactenum}
\end{exercise}
\begin{solution}
\begin{compactenum}
    \item 在曲线 $C$ 上取弧长微元 $\mathrm{d}s$。流速向量在法向 $\vec e_n$ 上的分量为 \(\displaystyle \vec v\cdot \vec e_n=P\cos(\vec e_n,x)+Q\cos(\vec e_n,y)\)，因此 \(\displaystyle \mathrm{d}q(x,y)=\left[P\cos(\vec e_n,x)+Q\cos(\vec e_n,y)\right]\mathrm{d}s\)。
    \item 在区域 $D$ 内取小矩形 $\mathrm{d}x\times \mathrm{d}y$。沿 $x$ 方向与 $y$ 方向的净流出量近似为
    \begin{align*}
    \big[P(x+\mathrm{d}x,y)-P(x,y)\big]\mathrm{d}y
    &\approx \frac{\partial P}{\partial x}\,\mathrm{d}x\,\mathrm{d}y,\\
    \big[Q(x,y+\mathrm{d}y)-Q(x,y)\big]\mathrm{d}x
    &\approx \frac{\partial Q}{\partial y}\,\mathrm{d}x\,\mathrm{d}y,\\
    \mathrm{d}q'(x,y)
    &=\left(\frac{\partial P}{\partial x}+\frac{\partial Q}{\partial y}\right)\mathrm{d}x\,\mathrm{d}y.
    \end{align*}
    这个量 $\frac{\partial P}{\partial x}+\frac{\partial Q}{\partial y}$ 后面会被正式命名为向量场 $(P,Q)$ 的\textbf{散度}（divergence），它度量的正是每一点的"净流出率"。
    \item 将边界流量对整个边界积分，将内部流量微元对整个区域积分。由于区域内部公共边上的流量相互抵消，最后只剩边界总流量，因此 \(\displaystyle \oint_C \left[P\cos(\vec e_n,x)+Q\cos(\vec e_n,y)\right]\mathrm{d}s=\iint_D\left(\frac{\partial P}{\partial x}+\frac{\partial Q}{\partial y}\right)\mathrm{d}x\,\mathrm{d}y\)。
    这就是平面通量形式的 Green 公式。
\end{compactenum}
\end{solution}


\subsection{第三节习题：第一类曲面积分}

\subsubsection{A组}

\begin{exercise}[基本概念]
回答下列问题：
\begin{shortexenum}
    \shortitem{1}{光滑曲面 $S$ 的面积怎样定义和计算？} &
    \shortitem{2}{第一类曲面积分是怎样定义的？}
\end{shortexenum}
\end{exercise}
\begin{solution}
\begin{compactenum}
    \item 光滑曲面的面积可看成曲面片面积和的极限。若曲面 $S$ 的参数方程为
    \[
    \vec{r}=\vec{r}(u,v),\qquad (u,v)\in D,
    \qquad
    \mathrm{d}S=\|\vec{r}_u\times \vec{r}_v\|\,\mathrm{d}u\,\mathrm{d}v,
    \qquad
    \operatorname{Area}(S)=\iint_D \|\vec{r}_u\times \vec{r}_v\|\,\mathrm{d}u\,\mathrm{d}v.
    \]
    特别地，当曲面表示为 $z=z(x,y)$ 时，有
    \(\displaystyle \mathrm{d}S=\sqrt{1+z_x^2+z_y^2}\,\mathrm{d}x\,\mathrm{d}y\)。
    \item 把曲面 $S$ 分成许多小曲面片，其面积为 $\Delta S_i$，在每片上任取一点 $\xi_i$，作和式 $\sum_{i=1}^{n} f(\xi_i)\Delta S_i$。当分割无限加细时，若该和式极限存在，且与分割方式和取点方式无关，则称该极限为函数 $f$ 在曲面 $S$ 上的第一类曲面积分，记作 $\iint_S f\,\mathrm{d}S$。若 $S$ 的参数方程为 $\vec{r}(u,v)$，则
    \[
    \iint_S f\,\mathrm{d}S=\iint_D f(\vec{r}(u,v))\|\vec{r}_u\times \vec{r}_v\|\,\mathrm{d}u\,\mathrm{d}v.
    \]
\end{compactenum}
\end{solution}

\begin{exercise}[第一类曲面积分的计算]
计算下列第一类曲面积分：
\begin{shortexenum}
    \shortitem{1}{$\displaystyle \iint_S (x^2+y^2+z^2)\,\mathrm{d}S$，其中 $S$ 是 $x=0$，$y=0$ 及 $x^2+y^2+z^2=a^2\ (x\ge 0,\ y\ge 0)$ 所围成的闭曲面；} &
    \shortitem{2}{$\displaystyle \iint_S (x^2+y^2)\,\mathrm{d}S$，其中 $S$ 为立体 $\sqrt{x^2+y^2}\le z\le 1$ 的边界曲面；}\\
    \shortitem{3}{$\displaystyle \iint_S \sqrt{x^2+y^2}\,\mathrm{d}S$，其中 $S$ 是锥面 $z=\sqrt{x^2+y^2}$ 被平面 $z=h\ (h>0)$ 所截部分；} &
    \shortitem{4}{$\displaystyle \iint_S (x+y+z)\,\mathrm{d}S$，其中 $S$ 为上半球面 $x^2+y^2+z^2=a^2,\ z\ge 0$；}\\
    \shortitem{5}{$\displaystyle \iint_S z\,\mathrm{d}S$，其中 $S$ 为曲面 $z=\dfrac{1}{2}(x^2+y^2)$ 被平面 $z=2$ 截下的有限部分；} &
    \shortitem{6}{$\displaystyle \iint_S y\,\mathrm{d}S$，其中 $S$ 是平面 $x+y+z=4$ 被圆柱面 $x^2+y^2=1$ 截出的有限部分。}
\end{shortexenum}
\end{exercise}
\begin{solution}
\begin{compactenum}
    \item 曲面 $S$ 由三部分组成：球面的四分之一部分 $S_1$，平面 $x=0$ 上的半圆盘 $S_2$，平面 $y=0$ 上的半圆盘 $S_3$。
    在球面 $S_1$ 上，$x^2+y^2+z^2=a^2$，故
    \(\displaystyle \iint_{S_1}(x^2+y^2+z^2)\,\mathrm{d}S=a^2\cdot \operatorname{Area}(S_1)=a^2\cdot \pi a^2=\pi a^4\)。
    在 $S_2$ 上，$x=0$，被积函数化为 $y^2+z^2$，改用平面极坐标得
    \(\displaystyle \iint_{S_2}(y^2+z^2)\,\mathrm{d}S=\int_0^\pi\int_0^a r^2\cdot r\,\mathrm{d}r\,\mathrm{d}\theta=\frac{\pi a^4}{4}\)。
    对于 $S_3$ 计算完全对应，因此
    \[
    \iint_{S_3}(x^2+z^2)\,\mathrm{d}S=\frac{\pi a^4}{4},\qquad
    \boxed{\iint_S (x^2+y^2+z^2)\,\mathrm{d}S=\frac{3\pi a^4}{2}}.
    \]
    \item 曲面边界由锥侧面 $S_1:\ z=\sqrt{x^2+y^2}$ 与顶圆盘 $S_2:\ z=1$ 组成。
    在顶圆盘上
    \(\displaystyle \iint_{S_2}(x^2+y^2)\,\mathrm{d}S=\int_0^{2\pi}\int_0^1 r^2\cdot r\,\mathrm{d}r\,\mathrm{d}\theta=\frac{\pi}{2}\)，
    在锥侧面上
    \begin{align*}
    \mathrm{d}S&=\sqrt{2}\,r\,\mathrm{d}r\,\mathrm{d}\theta,\qquad x^2+y^2=r^2,\\
    \iint_{S_1}(x^2+y^2)\,\mathrm{d}S&=\int_0^{2\pi}\int_0^1 r^2\cdot \sqrt{2}r\,\mathrm{d}r\,\mathrm{d}\theta=\frac{\pi\sqrt{2}}{2},\\
    \boxed{\iint_S (x^2+y^2)\,\mathrm{d}S}&=\frac{\pi}{2}(1+\sqrt{2}).
    \end{align*}
    \item 在锥面上取参数 $(r,\theta)$，其中 $0\le r\le h$，$0\le \theta\le 2\pi$。则 \(\displaystyle \sqrt{x^2+y^2}=r,\qquad \mathrm{d}S=\sqrt{2}\,r\,\mathrm{d}r\,\mathrm{d}\theta\)，
    故
    \[
    \iint_S \sqrt{x^2+y^2}\,\mathrm{d}S
    =\int_0^{2\pi}\int_0^h r\cdot \sqrt{2}r\,\mathrm{d}r\,\mathrm{d}\theta
    =\frac{2\pi\sqrt{2}}{3}h^3.
    \]
    \item 由对称性，
    \[
    \iint_S x\,\mathrm{d}S=\iint_S y\,\mathrm{d}S=0,\qquad
    \iint_S (x+y+z)\,\mathrm{d}S=\iint_S z\,\mathrm{d}S.
    \]
    在球坐标下
    \[
    z=a\cos\varphi,\qquad \mathrm{d}S=a^2\sin\varphi\,\mathrm{d}\varphi\,\mathrm{d}\theta,\qquad
    0\le \theta\le 2\pi,\quad 0\le \varphi\le \frac{\pi}{2}.
    \]
    于是
    \[
    \iint_S z\,\mathrm{d}S
    =\int_0^{2\pi}\int_0^{\pi/2} a\cos\varphi\cdot a^2\sin\varphi\,\mathrm{d}\varphi\,\mathrm{d}\theta
    =2\pi a^3\int_0^{\pi/2}\sin\varphi\cos\varphi\,\mathrm{d}\varphi
    =\pi a^3.
    \]
    因而 \(\displaystyle \boxed{\iint_S (x+y+z)\,\mathrm{d}S=\pi a^3}\)。
    \item 用极坐标表示曲面
    \(\displaystyle z=\frac{r^2}{2},\qquad 0\le r\le 2,\qquad 0\le \theta\le 2\pi,\qquad
    \mathrm{d}S=\sqrt{1+r^2}\,r\,\mathrm{d}r\,\mathrm{d}\theta\)，
    所以
    \[
    \iint_S z\,\mathrm{d}S
    =\int_0^{2\pi}\int_0^2 \frac{r^2}{2}\sqrt{1+r^2}\,r\,\mathrm{d}r\,\mathrm{d}\theta
    =\pi\int_0^2 r^3\sqrt{1+r^2}\,\mathrm{d}r=\frac{2\pi}{15}(25\sqrt{5}+1).
    \]
    \item 将平面写成
    \(\displaystyle z=4-x-y,\qquad \mathrm{d}S=\sqrt{3}\,\mathrm{d}x\,\mathrm{d}y\)。
    曲面在 $xOy$ 平面上的投影为单位圆盘 $x^2+y^2\le 1$。由于该区域关于 $x$ 轴对称，而函数 $y$ 关于 $x$ 轴是奇函数，所以
    \(\displaystyle \iint_S y\,\mathrm{d}S=\sqrt{3}\iint_{x^2+y^2\le 1} y\,\mathrm{d}x\,\mathrm{d}y=0\)。
\end{compactenum}
\end{solution}

\begin{exercise}[锥面带状部分的面积]
求曲面 $z=\sqrt{x^2+y^2}$ 夹在两曲面 $x^2+y^2=y$，$x^2+y^2=2y$ 之间的那部分的面积。
\end{exercise}
\begin{solution}
把两柱面化为极坐标方程 \(\displaystyle x^2+y^2=y\Longrightarrow r=\sin\theta,\quad x^2+y^2=2y\Longrightarrow r=2\sin\theta,\quad 0\le \theta\le \pi\)。所求曲面是锥面 $z=r$ 上对应于投影区域 \(\displaystyle \sin\theta\le r\le 2\sin\theta,\ 0\le \theta\le \pi\)，
而锥面面积元素为 $\mathrm{d}S=\sqrt{2}\,r\,\mathrm{d}r\,\mathrm{d}\theta$，故
\begin{align*}
A
&=\int_0^{\pi}\int_{\sin\theta}^{2\sin\theta}\sqrt{2}\,r\,\mathrm{d}r\,\mathrm{d}\theta
=\frac{\sqrt{2}}{2}\int_0^\pi \left(4\sin^2\theta-\sin^2\theta\right)\mathrm{d}\theta \\
&=\frac{3\sqrt{2}}{2}\int_0^\pi \sin^2\theta\,\mathrm{d}\theta
=\frac{3\pi\sqrt{2}}{4}.
\end{align*}
因此 \(\displaystyle \boxed{A=\frac{3\pi\sqrt{2}}{4}}\)。
\end{solution}

\subsubsection{B组}

\begin{exercise}[球面被柱面截下部分的面积]
求球面 $x^2+y^2+z^2=a^2$ 被柱面 $x^2+y^2=ax$ 截下部分的面积。
\end{exercise}
\begin{solution}
先判断为什么适合投影到 $xOy$ 平面：球面可以写成上下两张显式曲面 $z=\pm\sqrt{a^2-x^2-y^2}$，而柱面 $x^2+y^2=ax$ 恰好直接给出投影边界，所以从投影区域出发最自然。

所求部分在 $xOy$ 平面上的投影区域为 $D:\ x^2+y^2\le ax$，在极坐标下即 \(\displaystyle 0\le r\le a\cos\theta,\ -\frac{\pi}{2}\le \theta\le \frac{\pi}{2}\)。
由于柱面截下的是整条球带，球面的上、下两半都要算，而且二者面积相同。对上半球面 $z=\sqrt{a^2-r^2}$，有
\(\displaystyle z_x=-\frac{x}{\sqrt{a^2-r^2}},\quad z_y=-\frac{y}{\sqrt{a^2-r^2}}\)，
\(\displaystyle \mathrm{d}S=\frac{a}{\sqrt{a^2-r^2}}\,\mathrm{d}A\)。
下半球面的面积元与此完全相同，所以所求面积 \(\displaystyle A=2a\int_{-\pi/2}^{\pi/2}\int_0^{a\cos\theta}\frac{r}{\sqrt{a^2-r^2}}\,\mathrm{d}r\,\mathrm{d}\theta\)。
先算内层积分：\(\displaystyle \int_0^{a\cos\theta}\frac{r}{\sqrt{a^2-r^2}}\,\mathrm{d}r
=a-a|\sin\theta|\)，其中 \(\displaystyle \sqrt{a^2-a^2\cos^2\theta}=a|\sin\theta|\)，
而不是可以随手把 $\sin\theta$ 当成非负量；因为当前角度区间跨过了 $0$，左右两侧的符号不同，故 \(\displaystyle A
= 2a^2\int_{-\pi/2}^{\pi/2}\left(1-|\sin\theta|\right)\mathrm{d}\theta
= 2a^2(\pi-2)\)，即 \(\displaystyle \boxed{A=2a^2(\pi-2)}\)。
\end{solution}

\begin{exercise}[抛物面的质量]
求抛物面 $z=2-(x^2+y^2)$ 在 $Oxy$ 平面上方部分的质量，其密度为 $\rho(x,y,z)=x^2+y^2$。
\end{exercise}
\begin{solution}
曲面在 $Oxy$ 平面上方意味着 $2-(x^2+y^2)\ge 0$，即 $x^2+y^2\le 2$。记 $r^2=x^2+y^2$，则 \(\displaystyle z=2-r^2,\ z_x=-2x,\ z_y=-2y\)，且 \(\displaystyle \mathrm{d}S=\sqrt{1+z_x^2+z_y^2}\,\mathrm{d}x\,\mathrm{d}y
=\sqrt{1+4r^2}\,r\,\mathrm{d}r\,\mathrm{d}\theta\)。
于是质量。令 $u=1+4r^2$，则 $\mathrm{d}u=8r\,\mathrm{d}r,\ r^2=\frac{u-1}{4}$，从而
\begin{align*}
M
&=\iint_S \rho\,\mathrm{d}S
=\int_0^{2\pi}\int_0^{\sqrt{2}} r^2\sqrt{1+4r^2}\,r\,\mathrm{d}r\,\mathrm{d}\theta
=2\pi\int_0^{\sqrt{2}} r^3\sqrt{1+4r^2}\,\mathrm{d}r\\
&=\frac{\pi}{16}\int_1^9 (u-1)\sqrt{u}\,\mathrm{d}u
=\frac{\pi}{16}\left[\frac{2}{5}u^{5/2}-\frac{2}{3}u^{3/2}\right]_1^9
=\frac{149\pi}{30}.
\end{align*}
因此 \(\displaystyle \boxed{M=\frac{149\pi}{30}}\)。
\end{solution}

\begin{exercise}[上半球面的质心]
求均匀曲面 $z=\sqrt{a^2-x^2-y^2}$ 的质心坐标。
\end{exercise}
\begin{solution}
该曲面是半径为 $a$ 的上半球面。由对称性可知 $\bar{x}=0,\ \bar{y}=0$；又 \(\displaystyle S=2\pi a^2,\ \iint_S z\,\mathrm{d}S=\pi a^3,\ \bar{z}=\frac{1}{S}\iint_S z\,\mathrm{d}S=\frac{\pi a^3}{2\pi a^2}=\frac{a}{2}\)。因此质心坐标为 \(\displaystyle \boxed{\left(0,0,\frac{a}{2}\right)}\)。
\end{solution}


\subsection{第四节习题：第二类曲面积分}

\subsubsection{A组}

\begin{exercise}[基本概念]
回答下列问题：
\begin{shorttrienum}
    \shortitem{1}{曲面的方程为 $z=z(x,y)$ 时，曲面的上侧与下侧如何通过法向量来表示？} &
    \shortitem{2}{第二类曲面积分的概念是如何引进的？} &
    \shortitem{3}{第二类曲面积分有哪四种等价形式？}
\end{shorttrienum}
\end{exercise}
\begin{solution}
\begin{compactenum}
    \item 当曲面写成 $z=z(x,y)$ 时，可取 \(\displaystyle \vec{n}_{\uparrow}=(-z_x,-z_y,1),\quad \frac{\vec n_\uparrow}{\|\vec n_\uparrow\|}=\frac{(-z_x,-z_y,1)}{\sqrt{1+z_x^2+z_y^2}},\quad \vec{n}_{\downarrow}=(z_x,z_y,-1)\)，分别表示上侧的法向、上侧单位法向和下侧法向。
    \item 第二类曲面积分来源于流量问题。设向量场 $\vec{F}$ 穿过定向曲面 $S$，把曲面分成许多小曲面片 $\Delta S_i$，在第 $i$ 个小片上取一点 $\xi_i$，若记该点处单位法向量为 $\vec{n}_i$，则穿过该小片的流量近似为 $\vec{F}(\xi_i)\cdot \vec{n}_i\,\Delta S_i$。对所有小片求和并令分割无限加细，若极限存在，就得到向量场穿过曲面的通量，这就是第二类曲面积分。
    \item 若 $\vec{F}=(P,Q,R)$，则第二类曲面积分常写为以下四种等价形式：
    \begin{align*}
    &\iint_S \vec{F}\cdot \vec{n}\,\mathrm{d}S,\qquad
    \iint_S P\,\mathrm{d}y\,\mathrm{d}z+Q\,\mathrm{d}z\,\mathrm{d}x+R\,\mathrm{d}x\,\mathrm{d}y,\\
    &\iint_S (P\cos\alpha+Q\cos\beta+R\cos\gamma)\,\mathrm{d}S,\qquad
    \iint_D \bigl(-Pz_x-Qz_y+R\bigr)\,\mathrm{d}x\,\mathrm{d}y
    \end{align*}
    （最后一个公式是对上侧而言；下侧取相反号）。
\end{compactenum}
\end{solution}

\begin{exercise}[第二类曲面积分的计算]
计算下列第二类曲面积分：
\begin{shortsingleenum}
    \shortsingleitem{1}{$\displaystyle \iint_S x\,\mathrm{d}y\,\mathrm{d}z+y\,\mathrm{d}z\,\mathrm{d}x+z\,\mathrm{d}x\,\mathrm{d}y$，其中 $S$ 为曲面 $z=x^2+y^2$ 在第一卦限中、$0\le z\le 1$ 之间部分的上侧；}
    \shortsingleitem{2}{$\displaystyle \iint_S \frac{x^2}{x^2+y^2}\,\mathrm{d}x\,\mathrm{d}y$，其中 $S$ 为上半球面 $z=\sqrt{2ax-x^2-y^2}\ (a>0)$ 在圆柱面 $x^2+y^2=a^2$ 的外部部分的上侧；}
    \shortsingleitem{3}{$\displaystyle \iint_S x^2y^2z\,\mathrm{d}x\,\mathrm{d}y$，其中 $S$ 为锥面 $z=\sqrt{x^2+y^2}\ (0\le z\le R)$ 的下侧；}
    \shortsingleitem{4}{$\displaystyle \iint_S x^2\,\mathrm{d}y\,\mathrm{d}z+y^2\,\mathrm{d}z\,\mathrm{d}x+z^2\,\mathrm{d}x\,\mathrm{d}y$，其中 $S$ 是球面 $x^2+y^2+z^2=R^2\ (R>0)$ 的内侧。}
\end{shortsingleenum}
\end{exercise}
\begin{solution}
\begin{compactenum}
    \item 按定理~\ref{thm:surface-integral-type2-calc}，就把 $x,y$ 当参数：
    \[
    \vec r(x,y)=(x,y,x^2+y^2),\qquad
    D=\{(x,y)\mid x\ge0,\ y\ge0,\ x^2+y^2\le1\}.
    \]
    先列偏导：
    \[
    x_x=1,\quad x_y=0,\quad y_x=0,\quad y_y=1,\quad z_x=2x,\quad z_y=2y.
    \]
    逐个算雅可比行列式：
    \[
    \begin{aligned}
    \frac{\partial(y,z)}{\partial(x,y)}
    &=\begin{vmatrix}y_x&z_x\\y_y&z_y\end{vmatrix}
    =\begin{vmatrix}0&2x\\1&2y\end{vmatrix}
    =0\cdot2y-2x\cdot1=-2x,\\
    \frac{\partial(z,x)}{\partial(x,y)}
    &=\begin{vmatrix}z_x&x_x\\z_y&x_y\end{vmatrix}
    =\begin{vmatrix}2x&1\\2y&0\end{vmatrix}
    =2x\cdot0-1\cdot2y=-2y,\\
    \frac{\partial(x,y)}{\partial(x,y)}
    &=\begin{vmatrix}x_x&y_x\\x_y&y_y\end{vmatrix}
    =\begin{vmatrix}1&0\\0&1\end{vmatrix}
    =1\cdot1-0\cdot0=1.
    \end{aligned}
    \]
    由于 $\vec r_x\times\vec r_y=(-2x,-2y,1)$ 指向上侧，故 $\epsilon=1$。于是
    \begin{align*}
    I
    &=\iint_D \bigl[x(-2x)+y(-2y)+(x^2+y^2)\bigr]\,\mathrm{d}x\,\mathrm{d}y\\
    &=\iint_D \bigl(-2x^2-2y^2+x^2+y^2\bigr)\,\mathrm{d}x\,\mathrm{d}y\\
    &=-\iint_D (x^2+y^2)\,\mathrm{d}x\,\mathrm{d}y.
    \end{align*}
    用极坐标计算得 \(\displaystyle I=-\int_0^{\pi/2}\int_0^1 r^3\,\mathrm{d}r\,\mathrm{d}\theta=-\frac{\pi}{8}\)。
    \item 仍然把 $x,y$ 当参数：
    \[
    \vec r(x,y)=\left(x,y,\sqrt{2ax-x^2-y^2}\right).
    \]
    记 $z=\sqrt{2ax-x^2-y^2}$，则
    \[
    x_x=1,\quad x_y=0,\quad y_x=0,\quad y_y=1,\quad
    z_x=\frac{a-x}{z},\quad z_y=-\frac{y}{z}.
    \]
    三个雅可比行列式为
    \[
    \begin{aligned}
    \frac{\partial(y,z)}{\partial(x,y)}
    &=\begin{vmatrix}y_x&z_x\\y_y&z_y\end{vmatrix}
    =\begin{vmatrix}0&z_x\\1&z_y\end{vmatrix}
    =0\cdot z_y-z_x\cdot1
    =-\frac{a-x}{z},\\
    \frac{\partial(z,x)}{\partial(x,y)}
    &=\begin{vmatrix}z_x&x_x\\z_y&x_y\end{vmatrix}
    =\begin{vmatrix}z_x&1\\z_y&0\end{vmatrix}
    =z_x\cdot0-1\cdot z_y
    =\frac{y}{z},\\
    \frac{\partial(x,y)}{\partial(x,y)}
    &=\begin{vmatrix}x_x&y_x\\x_y&y_y\end{vmatrix}
    =\begin{vmatrix}1&0\\0&1\end{vmatrix}=1.
    \end{aligned}
    \]
    题目取上侧，$\epsilon=1$。本题只含 $\mathrm{d}x\,\mathrm{d}y$，所以真正进入积分的就是最后一个雅可比行列式。球面的投影满足 $x^2+y^2\le 2ax$，又要求在圆柱 $x^2+y^2=a^2$ 的外部，所以
    \[
    D=\{(r,\theta)\mid a\le r\le 2a\cos\theta,\ -\frac{\pi}{3}\le \theta\le \frac{\pi}{3}\}.
    \]
    于是
    \begin{align*}
    I
    &=\int_{-\pi/3}^{\pi/3}\int_a^{2a\cos\theta}\frac{r^2\cos^2\theta}{r^2}\,r\,\mathrm{d}r\,\mathrm{d}\theta \\
    &=\int_{-\pi/3}^{\pi/3}\cos^2\theta\left[\frac{r^2}{2}\right]_{a}^{2a\cos\theta}\mathrm{d}\theta \\
    &=\frac{a^2}{2}\int_{-\pi/3}^{\pi/3}\left(4\cos^4\theta-\cos^2\theta\right)\mathrm{d}\theta \\
    &=a^2\int_0^{\pi/3}\left(4\cos^4\theta-\cos^2\theta\right)\mathrm{d}\theta.
    \end{align*}
    再利用恒等式 \(\displaystyle 4\cos^4\theta-\cos^2\theta=1+\frac{3}{2}\cos 2\theta+\frac{1}{2}\cos 4\theta\)，得
    \begin{align*}
    I
    &=a^2\int_0^{\pi/3}\left(1+\frac{3}{2}\cos 2\theta+\frac{1}{2}\cos 4\theta\right)\mathrm{d}\theta
    =a^2\left[\theta+\frac{3}{4}\sin 2\theta+\frac{1}{8}\sin 4\theta\right]_0^{\pi/3} \\
    &=a^2\left(\frac{\pi}{3}+\frac{5\sqrt{3}}{16}\right).
    \end{align*}
    \item 取参数
    \[
    \vec r(x,y)=\left(x,y,\sqrt{x^2+y^2}\right),\qquad D:\ x^2+y^2\le R^2.
    \]
    记 $r=\sqrt{x^2+y^2}$，则
    \[
    x_x=1,\quad x_y=0,\quad y_x=0,\quad y_y=1,\quad z_x=\frac{x}{r},\quad z_y=\frac{y}{r}.
    \]
    于是
    \[
    \begin{aligned}
    \frac{\partial(y,z)}{\partial(x,y)}
    &=\begin{vmatrix}y_x&z_x\\y_y&z_y\end{vmatrix}
    =\begin{vmatrix}0&x/r\\1&y/r\end{vmatrix}
    =0\cdot\frac{y}{r}-\frac{x}{r}\cdot1
    =-\frac{x}{r},\\
    \frac{\partial(z,x)}{\partial(x,y)}
    &=\begin{vmatrix}z_x&x_x\\z_y&x_y\end{vmatrix}
    =\begin{vmatrix}x/r&1\\y/r&0\end{vmatrix}
    =\frac{x}{r}\cdot0-1\cdot\frac{y}{r}
    =-\frac{y}{r},\\
    \frac{\partial(x,y)}{\partial(x,y)}
    &=\begin{vmatrix}x_x&y_x\\x_y&y_y\end{vmatrix}
    =\begin{vmatrix}1&0\\0&1\end{vmatrix}=1.
    \end{aligned}
    \]
    上述参数化对应上侧，题目取下侧，故 $\epsilon=-1$。本题只含 $\mathrm{d}x\,\mathrm{d}y$，所以
    \[
    I=-\iint_D x^2y^2\sqrt{x^2+y^2}\,\mathrm{d}x\,\mathrm{d}y.
    \]
    改用极坐标 $x=r\cos\theta,\ y=r\sin\theta$，得
    \begin{align*}
    I&=-\int_0^{2\pi}\int_0^R r^2\cos^2\theta\cdot r^2\sin^2\theta\cdot r\cdot r\,\mathrm{d}r\,\mathrm{d}\theta\\
    &=-\int_0^{2\pi}\cos^2\theta\sin^2\theta\,\mathrm{d}\theta\int_0^R r^6\,\mathrm{d}r
    =-\frac{\pi}{4}\cdot\frac{R^7}{7}=-\frac{\pi R^7}{28}.
    \end{align*}
    \item 取球面参数化
    \[
    \vec r(\varphi,\theta)
    =(R\sin\varphi\cos\theta,\ R\sin\varphi\sin\theta,\ R\cos\varphi),
    \quad 0\le\varphi\le\pi,\quad 0\le\theta\le2\pi.
    \]
    直接算三个雅可比行列式：
    \[
    \begin{aligned}
    \frac{\partial(y,z)}{\partial(\varphi,\theta)}
    &=\begin{vmatrix}
    R\cos\varphi\sin\theta&-R\sin\varphi\\
    R\sin\varphi\cos\theta&0
    \end{vmatrix}
    =R\cos\varphi\sin\theta\cdot0-(-R\sin\varphi)R\sin\varphi\cos\theta
    =R^2\sin^2\varphi\cos\theta,\\
    \frac{\partial(z,x)}{\partial(\varphi,\theta)}
    &=\begin{vmatrix}
    -R\sin\varphi&R\cos\varphi\cos\theta\\
    0&-R\sin\varphi\sin\theta
    \end{vmatrix}
    \\&=(-R\sin\varphi)(-R\sin\varphi\sin\theta)
    -R\cos\varphi\cos\theta\cdot0\\
    &=R^2\sin^2\varphi\sin\theta,\\
    \frac{\partial(x,y)}{\partial(\varphi,\theta)}
    &=\begin{vmatrix}
    R\cos\varphi\cos\theta&R\cos\varphi\sin\theta\\
    -R\sin\varphi\sin\theta&R\sin\varphi\cos\theta
    \end{vmatrix}
    \\&=R^2\sin\varphi\cos\varphi(\cos^2\theta+\sin^2\theta)\\
    &=R^2\sin\varphi\cos\varphi.
    \end{aligned}
    \]
    这组参数的 $\vec r_\varphi\times\vec r_\theta$ 指向外侧，题目取内侧，故 $\epsilon=-1$，所以
    \[
    \begin{aligned}
    \mathrm{d}y\,\mathrm{d}z&=-R^2\sin^2\varphi\cos\theta\,\mathrm{d}\varphi\,\mathrm{d}\theta,\\
    \mathrm{d}z\,\mathrm{d}x&=-R^2\sin^2\varphi\sin\theta\,\mathrm{d}\varphi\,\mathrm{d}\theta,\\
    \mathrm{d}x\,\mathrm{d}y&=-R^2\sin\varphi\cos\varphi\,\mathrm{d}\varphi\,\mathrm{d}\theta.
    \end{aligned}
    \]
    其中
    \[
    x=R\sin\varphi\cos\theta,\qquad
    y=R\sin\varphi\sin\theta,\qquad
    z=R\cos\varphi.
    \]
    于是
    \[
    \begin{aligned}
    x^2\,\mathrm{d}y\,\mathrm{d}z+y^2\,\mathrm{d}z\,\mathrm{d}x+z^2\,\mathrm{d}x\,\mathrm{d}y
    &=-R^4\bigl(\sin^4\varphi\cos^3\theta+\sin^4\varphi\sin^3\theta\\
    &\qquad\quad+\sin\varphi\cos^3\varphi\bigr)\mathrm{d}\varphi\,\mathrm{d}\theta.
    \end{aligned}
    \]
    第一项、第二项对 $\theta\in[0,2\pi]$ 积分为 $0$，第三项对 $\varphi\in[0,\pi]$ 积分为 $0$，故 \(\displaystyle I=0\)。
\end{compactenum}
\end{solution}

\begin{exercise}[长方体外表面的通量]
设 $\Omega$ 是立体 \(\displaystyle 0\le x\le a,\ 0\le y\le b,\ 0\le z\le c\)，$S$ 为 $\Omega$ 的外表面的外侧，$f(x),g(y),h(z)$ 为连续函数。求 \(\displaystyle \iint_S f(x)\,\mathrm{d}y\,\mathrm{d}z+g(y)\,\mathrm{d}z\,\mathrm{d}x+h(z)\,\mathrm{d}x\,\mathrm{d}y\)。
\end{exercise}
\begin{solution}
仍按雅可比行列式逐面计算。先看两张 $x$ 为常数的面。取参数 $(u,v)=(y,z)$，则 $x=\text{常数},\ y=u,\ z=v$，因此
\[
\begin{aligned}
\frac{\partial(y,z)}{\partial(u,v)}
&=\begin{vmatrix}y_u&z_u\\y_v&z_v\end{vmatrix}
=\begin{vmatrix}1&0\\0&1\end{vmatrix}=1,\\
\frac{\partial(z,x)}{\partial(u,v)}
&=\begin{vmatrix}z_u&x_u\\z_v&x_v\end{vmatrix}
=\begin{vmatrix}0&0\\1&0\end{vmatrix}=0,\\
\frac{\partial(x,y)}{\partial(u,v)}
&=\begin{vmatrix}x_u&y_u\\x_v&y_v\end{vmatrix}
=\begin{vmatrix}0&1\\0&0\end{vmatrix}=0.
\end{aligned}
\]
在 $x=a$ 上，该参数方向指向外侧，贡献为 $\displaystyle \int_0^b\int_0^c f(a)\,\mathrm{d}v\,\mathrm{d}u=bcf(a)$；在 $x=0$ 上，外侧相反，贡献为 $-bcf(0)$。

在 $y=b$ 与 $y=0$ 上取 $(u,v)=(z,x)$，则 $x=v,\ y=\text{常数},\ z=u$，所以
\[
\begin{aligned}
\frac{\partial(y,z)}{\partial(u,v)}
&=\begin{vmatrix}0&1\\0&0\end{vmatrix}=0,\\
\frac{\partial(z,x)}{\partial(u,v)}
&=\begin{vmatrix}1&0\\0&1\end{vmatrix}=1,\\
\frac{\partial(x,y)}{\partial(u,v)}
&=\begin{vmatrix}0&0\\1&0\end{vmatrix}=0,
\end{aligned}
\]
两面的贡献为 $acg(b)-acg(0)$。在 $z=c$ 与 $z=0$ 上取 $(u,v)=(x,y)$，此时
\[
\begin{aligned}
\frac{\partial(y,z)}{\partial(u,v)}
&=\begin{vmatrix}0&0\\1&0\end{vmatrix}=0,\\
\frac{\partial(z,x)}{\partial(u,v)}
&=\begin{vmatrix}0&1\\0&0\end{vmatrix}=0,\\
\frac{\partial(x,y)}{\partial(u,v)}
&=\begin{vmatrix}1&0\\0&1\end{vmatrix}=1,
\end{aligned}
\]
两面的贡献为 $abh(c)-abh(0)$。
因此
\[
\boxed{
\iint_S f(x)\,\mathrm{d}y\,\mathrm{d}z+g(y)\,\mathrm{d}z\,\mathrm{d}x+h(z)\,\mathrm{d}x\,\mathrm{d}y
=bc\bigl(f(a)-f(0)\bigr)+ac\bigl(g(b)-g(0)\bigr)+ab\bigl(h(c)-h(0)\bigr)
}.
\]
\end{solution}

\subsubsection{B组}

\begin{exercise}[闭曲面的第二类曲面积分]
计算曲面积分 \(\displaystyle I=\iint_S xz\,\mathrm{d}x\,\mathrm{d}y+xy\,\mathrm{d}y\,\mathrm{d}z+yz\,\mathrm{d}z\,\mathrm{d}x\)，其中 $S$ 是平面 $x=0,\ y=0,\ z=0,\ x+y+z=1$ 所围成的空间区域的整个边界曲面的外侧。
\end{exercise}
\begin{solution}
仍用雅可比行列式逐片算。三个坐标面也按同一套路处理：
\[
\begin{array}{c|c|c|c}
\text{曲面片} & \frac{\partial(y,z)}{\partial(u,v)}
& \frac{\partial(z,x)}{\partial(u,v)}
& \frac{\partial(x,y)}{\partial(u,v)}\\ \hline
x=0,\ (u,v)=(y,z) &
\begin{vmatrix}1&0\\0&1\end{vmatrix}=1 &
\begin{vmatrix}0&0\\1&0\end{vmatrix}=0 &
\begin{vmatrix}0&1\\0&0\end{vmatrix}=0\\[3ex]
y=0,\ (u,v)=(z,x) &
\begin{vmatrix}0&1\\0&0\end{vmatrix}=0 &
\begin{vmatrix}1&0\\0&1\end{vmatrix}=1 &
\begin{vmatrix}0&0\\1&0\end{vmatrix}=0\\[3ex]
z=0,\ (u,v)=(x,y) &
\begin{vmatrix}0&0\\1&0\end{vmatrix}=0 &
\begin{vmatrix}0&1\\0&0\end{vmatrix}=0 &
\begin{vmatrix}1&0\\0&1\end{vmatrix}=1
\end{array}
\]
在 $x=0$ 上，只有 $\mathrm{d}y\,\mathrm{d}z$ 可能留下，但系数 $xy=0$；在 $y=0$ 上，只有 $\mathrm{d}z\,\mathrm{d}x$ 可能留下，但系数 $yz=0$；在 $z=0$ 上，只有 $\mathrm{d}x\,\mathrm{d}y$ 可能留下，但系数 $xz=0$。因此三个坐标面贡献都为 $0$。

只剩斜面 $x+y+z=1$。令
\[
z=1-x-y,\qquad D=\{(x,y)\mid x\ge0,\ y\ge0,\ x+y\le1\}.
\]
取参数 $\vec r(x,y)=(x,y,1-x-y)$，则
\[
\begin{aligned}
    \frac{\partial(y,z)}{\partial(x,y)}
    &=\begin{vmatrix}y_x&z_x\\y_y&z_y\end{vmatrix}
    =\begin{vmatrix}0&-1\\1&-1\end{vmatrix}
    =0\cdot(-1)-(-1)\cdot1=1,\\
    \frac{\partial(z,x)}{\partial(x,y)}
    &=\begin{vmatrix}z_x&x_x\\z_y&x_y\end{vmatrix}
    =\begin{vmatrix}-1&1\\-1&0\end{vmatrix}
    =(-1)\cdot0-1\cdot(-1)=1,\\
    \frac{\partial(x,y)}{\partial(x,y)}
    &=\begin{vmatrix}1&0\\0&1\end{vmatrix}=1.
\end{aligned}
\]
又 $\vec r_x\times\vec r_y=(1,1,1)$，正好指向四面体外侧，故 $\epsilon=1$。于是
\[
\begin{aligned}
I&=\iint_D\bigl[xy+yz+xz\bigr]\,\mathrm{d}x\,\mathrm{d}y\\
&=\iint_D\bigl[xy+y(1-x-y)+x(1-x-y)\bigr]\,\mathrm{d}x\,\mathrm{d}y\\
&=\iint_D(x+y-x^2-y^2-xy)\,\mathrm{d}x\,\mathrm{d}y\\
&=\frac16+\frac16-\frac1{12}-\frac1{12}-\frac1{24}
=\frac18.
\end{aligned}
\]
所以 \(\displaystyle \boxed{I=\frac18}\)。
\end{solution}

\begin{exercise}[锥面的外表面通量]
计算曲面积分 \(\displaystyle I=\iint_S (y-z)\,\mathrm{d}y\,\mathrm{d}z+(z-x)\,\mathrm{d}z\,\mathrm{d}x+(x-y)\,\mathrm{d}x\,\mathrm{d}y\)，其中 $S$ 为锥面 $x^2+y^2=z^2\ (0\le z\le h)$ 的外表面。
\end{exercise}
\begin{solution}
取锥面参数化
\[
\vec r(r,\theta)=(r\cos\theta,r\sin\theta,r),\qquad
0\le r\le h,\quad 0\le\theta\le2\pi.
\]
直接算雅可比行列式：
\[
\begin{aligned}
\frac{\partial(y,z)}{\partial(r,\theta)}
&=\begin{vmatrix}\sin\theta&1\\ r\cos\theta&0\end{vmatrix}
=\sin\theta\cdot0-1\cdot r\cos\theta
=-r\cos\theta,\\
\frac{\partial(z,x)}{\partial(r,\theta)}
&=\begin{vmatrix}1&\cos\theta\\0&-r\sin\theta\end{vmatrix}
=1\cdot(-r\sin\theta)-\cos\theta\cdot0
=-r\sin\theta,\\
\frac{\partial(x,y)}{\partial(r,\theta)}
&=\begin{vmatrix}\cos\theta&\sin\theta\\-r\sin\theta&r\cos\theta\end{vmatrix}
=r(\cos^2\theta+\sin^2\theta)
=r.
\end{aligned}
\]
这组参数给出的法向为 $(-r\cos\theta,-r\sin\theta,r)$，指向锥面内侧；题目取外表面，故 $\epsilon=-1$。因此
\[
\mathrm{d}y\,\mathrm{d}z=r\cos\theta\,\mathrm{d}r\,\mathrm{d}\theta,\qquad
\mathrm{d}z\,\mathrm{d}x=r\sin\theta\,\mathrm{d}r\,\mathrm{d}\theta,\qquad
\mathrm{d}x\,\mathrm{d}y=-r\,\mathrm{d}r\,\mathrm{d}\theta.
\]
代入 $x=r\cos\theta,\ y=r\sin\theta,\ z=r$，得
\[
\begin{aligned}
I&=\int_0^{2\pi}\int_0^h r^2\bigl[(\sin\theta-1)\cos\theta
+(1-\cos\theta)\sin\theta-(\cos\theta-\sin\theta)\bigr]\,\mathrm{d}r\,\mathrm{d}\theta\\
&=\int_0^{2\pi}\int_0^h 2r^2(\sin\theta-\cos\theta)\,\mathrm{d}r\,\mathrm{d}\theta=0.
\end{aligned}
\]
所以 \(\displaystyle \boxed{I=0}\)。
\end{solution}

\begin{exercise}[球面的径向通量]
计算曲面积分 \(\displaystyle I=\iint_S \frac{x\,\mathrm{d}y\,\mathrm{d}z+y\,\mathrm{d}z\,\mathrm{d}x+z\,\mathrm{d}x\,\mathrm{d}y}{(x^2+y^2+z^2)^{3/2}}\)，其中 $S$ 是球面 $x^2+y^2+z^2=a^2$ 的外侧表面。
\end{exercise}
\begin{solution}
取球面参数化
\[
\vec r(\varphi,\theta)=(a\sin\varphi\cos\theta,a\sin\varphi\sin\theta,a\cos\varphi),
\quad 0\le\varphi\le\pi,\quad 0\le\theta\le2\pi.
\]
三个雅可比行列式为
\[
\begin{aligned}
\frac{\partial(y,z)}{\partial(\varphi,\theta)}
&=\begin{vmatrix}
a\cos\varphi\sin\theta&-a\sin\varphi\\
a\sin\varphi\cos\theta&0
\end{vmatrix}
=a\cos\varphi\sin\theta\cdot0-(-a\sin\varphi)a\sin\varphi\cos\theta\\
&=a^2\sin^2\varphi\cos\theta,\\
\frac{\partial(z,x)}{\partial(\varphi,\theta)}
&=\begin{vmatrix}
-a\sin\varphi&a\cos\varphi\cos\theta\\
0&-a\sin\varphi\sin\theta
\end{vmatrix}
=(-a\sin\varphi)(-a\sin\varphi\sin\theta)-a\cos\varphi\cos\theta\cdot0\\
&=a^2\sin^2\varphi\sin\theta,\\
\frac{\partial(x,y)}{\partial(\varphi,\theta)}
&=\begin{vmatrix}
a\cos\varphi\cos\theta&a\cos\varphi\sin\theta\\
-a\sin\varphi\sin\theta&a\sin\varphi\cos\theta
\end{vmatrix}\\
&=a^2\sin\varphi\cos\varphi(\cos^2\theta+\sin^2\theta)
=a^2\sin\varphi\cos\varphi.
\end{aligned}
\]
这组参数对应外侧，故 $\epsilon=1$。又在球面上 $x^2+y^2+z^2=a^2$，所以
\[
\begin{aligned}
I&=\int_0^{2\pi}\int_0^\pi
\frac{x\dfrac{\partial(y,z)}{\partial(\varphi,\theta)}
+y\dfrac{\partial(z,x)}{\partial(\varphi,\theta)}
+z\dfrac{\partial(x,y)}{\partial(\varphi,\theta)}}{a^3}
\,\mathrm{d}\varphi\,\mathrm{d}\theta\\
&=\int_0^{2\pi}\int_0^\pi
\frac{a^3\sin^3\varphi\cos^2\theta+a^3\sin^3\varphi\sin^2\theta
+a^3\sin\varphi\cos^2\varphi}{a^3}
\,\mathrm{d}\varphi\,\mathrm{d}\theta\\
&=\int_0^{2\pi}\int_0^\pi \sin\varphi\,\mathrm{d}\varphi\,\mathrm{d}\theta=4\pi.
\end{aligned}
\]
故 \(\displaystyle \boxed{I=4\pi}\)。
\end{solution}


\subsection{第五节习题：格林公式及其应用}

\subsubsection{A组}

\begin{exercise}[基本概念]
回答下列问题：
\begin{shorttrienum}
    \shortitem{1}{什么叫单连通区域？它的边界曲线的正向、负向如何规定？} &
    \shortitem{2}{什么叫复连通区域？其边界曲线的正向又是怎样规定的？} &
    \shortitem{3}{格林公式成立的条件是什么？}
\end{shorttrienum}
\end{exercise}
\begin{solution}
\begin{compactenum}
    \item 若平面区域 $D$ 内任意一条简单闭曲线所围成的内部都完全包含于 $D$，则称 $D$ 为单连通区域。对单连通区域的边界曲线，通常规定逆时针方向为正向，顺时针方向为负向；等价地说，沿正向行进时区域总在左手边。
    \item 若区域内部存在“洞”，即不是单连通区域，则称为复连通区域。对复连通区域，外边界取正向（逆时针），内边界取负向（顺时针），这样沿各边界的规定方向行进时，区域始终位于左侧。
    \item 格林公式成立的基本条件是：区域 $D$ 为由分段光滑简单闭曲线围成的平面区域。

    函数 $P(x,y),Q(x,y)$ 在包含 $D$ 的开区域内具有连续一阶偏导数。此时有 \(\displaystyle \oint_{\partial D}P\,\mathrm{d}x+Q\,\mathrm{d}y=\iint_D\left(\frac{\partial Q}{\partial x}-\frac{\partial P}{\partial y}\right)\mathrm{d}x\,\mathrm{d}y\)。
\end{compactenum}
\end{solution}

\begin{exercise}[利用格林公式计算曲线积分]
试利用格林公式计算下列曲线积分：
\begin{shortsingleenum}
    \shortsingleitem{1}{$\displaystyle \oint_C \vec{F}\cdot \mathrm{d}\vec{s}$，其中 $\vec{F}=2xy\,\vec{i}+(x^2+8y^3)\vec{j}$，$C$ 为圆周 $x^2+y^2=a^2$，取正向；}
    \shortsingleitem{2}{$\displaystyle \oint_L \vec{F}\cdot \mathrm{d}\vec{s}$，其中 $\vec{F}=(x-y)\vec{i}+(y-x)\vec{j}$，$L$ 是椭圆 $\dfrac{x^2}{a^2}+\dfrac{y^2}{b^2}=1$，取负向；}
    \shortsingleitem{3}{$\displaystyle \oint_L (yx^3+\mathrm{e}^y)\,\mathrm{d}x+(xy^3+x\mathrm{e}^y-2y)\,\mathrm{d}y$，其中 $L$ 是椭圆 $\dfrac{x^2}{a^2}+\dfrac{y^2}{b^2}=1$，取正向；}
    \shortsingleitem{4}{$\displaystyle \oint_C x^2y\,\mathrm{d}x+xy^2\,\mathrm{d}y$，其中 $L:\ |x|+|y|=1$，取正向。}
\end{shortsingleenum}
\end{exercise}
\begin{solution}
\begin{compactenum}
    \item 记 $P=2xy,\ Q=x^2+8y^3$，则 $\dfrac{\partial Q}{\partial x}=2x,\ \dfrac{\partial P}{\partial y}=2x$，从而 $\dfrac{\partial Q}{\partial x}-\dfrac{\partial P}{\partial y}=0$。由格林公式，\(\displaystyle \oint_C \vec{F}\cdot \mathrm{d}\vec{s}=0\)。
    \item 这里 $P=x-y,\ Q=y-x$，故 $\dfrac{\partial Q}{\partial x}=-1,\ \dfrac{\partial P}{\partial y}=-1$，从而旋度仍为 $0$。积分沿负向进行，但由于对应的面积分为 $0$，故 \(\displaystyle \oint_L \vec{F}\cdot \mathrm{d}\vec{s}=0\)。
    \item 记 $P=yx^3+\mathrm{e}^y,\ Q=xy^3+x\mathrm{e}^y-2y$。则 $\dfrac{\partial Q}{\partial x}=y^3+\mathrm{e}^y,\ \dfrac{\partial P}{\partial y}=x^3+\mathrm{e}^y$，所以 $\dfrac{\partial Q}{\partial x}-\dfrac{\partial P}{\partial y}=y^3-x^3$。
    椭圆区域关于原点中心对称，而函数 $y^3-x^3$ 在该区域上是奇函数，因此其二重积分为零。由格林公式，\(\displaystyle \oint_L (yx^3+\mathrm{e}^y)\,\mathrm{d}x+(xy^3+x\mathrm{e}^y-2y)\,\mathrm{d}y=0\)。
    \item 记 $P=x^2y,\ Q=xy^2$，则 $\dfrac{\partial Q}{\partial x}=y^2,\ \dfrac{\partial P}{\partial y}=x^2$。因而 \(\displaystyle \oint_C x^2y\,\mathrm{d}x+xy^2\,\mathrm{d}y=\iint_D (y^2-x^2)\,\mathrm{d}x\,\mathrm{d}y\)，其中 \(\displaystyle D=\{(x,y)\mid |x|+|y|\le 1\}\)。
    该区域关于直线 $y=x$ 对称，而被积函数 $y^2-x^2$ 关于交换 $x,y$ 变号，所以积分为 $0$。故 \(\displaystyle \oint_C x^2y\,\mathrm{d}x+xy^2\,\mathrm{d}y=0\)。
\end{compactenum}
\end{solution}

\begin{exercise}[利用曲线积分求面积]
利用曲线积分，求下列曲线所围成的图形的面积：
\begin{shortexenum}
    \shortitem{1}{星形线 $x=a\cos^3 t,\ y=a\sin^3 t\ (0\le t\le 2\pi)$；} &
    \shortitem{2}{双纽线 $r^2=a^2\cos 2\varphi$。}
\end{shortexenum}
\end{exercise}
\begin{solution}
\begin{compactenum}
    \item 由面积公式 \(\displaystyle A=\frac12\oint_C (x\,\mathrm{d}y-y\,\mathrm{d}x)\)。
    这里 $x=a\cos^3 t,\ y=a\sin^3 t$，故 $\mathrm{d}x=-3a\cos^2 t\sin t\,\mathrm{d}t$，$\mathrm{d}y=3a\sin^2 t\cos t\,\mathrm{d}t$。
    于是
    \begin{align*}
    A
    &=\frac12\int_0^{2\pi}\left(a\cos^3 t\cdot 3a\sin^2 t\cos t-a\sin^3 t\cdot(-3a\cos^2 t\sin t)\right)\mathrm{d}t \\
    &=\frac{3a^2}{2}\int_0^{2\pi}\cos^2 t\sin^2 t\,\mathrm{d}t
    =\frac{3a^2}{2}\cdot \frac{\pi}{4}
    =\frac{3\pi a^2}{8}.
    \end{align*}
    \item 双纽线关于原点对称，可先求一瓣面积再乘以 $2$。由极坐标面积公式 \(\displaystyle A_{\text{一瓣}}=\frac12\int_{-\pi/4}^{\pi/4} r^2\,\mathrm{d}\varphi=\frac{a^2}{2}\int_{-\pi/4}^{\pi/4}\cos 2\varphi\,\mathrm{d}\varphi=\frac{a^2}{2}\)。
    所以总面积为 $A=2A_{\text{一瓣}}=a^2$，即 $\boxed{A=a^2}$。
\end{compactenum}
\end{solution}

\begin{exercise}[开曲线上的曲线积分]
计算曲线积分 \(\displaystyle \int_L (x+y)^2\,\mathrm{d}x-(x^2+y^2\sin y)\,\mathrm{d}y\)，其中 $L$ 是抛物线 $y=x^2$ 上从点 $(-1,1)$ 到点 $(1,1)$ 的那一段。
\end{exercise}
\begin{solution}
这条路径不是闭曲线，题目也没有给出保守场或势函数信号，所以 Green 公式和端点差都不该先上。最稳的做法就是直接按曲线参数化。
将曲线参数化为 $x=t,\ y=t^2,\ -1\le t\le 1$，则 $\mathrm{d}x=\mathrm{d}t,\ \mathrm{d}y=2t\,\mathrm{d}t$。
故积分为
\begin{align*}
I
&=\int_{-1}^{1}\left[(t+t^2)^2-2t\bigl(t^2+t^4\sin(t^2)\bigr)\right]\mathrm{d}t \\
&=\int_{-1}^{1}\left(t^2+2t^3+t^4-2t^3-2t^5\sin(t^2)\right)\mathrm{d}t \\
&=\int_{-1}^{1}(t^2+t^4)\,\mathrm{d}t
=2\int_0^1 (t^2+t^4)\,\mathrm{d}t
=\frac{16}{15}.
\end{align*}
因此 \(\displaystyle \boxed{I=\frac{16}{15}}\)。
\end{solution}

\begin{exercise}[半圆弧与线段构成的路径积分]
求 \(\displaystyle I=\int_{\widehat{ABO}} (x^2-\mathrm{e}^x\cos y)\,\mathrm{d}x+(\mathrm{e}^x\sin y+3x)\,\mathrm{d}y\)，其中 $\widehat{ABO}$ 是从点 $A(0,2)$ 沿右半圆周 $x=\sqrt{1-(y-1)^2}$ 到点 $O(0,0)$ 的弧段。
\end{exercise}
\begin{solution}
\textbf{思路：}原路径只是半圆弧，不能直接套 Green；但它的两个端点刚好可以用一条很简单的线段补成闭曲线，所以先补成右半圆盘的正向边界，再把补边贡献减回去最省力。
记 $P=x^2-\mathrm{e}^x\cos y,\ Q=\mathrm{e}^x\sin y+3x$，则 $\dfrac{\partial P}{\partial y}=\mathrm{e}^x\sin y,\ \dfrac{\partial Q}{\partial x}=\mathrm{e}^x\sin y+3$，故 $\dfrac{\partial Q}{\partial x}-\dfrac{\partial P}{\partial y}=3$。

为使用 Green 公式，取右半圆盘 $D$ 的正向边界 $C=\overrightarrow{AO}\cup\widehat{OA}$，其中 $\overrightarrow{AO}$ 表示线段 $AO$ 由上向下，$\widehat{OA}$ 表示右半圆弧由 $O$ 到 $A$。则 \(\displaystyle \oint_C P\,\mathrm{d}x+Q\,\mathrm{d}y=\iint_D 3\,\mathrm{d}x\,\mathrm{d}y=3\cdot \frac{\pi}{2}\)，
因为 $D$ 是半径为 $1$ 的半圆盘。

又在线段 $AO$ 上，$x=0$，从而 $P\,\mathrm{d}x+Q\,\mathrm{d}y=0$，因此 \(\displaystyle \oint_C P\,\mathrm{d}x+Q\,\mathrm{d}y=\int_{\widehat{OA}}P\,\mathrm{d}x+Q\,\mathrm{d}y\)。
而题目所求是反向弧段 $\widehat{ABO}$（即从 $A$ 到 $O$），故 \(\displaystyle I=-\int_{\widehat{OA}}P\,\mathrm{d}x+Q\,\mathrm{d}y
=-\dfrac{3\pi}{2}\)，
即 $\boxed{I=-\dfrac{3\pi}{2}}$。
\end{solution}

\subsubsection{B组}

\begin{exercise}[含奇点的曲线积分]
计算曲线积分 \(\displaystyle I=\oint_C \frac{-y\,\mathrm{d}x+x\,\mathrm{d}y}{4x^2+y^2}\)，其中 $C$ 是以点 $(1,0)$ 为圆心、以 $R$ 为半径的圆周（$R\neq 1$），取逆时针方向。
\end{exercise}
\begin{solution}
记 \(\displaystyle P=-\frac{y}{4x^2+y^2},\qquad Q=\frac{x}{4x^2+y^2}\)。当 $(x,y)\neq(0,0)$ 时，\(\displaystyle \frac{\partial Q}{\partial x}=\frac{y^2-4x^2}{(4x^2+y^2)^2}=\frac{\partial P}{\partial y}\)，故该微分形式只在原点有奇点。

若 $0<R<1$，则原点在圆 $C$ 外部，被积函数在 $C$ 所围圆盘内光滑，由 Green 公式得 $I=0$。

若 $R>1$，则原点在圆 $C$ 内部。取充分小的椭圆 \(\displaystyle C_\varepsilon:\ 4x^2+y^2=\varepsilon^2\)，并取逆时针方向。由挖洞后的 Green 公式，在环域上有 \(\displaystyle \oint_C=\oint_{C_\varepsilon}\)。参数化 \(\displaystyle x=\frac{\varepsilon}{2}\cos t,\qquad y=\varepsilon\sin t,\qquad 0\le t\le 2\pi\)，则 \(\displaystyle \oint_{C_\varepsilon}\frac{-y\,\mathrm{d}x+x\,\mathrm{d}y}{4x^2+y^2}=\int_0^{2\pi}\frac{\varepsilon^2/2}{\varepsilon^2}\,\mathrm{d}t=\int_0^{2\pi}\frac12\,\mathrm{d}t=\pi\)。
因此
\[
\boxed{
I=\begin{cases}
0,&0<R<1,\\
\pi,&R>1.
\end{cases}}
\]
\end{solution}

\begin{exercise}[不经过原点的简单闭曲线]
设 $C$ 为不经过原点的简单封闭曲线，取逆时针方向。计算曲线积分 \(\displaystyle I=\oint_C \frac{x\,\mathrm{d}y-y\,\mathrm{d}x}{x^2+y^2}\)。
\end{exercise}
\begin{solution}
设 $D$ 为 $C$ 所围成的区域。若 $0\notin D$，则被积函数在 $D$ 上光滑，且 \(\displaystyle \frac{\partial}{\partial x}\!\left(\frac{x}{x^2+y^2}\right)-\frac{\partial}{\partial y}\!\left(-\frac{y}{x^2+y^2}\right)=0\)，由 Green 公式得 $I=0$。

若 $0\in D$，则在原点处有奇点。取充分小的圆 $C_\varepsilon:\ x^2+y^2=\varepsilon^2$，并取逆时针方向。对挖去小圆后的环域应用 Green 公式，得 \(\displaystyle \oint_C \frac{x\,\mathrm{d}y-y\,\mathrm{d}x}{x^2+y^2}=\oint_{C_\varepsilon} \frac{x\,\mathrm{d}y-y\,\mathrm{d}x}{x^2+y^2}\)。对 $C_\varepsilon$ 参数化 \(\displaystyle x=\varepsilon\cos t,\qquad y=\varepsilon\sin t,\qquad 0\le t\le 2\pi\)，则 \(\displaystyle \oint_{C_\varepsilon} \frac{x\,\mathrm{d}y-y\,\mathrm{d}x}{x^2+y^2}=\int_0^{2\pi}\frac{\varepsilon^2}{\varepsilon^2}\,\mathrm{d}t=2\pi\)。
因此
\[
\boxed{
I=\begin{cases}
0,&0\notin D,\\
2\pi,&0\in D.
\end{cases}}
\]
\end{solution}

\begin{exercise}[用 Green 公式表示 Laplace 积分]
试利用 Green 公式证明
\(\displaystyle \iint_D \Delta f\,\mathrm{d}\sigma
=\oint_C\left(\frac{\partial f}{\partial x}\,\mathrm{d}y-\frac{\partial f}{\partial y}\,\mathrm{d}x\right)\)，
其中 $\Delta f=\nabla\cdot(\nabla f)$，$C$ 是平面闭区域 $D$ 的正向边界曲线，$f$ 在 $D$ 上具有二阶连续偏导数。
\end{exercise}
\begin{solution}
在 Green 公式中取 \(\displaystyle P=-\frac{\partial f}{\partial y},\qquad Q=\frac{\partial f}{\partial x}\)。则 \(\displaystyle \frac{\partial Q}{\partial x}-\frac{\partial P}{\partial y}=\frac{\partial^2 f}{\partial x^2}+\frac{\partial^2 f}{\partial y^2}=\nabla\cdot(\nabla f)\)。代入 Green 公式并整理，得 \(\displaystyle \oint_C\left(-\frac{\partial f}{\partial y}\,\mathrm{d}x+\frac{\partial f}{\partial x}\,\mathrm{d}y\right)=\iint_D \nabla\cdot(\nabla f)\,\mathrm{d}\sigma\)，也就是 \(\displaystyle \iint_D \Delta f\,\mathrm{d}\sigma=\oint_C\left(\frac{\partial f}{\partial x}\,\mathrm{d}y-\frac{\partial f}{\partial y}\,\mathrm{d}x\right)\)。
\end{solution}


\subsection{第六节习题：保守场与势函数}

\subsubsection{A组}

\begin{exercise}[基本概念]
回答下列问题：
\begin{shortexenum}
    \shortitem{1}{什么叫保守场？什么叫势函数？} &
    \shortitem{2}{保守场的重要性质是什么？} \\
    \shortitem{3}{如何判别一个向量场是否是保守场？} &
    \shortitem{4}{什么叫全微分方程？} \\
    \shortitem{5}{怎样求势函数？} &
\end{shortexenum}
\end{exercise}
\begin{solution}
\textbf{思路：}这组题主要是在回顾保守场、势函数和全微分方程的定义与常用判据，答题时按“定义—性质—判别—构造”的顺序最稳。
\begin{compactenum}
    \item 若向量场 $\vec{F}$ 沿任意两点间的曲线积分只与起点、终点有关，而与路径无关，则称 $\vec{F}$ 为保守场。若存在数量函数 $u$ 使得 $\vec{F}=\nabla u$，则称 $u$ 为该保守场的势函数。
    \item 保守场最重要的性质有：曲线积分与路径无关；任意闭曲线上的积分为零；若 $\vec{F}=\nabla u$，则 \(\displaystyle \int_A^B \vec{F}\cdot \mathrm{d}\vec{r}=u(B)-u(A)\)。
    \item 在单连通区域内，若平面向量场 $\vec{F}=(P,Q)$ 满足 $\dfrac{\partial P}{\partial y}=\dfrac{\partial Q}{\partial x}$，或空间向量场 $\vec{F}$ 满足 $\operatorname{rot}\vec{F}=\vec{0}$，则它是保守场。单连通条件不可忽略。
    \item 形如 $P(x,y)\,\mathrm{d}x+Q(x,y)\,\mathrm{d}y=0$，且左端恰是某个函数 $u(x,y)$ 的全微分 \(\displaystyle \mathrm{d}u=u_x\,\mathrm{d}x+u_y\,\mathrm{d}y\) 的方程，称为全微分方程。
    \item 求势函数通常有两种做法：先对一个分量积分，再用另一个分量校正“积分常数”；或者直接利用 $\displaystyle u(B)-u(A)=\int_A^B \vec{F}\cdot \mathrm{d}\vec{r}$ 构造势函数。
\end{compactenum}
\end{solution}

\begin{exercise}[向量场是否保守]
判别下列向量场是不是保守场，或对什么区域来说是保守场：
\begin{shortexenum}
    \shortitem{1}{$\vec{F}=\varphi(x)\vec{i}+\psi(y)\vec{j}$，其中 $\varphi,\psi$ 是可微函数；} &
    \shortitem{2}{$\vec{F}=\dfrac{1}{r^3}(x\vec{i}+y\vec{j})$，$r=\sqrt{x^2+y^2}$；} \\
    \shortitem{3}{$\vec{F}=f(x+y)(\vec{i}+\vec{j})$，其中 $f(u)$ 连续；} &
    \shortitem{4}{$\vec{F}=\dfrac{1}{x^2+y^2}(y\vec{i}-x\vec{j})$。}
\end{shortexenum}
\end{exercise}
\begin{solution}
\begin{compactenum}
    \item 这里 $P=\varphi(x),\ Q=\psi(y)$，故 $\dfrac{\partial P}{\partial y}=0,\ \dfrac{\partial Q}{\partial x}=0$，所以在整个平面上它都是保守场。对应势函数可取 $u(x,y)=\int \varphi(x)\,\mathrm{d}x+\int \psi(y)\,\mathrm{d}y$。
    \item 由 $\vec{F}=\left(\dfrac{x}{r^3},\dfrac{y}{r^3}\right)=\nabla\left(-\dfrac{1}{r}\right)$，可知它在整个去心平面 $\mathbb{R}^2\setminus\{(0,0)\}$ 上就是保守场，其势函数可取 $u(x,y)=-\dfrac{1}{r}$。
    \item 这里 $P=f(x+y),\ Q=f(x+y)$。不妨令 $F(u)=\int_0^u f(s)\,\mathrm{d}s$，则 $F'(u)=f(u)$。于是 $\nabla F(x+y)=F'(x+y)\nabla(x+y)=f(x+y)(1,1)$，因而该场在整个平面上是保守场，势函数可取 $u(x,y)=F(x+y)=\int_0^{x+y} f(s)\,\mathrm{d}s$。
    \item 这里 $P=\dfrac{y}{x^2+y^2},\ Q=-\dfrac{x}{x^2+y^2}$。在原点外有 \(\displaystyle \frac{\partial P}{\partial y}=\frac{x^2-y^2}{(x^2+y^2)^2}=\frac{\partial Q}{\partial x}\)，但该场在整个去心平面上不是保守场，因为沿包围原点的闭曲线积分不为零。它只在不包围原点的单连通区域内是保守场。
\end{compactenum}
\end{solution}

\begin{exercise}[求势函数]
求势函数 $u(x,y)$，使得
\begin{shortexenum}
    \shortitem{1}{$\mathrm{d}u=2xy\,\mathrm{d}x+x^2\,\mathrm{d}y$；} &
    \shortitem{2}{$\mathrm{d}u=(3x^2-3y^2+4)\,\mathrm{d}x-6xy\,\mathrm{d}y$；} \\
    \shortitem{3}{$\mathrm{d}u=(2x\cos y+y^2\cos x)\,\mathrm{d}x+(2y\sin x-x^2\sin y)\,\mathrm{d}y$；} &
    \shortitem{4}{$\mathrm{d}u=(\mathrm{e}^x\sin y+2xy^2)\,\mathrm{d}x+(\mathrm{e}^x\cos y+2x^2y)\,\mathrm{d}y$。}
\end{shortexenum}
\end{exercise}
\begin{solution}
\begin{compactenum}
    \item 由 $u_x=2xy$ 积分得 $u=x^2y+\varphi(y)$。再由 $u_y=x^2$ 得 $\varphi'(y)=0$，故 $\boxed{u=x^2y+C}$。
    \item 由 $u_x=3x^2-3y^2+4$ 积分得 $u=x^3-3xy^2+4x+\varphi(y)$。再由 $u_y=-6xy$ 得 $\varphi'(y)=0$，故 $\boxed{u=x^3-3xy^2+4x+C}$。
    \item 由 $u_x=2x\cos y+y^2\cos x$ 积分得 $u=x^2\cos y+y^2\sin x+\varphi(y)$。再求 $u_y=-x^2\sin y+2y\sin x+\varphi'(y)$，与已知比较可得 $\varphi'(y)=0$，故 $\boxed{u=x^2\cos y+y^2\sin x+C}$。
    \item 由 $u_x=\mathrm{e}^x\sin y+2xy^2$ 积分得 $u=\mathrm{e}^x\sin y+x^2y^2+\varphi(y)$。再求 $u_y=\mathrm{e}^x\cos y+2x^2y+\varphi'(y)$，与已知比较可得 $\varphi'(y)=0$，故 $\boxed{u=\mathrm{e}^x\sin y+x^2y^2+C}$。
\end{compactenum}
\end{solution}

\begin{exercise}[验证全微分并计算积分]
验证被积函数为全微分，并计算下列曲线积分：
\begin{shortexenum}
    \shortitem{1}{$\displaystyle \int_{(-1,2)}^{(2,3)} x\,\mathrm{d}y+y\,\mathrm{d}x$；} &
    \shortitem{2}{$\displaystyle \int_{(0,0)}^{(1,1)} (x-y)(\mathrm{d}x-\mathrm{d}y)$；} \\
    \shortitem{3}{$\displaystyle \int_{(2,1)}^{(1,2)} \frac{y\,\mathrm{d}x-x\,\mathrm{d}y}{x^2}$，沿不和 $y$ 轴相交的路径；} &
    \shortitem{4}{$\displaystyle \int_{(1,0)}^{(6,8)} \frac{x\,\mathrm{d}x+y\,\mathrm{d}y}{\sqrt{x^2+y^2}}$，沿不通过坐标原点的路径。}
\end{shortexenum}
\end{exercise}
\begin{solution}
\begin{compactenum}
    \item 因为 $x\,\mathrm{d}y+y\,\mathrm{d}x=\mathrm{d}(xy)$，所以 $\displaystyle \int_{(-1,2)}^{(2,3)} x\,\mathrm{d}y+y\,\mathrm{d}x=(xy)\Big|_{(-1,2)}^{(2,3)}=6-(-2)=8$。
    \item 注意到 $(x-y)(\mathrm{d}x-\mathrm{d}y)=\dfrac12\,\mathrm{d}(x-y)^2$，因而 $\displaystyle \int_{(0,0)}^{(1,1)} (x-y)(\mathrm{d}x-\mathrm{d}y)=\frac12\Big[(x-y)^2\Big]_{(0,0)}^{(1,1)}=0$。
    \item 被积表达式化为 $\dfrac{y\,\mathrm{d}x-x\,\mathrm{d}y}{x^2}=\mathrm{d}\left(-\dfrac{y}{x}\right)$，因此 $\displaystyle \int_{(2,1)}^{(1,2)} \frac{y\,\mathrm{d}x-x\,\mathrm{d}y}{x^2}=-\frac{y}{x}\Bigg|_{(2,1)}^{(1,2)}=-2+\frac12=-\frac32$。
    \item 注意 $\mathrm{d}\sqrt{x^2+y^2}=\dfrac{x\,\mathrm{d}x+y\,\mathrm{d}y}{\sqrt{x^2+y^2}}$，所以 $\displaystyle \int_{(1,0)}^{(6,8)} \frac{x\,\mathrm{d}x+y\,\mathrm{d}y}{\sqrt{x^2+y^2}}=\sqrt{x^2+y^2}\Big|_{(1,0)}^{(6,8)}=10-1=9$。
\end{compactenum}
\end{solution}

\begin{exercise}[保守场作功]
设有平面力场 \(\displaystyle \vec{F}(x,y)=\bigl(2xy^3-y^2\cos x\bigr)\vec{i}+\bigl(1-2y\sin x+3x^2y^2\bigr)\vec{j}\)，求一质点沿曲线 $L:\ 2x=\pi y^2$ 从点 $O(0,0)$ 运动到点 $A\left(\dfrac{\pi}{2},1\right)$ 时，变力 $\vec{F}$ 所作的功。
\end{exercise}
\begin{solution}
\textbf{思路：}题目虽然给了具体轨道，但最稳的入口不是参数化曲线，而是先判保守性；一旦发现它是保守场，做功就直接退化成势函数的端点差。

记 \( \displaystyle P=2xy^3-y^2\cos x,\ Q=1-2y\sin x+3x^2y^2 \)，则 \(\displaystyle \frac{\partial P}{\partial y}=6xy^2-2y\cos x,\quad \frac{\partial Q}{\partial x}=6xy^2-2y\cos x\)。
由于整个平面是单连通的，所以该场是保守场。

设势函数为 $u$。由 \(\displaystyle u_x=2xy^3-y^2\cos x\)，得
\(\displaystyle u=x^2y^3-y^2\sin x+\varphi(y)\)；再由 \(\displaystyle u_y=Q\)，得
\(\displaystyle 3x^2y^2-2y\sin x+\varphi'(y)=1-2y\sin x+3x^2y^2\Longrightarrow \varphi'(y)=1\)。
得 \(\displaystyle \varphi(y)=y+C\)，因此可取 \(\displaystyle u(x,y)=x^2y^3-y^2\sin x+y\)。所作的功为 \(\displaystyle W=u\left(\frac{\pi}{2},1\right)-u(0,0)=\frac{\pi^2}{4}-1+1=\frac{\pi^2}{4}\)。
故 \(\displaystyle \boxed{W=\frac{\pi^2}{4}}\)。
\end{solution}

\subsubsection{B组}

\begin{exercise}[上半椭圆弧上的积分]
设 $L$ 是从点 $A(-a,0)$ 经上半椭圆 \(\displaystyle \frac{x^2}{a^2}+\frac{y^2}{b^2}=1\ (y\ge 0)\) 到点 $B(a,0)$ 的弧段，计算曲线积分 \(\displaystyle I=\int_L \frac{x-y}{x^2+y^2}\,\mathrm{d}x+\frac{x+y}{x^2+y^2}\,\mathrm{d}y\)。
\end{exercise}
\begin{solution}
\textbf{思路：}这题若直接参数化椭圆会比较重。更好的入口是先把被积式拆成“全微分 $+$ 极角微分”，这样路径信息只剩下端点模长和极角变化。

将被积式拆开：
\[
\frac{x-y}{x^2+y^2}\,\mathrm{d}x+\frac{x+y}{x^2+y^2}\,\mathrm{d}y
=\frac{x\,\mathrm{d}x+y\,\mathrm{d}y}{x^2+y^2}
+\frac{x\,\mathrm{d}y-y\,\mathrm{d}x}{x^2+y^2}.
\]
第一项是 $\dfrac12\,\mathrm{d}\ln(x^2+y^2)$，第二项是极角微分 $\mathrm{d}\theta$。沿上半椭圆从 $A$ 到 $B$，极角从 $\pi$ 变化到 $0$，故 \(\displaystyle \int_L \frac{x\,\mathrm{d}y-y\,\mathrm{d}x}{x^2+y^2}=-\pi\)。
又端点满足 $x^2+y^2=a^2$，所以第一项积分为零。于是 $\boxed{I=-\pi}$。
\end{solution}

\begin{exercise}[含参数函数的全微分积分]
设 $f(x)$ 在 $(-\infty,+\infty)$ 内有连续导数，$L$ 是从点 $A\left(3,\dfrac{2}{3}\right)$ 到点 $B(1,2)$ 的直线段，计算曲线积分
\[
I=\int_L \frac{1+y^2f(xy)}{y}\,\mathrm{d}x
+\frac{x}{y^2}\bigl[y^2f(xy)-1\bigr]\,\mathrm{d}y.
\]
\end{exercise}
\begin{solution}
\textbf{思路：}含参函数看起来复杂，但它只依赖组合量 $xy$。因此先试着把被积式拼成 $F(xy)+x/y$ 的微分，比沿直线段参数化稳得多。

记
\[
P=\frac{1+y^2f(xy)}{y}=\frac{1}{y}+yf(xy),\qquad
Q=\frac{x}{y^2}\bigl[y^2f(xy)-1\bigr]=xf(xy)-\frac{x}{y^2}.
\]
设 $F(t)=\int_0^t f(s)\,\mathrm{d}s$，则 $F'(t)=f(t)$。取函数 $u(x,y)=F(xy)+\dfrac{x}{y}$，便有 \(\displaystyle u_x=yf(xy)+\frac1y=P,\qquad u_y=xf(xy)-\frac{x}{y^2}=Q\)。故该积分为全微分积分，且 $I=u(B)-u(A)$。
代入端点 $A\left(3,\frac23\right)$，$B(1,2)$，均有 $xy=2$，于是 \(\displaystyle u(B)=F(2)+\frac12,\quad u(A)=F(2)+\frac{9}{2}\)。
所以 $\boxed{I=-4}$。
\end{solution}

\begin{exercise}[保守场上的定值积分]
设 $\vec{F}(x,y)=xy^2\,\vec{i}+yf(x)\,\vec{j}$ 为保守场，其中 $f(x)$ 具有连续导数，且 $f(0)=0$。计算曲线积分 \(\displaystyle I=\int_{(0,0)}^{(1,1)} \vec{F}\cdot \mathrm{d}\vec{s}\) 的值。
\end{exercise}
\begin{solution}
\textbf{思路：}题目已经告诉我们这是保守场，所以先别想路径；先用保守性反推出 $f(x)$，再把积分降成势函数的端点差。

因为 $\vec{F}$ 是保守场，所以必有 \(\displaystyle \frac{\partial}{\partial y}(xy^2)=\frac{\partial}{\partial x}\bigl(yf(x)\bigr)\)，
即 $2xy=yf'(x)$。这对一切 $y$ 成立，因此 $f'(x)=2x$。再由 $f(0)=0$ 得 $f(x)=x^2$，从而 $\vec{F}=(xy^2,\ x^2y)$。

这是势函数 $u(x,y)=\dfrac{x^2y^2}{2}$ 的梯度场。因此 $I=u(1,1)-u(0,0)=\dfrac12$，即 $\boxed{I=\dfrac12}$。
\end{solution}

\begin{exercise}[闭曲线上的径向积分]
证明：若 $f(u)$ 为连续函数，$C$ 为分段光滑的简单闭曲线，则 \(\displaystyle \oint_C f(x^2+y^2)(x\,\mathrm{d}x+y\,\mathrm{d}y)=0\)。
\end{exercise}
\begin{solution}
\textbf{思路：}看到 $x\,\mathrm{d}x+y\,\mathrm{d}y$，先联想到 $\mathrm{d}(x^2+y^2)$。这类闭路上的径向积分通常本质上就是某个势函数的全微分。

注意到 \(x\,\mathrm{d}x+y\,\mathrm{d}y=\dfrac12\,\mathrm{d}(x^2+y^2)\)。令 \(U(t)=\dfrac12\int_0^t f(s)\,\mathrm{d}s\)，则
\(\mathrm{d}U(x^2+y^2)=f(x^2+y^2)(x\,\mathrm{d}x+y\,\mathrm{d}y)\)。
因此被积表达式是一个全微分，沿任意闭曲线积分恒为零，所以原积分为 \(\boxed{0}\)。
\end{solution}


\subsection{第七节习题：散度和Gauss公式}

\subsubsection{A组}

\begin{exercise}[基本概念]
回答下列问题：
\begin{shorttrienum}
    \shortitem{1}{向量场的散度怎样定义？} &
    \shortitem{2}{在直角坐标系下，散度的计算公式是怎样的？} &
    \shortitem{3}{高斯公式成立有什么条件？}
\end{shorttrienum}
\end{exercise}
\begin{solution}
\begin{compactanswerenum}
    \item 向量场在一点的散度，表示该点附近单位体积的净流出强度。若把向量场看成流体速度场，则散度反映该点是“源”还是“汇”。
    \item 在直角坐标系下，若 $\vec{F}=(P,Q,R)$，则
    \[
    \operatorname{div}\vec{F}
    =\frac{\partial P}{\partial x}+\frac{\partial Q}{\partial y}+\frac{\partial R}{\partial z}.
    \]
    \item 高斯公式成立的条件通常是：空间区域 $\Omega$ 的边界 $S=\partial\Omega$ 为分片光滑闭曲面，向量场 $\vec{F}$ 在包含 $\Omega$ 的区域内具有连续一阶偏导数。此时
    \[
    \iint_{\partial\Omega}\vec{F}\cdot\vec{n}\,\mathrm{d}S
    =\iiint_{\Omega}\operatorname{div}\vec{F}\,\mathrm{d}V.
    \]
\end{compactanswerenum}
\end{solution}

\begin{exercise}[散度的基本性质]
证明：
\begin{shortexenum}
    \shortitem{1}{$\operatorname{div}(\vec{F}+\vec{G})=\operatorname{div}\vec{F}+\operatorname{div}\vec{G}$；} &
    \shortitem{2}{$\operatorname{div}(u\operatorname{grad}u)=u\Delta u+|\nabla u|^2$，其中 $u$ 为标量函数。}
\end{shortexenum}
\end{exercise}
\begin{solution}
\begin{compactenum}
    \item 设 $\vec{F}=(P_1,Q_1,R_1),\ \vec{G}=(P_2,Q_2,R_2)$，则
    \[
    \operatorname{div}(\vec{F}+\vec{G})
    =\frac{\partial(P_1+P_2)}{\partial x}
    +\frac{\partial(Q_1+Q_2)}{\partial y}
    +\frac{\partial(R_1+R_2)}{\partial z}.
    \]
    按线性展开可得
    \(\displaystyle \operatorname{div}(\vec{F}+\vec{G})=\operatorname{div}\vec{F}+\operatorname{div}\vec{G}\)。
    \item 由乘积法则，
    \[
    \operatorname{div}(u\nabla u)
    =\nabla u\cdot \nabla u+u\,\operatorname{div}(\nabla u)
    =|\nabla u|^2+u\Delta u,
    \]
    故
    \(\displaystyle \boxed{\operatorname{div}(u\operatorname{grad}u)=u\Delta u+|\nabla u|^2}\)。
\end{compactenum}
\end{solution}

\begin{exercise}[计算散度]
计算：
\begin{shortexenum}
    \shortitem{1}{$\operatorname{div}(u\operatorname{grad}u)$；} &
    \shortitem{2}{$\operatorname{div}\vec{r}$，其中 $\vec{r}=x\vec{i}+y\vec{j}+z\vec{k}$。}
\end{shortexenum}
\end{exercise}
\begin{solution}
\begin{compactenum}
    \item 由乘积散度公式，
    \[
    \operatorname{div}(u\nabla u)
    =u\,\operatorname{div}(\nabla u)+\nabla u\cdot\nabla u
    =u\Delta u+|\nabla u|^2,
    \]
    所以 \(\displaystyle \boxed{\operatorname{div}(u\operatorname{grad}u)=u\Delta u+|\nabla u|^2}\)。
    \item 直接计算 \(\displaystyle \operatorname{div}\vec{r}=\frac{\partial x}{\partial x}+\frac{\partial y}{\partial y}+\frac{\partial z}{\partial z}=3\)，故 \(\displaystyle \boxed{\operatorname{div}\vec{r}=3}\)。
\end{compactenum}
\end{solution}

\begin{exercise}[利用高斯公式计算曲面积分]
利用高斯公式计算下列曲面积分：
\begin{shortsingleenum}
    \shortsingleitem{1}{$\displaystyle \oiint_S x\,\mathrm{d}y\,\mathrm{d}z+y\,\mathrm{d}z\,\mathrm{d}x+z\,\mathrm{d}x\,\mathrm{d}y$，其中 $S$ 是介于平面 $z=0$ 与 $z=3$ 之间的圆柱体 $x^2+y^2\le 9$ 的整个表面的外侧；}
    \shortsingleitem{2}{$\displaystyle \oiint_S (x-y^2+z^2)\,\mathrm{d}y\,\mathrm{d}z+(y-z^2+x^2)\,\mathrm{d}z\,\mathrm{d}x+(z-x^2+y^2)\,\mathrm{d}x\,\mathrm{d}y$，其中 $S$ 是球面 $(x-a)^2+(y-b)^2+(z-c)^2=R^2$ 的外侧；}
    \shortsingleitem{3}{$\displaystyle \oiint_S yz\,\mathrm{d}y\,\mathrm{d}z+zx\,\mathrm{d}z\,\mathrm{d}x+xy\,\mathrm{d}x\,\mathrm{d}y$，其中 $S$ 是空间区域 $\Omega=\{(x,y,z)\mid |x|\le 1,\ |y|\le 1,\ |z|\le 1\}$ 的边界面的外侧；}
    \shortsingleitem{4}{$\displaystyle \oiint_S \frac{x\,\mathrm{d}y\,\mathrm{d}z+y\,\mathrm{d}z\,\mathrm{d}x+z\,\mathrm{d}x\,\mathrm{d}y}{\sqrt{x^2+y^2+z^2}}$，其中 $S$ 为球面 $x^2+y^2+z^2=a^2$ 的外侧。}
\end{shortsingleenum}
\end{exercise}
\begin{solution}
\begin{compactenum}
    \item 令 $\vec{F}=(x,y,z)$，则 \(\displaystyle \operatorname{div}\vec{F}=3\)，而圆柱体体积 \(\displaystyle V=\pi\cdot 3^2\cdot 3=27\pi\)。所以 \(\displaystyle \oiint_S x\,\mathrm{d}y\,\mathrm{d}z+y\,\mathrm{d}z\,\mathrm{d}x+z\,\mathrm{d}x\,\mathrm{d}y=\iiint_{\Omega} 3\,\mathrm{d}V=81\pi\)。
    \item 令 \(\displaystyle \vec{F}=(x-y^2+z^2,\ y-z^2+x^2,\ z-x^2+y^2)\)，则 \(\displaystyle \operatorname{div}\vec{F}=1+1+1=3\)。
    球体体积为 $\dfrac{4}{3}\pi R^3$，故积分为 \(\displaystyle 3\cdot \frac{4}{3}\pi R^3=4\pi R^3\)。
    \item 令 \(\displaystyle \vec{F}=(yz,zx,xy)\)，则 \(\displaystyle \operatorname{div}\vec{F}=0+0+0=0\)，因而曲面积分为 \(\displaystyle 0\)。
    \item 令 \(\displaystyle \vec{F}=\frac{(x,y,z)}{\sqrt{x^2+y^2+z^2}}\)。
    这题虽然也可以先算散度再走体积分，但在球面上 $\vec F$ 本身就是径向场，$\vec F\cdot \vec n$ 立刻化成常数，因此直接按“通量 = 常数 $\times$ 面积”更快。
    在球面 $r=a$ 上，\(\displaystyle \vec{n}=\frac{(x,y,z)}{a}\)，且 \(\displaystyle \vec{F}\cdot \vec{n}=\frac{a^2}{a\cdot a}=1\)，因而积分等于球面面积 \(\displaystyle \boxed{4\pi a^2}\)。
\end{compactenum}
\end{solution}

\begin{exercise}[向量 $\vec{r}$ 的通量]
求向量 $\vec{r}=x\vec{i}+y\vec{j}+z\vec{k}$ 的通量：
\begin{shortexenum}
    \shortitem{1}{通过圆锥形 $x^2+y^2\le z^2\ (0\le z\le h)$ 的侧表面；} &
    \shortitem{2}{通过此圆锥形的底。}
\end{shortexenum}
\end{exercise}
\begin{solution}
题目只分别问侧面和底面的通量，但这两部分天然拼成一个闭曲面；而 $\operatorname{div}\vec r$ 又极简单，所以先封口再用 Gauss，比单独硬算侧面法向投影稳得多。
考虑整个圆锥体，其边界由侧面和底面组成。圆锥体的底半径为 $h$，高为 $h$，体积 \(\displaystyle V=\frac13\pi h^3\)。由高斯公式，
\(\displaystyle \Phi_{\partial\Omega}
=\iint_{\partial\Omega}\vec{r}\cdot \vec{n}\,\mathrm{d}S
=\iiint_{\Omega}\operatorname{div}\vec{r}\,\mathrm{d}V
=3V=\pi h^3\)。

这里圆锥体的底面是平面 $z=h$ 上的圆盘 $x^2+y^2\le h^2$，其外法向量为 $\vec{k}$。因此
\[
\Phi_{\text{底}}
=\iint_{\text{底}} \vec{r}\cdot \vec{n}\,\mathrm{d}S
=\iint_{x^2+y^2\le h^2} z\,\mathrm{d}A
=h\cdot \pi h^2=\pi h^3.
\]
从而 \(\displaystyle \Phi_{\text{侧}}=\pi h^3-\pi h^3=0\)。
故 \(\displaystyle \boxed{\text{(1) 侧面通量 }=0,\quad \text{(2) 底面通量 }=\pi h^3}\)。
\end{solution}

\begin{exercise}[向量场 $\vec{A}=yz\vec{i}+zx\vec{j}+xy\vec{k}$ 的通量]
求向量 $\vec{A}=yz\vec{i}+zx\vec{j}+xy\vec{k}$ 的通量：
\begin{shortexenum}
    \shortitem{1}{通过圆柱 $x^2+y^2\le a^2,\ 0\le z\le h$ 的全表面；} &
    \shortitem{2}{通过此圆柱的侧表面。}
\end{shortexenum}
\end{exercise}
\begin{solution}
对 \(\displaystyle \vec{A}=(yz,zx,xy)\) 有 \(\displaystyle \operatorname{div}\vec{A}=0\)，因此通过整个闭圆柱面的总通量为 \(\displaystyle \boxed{0}\)。

再分别看上下底面。上底面 $z=h$，外法向量为 $\vec{k}$，通量为 $\iint_{x^2+y^2\le a^2}xy\,\mathrm{d}A=0$（被积函数关于 $x$ 或 $y$ 为奇函数）。下底面 $z=0$，通量也为 $0$。故侧面的通量等于全表面通量减去上下底面通量，仍为 $\boxed{0}$。
\end{solution}

\subsubsection{B组}

\begin{exercise}[带方向余弦的曲面积分]
计算曲面积分 \(\displaystyle \iint_S (x^2\cos\alpha+y^2\cos\beta+z^2\cos\gamma)\,\mathrm{d}S\)，其中 $S$ 为圆锥面 $x^2+y^2=z^2$ 介于平面 $z=0$ 及 $z=h\ (h>0)$ 之间的部分的下侧，$\cos\alpha,\cos\beta,\cos\gamma$ 是 $S$ 在点 $(x,y,z)$ 处的法向量方向余弦。
\end{exercise}
\begin{solution}
先封口再用 Gauss。令 \(\vec F=(x^2,y^2,z^2)\)，则 \(\operatorname{div}\vec F=2x+2y+2z\)。
把锥侧面与顶面 \(z=h\) 合起来看成闭曲面 \(\partial\Omega\)，柱坐标范围为 \(\displaystyle 0\le z\le h,\qquad 0\le r\le z,\qquad 0\le \theta\le 2\pi\)。
由对称性，\(x,y\) 两项体积分为零，于是
\[
\iint_{\partial\Omega}\vec F\cdot \vec n\,\mathrm{d}S
=2\int_0^{2\pi}\int_0^h\int_0^z zr\,\mathrm{d}r\,\mathrm{d}z\,\mathrm{d}\theta
=\frac{\pi h^4}{2}.
\]
顶面 \(z=h\) 的外法向为 \(\vec k\)，在顶面上 \(\vec F\cdot\vec k=h^2\)，圆盘面积为 \(\pi h^2\)。故顶面外向通量为 \(h^2\cdot\pi h^2=\pi h^4\)。
底面退化成顶点，不贡献通量，所以侧面外向通量为 \(\displaystyle \frac{\pi h^4}{2}-\pi h^4=-\frac{\pi h^4}{2}\)。
\end{solution}

\begin{exercise}[穿过曲面 $z=1-\sqrt{x^2+y^2}$ 的通量]
求向径 $\vec{r}=x\vec{i}+y\vec{j}+z\vec{k}$ 穿过曲面 \(\displaystyle S:\ z=1-\sqrt{x^2+y^2}\ (0\le z\le 1)\) 的流量，其中 $S$ 取与其和平面 $z=0$ 围成的圆锥体外法向一致的一侧。
\end{exercise}
\begin{solution}
先封口再用 Gauss。题设取外侧，曲面 \(S\) 与平面 \(z=0\) 围成圆锥体 \(\Omega\)，而 \(\operatorname{div}\vec r=3\)。
因此 \(\displaystyle \iint_{\partial\Omega}\vec r\cdot \vec n\,\mathrm{d}S
=\iiint_{\Omega}3\,\mathrm{d}V
=3\cdot \frac{\pi}{3}
=\pi\)。
底面在 \(z=0\) 上，\(\vec r\cdot \vec n=0\)，故曲面 \(S\) 上的通量为 \(\boxed{\pi}\)。
\end{solution}


\subsection{第八节习题：旋度和Stokes公式}

\subsubsection{A组}

\begin{exercise}[基本概念]
回答下列问题：
\begin{shorttrienum}
    \shortitem{1}{什么叫方向旋量？} &
    \shortitem{2}{向量场的旋度是怎样定义的？} &
    \shortitem{3}{斯托克斯公式成立需要什么条件？}
\end{shorttrienum}
\end{exercise}
\begin{solution}
\begin{compactenum}
\item 方向旋量是指向量场沿某一给定单位法向的小闭曲线环量与所围面积之比在面积趋于零时的极限，用来刻画场在该方向上的旋转强弱。
    \item 若 $\vec{F}=(P,Q,R)$，则其旋度定义为 \(\displaystyle \operatorname{rot}\vec{F}=\nabla\times \vec{F}
    =\left(\frac{\partial R}{\partial y}-\frac{\partial Q}{\partial z},\ \frac{\partial P}{\partial z}-\frac{\partial R}{\partial x},\ \frac{\partial Q}{\partial x}-\frac{\partial P}{\partial y}\right)\)。
    \item 斯托克斯公式成立的条件通常是：曲面 $S$ 为分片光滑定向曲面，其边界曲线 $\partial S$ 为分段光滑闭曲线，向量场 $\vec{F}$ 在包含 $S$ 的区域内具有连续一阶偏导数。此时
    \(\displaystyle \oint_{\partial S}\vec{F}\cdot \mathrm{d}\vec{r}
    =\iint_S (\operatorname{rot}\vec{F})\cdot \vec{n}\,\mathrm{d}S\)。
\end{compactenum}
\end{solution}

\begin{exercise}[利用斯托克斯公式计算曲线积分]
利用斯托克斯公式计算下列曲线积分：
\begin{shortsingleenum}
    \shortsingleitem{1}{$\displaystyle \oint_L z\,\mathrm{d}x+x\,\mathrm{d}y+y\,\mathrm{d}z$，其中 $L$ 为平面 $x+y+z=1$ 被三个坐标面所截成的三角形的整个边界，它的正向与这个三角形上侧的法向量之间符合右手规则；}
    \shortsingleitem{2}{$\displaystyle \oint_L (y-z)\,\mathrm{d}x+(z-x)\,\mathrm{d}y+(x-y)\,\mathrm{d}z$，其中 $L$ 为柱面 $x^2+y^2=a^2$ 和平面 $\dfrac{x}{a}+\dfrac{z}{b}=1\ (a>0,b>0)$ 的交线，即 $L$ 是椭圆边界，从 $x$ 轴正向看去，椭圆按逆时针方向。}
\end{shortsingleenum}
\end{exercise}
\begin{solution}
\begin{compactenum}
    \item 记 $\vec{F}=(z,x,y)$，则 $\operatorname{rot}\vec{F}=(1,1,1)$。取三角形平面片 $S$ 为积分曲面。平面 $x+y+z=1$ 的上侧单位法向量为 \(\displaystyle \vec{n}=\dfrac{1}{\sqrt{3}}(1,1,1)\)，故 $(\operatorname{rot}\vec{F})\cdot \vec{n}=\sqrt{3}$。该三角形的三个顶点为 $(1,0,0),(0,1,0),(0,0,1)$，面积为 $\dfrac{\sqrt{3}}{2}$。所以
    \(\displaystyle \oint_L z\,\mathrm{d}x+x\,\mathrm{d}y+y\,\mathrm{d}z
    =\sqrt{3}\cdot \frac{\sqrt{3}}{2}
    =\frac{3}{2}\)。
    \item 记 $\vec{F}=(y-z,z-x,x-y)$，则 $\operatorname{rot}\vec{F}=(-2,-2,-2)$。
    取平面截得的椭圆盘为积分曲面。该平面为 $\dfrac{x}{a}+\dfrac{z}{b}=1$，
    从 $x$ 轴正向看去边界逆时针，对应的单位法向量应取
    \[
      \vec{n}=\frac{(b,0,a)}{\sqrt{a^2+b^2}}.
    \]
    于是
    \[
      (\operatorname{rot}\vec{F})\cdot \vec{n}
      =-\frac{2(a+b)}{\sqrt{a^2+b^2}}.
    \]
    将平面写成 \(z=b\left(1-\dfrac{x}{a}\right)\)，则椭圆盘在 \(xOy\) 平面上的投影为圆盘 \(x^2+y^2\le a^2\)。由投影公式，
    \[
      \mathrm{d}S
      =\sqrt{1+z_x^2+z_y^2}\,\mathrm{d}x\,\mathrm{d}y
      =\frac{\sqrt{a^2+b^2}}{a}\,\mathrm{d}x\,\mathrm{d}y.
    \]
    因而椭圆盘面积为
    \[
      \iint_D \mathrm{d}S
      =\frac{\sqrt{a^2+b^2}}{a}
       \iint_{x^2+y^2\le a^2}\mathrm{d}x\,\mathrm{d}y
      =\pi a\sqrt{a^2+b^2}.
    \]
    所以
    \[
      \oint_L (y-z)\,\mathrm{d}x+(z-x)\,\mathrm{d}y+(x-y)\,\mathrm{d}z
      =-\frac{2(a+b)}{\sqrt{a^2+b^2}}\cdot \pi a\sqrt{a^2+b^2}
      =-2\pi a(a+b).
    \]
\end{compactenum}
\end{solution}

\begin{exercise}[求向量场的旋度]
求下列向量场的旋度：
\begin{shortexenum}
    \shortitem{1}{$\vec{F}=2xy\,\vec{i}+\mathrm{e}^x\sin y\,\vec{j}+(x^2+y^2+z^2)\,\vec{k}$；} &
    \shortitem{2}{$\vec{F}=\operatorname{grad}u$，其中 $u=u(x,y,z)$ 是具有二阶连续偏导数的函数。}
\end{shortexenum}
\end{exercise}
\begin{solution}
\begin{compactenum}
    \item 设 $P=2xy,\ Q=\mathrm{e}^x\sin y,\ R=x^2+y^2+z^2$。则
    \(\displaystyle \operatorname{rot}\vec{F}
    =\left(2y-0,\ 0-2x,\ \mathrm{e}^x\sin y-2x\right)\)，即 \(\displaystyle \boxed{\operatorname{rot}\vec{F}=(2y,\ -2x,\ \mathrm{e}^x\sin y-2x)}\)。
    \item 因为 $\nabla\times(\nabla u)=\left(u_{zy}-u_{yz},\ u_{xz}-u_{zx},\ u_{yx}-u_{xy}\right)$，而二阶连续偏导满足混合偏导可交换，所以 $\boxed{\operatorname{rot}(\operatorname{grad}u)=\vec{0}}$。
\end{compactenum}
\end{solution}

\subsubsection{B组}

\begin{exercise}[空间曲线上的斯托克斯公式]
利用斯托克斯公式计算下列曲线积分：
\begin{shortsingleenum}
    \shortsingleitem{1}{计算 \(\displaystyle \oint_L yz\,\mathrm{d}x+3zx\,\mathrm{d}y-xy\,\mathrm{d}z\)，其中 $L$ 是曲线 \(\displaystyle x^2+y^2=4y,\ 3y-z+1=0\)，从 $z$ 轴正向看去，$L$ 是沿逆时针方向；}
    \shortsingleitem{2}{计算 \(\displaystyle \oint_L y\,\mathrm{d}x+z\,\mathrm{d}y+x\,\mathrm{d}z\)，其中 $L$ 是球面 $x^2+y^2+z^2=R^2$ 与平面 $x+z=R$ 的交线，方向由点 $(R,0,0)$ 出发，先经过 $x>0,y>0$ 部分，再经 $x>0,y<0$ 部分回到出发点。}
\end{shortsingleenum}
\end{exercise}
\begin{solution}
\begin{compactenum}
    \item \textbf{思路：}这是空间闭曲线上的第二类曲线积分，先用 Stokes 把线积分换成曲面积分；又因为边界本身就是平面截线，所以最省的一步是直接取同边界的平面圆盘。

    记 $\vec{F}=(yz,3zx,-xy)$，则 $\operatorname{rot}\vec{F}=(-4x,\ 2y,\ 2z)$。交线已经落在平面 $z=3y+1$ 内，所以直接取该平面内由圆周围成的圆盘作积分曲面。该圆在 $xOy$ 平面上的投影为 $x^2+(y-2)^2=4$。由题意“从 $z$ 轴正向看去为逆时针方向”，可知法向量应取 $z$ 分量为正的一侧。平面法向量可取 $\vec{n}_0=(0,-3,1)$，故可用 $\mathrm{d}\vec{S}=(0,-3,1)\,\mathrm{d}x\,\mathrm{d}y$。
    故
    \[
    \oint_L yz\,\mathrm{d}x+3zx\,\mathrm{d}y-xy\,\mathrm{d}z
    =\iint_D \operatorname{rot}\vec{F}\cdot (0,-3,1)\,\mathrm{d}x\,\mathrm{d}y.
    \]
    又因 \(z=3y+1\)，被积函数为 \(-6y+2z=2\)，所以原积分为 \(\displaystyle \iint_D 2\,\mathrm{d}x\,\mathrm{d}y\)。把 $D$ 看成圆心在 $(0,2)$、半径为 $2$ 的圆盘，其面积为 $4\pi$，故 \(\displaystyle \oint_L yz\,\mathrm{d}x+3zx\,\mathrm{d}y-xy\,\mathrm{d}z=2\cdot 4\pi=8\pi\)。
    \item \textbf{思路：}同样先看 Stokes。边界是球面与平面的交线，最合适的曲面不是球面帽，而是同边界的平面圆盘。

    记 $\vec{F}=(y,z,x)$，则 $\operatorname{rot}\vec{F}=(-1,-1,-1)$。取平面 $x+z=R$ 内的圆盘为积分曲面。由题给方向可知法向量取 $\vec{n}=\dfrac{1}{\sqrt{2}}(1,0,1)$，故 $(\operatorname{rot}\vec{F})\cdot \vec{n}=-\sqrt{2}$。该平面到球心的距离为 $\dfrac{R}{\sqrt{2}}$，所以交圆半径 $\rho=\sqrt{R^2-\dfrac{R^2}{2}}=\dfrac{R}{\sqrt{2}}$，圆盘面积为 $\pi \rho^2=\dfrac{\pi R^2}{2}$。
    从而
    \(\displaystyle \oint_L y\,\mathrm{d}x+z\,\mathrm{d}y+x\,\mathrm{d}z=-\sqrt{2}\cdot \frac{\pi R^2}{2}=-\frac{\sqrt{2}}{2}\pi R^2\)。
\end{compactenum}
\end{solution}

\begin{exercise}[证明 $\operatorname{rot}(\vec{F}+\vec{G})=\operatorname{rot}\vec{F}+\operatorname{rot}\vec{G}$]
证明旋度的线性性质 \(\displaystyle \operatorname{rot}(\vec{F}+\vec{G})=\operatorname{rot}\vec{F}+\operatorname{rot}\vec{G}\)。
\end{exercise}
\begin{solution}
设 $\vec{F}=(P_1,Q_1,R_1),\ \vec{G}=(P_2,Q_2,R_2)$，则 $\vec{F}+\vec{G}=(P_1+P_2,\ Q_1+Q_2,\ R_1+R_2)$。
由旋度定义，
\begin{align*}
\operatorname{rot}(\vec{F}+\vec{G})
&=\left(
\frac{\partial(R_1+R_2)}{\partial y}-\frac{\partial(Q_1+Q_2)}{\partial z},
\frac{\partial(P_1+P_2)}{\partial z}-\frac{\partial(R_1+R_2)}{\partial x},
\frac{\partial(Q_1+Q_2)}{\partial x}-\frac{\partial(P_1+P_2)}{\partial y}
\right) \\
&=\left(\frac{\partial R_1}{\partial y}-\frac{\partial Q_1}{\partial z},\ \frac{\partial P_1}{\partial z}-\frac{\partial R_1}{\partial x},\ \frac{\partial Q_1}{\partial x}-\frac{\partial P_1}{\partial y}\right) \\
&\quad+\left(\frac{\partial R_2}{\partial y}-\frac{\partial Q_2}{\partial z},\ \frac{\partial P_2}{\partial z}-\frac{\partial R_2}{\partial x},\ \frac{\partial Q_2}{\partial x}-\frac{\partial P_2}{\partial y}\right) \\
&=\operatorname{rot}\vec{F}+\operatorname{rot}\vec{G}.
\end{align*}
证毕。
\end{solution}


\subsection{第九节习题：梯度算子}

\begin{exercise}[基本公式证明]
证明公式：
\begin{shortexenum}
    \shortitem{1}{$\nabla\cdot(C\vec{A})=C\,\nabla\cdot\vec{A}$，其中 $C$ 为常数；} &
    \shortitem{2}{$\nabla\cdot(u\vec{C})=\nabla u\cdot \vec{C}$，其中 $\vec{C}$ 为常向量。}
\end{shortexenum}
\end{exercise}
\begin{solution}
\begin{compactenum}
    \item 设 $\vec{A}=(P,Q,R)$，则 $C\vec{A}=(CP,CQ,CR)$。所以
    \begin{align*}
    \nabla\cdot(C\vec{A})
    &=\frac{\partial(CP)}{\partial x}
      +\frac{\partial(CQ)}{\partial y}
      +\frac{\partial(CR)}{\partial z}\\
    &=C\left(\frac{\partial P}{\partial x}
      +\frac{\partial Q}{\partial y}
      +\frac{\partial R}{\partial z}\right)
    =C\,\nabla\cdot \vec{A}.
    \end{align*}
    \item 设 $\vec{C}=(a,b,c)$ 为常向量，则 $u\vec{C}=(au,bu,cu)$。因而
    \[
    \nabla\cdot(u\vec{C})
    =a\frac{\partial u}{\partial x}
    +b\frac{\partial u}{\partial y}
    +c\frac{\partial u}{\partial z}
    =\nabla u\cdot \vec{C}.
    \]
\end{compactenum}
\end{solution}

\begin{exercise}[证明 $\nabla\cdot \nabla u=\Delta u$]
证明 \(\displaystyle \nabla\cdot \nabla u=\Delta u\)，其中 $u$ 是数量场，$\Delta$ 为拉普拉斯算子。
\end{exercise}
\begin{solution}
由梯度定义，
\begin{align*}
\nabla u&=\left(\dfrac{\partial u}{\partial x},
\dfrac{\partial u}{\partial y},
\dfrac{\partial u}{\partial z}\right),\\
\nabla\cdot \nabla u
&=\frac{\partial^2 u}{\partial x^2}
+\frac{\partial^2 u}{\partial y^2}
+\frac{\partial^2 u}{\partial z^2}.
\end{align*}
而右端正是拉普拉斯算子 $\Delta u$ 的定义，因此 \(\displaystyle \boxed{\nabla\cdot \nabla u=\Delta u}\)。
\end{solution}

\begin{exercise}[梯度的旋度]
由直接计算验证，$f$ 的梯度的旋度是 $0$，即 \(\nabla\times(\nabla f)=0\)。
\end{exercise}
\begin{solution}
设 $f=f(x,y,z)$ 具有二阶连续偏导数，则 $\nabla f=(f_x,f_y,f_z)$。于是
\[
\nabla\times(\nabla f)=
\left(f_{zy}-f_{yz},\ f_{xz}-f_{zx},\ f_{yx}-f_{xy}\right).
\]
由混合偏导可交换，$f_{zy}=f_{yz},\ f_{xz}=f_{zx},\ f_{yx}=f_{xy}$，故三分量全为零，即 \(\displaystyle \boxed{\nabla\times(\nabla f)=\vec{0}}\)。
\end{solution}

\begin{exercise}[旋度的散度]
由直接计算验证，向量场 $\vec{A}=A_1\vec{i}+A_2\vec{j}+A_3\vec{k}$ 的旋度的散度为 $0$，即 \(\displaystyle \nabla\cdot(\nabla\times \vec{A})=0\)。
\end{exercise}
\begin{solution}
由定义，
\[
\nabla\times \vec{A}
=\left(\frac{\partial A_3}{\partial y}-\frac{\partial A_2}{\partial z},\ \frac{\partial A_1}{\partial z}-\frac{\partial A_3}{\partial x},\ \frac{\partial A_2}{\partial x}-\frac{\partial A_1}{\partial y}\right).
\]
再取散度：
\begin{align*}
\nabla\cdot(\nabla\times \vec{A})
&=\frac{\partial}{\partial x}\left(\frac{\partial A_3}{\partial y}-\frac{\partial A_2}{\partial z}\right)
+\frac{\partial}{\partial y}\left(\frac{\partial A_1}{\partial z}-\frac{\partial A_3}{\partial x}\right)
+\frac{\partial}{\partial z}\left(\frac{\partial A_2}{\partial x}-\frac{\partial A_1}{\partial y}\right) \\
&=A_{3yx}-A_{2zx}+A_{1zy}-A_{3xy}+A_{2xz}-A_{1yz}.
\end{align*}
由混合偏导可交换，三对项两两抵消，故 \(\displaystyle \boxed{\nabla\cdot(\nabla\times \vec{A})=0}\)。
\end{solution}

\begin{exercise}[证明 $\operatorname{div}(\vec{F}\times\vec{G})=\vec{G}\cdot \operatorname{rot}\vec{F}-\vec{F}\cdot \operatorname{rot}\vec{G}$]
设 $\vec{F}=P\vec{i}+Q\vec{j}+R\vec{k}$，$\vec{G}=L\vec{i}+M\vec{j}+N\vec{k}$，证明 \(\displaystyle \operatorname{div}(\vec{F}\times\vec{G})=\vec{G}\cdot \operatorname{rot}\vec{F}-\vec{F}\cdot \operatorname{rot}\vec{G}\)。
\end{exercise}
\begin{solution}
先写出叉积：\(\displaystyle \vec{F}\times\vec{G}=(QN-RM,\ RL-PN,\ PM-QL)\)。所以 \(\displaystyle \operatorname{div}(\vec{F}\times\vec{G}) =\frac{\partial(QN-RM)}{\partial x}+\frac{\partial(RL-PN)}{\partial y}+\frac{\partial(PM-QL)}{\partial z}\)。
逐项展开并按含 $\vec{F}$ 的偏导与含 $\vec{G}$ 的偏导分组，可得
\begin{align*}
\operatorname{div}(\vec{F}\times\vec{G})
&=L\left(\frac{\partial R}{\partial y}-\frac{\partial Q}{\partial z}\right)
+M\left(\frac{\partial P}{\partial z}-\frac{\partial R}{\partial x}\right)
+N\left(\frac{\partial Q}{\partial x}-\frac{\partial P}{\partial y}\right) \\
&\quad-\Bigg[
P\left(\frac{\partial N}{\partial y}-\frac{\partial M}{\partial z}\right)
+Q\left(\frac{\partial L}{\partial z}-\frac{\partial N}{\partial x}\right)
+R\left(\frac{\partial M}{\partial x}-\frac{\partial L}{\partial y}\right)
\Bigg].
\end{align*}
这正是 $\vec{G}\cdot \operatorname{rot}\vec{F}-\vec{F}\cdot \operatorname{rot}\vec{G}$，故命题成立。
\end{solution}


\subsection{本章总习题}

本章总习题不再只考某一个孤立公式，常见入口是混合出现的：看到 \(\mathrm{d}s\) 或 \(\mathrm{d}S\) 先按第一类积分处理；看到 \(P\,\mathrm{d}x+Q\,\mathrm{d}y(+R\,\mathrm{d}z)\) 先判开路保守场、平面闭路 Green、空间闭路 Stokes；看到 \(P\,\mathrm{d}y\,\mathrm{d}z+Q\,\mathrm{d}z\,\mathrm{d}x+R\,\mathrm{d}x\,\mathrm{d}y\) 先判投影、封口与 Gauss。若题目含奇点、洞、补边、补面或方向描述，先把合法性和方向账本写清，再进入计算。

% ============================================================
% 零、补充题目
% ============================================================

\subsubsection{补充题目}

\begin{exercise}[课堂补充题目]
\begin{shortsingleenum}
  \shortsingleitem{1}{（20分）计算 $\int_L \mathrm{e}^{\sqrt{x^2+y^2}}\,\mathrm{d}s$，其中 $L$ 为圆周 $x^2+y^2=a^2\ (a>0)$、直线 $y=x$ 与 $x$ 轴在第一象限内所围成的扇形的整个边界。}
  \shortsingleitem{2}{（20分）计算 $\iint_\Sigma x\,\mathrm{d}y\,\mathrm{d}z$，其中 $\Sigma$ 为 $x^2+y^2=1,\ z=x+2$ 和 $z=0$ 所围成立体的表面积，取外侧。}
  \shortsingleitem{3}{（20分）设曲线积分 $\int_L xy^2\,\mathrm{d}x+y\phi(x)\,\mathrm{d}y$ 与路径无关，其中 $\phi(x)$ 具有连续的导数，且 $\phi(0)=0$，计算 $\int_{(0,0)}^{(1,1)} xy^2\,\mathrm{d}x+y\phi(x)\,\mathrm{d}y$。}
  \shortsingleitem{4}{（20分）设 $\Sigma$ 是锥面 $z=\sqrt{x^2+y^2}$ 被平面 $z=0$ 及 $z=1$ 截下的部分的下侧。计算第二型曲面积分 \(\displaystyle I=\iint_\Sigma x\,\mathrm{d}y\,\mathrm{d}z+y\,\mathrm{d}z\,\mathrm{d}x+(z^2-2z)\,\mathrm{d}x\,\mathrm{d}y\)。}
  \shortsingleitem{5}{（10分）设曲线 $L$ 为球面 $x^2+y^2+z^2=1$ 与平面 $x+2y+z=0$ 的交线，计算积分 $\int_L (2x^2+x+y+y^2)\,\mathrm{d}s$。}
\end{shortsingleenum}
\end{exercise}
\begin{solution}
\begin{compactenum}
    \item 将边界 $L$ 分成 $x$ 轴上的线段 $L_1$、圆弧 $L_2$ 及直线 $y=x$ 上的线段 $L_3$。于是 \(\displaystyle \int_{L_1}\mathrm{e}^{\sqrt{x^2+y^2}}\,\mathrm{d}s=\int_0^a \mathrm{e}^x\,\mathrm{d}x=\mathrm{e}^a-1\)。
    在圆弧 $L_2$ 上有 $\sqrt{x^2+y^2}=a$，且弧长为 $a\cdot\dfrac{\pi}{4}$，故 \(\displaystyle \int_{L_2}\mathrm{e}^{\sqrt{x^2+y^2}}\,\mathrm{d}s=\mathrm{e}^a\cdot a\frac{\pi}{4}\)。
    对 $L_3$ 取参数 $x=t,\ y=t,\ 0\le t\le \dfrac{a}{\sqrt2}$，则 $\sqrt{x^2+y^2}=\sqrt2\,t,\ \mathrm{d}s=\sqrt2\,\mathrm{d}t$，从而 \(\displaystyle \int_{L_3}\mathrm{e}^{\sqrt{x^2+y^2}}\,\mathrm{d}s=\int_0^{a/\sqrt2}\mathrm{e}^{\sqrt2\,t}\sqrt2\,\mathrm{d}t=\mathrm{e}^a-1\)。
    所以 \(\displaystyle \boxed{\int_L \mathrm{e}^{\sqrt{x^2+y^2}}\,\mathrm{d}s=2(\mathrm{e}^a-1)+\frac{\pi a}{4}\mathrm{e}^a}\)。
    \item 这一问最适合用高斯公式。先把高斯公式写成第二型曲面积分常见的样子：
    \[
    \iint_\Sigma P\,\mathrm{d}y\,\mathrm{d}z+Q\,\mathrm{d}z\,\mathrm{d}x+R\,\mathrm{d}x\,\mathrm{d}y
    =\iiint_V \left(\frac{\partial P}{\partial x}+\frac{\partial Q}{\partial y}+\frac{\partial R}{\partial z}\right)\mathrm{d}V,
    \]
    其中 $\Sigma$ 是闭曲面，取外侧，$V$ 是它所围成的立体。
    本题里 \(\displaystyle P=x,\ Q=0,\ R=0\)，所以 \(\displaystyle \frac{\partial P}{\partial x}+\frac{\partial Q}{\partial y}+\frac{\partial R}{\partial z}=1\)。
    因而原积分立刻化为 \(\displaystyle \iint_\Sigma x\,\mathrm{d}y\,\mathrm{d}z=\iiint_V 1\,\mathrm{d}V\)。
    右边就是立体 $V$ 的体积。

    下面把立体 $V$ 写清楚。它由圆柱面 $x^2+y^2=1$、上平面 $z=x+2$ 和下平面 $z=0$ 围成。
    由于圆柱内有 $-1\le x\le 1$，所以 $x+2\ge 1>0$，也就是说上平面始终在下平面上方，故 \(\displaystyle V=\{(x,y,z)\mid x^2+y^2\le 1,\ 0\le z\le x+2\}\)。
    于是可先对 $z$ 积分，并把它拆开：记 $D=\{(x,y)\mid x^2+y^2\le 1\}$，则 \(\displaystyle \iiint_V 1\,\mathrm{d}V
    =\iint_D (x+2)\,\mathrm{d}x\,\mathrm{d}y
    =\iint_D x\,\mathrm{d}x\,\mathrm{d}y+2\iint_D\mathrm{d}x\,\mathrm{d}y\)。
    第一项为 $0$，因为单位圆盘关于 $y$ 轴对称，点 $(x,y)$ 与 $(-x,y)$ 对应的函数值互为相反数，积分彼此抵消。
    第二项等于 $2\times$ 单位圆面积，即 \(\displaystyle 2\iint_{x^2+y^2\le 1}\mathrm{d}x\,\mathrm{d}y=2\pi\)，
    所以 \(\displaystyle \boxed{\iint_\Sigma x\,\mathrm{d}y\,\mathrm{d}z=2\pi}\)。
    \item 设 $P=xy^2,\ Q=y\phi(x)$。因积分与路径无关，必有
    \(\displaystyle \frac{\partial P}{\partial y}=2xy=\frac{\partial Q}{\partial x}=y\phi'(x)\)，
    从而 \(\displaystyle \phi'(x)=2x,\ \phi(0)=0,\ \phi(x)=x^2\)。
    这时可取势函数 \(\displaystyle u(x,y)=\frac12 x^2y^2\)，因为 $u_x=xy^2,\ u_y=x^2y=y\phi(x)$。故
    \(\displaystyle \int_{(0,0)}^{(1,1)} xy^2\,\mathrm{d}x+y\phi(x)\,\mathrm{d}y
    =u(1,1)-u(0,0)=\frac12\)，即
    \(\displaystyle \boxed{\int_{(0,0)}^{(1,1)} xy^2\,\mathrm{d}x+y\phi(x)\,\mathrm{d}y=\frac12}\)。
    \item 记 \(\displaystyle \vec F=(x,y,z^2-2z)\)。
    设 $V$ 为锥面与平面 $z=1$ 围成的立体，顶点 $z=0$ 处只是一个点，不必另加底面。再记顶圆盘为 $D:\ z=1,\ x^2+y^2\le 1$。由于锥面的下侧正是闭曲面的外侧，故由高斯公式有 \(\displaystyle I+\iint_D \vec F\cdot \vec n\,\mathrm{d}S=\iiint_V \nabla\cdot \vec F\,\mathrm{d}V\)。
    而 \(\displaystyle \nabla\cdot \vec F=1+1+(2z-2)=2z\)。
    用柱坐标表示 $V:\ 0\le z\le 1,\ 0\le r\le z,\ 0\le \theta\le 2\pi$，则 \(\displaystyle \iiint_V 2z\,\mathrm{d}V=\int_0^{2\pi}\int_0^1\int_0^z 2zr\,\mathrm{d}r\,\mathrm{d}z\,\mathrm{d}\theta=\frac{\pi}{2}\)。
    在顶圆盘 $D$ 上，外法向量向上，故 \(\displaystyle \iint_D \vec F\cdot \vec n\,\mathrm{d}S=\iint_D (z^2-2z)\,\mathrm{d}x\,\mathrm{d}y=-\iint_D \mathrm{d}x\,\mathrm{d}y=-\pi\)。
    所以 \(\displaystyle \boxed{I=\frac{\pi}{2}-(-\pi)=\frac{3\pi}{2}}\)。
    \item 取平面 $x+2y+z=0$ 内的一组标准正交基 \(\displaystyle \vec e_1=\left(\frac{1}{\sqrt2},0,-\frac{1}{\sqrt2}\right),\ \vec e_2=\left(-\frac{1}{\sqrt3},\frac{1}{\sqrt3},-\frac{1}{\sqrt3}\right)\)。
    交线 $L$ 是单位圆，可参数化为 \(\displaystyle \vec r(t)=\vec e_1\cos t+\vec e_2\sin t,\ 0\le t\le 2\pi\)。于是 \(\displaystyle x=\frac{\cos t}{\sqrt2}-\frac{\sin t}{\sqrt3},\ y=\frac{\sin t}{\sqrt3},\ \mathrm{d}s=\mathrm{d}t\)。
    代入得
    \begin{align*}
    2x^2+x+y+y^2
    &=2\left(\frac{\cos t}{\sqrt2}-\frac{\sin t}{\sqrt3}\right)^2
    +\left(\frac{\cos t}{\sqrt2}-\frac{\sin t}{\sqrt3}\right)
    +\frac{\sin t}{\sqrt3}+\frac{\sin^2 t}{3} \\
    &=1-\frac{4}{\sqrt6}\sin t\cos t+\frac{\cos t}{\sqrt2}.
    \end{align*}
    故
    \[
    \int_L (2x^2+x+y+y^2)\,\mathrm{d}s
    =\int_0^{2\pi}\left(1-\frac{4}{\sqrt6}\sin t\cos t+\frac{\cos t}{\sqrt2}\right)\mathrm{d}t
    =2\pi.
    \]
    即 \(\displaystyle \boxed{\int_L (2x^2+x+y+y^2)\,\mathrm{d}s=2\pi}\)。
\end{compactenum}
\end{solution}

% ============================================================
% 一、第一类曲线积分与向量场基础
% ============================================================

\begin{exercise}[第一类曲线积分与做功]
\begin{shortsingleenum}
    \shortsingleitem{1}{计算第一类曲线积分 \(\displaystyle \int_L (x^2+y^2)\,\mathrm{d}s\)，其中 $L$ 为上半圆 $x^2+y^2=1,\ y\ge 0$；}
    \shortsingleitem{2}{题图中折线 $O\to A\to B$ 上的第一类曲线积分（按图中题面识别），答案为 \underline{\hspace{1cm}}；}
    \shortsingleitem{3}{计算第二类曲线积分 \(\displaystyle \int_L xy\,\mathrm{d}x+(y-x)\,\mathrm{d}y\)，其中 $L$ 为从 $O(0,0)$ 到 $A(1,1)$ 的线段；}
    \shortsingleitem{4}{求心形线 $r=a(1-\cos\varphi)$ 的质心（取线密度 $\rho=1$）；}
    \shortsingleitem{5}{两个定点质量分别位于 $A(1,0)$ 与 $B(0,1)$，质量为 $m$ 的质点沿以 $A$ 为圆心的上半圆弧从 $O$ 运动到 $D(2,0)$，求重力做功；}
    \shortsingleitem{6}{设螺旋线 $x=a\cos t,\ y=a\sin t,\ z=bt$，质点由 $A(a,0,0)$ 运动到 $B(a,0,2\pi b)$，所受力大小与到原点距离成正比、方向指向原点，求做功；}
    \shortsingleitem{7}{设向量场 $\vec{F}=(yz,zx,xy)$，质点沿从原点到椭球面 \(\displaystyle \frac{x^2}{a^2}+\frac{y^2}{b^2}+\frac{z^2}{c^2}=1\) 的第一卦限点的直线段运动，求最大做功。}
\end{shortsingleenum}
\end{exercise}
\begin{solution}
\begin{compactenum}
    \item 取参数 $x=\cos t,\ y=\sin t,\ 0\le t\le\pi$，则 $\mathrm{d}s=\mathrm{d}t$，故
    \(\displaystyle \int_L (x^2+y^2)\,\mathrm{d}s=\int_0^\pi 1\,\mathrm{d}t=\pi\)。
    \item 按题图所示结果为 $\boxed{\frac12+\sqrt2}$。
    \item 线段 $L:\ x=t,\ y=t,\ 0\le t\le1$，于是
    \(\displaystyle \int_L xy\,\mathrm{d}x+(y-x)\,\mathrm{d}y
    =\int_0^1 t^2\,\mathrm{d}t=\frac13\)。
    \item 心形线关于 $x$ 轴对称，故 $\bar y=0$。设参数为极角 $\varphi$，则
    \(\displaystyle r=a(1-\cos\varphi),\ \frac{\mathrm{d}r}{\mathrm{d}\varphi}=a\sin\varphi\)，因而弧长元
    \begin{align*}
    \mathrm{d}s
    &=\sqrt{r^2+\left(\frac{\mathrm{d}r}{\mathrm{d}\varphi}\right)^2}\,\mathrm{d}\varphi
    =a\sqrt{(1-\cos\varphi)^2+\sin^2\varphi}\,\mathrm{d}\varphi \\
    &=a\sqrt{2-2\cos\varphi}\,\mathrm{d}\varphi
    =2a\sin\frac{\varphi}{2}\,\mathrm{d}\varphi.
    \end{align*}
    于是曲线总长为
    \(\displaystyle L=\int_0^{2\pi}\mathrm{d}s
    =2a\int_0^{2\pi}\sin\frac{\varphi}{2}\,\mathrm{d}\varphi
    =8a\)。
    又 $x=r\cos\varphi=a(1-\cos\varphi)\cos\varphi$，所以 \(\displaystyle \int_L x\,\mathrm{d}s=2a^2\int_0^{2\pi}(1-\cos\varphi)\cos\varphi\sin\frac{\varphi}{2}\,\mathrm{d}\varphi=-\frac{32a^2}{5}\)。
    因此 \(\displaystyle \bar x=\frac{1}{L}\int_L x\,\mathrm{d}s
    =\frac{-32a^2/5}{8a}
    =-\frac{4a}{5}\)，故质心为 $\boxed{(\bar x,\bar y)=\left(-\frac{4a}{5},\,0\right)}$。
    \item 两个定点对质点的引力都是保守力，故总做功等于总势函数的端点差。
    点 $A(1,0)$ 到起点 $O(0,0)$ 和终点 $D(2,0)$ 的距离都等于 $1$，所以 $A$ 对应的势能改变量为 $0$。
    点 $B(0,1)$ 到两端点的距离分别为 $BO=1$ 与 $BD=\sqrt5$，因而只有点 $B$ 产生净做功：\(\displaystyle W=m\left(\frac{1}{BD}-\frac{1}{BO}\right)=\boxed{m\left(\frac1{\sqrt5}-1\right)}\)。
    \item 力场可写成 $\vec{F}=-k\|\vec r\|\,\vec r$，其势函数为 $\Phi=\dfrac{k}{3}\|\vec r\|^3$，故
    \(\displaystyle \boxed{W=\frac{k}{3}\left(a^3-(a^2+4\pi^2b^2)^{3/2}\right)}\)。
    \item 先验证该力场是保守场。若取势函数
    \(\displaystyle u(x,y,z)=xyz,\ \nabla u=(yz,zx,xy)=\vec F\)，
    因而从原点沿任意路径到点 $(\xi,\eta,\zeta)$ 的做功都等于
    \(\displaystyle W=u(\xi,\eta,\zeta)-u(0,0,0)=\xi\eta\zeta\)。
    于是问题化为在约束
    \(\displaystyle \frac{\xi^2}{a^2}+\frac{\eta^2}{b^2}+\frac{\zeta^2}{c^2}=1,\ \xi,\eta,\zeta\ge 0\)
    下求 $\xi\eta\zeta$ 的最大值。
    令 $X=\dfrac{\xi^2}{a^2},\ Y=\dfrac{\eta^2}{b^2},\ Z=\dfrac{\zeta^2}{c^2}$，则 $X+Y+Z=1$。由 AM-GM 不等式，\(\displaystyle XYZ\le \left(\frac{X+Y+Z}{3}\right)^3=\frac{1}{27}\)，
    等号当且仅当 $X=Y=Z=\dfrac13$ 成立。故当 $\xi=\dfrac{a}{\sqrt3},\ \eta=\dfrac{b}{\sqrt3},\ \zeta=\dfrac{c}{\sqrt3}$ 时取等号，且
    \(\displaystyle \boxed{W_{\max}=\frac{abc}{3\sqrt3}}\)。
\end{compactenum}
\end{solution}

% ============================================================
% 二、第一类曲面积分与第二类曲面积分
% ============================================================

\begin{exercise}[第一类与第二类曲面积分]
\begin{shortsingleenum}
    \shortsingleitem{1}{计算第一类曲面积分 \(\displaystyle \iint_{\Sigma}(x^2+y^2+z^2)\,\mathrm{d}S\)，其中 $\Sigma:x^2+y^2+z^2=2ax$；}
    \shortsingleitem{2}{计算 \(\displaystyle \iint_S (x+y+z)\,\mathrm{d}S\)，其中 $S$ 为平面 $y+z=5$ 被圆柱 $x^2+y^2=25$ 截出的部分；}
    \shortsingleitem{3}{计算 \(\displaystyle \iint_\Sigma |xyz|\,\mathrm{d}S\)，其中 $\Sigma:z=x^2+y^2,\ z\le 1$；}
    \shortsingleitem{4}{计算闭曲面上的第二类曲面积分 \(\displaystyle \oiint_{\Sigma}\frac{xz\,\mathrm{d}y\,\mathrm{d}z+z^2\,\mathrm{d}x\,\mathrm{d}y}{x^2+y^2+z^2}\)，其中 $\Sigma$ 为圆柱面 $x^2+y^2=R^2$ 与平面 $z=\pm R$ 围成的闭曲面。}
\end{shortsingleenum}
\end{exercise}
\begin{solution}
\begin{compactenum}
    \item 在球面 $\Sigma:x^2+y^2+z^2=2ax$ 上恒有 $x^2+y^2+z^2=2ax$，故
    \(\displaystyle \iint_{\Sigma}(x^2+y^2+z^2)\,\mathrm{d}S
    =2a\iint_\Sigma x\,\mathrm{d}S\)。
    该球面可写成 $(x-a)^2+y^2+z^2=a^2$，即球心在 $(a,0,0)$、半径为 $a$。由球面对称性，\(\displaystyle \frac{1}{\operatorname{Area}(\Sigma)}\iint_\Sigma x\,\mathrm{d}S=a,\ \iint_\Sigma x\,\mathrm{d}S=a\cdot 4\pi a^2=4\pi a^3\)。
    代回得 \(\displaystyle \iint_{\Sigma}(x^2+y^2+z^2)\,\mathrm{d}S=2a\cdot 4\pi a^3=8\pi a^4\)。
    \item 平面上 $z=5-y$，故 $\mathrm{d}S=\sqrt2\,\mathrm{d}x\,\mathrm{d}y$，投影为圆盘 $x^2+y^2\le25$。于是 \(\displaystyle \iint_S (x+y+z)\,\mathrm{d}S=\sqrt2\iint_{x^2+y^2\le25}(x+5)\,\mathrm{d}x\,\mathrm{d}y\)。
    其中 $\iint x\,\mathrm{d}x\,\mathrm{d}y=0$，因为积分区域关于 $y$ 轴对称；故
    \(\displaystyle \iint_S (x+y+z)\,\mathrm{d}S
    =5\sqrt2\iint_{x^2+y^2\le25}\mathrm{d}x\,\mathrm{d}y
    =125\sqrt2\,\pi\)。
    \item 取极坐标 $x=r\cos\theta,\ y=r\sin\theta,\ z=r^2$，则 \(\displaystyle \mathrm{d}S=\sqrt{1+4r^2}\,r\,\mathrm{d}r\,\mathrm{d}\theta\)，且 \(\displaystyle |xyz|=r^4|\cos\theta\sin\theta|\)。
    于是
    \begin{align*}
    \iint_\Sigma |xyz|\,\mathrm{d}S
    &=\int_0^{2\pi}\!\!\int_0^1
    r^5|\cos\theta\sin\theta|\sqrt{1+4r^2}\,\mathrm{d}r\,\mathrm{d}\theta \\
    &=\left(\int_0^{2\pi}|\cos\theta\sin\theta|\,\mathrm{d}\theta\right)
    \left(\int_0^1 r^5\sqrt{1+4r^2}\,\mathrm{d}r\right).
    \end{align*}
    角向积分为
    \(\displaystyle \int_0^{2\pi}|\cos\theta\sin\theta|\,\mathrm{d}\theta
    =\frac12\int_0^{2\pi}|\sin2\theta|\,\mathrm{d}\theta=2\)。
    对径向积分令 $u=1+4r^2$，则 $\mathrm{d}u=8r\,\mathrm{d}r$，从而
    \begin{align*}
    \int_0^1 r^5\sqrt{1+4r^2}\,\mathrm{d}r
    &=\frac{1}{128}\int_1^5 (u-1)^2u^{1/2}\,\mathrm{d}u
    =\frac{1}{128}\int_1^5 \left(u^{5/2}-2u^{3/2}+u^{1/2}\right)\mathrm{d}u \\
    &=\frac{125\sqrt5-1}{840}.
    \end{align*}
    因此 \(\displaystyle
    \iint_\Sigma |xyz|\,\mathrm{d}S
    =2\cdot \frac{125\sqrt5-1}{840}
    =\frac{125\sqrt5-1}{420}\)。
    \item \textbf{思路：}对象虽是闭曲面，但场在原点有奇点，而原点恰在圆柱内部，所以不要直接整块套 Gauss；这里改走“侧面与上下底分开，再看对称性”更稳。

    记 \(\displaystyle \vec F=\left(\frac{xz}{x^2+y^2+z^2},\ 0,\ \frac{z^2}{x^2+y^2+z^2}\right)\)。
    在侧面 $x^2+y^2=R^2$ 上，外法向量为 \(\displaystyle \vec n=\frac{1}{R}(x,y,0)\)，且 \(\displaystyle \vec F\cdot \vec n=\frac{x^2z}{R(x^2+y^2+z^2)}\)，
    它关于 $z$ 是奇函数，所以侧面通量为 $0$。上底 $z=R$ 与下底 $z=-R$ 上，$F_z$ 的大小相同，而外法向方向相反，因此两块底面的通量也恰好相消。故 $\boxed{0}$。
\end{compactenum}
\end{solution}

% ============================================================
% 三、曲线积分性质与 Green 公式
% ============================================================

\begin{exercise}[曲线积分性质与 Green 公式]
\begin{compactenum}
    \item 判断下列命题真伪：
    \begin{shortexenum}
        \shortitem{1}{$\displaystyle\int_C \vec F\cdot\mathrm{d}\vec s$ 是一个向量；} &
        \shortitem{2}{$\displaystyle\int_C \vec F\cdot\mathrm{d}\vec s=F(B)-F(A)$；} \\
        \shortitem{3}{若单位圆周上的一个第二类曲线积分为 $0$，则该向量场必为保守场；} &
        \shortitem{4}{$\displaystyle V=\frac13\oiint_S(x\cos\alpha+y\cos\beta+z\cos\gamma)\,\mathrm{d}S$。}
    \end{shortexenum}
    \item 计算 \(\displaystyle \oint_L xy^2\,\mathrm{d}y-x^2y\,\mathrm{d}x\)，其中 $L$ 为椭圆 $\dfrac{x^2}{a^2}+\dfrac{y^2}{b^2}=1$ 的正向边界；
    \item 计算 \(\displaystyle \oint_L \left(\mathrm{e}^y\sin x\,\mathrm{d}x+\mathrm{e}^{-x}\sin y\,\mathrm{d}y\right)\)，其中 $L$ 为矩形 $a\le x\le b,\ c\le y\le d$ 的正向边界；
    \item 证明
    \[
    \iint_D\left(P\frac{\partial R}{\partial x}+Q\frac{\partial R}{\partial y}\right)\,\mathrm{d}x\,\mathrm{d}y
    =\oint_C PR\,\mathrm{d}y-QR\,\mathrm{d}x
    -\iint_D R\left(\frac{\partial P}{\partial x}+\frac{\partial Q}{\partial y}\right)\,\mathrm{d}x\,\mathrm{d}y;
    \]
    \item 求 $\varphi(x)$，使 \(\displaystyle \int_L \left(\sin x-\varphi(x)\right)\frac{y}{x}\,\mathrm{d}x+\varphi(x)\,\mathrm{d}y\) 与路径无关，且 $\varphi(\pi)=1$；并求从 $A(1,0)$ 到 $B(\pi,\pi)$ 的积分值；
    \item 求微分 \(\displaystyle \mathrm{d}u=\frac{(x+2y)\,\mathrm{d}x+y\,\mathrm{d}y}{(x+y)^2}\) 所对应的原函数。
\end{compactenum}
\end{exercise}
\begin{solution}
\begin{compactenum}
    \item 依次为：假、假、假、真。
    其中：
    \begin{itemize}
        \item 第一类或第二类曲线积分的结果都是数量，不是向量；
        \item 一般只有在 $\vec F$ 为保守场且 $F$ 写成势函数时，才有“端点差”公式；
        \item 某一条闭曲线积分为 $0$ 只说明这条曲线上的积分恰好为零，不能推出对任意闭曲线都为零；
        \item 这是体积公式 $\displaystyle V=\frac13\oiint_S \vec r\cdot \vec n\,\mathrm{d}S$ 的常见处理。
    \end{itemize}
    \item \textbf{思路：}这是平面闭曲线围成的椭圆边界，先用 Green 把边界积分翻成区域积分；椭圆区域再用伸缩代换化成单位圆盘最顺。

    设 $P=-x^2y,\ Q=xy^2$，则 \(\displaystyle \frac{\partial Q}{\partial x}-\frac{\partial P}{\partial y}=x^2+y^2\)。由 Green 公式，
    \(\displaystyle \oint_L xy^2\,\mathrm{d}y-x^2y\,\mathrm{d}x
    =\iint_D (x^2+y^2)\,\mathrm{d}x\,\mathrm{d}y\)。
    再作椭圆区域的伸缩代换 $x=ar\cos\theta,\ y=br\sin\theta,\ 0\le r\le 1,\ 0\le \theta\le 2\pi$，其 Jacobian 为 $ab\,r$，故
    \begin{align*}
    \iint_D (x^2+y^2)\,\mathrm{d}x\,\mathrm{d}y
    &=ab\int_0^{2\pi}\!\!\int_0^1
    \left(a^2r^2\cos^2\theta+b^2r^2\sin^2\theta\right)r\,\mathrm{d}r\,\mathrm{d}\theta \\
    &=ab\left(\int_0^1 r^3\,\mathrm{d}r\right)
    \left(a^2\int_0^{2\pi}\cos^2\theta\,\mathrm{d}\theta
    +b^2\int_0^{2\pi}\sin^2\theta\,\mathrm{d}\theta\right) \\
    &=ab\cdot \frac14\cdot \pi(a^2+b^2)
    =\frac{\pi ab(a^2+b^2)}{4}.
    \end{align*}
    \item \textbf{思路：}边界是矩形闭曲线，若逐边参数化会比较碎；先用 Green 化成矩形上的二重积分，再按 $x,y$ 分离计算更稳定。

    由 Green 公式，
    \begin{align*}
    \oint_L \left(\mathrm{e}^y\sin x\,\mathrm{d}x+\mathrm{e}^{-x}\sin y\,\mathrm{d}y\right)
    &=\iint_D\left(
    \frac{\partial}{\partial x}\left(\mathrm{e}^{-x}\sin y\right)
    -\frac{\partial}{\partial y}\left(\mathrm{e}^y\sin x\right)
    \right)\mathrm{d}x\,\mathrm{d}y \\
    &=\iint_D\left(-\mathrm{e}^{-x}\sin y-\mathrm{e}^y\sin x\right)\mathrm{d}x\,\mathrm{d}y \\
    &=-(\mathrm{e}^{-a}-\mathrm{e}^{-b})(\cos c-\cos d)
    +(\cos b-\cos a)(\mathrm{e}^d-\mathrm{e}^c).
    \end{align*}
    \item 令 $M=-QR,\ N=PR$，则 \(\displaystyle \oint_C PR\,\mathrm{d}y-QR\,\mathrm{d}x
    =\oint_C M\,\mathrm{d}x+N\,\mathrm{d}y\)。由 Green 公式，
    \begin{align*}
    \oint_C M\,\mathrm{d}x+N\,\mathrm{d}y
    &=\iint_D\left(\frac{\partial N}{\partial x}-\frac{\partial M}{\partial y}\right)\mathrm{d}x\,\mathrm{d}y \\
    &=\iint_D\left(
    P\frac{\partial R}{\partial x}+R\frac{\partial P}{\partial x}
    +Q\frac{\partial R}{\partial y}+R\frac{\partial Q}{\partial y}
    \right)\mathrm{d}x\,\mathrm{d}y \\
    &=\iint_D\left(P\frac{\partial R}{\partial x}+Q\frac{\partial R}{\partial y}\right)\mathrm{d}x\,\mathrm{d}y
    +\iint_D R\left(\frac{\partial P}{\partial x}+\frac{\partial Q}{\partial y}\right)\mathrm{d}x\,\mathrm{d}y.
    \end{align*}
    移项即得所证公式。
    \item \textbf{思路：}题目先问“与路径无关”，入口就不是挑一条路径硬算，而是先做保守场判别：令交叉偏导相等求出 $\varphi$，再把积分降成势函数的端点差。

    路径无关条件为
    \[
    \frac{\partial}{\partial y}\!\left(\left(\sin x-\varphi(x)\right)\frac{y}{x}\right)
    =\frac{\partial \varphi(x)}{\partial x},\quad x\varphi'(x)+\varphi(x)=\sin x.
    \]
    解得 \(\displaystyle \varphi(x)=\frac{C-\cos x}{x}\)。由 $\varphi(\pi)=1$ 得 $C=\pi-1$，故 \(\displaystyle \boxed{\varphi(x)=\frac{\pi-1-\cos x}{x}}\)。
    取势函数 $u(x,y)=\varphi(x)y$，则 \(\displaystyle \int_A^B \left(\sin x-\varphi(x)\right)\frac{y}{x}\,\mathrm{d}x+\varphi(x)\,\mathrm{d}y=u(B)-u(A)=\pi\)。
    \item 设 $u(x,y)=\ln(x+y)-\dfrac{y}{x+y}$，则 \(\displaystyle \mathrm{d}u=\frac{(x+2y)\,\mathrm{d}x+y\,\mathrm{d}y}{(x+y)^2}\)。
\end{compactenum}
\end{solution}

\begin{exercise}[若干重要定理的证明]
\begin{compactenum}
    \item 证明：若 $\vec a$ 为固定单位向量，$C$ 为平面闭曲线，则 \(\displaystyle \oint_C \cos(\vec a,\vec n)\,\mathrm{d}s=0\)；
    \item 证明：若 $\vec a$ 为固定单位向量，$\Sigma$ 为闭曲面，则 \(\displaystyle \oiint_\Sigma \cos(\vec a,\vec n)\,\mathrm{d}S=0\)；
    \item 证明：\(\displaystyle \oiint_S \frac{\partial u}{\partial n}\,\mathrm{d}S=\iiint_\Omega \Delta u\,\mathrm{d}V\)；
    \item 证明 Green 第二公式：\(\displaystyle \int_\Omega (u\Delta v-v\Delta u)\,\mathrm{d}V
    =\int_{\partial\Omega}\left(u\frac{\partial v}{\partial n}-v\frac{\partial u}{\partial n}\right)\mathrm{d}S\)；
    \item 计算圆锥侧面 \(z=\sqrt{x^2+y^2},\ 0\le z\le 1\) 上的第二类曲面积分
    \(\displaystyle \iint_S x\,\mathrm{d}y\,\mathrm{d}z+y\,\mathrm{d}z\,\mathrm{d}x+(z^2-2z)\,\mathrm{d}x\,\mathrm{d}y\)；
    \item 计算旋转曲面 \(z=\sqrt{y-1},\ 1\le y\le 3\) 围成的旋转曲面上的第二类曲面积分
    \[
    \iint_S\left\{(8y+1)x\,\mathrm{d}y\,\mathrm{d}z+2(1-y^2)\,\mathrm{d}z\,\mathrm{d}x-4yz\,\mathrm{d}x\,\mathrm{d}y\right\}.
    \]
\end{compactenum}
\end{exercise}
\begin{solution}
\begin{compactenum}
    \item 取常向量场 $\vec F=\vec a$，由 Green 公式得 \(\displaystyle \oint_C \cos(\vec a,\vec n)\,\mathrm{d}s=\oint_C \vec a\cdot\vec n\,\mathrm{d}s=\iint_D \operatorname{div}\vec a\,\mathrm{d}x\,\mathrm{d}y=0\)。
    \item 取常向量场 $\vec F=\vec a$，由 Gauss 公式得 \(\displaystyle \oiint_\Sigma \cos(\vec a,\vec n)\,\mathrm{d}S=\iiint_\Omega \operatorname{div}\vec a\,\mathrm{d}V=0\)。
    \item 对向量场 $\nabla u$ 用 Gauss 公式即可：因为 \(\displaystyle \frac{\partial u}{\partial n}=\nabla u\cdot \vec n\)，且 \(\displaystyle \operatorname{div}(\nabla u)=\Delta u\)，故 \(\displaystyle \oiint_S \frac{\partial u}{\partial n}\,\mathrm{d}S=\iiint_\Omega \Delta u\,\mathrm{d}V\)。
    \item 分别对 $u\nabla v$ 与 $v\nabla u$ 用 Gauss 公式。先有 \(\displaystyle \operatorname{div}(u\nabla v)=u\Delta v+\nabla u\cdot \nabla v\)，故 \(\displaystyle
    \oiint_{\partial\Omega}u\frac{\partial v}{\partial n}\,\mathrm{d}S
    =\iiint_\Omega \left(u\Delta v+\nabla u\cdot \nabla v\right)\mathrm{d}V\)。
    再对 $v\nabla u$ 用同样的散度展开与高斯公式，并将两式相减，梯度内积项相消，得
    \begin{align*}
    \oiint_{\partial\Omega}v\frac{\partial u}{\partial n}\,\mathrm{d}S
    &=\iiint_\Omega \left(v\Delta u+\nabla u\cdot \nabla v\right)\mathrm{d}V,\\
    \iiint_\Omega (u\Delta v-v\Delta u)\,\mathrm{d}V
    &=\oiint_{\partial\Omega}\left(u\frac{\partial v}{\partial n}-v\frac{\partial u}{\partial n}\right)\mathrm{d}S.
    \end{align*}
    \item 记 $\vec F=(x,\ y,\ z^2-2z)$。
    将侧面与顶面 $z=1$ 补成闭曲面 $\partial\Omega$，其中
    \(\displaystyle \Omega=\{(x,y,z)\mid x^2+y^2\le z^2,\ 0\le z\le 1\}\)。
    由 Gauss 公式，\(\displaystyle
    \oiint_{\partial\Omega}\vec F\cdot \vec n\,\mathrm{d}S
    =\iiint_\Omega \operatorname{div}\vec F\,\mathrm{d}V\)，而
    \(\displaystyle \operatorname{div}\vec F
    =\frac{\partial x}{\partial x}+\frac{\partial y}{\partial y}
    +\frac{\partial(z^2-2z)}{\partial z}
    =2z\)。
    用柱坐标计算体积分：\(\displaystyle \iiint_\Omega 2z\,\mathrm{d}V=\int_0^1\int_0^{2\pi}\int_0^z 2z\,r\,\mathrm{d}r\,\mathrm{d}\theta\,\mathrm{d}z=2\pi\int_0^1 z^3\,\mathrm{d}z =\frac{\pi}{2}\)。
    在顶面 $z=1$ 上，外法向量为 $\vec k$，故顶面通量为
    \(\displaystyle \iint_{x^2+y^2\le 1}(1-2)\,\mathrm{d}x\,\mathrm{d}y=-\pi\)。
    于是侧面通量 \(\displaystyle
    \iint_S \vec F\cdot \vec n\,\mathrm{d}S
    =\frac{\pi}{2}-(-\pi)=\frac{3\pi}{2}\)。
    故 $\boxed{\iint_S x\,\mathrm{d}y\,\mathrm{d}z+y\,\mathrm{d}z\,\mathrm{d}x+(z^2-2z)\,\mathrm{d}x\,\mathrm{d}y=\frac{3\pi}{2}}$。
    \item 旋转后所得曲面为抛物面侧面 $y=x^2+z^2+1,\ 1\le y\le 3$。记
    \(\displaystyle \vec F=\bigl((8y+1)x,\ 2(1-y^2),\ -4yz\bigr)\)。
    将侧面与上端圆盘 $y=3$ 补成闭曲面。先算散度：
    \(\displaystyle \operatorname{div}\vec F=(8y+1)-4y-4y=1\)。
    因而闭曲面的总通量等于所围体积。该体积为 \(\displaystyle V=\int_{1}^{3}\pi(y-1)\,\mathrm{d}y=\pi\left[\frac{(y-1)^2}{2}\right]_{1}^{3}=2\pi\)。
    再算上端圆盘 $y=3$ 的通量。此时外法向量为 $\vec j$，故
    \(\displaystyle
    \iint_{x^2+z^2\le 2}2(1-y^2)\,\mathrm{d}x\,\mathrm{d}z
    =2(1-9)\cdot \pi(\sqrt2)^2
    =-32\pi\)。
    底端 $y=1$ 退化为一点，没有面积贡献，所以侧面通量为 $2\pi-(-32\pi)=34\pi$，即 $\boxed{34\pi}$。
\end{compactenum}
\end{solution}

\begin{exercise}[Stokes 公式与通量计算]
\begin{shortsingleenum}
    \shortsingleitem{1}{计算 \(\displaystyle \oint_C (z-y)\,\mathrm{d}x+(x-z)\,\mathrm{d}y+(x-y)\,\mathrm{d}z\)，其中 $C$ 为平面 \(x^2+y^2=1,\ x-y+z=2\) 的交线，方向按从 $+z$ 轴正向看去为顺时针；}
    \shortsingleitem{2}{计算闭曲面 $\Sigma$ 上的通量 \(\displaystyle \oiint_\Sigma yz\,\mathrm{d}x\,\mathrm{d}y+xz\,\mathrm{d}y\,\mathrm{d}z+xy\,\mathrm{d}z\,\mathrm{d}x\)，其中 $\Sigma$ 由圆柱面 $x^2+y^2=R^2$ 的第一卦限侧面、平面 $x+2y=R$、平面 $z=0$、$z=H$ 与 $x=0$ 围成。}
\end{shortsingleenum}
\end{exercise}
\begin{solution}
\begin{compactenum}
    \item \textbf{思路：}这是空间闭曲线上的第二类曲线积分，先看 Stokes；而边界已经是一个平面截线，所以直接换成同边界的平面椭圆盘最省。

    记 $\vec F=(z-y,\ x-z,\ x-y)$，则 $\operatorname{rot}\vec F=(0,0,2)$。
    由于旋度只剩 $z$ 分量，真正要算的就是该椭圆盘在 $xOy$ 平面上的有向投影面积。取平面 $x-y+z=2$ 内由交线围成的椭圆盘为积分曲面 $S$。由题给方向“从 $+z$ 轴正向看去为顺时针”，按右手法则可知 $S$ 的法向应取 $z$ 分量为负的一侧。
    由于 $\operatorname{rot}\vec F$ 只有 $z$ 分量，Stokes 公式化为
    \[
    \oint_C \vec F\cdot \mathrm{d}\vec r
    =\iint_S (\operatorname{rot}\vec F)\cdot \vec n\,\mathrm{d}S
    =2\iint_S n_z\,\mathrm{d}S.
    \]
    而 $\iint_S n_z\,\mathrm{d}S$ 正是 $S$ 在 $xOy$ 平面上的有向投影面积。交线投影满足 \(x^2+y^2=1\)，即单位圆盘，面积为 $\pi$；由于取负 $z$ 方向，投影面积应记为 $-\pi$。因此 \(\displaystyle \oint_C (z-y)\,\mathrm{d}x+(x-z)\,\mathrm{d}y+(x-y)\,\mathrm{d}z=2(-\pi)=-2\pi\)。
    $\boxed{-2\pi}$。
    \item \textbf{思路：}这一题对象本来就是闭曲面通量，优先直接用 Gauss；这样可以避免把圆柱侧面、平面侧面和上下底分别拆开。

    记 $\vec F=(xz,xy,yz)$，则 $\operatorname{div}\vec F=x+y+z$。
    对象既然已经封闭，后面就只需把通量改写为所围区域上的三重积分。设 $\Omega$ 为题中闭曲面所围成的立体，则由 Gauss 公式 \(\displaystyle I=\iiint_\Omega (x+y+z)\,\mathrm{d}V\)。
    由于 $0\le z\le H$，可分离为 \(\displaystyle
    I=H\iint_D (x+y)\,\mathrm{d}A+\frac{H^2}{2}\iint_D \mathrm{d}A\)，
    其中 $D$ 是 $xOy$ 平面上的底面区域。该区域可看作第一象限扇形 $x^2+y^2\le R^2,\ x\ge 0,\ y\ge 0$ 去掉直线 $x+2y=R$ 以下的三角形 $T=\{(x,y)\mid x\ge 0,\ y\ge 0,\ x+2y\le R\}$。
    因而 \(\displaystyle \operatorname{Area}(D)=\frac{\pi R^2}{4}-\frac{R^2}{4}
    =\frac{R^2(\pi-1)}{4}\)。
    再算 $\iint_D (x+y)\,\mathrm{d}A$。在扇形部分，\(\displaystyle \iint_{\text{扇形}}(x+y)\,\mathrm{d}A=\frac{2R^3}{3}\)；
    在三角形 $T$ 上，利用三角形形心 $\left(\dfrac{R}{3},\dfrac{R}{6}\right)$ 与面积 $\operatorname{Area}(T)=\dfrac{R^2}{4}$，可得 \(\displaystyle \iint_T (x+y)\,\mathrm{d}A=\frac{R^2}{4}\left(\frac{R}{3}+\frac{R}{6}\right)=\frac{R^3}{8}\)。
    所以 \(\displaystyle
    \iint_D (x+y)\,\mathrm{d}A
    =\frac{2R^3}{3}-\frac{R^3}{8}
    =\frac{13R^3}{24}\)。
    代回即得 \(\displaystyle I=H\cdot \frac{13R^3}{24}
    +\frac{H^2}{2}\cdot \frac{R^2(\pi-1)}{4}
    =\boxed{\frac{HR^2}{24}\bigl(13R+3H(\pi-1)\bigr)}\)。
\end{compactenum}
\end{solution}

\chapter{级数}

\section{知识讲解}

本章不是把若干种级数并排陈列，而是沿着“对象升级”的路线展开：先研究数项级数，解决无穷多个数能不能相加；再研究函数项级数，解决每个 \(x\) 处能相加以后，能否对整个区间统一控制；随后进入幂级数和 Taylor 级数，用幂函数系表示函数；最后进入 Fourier 级数，用三角函数系表示周期函数。

其中数项级数判敛最容易显得方法繁多，本章把它整理成一条判题链：
\[
\boxed{
\text{通项检查}\longrightarrow
\text{转正}\longrightarrow
\text{估阶归模型}\longrightarrow
\text{极限形式比较}\longrightarrow
\text{特殊结构辅助}\longrightarrow
\text{条件收敛}
}.
\]
这里的 Cauchy 准则是理论上最根本的收敛准则，但普通计算题通常不从它入手。考试中更常见、也最稳定的写法，是先把正项通项 \(a_n\) 估阶为某个标准模型 \(b_n\)，再证明
\[
\lim_{n\to\infty}\frac{a_n}{b_n}=1,
\]
然后使用比较判别法的极限形式说明 \(\sum a_n\) 与 \(\sum b_n\) 同敛散。

\subsection{数项级数判敛总流程：先排除，再转正，再估阶}

\subsubsection{第一关：通项是否趋于零}

判断级数前，先把对象和最基本的必要条件说清楚。部分和给出级数的定义，通项趋零给出第一道排除关口。

\begin{definition}[无穷级数与部分和]
设数列 $\{a_n\}$ 给定，称 \(\displaystyle \sum_{n=1}^{\infty} a_n=a_1+a_2+\cdots+a_n+\cdots\) 为无穷级数，其前 $n$ 项和 \(\displaystyle S_n=\sum_{k=1}^{n}a_k\) 称为部分和。若数列 $\{S_n\}$ 有有限极限 $S$，则称级数收敛，且记 \(\displaystyle \sum_{n=1}^{\infty}a_n=S\)；否则称该级数发散。
\end{definition}

\begin{theorem}[收敛级数通项趋于零]
若级数 $\sum_{n=1}^{\infty}a_n$ 收敛，则 \(\displaystyle \lim_{n\to\infty}a_n=0\)。
但其逆命题一般不成立。
\end{theorem}

\begin{example}[通项不趋于零时的发散]
级数 \(\displaystyle \sum_{n=0}^{\infty}(-1)^n\) 发散，因为其通项不趋于零。
\end{example}

\begin{theorem}[有限项不影响敛散性]
在级数前去掉有限项或添上有限项，不改变级数的收敛性或发散性。
\end{theorem}

第一关只给必要条件：若 \(a_n\not\to0\)，级数一定发散；若 \(a_n\to0\)，只是说明还没有被排除，后面仍要继续判断。

\subsubsection{第二关：先看绝对值，把一般级数转成正项级数}

带正负号的级数不要一开始就盯着符号变化。先看绝对值级数 \(\sum |a_n|\)，因为它是正项级数，能直接进入估阶和比较的流程；若绝对收敛已经成立，原级数自动收敛。

\begin{definition}[正项级数]
若对一切 $n$ 都有 $a_n\ge 0$，则 \(\displaystyle \sum_{n=1}^{\infty}a_n\) 称为正项级数。只有有限多个负项的级数，也可按正项级数的思想讨论其敛散性。
\end{definition}

\begin{definition}[绝对收敛]
若 \(\displaystyle \sum_{n=1}^{\infty}|a_n|\) 收敛，则称级数 $\sum a_n$ 绝对收敛。
\end{definition}

\begin{theorem}[绝对收敛推出收敛]
若 \(\displaystyle \sum_{n=1}^{\infty}|a_n|\) 收敛，则 \(\displaystyle \sum_{n=1}^{\infty}a_n\) 必收敛。
\end{theorem}

\begin{theorem}[正项级数收敛充要条件]
正项级数 $\sum a_n$ 收敛，当且仅当其部分和数列 \(\displaystyle S_n=\sum_{k=1}^{n}a_k\) 有界。
\end{theorem}

\subsubsection{第三关：估阶归模型}

正项级数进入真正计算时，先问通项最终像谁。教材和考试中最常见的目标模型是几何级数、$p$ 级数和对数边界型级数；复杂表达式通常只是换了一层外衣。

\begin{conclusion}[估阶判敛的标准落笔]
对正项级数 \(\sum a_n\)，常规判敛的落笔顺序是：
\[
\boxed{\text{找标准模型 }b_n\quad\longrightarrow\quad
\lim_{n\to\infty}\frac{a_n}{b_n}=1\quad\longrightarrow\quad
\text{由比较判别法的极限形式得 }\sum a_n\text{ 与 }\sum b_n\text{ 同敛散}.}
\]
这里 \(b_n\) 通常取 \(\displaystyle \frac{1}{n^p}\)、\(\displaystyle \frac{1}{n(\ln n)^q}\)、\(q^n\)、\(\displaystyle \frac{C}{n^p}\) 等已知模型。若极限不是 \(1\)，只要是有限正数 \(0<l<+\infty\)，同样可以用比较判别法的极限形式；写成 \(1\) 只是最清楚的等价估阶形式。
\end{conclusion}

\begin{example}[几何级数]
几何级数 \(\displaystyle \sum_{n=0}^{\infty}q^n\) 在 $|q|<1$ 时收敛且和为 \(\displaystyle \frac{1}{1-q}\)，在 $|q|\ge 1$ 时发散。
特别地，当 $q=1$ 时部分和趋于 $+\infty$；当 $q=-1$ 时部分和在 $0,1$ 间振荡，因此发散。
\end{example}

\begin{example}[Taylor 展开导出指数级数]
由 Taylor 公式 \(\displaystyle e^x=1+x+\frac{x^2}{2!}+\cdots+\frac{x^n}{n!}+\frac{x^{n+1}}{(n+1)!}e^{\theta x},\ 0<\theta<1\)，对固定 $x$ 有 \(\displaystyle \lim_{n\to\infty}\frac{x^{n+1}}{(n+1)!}e^{\theta x}=0\)，故 \(\displaystyle e^x=\sum_{n=0}^{\infty}\frac{x^n}{n!}\)。
\end{example}

这里先把指数级数当作一个阶乘型标准模型来用；至于为什么 Taylor 展开能够代表原函数，后面的 Taylor 专题会用余项趋零来说明。

\begin{conclusion}[常见估阶模型库]
设 \(a_n\ge 0\)。考场上常把 \(a_n\) 化到下面几类模型中，再用
\(\displaystyle \lim a_n/b_n=1\) 的比较判别法极限形式收尾。
\begin{compactenum}
\item 若取 \(\displaystyle b_n=\frac{C}{n^p}\ (C>0)\) 且 \(\displaystyle \lim_{n\to\infty}a_n/b_n=1\)，则 \(p>1\) 时收敛，\(p\le1\) 时发散。
\item 若取 \(\displaystyle b_n=\frac{C}{n(\ln n)^q}\ (C>0)\) 且 \(\displaystyle \lim_{n\to\infty}a_n/b_n=1\)，则 \(q>1\) 时收敛，\(q\le1\) 时发散。
\item 若 \(\displaystyle a_n=O(q^n)\ (0<q<1)\)，则 \(\sum a_n\) 收敛；指数衰减通常比任何多项式增长都强。
\item 若只差一个对数因子，常用 \(\ln n=o(n^\varepsilon)\)。例如充分大时 \(\ln n\le n^{1/4}\)，于是 \(\displaystyle \frac{\ln n}{n^{3/2}}\le \frac1{n^{5/4}}\)。
\item 根式差先有理化：
\[
\sqrt{n^2+1}-n=\frac1{\sqrt{n^2+1}+n}\sim \frac1{2n}.
\]
\item 小量函数先换成 \(x\to0\) 的等价无穷小：
\[
\sin x\sim x,\qquad
1-\cos x\sim \frac{x^2}{2},\qquad
\ln(1+x)\sim x,\qquad
e^x-1\sim x.
\]
\item 阶乘和二项式系数需要时可用 Stirling 型估阶，例如
\[
\binom{2n}{n}\sim \frac{4^n}{\sqrt{\pi n}}.
\]
\end{compactenum}
\end{conclusion}

\begin{conclusion}[标准模型的考试证明依据]
用估阶比较时，不能只写“与某标准级数比较”，还要知道标准级数本身为什么收敛或发散。下面这些写法可直接作为考试中的依据。
\begin{compactenum}
\item 几何级数：若 \(|q|<1\)，则
\[
\sum_{n=0}^{N}q^n=\frac{1-q^{N+1}}{1-q}\longrightarrow \frac1{1-q},
\]
所以 \(\sum q^n\) 收敛；若 \(|q|\ge1\)，通项 \(q^n\) 不趋于 \(0\) 或部分和振荡，故发散。

\item 调和级数：按块分组有
\[
1+\frac12+\left(\frac13+\frac14\right)
+\left(\frac15+\cdots+\frac18\right)+\cdots .
\]
从第 \(k\) 块开始，每块有 \(2^{k-1}\) 项，且每项不小于 \(1/2^k\)，所以每块和至少为 \(1/2\)。部分和无界，故 \(\sum 1/n\) 发散。

\item \(p\)-级数：当 \(p>1\) 时，把 \(n=2^k,\ldots,2^{k+1}-1\) 分成一块，则
\[
\sum_{n=2^k}^{2^{k+1}-1}\frac1{n^p}
\le \frac{2^k}{(2^k)^p}
=2^{k(1-p)}.
\]
右边是收敛的几何级数，所以 \(\sum 1/n^p\) 收敛。当 \(0<p\le1\) 时，有 \(1/n^p\ge1/n\)，故由调和级数发散知其发散；当 \(p\le0\) 时，通项不趋于 \(0\)，也发散。

\item 对数边界型：讨论
\(\displaystyle \sum_{n=3}^{\infty}\frac1{n(\ln n)^q}\)。仍按 \(n=2^k,\ldots,2^{k+1}-1\) 分块。若 \(q>1\)，则一块和满足
\[
\sum_{n=2^k}^{2^{k+1}-1}\frac1{n(\ln n)^q}
\le
\frac{2^k}{2^k(k\ln2)^q}
=\frac{C}{k^q},
\]
而 \(\sum 1/k^q\) 收敛，故原级数收敛。若 \(0<q\le1\)，则
\[
\sum_{n=2^k}^{2^{k+1}-1}\frac1{n(\ln n)^q}
\ge
\frac{2^k}{2^{k+1}((k+1)\ln2)^q}
=\frac{C}{(k+1)^q},
\]
对应块和发散，故原级数发散。若 \(q\le0\)，则充分大时 \((\ln n)^q\le1\)，于是
\[
\frac1{n(\ln n)^q}\ge \frac1n,
\]
由调和级数发散，原级数也发散。

\item 指数衰减：若 \(a_n=O(q^n)\) 且 \(0<q<1\)，则存在 \(M>0\)，使充分大时 \(0\le a_n\le Mq^n\)。由于 \(\sum q^n\) 收敛，比较判别法给出 \(\sum a_n\) 收敛。
\end{compactenum}
\end{conclusion}

\subsubsection{第四关：比较判别法的极限形式为主，普通比较为辅}

估阶完成后，真正落笔通常优先用比较判别法的极限形式。也就是说，先选标准模型 \(b_n\)，再算
\[
\lim_{n\to\infty}\frac{a_n}{b_n}.
\]
若极限为 \(1\)，就说明 \(a_n\) 的阶正是 \(b_n\)，二者对应级数同敛散。只有当题目天然给出明显不等式时，才直接用普通比较判别法。

\begin{theorem}[比较判别法的极限形式]
若 \(\displaystyle \lim_{n\to\infty}\frac{a_n}{b_n}=l\)，
则有：
\begin{itemize}
  \item 当 $0<l<\infty$ 时，$\sum a_n$ 与 $\sum b_n$ 同敛散；
  \item 当 $l=0$ 且 $\sum b_n$ 收敛时，$\sum a_n$ 收敛；
  \item 当 $l=\infty$ 且 $\sum b_n$ 发散时，$\sum a_n$ 发散。
\end{itemize}
\end{theorem}

\begin{theorem}[比较判别法]
设 $\sum a_n,\sum b_n$ 为正项级数。若从某项起有 $a_n\le b_n$，则 \(\displaystyle \sum b_n\ \text{收敛}\Rightarrow \sum a_n\ \text{收敛},\quad \sum a_n\ \text{发散}\Rightarrow \sum b_n\ \text{发散}\)。
\end{theorem}

\begin{example}[比较判别法]
讨论 \(\displaystyle \sum_{n=1}^{\infty}\frac{3+\sin^3(n+1)}{2^n+n}\)。
由 \(\displaystyle 0<\frac{3+\sin^3(n+1)}{2^n+n}<\frac4{2^n}\)，
而几何级数 $\sum \frac4{2^n}$ 收敛，所以原级数收敛。
\end{example}

\begin{example}[比较判别法的极限形式]
讨论 \(\displaystyle \sum_{n=1}^{\infty}\ln\left(1+\frac1{n^2}\right)\)。
取 \(\displaystyle b_n=\frac1{n^2}\)。因为
\[
\lim_{n\to\infty}\frac{\ln\left(1+\frac1{n^2}\right)}{\frac1{n^2}}
=\lim_{t\to0^+}\frac{\ln(1+t)}{t}=1,
\]
所以原通项的阶正是 \(b_n\)。又 \(\displaystyle \sum\frac1{n^2}\) 收敛，故由比较判别法的极限形式知原级数收敛。
\end{example}

\begin{example}[与调和级数比较]
讨论 \(\displaystyle \sum_{n=1}^{\infty}\sin\frac1n\)。
取 \(\displaystyle b_n=\frac1n\)。因为
\[
\lim_{n\to\infty}\frac{\sin(1/n)}{1/n}=1,
\]
而调和级数 \(\displaystyle \sum\frac1n\) 发散，故由比较判别法的极限形式知该级数发散。
\end{example}

\begin{example}[再证 $\sum \ln(1+\frac1n)$ 发散]
取 \(\displaystyle b_n=\frac1n\)。因为
\[
\lim_{n\to\infty}\frac{\ln\left(1+\frac1n\right)}{\frac1n}
=\lim_{t\to0^+}\frac{\ln(1+t)}{t}=1,
\]
而调和级数 \(\displaystyle \sum\frac1n\) 发散，所以由比较判别法的极限形式知
\(\displaystyle \sum_{n=1}^{\infty}\ln\left(1+\frac1n\right)\) 发散。
\end{example}

\subsubsection{第五关：比值、根值与积分判别作为估阶辅助}

比值判别、根值判别和积分判别不必抢在比较法前面。它们更像识别量级的工具：积分判别帮助确定 \(p\) 级数、对数边界型级数；比值和根值适合阶乘、指数、连乘、\(n\) 次幂放在 \(n\) 次方里的结构。

\begin{theorem}[积分判别法]
若函数 $f(x)$ 在 $[N,+\infty)$ 上非负且单调递减，并令 $a_n=f(n)$，则级数
\(\displaystyle \sum_{n=N}^{\infty}a_n\) 与反常积分 \(\displaystyle \int_{N}^{+\infty}f(x)\,\mathrm{d}x\) 同敛散。
\end{theorem}

\begin{example}[$\sum \frac1{n(\ln n)^q}$ 的敛散性]
对 $q>0$，讨论 \(\displaystyle \sum_{n=3}^{\infty}\frac1{n(\ln n)^q}\)。
由积分 \(\displaystyle \int_{3}^{+\infty}\frac{\mathrm{d}x}{x(\ln x)^q}\)
可知：
\begin{itemize}
  \item 当 $q>1$ 时级数收敛；
  \item 当 $q\le 1$ 时级数发散。
\end{itemize}
\end{example}

\begin{example}[$p$ 级数]
对 $\alpha>0$，级数 \(\displaystyle \sum_{n=1}^{\infty}\frac1{n^\alpha}\)
满足：
\begin{itemize}
  \item 当 $\alpha>1$ 时收敛；
  \item 当 $\alpha\le 1$ 时发散。
\end{itemize}
\end{example}

\begin{theorem}[D'Alembert 判别法]
若正项级数从某项起满足
\(\displaystyle \frac{a_{n+1}}{a_n}\le q<1\)，则收敛；若 \(\displaystyle \frac{a_{n+1}}{a_n}\ge 1\)，则发散。
\end{theorem}

\begin{theorem}[D'Alembert 判别法的极限形式]
若 \(\displaystyle \lim_{n\to\infty}\frac{a_{n+1}}{a_n}=q\)，则 \(q<1\Rightarrow \sum a_n\) 收敛，\(q>1\Rightarrow \sum a_n\) 发散。
当 $q=1$ 时，此法不能判定。
\end{theorem}

\begin{theorem}[Cauchy 根值判别法]
若从某项起满足
\(\displaystyle \sqrt[n]{a_n}\le q<1\)，则收敛；若 \(\displaystyle \sqrt[n]{a_n}\ge 1\) 对无穷多个 $n$ 成立，则级数发散。
\end{theorem}

\begin{theorem}[根值判别法的极限形式]
若 \(\displaystyle \lim_{n\to\infty}\sqrt[n]{a_n}=q\)，则 \(q<1\Rightarrow \sum a_n\) 收敛，\(q>1\Rightarrow \sum a_n\) 发散。
当 $q=1$ 时不能判定。
\end{theorem}

下面两个判别法主要处理比值法落在边界 \(q=1\) 时的少数题型，属于补充工具。常规判敛题仍优先回到通项估阶、标准模型和比较判别。

\begin{theorem}[Raabe 判别法]
若 \(\displaystyle \frac{a_n}{a_{n+1}}=1+\frac{l}{n}+o\left(\frac1n\right)\ (n\to\infty)\)，则 \(l>1\Rightarrow \sum a_n\) 收敛，\(l<1\Rightarrow \sum a_n\) 发散。
\end{theorem}

\begin{theorem}[Gauss 判别法]
若
\(\displaystyle \frac{a_n}{a_{n+1}}
=1+\frac1n+\frac{\beta}{n\ln n}+o\left(\frac1{n\ln n}\right)\ (n\to\infty)\)，
则 \(\beta>1\Rightarrow \sum a_n\) 收敛，\(\beta<1\Rightarrow \sum a_n\) 发散。
\end{theorem}

\begin{example}[比值判别法]
级数 \(\displaystyle \sum_{n=1}^{\infty}\frac1{n!}\) 中，若 $a_n=\frac1{n!}$，则 \(\displaystyle \frac{a_{n+1}}{a_n}=\frac1{n+1}\to 0<1\)，
故该级数收敛。
\end{example}

\begin{example}[含参数级数]
设 $a>0$，讨论
\(\displaystyle \sum_{n=1}^{\infty}\frac{a^n}{n!}\) 与 \(\displaystyle \sum_{n=1}^{\infty}na^n\)。结论为：
\begin{itemize}
  \item $\sum \dfrac{a^n}{n!}$ 对任意 $a>0$ 都收敛；
  \item $\sum na^n$ 在 $0\le a<1$ 时收敛，在 $a\ge 1$ 时发散。
\end{itemize}
\end{example}

\begin{example}[幂函数型参数级数]
设 $x>0$，讨论
\(\displaystyle \sum_{n=1}^{\infty}\frac{x^n}{n}\) 与 \(\displaystyle \sum_{n=1}^{\infty}\frac{x^n}{n^2}\)。则：
\begin{itemize}
  \item $\sum \frac{x^n}{n}$ 在 $0<x<1$ 时收敛，在 $x\ge 1$ 时发散；
  \item $\sum \frac{x^n}{n^2}$ 在 $0<x\le 1$ 时收敛，在 $x>1$ 时发散。
\end{itemize}
\end{example}

这类题在数项级数部分把参数看成固定常数；若进一步把 \(x\) 看成变量，研究所有 \(x\) 的收敛范围和和函数，就自然进入后面的幂级数专题。

\begin{example}[根值判别法]
讨论 \(\displaystyle \sum_{n=1}^{\infty}\frac1{(\ln n)^n}\)。
因 \(\displaystyle \sqrt[n]{a_n}=\frac1{\ln n}\to 0\)，
故原级数收敛。
\end{example}

\begin{example}[再看根值判别法]
设 $x\ge 0$，讨论 \(\displaystyle \sum_{n=1}^{\infty}2^n x^{2n}\)。
由 \(\displaystyle \sqrt[n]{a_n}=2x^2\)
得：
\begin{itemize}
  \item $0\le x<\frac1{\sqrt2}$ 时收敛；
  \item $x>\frac1{\sqrt2}$ 时发散；
  \item $x=\frac1{\sqrt2}$ 时级数化为 $\sum 1$，发散。
\end{itemize}
\end{example}

\begin{example}[欧拉常数的出现]
讨论 \(\displaystyle \sum_{n=1}^{\infty}\left(\frac1n-\ln\left(1+\frac1n\right)\right)\)。
由不等式 \(\displaystyle \frac1{n+1}<\ln\left(1+\frac1n\right)<\frac1n\)，可得
\(\displaystyle 0<\frac1n-\ln\left(1+\frac1n\right)<\frac1{n(n+1)}\)。
故该级数收敛。并可定义欧拉常数 \(\displaystyle \gamma=\lim_{n\to\infty}\left(\sum_{k=1}^{n}\frac1k-\ln n\right)\)。
\end{example}

\subsubsection{第六关：绝对收敛失败后再看交错与条件收敛}

若 \(\sum |a_n|\) 已经收敛，原级数就是绝对收敛，不必再讨论交错。只有绝对收敛失败，而符号又呈现稳定交错、单调趋零等抵消结构时，才进入条件收敛的判断。

\begin{definition}[交错级数]
形如 \(\displaystyle \sum_{n=1}^{\infty}(-1)^{n-1}a_n
=a_1-a_2+a_3-a_4+\cdots,\ a_n\ge 0\) 的级数称为交错级数。
\end{definition}

\begin{theorem}[Leibnitz 判别法]
若正数列 $\{a_n\}$ 单调递减且 $a_n\to 0$，则交错级数 \(\displaystyle \sum_{n=1}^{\infty}(-1)^{n-1}a_n\) 收敛。
\end{theorem}

\begin{theorem}[Leibnitz 型误差估计]
若交错级数 $\sum_{n=1}^{\infty}(-1)^{n-1}a_n$ 满足 Leibnitz 判别法条件，且其和为 $S$，则
\(\displaystyle |S-S_n|\le a_{n+1}\)。
\end{theorem}

\begin{example}
判定级数 \(\displaystyle \sum_{k=1}^{\infty}(-1)^{k-1}\frac{3^k}{k!}\)。
因为 $\frac{3^n}{n!}\to 0$，且当 $n$ 足够大时 $\frac{3^n}{n!}$ 单调递减，所以由 Leibnitz 判别法可知该级数收敛。
\end{example}

\begin{example}
讨论 \(\displaystyle \sum_{n=1}^{\infty}\frac{(-1)^{n-1}}{n^p},\ p>0\)。结论为：
\begin{itemize}
  \item $p>1$ 时绝对收敛；
  \item $0<p\le 1$ 时条件收敛。
\end{itemize}
\end{example}

\begin{definition}[条件收敛]
若 \(\displaystyle \sum_{n=1}^{\infty}a_n\) 收敛，而 \(\displaystyle \sum_{n=1}^{\infty}|a_n|\) 发散，则称级数 $\sum a_n$ 条件收敛。
\end{definition}

\begin{example}
下列级数都条件收敛：
\(\displaystyle \sum_{n=1}^{\infty}\frac{(-1)^{n-1}}{n}\) 与 \(\displaystyle \sum_{n=2}^{\infty}\frac{(-1)^n}{\ln n}\)。
它们都可由 Leibnitz 判别法证明收敛。再看绝对值级数：第一列的绝对值级数就是调和级数，发散；第二列中，当 \(n\) 足够大时 \(\ln n<n\)，所以
\[
\frac1{\ln n}>\frac1n,
\]
由普通比较判别法知 \(\sum 1/\ln n\) 发散。因此这两个级数都不绝对收敛，只能是条件收敛。
\end{example}

至此，数项级数的判敛流程已经闭合。接下来讨论求和、构造以及收敛后的合法运算；这些内容很重要，但不再作为普通判敛题的入口。

\subsection{数项级数的求和、构造与合法运算}

\subsubsection{裂项、几何拆分与构造}

判敛回答的是“能不能加”。这一小节不再作为判敛主流程的一部分，而是继续处理三个后续问题：已经知道可以加时，怎样算出和；怎样按要求构造一个级数；以及级数收敛后，哪些代数运算、分组和乘法是合法的。求和题仍然先识别结构：能裂项就写部分和，能拆成几何级数就拆，能转成已经熟悉的模型就转。

\begin{example}[裂项相消法]
设 $p\ge 0,\ q>0,\ s=p+q$，则
\(\displaystyle \frac{1}{(sn-p)(sn+q)}
=\frac{1}{s}\left(\frac{1}{sn-p}-\frac{1}{sn+q}\right)\)，从而
\(\displaystyle \sum_{n=1}^{\infty}\frac{1}{(sn-p)(sn+q)}=\frac{1}{s(s-p)}\)。
\end{example}

\begin{example}[求和 $\sum \frac1{n(n+1)}$]
由 \(\displaystyle \frac1{n(n+1)}=\frac1n-\frac1{n+1}\)，有
\(\displaystyle S_n=\sum_{k=1}^{n}\frac1{k(k+1)}=1-\frac1{n+1}\to 1\)。
故 $\sum_{n=1}^{\infty}\frac1{n(n+1)}=1$。
\end{example}

\begin{example}[对数裂项构造发散级数]
由 \(\displaystyle \ln\left(1+\frac1n\right)=\ln(n+1)-\ln n\)，有
\(\displaystyle S_n=\sum_{k=1}^{n}\ln\left(1+\frac1k\right)=\ln(n+1)\to+\infty\)。
故 $\sum_{n=1}^{\infty}\ln\left(1+\frac1n\right)$ 发散。
\end{example}

\begin{example}[求级数的和]
\[
\sum_{n=0}^{\infty}\frac{(-1)^n+2^n}{3^{n+2}}
=\frac19\sum_{n=0}^{\infty}\left(-\frac13\right)^n+\frac19\sum_{n=0}^{\infty}\left(\frac23\right)^n
=\frac1{12}+\frac13=\frac5{12}.
\]
\end{example}

\begin{example}[构造收敛级数]
设 $a_n>0$，$S_n=\sum_{k=1}^{n}a_k$。则 \(\displaystyle \sum_{n=1}^{\infty}\frac{a_n}{S_n^2}\) 收敛。事实上，对 $k\ge 2$，
\begin{align*}
\frac{a_k}{S_k^2}
&=\frac{S_k-S_{k-1}}{S_k^2}
\le \frac{S_k-S_{k-1}}{S_kS_{k-1}}
=\frac1{S_{k-1}}-\frac1{S_k},\\
\sum_{k=2}^{n}\frac{a_k}{S_k^2}
&\le \frac1{S_1}-\frac1{S_n}
<\frac1{a_1}.
\end{align*}
\end{example}

\begin{example}[不存在“收敛得最慢”的正项级数]
设 $\sum a_n$ 为收敛正项级数，记 \(\displaystyle \sum_{n=1}^{\infty}a_n=S,\qquad S_n=\sum_{k=1}^{n}a_k,\qquad \beta_n=S-S_{n-1}>0\)。
则 $\beta_n\downarrow 0$。令 \(\displaystyle b_n=\sqrt{\beta_n}-\sqrt{\beta_{n+1}}>0\)，则 $\sum b_n$ 收敛，且 \(\displaystyle \frac{a_n}{b_n}=\sqrt{\beta_n}+\sqrt{\beta_{n+1}}\to 0\)。
这表明总可以构造一个“更慢收敛”的正项级数，因此不存在绝对意义上的“收敛得最慢”的正项级数。
\end{example}

\subsubsection{收敛后的合法运算：正负部分、重排与 Cauchy 乘积}

下面这些内容说明级数已经收敛以后还能做什么、不能做什么。判敛时先用前面的流程；一旦进入求和、重排或级数相乘，就要额外检查绝对收敛等条件。

\begin{definition}[正负部分]
定义 \(\displaystyle a_n^+=\max\{a_n,0\},\qquad a_n^-=\max\{-a_n,0\}\)。
则 \(\displaystyle |a_n|=a_n^++a_n^-,\qquad a_n=a_n^+-a_n^-\)。
\end{definition}

\begin{theorem}[绝对收敛的等价刻画]
级数 $\sum a_n$ 绝对收敛，当且仅当 \(\displaystyle \sum a_n^+,\ \sum a_n^-\) 都收敛，且此时 \(\displaystyle \sum_{n=1}^{\infty}a_n=\sum_{n=1}^{\infty}a_n^+-\sum_{n=1}^{\infty}a_n^-\)。
\end{theorem}

\begin{theorem}[条件收敛时正负部分都发散到 $+\infty$]
若 $\sum a_n$ 条件收敛，则
\(\displaystyle \sum_{n=1}^{\infty}a_n^+=\sum_{n=1}^{\infty}a_n^-=+\infty\)。
\end{theorem}

\begin{theorem}[绝对收敛级数的重排不变性]
若级数 $\sum a_n$ 绝对收敛，则对其作任意重排后所得新级数仍绝对收敛，且和不变。
\end{theorem}

\begin{theorem}[Riemann 重排定理]
若级数 $\sum a_n$ 条件收敛，则适当重排其项后，可以使级数收敛到任意预先指定的实数，也可以使其发散。
\end{theorem}

\begin{definition}[Cauchy 乘积]
设
\(\displaystyle \sum_{n=1}^{\infty}a_n=A,\quad
\sum_{n=1}^{\infty}b_n=B\)，定义 \(\displaystyle c_n=a_1b_n+a_2b_{n-1}+\cdots+a_nb_1
=\sum_{k=1}^{n}a_kb_{n-k+1}\)。则级数 \(\displaystyle \sum_{n=1}^{\infty}c_n\) 称为级数 $\sum a_n$ 与 $\sum b_n$ 的 Cauchy 乘积。
\end{definition}

\begin{theorem}[绝对收敛级数的 Cauchy 乘积]
若 $\sum a_n$ 与 $\sum b_n$ 都绝对收敛，且和分别为 $A,B$，则其 Cauchy 乘积也绝对收敛，且
\(\displaystyle \sum_{n=1}^{\infty}c_n=AB\)。
\end{theorem}

\begin{example}[条件收敛重排改变和]
对交错调和级数 \(\displaystyle \sum_{n=1}^{\infty}\frac{(-1)^{n-1}}n=S\)，适当重排后，可以得到和为 \(\displaystyle \frac32 S\) 的新级数。这说明条件收敛级数重排后和会发生改变。
\end{example}

\begin{example}[一般收敛级数相乘未必收敛]
取
\[
a_n=b_n=\frac{(-1)^{n-1}}{\sqrt n}\qquad(n=1,2,\dots).
\]
由 Leibnitz 判别法，\(\sum a_n\) 与 \(\sum b_n\) 都收敛，但都不是绝对收敛。其 Cauchy 乘积通项为
\[
c_n=\sum_{k=1}^{n}a_kb_{n-k+1}
=(-1)^{n-1}\sum_{k=1}^{n}\frac1{\sqrt{k(n-k+1)}}.
\]
而
\[
\sum_{k=1}^{n}\frac1{\sqrt{k(n-k+1)}}
\ge \sum_{k=1}^{n}\frac{2}{n+1}
=\frac{2n}{n+1}.
\]
这里用到的是 \(k+(n-k+1)=n+1\)，所以
\[
k(n-k+1)\le \left(\frac{n+1}{2}\right)^2.
\]
于是 \(|c_n|\ge \dfrac{2n}{n+1}\)，从而 \(c_n\not\to0\)，Cauchy 乘积级数发散。这说明“两个收敛级数的 Cauchy 乘积仍收敛”在一般情形下不成立，必须加上绝对收敛等条件。
\end{example}

\subsubsection{理论补充：Cauchy 准则、线性与分组}

Cauchy 准则从尾和角度刻画收敛，适合证明理论结论，也适合处理少数需要直接估计尾和的题。但在常规判敛题中，仍优先走前面的估阶比较流程。

\begin{theorem}[Cauchy 准则]
级数 $\sum_{n=1}^{\infty}a_n$ 收敛的充要条件是：对任意 $\varepsilon>0$，存在正整数 $N$，使得当 $m>n>N$ 时，
\(\displaystyle \left|a_{n+1}+a_{n+2}+\cdots+a_m\right|<\varepsilon\)。
\end{theorem}

\begin{example}[用 Cauchy 准则证明 $\sum \frac1{n^2}$ 收敛]
对 $m>n$ 有估计
\(\displaystyle \sum_{k=n+1}^{m}\frac1{k^2}\le
\sum_{k=n+1}^{m}\frac{1}{k(k-1)}=\frac1n-\frac1m<\frac1n\)。
故当 $n$ 足够大时尾和可任意小，由 Cauchy 准则可知 $\sum_{n=1}^{\infty}\frac1{n^2}$ 收敛。
\end{example}

\begin{theorem}[线性性质]
若 $\sum a_n$ 与 $\sum b_n$ 都收敛，则对任意常数 $\alpha,\beta\in\mathbb{R}$，有 \(\displaystyle \sum_{n=1}^{\infty}(\alpha a_n+\beta b_n)=\alpha\sum_{n=1}^{\infty}a_n+\beta\sum_{n=1}^{\infty}b_n\)。
\end{theorem}

\begin{theorem}[分组的可结合性]
若收敛级数 $\sum a_n$ 在不改变项的先后顺序的前提下作任意分组，则所得新级数仍收敛，且和不变。
\end{theorem}

\begin{theorem}[分组结论的可逆条件]
若把级数 $\sum a_n$ 按原顺序分组后得到的新级数收敛，并且每一组内各项同号，则原级数也收敛，且和与分组后的级数相同。
\end{theorem}

\begin{remark}[补充说明]
分组不换序的结论不能反过来用。比如 \(\displaystyle 1-1+1-1+\cdots\) 两项一组后形式上可得 $0$，但原级数本身发散。
\end{remark}

至此，数项级数部分收束。后面函数项级数、幂级数、Taylor 级数中遇到端点或逐点问题时，最终仍会回到这里：先把固定的 \(x\) 代入，得到一个数项级数，再按前面的流程判断。

\subsection{函数项级数专题：逐点收敛、一致收敛与交换定理}

从这里开始，级数的每一项不再只是数，而是函数。固定一个 \(x\) 后，它仍然退回数项级数；但若要让极限、积分、求导在整个区间上可靠交换，就必须控制这种收敛对 \(x\) 是否统一。

\subsubsection{函数项级数与和函数}

\begin{definition}[函数项级数]
设 $u_1(x),u_2(x),\dots$ 定义在区间 $I$ 上，则 \(\displaystyle \sum_{n=1}^{\infty}u_n(x)\) 称为定义在 $I$ 上的函数项级数。
\end{definition}

\begin{definition}[收敛点、发散点与和函数]
对固定的 $x_0\in I$，若数项级数 \(\displaystyle \sum_{n=1}^{\infty}u_n(x_0)\) 收敛，则称函数项级数在 $x_0$ 处收敛，否则称在 $x_0$ 处发散。若它在区间内每一点都收敛，则称在该区间上逐点收敛。此时定义和函数 \(\displaystyle S(x)=\sum_{n=1}^{\infty}u_n(x)\)。
\end{definition}

\begin{remark}[三个基本问题]
若每一项 $u_n(x)$ 都在区间 $[a,b]$ 上连续、可积、可导，则其和函数是否仍连续、可积、可导？极限与积分、极限与求导能否交换？这些问题都需要通过一致收敛来控制。
\end{remark}

\subsubsection{一致收敛的定义与判别}

\begin{definition}[一致收敛]
设函数项级数 $\sum_{n=1}^{\infty}u_n(x)$ 在 $I$ 上的和函数为 $S(x)$。若对任意 $\varepsilon>0$，存在与 $x$ 无关的自然数 $N_\varepsilon$，使得当 $n>N_\varepsilon$ 时，
\(\displaystyle |S_n(x)-S(x)|<\varepsilon,\qquad \forall x\in I\)，
则称该级数在 $I$ 上一致收敛于 $S(x)$，记作 $S_n(x)\rightrightarrows S(x)$。
\end{definition}

\begin{definition}[函数列的一致收敛]
设函数列 $\{f_n(x)\}$ 定义在区间 $I$ 上。若存在函数 $f(x)$，使得对任意 $\varepsilon>0$，存在与 $x$ 无关的自然数 $N_\varepsilon$，使得当 $n>N_\varepsilon$ 时，
\(\displaystyle |f_n(x)-f(x)|<\varepsilon,\qquad \forall x\in I\)，
则称函数列 $\{f_n(x)\}$ 在 $I$ 上一致收敛于 $f(x)$，记作 $f_n(x)\rightrightarrows f(x)$。
\end{definition}

\begin{theorem}[一致收敛的 Cauchy 准则]
函数项级数 $\sum_{n=1}^{\infty}u_n(x)$ 在 $I$ 上一致收敛的充要条件是：对任意 $\varepsilon>0$，存在与 $x$ 无关的自然数 $N_\varepsilon$，使得当 $m>n>N_\varepsilon$ 时，
\(\displaystyle |u_{n+1}(x)+u_{n+2}(x)+\cdots+u_m(x)|<\varepsilon,\qquad \forall x\in I\)。
\end{theorem}

\begin{theorem}[Weierstrass 判别法]
若在区间 $I$ 上满足
\(\displaystyle |u_n(x)|\le M_n,\qquad x\in I\)，
且正项级数 $\sum_{n=1}^{\infty}M_n$ 收敛，则函数项级数
\(\displaystyle \sum_{n=1}^{\infty}u_n(x)\)
在 $I$ 上一致收敛。
\end{theorem}

\begin{theorem}[一致收敛的等价刻画]
函数列 $\{S_n(x)\}$ 在 $I$ 上一致收敛于 $S(x)$ 当且仅当
\(\displaystyle \sup_{x\in I}|S_n(x)-S(x)|\to 0\)。
等价地，对任意数列 $\{x_n\}\subset I$，都有
\(\displaystyle S_n(x_n)-S(x_n)\to 0\)。
\end{theorem}

\begin{theorem}[一致收敛的必要条件]
若函数项级数 $\sum_{n=1}^{\infty}u_n(x)$ 在区间 $I$ 上一致收敛，则 \(u_n(x)\rightrightarrows 0\)。若函数列 $\{f_n(x)\}$ 在 $I$ 上一致收敛于 $f(x)$，则 \(f_n(x)-f(x)\rightrightarrows 0\)。
\end{theorem}

\subsubsection{一致收敛的性质}

\begin{theorem}[连续性]
若级数 $\sum_{n=1}^{\infty}u_n(x)$ 在 $[a,b]$ 上一致收敛于 $S(x)$，且每一项都在 $[a,b]$ 上连续，则 $S(x)$ 也在 $[a,b]$ 上连续。
\end{theorem}

\begin{theorem}[逐项积分]
若级数 $\sum_{n=1}^{\infty}u_n(x)$ 在 $[a,b]$ 上一致收敛于 $S(x)$，且每一项都在 $[a,b]$ 上可积，则
\(\displaystyle \int_a^b S(x)\,\mathrm{d}x=\sum_{n=1}^{\infty}\int_a^b u_n(x)\,\mathrm{d}x\)。
\end{theorem}

\begin{theorem}[逐项求导]
若函数项级数 $\sum_{n=1}^{\infty}u_n(x)$ 在 $[a,b]$ 上收敛于 $S(x)$，每个 $u_n$ 在 $[a,b]$ 上有连续导数，且导函数级数
 \(\displaystyle \sum_{n=1}^{\infty}u_n'(x)\) 在 $[a,b]$ 上一致收敛，则
\(\displaystyle S'(x)=\sum_{n=1}^{\infty}u_n'(x)\)。
\end{theorem}

\subsubsection{典型例题：一致收敛与交换定理}

函数项级数本质上先形成一列部分和函数 \(S_n(x)=\sum_{k=1}^{n}u_k(x)\)。下面的例题都围绕同一个问题：\(S_n(x)\) 对每个固定 \(x\) 收敛，并不等于它在整个区间上收敛得同样稳定。

\begin{example}[逐点极限不保连续]
在区间 $[0,1]$ 上，考虑级数
\(\displaystyle x+\sum_{n=2}^{\infty}(x^n-x^{n-1})\)。
其部分和为
\(\displaystyle S_n(x)=x^n\)，并且当 $0\le x<1$ 时 \(\displaystyle \lim_{n\to\infty}S_n(x)=0\)，当 $x=1$ 时 \(\displaystyle \lim_{n\to\infty}S_n(x)=1\)。
因此和函数在 $x=1$ 处不连续，这说明逐点收敛不足以保证连续性。
\end{example}

\begin{example}[极限与积分不可随意交换]
设 \(\displaystyle S_n(x)=nx(1-x^2)^n,\qquad 0\le x\le 1\)。
对每个固定 $x$，有 $S_n(x)\to 0$，从而
\(\displaystyle \int_0^1 S(x)\,\mathrm{d}x=0\)。
但 \(\displaystyle \int_0^1 S_n(x)\,\mathrm{d}x
=\int_0^1 nx(1-x^2)^n\,\mathrm{d}x
=\frac{n}{2(n+1)}\to \frac12\)。
故极限与积分不能随意交换。
\end{example}

\begin{example}[极限与求导不可随意交换]
设 \(\displaystyle S_n(x)=e^{-n^2x^2},\qquad 0\le x\le 1\)，则 \(\displaystyle S_n'(x)=-2n^2xe^{-n^2x^2}\)。
对每个固定 $x$，当 $x=0$ 时 \(\displaystyle S(x)=\lim_{n\to\infty}S_n(x)=1\)，当 $0<x\le1$ 时 \(\displaystyle S(x)=0\)。
所以 $S(x)$ 在 $x=0$ 处不可导，但 $S_n$ 都可导。故仅有逐点极限时，不能交换极限与求导。
\end{example}

\begin{example}[在 $(0,1)$ 上不一致收敛]
对同一组部分和 $S_n(x)=x^n$，其在 $(0,1)$ 上的和函数为 $S(x)=0$。若一致收敛，则应有
\(\displaystyle |x^n|<\frac14,\qquad \forall x\in(0,1),\ n>N\)。
但对任意固定 $n$，当 $x$ 充分接近 $1$ 时，$x^n$ 可任意接近 $1$，矛盾。故它在 $(0,1)$ 上不一致收敛。
\end{example}

\begin{example}[在 $(0,1)$ 上一致收敛]
考虑级数 \(\displaystyle \frac1{1+x}+\sum_{n=2}^{\infty}\left(\frac1{n+x}-\frac1{n-1+x}\right)\)。
其部分和为
\(\displaystyle S_n(x)=\frac1{n+x}\)。
当 $x\in(0,1)$ 时，
\(\displaystyle |S_n(x)|=\frac1{n+x}<\frac1n\)。
故对任意 $\varepsilon>0$，取 $N>\frac1\varepsilon$ 即可推出
\(\displaystyle |S_n(x)|<\varepsilon,\qquad \forall x\in(0,1),\ n>N\)。
因此该级数在 $(0,1)$ 上一致收敛于 $0$。
\end{example}

\begin{example}[在闭区间上的一致收敛]
设 $a\in(0,1)$，考虑函数列 $S_n(x)=x^n$。对 $x\in[0,a]$，有 \(\displaystyle |S_n(x)|=x^n\le a^n\to 0\)。
故 $\{x^n\}$ 在 $[0,a]$ 上一致收敛于 $0$。
\end{example}

\begin{example}[用上界判别一致收敛]
设 \(\displaystyle f_n(x)=\frac{x^2}{1+nx^2}\)。则 \(\displaystyle |f_n(x)|\le \frac1n,\qquad \forall x\in\mathbb R\)。
故 $f_n(x)\rightrightarrows 0$ 于 $\mathbb R$。
\end{example}

\begin{example}[求最大值证明一致收敛]
设 \(\displaystyle S_n(x)=(1-x)x^n,\ x\in[0,1]\)。其最大值在 $x=\frac{n}{n+1}$ 处取到，且 \(\displaystyle \max_{x\in[0,1]}|(1-x)x^n|
=\frac1{n+1}\left(\frac{n}{n+1}\right)^n\to 0\)。故 $\{S_n\}$ 在 $[0,1]$ 上一致收敛于 $0$。
\end{example}

\begin{example}[在有界区间上一致收敛到 $e^x$]
设 \(\displaystyle S_n(x)=\left(1+\frac{x}{n}\right)^n\)。
对任意固定的 $a>0$，该函数列在 $[0,a]$ 上一致收敛于 $e^x$。
\end{example}

\begin{example}[构造反例判定不一致收敛]
设 \(\displaystyle S_n(x)=nx(1-x^2)^n,\ x\in[0,1]\)。取 $x_n=\frac1n$，则 \(\displaystyle |S_n(x_n)|=\left(1-\frac1{n^2}\right)^n\to 1\)。故该函数列在 $[0,1]$ 上不一致收敛于 $0$。
\end{example}

\begin{example}[在无界区间上不一致收敛]
设 \(\displaystyle S_n(x)=\frac{nx}{1+n^2x^2},\ x\in[0,+\infty)\)。对每个固定 $x$ 有 $S_n(x)\to 0$，但 \(\displaystyle \left|S_n\!\left(\frac1n\right)\right|=\frac12\)。故它在 $[0,+\infty)$ 上不一致收敛于 $0$。
\end{example}

\begin{example}[再看无界区间]
设 \(\displaystyle S_n(x)=\left(1+\frac{x}{n}\right)^n,\ x\in[0,+\infty)\)。虽然对每个固定 $x$ 有 $S_n(x)\to e^x$，但令 $x_n=n$，则 \(\displaystyle S_n(x_n)-e^{x_n}=2^n-e^n\not\to 0\)，故它在 $[0,+\infty)$ 上不一致收敛。
\end{example}

\begin{example}[Weierstrass 判别法]
证明 \(\displaystyle \sum_{n=1}^{\infty}\frac{x^n}{n^2}\) 在 $[-1,1]$ 上一致收敛。事实上，\(\displaystyle \left|\frac{x^n}{n^2}\right|\le \frac1{n^2},\ x\in[-1,1]\)，而 $\sum \frac1{n^2}$ 收敛，因此由 Weierstrass 判别法知原级数一致收敛。
\end{example}

\begin{example}[一致收敛保证连续]
设 \(\displaystyle S(x)=\sum_{n=0}^{\infty}\frac{x^n}{3^n}\cos(n\pi x)\)。在区间 $[-2,2]$ 上，\(\displaystyle \left|\frac{x^n}{3^n}\cos(n\pi x)\right|\le \left(\frac23\right)^n\)。故该级数在 $[-2,2]$ 上一致收敛，从而和函数连续。于是 \(\displaystyle \lim_{x\to 1}S(x)=S(1)=\sum_{n=0}^{\infty}\left(-\frac13\right)^n=\frac34\)。
\end{example}

\begin{example}[一致收敛保证逐项积分]
设 \(\displaystyle f(x)=\sum_{n=1}^{\infty}\frac{\cos nx}{n^2}\)。则 \(\displaystyle \int_0^\pi f(x)\,\mathrm{d}x
=\sum_{n=1}^{\infty}\frac1{n^2}\int_0^\pi \cos nx\,\mathrm{d}x=0\)。
\end{example}

\begin{example}[一致收敛保证逐项求导]
设 \(\displaystyle f(x)=\sum_{n=1}^{\infty}\frac{\sin nx}{n^4}\)。则它在 $\mathbb R$ 上有二阶连续导数，且 \(\displaystyle f'(x)=\sum_{n=1}^{\infty}\frac{\cos nx}{n^3},\qquad
f''(x)=-\sum_{n=1}^{\infty}\frac{\sin nx}{n^2}\)。
\end{example}

\subsection{幂级数专题：收敛半径、端点与和函数}

幂级数是函数项级数中最稳定的一类。它的收敛区域由收敛半径统一控制，内部可以逐项求导、逐项积分；端点一旦代入，又回到前面的数项级数判敛流程。

\subsubsection{幂级数与收敛半径}

\begin{definition}[幂级数]
形如 \(\displaystyle \sum_{n=0}^{\infty}a_nx^n
=a_0+a_1x+a_2x^2+\cdots\) 的函数项级数称为幂级数。
\end{definition}

\begin{theorem}[根值公式]
若 \(\displaystyle \alpha=\lim_{n\to\infty}\sqrt[n]{|a_n|},\qquad R=\frac1\alpha\)，
则幂级数 $\sum a_nx^n$ 满足：
\begin{itemize}
  \item 当 $R=0$ 时，只在 $x=0$ 处收敛；
  \item 当 $R=+\infty$ 时，在整个实轴上绝对收敛；
  \item 当 $0<R<+\infty$ 时，在 $(-R,R)$ 内绝对收敛，在 $|x|>R$ 时发散。
\end{itemize}
\end{theorem}

\begin{theorem}[比值公式]
若 \(\displaystyle \alpha=\lim_{n\to\infty}\left|\frac{a_{n+1}}{a_n}\right|,\qquad R=\frac1\alpha\)，
则可得到与根值公式相同的结论。
\end{theorem}

\begin{remark}[收敛区间与端点]
幂级数的收敛点集总是某个以原点为中心的区间 $(-R,R)$。端点 $x=\pm R$ 的敛散性必须单独讨论。
端点一代入，就不再是幂级数问题，而是普通数项级数问题；因此仍按前面的流程检查通项、绝对值、估阶模型和条件收敛。
\end{remark}

\begin{theorem}[Abel 定理]
若幂级数 $\sum a_nx^n$ 在 $x=x_0\ne 0$ 处收敛，则它在所有满足 $|x|<|x_0|$ 的点绝对收敛；若它在 $x=x_1$ 处发散，则它在所有满足 $|x|>|x_1|$ 的点必发散。
\end{theorem}

\begin{theorem}[收敛区间内部的一致收敛]
若幂级数 $\sum a_nx^n$ 的收敛半径为 $R$，则对任意 $r\in(0,R)$，它在闭区间 $[-r,r]$ 上一致收敛。
\end{theorem}

\subsubsection{和函数的性质}

\begin{theorem}[连续、可导、可积]
设幂级数 $\sum_{n=0}^{\infty}a_nx^n$ 的收敛半径为 $R$，和函数为 $S(x)$，则在 $(-R,R)$ 内：
\begin{itemize}
  \item $S(x)$ 连续；
  \item $S(x)$ 可逐项求导，且 \(\displaystyle S'(x)=\sum_{n=1}^{\infty}na_nx^{n-1}\)；
  \item $S(x)$ 可逐项积分，且 \(\displaystyle \int_0^x\left(\sum_{n=0}^{\infty}a_nt^n\right)\mathrm{d}t=\sum_{n=0}^{\infty}\frac{a_n}{n+1}x^{n+1}\)。
\end{itemize}
导出级数与积分后所得级数的收敛半径仍为 $R$。
\end{theorem}

\begin{theorem}[端点处的连续性]
若幂级数 $\sum_{n=0}^{\infty}a_nx^n$ 的收敛半径为 $R$，并且它在 $x=R$ 处收敛，则其和函数在 $x=R$ 处左连续；若它在 $x=-R$ 处收敛，则其和函数在 $x=-R$ 处右连续。
\end{theorem}

\begin{theorem}[幂级数的加法与乘法]
设幂级数 \(\displaystyle \sum_{n=0}^{\infty}a_nx^n\) 与 \(\displaystyle \sum_{n=0}^{\infty}b_nx^n\) 在公共收敛区间内都有意义，则它们可逐项加减。若两者收敛半径同为 $R$，则在 $|x|<R$ 内可作乘法，并有
\[
\left(\sum_{n=0}^{\infty}a_nx^n\right)\left(\sum_{n=0}^{\infty}b_nx^n\right)
=\sum_{n=0}^{\infty}c_nx^n,
\qquad
c_n=\sum_{i=0}^{n}a_ib_{n-i}.
\]
\end{theorem}

\begin{conclusion}[常见幂级数的和函数]
下面这些式子先按收敛区间内部使用；若题目要取端点值，必须把端点级数单独代回去判敛，再借端点连续性补值。

\begin{compactenum}
\item 几何级数及其求导变形：
\[
\sum_{n=0}^{\infty}x^n=\frac1{1-x},\qquad
\sum_{n=1}^{\infty}nx^{n-1}=\frac1{(1-x)^2},\qquad |x|<1.
\]
因此
\[
\sum_{n=1}^{\infty}nx^n=\frac{x}{(1-x)^2},\qquad
\sum_{n=0}^{\infty}(n+1)x^n=\frac1{(1-x)^2},\qquad
\sum_{n=1}^{\infty}n^2x^n=\frac{x(1+x)}{(1-x)^3}.
\]

\item 对几何级数逐项积分，常得到对数型和函数：
\[
\sum_{n=1}^{\infty}\frac{x^n}{n}=-\ln(1-x),\qquad |x|<1,
\]
\[
\sum_{n=1}^{\infty}(-1)^{n-1}\frac{x^n}{n}=\ln(1+x),\qquad |x|<1.
\]
特别地，后一式在 $x=1$ 处可补成 \(\displaystyle \sum_{n=1}^{\infty}(-1)^{n-1}\frac1n=\ln2\)。

\item 只含偶次或奇次幂时，先把变量换成 $x^2$：
\[
\sum_{n=0}^{\infty}x^{2n}=\frac1{1-x^2},\qquad
\sum_{n=0}^{\infty}x^{2n+1}=\frac{x}{1-x^2},\qquad |x|<1.
\]
再逐项积分可得
\[
\sum_{n=0}^{\infty}\frac{x^{2n+1}}{2n+1}
=\frac12\ln\frac{1+x}{1-x},\qquad |x|<1,
\]
\[
\sum_{n=0}^{\infty}(-1)^n\frac{x^{2n+1}}{2n+1}
=\arctan x,\qquad |x|\le 1.
\]

\item 阶乘型和三角型常直接认 Taylor 级数：
\[
\sum_{n=0}^{\infty}\frac{x^n}{n!}=e^x,\qquad
\sum_{n=0}^{\infty}\frac{nx^n}{n!}=xe^x,\qquad
\sum_{n=0}^{\infty}\frac{n^2x^n}{n!}=(x^2+x)e^x.
\]
\[
\sum_{n=0}^{\infty}(-1)^n\frac{x^{2n}}{(2n)!}=\cos x,\qquad
\sum_{n=0}^{\infty}(-1)^n\frac{x^{2n+1}}{(2n+1)!}=\sin x.
\]
\[
\sum_{n=0}^{\infty}\frac{x^{2n}}{(2n)!}=\cosh x,\qquad
\sum_{n=0}^{\infty}\frac{x^{2n+1}}{(2n+1)!}=\sinh x.
\]

\item 二项式型可记为
\[
\sum_{n=0}^{\infty}\binom{\alpha}{n}x^n=(1+x)^\alpha,\qquad |x|<1,
\]
其中 \(\displaystyle \binom{\alpha}{n}=\frac{\alpha(\alpha-1)\cdots(\alpha-n+1)}{n!}\)。当 $\alpha$ 是非负整数时，右边退化为有限项多项式。
\end{compactenum}
\end{conclusion}

\subsubsection{典型例题：收敛半径与和函数}

\begin{example}[求收敛半径]
\begin{compactenum}
  \item 对 \(\displaystyle \sum_{n=0}^{\infty}\frac{x^n}{n!}\)，有 \(\displaystyle \lim_{n\to\infty}\left|\frac{1/(n+1)!}{1/n!}\right|=0\)，故 $R=+\infty$。
  \item 对 \(\displaystyle \sum_{n=1}^{\infty}n^n x^n\)，有 \(\displaystyle \lim_{n\to\infty}\sqrt[n]{n^n}=+\infty\)，故 $R=0$。
  \item 对 \(\displaystyle \sum_{n=1}^{\infty}nx^n\)，有 \(\displaystyle \lim_{n\to\infty}\sqrt[n]{n}=1\)，故 $R=1$。
\end{compactenum}
\end{example}

\begin{example}[端点讨论]
\begin{compactenum}
  \item \(\displaystyle \sum_{n=1}^{\infty}\frac{x^n}{n}\) 的收敛半径为 $1$，在 $x=-1$ 处条件收敛，在 $x=1$ 处发散。
  \item \(\displaystyle \sum_{n=1}^{\infty}\frac{x^n}{n^2}\) 的收敛半径为 $1$，在 $x=\pm1$ 处都绝对收敛。
  \item \(\displaystyle \sum_{n=1}^{\infty}nx^n\) 的收敛半径为 $1$，在 $x=\pm1$ 处都发散。
\end{compactenum}
\end{example}

\begin{example}[求收敛域]
求 \(\displaystyle \sum_{n=1}^{\infty}(-1)^{n-1}\frac{(x-1)^n}{n}\) 的收敛域。令 $y=x-1$，则化为 \(\displaystyle \sum_{n=1}^{\infty}(-1)^{n-1}\frac{y^n}{n}\)。
其收敛半径为 $1$，且在 $y=1$ 处收敛、在 $y=-1$ 处发散，因此原级数的收敛域为
\(\displaystyle 0<x\le 2\)。
\end{example}

\begin{example}[再求收敛域]
求 \(\displaystyle \sum_{n=1}^{\infty}\frac{x^n}{n2^n}\) 的收敛域。由 \(\displaystyle \lim_{n\to\infty}\sqrt[n]{\frac1{n2^n}}=\frac12\) 得收敛半径 $R=2$。当 $x=2$ 时，级数化为调和级数而发散；当 $x=-2$ 时，化为交错调和级数而收敛，故收敛域为 \(\displaystyle [-2,2)\)。
\end{example}

\begin{example}[含偶次幂的收敛域]
求 \(\displaystyle \sum_{n=1}^{\infty}\frac{x^{2n}}{n^2}\) 的收敛域。令 $y=x^2$，则化为 \(\displaystyle \sum_{n=1}^{\infty}\frac{y^n}{n^2}\)。该级数在 $|y|\le 1$ 时收敛，而 $y=x^2\ge 0$，故原级数的收敛域为 \(\displaystyle -1\le x\le 1\)。
\end{example}

\begin{example}[已知一点判断另一点]
若级数 \(\displaystyle \sum_{n=1}^{\infty}a_n(x-1)^n\) 在 $x=-1$ 处收敛，则令 $y=x-1$ 可知其在 $y=-2$ 处收敛。由 Abel 定理，凡满足 $|y|<2$ 的点都绝对收敛。故在 $x_0=2$ 处，即 $y_0=1$ 处，原级数绝对收敛。
\end{example}

\begin{example}[求和函数]
由几何级数 \(\displaystyle \sum_{n=0}^{\infty}x^n=\frac1{1-x}\ (|x|<1)\)，逐项求导得
\(\displaystyle \sum_{n=1}^{\infty}nx^{n-1}=\frac1{(1-x)^2}\)，再乘以 $x$ 得
\(\displaystyle \sum_{n=1}^{\infty}nx^n=\frac{x}{(1-x)^2}\ (|x|<1)\)。
特别地，
\(\displaystyle \sum_{n=1}^{\infty}\frac{n}{2^n}=2,\qquad
\sum_{n=1}^{\infty}\frac{n}{3^n}=\frac34\)。
\end{example}

\begin{example}[再求和函数]
设 \(\displaystyle S(x)=\sum_{n=1}^{\infty}\frac{x^n}{n}\)。则 \(\displaystyle S'(x)=\sum_{n=1}^{\infty}x^{n-1}=\frac1{1-x}\ (|x|<1)\)，从而 \(\displaystyle S(x)=-\ln(1-x)\ (x\in[-1,1))\)。
代入 $x=-1$ 得
\(\displaystyle \ln 2=1-\frac12+\frac13-\frac14+\cdots\)。
\end{example}

\begin{example}[积分法求和函数]
设 \(\displaystyle S(x)=\sum_{n=0}^{\infty}(n+1)x^n\)。则
\(\displaystyle \int_0^x S(t)\,\mathrm{d}t
=\sum_{n=0}^{\infty}\int_0^x (n+1)t^n\,\mathrm{d}t
=\sum_{n=0}^{\infty}x^{n+1}
=\frac{x}{1-x}\)，故
\(\displaystyle S(x)=\left(\frac{x}{1-x}\right)'=\frac1{(1-x)^2},\quad |x|<1\)。
\end{example}

\begin{example}[用和函数证明 Leibniz 公式]
由 \(\displaystyle \frac1{1+x^2}=\sum_{n=0}^{\infty}(-1)^nx^{2n},\ |x|<1\)，逐项积分得
\(\displaystyle \arctan x=\sum_{n=0}^{\infty}(-1)^n\frac{x^{2n+1}}{2n+1},\ |x|<1\)。
令 $x\to 1^-$，得到 \(\displaystyle \frac{\pi}{4}=1-\frac13+\frac15-\frac17+\cdots\)。
\end{example}

\begin{example}[数项级数求和]
\begin{compactenum}
  \item 由 \(\displaystyle 1-x^3+x^6-x^9+\cdots=\frac1{1+x^3}\)，逐项积分可得
  \(\displaystyle 1-\frac14+\frac17-\frac1{10}+\cdots=\frac13\ln 2+\frac{\pi}{3\sqrt3}\)。
  \item 对 \(\displaystyle 1+x-x^2-x^3+x^4+x^5-\cdots\) 分组并逐项积分，可得
  \(\displaystyle 1+\frac12-\frac13-\frac14+\frac15+\frac16-\frac17-\frac18+\cdots=\frac{\pi}{4}+\frac12\ln 2\)。
\end{compactenum}
\end{example}

\begin{example}[恒等式]
证明 \(\displaystyle \sum_{n=0}^{\infty}(n+1)(n+2)x^n=\frac{2}{(1-x)^3},\ |x|<1\)。
只需写成 \(\displaystyle (n+1)(n+2)=n^2+3n+2\)，并结合
\(\displaystyle \sum_{n=0}^{\infty}nx^n=\frac{x}{(1-x)^2}\) 与
\(\displaystyle \sum_{n=0}^{\infty}n^2x^n=\frac{x(x+1)}{(1-x)^3}\)，即可整理得到结论。
\end{example}

\subsection{泰勒展开专题：经典展开、余项与应用}

前面是“给出一个幂级数，研究它的和函数”。Taylor 展开反过来问：给定一个函数，能不能把它写成幂级数？关键不只是写出形式，而是证明余项趋于零，使级数真的等于原函数。

\subsubsection{Taylor 级数与 Maclaurin 级数}

\begin{definition}[Taylor 级数]
若函数 $f(x)$ 在 $x=a$ 附近可展开成幂级数
\(\displaystyle f(x)=\sum_{n=0}^{\infty}a_n(x-a)^n,\ a_n=\frac{f^{(n)}(a)}{n!}\)，
因此其展开式只能是
\(\displaystyle f(x)=\sum_{n=0}^{\infty}\frac{f^{(n)}(a)}{n!}(x-a)^n\)。
这称为 $f(x)$ 在点 $a$ 附近的 Taylor 展开。
\end{definition}

\begin{definition}[Maclaurin 级数]
当 $a=0$ 时，Taylor 级数 \(\displaystyle \sum_{n=0}^{\infty}\frac{f^{(n)}(0)}{n!}x^n\) 称为 Maclaurin 级数。
\end{definition}

\begin{theorem}[Taylor 公式]
若 $f$ 有足够阶导数，则
\[
f(x)=f(a)+f'(a)(x-a)+\frac{f''(a)}{2!}(x-a)^2+\cdots+\frac{f^{(n)}(a)}{n!}(x-a)^n+R_n(x),
\]
其中 Lagrange 余项为 \(\displaystyle R_n(x)=\frac{f^{(n+1)}(\xi)}{(n+1)!}(x-a)^{n+1}\)，
$\xi$ 介于 $a$ 与 $x$ 之间。
\end{theorem}

\begin{theorem}[展开成 Taylor 级数的充要条件]
设 $f(x)$ 在区间 $(a-R,a+R)$ 上有任意阶导数，则 $f(x)$ 在该区间内等于其 Taylor 级数的充要条件是：对任意 $x\in(a-R,a+R)$，\(\displaystyle \lim_{n\to\infty}R_n(x)=0\)。
\end{theorem}

\begin{theorem}[一个常用充分条件]
若存在常数 $M$，使对 $(a-R,a+R)$ 内任意 $x$ 和任意自然数 $n$ 都有 \(\displaystyle |f^{(n)}(x)|\le M\)，则 $f(x)$ 在该区间内可以展开成 Taylor 级数。
\end{theorem}

\subsubsection{典型展开：基础函数与换中心}

\begin{example}[Taylor 级数收敛但不等于原函数]
设 \(f(x)=e^{-1/x^2}\ (x\ne0)\)，且 \(f(0)=0\)。
可证明 \(\displaystyle f^{(n)}(0)=0,\qquad n=0,1,2,\dots\)
故其 Taylor 级数恒为 $0$。虽然该级数收敛，但它的和不是原函数，这说明“有任意阶导数”并不足以保证 Taylor 级数收敛到原函数。
\end{example}

\begin{example}[$e^x$ 的 Maclaurin 展开]
因为 \(\displaystyle (e^x)^{(n)}=e^x,\ (e^x)^{(n)}\big|_{x=0}=1\)，故
\(\displaystyle e^x=\sum_{n=0}^{\infty}\frac{x^n}{n!},\quad x\in\mathbb R\)。
\end{example}

\begin{example}[$\sin x,\cos x$ 的 Maclaurin 展开]
\(\displaystyle \sin x=\sum_{k=0}^{\infty}(-1)^k\frac{x^{2k+1}}{(2k+1)!},\qquad \cos x=\sum_{k=0}^{\infty}(-1)^k\frac{x^{2k}}{(2k)!},\qquad x\in\mathbb R\)。
\end{example}

\begin{example}[$\ln(1+x)$ 的 Maclaurin 展开]
由 \(\displaystyle \frac1{1+x}=\sum_{n=0}^{\infty}(-1)^nx^n,\ |x|<1\)，逐项积分得
\(\displaystyle \ln(1+x)=\sum_{n=0}^{\infty}(-1)^n\frac{x^{n+1}}{n+1},\ |x|<1\)。
该级数在 $x=1$ 处也收敛，因此
\(\displaystyle \ln(1+x)=\sum_{n=0}^{\infty}(-1)^n\frac{x^{n+1}}{n+1},\qquad x\in(-1,1]\)。
\end{example}

\begin{example}[$(1+x)^\alpha$ 的 Maclaurin 展开]
设 $\alpha\in\mathbb R$，则
\[
(1+x)^\alpha
\sim
1+\alpha x+\frac{\alpha(\alpha-1)}{2!}x^2+\cdots+\frac{\alpha(\alpha-1)\cdots(\alpha-n+1)}{n!}x^n+\cdots
\]
其收敛半径为 $1$，且在 $(-1,1)$ 内有级数恒等式成立。端点敛散性为：
\begin{itemize}
  \item $\alpha\le -1$ 时，收敛区间为 $(-1,1)$；
  \item $-1<\alpha<0$ 时，收敛区间为 $(-1,1]$；
  \item $\alpha>0$ 时，收敛区间为 $[-1,1]$。
\end{itemize}
\end{example}

\begin{example}[常用熟知展开]
\begin{align*}
\frac1{1+x}&=1-x+x^2-x^3+\cdots,\qquad x\in(-1,1),\\
\sqrt{1+x}
&=1+\frac12x+\sum_{n=2}^{\infty}(-1)^{n-1}\frac{(2n-3)!!}{(2n)!!}x^n,\qquad x\in[-1,1],\\
\frac1{\sqrt{1+x}}
&=1-\frac12x+\sum_{n=2}^{\infty}(-1)^n\frac{(2n-1)!!}{(2n)!!}x^n,\qquad x\in(-1,1].
\end{align*}
\end{example}

\begin{example}[$\arctan x$ 的 Maclaurin 展开]
由 \(\displaystyle (\arctan x)'=\frac1{1+x^2}=1-x^2+x^4-\cdots,\qquad |x|<1\)，积分得 \(\displaystyle \arctan x=x-\frac{x^3}{3}+\frac{x^5}{5}-\cdots+(-1)^n\frac{x^{2n+1}}{2n+1}+\cdots,\qquad x\in[-1,1]\)。
\end{example}

\begin{example}[$\arcsin x$ 的 Maclaurin 展开]
由
\begin{align*}
(\arcsin x)'&=\frac1{\sqrt{1-x^2}}
=1+\frac12x^2+\frac{1\cdot3}{2\cdot4}x^4+\cdots+\frac{(2n-1)!!}{(2n)!!}x^{2n}+\cdots,\\
\arcsin x
&=x+\sum_{n=1}^{\infty}\frac{(2n-1)!!}{(2n)!!(2n+1)}x^{2n+1},\qquad x\in(-1,1].
\end{align*}
再由 \(\displaystyle \arccos x=\frac\pi2-\arcsin x,\qquad
\operatorname{arccot}x=\frac\pi2-\arctan x\) 可得到相关展开。
\end{example}

\begin{example}[换中心展开：$\ln x$]
写成 \(\displaystyle x=2\left(1+\frac{x-2}{2}\right),\ \ln x=\ln 2+\ln\left(1+\frac{x-2}{2}\right)\)。
代入 $\ln(1+y)$ 的展开式，得
\(\displaystyle \ln x=\ln 2+\sum_{n=1}^{\infty}(-1)^{n-1}\frac{(x-2)^n}{n2^n},\quad x\in(0,4]\)。
\end{example}

\begin{example}[换中心展开：$\sin x$]
由和角公式 \(\displaystyle
\sin x=\sin\left(\frac\pi4+x-\frac\pi4\right)
=\frac{\sqrt2}{2}\cos\left(x-\frac\pi4\right)+\frac{\sqrt2}{2}\sin\left(x-\frac\pi4\right)\)，
再分别展开 $\sin$ 与 $\cos$，可得关于 $x-\frac\pi4$ 的幂级数展开。
\end{example}

\begin{example}[幂级数法求微分方程的级数解]
设微分方程 \(\displaystyle xy''+y'+xy=0,\qquad y=\sum_{n=0}^{\infty}a_nx^n\)，则
\[
y'=a_1+2a_2x+3a_3x^2+\cdots+na_nx^{n-1}+\cdots,\qquad
y''=2a_2+3\cdot2a_3x+4\cdot3a_4x^2+\cdots+n(n-1)a_nx^{n-2}+\cdots.
\]
代回方程后得 \(\displaystyle 0=xy''+y'+xy=a_1+(a_0+4a_2)x+(a_1+9a_3)x^2+(a_2+16a_4)x^3+\cdots\)。
比较系数并令 $a_0=c$，得到
\begin{align*}
a_1&=0,\quad a_0+4a_2=0,\quad a_1+9a_3=0,\quad a_2+16a_4=0,\\
a_{2k-1}+(2k+1)^2a_{2k+1}&=0,\qquad
a_{2k}+(2k+2)^2a_{2k+2}=0\qquad (k=1,2,\cdots),\\
a_1&=a_3=\cdots=0,\\
a_{2k}&=(-1)^k\frac{c}{2^{2k}(k!)^2},\\
y&=c\sum_{k=0}^{\infty}(-1)^k\frac{1}{(k!)^2}\left(\frac{x}{2}\right)^{2k}.
\end{align*}
这是一个 Bessel 函数型解。
\end{example}

\begin{example}[幂级数法求微分方程]
求解 \(\displaystyle (1-x^2)y''=-2y\)。设 \(\displaystyle y=\sum_{n=0}^{\infty}a_nx^n\)，则
\(\displaystyle y'=\sum_{n=1}^{\infty}na_nx^{n-1},\qquad
y''=\sum_{n=2}^{\infty}n(n-1)a_nx^{n-2}\)。
代入并比较同次幂系数，得到
\begin{align*}
a_2&=-a_0,\qquad
a_3=-\frac{a_1}{3},\qquad
a_4=0,\qquad
a_5=-\frac{a_1}{5\cdot 3},\\
a_{n+2}&=\frac{n-2}{n+2}a_n,\qquad
a_4=a_6=\cdots=0,\qquad
a_{2n+1}=-\frac{a_1}{(2n+1)(2n-1)}.
\end{align*}
故方程的解可写为 \(\displaystyle y=a_0(1-x^2)+a_1x-a_1\sum_{n=1}^{\infty}\frac{x^{2n+1}}{(2n+1)(2n-1)},\ x\in(-1,1)\)。
\end{example}

\subsection{傅里叶展开专题一：周期函数的三角级数}

Taylor 展开用 \(1,x,x^2,\dots\) 这组幂函数刻画函数的局部或幂级数结构。对于周期函数，更自然的基底是 \(1,\cos nx,\sin nx\)。Fourier 级数就是把函数展开到这组三角函数系上。

\subsubsection{三角级数、正交性与系数公式}

最简单的周期函数是简谐波 $x(t)=a\sin(\omega t+\varphi)$，其周期为 $T=\frac{2\pi}{\omega}$。若原函数以 $T$ 为周期，可先作代换
\(\displaystyle x=\frac{2\pi}{T}t,\quad t=\frac{T}{2\pi}x,\quad f(x)=g\!\left(\frac{T}{2\pi}x\right)\)，
把它化为周期 $2\pi$ 的问题来处理；若题目已经指明按某个周期展开，就以题目指定周期为准，不要自行改用函数的最小正周期。对周期为 $2\pi$ 的函数，最常用的展开形式是
\(\displaystyle f(x)\sim \frac{a_0}{2}
+\sum_{n=1}^{\infty}\bigl(a_n\cos nx+b_n\sin nx\bigr)\)。
这称为 $f(x)$ 的傅里叶级数。做题时应先看周期、奇偶性和定义区间，再决定是否直接套系数公式。

\begin{theorem}[三角函数系的正交性]
在区间 $[-\pi,\pi]$ 上，
\begin{align*}
\int_{-\pi}^{\pi}1^2\,\mathrm{d}x&=2\pi,\qquad
\int_{-\pi}^{\pi}\sin nx\,\mathrm{d}x=0,\qquad
\int_{-\pi}^{\pi}\cos nx\,\mathrm{d}x=0\ (n\ge 1),\\
\int_{-\pi}^{\pi}\cos mx\cos nx\,\mathrm{d}x&=
\begin{cases}
0,& m\neq n,\\
\pi,& m=n\ge 1,
\end{cases}
\qquad
\int_{-\pi}^{\pi}\sin mx\sin nx\,\mathrm{d}x=
\begin{cases}
0,& m\neq n,\\
\pi,& m=n\ge 1.
\end{cases}\\
\int_{-\pi}^{\pi}\sin mx\cos nx\,\mathrm{d}x&=0,\\
\int_{-\pi}^{\pi}\cos mx\cos nx\,\mathrm{d}x&=\pi\delta_{mn},\qquad
\int_{-\pi}^{\pi}\sin mx\sin nx\,\mathrm{d}x=\pi\delta_{mn}\quad(m,n\ge 1).
\end{align*}
其中 $\delta_{mn}$ 为克罗内克符号，即 \(\displaystyle \delta_{mn}=0\ (m\ne n),\ \delta_{mn}=1\ (m=n)\)。相应的标准正交系可写成 \(\displaystyle \frac{1}{\sqrt{2\pi}},\ \frac{\cos x}{\sqrt{\pi}},\ \frac{\sin x}{\sqrt{\pi}},\dots\)。
若定义内积 $\langle f,g\rangle=\int_{-\pi}^{\pi}f(x)g(x)\,\mathrm{d}x$，则 $\langle f,g\rangle=0$ 就表示 $f,g$ 在 $[-\pi,\pi]$ 上正交。
\end{theorem}

\begin{theorem}[傅里叶系数公式]
若在 $[-\pi,\pi]$ 上有展开
\(\displaystyle f(x)\sim \frac{a_0}{2}+\sum_{n=1}^{\infty}\bigl(a_n\cos nx+b_n\sin nx\bigr)\)
并收敛到 $f(x)$，则
\(\displaystyle a_n=\frac{1}{\pi}\int_{-\pi}^{\pi}f(x)\cos nx\,\mathrm{d}x\ (n=0,1,2,\dots)\)，
\(\displaystyle b_n=\frac{1}{\pi}\int_{-\pi}^{\pi}f(x)\sin nx\,\mathrm{d}x\ (n=1,2,\dots)\)。
若 $f$ 为奇函数，则 $a_0=a_n=0$，且 $b_n=\frac{2}{\pi}\int_{0}^{\pi}f(x)\sin nx\,\mathrm{d}x$；若 $f$ 为偶函数，则 $b_n=0$，且 $a_n=\frac{2}{\pi}\int_{0}^{\pi}f(x)\cos nx\,\mathrm{d}x$。
\end{theorem}

\begin{theorem}[Dirichlet 收敛定理]
若函数 $f(x)$ 以 $2\pi$ 为周期，并且在一个周期内连续或只有有限个第一类间断点，同时逐段单调，则它的傅里叶级数收敛。连续点 $x$ 处收敛到 $f(x)$，第一类间断点 $x$ 处收敛到 $\frac{f(x-0)+f(x+0)}{2}$。
\end{theorem}

\begin{note}[收敛值陷阱]
把展开式代入特殊点求数项级数和时，必须先判断该点是不是间断点。若是间断点，代入的应是左右极限的平均值，而不是原函数在该点的写法。
\end{note}

\subsubsection{典型展开：奇偶函数与分段函数}

\begin{example}[分段函数的傅里叶展开]
设 $f(x)$ 是以 $2\pi$ 为周期的函数，在 $(-\pi,\pi]$ 上
\[
f(x)=
\begin{cases}
0, & x\in(-\pi,0],\\
x, & x\in(0,\pi].
\end{cases}
\]
则 \(\displaystyle a_0=\frac{\pi}{2},\ a_n=\frac{(-1)^n-1}{n^2\pi},\ b_n=\frac{(-1)^{n+1}}{n}\)。故其傅里叶级数为
\[
f(x)\sim \frac{\pi}{4}
+\sum_{n=1}^{\infty}\left[
\frac{(-1)^n-1}{n^2\pi}\cos nx+\frac{(-1)^{n+1}}{n}\sin nx
\right].
\]
当 $x=(2k-1)\pi$ 时，函数有跳跃间断，级数收敛到 $\frac{f(x-0)+f(x+0)}{2}=\frac{\pi}{2}$。
由 $x=0$ 或 $x=\pi$ 等特殊值，还可得到
\(\displaystyle \sum_{n=1}^{\infty}\frac1{n^2}=\frac{\pi^2}{6},\qquad
\sum_{n=1}^{\infty}\frac{1}{(2n-1)^2}=\frac{\pi^2}{8},\qquad
\sum_{n=1}^{\infty}\frac{(-1)^{n-1}}{n^2}=\frac{\pi^2}{12}\)。
\end{example}

\begin{example}[函数 $f(x)=x$ 的傅里叶展开]
设 $f(x)=x$ 在 $(-\pi,\pi]$ 上定义，并以 $2\pi$ 为周期延拓。由于 $f$ 为奇函数，有 $a_0=a_n=0$，且 \(\displaystyle b_n=\frac{1}{\pi}\int_{-\pi}^{\pi}x\sin nx\,\mathrm{d}x=\frac{2(-1)^{n+1}}{n}\)。
因此 \(\displaystyle x\sim 2\sum_{n=1}^{\infty}\frac{(-1)^{n+1}}{n}\sin nx,\quad x\in(-\pi,\pi)\)。
当 $x=(2k+1)\pi$ 时，级数收敛到 $0$。取 $x=\dfrac{\pi}{2}$ 得 $\frac{\pi}{4}=1-\frac13+\frac15-\frac17+\cdots$。
\end{example}

\begin{example}[方波的傅里叶展开]
设 $f(x)$ 以 $2\pi$ 为周期，在 $(-\pi,\pi]$ 上满足 \(f(x)=1\ (x\in(-\pi,0))\)，且 \(f(x)=-1\ (x\in(0,\pi))\)。
它是奇函数，因此
\(\displaystyle a_0=a_n=0,\qquad b_{2k}=0,\qquad b_{2k-1}=-\frac{4}{(2k-1)\pi}\)。
\(\displaystyle f(x)\sim -\frac{4}{\pi}\sum_{k=1}^{\infty}\frac{\sin(2k-1)x}{2k-1}\)。
在间断点 $x=k\pi$ 处，级数收敛到 $0$。
\end{example}

\begin{example}[偶函数的余弦展开]
设 $f(x)$ 以 $2\pi$ 为周期，在 $(-\pi,\pi]$ 上满足 \(f(x)=1+\frac{2x}{\pi}\ (-\pi<x\le 0)\)，且 \(f(x)=1-\frac{2x}{\pi}\ (0<x\le \pi)\)。
它是偶函数，因此 $b_n=0$。计算得
\(\displaystyle a_0=0,\ a_{2k-1}=\frac{8}{(2k-1)^2\pi^2},\ a_{2k}=0\)。
\(\displaystyle f(x)\sim \frac{8}{\pi^2}\sum_{k=0}^{\infty}\frac{\cos(2k+1)x}{(2k+1)^2}\)。
\end{example}

\subsection{傅里叶展开专题二：半区间、任意区间与缩放}

\subsubsection{周期延拓、半区间展开与区间缩放}

若函数只在有限区间上给出，做题时要先判断是直接作周期延拓，还是先作奇延拓、偶延拓，再决定用整区间公式、正弦级数还是余弦级数。

\begin{theorem}[区间 $[-\pi,\pi]$ 上的周期延拓]
若函数 $f(x)$ 在 $[-\pi,\pi]$ 上有定义，可把它以 $2\pi$ 为周期延拓到整个实轴。若延拓后的函数满足收敛条件，则在连续点处，傅里叶级数收敛于 $f(x)$；在第一类间断点处，收敛于 $\frac{f(x-0)+f(x+0)}{2}$。特别地，若 $f(\pi)\ne f(-\pi)$，则在端点 $x=\pm\pi$ 处收敛到 $\frac{f(\pi-0)+f(-\pi+0)}{2}$。
\end{theorem}

\begin{theorem}[半区间上的正弦级数与余弦级数]
若函数只在 $[0,\pi]$ 上给出，则：
\begin{itemize}
\item 作偶延拓，可得到余弦级数 $f(x)\sim \frac{a_0}{2}+\sum_{n=1}^{\infty}a_n\cos nx$，其中 $a_n=\frac{2}{\pi}\int_{0}^{\pi}f(x)\cos nx\,\mathrm{d}x$；
\item 作奇延拓，可得到正弦级数 $f(x)\sim \sum_{n=1}^{\infty}b_n\sin nx$，其中 $b_n=\frac{2}{\pi}\int_{0}^{\pi}f(x)\sin nx\,\mathrm{d}x$。
\end{itemize}
若区间改为 $[0,l]$，则相应地有
\begin{align*}
f(x)&\sim \frac{a_0}{2}+\sum_{n=1}^{\infty}a_n\cos\frac{n\pi x}{l},
&a_n&=\frac{2}{l}\int_{0}^{l}f(x)\cos\frac{n\pi x}{l}\,\mathrm{d}x,\\
f(x)&\sim \sum_{n=1}^{\infty}b_n\sin\frac{n\pi x}{l},
&b_n&=\frac{2}{l}\int_{0}^{l}f(x)\sin\frac{n\pi x}{l}\,\mathrm{d}x.
\end{align*}
在 $[0,l]$ 上的正交关系为
\[
\int_{0}^{l}\sin\frac{m\pi x}{l}\sin\frac{n\pi x}{l}\,\mathrm{d}x
=\int_{0}^{l}\cos\frac{m\pi x}{l}\cos\frac{n\pi x}{l}\,\mathrm{d}x
=
\begin{cases}
0,& m\neq n,\\[4pt]
\dfrac{l}{2},& m=n\ge 1.
\end{cases}
\]
\end{theorem}

\begin{theorem}[区间 $[-l,l]$ 上的傅里叶级数]
令 $y=\frac{\pi x}{l}$，则定义在 $[-l,l]$ 上的函数可缩放到 $[-\pi,\pi]$ 上，从而得到
\begin{align*}
f(x)&\sim \frac{a_0}{2}
+\sum_{n=1}^{\infty}\left(a_n\cos\frac{n\pi x}{l}+b_n\sin\frac{n\pi x}{l}\right),\\
a_n&=\frac{1}{l}\int_{-l}^{l}f(x)\cos\frac{n\pi x}{l}\,\mathrm{d}x,\qquad
b_n=\frac{1}{l}\int_{-l}^{l}f(x)\sin\frac{n\pi x}{l}\,\mathrm{d}x.
\end{align*}
\end{theorem}

若原区间是 $[a,b]$ 而不是关于原点对称的区间，先取 $c=\frac{a+b}{2},\ l=\frac{b-a}{2},\ t=x-c$，把区间平移成 $[-l,l]$，再对变量 $t$ 使用上面的公式；也就是说，一般区间题应先平移，再缩放。

\begin{note}[先判延拓方式]
半区间题最容易错的不是积分，而是选错延拓。题目若要求正弦级数，就必须作奇延拓；若要求余弦级数，就必须作偶延拓；若题目已指定周期，就必须按指定周期展开，不能私自换成别的周期；若题目只说“作傅里叶展开”，则要先看原区间是 $[-l,l]$ 还是 $[0,l]$，不要把几种情形混在一起。
\end{note}

\subsubsection{典型展开：半区间与区间缩放}

\begin{example}[二次函数在对称区间上的傅里叶展开]
因为 $x^2$ 为偶函数，所以 $b_n=0$。计算得
\begin{align*}
a_0&=\frac{2\pi^2}{3},\qquad a_n=\frac{4(-1)^n}{n^2},\\
x^2&=\frac{\pi^2}{3}+4\sum_{n=1}^{\infty}\frac{(-1)^n}{n^2}\cos nx,\qquad -\pi\le x\le \pi.
\end{align*}
取 $x=\pi$ 与 $x=0$ 可得到
\[
\sum_{n=1}^{\infty}\frac1{n^2}=\frac{\pi^2}{6},\qquad
\sum_{n=1}^{\infty}\frac{(-1)^{n-1}}{n^2}=\frac{\pi^2}{12},\qquad
\sum_{n=1}^{\infty}\frac{1}{(2n-1)^2}=\frac{\pi^2}{8}.
\]
\end{example}

\begin{example}[半区间函数分别展开成余弦级数与正弦级数]
设 $f(x)=\frac{x^2}{4}-\frac{\pi x}{2}$。
\begin{compactenum}
\item 作偶延拓，得到余弦级数。计算得
\[
a_0=\frac{2}{\pi}\int_{0}^{\pi}\left(\frac{x^2}{4}-\frac{\pi x}{2}\right)\,\mathrm{d}x=-\frac{\pi^2}{3},\qquad
a_n=\frac{2}{\pi}\int_{0}^{\pi}\left(\frac{x^2}{4}-\frac{\pi x}{2}\right)\cos nx\,\mathrm{d}x=\frac{1}{n^2}.
\]
因此 \(\displaystyle \frac{x^2}{4}-\frac{\pi x}{2}=-\frac{\pi^2}{6}+\sum_{n=1}^{\infty}\frac{\cos nx}{n^2},\quad x\in[0,\pi]\)。
    \item 作奇延拓，得到正弦级数：
    \(\displaystyle b_n=\frac{2}{\pi}\int_{0}^{\pi}\left(\frac{x^2}{4}-\frac{\pi x}{2}\right)\sin nx\,\mathrm{d}x=\frac{(-1)^n\pi}{2n}+\frac{(-1)^n-1}{n^3\pi}\)。于是
    \[
    \frac{x^2}{4}-\frac{\pi x}{2}
    =\sum_{n=1}^{\infty}\left[\frac{(-1)^n\pi}{2n}+\frac{(-1)^n-1}{n^3\pi}\right]\sin nx,\qquad x\in[0,\pi).
    \]
\end{compactenum}
\end{example}

\begin{example}[绝对值函数按周期 $2$ 展开]
把 $f(x)=2+|x|$ 展开成以 $2$ 为周期的傅里叶级数。函数为偶函数，所以只含余弦项。计算得 \(\displaystyle a_0=5,\quad b_n=0,\quad a_n=\frac{2(\cos n\pi-1)}{n^2\pi^2}\)。于是
\[
2+|x|
=\frac{5}{2}+\sum_{n=1}^{\infty}\frac{2(\cos n\pi-1)}{n^2\pi^2}\cos n\pi x
=\frac{5}{2}-\frac{4}{\pi^2}\sum_{k=0}^{\infty}\frac{\cos(2k+1)\pi x}{(2k+1)^2},
\]
其中 \(x\in[-1,1]\)。
\end{example}

\begin{example}[分段函数在区间 $(0,4)$ 上展开成余弦级数]
设 \(f(x)=x\ (0<x\le 2)\)，且 \(f(x)=2-\dfrac{x}{2}\ (2<x<4)\)。
作偶延拓并取最小正周期 $4$，则 $l=2$。计算得 \(\displaystyle a_0=1,\quad a_{2k}=0,\quad a_{2k-1}=-\frac{4}{(2k-1)^2\pi^2}\)。故
\[
f(x)=\frac12-\frac{4}{\pi^2}\sum_{n=1}^{\infty}\frac{1}{(2n-1)^2}\cos\frac{(2n-1)\pi x}{2},\qquad 0<x<4.
\]
\end{example}

\begin{example}[线性函数按周期 $\pi$ 延拓]
可把它视作区间 $\left[-\dfrac{\pi}{2},\dfrac{\pi}{2}\right]$ 上的函数
\(g(x)=x+\pi\ (-\dfrac{\pi}{2}\le x<0)\)，且 \(g(x)=x\ (0<x\le \dfrac{\pi}{2})\)，
再以周期 $\pi$ 延拓。计算得 $a_0=\pi,\ a_n=0,\ b_n=-\frac{1}{n}\ (n\ge 1)$，从而
\(\displaystyle x=\frac{\pi}{2}-\sum_{n=1}^{\infty}\frac{\sin 2nx}{n},\qquad x\in(0,\pi)\)，
在 $x=0,\pi$ 处级数收敛到 $\dfrac{\pi}{2}$。
\end{example}

\subsection{傅里叶展开专题三：复数形式与系数互化}

\subsubsection{复数形式与复系数公式}

把三角形式中的余弦、正弦用欧拉公式改写，就得到傅里叶级数的复数形式。它把两套实系数统一成一套复系数，常用于推导公式和处理周期为 $2l$ 的情形。

\begin{theorem}[欧拉公式]
\(\displaystyle \cos nx=\frac{\mathrm{e}^{inx}+\mathrm{e}^{-inx}}{2},\qquad
\sin nx=\frac{\mathrm{e}^{inx}-\mathrm{e}^{-inx}}{2i}\)。
\end{theorem}

\begin{theorem}[傅里叶级数的复数形式]
若 \(\displaystyle f(x)\sim \frac{a_0}{2}+\sum_{n=1}^{\infty}(a_n\cos nx+b_n\sin nx)\)，
则可写成 $f(x)\sim \sum_{n=-\infty}^{+\infty}c_n\mathrm{e}^{inx}$，
其中 \(\displaystyle c_0=\frac{a_0}{2},\qquad
c_n=\frac{a_n-ib_n}{2},\qquad
c_{-n}=\frac{a_n+ib_n}{2}\quad(n\ge 1)\)。
反过来 $a_n=c_n+c_{-n},\ b_n=i(c_n-c_{-n})\ (n\ge 1)$。若 $f$ 为实值函数，还满足 $c_{-n}=\overline{c_n}$。
进一步有 \(\displaystyle c_n=\frac{1}{2\pi}\int_{-\pi}^{\pi}f(x)\mathrm{e}^{-inx}\,\mathrm{d}x,\qquad n\in\mathbb Z\)。
若周期改为 $2l$，则 \(\displaystyle f(x)\sim \sum_{n=-\infty}^{+\infty}c_n\mathrm{e}^{i n\pi x/l},
\qquad
c_n=\frac{1}{2l}\int_{-l}^{l}f(x)\mathrm{e}^{-i n\pi x/l}\,\mathrm{d}x\)。
\end{theorem}

\subsubsection{典型展开：复傅里叶级数}

\begin{example}[矩形波的复傅里叶展开]
设矩形波在一个周期 $[-l,l]$ 内满足 \(f(x)=E\ (-\tau<x<\tau)\)，其余部分 \(f(x)=0\)，其中 \(l>\tau>0\)。
由复系数公式得 \(\displaystyle c_0=\frac{1}{2l}\int_{-\tau}^{\tau}E\,\mathrm{d}x=\frac{E\tau}{l}\)，且 \(\displaystyle c_n=\frac{1}{2l}\int_{-\tau}^{\tau}E\mathrm{e}^{-i n\pi x/l}\,\mathrm{d}x=\frac{E}{n\pi}\sin\frac{n\pi\tau}{l}\ (n\ne 0)\)。因此
\[
f(x)=\frac{E\tau}{l}
+\sum_{\substack{n=-\infty\\ n\ne 0}}^{+\infty}
\frac{E}{n\pi}\sin\frac{n\pi\tau}{l}\,
\mathrm{e}^{i n\pi x/l},
\qquad -l\le x\le l,\ x\ne \pm\tau.
\]
在间断点 $x=\pm\tau$ 处，级数收敛到 $\frac{E}{2}$。
\end{example}

\section{习题}

\subsection{第一节习题：数项级数的收敛与发散}
\subsubsection*{A组}

\begin{exercise}[基本概念]
回答下列问题：
\begin{shortexenum}
    \shortitem{1}{数项级数的部分和是什么？} &
    \shortitem{2}{数项级数的收敛与发散怎样定义？} \\
    \shortitem{3}{若级数为 $\sum_{n=1}^{\infty}a_n$，其部分和 $S_n=\sum_{k=1}^{n}a_k\ (n=1,2,\cdots)$ 是一个数列。反过来，你能用 $S_n$ 把原级数表示出来吗？} &
    \shortitem{4}{一般项 $a_n\to 0\ (n\to\infty)$ 是不是级数 $\sum_{n=1}^{\infty}a_n$ 收敛的充分条件？若不是，请举一例说明。}
\end{shortexenum}
\end{exercise}
\begin{solution}
\begin{shortanswerenum}
    \shortansweritem{1}{级数 $\sum_{n=1}^{\infty}a_n$ 的前 $n$ 项和 \(\displaystyle S_n=a_1+a_2+\cdots+a_n\) 称为它的部分和。}
    \shortansweritem{2}{若部分和数列 $\{S_n\}$ 有极限 $S$，则称级数收敛，且和为 $S$；若 $\{S_n\}$ 无极限，则称级数发散。}
    \shortansweritem{3}{有 \(\displaystyle a_1=S_1,\qquad a_n=S_n-S_{n-1}\ (n\ge 2)\)，因而原级数完全由部分和数列决定。}
    \shortansweritem{4}{不是。反例是调和级数 \(\displaystyle \sum_{n=1}^{\infty}\frac{1}{n}\)，它满足 $\frac1n\to 0$，但级数发散。}
\end{shortanswerenum}
\end{solution}

\begin{exercise}[裂项求和与几何级数]
\begin{compactenum}
    \item 求下列级数的和：
    \begin{shortexenum}
        \shortitem{1}{$\displaystyle \sum_{n=1}^{\infty}\frac{1}{4n^2-1}$；} &
        \shortitem{2}{$\displaystyle \sum_{n=1}^{\infty}\frac{1}{(3n-2)(3n+1)}$。}
    \end{shortexenum}
    \item 证明级数 \(\displaystyle \left(\frac12+\frac13\right)+\left(\frac1{2^2}+\frac1{3^2}\right)+\cdots+\left(\frac1{2^n}+\frac1{3^n}\right)+\cdots\) 收敛，并求其和。
\end{compactenum}
\end{exercise}
\begin{solution}
\begin{compactanswerenum}
    \item
    \begin{compactenum}
        \item
        \(\displaystyle \frac{1}{4n^2-1}=\frac{1}{(2n-1)(2n+1)}=\frac12\left(\frac{1}{2n-1}-\frac{1}{2n+1}\right)\)，
        故 \(\displaystyle \sum_{n=1}^{N}\frac{1}{4n^2-1}=\frac12\left(1-\frac{1}{2N+1}\right)\)。
        令 $N\to\infty$，得 $\sum_{n=1}^{\infty}\frac{1}{4n^2-1}=\frac12$。
        \item
        \(\displaystyle \frac{1}{(3n-2)(3n+1)}
        =\frac13\left(\frac{1}{3n-2}-\frac{1}{3n+1}\right),\quad
        \sum_{n=1}^{N}\frac{1}{(3n-2)(3n+1)}
        =\frac13\left(1-\frac{1}{3N+1}\right)\)。
        于是 $\sum_{n=1}^{\infty}\frac{1}{(3n-2)(3n+1)}=\frac13$。
    \end{compactenum}
    \item 该级数可拆成两个等比级数：
    \(\displaystyle \sum_{n=1}^{\infty}\left(\frac1{2^n}+\frac1{3^n}\right)
    =\sum_{n=1}^{\infty}\frac1{2^n}+\sum_{n=1}^{\infty}\frac1{3^n}
    =1+\frac12=\frac32\)。
\end{compactanswerenum}
\end{solution}

\begin{exercise}[和式的线性关系]
\begin{shortsingleenum}
    \shortsingleitem{1}{证明：若 $\sum_{n=1}^{\infty}a_n$ 收敛，但 $\sum_{n=1}^{\infty}b_n$ 发散，则 $\sum_{n=1}^{\infty}(a_n+b_n)$ 发散。}
    \shortsingleitem{2}{如果 $\sum_{n=1}^{\infty}a_n$ 与 $\sum_{n=1}^{\infty}b_n$ 都发散，试问 $\sum_{n=1}^{\infty}(a_n+b_n)$ 一定发散吗？如果这里的 $a_n,b_n\ (n=1,2,\cdots)$ 都是非负数，则能得出什么结论？}
\end{shortsingleenum}
\end{exercise}
\begin{solution}
\begin{compactanswerenum}
    \item 反设 $\sum(a_n+b_n)$ 收敛。由于 $\sum a_n$ 已收敛，两式相减可得
    \(\displaystyle \sum_{n=1}^{\infty}b_n=\sum_{n=1}^{\infty}(a_n+b_n)-\sum_{n=1}^{\infty}a_n\)
    也收敛，这与题设矛盾，故 $\sum(a_n+b_n)$ 必发散。
    \item 不一定。取反例 \(\displaystyle a_n=1,\ b_n=-1\)。则 $\sum a_n,\sum b_n$ 都发散，但 \(\displaystyle \sum_{n=1}^{\infty}(a_n+b_n)=\sum_{n=1}^{\infty}0\)
    收敛。若 $a_n,b_n\ge 0$，则两个部分和数列都单调增加且无界，于是 \(\displaystyle \sum_{n=1}^{\infty}(a_n+b_n)\)
    的部分和是两个无界增数列之和，因此也无界，从而必发散。
\end{compactanswerenum}
\end{solution}

\subsubsection*{B组}

\begin{exercise}[进一步的求和]
求下列级数的和：
\begin{shortexenum}
    \shortitem{1}{$\displaystyle \sum_{n=1}^{\infty}\frac{(-1)^{n-1}}{n(n+2)}$；} &
    \shortitem{2}{$\displaystyle \sum_{n=1}^{\infty}\arctan\frac{1}{2n^2}$。}
\end{shortexenum}
\end{exercise}
\begin{solution}
\begin{compactanswerenum}
    \item \(\displaystyle \frac{1}{n(n+2)}=\frac12\left(\frac1n-\frac1{n+2}\right)\)，所以
    从第三项起逐项抵消，故 \(\displaystyle \sum_{n=1}^{\infty}\frac{(-1)^{n-1}}{n(n+2)}
    =\frac12\left(1-\frac12\right)=\frac14\)。
    \item 利用反正切差角公式，
    \(\displaystyle \arctan\frac{1}{2n^2}
    =\arctan\frac{1}{2n-1}-\arctan\frac{1}{2n+1}\)，故
    \(\displaystyle \sum_{n=1}^{N}\arctan\frac{1}{2n^2}
    =\arctan 1-\arctan\frac{1}{2N+1}\)。
    令 $N\to\infty$，得 \(\displaystyle \sum_{n=1}^{\infty}\arctan\frac{1}{2n^2}=\frac{\pi}{4}\)。
\end{compactanswerenum}
\end{solution}

\begin{exercise}[用 Cauchy 准则判别收敛性]
利用柯西准则判别下列级数的收敛性：
\begin{shortexenum}
    \shortitem{1}{$\displaystyle \frac{\sin x}{2}+\frac{\sin 2x}{2^2}+\cdots+\frac{\sin nx}{2^n}+\cdots$；} &
    \shortitem{2}{$\displaystyle 1+\frac12-\frac13+\frac14+\frac15-\frac16+\cdots$。}
\end{shortexenum}
\end{exercise}
\begin{solution}
\begin{compactenum}
    \item 设 \(\displaystyle u_n=\frac{\sin nx}{2^n}\)，则
    \(\displaystyle \left|\sum_{k=n+1}^{m}u_k\right|
    \le \sum_{k=n+1}^{m}\frac{1}{2^k}<\frac{1}{2^n}\)。
    当 $n\to\infty$ 时右端趋于 $0$，故满足 Cauchy 准则，所以该级数对任意实数 $x$ 都收敛。
    \item 把该级数按三项一组写成
    \[
    \left(1+\frac12-\frac13\right)+\left(\frac14+\frac15-\frac16\right)+\cdots
    +\left(\frac{1}{3k-2}+\frac{1}{3k-1}-\frac{1}{3k}\right)+\cdots
    \]
    第 $k$ 组满足 \(\displaystyle \frac{1}{3k-2}+\frac{1}{3k-1}-\frac{1}{3k}
    >\frac{1}{3k}\)。
    因而任取大尾和，都可以从下面估计为一段正的调和级数尾和，不可能趋于 $0$，故不满足 Cauchy 准则。于是原级数发散。
\end{compactenum}
\end{solution}

\subsection{第二节习题：正项级数}
\subsubsection*{A组}

\begin{exercise}[基本概念]
回答下列问题：
\begin{shortexenum}
    \shortitem{1}{什么叫正项级数？它的部分和序列有什么特点？} &
    \shortitem{2}{正项级数收敛的充要条件是什么？} \\
    \shortitem{3}{$p$ 级数 $\displaystyle \sum_{n=1}^{\infty}\frac1{n^p}$ 在什么条件下收敛？在什么条件下发散？} &
    \shortitem{4}{设 $a_n\le b_n\ (n=1,2,\cdots)$，如果 $\sum_{n=1}^{\infty}b_n$ 收敛，能否断言 $\sum_{n=1}^{\infty}a_n$ 也收敛？试举例说明。}
\end{shortexenum}
\end{exercise}
\begin{solution}
\begin{shortanswerenum}
    \shortansweritem{1}{若 $a_n\ge 0$，则 $\sum a_n$ 为正项级数。它的部分和 \(\displaystyle S_n=\sum_{k=1}^{n}a_k\) 单调不减。}
    \shortansweritem{2}{正项级数收敛当且仅当其部分和数列有上界。}
    \shortansweritem{3}{\(\displaystyle \sum_{n=1}^{\infty}\frac1{n^p}\) 在 $p>1$ 时收敛，在 $p\le 1$ 时发散。}
    \shortansweritem{4}{不能。比较判别法要求被比较的级数也从某项起非负。反例：\(\displaystyle a_n=-1,\qquad b_n=\frac1{n^2}\)。虽然 $a_n\le b_n$ 且 $\sum b_n$ 收敛，但 $\sum a_n$ 发散。}
\end{shortanswerenum}
\end{solution}

\begin{exercise}[估阶与比较判别法的极限形式]
讨论下列级数的敛散性：
\begin{compactenum}
    \item 可直接普通比较，或先估阶再用比较判别法的极限形式：
    \begin{shortexenum}
        \shortitem{1}{$\displaystyle \sum_{n=1}^{\infty}\frac{1}{n^2+3}$；} &
        \shortitem{2}{$\displaystyle \sum_{n=1}^{\infty}\frac{\sin^2 n}{n^2}$；} \\
        \shortitem{3}{$\displaystyle \sum_{n=1}^{\infty}\frac{n+2}{(n+1)\sqrt n}$；} &
        \shortitem{4}{$\displaystyle \sum_{n=1}^{\infty}\frac1{n2^n}$。}
    \end{shortexenum}
    \item 先估阶为标准模型 \(b_n\)，再用比较判别法的极限形式：
    \begin{shortexenum}
        \shortitem{1}{$\displaystyle \sum_{n=1}^{\infty}\frac{5n+1}{(n+2)n^2}$；} &
        \shortitem{2}{$\displaystyle \sum_{n=1}^{\infty}\frac{2^n+n}{3^n}$；} \\
        \shortitem{3}{$\displaystyle \sum_{n=1}^{\infty}\sin\frac{\pi}{2^n}$；} &
        \shortitem{4}{$\displaystyle \sum_{n=1}^{\infty}\frac{\left(1+\frac1n\right)^n}{n^2}$；} \\
        \shortitem{5}{$\displaystyle \sum_{n=1}^{\infty}\frac1{1+a^n}\ (a>0)$；} &
        \shortitem{6}{$\displaystyle \sum_{n=1}^{\infty}\frac1{n\sqrt n}$。}
    \end{shortexenum}
\end{compactenum}
\end{exercise}
\begin{solution}
\begin{compactanswerenum}
    \item
    \begin{compactenum}
        \item 取 \(b_n=\dfrac1{n^2}\)，则
        \[
        \lim_{n\to\infty}\frac{\frac1{n^2+3}}{\frac1{n^2}}
        =\lim_{n\to\infty}\frac{n^2}{n^2+3}=1.
        \]
        因为 \(\sum 1/n^2\) 收敛，所以由比较判别法的极限形式知原级数收敛。
        \item 因为 \(\displaystyle 0\le \frac{\sin^2 n}{n^2}\le \frac1{n^2}\)，所以原级数收敛。
        \item 取 \(b_n=\dfrac1{\sqrt n}\)，则
        \[
        \lim_{n\to\infty}\frac{\frac{n+2}{(n+1)\sqrt n}}{\frac1{\sqrt n}}
        =\lim_{n\to\infty}\frac{n+2}{n+1}=1.
        \]
        而 $\sum 1/\sqrt n$ 发散，所以由比较判别法的极限形式知原级数发散。
        \item \(\displaystyle 0<\frac1{n2^n}\le \frac1{2^n}\)，故原级数收敛。
    \end{compactenum}
    \item
    \begin{compactenum}
        \item 取 \(b_n=\dfrac1{n^2}\)，则
        \[
        \lim_{n\to\infty}\frac{\frac{5n+1}{(n+2)n^2}}{\frac1{n^2}}
        =\lim_{n\to\infty}\frac{5n+1}{n+2}=5.
        \]
        因为 \(\sum 1/n^2\) 收敛，所以由比较判别法的极限形式知原级数收敛。
        \item 取 \(b_n=\left(\dfrac23\right)^n\)，则
        \[
        \lim_{n\to\infty}\frac{\frac{2^n+n}{3^n}}{\left(\frac23\right)^n}
        =\lim_{n\to\infty}\left(1+\frac{n}{2^n}\right)=1.
        \]
        因为几何级数 \(\sum (2/3)^n\) 收敛，所以由比较判别法的极限形式知原级数收敛。
        \item 取 \(b_n=\dfrac{\pi}{2^n}\)，则
        \[
        \lim_{n\to\infty}\frac{\sin\frac{\pi}{2^n}}{\frac{\pi}{2^n}}=1.
        \]
        因为 \(\sum \pi/2^n\) 收敛，所以由比较判别法的极限形式知原级数收敛。
        \item 取 \(b_n=\dfrac1{n^2}\)。因为 \(\displaystyle \left(1+\frac1n\right)^n\to e\)，所以
        \[
        \lim_{n\to\infty}\frac{\frac{(1+1/n)^n}{n^2}}{\frac1{n^2}}=e.
        \]
        因为 \(\sum 1/n^2\) 收敛，所以由比较判别法的极限形式知原级数收敛。
        \item 若 $a>1$，取 \(b_n=\dfrac1{a^n}\)，则
        \[
        \lim_{n\to\infty}\frac{\frac1{1+a^n}}{\frac1{a^n}}
        =\lim_{n\to\infty}\frac{a^n}{1+a^n}=1,
        \]
        因为几何级数 \(\sum 1/a^n\) 收敛，所以由比较判别法的极限形式知原级数收敛；若 $0<a\le 1$，则通项不趋于 $0$，故发散。即原级数当且仅当 $a>1$ 时收敛。
        \item 取 \(b_n=\dfrac1{n^{3/2}}\)。因为
        \[
        \frac1{n\sqrt n}=\frac1{n^{3/2}},\qquad
        \lim_{n\to\infty}\frac{\frac1{n\sqrt n}}{\frac1{n^{3/2}}}=1,
        \]
        而 \(p=\frac32>1\) 的 \(p\) 级数 \(\sum 1/n^{3/2}\) 收敛，所以由比较判别法的极限形式知原级数收敛。
    \end{compactenum}
\end{compactanswerenum}
\end{solution}

\begin{exercise}[比值判别法与根值判别法]
判别下列级数的敛散性：
\begin{compactenum}
    \item 用比值判别法：
    \begin{shortexenum}
        \shortitem{1}{$\displaystyle \sum_{n=1}^{\infty}np^n\ (0<p<1)$；} &
        \shortitem{2}{$\displaystyle \sum_{n=1}^{\infty}\frac{x^n}{n!}\ (x>0)$；} \\
        \shortitem{3}{$\displaystyle \sum_{n=1}^{\infty}\frac{n^2}{n!}$；} &
        \shortitem{4}{$\displaystyle \sum_{n=1}^{\infty}\frac{n!}{n^n}$；} \\
        \shortitem{5}{$\displaystyle \sum_{n=1}^{\infty}\frac{2^nn!}{n^n}$；} &
        \shortitem{6}{$\displaystyle \sum_{n=1}^{\infty}\frac{3^nn!}{n^n}$。}
    \end{shortexenum}
    \item 用根值判别法：
    \begin{shortexenum}
        \shortitem{1}{$\displaystyle \sum_{n=1}^{\infty}\frac1{n^n}$；} &
        \shortitem{2}{$\displaystyle \sum_{n=1}^{\infty}\frac{2^n}{(n+1)^n}$；} \\
        \shortitem{3}{$\displaystyle \sum_{n=1}^{\infty}\frac{5^n}{n^{n+1}}$；} &
        \shortitem{4}{$\displaystyle \sum_{n=1}^{\infty}\frac{n^2}{\left(n+\frac1n\right)^n}$。}
    \end{shortexenum}
\end{compactenum}
\end{exercise}
\begin{solution}
\begin{compactanswerenum}
    \item
    \begin{compactenum}
        \item \(\displaystyle \frac{a_{n+1}}{a_n}=\frac{(n+1)p^{n+1}}{np^n}=p\left(1+\frac1n\right)\to p<1\)，故收敛。
        \item \(\displaystyle \frac{a_{n+1}}{a_n}=\frac{x^{n+1}/(n+1)!}{x^n/n!}=\frac{x}{n+1}\to 0\)，故收敛。
        \item \(\displaystyle \frac{a_{n+1}}{a_n}=\frac{(n+1)^2/(n+1)!}{n^2/n!}=\frac{n+1}{n^2}\to 0\)，故收敛。
        \item \(\displaystyle \frac{a_{n+1}}{a_n}
        =\frac{(n+1)!}{(n+1)^{n+1}}\cdot \frac{n^n}{n!}
        =\left(\frac{n}{n+1}\right)^n\to \frac1e<1\)，故收敛。
        \item \(\displaystyle \frac{a_{n+1}}{a_n}=2\left(\frac{n}{n+1}\right)^n\to \frac2e<1\)，故收敛。
        \item \(\displaystyle \frac{a_{n+1}}{a_n}=3\left(\frac{n}{n+1}\right)^n\to \frac3e>1\)，故发散。
    \end{compactenum}
    \item
    \begin{compactenum}
        \item
        \(\displaystyle \sqrt[n]{a_n}=\frac1n\to 0\)，故收敛。
        \item
        \(\displaystyle \sqrt[n]{a_n}=\frac{2}{n+1}\to 0\)，故收敛。
        \item
        \(\displaystyle \sqrt[n]{a_n}=\frac{5}{n^{1+1/n}}\to 0\)，故收敛。
        \item
        \(\displaystyle \sqrt[n]{a_n}=\frac{n^{2/n}}{n+\frac1n}\to 0\)，故收敛。
    \end{compactenum}
\end{compactanswerenum}
\end{solution}

\subsubsection*{B组}

\begin{exercise}[综合判别]
用所学过的判别法判别下列级数的敛散性：
\begin{shortexenum}
    \shortitem{1}{$\displaystyle \sum_{n=1}^{\infty}n\left(\frac34\right)^n$；} &
    \shortitem{2}{$\displaystyle \sum_{n=1}^{\infty}\frac{1+\cos n}{n^2}$；} \\
    \shortitem{3}{$\displaystyle \sum_{n=2}^{\infty}\frac1{(\ln n)^k}\ (k>0)$；} &
    \shortitem{4}{$\displaystyle \sum_{n=2}^{\infty}\frac1{n(\ln n)^2}$；} \\
    \shortitem{5}{$\displaystyle \sum_{n=1}^{\infty}\frac{(n!)^2}{2n^2}$；} &
    \shortitem{6}{$\displaystyle \sum_{n=1}^{\infty}\sqrt{\frac{n+1}{n}}$；} \\
    \shortitem{7}{$\displaystyle \sum_{n=1}^{\infty}\frac1{n^{1/\ln n}}$；} &
    \shortitem{8}{$\displaystyle \sum_{n=2}^{\infty}\frac1{(\ln n)^{\ln n}}$；} \\
    \shortitem{9}{$\displaystyle \sum_{n=1}^{\infty}\frac{a^n}{1+a^{2n}}\ (a\ge 0)$；} &
    \shortitem{10}{$\displaystyle \sum_{n=1}^{\infty}\left(\frac{b}{a_n}\right)^n$，其中 $a_n\to a\ (n\to\infty)$，且 $a_n,b,a$ 均为正数。}
\end{shortexenum}
\end{exercise}
\begin{solution}
\begin{compactanswerenum}
    \item 收敛。可用比值判别法，极限为 $\frac34$。
    \item 收敛，因为 \(\displaystyle 0\le \frac{1+\cos n}{n^2}\le \frac2{n^2}\)。
    \item 发散。对充分大的 $n$ 有 $(\ln n)^k<n$，故 \(\displaystyle \frac1{(\ln n)^k}>\frac1n\)，与调和级数比较可知发散。
    \item 收敛。它是 Bertrand 型级数，亦可由积分判别法得知收敛。
    \item 发散，因为通项
    \(\displaystyle \frac{(n!)^2}{2n^2}\not\to 0\)。
    \item 发散，因为
    \(\displaystyle \sqrt{\frac{n+1}{n}}\to 1\neq 0\)。
    \item 发散，因为
    \(\displaystyle n^{1/\ln n}=e,\qquad \frac1{n^{1/\ln n}}=\frac1e\)。
    \item 收敛，因为
    \(\displaystyle (\ln n)^{\ln n}=e^{\ln n\cdot \ln\ln n}=n^{\ln\ln n}\)，
    当 $n$ 足够大时有 $n^{\ln\ln n}>n^2$，故与 $\sum 1/n^2$ 比较收敛。
    \item 当 $a=1$ 时通项恒为 $1/2$，发散；当 $a=0$ 时通项恒为 $0$，收敛。若 \(0<a<1\)，取 \(b_n=a^n\)，则
    \[
    \lim_{n\to\infty}\frac{\frac{a^n}{1+a^{2n}}}{a^n}
    =\lim_{n\to\infty}\frac1{1+a^{2n}}=1,
    \]
    而几何级数 \(\sum a^n\) 收敛，故原级数收敛。若 \(a>1\)，取 \(b_n=\dfrac1{a^n}\)，则
    \[
    \lim_{n\to\infty}\frac{\frac{a^n}{1+a^{2n}}}{\frac1{a^n}}
    =\lim_{n\to\infty}\frac{a^{2n}}{1+a^{2n}}=1,
    \]
    而几何级数 \(\sum 1/a^n\) 收敛，故原级数收敛。故当且仅当 $a\neq 1$ 时收敛。
    \item 用根值判别法：
    \(\displaystyle \sqrt[n]{\left(\frac{b}{a_n}\right)^n}=\frac{b}{a_n}\to \frac{b}{a}\)。
    所以当 $b<a$ 时收敛，当 $b>a$ 时发散；当 $b=a$ 时通项趋于 $1$，也发散。
\end{compactanswerenum}
\end{solution}

\begin{exercise}[用 Taylor 公式估计无穷小的阶]
利用 Taylor 公式估计下列级数 $\sum_{n=1}^{\infty}a_n$ 中无穷小量 $a_n$ 的阶，从而判别级数的敛散性：
\begin{shortexenum}
    \shortitem{1}{$\displaystyle \sum_{n=1}^{\infty}2^n\sin\frac{\pi}{3^n}$；} &
    \shortitem{2}{$\displaystyle \sum_{n=1}^{\infty}\left(\sqrt{n+1}-\sqrt[4]{n^2+n+1}\right)$。}
\end{shortexenum}
\end{exercise}
\begin{solution}
\begin{compactanswerenum}
    \item 取 \(\displaystyle b_n=\pi\left(\frac23\right)^n\)。由 Taylor 公式 \(\sin t=t+O(t^3)\) 得
    \[
    \lim_{n\to\infty}
    \frac{2^n\sin\frac{\pi}{3^n}}{\pi\left(\frac23\right)^n}
    =
    \lim_{n\to\infty}
    \frac{\sin\frac{\pi}{3^n}}{\frac{\pi}{3^n}}
    =1.
    \]
    因为几何级数 \(\displaystyle \sum \pi(2/3)^n\) 收敛，所以由比较判别法的极限形式知原级数收敛。
    \item 展开得
\begin{align*}
\sqrt{n+1}
&=\sqrt n\left(1+\frac1n\right)^{1/2}
=\sqrt n+\frac1{2\sqrt n}+O\left(\frac1{n^{3/2}}\right),\\
\sqrt[4]{n^2+n+1}
&=\sqrt n\left(1+\frac1n+\frac1{n^2}\right)^{1/4}
=\sqrt n+\frac1{4\sqrt n}+O\left(\frac1{n^{3/2}}\right),\\
\sqrt{n+1}-\sqrt[4]{n^2+n+1}
&=\frac1{4\sqrt n}+O\left(\frac1{n^{3/2}}\right).
    \end{align*}
    取 \(\displaystyle b_n=\frac1{4\sqrt n}\)，则
    \[
    \lim_{n\to\infty}
    \frac{\sqrt{n+1}-\sqrt[4]{n^2+n+1}}{\frac1{4\sqrt n}}
    =1.
    \]
    因为 \(\displaystyle \sum 1/\sqrt n\) 发散，所以由比较判别法的极限形式知原级数发散。
\end{compactanswerenum}
\end{solution}

\begin{exercise}[证明题]
\begin{shortsingleenum}
    \shortsingleitem{1}{若正项级数 $\displaystyle \sum_{n=1}^{\infty}a_n$ 收敛，证明 $\displaystyle \sum_{n=1}^{\infty}a_n^2$ 也收敛；但反之不然，举例说明。}
    \shortsingleitem{2}{设级数 $\displaystyle \sum_{n=1}^{\infty}a_n^2$ 与 $\displaystyle \sum_{n=1}^{\infty}b_n^2$ 都收敛，证明 \(\displaystyle \sum_{n=1}^{\infty}|a_nb_n|\) 与 \(\displaystyle \sum_{n=1}^{\infty}(a_n+b_n)^2\) 也收敛。}
    \shortsingleitem{3}{设正数序列 $\{x_n\}$ 单调上升且有界，证明级数 \(\displaystyle \sum_{n=1}^{\infty}\left(1-\frac{x_n}{x_{n+1}}\right)\) 收敛。}
    \shortsingleitem{4}{若正项级数 $\displaystyle \sum_{n=1}^{\infty}a_n$ 收敛，证明正项级数 \(\displaystyle \sum_{n=1}^{\infty}\frac{\sqrt{a_n}}{n}\) 也收敛。}
\end{shortsingleenum}
\end{exercise}
\begin{solution}
\begin{compactanswerenum}
    \item 因 $\sum a_n$ 收敛，所以 $a_n\to 0$。从某项起 \(0<a_n<1\)，于是 \(\displaystyle 0<a_n^2\le a_n\)。用比较判别法知 $\sum a_n^2$ 收敛。反例可取 \(\displaystyle a_n=\frac1n\)，则 $\sum a_n^2=\sum 1/n^2$ 收敛，而 $\sum a_n$ 发散。
    \item 利用不等式 \(\displaystyle 2|a_nb_n|\le a_n^2+b_n^2,\qquad (a_n+b_n)^2\le 2(a_n^2+b_n^2)\)。右端对应的级数都收敛，因此左端对应的级数也收敛。
    \item 因 $\{x_n\}$ 单调上升且有界，设其上界极限为 $L$。又因 $x_n>0$，故 \(\displaystyle 0<1-\frac{x_n}{x_{n+1}}=\frac{x_{n+1}-x_n}{x_{n+1}}\le \frac{x_{n+1}-x_n}{x_1}\)，且 \(\displaystyle \sum_{n=1}^{\infty}(x_{n+1}-x_n)=L-x_1<\infty\)，
    故由比较判别法可知原级数收敛。
    \item 由均值不等式 \(\displaystyle \frac{\sqrt{a_n}}{n}\le \frac12\left(a_n+\frac1{n^2}\right)\)。因为 $\sum a_n$ 与 $\sum 1/n^2$ 都收敛，所以 $\sum \sqrt{a_n}/n$ 也收敛。
\end{compactanswerenum}
\end{solution}

\subsection{第三节习题：一般级数}
\subsubsection*{A组}

\begin{exercise}[基本概念]
回答下列问题：
\begin{shorttrienum}
    \shortitem{1}{什么叫交错级数？什么叫莱布尼兹型级数？} &
    \shortitem{2}{什么是级数的绝对收敛性和条件收敛性？} &
    \shortitem{3}{两级数的乘积是怎样定义的？在什么条件下两级数的乘积级数必定收敛？}
\end{shorttrienum}
\end{exercise}
\begin{solution}
\begin{compactanswerenum}
    \item 形如 \(\displaystyle a_1-a_2+a_3-a_4+\cdots=\sum_{n=1}^{\infty}(-1)^{n-1}a_n,\ a_n\ge 0\) 的级数称为交错级数。若其中 $a_n$ 单调下降且 $a_n\to 0$，就称它为莱布尼兹型级数。
    \item 若 $\sum |a_n|$ 收敛，则称 $\sum a_n$ 绝对收敛；若 $\sum a_n$ 收敛但 $\sum |a_n|$ 发散，则称其条件收敛。
    \item 设 \(\displaystyle \sum_{n=1}^{\infty}a_n\) 与 \(\displaystyle \sum_{n=1}^{\infty}b_n\)，则其 Cauchy 乘积定义为 \(\displaystyle \sum_{n=1}^{\infty}c_n,\ c_n=a_1b_n+a_2b_{n-1}+\cdots+a_nb_1\)。若两个级数都绝对收敛，则乘积级数绝对收敛，且其和等于两级数和的乘积。
\end{compactanswerenum}
\end{solution}

\begin{exercise}[判断正误]
对下列命题，正确的给出证明，错误的举出反例：
\begin{shortexenum}
    \shortitem{1}{若 $a_n>0$，则 $a_1-a_1+a_2-a_2+a_3-a_3+\cdots$ 收敛；} &
    \shortitem{2}{若 $\sum_{n=1}^{\infty}a_n$ 收敛，则 $\sum_{n=1}^{\infty}(-1)^na_n$ 也收敛；}\\
    \shortitem{3}{若 $\sum_{n=1}^{\infty}a_n^2$ 收敛，则 $\sum_{n=1}^{\infty}a_n^3$ 绝对收敛；} &
    \shortitem{4}{若 $\sum_{n=1}^{\infty}a_n$ 收敛，则 $\sum_{n=1}^{\infty}a_n^2$ 收敛；}\\
    \shortitem{5}{若 $\sum_{n=1}^{\infty}a_n$ 不收敛，则 $\sum_{n=1}^{\infty}a_n$ 不绝对收敛；} &
    \shortitem{6}{绝对收敛级数也是条件收敛的；}\\
    \shortitem{7}{若 $\sum_{n=1}^{\infty}a_n$ 不是条件收敛的，则 $\sum_{n=1}^{\infty}a_n$ 发散；} &
    \shortitem{8}{若 $\sum_{n=1}^{\infty}a_n$ 收敛，且 $\displaystyle \lim_{n\to\infty}\frac{b_n}{a_n}=1$，则 $\sum_{n=1}^{\infty}b_n$ 也收敛。}
\end{shortexenum}
\end{exercise}
\begin{solution}
\begin{compactanswerenum}
    \item 错。其部分和满足 \(\displaystyle S_{2n}=0,\qquad S_{2n-1}=a_n\)。若取 $a_n\equiv 1$，则 $S_{2n-1}=1$、$S_{2n}=0$，不收敛。
    \item 错。取 \(\displaystyle a_n=\frac{(-1)^{n-1}}{\sqrt n}\)，则 $\sum a_n$ 收敛，而 \(\displaystyle \sum (-1)^na_n=-\sum \frac1{\sqrt n}\) 发散。
    \item 对。由 $\sum a_n^2$ 收敛知 $a_n\to 0$，故从某项起 $|a_n|<1$，于是 \(\displaystyle |a_n|^3\le a_n^2\)，所以 $\sum |a_n|^3$ 收敛。
    \item 错。仍取 \(\displaystyle a_n=\frac{(-1)^{n-1}}{\sqrt n}\)，则 $\sum a_n$ 收敛，而 \(\displaystyle \sum a_n^2=\sum \frac1n\) 发散。
    \item 对。绝对收敛必推出收敛，所以若原级数不收敛，就不可能绝对收敛。
    \item 错。绝对收敛与条件收敛互斥。
    \item 错。不是条件收敛还可能是绝对收敛。
    \item 错。取 \(\displaystyle a_n=\frac{(-1)^n}{\sqrt n},\qquad b_n=a_n+\frac1n\)。则 \(\displaystyle \frac{b_n}{a_n}=1+\frac{(-1)^n}{\sqrt n}\to 1\)，但 \(\displaystyle \sum b_n=\sum a_n+\sum \frac1n\) 发散。
\end{compactanswerenum}
\end{solution}

\begin{exercise}[判别绝对收敛、条件收敛与发散]
讨论下列级数的敛散性（包括绝对收敛、条件收敛、发散）：
\begin{shortexenum}
    \shortitem{1}{$\displaystyle \sum_{n=1}^{\infty}(-1)^{n-1}\frac1{\sqrt n}$；} &
    \shortitem{2}{$\displaystyle \sum_{n=1}^{\infty}(-1)^{n-1}\frac1{3^n}$；} \\
    \shortitem{3}{$\displaystyle \sum_{n=1}^{\infty}(-1)^{n-1}\frac1{n(n+1)}$；} &
    \shortitem{4}{$\displaystyle \sum_{n=1}^{\infty}(-1)^{n-1}\frac{n}{2n-1}$；} \\
    \shortitem{5}{$\displaystyle \sum_{n=2}^{\infty}(-1)^n\frac1{\ln n}$；} &
    \shortitem{6}{$\displaystyle \sum_{n=1}^{\infty}(-1)^{n-1}\frac{\ln n}{\sqrt n}$。}
\end{shortexenum}
\end{exercise}
\begin{solution}
\begin{compactanswerenum}
    \item 条件收敛。满足 Leibniz 判别法，而绝对值级数 $\sum 1/\sqrt n$ 发散。
    \item 绝对收敛。因为 \(\displaystyle \sum \left|\frac{(-1)^{n-1}}{3^n}\right|=\sum \frac1{3^n}\) 为几何级数。
    \item 绝对收敛。取 \(\displaystyle b_n=\frac1{n^2}\)，则
    \[
    \lim_{n\to\infty}\frac{\frac1{n(n+1)}}{\frac1{n^2}}
    =\lim_{n\to\infty}\frac{n}{n+1}=1.
    \]
    因为 \(\sum 1/n^2\) 收敛，所以绝对值级数收敛。
    \item 发散。其通项
    \(\displaystyle \frac{n}{2n-1}\to \frac12\neq 0\)。
    \item 条件收敛。因为 $1/\ln n$ 单调趋于 $0$，故交错级数收敛；但 $\sum 1/\ln n$ 发散。
    \item 条件收敛。因为 $\ln n/\sqrt n\downarrow 0$（充分大时成立），由 Leibniz 判别法知收敛；而 \(\displaystyle \sum \frac{\ln n}{\sqrt n}\)
    发散，所以不是绝对收敛。
\end{compactanswerenum}
\end{solution}

\begin{exercise}[证明与讨论]
\begin{shortsingleenum}
    \shortsingleitem{1}{证明：若级数 $\displaystyle \sum_{n=1}^{\infty}a_n$ 及 $\displaystyle \sum_{n=1}^{\infty}b_n$ 皆收敛，且 $a_n\le c_n\le b_n\ (n=1,2,\cdots)$，则 $\displaystyle \sum_{n=1}^{\infty}c_n$ 也收敛。若 $\sum a_n$ 与 $\sum b_n$ 皆发散，试问级数 $\sum c_n$ 的收敛性如何？}
    \shortsingleitem{2}{已知 $\displaystyle \sum_{n=1}^{\infty}a_n^2$ 与 $\displaystyle \sum_{n=1}^{\infty}b_n^2$ 皆收敛，证明级数 $\displaystyle \sum_{n=1}^{\infty}a_nb_n$ 绝对收敛。}
    \shortsingleitem{3}{设常数 $k>0$，讨论级数 $\displaystyle \sum_{n=1}^{\infty}(-1)^n\frac{k+n}{n^2}$ 的敛散性（包括绝对收敛与条件收敛）。}
\end{shortsingleenum}
\end{exercise}
\begin{solution}
\begin{compactanswerenum}
    \item 因为 \(\displaystyle \sum a_n,\ \sum b_n\) 都收敛，所以它们的部分和数列都有极限。对任意 $N$，有 \(\displaystyle \sum_{n=1}^{N}a_n\le \sum_{n=1}^{N}c_n\le \sum_{n=1}^{N}b_n\)。
    由夹逼原理可知 $\sum c_n$ 收敛。若 $\sum a_n,\sum b_n$ 都发散，则对 $\sum c_n$ 不能作统一结论。例如取 \(\displaystyle a_n=-1,\ b_n=1,\ c_n=0\)，
    则 $\sum c_n$ 收敛；若取 $c_n=1$，则 $\sum c_n$ 发散。
    \item 由 \(\displaystyle 2|a_nb_n|\le a_n^2+b_n^2\) 可得
    \[
    |a_nb_n|\le \frac12(a_n^2+b_n^2).
    \]
    右端对应级数收敛，故由普通比较判别法知 \(\displaystyle \sum |a_nb_n|\) 收敛。
    \item 原级数可写为 \(\displaystyle \sum_{n=1}^{\infty}(-1)^n\left(\frac1n+\frac{k}{n^2}\right)\)。
    因为 \(\displaystyle \frac1n+\frac{k}{n^2}\downarrow 0\)，所以按 Leibniz 判别法，原级数对任意 $k>0$ 都收敛。再看绝对值级数，取 \(b_n=\dfrac1n\)，则
    \[
    \lim_{n\to\infty}
    \frac{\frac1n+\frac{k}{n^2}}{\frac1n}
    =\lim_{n\to\infty}\left(1+\frac{k}{n}\right)=1.
    \]
    而调和级数 \(\sum 1/n\) 发散，故绝对值级数发散。原级数对一切 $k>0$ 都是条件收敛。
\end{compactanswerenum}
\end{solution}

\subsubsection*{B组}

\begin{exercise}[一个判别定理的应用]
设 $\displaystyle \lim_{n\to\infty}n^pu_n=A$，证明：
\begin{shortexenum}
    \shortitem{1}{当 $p>1$ 时，$\displaystyle \sum_{n=1}^{\infty}u_n$ 绝对收敛；} &
    \shortitem{2}{当 $p=1$ 且 $A\neq 0$ 时，$\displaystyle \sum_{n=1}^{\infty}u_n$ 发散；} \\
    \shortitem{3}{当 $p=1$ 且 $A=0$ 时，$\displaystyle \sum_{n=1}^{\infty}u_n$ 能否收敛？} &
\end{shortexenum}
\end{exercise}
\begin{solution}
\begin{shortanswerenum}
    \shortansweritem{1}{取 \(b_n=1/n^p\)。由 \(n^pu_n\to A\) 得
    \(\displaystyle \lim_{n\to\infty}\frac{|u_n|}{b_n}=|A|\)。若 \(A\ne0\)，由比较判别法的极限形式知 \(\sum |u_n|\) 与 \(\sum1/n^p\) 同敛散；若 \(A=0\)，则 \(|u_n|/b_n\to0\)，仍由比较判别法的极限形式知 \(\sum |u_n|\) 收敛。因 \(p>1\)，故原级数绝对收敛。}
    \shortansweritem{2}{当 $p=1$ 且 $A\neq 0$ 时，取 \(b_n=1/n\)，则
    \(\displaystyle \lim_{n\to\infty}\frac{|u_n|}{b_n}=|A|>0\)。又 \(u_n\) 从某项起与 \(A\) 同号，而调和级数发散，所以 \(\sum u_n\) 发散。}
    \shortansweritem{3}{不能确定。反例一：\(\displaystyle u_n=\frac1{n\ln n}\ (n\ge 2)\)，则 $nu_n\to 0$，但级数发散。反例二：\(\displaystyle u_n=\frac1{n^2}\)，也有 $nu_n\to 0$，但级数收敛。}
\end{shortanswerenum}
\end{solution}

\begin{exercise}[利用比较判别法证明绝对收敛定理]
试利用比较判别法证明：若级数 $\displaystyle \sum_{n=1}^{\infty}|a_n|$ 收敛，则级数 $\displaystyle \sum_{n=1}^{\infty}a_n$ 收敛。
\end{exercise}
\begin{solution}
定义 \(\displaystyle a_n^+=\frac{|a_n|+a_n}{2},\qquad
a_n^-=\frac{|a_n|-a_n}{2}\)，则 \(\displaystyle 0\le a_n^+\le |a_n|,\qquad 0\le a_n^-\le |a_n|\)。
由比较判别法知
\(\displaystyle \sum a_n^+,\ \sum a_n^-\) 都收敛，并且 \(\displaystyle a_n=a_n^+-a_n^-,\quad
\sum a_n=\sum a_n^+-\sum a_n^-\)。
故 $\sum a_n$ 收敛。
\end{solution}

\subsection{第四节习题：函数项级数的基本概念}
\subsubsection*{A组}

\begin{exercise}[基本概念]
回答下列问题：
\begin{shortexenum}
    \shortitem{1}{级数 $\displaystyle \sum_{n=1}^{\infty}u_n(x)$ 在 $I$ 上点态收敛于 $S(x)$ 是什么意思？} &
    \shortitem{2}{什么叫级数 $\displaystyle \sum_{n=1}^{\infty}u_n(x)$ 的收敛域和发散域？} \\
    \shortitem{3}{函数列 $\{f_n(x)\}$ 在 $I$ 上一致收敛于 $f(x)$ 的定义是怎样的？} &
    \shortitem{4}{级数 $\displaystyle \sum_{n=1}^{\infty}u_n(x)$ 在 $I$ 上一致收敛于 $S(x)$ 的定义是怎样的？} \\
    \shortitem{5}{一致收敛性与点态收敛性有什么不同？} &
\end{shortexenum}
\end{exercise}
\begin{solution}
\begin{compactenum}
    \item 对每个固定的 $x\in I$，数项级数 $\sum u_n(x)$ 收敛，且其和为 $S(x)$。
    \item 使 $\sum u_n(x)$ 收敛的点 $x$ 全体构成收敛域；使其发散的点 $x$ 全体构成发散域。
    \item 对任意 $\varepsilon>0$，存在与 $x$ 无关的整数 $N$，使得当 $n>N$ 时，\(\displaystyle |f_n(x)-f(x)|<\varepsilon,\ \forall x\in I\)。
    \item 设部分和为 $S_n(x)$。若对任意 $\varepsilon>0$，存在与 $x$ 无关的整数 $N$，使得当 $n>N$ 时，\(\displaystyle |S_n(x)-S(x)|<\varepsilon,\ \forall x\in I\)，就称级数在 $I$ 上一致收敛于 $S(x)$。
    \item 点态收敛允许 $N$ 依赖于 $x$；一致收敛要求同一个 $N$ 对整个区间都有效，所以一致收敛更强。
\end{compactenum}
\end{solution}

\begin{exercise}[Weierstrass $M$ 判别法]
试利用维尔斯特拉斯 $M$ 判别法证明下列级数在指定区间上的一致收敛性：
\begin{shortexenum}
    \shortitem{1}{$\displaystyle \sum_{n=1}^{\infty}\frac{\cos nx}{2^n},\quad -\infty<x<+\infty$；} &
    \shortitem{2}{$\displaystyle \sum_{n=1}^{\infty}\frac{\sin nx}{n^2},\quad -\infty<x<+\infty$；} \\
    \shortitem{3}{$\displaystyle \sum_{n=1}^{\infty}\frac{x^n}{n^{3/2}},\quad x\in[-1,1]$；} &
    \shortitem{4}{$\displaystyle \sum_{n=1}^{\infty}\frac{(-1)^n(1-\mathrm{e}^{-nx})}{n^2+x^2},\quad 0\le x<+\infty$。}
\end{shortexenum}
\end{exercise}
\begin{solution}
\begin{compactenum}
    \item \(\displaystyle \left|\frac{\cos nx}{2^n}\right|\le \frac1{2^n}\)，而 $\sum 2^{-n}$ 收敛，所以一致收敛。
    \item \(\displaystyle \left|\frac{\sin nx}{n^2}\right|\le \frac1{n^2}\)，且 $\sum 1/n^2$ 收敛，所以一致收敛。
    \item 对 $x\in[-1,1]$，\(\displaystyle \left|\frac{x^n}{n^{3/2}}\right|\le \frac1{n^{3/2}}\)，而 $\sum 1/n^{3/2}$ 收敛，所以一致收敛。
    \item 对任意 $x\ge 0$，\(\displaystyle \left|\frac{(-1)^n(1-\mathrm{e}^{-nx})}{n^2+x^2}\right|\le \frac1{n^2}\)，由 $M$ 判别法知一致收敛。
\end{compactenum}
\end{solution}

\begin{exercise}[按定义讨论一致收敛性]
按定义讨论下列级数在指定区间上的一致收敛性：
\begin{shortexenum}
    \shortitem{1}{$\displaystyle \sum_{n=1}^{\infty}(-1)^{n-1}\frac{x^2}{(1+x^2)^n},\quad -\infty<x<+\infty$；} &
    \shortitem{2}{$\displaystyle \sum_{n=1}^{\infty}(1-x)x^n,\quad 0<x<1$。}
\end{shortexenum}
\end{exercise}
\begin{solution}
\begin{compactenum}
    \item 记 \(\displaystyle a_n(x)=\frac{x^2}{(1+x^2)^n}\ge 0\)。
    对每个固定的 $x$，$\{a_n(x)\}$ 单调下降且趋于 $0$，所以原级数按 Leibniz 判别法逐点收敛。又
    \[
    \sup_{x\in\mathbb R}a_n(x)=\sup_{t\ge 0}\frac{t}{(1+t)^n}
    =\frac{(n-1)^{n-1}}{n^n}\to 0,\qquad
    \sup_{x\in\mathbb R}|R_n(x)|\le \sup_{x\in\mathbb R}a_{n+1}(x)\to 0,
    \]
    所以该级数一致收敛。
    \item 部分和为 \(\displaystyle S_N(x)=(1-x)\sum_{n=1}^{N}x^n=x(1-x^N)\)，故和函数为 $S(x)=x$，余项为 \(\displaystyle S(x)-S_N(x)=x^{N+1}\)。
    于是 \(\displaystyle \sup_{0<x<1}|S(x)-S_N(x)|=\sup_{0<x<1}x^{N+1}=1\)，
    不趋于 $0$，因此不是一致收敛。
\end{compactenum}
\end{solution}

\subsubsection*{B组}

\begin{exercise}[一致收敛的证明]
证明级数 \(\displaystyle \sum_{n=1}^{\infty}x^2\mathrm{e}^{-nx}\) 在 $[0,+\infty)$ 上一致收敛。
\end{exercise}
\begin{solution}
设 \(\displaystyle u_n(x)=x^2\mathrm{e}^{-nx}\)。对固定的 $n$，令 \(\displaystyle \varphi(x)=x^2\mathrm{e}^{-nx}\)，则 \(\displaystyle \varphi'(x)=x\mathrm{e}^{-nx}(2-nx)\)。
所以最大值在 $x=\frac2n$ 处取得，且 \(\displaystyle \sup_{x\ge 0}u_n(x)=\frac{4}{n^2\mathrm{e}^2}\)。而 \(\displaystyle \sum_{n=1}^{\infty}\frac{4}{n^2\mathrm{e}^2}\)
收敛，因此由 $M$ 判别法知原级数在 $[0,+\infty)$ 上一致收敛。
\end{solution}

\begin{exercise}[逐点收敛而不一致收敛]
证明函数列 \(\displaystyle f_n(x)=\frac{nx}{1+n^2x^2}\ (n=1,2,\cdots)\) 在 $[0,+\infty)$ 上逐点收敛于 $0$，但不一致收敛。
\end{exercise}
\begin{solution}
对每个固定的 $x\ge 0$：
\begin{itemize}
    \item 当 $x=0$ 时，$f_n(0)=0$；
    \item 当 $x>0$ 时，\(\displaystyle f_n(x)=\frac{1}{nx+\frac1{nx}}\to 0\)。
\end{itemize}
所以 $f_n(x)\to 0$，即逐点收敛于零函数。

但取 \(\displaystyle x_n=\frac1n\)，则 \(\displaystyle f_n(x_n)=\frac{n\cdot \frac1n}{1+n^2\cdot \frac1{n^2}}=\frac12,\qquad
\sup_{x\ge 0}|f_n(x)-0|\ge \frac12\)，
不趋于 $0$，故不一致收敛。
\end{solution}

\subsection{第五节习题：幂级数及其收敛性}
\subsubsection*{A组}

\begin{exercise}[基本概念]
回答下列问题：
\begin{shortexenum}
    \shortitem{1}{什么是幂级数的收敛半径和收敛域？} &
    \shortitem{2}{幂级数的收敛域有什么特点？} \\
    \shortitem{3}{什么叫做幂级数在其收敛区间中的内闭一致收敛性？} &
    \shortitem{4}{幂级数在其收敛区间内，和函数有哪些重要性质？}
\end{shortexenum}
\end{exercise}
\begin{solution}
\begin{compactenum}
    \item 对幂级数 \(\displaystyle \sum_{n=0}^{\infty}a_n(x-x_0)^n\)，
    存在一个 $R\in[0,+\infty]$，使得当 $|x-x_0|<R$ 时绝对收敛，当 $|x-x_0|>R$ 时发散。这个 $R$ 叫收敛半径；全体收敛点构成收敛域。
    \item 收敛域关于展开中心是一个区间，内部必收敛，外部必发散，两个端点要另行判别。
    \item 对任意 $r$ 满足 $0<r<R$，幂级数在闭区间 $[x_0-r,x_0+r]$ 上一致收敛，这就叫内闭一致收敛性。
    \item 在收敛半径内，和函数连续、可逐项积分、可逐项求导，而且导数级数与积分级数具有同样的收敛半径。
\end{compactenum}
\end{solution}

\begin{exercise}[求收敛半径与收敛域]
求下列幂级数的收敛半径与收敛域：
\begin{shortexenum}
    \shortitem{1}{$\displaystyle \sum_{n=1}^{\infty}nx^n$；} &
    \shortitem{2}{$\displaystyle \sum_{n=0}^{\infty}\frac{x^n}{n!}$；} \\
    \shortitem{3}{$\displaystyle \sum_{n=1}^{\infty}\frac{x^n}{\sqrt n}$；} &
    \shortitem{4}{$\displaystyle \sum_{n=1}^{\infty}\frac{(x-3)^n}{n3^n}$；} \\
    \shortitem{5}{$\displaystyle \sum_{n=0}^{\infty}\frac{2n^2+1}{n^2-5}x^n$；} &
    \shortitem{6}{$\displaystyle \sum_{n=1}^{\infty}\frac{2^nx^n}{n!}$；} \\
    \shortitem{7}{$\displaystyle \sum_{n=0}^{\infty}\frac{(x-4)^n}{2n+1}$；} &
    \shortitem{8}{$\displaystyle \sum_{n=0}^{\infty}\frac{(x+3)^n}{5^n}$；} \\
    \shortitem{9}{$\displaystyle \sum_{n=2}^{\infty}\frac{(x-5)^n}{n\ln n}$；} &
    \shortitem{10}{$\displaystyle \sum_{n=0}^{\infty}n\left(\frac{x-3}{2}\right)^n$；} \\
    \shortitem{11}{$\displaystyle \sum_{n=1}^{\infty}\left(1+\frac1n\right)^{-n^2}x^n$；} &
    \shortitem{12}{$\displaystyle \sum_{n=0}^{\infty}(-1)^n\frac{1\cdot 3\cdots (2n-1)}{1\cdot 2\cdots n}x^{2n}$。}
\end{shortexenum}
\end{exercise}
\begin{solution}
\begin{compactenum}
    \item $\lim\sqrt[n]{n}=1$，故 $R=1$，收敛域为 $(-1,1)$。
    \item $R=+\infty$，在全体实数上收敛。
    \item $\lim\sqrt[n]{1/\sqrt n}=1$，故 $R=1$。端点 $x=1$ 时为 $\sum 1/\sqrt n$ 发散，$x=-1$ 时为交错级数收敛，所以收敛域为 $[-1,1)$？注意右端点是 $1$ 发散，左端点 $-1$ 收敛，故收敛域为 $[-1,1)$。
    \item 视作关于 $x-3$ 的幂级数，$R=3$。端点 $x=6$ 时为 $\sum 1/n$ 发散，$x=0$ 时为 $\sum (-1)^n/n$ 收敛，故收敛域为 $[0,6)$。
    \item 系数极限非零，故 $R=1$。端点需另判；当 $x=\pm 1$ 时通项不趋零，因此收敛域为 $(-1,1)$。
    \item $R=+\infty$，全体实数收敛。
    \item $R=1$。当 $x=5$ 时为 $\sum 1/(2n+1)$ 发散；当 $x=3$ 时为交错级数 $\sum (-1)^n/(2n+1)$ 收敛，故收敛域为 $[3,5)$。
    \item $R=5$。端点 $x=2$ 时为 $\sum (-1)^n$ 发散，$x=-8$ 时为 $\sum 1$ 发散，所以收敛域为 $(-8,2)$。
    \item $R=1$。端点 $x=6$ 时为 $\sum 1/(n\ln n)$ 发散，$x=4$ 时为交错级数 $\sum (-1)^n/(n\ln n)$ 收敛，所以收敛域为 $[4,6)$。
    \item 令 $y=(x-3)/2$，则为 $\sum ny^n$，故 $|y|<1$，即 $|x-3|<2$。端点均因通项不趋零而发散，所以收敛域为 $(1,5)$。
    \item \(\displaystyle \sqrt[n]{a_n}=\left(1+\frac1n\right)^{-n}\to \frac1e\)，故 $R=e$。端点由题面不要求细判，可写收敛半径为 $e$，收敛区间至少为 $(-e,e)$。
    \item 利用 \(\displaystyle \frac{1\cdot 3\cdots (2n-1)}{1\cdot 2\cdots n}
    =\frac{(2n)!}{2^{2n}(n!)^2}\)。由 Stirling 公式可得
    \[
    \lim_{n\to\infty}
    \frac{\frac{(2n)!}{2^{2n}(n!)^2}}{\frac1{\sqrt{\pi n}}}=1.
    \]
    因而关于 \(x^2\) 的收敛半径为 \(1\)，故 \(R=1\)。端点 \(x=\pm 1\) 时为交错级数或正项级数边界情形，可得 \(x=1\) 条件收敛，\(x=-1\) 也收敛，故收敛域为 \([-1,1]\)。
\end{compactenum}
\end{solution}

\begin{exercise}[求收敛域及和函数]
求下列幂级数的收敛域及和函数：
\begin{shortexenum}
    \shortitem{1}{$\displaystyle \sum_{n=1}^{\infty}\frac{x^{n-1}}{n2^n}$；} &
    \shortitem{2}{$\displaystyle \sum_{n=1}^{\infty}n(n+1)x^n$；} \\
    \shortitem{3}{$\displaystyle \sum_{n=0}^{\infty}(2n+1)x^n$；} &
    \shortitem{4}{$\displaystyle \sum_{n=1}^{\infty}\frac{n}{n+1}x^n$；} \\
    \shortitem{5}{$\displaystyle \sum_{n=1}^{\infty}nx^{2n}$；} &
    \shortitem{6}{$\displaystyle \sum_{n=0}^{\infty}\frac{x^{4n+1}}{4n+1}$。}
\end{shortexenum}
\end{exercise}
\begin{solution}
\begin{compactenum}
    \item 由 \(\displaystyle -\ln(1-t)=\sum_{n=1}^{\infty}\frac{t^n}{n},\ |t|<1\)，取 $t=x/2$，得
    \(\displaystyle \sum_{n=1}^{\infty}\frac{x^{n-1}}{n2^n}
    =\frac1x\sum_{n=1}^{\infty}\frac{(x/2)^n}{n}
    =-\frac1x\ln\left(1-\frac{x}{2}\right)\)，
    收敛域为 $[-2,2)$。
    \item 由 \(\displaystyle \sum_{n=1}^{\infty}nx^n=\frac{x}{(1-x)^2}\) 和
    \(\displaystyle \sum_{n=1}^{\infty}n^2x^n=\frac{x(x+1)}{(1-x)^3}\)，得
    \(\displaystyle \sum_{n=1}^{\infty}n(n+1)x^n=\sum n^2x^n+\sum nx^n=\frac{2x}{(1-x)^3},\quad |x|<1\)。
    收敛域为 $(-1,1)$。
    \item \(\displaystyle \sum_{n=0}^{\infty}(2n+1)x^n
    =2\sum_{n=1}^{\infty}nx^n+\sum_{n=0}^{\infty}x^n
    =\frac{1+x}{(1-x)^2},\quad |x|<1\)。
    收敛域为 $(-1,1)$。
    \item 由 \(\displaystyle \frac{n}{n+1}=1-\frac1{n+1}\)，可得
    \(\displaystyle \sum_{n=1}^{\infty}\frac{n}{n+1}x^n
    =\sum_{n=1}^{\infty}x^n-\sum_{n=1}^{\infty}\frac{x^n}{n+1}
    =\frac{x}{1-x}-\frac1x\left[-\ln(1-x)-x\right]
    =\frac{x}{1-x}+\frac{\ln(1-x)+x}{x},\quad |x|<1\)。
    收敛域为 $(-1,1)$。
    \item 设 $t=x^2$，则
    \(\displaystyle \sum_{n=1}^{\infty}nx^{2n}=\sum_{n=1}^{\infty}nt^n=\frac{t}{(1-t)^2}
    =\frac{x^2}{(1-x^2)^2},\quad |x|<1\)。
    收敛域为 $(-1,1)$。
    \item 这是 \(\displaystyle \int_0^x \sum_{n=0}^{\infty}t^{4n}\,\mathrm{d}t
    =\int_0^x \frac1{1-t^4}\,\mathrm{d}t,\quad |x|<1\)，
    所以和函数就是该积分；收敛域为 $[-1,1]$ 中需另判端点，但题面只要求写出和函数即可，标准收敛区间为 $(-1,1)$。
\end{compactenum}
\end{solution}

\begin{exercise}[求和与收敛半径]
\begin{compactenum}
    \item 求下列级数的和：
    \begin{shortexenum}
        \shortitem{1}{$\displaystyle \sum_{n=1}^{\infty}(-1)^n\frac{n(n+1)}{2^n}$；} &
        \shortitem{2}{$\displaystyle \sum_{n=1}^{\infty}\frac{2n-1}{2^n}$；} \\
        \shortitem{3}{$\displaystyle \sum_{n=0}^{\infty}\frac1{(4n+1)(4n+3)}$；} &
        \shortitem{4}{$\displaystyle \sum_{n=0}^{\infty}(-1)^n\frac{n^2-n+1}{2^n}$。}
    \end{shortexenum}
    \item 若 $\displaystyle \sum_{n=0}^{\infty}a_nx^n$ 的收敛半径为 $3$，$\displaystyle \sum_{n=0}^{\infty}b_nx^n$ 的收敛半径为 $5$，试问 $\displaystyle \sum_{n=0}^{\infty}(a_n+b_n)x^n$ 的收敛半径是多少？
\end{compactenum}
\end{exercise}
\begin{solution}
\begin{compactenum}
    \item
    \begin{compactenum}
        \item 取 $x=-\frac12$ 代入
        \(\displaystyle \sum_{n=1}^{\infty}n(n+1)x^n=\frac{2x}{(1-x)^3}\)，得
        \(\displaystyle \sum_{n=1}^{\infty}(-1)^n\frac{n(n+1)}{2^n}
        =\frac{2(-1/2)}{(1+1/2)^3}=-\frac{8}{27}\)。
        \item \(\displaystyle \sum_{n=1}^{\infty}\frac{2n-1}{2^n}
        =2\sum_{n=1}^{\infty}\frac{n}{2^n}-\sum_{n=1}^{\infty}\frac1{2^n}
        =2\cdot 2-1=3\)。
        \item \(\displaystyle \frac1{(4n+1)(4n+3)}=\frac12\left(\frac1{4n+1}-\frac1{4n+3}\right)\)，故级数裂项相消，和为 \(\displaystyle \frac12\)。
        \item 由
        \(\displaystyle \sum n^2x^n=\frac{x(x+1)}{(1-x)^3},\quad
        \sum nx^n=\frac{x}{(1-x)^2},\quad
        \sum x^n=\frac1{1-x}\)，
        代入 $x=-\frac12$，得和为 \(\displaystyle \frac{14}{27}\)。
    \end{compactenum}
    \item 新级数的收敛半径至少不小于 $\min\{3,5\}=3$，一般只能断言 \(\displaystyle R\ge 3\)。
    若无进一步条件，不能一定等于 $3$。
\end{compactenum}
\end{solution}

\subsubsection*{B组}

\begin{exercise}[函数项级数与幂级数的收敛域]
\begin{compactenum}
    \item 求下列函数项级数的收敛域：
    \begin{shortexenum}
        \shortitem{1}{$\displaystyle \sum_{n=1}^{\infty}\frac{3^n}{n!x^n}$；} &
        \shortitem{2}{$\displaystyle \sum_{n=1}^{\infty}\frac1n\left(\frac{x-1}{x}\right)^n$；} \\
        \shortitem{3}{$\displaystyle \sum_{n=0}^{\infty}\frac{n!}{x^n}$；} &
        \shortitem{4}{$\displaystyle \sum_{n=1}^{\infty}(-1)^n\frac{(\ln x)^n}{n}$。}
    \end{shortexenum}
    \item
    \begin{shortexenum}
        \shortitem{1}{如果幂级数 $\displaystyle \sum_{n=1}^{\infty}a_nx^n$ 在 $x=3$ 处发散，那么它在哪些 $x$ 处必然发散？} &
        \shortitem{2}{如果幂级数 $\displaystyle \sum_{n=0}^{\infty}a_n(x+5)^n$ 在 $x=-2$ 处发散，那么它在哪些 $x$ 处必然发散？}
    \end{shortexenum}
    \item 如果幂级数 $\displaystyle \sum_{n=0}^{\infty}a_n(x-3)^n$ 在 $x=7$ 处收敛，那么它在 $x$ 的什么值处必然收敛？
    \item 如果正项级数 $\displaystyle \sum_{n=0}^{\infty}a_n$ 收敛，证明 $f(x)=\displaystyle \sum_{n=0}^{\infty}a_nx^n$ 在 $(-1,1)$ 上连续。
\end{compactenum}
\end{exercise}
\begin{solution}
\begin{compactenum}
    \item
    \begin{compactenum}
        \item \(\displaystyle \sum_{n=1}^{\infty}\frac{3^n}{n!x^n}
        =\sum_{n=1}^{\infty}\frac1{n!}\left(\frac3x\right)^n\)
        对一切 $x\neq 0$ 收敛，所以收敛域为 $\mathbb R\setminus\{0\}$。
        \item 相当于 \(\displaystyle \sum_{n=1}^{\infty}\frac{t^n}{n},\ t=\frac{x-1}{x}\)。
        由对数级数收敛域 $[-1,1)$ 得 \(\displaystyle -1\le \frac{x-1}{x}<1\)，
        解得收敛域为 $[1/2,+\infty)$。
        \item 因 $n!/x^n$ 的通项对任意有限 $x$ 都不趋于 $0$，故处处发散。
        \item 等价于 \(\displaystyle -\sum_{n=1}^{\infty}\frac{(-\ln x)^n}{n}\)，收敛条件为 $-1\le \ln x<1$，即 \(\displaystyle e^{-1}\le x<e\)。
    \end{compactenum}
    \item
    \begin{compactenum}
        \item 若在 $x=3$ 发散，则收敛半径 $R\le 3$，因此当 $|x|>3$ 时必发散。
        \item 因为中心是 $-5$，且 $x=-2$ 满足 $|x+5|=3$。若该点发散，则当 $|x+5|>3$ 时必发散。
    \end{compactenum}
    \item 因为 $|7-3|=4$，在该点收敛说明收敛半径 $R\ge 4$，故凡满足 \(\displaystyle |x-3|<4\) 的点都必收敛。
    \item 对任意 $r\in(0,1)$ 与 $|x|\le r$，有 \(\displaystyle 0\le a_n|x|^n\le a_n\)。
    由 $M$ 判别法，$\sum a_nx^n$ 在 $[-r,r]$ 上一致收敛。每一项 $a_nx^n$ 连续，所以和函数在 $[-r,r]$ 上连续。由于任意 $x_0\in(-1,1)$ 都可落在某个这样的闭区间中，故 $f(x)$ 在 $(-1,1)$ 上连续。
\end{compactenum}
\end{solution}

\subsection{第六节习题：泰勒级数}
\subsubsection*{A组}

\begin{exercise}[基本概念]
回答下列问题：
\begin{shortexenum}
    \shortitem{1}{函数 $f(x)$ 能在区间 $(x_0-R,x_0+R)$ 中展开为幂级数的充要条件是什么？} &
    \shortitem{2}{如果函数 $f(x)$ 在区间 $(x_0-R,x_0+R)$ 中有任意阶导数，$f(x)$ 是否能在这区间上展开幂级数？}
\end{shortexenum}
\end{exercise}
\begin{solution}
\begin{compactenum}
    \item 充要条件是：$f(x)$ 在该区间内的 Taylor 余项满足
    \(\displaystyle \lim_{n\to\infty}R_n(x)=0\)，从而 \(\displaystyle f(x)=\sum_{n=0}^{\infty}\frac{f^{(n)}(x_0)}{n!}(x-x_0)^n\)。
    \item 不一定。存在处处有任意阶导数但不等于其 Taylor 级数的函数，例如 \(f(x)=\mathrm{e}^{-1/x^2}\ (x\neq0)\)，且 \(f(0)=0\)，
    在 $x=0$ 处的 Taylor 级数恒为 $0$，但函数本身不恒为 $0$。
\end{compactenum}
\end{solution}

\begin{exercise}[在 $x=0$ 处展开]
将下列函数展开成 $x$ 的幂级数，并指出展开式成立的范围：
\begin{shortexenum}
    \shortitem{1}{$\displaystyle f(x)=\frac1{1+x+x^2}$；} &
    \shortitem{2}{$\displaystyle f(x)=\sin^2x$；} \\
    \shortitem{3}{$\displaystyle f(x)=\frac{x}{\sqrt{1+x^2}}$；} &
    \shortitem{4}{$\displaystyle f(x)=\int_0^x\mathrm{e}^{-t^2}\,\mathrm{d}t$。}
\end{shortexenum}
\end{exercise}
\begin{solution}
\begin{compactenum}
    \item 先拆成 \(\displaystyle \frac1{1+x+x^2}=\frac{1-x}{1-x^3}=(1-x)\sum_{n=0}^{\infty}x^{3n},\quad |x|<1\)。
    所以 \(\displaystyle \frac1{1+x+x^2}=1-x+x^3-x^4+x^6-x^7+\cdots,\quad |x|<1\)。
    \item \(\displaystyle \sin^2x=\frac{1-\cos 2x}{2}
    =x^2-\frac{x^4}{3}+\frac{2x^6}{45}-\cdots,\qquad x\in\mathbb R\)。
    \item 利用 \(\displaystyle (1+x^2)^{-1/2}=1-\frac12x^2+\frac{1\cdot 3}{2\cdot 4}x^4-\cdots,\quad |x|<1\)，
    故 \(\displaystyle \frac{x}{\sqrt{1+x^2}}=x-\frac12x^3+\frac{3}{8}x^5-\cdots,\quad |x|<1\)。
    \item 由 \(\displaystyle \mathrm{e}^{-t^2}=\sum_{n=0}^{\infty}\frac{(-1)^nt^{2n}}{n!}\)，逐项积分得 \(\displaystyle \int_0^x\mathrm{e}^{-t^2}\,\mathrm{d}t=\sum_{n=0}^{\infty}\frac{(-1)^n}{(2n+1)n!}x^{2n+1},\ x\in\mathbb R\)。
\end{compactenum}
\end{solution}

\begin{exercise}[在指定点展开]
将下列函数在指定点展开成幂级数，并指出展开式成立的范围：
\begin{shortexenum}
    \shortitem{1}{$\displaystyle f(x)=\frac1{x^2+4x+3},\quad x=1$；} &
    \shortitem{2}{$\displaystyle f(x)=\frac{2x+1}{x^2+x-2},\quad x=2$；} \\
    \shortitem{3}{$\displaystyle f(x)=\ln\frac1{2+2x+x^2},\quad x=-1$；} &
    \shortitem{4}{$\displaystyle f(x)=\cos x,\quad x=-\frac{\pi}{3}$。}
\end{shortexenum}
\end{exercise}
\begin{solution}
\begin{compactenum}
    \item 设 $t=x-1$，则 \(\displaystyle \frac1{x^2+4x+3}=\frac1{(x+1)(x+3)}
    =\frac12\left(\frac1{x+1}-\frac1{x+3}\right)\)。
    又有 \(\displaystyle x+1=2+t,\qquad x+3=4+t\)，
    所以可分别按 $t$ 展开，得在 $|t|<2$ 内的幂级数。
    \item 因 \(\displaystyle \frac{2x+1}{x^2+x-2}=\frac{2x+1}{(x-1)(x+2)}\)，令 $t=x-2$，再作部分分式展开，可得到关于 $t$ 的几何级数展开，收敛半径由最近奇点 $x=1$ 决定，为 $1$。
    \item 令 $t=x+1$，则 \(\displaystyle 2+2x+x^2=1+t^2,\qquad \ln\frac1{2+2x+x^2}=-\ln(1+t^2)\)。
    故 \(\displaystyle \ln\frac1{2+2x+x^2}
    =-\sum_{n=1}^{\infty}\frac{(-1)^{n-1}}{n}t^{2n},\ |t|<1\)。
    即 \(\displaystyle \ln\frac1{2+2x+x^2}
    =-(x+1)^2+\frac{(x+1)^4}{2}-\frac{(x+1)^6}{3}+\cdots,\ |x+1|<1\)。
    \item 设 $t=x+\frac{\pi}{3}$，则 \(\displaystyle \cos x=\cos\left(t-\frac{\pi}{3}\right)=\frac12\cos t+\frac{\sqrt3}{2}\sin t\)。
    再把 $\cos t,\sin t$ 展开即可，且对一切实数成立。
\end{compactenum}
\end{solution}

\subsubsection*{B组}

\begin{exercise}[由已知展开推出新恒等式]
将函数 \(\displaystyle \frac{\mathrm{d}}{\mathrm{d}x}\left(\frac{\mathrm{e}^x-1}{x}\right)\)
展开成 $x$ 的幂级数，并证明 \(\displaystyle \sum_{n=1}^{\infty}\frac{n}{(n+1)!}=1\)。
\end{exercise}
\begin{solution}
先有 \(\displaystyle \frac{\mathrm{e}^x-1}{x}=\frac1x\sum_{n=1}^{\infty}\frac{x^n}{n!}=\sum_{n=0}^{\infty}\frac{x^n}{(n+1)!}\)。
逐项求导得 \(\displaystyle \frac{\mathrm{d}}{\mathrm{d}x}\left(\frac{\mathrm{e}^x-1}{x}\right)=\sum_{n=1}^{\infty}\frac{n}{(n+1)!}x^{n-1}\)。
令 $x=1$，左端化为 \(\displaystyle \frac{\mathrm{e}(1)-(\mathrm{e}-1)}{1^2}=1\)，故 \(\displaystyle \sum_{n=1}^{\infty}\frac{n}{(n+1)!}=1\)。
\end{solution}

\begin{exercise}[由逐项求导与逐项积分求展开]
利用逐项求导和逐项积分的方法，将函数
\(\displaystyle f(x)=\arctan\frac{4+x^2}{4-x^2}\)
展开成 $x$ 的幂级数。
\end{exercise}
\begin{solution}
先用反正切恒等式化简，再代入
\(\displaystyle \arctan t=\sum_{n=0}^{\infty}(-1)^n\frac{t^{2n+1}}{2n+1}\)：
\(\displaystyle \arctan\frac{4+x^2}{4-x^2}=\arctan 1+\arctan\frac{x^2}{2}=\frac{\pi}{4}+\arctan\frac{x^2}{2}\)。
因此 \(\displaystyle \arctan\frac{4+x^2}{4-x^2}=\frac{\pi}{4}+\sum_{n=0}^{\infty}(-1)^n\frac{x^{4n+2}}{(2n+1)2^{2n+1}},\quad |x|\le \sqrt2\)。
\end{solution}

\begin{exercise}[微分方程的幂级数解]
求下列微分方程的幂级数解：
\begin{shortexenum}
    \shortitem{1}{$y''=xy$；} &
    \shortitem{2}{$xy''+y'+xy=0$。}
\end{shortexenum}
\end{exercise}
\begin{solution}
\begin{compactenum}
    \item 设 \(\displaystyle y=\sum_{n=0}^{\infty}a_nx^n,\qquad
    y''=\sum_{n=0}^{\infty}(n+2)(n+1)a_{n+2}x^n\)。
    代入 $y''=xy$，得
    \(\displaystyle (n+2)(n+1)a_{n+2}=a_{n-1}\quad (n\ge 1),\qquad 2a_2=0\)。
    所以 \(\displaystyle a_2=0,\qquad
    a_{n+2}=\frac{a_{n-1}}{(n+2)(n+1)}\ (n\ge 1)\)。
    从而系数按模 $3$ 分成三链，得到两组独立常数决定的幂级数解。
    \item 设 \(\displaystyle y=\sum_{n=0}^{\infty}a_nx^n\)。由正文中递推关系可得 \(\displaystyle a_{2k+1}=0,\qquad a_{2k}=(-1)^k\frac{a_0}{2^{2k}(k!)^2}\)。
    因此
    \[
    y=a_0\sum_{k=0}^{\infty}(-1)^k\frac1{(k!)^2}\left(\frac{x}{2}\right)^{2k}.
    \]
\end{compactenum}
\end{solution}

\subsection{第七节习题：周期函数的傅里叶级数}
\subsubsection*{A组}

\begin{exercise}[基本概念]
回答下列问题：
\begin{shortexenum}
    \shortitem{1}{三角函数系的正交性是指什么？} &
    \shortitem{2}{收敛定理的条件有哪些？} \\
    \shortitem{3}{周期为 $2\pi$ 的函数 $f(x)$ 的傅里叶级数是否一定收敛？如果收敛，是否一定收敛到自身，即收敛到 $f(x)$？} &
    \shortitem{4}{奇函数和偶函数的傅里叶系数有什么特点？}
\end{shortexenum}
\end{exercise}
\begin{solution}
\begin{compactenum}
    \item 在 $[-\pi,\pi]$ 上，\(\displaystyle \int_{-\pi}^{\pi}\cos mx\cos nx\,\mathrm{d}x\) 与 \(\displaystyle \int_{-\pi}^{\pi}\sin mx\sin nx\,\mathrm{d}x\) 均满足：当 $m\ne n$ 时为 $0$，当 $m=n\ge1$ 时为 $\pi$；并且 \(\displaystyle \int_{-\pi}^{\pi}\sin mx\cos nx\,\mathrm{d}x=0\)。这就叫三角函数系的正交性。
    \item 常用的充分条件是：函数在一个周期内分段光滑，只有有限个极值点和第一类间断点。
    \item 不一定。傅里叶级数是否收敛取决于函数的光滑性。即使收敛，在间断点处也不收敛到函数值，而收敛到左右极限的平均值。
    \item 若 $f$ 为偶函数，则 $b_n=0$，并且 $f(x)\sim \frac{a_0}{2}+\sum_{n=1}^{\infty}a_n\cos nx$。若 $f$ 为奇函数，则 $a_0=0,\ a_n=0$，并且 $f(x)\sim \sum_{n=1}^{\infty}b_n\sin nx$。
\end{compactenum}
\end{solution}

\begin{exercise}[展开成傅里叶级数]
试将下列以 $2\pi$ 为周期的函数 $f(x)$ 展开成傅里叶级数：
\begin{shortsingleenum}
    \shortsingleitem{1}{$\displaystyle f(x)=\begin{cases}
        -\pi, & -\pi\le x<0,\\
        x, & 0\le x<\pi;
    \end{cases}$}
    \shortsingleitem{2}{$\displaystyle f(x)=\begin{cases}
        -1, & -\pi\le x<0,\\
        1, & 0\le x<\pi;
    \end{cases}$}
    \shortsingleitem{3}{$f(x)=|\sin x|,\ -\pi\le x<\pi$；}
    \shortsingleitem{4}{$\displaystyle f(x)=\begin{cases}
        \mathrm{e}^x, & -\pi\le x<0,\\
        0, & 0\le x<\pi;
    \end{cases}$}
    \shortsingleitem{5}{$f(x)=\dfrac{\pi-x}{2},\ 0<x<2\pi$；}
    \shortsingleitem{6}{$f(x)=2\sin\dfrac{x}{3},\ -\pi\le x\le\pi$。}
\end{shortsingleenum}
\end{exercise}
\begin{solution}
\begin{compactenum}
    \item 计算得 \(\displaystyle a_0=-\frac{\pi}{2},\ a_n=\frac{(-1)^n-1}{\pi n^2},\ b_n=\frac{1-2(-1)^n}{n}\)，因而 \(\displaystyle f(x)\sim -\frac{\pi}{4}+\sum_{n=1}^{\infty}\frac{(-1)^n-1}{\pi n^2}\cos nx+\sum_{n=1}^{\infty}\frac{1-2(-1)^n}{n}\sin nx\)。
    \item 这是奇函数，故 $a_0=a_n=0$。又 \(\displaystyle b_n=\frac{2}{\pi}\int_0^\pi \sin nx\,\mathrm{d}x=\frac{2}{n\pi}\bigl(1-(-1)^n\bigr)\)，所以 \(\displaystyle f(x)\sim \frac{4}{\pi}\left(\sin x+\frac{\sin 3x}{3}+\frac{\sin 5x}{5}+\cdots\right)\)。
    \item 这是偶函数，故 $b_n=0$。计算得 \(\displaystyle a_0=\frac{4}{\pi},\ a_{2n}=-\frac{4}{\pi(4n^2-1)},\ a_{2n-1}=0\)，故 \(\displaystyle |\sin x|=\frac{2}{\pi}-\frac{4}{\pi}\sum_{n=1}^{\infty}\frac{\cos 2nx}{4n^2-1}\)。
    \item 计算得 \(\displaystyle a_0=\frac{1-\mathrm{e}^{-\pi}}{\pi},\ a_n=\frac{1-(-1)^n\mathrm{e}^{-\pi}}{\pi(1+n^2)},\ b_n=-\frac{n\bigl(1-(-1)^n\mathrm{e}^{-\pi}\bigr)}{\pi(1+n^2)}\)，因此 \(\displaystyle f(x)\sim \frac{1-\mathrm{e}^{-\pi}}{2\pi}+\sum_{n=1}^{\infty}\frac{1-(-1)^n\mathrm{e}^{-\pi}}{\pi(1+n^2)}(\cos nx-n\sin nx)\)。
    \item 这是经典锯齿波，\(\displaystyle \frac{\pi-x}{2}=\sum_{n=1}^{\infty}\frac{\sin nx}{n}\ (0<x<2\pi)\)。
    \item 这是奇函数，故 $a_0=a_n=0$。由 \(\displaystyle b_n=\frac{2}{\pi}\int_0^\pi 2\sin\frac{x}{3}\sin nx\,\mathrm{d}x=(-1)^{n+1}\frac{18\sqrt3\,n}{\pi(9n^2-1)}\)，得 \(\displaystyle 2\sin\frac{x}{3}\sim\sum_{n=1}^{\infty}(-1)^{n+1}\frac{18\sqrt3\,n}{\pi(9n^2-1)}\sin nx\)。
\end{compactenum}
\end{solution}

\begin{exercise}[和函数在特殊点的取值]
设 $f(x)$ 是以 $2\pi$ 为周期的函数，它在 $[-\pi,\pi]$ 上满足 \(f(x)=x+1\ (-\pi\le x<0)\)，且 \(f(x)=x^2\ (0\le x<\pi)\)。
若它的傅里叶级数的和函数为 $S(x)$，试问 $S(-\pi),S(0)$ 和 $S(\pi)$ 的值各为多少？
\end{exercise}
\begin{solution}
由傅里叶级数收敛定理，在 $x=0$ 处，\(\displaystyle S(0)=\frac{f(0-)+f(0+)}{2}=\frac12\)。在端点 $x=\pm\pi$ 处，应取一个周期两端邻值平均。这里 \(f(-\pi+)=1-\pi,\ f(\pi-)=\pi^2\)，所以 \(\displaystyle S(-\pi)=S(\pi)=\frac{\pi^2+1-\pi}{2}\)。
\end{solution}

\subsubsection*{B组}

\begin{exercise}[系数之间的关系]
设 $\varphi$ 和 $\psi$ 都是周期为 $2\pi$ 的函数。
\begin{shortexenum}
    \shortitem{1}{若函数 $\varphi(-x)=\psi(x),\ -\pi\le x<\pi$，问 $\varphi(x)$ 和 $\psi(x)$ 的傅里叶系数 $a_n,b_n$ 与 $\alpha_n,\beta_n\ (n=0,1,2,\cdots)$ 之间有何关系？} &
    \shortitem{2}{若函数 $\varphi(-x)=-\psi(x),\ -\pi\le x<\pi$，问 $\varphi(x)$ 和 $\psi(x)$ 的傅里叶系数 $a_n,b_n$ 与 $\alpha_n,\beta_n\ (n=0,1,2,\cdots)$ 之间有何关系？}
\end{shortexenum}
\end{exercise}
\begin{solution}
\begin{compactenum}
    \item 由变量代换 $x\mapsto -x$，有 \(\displaystyle \alpha_n=\frac1\pi\int_{-\pi}^{\pi}\psi(x)\cos nx\,\mathrm{d}x
    =\frac1\pi\int_{-\pi}^{\pi}\varphi(-x)\cos nx\,\mathrm{d}x=a_n\)，以及 \(\displaystyle \beta_n=\frac1\pi\int_{-\pi}^{\pi}\psi(x)\sin nx\,\mathrm{d}x
    =\frac1\pi\int_{-\pi}^{\pi}\varphi(-x)\sin nx\,\mathrm{d}x=-b_n\)。
    因而 $\alpha_n=a_n,\ \beta_n=-b_n$。
    \item 同理，$\alpha_n=-a_n,\ \beta_n=b_n$。
\end{compactenum}
\end{solution}

\begin{exercise}[利用半周期对称性]
设 $f(x)$ 是以 $2\pi$ 为周期的函数。证明：
\begin{shortexenum}
    \shortitem{1}{若函数 $f(x-\pi)=f(x)$，则 $f(x)$ 的傅里叶系数满足 $a_{2k+1}=0,\ b_{2k+1}=0\ (k=0,1,2,\cdots)$。} &
    \shortitem{2}{若函数 $f(x-\pi)=-f(x)$，则 $f(x)$ 的傅里叶系数满足 $a_0=0,\ a_{2k}=0,\ b_{2k}=0\ (k=1,2,\cdots)$。}
\end{shortexenum}
\end{exercise}
\begin{solution}
\begin{compactenum}
    \item 以 $a_n$ 为例，$a_n=\frac1\pi\int_{-\pi}^{\pi}f(x)\cos nx\,\mathrm{d}x$。
    把积分区间平移 $\pi$，并用 $f(x-\pi)=f(x)$，得 \(\displaystyle a_n=\frac1\pi\int_{-\pi}^{\pi}f(x-\pi)\cos n(x-\pi)\,\mathrm{d}x
    =(-1)^na_n\)。
    因而当 $n$ 为奇数时 $a_n=-a_n$，故 $a_n=0$。同理 $b_n=0$。
    \item 同样地，\(\displaystyle a_n=\frac1\pi\int_{-\pi}^{\pi}f(x-\pi)\cos n(x-\pi)\,\mathrm{d}x
    =-(-1)^na_n\)。
    当 $n$ 为偶数时有 $a_n=-a_n$，所以 $a_{2k}=0$，特别地 $a_0=0$。对 $b_n$ 同理可得 $b_{2k}=0$。
\end{compactenum}
\end{solution}

\subsection{第八节习题：任意区间上的傅里叶级数}
\subsubsection*{A组}

\begin{exercise}[周期延拓与半区间展开]
试将下列函数展开成傅里叶级数：
\begin{shortexenum}
    \shortitem{1}{周期函数在一个周期内满足 \(f(x)=0\ (-2\le x<0)\)，且 \(f(x)=1\ (0\le x<2)\)；} &
    \shortitem{2}{周期函数在一个周期内满足 \(f(x)=2x+1\ (-3\le x<0)\)，且 \(f(x)=1\ (0\le x<3)\)；}\\
    \shortitem{3}{$f(x)=x\cos x,\ -\dfrac{\pi}{2}\le x\le \dfrac{\pi}{2}$；} &
    \shortitem{4}{周期函数在一个周期内满足 \(f(x)=2-x\ (0\le x\le 4)\)，且 \(f(x)=x-6\ (4<x\le 8)\)；}\\
    \shortitem{5}{将 $f(x)=\dfrac{x^2}{2}\ (0\le x\le \pi)$ 展开成正弦级数；} &
    \shortitem{6}{证明：当 $0<x<\pi$ 时，$\sin x+\frac13\sin 3x+\frac15\sin 5x+\cdots=\frac{\pi}{4}$。}\\
    \shortitem{7}{将 $f(x)=x-1\ (0\le x\le 2)$ 展开成周期为 $4$ 的余弦级数；} &
    \shortitem{8}{试将 $f(x)=\dfrac{\pi-x}{2}\ (0\le x\le \pi)$ 展开成正弦级数。}
\end{shortexenum}
\end{exercise}
\begin{solution}
\begin{compactenum}
    \item 这里周期为 $4$，取 $l=2$。计算得 \(a_0=1,\ a_n=0\)，且 \(\displaystyle b_n=\frac{1-(-1)^n}{n\pi}\)。故 \(\displaystyle f(x)\sim \frac12+\frac{2}{\pi}\sum_{k=1}^{\infty}\frac{1}{2k-1}\sin\frac{(2k-1)\pi x}{2}\)。
    \item 这里周期为 $6$，取 $l=3$。计算得 \(\displaystyle a_0=-1,\ a_n=\frac{6(1-(-1)^n)}{n^2\pi^2}\)，且 \(\displaystyle b_n=\frac{6(-1)^{n+1}}{n\pi}\)。所以 \(\displaystyle f(x)\sim -\frac12+\sum_{n=1}^{\infty}\left[\frac{6(1-(-1)^n)}{n^2\pi^2}\cos\frac{n\pi x}{3}+\frac{6(-1)^{n+1}}{n\pi}\sin\frac{n\pi x}{3}\right]\)。
    \item 一个周期长度为 $\pi$，故用基函数 $\sin 2nx,\cos 2nx$。由于 $x\cos x$ 为奇函数，所以只有正弦项。计算得 \(\displaystyle x\cos x=\sum_{n=1}^{\infty}\frac{16n(-1)^{n-1}}{\pi(2n-1)^2(2n+1)^2}\sin 2nx\)。
    \item 令 $t=x-4$，则一周期内函数变为 \(g(t)=|t|-2\ (-4\le t\le 4)\)。它是偶函数，故只有余弦项，且 \(\displaystyle g(t)=-\frac{16}{\pi^2}\sum_{k=1}^{\infty}\frac{\cos\frac{(2k-1)\pi t}{4}}{(2k-1)^2}\)。换回 $x$ 得 \(\displaystyle f(x)=\frac{16}{\pi^2}\sum_{k=1}^{\infty}\frac{\cos\frac{(2k-1)\pi x}{4}}{(2k-1)^2}\)。
    \item 正弦级数系数为 \(\displaystyle b_n=\frac{2}{\pi}\int_0^\pi \frac{x^2}{2}\sin nx\,\mathrm{d}x\)，计算得 \(\displaystyle b_n=-\frac{\pi(-1)^n}{n}+\frac{2\bigl((-1)^n-1\bigr)}{\pi n^3}\)。因此 \(\displaystyle \frac{x^2}{2}\sim\sum_{n=1}^{\infty}\left[-\frac{\pi(-1)^n}{n}+\frac{2\bigl((-1)^n-1\bigr)}{\pi n^3}\right]\sin nx\)。
    \item 取第 10.7 节中的方波 \(f(x)=-1\ (-\pi\le x<0)\)，且 \(f(x)=1\ (0\le x<\pi)\)。它的傅里叶级数为 \(\displaystyle \frac{4}{\pi}\left(\sin x+\frac13\sin 3x+\frac15\sin 5x+\cdots\right)\)。当 $0<x<\pi$ 时，和函数等于 $1$，故结论成立。
    \item 周期为 $4$，故 $l=2$。余弦级数为 \(\displaystyle x-1\sim \sum_{n=1}^{\infty}a_n\cos\frac{n\pi x}{2}\)，其中 \(\displaystyle a_n=\int_0^2 (x-1)\cos\frac{n\pi x}{2}\,\mathrm{d}x=\frac{4\bigl((-1)^n-1\bigr)}{n^2\pi^2}\)。所以 \(\displaystyle x-1=-\frac{8}{\pi^2}\sum_{k=1}^{\infty}\frac{\cos\frac{(2k-1)\pi x}{2}}{(2k-1)^2}\)。
    \item \(\displaystyle \frac{\pi-x}{2}=\sum_{n=1}^{\infty}\frac{\sin nx}{n}\qquad (0<x<\pi)\)。
\end{compactenum}
\end{solution}

\begin{exercise}[同一函数的不同展开]
将 $f(x)=x$ 在 $[0,\pi]$ 上分别展开成余弦级数和正弦级数。
\end{exercise}
\begin{solution}
\begin{itemize}
    \item 余弦级数：\(\displaystyle a_0=\frac{2}{\pi}\int_0^\pi x\,\mathrm{d}x=\pi,\ a_n=\frac{2}{\pi}\int_0^\pi x\cos nx\,\mathrm{d}x=\frac{2\bigl((-1)^n-1\bigr)}{\pi n^2}\)，因此 \(\displaystyle x=\frac{\pi}{2}-\frac{4}{\pi}\sum_{k=1}^{\infty}\frac{\cos(2k-1)x}{(2k-1)^2},\ 0\le x\le \pi\)。
    \item 正弦级数：\(\displaystyle b_n=\frac{2}{\pi}\int_0^\pi x\sin nx\,\mathrm{d}x=\frac{2(-1)^{n+1}}{n}\)，因此 \(\displaystyle x=2\sum_{n=1}^{\infty}\frac{(-1)^{n+1}}{n}\sin nx,\ 0<x<\pi\)。
\end{itemize}
\end{solution}

\begin{exercise}[余弦级数与 Fourier 级数应用]
\begin{shortexenum}
    \shortitem{1}{将函数 \(f(x)=\sin x\ (0\le x\le \frac{\pi}{2})\)，且 \(f(x)=0\ (\frac{\pi}{2}<x\le \pi)\) 展开成余弦级数。} &
    \shortitem{2}{将 $f(x)=x\ (1<x<3)$ 展开成傅里叶级数，并用它证明
    \(\displaystyle \sum_{n=1}^{\infty}\frac{(-1)^{n-1}}{2n-1}=\frac{\pi}{4}\)。}
\end{shortexenum}
\end{exercise}
\begin{solution}
\begin{compactenum}
    \item 余弦级数系数为 \(\displaystyle a_0=\frac{2}{\pi}\int_0^{\pi/2}\sin x\,\mathrm{d}x=\frac{2}{\pi}\)，且 \(\displaystyle a_n=\frac{2}{\pi}\int_0^{\pi/2}\sin x\cos nx\,\mathrm{d}x\)。由积化和差公式，\(\displaystyle a_n=\frac1\pi\left[\frac{1-\cos\frac{(n+1)\pi}{2}}{n+1}-\frac{1-\cos\frac{(n-1)\pi}{2}}{n-1}\right]\quad (n\ne 1)\)，且 \(a_1=\frac{1}{\pi}\)。故 \(\displaystyle f(x)\sim \frac1\pi+\sum_{n=1}^{\infty}a_n\cos nx\)。
    \item 令 \(t=x-2\)，则 \(t\in(-1,1)\)，且 \(f(x)=t+2\)。将其作周期 $2$ 延拓，有 \(\displaystyle t=2\sum_{n=1}^{\infty}\frac{(-1)^{n+1}}{n\pi}\sin n\pi t\)。所以 \(\displaystyle x=2+2\sum_{n=1}^{\infty}\frac{(-1)^{n+1}}{n\pi}\sin n\pi(x-2),\qquad 1<x<3\)。取 \(x=\dfrac32\)，则 \(t=-\dfrac12\)，从而 \(\displaystyle -\frac12=2\sum_{n=1}^{\infty}\frac{(-1)^{n+1}}{n\pi}\sin\frac{n\pi}{2}\)。只剩奇数项，整理得 \(\displaystyle \sum_{n=1}^{\infty}\frac{(-1)^{n-1}}{2n-1}=\frac{\pi}{4}\)。
\end{compactenum}
\end{solution}

\subsubsection*{B组}

\begin{exercise}[由 Fourier 展开推出恒等式]
证明在 $[0,\pi]$ 上下式成立：
\begin{shortexenum}
    \shortitem{1}{$\displaystyle x(\pi-x)=\frac{\pi^2}{6}-\sum_{n=1}^{\infty}\frac{\cos 2nx}{n^2}$；} &
    \shortitem{2}{$\displaystyle x(\pi-x)=\frac{8}{\pi}\sum_{n=1}^{\infty}\frac{\sin(2n-1)x}{(2n-1)^3}$。}
\end{shortexenum}
并利用以上结果证明
\begin{shortexenum}
    \shortitem{1}{$\displaystyle \sum_{n=1}^{\infty}\frac{(-1)^{n-1}}{n^2}=\frac{\pi^2}{12}$；} &
    \shortitem{2}{$\displaystyle \sum_{n=1}^{\infty}\frac{(-1)^{n-1}}{(2n-1)^3}=\frac{\pi^3}{32}$。}
\end{shortexenum}
\end{exercise}
\begin{solution}
\begin{compactenum}
    \item 这是 \(x(\pi-x)\) 在 \([0,\pi]\) 上作余弦级数展开所得。计算
    \[
    a_0=\frac{2}{\pi}\int_0^\pi x(\pi-x)\,\mathrm{d}x=\frac{\pi^2}{3},
    \qquad a_{2n}=-\frac{2}{n^2},\quad a_{2n-1}=0.
    \]
    所以
    \[
    x(\pi-x)=\frac{\pi^2}{6}-\sum_{n=1}^{\infty}\frac{\cos 2nx}{n^2}.
    \]
    \item 这是同一函数的正弦级数展开。计算得 \(b_{2n}=0\)，且
    \[
    b_{2n-1}=\frac{8}{\pi(2n-1)^3}.
    \]
    故
    \[
    x(\pi-x)=\frac{8}{\pi}\sum_{n=1}^{\infty}\frac{\sin(2n-1)x}{(2n-1)^3}.
    \]
    \item 在第一式中取 \(x=\dfrac{\pi}{2}\)，有 \(\displaystyle \frac{\pi^2}{4}=\frac{\pi^2}{6}-\sum_{n=1}^{\infty}\frac{(-1)^n}{n^2}\)，整理得 \(\displaystyle \sum_{n=1}^{\infty}\frac{(-1)^{n-1}}{n^2}=\frac{\pi^2}{12}\)。
    \item 在第二式中取 \(x=\dfrac{\pi}{2}\)，有 \(\displaystyle \frac{\pi^2}{4}=\frac{8}{\pi}\sum_{n=1}^{\infty}\frac{(-1)^{n-1}}{(2n-1)^3}\)，从而 \(\displaystyle \sum_{n=1}^{\infty}\frac{(-1)^{n-1}}{(2n-1)^3}=\frac{\pi^3}{32}\)。
\end{compactenum}
\end{solution}

\subsection{第九节习题：傅里叶级数的复数形式}
\begin{exercise}[由复系数写出实数形式]
在例 10.9.1 所得到的式 (10.9.8) 中，取 $\tau=\dfrac{l}{3}$，试将复数形式的傅里叶级数 (10.9.8) 的实数形式写出来。
\end{exercise}
\begin{solution}
矩形波在 $[-l,l]$ 内满足 \(f(x)=E\ (|x|<\dfrac{l}{3})\)，且 \(f(x)=0\ (\dfrac{l}{3}<|x|\le l)\)。
由正文复系数公式，
\[
\begin{array}{c}
c_0=\frac{E}{3},\qquad
c_n=\frac{E}{n\pi}\sin\frac{n\pi}{3}\quad (n\ne 0),\\[2pt]
a_n=c_n+c_{-n}=2c_n,\qquad b_n=0.
\end{array}
\]
该函数是偶函数，所以实数形式中只有余弦项，
从而
\[
f(x)\sim \frac{E}{3}
+\frac{2E}{\pi}\sum_{n=1}^{\infty}\frac{\sin\frac{n\pi}{3}}{n}\cos\frac{n\pi x}{l}.
\]
这就是所求的实数形式。
\end{solution}

\begin{exercise}[周期为 $2$ 的函数的复 Fourier 展开]
设 $f(x)$ 是周期为 $2$ 的周期函数，它在 $[-1,1]$ 上的表达式为 $f(x)=x$。试将 $f(x)$ 展开为复数形式的傅里叶级数。
\end{exercise}
\begin{solution}
这里周期 $2l=2$，故 $l=1$，复数形式应写为
\[
f(x)\sim \sum_{n=-\infty}^{+\infty}c_n\mathrm{e}^{in\pi x},\qquad
c_n=\frac12\int_{-1}^{1}x\,\mathrm{e}^{-in\pi x}\,\mathrm{d}x.
\]
显然 $c_0=0$；当 $n\ne0$ 时，
\[
c_n=-i\int_0^1 x\sin(n\pi x)\,\mathrm{d}x
=-i\left[-\frac{x\cos(n\pi x)}{n\pi}+\frac{\sin(n\pi x)}{n^2\pi^2}\right]_0^1
=\frac{i(-1)^n}{n\pi}.
\]
所以 \(\displaystyle x\sim \sum_{\substack{n=-\infty\\ n\ne 0}}^{+\infty}\frac{i(-1)^n}{n\pi}\mathrm{e}^{in\pi x}\)。
\end{solution}

\subsection{本章总习题}

\begin{exercise}[填空题]
填写下列各题：
\begin{shortexenum}
    \shortitem{1}{级数 $\displaystyle \sum_{n=0}^{\infty}\frac{(\ln 3)^n}{2^n}$ 的和为\underline{\hspace{5em}}。} &
    \shortitem{2}{若级数 $\displaystyle \sum_{n=1}^{\infty}u_n$ 收敛于 $A$，则级数 $\displaystyle \sum_{n=1}^{\infty}(u_n+u_{n+1})$ 收敛于\underline{\hspace{5em}}。} \\
    \shortitem{3}{若级数 $\displaystyle \sum_{n=1}^{\infty}a_n$ 的部分和序列为 $S_n=\dfrac{2n}{n+1}$，则 $a_n=$\underline{\hspace{5em}}，$\displaystyle \sum_{n=2}^{\infty}a_n=$\underline{\hspace{5em}}。} &
    \shortitem{4}{若级数 $\displaystyle \sum_{n=1}^{\infty}\frac{(-1)^n+a}{n}$ 收敛，则 $a$ 的取值范围为\underline{\hspace{5em}}。} \\
    \shortitem{5}{级数 $\displaystyle \sum_{n=2}^{\infty}\frac{(-1)^n\ln n}{n^p}$ 在 $p=$\underline{\hspace{5em}}时收敛。} &
    \shortitem{6}{设有级数 $\displaystyle \sum_{n=0}^{\infty}a_n\left(\frac{x+1}{2}\right)^n$，若 $\displaystyle \lim_{n\to\infty}\left|\frac{a_n}{a_{n+1}}\right|=\frac13$，则该级数的收敛半径等于\underline{\hspace{5em}}。} \\
    \shortitem{7}{设幂级数 $\displaystyle \sum_{n=0}^{\infty}a_nx^n$ 的收敛半径为 $3$，则幂级数 $\displaystyle \sum_{n=1}^{\infty}na_n(x-1)^{n+1}$ 的收敛区间为\underline{\hspace{5em}}。} &
    \shortitem{8}{设 $f(x)$ 是周期为 $2$ 的周期函数，它在区间 $(-1,1]$ 上满足 \(f(x)=2\ (-1<x\le0)\)，且 \(f(x)=x^3\ (0<x\le1)\)，则 $f(x)$ 的傅里叶级数在 $x=1$ 处收敛于\underline{\hspace{5em}}。} \\
    \shortitem{9}{设函数 $f(x)=\pi x+x^2\ (-\pi<x<\pi)$ 的傅里叶级数展开式为 \(\displaystyle \frac{a_0}{2}+\sum_{n=1}^{\infty}(a_n\cos nx+b_n\sin nx)\)，则其中系数 $b_3$ 的值为\underline{\hspace{5em}}。} &
    \shortitem{10}{幂级数 $\displaystyle \sum_{n=1}^{\infty}\frac{(x-2)^n}{n4^n}$ 的收敛域为\underline{\hspace{5em}}。}
\end{shortexenum}
\end{exercise}
\begin{solution}
\begin{compactanswerenum}
    \item $\sum_{n=0}^{\infty}\left(\frac{\ln 3}{2}\right)^n=\frac{1}{1-\frac{\ln 3}{2}}=\frac{2}{2-\ln 3}$。
    \item $\sum_{n=1}^{\infty}(u_n+u_{n+1})=\sum_{n=1}^{\infty}u_n+\sum_{n=2}^{\infty}u_n=A+(A-u_1)=2A-u_1$。
    \item $a_n=S_n-S_{n-1}=\frac{2}{n(n+1)}$，且 $\sum_{n=2}^{\infty}a_n=2-1=1$。
    \item 只能取 $a=0$。
    \item 当且仅当 $p>0$ 时收敛。
    \item 设 $c_n=\dfrac{a_n}{2^n}$，则 $\lim_{n\to\infty}\left|\frac{c_n}{c_{n+1}}\right|=2\lim_{n\to\infty}\left|\frac{a_n}{a_{n+1}}\right|=\frac23$。
    所以收敛半径为 $\dfrac23$。
    \item 它关于 $x-1$ 的收敛半径仍为 $3$，所以开区间是 $(-2,4)$。
    \item 端点取左右极限平均值，故 $\dfrac{f(1-0)+f(-1+0)}{2}=\dfrac{1+2}{2}=\dfrac32$。
    \item 由
    \(\displaystyle b_n=\frac1\pi\int_{-\pi}^{\pi}(\pi x+x^2)\sin nx\,\mathrm{d}x
    =\int_{-\pi}^{\pi}x\sin nx\,\mathrm{d}x
    =-\frac{2\pi(-1)^n}{n},\qquad
    b_3=\frac{2\pi}{3}\)。
    \item 收敛半径为 $4$，端点 $x=6$ 发散、$x=-2$ 收敛，所以收敛域为 $[-2,6)$。
\end{compactanswerenum}
\end{solution}

\begin{exercise}[选择题]
选择唯一正确的选项：
\begin{compactenum}
    \item 已知级数 $\displaystyle \sum_{n=1}^{\infty}(-1)^{n-1}a_n=2$，$\displaystyle \sum_{n=1}^{\infty}a_{2n-1}=5$，则级数 $\displaystyle \sum_{n=1}^{\infty}a_n$ 等于
    
    \begin{tabular*}{\linewidth}{@{\extracolsep{\fill}}llll@{}}
    \text{(A) }3 & \text{(B) }7 & \text{(C) }8 & \text{(D) }9
    \end{tabular*}
    \item 设 $\alpha$ 为常数，则级数
    \(\displaystyle \sum_{n=1}^{\infty}\left[\frac{\sin(n\alpha)}{n^2}-\frac1{\sqrt n}\right]\)
    的敛散性为
    
    \begin{tabular*}{\linewidth}{@{\extracolsep{\fill}}ll@{}}
    \text{(A) 绝对收敛} & \text{(B) 发散} \\
    \text{(C) 条件收敛} & \text{(D) 收敛性与 $\alpha$ 有关}
    \end{tabular*}
    \item 设 $0\le a_n<\dfrac1n\ (n=1,2,\cdots)$，则下列级数中肯定收敛的是
    
    \begin{tabular*}{\linewidth}{@{\extracolsep{\fill}}llll@{}}
    \text{(A) }\(\displaystyle \sum_{n=1}^{\infty}a_n\) &
    \text{(B) }\(\displaystyle \sum_{n=1}^{\infty}(-1)^na_n\) &
    \text{(C) }\(\displaystyle \sum_{n=1}^{\infty}\sqrt{a_n}\) &
    \text{(D) }\(\displaystyle \sum_{n=1}^{\infty}(-1)^na_n^2\)
    \end{tabular*}
    \item 下列各选项中正确的是
    
    \begin{tabular*}{\linewidth}{@{\extracolsep{\fill}}ll@{}}
    \makecell[l]{\text{(A) 若 }\(\sum a_n^2\)\text{ 和 }\(\sum b_n^2\)\text{ 都收敛，}\\\text{则 }\(\sum (a_n+b_n)^2\)\text{ 收敛；}} &
    \makecell[l]{\text{(B) 若 }\(\sum |a_nb_n|\)\text{ 收敛，}\\\text{则 }\(\sum a_n^2\)\text{ 与 }\(\sum b_n^2\)\text{ 都收敛；}} \\
    \makecell[l]{\text{(C) 若正项级数 }\(\sum a_n\)\text{ 发散，}\\\text{则 }$a_n>\frac1n$\text{；}} &
    \makecell[l]{\text{(D) 若 }\(\sum a_n\)\text{ 收敛，且 }$a_n\ge b_n\ (n=1,2,\cdots)$\text{，}\\\text{则 }\(\sum b_n\)\text{ 也收敛。}}
    \end{tabular*}
    \item 设 $f(x)=x^2\ (0\le x<1)$，而 $S(x)=\sum_{n=1}^{\infty}b_n\sin n\pi x\ (-\infty<x<+\infty)$，其中 $b_n=2\int_0^1 f(x)\sin n\pi x\,\mathrm{d}x$。
    则 $S\!\left(-\dfrac12\right)$ 等于
    
    \begin{tabular*}{\linewidth}{@{\extracolsep{\fill}}llll@{}}
    \text{(A) }$-\frac12$ & \text{(B) }$-\frac14$ & \text{(C) }$\frac14$ & \text{(D) }$\frac12$
    \end{tabular*}
    \item 设常数 $p>0$，则幂级数 $\displaystyle \sum_{n=1}^{\infty}(-1)^{n-1}\frac{x^n}{n^p}$ 在其收敛区间的右端点处是
    
    \begin{tabular*}{\linewidth}{@{\extracolsep{\fill}}ll@{}}
    \text{(A) 条件收敛的} & \text{(B) 绝对收敛的} \\
    \makecell[l]{\text{(C) 当 }$0<p\le 1$\text{ 时为条件收敛，}\\\text{当 }$p>1$\text{ 时为绝对收敛}} &
    \makecell[l]{\text{(D) 当 }$0<p\le 1$\text{ 时为绝对收敛，}\\\text{当 }$p>1$\text{ 时为条件收敛。}}
    \end{tabular*}
    \item 幂级数 $\displaystyle \sum_{n=1}^{\infty}\frac{\ln(n+1)}{n}x^n$ 的收敛域为
    
    \begin{tabular*}{\linewidth}{@{\extracolsep{\fill}}llll@{}}
    \text{(A) }$(-1,1)$ & \text{(B) }$[-1,1]$ & \text{(C) }$(-1,1]$ & \text{(D) }$[-1,1)$
    \end{tabular*}
    \item 幂级数 $\displaystyle \sum_{n=0}^{\infty}\frac{3n+1}{n!}x^{3n}$ 的和函数为
    
    \begin{tabular*}{\linewidth}{@{\extracolsep{\fill}}llll@{}}
    \text{(A) }$xe^{x^3}$ &
    \text{(B) }$(1+3x^3)e^{x^3}$ &
    \text{(C) }$3x^3e^{x^3}$ &
    \text{(D) }$(2+3x^3)e^{x^3}$
    \end{tabular*}
\end{compactenum}
\end{exercise}
\begin{solution}
\begin{compactenum}
    \item \textbf{(C)}。记奇项和为
    \(\displaystyle O=\sum_{n=1}^{\infty}a_{2n-1}=5\)，偶项和为
    \(\displaystyle E=\sum_{n=1}^{\infty}a_{2n}\)。已知交错和为
    \[
    \sum_{n=1}^{\infty}(-1)^{n-1}a_n=O-E=2,
    \]
    所以 $E=3$。原级数的和为 $O+E=5+3=8$。

    \item \textbf{(B)}。其中
    \[
    \sum_{n=1}^{\infty}\frac{\sin(n\alpha)}{n^2}
    \]
    绝对收敛，因为 \(\left|\sin(n\alpha)\right|\le 1\)，可与 \(\sum 1/n^2\) 比较。而
    \[
    \sum_{n=1}^{\infty}\frac1{\sqrt n}
    \]
    是发散的 $p$ 级数。一个收敛级数减去一个发散到 $+\infty$ 的正项级数，结果必发散。

    \item \textbf{(D)}。由 $0\le a_n<1/n$ 得
    \[
    0\le a_n^2<\frac1{n^2}.
    \]
    因此 \(\sum a_n^2\) 收敛，从而 \(\sum (-1)^na_n^2\) 绝对收敛。选项 (A)、(C) 都可由比较例子排除；选项 (B) 缺少单调性条件，不能直接套 Leibniz 判别法。

    \item \textbf{(A)}。若 \(\sum a_n^2\) 与 \(\sum b_n^2\) 都收敛，则
    \[
    (a_n+b_n)^2\le 2a_n^2+2b_n^2,
    \]
    故由比较判别法知 \(\sum(a_n+b_n)^2\) 收敛。其余选项均不成立：例如 (B) 中取 $a_n=0,\ b_n=1$，则 \(\sum |a_nb_n|\) 收敛而 \(\sum b_n^2\) 发散；(C) 中取 \(a_n=1/(n\ln n)\) 可见发散正项级数不必满足 \(a_n>1/n\)；(D) 中取 $a_n=0,\ b_n=-1$ 即可排除。

    \item \textbf{(B)}。由题设
    \[
    b_n=2\int_0^1x^2\sin n\pi x\,\mathrm{d}x
    \]
    可知 \(S(x)\) 是函数 \(f(x)=x^2\ (0<x<1)\) 的半区间正弦级数，也就是把 \(x^2\) 作奇延拓后再展开。故在连续点处，
    \[
    S(-x)=-S(x).
    \]
    因为 \(1/2\in(0,1)\)，有 \(S(1/2)=(1/2)^2=1/4\)，所以
    \[
    S\!\left(-\frac12\right)=-\frac14.
    \]

    \item \textbf{(C)}。该幂级数的收敛半径为 $1$。在右端点 \(x=1\) 处，级数变为
    \[
    \sum_{n=1}^{\infty}(-1)^{n-1}\frac1{n^p}.
    \]
    当 \(p>0\) 时它按交错级数判别法收敛；其绝对值级数为 \(\sum 1/n^p\)，当且仅当 \(p>1\) 收敛。因此 \(0<p\le 1\) 时条件收敛，\(p>1\) 时绝对收敛。

    \item \textbf{(D)}。由根值判别可知收敛半径为 $1$。当 \(x=1\) 时，
    \[
    \sum_{n=1}^{\infty}\frac{\ln(n+1)}n
    \]
    发散，因为其项大于调和级数的一个常数倍（例如 $n$ 足够大时 \(\ln(n+1)>1\)）。当 \(x=-1\) 时，
    \[
    \sum_{n=1}^{\infty}(-1)^n\frac{\ln(n+1)}n
    \]
    条件收敛：\(\ln(n+1)/n\to0\)，且从某项起单调递减；但绝对值级数与 $x=1$ 处相同，发散。故收敛域为 \([-1,1)\)。

    \item \textbf{(B)}。令 \(u=x^3\)。利用
    \[
    \sum_{n=0}^{\infty}\frac{u^n}{n!}=e^u,\qquad
    \sum_{n=0}^{\infty}\frac{n u^n}{n!}=u e^u,
    \]
    得
    \[
    \sum_{n=0}^{\infty}\frac{3n+1}{n!}u^n
    =3ue^u+e^u=(1+3u)e^u.
    \]
    代回 \(u=x^3\)，和函数为 \((1+3x^3)e^{x^3}\)。
\end{compactenum}
\end{solution}

\begin{exercise}[判定敛散性]
判定下列级数的敛散性、绝对收敛性与条件收敛性：
\begin{shortexenum}
    \shortitem{1}{$\displaystyle \sum_{n=1}^{\infty}(-1)^n\left(1-\cos\frac{\alpha}{n}\right)\qquad (\alpha>0)$；} &
    \shortitem{2}{$\displaystyle \sum_{n=1}^{\infty}\frac{n}{\mathrm{e}^n-1}$；} \\
    \shortitem{3}{$\displaystyle \sum_{n=2}^{\infty}\frac1{\sqrt[n]{\ln n}}$；} &
    \shortitem{4}{$\displaystyle \sum_{n=1}^{\infty}a^{\ln n}\qquad (a>0)$；} \\
    \shortitem{5}{$\displaystyle \sum_{n=1}^{\infty}\left(\sqrt[n]{a}-\sqrt{1+\frac1n}\right)\qquad (a>0)$；} &
    \shortitem{6}{$\displaystyle \sum_{n=1}^{\infty}\left[\frac1n-\ln\left(1+\frac1n\right)\right]$；} \\
    \shortitem{7}{$\displaystyle \sum_{n=1}^{\infty}(-1)^{n-1}\frac1{n-\ln n}$；} &
    \shortitem{8}{$\displaystyle \sum_{n=1}^{\infty}(-1)^n\int_n^{n+1}\frac{\mathrm{e}^{-x}}{x}\,\mathrm{d}x$；} \\
    \shortitem{9}{$\displaystyle \sum_{n=2}^{\infty}\frac{(-1)^n}{(-1)^n+\sqrt n}$；} &
    \shortitem{10}{$\displaystyle \sum_{n=2}^{\infty}\frac1{\ln(n!)}$。}
\end{shortexenum}
\end{exercise}
\begin{solution}
\begin{compactenum}
    \item 因
    \[
    1-\cos t=\frac{t^2}{2}+O(t^4)\quad(t\to0),
    \]
    取 \(t=\alpha/n\)，得
    \[
    1-\cos\frac{\alpha}{n}
    =\frac{\alpha^2}{2n^2}+O\left(\frac1{n^4}\right).
    \]
    取 \(\displaystyle b_n=\frac{\alpha^2}{2n^2}\)，则
    \[
    \lim_{n\to\infty}
    \frac{1-\cos\frac{\alpha}{n}}{\frac{\alpha^2}{2n^2}}=1.
    \]
    因为 \(\sum b_n\) 收敛，所以由比较判别法的极限形式知绝对值级数收敛，原级数绝对收敛。

    \item 取 \(\displaystyle b_n=\frac{n}{\mathrm e^n}\)。因为
    \[
    \lim_{n\to\infty}\frac{\frac{n}{\mathrm e^n-1}}{\frac{n}{\mathrm e^n}}
    =\lim_{n\to\infty}\frac{\mathrm e^n}{\mathrm e^n-1}=1.
    \]
    再看模型级数 \(\sum b_n=\sum n\mathrm e^{-n}\)，用比值判别法：
    \[
    \lim_{n\to\infty}\frac{(n+1)\mathrm e^{-(n+1)}}{n\mathrm e^{-n}}=\frac1{\mathrm e}<1.
    \]
    故 \(\sum b_n\) 收敛，从而原级数收敛；它本身为正项级数，所以也可说绝对收敛。

    \item 因为
    \[
    \sqrt[n]{\ln n}=\exp\left(\frac{\ln\ln n}{n}\right)\to 1,
    \]
    所以通项
    \[
    \frac1{\sqrt[n]{\ln n}}\to 1\ne0.
    \]
    级数通项不趋于 $0$，故发散。

    \item 有
    \[
    a^{\ln n}=\mathrm e^{(\ln a)(\ln n)}=n^{\ln a}.
    \]
    于是原级数是 \(p\) 型幂级数 \(\sum n^{\ln a}\)。它收敛当且仅当指数小于 \(-1\)，即
    \[
    \ln a<-1\quad\Longleftrightarrow\quad 0<a<\mathrm e^{-1}.
    \]
    当 \(a\ge \mathrm e^{-1}\) 时发散。

    \item 利用展开
    \[
    a^{1/n}=\mathrm e^{(\ln a)/n}
    =1+\frac{\ln a}{n}+O\left(\frac1{n^2}\right),
    \]
    以及
    \[
    \sqrt{1+\frac1n}
    =1+\frac1{2n}+O\left(\frac1{n^2}\right).
    \]
    两式相减得
    \[
    \sqrt[n]{a}-\sqrt{1+\frac1n}
    =\frac{\ln a-\frac12}{n}+O\left(\frac1{n^2}\right).
    \]
    若 \(c=\ln a-\frac12\ne0\)，则通项从某项起与 \(c\) 同号，且取 \(b_n=\dfrac{|c|}{n}\) 有
    \[
    \lim_{n\to\infty}
    \frac{\left|\sqrt[n]{a}-\sqrt{1+\frac1n}\right|}{\frac{|c|}{n}}=1.
    \]
    因为调和级数发散，所以原级数发散。若 \(\ln a=1/2\)，即 \(a=\sqrt{\mathrm e}\)，继续展开到二阶：
    \[
    a^{1/n}=1+\frac1{2n}+\frac1{8n^2}+O\left(\frac1{n^3}\right),
    \qquad
    \sqrt{1+\frac1n}=1+\frac1{2n}-\frac1{8n^2}+O\left(\frac1{n^3}\right).
    \]
    因而
    \[
    \sqrt[n]{a}-\sqrt{1+\frac1n}
    =\frac1{4n^2}+O\left(\frac1{n^3}\right).
    \]
    取 \(b_n=\dfrac1{4n^2}\)，则通项与 \(b_n\) 的比值趋于 \(1\)，故原级数绝对收敛。

    \item 由 Taylor 公式
    \[
    \ln(1+t)=t-\frac{t^2}{2}+O(t^3)\quad(t\to0),
    \]
    取 \(t=1/n\)，得
    \[
    \frac1n-\ln\left(1+\frac1n\right)
    =\frac1{2n^2}+O\left(\frac1{n^3}\right).
    \]
    取 \(\displaystyle b_n=\frac1{2n^2}\)，则
    \[
    \lim_{n\to\infty}
    \frac{\frac1n-\ln\left(1+\frac1n\right)}{\frac1{2n^2}}=1.
    \]
    因为 \(\sum b_n\) 收敛，所以原级数收敛。

    \item 先看绝对值级数，取 \(\displaystyle b_n=\frac1n\)，则
    \[
    \lim_{n\to\infty}
    \frac{\frac1{n-\ln n}}{\frac1n}
    =\lim_{n\to\infty}\frac{n}{n-\ln n}=1.
    \]
    因为调和级数发散，所以绝对值级数发散。另一方面，\(n-\ln n\) 从 \(n\ge2\) 起严格增大，所以 \(1/(n-\ln n)\) 从某项起单调递减并趋于 \(0\)。由 Leibniz 判别法，原交错级数收敛但不绝对收敛，即条件收敛。

    \item 记
    \[
    A_n=\int_n^{n+1}\frac{\mathrm e^{-x}}{x}\,\mathrm{d}x>0.
    \]
    则
    \[
    0<A_n\le \int_n^{n+1}\mathrm e^{-x}\,\mathrm{d}x.
    \]
    对 \(n\) 求和后，右端为收敛的指数尾和，所以 \(\sum A_n\) 收敛。故原级数绝对收敛。

    \item 先把通项按 \(1/\sqrt n\) 展开：
    \(\displaystyle \frac{(-1)^n}{(-1)^n+\sqrt n}
    =\frac{(-1)^n}{\sqrt n}-\frac1n+O\left(\frac1{n^{3/2}}\right)\)。
    其中 \(\sum (-1)^n/\sqrt n\) 收敛，\(\sum O(n^{-3/2})\) 绝对收敛，但 \(-\sum1/n\) 发散。因此原级数发散。

    \item 由 Stirling 公式，
    \[
    \ln(n!)=n\ln n-n+O(\ln n),
    \qquad
    \lim_{n\to\infty}\frac{\ln(n!)}{n\ln n}=1.
    \]
    故取 \(\displaystyle b_n=\frac1{n\ln n}\) 有
    \[
    \lim_{n\to\infty}
    \frac{\frac1{\ln(n!)}}{\frac1{n\ln n}}
    =\lim_{n\to\infty}\frac{n\ln n}{\ln(n!)}=1.
    \]
    而 \(\sum_{n=2}^{\infty}1/(n\ln n)\) 发散，所以由比较判别法的极限形式知原级数发散。
\end{compactenum}
\end{solution}

\begin{exercise}[证明与估阶]
\begin{compactenum}
    \item 设 $a_n\ne 0\ (n=1,2,\cdots)$，且 $\displaystyle \lim_{n\to\infty}a_n=l\ne 0$。求证：级数
    \(\displaystyle \sum_{n=1}^{\infty}|a_{n+1}-a_n|\)
    与
    \(\displaystyle \sum_{n=1}^{\infty}\left|\frac1{a_{n+1}}-\frac1{a_n}\right|\)
    同敛散。
    \item 利用 Taylor 公式估计无穷小量的阶，从而判别下列级数的敛散性：
    \begin{shortexenum}
        \shortitem{1}{$\displaystyle \sum_{n=1}^{\infty}\frac1{\ln(n+1)}\sin\frac1n$；} &
        \shortitem{2}{$\displaystyle \sum_{n=1}^{\infty}\left(\sqrt{n+a}-\sqrt[4]{n^2+n+b}\right)$。}
    \end{shortexenum}
\end{compactenum}
\end{exercise}
\begin{solution}
\begin{compactenum}
    \item 因为 $a_n\to l\ne0$，所以存在正整数 $N$，当 $n\ge N$ 时
    \[
    \frac{|l|}{2}\le |a_n|\le \frac{3|l|}{2}.
    \]
    于是当 $n\ge N$ 时，
    \[
    \left|\frac1{a_{n+1}}-\frac1{a_n}\right|
    =
    \frac{|a_{n+1}-a_n|}{|a_na_{n+1}|}.
    \]
    取
    \[
    b_n=\frac{|a_{n+1}-a_n|}{|l|^2}.
    \]
    因为
    \[
    \lim_{n\to\infty}
    \frac{\left|\frac1{a_{n+1}}-\frac1{a_n}\right|}{b_n}
    =
    \lim_{n\to\infty}\frac{|l|^2}{|a_na_{n+1}|}=1,
    \]
    所以第二个级数的通项正好估阶为 \(b_n\)。又 \(b_n\) 只比 \(|a_{n+1}-a_n|\) 多一个正常数因子，不影响敛散性。等价地，也可写成：由于 \(|a_na_{n+1}|\) 从某项起被夹在两个正常数之间，存在常数 $C_1,C_2>0$，使得
    \[
    C_1|a_{n+1}-a_n|
    \le
    \left|\frac1{a_{n+1}}-\frac1{a_n}\right|
    \le
    C_2|a_{n+1}-a_n|.
    \]
    因此两个级数只差有限多个前项，尾项又由比较判别法的极限形式相互对应，所以它们同敛散。

    \item
    \begin{compactenum}
        \item 取 \(\displaystyle b_n=\frac1{n\ln n}\)。直接计算比值：
        \[
        \lim_{n\to\infty}
        \frac{\frac1{\ln(n+1)}\sin\frac1n}{\frac1{n\ln n}}
        =
        \lim_{n\to\infty}\frac{\sin(1/n)}{1/n}\cdot\frac{\ln n}{\ln(n+1)}
        =1.
        \]
        由于 \(\sum_{n=2}^{\infty}1/(n\ln n)\) 发散，故由比较判别法的极限形式知原级数发散。

        \item
        先把两项都提出主量 \(\sqrt n\)：
        \[
        \sqrt{n+a}
        =\sqrt n\left(1+\frac an\right)^{1/2}
        =\sqrt n+\frac{a}{2\sqrt n}+O\left(\frac1{n^{3/2}}\right),
        \]
        \[
        \sqrt[4]{n^2+n+b}
        =\sqrt n\left(1+\frac1n+\frac b{n^2}\right)^{1/4}
        =\sqrt n+\frac1{4\sqrt n}+O\left(\frac1{n^{3/2}}\right).
        \]
        从而
        \[
        \sqrt{n+a}-\sqrt[4]{n^2+n+b}
        =
        \frac{a-\frac12}{2\sqrt n}
        +O\left(\frac1{n^{3/2}}\right).
        \]
        若 \(a\ne1/2\)，令 \(c=\dfrac{a-\frac12}{2}\)，则通项从某项起与 \(c\) 同号，且取 \(\displaystyle b_n=\frac{|c|}{\sqrt n}\) 有
        \[
        \lim_{n\to\infty}
        \frac{\left|\sqrt{n+a}-\sqrt[4]{n^2+n+b}\right|}{\frac{|c|}{\sqrt n}}=1.
        \]
        因为 \(\sum1/\sqrt n\) 发散，故原级数发散。若 \(a=1/2\)，则首项抵消，通项为 \(O(n^{-3/2})\)，可与收敛的 \(p\) 级数 \(\sum n^{-3/2}\) 比较，故绝对收敛。注意参数 \(b\) 只进入更高阶项，不影响敛散分界。
    \end{compactenum}
\end{compactenum}
\end{solution}

\begin{exercise}[证明题]
\begin{shortsingleenum}
    \shortsingleitem{1}{设偶函数 $f(x)$ 的二阶导数 $f''(x)$ 在 $x=0$ 的一个邻域内连续，且 $f(0)=1,\ f''(0)=2$。试证明级数 $\sum_{n=1}^{\infty}\left[f\left(\frac1n\right)-1\right]$ 绝对收敛。}
    \shortsingleitem{2}{设 $f(x)$ 在点 $x=0$ 的某邻域内有二阶连续导数，且 $\displaystyle \lim_{x\to 0}\frac{f(x)}{x}=0$。证明级数 $\displaystyle \sum_{n=1}^{\infty}f\left(\frac1n\right)$ 绝对收敛。}
    \shortsingleitem{3}{设 $a_n>0$，级数 $\displaystyle \sum_{n=1}^{\infty}a_n$ 收敛，$b_n=1-\dfrac{\ln(1+a_n)}{a_n}$。证明级数 $\displaystyle \sum_{n=1}^{\infty}b_n$ 收敛。}
    \shortsingleitem{4}{设数列 $\{na_n\}$ 收敛，且级数 $\displaystyle \sum_{n=1}^{\infty}n(a_n-a_{n-1})$ 收敛，证明级数 $\displaystyle \sum_{n=1}^{\infty}a_n$ 收敛。}
\end{shortsingleenum}
\end{exercise}
\begin{solution}
\begin{compactenum}
    \item 因 $f$ 为偶函数，$f(-x)=f(x)$。两边对 $x$ 求导后令 $x=0$，得
    \[
    -f'(0)=f'(0),
    \]
    所以 $f'(0)=0$。在 $0$ 附近用 Taylor 公式，有
    \[
    f(x)=f(0)+f'(0)x+\frac{f''(\xi_x)}2x^2
    =1+\frac{f''(\xi_x)}2x^2,
    \]
    其中 \(\xi_x\) 介于 $0$ 与 $x$ 之间。由于 $f''$ 在 $0$ 的邻域内连续，故在较小邻域内有界，记 \(|f''|\le M\)。取 \(x=1/n\)，当 $n$ 充分大时
    \[
    \left|f\left(\frac1n\right)-1\right|
    \le \frac{M}{2n^2}.
    \]
    右端构成收敛的 \(p\) 级数，有限多个前项不影响敛散，因此原级数绝对收敛。

  \item 由
    \[
    \lim_{x\to0}\frac{f(x)}x=0
    \]
    可知 $f(0)=0$，否则分子趋于非零常数，商不可能趋于 $0$。又
    \[
    f'(0)=\lim_{x\to 0}\frac{f(x)-f(0)}{x}=0.
    \]
    因 $f$ 在 $0$ 附近有二阶连续导数，由 Taylor 公式得
    \[
    f(x)=f(0)+f'(0)x+\frac{f''(\xi_x)}2x^2
    =\frac{f''(\xi_x)}2x^2.
    \]
    同样由 $f''$ 局部有界，取 \(x=1/n\) 得
    \[
    \left|f\left(\frac1n\right)\right|\le \frac{C}{n^2}.
    \]
    故 \(\sum |f(1/n)|\) 收敛，即原级数绝对收敛。

    \item 因 \(\sum a_n\) 收敛且 \(a_n>0\)，必有 \(a_n\to0\)。对 \(t\to0^+\)，
    \[
    \ln(1+t)=t-\frac{t^2}{2}+O(t^3).
    \]
    代入 \(t=a_n\)，得
    \[
    \frac{\ln(1+a_n)}{a_n}
    =1-\frac{a_n}{2}+O(a_n^2),
    \]
    因此
    \[
    b_n=1-\frac{\ln(1+a_n)}{a_n}
    =\frac{a_n}{2}+O(a_n^2).
    \]
    又因为 \(a_n\to0\)，从某项起 \(a_n^2\le a_n\)，所以 \(\sum a_n^2\) 收敛。于是 \(\sum b_n\) 与收敛正项级数 \(\sum a_n\)、\(\sum a_n^2\) 控制在一起，故 \(\sum b_n\) 收敛。

    \item 记
    \[
    c_n=n(a_n-a_{n-1}).
    \]
    题设说明 \(\sum c_n\) 收敛，且数列 \(\{na_n\}\) 收敛。利用恒等式
    \[
    na_n-(n-1)a_{n-1}
    =n(a_n-a_{n-1})+a_{n-1}
    =c_n+a_{n-1},
    \]
    得
    \[
    a_{n-1}=na_n-(n-1)a_{n-1}-c_n.
    \]
    对 $n=1,\dots,N$ 求和：
    \[
    \sum_{n=1}^{N}a_{n-1}
    =
    \sum_{n=1}^{N}\bigl[na_n-(n-1)a_{n-1}\bigr]
    -\sum_{n=1}^{N}c_n.
    \]
    第一项是望远镜和：
    \[
    \sum_{n=1}^{N}\bigl[na_n-(n-1)a_{n-1}\bigr]=Na_N.
    \]
    当 \(N\to\infty\) 时，\(Na_N\) 有极限，\(\sum_{n=1}^{N}c_n\) 也有极限，所以 \(\sum_{n=1}^{N}a_{n-1}\) 有极限。故 \(\sum a_{n-1}\) 收敛，从而去掉或平移有限项后，\(\sum a_n\) 也收敛。
\end{compactenum}
\end{solution}

\begin{exercise}[幂级数综合]
\begin{compactenum}
    \item 求下列幂级数的收敛域：
    \begin{shortexenum}
        \shortitem{1}{$\displaystyle \sum_{n=1}^{\infty}\frac{x^n}{(n+1)^p}\qquad (p>0)$；} &
        \shortitem{2}{$\displaystyle \sum_{n=1}^{\infty}\frac{x^{2n}}{(2n-1)2^n}$；} \\
        \shortitem{3}{$\displaystyle \sum_{n=0}^{\infty}\frac{\ln(n+1)}{n+1}x^{n+1}$；} &
        \shortitem{4}{$\displaystyle \sum_{n=1}^{\infty}\frac{x^n}{a^n+b^n}\qquad (a>0,b>0)$。}
    \end{shortexenum}
    \item 求下列函数项级数的收敛域：
    \begin{shortexenum}
        \shortitem{1}{$\displaystyle \sum_{n=1}^{\infty}\frac{x^n}{1+x^{2n}}$；} &
        \shortitem{2}{$\displaystyle \sum_{n=1}^{\infty}\frac{(n+x)^n}{n^{n+x}}$。}
    \end{shortexenum}
    \item 求下列幂级数的和函数，并指出其收敛域：
    \begin{shortexenum}
        \shortitem{1}{$\displaystyle \sum_{n=1}^{\infty}\frac{x^{n+1}}{n(n+1)}$；} &
        \shortitem{2}{$\displaystyle \sum_{n=1}^{\infty}\frac{2n+1}{n!}x^{2n}$。}
    \end{shortexenum}
    \item 求下列级数的和：
    \begin{shortexenum}
        \shortitem{1}{$\displaystyle \sum_{n=1}^{\infty}\frac{(-1)^{n-1}}{n(2n-1)}$；} &
        \shortitem{2}{$\displaystyle \sum_{n=1}^{\infty}\frac{(-1)^nn}{(2n+1)!}$。}
    \end{shortexenum}
\end{compactenum}
\end{exercise}
\begin{solution}
\begin{compactenum}
    \item
    \begin{compactenum}
        \item 先看开区间。系数 \((n+1)^{-p}\) 的 \(n\) 次根趋于 $1$，所以收敛半径为 $1$。当 \(|x|<1\) 时收敛，当 \(|x|>1\) 时通项不趋于 $0$。

        再查端点。$x=1$ 时为
        \[
        \sum_{n=1}^{\infty}\frac1{(n+1)^p},
        \]
        它当且仅当 \(p>1\) 收敛。$x=-1$ 时为
        \[
        \sum_{n=1}^{\infty}\frac{(-1)^n}{(n+1)^p},
        \]
        对任意 \(p>0\) 都由 Leibniz 判别法收敛。因此
        \[
        \boxed{0<p\le1:\ [-1,1),\qquad p>1:\ [-1,1]}.
        \]

        \item 令
        \[
        y=\frac{x^2}{2},
        \]
        原级数化为
        \[
        \sum_{n=1}^{\infty}\frac{y^n}{2n-1}.
        \]
        关于 \(y\) 的收敛半径为 $1$，所以 \(|x^2/2|<1\)，即 \(|x|<\sqrt2\)。当 \(x=\pm\sqrt2\) 时，\(y=1\)，级数变为
        \[
        \sum_{n=1}^{\infty}\frac1{2n-1},
        \]
        取 \(b_n=\dfrac1{2n}\)，则
        \[
        \lim_{n\to\infty}\frac{\frac1{2n-1}}{\frac1{2n}}=1.
        \]
        因为 \(\sum 1/(2n)\) 发散，所以该端点级数发散。故收敛域为
        \[
        \boxed{(-\sqrt2,\sqrt2)}.
        \]

        \item 改写指标 \(m=n+1\)，得
        \[
        \sum_{n=0}^{\infty}\frac{\ln(n+1)}{n+1}x^{n+1}
        =
        \sum_{m=1}^{\infty}\frac{\ln m}{m}x^m.
        \]
        由根值判别，\(\sqrt[m]{\ln m/m}\to1\)，所以半径为 $1$。当 \(x=1\) 时，\(\sum \ln m/m\) 发散；当 \(x=-1\) 时，
        \[
        \sum_{m=1}^{\infty}(-1)^m\frac{\ln m}{m}
        \]
        从某项起满足 Leibniz 判别法，条件收敛。故收敛域为
        \[
        \boxed{[-1,1)}.
        \]

        \item 设 \(M=\max\{a,b\}\)。当 \(a\ne b\) 时，取 \(b_n=M^n\)，则
        \[
        \lim_{n\to\infty}\frac{a^n+b^n}{M^n}=1;
        \]
        当 \(a=b\) 时，\(\displaystyle (a^n+b^n)/M^n=2\)。因此系数 \(\dfrac1{a^n+b^n}\) 的 \(n\) 次根都趋于 \(\dfrac1M\)，
        \[
        \lim_{n\to\infty}\sqrt[n]{\frac1{a^n+b^n}}=\frac1M,
        \]
        所以收敛半径为 \(M\)。当 \(x=M\) 或 \(x=-M\) 时，通项的绝对值趋于 $1$ 或 $1/2$，不趋于 $0$。故收敛域为
        \[
        \boxed{(-M,M)},\qquad M=\max\{a,b\}.
        \]
    \end{compactenum}

    \item
    \begin{compactenum}
        \item 当 \(|x|<1\) 时，取 \(b_n=|x|^n\)，则
        \[
        \lim_{n\to\infty}
        \frac{\left|\frac{x^n}{1+x^{2n}}\right|}{|x|^n}
        =\lim_{n\to\infty}\frac1{1+x^{2n}}=1,
        \]
        故按几何级数收敛。当 \(|x|>1\) 时，取 \(b_n=|x|^{-n}\)，则
        \[
        \lim_{n\to\infty}
        \frac{\left|\frac{x^n}{1+x^{2n}}\right|}{|x|^{-n}}
        =\lim_{n\to\infty}\frac{x^{2n}}{1+x^{2n}}=1,
        \]
        仍按几何级数收敛。当 \(x=1\) 时通项为 \(1/2\)，当 \(x=-1\) 时通项为 \((-1)^n/2\)，都不趋于 $0$。故收敛域为
        \[
        \boxed{\mathbb R\setminus\{-1,1\}}.
        \]

        \item 对固定的实数 \(x\)，当 \(n\) 充分大时 \(n+x>0\)，且
        \[
        \frac{(n+x)^n}{n^{n+x}}
        =
        \frac{\left(1+\frac{x}{n}\right)^n}{n^x}.
        \]
        取 \(\displaystyle b_n=\frac{\mathrm e^x}{n^x}\)，则
        \[
        \lim_{n\to\infty}
        \frac{\frac{(n+x)^n}{n^{n+x}}}{\frac{\mathrm e^x}{n^x}}
        =
        \lim_{n\to\infty}\frac{\left(1+\frac{x}{n}\right)^n}{\mathrm e^x}=1.
        \]
        因而由比较判别法的极限形式可知，它与 \(p\) 级数 \(\sum 1/n^x\) 同敛散，故当且仅当
        \[
        \boxed{x>1}
        \]
        时收敛。
    \end{compactenum}

    \item
    \begin{compactenum}
        \item 先在 \(|x|<1\) 内求和。利用
        \[
        \frac1{n(n+1)}=\frac1n-\frac1{n+1},
        \]
        得
        \begin{align*}
        \sum_{n=1}^{\infty}\frac{x^{n+1}}{n(n+1)}
        &=
        x\sum_{n=1}^{\infty}\frac{x^n}{n}
        -\sum_{n=1}^{\infty}\frac{x^{n+1}}{n+1}\\
        &=
        -x\ln(1-x)-\left[-\ln(1-x)-x\right]\\
        &=x+(1-x)\ln(1-x).
        \end{align*}
        收敛半径为 $1$；在 \(x=1\) 时为 \(\sum1/[n(n+1)]\)，收敛；在 \(x=-1\) 时绝对值仍为 \(\sum1/[n(n+1)]\)。取 \(b_n=1/n^2\)，有
        \[
        \lim_{n\to\infty}\frac{\frac1{n(n+1)}}{\frac1{n^2}}=1,
        \]
        故端点 \(x=-1\) 也绝对收敛。因此收敛域为
        \[
        \boxed{[-1,1]}.
        \]
        和函数在开区间内为 \(x+(1-x)\ln(1-x)\)，端点处可取相应极限值。

        \item 令 \(u=x^2\)。由
        \[
        \sum_{n=0}^{\infty}\frac{u^n}{n!}=e^u,\qquad
        \sum_{n=1}^{\infty}\frac{n u^n}{n!}=u e^u,
        \]
        得
        \begin{align*}
        \sum_{n=1}^{\infty}\frac{2n+1}{n!}u^n
        &=2\sum_{n=1}^{\infty}\frac{n u^n}{n!}
        +\sum_{n=1}^{\infty}\frac{u^n}{n!}\\
        &=2u e^u+(e^u-1)
        =(2u+1)e^u-1.
        \end{align*}
        代回 \(u=x^2\)，得
        \[
        \boxed{\sum_{n=1}^{\infty}\frac{2n+1}{n!}x^{2n}
        =(2x^2+1)e^{x^2}-1},
        \]
        收敛域为 \(\boxed{\mathbb R}\)。
    \end{compactenum}

    \item
    \begin{compactenum}
        \item 先拆分
        \[
        \frac1{n(2n-1)}=-\frac1n+\frac{2}{2n-1}.
        \]
        因而
        \begin{align*}
        \sum_{n=1}^{\infty}\frac{(-1)^{n-1}}{n(2n-1)}
        &=
        -\sum_{n=1}^{\infty}\frac{(-1)^{n-1}}n
        +2\sum_{n=1}^{\infty}\frac{(-1)^{n-1}}{2n-1}\\
        &=-\ln2+2\cdot\frac{\pi}{4}
        =\boxed{\frac{\pi}{2}-\ln2}.
        \end{align*}

        \item 由
        \[
        n=\frac{(2n+1)-1}{2},
        \]
        得
        \begin{align*}
        \sum_{n=1}^{\infty}\frac{(-1)^n n}{(2n+1)!}
        &=
        \frac12\sum_{n=1}^{\infty}\frac{(-1)^n}{(2n)!}
        -\frac12\sum_{n=1}^{\infty}\frac{(-1)^n}{(2n+1)!}.
        \end{align*}
        又
        \[
        \sum_{n=1}^{\infty}\frac{(-1)^n}{(2n)!}=\cos1-1,\qquad
        \sum_{n=1}^{\infty}\frac{(-1)^n}{(2n+1)!}=\sin1-1.
        \]
        所以
        \[
        \boxed{
        \sum_{n=1}^{\infty}\frac{(-1)^n n}{(2n+1)!}
        =\frac{\cos1-\sin1}{2}}.
        \]
    \end{compactenum}
\end{compactenum}
\end{solution}

\begin{exercise}[幂级数展开与 Fourier 综合]
\begin{compactenum}
    \item 将下列函数展开成 $x$ 的幂级数，并指出其收敛域：
    \begin{shortexenum}
        \shortitem{1}{$\ln(a+x)\ (a>0)$；} &
        \shortitem{2}{$\dfrac{x^2+1}{(x^2-1)^2}$；} \\
        \shortitem{3}{$\dfrac{x}{\sqrt{1-2x}}$；} &
        \shortitem{4}{$\arctan\dfrac{1+x}{1-x}$。}
    \end{shortexenum}
    \item 将下列函数在指定点展开成幂级数：
    \begin{shortexenum}
        \shortitem{1}{$\displaystyle f(x)=\frac1{2x^2+x-3},\ x_0=3$；} &
        \shortitem{2}{$\displaystyle f(x)=(x-2)\mathrm{e}^{-x},\ x_0=1$；} \\
        \shortitem{3}{$\displaystyle f(x)=\frac{\mathrm{d}}{\mathrm{d}x}\left(\frac{\mathrm{e}^x-\mathrm{e}}{x-1}\right),\ x_0=1$。} &
    \end{shortexenum}
    \item 试将函数 $f(x)=10-x\ (5\le x\le 15)$ 展开成以 $10$ 为周期的傅里叶级数。
    \item 试利用 $\dfrac{\pi-x}{2}$ 在 $[0,\pi]$ 上的正弦级数，对于 $-\pi\le \alpha<0<\beta\le \pi$，求极限
    \[
    \lim_{n\to\infty}\int_{\alpha}^{\beta}\frac1\pi
    \left[\frac12+\sum_{k=1}^{n}\cos kx\right]\mathrm{d}x.
    \]
    \item 在区间 $(0,2)$ 上将函数
    \[
    f(x)=
    \begin{cases}
      x, & 0\le x<1,\\
      2-x, & 1\le x\le2
    \end{cases}
    \]
    展开成余弦级数。
    \item 设
    \[
    S(x)=\sum_{n=1}^{\infty}b_n\sin nx\quad (-\pi<x<\pi),
    \qquad
    \frac{\pi-x}{2}=\sum_{n=1}^{\infty}b_n\sin nx\quad (0<x<\pi).
    \]
    试求 $b_n$ 及 $S(x)$。
  \item 证明：在区间 $[-\pi,\pi]$ 上等式 \(\displaystyle \sum_{n=1}^{\infty}\frac{(-1)^{n-1}}{n^2}\cos nx=\frac{\pi^2}{12}-\frac{x^2}{4}\) 成立，并求级数 $\displaystyle \sum_{n=1}^{\infty}\frac{(-1)^{n-1}}{n^2}$ 的和。
    \item 证明：当 $0\le x\le \pi$ 时，
    \begin{align*}
    \mathrm{e}^{2x}
    &=\frac{\mathrm{e}^{2\pi}-1}{2\pi}
    +\frac{4}{\pi}\sum_{n=1}^{\infty}
    \frac{(-1)^n\mathrm{e}^{2\pi}-1}{4+n^2}\cos nx.
    \end{align*}
\end{compactenum}
\end{exercise}
\begin{solution}
\begin{compactenum}
    \item
    \begin{compactenum}
        \item 先提出常数：
        \[
        \ln(a+x)=\ln a+\ln\left(1+\frac xa\right).
        \]
        利用
        \[
        \ln(1+t)=\sum_{n=1}^{\infty}(-1)^{n-1}\frac{t^n}{n}
        \quad(-1<t\le1),
        \]
        取 \(t=x/a\)，得
        \[
        \boxed{
        \ln(a+x)=\ln a+\sum_{n=1}^{\infty}(-1)^{n-1}\frac{x^n}{na^n}
        }.
        \]
        收敛域为 \(\boxed{-a<x\le a}\)。

        \item 因为 \((x^2-1)^2=(1-x^2)^2\)，所以
        \[
        \frac{x^2+1}{(x^2-1)^2}
        =(1+x^2)\frac1{(1-x^2)^2}.
        \]
        由
        \[
        \frac1{(1-t)^2}=\sum_{n=0}^{\infty}(n+1)t^n\quad(|t|<1),
        \]
        取 \(t=x^2\)，再乘以 \(1+x^2\)，得
        \[
        \frac{x^2+1}{(x^2-1)^2}
        =
        1+\sum_{n=1}^{\infty}(2n+1)x^{2n}.
        \]
        收敛域为 \(\boxed{|x|<1}\)，端点 \(x=\pm1\) 处原函数本身也无定义。

        \item 利用二项式展开
        \[
        (1-t)^{-1/2}
        =
        \sum_{n=0}^{\infty}\frac{(2n)!}{4^n(n!)^2}t^n
        \quad(|t|<1).
        \]
        取 \(t=2x\)，再乘以 \(x\)，得
        \[
        \boxed{
        \frac{x}{\sqrt{1-2x}}
        =
        x\sum_{n=0}^{\infty}\frac{(2n)!}{2^n(n!)^2}x^n
        }.
        \]
        半径为 \(1/2\)。当 \(x=-1/2\) 时，级数为交错型且项量级为 \(1/\sqrt n\)，收敛；当 \(x=1/2\) 时发散。故收敛域为
        \[
        \boxed{\left[-\frac12,\frac12\right)}.
        \]

        \item 在 \(|x|<1\) 内有恒等式
        \[
        \arctan\frac{1+x}{1-x}=\frac{\pi}{4}+\arctan x.
        \]
        又
        \[
        \arctan x=\sum_{n=0}^{\infty}(-1)^n\frac{x^{2n+1}}{2n+1}
        \quad(|x|\le1).
        \]
        因此
        \[
        \boxed{
        \arctan\frac{1+x}{1-x}
        =
        \frac{\pi}{4}
        +\sum_{n=0}^{\infty}(-1)^n\frac{x^{2n+1}}{2n+1}
        }.
        \]
        幂级数本身的收敛域为 \([-1,1]\)；作为原函数展开，\(x=1\) 处原式无定义，只能理解为左极限 \(\pi/2\)。
    \end{compactenum}

    \item
    \begin{compactenum}
    \item 令 \(t=x-3\)。先作部分分式分解：
    \[
    \frac1{2x^2+x-3}
    =
    \frac1{(x-1)(2x+3)}
    =
    \frac1{5(x-1)}-\frac2{5(2x+3)}.
    \]
    由于 \(x-1=t+2,\ 2x+3=2t+9\)，所以
    \[
    \frac1{5(x-1)}
    =\frac1{10}\frac1{1+t/2}
    =\frac1{10}\sum_{n=0}^{\infty}\left(-\frac t2\right)^n,
    \]
    \[
    -\frac2{5(2x+3)}
    =
    -\frac2{45}\frac1{1+2t/9}
    =
    -\frac2{45}\sum_{n=0}^{\infty}\left(-\frac{2t}{9}\right)^n.
    \]
    因此
    \[
    \boxed{
    \frac1{2x^2+x-3}
    =
    \frac1{10}\sum_{n=0}^{\infty}\left(-\frac{x-3}{2}\right)^n
    -\frac2{45}\sum_{n=0}^{\infty}\left(-\frac{2(x-3)}9\right)^n
    }.
    \]
    收敛条件取两者交集，得 \(|x-3|<2\)。

        \item 令 \(t=x-1\)，则 \(x-2=t-1,\ e^{-x}=e^{-1}e^{-t}\)。于是
        \[
        (x-2)e^{-x}
        =(t-1)e^{-1}\sum_{n=0}^{\infty}\frac{(-1)^n}{n!}t^n.
        \]
        也可把乘积写开为
        \[
        \boxed{
        (x-2)e^{-x}
        =
        e^{-1}\left[
        -1+\sum_{n=1}^{\infty}
        (-1)^{n-1}\left(\frac1{(n-1)!}+\frac1{n!}\right)(x-1)^n
        \right]}.
        \]
        指数级数对一切 \(t\) 收敛，所以收敛域为 \(\mathbb R\)。

        \item 令 $t=x-1$，则
        \[
        \frac{\mathrm{e}^x-\mathrm{e}}{x-1}
        =\mathrm e\,\frac{\mathrm{e}^t-1}{t}
        =\mathrm e\sum_{n=0}^{\infty}\frac{t^n}{(n+1)!}.
        \]
        并且
        \[
        \frac{\mathrm{d}}{\mathrm{d}x}\left(\frac{\mathrm{e}^x-\mathrm{e}}{x-1}\right)
        =\mathrm e\sum_{n=0}^{\infty}\frac{n+1}{(n+2)!}(x-1)^n.
        \]
        这里 \(\mathrm{d}t/\mathrm{d}x=1\)，所以对 \(x\) 求导与对 \(t\) 求导相同；收敛域仍为 \(\mathbb R\)。
    \end{compactenum}

  \item 令 \(t=x-10\)，则 \(t\in[-5,5]\)，且
  \[
  f(x)=10-x=-t.
  \]
  这是以 \(2L=10\) 为周期的函数，故 \(L=5\)。在变量 \(t\) 下它是奇函数，所以只含正弦项：
  \[
  -t\sim \sum_{n=1}^{\infty}b_n\sin\frac{n\pi t}{5}.
  \]
  系数为
  \begin{align*}
  b_n
  &=\frac{2}{5}\int_0^5(-t)\sin\frac{n\pi t}{5}\,\mathrm{d}t\\
  &=\frac{10}{\pi}\frac{(-1)^n}{n}.
  \end{align*}
  又 \(\sin\frac{n\pi(x-10)}5=\sin\frac{n\pi x}{5}\)，故
  \[
  \boxed{
  10-x=\frac{10}{\pi}\sum_{n=1}^{\infty}
  \frac{(-1)^n}{n}\sin\frac{n\pi x}{5}}
  \]
  在连续点处成立，端点处取周期延拓后的左右极限平均值。

    \item 记
    \[
    I_n=\frac1\pi\int_{\alpha}^{\beta}
    \left[\frac12+\sum_{k=1}^{n}\cos kx\right]\mathrm{d}x.
    \]
    逐项积分得
    \[
    I_n=
    \frac1\pi\left[
    \frac{x}{2}+\sum_{k=1}^{n}\frac{\sin kx}{k}
    \right]_{\alpha}^{\beta}.
    \]
    已知
    \[
    \sum_{k=1}^{\infty}\frac{\sin kx}{k}
    =
    \begin{cases}
        \dfrac{\pi-x}{2},&0<x<\pi,\\[4pt]
        -\dfrac{\pi+x}{2},&-\pi<x<0.
    \end{cases}
    \]
    因而对 \(0<\beta\le\pi\)，括号内的极限为 \(\pi/2\)；对 \(-\pi\le\alpha<0\)，括号内的极限为 \(-\pi/2\)。所以
    \[
    \boxed{
    \lim_{n\to\infty}I_n
    =\frac1\pi\left(\frac\pi2-\left(-\frac\pi2\right)\right)=1}.
    \]

    \item 这是区间 \((0,2)\) 上的半区间余弦展开，取 \(L=2\)：
    \[
    f(x)\sim \frac{a_0}{2}+\sum_{n=1}^{\infty}a_n\cos\frac{n\pi x}{2}.
    \]
    先算常数项：
    \[
    a_0=\int_0^2 f(x)\,\mathrm{d}x
    =\int_0^1x\,\mathrm{d}x+\int_1^2(2-x)\,\mathrm{d}x=1,
    \]
    故 \(a_0/2=1/2\)。再算 \(a_n\)：
    \[
    a_n=\int_0^1x\cos\frac{n\pi x}{2}\,\mathrm{d}x
    +\int_1^2(2-x)\cos\frac{n\pi x}{2}\,\mathrm{d}x.
    \]
    第二个积分令 \(u=2-x\)，可得
    \[
    a_n=\bigl(1+(-1)^n\bigr)
    \int_0^1x\cos\frac{n\pi x}{2}\,\mathrm{d}x.
    \]
    所以 \(n\) 为奇数时 \(a_n=0\)。令 \(n=2m\)，则
    \[
    a_{2m}=2\int_0^1x\cos(m\pi x)\,\mathrm{d}x
    =\frac{2\bigl((-1)^m-1\bigr)}{m^2\pi^2}.
    \]
    当 \(m\) 为偶数时仍为 $0$；当 \(m=2k+1\) 时，
    \[
    a_{4k+2}=-\frac{4}{(2k+1)^2\pi^2}.
    \]
    因此
    \[
    \boxed{
    f(x)=\frac12-\frac4{\pi^2}
    \sum_{k=0}^{\infty}
    \frac{\cos(2k+1)\pi x}{(2k+1)^2}},
    \qquad 0<x<2.
    \]

    \item 由题给的正弦级数
    \[
    \frac{\pi-x}{2}=\sum_{n=1}^{\infty}b_n\sin nx\quad(0<x<\pi)
    \]
    与标准公式对比，可得
    \[
    \boxed{b_n=\frac1n}.
    \]
    因为正弦级数表示的是奇延拓，所以在 \((-\pi,\pi)\) 内
    \[
    \boxed{
    S(x)=
    \begin{cases}
    -\dfrac{\pi+x}{2},&-\pi<x<0,\\[4pt]
    0,&x=0,\\[4pt]
    \dfrac{\pi-x}{2},&0<x<\pi.
    \end{cases}}
    \]
    其中 \(x=0\) 处取左右极限平均值。

    \item 从函数 \(x^2\) 的 Fourier 展开出发。它是偶函数，所以只含余弦项：
    \[
    x^2\sim \frac{a_0}{2}+\sum_{n=1}^{\infty}a_n\cos nx.
    \]
    其中
    \[
    a_0=\frac1\pi\int_{-\pi}^{\pi}x^2\,\mathrm{d}x
    =\frac{2\pi^2}{3},
    \qquad \frac{a_0}{2}=\frac{\pi^2}{3}.
    \]
    对 \(n\ge1\)，分部积分两次可得
    \[
    a_n=\frac2\pi\int_0^\pi x^2\cos nx\,\mathrm{d}x
    =\frac{4(-1)^n}{n^2}.
    \]
    因此
    \[
    x^2=\frac{\pi^2}{3}
    +4\sum_{n=1}^{\infty}\frac{(-1)^n}{n^2}\cos nx
    \quad(-\pi\le x\le\pi).
    \]
    移项得
    \[
    \boxed{
    \sum_{n=1}^{\infty}\frac{(-1)^{n-1}}{n^2}\cos nx
    =
    \frac{\pi^2}{12}-\frac{x^2}{4}}.
    \]
    取 \(x=0\)，即得
    \[
    \boxed{\sum_{n=1}^{\infty}\frac{(-1)^{n-1}}{n^2}
    =\frac{\pi^2}{12}}.
    \]

  \item 作 \(e^{2x}\) 在 \([0,\pi]\) 上的余弦级数展开：
    \[
    \mathrm{e}^{2x}\sim \frac{a_0}{2}+\sum_{n=1}^{\infty}a_n\cos nx.
    \]
    常数项为
    \[
    a_0=\frac{2}{\pi}\int_0^\pi \mathrm{e}^{2x}\,\mathrm{d}x
    =\frac{\mathrm{e}^{2\pi}-1}{\pi},
    \qquad
    \frac{a_0}{2}=\frac{\mathrm{e}^{2\pi}-1}{2\pi}.
    \]
    对 \(n\ge1\)，利用
    \[
    \int e^{2x}\cos nx\,\mathrm{d}x
    =
    \frac{e^{2x}(2\cos nx+n\sin nx)}{4+n^2},
    \]
    得
    \begin{align*}
    a_n
    &=\frac{2}{\pi}\int_0^\pi e^{2x}\cos nx\,\mathrm{d}x\\
    &=\frac{2}{\pi}
    \left[
    \frac{e^{2x}(2\cos nx+n\sin nx)}{4+n^2}
    \right]_0^\pi\\
    &=\frac4\pi\frac{(-1)^n e^{2\pi}-1}{4+n^2}.
    \end{align*}
    故
    \[
    \boxed{
    \mathrm{e}^{2x}
    =
    \frac{\mathrm{e}^{2\pi}-1}{2\pi}
    +\frac{4}{\pi}\sum_{n=1}^{\infty}
    \frac{(-1)^n\mathrm{e}^{2\pi}-1}{4+n^2}\cos nx
    }\quad(0\le x\le\pi).
    \]
\end{compactenum}
\end{solution}

\chapter{往年真题整理}

\section{2006--2007 第二学期 A 卷}

\par\medskip
\noindent\textbf{一、单项选择题（每小题 3 分，共 15 分）}\par\nobreak\smallskip
\begin{exercise}
\par
设有直线
  \[
    L:\begin{cases}
      x+3y+2z+1=0,\\
      2x-y-10z+3=0
    \end{cases}
  \]
  及平面 $\pi:4x-2y+z-2=0$，则直线 $L$（\quad）。

  A. 平行于平面 $\pi$；\quad
  B. 在平面 $\pi$ 上；\quad
  C. 垂直于平面 $\pi$；\quad
  D. 与平面 $\pi$ 斜交。
\end{exercise}

\begin{solution}
\par

直线 $L$ 是两个平面的交线，因此先取这两个平面的法向量
\[
  \boldsymbol n_1=(1,3,2),\qquad \boldsymbol n_2=(2,-1,-10).
\]
交线的方向向量应同时垂直于 $\boldsymbol n_1,\boldsymbol n_2$，可由叉乘得到
\[
  \boldsymbol s=\boldsymbol n_1\times \boldsymbol n_2=(-28,14,-7),
\]
故方向向量可化简取为
\[
  \boldsymbol s=(4,-2,1).
\]
平面 $\pi:4x-2y+z-2=0$ 的法向量为
\[
  \boldsymbol n_\pi=(4,-2,1).
\]
可见直线的方向向量与平面的法向量平行。直线方向若平行于平面法向量，则直线垂直于该平面。因此选 C。
\end{solution}
\begin{exercise}
\par
二元函数
  \[
    f(x,y)=
    \begin{cases}
      \dfrac{x^2y}{x^2+y^2}, & (x,y)\ne(0,0),\\
      0, & (x,y)=(0,0)
    \end{cases}
  \]
  在点 $(0,0)$ 处（\quad）。

  A. 偏导数存在，但不可微；\quad
  B. 连续、偏导数不存在；\quad
  C. 不连续、偏导数不存在；\quad
  D. 偏导数不存在。
\end{exercise}

\begin{solution}
\par

先判连续性。令 $x=r\cos\theta,\ y=r\sin\theta$，则当 $(x,y)\ne(0,0)$ 时
\[
  f(x,y)=\frac{x^2y}{x^2+y^2}
  =r\cos^2\theta\sin\theta .
\]
所以 $|f(x,y)|\le r\to0$，故 $f$ 在 $(0,0)$ 连续。

再判偏导数。按定义，
\[
  f_x(0,0)=\lim_{h\to0}\frac{f(h,0)-f(0,0)}{h}=0,\qquad
  f_y(0,0)=\lim_{k\to0}\frac{f(0,k)-f(0,0)}{k}=0.
\]
因此两个偏导数都存在。

最后判可微性。若可微，则由于两个偏导数均为 $0$，应有
\[
  \frac{f(x,y)-f(0,0)}{\sqrt{x^2+y^2}}\to0.
\]
但沿 $y=x$ 有
\[
  \frac{f(x,x)-f(0,0)}{\sqrt{x^2+x^2}}
  =\frac{x/2}{\sqrt2|x|},
\]
该式在 $x\to0^+$ 与 $x\to0^-$ 时极限不同，特别不趋于 $0$。故函数在原点不可微。选 A。
\end{solution}
\begin{exercise}
\par
设平面区域 \(\displaystyle D=\{(x,y)\mid -a\le x\le a,\ x\le y\le a\},\qquad
    D_1=\{(x,y)\mid 0\le x\le a,\ x\le y\le a\}\)，
  则
  \(\displaystyle \iint_D (xy+\cos x\sin y)\,\mathrm{d}x\,\mathrm{d}y\)
  等于（\quad）。

  A. $\displaystyle 2\iint_{D_1}\cos x\sin y\,\mathrm{d}x\,\mathrm{d}y$；\quad
  B. $\displaystyle 2\iint_{D_1}xy\,\mathrm{d}x\,\mathrm{d}y$；\quad
  C. $\displaystyle 4\iint_{D_1}(xy+\cos x\sin y)\,\mathrm{d}x\,\mathrm{d}y$；\quad
  D. $0$。
\end{exercise}

\begin{solution}
\par

把积分拆成两部分：
\[
  \iint_D(xy+\cos x\sin y)\,\mathrm{d}x\,\mathrm{d}y
  =\iint_Dxy\,\mathrm{d}x\,\mathrm{d}y+\iint_D\cos x\sin y\,\mathrm{d}x\,\mathrm{d}y .
\]
对第一部分，按题给区域写成累次积分：
\[
  \iint_Dxy\,\mathrm{d}x\,\mathrm{d}y
  =\int_{-a}^{a}x\left(\int_x^a y\,\mathrm{d}y\right)\mathrm{d}x
  =\frac12\int_{-a}^{a}x(a^2-x^2)\,\mathrm{d}x=0,
\]
因为被积函数 $x(a^2-x^2)$ 是奇函数。

对第二部分，先对 $y$ 积分：
\[
  \iint_D\cos x\sin y\,\mathrm{d}x\,\mathrm{d}y
  =\int_{-a}^{a}\cos x\left(\int_x^a\sin y\,\mathrm{d}y\right)\mathrm{d}x
  =\int_{-a}^{a}\cos x(\cos x-\cos a)\,\mathrm{d}x.
\]
右端关于 $x$ 为偶函数，所以等于
\[
  2\int_0^a\cos x(\cos x-\cos a)\,\mathrm{d}x
  =2\iint_{D_1}\cos x\sin y\,\mathrm{d}x\,\mathrm{d}y .
\]
因此原积分等于选项 A。
\end{solution}
\begin{exercise}
\par
若级数 $\displaystyle \sum_{n=1}^{\infty}a_n^2$，$\displaystyle \sum_{n=1}^{\infty}b_n^2$ 都收敛，则下列结论不成立的是（\quad）。

  A. $\displaystyle \sum_{n=1}^{\infty}(a_n+b_n)^2$ 收敛；\quad
  B. $\displaystyle \sum_{n=1}^{\infty}\dfrac{|a_n|}{n}$ 收敛；\quad
  C. $\displaystyle \sum_{n=1}^{\infty}a_nb_n$ 收敛；\quad
  D. $\displaystyle \sum_{n=1}^{\infty}a_nb_n$ 发散。
\end{exercise}

\begin{solution}
\par

已知
\[
  \sum a_n^2,\qquad \sum b_n^2
\]
都收敛，说明 $\{a_n\},\{b_n\}$ 都属于平方可和序列。逐项判断：

A 项中
\[
  (a_n+b_n)^2\le 2a_n^2+2b_n^2,
\]
右端级数收敛，所以 $\sum(a_n+b_n)^2$ 收敛，A 正确。

B 项用 Cauchy 不等式：
\[
  \sum_{n=1}^{\infty}\frac{|a_n|}{n}
  \le
  \left(\sum_{n=1}^{\infty}a_n^2\right)^{1/2}
  \left(\sum_{n=1}^{\infty}\frac1{n^2}\right)^{1/2}<\infty,
\]
所以 B 正确。

C 项同样用 Cauchy 不等式：
\[
  \sum_{n=1}^{\infty}|a_nb_n|
  \le
  \left(\sum_{n=1}^{\infty}a_n^2\right)^{1/2}
  \left(\sum_{n=1}^{\infty}b_n^2\right)^{1/2}<\infty.
\]
因此 $\sum a_nb_n$ 绝对收敛，C 正确。D 声称同一个级数发散，与 C 的结论矛盾，故不成立的是 D。
\end{solution}
\begin{exercise}
\par
微分方程 \(\displaystyle y'+\frac{1}{x}y=\frac{\sin x}{x}\)
  的通解为（\quad）。

  A. $\displaystyle y=\dfrac{1}{x}(-\cos x+C)$；\quad
  B. $y=x(-\cos x+C)$；\quad
  C. $\displaystyle y=\dfrac{1}{x}(-\sin x+C)$；\quad
  D. $\displaystyle y=\dfrac{1}{x}(\cos x+C)$。
\end{exercise}

\begin{solution}
\par

这是标准一阶线性方程
\[
  y'+P(x)y=Q(x),\qquad P(x)=\frac1x,\quad Q(x)=\frac{\sin x}{x}.
\]
在不跨过 $x=0$ 的区间内，积分因子可取
\[
  \mu(x)=\mathrm e^{\int \frac1x\,\mathrm{d}x}=x.
\]
方程两边同乘以 $x$，得
\[
  xy'+y=\sin x.
\]
左端正好是乘积导数：
\[
  (xy)'=\sin x.
\]
积分得
\[
  xy=-\cos x+C,
\]
所以
\[
  y=\frac{-\cos x+C}{x}.
\]
与选项 A 相同。
\end{solution}

\par\medskip
\noindent\textbf{二、填空题（每小题 3 分，共 15 分）}\par\nobreak\smallskip
\begin{exercise}
\par
过点 $(0,2,4)$ 且同时平行于平面 $x+2z=1$ 和 $y-3z=2$ 的直线方程是 \underline{\hspace{4cm}}。
\end{exercise}

\begin{solution}
\par

所求直线同时平行于两个平面，因此它的方向向量必须同时垂直于这两个平面的法向量。两个平面
\[
  x+2z=1,\qquad y-3z=2
\]
的法向量分别为
\[
  \boldsymbol n_1=(1,0,2),\qquad \boldsymbol n_2=(0,1,-3).
\]
故直线方向向量可取
\[
  \boldsymbol s=\boldsymbol n_1\times\boldsymbol n_2=(-2,3,1).
\]
又直线过点 $(0,2,4)$，所以点向式为
\[
  (x,y,z)=(0,2,4)+t(-2,3,1).
\]
写成对称式即
\[
  \frac{x}{-2}=\frac{y-2}{3}=z-4.
\]
\end{solution}
\begin{exercise}
\par
将积分 \(\displaystyle \int_1^e \mathrm{d}x\int_0^{\ln x} f(x,y)\,\mathrm{d}y\)
  改变积分顺序为 \underline{\hspace{4cm}}。
  \textit{（校勘：原卷内层微分疑似误印为 $\mathrm{d}x$；按二重积分换序题意应为 $\mathrm{d}y$。）}
\end{exercise}

\begin{solution}
\par

先把原累次积分翻译成平面区域：
\[
  1\le x\le e,\qquad 0\le y\le \ln x.
\]
边界 $y=\ln x$ 等价于 $x=e^y$。当 $x$ 从 $1$ 到 $e$ 时，$y$ 的范围从 $0$ 到 $1$，所以换成先取 $y$ 时有
\[
  0\le y\le1.
\]
固定 $y$ 后，区域内的 $x$ 从曲线 $x=e^y$ 到右边界 $x=e$，即
\[
  e^y\le x\le e.
\]
因此换序后为
\[
  \int_1^e \mathrm{d}x\int_0^{\ln x}f(x,y)\,\mathrm{d}y
  =\int_0^1\mathrm{d}y\int_{e^y}^{e}f(x,y)\,\mathrm{d}x.
\]
\end{solution}
\begin{exercise}
\par
级数 \(\displaystyle \sum_{n=1}^{\infty}\frac{3^n+(-2)^n}{n}(x+1)^n\)
  的收敛域为 \underline{\hspace{4cm}}。
\end{exercise}

\begin{solution}
\par

这是以 $x+1$ 为变量的幂级数。设
\[
  c_n=\frac{3^n+(-2)^n}{n}.
\]
因为 $3^n$ 的增长速度快于 $2^n$，所以
\[
  \limsup_{n\to\infty}|c_n|^{1/n}=3,
\]
从而收敛半径为
\[
  R=\frac13.
\]
中心是 $x=-1$，所以先得开区间
\[
  -1-\frac13<x<-1+\frac13,
  \qquad\text{即}\qquad
  -\frac43<x<-\frac23.
\]
再判端点。

当 $x+1=\frac13$，级数化为
\[
  \sum_{n=1}^{\infty}\frac{1+(-2/3)^n}{n},
\]
其中 $\sum 1/n$ 发散，而 $\sum (-2/3)^n/n$ 收敛，故此端点发散。

当 $x+1=-\frac13$，级数化为
\[
  \sum_{n=1}^{\infty}\frac{(-1)^n+(2/3)^n}{n},
\]
其中交错调和级数收敛，$\sum (2/3)^n/n$ 也收敛，故此端点收敛。
因此收敛域为
\[
  \left[-\frac43,-\frac23\right).
\]
\end{solution}
\begin{exercise}
\par
设函数 $u=xyz$，则 $u$ 沿方向 $\ell=\{2,-1,3\}$ 的方向导数为 \underline{\hspace{4cm}}。
\end{exercise}

\begin{solution}
\par

方向导数等于梯度与单位方向向量的点积。先求梯度：
\[
  \nabla u=(u_x,u_y,u_z)=(yz,xz,xy).
\]
方向 $\ell=\{2,-1,3\}$ 的长度为
\[
  |\ell|=\sqrt{2^2+(-1)^2+3^2}=\sqrt{14},
\]
所以对应单位方向向量为
\[
  \boldsymbol e_\ell=\frac{(2,-1,3)}{\sqrt{14}}.
\]
于是
\[
  D_\ell u=\nabla u\cdot \boldsymbol e_\ell
  =\frac{2yz-xz+3xy}{\sqrt{14}}.
\]
\end{solution}
\begin{exercise}
\par
设 $L$ 为取逆时针方向的圆周 $x^2+y^2=9$，求曲线积分 \(\displaystyle \oint_L (2xy-2y)\,\mathrm{d}x+(x^2-4x)\,\mathrm{d}y\)
  的值为 \underline{\hspace{4cm}}。
\end{exercise}

\begin{solution}
\par

记
\[
  P=2xy-2y,\qquad Q=x^2-4x.
\]
曲线 $L:x^2+y^2=9$ 取逆时针方向，正好是 Green 公式的正向边界。因此
\[
  \oint_LP\,\mathrm{d}x+Q\,\mathrm{d}y
  =\iint_D(Q_x-P_y)\,\mathrm{d}x\,\mathrm{d}y,
\]
其中 $D$ 是半径为 $3$ 的圆盘。计算偏导：
\[
  Q_x=2x-4,\qquad P_y=2x-2,
\]
所以
\[
  Q_x-P_y=(2x-4)-(2x-2)=-2.
\]
于是
\[
  \oint_L(2xy-2y)\,\mathrm{d}x+(x^2-4x)\,\mathrm{d}y
  =\iint_D(-2)\,\mathrm{d}x\,\mathrm{d}y
  =-2\cdot 9\pi=-18\pi.
\]
\end{solution}

\par\medskip
\noindent\textbf{三、计算题（每小题 7 分，共 35 分）}\par\nobreak\smallskip
\begin{exercise}
\par
将函数 \(\displaystyle f(x)=\ln\frac{1}{2+2x+x^2}\)
  按 $(x+1)$ 的正整数次幂展开成幂级数，即在 $x_0=-1$ 处展开成泰勒级数。
\end{exercise}

\begin{solution}
\par

因为要按 $(x+1)$ 的正整数次幂展开，令
\[
  t=x+1.
\]
原分母可化为
\[
  2+2x+x^2=(x+1)^2+1=t^2+1.
\]
因此
\[
  f(x)=\ln\frac1{1+t^2}=-\ln(1+t^2).
\]
利用基本展开式
\[
  \ln(1+u)=\sum_{n=1}^{\infty}\frac{(-1)^{n-1}}{n}u^n,\qquad |u|<1,
\]
取 $u=t^2$，得
\[
  -\ln(1+t^2)=\sum_{n=1}^{\infty}\frac{(-1)^n}{n}t^{2n}.
\]
代回 $t=x+1$，得到
\[
  f(x)=\sum_{n=1}^{\infty}\frac{(-1)^n}{n}(x+1)^{2n},
  \qquad |x+1|<1.
\]
\end{solution}
\begin{exercise}
\par
利用高斯公式计算曲面积分 \(\displaystyle \iint_{\Sigma}x^2\,\mathrm{d}y\,\mathrm{d}z+2y^2\,\mathrm{d}z\,\mathrm{d}x+(3z^2-4x^2y^2)\,\mathrm{d}x\,\mathrm{d}y\)，
  其中 $\Sigma$ 为 $z=\sqrt{x^2+y^2}$ 与 $z=2$ 围成的立体表面，取外侧。
\end{exercise}

\begin{solution}
\par

将第二类曲面积分看成向量场通量。设
\[
  \mathbf F=(P,Q,R)=(x^2,2y^2,3z^2-4x^2y^2).
\]
曲面 $\Sigma$ 是由圆锥面 $z=\sqrt{x^2+y^2}$ 与平面 $z=2$ 围成的闭曲面，并取外侧方向，正可使用 Gauss 公式：
\[
  \iint_{\Sigma}P\,\mathrm{d}y\,\mathrm{d}z+Q\,\mathrm{d}z\,\mathrm{d}x+R\,\mathrm{d}x\,\mathrm{d}y
  =\iiint_{\Omega}\nabla\cdot\mathbf F\,\mathrm{d}V.
\]
计算散度：
\[
  \nabla\cdot\mathbf F
  =\frac{\partial x^2}{\partial x}
   +\frac{\partial(2y^2)}{\partial y}
   +\frac{\partial(3z^2-4x^2y^2)}{\partial z}
  =2x+4y+6z.
\]
区域 $\Omega$ 关于 $x$、$y$ 都对称，所以 $2x$ 与 $4y$ 的体积分为 $0$，只剩
\[
  \iiint_{\Omega}6z\,\mathrm{d}V.
\]
用柱坐标 $x=r\cos\theta,\ y=r\sin\theta$。圆锥面为 $z=r$，平面为 $z=2$，故可取
\[
  0\le\theta\le2\pi,\qquad 0\le z\le2,\qquad 0\le r\le z.
\]
于是
\[
  \iint_{\Sigma}\mathbf F\cdot d\mathbf S
  =\int_0^{2\pi}\int_0^2\int_0^z 6z\,r\,\mathrm{d}r\,\mathrm{d}z\,\mathrm{d}\theta.
\]
先对 $r$ 积分得
\[
  \int_0^z6zr\,\mathrm{d}r=3z^3,
\]
再积分：
\[
  \int_0^{2\pi}\mathrm{d}\theta\int_0^2 3z^3\,\mathrm{d}z
  =2\pi\cdot 12=24\pi.
\]
\end{solution}
\begin{exercise}
\par
设二阶可微函数 $\varphi(x)$ 满足方程 \(\displaystyle \varphi(x)=e-\int_0^x (x-u)\varphi(u)\,\mathrm{d}u\)，
  求 $\varphi(x)$。
\end{exercise}

\begin{solution}
\par

对方程
\[
  \varphi(x)=e-\int_0^x(x-u)\varphi(u)\,\mathrm{d}u
\]
求导，得
\[
  \varphi'(x)=-\int_0^x\varphi(u)\,\mathrm{d}u .
\]
再求导得
\[
  \varphi''(x)=-\varphi(x),\qquad \varphi(0)=e,\quad \varphi'(0)=0 .
\]
即
\[
  \varphi''+\varphi=0.
\]
其通解为
\[
  \varphi(x)=C_1\cos x+C_2\sin x.
\]
由 $\varphi(0)=e$ 得 $C_1=e$；由 $\varphi'(0)=0$ 得 $C_2=0$。因此
\[
  \varphi(x)=e\cos x.
\]
\end{solution}
\begin{exercise}
\par
设 \(\displaystyle z=xf\left(2x,\frac{y^2}{x}\right)\)，
  其中 $f$ 具有二阶连续偏导数，求 $\displaystyle \dfrac{\partial^2 z}{\partial x\partial y}$。
\end{exercise}

\begin{solution}
\par

设
\[
  u=2x,\qquad v=\frac{y^2}{x},
\]
则
\[
  z=x f(u,v).
\]
先对 $y$ 求偏导。此时 $x$ 看作常数，且
\[
  u_y=0,\qquad v_y=\frac{2y}{x}.
\]
因此
\[
  z_y=x\left(f_1(u,v)u_y+f_2(u,v)v_y\right)
  =x f_2(u,v)\frac{2y}{x}=2y f_2(u,v).
\]
再对 $x$ 求偏导：
\[
  u_x=2,\qquad v_x=-\frac{y^2}{x^2}.
\]
所以
\[
  z_{yx}
  =2y\left(f_{21}(u,v)u_x+f_{22}(u,v)v_x\right)
  =2y\left(2f_{21}(u,v)-\frac{y^2}{x^2}f_{22}(u,v)\right).
\]
由于 $f$ 具有二阶连续偏导数，混合偏导可交换，即 $f_{21}=f_{12}$。故
\[
  \frac{\partial^2z}{\partial x\partial y}
  =4y f_{12}\left(2x,\frac{y^2}{x}\right)
  -\frac{2y^3}{x^2}f_{22}\left(2x,\frac{y^2}{x}\right).
\]
\end{solution}
\begin{exercise}
\par
设三重积分 \(\displaystyle I=\iiint_V (x^2+y^2+z^2)\,\mathrm{d}x\,\mathrm{d}y\,\mathrm{d}z\)，
  其中积分区域 $V$ 是由球面 $x^2+y^2+z^2=R^2$（$z\ge 0$）及锥面 $x^2+y^2=z^2$ 所围成的区域。
\end{exercise}

\begin{solution}
\par

该区域由上半球面与圆锥面围成，适合用球坐标
\[
  x=\rho\sin\varphi\cos\theta,\quad
  y=\rho\sin\varphi\sin\theta,\quad
  z=\rho\cos\varphi.
\]
球面给出
\[
  0\le \rho\le R.
\]
锥面 $x^2+y^2=z^2$ 在球坐标中为
\[
  \rho^2\sin^2\varphi=\rho^2\cos^2\varphi,
\]
且 $z\ge0$，所以
\[
  \varphi=\frac{\pi}{4}.
\]
所围成的是靠近 $z$ 轴的锥内部分，因此
\[
  0\le\varphi\le\frac{\pi}{4},\qquad 0\le\theta\le2\pi.
\]
又
\[
  x^2+y^2+z^2=\rho^2,\qquad \mathrm{d}V=\rho^2\sin\varphi\,\mathrm{d}\rho \mathrm{d}\varphi \mathrm{d}\theta.
\]
于是
\[
  I=\int_0^{2\pi}\int_0^{\pi/4}\int_0^R
  \rho^2\cdot\rho^2\sin\varphi\,\mathrm{d}\rho \mathrm{d}\varphi \mathrm{d}\theta.
\]
逐次积分：
\[
  I=2\pi\cdot \frac{R^5}{5}\cdot
  \left(1-\cos\frac{\pi}{4}\right)
  =\frac{\pi R^5}{5}(2-\sqrt2).
\]
\end{solution}

\par\medskip
\noindent\textbf{四、证明题（第 1 题 7 分，第 2 题 8 分，共 15 分）}\par\nobreak\smallskip
\begin{exercise}
\par
证明：一般项级数 \(\displaystyle \sum_{n=1}^{\infty}\sin\left(\pi\sqrt{n^2+1}\right)\)
  条件收敛。
\end{exercise}

\begin{solution}
\par

把根号中的主项 $n$ 分离出来。令
\[
  \delta_n=\pi\left(\sqrt{n^2+1}-n\right)>0.
\]
由于
\[
  \sqrt{n^2+1}-n=\frac{1}{\sqrt{n^2+1}+n}\sim \frac1{2n},
\]
所以
\[
  \delta_n\sim \frac{\pi}{2n}.
\]
又
\[
  \pi\sqrt{n^2+1}=n\pi+\delta_n,
\]
故
\[
  \sin\left(\pi\sqrt{n^2+1}\right)
  =\sin(n\pi+\delta_n)=(-1)^n\sin\delta_n.
\]
当 $n$ 充分大时，$\delta_n\downarrow0$，从而 $\sin\delta_n$ 也趋于 $0$ 且最终单调递减。因此由交错级数判别法，
\[
  \sum_{n=1}^{\infty}\sin\left(\pi\sqrt{n^2+1}\right)
\]
收敛。

再看绝对收敛性。因为 $\sin\delta_n\sim\delta_n$，所以
\[
  \left|\sin\left(\pi\sqrt{n^2+1}\right)\right|
  =\sin\delta_n\sim \frac{\pi}{2n}.
\]
而 $\sum 1/n$ 发散，故绝对值级数发散。综上，原级数条件收敛。
\end{solution}
\begin{exercise}
\par
证明：曲面 \(\displaystyle F\left(\frac{x-a}{z-c},\frac{y-b}{z-c}\right)=0\)
  上任意一点的切平面都通过一定点。
\end{exercise}

\begin{solution}
\par

记
\[
  u=\frac{x-a}{z-c},\qquad v=\frac{y-b}{z-c}.
\]
曲面方程为 $F(u,v)=0$。若在曲面上取定一组满足 $F(u,v)=0$ 的常数 $u,v$，则由
\[
  \frac{x-a}{z-c}=u,\qquad \frac{y-b}{z-c}=v
\]
可得
\[
  x=a+u(z-c),\qquad y=b+v(z-c).
\]
令 $t=z-c$，曲面上的这一条曲线可写为
\[
  (x,y,z)=(a,b,c)+t(u,v,1).
\]
这说明曲面由一族都经过定点 $(a,b,c)$ 的直线生成。

现在取曲面上一点 $P$。经过 $P$ 的那条生成直线完全位于曲面上，因此这条直线在 $P$ 处的方向属于曲面的切平面。切平面既经过点 $P$，又含有从 $P$ 指向定点 $(a,b,c)$ 的生成直线方向，所以该切平面必经过定点 $(a,b,c)$。
\end{solution}

\par\medskip
\noindent\textbf{五、应用题（每小题 10 分，共 20 分）}\par\nobreak\smallskip
\begin{exercise}
\par
设函数 $y(x)$（$x\ge 0$）二阶可导，且 $y'(x)>0$，$y(0)=1$。
  过曲线 $y=y(x)$ 上任意一点 $P(x,y)$ 作该曲线的切线及 $x$ 轴的垂线，上述两直线与 $x$ 轴所围成的三角形面积记为 $S_1$。
  区间 $(0,x]$ 上以 $y=y(x)$ 为曲边的曲边梯形面积记为 $S_2$。
  若 $2S_1-S_2$ 恒为 $1$，求此曲线 $y=y(x)$ 的方程。
\end{exercise}

\begin{solution}
\par

设曲线上点为 $P(x,y)$。该点切线方程为
\[
  Y-y=y'(X-x).
\]
令 $Y=0$，得切线与 $x$ 轴交点的横坐标
\[
  X=x-\frac{y}{y'}.
\]
由于 $y'(x)>0$ 且 $y(0)=1$，曲线在 $x\ge0$ 上有 $y>0$，因此三角形的水平底边长为 $y/y'$，高为 $y$。于是
\[
  S_1=\frac12\cdot y\cdot\frac{y}{y'}=\frac{y^2}{2y'}.
\]
曲边梯形面积为
\[
  S_2=\int_0^x y(t)\,\mathrm{d}t.
\]
条件 $2S_1-S_2=1$ 化为
\[
  \frac{y^2}{y'}-\int_0^x y(t)\,\mathrm{d}t=1.
\]
对 $x$ 求导，得
\[
  \frac{\mathrm{d}}{\mathrm{d}x}\left(\frac{y^2}{y'}\right)-y=0.
\]
而
\[
  \frac{\mathrm{d}}{\mathrm{d}x}\left(\frac{y^2}{y'}\right)
  =\frac{2y(y')^2-y^2y''}{(y')^2}.
\]
代入并化简：
\[
  \frac{2y(y')^2-y^2y''}{(y')^2}-y=0
  \quad\Longrightarrow\quad
  yy''=(y')^2.
\]
这等价于
\[
  \left(\frac{y'}{y}\right)'=\frac{yy''-(y')^2}{y^2}=0.
\]
所以
\[
  \frac{y'}{y}=k,
\]
其中 $k$ 为常数。解得
\[
  y=Ce^{kx}.
\]
由 $y(0)=1$ 得 $C=1$，故 $y=e^{kx}$。再把 $x=0$ 代回原条件：
\[
  \frac{y(0)^2}{y'(0)}-\int_0^0y(t)\,\mathrm{d}t=1,
\]
即
\[
  \frac1{y'(0)}=1.
\]
所以 $y'(0)=1$。而 $y'=ke^{kx}$，故 $k=1$。最终曲线为
\[
  y=e^x.
\]
\end{solution}
\begin{exercise}
\par
将长为 $a$ 的线段截成两段，一段围成一个正方形，另一段围成一个圆。问这两段的长度分别为多少时，它们所围成的正方形和圆的面积之和最大？
\end{exercise}

\begin{solution}
\par

设用来围正方形的线段长为 $s$，则用来围圆的线段长为 $a-s$，其中
\[
  0\le s\le a.
\]
若正方形周长为 $s$，则边长为 $s/4$，面积为
\[
  A_1=\left(\frac{s}{4}\right)^2=\frac{s^2}{16}.
\]
若圆周长为 $a-s$，则半径为
\[
  r=\frac{a-s}{2\pi},
\]
圆面积为
\[
  A_2=\pi r^2=\frac{(a-s)^2}{4\pi}.
\]
面积和为
\[
  A(s)=\frac{s^2}{16}+\frac{(a-s)^2}{4\pi}.
\]
这是关于 $s$ 的二次函数，并且二次项系数
\[
  \frac1{16}+\frac1{4\pi}>0,
\]
所以图像开口向上。开口向上的二次函数在闭区间上的最大值只能在端点取得。

比较两个端点：
\[
  A(0)=\frac{a^2}{4\pi},\qquad
  A(a)=\frac{a^2}{16}.
\]
因为
\[
  \frac1{4\pi}>\frac1{16}\quad(\text{即 }16>4\pi),
\]
所以 $A(0)>A(a)$。因此允许某一段长度为 $0$ 时，面积和最大应取
\[
  s=0,\qquad a-s=a,
\]
即全部线段用于围圆，正方形那段长度为 $0$，圆那段长度为 $a$。

若题意强制“两段”都必须为正长度，则 $0<s<a$，此时最大值不在开区间内取得；只能让 $s\to0^+$，面积和趋近于上确界 $\dfrac{a^2}{4\pi}$，但没有真正最大值。
\end{solution}

\section{2006--2007 第二学期 B 卷}

\par\medskip
\noindent\textbf{一、单项选择题（每小题 3 分，共 15 分）}\par\nobreak\smallskip
\begin{exercise}
\par
下列说法不正确的是（\quad）。
  A. 如果函数 $f(x,y)$ 在区域 $D$ 中连续，则 $|f(x,y)|$ 也在 $D$ 中连续；\quad
  B. 函数 $f(x,y)$ 沿射线 $\ell$ 的方向导数，等于梯度在 $\ell$ 上的投影；\quad
  C. 如果函数 $f(x,y)$ 在点 $P(x_0,y_0)$ 取极值，则必有 $f_x(x_0,y_0)=0,\ f_y(x_0,y_0)=0$；\quad
  D. 如果函数 $f(x,y)$ 在点 $P(x_0,y_0)$ 存在两个连续的偏导数，则在该点一定可微。
\end{exercise}

\begin{solution}
\par

逐项判断。A 项中，连续函数经过绝对值函数复合后仍连续，因此 $|f(x,y)|$ 在 $D$ 中连续。B 项中，若函数在点处可微，沿单位方向 $\boldsymbol e$ 的方向导数为
\[
  \frac{\partial f}{\partial \boldsymbol e}=\nabla f\cdot \boldsymbol e,
\]
也就是梯度在该方向上的投影。D 项中，两个一阶偏导数在点的邻域内存在并在该点连续，是函数在该点可微的充分条件，因此 D 可作为常用充分判据。

C 项的问题在于：内点极值的必要条件 $f_x(x_0,y_0)=f_y(x_0,y_0)=0$ 还要求这两个偏导数存在。若极值点处偏导数不存在，就不能推出偏导数等于零。例如 $f(x,y)=\sqrt{x^2+y^2}$ 在原点取极小值，但偏导数判别不能直接按 $f_x=f_y=0$ 使用。因此不正确的是 C。
\end{solution}
\begin{exercise}
\par
函数 $z=\ln(1+x^2+y^2)$ 当 $x=1,\ y=2$ 时的全微分为（\quad）。
  A. $\displaystyle \dfrac13\mathrm{d}x+\dfrac13\mathrm{d}y$；\quad
  B. $\displaystyle \dfrac23\mathrm{d}x+\dfrac13\mathrm{d}y$；\quad
  C. $\displaystyle \dfrac23\mathrm{d}x+\dfrac23\mathrm{d}y$；\quad
  D. $\displaystyle \dfrac13\mathrm{d}x+\dfrac23\mathrm{d}y$。
\end{exercise}

\begin{solution}
\par

全微分按
\[
  \mathrm{d}z=z_x\,\mathrm{d}x+z_y\,\mathrm{d}y
\]
计算。这里
\[
  z_x=\frac{2x}{1+x^2+y^2},\qquad
  z_y=\frac{2y}{1+x^2+y^2}.
\]
代入 $x=1,\ y=2$，分母为 $1+1^2+2^2=6$，所以
\[
  z_x(1,2)=\frac{2}{6}=\frac13,\qquad
  z_y(1,2)=\frac{4}{6}=\frac23.
\]
于是
\[
  \mathrm{d}z=\frac13\,\mathrm{d}x+\frac23\,\mathrm{d}y.
\]
故选 D。
\end{solution}
\begin{exercise}
\par
设 $\Omega$ 是曲面 $z=x^2+y^2$ 与平面 $z=1$ 所围成的闭区域，则 $\displaystyle \iiint_{\Omega}f(x,y,z)\,\mathrm{d}V$ 化为三次积分为（\quad）。
\begin{align*}
  \text{A. }&\int_{-1}^{1}\mathrm{d}y\int_{-\sqrt{1-y^2}}^{\sqrt{1-y^2}}\mathrm{d}x\int_1^{x^2+y^2}f(x,y,z)\,\mathrm{d}z;\\
  \text{B. }&4\int_0^1\mathrm{d}x\int_0^{\sqrt{1-x^2}}\mathrm{d}y\int_{x^2+y^2}^{1}f(x,y,z)\,\mathrm{d}z;\\
  \text{C. }&4\int_0^1\mathrm{d}x\int_0^{\sqrt{1-x^2}}\mathrm{d}y\int_0^{x^2+y^2}f(x,y,z)\,\mathrm{d}z;\\
  \text{D. }&\int_{-1}^{1}\mathrm{d}x\int_{-\sqrt{1-x^2}}^{\sqrt{1-x^2}}\mathrm{d}y\int_{x^2+y^2}^{1}f(x,y,z)\,\mathrm{d}z.
\end{align*}
\end{exercise}

\begin{solution}
\par

曲面 $z=x^2+y^2$ 与平面 $z=1$ 相交时满足
\[
  x^2+y^2=1,
\]
因此 $\Omega$ 在 $xOy$ 平面上的投影是单位圆盘
\[
  x^2+y^2\le1.
\]
对投影内任一点 $(x,y)$，下方是抛物面 $z=x^2+y^2$，上方是平面 $z=1$，故
\[
  x^2+y^2\le z\le 1.
\]
若用直角坐标写成三次积分，可取
\[
  -1\le x\le1,\qquad
  -\sqrt{1-x^2}\le y\le\sqrt{1-x^2},\qquad
  x^2+y^2\le z\le1.
\]
所以
\[
  \iiint_{\Omega}f(x,y,z)\,\mathrm{d}V
  =
  \int_{-1}^{1}\mathrm{d}x
  \int_{-\sqrt{1-x^2}}^{\sqrt{1-x^2}}\mathrm{d}y
  \int_{x^2+y^2}^{1}f(x,y,z)\,\mathrm{d}z.
\]
B 项虽然用了第一象限乘 $4$，但 $f(x,y,z)$ 是任意函数，不能默认关于 $x,y$ 对称；A 项上下限反向，C 项下界错误。因此选 D。
\end{solution}
\begin{exercise}
\par
设曲线 $L$ 为沿顺时针方向的圆周 $x^2+y^2=9$，则第二型闭路曲线积分 \(\displaystyle \oint_L (2xy-2y)\,\mathrm{d}x+(x^2-4x)\,\mathrm{d}y\)
  等于（\quad）。
  A. $0$；\quad B. $2\pi$；\quad C. $6\pi$；\quad D. $18\pi$。
\end{exercise}

\begin{solution}
\par

记
\[
  P=2xy-2y,\qquad Q=x^2-4x.
\]
先按正向，即逆时针方向使用 Green 公式：
\[
  \oint_{\partial D}^{+}P\,\mathrm{d}x+Q\,\mathrm{d}y
  =
  \iint_D\left(\frac{\partial Q}{\partial x}-\frac{\partial P}{\partial y}\right)\,\mathrm{d}\sigma.
\]
这里
\[
  \frac{\partial Q}{\partial x}=2x-4,\qquad
  \frac{\partial P}{\partial y}=2x-2,
\]
所以
\[
  Q_x-P_y=(2x-4)-(2x-2)=-2.
\]
圆盘 $D:x^2+y^2\le9$ 的面积为 $9\pi$，因此逆时针方向积分为
\[
  \iint_D(-2)\,\mathrm{d}\sigma=-18\pi.
\]
题设 $L$ 为顺时针方向，与正向相反，积分变号，所以所求值为
\[
  18\pi.
\]
故选 D。
\end{solution}
\begin{exercise}
\par
已知幂级数 $\displaystyle \sum_{n=1}^{\infty}a_nx^n$ 的收敛域为 $[-8,8]$，则 \(\displaystyle \sum_{n=2}^{\infty}\frac{a_nx^n}{n(n-1)}\)
  的收敛半径为（\quad）。
  \textit{（校勘：原卷下标疑似为 $n=1$，但 $n(n-1)$ 使首项无定义；此处按 $n=2$ 处理。）}
  A. $R>8$；\quad B. $R<8$；\quad C. $R=8$；\quad D. 不能确定。
\end{exercise}

\begin{solution}
\par

设原幂级数 $\sum a_nx^n$ 的收敛半径为 $R_0$。题设收敛域为 $[-8,8]$，所以原级数的收敛半径为
\[
  R_0=8.
\]
新级数的系数为
\[
  b_n=\frac{a_n}{n(n-1)}\qquad(n\ge2).
\]
收敛半径只由系数的指数级增长速度决定，而
\[
  \sqrt[n]{n(n-1)}\to1.
\]
因此
\[
  \limsup_{n\to\infty}\sqrt[n]{|b_n|}
  =
  \limsup_{n\to\infty}\sqrt[n]{|a_n|},
\]
收敛半径不变。故新级数的收敛半径仍为
\[
  R=8.
\]
选 C。注意题目只问收敛半径，不需要重新讨论端点。
\end{solution}

\par\medskip
\noindent\textbf{二、填空题（每小题 3 分，共 15 分）}\par\nobreak\smallskip
\begin{exercise}
\par
点 $(2,1,0)$ 到平面 $3x+4y+5z=0$ 的距离为 \underline{\hspace{4cm}}。
\end{exercise}

\begin{solution}
\par

点 $(x_0,y_0,z_0)$ 到平面 $Ax+By+Cz+D=0$ 的距离为
\[
  d=\frac{|Ax_0+By_0+Cz_0+D|}{\sqrt{A^2+B^2+C^2}}.
\]
本题中平面为 $3x+4y+5z=0$，即 $A=3,\ B=4,\ C=5,\ D=0$。代入点 $(2,1,0)$ 得
\[
  d=\frac{|3\cdot2+4\cdot1+5\cdot0|}{\sqrt{3^2+4^2+5^2}}
  =\frac{10}{\sqrt{50}}
  =\sqrt2.
\]
故填 $\sqrt2$。
\end{solution}
\begin{exercise}
\par
微分方程 \(\displaystyle \frac{1}{y}\frac{\mathrm{d}y}{\mathrm{d}x}=2x+\frac{x-x^3}{y}\)
  （$y\ne0$）的通解为 \underline{\hspace{4cm}}。
\end{exercise}

\begin{solution}
\par

由于 $y\ne0$，可将方程两边乘以 $y$：
\[
  y'=2xy+x-x^3.
\]
整理成一阶线性方程标准形式
\[
  y'-2xy=x-x^3.
\]
这里 $P(x)=-2x$，积分因子为
\[
  \mu(x)=\mathrm e^{\int -2x\,\mathrm{d}x}=\mathrm e^{-x^2}.
\]
两边同乘 $\mathrm e^{-x^2}$，得
\[
  \left(y\mathrm e^{-x^2}\right)'=(x-x^3)\mathrm e^{-x^2}.
\]
计算右端积分，令 $t=x^2$，则
\[
  \int (x-x^3)\mathrm e^{-x^2}\,\mathrm{d}x
  =
  \frac12\int(1-t)\mathrm e^{-t}\,\mathrm{d}t
  =
  \frac{x^2}{2}\mathrm e^{-x^2}+C.
\]
因此
\[
  y\mathrm e^{-x^2}=\frac{x^2}{2}\mathrm e^{-x^2}+C,
\]
从而
\[
  y=\frac{x^2}{2}+C\mathrm e^{x^2}.
\]
这就是所求通解。
\end{solution}
\begin{exercise}
\par
设 \(\displaystyle f(x,y)=\ln\left(x+\frac{y}{2x}\right)\)，
  则 $\displaystyle \left.\dfrac{\partial f}{\partial y}\right|_{x=1,\ y=0}=$ \underline{\hspace{4cm}}。
\end{exercise}

\begin{solution}
\par

对 $y$ 求偏导时，把 $x$ 看成常数。由链式法则，
\[
  f_y(x,y)
  =
  \frac{1}{x+\frac{y}{2x}}\cdot\frac{\partial}{\partial y}\left(x+\frac{y}{2x}\right)
  =
  \frac{1}{x+\frac{y}{2x}}\cdot\frac1{2x}.
\]
代入 $(x,y)=(1,0)$，有
\[
  \left.f_y\right|_{(1,0)}
  =
  \frac{1}{1+0}\cdot\frac12
  =
  \frac12.
\]
故填 $\dfrac12$。
\end{solution}
\begin{exercise}
\par
设
  \[
    f(x)=
    \begin{cases}
      x^2, & -2\le x<0,\\
      4-x^2, & 0\le x<2
    \end{cases}
  \]
  则 $f(x)$ 以 $4$ 为周期的傅里叶级数的和函数 $s(x)$ 在 $[-2,2)$ 上的表达式为 \underline{\hspace{4cm}}。
\end{exercise}

\begin{solution}
\par

函数以 $4$ 为周期，傅里叶级数的和函数满足：在连续点收敛到函数值，在跳跃点收敛到左右极限的平均值。

在区间 $(-2,0)$ 上函数连续，故
\[
  s(x)=x^2,\qquad -2<x<0.
\]
在区间 $(0,2)$ 上函数也连续，故
\[
  s(x)=4-x^2,\qquad 0<x<2.
\]
在 $x=0$ 处，
\[
  f(0-)=0,\qquad f(0+)=4,
\]
所以
\[
  s(0)=\frac{0+4}{2}=2.
\]
在端点 $x=-2$ 处要按周期延拓理解。右极限为 $f(-2+)=4$，左极限等于周期延拓后 $x=2-$ 处的极限，即 $0$，故
\[
  s(-2)=\frac{4+0}{2}=2.
\]
因此在 $[-2,2)$ 上
\[
  s(x)=
  \begin{cases}
  x^2,&-2<x<0,\\
  4-x^2,&0<x<2,\\
  2,&x=-2,0.
  \end{cases}
\]
\end{solution}
\begin{exercise}
\par
设曲线积分 \(\displaystyle \int_L [f(x)-e^x]\sin y\,\mathrm{d}x-f(x)\cos y\,\mathrm{d}y\)
  与积分路径无关，其中 $f(x)$ 具有一阶连续导数，且 $f(0)=0$，则 $f(x)=$ \underline{\hspace{4cm}}。
\end{exercise}

\begin{solution}
\par

记
\[
  P(x,y)=[f(x)-\mathrm e^x]\sin y,\qquad
  Q(x,y)=-f(x)\cos y.
\]
曲线积分与路径无关时，应满足
\[
  P_y=Q_x.
\]
分别求偏导：
\[
  P_y=[f(x)-\mathrm e^x]\cos y,\qquad
  Q_x=-f'(x)\cos y.
\]
于是
\[
  [f(x)-\mathrm e^x]\cos y=-f'(x)\cos y.
\]
该等式对 $y$ 恒成立，因此得到关于 $f$ 的微分方程
\[
  f'(x)+f(x)=\mathrm e^x.
\]
这是线性方程。乘以积分因子 $\mathrm e^x$，得
\[
  (\mathrm e^x f(x))'=\mathrm e^{2x}.
\]
积分得
\[
  \mathrm e^x f(x)=\frac12\mathrm e^{2x}+C,
\]
即
\[
  f(x)=\frac12\mathrm e^x+C\mathrm e^{-x}.
\]
由 $f(0)=0$ 得 $\frac12+C=0$，所以 $C=-\frac12$。故
\[
  f(x)=\frac{\mathrm e^x-\mathrm e^{-x}}2.
\]
\end{solution}

\par\medskip
\noindent\textbf{三、计算题（每小题 7 分，共 35 分）}\par\nobreak\smallskip
\begin{exercise}
\par
判断交错级数 \(\displaystyle \sum_{n=2}^{\infty}\frac{(-1)^n}{\sqrt n+(-1)^n}\)
  的敛散性。
  \textit{（校勘：原卷下标疑似为 $n=1$，但首项分母为 $0$；此处按 $n=2$ 处理。）}
\end{exercise}

\begin{solution}
\par

对通项先作有理化分解。由于 $n\ge2$，
\[
  \frac{(-1)^n}{\sqrt n+(-1)^n}
  =
  \frac{(-1)^n(\sqrt n-(-1)^n)}{n-1}
  =
  \frac{(-1)^n\sqrt n}{n-1}-\frac1{n-1}.
\]
因此原级数可看作两部分之和：
\[
  \sum_{n=2}^{\infty}\frac{(-1)^n\sqrt n}{n-1}
  -
  \sum_{n=2}^{\infty}\frac1{n-1}.
\]
第一部分中
\[
  \frac{\sqrt n}{n-1}\sim\frac1{\sqrt n}\to0,
\]
且从某项以后单调下降，所以由交错级数判别法可知
\[
  \sum_{n=2}^{\infty}\frac{(-1)^n\sqrt n}{n-1}
\]
收敛。第二部分是调和级数
\[
  \sum_{n=2}^{\infty}\frac1{n-1}=\sum_{m=1}^{\infty}\frac1m,
\]
发散到 $+\infty$。原级数等于“一个收敛级数减去发散到 $+\infty$ 的级数”，故发散，并且部分和趋于 $-\infty$。
\end{solution}
\begin{exercise}
\par
利用高斯公式计算曲面积分 \(\displaystyle \iint_{\Sigma}x^2\,\mathrm{d}y\,\mathrm{d}z+2y^2\,\mathrm{d}z\,\mathrm{d}x+(3z^2-4x^2y^2)\,\mathrm{d}x\,\mathrm{d}y\)，
  其中 $\Sigma$ 为 $z=\sqrt{x^2+y^2}$ 与 $z=2$ 围成的立体表面，取外侧。
\end{exercise}

\begin{solution}
\par

这是第二类曲面积分。记
\[
  P=x^2,\qquad Q=2y^2,\qquad R=3z^2-4x^2y^2.
\]
曲面 $\Sigma$ 是由圆锥面 $z=\sqrt{x^2+y^2}$ 与平面 $z=2$ 围成的闭曲面，取外侧，正好可以用 Gauss 公式：
\[
  \iint_{\Sigma}P\,\mathrm{d}y\,\mathrm{d}z+Q\,\mathrm{d}z\,\mathrm{d}x+R\,\mathrm{d}x\,\mathrm{d}y
  =
  \iiint_{\Omega}(P_x+Q_y+R_z)\,\mathrm{d}V.
\]
先算散度：
\[
  P_x=2x,\qquad Q_y=4y,\qquad R_z=6z,
\]
故
\[
  P_x+Q_y+R_z=2x+4y+6z.
\]
区域 $\Omega$ 关于 $x,y$ 对称，所以 $\iiint_\Omega 2x\,\mathrm{d}V=0$，$\iiint_\Omega 4y\,\mathrm{d}V=0$，只剩
\[
  \iiint_\Omega 6z\,\mathrm{d}V.
\]
用柱坐标描述圆锥体：
\[
  0\le z\le2,\qquad 0\le r\le z,\qquad 0\le\theta\le2\pi.
\]
于是
\[
  \iiint_\Omega 6z\,\mathrm{d}V
  =
  \int_0^2\int_0^{2\pi}\int_0^z 6z\,r\,\mathrm{d}r\,\mathrm{d}\theta\,\mathrm{d}z
  =
  6\pi\int_0^2z^3\,\mathrm{d}z
  =
  24\pi.
\]
所以所求曲面积分为 $24\pi$。
\end{solution}
\begin{exercise}
\par
已知 $y_1(x)=e^x$ 是齐次线性微分方程 \(\displaystyle (2x-1)y''-(2x+1)y'+2y=0\)
  的一个解，求此方程的通解。
\end{exercise}

\begin{solution}
\par

已知齐次方程的一个非零解 $y_1=e^x$，可用降阶法。令
\[
  y=v(x)e^x.
\]
则
\[
  y'=e^x(v'+v),\qquad
  y''=e^x(v''+2v'+v).
\]
代入原方程并约去 $e^x$，得
\[
  (2x-1)(v''+2v'+v)-(2x+1)(v'+v)+2v=0.
\]
合并同类项，$v$ 项系数为
\[
  (2x-1)-(2x+1)+2=0,
\]
所以方程化为
\[
  (2x-1)v''+(2x-3)v'=0.
\]
设 $w=v'$，得到一阶方程
\[
  (2x-1)w'+(2x-3)w=0.
\]
当 $w\ne0$ 时，
\[
  \frac{w'}{w}=-\frac{2x-3}{2x-1}
  =
  -1+\frac{2}{2x-1}.
\]
积分得
\[
  \ln|w|=-x+\ln|2x-1|+C,
\]
故
\[
  w=C(2x-1)\mathrm e^{-x}.
\]
于是
\[
  v=\int C(2x-1)\mathrm e^{-x}\,\mathrm{d}x
  =
  -C(2x+1)\mathrm e^{-x}+C_1.
\]
乘回 $e^x$，得到
\[
  y=C_1e^x+C_2(2x+1).
\]
因此通解为
\[
  y=C_1e^x+C_2(2x+1).
\]
\end{solution}
\begin{exercise}
\par
设 \(\displaystyle z=xf\left(\frac{y}{x}\right)+yg\left(x,\frac{x}{y}\right)\)，
  其中 $f,g$ 均为二阶可微函数，求 $\displaystyle \dfrac{\partial^2 z}{\partial x\partial y}$。
\end{exercise}

\begin{solution}
\par

记
\[
  u=\frac yx,\qquad v=x,\qquad w=\frac xy.
\]
先对 $x$ 求偏导。第一项为 $xf(u)$，其中 $u_x=-y/x^2=-u/x$，所以
\[
  \frac{\partial}{\partial x}\bigl[xf(u)\bigr]
  =
  f(u)+x f'(u)u_x
  =
  f(u)-u f'(u).
\]
第二项为 $yg(v,w)$，其中 $v_x=1,\ w_x=1/y$，所以
\[
  \frac{\partial}{\partial x}\bigl[yg(v,w)\bigr]
  =
  y\left(g_1(v,w)v_x+g_2(v,w)w_x\right)
  =
  yg_1(v,w)+g_2(v,w).
\]
因此
\[
  z_x=f(u)-u f'(u)+yg_1(v,w)+g_2(v,w).
\]
再对 $y$ 求偏导。对第一部分，
\[
  \frac{\partial}{\partial y}[f(u)-u f'(u)]
  =
  \bigl(f'(u)-f'(u)-u f''(u)\bigr)u_y
  =
  -u f''(u)\cdot\frac1x
  =
  -\frac{y}{x^2}f''\left(\frac yx\right).
\]
对第二部分，注意 $w_y=-x/y^2$，于是
\[
  \frac{\partial}{\partial y}\bigl[yg_1(v,w)\bigr]
  =
  g_1(v,w)+y g_{12}(v,w)w_y
  =
  g_1(v,w)-\frac{x}{y}g_{12}(v,w),
\]
并且
\[
  \frac{\partial}{\partial y}g_2(v,w)
  =
  g_{22}(v,w)w_y
  =
  -\frac{x}{y^2}g_{22}(v,w).
\]
合并得到
\[
  \frac{\partial^2z}{\partial y\partial x}
  =
  -\frac{y}{x^2}f''\left(\frac yx\right)
  +g_1\left(x,\frac xy\right)
  -\frac{x}{y}g_{12}\left(x,\frac xy\right)
  -\frac{x}{y^2}g_{22}\left(x,\frac xy\right).
\]
由于题设 $f,g$ 二阶可微，在相应连续条件下混合偏导可按题目记号写为
\[
  \frac{\partial^2z}{\partial x\partial y}
  =
  -\frac{y}{x^2}f''\left(\frac yx\right)
  +g_1\left(x,\frac xy\right)
  -\frac{x}{y}g_{12}\left(x,\frac xy\right)
  -\frac{x}{y^2}g_{22}\left(x,\frac xy\right).
\]
\end{solution}
\begin{exercise}
\par
计算二重积分 \(\displaystyle \iint_D (x^2+y^2)\,\mathrm{d}\sigma\)，
  其中 $D$ 是由 $y=\sqrt{1-x^2}$ 和 $x$ 轴所围成的平面区域。
\end{exercise}

\begin{solution}
\par

曲线 $y=\sqrt{1-x^2}$ 是单位圆 $x^2+y^2=1$ 的上半圆，和 $x$ 轴围成的区域为上半单位圆盘。取极坐标
\[
  x=r\cos\theta,\qquad y=r\sin\theta.
\]
区域对应
\[
  0\le r\le1,\qquad 0\le\theta\le\pi.
\]
又
\[
  x^2+y^2=r^2,\qquad \mathrm{d}\sigma=r\,\mathrm{d}r\,\mathrm{d}\theta.
\]
因此
\[
  \iint_D(x^2+y^2)\,\mathrm{d}\sigma
  =
  \int_0^\pi\int_0^1 r^2\cdot r\,\mathrm{d}r\,\mathrm{d}\theta
  =
  \int_0^\pi\left[\frac{r^4}{4}\right]_0^1\mathrm{d}\theta
  =
  \frac{\pi}{4}.
\]
\end{solution}

\par\medskip
\noindent\textbf{四、证明题（第 1 题 7 分，第 2 题 8 分，共 15 分）}\par\nobreak\smallskip
\begin{exercise}
\par
证明：级数 $\displaystyle \sum_{n=1}^{\infty}(a_{n+1}-a_n)$ 收敛的充分必要条件是数列 $\{a_n\}$ 收敛。
\end{exercise}

\begin{solution}
\par

记级数的第 $N$ 个部分和为
\[
  S_N=\sum_{n=1}^N(a_{n+1}-a_n).
\]
这是望远镜求和，逐项相消后
\[
  S_N=a_{N+1}-a_1.
\]

先证必要性。若级数 $\sum_{n=1}^{\infty}(a_{n+1}-a_n)$ 收敛，则部分和 $S_N$ 收敛。由
\[
  a_{N+1}=S_N+a_1
\]
可知 $\{a_{N+1}\}$ 收敛，因此原数列 $\{a_n\}$ 也收敛。

再证充分性。若 $\{a_n\}$ 收敛，设 $a_n\to A$，则
\[
  S_N=a_{N+1}-a_1\to A-a_1.
\]
故部分和数列 $\{S_N\}$ 收敛，于是级数
\[
  \sum_{n=1}^{\infty}(a_{n+1}-a_n)
\]
收敛。综上，命题成立。
\end{solution}
\begin{exercise}
\par
证明：曲面 \(\displaystyle F\left(\frac{x-a}{z-c},\frac{y-b}{z-c}\right)=0\)
  上任意一点的切平面都通过一定点。
\end{exercise}

\begin{solution}
\par

设曲面上一点为 $P_0(x_0,y_0,z_0)$，且 $z_0\ne c$。令
\[
  u_0=\frac{x_0-a}{z_0-c},\qquad
  v_0=\frac{y_0-b}{z_0-c}.
\]
因为 $P_0$ 在曲面上，所以
\[
  F(u_0,v_0)=0.
\]

考虑过点 $(a,b,c)$ 与 $P_0$ 的直线
\[
  (x,y,z)=(a,b,c)+\lambda(x_0-a,\ y_0-b,\ z_0-c).
\]
当 $\lambda\ne0$ 时，
\[
  \frac{x-a}{z-c}
  =
  \frac{\lambda(x_0-a)}{\lambda(z_0-c)}
  =
  u_0,\qquad
  \frac{y-b}{z-c}
  =
  \frac{\lambda(y_0-b)}{\lambda(z_0-c)}
  =
  v_0.
\]
因此这条直线上的点都满足
\[
  F\left(\frac{x-a}{z-c},\frac{y-b}{z-c}\right)=F(u_0,v_0)=0,
\]
也就是说曲面包含一条过固定点 $(a,b,c)$ 的直线族。对曲面上任一点 $P_0$，这条通过 $P_0$ 的母线就在曲面上，其方向向量属于该点的切平面。切平面既过 $P_0$，又含有这条母线方向，所以它包含整条母线，从而通过固定点
\[
  (a,b,c).
\]
故曲面上任意一点的切平面都通过一定点 $(a,b,c)$。
\end{solution}

\par\medskip
\noindent\textbf{五、应用题（每小题 10 分，共 20 分）}\par\nobreak\smallskip
\begin{exercise}
\par
设有力场 \(\displaystyle \mathbf F=\left(-\frac1y,\frac{x}{y^2}\right)\)，
  证明力 $\mathbf F$ 在上半平面内对质点所做的功与路径无关，并计算质点从点 $A(1,2)$ 移动到点 $B(2,1)$ 时力 $\mathbf F$ 所做的功。
\end{exercise}

\begin{solution}
\par

把力场写成
\[
  \mathbf F=(P,Q),\qquad
  P(x,y)=-\frac1y,\qquad Q(x,y)=\frac{x}{y^2}.
\]
在上半平面 $y>0$ 内，$P,Q$ 具有连续一阶偏导数，且该区域单连通。计算
\[
  P_y=\frac1{y^2},\qquad Q_x=\frac1{y^2}.
\]
于是 $P_y=Q_x$，所以第二类曲线积分
\[
  \int_L P\,\mathrm{d}x+Q\,\mathrm{d}y
\]
在上半平面内与路径无关。

接着求势函数 $\phi$，使
\[
  \phi_x=P=-\frac1y,\qquad \phi_y=Q=\frac{x}{y^2}.
\]
由 $\phi_x=-1/y$ 对 $x$ 积分，得
\[
  \phi(x,y)=-\frac{x}{y}+C(y).
\]
再对 $y$ 求偏导：
\[
  \phi_y=\frac{x}{y^2}+C'(y).
\]
与 $Q=x/y^2$ 比较，得到 $C'(y)=0$，故可取
\[
  \phi(x,y)=-\frac{x}{y}.
\]
因此从 $A(1,2)$ 到 $B(2,1)$ 的功为势函数增量
\[
  W=\phi(B)-\phi(A)
  =
  \left(-\frac21\right)-\left(-\frac12\right)
  =
  -\frac32.
\]
\end{solution}
\begin{exercise}
\par
设周长为 $2p$ 的矩形绕它的一边旋转构成圆柱形，求矩形的边长各为多少时，圆柱的体积最大。
\end{exercise}

\begin{solution}
\par

设矩形两边长分别为 $x$ 与 $p-x$，其中 $0<x<p$，因为周长为 $2p$，所以两边之和为 $p$。

若绕长度为 $x$ 的边旋转，则圆柱的高为 $x$，半径为 $p-x$，体积为
\[
  V(x)=\pi(p-x)^2x.
\]
求导：
\[
  V'(x)=\pi\bigl[(p-x)^2-2x(p-x)\bigr]
  =
  \pi(p-x)(p-3x).
\]
在 $0<x<p$ 内的驻点为
\[
  x=\frac p3.
\]
端点处体积趋于 $0$，而驻点处
\[
  V\left(\frac p3\right)
  =
  \pi\left(\frac{2p}{3}\right)^2\frac p3
  =
  \frac{4\pi p^3}{27},
\]
故这里取得最大值。

此时矩形两边为
\[
  \frac p3,\qquad \frac{2p}{3}.
\]
为了使体积最大，应绕较短边 $\dfrac p3$ 旋转，使较长边 $\dfrac{2p}{3}$ 成为圆柱半径。
\end{solution}

\section{2007--2008 第二学期 A 卷}

\par\medskip
\noindent\textbf{一、填空题（共 5 小题，每小题 3 分，共 15 分）}\par\nobreak\smallskip
\begin{exercise}
\par
函数 $\displaystyle z=\arctan\dfrac{y}{x}$ 在点 $(1,0)$ 处的梯度为 \underline{\hspace{3cm}}。
\end{exercise}

\begin{solution}
\par
梯度由两个偏导数组成：
\[
  \operatorname{grad}z=(z_x,z_y).
\]
设
\[
  z=\arctan\frac yx.
\]
对 $x$ 求偏导时把 $y$ 看成常数，由链式法则
\[
  z_x=\frac{1}{1+(y/x)^2}\left(-\frac{y}{x^2}\right)
  =
  -\frac{y}{x^2+y^2}.
\]
对 $y$ 求偏导时把 $x$ 看成常数，
\[
  z_y=\frac{1}{1+(y/x)^2}\cdot\frac1x
  =
  \frac{x}{x^2+y^2}.
\]
代入 $(1,0)$ 得
\[
  z_x(1,0)=0,\qquad z_y(1,0)=1.
\]
因此
\[
  \operatorname{grad}z(1,0)=(0,1).
\]
\end{solution}
\begin{exercise}
\par
积分 \(\displaystyle \int_0^2 \mathrm{d}x\int_x^2 e^{-y^2}\,\mathrm{d}y=\underline{\hspace{3cm}}\)。
\end{exercise}

\begin{solution}
\par
先把积分区域写出来：
\[
  0\le x\le2,\qquad x\le y\le2.
\]
该区域等价于
\[
  0\le y\le2,\qquad 0\le x\le y.
\]
因此换序后
\[
  \int_0^2 \mathrm{d}x\int_x^2 e^{-y^2}\,\mathrm{d}y
  =
  \int_0^2 \mathrm{d}y\int_0^y e^{-y^2}\,\mathrm{d}x.
\]
内层对 $x$ 积分时，$e^{-y^2}$ 与 $x$ 无关，所以
\[
  \int_0^2 \mathrm{d}y\int_0^y e^{-y^2}\,\mathrm{d}x
  =
  \int_0^2 ye^{-y^2}\,\mathrm{d}y.
\]
令 $u=y^2$，$\mathrm{d}u=2y\,\mathrm{d}y$，得
\[
  \int_0^2 ye^{-y^2}\,\mathrm{d}y
  =
  \frac12\int_0^4 e^{-u}\,\mathrm{d}u
  =
  \frac{1-e^{-4}}2.
\]
故填 $\dfrac{1-e^{-4}}2$。

\textit{（校勘：答案版给出的数值对应 $e^{y^2}$，但试题 PDF 图像为 $e^{-y^2}$；此处按题面计算。）}
\end{solution}
\begin{exercise}
\par
幂级数 \(\displaystyle \sum_{n=1}^{\infty}\frac{(x-2)^n}{n4^n}\) 的收敛域为 \underline{\hspace{3cm}}。
\end{exercise}

\begin{solution}
\par
令
\[
  t=\frac{x-2}{4}.
\]
原级数化为
\[
  \sum_{n=1}^{\infty}\frac{t^n}{n}.
\]
已知幂级数 $\sum t^n/n$ 的收敛半径为 $1$。因此先有
\[
  |t|<1
  \quad\Longleftrightarrow\quad
  \left|\frac{x-2}{4}\right|<1
  \quad\Longleftrightarrow\quad
  -2<x<6.
\]
端点必须另判。若 $x=-2$，则 $t=-1$，级数为
\[
  \sum_{n=1}^{\infty}\frac{(-1)^n}{n},
\]
由交错级数判别法收敛。若 $x=6$，则 $t=1$，级数为
\[
  \sum_{n=1}^{\infty}\frac1n,
\]
调和级数发散。故收敛域为
\[
  [-2,6).
\]
\end{solution}
\begin{exercise}
\par
设数量场 $u=\ln\sqrt{x^2+y^2}$，则 \(\displaystyle \operatorname{div}(\operatorname{grad}u)=\underline{\hspace{3cm}}\)。
\end{exercise}

\begin{solution}
\par
这里
\[
  u=\ln\sqrt{x^2+y^2}
  =
  \frac12\ln(x^2+y^2),
\]
定义在 $(x,y)\ne(0,0)$ 的区域内。先求梯度：
\[
  u_x=\frac{x}{x^2+y^2},\qquad
  u_y=\frac{y}{x^2+y^2}.
\]
再求散度，即二维 Laplace 算子：
\[
  \operatorname{div}(\operatorname{grad}u)=u_{xx}+u_{yy}.
\]
分别计算
\[
  u_{xx}=\frac{y^2-x^2}{(x^2+y^2)^2},\qquad
  u_{yy}=\frac{x^2-y^2}{(x^2+y^2)^2}.
\]
相加得
\[
  u_{xx}+u_{yy}=0.
\]
因此
\[
  \operatorname{div}(\operatorname{grad}u)=0\qquad ((x,y)\ne(0,0)).
\]
\end{solution}
\begin{exercise}
\par
函数 $z=xy$ 在条件 $x+y=1$ 下的极大值为 \underline{\hspace{3cm}}。
\end{exercise}

\begin{solution}
\par
由约束 $x+y=1$ 可令
\[
  y=1-x.
\]
于是
\[
  z=xy=x(1-x)=-x^2+x
  =
  -\left(x-\frac12\right)^2+\frac14.
\]
这是开口向下的二次函数，因此最大值在
\[
  x=\frac12,\qquad y=\frac12
\]
处取得，极大值为
\[
  z_{\max}=\frac12\cdot\frac12=\frac14.
\]
由于约束 \(x+y=1\) 是一条直线，把二维约束问题化为一元二次函数后已经覆盖了全部可行点；二次函数开口向下，所以该驻点就是全局最大点。
\end{solution}

\par\medskip
\noindent\textbf{二、选择题（共 5 小题，每小题 3 分，共 15 分）}\par\nobreak\smallskip
\begin{exercise}
\par
抛物线 $y^2=4x$ 上与直线 $x-y+4=0$ 相距最近的点是（\quad）。
  A. $(0,0)$；\quad B. $(1,1)$；\quad C. $(1,2)$；\quad D. $(2,2\sqrt2)$。
\end{exercise}

\begin{solution}
\par
抛物线 $y^2=4x$ 上的点可用参数 $y$ 表示为
\[
  \left(\frac{y^2}{4},\,y\right).
\]
点 $(x,y)$ 到直线 $x-y+4=0$ 的距离为
\[
  d=\frac{|x-y+4|}{\sqrt{1^2+(-1)^2}}
  =
  \frac{|x-y+4|}{\sqrt2}.
\]
代入抛物线参数点，得到
\[
  x-y+4=\frac{y^2}{4}-y+4
  =
  \frac{(y-2)^2}{4}+3.
\]
该表达式恒为正，因此距离最小等价于使
\[
  \frac{(y-2)^2}{4}+3
\]
最小。显然 $y=2$ 时最小，此时
\[
  x=\frac{y^2}{4}=1.
\]
故最近点为 $(1,2)$，选 C。
\end{solution}
\begin{exercise}
\par
设区域 $D:0\le y\le x,\ 0\le x\le \pi$，则 \(\displaystyle \iint_D \sqrt{1-\cos^2 x}\,\mathrm{d}x\,\mathrm{d}y\)
  等于（\quad）。
  A. $0$；\quad B. $\displaystyle \dfrac{\pi}{2}$；\quad C. $2$；\quad D. $\pi$。
\end{exercise}

\begin{solution}
\par
区域为
\[
  0\le x\le\pi,\qquad 0\le y\le x.
\]
被积函数只与 $x$ 有关，并且在 $0\le x\le\pi$ 上
\[
  \sqrt{1-\cos^2x}=|\sin x|=\sin x.
\]
于是
\[
  \iint_D \sqrt{1-\cos^2 x}\,\mathrm{d}x\,\mathrm{d}y
  =
  \int_0^\pi\int_0^x \sin x\,\mathrm{d}y\,\mathrm{d}x.
\]
内层对 $y$ 积分得到长度 $x$，所以
\[
  \int_0^\pi\int_0^x \sin x\,\mathrm{d}y\,\mathrm{d}x
  =
  \int_0^\pi x\sin x\,\mathrm{d}x.
\]
分部积分：
\[
  \int_0^\pi x\sin x\,\mathrm{d}x
  =
  \left[-x\cos x\right]_0^\pi+\int_0^\pi\cos x\,\mathrm{d}x
  =
  \pi.
\]
故选 D。
\end{solution}
\begin{exercise}
\par
设 $\alpha$ 是常数，则级数 \(\displaystyle \sum_{n=1}^{\infty}\left[\frac{\sin(n\alpha)}{n^2}-\frac1{\sqrt n}\right]\)
  是（\quad）。
  A. 绝对收敛；\quad B. 发散；\quad C. 条件收敛；\quad D. 敛散性与 $\alpha$ 取值有关。
\end{exercise}

\begin{solution}
\par
把级数拆成两部分：
\[
  \sum_{n=1}^{\infty}\frac{\sin(n\alpha)}{n^2}
  -
  \sum_{n=1}^{\infty}\frac1{\sqrt n}.
\]
第一部分绝对收敛，因为
\[
  \left|\frac{\sin(n\alpha)}{n^2}\right|
  \le
  \frac1{n^2},
\]
而 $\sum 1/n^2$ 收敛。第二部分是 $p$ 级数
\[
  \sum_{n=1}^{\infty}\frac1{\sqrt n}
  =
  \sum_{n=1}^{\infty}\frac1{n^{1/2}},
\]
其中 $p=\frac12<1$，所以发散到 $+\infty$。因此原级数等于“收敛级数减去发散到 $+\infty$ 的级数”，必发散，且与 $\alpha$ 的取值无关。选 B。
\end{solution}
\begin{exercise}
\par
若 $\displaystyle 0\le a_n<\dfrac1n\ (n=1,2,\cdots)$，则下列级数中肯定收敛的是（\quad）。
  A. $\displaystyle \sum_{n=1}^{\infty}a_n$；\quad
  B. $\displaystyle \sum_{n=1}^{\infty}(-1)^na_n$；\quad
  C. $\displaystyle \sum_{n=1}^{\infty}\sqrt{a_n}$；\quad
  D. $\displaystyle \sum_{n=1}^{\infty}(-1)^na_n^2$。
\end{exercise}

\begin{solution}
\par
由条件
\[
  0\le a_n<\frac1n
\]
可得
\[
  0\le a_n^2<\frac1{n^2}.
\]
因此
\[
  \sum_{n=1}^{\infty}a_n^2
\]
由比较判别法收敛。于是
\[
  \sum_{n=1}^{\infty}(-1)^n a_n^2
\]
绝对收敛，必定收敛，所以 D 正确。

其他选项不能保证。A 中可取 $a_n=\frac1{2n}$，则 $\sum a_n$ 发散；B 中虽然是交错形式，但题设没有给出 $a_n$ 单调下降，不能直接保证交错级数收敛；C 中可取 $a_n$ 接近 $1/n$，则 $\sqrt{a_n}$ 接近 $1/\sqrt n$，不能保证收敛。故肯定收敛的是 D。
\end{solution}
\begin{exercise}
\par
设 $y_1$ 与 $y_2$ 是非齐次线性微分方程 \(\displaystyle y''+ay'+by=f(x)\)
  的两个特解，则下列说法正确的是（\quad）。
  A. $y_1+y_2$ 是非齐次方程的解；\quad
  B. $y_1-y_2$ 是非齐次方程的解；\quad
  C. $y_1+y_2$ 是对应齐次方程的解；\quad
  D. $y_1-y_2$ 是对应齐次方程的解。
\end{exercise}

\begin{solution}
\par
记线性微分算子
\[
  L[y]=y''+ay'+by.
\]
因为 $y_1,y_2$ 都是非齐次方程的特解，所以
\[
  L[y_1]=f(x),\qquad L[y_2]=f(x).
\]
利用线性性，
\[
  L[y_1-y_2]=L[y_1]-L[y_2]=f(x)-f(x)=0.
\]
这说明 $y_1-y_2$ 满足对应齐次方程
\[
  y''+ay'+by=0.
\]
而
\[
  L[y_1+y_2]=L[y_1]+L[y_2]=2f(x),
\]
一般既不是非齐次方程的右端 $f(x)$，也不是齐次方程的右端 $0$。故正确的是 D。
\end{solution}

\par\medskip
\noindent\textbf{三、计算题（共 2 小题，每小题 7 分，共 14 分）}\par\nobreak\smallskip
\begin{exercise}
\par
设 $u(x,y,z)=x+yz$，其中 $z=z(x,y)$ 由 $x+y+z=0$ 确定，求 \(\displaystyle \frac{\partial^2u}{\partial x\partial y}\)。
\end{exercise}

\begin{solution}
\par
由约束
\[
  x+y+z=0
\]
可直接解出
\[
  z=-x-y.
\]
代入
\[
  u=x+yz
\]
得
\[
  u=x+y(-x-y)=x-xy-y^2.
\]
先对 $y$ 求偏导：
\[
  u_y=-x-2y.
\]
再对 $x$ 求偏导：
\[
  (u_y)_x=-1.
\]
因此
\[
  \frac{\partial^2u}{\partial x\partial y}=-1.
\]
也可以先由隐函数求出 $z_x=z_y=-1$ 再代入，结果相同。
\end{solution}
\begin{exercise}
\par
求曲面 $x^2+2y^2+3z^2=21$ 上平行于平面 $x+4y+6z=0$ 的各个切平面。
\end{exercise}

\begin{solution}
\par
设切点为 $(x_0,y_0,z_0)$。曲面
\[
  x^2+2y^2+3z^2=21
\]
的法向量为梯度
\[
  \nabla F(x_0,y_0,z_0)=(2x_0,4y_0,6z_0).
\]
要求切平面平行于平面
\[
  x+4y+6z=0,
\]
就要求两个平面的法向量平行，即
\[
  (2x_0,4y_0,6z_0)=\lambda(1,4,6).
\]
于是
\[
  x_0=\frac{\lambda}{2},\qquad y_0=\lambda,\qquad z_0=\lambda.
\]
代入椭球面方程：
\[
  \left(\frac{\lambda}{2}\right)^2+2\lambda^2+3\lambda^2=21,
\]
即
\[
  \frac{21}{4}\lambda^2=21,
\]
所以
\[
  \lambda=\pm2.
\]
当 $\lambda=2$ 时，切点为 $(1,2,2)$；当 $\lambda=-2$ 时，切点为 $(-1,-2,-2)$。

椭球面在 $(x_0,y_0,z_0)$ 处的切平面为
\[
  x_0x+2y_0y+3z_0z=21.
\]
代入两个切点，得
\[
  x+4y+6z=21,\qquad x+4y+6z=-21.
\]
故所求切平面为
\[
  x+4y+6z=\pm21.
\]
\end{solution}

\par\medskip
\noindent\textbf{四、计算题（共 2 小题，每小题 7 分，共 14 分）}\par\nobreak\smallskip
\begin{exercise}
\par
计算曲线积分 \(\displaystyle \int_L xy\,\mathrm{d}x+(y-x)\,\mathrm{d}y\)，
  其中 $L$ 为 $y=x^2$ 上由点 $O(0,0)$ 到点 $A(1,1)$ 的一段。
\end{exercise}

\begin{solution}
\par
曲线 $L$ 由 $O(0,0)$ 到 $A(1,1)$，可取参数
\[
  x=t,\qquad y=t^2,\qquad 0\le t\le1.
\]
于是
\[
  \mathrm{d}x=\mathrm{d}t,\qquad \mathrm{d}y=2t\,\mathrm{d}t.
\]
代入第二类曲线积分：
\[
  \int_L xy\,\mathrm{d}x+(y-x)\,\mathrm{d}y
  =
  \int_0^1 \left[t\cdot t^2\,\mathrm{d}t+(t^2-t)\cdot 2t\,\mathrm{d}t\right].
\]
整理被积表达式：
\[
  t^3+2t(t^2-t)=t^3+2t^3-2t^2=3t^3-2t^2.
\]
因此
\[
  \int_0^1(3t^3-2t^2)\,\mathrm{d}t
  =
  \left[\frac{3t^4}{4}-\frac{2t^3}{3}\right]_0^1
  =
  \frac34-\frac23
  =
  \frac1{12}.
\]
\end{solution}
\begin{exercise}
\par
设 $\Sigma$ 是锥面 $z=\sqrt{x^2+y^2}$ 被平面 $z=1$ 所截下的有限部分下侧，计算 \(\displaystyle \iint_{\Sigma}x\,\mathrm{d}y\,\mathrm{d}z+y\,\mathrm{d}z\,\mathrm{d}x+(4-2z)\,\mathrm{d}x\,\mathrm{d}y\)。
\end{exercise}

\begin{solution}
\par
记向量场
\[
  \mathbf F=(P,Q,R)=(x,y,4-2z).
\]
题中第二类曲面积分就是 $\iint_\Sigma \mathbf F\cdot \boldsymbol n\,\mathrm{d}S$。锥面
\[
  z=\sqrt{x^2+y^2}
\]
被 $z=1$ 截下的有限部分，与顶面
\[
  \Sigma_0:\ z=1,\quad x^2+y^2\le1
\]
共同围成一个圆锥体。题设取锥面下侧，正好对应圆锥体的外侧方向。

添上顶面 $\Sigma_0$ 后用 Gauss 公式。散度为
\[
  \operatorname{div}\mathbf F
  =
  \frac{\partial x}{\partial x}
  +\frac{\partial y}{\partial y}
  +\frac{\partial(4-2z)}{\partial z}
  =
  1+1-2=0.
\]
所以闭曲面总通量为
\[
  \iint_{\Sigma\cup\Sigma_0}\mathbf F\cdot\boldsymbol n\,\mathrm{d}S=0.
\]
顶面 $\Sigma_0$ 的外法向量向上，故在顶面上
\[
  \mathbf F\cdot\boldsymbol n=R|_{z=1}=4-2=2.
\]
顶面投影是单位圆盘，面积为 $\pi$，因此
\[
  \iint_{\Sigma_0}\mathbf F\cdot\boldsymbol n\,\mathrm{d}S
  =
  \iint_{x^2+y^2\le1}2\,\mathrm{d}x\,\mathrm{d}y
  =
  2\pi.
\]
闭曲面总通量为 $0$，所以锥面下侧的积分为
\[
  \iint_{\Sigma}\mathbf F\cdot\boldsymbol n\,\mathrm{d}S
  =
  -2\pi.
\]
\end{solution}

\par\medskip
\noindent\textbf{五、证明题（8 分）}\par\nobreak\smallskip
\begin{exercise}
\par
将函数 $f(x)=x$ 在 $[0,\pi]$ 上展开成余弦级数，并证明
\(\displaystyle \sum_{n=1}^{\infty}\frac1{(2n-1)^2}=\frac{\pi^2}{8}\)。
\end{exercise}

\begin{solution}
\par
在 $[0,\pi]$ 上展开成余弦级数，就是把 $f(x)=x$ 作偶延拓。余弦级数形式为
\[
  f(x)\sim \frac{a_0}{2}+\sum_{n=1}^{\infty}a_n\cos nx,
\]
其中
\[
  a_0=\frac2\pi\int_0^\pi x\,\mathrm{d}x=\pi,
\]
所以常数项为 $a_0/2=\pi/2$。对 $n\ge1$，
\[
  a_n=\frac2\pi\int_0^\pi x\cos nx\,\mathrm{d}x.
\]
分部积分得
\[
  \int_0^\pi x\cos nx\,\mathrm{d}x
  =
  \left.\frac{x\sin nx}{n}\right|_0^\pi
  -\frac1n\int_0^\pi \sin nx\,\mathrm{d}x
  =
  \frac{(-1)^n-1}{n^2}.
\]
因此
\[
  a_n=\frac{2\bigl((-1)^n-1\bigr)}{\pi n^2}.
\]
当 $n$ 为偶数时 $a_n=0$，当 $n$ 为奇数时 $a_n=-4/(\pi n^2)$。于是
\[
  x\sim \frac{\pi}{2}
  +\sum_{n=1}^{\infty}\frac{2\big((-1)^n-1\big)}{\pi n^2}\cos nx
  =
  \frac{\pi}{2}
  -\frac{4}{\pi}\sum_{n=1}^{\infty}\frac{\cos(2n-1)x}{(2n-1)^2}.
\]
在 $x=0$ 处，偶延拓函数连续且函数值为 $0$，所以傅里叶级数收敛到 $0$。代入 $x=0$：
\[
  0=\frac{\pi}{2}
  -\frac{4}{\pi}\sum_{n=1}^{\infty}\frac1{(2n-1)^2}.
\]
移项得
\[
  \sum_{n=1}^{\infty}\frac1{(2n-1)^2}
  =
  \frac{\pi^2}{8}.
\]
\end{solution}

\par\medskip
\noindent\textbf{六、计算题（8 分）}\par\nobreak\smallskip
\begin{exercise}
\par
求幂级数 \(\displaystyle \sum_{n=1}^{\infty}n x^{2n}\) 的收敛域及和函数。
\end{exercise}

\begin{solution}
\par
先判收敛域。令 $u=x^2$，原级数变为
\[
  \sum_{n=1}^{\infty}n u^n.
\]
对 $\sum n u^n$ 用比值判别：
\[
  \lim_{n\to\infty}\left|\frac{(n+1)u^{n+1}}{n u^n}\right|
  =
  |u|.
\]
所以当 $|u|<1$ 时收敛，即
\[
  x^2<1\quad\Longleftrightarrow\quad -1<x<1.
\]
当 $x=1$ 或 $x=-1$ 时，通项为 $n$，不趋于 $0$，故发散。因此收敛域为
\[
  (-1,1).
\]

在 $|x|<1$ 内求和函数。由几何级数
\[
  \sum_{n=0}^{\infty}u^n=\frac1{1-u}\qquad(|u|<1)
\]
两边对 $u$ 求导，得
\[
  \sum_{n=1}^{\infty}n u^{n-1}=\frac1{(1-u)^2}.
\]
再乘以 $u$：
\[
  \sum_{n=1}^{\infty}n u^n=\frac{u}{(1-u)^2}.
\]
代回 $u=x^2$，得到
\[
  \sum_{n=1}^{\infty}nx^{2n}
  =
  \frac{u}{(1-u)^2}
  =
  \frac{x^2}{(1-x^2)^2}.
\]
故和函数为
\[
  S(x)=\frac{x^2}{(1-x^2)^2},\qquad -1<x<1.
\]
\end{solution}

\par\medskip
\noindent\textbf{七、证明题（8 分）}\par\nobreak\smallskip
\begin{exercise}
\par
设 $f(x)\in C[0,2]$，证明
\(\displaystyle \int_0^2 \mathrm{d}x\int_x^2 f(x)f(y)\,\mathrm{d}y
  =\frac12\left[\int_0^2 f(x)\,\mathrm{d}x\right]^2\)。
\end{exercise}

\begin{solution}
\par
记
\[
  I=\int_0^2 \mathrm{d}x\int_x^2 f(x)f(y)\,\mathrm{d}y .
\]
对应的积分区域为
\[
  0\le x\le2,\qquad x\le y\le2,
\]
即正方形 $[0,2]\times[0,2]$ 中直线 $y=x$ 上方的三角形。由于 $f(x)f(y)$ 关于 $x,y$ 对称，交换变量可得下方三角形上的积分：
\[
  I=\int_0^2 \mathrm{d}y\int_0^y f(x)f(y)\,\mathrm{d}x.
\]
原积分与交换后的积分相加，刚好覆盖整个正方形 $[0,2]\times[0,2]$，对角线面积为零，不影响积分值。因此
\[
  2I=
  \int_0^2\int_0^2 f(x)f(y)\,\mathrm{d}x\,\mathrm{d}y.
\]
因为 $f(x)$ 只含 $x$，$f(y)$ 只含 $y$，二重积分可分离：
\[
  \int_0^2\int_0^2 f(x)f(y)\,\mathrm{d}x\,\mathrm{d}y
  =
  \left(\int_0^2 f(x)\,\mathrm{d}x\right)
  \left(\int_0^2 f(y)\,\mathrm{d}y\right)
  =
  \left[\int_0^2 f(x)\,\mathrm{d}x\right]^2.
\]
故
\[
  I=\frac12\left[\int_0^2 f(x)\,\mathrm{d}x\right]^2,
\]
命题得证。
\end{solution}

\par\medskip
\noindent\textbf{八、证明题（8 分）}\par\nobreak\smallskip
\begin{exercise}
\par
证明函数 \(\displaystyle F(y)=\int_1^{+\infty}\frac{\mathrm{d}x}{\sqrt[3]{x^4+y^2}}\) 在 $-\infty<y<+\infty$ 上连续。
\end{exercise}

\begin{solution}
\par
对任意固定的 $x\ge1$，函数
\[
  \frac{1}{\sqrt[3]{x^4+y^2}}
\]
关于 $y$ 连续。要证明含参广义积分定义的 $F(y)$ 连续，只需在任意有限区间上证明积分关于 $y$ 一致收敛。

任取 $N>0$，在 $y\in[-N,N]$ 且 $x\ge1$ 时，有
\[
  0\le
  \frac1{\sqrt[3]{x^4+y^2}}
  \le
  \frac1{\sqrt[3]{x^4}}
  =
  \frac1{x^{4/3}}.
\]
而
\[
  \int_1^{+\infty}x^{-4/3}\,\mathrm{d}x
\]
收敛，因为 $4/3>1$。由一致收敛判别法，含参广义积分
\[
  \int_1^{+\infty}\frac{\mathrm{d}x}{\sqrt[3]{x^4+y^2}}
\]
在 $[-N,N]$ 上一致收敛。

被积函数关于 $y$ 连续，且广义积分在 $[-N,N]$ 上一致收敛，所以 $F(y)$ 在 $[-N,N]$ 上连续。由于 $N>0$ 任意，任意实数点都包含在某个 $[-N,N]$ 内，因此
\[
  F(y)
\]
在 $(-\infty,+\infty)$ 上连续。
\end{solution}

\par\medskip
\noindent\textbf{九、应用题（10 分）}\par\nobreak\smallskip
\begin{exercise}
\par
设 \(\displaystyle f(x)=\sin x-\int_0^x (x-t)f(t)\,\mathrm{d}t\)，其中 $f(x)$ 为连续函数，求 $f(x)$。
\end{exercise}

\begin{solution}
\par
这是含变上限积分的方程。设
\[
  f(x)=\sin x-\int_0^x (x-t)f(t)\,\mathrm{d}t.
\]
先令 $x=0$，积分区间退化，得
\[
  f(0)=0.
\]

对等式两边关于 $x$ 求导。由 Leibniz 公式，
\[
  \frac{\mathrm{d}}{\mathrm{d}x}\int_0^x (x-t)f(t)\,\mathrm{d}t
  =
  (x-x)f(x)+\int_0^x f(t)\,\mathrm{d}t
  =
  \int_0^x f(t)\,\mathrm{d}t.
\]
因此
\[
  f'(x)=\cos x-\int_0^x f(t)\,\mathrm{d}t.
\]
再令 $x=0$，得
\[
  f'(0)=1.
\]
继续求导：
\[
  f''(x)=-\sin x-f(x),
\]
即
\[
  f''+f=-\sin x.
\]

齐次方程 $f''+f=0$ 的通解为
\[
  C_1\cos x+C_2\sin x.
\]
由于右端 $-\sin x$ 与齐次解共振，取特解
\[
  f_p(x)=\frac{x}{2}\cos x.
\]
直接代入可得
\[
  f_p''+f_p=-\sin x.
\]
故通解为
\[
  f(x)=C_1\cos x+C_2\sin x+\frac{x}{2}\cos x.
\]
由 $f(0)=0$ 得 $C_1=0$。再求导：
\[
  f'(x)=C_2\cos x+\frac12\cos x-\frac{x}{2}\sin x,
\]
代入 $f'(0)=1$ 得
\[
  C_2+\frac12=1,\qquad C_2=\frac12.
\]
因此
\[
  f(x)=\frac12\sin x+\frac{x}{2}\cos x.
\]
\end{solution}

\section{2009级工科数学分析下 A卷}

\subsubsection*{一、填空题（共 5 小题，每小题 4 分，共 20 分）}
\begin{exercise}
\par
$(y''')^3+(y')^4+x^2+x+1=0$ 是 \underline{\hspace{2cm}} 阶微分方程。
\end{exercise}

\begin{solution}
\par
微分方程的阶数由方程中出现的最高阶导数决定。题中含有
\[
  y''',\qquad y',
\]
其中最高阶导数是 $y'''$，它是三阶导数。虽然 $y'''$ 外面还有三次方，但这不改变“导数阶数”的判断。因此该方程是
\[
  3
\]
阶微分方程。这里要特别区分“导数的阶”和“导数的幂”：\((y''')^3\) 只是三阶导数的三次幂，不会把方程变成九阶；只要最高出现到 \(y'''\)，阶数就固定为三阶。
因此空格中应填 \(3\)。若题目问“次数”，才需要讨论最高阶导数的幂次；本题只问“阶”。
\end{solution}
\begin{exercise}
\par
已知 $(a\times b)\cdot c=2$，则 $[(a+b)\times(b+c)]\cdot(c+a)=$ \underline{\hspace{3cm}}。
\end{exercise}

\begin{solution}
\par
先展开叉乘：
\[
  (a+b)\times(b+c)
  =
  a\times b+a\times c+b\times b+b\times c.
\]
其中 $b\times b=0$，所以
\[
  (a+b)\times(b+c)=a\times b+a\times c+b\times c.
\]
再与 $c+a$ 作数量积：
\[
  [(a+b)\times(b+c)]\cdot(c+a)
  =
  (a\times b)\cdot(c+a)
  +(a\times c)\cdot(c+a)
  +(b\times c)\cdot(c+a).
\]
利用 $a\times b$ 同时垂直于 $a,b$，可知
\[
  (a\times b)\cdot a=0,
\]
同理
\[
  (a\times c)\cdot c=0,\qquad (a\times c)\cdot a=0,\qquad (b\times c)\cdot c=0.
\]
于是只剩
\[
  (a\times b)\cdot c+(b\times c)\cdot a.
\]
数量三重积在循环置换下不变，所以
\[
  (b\times c)\cdot a=(a\times b)\cdot c=2.
\]
故原式为
\[
  2+2=4.
\]
\end{solution}
\begin{exercise}
\par
函数 $\displaystyle z=\arctan\dfrac{y}{x}+\ln\sqrt{x^2+y^2}$ 在点 $(1,0)$ 处的梯度 $\operatorname{grad}z\big|_{(1,0)}=$ \underline{\hspace{3cm}}。
\end{exercise}

\begin{solution}
\par
把函数分成两部分：
\[
  z=\arctan\frac yx+\ln\sqrt{x^2+y^2}.
\]
第一部分的偏导数为
\[
  \left(\arctan\frac yx\right)_x=-\frac{y}{x^2+y^2},\qquad
  \left(\arctan\frac yx\right)_y=\frac{x}{x^2+y^2}.
\]
第二部分可写成
\[
  \ln\sqrt{x^2+y^2}=\frac12\ln(x^2+y^2),
\]
所以
\[
  \left(\ln\sqrt{x^2+y^2}\right)_x=\frac{x}{x^2+y^2},\qquad
  \left(\ln\sqrt{x^2+y^2}\right)_y=\frac{y}{x^2+y^2}.
\]
合并得
\[
  z_x=\frac{x-y}{x^2+y^2},\qquad
  z_y=\frac{x+y}{x^2+y^2}.
\]
代入 $(1,0)$：
\[
  z_x(1,0)=1,\qquad z_y(1,0)=1.
\]
因此
\[
  \operatorname{grad}z\big|_{(1,0)}=(1,1).
\]
\end{solution}
\begin{exercise}
\par
交换积分顺序，则 $\displaystyle\int_1^e \mathrm{d}x\int_0^{\ln x}f(x,y)\,\mathrm{d}y=$ \underline{\hspace{3cm}}。
\end{exercise}

\begin{solution}
\par
先把原积分对应的区域写成不等式：
\[
  1\le x\le e,\qquad 0\le y\le\ln x.
\]
边界 $y=\ln x$ 等价于
\[
  x=e^y.
\]
当 $x$ 从 $1$ 到 $e$ 时，$y$ 的最小值为 $0$，最大值为 $\ln e=1$，所以换序后
\[
  0\le y\le1.
\]
固定 $y$ 后，区域中的 $x$ 从曲线 $x=e^y$ 到右边界 $x=e$，即
\[
  e^y\le x\le e.
\]
因此
\[
  \int_1^e \mathrm{d}x\int_0^{\ln x}f(x,y)\,\mathrm{d}y
  =
  \int_0^1\mathrm{d}y\int_{e^y}^{e}f(x,y)\,\mathrm{d}x.
\]
\end{solution}
\begin{exercise}
\par
设 $f(x)$ 是以 $2\pi$ 为周期的函数，它在 $[-\pi,\pi]$ 上的表达式是
  \[
  f(x)=
  \begin{cases}
  x+1,&-\pi\le x<0,\\
  x^2,&0\le x<\pi.
  \end{cases}
  \]
  若它的傅里叶级数的和函数为 $S(x)$，则 $S(0)=$ \underline{\hspace{3cm}}。
\end{exercise}

\begin{solution}
\par
傅里叶级数在连续点收敛到函数值；在跳跃间断点收敛到左右极限的平均值。这里 $x=0$ 处左右表达式不同：
\[
  f(0-)=0+1=1,\qquad f(0+)=0^2=0.
\]
因此
\[
  S(0)=\frac{f(0-)+f(0+)}2
  =
  \frac{1+0}{2}
  =
  \frac12.
\]
这里不能直接把 \(x=0\) 代入某一侧表达式，因为 \(0\) 正是分段点；必须把周期延拓后的左、右极限分别算出，再按傅里叶收敛定理取平均。
最终空格应填 \(\dfrac12\)，而不是原函数在分段点处某一侧的定义值。
所以本题不需要计算任何傅里叶系数，核心只是在分段点正确识别和函数值。
最终空格应填 \(\dfrac12\)，而不是 \(f(0)\) 的某个单侧定义值。
本题的关键不是求傅里叶系数，而是使用收敛定理判断和函数取值。只要题目问的是分段点处的 \(S(0)\)，就先取左右极限，再取平均值；原函数在分段点处写成哪一侧不影响傅里叶级数的收敛值。
这里不能直接把 \(x=0\) 代入某一侧表达式，因为 \(0\) 正是分段点；必须取左右极限平均值。
\end{solution}

\subsubsection*{二、选择题（共 5 小题，每小题 4 分，共 20 分）}
\begin{exercise}
\par
抛物线 $y^2=4x$ 上与直线 $x-y+4=0$ 相距最近的点是（\quad）。
  A. $(0,0)$；\quad B. $(1,1)$；\quad C. $(1,2)$；\quad D. $(2,2\sqrt2)$。
\end{exercise}

\begin{solution}
\par
抛物线 $y^2=4x$ 上的点可设为
\[
  y=t,\qquad x=\frac{t^2}{4}.
\]
点到直线 $x-y+4=0$ 的距离为
\[
  d=\frac{|x-y+4|}{\sqrt2}.
\]
代入参数点：
\[
  x-y+4=\frac{t^2}{4}-t+4
  =
  \frac{(t-2)^2}{4}+3.
\]
该式恒为正，所以距离最小等价于使
\[
  \frac{(t-2)^2}{4}+3
\]
最小。显然 $t=2$ 时最小，此时
\[
  y=2,\qquad x=\frac{2^2}{4}=1.
\]
故最近点为 $(1,2)$，选 C。
\end{solution}
\begin{exercise}
\par
直线 $\displaystyle L:\dfrac{x-2}{2}=\dfrac{y+1}{-2}=\dfrac{z-3}{1}$ 与平面 $\pi:x+2y-2z=6$ 的关系是（\quad）。
  A. 平行；\quad B. 垂直；\quad C. 相交但不垂直；\quad D. 重合。
\end{exercise}

\begin{solution}
\par
由直线的对称式可读出方向向量
\[
  \boldsymbol v=(2,-2,1).
\]
平面
\[
  x+2y-2z=6
\]
的法向量为
\[
  \boldsymbol n=(1,2,-2).
\]
先判断是否平行于平面。若直线平行于平面，则应有 $\boldsymbol v\cdot\boldsymbol n=0$。但
\[
  \boldsymbol v\cdot\boldsymbol n
  =
  2\cdot1+(-2)\cdot2+1\cdot(-2)
  =
  -4\ne0.
\]
所以直线不平行于平面，必与平面相交。再判断是否垂直于平面：若直线垂直于平面，则 $\boldsymbol v$ 应与 $\boldsymbol n$ 平行。显然
\[
  (2,-2,1)
\]
不是 $(1,2,-2)$ 的数倍，因此不垂直。故关系为相交但不垂直，选 C。
\end{solution}
\begin{exercise}
\par
设 $f(x,y)$ 在点 $(x_0,y_0)$ 处的两个偏导数都存在，则（\quad）。
  A. $f(x,y)$ 在点 $(x_0,y_0)$ 必连续；\quad
  B. $f(x,y)$ 在点 $(x_0,y_0)$ 必可微；\quad
  C. $\displaystyle\lim_{x\to x_0}f(x,y_0)$ 与 $\displaystyle\lim_{y\to y_0}f(x_0,y)$ 都存在；\quad
  D. $\displaystyle\lim_{\substack{x\to x_0\\y\to y_0}}f(x,y)$ 存在。
\end{exercise}

\begin{solution}
\par
两个偏导数在 $(x_0,y_0)$ 处存在，意思是
\[
  f_x(x_0,y_0)=\lim_{h\to0}\frac{f(x_0+h,y_0)-f(x_0,y_0)}{h}
\]
和
\[
  f_y(x_0,y_0)=\lim_{k\to0}\frac{f(x_0,y_0+k)-f(x_0,y_0)}{k}
\]
都存在。这说明一元函数 $x\mapsto f(x,y_0)$ 在 $x_0$ 处可导，从而连续；一元函数 $y\mapsto f(x_0,y)$ 在 $y_0$ 处可导，从而连续。因此
\[
  \lim_{x\to x_0}f(x,y_0)=f(x_0,y_0),\qquad
  \lim_{y\to y_0}f(x_0,y)=f(x_0,y_0),
\]
C 正确。

但偏导数存在并不能推出二元函数连续，更不能推出可微，也不能推出二重极限存在。因此 A、B、D 都不能保证。选 C。
\end{solution}
\begin{exercise}
\par
设区域 $D:0\le y\le x,\ 0\le x\le\pi$，则二重积分 $\displaystyle\iint_D\sqrt{1-\sin^2x}\,\mathrm{d}x\,\mathrm{d}y$（\quad）。
  A. $0$；\quad B. $\displaystyle \dfrac{\pi}{2}$；\quad C. $2$；\quad D. $\pi$。
\end{exercise}

\begin{solution}
\par
在区域
\[
  D:\quad 0\le x\le\pi,\qquad 0\le y\le x
\]
上，被积函数只与 $x$ 有关。由于
\[
  \sqrt{1-\sin^2x}=|\cos x|,
\]
所以
\[
  \iint_D\sqrt{1-\sin^2x}\,\mathrm{d}x\,\mathrm{d}y
  =
  \int_0^\pi\int_0^x|\cos x|\,\mathrm{d}y\,\mathrm{d}x
  =
  \int_0^\pi x|\cos x|\,\mathrm{d}x.
\]
在 $[0,\pi/2]$ 上 $|\cos x|=\cos x$，在 $[\pi/2,\pi]$ 上 $|\cos x|=-\cos x$，故
\[
  \int_0^\pi x|\cos x|\,\mathrm{d}x
  =
  \int_0^{\pi/2}x\cos x\,\mathrm{d}x
  -
  \int_{\pi/2}^{\pi}x\cos x\,\mathrm{d}x.
\]
分部积分计算可得
\[
  \int_0^{\pi/2}x\cos x\,\mathrm{d}x=\frac\pi2-1,
\]
以及
\[
  -\int_{\pi/2}^{\pi}x\cos x\,\mathrm{d}x=1+\frac\pi2.
\]
相加得
\[
  \pi.
\]
故选 D。
\end{solution}
\begin{exercise}
\par
设 $\alpha$ 为常数，则级数 $\displaystyle\sum_{n=1}^{\infty}\left[\dfrac{\sin(n\alpha)}{n^2}-\dfrac1n\right]$（\quad）。
  A. 绝对收敛；\quad B. 发散；\quad C. 条件收敛；\quad D. 收敛性与 $\alpha$ 的取值有关。
\end{exercise}

\begin{solution}
\par
把原级数拆成
\[
  \sum_{n=1}^{\infty}\frac{\sin(n\alpha)}{n^2}
  -
  \sum_{n=1}^{\infty}\frac1n.
\]
第一部分绝对收敛，因为
\[
  \left|\frac{\sin(n\alpha)}{n^2}\right|\le\frac1{n^2},
\]
而 $\sum1/n^2$ 收敛。第二部分是调和级数，发散到 $+\infty$。因此原级数等于一个收敛级数减去发散级数，必然发散。这个结论与 $\alpha$ 的取值无关，故选 B。
\end{solution}

\subsubsection*{三、计算题（本大题共 4 小题，每小题 8 分，共 32 分）}
\begin{exercise}
\par
求微分方程 $\displaystyle xy''=y'\ln\dfrac{y'}{x}$ 的通解。
\end{exercise}

\begin{solution}
\par
方程中不显含 $y$，令
\[
  p=y',
\]
则 $y''=p'$，原方程化为
\[
  xp'=p\ln\frac px.
\]
为处理对数项，令
\[
  q=\ln\frac px.
\]
于是
\[
  p=xe^q.
\]
对 $x$ 求导：
\[
  p'=e^q+xe^q q'=e^q(1+xq').
\]
代入 $xp'=p q$，得
\[
  x e^q(1+xq')=x e^q q.
\]
约去 $xe^q$，得到
\[
  1+xq'=q,
\]
即
\[
  q'=\frac{q-1}{x}.
\]
分离变量：
\[
  \frac{\mathrm{d}q}{q-1}=\frac{\mathrm{d}x}{x}.
\]
积分得
\[
  \ln|q-1|=\ln|x|+C,
\]
故
\[
  q=1+C_1x.
\]
于是
\[
  y'=p=xe^q=ex\,e^{C_1x}.
\]
再积分：
\[
  y=C_2+e\int xe^{C_1x}\,\mathrm{d}x.
\]
若 $C_1\ne0$，则
\[
  \int xe^{C_1x}\,\mathrm{d}x=\frac{e^{C_1x}(C_1x-1)}{C_1^2},
\]
从而
\[
  y=C_2+\frac{e\,e^{C_1x}(C_1x-1)}{C_1^2}.
\]
若 $C_1=0$，则 $y'=ex$，所以
\[
  y=C_2+\frac e2x^2.
\]
以上即为通解形式。
\end{solution}
\begin{exercise}
\par
计算三重积分 $\displaystyle I=\iiint_{\Omega}(x^2+y^2+z^2)\,\mathrm{d}V$，其中 $\Omega:x^2+y^2+z^2\le 2z$。
\end{exercise}

\begin{solution}
\par
先整理区域：
\[
  x^2+y^2+z^2\le2z
  \quad\Longleftrightarrow\quad
  x^2+y^2+(z-1)^2\le1.
\]
这是球心为 $(0,0,1)$、半径为 $1$ 的球。令
\[
  w=z-1,
\]
则区域变成单位球
\[
  x^2+y^2+w^2\le1.
\]
被积函数变为
\[
  x^2+y^2+z^2
  =
  x^2+y^2+(w+1)^2
  =
  x^2+y^2+w^2+1+2w.
\]
在关于 $w=0$ 对称的单位球内，奇函数 $2w$ 的积分为 $0$，所以
\[
  I=\iiint_{x^2+y^2+w^2\le1}(x^2+y^2+w^2)\,\mathrm{d}V
  +
  \iiint_{x^2+y^2+w^2\le1}1\,\mathrm{d}V.
\]
用球坐标，第一项为
\[
  \int_0^{2\pi}\int_0^\pi\int_0^1 \rho^2\cdot \rho^2\sin\varphi\,\mathrm{d}\rho\,\mathrm{d}\varphi\,\mathrm{d}\theta
  =
  4\pi\int_0^1\rho^4\,\mathrm{d}\rho
  =
  \frac{4\pi}{5}.
\]
第二项是单位球体积：
\[
  \frac{4\pi}{3}.
\]
因此
\[
  I=\frac{4\pi}{5}+\frac{4\pi}{3}
  =
  \frac{32\pi}{15}.
\]
\end{solution}
\begin{exercise}
\par
计算曲线积分 \(\displaystyle I=\int_L (x+y)^2\,\mathrm{d}x-(x^2+y^2\sin y)\,\mathrm{d}y\)，
  其中 $L$ 是抛物线 $y=x^2$ 上从点 $(-1,1)$ 到点 $(1,1)$ 的那一段。
\end{exercise}

\begin{solution}
\par
曲线为抛物线 $y=x^2$ 从 $(-1,1)$ 到 $(1,1)$ 的一段，方向对应 $x$ 从 $-1$ 增加到 $1$。取参数
\[
  x=t,\qquad y=t^2,\qquad -1\le t\le1.
\]
于是
\[
  \mathrm{d}x=\mathrm{d}t,\qquad \mathrm{d}y=2t\,\mathrm{d}t.
\]
代入积分：
\[
  I=
  \int_{-1}^{1}
  \left[(t+t^2)^2\,\mathrm{d}t-\bigl(t^2+t^4\sin(t^2)\bigr)\,2t\,\mathrm{d}t\right].
\]
整理被积函数：
\[
  (t+t^2)^2-2t\bigl(t^2+t^4\sin(t^2)\bigr)
  =
  t^2+t^4-2t^5\sin(t^2).
\]
其中 $-2t^5\sin(t^2)$ 是奇函数，在对称区间 $[-1,1]$ 上积分为 $0$，故
\[
  I=\int_{-1}^{1}(t^2+t^4)\,\mathrm{d}t.
\]
计算得
\[
  I=2\int_0^1(t^2+t^4)\,\mathrm{d}t
  =
  2\left(\frac13+\frac15\right)
  =
  \frac{16}{15}.
\]
\end{solution}
\begin{exercise}
\par
求幂级数 $\displaystyle\sum_{n=0}^{\infty}(2n+1)x^{2n}$ 的收敛域及和函数。
\end{exercise}

\begin{solution}
\par
令
\[
  r=x^2.
\]
原级数变为
\[
  \sum_{n=0}^{\infty}(2n+1)r^n.
\]
先判收敛。由于 $\sum (2n+1)r^n$ 的收敛半径为 $1$，所以
\[
  |r|<1.
\]
而 $r=x^2\ge0$，故
\[
  x^2<1\quad\Longleftrightarrow\quad -1<x<1.
\]
当 $x=\pm1$ 时，通项为 $2n+1$，不趋于 $0$，所以端点发散。因此收敛域为
\[
  (-1,1).
\]

在 $|r|<1$ 内求和：
\[
  \sum_{n=0}^{\infty}(2n+1)r^n
  =
  2\sum_{n=0}^{\infty}nr^n+\sum_{n=0}^{\infty}r^n.
\]
由
\[
  \sum_{n=0}^{\infty}r^n=\frac1{1-r},\qquad
  \sum_{n=1}^{\infty}nr^n=\frac{r}{(1-r)^2},
\]
得到
\[
  S=2\frac{r}{(1-r)^2}+\frac1{1-r}
  =
  \frac{2r+1-r}{(1-r)^2}
  =
  \frac{1+r}{(1-r)^2}.
\]
代回 $r=x^2$，和函数为
\[
  S(x)=\frac{1+x^2}{(1-x^2)^2},\qquad -1<x<1.
\]
\end{solution}

\subsubsection*{四、应用题（本大题共 2 小题，每小题 7 分，共 14 分）}
\begin{exercise}
\par
造一个容积为 $V$ 的长方形无盖铝盒，怎样设计尺寸才能使它的表面积最小？
\end{exercise}

\begin{solution}
\par
设无盖铝盒底面长、宽分别为 $x,y$，高为 $h$，其中 $x,y,h>0$。体积约束为
\[
  xyh=V.
\]
无盖表面积由底面和四个侧面组成：
\[
  S=xy+2xh+2yh.
\]
由约束得
\[
  h=\frac{V}{xy}.
\]
代入表面积：
\[
  S(x,y)=xy+\frac{2V}{y}+\frac{2V}{x}.
\]
求偏导：
\[
  S_x=y-\frac{2V}{x^2},\qquad
  S_y=x-\frac{2V}{y^2}.
\]
令 $S_x=S_y=0$，得
\[
  x^2y=2V,\qquad xy^2=2V.
\]
两式相除得
\[
  x=y.
\]
设 $x=y=a$，代入 $a^3=2V$，所以
\[
  a=(2V)^{1/3}.
\]
此时
\[
  h=\frac{V}{a^2}=\frac{a}{2}=\left(\frac V4\right)^{1/3}.
\]
当 $x,y,h$ 趋近边界时表面积趋于无穷大，因此该驻点给出最小表面积。故底面应设计为正方形，边长
\[
  (2V)^{1/3},
\]
高为
\[
  \left(\frac V4\right)^{1/3}.
\]
\end{solution}
\begin{exercise}
\par
求密度 $\rho\equiv1$ 的均匀球壳 $x^2+y^2+z^2=a^2\ (z\ge0)$ 对于 $Oz$ 轴的转动惯量。
\end{exercise}

\begin{solution}
\par
对于面密度 $\rho\equiv1$ 的曲面，关于 $Oz$ 轴的转动惯量为
\[
  I_{Oz}=\iint_{\Sigma}(x^2+y^2)\rho\,\mathrm{d}S
  =
  \iint_{\Sigma}(x^2+y^2)\,\mathrm{d}S.
\]
上半球面用球坐标参数化：
\[
  x=a\sin\varphi\cos\theta,\qquad
  y=a\sin\varphi\sin\theta,\qquad
  z=a\cos\varphi,
\]
其中
\[
  0\le\varphi\le\frac\pi2,\qquad 0\le\theta\le2\pi.
\]
在球面上
\[
  x^2+y^2=a^2\sin^2\varphi,
\]
面积元为
\[
  \mathrm{d}S=a^2\sin\varphi\,\mathrm{d}\varphi \mathrm{d}\theta.
\]
因此
\[
  I_{Oz}
  =
  \int_0^{2\pi}\int_0^{\pi/2}
  a^2\sin^2\varphi\cdot a^2\sin\varphi\,\mathrm{d}\varphi \mathrm{d}\theta.
\]
即
\[
  I_{Oz}
  =
  2\pi a^4\int_0^{\pi/2}\sin^3\varphi\,\mathrm{d}\varphi.
\]
计算
\[
  \int_0^{\pi/2}\sin^3\varphi\,\mathrm{d}\varphi
  =
  \int_0^{\pi/2}\sin\varphi(1-\cos^2\varphi)\,\mathrm{d}\varphi
  =
  \frac23.
\]
故
\[
  I_{Oz}=\frac{4\pi a^4}{3}.
\]
\end{solution}

\subsubsection*{五、证明题（本大题共 2 小题，每小题 7 分，共 14 分）}
\begin{exercise}
\par
设函数 $f(x)$ 在区间 $[a,b]$ 上连续，证明：
  \(\displaystyle \int_a^b \mathrm{d}x\int_a^x f(y)\,\mathrm{d}y=\int_a^b f(y)(b-y)\,\mathrm{d}y\)。
\end{exercise}

\begin{solution}
\par
左端是二重积分
\[
  \int_a^b \mathrm{d}x\int_a^x f(y)\,\mathrm{d}y.
\]
它对应的积分区域为
\[
  a\le x\le b,\qquad a\le y\le x.
\]
也就是
\[
  a\le y\le x\le b.
\]
若改为先取 $y$，则 $y$ 从 $a$ 到 $b$；固定 $y$ 后，$x$ 从 $y$ 到 $b$。因此
\[
  \int_a^b \mathrm{d}x\int_a^x f(y)\,\mathrm{d}y
  =
  \int_a^b \mathrm{d}y\int_y^b f(y)\,\mathrm{d}x.
\]
由于 $f(y)$ 与 $x$ 无关，
\[
  \int_y^b f(y)\,\mathrm{d}x=f(y)(b-y).
\]
故
\[
  \int_a^b \mathrm{d}x\int_a^x f(y)\,\mathrm{d}y
  =
  \int_a^b f(y)(b-y)\,\mathrm{d}y.
\]
命题得证。
\end{solution}
\begin{exercise}
\par
设正数序列 $\{x_n\}$ 单调上升且有界，证明级数 \(\displaystyle \sum_{n=1}^{\infty}\left(1-\frac{x_n}{x_{n+1}}\right)\)
  收敛。
\end{exercise}

\begin{solution}
\par
因为 $\{x_n\}$ 是正数序列且单调上升，所以
\[
  0<x_1\le x_n\le x_{n+1}.
\]
于是每一项非负：
\[
  1-\frac{x_n}{x_{n+1}}
  =
  \frac{x_{n+1}-x_n}{x_{n+1}}
  \ge0.
\]
同时，由 $x_{n+1}\ge x_1$，有
\[
  1-\frac{x_n}{x_{n+1}}
  =
  \frac{x_{n+1}-x_n}{x_{n+1}}
  \le
  \frac{x_{n+1}-x_n}{x_1}.
\]
取前 $N$ 项部分和：
\[
  \sum_{n=1}^N\left(1-\frac{x_n}{x_{n+1}}\right)
  \le
  \frac1{x_1}\sum_{n=1}^N(x_{n+1}-x_n).
\]
右端望远镜相消：
\[
  \frac1{x_1}\sum_{n=1}^N(x_{n+1}-x_n)
  =
  \frac{x_{N+1}-x_1}{x_1}.
\]
由于 $\{x_n\}$ 单调上升且有界，存在常数 $M$ 使 $x_{N+1}\le M$，于是
\[
  \frac{x_{N+1}-x_1}{x_1}\le\frac{M-x_1}{x_1}.
\]
因此原级数是正项级数，且部分和有上界，所以收敛。
\end{solution}

\section{2009级工科数学分析下 B卷}

\subsubsection*{一、填空题（共 5 小题，每小题 4 分，共 20 分）}
\begin{exercise}
\par
方程 $y''+5y'+4=0$ 的通解为 \underline{\hspace{4cm}}。
\end{exercise}

\begin{solution}
\par
方程中不显含 $y$，令
\[
  p=y',
\]
则 $y''=p'$，原方程化为一阶线性方程
\[
  p'+5p=-4.
\]
先求齐次方程 $p'+5p=0$，得
\[
  p_h=C e^{-5x}.
\]
再取常数特解 $p=A$，代入
\[
  5A=-4,
\]
所以
\[
  A=-\frac45.
\]
于是
\[
  p=y'=C_1e^{-5x}-\frac45.
\]
对 $x$ 积分，常数倍可以重新记号，得到
\[
  y=C_2+C_1e^{-5x}-\frac45x.
\]
这就是通解。
\end{solution}
\begin{exercise}
\par
原点到平面 $\displaystyle x+\dfrac{y}{-4}+\dfrac{z}{3}=1$ 的距离为 \underline{\hspace{3cm}}。
\end{exercise}

\begin{solution}
\par
先把平面写成一般式：
\[
  x+\frac{y}{-4}+\frac z3=1
  \quad\Longleftrightarrow\quad
  x-\frac14y+\frac13z-1=0.
\]
点 $(x_0,y_0,z_0)$ 到平面 $Ax+By+Cz+D=0$ 的距离为
\[
  d=\frac{|Ax_0+By_0+Cz_0+D|}{\sqrt{A^2+B^2+C^2}}.
\]
代入原点 $(0,0,0)$，有
\[
  d=
  \frac{|-1|}{\sqrt{1+\frac1{16}+\frac1{9}}}.
\]
分母为
\[
  \sqrt{1+\frac1{16}+\frac1{9}}
  =
  \sqrt{\frac{144+9+16}{144}}
  =
  \frac{13}{12}.
\]
因此
\[
  d=\frac{12}{13}.
\]
\end{solution}
\begin{exercise}
\par
函数 $\displaystyle z=\arctan\dfrac{y}{x}+\ln\sqrt{x^2+y^2}$ 在点 $(1,0)$ 处的梯度 $\operatorname{grad}z\big|_{(1,0)}=$ \underline{\hspace{3cm}}。
\end{exercise}

\begin{solution}
\par
分开求导。第一部分
\[
  \arctan\frac yx
\]
的偏导数为
\[
  \left(\arctan\frac yx\right)_x=-\frac{y}{x^2+y^2},\qquad
  \left(\arctan\frac yx\right)_y=\frac{x}{x^2+y^2}.
\]
第二部分
\[
  \ln\sqrt{x^2+y^2}=\frac12\ln(x^2+y^2)
\]
的偏导数为
\[
  \left(\ln\sqrt{x^2+y^2}\right)_x=\frac{x}{x^2+y^2},\qquad
  \left(\ln\sqrt{x^2+y^2}\right)_y=\frac{y}{x^2+y^2}.
\]
故
\[
  z_x=\frac{x-y}{x^2+y^2},\qquad
  z_y=\frac{x+y}{x^2+y^2}.
\]
代入 $(1,0)$：
\[
  z_x(1,0)=1,\qquad z_y(1,0)=1.
\]
所以
\[
  \operatorname{grad}z\big|_{(1,0)}=(1,1).
\]
\end{solution}
\begin{exercise}
\par
设 $L$ 是 $xOy$ 平面上任意一条分段光滑的闭曲线，则第二类曲线积分
  \(\displaystyle \oint_L(2xy+\cos x)\,\mathrm{d}x+(x^2+\sin y)\,\mathrm{d}y=\underline{\hspace{3cm}}\)。
\end{exercise}

\begin{solution}
\par
记
\[
  P=2xy+\cos x,\qquad Q=x^2+\sin y.
\]
检查是否为全微分：
\[
  P_y=2x,\qquad Q_x=2x.
\]
在整个 $xOy$ 平面上 $P,Q$ 都有连续偏导数，且平面单连通，因此该曲线积分与路径无关。也可以直接找势函数 $\varphi$：
\[
  \varphi_x=P=2xy+\cos x.
\]
对 $x$ 积分得
\[
  \varphi=x^2y+\sin x+C(y).
\]
再对 $y$ 求偏导：
\[
  \varphi_y=x^2+C'(y).
\]
与 $Q=x^2+\sin y$ 比较，得
\[
  C'(y)=\sin y,
\]
可取
\[
  C(y)=-\cos y.
\]
故势函数为
\[
  \varphi=x^2y+\sin x-\cos y.
\]
闭曲线上的第二类曲线积分等于势函数终点值减起点值，而闭曲线起点终点相同，所以积分为
\[
  0.
\]
\end{solution}
\begin{exercise}
\par
设 $f(x)$ 是以 $2\pi$ 为周期的函数，它在 $[-\pi,\pi]$ 上的表达式是
  \[
  f(x)=
  \begin{cases}
  x+1,&-\pi\le x<0,\\
  x^2,&0\le x<\pi.
  \end{cases}
  \]
  若它的傅里叶级数的和函数为 $S(x)$，则 $S(0)=$ \underline{\hspace{3cm}}。
\end{exercise}

\begin{solution}
\par
傅里叶级数在跳跃间断点收敛到左右极限的平均值。这里
\[
  f(0-)=0+1=1,\qquad f(0+)=0^2=0.
\]
因此
\[
  S(0)=\frac{f(0-)+f(0+)}2
  =
  \frac{1+0}{2}
  =
  \frac12.
\]
这里不能直接把 \(x=0\) 代入某一侧表达式，因为 \(0\) 正是分段点；必须把周期延拓后的左、右极限分别算出，再按傅里叶收敛定理取平均。
最终空格应填 \(\dfrac12\)，而不是原函数在分段点处某一侧的定义值。
\end{solution}

\subsubsection*{二、选择题（共 5 小题，每小题 4 分，共 20 分）}
\begin{exercise}
\par
在 $R^3$ 中，方程 $x^2=4y$ 的图形是（\quad）。
  A. 抛物线；\quad B. 抛物柱面；\quad C. 椭圆抛物面；\quad D. 旋转抛物面。
\end{exercise}

\begin{solution}
\par
在三维空间中，方程
\[
  x^2=4y
\]
只限制 $x$ 与 $y$ 的关系，对 $z$ 没有任何限制。也就是说，在每一个固定的 $z$ 平面内，截线都是同一条抛物线
\[
  x^2=4y.
\]
这条抛物线沿 $z$ 轴方向平移生成曲面，因此图形是抛物柱面。选 B。
判断柱面时可看是否有某个变量完全没有出现；本题缺少 \(z\)，说明母线方向平行于 \(z\) 轴。
因此它不是平面内的一条抛物线，而是在空间中由抛物线沿 \(z\) 轴平行移动得到的一整张曲面。
\end{solution}
\begin{exercise}
\par
曲面 $3x^2+y^2+z^2=12$ 上点 $M(-1,0,3)$ 处的切平面与平面 $z=0$ 的夹角是（\quad）。
  A. $\displaystyle \dfrac{\pi}{6}$；\quad B. $\displaystyle \dfrac{\pi}{4}$；\quad C. $\displaystyle \dfrac{\pi}{3}$；\quad D. $\displaystyle \dfrac{\pi}{2}$。
\end{exercise}

\begin{solution}
\par
设
\[
  F(x,y,z)=3x^2+y^2+z^2.
\]
曲面在点 $M(-1,0,3)$ 处的法向量为
\[
  \nabla F(M)=(6x,2y,2z)\big|_{(-1,0,3)}=(-6,0,6).
\]
平面 $z=0$ 的法向量为
\[
  \boldsymbol k=(0,0,1).
\]
两平面的夹角等于它们法向量夹角的锐角，因此
\[
  \cos\theta
  =
  \frac{|(-6,0,6)\cdot(0,0,1)|}{|(-6,0,6)|\,|(0,0,1)|}.
\]
计算得
\[
  |(-6,0,6)|=6\sqrt2,\qquad
  |(0,0,1)|=1,
\]
且点积为 $6$，所以
\[
  \cos\theta=\frac6{6\sqrt2}=\frac1{\sqrt2}.
\]
故
\[
  \theta=\frac\pi4.
\]
选 B。
\end{solution}
\begin{exercise}
\par
设 $f(x,y)$ 在点 $(x_0,y_0)$ 处的两个偏导数都存在，则（\quad）。
  A. $f(x,y)$ 在点 $(x_0,y_0)$ 必连续；\quad
  B. $f(x,y)$ 在点 $(x_0,y_0)$ 必可微；\quad
  C. $\displaystyle\lim_{x\to x_0}f(x,y_0)$ 与 $\displaystyle\lim_{y\to y_0}f(x_0,y)$ 都存在；\quad
  D. $\displaystyle\lim_{\substack{x\to x_0\\y\to y_0}}f(x,y)$ 存在。
\end{exercise}

\begin{solution}
\par
偏导数存在的含义是沿坐标轴方向的一元导数存在：
\[
  f_x(x_0,y_0)=\lim_{h\to0}\frac{f(x_0+h,y_0)-f(x_0,y_0)}{h},
\]
\[
  f_y(x_0,y_0)=\lim_{k\to0}\frac{f(x_0,y_0+k)-f(x_0,y_0)}{k}.
\]
一元函数可导必连续，所以
\[
  \lim_{x\to x_0}f(x,y_0)=f(x_0,y_0),\qquad
  \lim_{y\to y_0}f(x_0,y)=f(x_0,y_0).
\]
因此 C 正确。

但偏导数存在不保证二元函数连续，也不保证可微，更不保证二重极限存在，所以 A、B、D 都不能必然成立。选 C。
\end{solution}
\begin{exercise}
\par
设平面曲线 $L$ 为下半圆周 $y=-\sqrt{1-x^2}$，则曲线积分 $\displaystyle\int_L(x^2+y^2)\,\mathrm{d}s$（\quad）。
  A. $0$；\quad B. $-1$；\quad C. $2$；\quad D. $\pi$。
\end{exercise}

\begin{solution}
\par
曲线
\[
  y=-\sqrt{1-x^2}
\]
是单位圆 $x^2+y^2=1$ 的下半圆周。因此在 $L$ 上恒有
\[
  x^2+y^2=1.
\]
第一类曲线积分与方向无关，所以
\[
  \int_L(x^2+y^2)\,\mathrm{d}s
  =
  \int_L1\,\mathrm{d}s.
\]
右端就是下半单位圆的弧长。单位圆半周长为
\[
  \pi.
\]
故积分值为 $\pi$，选 D。
如果按参数计算，可取 \(x=\cos t,\ y=-\sin t,\ 0\le t\le\pi\)，此时 \(\mathrm{d}s=\mathrm{d}t\)，被积函数恒为 \(1\)，积分就是 \(\int_0^\pi 1\,\mathrm{d}t\)。这也说明第一类曲线积分只看弧长，不看行进方向。
\end{solution}
\begin{exercise}
\par
幂级数 $\displaystyle\sum_{n=1}^{\infty}\frac{\ln(n+1)}{n}x^n$ 的收敛域为（\quad）。
  A. $(-1,1)$；\quad B. $[-1,1]$；\quad C. $[-1,1)$；\quad D. $(-1,1]$。
\end{exercise}

\begin{solution}
\par
先求收敛半径。系数为
\[
  c_n=\frac{\ln(n+1)}{n}.
\]
由于
\[
  \sqrt[n]{c_n}\to1,
\]
所以幂级数的收敛半径为
\[
  R=1.
\]
因此在 $|x|<1$ 时收敛，在 $|x|>1$ 时发散，端点另判。

当 $x=1$ 时，级数为
\[
  \sum_{n=1}^{\infty}\frac{\ln(n+1)}{n}.
\]
由于对充分大的 $n$ 有 $\ln(n+1)>1$，该级数逐项大于调和级数的一部分，故发散。

当 $x=-1$ 时，级数为
\[
  \sum_{n=1}^{\infty}(-1)^n\frac{\ln(n+1)}{n}.
\]
此时
\[
  \frac{\ln(n+1)}{n}\to0,
\]
且从某项以后单调递减，因此由交错级数判别法收敛。

故收敛域为
\[
  [-1,1).
\]
选 C。
\end{solution}

\subsubsection*{三、计算题（本大题共 4 小题，每小题 8 分，共 32 分）}
\begin{exercise}
\par
求微分方程 $y''+y'+2y=0$ 的通解。
\end{exercise}

\begin{solution}
\par
这是二阶常系数齐次线性方程。写出特征方程：
\[
  r^2+r+2=0.
\]
判别式为
\[
  \Delta=1-8=-7,
\]
所以特征根为
\[
  r=\frac{-1\pm i\sqrt7}{2}.
\]
若特征根为 $\alpha\pm i\beta$，通解为
\[
  y=e^{\alpha x}(C_1\cos\beta x+C_2\sin\beta x).
\]
这里
\[
  \alpha=-\frac12,\qquad \beta=\frac{\sqrt7}{2}.
\]
因此
\[
  y=e^{-x/2}\left(C_1\cos\frac{\sqrt7}{2}x+C_2\sin\frac{\sqrt7}{2}x\right).
\]
\end{solution}
\begin{exercise}
\par
计算二重积分 $\displaystyle I=\iint_D\frac{x}{x^2+y^2}\,\mathrm{d}x\,\mathrm{d}y$，其中 $D$ 由曲线 $\displaystyle y=\dfrac{x^2}{2}$ 与直线 $y=x$ 围成。
\end{exercise}

\begin{solution}
\par
先求区域。两条边界为
\[
  y=\frac{x^2}{2},\qquad y=x.
\]
交点由
\[
  \frac{x^2}{2}=x
\]
得
\[
  x=0,\quad x=2.
\]
在 $0\le x\le2$ 上，有
\[
  \frac{x^2}{2}\le y\le x.
\]
因此
\[
  I=\int_0^2\int_{x^2/2}^{x}\frac{x}{x^2+y^2}\,\mathrm{d}y\,\mathrm{d}x.
\]
对 $y$ 积分时 $x$ 为常数：
\[
  \int \frac{x}{x^2+y^2}\,\mathrm{d}y
  =
  \arctan\frac{y}{x}.
\]
所以
\[
  I=\int_0^2
  \left[
  \arctan\frac{y}{x}
  \right]_{y=x^2/2}^{y=x}
  \mathrm{d}x.
\]
当 $y=x$ 时，$\arctan(y/x)=\arctan1=\pi/4$；当 $y=x^2/2$ 时，$\arctan(y/x)=\arctan(x/2)$。故
\[
  I=\int_0^2\left(\frac\pi4-\arctan\frac{x}{2}\right)\mathrm{d}x.
\]
计算
\[
  \int_0^2\arctan\frac{x}{2}\,\mathrm{d}x
  =
  \left[x\arctan\frac{x}{2}-\ln(x^2+4)\right]_0^2
  =
  \frac\pi2-\ln2.
\]
因此
\[
  I=\frac\pi2-\left(\frac\pi2-\ln2\right)=\ln2.
\]
\end{solution}
\begin{exercise}
\par
计算曲线积分 \(\displaystyle I=\int_L (x+y)^2\,\mathrm{d}x-(x^2+y^2\sin y)\,\mathrm{d}y\)，
  其中 $L$ 是抛物线 $y=x^2$ 上从点 $(-1,1)$ 到点 $(1,1)$ 的那一段。
\end{exercise}

\begin{solution}
\par
曲线方向从 $(-1,1)$ 到 $(1,1)$，取参数
\[
  x=t,\qquad y=t^2,\qquad -1\le t\le1.
\]
于是
\[
  \mathrm{d}x=\mathrm{d}t,\qquad \mathrm{d}y=2t\,\mathrm{d}t.
\]
代入积分：
\[
  I=
  \int_{-1}^{1}
  \left[(t+t^2)^2-\bigl(t^2+t^4\sin(t^2)\bigr)2t\right]\mathrm{d}t.
\]
化简被积函数：
\[
  (t+t^2)^2-2t(t^2+t^4\sin(t^2))
  =
  t^2+t^4-2t^5\sin(t^2).
\]
其中 $-2t^5\sin(t^2)$ 是奇函数，在 $[-1,1]$ 上积分为 $0$。故
\[
  I=\int_{-1}^{1}(t^2+t^4)\,\mathrm{d}t
  =
  2\int_0^1(t^2+t^4)\,\mathrm{d}t
  =
  \frac{16}{15}.
\]
\end{solution}
\begin{exercise}
\par
求幂级数 $\displaystyle\sum_{n=0}^{\infty}\frac{n}{n+1}x^n$ 的收敛域及和函数。
\end{exercise}

\begin{solution}
\par
先判收敛域。幂级数系数为
\[
  \frac{n}{n+1}\to1,
\]
所以收敛半径为
\[
  R=1.
\]
当 $|x|<1$ 时收敛。当 $x=1$ 时，通项为
\[
  \frac{n}{n+1}\to1\ne0,
\]
故发散。当 $x=-1$ 时，通项为
\[
  \frac{n}{n+1}(-1)^n,
\]
仍不趋于 $0$，故也发散。因此收敛域为
\[
  (-1,1).
\]

在 $|x|<1$ 内求和。利用
\[
  \frac{n}{n+1}=1-\frac1{n+1},
\]
有
\[
  \sum_{n=0}^{\infty}\frac{n}{n+1}x^n
  =
  \sum_{n=0}^{\infty}x^n
  -
  \sum_{n=0}^{\infty}\frac{x^n}{n+1}.
\]
第一项为
\[
  \sum_{n=0}^{\infty}x^n=\frac1{1-x}.
\]
第二项由
\[
  -\ln(1-x)=\sum_{n=1}^{\infty}\frac{x^n}{n}
\]
得到
\[
  \sum_{n=0}^{\infty}\frac{x^n}{n+1}
  =
  \frac{-\ln(1-x)}{x}\qquad (x\ne0).
\]
因此
\[
  S(x)=\frac1{1-x}+\frac{\ln(1-x)}{x},\qquad 0<|x|<1.
\]
当 $x=0$ 时，原级数每项除首项外含 $x^n$，且首项系数为 $0$，所以
\[
  S(0)=0.
\]
上式在 $x\to0$ 时也趋于 $0$，故可把 $S(0)=0$ 作为补充定义。
\end{solution}

\subsubsection*{四、应用题（本大题共 2 小题，每小题 7 分，共 14 分）}
\begin{exercise}
\par
欲造一个无盖的长方体容器，已知底部造价为 $3$ 元/平方米，侧面造价为 $1$ 元/平方米。现想用 36 元造一个容积最大的容器，问应该怎样设计尺寸？
\end{exercise}

\begin{solution}
\par
设底面长、宽分别为 $x,y$，高为 $h$。无盖容器的费用由底面和四个侧面组成：
\[
  3xy+2xh+2yh=36.
\]
目标是最大化体积
\[
  V=xyh.
\]
由于底面长和宽在费用与体积中地位对称，极值点应满足 $x=y$；也可由 Lagrange 乘子法推出这一点。设
\[
  x=y=a.
\]
费用约束变为
\[
  3a^2+4ah=36,
\]
所以
\[
  h=\frac{36-3a^2}{4a}.
\]
体积为
\[
  V(a)=a^2h
  =
  \frac{a(36-3a^2)}4
  =
  9a-\frac34a^3.
\]
求导：
\[
  V'(a)=9-\frac94a^2.
\]
令 $V'(a)=0$，得
\[
  a^2=4.
\]
因 $a>0$，故
\[
  a=2.
\]
代回费用约束：
\[
  3\cdot 2^2+4\cdot2\cdot h=36,
\]
即
\[
  12+8h=36,\qquad h=3.
\]
因此应设计为底面 $2\times2$ 的正方形，高为 $3$。
\end{solution}
\begin{exercise}
\par
设 $\mathbf F(x,y)=xy^2\mathbf i+yf(x)\mathbf j$ 为保守场，其中 $f(x)$ 具有连续导数，且 $f(0)=0$。求曲线积分
  \(\displaystyle I=\int_{(0,0)}^{(1,1)}\mathbf F\cdot d\mathbf s\)。
\end{exercise}

\begin{solution}
\par
写成
\[
  \mathbf F=(P,Q),\qquad P=xy^2,\qquad Q=yf(x).
\]
若 $\mathbf F$ 为保守场，则在单连通区域内必须满足
\[
  P_y=Q_x.
\]
分别计算：
\[
  P_y=2xy,\qquad Q_x=yf'(x).
\]
因此
\[
  2xy=yf'(x).
\]
作为恒等式，应有
\[
  f'(x)=2x.
\]
积分得
\[
  f(x)=x^2+C.
\]
由 $f(0)=0$ 得 $C=0$，所以
\[
  f(x)=x^2.
\]
此时
\[
  \mathbf F=(xy^2,x^2y).
\]
求势函数 $\varphi$，使
\[
  \varphi_x=xy^2,\qquad \varphi_y=x^2y.
\]
由 $\varphi_x=xy^2$ 对 $x$ 积分：
\[
  \varphi=\frac12x^2y^2+C(y).
\]
再对 $y$ 求偏导：
\[
  \varphi_y=x^2y+C'(y).
\]
与 $Q=x^2y$ 比较，得 $C'(y)=0$，故可取
\[
  \varphi=\frac12x^2y^2.
\]
曲线积分与路径无关，等于势函数增量：
\[
  I=\varphi(1,1)-\varphi(0,0)
  =
  \frac12-0
  =
  \frac12.
\]
\end{solution}

\subsubsection*{五、证明题（本大题共 2 小题，每小题 7 分，共 14 分）}
\begin{exercise}
\par
证明：圆柱螺线 $x=a\cos t,\ y=a\sin t,\ z=bt$ 的切线与 $z$ 轴的夹角为常数。
\end{exercise}

\begin{solution}
\par
圆柱螺线的向量参数方程为
\[
  \boldsymbol r(t)=(a\cos t,\ a\sin t,\ bt).
\]
它的切向量为
\[
  \boldsymbol r'(t)=(-a\sin t,\ a\cos t,\ b).
\]
$z$ 轴的方向向量可取
\[
  \boldsymbol k=(0,0,1).
\]
设切线与 $z$ 轴的夹角为 $\theta$，则
\[
  \cos\theta
  =
  \frac{|\boldsymbol r'(t)\cdot\boldsymbol k|}
  {|\boldsymbol r'(t)|\,|\boldsymbol k|}.
\]
计算点积：
\[
  \boldsymbol r'(t)\cdot\boldsymbol k=b.
\]
再计算切向量长度：
\[
  |\boldsymbol r'(t)|
  =
  \sqrt{a^2\sin^2t+a^2\cos^2t+b^2}
  =
  \sqrt{a^2+b^2}.
\]
因此
\[
  \cos\theta=\frac{|b|}{\sqrt{a^2+b^2}},
\]
右端与 $t$ 无关，所以 $\theta$ 为常数。命题得证。
\end{solution}
\begin{exercise}
\par
若正项级数 $\displaystyle\sum_{n=1}^{\infty}a_n$ 收敛，证明正项级数 \(\displaystyle \sum_{n=1}^{\infty}\frac{\sqrt{a_n}}{n}\)
  也收敛。
\end{exercise}

\begin{solution}
\par
因为 $a_n>0$，可使用基本不等式
\[
  2\sqrt{uv}\le u+v\qquad(u,v\ge0).
\]
取
\[
  u=a_n,\qquad v=\frac1{n^2},
\]
则
\[
  2\sqrt{a_n\cdot\frac1{n^2}}
  \le
  a_n+\frac1{n^2}.
\]
也就是
\[
  \frac{\sqrt{a_n}}{n}
  \le
  \frac12a_n+\frac1{2n^2}.
\]
题设 $\sum a_n$ 收敛，而
\[
  \sum_{n=1}^{\infty}\frac1{n^2}
\]
也收敛，所以
\[
  \sum_{n=1}^{\infty}\left(\frac12a_n+\frac1{2n^2}\right)
\]
收敛。由比较判别法，正项级数
\[
  \sum_{n=1}^{\infty}\frac{\sqrt{a_n}}{n}
\]
收敛。
\end{solution}

\section{2012级工科数学分析下 A卷}

\subsubsection*{一、单项选择题（每小题 3 分，共 15 分）}
\begin{exercise}
\par
下述函数中有可能成为二阶微分方程 \(\displaystyle f(x,y,y'')=0\) 的通解的是（\quad）。

  A. $C_1x+C_2y=0$；\quad B. $y-C_1^2=e^x+C_2$；\quad
  C. $C_1y=C_2e^{C_3x}$；\quad D. $y=C_1\ln x+C_2x+C_3$。
\end{exercise}

\begin{solution}
\par
二阶微分方程的通解通常应含有两个任意常数，并且要能表示出未知函数 $y$ 的一般形式。逐项看：

A 项
\[
  C_1x+C_2y=0
\]
若解出 $y$，只得到 $y=-\dfrac{C_1}{C_2}x$，本质上只有一个任意常数。C、D 两项含有三个任意常数，超出了二阶方程通解的常规个数。按本题的选择题口径，B 项能写成
\[
  y=e^x+C_1^2+C_2,
\]
形式上含两个任意常数并能表示 $y$，故选 B。
这里的“有可能”按通解形式筛选，不要求反推出具体方程；先看能否表示 \(y\)，再看独立任意常数的个数。
\end{solution}
\begin{exercise}
\par
设 $f(x,y)$ 是可微函数，且 \(\displaystyle f(0,0)=0,\ f_x(0,0)=a,\ f_y(0,0)=b\)，令 $\varphi(t)=f(t,f(t,t))$，则 $\varphi'(0)=$（\quad）。

  A. $a$；\quad B. $a+b(a+b)$；\quad C. $a+1$；\quad D. $\displaystyle \dfrac{a}{1-b}$。
\end{exercise}

\begin{solution}
\par
设
\[
  \varphi(t)=f(t,f(t,t)).
\]
这是复合函数。外层 $f$ 的两个自变量分别为
\[
  X(t)=t,\qquad Y(t)=f(t,t).
\]
由链式法则
\[
  \varphi'(t)=f_x(X(t),Y(t))X'(t)+f_y(X(t),Y(t))Y'(t).
\]
先求 $t=0$ 时的内层点。题设 $f(0,0)=0$，所以
\[
  X(0)=0,\qquad Y(0)=f(0,0)=0.
\]
又
\[
  X'(t)=1.
\]
对 $Y(t)=f(t,t)$ 求导：
\[
  Y'(t)=f_x(t,t)+f_y(t,t).
\]
因此
\[
  Y'(0)=f_x(0,0)+f_y(0,0)=a+b.
\]
代入链式法则：
\[
  \varphi'(0)=f_x(0,0)+f_y(0,0)(a+b)
  =
  a+b(a+b).
\]
故选 B。
\end{solution}
\begin{exercise}
\par
\(\displaystyle \int_0^1\mathrm{d}y\int_{-y}^{\sqrt{2y-y^2}} f(x,y)\,\mathrm{d}x\)
  交换积分顺序后，正确的结果是（\quad）。

  A. $\displaystyle \int_{-1}^0\mathrm{d}x\int_{-x}^1f(x,y)\,\mathrm{d}y+\int_0^1\mathrm{d}x\int_{1-\sqrt{1-x^2}}^1f(x,y)\,\mathrm{d}y$；

  B. $\displaystyle \int_0^1\mathrm{d}x\int_{1-\sqrt{1-x^2}}^1f(x,y)\,\mathrm{d}y$；

  C. $\displaystyle \int_{-1}^0\mathrm{d}x\int_{-x}^1f(x,y)\,\mathrm{d}y$；

  D. $\displaystyle \int_{-1}^1\mathrm{d}x\int_{1-\sqrt{1-x^2}}^1f(x,y)\,\mathrm{d}y$。
\end{exercise}

\begin{solution}
\par
原积分区域由
\[
  0\le y\le1,\qquad -y\le x\le \sqrt{2y-y^2}
\]
给出。右边界
\[
  x=\sqrt{2y-y^2}
\]
等价于
\[
  x^2=2y-y^2
  \quad\Longleftrightarrow\quad
  x^2+(y-1)^2=1,
\]
且取的是右半圆弧。左边界 $x=-y$ 等价于
\[
  y=-x.
\]

换成先取 $x$ 时，要按 $x$ 的正负分段。左侧部分为
\[
  -1\le x\le0,\qquad -x\le y\le1.
\]
右侧部分来自圆
\[
  x^2+(y-1)^2\le1,
\]
即
\[
  0\le x\le1,\qquad 1-\sqrt{1-x^2}\le y\le1.
\]
所以换序结果为
\[
  \int_{-1}^0\mathrm{d}x\int_{-x}^1f(x,y)\,\mathrm{d}y
  +
  \int_0^1\mathrm{d}x\int_{1-\sqrt{1-x^2}}^1f(x,y)\,\mathrm{d}y.
\]
故选 A。
\end{solution}
\begin{exercise}
\par
设 $C$ 是右半圆周 \(\displaystyle x^2+y^2=a^2,\ x\ge0,\ a>0\)，则曲线积分 \(\displaystyle \int_C(x+y)\,\mathrm{d}s\)（\quad）。

  A. $0$；\quad B. $2a^2$；\quad C. $a^2$；\quad D. $4a$。
\end{exercise}

\begin{solution}
\par
右半圆周可参数化为
\[
  x=a\cos t,\qquad y=a\sin t,\qquad -\frac\pi2\le t\le\frac\pi2.
\]
此时
\[
  \mathrm{d}s=\sqrt{(-a\sin t)^2+(a\cos t)^2}\,\mathrm{d}t=a\,\mathrm{d}t.
\]
所以
\[
  \int_C(x+y)\,\mathrm{d}s
  =
  \int_{-\pi/2}^{\pi/2}(a\cos t+a\sin t)\,a\,\mathrm{d}t.
\]
即
\[
  a^2\int_{-\pi/2}^{\pi/2}(\cos t+\sin t)\,\mathrm{d}t.
\]
其中 $\sin t$ 在对称区间上积分为 $0$，而
\[
  \int_{-\pi/2}^{\pi/2}\cos t\,\mathrm{d}t=2.
\]
故积分值为
\[
  2a^2.
\]
选 B。
\end{solution}
\begin{exercise}
\par
设 $f(x)=x^2\ (0\le x<1)$，而
  \[
  S(x)=\sum_{n=1}^{\infty}b_n\sin n\pi x\quad(-\infty<x<+\infty),
  \qquad b_n=2\int_0^1f(x)\sin n\pi x\,\mathrm{d}x,
  \]
  则 $\displaystyle S\!\left(-\dfrac12\right)$（\quad）。

  A. $\displaystyle -\dfrac12$；\quad B. $\displaystyle -\dfrac14$；\quad C. $\displaystyle \dfrac14$；\quad D. $\displaystyle \dfrac12$。
\end{exercise}

\begin{solution}
\par
这是 $[0,1]$ 上的正弦级数展开。正弦级数对应把 $f(x)$ 作奇延拓，并且周期为 $2$。因此在连续点处，
\[
  S(-x)=-S(x).
\]
又 $x=1/2$ 是原函数的连续点，所以
\[
  S\left(\frac12\right)=f\left(\frac12\right)=\left(\frac12\right)^2=\frac14.
\]
由奇延拓性质，
\[
  S\left(-\frac12\right)=-S\left(\frac12\right)=-\frac14.
\]
故选 B。
\end{solution}

\subsubsection*{二、填空题（每小题 3 分，共 15 分）}
\begin{exercise}
\par
微分方程 \(\displaystyle yy''-y'^2=0\) 的通解为 \underline{\hspace{4cm}}。
\end{exercise}

\begin{solution}
\par
注意
\[
  \left(\frac{y'}{y}\right)'
  =
  \frac{y''y-(y')^2}{y^2}.
\]
原方程
\[
  yy''-(y')^2=0
\]
等价于
\[
  \left(\frac{y'}{y}\right)'=0
\]
在 $y\ne0$ 的区间上成立。于是
\[
  \frac{y'}{y}=k,
\]
其中 $k$ 为常数。分离变量：
\[
  \frac{\mathrm{d}y}{y}=k\,\mathrm{d}x.
\]
积分得
\[
  \ln|y|=kx+C,
\]
所以
\[
  y=Ce^{kx}.
\]
当 $C=0$ 时也包含零解，因此通解可写为
\[
  y=Ce^{kx}.
\]
\end{solution}
\begin{exercise}
\par
设 \(\displaystyle u=x+(y-1)\arcsin\frac{x}{y}\)，则 $\displaystyle \left.\dfrac{\partial u}{\partial x}\right|_{(1,2)}=$ \underline{\hspace{3cm}}。
\end{exercise}

\begin{solution}
\par
对 $x$ 求偏导时把 $y$ 看成常数。由
\[
  u=x+(y-1)\arcsin\frac{x}{y}
\]
得
\[
  u_x=1+(y-1)\cdot
  \frac{1}{\sqrt{1-(x/y)^2}}\cdot\frac1y.
\]
注意这里第一项 $x$ 对 $x$ 的偏导为 $1$，不能漏掉。化简第二项：
\[
  \frac{1}{\sqrt{1-(x/y)^2}}\cdot\frac1y
  =
  \frac1{\sqrt{y^2-x^2}}
\]
在 $y>0$ 的点 $(1,2)$ 附近成立。故
\[
  u_x=1+\frac{y-1}{\sqrt{y^2-x^2}}.
\]
代入 $(1,2)$：
\[
  u_x(1,2)=1+\frac{1}{\sqrt{4-1}}
  =
  1+\frac1{\sqrt3}.
\]
故填
\[
  1+\frac1{\sqrt3}.
\]
\end{solution}
\begin{exercise}
\par
设区域 \(\displaystyle D:0\le y\le x,\ 0\le x\le\pi\)，则二重积分
  \(\displaystyle \iint_D \sqrt{1-\sin^2x}\,\mathrm{d}x\,\mathrm{d}y=\) \underline{\hspace{3cm}}。
\end{exercise}

\begin{solution}
\par
区域为
\[
  0\le x\le\pi,\qquad 0\le y\le x.
\]
在该区间上
\[
  \sqrt{1-\sin^2x}=|\cos x|.
\]
因此
\[
  \iint_D \sqrt{1-\sin^2x}\,\mathrm{d}x\,\mathrm{d}y
  =
  \int_0^\pi\int_0^x|\cos x|\,\mathrm{d}y\,\mathrm{d}x
  =
  \int_0^\pi x|\cos x|\,\mathrm{d}x.
\]
分段计算：
\[
  \int_0^\pi x|\cos x|\,\mathrm{d}x
  =
  \int_0^{\pi/2}x\cos x\,\mathrm{d}x
  -
  \int_{\pi/2}^{\pi}x\cos x\,\mathrm{d}x.
\]
由分部积分得
\[
  \int_0^{\pi/2}x\cos x\,\mathrm{d}x=\frac\pi2-1,
\]
且
\[
  -\int_{\pi/2}^{\pi}x\cos x\,\mathrm{d}x=1+\frac\pi2.
\]
相加得
\[
  \pi.
\]
\end{solution}
\begin{exercise}
\par
已知曲线 $C$ 为 \(\displaystyle x^2+y^2=ax\)，则曲线积分
  \(\displaystyle \int_C\sqrt{x^2+y^2}\,\mathrm{d}s=\) \underline{\hspace{3cm}}。
\end{exercise}

\begin{solution}
\par
曲线
\[
  x^2+y^2=ax
\]
化为极坐标。由于 $x=r\cos\theta$，$x^2+y^2=r^2$，故
\[
  r^2=ar\cos\theta.
\]
除原点外，曲线为
\[
  r=a\cos\theta,\qquad -\frac\pi2\le\theta\le\frac\pi2.
\]
这是圆心 $(a/2,0)$、半径 $a/2$ 的圆。被积函数
\[
  \sqrt{x^2+y^2}=r=a\cos\theta.
\]
极坐标下曲线弧长元为
\[
  \mathrm{d}s=\sqrt{r^2+\left(\frac{\mathrm{d}r}{\mathrm{d}\theta}\right)^2}\,\mathrm{d}\theta.
\]
这里
\[
  r=a\cos\theta,\qquad \frac{\mathrm{d}r}{\mathrm{d}\theta}=-a\sin\theta,
\]
所以
\[
  \mathrm{d}s=a\,\mathrm{d}\theta.
\]
于是
\[
  \int_C\sqrt{x^2+y^2}\,\mathrm{d}s
  =
  \int_{-\pi/2}^{\pi/2}a\cos\theta\cdot a\,\mathrm{d}\theta
  =
  a^2\left[\sin\theta\right]_{-\pi/2}^{\pi/2}
  =
  2a^2.
\]
\end{solution}
\begin{exercise}
\par
幂级数 \(\displaystyle \sum_{n=1}^{\infty}\frac{(x-2)^n}{n4^n}\) 的收敛域为 \underline{\hspace{4cm}}。
\end{exercise}

\begin{solution}
\par
令
\[
  t=\frac{x-2}{4}.
\]
原级数变为
\[
  \sum_{n=1}^{\infty}\frac{t^n}{n}.
\]
它的收敛半径为 $1$，所以先有
\[
  |t|<1
  \quad\Longleftrightarrow\quad
  \left|\frac{x-2}{4}\right|<1
  \quad\Longleftrightarrow\quad
  -2<x<6.
\]
端点另判。若 $x=6$，则 $t=1$，级数为
\[
  \sum_{n=1}^{\infty}\frac1n,
\]
发散。若 $x=-2$，则 $t=-1$，级数为
\[
  \sum_{n=1}^{\infty}\frac{(-1)^n}{n},
\]
由交错级数判别法收敛。故收敛域为
\[
  [-2,6).
\]
\end{solution}

\subsubsection*{三、计算题（每小题 10 分，共 50 分）}
\begin{exercise}
\par
计算三重积分 \(\displaystyle I=\iiint_\Omega (x^2+y^2+z^2)\,\mathrm{d}V\)，
  其中 $\Omega:x^2+y^2+z^2\le2z$。
\end{exercise}

\begin{solution}
\par
先配方：
\[
  x^2+y^2+z^2\le2z
  \quad\Longleftrightarrow\quad
  x^2+y^2+(z-1)^2\le1.
\]
因此 $\Omega$ 是球心 $(0,0,1)$、半径 $1$ 的球。令
\[
  w=z-1,
\]
则区域变成单位球
\[
  x^2+y^2+w^2\le1.
\]
被积函数为
\[
  x^2+y^2+z^2
  =
  x^2+y^2+(w+1)^2
  =
  x^2+y^2+w^2+1+2w.
\]
单位球关于 $w=0$ 对称，故奇函数 $2w$ 的积分为 $0$。于是
\[
  I=
  \iiint_{x^2+y^2+w^2\le1}(x^2+y^2+w^2)\,\mathrm{d}V
  +
  \iiint_{x^2+y^2+w^2\le1}1\,\mathrm{d}V.
\]
用球坐标计算第一项：
\[
  \int_0^{2\pi}\int_0^\pi\int_0^1
  \rho^2\cdot \rho^2\sin\varphi\,\mathrm{d}\rho\,\mathrm{d}\varphi\,\mathrm{d}\theta
  =
  4\pi\int_0^1\rho^4\,\mathrm{d}\rho
  =
  \frac{4\pi}{5}.
\]
第二项为单位球体积：
\[
  \frac{4\pi}{3}.
\]
所以
\[
  I=\frac{4\pi}{5}+\frac{4\pi}{3}
  =
  \frac{32\pi}{15}.
\]
\end{solution}
\begin{exercise}
\par
设 $\Sigma$ 是锥面 $z=\sqrt{x^2+y^2}$ 被平面 $z=0$ 及 $z=1$ 所截部分的外侧，计算第二类曲面积分 \(\displaystyle I=\iint_\Sigma xy\,\mathrm{d}y\,\mathrm{d}z+y\,\mathrm{d}z\,\mathrm{d}x+(z^2-2z)\,\mathrm{d}x\,\mathrm{d}y\)。
\end{exercise}

\begin{solution}
\par
记
\[
  P=xy,\qquad Q=y,\qquad R=z^2-2z.
\]
锥面为 $z=\sqrt{x^2+y^2}=r$，被 $z=0$ 与 $z=1$ 截下后，其在 $xOy$ 平面上的投影为
\[
  D:\ x^2+y^2\le1.
\]
对圆锥体而言，外侧法向远离 $z$ 轴，且 $z$ 分量为负。把锥面写成 $z=f(x,y)=r$，下侧方向对应的有向面积元为
\[
  (f_x,f_y,-1)\,\mathrm{d}x\,\mathrm{d}y
  =
  \left(\frac{x}{r},\frac{y}{r},-1\right)\,\mathrm{d}x\,\mathrm{d}y.
\]
因此第二类曲面积分化为
\[
  I=\iint_D\left(P\frac{x}{r}+Q\frac{y}{r}-R\right)\,\mathrm{d}x\,\mathrm{d}y.
\]
在锥面上 $z=r$，所以
\[
  R=r^2-2r.
\]
代入 $P,Q$：
\[
  I=\iint_D\left(\frac{x^2y}{r}+\frac{y^2}{r}-r^2+2r\right)\,\mathrm{d}x\,\mathrm{d}y.
\]
转为极坐标。第一项 $\dfrac{x^2y}{r}$ 关于 $y$ 为奇函数，在圆盘 $D$ 上积分为 $0$。第二项
\[
  \iint_D\frac{y^2}{r}\,\mathrm{d}x\,\mathrm{d}y
  =
  \int_0^{2\pi}\int_0^1
  \frac{r^2\sin^2\theta}{r}\,r\,\mathrm{d}r\,\mathrm{d}\theta
  =
  \left(\int_0^{2\pi}\sin^2\theta\,\mathrm{d}\theta\right)
  \left(\int_0^1r^2\,\mathrm{d}r\right)
  =
  \frac{\pi}{3}.
\]
第三项
\[
  \iint_D(-r^2)\,\mathrm{d}x\,\mathrm{d}y
  =
  -\int_0^{2\pi}\int_0^1r^3\,\mathrm{d}r\,\mathrm{d}\theta
  =
  -\frac{\pi}{2}.
\]
第四项
\[
  \iint_D2r\,\mathrm{d}x\,\mathrm{d}y
  =
  \int_0^{2\pi}\int_0^1 2r^2\,\mathrm{d}r\,\mathrm{d}\theta
  =
  \frac{4\pi}{3}.
\]
合并：
\[
  I=0+\frac{\pi}{3}-\frac{\pi}{2}+\frac{4\pi}{3}
  =
  \frac{7\pi}{6}.
\]
\end{solution}
\begin{exercise}
\par
求幂级数 \(\displaystyle \sum_{n=0}^{\infty}(2n+1)x^n\) 的和函数，并据此求数项级数 \(\displaystyle \sum_{n=0}^{\infty}\frac{2n+1}{2^n}\) 的值。
\end{exercise}

\begin{solution}
\par
先由几何级数出发：
\[
  \sum_{n=0}^{\infty}x^n=\frac1{1-x},\qquad |x|<1.
\]
两边对 $x$ 求导，得
\[
  \sum_{n=1}^{\infty}n x^{n-1}=\frac1{(1-x)^2}.
\]
乘以 $x$：
\[
  \sum_{n=1}^{\infty}n x^n=\frac{x}{(1-x)^2}.
\]
因此
\[
  \sum_{n=0}^{\infty}(2n+1)x^n
  =
  2\sum_{n=0}^{\infty}n x^n+\sum_{n=0}^{\infty}x^n
  =
  \frac{2x}{(1-x)^2}+\frac1{1-x}.
\]
通分：
\[
  \frac{2x}{(1-x)^2}+\frac1{1-x}
  =
  \frac{2x+1-x}{(1-x)^2}
  =
  \frac{1+x}{(1-x)^2}.
\]
所以和函数为
\[
  S(x)=\frac{1+x}{(1-x)^2},\qquad |x|<1.
\]
要求的数项级数对应 $x=1/2$：
\[
  \sum_{n=0}^{\infty}\frac{2n+1}{2^n}
  =
  S\left(\frac12\right)
  =
  \frac{1+\frac12}{\left(1-\frac12\right)^2}
  =
  6.
\]
\end{solution}
\begin{exercise}
\par
求 $a,b$ 的值，使得包含圆周 \(\displaystyle (x-1)^2+y^2=1\)
  在其内部的椭圆 \(\displaystyle \frac{x^2}{a^2}+\frac{y^2}{b^2}=1\quad(a>0,\ b>0,\ a\ne b)\) 有最小的面积。
\end{exercise}

\begin{solution}
\par
令
\[
  \alpha=\frac1{a^2},\qquad \beta=\frac1{b^2}.
\]
椭圆面积为 $\pi ab=\dfrac{\pi}{\sqrt{\alpha\beta}}$，所以面积最小等价于 $\alpha\beta$ 最大。圆周参数化为
\[
  x=1+\cos t,\qquad y=\sin t.
\]
令 $u=\cos t$，则 $u\in[-1,1]$，包含条件为
\[
  \alpha(1+u)^2+\beta(1-u^2)\le1.
\]
再令 $s=1+u$，则 $0\le s\le2$，且 $1-u^2=s(2-s)$。条件变成
\[
  \alpha s^2+\beta s(2-s)\le1.
\]
对 $s>0$，这等价于
\[
  \beta\le \frac{1-\alpha s^2}{s(2-s)}.
\]
最优情形下右端在某个内点达到最小值。记
\[
  q(s)=\frac{1-\alpha s^2}{s(2-s)}.
\]
由 $q'(s)=0$ 得
\[
  \alpha=\frac{s-1}{s^2}.
\]
此时
\[
  \beta=q(s)=\frac1s,
\]
所以
\[
  \alpha\beta=\frac{s-1}{s^3}.
\]
对 $1<s\le2$ 最大化该式：
\[
  \frac{\mathrm{d}}{\mathrm{d}s}\left(\frac{s-1}{s^3}\right)
  =
  \frac{3-2s}{s^4}.
\]
故最优点为
\[
  s=\frac32.
\]
于是
\[
  \alpha=\frac{2}{9},\qquad \beta=\frac23.
\]
换回 $a,b$：
\[
  a=\frac{3}{\sqrt2},\qquad b=\sqrt{\frac32}.
\]
此时面积最小，最小面积为
\[
  \pi ab=\frac{3\sqrt3}{2}\pi.
\]
\end{solution}
\begin{exercise}
\par
求微分方程 \(\displaystyle y''-y=2xe^x\) 的通解。
\end{exercise}

\begin{solution}
\par
先解齐次方程
\[
  y''-y=0.
\]
特征方程为
\[
  r^2-1=0,
\]
故
\[
  r=\pm1.
\]
齐次解为
\[
  y_h=C_1e^x+C_2e^{-x}.
\]

右端为 $2xe^x$，其中 $e^x$ 对应的指数 $1$ 是齐次方程的一重特征根，所以特解要多乘一个 $x$。设
\[
  y_p=e^x(Ax^2+Bx).
\]
令 $u=Ax^2+Bx$，则
\[
  y_p=e^x u.
\]
对算子 $D^2-1$ 有
\[
  (D^2-1)(e^x u)=e^x(u''+2u').
\]
计算
\[
  u'=2Ax+B,\qquad u''=2A.
\]
因此
\[
  u''+2u'=2A+4Ax+2B.
\]
要求
\[
  e^x(2A+4Ax+2B)=2xe^x,
\]
所以比较系数：
\[
  4A=2,\qquad 2A+2B=0.
\]
解得
\[
  A=\frac12,\qquad B=-\frac12.
\]
因此
\[
  y_p=\frac12e^x(x^2-x).
\]
通解为
\[
  y=C_1e^x+C_2e^{-x}+\frac12e^x(x^2-x).
\]
\end{solution}

\subsubsection*{四、证明题（本题 10 分）}
\begin{exercise}
\par
证明数项级数 \(\displaystyle \sum_{n=1}^{\infty}\sin(\pi\sqrt{n^2+1})\) 条件收敛。
\end{exercise}

\begin{solution}
\par
先提取通项的主部。由于
\[
  \sqrt{n^2+1}
  =
  n\sqrt{1+\frac1{n^2}}
  =
  n+\frac1{2n}+O\left(\frac1{n^3}\right),
\]
所以
\[
  \pi\sqrt{n^2+1}
  =
  n\pi+\frac{\pi}{2n}+O\left(\frac1{n^3}\right).
\]
利用
\[
  \sin(n\pi+\alpha)=(-1)^n\sin\alpha,
\]
得到
\[
  \sin(\pi\sqrt{n^2+1})
  =
  (-1)^n
  \sin\left(\frac{\pi}{2n}+O\left(\frac1{n^3}\right)\right).
\]
当 $n\to\infty$ 时，括号内趋于 $0$，故
\[
  \sin(\pi\sqrt{n^2+1})
  \sim
  (-1)^n\frac{\pi}{2n}.
\]

这说明原级数具有交错调和型主部，通项绝对值趋于 $0$ 且从某项后单调下降，因此原级数由交错级数判别法收敛。

再看绝对收敛性：
\[
  \left|\sin(\pi\sqrt{n^2+1})\right|
  \sim
  \frac{\pi}{2n}.
\]
于是绝对值级数与
\[
  \sum_{n=1}^{\infty}\frac1n
\]
同敛散，而调和级数发散，所以原级数不绝对收敛。综上，原级数条件收敛。
\end{solution}

\subsubsection*{五、应用题（本题 10 分）}
\begin{exercise}
\par
设某山峰可由曲面 \(\displaystyle z=5-x^2-2y^2\) 表示。位于点
\(\displaystyle \left(\frac12,-\frac12,\frac{17}{4}\right)\) 处的登山者发现其供氧面具漏气，必须返回。他应该沿哪个方向才能最快到达山底？
\end{exercise}

\begin{solution}
\par
山峰可看成高度函数
\[
  z=z(x,y)=5-x^2-2y^2.
\]
在水平面内，函数增长最快的方向是梯度方向
\[
  \nabla z=(z_x,z_y).
\]
下山要让高度下降最快，因此应取负梯度方向。先求梯度：
\[
  z_x=-2x,\qquad z_y=-4y,
\]
所以
\[
  -\nabla z=(2x,4y).
\]
登山者所在点的水平坐标为
\[
  \left(\frac12,-\frac12\right).
\]
代入得
\[
  -\nabla z\left(\frac12,-\frac12\right)
  =
  \left(1,-2\right).
\]
因此应沿水平面内方向
\[
  (1,-2)
\]
前进。若要求单位方向，则为
\[
  \frac{(1,-2)}{\sqrt{1^2+(-2)^2}}
  =
  \frac{(1,-2)}{\sqrt5}.
\]
\end{solution}

\section{2012级工科数学分析下 B卷}

\subsubsection*{一、单项选择题（每小题 3 分，共 15 分）}
\begin{exercise}
\par
初值问题 \(\displaystyle y''+4y'+4y=0,\qquad y(0)=1,\quad y'(0)=3\)
  的解为（\quad）。
  A. $y=(C_1+C_2x)e^{-2x}$；\quad B. $y=(1+5x)e^{-2x}$；\quad
  C. $y=(1+3x)e^{-2x}$；\quad D. $y=e^{-2x}+5e^{2x}$。
\end{exercise}

\begin{solution}
\par
这是二阶常系数齐次方程。特征方程为
\[
  r^2+4r+4=0,\qquad (r+2)^2=0,
\]
有二重根 $r=-2$，所以通解为
\[
  y=(C_1+C_2x)e^{-2x}.
\]
由 $y(0)=1$ 得 $C_1=1$。再求导：
\[
  y'=\bigl(C_2-2C_1-2C_2x\bigr)e^{-2x},
\]
代入 $x=0$ 得 $y'(0)=C_2-2C_1=3$，于是 $C_2=5$。故
\[
  y=(1+5x)e^{-2x}.
\]
选 B。
\end{solution}
\begin{exercise}
\par
对于二元函数 $z=f(x,y)$，下列说法正确的是（\quad）。
  A. 若函数 $z=f(x,y)$ 在 $(x_0,y_0)$ 连续，则函数在该点一定存在偏导数；
  B. 若函数 $z=f(x,y)$ 在 $(x_0,y_0)$ 存在偏导数，则函数在该点一定连续；
  C. 若函数 $z=f(x,y)$ 在 $(x_0,y_0)$ 的偏导数连续，则函数在该点一定可微；
  D. 若函数 $z=f(x,y)$ 在 $(x_0,y_0)$ 可微，则函数在该点的偏导数一定连续。
\end{exercise}

\begin{solution}
\par
逐项判断。连续只说明函数值随点变化连续，并不能保证沿坐标方向的导数存在，因此 A 错。偏导数存在只考察沿 $x$ 轴、$y$ 轴两个方向的变化，不能控制任意方向趋近，故也不能推出连续，B 错。若 $f_x,f_y$ 在点的某邻域内存在并且在该点连续，则全增量可写成
\[
  \Delta z=f_x(x_0,y_0)\Delta x+f_y(x_0,y_0)\Delta y+o(\rho),
  \qquad \rho=\sqrt{(\Delta x)^2+(\Delta y)^2},
\]
所以函数在该点可微，C 对。可微只能推出偏导数存在，不能反推偏导数连续，D 错。选 C。
\end{solution}
\begin{exercise}
\par
曲面 \(\displaystyle 3x^2+y^2+z^2=12\)
  上点 $M(-1,0,3)$ 处的切平面与平面 $z=0$ 的夹角是（\quad）。
  A. $\displaystyle \dfrac{\pi}{6}$；\quad B. $\displaystyle \dfrac{\pi}{4}$；\quad C. $\displaystyle \dfrac{\pi}{3}$；\quad D. $\displaystyle \dfrac{\pi}{2}$。
\end{exercise}

\begin{solution}
\par
设
\[
  F(x,y,z)=3x^2+y^2+z^2-12.
\]
曲面在一点处的法向量为 $\nabla F=(6x,2y,2z)$。在 $M(-1,0,3)$ 处，
\[
  \nabla F(M)=(-6,0,6).
\]
平面 $z=0$ 的法向量为 $\mathbf k=(0,0,1)$。两平面的夹角等于两法向量夹角的锐角，因此
\[
  \cos\theta
  =
  \frac{|\nabla F(M)\cdot \mathbf k|}{|\nabla F(M)|\,|\mathbf k|}
  =
  \frac{6}{\sqrt{36+36}}
  =
  \frac1{\sqrt2}.
\]
故 $\theta=\dfrac{\pi}{4}$。选 B。
\end{solution}
\begin{exercise}
\par
已知 $L$ 是球面 \(\displaystyle x^2+y^2+z^2=1\)
  与平面 $x+y+z=0$ 的交线，则第一类曲线积分 \(\displaystyle \int_Lx^2\,\mathrm{d}s\)
  （\quad）。
  A. $0$；\quad B. $\displaystyle \dfrac{2\pi}{3}$；\quad C. $\displaystyle \dfrac{\pi}{3}$；\quad D. $\displaystyle \dfrac{4\pi}{3}$。
\end{exercise}

\begin{solution}
\par
平面 $x+y+z=0$ 过球心，所以交线 $L$ 是单位球上的大圆，半径为 $1$，弧长为 $2\pi$。在这个圆上有
\[
  x^2+y^2+z^2=1.
\]
由于平面方程与 $x,y,z$ 对称，沿该大圆积分时
\[
  \int_L x^2\,\mathrm{d}s=\int_L y^2\,\mathrm{d}s=\int_L z^2\,\mathrm{d}s.
\]
三式相加得
\[
  3\int_L x^2\,\mathrm{d}s
  =
  \int_L (x^2+y^2+z^2)\,\mathrm{d}s
  =
  \int_L 1\,\mathrm{d}s
  =
  2\pi.
\]
故 $\displaystyle \int_Lx^2\,\mathrm{d}s=\dfrac{2\pi}{3}$。选 B。
\end{solution}
\begin{exercise}
\par
幂级数 \(\displaystyle \sum_{n=1}^{\infty}\frac{x^n}{n}\)
  的和函数是（\quad）。
  A. $s(x)=-\ln(1-x),\ (-1\le x<1)$；\quad
  B. $s(x)=-\ln(x-1),\ (-1<x<1)$；
  C. $s(x)=-\ln(1+x),\ (-1\le x<1)$；\quad
  D. $s(x)=\ln(1-x),\ (-1\le x<1)$。
\end{exercise}

\begin{solution}
\par
从几何级数
\[
  \frac1{1-x}=\sum_{n=0}^{\infty}x^n,\qquad |x|<1
\]
逐项积分，得到
\[
  -\ln(1-x)=\sum_{n=1}^{\infty}\frac{x^n}{n},\qquad |x|<1.
\]
端点要另判：当 $x=-1$ 时为 $\sum_{n=1}^{\infty}(-1)^n/n$，收敛；当 $x=1$ 时为调和级数 $\sum 1/n$，发散。因此和函数为
\[
  s(x)=-\ln(1-x),\qquad -1\le x<1.
\]
选 A。
\end{solution}

\subsubsection*{二、填空题（每小题 3 分，共 15 分）}
\begin{exercise}
\par
已知 \(\displaystyle f(x,y)=\sin x\ln(y+1)+\cos y\ln(1-x)\)，
  则 $\displaystyle \left.\dfrac{\partial f}{\partial x}\right|_{(0,0)}=$ \underline{\hspace{3cm}}。
\end{exercise}

\begin{solution}
\par
对 $x$ 求偏导时把 $y$ 看作常数：
\[
  \frac{\partial}{\partial x}\bigl(\sin x\ln(y+1)\bigr)=\cos x\ln(y+1),
\]
\[
  \frac{\partial}{\partial x}\bigl(\cos y\ln(1-x)\bigr)
  =
  \cos y\cdot\frac{-1}{1-x}.
\]
所以
\[
  f_x=\cos x\ln(y+1)-\frac{\cos y}{1-x}.
\]
代入 $(0,0)$，有 $\ln1=0,\ \cos0=1$，故
\[
  f_x(0,0)=0-1=-1.
\]
\end{solution}
\begin{exercise}
\par
微分方程 \(\displaystyle y^{(4)}-y=0\)
  的通解为 \underline{\hspace{4cm}}。
\end{exercise}

\begin{solution}
\par
这是常系数线性齐次方程。设 $y=e^{rx}$，得特征方程
\[
  r^4-1=0,\qquad (r^2-1)(r^2+1)=0.
\]
故特征根为
\[
  r=1,\quad r=-1,\quad r=i,\quad r=-i.
\]
实根 $1,-1$ 给出 $e^x,e^{-x}$，共轭虚根 $\pm i$ 给出 $\cos x,\sin x$。因此通解为
\[
  y=C_1e^x+C_2e^{-x}+C_3\cos x+C_4\sin x.
\]
\end{solution}
\begin{exercise}
\par
设 $D=[0,1;0,2]$，则二重积分
  \(\displaystyle \iint_D\sqrt{x+y}\,\mathrm{d}x\,\mathrm{d}y=\underline{\hspace{3cm}}\)。
\end{exercise}

\begin{solution}
\par
这里 $D=[0,1;0,2]$ 表示矩形区域 $0\le x\le1,\ 0\le y\le2$。先对 $x$ 积分：
\[
  \iint_D\sqrt{x+y}\,\mathrm{d}x\,\mathrm{d}y
  =
  \int_0^2\int_0^1 (x+y)^{1/2}\,\mathrm{d}x\,\mathrm{d}y.
\]
内层积分为
\[
  \int_0^1 (x+y)^{1/2}\,\mathrm{d}x
  =
  \frac23\bigl((y+1)^{3/2}-y^{3/2}\bigr).
\]
于是
\[
  I
  =
  \frac23\int_0^2\bigl((y+1)^{3/2}-y^{3/2}\bigr)\,\mathrm{d}y
  =
  \frac23\cdot\frac25
  \left[(y+1)^{5/2}-y^{5/2}\right]_{0}^{2}.
\]
代入上下限得
\[
  I
  =
  \frac4{15}\left(3^{5/2}-2^{5/2}-1\right).
\]
\end{solution}
\begin{exercise}
\par
设 $L$ 是任意一条光滑闭曲线，则
  \(\displaystyle \oint_L(2xy+\cos x)\,\mathrm{d}x+(x^2+\sin y)\,\mathrm{d}y=\underline{\hspace{3cm}}\)。
\end{exercise}

\begin{solution}
\par
记
\[
  P=2xy+\cos x,\qquad Q=x^2+\sin y.
\]
有
\[
  P_y=2x,\qquad Q_x=2x,
\]
因此 $P_y=Q_x$。由于 $P,Q$ 在整个平面上有连续偏导数，且平面单连通，所以该微分形式为全微分。也可以直接求势函数：
\[
  \Phi_x=P
  \quad\Longrightarrow\quad
  \Phi=x^2y+\sin x+h(y).
\]
再由 $\Phi_y=Q$ 得 $x^2+h'(y)=x^2+\sin y$，故 $h'(y)=\sin y$。所以闭曲线积分只与端点有关，而闭曲线起点终点相同，故
\[
  \oint_L(2xy+\cos x)\,\mathrm{d}x+(x^2+\sin y)\,\mathrm{d}y=0.
\]
\end{solution}
\begin{exercise}
\par
函数项级数 \(\displaystyle \sum_{n=1}^{\infty}\frac1n\left(\frac{x-1}{x}\right)^n\)
  的收敛域为 \underline{\hspace{4cm}}。
\end{exercise}

\begin{solution}
\par
令
\[
  t=\frac{x-1}{x}=1-\frac1x,\qquad x\ne0.
\]
原级数化为
\[
  \sum_{n=1}^{\infty}\frac{t^n}{n}.
\]
已知 $\sum t^n/n$ 在 $-1\le t<1$ 时收敛，在 $t=1$ 时发散。因此需要
\[
  -1\le 1-\frac1x<1.
\]
先由 $1-\dfrac1x<1$ 得 $-\dfrac1x<0$，故 $x>0$。在 $x>0$ 下，再由
\[
  -1\le1-\frac1x
  \quad\Longleftrightarrow\quad
  \frac1x\le2
  \quad\Longleftrightarrow\quad
  x\ge\frac12.
\]
所以收敛域为
\[
  \left[\frac12,+\infty\right).
\]
\end{solution}

\subsubsection*{三、计算题（每小题 10 分，共 50 分）}
\begin{exercise}
\par
求曲面 \(\displaystyle x^2+y^2+z^2=x\)
  的切平面，使它垂直于平面 $x-y-z=2$ 和 $\displaystyle x-y-\dfrac12z=2$。
\end{exercise}

\begin{solution}
\par
设所求切点为 $(x_0,y_0,z_0)$。曲面可写成
\[
  F(x,y,z)=x^2+y^2+z^2-x=0,
\]
所以切平面的法向量为
\[
  \nabla F(x_0,y_0,z_0)=(2x_0-1,\,2y_0,\,2z_0).
\]
给定两个平面的法向量分别为
\[
  \mathbf n_1=(1,-1,-1),\qquad
  \mathbf n_2=\left(1,-1,-\frac12\right).
\]
若所求切平面同时垂直于这两个平面，则它的法向量必须同时垂直于 $\mathbf n_1,\mathbf n_2$，故平行于
\[
  \mathbf n_1\times\mathbf n_2
  =
  \left(-\frac12,-\frac12,0\right),
\]
也就是平行于 $(1,1,0)$。于是
\[
  (2x_0-1,2y_0,2z_0)=\lambda(1,1,0),
\]
从而
\[
  z_0=0,\qquad 2x_0-1=2y_0,\qquad y_0=x_0-\frac12.
\]
代回曲面方程：
\[
  x_0^2+\left(x_0-\frac12\right)^2=x_0
  \quad\Longleftrightarrow\quad
  8x_0^2-8x_0+1=0.
\]
解得
\[
  x_0=\frac{2+\sqrt2}{4}\quad\text{或}\quad x_0=\frac{2-\sqrt2}{4},
\]
相应地
\[
  y_0=\frac{\sqrt2}{4}\quad\text{或}\quad y_0=-\frac{\sqrt2}{4},\qquad z_0=0.
\]
由于切平面法向量可取 $(1,1,0)$，切平面为
\[
  x+y=x_0+y_0.
\]
所以所求两个切平面为
\[
  x+y-\frac{1+\sqrt2}{2}=0,\qquad
  x+y-\frac{1-\sqrt2}{2}=0.
\]
\end{solution}
\begin{exercise}
\par
求微分方程 \(\displaystyle y''+(4x+e^{2y})(y')^3=0\)
  的通解。
\end{exercise}

\begin{solution}
\par
方程中没有显含自变量的导数项组合，且 $(y')^3$ 出现，适合把 $x$ 看成 $y$ 的函数。设
\[
  x=x(y),\qquad x'=\frac{\mathrm{d}x}{\mathrm{d}y}.
\]
当 $y'\ne0$ 时，有
\[
  y'=\frac1{x'},\qquad
  y''
  =
  \frac{\mathrm{d}}{\mathrm{d}x}\left(\frac1{x'}\right)
  =
  \frac{\mathrm{d}}{\mathrm{d}y}\left(\frac1{x'}\right)\frac{\mathrm{d}y}{\mathrm{d}x}
  =
  -\frac{x''}{(x')^3}.
\]
代入原方程：
\[
  -\frac{x''}{(x')^3}
  +(4x+e^{2y})\frac1{(x')^3}=0.
\]
两边乘以 $(x')^3$，得
\[
  x''-4x=e^{2y}.
\]
这是以 $y$ 为自变量的二阶常系数非齐次线性方程。齐次方程 $x''-4x=0$ 的通解为
\[
  x_h=C_1e^{2y}+C_2e^{-2y}.
\]
右端 $e^{2y}$ 与齐次解中的 $e^{2y}$ 重合，取特解 $x_p=Ay e^{2y}$，代入得
\[
  x_p''-4x_p=4A e^{2y}.
\]
令 $4A=1$，得 $A=\dfrac14$。因此
\[
  x=C_1e^{2y}+C_2e^{-2y}+\frac14\,y e^{2y}
\]
给出原方程的一般积分形式。还要注意，上面的变换默认 $y'\ne0$；若 $y'\equiv0$，则 $y\equiv C$ 也满足原方程，应作为常值解补上。
\end{solution}
\begin{exercise}
\par
计算三重积分 \(\displaystyle I=\iiint_D(x^2+y^2)\,\mathrm{d}x\,\mathrm{d}y\,\mathrm{d}z\)，
  其中 $D$ 是由圆锥面 $x^2+y^2=z^2$ 与平面 $z=1$ 所围成的区域。
\end{exercise}

\begin{solution}
\par
区域 $D$ 是上圆锥 $r=z$ 与平面 $z=1$ 围成的实心圆锥。用柱坐标
\[
  x=r\cos\theta,\qquad y=r\sin\theta,\qquad \mathrm{d}x\,\mathrm{d}y\,\mathrm{d}z=r\,\mathrm{d}r\,\mathrm{d}\theta\,\mathrm{d}z.
\]
在该区域内
\[
  0\le z\le1,\qquad 0\le r\le z,\qquad 0\le\theta\le2\pi,
\]
并且 $x^2+y^2=r^2$。于是
\[
  I
  =
  \int_0^{2\pi}\int_0^1\int_0^z r^2\cdot r\,\mathrm{d}r\,\mathrm{d}z\,\mathrm{d}\theta
  =
  2\pi\int_0^1\left[\frac{r^4}{4}\right]_{0}^{z}\,\mathrm{d}z.
\]
继续计算：
\[
  I
  =
  \frac{\pi}{2}\int_0^1 z^4\,\mathrm{d}z
  =
  \frac{\pi}{2}\cdot\frac15
  =
  \frac{\pi}{10}.
\]
\end{solution}
\begin{exercise}
\par
设 $\Sigma$ 为曲面 \(\displaystyle x^2+y^2+z^2=2ax\quad(a>0)\)，
  计算曲面积分 \(\displaystyle I=\iint_\Sigma(x^2+y^2+z^2)\,\mathrm{d}S\)。
\end{exercise}

\begin{solution}
\par
将曲面方程配方：
\[
  x^2-2ax+y^2+z^2=0
  \quad\Longleftrightarrow\quad
  (x-a)^2+y^2+z^2=a^2.
\]
所以 $\Sigma$ 是以 $(a,0,0)$ 为球心、半径为 $a$ 的球面。在曲面上，
\[
  x^2+y^2+z^2=2ax,
\]
因此
\[
  I=\iint_\Sigma (x^2+y^2+z^2)\,\mathrm{d}S
  =
  2a\iint_\Sigma x\,\mathrm{d}S.
\]
设 $X=x-a$，则球面对球心对称，有
\[
  \iint_\Sigma X\,\mathrm{d}S=0.
\]
于是
\[
  \iint_\Sigma x\,\mathrm{d}S
  =
  \iint_\Sigma (a+X)\,\mathrm{d}S
  =
  a\cdot 4\pi a^2.
\]
故
\[
  I=2a\cdot a\cdot4\pi a^2=8\pi a^4.
\]
\end{solution}
\begin{exercise}
\par
将函数 \(\displaystyle f(x)=\frac1{x^2+4x+3}\)
  在 $x=1$ 处展成幂级数。
\end{exercise}

\begin{solution}
\par
先分解为部分分式：
\[
  \frac1{x^2+4x+3}
  =
  \frac1{(x+1)(x+3)}
  =
  \frac12\left(\frac1{x+1}-\frac1{x+3}\right).
\]
在 $x=1$ 处展开，令 $t=x-1$，则
\[
  x+1=t+2,\qquad x+3=t+4.
\]
于是
\[
  f(x)
  =
  \frac12\left(\frac1{t+2}-\frac1{t+4}\right).
\]
分别化成几何级数形式：
\[
  \frac1{t+2}
  =
  \frac12\cdot\frac1{1+t/2}
  =
  \sum_{n=0}^{\infty}\frac{(-1)^n t^n}{2^{n+1}},
  \qquad |t|<2,
\]
\[
  \frac1{t+4}
  =
  \frac14\cdot\frac1{1+t/4}
  =
  \sum_{n=0}^{\infty}\frac{(-1)^n t^n}{4^{n+1}},
  \qquad |t|<4.
\]
两式相减后再乘 $\dfrac12$，得
\[
  f(x)
  =
  \sum_{n=0}^{\infty}
  (-1)^n
  \left(\frac1{2^{n+2}}-\frac1{2^{2n+3}}\right)
  (x-1)^n.
\]
共同收敛条件取较小者，即
\[
  |x-1|<2.
\]
\end{solution}

\subsubsection*{四、证明题（本题 10 分）}
\begin{exercise}
\par
设 \(\displaystyle x=x(y,z),\qquad y=y(x,z),\qquad z=z(x,y)\)
都是由方程 \(\displaystyle F(x,y,z)=0\)
所确定的具有连续偏导数的函数，证明
\(\displaystyle \frac{\partial x}{\partial y}\frac{\partial y}{\partial z}\frac{\partial z}{\partial x}=-1\)。
\end{exercise}

\begin{solution}
\par
把同一个方程 $F(x,y,z)=0$ 分别看成确定 $x=x(y,z)$、$y=y(x,z)$、$z=z(x,y)$ 的隐函数。为使下列商有意义，在相应点需有 $F_xF_yF_z\ne0$。

先把 $x$ 看作 $y,z$ 的函数，对 $y$ 求偏导而保持 $z$ 不变：
\[
  F_x\frac{\partial x}{\partial y}+F_y=0,
  \qquad
  \frac{\partial x}{\partial y}=-\frac{F_y}{F_x}.
\]
再把 $y$ 看作 $x,z$ 的函数，对 $z$ 求偏导而保持 $x$ 不变：
\[
  F_y\frac{\partial y}{\partial z}+F_z=0,
  \qquad
  \frac{\partial y}{\partial z}=-\frac{F_z}{F_y}.
\]
最后把 $z$ 看作 $x,y$ 的函数，对 $x$ 求偏导而保持 $y$ 不变：
\[
  F_z\frac{\partial z}{\partial x}+F_x=0,
  \qquad
  \frac{\partial z}{\partial x}=-\frac{F_x}{F_z}.
\]
三式相乘：
\[
  \frac{\partial x}{\partial y}
  \frac{\partial y}{\partial z}
  \frac{\partial z}{\partial x}
  =
  \left(-\frac{F_y}{F_x}\right)
  \left(-\frac{F_z}{F_y}\right)
  \left(-\frac{F_x}{F_z}\right)
  =
  -1.
\]
命题得证。
\end{solution}

\subsubsection*{五、应用题（本题 10 分）}
\begin{exercise}
\par
设有平面力场
\[
\mathbf F(x,y)=(2xy^3-y^2\cos x)\mathbf i+(1-2y\sin x+3x^2y^2)\mathbf j,
\]
求一质点沿曲线 \(\displaystyle L:2x=\pi y^2\)
从点 $O(0,0)$ 运动到点 $\displaystyle A\!\left(\dfrac\pi2,1\right)$ 时变力 $\mathbf F$ 所作的功。
\end{exercise}

\begin{solution}
\par
记
\[
  P(x,y)=2xy^3-y^2\cos x,\qquad
  Q(x,y)=1-2y\sin x+3x^2y^2.
\]
先判断是否与路径有关：
\[
  P_y=6xy^2-2y\cos x,\qquad
  Q_x=-2y\cos x+6xy^2.
\]
两者相等，且 $P,Q$ 在全平面上有连续偏导数，因此该力场为保守场，沿曲线所作功只与起点和终点有关。

求势函数 $\Phi$。由 $\Phi_x=P$，对 $x$ 积分：
\[
  \Phi(x,y)
  =
  \int (2xy^3-y^2\cos x)\,\mathrm{d}x
  =
  x^2y^3-y^2\sin x+h(y).
\]
再对 $y$ 求偏导：
\[
  \Phi_y=3x^2y^2-2y\sin x+h'(y).
\]
令 $\Phi_y=Q=1-2y\sin x+3x^2y^2$，得 $h'(y)=1$，故可取 $h(y)=y$。于是
\[
  \Phi(x,y)=x^2y^3-y^2\sin x+y.
\]
从 $O(0,0)$ 到 $A\!\left(\dfrac\pi2,1\right)$ 的功为
\[
  W
  =
  \Phi\!\left(\frac\pi2,1\right)-\Phi(0,0)
  =
  \left(\frac{\pi^2}{4}-1+1\right)-0
  =
  \frac{\pi^2}{4}.
\]
\end{solution}

\section{2013级工科数学分析下 A卷}

\par\medskip
\noindent\textbf{一、单项选择题（每小题3分，共15分）}\par\nobreak\smallskip
\begin{exercise}
\par
若 $y_1,y_2$ 是非齐次线性方程 $y''+ay'+by=f(x)$ 的两个特解，则下面结论中正确的是（\quad）。
  \par\noindent\begin{tabular}{@{}>{\raggedright\arraybackslash}p{0.46\linewidth}>{\raggedright\arraybackslash}p{0.46\linewidth}@{}}
  A. $y_1+y_2$ 是非齐次线性方程的解。 & B. $y_1-y_2$ 是非齐次线性方程的解。\\
  C. $y_1+y_2$ 是 $y''+ay'+by=0$ 的解。 & D. $y_1-y_2$ 是 $y''+ay'+by=0$ 的解。
  \end{tabular}
\end{exercise}

\begin{solution}
\par
记线性微分算子
\[
  L[y]=y''+ay'+by.
\]
因为 $y_1,y_2$ 都是非齐次方程的特解，所以
\[
  L[y_1]=f(x),\qquad L[y_2]=f(x).
\]
利用线性性，
\[
  L[y_1-y_2]=L[y_1]-L[y_2]=f(x)-f(x)=0.
\]
因此 $y_1-y_2$ 是对应齐次方程 $y''+ay'+by=0$ 的解。一般地，$y_1+y_2$ 满足 $L[y_1+y_2]=2f(x)$，不满足原非齐次方程；$y_1-y_2$ 也不是非齐次方程的解。选 D。
\end{solution}
\begin{exercise}
\par
设 $z=f(x,y)$ 可微，且 $\displaystyle f(x,x^2)=\frac12,\ \left.f_x(x,y)\right|_{y=x^2}=\frac1x$，则 $\left.f_y(x,y)\right|_{y=x^2}$（\quad）。
  \par\noindent\begin{tabular}{@{}>{\raggedright\arraybackslash}p{0.22\linewidth}>{\raggedright\arraybackslash}p{0.22\linewidth}>{\raggedright\arraybackslash}p{0.22\linewidth}>{\raggedright\arraybackslash}p{0.22\linewidth}@{}}
  A. $\displaystyle -\dfrac{1}{x^2}$ & B. $\displaystyle \dfrac{1}{x^2}$ & C. $\displaystyle \dfrac{1}{2x^2}$ & D. $\displaystyle -\dfrac{1}{2x^2}$
  \end{tabular}
\end{exercise}

\begin{solution}
\par
把
\[
  f(x,x^2)=\frac12
\]
看成关于 $x$ 的恒等式。由于右端为常数，对 $x$ 求导得
\[
  f_x(x,x^2)\cdot 1+f_y(x,x^2)\cdot 2x=0.
\]
题设给出 $\left.f_x(x,y)\right|_{y=x^2}=\dfrac1x$，代入上式：
\[
  \frac1x+2x\,f_y(x,x^2)=0.
\]
所以
\[
  f_y(x,x^2)=-\frac1{2x^2}.
\]
选 D。
\end{solution}
\begin{exercise}
\par
设 $C:x^2+y^2=a^2$ 取逆时针方向，则 \(\displaystyle \oint_C \frac{(x+y)\,\mathrm{d}x-(x-y)\,\mathrm{d}y}{x^2+y^2}\)（\quad）。
  \par\noindent\begin{tabular}{@{}>{\raggedright\arraybackslash}p{0.22\linewidth}>{\raggedright\arraybackslash}p{0.22\linewidth}>{\raggedright\arraybackslash}p{0.22\linewidth}>{\raggedright\arraybackslash}p{0.22\linewidth}@{}}
  A. $-2\pi$ & B. $2\pi$ & C. $0$ & D. $-\pi$
  \end{tabular}
\end{exercise}

\begin{solution}
\par
将分子拆开：
\[
  (x+y)\,\mathrm{d}x-(x-y)\,\mathrm{d}y
  =
  x\,\mathrm{d}x+y\,\mathrm{d}y+y\,\mathrm{d}x-x\,\mathrm{d}y.
\]
因此
\[
  \frac{(x+y)\,\mathrm{d}x-(x-y)\,\mathrm{d}y}{x^2+y^2}
  =
  \frac{x\,\mathrm{d}x+y\,\mathrm{d}y}{x^2+y^2}
  +
  \frac{y\,\mathrm{d}x-x\,\mathrm{d}y}{x^2+y^2}.
\]
第一项为 $\mathrm{d}\ln r$，其中 $r=\sqrt{x^2+y^2}$；第二项为 $-\mathrm{d}\theta$，因为
\[
  \mathrm{d}\theta=\frac{x\,\mathrm{d}y-y\,\mathrm{d}x}{x^2+y^2}.
\]
在圆 $C:x^2+y^2=a^2$ 上，$r=a$ 为常数，所以 $\oint_C \mathrm{d}\ln r=0$。曲线逆时针绕原点一周，$\theta$ 增加 $2\pi$，故
\[
  \oint_C \frac{(x+y)\,\mathrm{d}x-(x-y)\,\mathrm{d}y}{x^2+y^2}
  =
  -\oint_C \mathrm{d}\theta
  =
  -2\pi.
\]
选 A。
\end{solution}
\begin{exercise}
\par
幂级数 $\displaystyle\sum_{n=1}^{\infty}\frac{(-1)^{n-1}}{n}(x-1)^n$ 的收敛域是（\quad）。
  \par\noindent\begin{tabular}{@{}>{\raggedright\arraybackslash}p{0.22\linewidth}>{\raggedright\arraybackslash}p{0.22\linewidth}>{\raggedright\arraybackslash}p{0.22\linewidth}>{\raggedright\arraybackslash}p{0.22\linewidth}@{}}
  A. $(0,2)$ & B. $(0,2]$ & C. $[0,2)$ & D. $[0,2]$
  \end{tabular}
\end{exercise}

\begin{solution}
\par
令 $t=x-1$，原级数变为
\[
  \sum_{n=1}^{\infty}\frac{(-1)^{n-1}}{n}t^n.
\]
这是 $\ln(1+t)$ 的幂级数展开。其半径为 $1$，先得 $|t|<1$。端点另判：
\[
  t=1:\quad \sum_{n=1}^{\infty}\frac{(-1)^{n-1}}{n}
  \text{ 收敛},
\]
\[
  t=-1:\quad -\sum_{n=1}^{\infty}\frac1n
  \text{ 发散}.
\]
所以 $-1<t\le1$。换回 $t=x-1$，得到
\[
  0<x\le2.
\]
选 B。
\end{solution}
\begin{exercise}
\par
设 \(f(x)=x^2\ (0\le x<1)\)，而
  \[
  S(x)=\sum_{n=1}^{\infty}b_n\sin n\pi x\quad(-\infty<x<+\infty),
  \qquad
  b_n=2\int_0^1 f(x)\sin n\pi x\,\mathrm{d}x .
  \]
  则 $\displaystyle S\!\left(-\dfrac12\right)$（\quad）。
  \par\noindent\begin{tabular}{@{}>{\raggedright\arraybackslash}p{0.22\linewidth}>{\raggedright\arraybackslash}p{0.22\linewidth}>{\raggedright\arraybackslash}p{0.22\linewidth}>{\raggedright\arraybackslash}p{0.22\linewidth}@{}}
  A. $\displaystyle -\dfrac12$ & B. $\displaystyle \dfrac14$ & C. $\displaystyle -\dfrac14$ & D. $\displaystyle \dfrac12$
  \end{tabular}
\end{exercise}

\begin{solution}
\par
给出的系数
\[
  b_n=2\int_0^1 f(x)\sin n\pi x\,\mathrm{d}x
\]
是 $[0,1]$ 上的 Fourier 正弦级数系数。正弦级数对应把 $f(x)=x^2$ 作奇延拓，并以周期 $2$ 延拓到全直线。于是当 $x=-\dfrac12$ 时，奇延拓函数值为
\[
  S\!\left(-\frac12\right)
  =
  -f\!\left(\frac12\right)
  =
  -\left(\frac12\right)^2
  =
  -\frac14.
\]
这里 \(x=-\frac12\) 先由奇延拓对应到区间内的 \(\frac12\)，再加负号；因为该点不是延拓后的跳跃点，所以不需要再取左右极限平均。
这里 \(x=-\frac12\) 先由奇延拓对应到区间内的 \(\frac12\)，再加负号；因为该点不是延拓后的跳跃点，所以不需要再取左右极限平均。
该点不是跳跃点，因此级数收敛到上述函数值。选 C。
\end{solution}

\par\medskip
\noindent\textbf{二、填空题（每小题3分，共15分）}\par\nobreak\smallskip
\begin{exercise}
\par
微分方程 $y^{(4)}-y=0$ 的通解为 \underline{\hspace{5cm}}。
\end{exercise}

\begin{solution}
\par
设 $y=e^{rx}$，得到特征方程
\[
  r^4-1=0
  \quad\Longleftrightarrow\quad
  (r^2-1)(r^2+1)=0.
\]
因此特征根为 $1,-1,i,-i$。两个实根给出 $e^x,e^{-x}$，一对共轭虚根给出 $\cos x,\sin x$，故通解为
\[
  y=C_1e^x+C_2e^{-x}+C_3\cos x+C_4\sin x.
\]
四阶齐次方程的通解应含四个线性无关解，这里两个实根和一对共轭虚根正好给出四个基本解。
\end{solution}
\begin{exercise}
\par
函数 $z=2x^2+y^2$ 在点 $(1,1)$ 处沿该点梯度方向的方向导数为 \underline{\hspace{5cm}}。
\end{exercise}

\begin{solution}
\par
先求梯度：
\[
  \nabla z=(z_x,z_y)=(4x,2y).
\]
在点 $(1,1)$ 处，
\[
  \nabla z(1,1)=(4,2),\qquad |\nabla z(1,1)|=\sqrt{4^2+2^2}=2\sqrt5.
\]
沿梯度方向的单位向量为 $\dfrac{\nabla z}{|\nabla z|}$，方向导数等于梯度在该方向上的投影：
\[
  D_{\nabla z/|\nabla z|}z
  =
  \nabla z\cdot \frac{\nabla z}{|\nabla z|}
  =
  |\nabla z|
  =
  2\sqrt5.
\]
\end{solution}
\begin{exercise}
\par
交换积分的顺序，则 $\displaystyle\int_1^e \mathrm{d}x\int_0^{\ln x} f(x,y)\,\mathrm{d}y=$ \underline{\hspace{5cm}}。
\end{exercise}

\begin{solution}
\par
原积分区域由
\[
  1\le x\le e,\qquad 0\le y\le \ln x
\]
描述。将不等式 $y\le\ln x$ 改写为 $x\ge e^y$。当 $x$ 从 $1$ 到 $e$ 时，$y$ 的范围为 $0\le y\le1$。固定 $y$ 后，$x$ 从曲线 $x=e^y$ 到右边界 $x=e$。因此
\[
  \int_1^e \mathrm{d}x\int_0^{\ln x} f(x,y)\,\mathrm{d}y
  =
  \int_0^1\mathrm{d}y\int_{e^y}^{e}f(x,y)\,\mathrm{d}x.
\]
\end{solution}
\begin{exercise}
\par
设 $\Sigma:x^2+y^2+z^2=a^2$，则 $\displaystyle\iint_{\Sigma}(x^2+y^2+z^2)\,\mathrm{d}S=$ \underline{\hspace{5cm}}。
\end{exercise}

\begin{solution}
\par
在球面 $\Sigma:x^2+y^2+z^2=a^2$ 上，被积函数恒等于 $a^2$。因此
\[
  \iint_{\Sigma}(x^2+y^2+z^2)\,\mathrm{d}S
  =
  \iint_{\Sigma}a^2\,\mathrm{d}S
  =
  a^2\cdot \operatorname{Area}(\Sigma).
\]
半径为 $a$ 的球面面积为 $4\pi a^2$，所以结果为
\[
  4\pi a^4.
\]
这是第一类曲面积分，积分元是 \(\mathrm{d}S\)，没有方向问题；只需把球面上的被积函数常值乘以球面面积。
\end{solution}
\begin{exercise}
\par
设幂级数 \(\displaystyle\sum_{n=1}^{\infty}a_nx^n\) 的收敛半径为 \(3\)。则
  \[
    \sum_{n=1}^{\infty}n a_n(x-1)^{n+1}
  \]
  的收敛区间为 \underline{\hspace{5cm}}。
\end{exercise}

\begin{solution}
\par
令 $t=x-1$，则待求级数为
\[
  \sum_{n=1}^{\infty}n a_n t^{n+1}
  =
  t\sum_{n=1}^{\infty}n a_n t^n.
\]
已知 $\sum a_nx^n$ 的收敛半径为 $3$。对幂级数逐项求导或给系数乘以 $n$ 都不改变收敛半径，所以 $\sum n a_n t^n$ 的收敛半径仍为 $3$。再乘上一个 $t$ 也不改变半径。故
\[
  |t|<3
  \quad\Longleftrightarrow\quad
  |x-1|<3.
\]
因此可确定的收敛区间为
\[
  (-2,4).
\]
端点是否收敛取决于具体的 $a_n$，题设信息不足，不能加入端点。
\end{solution}

\par\medskip
\noindent\textbf{三、计算题（每小题10分，共50分）}\par\nobreak\smallskip
\begin{exercise}
\par
设函数 \(f(u)\) 具有二阶连续导数，而 \(z=f(e^x\sin y)\) 满足方程 \(\displaystyle \frac{\partial^2 z}{\partial x^2}+\frac{\partial^2 z}{\partial y^2}=e^{2x}z\)，求 \(f(u)\)。
\end{exercise}

\begin{solution}
\par
令
\[
  u=e^x\sin y,\qquad z=f(u).
\]
先求关于 $x$ 的二阶偏导。因为 $u_x=u$，所以
\[
  z_x=f'(u)u,\qquad
  z_{xx}=f''(u)u^2+f'(u)u.
\]
再求关于 $y$ 的二阶偏导。由
\[
  u_y=e^x\cos y,\qquad u_{yy}=-e^x\sin y=-u,
\]
得
\[
  z_y=f'(u)e^x\cos y,
\]
\[
  z_{yy}=f''(u)e^{2x}\cos^2y-f'(u)u.
\]
两式相加，$f'(u)u$ 与 $-f'(u)u$ 抵消：
\[
  z_{xx}+z_{yy}
  =
  f''(u)\bigl(u^2+e^{2x}\cos^2y\bigr).
\]
而
\[
  u^2+e^{2x}\cos^2y
  =
  e^{2x}\sin^2y+e^{2x}\cos^2y
  =
  e^{2x}.
\]
故方程化为
\[
  e^{2x}f''(u)=e^{2x}f(u).
\]
由于 $e^{2x}\ne0$，得到常微分方程
\[
  f''(u)=f(u).
\]
解得
\[
  f(u)=C_1e^u+C_2e^{-u}.
\]
\end{solution}
\begin{exercise}
\par
计算三重积分 $I=\displaystyle\iiint_{\Omega}(x+z)\,\mathrm{d}V$，其中 $\Omega$ 是由曲面 $z=x^2+y^2$ 与曲面 $z=1-x^2-y^2$ 所围成的区域。
\end{exercise}

\begin{solution}
\par
区域由下方抛物面 $z=x^2+y^2$ 与上方抛物面 $z=1-x^2-y^2$ 围成。两曲面交线满足
\[
  r^2=1-r^2,\qquad r=\frac1{\sqrt2}.
\]
区域关于平面 $x=0$ 对称，而被积函数中的 $x$ 关于 $x$ 为奇函数，所以
\[
  \iiint_{\Omega}x\,\mathrm{d}V=0.
\]
因此只需计算 $\iiint_{\Omega}z\,\mathrm{d}V$。用柱坐标，有
\[
  0\le \theta\le2\pi,\qquad
  0\le r\le\frac1{\sqrt2},\qquad
  r^2\le z\le1-r^2.
\]
于是
\[
  I
  =
  \int_0^{2\pi}\int_0^{1/\sqrt2}\int_{r^2}^{1-r^2} z\,r\,\mathrm{d}z\,\mathrm{d}r\,\mathrm{d}\theta.
\]
先对 $z$ 积分：
\[
  I
  =
  2\pi\int_0^{1/\sqrt2}
  \frac12\bigl((1-r^2)^2-r^4\bigr)r\,\mathrm{d}r.
\]
化简被积函数：
\[
  (1-r^2)^2-r^4=1-2r^2.
\]
所以
\[
  I
  =
  \pi\int_0^{1/\sqrt2}(r-2r^3)\,\mathrm{d}r
  =
  \pi\left[\frac{r^2}{2}-\frac{r^4}{2}\right]_{0}^{1/\sqrt2}
  =
  \frac{\pi}{8}.
\]
\end{solution}
\begin{exercise}
\par
设 $f(x)$ 在 $(-\infty,+\infty)$ 内有连续导数，$L$ 是从点 $\displaystyle A(3,\frac23)$ 到点 $B(1,2)$ 的直线段，计算曲线积分 \(\displaystyle I=\int_L \frac{1+y^2f(xy)}{y}\,\mathrm{d}x+\frac{x}{y^2}\bigl[y^2f(xy)-1\bigr]\,\mathrm{d}y\)。
\end{exercise}

\begin{solution}
\par
取一个原函数 $F$，使
\[
  F'(t)=f(t).
\]
考察函数
\[
  \Phi(x,y)=F(xy)+\frac{x}{y}.
\]
它的全微分为
\[
  d\Phi
  =
  f(xy)(y\,\mathrm{d}x+x\,\mathrm{d}y)+\frac1y\,\mathrm{d}x-\frac{x}{y^2}\,\mathrm{d}y.
\]
整理 $\mathrm{d}x,\mathrm{d}y$ 的系数：
\[
  d\Phi
  =
  \left(yf(xy)+\frac1y\right)\mathrm{d}x
  +
  \left(xf(xy)-\frac{x}{y^2}\right)\mathrm{d}y
\]
\[
  =
  \frac{1+y^2f(xy)}{y}\,\mathrm{d}x
  +
  \frac{x}{y^2}\bigl[y^2f(xy)-1\bigr]\,\mathrm{d}y.
\]
这正是题中曲线积分的微分形式，所以积分与路径无关，只需代入端点：
\[
  I=\Phi(1,2)-\Phi\!\left(3,\frac23\right).
\]
两端点都有 $xy=2$，因此 $F(xy)$ 部分相消。于是
\[
  I
  =
  \left(F(2)+\frac12\right)
  -
  \left(F(2)+\frac{3}{2/3}\right)
  =
  \frac12-\frac92
  =
  -4.
\]
\end{solution}
\begin{exercise}
\par
设 \(\Sigma\) 是锥面 \(z=\sqrt{x^2+y^2}\) 被平面 \(z=0\) 及 \(z=1\) 所截部分的下侧，计算第二类曲面积分
  \(\displaystyle I=\iint_{\Sigma}x\,\mathrm{d}y\,\mathrm{d}z+y\,\mathrm{d}z\,\mathrm{d}x+(z^2-2z)\,\mathrm{d}x\,\mathrm{d}y\)。
\end{exercise}

\begin{solution}
\par
将第二类曲面积分写成通量形式，其中
\[
  P=x,\qquad Q=y,\qquad R=z^2-2z.
\]
锥面为 $z=\sqrt{x^2+y^2}$，可用参数
\[
  \mathbf r(\theta,z)=(z\cos\theta,z\sin\theta,z),
  \qquad 0\le z\le1,\quad 0\le\theta\le2\pi.
\]
计算
\[
  \mathbf r_\theta=(-z\sin\theta,z\cos\theta,0),
  \qquad
  \mathbf r_z=(\cos\theta,\sin\theta,1).
\]
下侧要求法向量的 $z$ 分量为负，故取
\[
  \mathbf r_\theta\times\mathbf r_z
  =
  (z\cos\theta,z\sin\theta,-z).
\]
在曲面上
\[
  (P,Q,R)=(z\cos\theta,z\sin\theta,z^2-2z).
\]
所以被积表达式为
\[
  (P,Q,R)\cdot(\mathbf r_\theta\times\mathbf r_z)
  =
  z^2\cos^2\theta+z^2\sin^2\theta-(z^2-2z)z
  =
  3z^2-z^3.
\]
于是
\[
  I
  =
  \int_0^{2\pi}\int_0^1(3z^2-z^3)\,\mathrm{d}z\,\mathrm{d}\theta
  =
  2\pi\left[z^3-\frac{z^4}{4}\right]_0^1
  =
  \frac{3\pi}{2}.
\]
\end{solution}
\begin{exercise}
\par
求幂级数 $\displaystyle\sum_{n=1}^{\infty}\frac{2n+1}{n!}x^{2n}$ 的和函数。
\end{exercise}

\begin{solution}
\par
令 $t=x^2$，则
\[
  S(x)=\sum_{n=1}^{\infty}\frac{2n+1}{n!}t^n
  =
  2\sum_{n=1}^{\infty}\frac{n t^n}{n!}
  +
  \sum_{n=1}^{\infty}\frac{t^n}{n!}.
\]
第二个级数直接由指数级数得
\[
  \sum_{n=1}^{\infty}\frac{t^n}{n!}=e^t-1.
\]
第一个级数中
\[
  \frac{n}{n!}=\frac1{(n-1)!},
\]
所以
\[
  \sum_{n=1}^{\infty}\frac{n t^n}{n!}
  =
  t\sum_{n=1}^{\infty}\frac{t^{n-1}}{(n-1)!}
  =
  te^t.
\]
因此
\[
  S(x)=2te^t+e^t-1=(2t+1)e^t-1.
\]
换回 $t=x^2$，得
\[
  S(x)=(2x^2+1)e^{x^2}-1.
\]
\end{solution}

\par\medskip
\noindent\textbf{四、证明题（本题10分）}\par\nobreak\smallskip
\begin{exercise}
\par
证明级数 $\displaystyle\sum_{n=1}^{\infty}\frac{(-1)^n(1-e^{-nx})}{n^2+x^2}$ 在 $[0,+\infty)$ 上一致收敛。
\end{exercise}

\begin{solution}
\par
对任意 $x\in[0,+\infty)$，有
\[
  0\le e^{-nx}\le1,
  \qquad
  0\le1-e^{-nx}\le1.
\]
同时
\[
  n^2+x^2\ge n^2.
\]
因此每一项满足一致估计
\[
  \left|
  \frac{(-1)^n(1-e^{-nx})}{n^2+x^2}
  \right|
  =
  \frac{1-e^{-nx}}{n^2+x^2}
  \le
  \frac1{n^2}.
\]
而数项级数
\[
  \sum_{n=1}^{\infty}\frac1{n^2}
\]
收敛。由 Weierstrass 判别法，原函数项级数在 $[0,+\infty)$ 上一致收敛。
\end{solution}

\par\medskip
\noindent\textbf{五、应用题（本题10分）}\par\nobreak\smallskip
\begin{exercise}
\par
在第一象限内作椭球面 $\displaystyle \dfrac{x^2}{a^2}+\dfrac{y^2}{b^2}+\dfrac{z^2}{c^2}=1$ 的切平面，使切平面与三个坐标面所围成的四面体的体积最小，求切点的坐标。
\end{exercise}

\begin{solution}
\par
设切点为
\[
  P(x_0,y_0,z_0),\qquad x_0>0,\ y_0>0,\ z_0>0.
\]
它在椭球面上，所以
\[
  \frac{x_0^2}{a^2}+\frac{y_0^2}{b^2}+\frac{z_0^2}{c^2}=1.
\]
椭球面
\[
  \frac{x^2}{a^2}+\frac{y^2}{b^2}+\frac{z^2}{c^2}=1
\]
在 $P$ 处的切平面为
\[
  \frac{x_0x}{a^2}+\frac{y_0y}{b^2}+\frac{z_0z}{c^2}=1.
\]
写成截距式：
\[
  \frac{x}{a^2/x_0}+\frac{y}{b^2/y_0}+\frac{z}{c^2/z_0}=1.
\]
因此它与三个坐标面围成的四面体三个截距分别为
\[
  \frac{a^2}{x_0},\qquad \frac{b^2}{y_0},\qquad \frac{c^2}{z_0}.
\]
四面体体积为
\[
  V
  =
  \frac16\cdot\frac{a^2}{x_0}\cdot\frac{b^2}{y_0}\cdot\frac{c^2}{z_0}
  =
  \frac{a^2b^2c^2}{6x_0y_0z_0}.
\]
要使 $V$ 最小，等价于在约束条件下使 $x_0y_0z_0$ 最大。令
\[
  X=\frac{x_0}{a},\qquad Y=\frac{y_0}{b},\qquad Z=\frac{z_0}{c}.
\]
约束变为
\[
  X^2+Y^2+Z^2=1,\qquad X,Y,Z>0,
\]
而 $x_0y_0z_0=abc\,XYZ$，所以只需最大化 $XYZ$。用拉格朗日乘数法，设
\[
  \mathcal L=\ln X+\ln Y+\ln Z-\lambda(X^2+Y^2+Z^2-1).
\]
驻点方程为
\[
  \frac1X=2\lambda X,\qquad
  \frac1Y=2\lambda Y,\qquad
  \frac1Z=2\lambda Z.
\]
因此
\[
  X^2=Y^2=Z^2.
\]
在第一象限 $X,Y,Z>0$，故 $X=Y=Z$。再由约束得
\[
  3X^2=1,\qquad X=Y=Z=\frac1{\sqrt3}.
\]
于是
\[
  x_0=\frac{a}{\sqrt3},\qquad
  y_0=\frac{b}{\sqrt3},\qquad
  z_0=\frac{c}{\sqrt3}.
\]
由于边界上至少有一个变量趋于 $0$ 时体积趋于无穷大，而第一象限内驻点唯一，所以该点给出最小体积。所求切点为
\[
  \left(\frac{a}{\sqrt3},\frac{b}{\sqrt3},\frac{c}{\sqrt3}\right).
\]
\end{solution}

\section{2013级工科数学分析下 B卷}

\par\medskip
\noindent\textbf{一、单项选择题（每小题3分，共15分）}\par\nobreak\smallskip
\begin{exercise}
\par
方程 \((y-\ln x)\,\mathrm{d}x+x\,\mathrm{d}y=0\) 是（\quad）。
  \par\noindent\begin{tabular}{@{}>{\raggedright\arraybackslash}p{0.22\linewidth}>{\raggedright\arraybackslash}p{0.22\linewidth}>{\raggedright\arraybackslash}p{0.22\linewidth}>{\raggedright\arraybackslash}p{0.22\linewidth}@{}}
  A. 可分离变量微分方程。 & B. 一阶线性非齐次方程。 & C. 一阶线性齐次方程。 & D. 非线性方程。
  \end{tabular}
\end{exercise}

\begin{solution}
\par
把原方程除以 $\mathrm{d}x$，得
\[
  y-\ln x+x\frac{\mathrm{d}y}{\mathrm{d}x}=0.
\]
整理为
\[
  y'+\frac1x y=\frac{\ln x}{x}.
\]
这是形如 $y'+P(x)y=Q(x)$ 的一阶线性方程，并且右端 $Q(x)=\dfrac{\ln x}{x}$ 不恒为零，所以是一阶线性非齐次方程。这里默认在 \(x>0\) 的区间上讨论，保证 \(\ln x\) 有意义；未知函数 \(y\) 与 \(y'\) 都只一次出现，因此不属于非线性方程。选 B。
\end{solution}
\begin{exercise}
\par
设 \(f(x,y)\) 是可微函数，且 \(f(0,0)=0\)，\(f_x(0,0)=a\)，\(f_y(0,0)=b\)，令
  \(\varphi(t)=f(t,f(t,t))\)，则 \(\varphi'(0)=\)（\quad）。
  \par\noindent\begin{tabular}{@{}>{\raggedright\arraybackslash}p{0.22\linewidth}>{\raggedright\arraybackslash}p{0.22\linewidth}>{\raggedright\arraybackslash}p{0.22\linewidth}>{\raggedright\arraybackslash}p{0.22\linewidth}@{}}
  A. \(a\) & B. \(a+b(a+b)\) & C. \(a+1\) & D. \(\displaystyle \dfrac{a}{1-b}\)
  \end{tabular}
\end{exercise}

\begin{solution}
\par
设内层函数
\[
  g(t)=f(t,t),
\]
则
\[
  \varphi(t)=f(t,g(t)).
\]
由链式法则，
\[
  \varphi'(t)=f_x(t,g(t))+f_y(t,g(t))g'(t).
\]
又
\[
  g'(t)=f_x(t,t)+f_y(t,t),
\]
所以在 $t=0$ 时，由 $f(0,0)=0$ 得 $g(0)=0$，并且
\[
  g'(0)=f_x(0,0)+f_y(0,0)=a+b.
\]
于是
\[
  \varphi'(0)
  =
  f_x(0,0)+f_y(0,0)\,g'(0)
  =
  a+b(a+b).
\]
选 B。
\end{solution}
\begin{exercise}
\par
设 \(C:x^2+y^2=a^2\)，取逆时针方向，则 \(\displaystyle \oint_C xy^2\,\mathrm{d}y-x^2y\,\mathrm{d}x=\)（\quad）。
  \par\noindent\begin{tabular}{@{}>{\raggedright\arraybackslash}p{0.22\linewidth}>{\raggedright\arraybackslash}p{0.22\linewidth}>{\raggedright\arraybackslash}p{0.22\linewidth}>{\raggedright\arraybackslash}p{0.22\linewidth}@{}}
  A. \(\displaystyle \dfrac{\pi a^4}{2}\) & B. \(\displaystyle -\dfrac{\pi a^4}{2}\) & C. \(\pi a^4\) & D. \(-\pi a^4\)
  \end{tabular}
\end{exercise}

\begin{solution}
\par
把积分写成
\[
  \oint_C P\,\mathrm{d}x+Q\,\mathrm{d}y,
  \qquad
  P=-x^2y,\quad Q=xy^2.
\]
曲线 $C$ 取逆时针方向，正好是圆盘 $D:x^2+y^2\le a^2$ 的正向边界。由 Green 公式，
\[
  \oint_C P\,\mathrm{d}x+Q\,\mathrm{d}y
  =
  \iint_D(Q_x-P_y)\,\mathrm{d}A.
\]
计算偏导：
\[
  Q_x=y^2,\qquad P_y=-x^2,
\]
故
\[
  Q_x-P_y=x^2+y^2.
\]
改用极坐标，
\[
  \iint_D(x^2+y^2)\,\mathrm{d}A
  =
  \int_0^{2\pi}\int_0^a r^2\cdot r\,\mathrm{d}r\,\mathrm{d}\theta
  =
  2\pi\cdot\frac{a^4}{4}
  =
  \frac{\pi a^4}{2}.
\]
选 A。
\end{solution}
\begin{exercise}
\par
幂级数 \(\displaystyle \sum_{n=1}^{\infty}\frac{x^n}{n2^n}\) 的收敛域是（\quad）。
  \par\noindent\begin{tabular}{@{}>{\raggedright\arraybackslash}p{0.22\linewidth}>{\raggedright\arraybackslash}p{0.22\linewidth}>{\raggedright\arraybackslash}p{0.22\linewidth}>{\raggedright\arraybackslash}p{0.22\linewidth}@{}}
  A. \((-2,2)\) & B. \([-2,2]\) & C. \((-2,2]\) & D. \([-2,2)\)
  \end{tabular}
\end{exercise}

\begin{solution}
\par
原级数可写为
\[
  \sum_{n=1}^{\infty}\frac1n\left(\frac{x}{2}\right)^n.
\]
先看幂级数内部：当
\[
  \left|\frac{x}{2}\right|<1
\]
时收敛，所以收敛半径为 $R=2$，开区间为 $(-2,2)$。端点必须另判。当 $x=2$ 时，
\[
  \sum_{n=1}^{\infty}\frac{2^n}{n2^n}
  =
  \sum_{n=1}^{\infty}\frac1n
\]
为调和级数，发散。当 $x=-2$ 时，
\[
  \sum_{n=1}^{\infty}\frac{(-2)^n}{n2^n}
  =
  \sum_{n=1}^{\infty}\frac{(-1)^n}{n}
\]
为交错调和级数，收敛。因此收敛域为
\[
  [-2,2).
\]
选 D。
\end{solution}
\begin{exercise}
\par
设 \(f(x)\) 是周期为 \(2\) 的周期函数，它在区间 \((-1,1]\) 上的表达式为
  \[
    f(x)=
    \begin{cases}
      2, & -1<x\le 0,\\
      x^3, & 0<x\le 1,
    \end{cases}
  \]
  则 \(f(x)\) 的 Fourier 级数在 \(x=1\) 处收敛于（\quad）。
  \par\noindent\begin{tabular}{@{}>{\raggedright\arraybackslash}p{0.22\linewidth}>{\raggedright\arraybackslash}p{0.22\linewidth}>{\raggedright\arraybackslash}p{0.22\linewidth}>{\raggedright\arraybackslash}p{0.22\linewidth}@{}}
  A. \(1\) & B. \(0\) & C. \(\displaystyle \dfrac32\) & D. \(\displaystyle -\dfrac32\)
  \end{tabular}
\end{exercise}

\begin{solution}
\par
Fourier 级数在跳跃点处收敛到左右极限的平均值。这里函数周期为 $2$，在基本区间 $(-1,1]$ 上，
\[
  f(1-0)=1^3=1.
\]
而 $1+0$ 按周期 $2$ 折回到 $-1+0$，此时
\[
  f(1+0)=f(-1+0)=2.
\]
因此 Fourier 级数在 $x=1$ 处的和为
\[
  \frac{f(1-0)+f(1+0)}2
  =
  \frac{1+2}{2}
  =
  \frac32.
\]
选 C。
\end{solution}

\par\medskip
\noindent\textbf{二、填空题（每小题3分，共15分）}\par\nobreak\smallskip
\begin{exercise}
\par
微分方程 \(y''-2y'+y=0\) 的通解为 \underline{\hspace{5cm}}。
\end{exercise}

\begin{solution}
\par
设 $y=e^{rx}$，得到特征方程
\[
  r^2-2r+1=0,\qquad (r-1)^2=0.
\]
特征根 $r=1$ 是二重根。二阶常系数齐次方程遇到二重根时，对应两个线性无关解为 $e^x$ 与 $xe^x$，故通解为
\[
  y=(C_1+C_2x)e^x.
\]
这里的 \(xe^x\) 不能漏掉：重根只给出一个指数函数 \(e^x\) 还不够二阶方程的两个独立任意常数，必须再乘一个 \(x\) 构造第二个线性无关解。
\end{solution}
\begin{exercise}
\par
函数 \(\displaystyle z=\arctan\frac{y}{x}+\ln\sqrt{x^2+y^2}\) 在点 \((1,2)\) 处的梯度为 \underline{\hspace{5cm}}。
\end{exercise}

\begin{solution}
\par
把
\[
  z=\arctan\frac{y}{x}+\ln\sqrt{x^2+y^2}
  =
  \arctan\frac{y}{x}+\frac12\ln(x^2+y^2)
\]
分别对 $x,y$ 求偏导。对 $x$ 求导：
\[
  z_x
  =
  \frac{1}{1+(y/x)^2}\left(-\frac{y}{x^2}\right)
  +
  \frac{x}{x^2+y^2}
  =
  -\frac{y}{x^2+y^2}+\frac{x}{x^2+y^2}.
\]
对 $y$ 求导：
\[
  z_y
  =
  \frac{1}{1+(y/x)^2}\cdot\frac1x
  +
  \frac{y}{x^2+y^2}
  =
  \frac{x}{x^2+y^2}+\frac{y}{x^2+y^2}.
\]
代入 $(1,2)$，得
\[
  z_x(1,2)=\frac{1-2}{5}=-\frac15,\qquad
  z_y(1,2)=\frac{1+2}{5}=\frac35.
\]
因此
\[
  \nabla z(1,2)=\left(-\frac15,\frac35\right).
\]
\end{solution}
\begin{exercise}
\par
交换积分的顺序，则 \(\displaystyle \int_0^2\mathrm{d}x\int_x^2 \mathrm{e}^{-y^2}\,\mathrm{d}y=\) \underline{\hspace{5cm}}。
\end{exercise}

\begin{solution}
\par
原积分区域由
\[
  0\le x\le2,\qquad x\le y\le2
\]
给出，也就是三角形区域
\[
  0\le y\le2,\qquad 0\le x\le y.
\]
因此交换积分顺序后，
\[
  \int_0^2\mathrm{d}x\int_x^2 e^{-y^2}\,\mathrm{d}y
  =
  \int_0^2\mathrm{d}y\int_0^y e^{-y^2}\,\mathrm{d}x.
\]
由于内层积分中 $e^{-y^2}$ 与 $x$ 无关，
\[
  \int_0^y e^{-y^2}\,\mathrm{d}x=ye^{-y^2}.
\]
于是
\[
  \int_0^2 y e^{-y^2}\,\mathrm{d}y
  =
  -\frac12 e^{-y^2}\Big|_0^2
  =
  \frac{1-e^{-4}}2.
\]
\end{solution}
\begin{exercise}
\par
设 \(\Sigma:x^2+y^2+z^2=a^2\)，则
  \(\displaystyle \iint_\Sigma (x^2+y^2+z^2)\,\mathrm{d}S=\) \underline{\hspace{5cm}}。
\end{exercise}

\begin{solution}
\par
在球面
\[
  \Sigma:x^2+y^2+z^2=a^2
\]
上，被积函数 $x^2+y^2+z^2$ 恒等于 $a^2$。因此
\[
  \iint_\Sigma (x^2+y^2+z^2)\,\mathrm{d}S
  =
  a^2\iint_\Sigma \mathrm{d}S.
\]
半径为 $a$ 的球面面积为 $4\pi a^2$，所以
\[
  \iint_\Sigma (x^2+y^2+z^2)\,\mathrm{d}S
  =
  a^2\cdot4\pi a^2
  =
  4\pi a^4.
\]
\end{solution}
\begin{exercise}
\par
设幂级数
  \[
    \sum_{n=1}^{\infty}a_n(x-1)^n
  \]
  在 \(x=-1\) 处收敛。则此级数在 \(x=2\) 的敛散性是
  \underline{\hspace{5cm}}。
\end{exercise}

\begin{solution}
\par
这是以 $x=1$ 为中心的幂级数。若它在 $x=-1$ 处收敛，则点 $x=-1$ 到中心 $1$ 的距离为
\[
  |-1-1|=2.
\]
幂级数有这样的性质：若在某点收敛，则在所有距离中心更近的点绝对收敛。因此收敛半径至少为 $2$。而点 $x=2$ 到中心 $1$ 的距离为
\[
  |2-1|=1<2.
\]
所以该级数在 $x=2$ 处绝对收敛。
这里不能反过来断言收敛半径恰好等于 \(2\)；题目只给出 \(x=-1\) 收敛，所以能确定的是半径至少覆盖到距离更近的 \(x=2\)。
\end{solution}

\par\medskip
\noindent\textbf{三、计算题（每小题10分，共50分）}\par\nobreak\smallskip
\begin{exercise}
\par
设 \(f(\xi,\eta)\) 具有连续的二阶偏导数，且满足 Laplace 方程
  \(\displaystyle \frac{\partial^2 f}{\partial \xi^2}+\frac{\partial^2 f}{\partial \eta^2}=0\)，令
  \(z=f(x^2-y^2,2xy)\)，求
  \(\displaystyle \frac{\partial^2z}{\partial x^2}+\frac{\partial^2z}{\partial y^2}\)。
\end{exercise}

\begin{solution}
\par
记
\[
  \xi=x^2-y^2,\qquad \eta=2xy,\qquad z=f(\xi,\eta).
\]
先列出一阶偏导：
\[
  \xi_x=2x,\quad \eta_x=2y,\qquad
  \xi_y=-2y,\quad \eta_y=2x.
\]
由链式法则，
\[
  z_x=2x f_\xi+2y f_\eta,\qquad
  z_y=-2y f_\xi+2x f_\eta.
\]
继续求二阶偏导并合并。对 $z_x$ 再对 $x$ 求导：
\[
  z_{xx}
  =
  2f_\xi
  +4x^2f_{\xi\xi}
  +8xyf_{\xi\eta}
  +4y^2f_{\eta\eta}.
\]
对 $z_y$ 再对 $y$ 求导：
\[
  z_{yy}
  =
  -2f_\xi
  +4y^2f_{\xi\xi}
  -8xyf_{\xi\eta}
  +4x^2f_{\eta\eta}.
\]
相加后一次项与混合项抵消：
\[
  z_{xx}+z_{yy}
  =
  4(x^2+y^2)(f_{\xi\xi}+f_{\eta\eta}).
\]
题设 $f$ 满足 Laplace 方程
\[
  f_{\xi\xi}+f_{\eta\eta}=0,
\]
所以
\[
  \frac{\partial^2z}{\partial x^2}
  +
  \frac{\partial^2z}{\partial y^2}
  =
  0.
\]
\end{solution}
\begin{exercise}
\par
计算三重积分
  \[
    I=\iiint_\Omega z\mathrm{e}^{-(x^2+y^2+z^2)}\,\mathrm{d}V,
    \qquad
    \Omega:\begin{cases}
      1\le x^2+y^2+z^2\le 4,\\
      x\ge 0,\ y\ge 0,\ z\ge 0.
    \end{cases}
  \]
\end{exercise}

\begin{solution}
\par
用球坐标
\[
  x=r\sin\varphi\cos\theta,\qquad
  y=r\sin\varphi\sin\theta,\qquad
  z=r\cos\varphi.
\]
区域在第一卦限且夹在球面 $r=1$ 与 $r=2$ 之间，所以
\[
  1\le r\le2,\qquad
  0\le\varphi\le\frac{\pi}{2},\qquad
  0\le\theta\le\frac{\pi}{2}.
\]
体积元为 $\mathrm{d}V=r^2\sin\varphi\,\mathrm{d}r\,\mathrm{d}\varphi\,\mathrm{d}\theta$，且
\[
  z e^{-(x^2+y^2+z^2)}
  =
  r\cos\varphi\,e^{-r^2}.
\]
因此
\[
  I
  =
  \int_0^{\pi/2}\int_0^{\pi/2}\int_1^2
  r^3e^{-r^2}\cos\varphi\sin\varphi
  \,\mathrm{d}r\,\mathrm{d}\varphi\,\mathrm{d}\theta.
\]
三个变量可以分离：
\[
  \int_0^{\pi/2}\mathrm{d}\theta=\frac{\pi}{2},
  \qquad
  \int_0^{\pi/2}\cos\varphi\sin\varphi\,\mathrm{d}\varphi=\frac12.
\]
径向积分令 $u=r^2$，则 $r^3\mathrm{d}r=\dfrac12u\,\mathrm{d}u$，所以
\[
  \int_1^2 r^3e^{-r^2}\,\mathrm{d}r
  =
  \frac12\int_1^4 u e^{-u}\,\mathrm{d}u
  =
  \frac12\left[-(u+1)e^{-u}\right]_{1}^{4}
  =
  \frac12\left(\frac2e-\frac5{e^4}\right).
\]
于是
\[
  I
  =
  \frac{\pi}{2}\cdot\frac12\cdot
  \frac12\left(\frac2e-\frac5{e^4}\right)
  =
  \frac{\pi}{8}\left(\frac2e-\frac5{e^4}\right).
\]
\end{solution}
\begin{exercise}
\par
已知 \(\varphi(\pi)=1\)，试确定 \(\varphi(x)\)，使曲线积分
  \(\displaystyle I=\int_A^B\left[\sin x-\varphi(x)\right]\frac{y}{x}\,\mathrm{d}x+\varphi(x)\,\mathrm{d}y\)
  与路径无关，并求当 \(A,B\) 两点分别为 \((1,0),(\pi,\pi)\) 时积分 \(I\) 的值。
\end{exercise}

\begin{solution}
\par
记
\[
  P(x,y)=\left[\sin x-\varphi(x)\right]\frac{y}{x},
  \qquad
  Q(x,y)=\varphi(x).
\]
曲线积分与路径无关的条件是
\[
  P_y=Q_x.
\]
计算得
\[
  P_y=\frac{\sin x-\varphi(x)}{x},
  \qquad
  Q_x=\varphi'(x).
\]
所以
\[
  \varphi'(x)=\frac{\sin x-\varphi(x)}{x}.
\]
整理为
\[
  x\varphi'(x)+\varphi(x)=\sin x,
  \qquad
  (x\varphi(x))'=\sin x.
\]
积分得
\[
  x\varphi(x)=-\cos x+C.
\]
由 $\varphi(\pi)=1$，得
\[
  \pi\cdot1=-\cos\pi+C=1+C,
\]
故 $C=\pi-1$。于是
\[
  \varphi(x)=\frac{\pi-1-\cos x}{x}.
\]
下面求积分值。由上面的条件可知该微分形式为全微分。取势函数
\[
  \Phi(x,y)=y\varphi(x).
\]
则
\[
  \Phi_y=\varphi(x)=Q,
\]
且
\[
  \Phi_x=y\varphi'(x)
  =
  y\frac{\sin x-\varphi(x)}{x}
  =
  P.
\]
因此
\[
  I=\Phi(\pi,\pi)-\Phi(1,0).
\]
代入端点：
\[
  \Phi(\pi,\pi)=\pi\varphi(\pi)=\pi,\qquad
  \Phi(1,0)=0.
\]
故
\[
  I=\pi.
\]
\end{solution}
\begin{exercise}
\par
计算曲面积分
  \[
    I=\iint_\Sigma (8y+1)x\,\mathrm{d}y\,\mathrm{d}z+2(1-y^2)\,\mathrm{d}z\,\mathrm{d}x-4yz\,\mathrm{d}x\,\mathrm{d}y,
  \]
  其中 \(\Sigma:y=1+x^2+z^2\)，介于 \(1\le y\le 3\) 之间，且其法向量与 \(y\) 轴正向的夹角恒大于 \(\displaystyle \dfrac{\pi}{2}\)。
\end{exercise}

\begin{solution}
\par
把曲面写成参数形式
\[
  \mathbf r(x,z)=(x,\,1+x^2+z^2,\,z).
\]
题设要求法向量与 $y$ 轴正向夹角大于 $\dfrac{\pi}{2}$，即法向量的 $y$ 分量为负。计算
\[
  \mathbf r_x=(1,2x,0),\qquad
  \mathbf r_z=(0,2z,1),
\]
\[
  \mathbf r_x\times\mathbf r_z=(2x,-1,2z),
\]
正好满足 $y$ 分量为负。投影区域由 $1\le y\le3$ 给出：
\[
  1\le1+x^2+z^2\le3
  \quad\Longleftrightarrow\quad
  x^2+z^2\le2.
\]
记 $D:x^2+z^2\le2$。第二类曲面积分为通量形式，向量场为
\[
  (P,Q,R)=\bigl((8y+1)x,\ 2(1-y^2),\ -4yz\bigr).
\]
因此
\[
  I=\iint_D (P,Q,R)\cdot(2x,-1,2z)\,\mathrm{d}x\,\mathrm{d}z,
  \qquad y=1+x^2+z^2.
\]
即
\[
  I=\iint_D\left[(8y+1)x\cdot2x-2(1-y^2)-8yz^2\right]\,\mathrm{d}x\,\mathrm{d}z.
\]
令 $r^2=x^2+z^2$，代入 $y=1+r^2$ 并整理。利用圆盘关于 $x,z$ 的对称性，可写成
\[
  I
  =
  10\iint_Dx^2\,\mathrm{d}A
  +4\iint_Dr^2\,\mathrm{d}A
  +6\iint_Dr^4\,\mathrm{d}A.
\]
下面分别计算三个积分。对半径为 $\sqrt2$ 的圆盘，
\[
  \iint_Dx^2\,\mathrm{d}A
  =
  \frac12\iint_Dr^2\,\mathrm{d}A
  =
  \frac12\int_0^{2\pi}\int_0^{\sqrt2}r^2\cdot r\,\mathrm{d}r\,\mathrm{d}\theta
  =
  \pi.
\]
同时
\[
  \iint_Dr^2\,\mathrm{d}A
  =
  \int_0^{2\pi}\int_0^{\sqrt2}r^3\,\mathrm{d}r\,\mathrm{d}\theta
  =
  2\pi,
\]
\[
  \iint_Dr^4\,\mathrm{d}A
  =
  \int_0^{2\pi}\int_0^{\sqrt2}r^5\,\mathrm{d}r\,\mathrm{d}\theta
  =
  \frac{8\pi}{3}.
\]
故
\[
  I=10\pi+4\cdot2\pi+6\cdot\frac{8\pi}{3}
  =
  34\pi.
\]
\end{solution}
\begin{exercise}
\par
求级数 \(\displaystyle \sum_{n=1}^{\infty}\frac{(-1)^n n}{(2n+1)!}\) 的和。
\end{exercise}

\begin{solution}
\par
利用
\[
  n=\frac{(2n+1)-1}{2},
\]
原级数可拆为
\[
  \sum_{n=1}^{\infty}\frac{(-1)^n n}{(2n+1)!}
  =
  \frac12\sum_{n=1}^{\infty}\frac{(-1)^n(2n+1)}{(2n+1)!}
  -
  \frac12\sum_{n=1}^{\infty}\frac{(-1)^n}{(2n+1)!}.
\]
第一项中
\[
  \frac{2n+1}{(2n+1)!}=\frac1{(2n)!},
\]
所以
\[
  \sum_{n=1}^{\infty}\frac{(-1)^n(2n+1)}{(2n+1)!}
  =
  \sum_{n=1}^{\infty}\frac{(-1)^n}{(2n)!}
  =
  \cos1-1.
\]
第二项由正弦展开
\[
  \sin1=\sum_{n=0}^{\infty}\frac{(-1)^n}{(2n+1)!}
\]
可得
\[
  \sum_{n=1}^{\infty}\frac{(-1)^n}{(2n+1)!}
  =
  \sin1-1.
\]
代回：
\[
  \sum_{n=1}^{\infty}\frac{(-1)^n n}{(2n+1)!}
  =
  \frac12(\cos1-1)-\frac12(\sin1-1)
  =
  \frac{\cos1-\sin1}{2}.
\]
\end{solution}

\par\medskip
\noindent\textbf{四、证明题（本题10分）}\par\nobreak\smallskip
\begin{exercise}
\par
证明级数 \(\displaystyle \sum_{n=1}^{\infty}\frac{\sin(nx)}{n^2}\) 在 \((-\infty,+\infty)\) 上一致收敛。
\end{exercise}

\begin{solution}
\par
对任意实数 $x$，都有 $|\sin(nx)|\le1$，因此
\[
  \left|\frac{\sin(nx)}{n^2}\right|\le \frac1{n^2}.
\]
右端与 $x$ 无关，且数项级数
\[
  \sum_{n=1}^{\infty}\frac1{n^2}
\]
收敛。由 Weierstrass 判别法，函数项级数
\[
  \sum_{n=1}^{\infty}\frac{\sin(nx)}{n^2}
\]
在整个实轴上，也就是在 $(-\infty,+\infty)$ 上一致收敛。
\end{solution}

\par\medskip
\noindent\textbf{五、应用题（本题10分）}\par\nobreak\smallskip
\begin{exercise}
\par
求内接于半径为 \(a\) 的球且有最大体积的长方体。
\end{exercise}

\begin{solution}
\par
最大体积的长方体可取球心为中心，棱分别平行于三坐标轴。设三个半边长为
\[
  x>0,\qquad y>0,\qquad z>0.
\]
长方体顶点在球面上时体积才可能最大，因此约束为
\[
  x^2+y^2+z^2=a^2.
\]
体积为
\[
  V=8xyz.
\]
在约束下最大化 $xyz$。用拉格朗日乘数法，令
\[
  \mathcal L=\ln x+\ln y+\ln z-\lambda(x^2+y^2+z^2-a^2).
\]
驻点方程为
\[
  \frac1x=2\lambda x,\qquad
  \frac1y=2\lambda y,\qquad
  \frac1z=2\lambda z.
\]
于是
\[
  x^2=y^2=z^2.
\]
由于 $x,y,z>0$，得到 $x=y=z$。代入约束：
\[
  3x^2=a^2,\qquad x=y=z=\frac{a}{\sqrt3}.
\]
因此长方体三条棱长均为
\[
  2x=2y=2z=\frac{2a}{\sqrt3},
\]
也就是说最大体积长方体是内接正方体。最大体积为
\[
  V_{\max}
  =
  8\left(\frac{a}{\sqrt3}\right)^3
  =
  \frac{8a^3}{3\sqrt3}.
\]
\end{solution}

\section{2014级工科数学分析下 A卷}

\subsubsection*{一、填空题（每小题 4 分，共 20 分）}
\begin{exercise}
\par
微分方程 \(\displaystyle y''-4y'=e^{4x}\) 的特解形式为 \underline{\hspace{4cm}}。
\end{exercise}

\begin{solution}
\par
先看对应齐次方程的特征方程：
\[
  r^2-4r=0,\qquad r(r-4)=0.
\]
右端为 $e^{4x}$，而 $r=4$ 是齐次方程的一重特征根，所以待定特解要在通常形式 $Ae^{4x}$ 前多乘一个 $x$。因此特解形式应取
\[
  y^*=Axe^{4x}.
\]
若不乘 \(x\)，\(Ae^{4x}\) 已经落在齐次解空间内，代入左端会被齐次算子消掉，无法产生右端的 \(e^{4x}\)。本题只问特解形式，因此常数 \(A\) 不必继续求出。
\end{solution}
\begin{exercise}
\par
设 \(\displaystyle xyz+\sqrt{x^2+y^2+z^2}=\sqrt2\)，则 $\left.\mathrm{d}z\right|_{(1,0,-1)}=$ \underline{\hspace{3cm}}。
\end{exercise}

\begin{solution}
\par
令
\[
  F(x,y,z)=xyz+\sqrt{x^2+y^2+z^2}-\sqrt2.
\]
隐函数关系为 $F(x,y,z)=0$，所以
\[
  F_x\,\mathrm{d}x+F_y\,\mathrm{d}y+F_z\,\mathrm{d}z=0.
\]
先求偏导：
\[
  F_x=yz+\frac{x}{\sqrt{x^2+y^2+z^2}},
\]
\[
  F_y=xz+\frac{y}{\sqrt{x^2+y^2+z^2}},
  \qquad
  F_z=xy+\frac{z}{\sqrt{x^2+y^2+z^2}}.
\]
在点 $(1,0,-1)$ 处，$\sqrt{x^2+y^2+z^2}=\sqrt2$，故
\[
  F_x=\frac1{\sqrt2},\qquad
  F_y=-1,\qquad
  F_z=-\frac1{\sqrt2}.
\]
代入全微分方程：
\[
  \frac1{\sqrt2}\mathrm{d}x-\mathrm{d}y-\frac1{\sqrt2}\mathrm{d}z=0.
\]
解得
\[
  \mathrm{d}z=\mathrm{d}x-\sqrt2\,\mathrm{d}y.
\]
\end{solution}
\begin{exercise}
\par
交换积分的顺序，则
  \(\displaystyle \int_0^1\mathrm{d}y\int_y^{\sqrt{2y-y^2}}f(x,y)\,\mathrm{d}x=\) \underline{\hspace{3cm}}。
\end{exercise}

\begin{solution}
\par
原区域满足
\[
  0\le y\le1,\qquad y\le x\le\sqrt{2y-y^2}.
\]
由 $x\le\sqrt{2y-y^2}$ 得
\[
  x^2\le2y-y^2
  \quad\Longleftrightarrow\quad
  x^2+(y-1)^2\le1.
\]
所以区域是圆 $x^2+(y-1)^2=1$ 内、直线 $y=x$ 下方的部分。两条边界在第一象限内交于 $(0,0)$ 与 $(1,1)$，故交换顺序后
\[
  0\le x\le1.
\]
固定 $x$，下边界来自圆的下半支
\[
  y=1-\sqrt{1-x^2},
\]
上边界为直线 $y=x$。因此
\[
  \int_0^1\mathrm{d}y\int_y^{\sqrt{2y-y^2}}f(x,y)\,\mathrm{d}x
  =
  \int_0^1\mathrm{d}x\int_{1-\sqrt{1-x^2}}^x f(x,y)\,\mathrm{d}y.
\]
\end{solution}
\begin{exercise}
\par
设 $f(x)=x^2\ (0\le x<1)$，而
  \[
  S(x)=\sum_{n=1}^{\infty}b_n\sin n\pi x\quad(-\infty<x<+\infty),
  \qquad
  b_n=2\int_0^1 f(x)\sin n\pi x\,\mathrm{d}x.
  \]
  则 $\displaystyle S\!\left(-\dfrac12\right)$ \underline{\hspace{3cm}}。
\end{exercise}

\begin{solution}
\par
这是 $[0,1]$ 上函数 $f(x)=x^2$ 的 Fourier 正弦级数。正弦级数对应先作奇延拓，再作周期为 $2$ 的周期延拓。因此在 $x=-\dfrac12$ 处，
\[
  S\!\left(-\frac12\right)
  =
  -f\!\left(\frac12\right)
  =
  -\left(\frac12\right)^2
  =
  -\frac14.
\]
这里 \(x=-\frac12\) 先由奇延拓对应到区间内的 \(\frac12\)，再加负号；因为该点不是延拓后的跳跃点，所以不需要再取左右极限平均。
\end{solution}
\begin{exercise}
\par
设 \(\displaystyle L:\frac{x^2}{a^2}+\frac{y^2}{b^2}=1\) 取逆时针方向，则
  \(\displaystyle \int_L(x-y)\,\mathrm{d}x+(y-x)\,\mathrm{d}y=\) \underline{\hspace{3cm}}。
\end{exercise}

\begin{solution}
\par
记
\[
  P=x-y,\qquad Q=y-x.
\]
椭圆 $L$ 取逆时针方向，是其内部区域 $D$ 的正向边界。由 Green 公式，
\[
  \int_L P\,\mathrm{d}x+Q\,\mathrm{d}y
  =
  \iint_D(Q_x-P_y)\,\mathrm{d}A.
\]
计算
\[
  Q_x=-1,\qquad P_y=-1,
\]
所以
\[
  Q_x-P_y=0.
\]
故
\[
  \int_L(x-y)\,\mathrm{d}x+(y-x)\,\mathrm{d}y=0.
\]
由于 \(P,Q\) 在整个椭圆区域内连续可偏导，且 Green 公式中的面积分密度恒为零，所以结果与椭圆半轴长无关。
\end{solution}

\subsubsection*{二、计算题（每小题 10 分，共 40 分）}
\begin{exercise}
\par
设 \(\displaystyle u=yf\!\left(\frac{x}{y}\right)+xg\!\left(\frac{y}{x}\right)\)，
  其中 $f,g$ 具有二阶连续偏导数，求
  \(\displaystyle x\frac{\partial^2u}{\partial x^2}+y\frac{\partial^2u}{\partial x\partial y}\)。
\end{exercise}

\begin{solution}
\par
注意
\[
  u=yf\!\left(\frac{x}{y}\right)+xg\!\left(\frac{y}{x}\right)
\]
是关于 $x,y$ 的一阶齐次函数。事实上，对任意 $\lambda>0$，
\[
  u(\lambda x,\lambda y)
  =
  \lambda y f\!\left(\frac{x}{y}\right)
  +
  \lambda x g\!\left(\frac{y}{x}\right)
  =
  \lambda u(x,y).
\]
由 Euler 齐次函数公式，
\[
  xu_x+yu_y=u.
\]
对上式两边关于 $x$ 求偏导，注意 $y$ 视为常数：
\[
  u_x+xu_{xx}+yu_{yx}=u_x.
\]
两边消去 $u_x$，且 $u_{yx}=u_{xy}$，得到
\[
  xu_{xx}+yu_{xy}=0.
\]
\end{solution}
\begin{exercise}
\par
计算三重积分 \(\displaystyle I=\iiint_\Omega (x^2+y^2)\,\mathrm{d}x\,\mathrm{d}y\,\mathrm{d}z\)，
  其中 $\Omega$ 是由圆锥面 $z=\sqrt{x^2+y^2}$ 与平面 $z=1$ 所围成的区域。
\end{exercise}

\begin{solution}
\par
用柱坐标
\[
  x=r\cos\theta,\qquad y=r\sin\theta.
\]
圆锥面 $z=\sqrt{x^2+y^2}$ 即 $z=r$。由圆锥面与平面 $z=1$ 围成的实心圆锥可写为
\[
  0\le z\le1,\qquad 0\le r\le z,\qquad 0\le\theta\le2\pi.
\]
又
\[
  x^2+y^2=r^2,\qquad \mathrm{d}x\,\mathrm{d}y\,\mathrm{d}z=r\,\mathrm{d}r\,\mathrm{d}\theta\,\mathrm{d}z.
\]
因此
\[
  I
  =
  \int_0^{2\pi}\int_0^1\int_0^z r^2\cdot r\,\mathrm{d}r\,\mathrm{d}z\,\mathrm{d}\theta
  =
  2\pi\int_0^1\frac{z^4}{4}\,\mathrm{d}z
  =
  \frac{\pi}{10}.
\]
\end{solution}
\begin{exercise}
\par
计算曲线积分 \(\displaystyle I=\int_L \cos(x+y^2)\,\mathrm{d}x+\left[2y\cos(x+y^2)-\frac1{\sqrt{1+y^4}}\right]\,\mathrm{d}y\)，
  其中
  \[
  L:\begin{cases}
  x=a(t-\sin t),\\
  y=a(1-\cos t),
  \end{cases}
  \]
  上从 $O(0,0)$ 到 $A(2\pi a,0)$ 的有向曲线弧段。
\end{exercise}

\begin{solution}
\par
先把前两项识别为全微分：
\[
  d\sin(x+y^2)
  =
  \cos(x+y^2)\,\mathrm{d}x+2y\cos(x+y^2)\,\mathrm{d}y.
\]
因此
\[
  I
  =
  \int_L d\sin(x+y^2)
  -
  \int_L\frac{\mathrm{d}y}{\sqrt{1+y^4}}.
\]
第一项只看端点：
\[
  \int_L d\sin(x+y^2)
  =
  \sin(x+y^2)\Big|_{O}^{A}.
\]
在 $O(0,0)$ 处为 $\sin0=0$；在 $A(2\pi a,0)$ 处为 $\sin(2\pi a)$，所以第一项为 $\sin(2\pi a)$。

第二项是仅关于 $y$ 的全微分。设
\[
  H'(y)=\frac1{\sqrt{1+y^4}},
\]
则
\[
  \int_L\frac{\mathrm{d}y}{\sqrt{1+y^4}}
  =
  H(y)\Big|_{0}^{0}=0,
\]
因为路径起点与终点的 $y$ 坐标都为 $0$。因此
\[
  I=\sin(2\pi a).
\]
\end{solution}
\begin{exercise}
\par
设 $\Sigma$ 是抛物面 $z=x^2+y^2$ 被平面 $z=0$ 及 $z=1$ 所截部分的下侧，计算曲面积分 \(\displaystyle I=\iint_\Sigma xy\,\mathrm{d}y\,\mathrm{d}z+2y\,\mathrm{d}z\,\mathrm{d}x+(z^2-3z)\,\mathrm{d}x\,\mathrm{d}y\)。
\end{exercise}

\begin{solution}
\par
曲面为图形 $z=x^2+y^2$。下侧取向的面积向量为
\[
  (z_x,z_y,-1)\,\mathrm{d}x\,\mathrm{d}y=(2x,2y,-1)\,\mathrm{d}x\,\mathrm{d}y.
\]
投影区域由 $0\le z\le1$ 给出：
\[
  D:x^2+y^2\le1.
\]
向量场为
\[
  (P,Q,R)=(xy,\,2y,\,z^2-3z).
\]
因此
\[
  I=\iint_D (P,Q,R)\cdot(2x,2y,-1)\,\mathrm{d}x\,\mathrm{d}y.
\]
代入 $z=x^2+y^2$：
\[
  I
  =
  \iint_D\bigl(2x^2y+4y^2-(z^2-3z)\bigr)\,\mathrm{d}x\,\mathrm{d}y.
\]
其中 $2x^2y$ 关于 $y$ 为奇函数，在圆盘 $D$ 上积分为 $0$。令 $r^2=x^2+y^2$，则
\[
  I=\iint_D(4y^2-r^4+3r^2)\,\mathrm{d}A.
\]
分别计算：
\[
  \iint_D4y^2\,\mathrm{d}A=4\cdot\frac{\pi}{4}=\pi,
\]
\[
  \iint_Dr^4\,\mathrm{d}A
  =
  \int_0^{2\pi}\int_0^1 r^4\cdot r\,\mathrm{d}r\,\mathrm{d}\theta
  =
  \frac{\pi}{3},
\]
\[
  \iint_D3r^2\,\mathrm{d}A
  =
  3\int_0^{2\pi}\int_0^1 r^2\cdot r\,\mathrm{d}r\,\mathrm{d}\theta
  =
  \frac{3\pi}{2}.
\]
故
\[
  I=\pi-\frac{\pi}{3}+\frac{3\pi}{2}
  =
  \frac{13\pi}{6}.
\]
\end{solution}

\subsubsection*{三、解答题（每小题 10 分，共 20 分）}
\begin{exercise}
\par
求幂级数 \(\displaystyle \sum_{n=1}^{\infty}\frac{(-1)^n n}{(2n+1)!}x^{2n+1}\) 的和函数。
\end{exercise}

\begin{solution}
\par
记
\[
  S(x)=\sum_{n=1}^{\infty}\frac{(-1)^n n}{(2n+1)!}x^{2n+1}.
\]
从正弦展开出发：
\[
  \sin x=\sum_{n=0}^{\infty}\frac{(-1)^n x^{2n+1}}{(2n+1)!}.
\]
两边求导得
\[
  \cos x=\sum_{n=0}^{\infty}\frac{(-1)^n x^{2n}}{(2n)!}.
\]
于是
\[
  x\cos x
  =
  \sum_{n=0}^{\infty}\frac{(-1)^n x^{2n+1}}{(2n)!}
  =
  \sum_{n=0}^{\infty}\frac{(-1)^n(2n+1)x^{2n+1}}{(2n+1)!}.
\]
相减：
\[
  x\cos x-\sin x
  =
  \sum_{n=1}^{\infty}
  \frac{(-1)^n\,2n}{(2n+1)!}x^{2n+1}.
\]
因此
\[
  S(x)=\frac{x\cos x-\sin x}{2}.
\]
\end{solution}
\begin{exercise}
\par
设 \(\displaystyle f(0)=0,\qquad f'(x)=1+\int_0^x(6\sin^2t-f(t))\,\mathrm{d}t\)，
  其中 $f(x)$ 二次可微，求 $f(x)$。
\end{exercise}

\begin{solution}
\par
对题设
\[
  f'(x)=1+\int_0^x(6\sin^2t-f(t))\,\mathrm{d}t
\]
两边求导，得
\[
  f''(x)=6\sin^2x-f(x).
\]
移项：
\[
  f''+f=6\sin^2x.
\]
利用
\[
  6\sin^2x=3-3\cos2x,
\]
得到
\[
  f''+f=3-3\cos2x.
\]
对应齐次解为 $C_1\cos x+C_2\sin x$。右端常数 $3$ 的特解取 $3$；右端 $-3\cos2x$ 的特解取 $A\cos2x$，代入 $f''+f$ 得 $-3A\cos2x$，所以 $A=1$。故
\[
  f(x)=C_1\cos x+C_2\sin x+3+\cos2x.
\]
由题设 $f(0)=0$ 得
\[
  C_1+3+1=0,\qquad C_1=-4.
\]
再由原式代入 $x=0$ 得 $f'(0)=1$。而
\[
  f'(x)=-C_1\sin x+C_2\cos x-2\sin2x,
\]
故 $f'(0)=C_2=1$。因此
\[
  f(x)=-4\cos x+\sin x+3+\cos2x.
\]
\end{solution}

\subsubsection*{四、证明题（本题 10 分）}
\begin{exercise}
\par
设正项级数 $\displaystyle \sum_{n=1}^{\infty}a_n$ 收敛，证明
\(\displaystyle f(x)=\sum_{n=1}^{\infty}a_nx^n\)
在 $(-1,1)$ 上连续。
\end{exercise}

\begin{solution}
\par
任取 $0<r<1$，在闭区间 $[-r,r]$ 上有
\[
  |a_nx^n|\le a_nr^n\le a_n.
\]
由于 $\sum_{n=1}^{\infty}a_n$ 是收敛的正项级数，所以由 Weierstrass 判别法，
\[
  \sum_{n=1}^{\infty}a_nx^n
\]
在 $[-r,r]$ 上一致收敛。每一项 $a_nx^n$ 都是连续函数，一致收敛的连续函数项级数的和函数在 $[-r,r]$ 上连续。

对任意 $x_0\in(-1,1)$，总可以取 $r$ 使 $|x_0|<r<1$，于是 $x_0$ 落在某个 $[-r,r]$ 内。故 $f(x)$ 在 $(-1,1)$ 上连续。
\end{solution}

\subsubsection*{五、应用题（本题 10 分）}
\begin{exercise}
\par
求 $a,b$ 的值，使得包含圆周 $(x-1)^2+y^2=1$ 在其内部的椭圆 \(\displaystyle \frac{x^2}{a^2}+\frac{y^2}{b^2}=1\quad(a>0,\ b>0,\ a\ne b)\) 有最小的面积。
\end{exercise}

\begin{solution}
\par
椭圆面积为 $\pi ab$。令
\[
  \alpha=\frac1{a^2},\qquad \beta=\frac1{b^2},
\]
则最小化 $\pi ab$ 等价于最大化 $\alpha\beta$。圆周可参数化为
\[
  x=1+\cos t,\qquad y=\sin t.
\]
令 $u=\cos t$，则 $u\in[-1,1]$，圆周在椭圆内要求
\[
  \alpha(1+u)^2+\beta(1-u^2)\le1
  \qquad (-1\le u\le1).
\]
令 $s=1+u$，则 $0\le s\le2$，并且 $1-u^2=s(2-s)$。条件变为
\[
  \alpha s^2+\beta s(2-s)\le1
  \qquad (0\le s\le2).
\]
当 $s>0$ 时，
\[
  \beta\le
  \frac{1-\alpha s^2}{s(2-s)}.
\]
为了使 $\alpha\beta$ 最大，临界时右端对某个内点 $s$ 取最小值。对
\[
  q(s)=\frac{1-\alpha s^2}{s(2-s)}
\]
求极值，得
\[
  \alpha=\frac{s-1}{s^2}.
\]
此时
\[
  \beta=q(s)=\frac1s.
\]
所以
\[
  \alpha\beta=\frac{s-1}{s^3},\qquad 1<s\le2.
\]
再对它求最大值：
\[
  \frac{\mathrm{d}}{\mathrm{d}s}\left(\frac{s-1}{s^3}\right)
  =
  \frac{3-2s}{s^4}.
\]
故最大点为
\[
  s=\frac32.
\]
于是
\[
  \alpha=\frac{(3/2)-1}{(3/2)^2}=\frac{2}{9},
  \qquad
  \beta=\frac{2}{3}.
\]
换回 $a,b$：
\[
  a=\frac1{\sqrt\alpha}=\frac{3}{\sqrt2},
  \qquad
  b=\frac1{\sqrt\beta}=\sqrt{\frac32}.
\]
此时最小面积为
\[
  \pi ab=\frac{3\sqrt3}{2}\pi.
\]
\end{solution}

\section{2014级工科数学分析下 B卷}

\subsubsection*{一、填空题（每小题 4 分，共 20 分）}
\begin{exercise}
\par
微分方程 \(\displaystyle y''-y=2xe^x\) 的特解形式为 \underline{\hspace{4cm}}。
\end{exercise}

\begin{solution}
\par
对应齐次方程的特征方程为
\[
  r^2-1=0,\qquad r=\pm1.
\]
右端 $2xe^x$ 是一次多项式乘 $e^x$，通常应设特解为 $e^x(Ax+B)$。但 $r=1$ 已经是齐次方程的一重特征根，发生共振，需要再乘一个 $x$。因此特解形式为
\[
  y^*=xe^x(Ax+B).
\]
这里的多项式次数由右端 \(2x\) 决定，所以括号内保留一次多项式 \(Ax+B\)；前面的 \(x\) 只负责避开齐次解中的 \(e^x\)，不改变待定多项式的次数。
\end{solution}
\begin{exercise}
\par
设由方程 \(\displaystyle f(x-y,y-z,z-x)=0\)
  确定了 $z=z(x,y)$，其中 $f$ 有连续的偏导数，则 $\mathrm{d}z=$ \underline{\hspace{3cm}}。
\end{exercise}

\begin{solution}
\par
记
\[
  u=x-y,\qquad v=y-z,\qquad w=z-x,
\]
并记 $f_1,f_2,f_3$ 分别为 $f$ 对第一、第二、第三个自变量的偏导数。由
\[
  f(u,v,w)=0
\]
取全微分：
\[
  f_1\,\mathrm{d}u+f_2\,\mathrm{d}v+f_3\,\mathrm{d}w=0.
\]
而
\[
  \mathrm{d}u=\mathrm{d}x-\mathrm{d}y,\qquad \mathrm{d}v=\mathrm{d}y-\mathrm{d}z,\qquad \mathrm{d}w=\mathrm{d}z-\mathrm{d}x.
\]
代入得
\[
  f_1(\mathrm{d}x-\mathrm{d}y)+f_2(\mathrm{d}y-\mathrm{d}z)+f_3(\mathrm{d}z-\mathrm{d}x)=0.
\]
合并 $\mathrm{d}x,\mathrm{d}y,\mathrm{d}z$：
\[
  (f_1-f_3)\mathrm{d}x+(f_2-f_1)\mathrm{d}y+(f_3-f_2)\mathrm{d}z=0.
\]
若 $f_3\ne f_2$，则
\[
  \mathrm{d}z=\frac{(f_3-f_1)\,\mathrm{d}x+(f_1-f_2)\,\mathrm{d}y}{f_3-f_2}.
\]
\end{solution}
\begin{exercise}
\par
交换积分的顺序，则
  \(\displaystyle \int_0^1\mathrm{d}y\int_{y^2}^{\sqrt y}f(x,y)\,\mathrm{d}x=\) \underline{\hspace{3cm}}。
\end{exercise}

\begin{solution}
\par
原区域满足
\[
  0\le y\le1,\qquad y^2\le x\le \sqrt y.
\]
由 $y^2\le x$ 得 $y\le\sqrt x$；由 $x\le\sqrt y$ 得 $x^2\le y$。在 $0\le x,y\le1$ 中，固定 $x$ 后，
\[
  x^2\le y\le\sqrt x.
\]
因此交换顺序后为
\[
  \int_0^1\mathrm{d}y\int_{y^2}^{\sqrt y}f(x,y)\,\mathrm{d}x
  =
  \int_0^1\mathrm{d}x\int_{x^2}^{\sqrt x}f(x,y)\,\mathrm{d}y.
\]
\end{solution}
\begin{exercise}
\par
设 $f(x)$ 是周期为 $2$ 的周期函数，它在区间 $(-1,1]$ 上的表达式为
  \[
  f(x)=
  \begin{cases}
  3,&-1<x\le0,\\
  x^3,&0<x\le1,
  \end{cases}
  \]
  则 $f(x)$ 的 Fourier 级数在 $x=1$ 处收敛于 \underline{\hspace{3cm}}。
\end{exercise}

\begin{solution}
\par
傅里叶级数在跳跃点处收敛到左右极限平均值。这里周期为 $2$。当 $x\to1-0$ 时，
\[
  f(x)=x^3\to1.
\]
当 $x\to1+0$ 时，按周期 $2$ 折回到 $x\to-1+0$，而在 $(-1,0]$ 上 $f(x)=3$，所以
\[
  f(1+0)=3.
\]
因此傅里叶级数在 $x=1$ 处收敛于
\[
  \frac{1+3}{2}=2.
\]
这里 \(x=1\) 是基本区间的右端点，右极限必须按周期性折回到 \(-1+0\)，不能只看左侧的 \(x^3\)。
\end{solution}
\begin{exercise}
\par
设 \(\displaystyle L:|x|+|y|=1\)
  取逆时针方向，则
  \(\displaystyle \int_L xy^2\,\mathrm{d}x+x^2y\,\mathrm{d}y=\) \underline{\hspace{3cm}}。
\end{exercise}

\begin{solution}
\par
记
\[
  P=xy^2,\qquad Q=x^2y.
\]
曲线 $L:|x|+|y|=1$ 取逆时针方向，是菱形区域 $D$ 的正向边界。由 Green 公式，
\[
  \int_L P\,\mathrm{d}x+Q\,\mathrm{d}y
  =
  \iint_D(Q_x-P_y)\,\mathrm{d}A.
\]
计算
\[
  Q_x=2xy,\qquad P_y=2xy,
\]
故
\[
  Q_x-P_y=0.
\]
所以
\[
  \int_L xy^2\,\mathrm{d}x+x^2y\,\mathrm{d}y=0.
\]
\end{solution}

\subsubsection*{二、计算题（每小题 10 分，共 40 分）}
\begin{exercise}
\par
设
  \(\displaystyle z=f\!\left(2x,\frac{y}{x}\right)\)，
  其中 $f$ 具有二阶连续的偏导数，求 $\displaystyle \dfrac{\partial^2z}{\partial x\partial y}$。
\end{exercise}

\begin{solution}
\par
令
\[
  \xi=2x,\qquad \eta=\frac{y}{x},
  \qquad z=f(\xi,\eta).
\]
先对 $y$ 求偏导：
\[
  z_y=f_\xi \xi_y+f_\eta \eta_y.
\]
由于 $\xi_y=0,\ \eta_y=\dfrac1x$，所以
\[
  z_y=\frac1x f_\eta.
\]
再对 $x$ 求偏导。注意 $f_\eta$ 是复合函数：
\[
  \frac{\partial}{\partial x}f_\eta(\xi,\eta)
  =
  f_{\eta\xi}\xi_x+f_{\eta\eta}\eta_x
  =
  2f_{\xi\eta}-\frac{y}{x^2}f_{\eta\eta}.
\]
于是
\[
  z_{xy}
  =
  -\frac1{x^2}f_\eta
  +
  \frac1x\left(2f_{\xi\eta}-\frac{y}{x^2}f_{\eta\eta}\right).
\]
即
\[
  \frac{\partial^2z}{\partial x\partial y}
  =
  \frac2x f_{\xi\eta}
  -
  \frac{y}{x^3}f_{\eta\eta}
  -
  \frac1{x^2}f_\eta,
\]
其中所有偏导均在
\[
  \left(2x,\frac{y}{x}\right)
\]
处取值。
\end{solution}
\begin{exercise}
\par
计算三重积分 \(\displaystyle I=\iiint_\Omega z\,\mathrm{d}x\,\mathrm{d}y\,\mathrm{d}z\)，
  其中 $\Omega$ 是由球面 $x^2+y^2+z^2=4$ 与抛物面 $x^2+y^2=3z$ 所围成的区域。
\end{exercise}

\begin{solution}
\par
用柱坐标。球面为
\[
  r^2+z^2=4,
\]
抛物面为
\[
  r^2=3z.
\]
求交线：
\[
  3z+z^2=4
  \quad\Longleftrightarrow\quad
  z=1\quad(\text{取上方交线}),
\]
于是 $r=\sqrt3$。所围区域为
\[
  0\le r\le\sqrt3,\qquad
  \frac{r^2}{3}\le z\le\sqrt{4-r^2},\qquad
  0\le\theta\le2\pi.
\]
因此
\[
  I
  =
  \int_0^{2\pi}\int_0^{\sqrt3}\int_{r^2/3}^{\sqrt{4-r^2}}
  z\,r\,\mathrm{d}z\,\mathrm{d}r\,\mathrm{d}\theta.
\]
先对 $z$ 积分：
\[
  I
  =
  2\pi\int_0^{\sqrt3}
  \frac12\left((4-r^2)-\frac{r^4}{9}\right)r\,\mathrm{d}r.
\]
即
\[
  I
  =
  \pi\int_0^{\sqrt3}
  \left(4r-r^3-\frac{r^5}{9}\right)\mathrm{d}r.
\]
计算得
\[
  I
  =
  \pi\left[
  2r^2-\frac{r^4}{4}-\frac{r^6}{54}
  \right]_0^{\sqrt3}
  =
  \pi\left(6-\frac94-\frac12\right)
  =
  \frac{13\pi}{4}.
\]
\end{solution}
\begin{exercise}
\par
计算曲线积分
  \(\displaystyle I=\int_L\frac{-y\,\mathrm{d}x+x\,\mathrm{d}y}{x^2+4y^2}\)，
  其中 $L$ 是从 $A(-\pi,0)$ 沿曲线 $\displaystyle y=\cos\dfrac{x}{2}$ 到 $B(\pi,0)$ 的有向曲线弧段。
\end{exercise}

\begin{solution}
\par
作线性代换
\[
  u=x,\qquad v=2y.
\]
则
\[
  \mathrm{d}u=\mathrm{d}x,\qquad \mathrm{d}v=2\mathrm{d}y,
\]
并且
\[
  -y\,\mathrm{d}x+x\,\mathrm{d}y
  =
  -\frac v2\,\mathrm{d}u+\frac u2\,\mathrm{d}v
  =
  \frac12(-v\,\mathrm{d}u+u\,\mathrm{d}v).
\]
分母为
\[
  x^2+4y^2=u^2+v^2.
\]
所以
\[
  \frac{-y\,\mathrm{d}x+x\,\mathrm{d}y}{x^2+4y^2}
  =
  \frac12\frac{-v\,\mathrm{d}u+u\,\mathrm{d}v}{u^2+v^2}
  =
  \frac12\,\mathrm{d}\theta,
\]
其中 $\theta$ 是向量 $(u,v)$ 的极角。沿曲线从 $A(-\pi,0)$ 到 $B(\pi,0)$ 时，在 $(u,v)$ 平面中从 $(-\pi,0)$ 经上半平面到 $(\pi,0)$，极角从 $\pi$ 变到 $0$。因此
\[
  I=\frac12\int_{\pi}^{0}\mathrm{d}\theta
  =
  -\frac{\pi}{2}.
\]
\end{solution}
\begin{exercise}
设 $\Sigma$ 是锥面 $z=\sqrt{x^2+y^2}$ 被平面 $z=0$ 及 $z=1$ 所截部分的下侧，计算曲面积分
\[
  I=\iint_\Sigma
  2x\,\mathrm{d}y\,\mathrm{d}z
  +3y\,\mathrm{d}z\,\mathrm{d}x
  +(z^2-5z)\,\mathrm{d}x\,\mathrm{d}y.
\]
\end{exercise}

\begin{solution}
\par
锥面为 $z=r=\sqrt{x^2+y^2}$。作为图形 $z=z(x,y)$，下侧取向的面积向量为
\[
  (z_x,z_y,-1)\,\mathrm{d}x\,\mathrm{d}y
  =
  \left(\frac{x}{r},\frac{y}{r},-1\right)\mathrm{d}x\,\mathrm{d}y.
\]
投影区域为
\[
  D:x^2+y^2\le1.
\]
向量场为
\[
  (P,Q,R)=(2x,\,3y,\,z^2-5z).
\]
因此
\[
  I
  =
  \iint_D
  \left(
  2x\frac{x}{r}+3y\frac{y}{r}-(r^2-5r)
  \right)\mathrm{d}x\,\mathrm{d}y.
\]
换成极坐标 $x=r\cos\theta,\ y=r\sin\theta$，得到
\[
  I
  =
  \int_0^{2\pi}\int_0^1
  \left(2r\cos^2\theta+3r\sin^2\theta-r^2+5r\right)r\,\mathrm{d}r\,\mathrm{d}\theta.
\]
对 $\theta$ 积分：
\[
  \int_0^{2\pi}2r\cos^2\theta\,\mathrm{d}\theta=2\pi r,
  \qquad
  \int_0^{2\pi}3r\sin^2\theta\,\mathrm{d}\theta=3\pi r,
\]
\[
  \int_0^{2\pi}5r\,\mathrm{d}\theta=10\pi r,
  \qquad
  \int_0^{2\pi}(-r^2)\,\mathrm{d}\theta=-2\pi r^2.
\]
再乘外面的 $r$ 并积分：
\[
  I=\int_0^1(15\pi r^2-2\pi r^3)\,\mathrm{d}r
  =
  5\pi-\frac{\pi}{2}
  =
  \frac{9\pi}{2}.
\]
\end{solution}

\subsubsection*{三、解答题（每小题 10 分，共 20 分）}
\begin{exercise}
\par
求幂级数 \(\displaystyle \sum_{n=1}^{\infty}\frac{2n-1}{n!}x^{2n}\) 的和函数。
\end{exercise}

\begin{solution}
\par
令 $t=x^2$。原级数化为
\[
  \sum_{n=1}^{\infty}\frac{2n-1}{n!}t^n
  =
  2\sum_{n=1}^{\infty}\frac{nt^n}{n!}
  -
  \sum_{n=1}^{\infty}\frac{t^n}{n!}.
\]
由指数级数，
\[
  \sum_{n=1}^{\infty}\frac{t^n}{n!}=e^t-1.
\]
又因为 $\dfrac{n}{n!}=\dfrac1{(n-1)!}$，
\[
  \sum_{n=1}^{\infty}\frac{nt^n}{n!}
  =
  t\sum_{n=1}^{\infty}\frac{t^{n-1}}{(n-1)!}
  =
  te^t.
\]
所以
\[
  S=2te^t-(e^t-1)=(2t-1)e^t+1.
\]
换回 $t=x^2$，
\[
  S(x)=(2x^2-1)e^{x^2}+1.
\]
\end{solution}
\begin{exercise}
\par
设 \(\displaystyle f(x)=\sin x-\int_0^x(x-t)f(t)\,\mathrm{d}t\)，其中 $f(x)$ 为连续函数，求 $f(x)$。
\end{exercise}

\begin{solution}
\par
对积分方程
\[
  f(x)=\sin x-\int_0^x(x-t)f(t)\,\mathrm{d}t
\]
求导。由含参上限积分求导法，
\[
  \frac{\mathrm{d}}{\mathrm{d}x}\int_0^x(x-t)f(t)\,\mathrm{d}t
  =
  \int_0^x f(t)\,\mathrm{d}t.
\]
所以
\[
  f'(x)=\cos x-\int_0^x f(t)\,\mathrm{d}t.
\]
再求导：
\[
  f''(x)=-\sin x-f(x),
\]
即
\[
  f''+f=-\sin x.
\]
由原方程代入 $x=0$ 得 $f(0)=0$；由 $f'(x)$ 式代入 $x=0$ 得 $f'(0)=1$。

齐次方程 $f''+f=0$ 的解为 $C_1\cos x+C_2\sin x$。由于右端 $-\sin x$ 与齐次解共振，取特解
\[
  f_p=A x\cos x.
\]
代入得
\[
  f_p''+f_p=-2A\sin x.
\]
令 $-2A=-1$，得 $A=\dfrac12$。因此
\[
  f=C_1\cos x+C_2\sin x+\frac12 x\cos x.
\]
由 $f(0)=0$ 得 $C_1=0$。再求导：
\[
  f'(x)=C_2\cos x+\frac12\cos x-\frac12x\sin x.
\]
代入 $f'(0)=1$，得
\[
  C_2+\frac12=1,\qquad C_2=\frac12.
\]
故
\[
  f(x)=\frac12\sin x+\frac12x\cos x.
\]
\end{solution}

\subsubsection*{四、证明题（本题 10 分）}
\begin{exercise}
\par
证明：级数 \(\displaystyle \sum_{n=1}^{\infty}\frac{\sin(nx)}{n^2}\) 在 $(-\infty,+\infty)$ 上一致收敛。
\end{exercise}

\begin{solution}
\par
对任意 $x\in(-\infty,+\infty)$，都有 $|\sin(nx)|\le1$，所以
\[
  \left|\frac{\sin(nx)}{n^2}\right|\le\frac1{n^2}.
\]
右端与 $x$ 无关，且
\[
  \sum_{n=1}^{\infty}\frac1{n^2}
\]
收敛。由 Weierstrass 判别法，级数
\[
  \sum_{n=1}^{\infty}\frac{\sin(nx)}{n^2}
\]
在 $(-\infty,+\infty)$ 上一致收敛。
\end{solution}

\subsubsection*{五、应用题（本题 10 分）}
\begin{exercise}
\par
求椭球面 \(\displaystyle x^2+y^2+\frac{z^2}{4}=1\) 在第一卦限的一点，使该点处的切平面在三个坐标轴上的截距平方和最小。
\end{exercise}

\begin{solution}
\par
设所求点为 $(x_0,y_0,z_0)$，且在第一卦限，所以三者均为正。椭球面方程为
\[
  x^2+y^2+\frac{z^2}{4}=1.
\]
在 $(x_0,y_0,z_0)$ 处的切平面为
\[
  x_0x+y_0y+\frac{z_0}{4}z=1.
\]
因此三个坐标轴截距分别为
\[
  \frac1{x_0},\qquad \frac1{y_0},\qquad \frac4{z_0}.
\]
要最小化截距平方和
\[
  S=\frac1{x_0^2}+\frac1{y_0^2}+\frac{16}{z_0^2},
\]
同时满足约束
\[
  x_0^2+y_0^2+\frac{z_0^2}{4}=1.
\]
令
\[
  u=x_0^2,\qquad v=y_0^2,\qquad w=\frac{z_0^2}{4}.
\]
则
\[
  u+v+w=1,\qquad
  S=\frac1u+\frac1v+\frac4w.
\]
用拉格朗日乘数法，令
\[
  \mathcal L=\frac1u+\frac1v+\frac4w+\lambda(u+v+w-1).
\]
驻点方程为
\[
  -\frac1{u^2}+\lambda=0,\qquad
  -\frac1{v^2}+\lambda=0,\qquad
  -\frac4{w^2}+\lambda=0.
\]
因此
\[
  u=v,\qquad w=2u.
\]
代入 $u+v+w=1$ 得
\[
  4u=1,\qquad u=v=\frac14,\qquad w=\frac12.
\]
于是
\[
  x_0=\frac12,\qquad y_0=\frac12,\qquad
  \frac{z_0^2}{4}=\frac12.
\]
第一卦限内 $z_0>0$，故
\[
  z_0=\sqrt2.
\]
所求点为
\[
  \left(\frac12,\frac12,\sqrt2\right).
\]
\end{solution}

\section{2015级工科数学分析下 A卷}

\subsubsection*{一、填空题（每小题 3 分，共 15 分）}
\begin{exercise}
\par
设 $y_1=3,\ y_2=3+x^2,\ y_3=3+x^2+e^x$ 都是方程
  \(\displaystyle (x^2-2x)y''-(x^2-2)y'+(2x-2)y=6x-6\)
  的解，则该方程的通解为 \underline{\hspace{4cm}}。
\end{exercise}

\begin{solution}
\par
这是线性非齐次方程。若 $y_1,y_2,y_3$ 都是同一个非齐次方程的解，则任意两个解的差是对应齐次方程的解。这里
\[
  y_2-y_1=x^2,\qquad y_3-y_2=e^x.
\]
两者显然线性无关，所以它们构成对应齐次方程的一组基本解。又 $y_1=3$ 是一个非齐次方程特解，因此原方程通解为
\[
  y=3+C_1x^2+C_2e^x.
\]
这类题的基本处理是“一个非齐次特解 + 齐次通解”。多个非齐次解本身不能直接线性组合，但它们的差一定回到齐次方程。
\end{solution}
\begin{exercise}
\par
设由方程 $F(x-y,y-z,z-x)=0$ 确定了 $z=z(x,y)$，其中 $F$ 可微，则 $\mathrm{d}z=$ \underline{\hspace{3cm}}。
\end{exercise}

\begin{solution}
\par
令
\[
  u=x-y,\qquad v=y-z,\qquad w=z-x.
\]
记 $F_1,F_2,F_3$ 分别为 $F$ 对第一、第二、第三个自变量的偏导数。由
\[
  F(u,v,w)=0
\]
取全微分：
\[
  F_1\,\mathrm{d}u+F_2\,\mathrm{d}v+F_3\,\mathrm{d}w=0.
\]
又
\[
  \mathrm{d}u=\mathrm{d}x-\mathrm{d}y,\qquad \mathrm{d}v=\mathrm{d}y-\mathrm{d}z,\qquad \mathrm{d}w=\mathrm{d}z-\mathrm{d}x.
\]
代入并合并同类项：
\[
  (F_1-F_3)\mathrm{d}x+(F_2-F_1)\mathrm{d}y+(F_3-F_2)\mathrm{d}z=0.
\]
当 $F_3\ne F_2$ 时，
\[
  \mathrm{d}z=\frac{(F_3-F_1)\,\mathrm{d}x+(F_1-F_2)\,\mathrm{d}y}{F_3-F_2}.
\]
\end{solution}
\begin{exercise}
\par
设函数
  \[
    f(x,y,z)=\ln(x^2+y^2+z^2).
  \]
  其在点 \(P(2,0,1)\) 处沿梯度方向的方向导数是 \underline{\hspace{4cm}}。
\end{exercise}

\begin{solution}
\par
先求梯度：
\[
  \nabla f
  =
  \left(
  \frac{2x}{x^2+y^2+z^2},
  \frac{2y}{x^2+y^2+z^2},
  \frac{2z}{x^2+y^2+z^2}
  \right).
\]
在 $P(2,0,1)$ 处，
\[
  x^2+y^2+z^2=5,
\]
所以
\[
  \nabla f(P)=\left(\frac45,0,\frac25\right).
\]
沿梯度方向的方向导数等于梯度模长：
\[
  |\nabla f(P)|
  =
  \sqrt{\left(\frac45\right)^2+\left(\frac25\right)^2}
  =
  \frac{2\sqrt5}{5}.
\]
\end{solution}
\begin{exercise}
\par
设 $\displaystyle L:\dfrac{x^2}{a^2}+\dfrac{y^2}{b^2}=1$ 取顺时针方向，则 \(\displaystyle \int_L(x-y)\,\mathrm{d}x+(x+y)\,\mathrm{d}y=\) \underline{\hspace{3cm}}。
\end{exercise}

\begin{solution}
\par
记
\[
  P=x-y,\qquad Q=x+y.
\]
若 $L$ 取逆时针方向，则由 Green 公式
\[
  \int_L P\,\mathrm{d}x+Q\,\mathrm{d}y
  =
  \iint_D(Q_x-P_y)\,\mathrm{d}A.
\]
计算
\[
  Q_x=1,\qquad P_y=-1,
\]
故
\[
  Q_x-P_y=2.
\]
椭圆内部面积为 $\pi ab$，所以逆时针方向积分为
\[
  2\pi ab.
\]
题中 $L$ 取顺时针方向，方向相反，积分变号，故结果为
\[
  -2\pi ab.
\]
\end{solution}
\begin{exercise}
\par
设 \(f(x)=\sin(\pi x^2)\ (0\le x<1)\)，而
  \[
  S(x)=\sum_{n=1}^{\infty}b_n\sin(n\pi x)\quad(-\infty<x<+\infty),
  \qquad
  b_n=2\int_0^1 f(x)\sin(n\pi x)\,\mathrm{d}x\ (n=1,2,\cdots).
  \]
  则 \(\displaystyle S\!\left(-\dfrac12\right)\) \underline{\hspace{3cm}}。
\end{exercise}

\begin{solution}
\par
这是 $[0,1]$ 上函数 $f(x)=\sin(\pi x^2)$ 的 Fourier 正弦级数。正弦级数对应奇延拓，并以周期 $2$ 延拓到全直线。点 $-\dfrac12$ 是连续点，因此
\[
  S\!\left(-\frac12\right)
  =
  -f\!\left(\frac12\right)
  =
  -\sin\left(\pi\cdot\frac14\right)
  =
  -\frac{\sqrt2}{2}.
\]
这里取负号来自奇延拓；因为 \(-\frac12\) 不是端点或跳跃点，所以和函数就等于延拓后的函数值。
\end{solution}

\subsubsection*{二、计算题（每小题 8 分，共 40 分）}
\begin{exercise}
\par
设函数 $z=f(x^2-y^2,2xy)$，其中 $f(\xi,\eta)$ 具有连续的二阶偏导数，且满足
  \(\displaystyle \frac{\partial^2 f}{\partial \xi^2}+\frac{\partial^2 f}{\partial \eta^2}=0\)，
  求 $\displaystyle \dfrac{\partial^2 z}{\partial x^2}+\dfrac{\partial^2 z}{\partial y^2}$。
\end{exercise}

\begin{solution}
\par
令
\[
  \xi=x^2-y^2,\qquad \eta=2xy,\qquad z=f(\xi,\eta).
\]
先求一阶偏导：
\[
  \xi_x=2x,\quad \eta_x=2y,\qquad
  \xi_y=-2y,\quad \eta_y=2x.
\]
由链式法则，
\[
  z_x=2x f_\xi+2y f_\eta,
  \qquad
  z_y=-2y f_\xi+2x f_\eta.
\]
继续求二阶偏导并相加，可得
\[
  z_{xx}+z_{yy}
  =
  4(x^2+y^2)(f_{\xi\xi}+f_{\eta\eta}).
\]
题设给出
\[
  f_{\xi\xi}+f_{\eta\eta}=0,
\]
所以
\[
  \frac{\partial^2 z}{\partial x^2}
  +
  \frac{\partial^2 z}{\partial y^2}
  =
  0.
\]
\end{solution}
\begin{exercise}
\par
计算曲线积分 \(\displaystyle I=\int_L\sqrt{x^2+y^2}\,\mathrm{d}s\)，
  其中 $L$ 是圆周 $x^2+y^2=ax$。
\end{exercise}

\begin{solution}
\par
圆周方程可配方为
\[
  \left(x-\frac a2\right)^2+y^2=\left(\frac a2\right)^2.
\]
取参数
\[
  x=\frac a2(1+\cos t),\qquad
  y=\frac a2\sin t,\qquad 0\le t\le2\pi.
\]
则
\[
  \mathrm{d}s=\sqrt{x'(t)^2+y'(t)^2}\,\mathrm{d}t=\frac a2\,\mathrm{d}t.
\]
又
\[
  x^2+y^2=ax
  =
  \frac{a^2}{2}(1+\cos t)
  =
  a^2\cos^2\frac t2,
\]
故
\[
  \sqrt{x^2+y^2}=a\left|\cos\frac t2\right|.
\]
于是
\[
  I
  =
  \frac{a^2}{2}
  \int_0^{2\pi}\left|\cos\frac t2\right|\mathrm{d}t.
\]
令 $u=\dfrac t2$，则
\[
  \int_0^{2\pi}\left|\cos\frac t2\right|\mathrm{d}t
  =
  2\int_0^\pi |\cos u|\,\mathrm{d}u
  =
  4.
\]
所以
\[
  I=2a^2.
\]
\end{solution}
\begin{exercise}
\par
设 $f(x)$ 在 $(-\infty,+\infty)$ 内有连续导数，$L$ 是从点 $\displaystyle A\!\left(3,\dfrac23\right)$ 到点 $B(1,2)$ 的直线段，计算曲线积分 \(\displaystyle I=\int_L\frac{1+y^2f(x)}{y}\,\mathrm{d}x+\frac{x}{y^2}\bigl[y^2f(x)-1\bigr]\,\mathrm{d}y\)。
\end{exercise}

\begin{solution}
\par
先把微分形式整理：
\[
  \frac{1+y^2f(x)}{y}\,\mathrm{d}x+\frac{x}{y^2}\bigl[y^2f(x)-1\bigr]\,\mathrm{d}y
  =
  \left(\frac1y+yf(x)\right)\mathrm{d}x+\left(xf(x)-\frac{x}{y^2}\right)\mathrm{d}y.
\]
注意
\[
  \mathrm{d}\left(\frac{x}{y}\right)=\frac1y\,\mathrm{d}x-\frac{x}{y^2}\,\mathrm{d}y,
\]
且
\[
  f(x)\,\mathrm{d}(xy)=f(x)(y\,\mathrm{d}x+x\,\mathrm{d}y).
\]
所以原积分为
\[
  I=\int_L \mathrm{d}\left(\frac{x}{y}\right)+\int_L f(x)\,\mathrm{d}(xy).
\]
直线段从 $A\!\left(3,\dfrac23\right)$ 到 $B(1,2)$，其方程为
\[
  y=\frac{8-2x}{3},\qquad x:3\to1.
\]
于是
\[
  xy=\frac{x(8-2x)}3,
  \qquad
  \mathrm{d}(xy)=\frac{8-4x}{3}\,\mathrm{d}x.
\]
第一项为端点值：
\[
  \left.\frac{x}{y}\right|_A^B
  =
  \frac12-\frac{3}{2/3}
  =
  -4.
\]
第二项化为单变量积分：
\[
  \int_L f(x)\,\mathrm{d}(xy)
  =
  \int_3^1 f(x)\frac{8-4x}{3}\,\mathrm{d}x.
\]
故
\[
  I=-4+\frac13\int_3^1(8-4x)f(x)\,\mathrm{d}x.
\]
\end{solution}
\begin{exercise}
\par
设 $\Sigma$ 为曲面 $x^2+y^2+z^2=2ax\ (a>0)$，计算曲面积分
  \(\displaystyle I=\iint_\Sigma (x^2+y^2+z^2)\,\mathrm{d}S\)。
\end{exercise}

\begin{solution}
\par
将曲面方程配方：
\[
  x^2-2ax+y^2+z^2=0
  \quad\Longleftrightarrow\quad
  (x-a)^2+y^2+z^2=a^2.
\]
所以 $\Sigma$ 是以 $(a,0,0)$ 为球心、半径为 $a$ 的球面。在曲面上有
\[
  x^2+y^2+z^2=2ax.
\]
因此
\[
  I=2a\iint_\Sigma x\,\mathrm{d}S.
\]
令 $X=x-a$，球面对球心对称，故
\[
  \iint_\Sigma X\,\mathrm{d}S=0.
\]
于是
\[
  \iint_\Sigma x\,\mathrm{d}S
  =
  \iint_\Sigma(a+X)\,\mathrm{d}S
  =
  a\cdot4\pi a^2.
\]
所以
\[
  I=2a\cdot a\cdot4\pi a^2=8\pi a^4.
\]
\end{solution}
\begin{exercise}
\par
设 $\Sigma$ 是上半球面 $z=\sqrt{a^2-x^2-y^2}\ (a>0)$，取上侧，计算曲面积分 \(\displaystyle I=\iint_\Sigma x^3\,\mathrm{d}y\,\mathrm{d}z+y^3\,\mathrm{d}z\,\mathrm{d}x+z^3\,\mathrm{d}x\,\mathrm{d}y\)。
\end{exercise}

\begin{solution}
\par
把上半球面 $\Sigma$ 与底面圆盘 $D:z=0$ 合成上半球体 $\Omega$ 的闭曲面。向量场为
\[
  \mathbf F=(x^3,y^3,z^3).
\]
其散度为
\[
  \nabla\cdot\mathbf F=3x^2+3y^2+3z^2=3(x^2+y^2+z^2).
\]
底面上 $z=0$，且底面外法向量为 $(0,0,-1)$，底面通量为 $-z^3=0$。因此上半球面的通量等于闭曲面通量。由 Gauss 公式，
\[
  I=\iiint_\Omega 3(x^2+y^2+z^2)\,\mathrm{d}V.
\]
用球坐标，在上半球体中角向面积为 $2\pi$，所以
\[
  I
  =
  3\int_0^a r^2\cdot 2\pi r^2\,\mathrm{d}r
  =
  6\pi\int_0^a r^4\,\mathrm{d}r
  =
  \frac{6\pi a^5}{5}.
\]
\end{solution}

\subsubsection*{三、解答题（每小题 9 分，共 18 分）}
\begin{exercise}
\par
求幂级数 \(\displaystyle \sum_{n=1}^{\infty}\frac{2n+1}{n!}x^{2n}\) 的和函数。
\end{exercise}

\begin{solution}
\par
令 $t=x^2$。则
\[
  S=\sum_{n=1}^{\infty}\frac{2n+1}{n!}t^n
  =
  2\sum_{n=1}^{\infty}\frac{nt^n}{n!}
  +
  \sum_{n=1}^{\infty}\frac{t^n}{n!}.
\]
由指数级数，
\[
  \sum_{n=1}^{\infty}\frac{t^n}{n!}=e^t-1.
\]
又
\[
  \sum_{n=1}^{\infty}\frac{nt^n}{n!}
  =
  t\sum_{n=1}^{\infty}\frac{t^{n-1}}{(n-1)!}
  =
  te^t.
\]
所以
\[
  S=2te^t+e^t-1=(2t+1)e^t-1.
\]
换回 $t=x^2$，得
\[
  S(x)=(2x^2+1)e^{x^2}-1.
\]
\end{solution}
\begin{exercise}
\par
求微分方程 \(\displaystyle xy'-2y=x^3\sin x\) 的通解。
\end{exercise}

\begin{solution}
\par
当 $x\ne0$ 时，将方程除以 $x$：
\[
  y'-\frac2x y=x^2\sin x.
\]
这是一阶线性方程。积分因子为
\[
  \mu(x)=e^{\int -2/x\,\mathrm{d}x}=x^{-2}.
\]
两边乘以 $x^{-2}$：
\[
  x^{-2}y'-2x^{-3}y=\sin x.
\]
左端正是
\[
  \left(\frac{y}{x^2}\right)'.
\]
因此
\[
  \left(\frac{y}{x^2}\right)'=\sin x.
\]
积分得
\[
  \frac{y}{x^2}=-\cos x+C.
\]
故通解为
\[
  y=Cx^2-x^2\cos x.
\]
\end{solution}

\subsubsection*{四、证明题（每小题 9 分，共 18 分）}
\begin{exercise}
\par
设 $f(x)\in C[0,1]$，求证：\(\displaystyle \int_0^1 \mathrm{d}x\int_x^1 f(x)f(y)\,\mathrm{d}y=\frac12\left[\int_0^1 f(x)\,\mathrm{d}x\right]^2\)。
\end{exercise}

\begin{solution}
\par
记
\[
  I=\int_0^1 \mathrm{d}x\int_x^1 f(x)f(y)\,\mathrm{d}y.
\]
积分区域为
\[
  0\le x\le y\le1.
\]
由于 $f(x)f(y)$ 关于 $x,y$ 对称，交换变量 $x,y$ 后，下三角区域上的积分与上三角区域上的积分相同：
\[
  I=\int_0^1\mathrm{d}y\int_0^y f(x)f(y)\,\mathrm{d}x.
\]
两个三角区域合起来就是单位正方形，除去的对角线面积为零。因此
\[
  2I
  =
  \int_0^1\int_0^1 f(x)f(y)\,\mathrm{d}x\,\mathrm{d}y.
\]
右端可分离：
\[
  \int_0^1\int_0^1 f(x)f(y)\,\mathrm{d}x\,\mathrm{d}y
  =
  \left(\int_0^1 f(x)\,\mathrm{d}x\right)
  \left(\int_0^1 f(y)\,\mathrm{d}y\right).
\]
故
\[
  I=\frac12\left[\int_0^1 f(x)\,\mathrm{d}x\right]^2.
\]
结论得证。
\end{solution}
\begin{exercise}
\par
证明：函数项级数
  \(\displaystyle \sum_{n=1}^{\infty}\frac{\cos(nx)}{n^2}\)
  在 $(-\infty,+\infty)$ 上一致收敛。
\end{exercise}

\begin{solution}
\par
对任意实数 $x$，有 $|\cos(nx)|\le1$，所以
\[
  \left|\frac{\cos(nx)}{n^2}\right|
  \le
  \frac1{n^2}.
\]
右端与 $x$ 无关，且数项级数
\[
  \sum_{n=1}^{\infty}\frac1{n^2}
\]
收敛。由 Weierstrass 判别法，函数项级数
\[
  \sum_{n=1}^{\infty}\frac{\cos(nx)}{n^2}
\]
在 $(-\infty,+\infty)$ 上一致收敛。
\end{solution}

\subsubsection*{五、应用题（本题 9 分）}
\begin{exercise}
\par
求内接于半径为 $a$ 的球且有最大体积的长方体。
\end{exercise}

\begin{solution}
\par
最大体积的内接长方体可取球心为中心。设三个半棱长为
\[
  x>0,\qquad y>0,\qquad z>0.
\]
顶点在球面上时体积最大，因此约束为
\[
  x^2+y^2+z^2=a^2.
\]
体积为
\[
  V=8xyz.
\]
用拉格朗日乘数法最大化 $\ln x+\ln y+\ln z$，令
\[
  \mathcal L=\ln x+\ln y+\ln z-\lambda(x^2+y^2+z^2-a^2).
\]
驻点方程给出
\[
  \frac1x=2\lambda x,\qquad
  \frac1y=2\lambda y,\qquad
  \frac1z=2\lambda z.
\]
因此
\[
  x=y=z.
\]
代入约束：
\[
  3x^2=a^2,\qquad x=y=z=\frac a{\sqrt3}.
\]
所以最大体积长方体为正方体，棱长为
\[
  \frac{2a}{\sqrt3},
\]
最大体积为
\[
  V_{\max}
  =
  \left(\frac{2a}{\sqrt3}\right)^3
  =
  \frac{8a^3}{3\sqrt3}.
\]
\end{solution}

\section{2015级工科数学分析下 B卷}

\subsubsection*{一、填空题（每小题 3 分，共 15 分）}
\begin{exercise}
\par
微分方程 \(\displaystyle y''-4y=xe^{2x}\) 的特解形式为 \underline{\hspace{4cm}}。
\end{exercise}

\begin{solution}
\par
对应齐次方程的特征方程为
\[
  r^2-4=0,\qquad r=\pm2.
\]
右端 $xe^{2x}$ 是一次多项式乘 $e^{2x}$。若没有共振，特解形式应为 $e^{2x}(Ax+B)$；但 $r=2$ 是齐次方程的一重特征根，所以要再乘一个 $x$。因此特解形式为
\[
  y^*=xe^{2x}(Ax+B).
\]
这里括号内仍取一次多项式 \(Ax+B\)，因为右端的多项式因子是 \(x\)；前面的 \(x\) 只是因为 \(r=2\) 与齐次解共振而额外补上的因子。
\end{solution}
\begin{exercise}
\par
设 $z=z(x,y)$ 满足方程
  \(\displaystyle x-az=\varphi(y-bz)\)，
  其中 $\varphi$ 可微，$a,b$ 为常数，则
  \(\displaystyle a\frac{\partial z}{\partial x}+b\frac{\partial z}{\partial y}=\) \underline{\hspace{3cm}}。
\end{exercise}

\begin{solution}
\par
令
\[
  u=y-bz.
\]
题设方程为
\[
  x-az=\varphi(u).
\]
对 $x$ 求偏导时，$y$ 为常数：
\[
  1-a z_x=\varphi'(u)\cdot(-b z_x).
\]
即
\[
  1=(a-b\varphi'(u))z_x.
\]
对 $y$ 求偏导时，$x$ 为常数：
\[
  -a z_y=\varphi'(u)(1-bz_y).
\]
整理得
\[
  -\varphi'(u)=(a-b\varphi'(u))z_y.
\]
于是
\[
  a z_x+b z_y
  =
  \frac{a}{a-b\varphi'(u)}
  -
  \frac{b\varphi'(u)}{a-b\varphi'(u)}
  =
  1.
\]
\end{solution}
\begin{exercise}
\par
设
\[
  f(x,y,z)=\cos(xyz).
\]
其在点 $\displaystyle P\!\left(\dfrac13,1,\pi\right)$ 处使方向导数取得最大的方向是 \underline{\hspace{4cm}}。
\end{exercise}

\begin{solution}
\par
方向导数在梯度方向取得最大值。先求梯度：
\[
  \nabla f
  =
  -\sin(xyz)\,(yz,xz,xy).
\]
在
\[
  P\!\left(\frac13,1,\pi\right)
\]
处，
\[
  xyz=\frac{\pi}{3},\qquad
  \sin(xyz)=\sin\frac{\pi}{3}=\frac{\sqrt3}{2}.
\]
并且
\[
  (yz,xz,xy)=\left(\pi,\frac{\pi}{3},\frac13\right)
  =
  \frac13(3\pi,\pi,1).
\]
所以
\[
  \nabla f(P)
  =
  -\frac{\sqrt3}{6}(3\pi,\pi,1).
\]
最大方向为其单位方向，即可取
\[
  -\frac{(3\pi,\pi,1)}{\sqrt{10\pi^2+1}}.
\]
\end{solution}
\begin{exercise}
\par
设 $L:x^2+y^2=a^2$ 取逆时针方向，则
  \(\displaystyle \int_L 2xy\,\mathrm{d}x+(x^2+8y^2)\,\mathrm{d}y=\) \underline{\hspace{3cm}}。
\end{exercise}

\begin{solution}
\par
记
\[
  P=2xy,\qquad Q=x^2+8y^2.
\]
圆 $L$ 取逆时针方向，是圆盘 $D:x^2+y^2\le a^2$ 的正向边界。由 Green 公式，
\[
  \int_LP\,\mathrm{d}x+Q\,\mathrm{d}y
  =
  \iint_D(Q_x-P_y)\,\mathrm{d}A.
\]
计算
\[
  Q_x=2x,\qquad P_y=2x.
\]
故
\[
  Q_x-P_y=0,
\]
所以积分为
\[
  0.
\]
这里圆盘内没有奇点，且被积形式的 Green 旋度处处为零，所以闭曲线积分直接为零。
\end{solution}
\begin{exercise}
\par
设幂级数 \(\displaystyle \sum_{n=0}^{\infty}a_nx^n\)
  在 $x=-3$ 条件收敛，则该幂级数的收敛半径为 $R=$ \underline{\hspace{2cm}}。
\end{exercise}

\begin{solution}
\par
幂级数
\[
  \sum_{n=0}^{\infty}a_nx^n
\]
以 $0$ 为中心。幂级数在收敛区间内部一定绝对收敛；只有端点才可能出现条件收敛。题设在 $x=-3$ 条件收敛，说明 $x=-3$ 是收敛区间端点。它到中心 $0$ 的距离为 $3$，因此收敛半径为
\[
  R=3.
\]
若 \(R>3\)，则 \(x=-3\) 位于收敛区间内部，应绝对收敛；若 \(R<3\)，则 \(x=-3\) 应发散。题设为条件收敛，所以只能是端点 \(R=3\)。
\end{solution}

\subsubsection*{二、计算题（每小题 8 分，共 40 分）}
\begin{exercise}
\par
设函数 \(\displaystyle z=f\!\left(xy,\frac{x}{y}\right)\)，
  其中 $f(\xi,\eta)$ 具有连续的二阶偏导数，求 $\displaystyle \dfrac{\partial^2z}{\partial x\partial y}$。
\end{exercise}

\begin{solution}
\par
令
\[
  \xi=xy,\qquad \eta=\frac{x}{y},\qquad z=f(\xi,\eta).
\]
先对 $x$ 求偏导：
\[
  z_x=f_\xi\xi_x+f_\eta\eta_x
  =
  yf_\xi+\frac1y f_\eta.
\]
再对 $y$ 求偏导。第一项：
\[
  \frac{\partial}{\partial y}(yf_\xi)
  =
  f_\xi
  +
  y(f_{\xi\xi}\xi_y+f_{\xi\eta}\eta_y).
\]
其中
\[
  \xi_y=x,\qquad \eta_y=-\frac{x}{y^2}.
\]
所以第一项为
\[
  f_\xi+xyf_{\xi\xi}-\frac{x}{y}f_{\xi\eta}.
\]
第二项：
\[
  \frac{\partial}{\partial y}\left(\frac1y f_\eta\right)
  =
  -\frac1{y^2}f_\eta
  +
  \frac1y(f_{\eta\xi}\xi_y+f_{\eta\eta}\eta_y)
\]
\[
  =
  -\frac1{y^2}f_\eta
  +
  \frac{x}{y}f_{\xi\eta}
  -
  \frac{x}{y^3}f_{\eta\eta}.
\]
混合偏导项相消，故
\[
  z_{xy}
  =
  f_\xi+xyf_{\xi\xi}
  -
  \frac{x}{y^3}f_{\eta\eta}
  -
  \frac1{y^2}f_\eta.
\]
各偏导均在
\[
  \left(xy,\frac{x}{y}\right)
\]
处取值。
\end{solution}
\begin{exercise}
\par
计算曲线积分 \(\displaystyle I=\int_\Gamma x^2\,\mathrm{d}s\)，
  其中 $\Gamma$ 是球面 $x^2+y^2+z^2=a^2$ 与平面 $x+y+z=0$ 的交线。
\end{exercise}

\begin{solution}
\par
平面 $x+y+z=0$ 过球心，所以交线 $\Gamma$ 是半径为 $a$ 的大圆，弧长为 $2\pi a$。在该大圆上
\[
  x^2+y^2+z^2=a^2.
\]
由于方程对 $x,y,z$ 对称，沿这条大圆积分时
\[
  \int_\Gamma x^2\,\mathrm{d}s
  =
  \int_\Gamma y^2\,\mathrm{d}s
  =
  \int_\Gamma z^2\,\mathrm{d}s.
\]
三式相加：
\[
  3\int_\Gamma x^2\,\mathrm{d}s
  =
  \int_\Gamma (x^2+y^2+z^2)\,\mathrm{d}s
  =
  a^2\int_\Gamma \mathrm{d}s
  =
  a^2\cdot2\pi a.
\]
因此
\[
  I=\int_\Gamma x^2\,\mathrm{d}s=\frac{2\pi a^3}{3}.
\]
\end{solution}
\begin{exercise}
\par
设曲线积分 \(\displaystyle \int_L\bigl(f(x)-e^x\bigr)\sin y\,\mathrm{d}x-f'(x)\cos y\,\mathrm{d}y\) 与路径无关，其中 $f'(x)$ 有一阶的连续导数，且 $f'(0)=0$。

  （1）求 $f'(x)$；（2）计算曲线积分 \(\displaystyle I=\int_{(0,0)}^{(1,1)}\bigl(f(x)-e^x\bigr)\sin y\,\mathrm{d}x-f'(x)\cos y\,\mathrm{d}y\)。
\end{exercise}

\begin{solution}
\par
记
\[
  P=(f(x)-e^x)\sin y,\qquad Q=-f'(x)\cos y.
\]
曲线积分与路径无关要求
\[
  P_y=Q_x.
\]
计算得
\[
  P_y=(f(x)-e^x)\cos y,\qquad
  Q_x=-f''(x)\cos y.
\]
所以
\[
  f(x)-e^x=-f''(x),
  \qquad
  f''(x)+f(x)=e^x.
\]
解这个常系数非齐次方程：
\[
  f(x)=C_1\cos x+C_2\sin x+\frac12e^x.
\]
因此
\[
  f'(x)=-C_1\sin x+C_2\cos x+\frac12e^x.
\]
由 $f'(0)=0$ 得
\[
  C_2+\frac12=0,\qquad C_2=-\frac12.
\]
所以
\[
  f'(x)=-C_1\sin x-\frac12\cos x+\frac12e^x.
\]
再求势函数。注意
\[
  d[-f'(x)\sin y]
  =
  -f''(x)\sin y\,\mathrm{d}x-f'(x)\cos y\,\mathrm{d}y.
\]
由 $f''=e^x-f$，有
\[
  -f''(x)\sin y=(f(x)-e^x)\sin y.
\]
故原微分形式就是
\[
  d[-f'(x)\sin y].
\]
于是
\[
  I
  =
  [-f'(x)\sin y]_{(0,0)}^{(1,1)}
  =
  -f'(1)\sin1.
\]
代入 $f'(1)$：
\[
  I
  =
  \left(C_1\sin1+\frac12\cos1-\frac12e\right)\sin1
  =
  C_1\sin^2 1+\frac12(\cos1-e)\sin1.
\]
按当前题面，只给出 $f'(0)=0$，只能确定 $C_2$，不能确定 $C_1$。因此积分值不能唯一确定；若原题需要唯一数值，还必须再给一个关于 $f$ 或 $f'$ 的独立条件。
\end{solution}
\begin{exercise}
\par
设 $\Sigma$ 为平面 $y+z=5$ 被柱面 $x^2+y^2=25$ 所截得的部分，计算曲面积分
  \(\displaystyle I=\iint_\Sigma (x+y+z)\,\mathrm{d}S\)。
\end{exercise}

\begin{solution}
\par
将平面写成图形
\[
  z=5-y.
\]
则
\[
  z_x=0,\qquad z_y=-1,
\]
所以
\[
  \mathrm{d}S=\sqrt{1+z_x^2+z_y^2}\,\mathrm{d}x\,\mathrm{d}y=\sqrt2\,\mathrm{d}x\,\mathrm{d}y.
\]
该平面在 $xy$ 平面上的投影区域由柱面给出：
\[
  D:x^2+y^2\le25.
\]
在平面上
\[
  x+y+z=x+y+5-y=x+5.
\]
因此
\[
  I
  =
  \sqrt2\iint_D(x+5)\,\mathrm{d}x\,\mathrm{d}y.
\]
圆盘 $D$ 关于 $y$ 轴对称，故
\[
  \iint_D x\,\mathrm{d}x\,\mathrm{d}y=0.
\]
而 $D$ 的面积为 $25\pi$，所以
\[
  I=\sqrt2\cdot5\cdot25\pi=125\sqrt2\,\pi.
\]
\end{solution}
\begin{exercise}
\par
设 $\Sigma$ 是上半球面 $z=\sqrt{a^2-x^2-y^2}\ (a>0)$，取上侧，计算曲面积分
  \[
  I=\iint_\Sigma (x^3-xy^2)\,\mathrm{d}y\,\mathrm{d}z+(y^3-yz^2)\,\mathrm{d}z\,\mathrm{d}x+(z^3-zx^2)\,\mathrm{d}x\,\mathrm{d}y.
  \]
\end{exercise}

\begin{solution}
\par
令
\[
  \mathbf F=(x^3-xy^2,\ y^3-yz^2,\ z^3-zx^2).
\]
则题中第二类曲面积分就是 $\mathbf F$ 通过上半球面的通量。先求散度：
\[
  \nabla\cdot\mathbf F
  =
  (3x^2-y^2)+(3y^2-z^2)+(3z^2-x^2)
  =
  2(x^2+y^2+z^2).
\]
把上半球面与底面圆盘合成上半球体 $\Omega$ 的闭曲面。底面 $z=0$ 的外法向量为 $(0,0,-1)$，底面通量只涉及第三分量；而第三分量
\[
  z^3-zx^2
\]
在 $z=0$ 上为 $0$，所以底面通量为 $0$。由 Gauss 公式，
\[
  I=\iiint_\Omega 2(x^2+y^2+z^2)\,\mathrm{d}V.
\]
用球坐标，在上半球体内角向面积为 $2\pi$，故
\[
  I
  =
  2\int_0^a r^2\cdot2\pi r^2\,\mathrm{d}r
  =
  4\pi\int_0^a r^4\,\mathrm{d}r
  =
  \frac{4\pi a^5}{5}.
\]
\end{solution}

\subsubsection*{三、解答题（每小题 9 分，共 18 分）}
\begin{exercise}
\par
将函数 \(\displaystyle \frac{\mathrm{d}}{\mathrm{d}x}\left(\frac{e^x-1}{x}\right)\)
  展开成 $x$ 的幂级数，并求数项级数
  \(\displaystyle \sum_{n=1}^{\infty}\frac{n}{(n+1)!}\) 的和。
\end{exercise}

\begin{solution}
\par
由指数函数展开
\[
  e^x=\sum_{n=0}^{\infty}\frac{x^n}{n!},
\]
得
\[
  \frac{e^x-1}{x}
  =
  \sum_{n=1}^{\infty}\frac{x^{n-1}}{n!}
  =
  \sum_{n=0}^{\infty}\frac{x^n}{(n+1)!}.
\]
逐项求导：
\[
  \frac{\mathrm{d}}{\mathrm{d}x}\left(\frac{e^x-1}{x}\right)
  =
  \sum_{n=1}^{\infty}\frac{n}{(n+1)!}x^{n-1}.
\]
另一方面，直接对闭式求导：
\[
  \frac{\mathrm{d}}{\mathrm{d}x}\left(\frac{e^x-1}{x}\right)
  =
  \frac{xe^x-(e^x-1)}{x^2}
  =
  \frac{(x-1)e^x+1}{x^2}.
\]
令 $x=1$，左边级数变为
\[
  \sum_{n=1}^{\infty}\frac{n}{(n+1)!},
\]
右边为 $1$。因此
\[
  \sum_{n=1}^{\infty}\frac{n}{(n+1)!}=1.
\]
\end{solution}
\begin{exercise}
\par
求微分方程 \(\displaystyle y'\tan x=y\ln y\) 的通解。
\end{exercise}

\begin{solution}
\par
因方程中含有 $\ln y$，先在 $y>0$ 的区间内讨论。原方程为
\[
  y'\tan x=y\ln y.
\]
当 $\tan x\ne0$ 且 $y\ln y\ne0$ 时分离变量：
\[
  \frac{\mathrm{d}y}{y\ln y}=\cot x\,\mathrm{d}x.
\]
积分：
\[
  \int\frac{\mathrm{d}y}{y\ln y}
  =
  \ln|\ln y|,
  \qquad
  \int\cot x\,\mathrm{d}x=\ln|\sin x|.
\]
所以
\[
  \ln|\ln y|=\ln|\sin x|+C.
\]
指数化后可写成
\[
  \ln y=C\sin x.
\]
因此
\[
  y=e^{C\sin x}.
\]
当 $C=0$ 时得到常值解 $y=1$，它也满足原方程。
\end{solution}

\subsubsection*{四、证明题（每小题 9 分，共 18 分）}
\begin{exercise}
\par
设 $f(x)\in C[a,b]$，证明：
  \(\displaystyle \int_a^b \mathrm{d}x\int_a^x f(y)\,\mathrm{d}y=\int_a^b f(y)(b-y)\,\mathrm{d}y\)。
\end{exercise}

\begin{solution}
\par
左端积分区域为
\[
  a\le y\le x\le b.
\]
若改为先固定 $y$，则 $y$ 从 $a$ 到 $b$，而 $x$ 从 $y$ 到 $b$。因此
\[
  \int_a^b \mathrm{d}x\int_a^x f(y)\,\mathrm{d}y
  =
  \int_a^b \mathrm{d}y\int_y^b f(y)\,\mathrm{d}x.
\]
由于内层积分中 $f(y)$ 与 $x$ 无关，
\[
  \int_y^b f(y)\,\mathrm{d}x=f(y)(b-y).
\]
故
\[
  \int_a^b \mathrm{d}x\int_a^x f(y)\,\mathrm{d}y
  =
  \int_a^b f(y)(b-y)\,\mathrm{d}y.
\]
结论得证。
\end{solution}
\begin{exercise}
\par
证明：函数项级数
  \(\displaystyle \sum_{n=1}^{\infty}\frac{(-1)^n(1-e^{-nx})}{n^2+x^2}\)
  在 $[0,+\infty)$ 上一致收敛。
\end{exercise}

\begin{solution}
\par
对任意 $x\in[0,+\infty)$，有
\[
  0\le e^{-nx}\le1,
  \qquad
  0\le1-e^{-nx}\le1.
\]
又
\[
  n^2+x^2\ge n^2.
\]
因此
\[
  \left|
  \frac{(-1)^n(1-e^{-nx})}{n^2+x^2}
  \right|
  =
  \frac{1-e^{-nx}}{n^2+x^2}
  \le
  \frac1{n^2}.
\]
数项级数 $\sum_{n=1}^{\infty}\dfrac1{n^2}$ 收敛。由 Weierstrass 判别法，原函数项级数在 $[0,+\infty)$ 上一致收敛。
\end{solution}

\subsubsection*{五、应用题（本题 9 分）}
\begin{exercise}
\par
从斜边之长为 $a$ 的一切直角三角形中，求有最大周长的直角三角形。
\end{exercise}

\begin{solution}
\par
设直角三角形两条直角边为 $x,y>0$，斜边为 $a$。则
\[
  x^2+y^2=a^2.
\]
周长为
\[
  P=a+x+y.
\]
由于 $a$ 固定，最大化周长等价于最大化 $x+y$。由
\[
  (x-y)^2\ge0
\]
得
\[
  x^2+y^2\ge2xy.
\]
于是
\[
  (x+y)^2=x^2+y^2+2xy\le a^2+a^2=2a^2,
\]
所以
\[
  x+y\le\sqrt2\,a.
\]
等号成立当且仅当 $x=y$。代入约束得
\[
  2x^2=a^2,\qquad x=y=\frac{a}{\sqrt2}.
\]
因此所求三角形为等腰直角三角形，最大周长为
\[
  P_{\max}=a+\sqrt2\,a=(1+\sqrt2)a.
\]
\end{solution}

\section{2016级工科数学分析下 A卷}

\par\medskip
\noindent\textbf{一、填空题（共5题，每题2分，共10分）}\par\nobreak\smallskip
\begin{exercise}
\par
设函数
  \[
    f(x,y)=2x^2+y^2.
  \]
  其在点 \((1,1)\) 处沿该点的梯度方向的方向导数为 \underline{\hspace{5cm}}。
\end{exercise}

\begin{solution}
\par
先求梯度：
\[
  \nabla f=(f_x,f_y)=(4x,2y).
\]
在点 $(1,1)$ 处，
\[
  \nabla f(1,1)=(4,2).
\]
沿梯度方向的单位向量为 $\dfrac{\nabla f}{|\nabla f|}$，方向导数等于
\[
  \nabla f(1,1)\cdot\frac{\nabla f(1,1)}{|\nabla f(1,1)|}
  =
  |\nabla f(1,1)|
  =
  \sqrt{4^2+2^2}
  =
  2\sqrt5.
\]
\end{solution}
\begin{exercise}
\par
向量场 \((2x-3y)\mathbf i+(3x-z)\mathbf j+(y-2x)\mathbf k\) 的旋度向量为 \underline{\hspace{5cm}}。
\end{exercise}

\begin{solution}
\par
记
\[
  \mathbf F=(P,Q,R)=(2x-3y,\ 3x-z,\ y-2x).
\]
旋度为
\[
  \nabla\times\mathbf F
  =
  (R_y-Q_z,\ P_z-R_x,\ Q_x-P_y).
\]
逐项计算：
\[
  R_y-Q_z=1-(-1)=2,
\]
\[
  P_z-R_x=0-(-2)=2,
\]
\[
  Q_x-P_y=3-(-3)=6.
\]
故
\[
  \nabla\times\mathbf F=2\mathbf i+2\mathbf j+6\mathbf k.
\]
\end{solution}
\begin{exercise}
\par
设 \(\Sigma\) 表示球面 \(x^2+y^2+z^2=R^2\) 的外侧，则第二类曲面积分
  \(\displaystyle \iint_\Sigma x\,\mathrm{d}y\,\mathrm{d}z+y\,\mathrm{d}z\,\mathrm{d}x+z\,\mathrm{d}x\,\mathrm{d}y=\)
  \underline{\hspace{5cm}}。
\end{exercise}

\begin{solution}
\par
这是向量场
\[
  \mathbf F=(x,y,z)
\]
通过球面外侧的通量。其散度为
\[
  \nabla\cdot\mathbf F=1+1+1=3.
\]
由 Gauss 公式，
\[
  \iint_\Sigma x\,\mathrm{d}y\,\mathrm{d}z+y\,\mathrm{d}z\,\mathrm{d}x+z\,\mathrm{d}x\,\mathrm{d}y
  =
  \iiint_V3\,\mathrm{d}V.
\]
球体体积为 $\dfrac43\pi R^3$，所以积分为
\[
  3\cdot\frac43\pi R^3=4\pi R^3.
\]
\end{solution}
\begin{exercise}
\par
初值问题
  \[
    \begin{cases}
      y'+y\cos x=\mathrm{e}^{-\sin x},\\
      y|_{x=0}=1
    \end{cases}
  \]
  的解为 \underline{\hspace{5cm}}。
\end{exercise}

\begin{solution}
\par
方程是
\[
  y'+(\cos x)y=e^{-\sin x}.
\]
积分因子为
\[
  \mu(x)=e^{\int\cos x\,\mathrm{d}x}=e^{\sin x}.
\]
两边乘以积分因子：
\[
  e^{\sin x}y'+e^{\sin x}(\cos x)y=1.
\]
左边正是
\[
  (ye^{\sin x})'.
\]
所以
\[
  (ye^{\sin x})'=1.
\]
积分得
\[
  ye^{\sin x}=x+C.
\]
由 $y(0)=1$ 得 $C=1$，故
\[
  y=(x+1)e^{-\sin x}.
\]
\end{solution}
\begin{exercise}
\par
设 \(f(x)\) 是周期为 \(2\) 的周期函数，它在 \((-1,1]\) 上的表达式为
  \[
    f(x)=
    \begin{cases}
      2, & -1<x\le 0,\\
      x^3+1, & 0<x\le 1,
    \end{cases}
  \]
  则 \(f(x)\) 的傅里叶级数在 \(x=0\) 处收敛于 \underline{\hspace{5cm}}。
\end{exercise}

\begin{solution}
\par
傅里叶级数在第一类间断点处收敛到左右极限的平均值。这里
\[
  f(0-0)=2,
\]
而
\[
  f(0+0)=0^3+1=1.
\]
因此在 $x=0$ 处的 Fourier 级数和为
\[
  \frac{f(0-0)+f(0+0)}2
  =
  \frac{2+1}{2}
  =
  \frac32.
\]
注意题目给出的区间是 \((-1,1]\)，但 \(0\) 是内部跳跃点，不是周期端点；处理方法仍然是取同一点的左、右极限平均值，而不是直接采用 \(f(0)=2\) 或右侧表达式的值。
\end{solution}

\par\medskip
\noindent\textbf{二、单选题（每题只有一个正确选项，共5题，每题2分，共10分）}\par\nobreak\smallskip
\begin{exercise}
\par
二元函数
  \[
    f(x,y)=
    \begin{cases}
      \dfrac{xy}{x^2+y^2}, & (x,y)\ne(0,0),\\
      0, & (x,y)=(0,0)
    \end{cases}
  \]
  在点 \((0,0)\) 处（\quad）。
  \par\noindent\begin{tabular}{@{}>{\raggedright\arraybackslash}p{0.22\linewidth}>{\raggedright\arraybackslash}p{0.22\linewidth}>{\raggedright\arraybackslash}p{0.22\linewidth}>{\raggedright\arraybackslash}p{0.22\linewidth}@{}}
  A. 连续，偏导数存在。 & B. 不连续，偏导数存在。 & C. 连续，偏导数不存在。 & D. 不连续，偏导数不存在。
  \end{tabular}
\end{exercise}

\begin{solution}
\par
沿直线 $y=x$ 趋于原点时，
\[
  f(x,x)=\frac{x^2}{2x^2}=\frac12,
\]
而 $f(0,0)=0$，所以函数在原点不连续。再看偏导数：
\[
  f_x(0,0)=\lim_{h\to0}\frac{f(h,0)-f(0,0)}{h}=0,
\]
\[
  f_y(0,0)=\lim_{h\to0}\frac{f(0,h)-f(0,0)}{h}=0.
\]
因此在 $(0,0)$ 处不连续，但两个偏导数存在。选 B。
\end{solution}
\begin{exercise}
\par
曲线
  \[
    \begin{cases}
      x^2+y^2+z^2=9,\\
      x-y+z=1
    \end{cases}
  \]
  在点 \(M(1,2,2)\) 处的切线一定平行于（\quad）。
  \par\noindent\begin{tabular}{@{}>{\raggedright\arraybackslash}p{0.22\linewidth}>{\raggedright\arraybackslash}p{0.22\linewidth}>{\raggedright\arraybackslash}p{0.22\linewidth}>{\raggedright\arraybackslash}p{0.22\linewidth}@{}}
  A. \(Oxy\) 平面。 & B. \(Oyz\) 平面。 & C. \(Ozx\) 平面。 & D. 平面 \(x-y+z=1\)。
  \end{tabular}
\end{exercise}

\begin{solution}
\par
交线的切向量同时垂直于两个曲面的法向量。球面的法向量为
\[
  \nabla(x^2+y^2+z^2)=(2x,2y,2z),
\]
在 $M(1,2,2)$ 处为 $(2,4,4)$。平面 $x-y+z=1$ 的法向量为
\[
  (1,-1,1).
\]
故切向量可取
\[
  (2,4,4)\times(1,-1,1)=(8,2,-6).
\]
该向量与平面 $x-y+z=1$ 的法向量点乘为
\[
  (8,2,-6)\cdot(1,-1,1)=8-2-6=0,
\]
说明切线平行于平面 $x-y+z=1$。选 D。
\end{solution}
\begin{exercise}
\par
设 \(D\) 是一个有界的平面闭区域，其边界曲线 \(\Gamma\) 分段光滑，则下列积分值不等于区域 \(D\) 的面积的是（\quad）。
  \par\noindent\begin{tabular}{@{}>{\raggedright\arraybackslash}p{0.46\linewidth}>{\raggedright\arraybackslash}p{0.46\linewidth}@{}}
  A. \(\displaystyle \oint_\Gamma y\,\mathrm{d}x\) & B. \(\displaystyle \frac12\oint_\Gamma x\,\mathrm{d}y-y\,\mathrm{d}x\)\\
  C. \(\displaystyle \oint_\Gamma x\,\mathrm{d}y\) & D. \(\displaystyle \iint_D 1\,\mathrm{d}x\,\mathrm{d}y\)
  \end{tabular}
\end{exercise}

\begin{solution}
\par
设边界 $\Gamma$ 为正向，即逆时针方向。由 Green 公式，
\[
  \oint_\Gamma y\,\mathrm{d}x
  =
  \iint_D(0-1)\,\mathrm{d}A
  =
  -S_D,
\]
所以 A 不等于面积。又
\[
  \oint_\Gamma x\,\mathrm{d}y
  =
  \iint_D(1-0)\,\mathrm{d}A
  =
  S_D,
\]
\[
  \frac12\oint_\Gamma x\,\mathrm{d}y-y\,\mathrm{d}x
  =
  \frac12(S_D-(-S_D))=S_D,
\]
且
\[
  \iint_D1\,\mathrm{d}x\,\mathrm{d}y=S_D.
\]
因此不等于区域面积的是 A。
\end{solution}
\begin{exercise}
\par
关于未知函数 \(y\) 的微分方程 \((y-\sin x)\,\mathrm{d}x+x\,\mathrm{d}y=0\) 是（\quad）。
  \par\noindent\begin{tabular}{@{}>{\raggedright\arraybackslash}p{0.22\linewidth}>{\raggedright\arraybackslash}p{0.22\linewidth}>{\raggedright\arraybackslash}p{0.22\linewidth}>{\raggedright\arraybackslash}p{0.22\linewidth}@{}}
  A. 可分离变量方程。 & B. 一阶非齐次线性方程。 & C. 一阶齐次线性方程。 & D. 非线性方程。
  \end{tabular}
\end{exercise}

\begin{solution}
\par
将
\[
  (y-\sin x)\,\mathrm{d}x+x\,\mathrm{d}y=0
\]
除以 $\mathrm{d}x$，得
\[
  y-\sin x+xy'=0.
\]
整理：
\[
  y'+\frac1x y=\frac{\sin x}{x}.
\]
这正是 $y'+P(x)y=Q(x)$ 的一阶线性方程，且右端不恒为 $0$，所以是一阶非齐次线性方程。判断时只看未知函数 \(y\) 与 \(y'\) 是否一次出现；这里的 \(\sin x\) 和 \(1/x\) 都只是关于自变量 \(x\) 的系数或右端项，不影响线性性。选 B。
\end{solution}
\begin{exercise}
\par
使得级数 \(\displaystyle \sum_{n=1}^{\infty}\frac{(-1)^n}{n^p}\) 条件收敛的常数 \(p\) 的取值范围是（\quad）。
  \par\noindent\begin{tabular}{@{}>{\raggedright\arraybackslash}p{0.22\linewidth}>{\raggedright\arraybackslash}p{0.22\linewidth}>{\raggedright\arraybackslash}p{0.22\linewidth}>{\raggedright\arraybackslash}p{0.22\linewidth}@{}}
  A. \(p\le 0\) & B. \(0<p<1\) & C. \(0<p\le 1\) & D. \(p>1\)
  \end{tabular}
\end{exercise}

\begin{solution}
\par
对交错级数
\[
  \sum_{n=1}^{\infty}\frac{(-1)^n}{n^p},
\]
若 $p>0$，则 $1/n^p$ 单调趋于 $0$，由交错级数判别法收敛；若 $p\le0$，通项不趋于 $0$，发散。绝对收敛对应
\[
  \sum_{n=1}^{\infty}\frac1{n^p}
\]
收敛，即 $p>1$。所以条件收敛为
\[
  0<p\le1.
\]
选 C。
这里 \(p=1\) 时原级数是交错调和型，仍然收敛，但绝对值级数是调和级数发散，所以应包含端点 \(p=1\)。
\end{solution}

\par\medskip
\noindent\textbf{三、计算题（共4题，每题7分，共28分）}\par\nobreak\smallskip
\begin{exercise}
\par
设函数 \(u=f(x,y,z)\) 有连续的偏导数，函数 \(y=y(x)\) 和 \(z=z(x)\) 分别由方程
  \(\mathrm{e}^{xy}=y\) 和 \(\mathrm{e}^z=xz\) 确定，计算 \(\displaystyle \dfrac{\mathrm{d}u}{\mathrm{d}x}\)。
\end{exercise}

\begin{solution}
\par
把 $u=f(x,y,z)$ 看成 $x$ 的复合函数：
\[
  \frac{\mathrm{d}u}{\mathrm{d}x}=f_x+f_y\,y'(x)+f_z\,z'(x).
\]
先由
\[
  e^{xy}=y
\]
对 $x$ 求导：
\[
  e^{xy}(y+xy')=y'.
\]
利用 $e^{xy}=y$，得到
\[
  y(y+xy')=y',
\]
即
\[
  y'(1-xy)=y^2,\qquad y'=\frac{y^2}{1-xy}.
\]
再由
\[
  e^z=xz
\]
求导：
\[
  e^z z'=z+xz'.
\]
故
\[
  z'(e^z-x)=z.
\]
又 $e^z=xz$，所以 $e^z-x=x(z-1)$，从而
\[
  z'=\frac{z}{x(z-1)}.
\]
代回链式法则：
\[
  \frac{\mathrm{d}u}{\mathrm{d}x}
  =
  f_x+f_y\frac{y^2}{1-xy}+f_z\frac{z}{x(z-1)}.
\]
\end{solution}
\begin{exercise}
\par
计算累次积分
  \[
    \int_1^2\mathrm{d}x\int_{\sqrt{x}}^x\sin\frac{\pi x}{2y}\,\mathrm{d}y
    +\int_2^4\mathrm{d}x\int_{\sqrt{x}}^2\sin\frac{\pi x}{2y}\,\mathrm{d}y.
  \]
\end{exercise}

\begin{solution}
\par
先把两块区域合并。第一块为
\[
  1\le x\le2,\qquad \sqrt{x}\le y\le x.
\]
第二块为
\[
  2\le x\le4,\qquad \sqrt{x}\le y\le2.
\]
合并后按 $y$ 描述，得到
\[
  1\le y\le2,\qquad y\le x\le y^2.
\]
因此
\[
  I=\int_1^2\mathrm{d}y\int_y^{y^2}\sin\frac{\pi x}{2y}\,\mathrm{d}x.
\]
固定 $y$，令 $u=\dfrac{\pi x}{2y}$，则
\[
  \int \sin\frac{\pi x}{2y}\,\mathrm{d}x
  =
  -\frac{2y}{\pi}\cos\frac{\pi x}{2y}.
\]
代入上下限：
\[
  \int_y^{y^2}\sin\frac{\pi x}{2y}\,\mathrm{d}x
  =
  -\frac{2y}{\pi}
  \left(
  \cos\frac{\pi y}{2}-\cos\frac{\pi}{2}
  \right)
  =
  -\frac{2y}{\pi}\cos\frac{\pi y}{2}.
\]
于是
\[
  I=-\frac2\pi\int_1^2 y\cos\frac{\pi y}{2}\,\mathrm{d}y.
\]
分部积分或直接套公式
\[
  \int y\cos ay\,\mathrm{d}y=\frac{y\sin ay}{a}+\frac{\cos ay}{a^2}
\]
并取 $a=\dfrac\pi2$，得
\[
  I=\frac4{\pi^2}+\frac8{\pi^3}.
\]
\end{solution}
\begin{exercise}
\par
计算三重积分 \(\displaystyle \iiint_\Omega z\,\mathrm{d}x\,\mathrm{d}y\,\mathrm{d}z\)，其中 \(\Omega\) 是由上半球面
  \(x^2+y^2+z^2=R^2\)，\(z\ge 0\)，\(R>0\) 和锥面 \(z=\sqrt{x^2+y^2}\) 所围成的闭区域。
\end{exercise}

\begin{solution}
\par
用球坐标
\[
  x=r\sin\varphi\cos\theta,\quad
  y=r\sin\varphi\sin\theta,\quad
  z=r\cos\varphi.
\]
锥面 $z=\sqrt{x^2+y^2}$ 对应
\[
  r\cos\varphi=r\sin\varphi,
  \qquad \varphi=\frac{\pi}{4}.
\]
所围区域在锥面内侧且在上半球内，因此
\[
  0\le r\le R,\qquad
  0\le\varphi\le\frac{\pi}{4},\qquad
  0\le\theta\le2\pi.
\]
又 $z=r\cos\varphi$，体积元为 $r^2\sin\varphi\,\mathrm{d}r\,\mathrm{d}\varphi\,\mathrm{d}\theta$，故
\[
  \iiint_\Omega z\,\mathrm{d}V
  =
  \int_0^{2\pi}\int_0^{\pi/4}\int_0^R
  r^3\cos\varphi\sin\varphi\,\mathrm{d}r\,\mathrm{d}\varphi\,\mathrm{d}\theta.
\]
分别计算：
\[
  \int_0^Rr^3\,\mathrm{d}r=\frac{R^4}{4},\qquad
  \int_0^{\pi/4}\cos\varphi\sin\varphi\,\mathrm{d}\varphi=\frac14.
\]
所以
\[
  \iiint_\Omega z\,\mathrm{d}V
  =
  2\pi\cdot\frac14\cdot\frac{R^4}{4}
  =
  \frac{\pi R^4}{8}.
\]
\end{solution}
\begin{exercise}
\par
计算第二类曲线积分
  \(\displaystyle \oint_\Gamma xy^2\,\mathrm{d}y-x^2y\,\mathrm{d}x\)，
  其中 \(\Gamma\) 为椭圆 \(\displaystyle \dfrac{x^2}{4}+y^2=1\)，取逆时针方向。
\end{exercise}

\begin{solution}
\par
写成 Green 公式的形式：
\[
  P=-x^2y,\qquad Q=xy^2.
\]
椭圆取逆时针方向，故
\[
  \oint_\Gamma P\,\mathrm{d}x+Q\,\mathrm{d}y
  =
  \iint_D(Q_x-P_y)\,\mathrm{d}A.
\]
计算
\[
  Q_x=y^2,\qquad P_y=-x^2,
\]
所以
\[
  Q_x-P_y=x^2+y^2.
\]
对椭圆区域作变换
\[
  x=2r\cos\theta,\qquad y=r\sin\theta,
  \qquad 0\le r\le1,\quad 0\le\theta\le2\pi.
\]
Jacobian 为
\[
  \left|\frac{\partial(x,y)}{\partial(r,\theta)}\right|=2r.
\]
因此
\[
  I=\int_0^{2\pi}\int_0^1
  (4r^2\cos^2\theta+r^2\sin^2\theta)\,2r\,\mathrm{d}r\,\mathrm{d}\theta.
\]
先对 $r$ 积分得 $\int_0^1 2r^3\,\mathrm{d}r=\dfrac12$，于是
\[
  I
  =
  \frac12\int_0^{2\pi}(4\cos^2\theta+\sin^2\theta)\,\mathrm{d}\theta
  =
  \frac12(4\pi+\pi)
  =
  \frac{5\pi}{2}.
\]
\end{solution}

\par\medskip
\noindent\textbf{四、解答题（共4题，每题8分，共32分）}\par\nobreak\smallskip
\begin{exercise}
\par
设二阶常系数线性微分方程
  \(\displaystyle y''+\alpha y'+\beta y=-\mathrm{e}^x\)
  的一个特解为 \(y=(1+x)\mathrm{e}^x\)，试确定常数 \(\alpha,\beta\)，并求该方程的通解。
\end{exercise}

\begin{solution}
\par
已知特解
\[
  y=(1+x)e^x.
\]
先求导：
\[
  y'=(2+x)e^x,\qquad y''=(3+x)e^x.
\]
代入方程
\[
  y''+\alpha y'+\beta y=-e^x,
\]
得
\[
  \bigl[(3+x)+\alpha(2+x)+\beta(1+x)\bigr]e^x=-e^x.
\]
比较 $x$ 的系数与常数项：
\[
  1+\alpha+\beta=0,
\]
\[
  3+2\alpha+\beta=-1.
\]
解得
\[
  \alpha=-3,\qquad \beta=2.
\]
原方程为
\[
  y''-3y'+2y=-e^x.
\]
对应齐次方程的特征方程为
\[
  r^2-3r+2=0,\qquad r=1,2.
\]
齐次解为
\[
  y_h=C_1e^x+C_2e^{2x}.
\]
题目给的特解 $(1+x)e^x$ 中的 $e^x$ 部分可并入齐次解，所以通解可写成
\[
  y=C_1e^x+C_2e^{2x}+xe^x.
\]
\end{solution}
\begin{exercise}
\par
已知螺旋形弹簧一圈的方程为
  \(\displaystyle x=a\cos t,\ y=a\sin t,\ z=bt,\ 0\le t\le 2\pi\)，
  其中 \(a,b\) 为大于零的常数，且弹簧上各点处的线密度等于该点到 \(Oxy\) 平面的距离，求此弹簧的质心坐标。
\end{exercise}

\begin{solution}
\par
曲线参数为
\[
  \mathbf r(t)=(a\cos t,a\sin t,bt),\qquad 0\le t\le2\pi.
\]
弧长微元为
\[
  \mathrm{d}s=\sqrt{(-a\sin t)^2+(a\cos t)^2+b^2}\,\mathrm{d}t
  =
  \sqrt{a^2+b^2}\,\mathrm{d}t.
\]
线密度等于到 $Oxy$ 平面的距离，即
\[
  \rho=z=bt.
\]
总质量
\[
  m=\int_0^{2\pi}\rho\,\mathrm{d}s
  =
  \int_0^{2\pi}bt\sqrt{a^2+b^2}\,\mathrm{d}t
  =
  2\pi^2b\sqrt{a^2+b^2}.
\]
质心坐标为
\[
  \bar x=\frac1m\int_0^{2\pi}x\rho\,\mathrm{d}s,\qquad
  \bar y=\frac1m\int_0^{2\pi}y\rho\,\mathrm{d}s,\qquad
  \bar z=\frac1m\int_0^{2\pi}z\rho\,\mathrm{d}s.
\]
先算 $\bar x$：
\[
  \int_0^{2\pi}a\cos t\cdot bt\sqrt{a^2+b^2}\,\mathrm{d}t
  =
  ab\sqrt{a^2+b^2}\int_0^{2\pi}t\cos t\,\mathrm{d}t=0.
\]
故 $\bar x=0$。再算 $\bar y$：
\[
  \int_0^{2\pi}a\sin t\cdot bt\sqrt{a^2+b^2}\,\mathrm{d}t
  =
  ab\sqrt{a^2+b^2}\int_0^{2\pi}t\sin t\,\mathrm{d}t.
\]
分部积分得
\[
  \int_0^{2\pi}t\sin t\,\mathrm{d}t=-2\pi,
\]
所以
\[
  \bar y=\frac{-2\pi ab\sqrt{a^2+b^2}}{2\pi^2b\sqrt{a^2+b^2}}
  =
  -\frac a\pi.
\]
最后
\[
  \int_0^{2\pi}bt\cdot bt\sqrt{a^2+b^2}\,\mathrm{d}t
  =
  b^2\sqrt{a^2+b^2}\int_0^{2\pi}t^2\,\mathrm{d}t
  =
  b^2\sqrt{a^2+b^2}\frac{8\pi^3}{3}.
\]
因此
\[
  \bar z=\frac{4\pi b}{3}.
\]
质心为
\[
  \left(0,-\frac a\pi,\frac{4\pi b}{3}\right).
\]
\end{solution}
\begin{exercise}
\par
求上半球面 \(x^2+y^2+z^2=a^2,\ z\ge 0\) 被柱面 \(x^2+y^2=ax\) 截下的部分的面积。
\end{exercise}

\begin{solution}
\par
柱面 $x^2+y^2=ax$ 在极坐标下为
\[
  r^2=ar\cos\theta,\qquad r=a\cos\theta.
\]
其内部投影区域为
\[
  -\frac{\pi}{2}\le\theta\le\frac{\pi}{2},\qquad
  0\le r\le a\cos\theta.
\]
上半球面为
\[
  z=\sqrt{a^2-r^2}.
\]
球面投影到 $xy$ 平面时，
\[
  \mathrm{d}S=\frac{a}{\sqrt{a^2-r^2}}\,\mathrm{d}x\,\mathrm{d}y
  =
  \frac{a}{\sqrt{a^2-r^2}}r\,\mathrm{d}r\,\mathrm{d}\theta.
\]
因此面积为
\[
  S=\int_{-\pi/2}^{\pi/2}\int_0^{a\cos\theta}
  \frac{ar}{\sqrt{a^2-r^2}}\,\mathrm{d}r\,\mathrm{d}\theta.
\]
内层积分：
\[
  \int \frac{ar}{\sqrt{a^2-r^2}}\,\mathrm{d}r
  =
  -a\sqrt{a^2-r^2}.
\]
所以
\[
  S
  =
  \int_{-\pi/2}^{\pi/2}
  \left(a^2-a\sqrt{a^2-a^2\cos^2\theta}\right)\mathrm{d}\theta.
\]
在该区间内 $\sqrt{a^2\sin^2\theta}=a|\sin\theta|$，故
\[
  S=a^2\int_{-\pi/2}^{\pi/2}(1-|\sin\theta|)\,\mathrm{d}\theta
  =
  a^2(\pi-2).
\]
\end{solution}
\begin{exercise}
\par
将函数 \(f(x)=\arctan x\) 展开为 \(x\) 的幂级数。
\end{exercise}

\begin{solution}
\par
先用几何级数：
\[
  \frac1{1+x^2}
  =
  \frac1{1-(-x^2)}
  =
  \sum_{n=0}^{\infty}(-1)^n x^{2n},
  \qquad |x|<1.
\]
因为
\[
  (\arctan x)'=\frac1{1+x^2},\qquad \arctan0=0,
\]
对上式从 $0$ 到 $x$ 逐项积分：
\[
  \arctan x
  =
  \int_0^x\frac{\mathrm{d}t}{1+t^2}
  =
  \sum_{n=0}^{\infty}(-1)^n\frac{x^{2n+1}}{2n+1}.
\]
在 $x=\pm1$ 时该级数也由交错级数判别法收敛，因此展开式可写为
\[
  \arctan x
  =
  \sum_{n=0}^{\infty}\frac{(-1)^n x^{2n+1}}{2n+1},
  \qquad |x|\le1.
\]
\end{solution}

\par\medskip
\noindent\textbf{五、证明题（共2题，每题6分，共12分）}\par\nobreak\smallskip
\begin{exercise}
\par
设函数 \(f(\xi,\eta)\) 具有连续的二阶偏导数，且满足 \(\displaystyle \frac{\partial^2f}{\partial \xi^2}+\frac{\partial^2f}{\partial \eta^2}=0\)。证明：函数 \(z=f(x^2-y^2,2xy)\) 满足 \(\displaystyle \frac{\partial^2z}{\partial x^2}+\frac{\partial^2z}{\partial y^2}=0\)。
\end{exercise}

\begin{solution}
\par
令
\[
  \xi=x^2-y^2,\qquad \eta=2xy,\qquad z=f(\xi,\eta).
\]
由链式法则，
\[
  z_x=2x f_\xi+2y f_\eta,\qquad
  z_y=-2y f_\xi+2x f_\eta.
\]
继续求二阶偏导并相加，混合偏导项相互抵消，得到
\[
  z_{xx}+z_{yy}
  =
  4(x^2+y^2)(f_{\xi\xi}+f_{\eta\eta}).
\]
题设
\[
  f_{\xi\xi}+f_{\eta\eta}=0,
\]
故
\[
  z_{xx}+z_{yy}=0.
\]
命题得证。
\end{solution}
\begin{exercise}
\par
证明：函数项级数 \(\displaystyle \sum_{n=0}^{\infty}(1-x)x^n\) 在区间 \((0,1)\) 上点态收敛，但不一致收敛。
\end{exercise}

\begin{solution}
\par
对固定 $x\in(0,1)$，这是首项为 $1-x$、公比为 $x$ 的几何级数：
\[
  \sum_{n=0}^{\infty}(1-x)x^n
  =
  (1-x)\frac1{1-x}
  =
  1.
\]
所以它在 $(0,1)$ 上点态收敛，和函数为 $1$。

第 $N$ 项部分和为
\[
  S_N(x)=(1-x)\sum_{n=0}^{N}x^n=1-x^{N+1}.
\]
余项为
\[
  1-S_N(x)=x^{N+1}.
\]
若一致收敛，则应有
\[
  \sup_{0<x<1}x^{N+1}\to0.
\]
但对每个固定 $N$，
\[
  \sup_{0<x<1}x^{N+1}=1,
\]
因为 $x$ 可以任意接近 $1$。余项不一致趋于 $0$，故该函数项级数在 $(0,1)$ 上不一致收敛。
\end{solution}

\par\medskip
\noindent\textbf{六、应用题（共1题，共8分）}\par\nobreak\smallskip
\begin{exercise}
\par
求曲线
\[
  \begin{cases}
    z=x^2+y^2,\\
    y=\dfrac1x
  \end{cases}
\]
上到 \(Oxy\) 平面距离最近的点。
\end{exercise}

\begin{solution}
\par
点到 $Oxy$ 平面的距离为
\[
  |z|.
\]
曲线上满足
\[
  z=x^2+y^2,\qquad y=\frac1x,\qquad x\ne0.
\]
因此
\[
  z=x^2+\frac1{x^2}>0,
\]
距离就是 $z$。令
\[
  \phi(x)=x^2+\frac1{x^2}.
\]
求导：
\[
  \phi'(x)=2x-\frac2{x^3}.
\]
令 $\phi'(x)=0$，得
\[
  2x-\frac2{x^3}=0
  \quad\Longleftrightarrow\quad
  x^4=1.
\]
所以
\[
  x=\pm1.
\]
当 $x=1$ 时，$y=1,z=2$；当 $x=-1$ 时，$y=-1,z=2$。并且由
\[
  x^2+\frac1{x^2}\ge2
\]
可知这是最小值。故最近点为
\[
  (1,1,2),\qquad (-1,-1,2).
\]
\end{solution}

\section{2016级工科数学分析下 B卷}

\par\medskip
\noindent\textbf{一、填空题（共 5 题，每题 2 分，共 10 分）}\par\nobreak\smallskip
\begin{exercise}
\par
由方程
  \[
    xyz+\sqrt{x^2+y^2+z^2}=\sqrt2
  \]
  所确定的函数 $z=z(x,y)$ 在点 $(1,0,-1)$ 处的全微分 $\mathrm{d}z=$\underline{\hspace{4cm}}。
\end{exercise}

\begin{solution}
\par
设
\[
  F(x,y,z)=xyz+\sqrt{x^2+y^2+z^2}-\sqrt2.
\]
由隐函数方程 $F(x,y,z)=0$，有
\[
  F_x\,\mathrm{d}x+F_y\,\mathrm{d}y+F_z\,\mathrm{d}z=0.
\]
计算偏导：
\[
  F_x=yz+\frac{x}{\sqrt{x^2+y^2+z^2}},
\]
\[
  F_y=xz+\frac{y}{\sqrt{x^2+y^2+z^2}},
  \qquad
  F_z=xy+\frac{z}{\sqrt{x^2+y^2+z^2}}.
\]
在 $(1,0,-1)$ 处，
\[
  F_x=\frac1{\sqrt2},\qquad
  F_y=-1,\qquad
  F_z=-\frac1{\sqrt2}.
\]
所以
\[
  \frac1{\sqrt2}\mathrm{d}x-\mathrm{d}y-\frac1{\sqrt2}\mathrm{d}z=0,
\]
解得
\[
  \mathrm{d}z=\mathrm{d}x-\sqrt2\,\mathrm{d}y.
\]
\end{solution}
\begin{exercise}
\par
设 $\Gamma$ 为椭圆
  \[
    \frac{x^2}{9}+\frac{y^2}{4}=1,
  \]
  取逆时针方向，则第二类曲线积分 \(\displaystyle \oint_\Gamma x\,\mathrm{d}y-y\,\mathrm{d}x\) 的值为 \underline{\hspace{4cm}}。
\end{exercise}

\begin{solution}
\par
记
\[
  P=-y,\qquad Q=x.
\]
椭圆取逆时针方向，由 Green 公式
\[
  \oint_\Gamma x\,\mathrm{d}y-y\,\mathrm{d}x
  =
  \iint_D(Q_x-P_y)\,\mathrm{d}A.
\]
因为
\[
  Q_x=1,\qquad P_y=-1,
\]
所以
\[
  Q_x-P_y=2.
\]
椭圆半轴长分别为 $3$ 与 $2$，面积为
\[
  S=6\pi.
\]
故积分为
\[
  2S=12\pi.
\]
方向为逆时针，正好是椭圆内部区域的正向边界；若方向反过来，最后结果才需要取负号。
\end{solution}
\begin{exercise}
\par
向量场 \(\displaystyle x^2yz\,\mathbf i+xy^2z\,\mathbf j+xyz^2\,\mathbf k\) 在点 $(1,1,1)$ 的散度为\underline{\hspace{4cm}}。
\end{exercise}

\begin{solution}
\par
散度为
\[
  \operatorname{div}\mathbf F
  =
  \frac{\partial}{\partial x}(x^2yz)
  +
  \frac{\partial}{\partial y}(xy^2z)
  +
  \frac{\partial}{\partial z}(xyz^2).
\]
逐项计算：
\[
  \frac{\partial}{\partial x}(x^2yz)=2xyz,
\]
\[
  \frac{\partial}{\partial y}(xy^2z)=2xyz,
  \qquad
  \frac{\partial}{\partial z}(xyz^2)=2xyz.
\]
所以
\[
  \operatorname{div}\mathbf F=6xyz.
\]
在 $(1,1,1)$ 处取值为
\[
  6.
\]
\end{solution}
\begin{exercise}
\par
初值问题
  \[
    \begin{cases}
      y'+\dfrac{3}{x}y=\dfrac{2}{x^3},\\
      y|_{x=1}=1
    \end{cases}
  \]
  的解为\underline{\hspace{4cm}}。
\end{exercise}

\begin{solution}
\par
这是线性方程
\[
  y'+\frac3x y=\frac2{x^3}.
\]
积分因子为
\[
  \mu(x)=e^{\int 3/x\,\mathrm{d}x}=x^3.
\]
两边乘以 $x^3$：
\[
  x^3y'+3x^2y=2.
\]
左端正是
\[
  (x^3y)'.
\]
故
\[
  (x^3y)'=2.
\]
积分得
\[
  x^3y=2x+C.
\]
由 $y(1)=1$ 得
\[
  1=2+C,\qquad C=-1.
\]
所以
\[
  y=\frac{2x-1}{x^3}
  =
  \frac2{x^2}-\frac1{x^3}.
\]
\end{solution}
\begin{exercise}
\par
设 $f(x)=\pi x+x^2,\ -\pi<x<\pi$ 的傅里叶级数为 \(\displaystyle \frac{a_0}{2}+\sum_{n=1}^{\infty}(a_n\cos nx+b_n\sin nx)\)，则其中系数 $b_2=$\underline{\hspace{4cm}}。
\end{exercise}

\begin{solution}
\par
Fourier 正弦系数为
\[
  b_n=\frac1\pi\int_{-\pi}^{\pi}(\pi x+x^2)\sin nx\,\mathrm{d}x.
\]
其中 $x^2\sin nx$ 是奇函数，在对称区间上积分为 $0$。而 $\pi x\sin nx$ 是偶函数，因此
\[
  b_n=\frac1\pi\cdot 2\int_0^\pi \pi x\sin nx\,\mathrm{d}x
  =
  2\int_0^\pi x\sin nx\,\mathrm{d}x.
\]
分部积分：
\[
  \int_0^\pi x\sin nx\,\mathrm{d}x
  =
  \left[-\frac{x\cos nx}{n}\right]_0^\pi
  +
  \frac1n\int_0^\pi\cos nx\,\mathrm{d}x.
\]
第二项为 $0$，所以
\[
  \int_0^\pi x\sin nx\,\mathrm{d}x
  =
  -\frac{\pi\cos n\pi}{n}
  =
  -\frac{\pi(-1)^n}{n}.
\]
于是
\[
  b_n=-\frac{2\pi(-1)^n}{n}.
\]
取 $n=2$，
\[
  b_2=-\pi.
\]
\end{solution}

\par\medskip
\noindent\textbf{二、单选题（共 5 题，每题 2 分，只有一个正确选项，共 10 分）}\par\nobreak\smallskip
\begin{exercise}
\par
二元函数 $f(x,y)$ 在点 $(a,b)$ 处的两个偏导数 \(\displaystyle \frac{\partial f}{\partial x}(a,b),\qquad \frac{\partial f}{\partial y}(a,b)\) 存在，则（\quad）。
  \par\noindent\begin{tabular}{@{}>{\raggedright\arraybackslash}p{0.46\linewidth}>{\raggedright\arraybackslash}p{0.46\linewidth}@{}}
  A. $f(x,y)$ 在点 $(a,b)$ 处连续； & B. $f(x,y)$ 在点 $(a,b)$ 处可微；\\
  C. \(\displaystyle \lim_{x\to a}f(x,b)\) 和 \(\displaystyle \lim_{y\to b}f(a,y)\) 都存在； & D. \(\displaystyle \lim_{(x,y)\to(a,b)}f(x,y)\) 存在。
  \end{tabular}
\end{exercise}

\begin{solution}
\par
偏导数存在的含义是
\[
  \lim_{x\to a}\frac{f(x,b)-f(a,b)}{x-a}
\]
和
\[
  \lim_{y\to b}\frac{f(a,y)-f(a,b)}{y-b}
\]
存在。由第一个极限存在可知
\[
  f(x,b)-f(a,b)=(x-a)\cdot\frac{f(x,b)-f(a,b)}{x-a}\to0,
\]
故 $\lim_{x\to a}f(x,b)$ 存在并等于 $f(a,b)$；同理 $\lim_{y\to b}f(a,y)$ 也存在。但是偏导存在不能推出二元连续，更不能推出可微，也不能推出二元极限存在。选 C。
\end{solution}
\begin{exercise}
\par
已知曲面 $z=4-x^2-y^2$ 上点 $P$ 处的切平面平行于平面 $2x+2y+z=1$，则 $P$ 点的坐标是（\quad）。
  \par\noindent\begin{tabular}{@{}>{\raggedright\arraybackslash}p{0.22\linewidth}>{\raggedright\arraybackslash}p{0.22\linewidth}>{\raggedright\arraybackslash}p{0.22\linewidth}>{\raggedright\arraybackslash}p{0.22\linewidth}@{}}
  A. $(1,-1,2)$； & B. $(-1,1,2)$； & C. $(1,1,2)$； & D. $(-1,-1,2)$。
  \end{tabular}
\end{exercise}

\begin{solution}
\par
把曲面写成
\[
  F(x,y,z)=x^2+y^2+z-4=0.
\]
曲面法向量为
\[
  \nabla F=(2x,2y,1).
\]
切平面平行于平面 $2x+2y+z=1$，说明两个平面的法向量平行。平面 $2x+2y+z=1$ 的法向量是
\[
  (2,2,1).
\]
由于第三个分量同为 $1$，只能有
\[
  (2x,2y,1)=(2,2,1),
\]
故
\[
  x=1,\qquad y=1.
\]
代回曲面 $z=4-x^2-y^2$，得
\[
  z=2.
\]
所以
\[
  P=(1,1,2).
\]
选 C。
\end{solution}
\begin{exercise}
\par
一个均匀物体由曲面 $z=x^2+y^2$ 及 $z=1$ 围成，则该物体的质心坐标为（\quad）。
  \par\noindent\begin{tabular}{@{}>{\raggedright\arraybackslash}p{0.22\linewidth}>{\raggedright\arraybackslash}p{0.22\linewidth}>{\raggedright\arraybackslash}p{0.22\linewidth}>{\raggedright\arraybackslash}p{0.22\linewidth}@{}}
  A. \(\displaystyle \left(\dfrac23,\dfrac23,\dfrac23\right)\)； & B. \(\displaystyle \left(0,0,\dfrac23\right)\)； & C. $(0,0,1)$； & D. \(\displaystyle \left(\dfrac23,0,0\right)\)。
  \end{tabular}
\end{exercise}

\begin{solution}
\par
物体由 $z=r^2$ 与 $z=1$ 围成，其中 $r^2=x^2+y^2$。区域关于 $z$ 轴旋转对称，所以
\[
  \bar x=\bar y=0.
\]
用柱坐标计算体积：
\[
  V=\int_0^{2\pi}\int_0^1\int_{r^2}^{1}r\,\mathrm{d}z\,\mathrm{d}r\,\mathrm{d}\theta
  =
  2\pi\int_0^1(1-r^2)r\,\mathrm{d}r
  =
  \frac{\pi}{2}.
\]
关于 $xOy$ 平面的静矩为
\[
  M_{xy}=\iiint_\Omega z\,\mathrm{d}V
  =
  \int_0^{2\pi}\int_0^1\int_{r^2}^{1} z\,r\,\mathrm{d}z\,\mathrm{d}r\,\mathrm{d}\theta.
\]
先对 $z$ 积分：
\[
  M_{xy}
  =
  2\pi\int_0^1\frac12(1-r^4)r\,\mathrm{d}r
  =
  \pi\left[\frac{r^2}{2}-\frac{r^6}{6}\right]_0^1
  =
  \frac{\pi}{3}.
\]
因此
\[
  \bar z=\frac{M_{xy}}{V}=\frac{\pi/3}{\pi/2}=\frac23.
\]
质心为
\[
  \left(0,0,\frac23\right).
\]
选 B。
\end{solution}
\begin{exercise}
\par
微分方程 \(\displaystyle (y-\ln x)\,\mathrm{d}x+x\,\mathrm{d}y=0\) 是（\quad）。
  \par\noindent\begin{tabular}{@{}>{\raggedright\arraybackslash}p{0.22\linewidth}>{\raggedright\arraybackslash}p{0.22\linewidth}>{\raggedright\arraybackslash}p{0.22\linewidth}>{\raggedright\arraybackslash}p{0.22\linewidth}@{}}
  A. 可分离变量方程； & B. 一阶非齐次线性方程； & C. 一阶齐次线性方程； & D. 非线性方程。
  \end{tabular}
\end{exercise}

\begin{solution}
\par
将方程
\[
  (y-\ln x)\,\mathrm{d}x+x\,\mathrm{d}y=0
\]
除以 $\mathrm{d}x$：
\[
  y-\ln x+xy'=0.
\]
整理为
\[
  y'+\frac1x y=\frac{\ln x}{x}.
\]
这是一阶线性方程，且右端不为零，所以是一阶非齐次线性方程。由于 \(\ln x\) 出现在右端，方程通常在 \(x>0\) 的区间上讨论；未知函数 \(y\) 没有平方、乘积或复合函数形式，因此不属于非线性方程。选 B。
因此选项判断的核心是把微分形式先化成 \(y'+P(x)y=Q(x)\)，再看 \(Q(x)\) 是否为零。
\end{solution}
\begin{exercise}
\par
下列级数条件收敛的是（\quad）。
  \par\noindent\begin{tabular}{@{}>{\raggedright\arraybackslash}p{0.46\linewidth}>{\raggedright\arraybackslash}p{0.46\linewidth}@{}}
  A. \(\displaystyle \sum_{n=1}^{\infty}\frac{(-1)^n}{n^2}\)； & B. \(\displaystyle \sum_{n=1}^{\infty}\left(1-\cos\frac1n\right)\)；\\
  C. \(\displaystyle \sum_{n=1}^{\infty}\frac{n!}{n^n}\)； & D. \(\displaystyle \sum_{n=1}^{\infty}\frac{(-1)^n}{n}\)。
  \end{tabular}
\end{exercise}

\begin{solution}
\par
A 项的绝对值级数为
\[
  \sum_{n=1}^{\infty}\frac1{n^2},
\]
收敛，所以 A 绝对收敛。B 项中
\[
  1-\cos\frac1n\sim \frac1{2n^2},
\]
为正项收敛级数。C 项用根值判别法：
\[
  \sqrt[n]{\frac{n!}{n^n}}\sim \frac1e<1,
\]
故收敛。D 项
\[
  \sum_{n=1}^{\infty}\frac{(-1)^n}{n}
\]
是交错调和级数，收敛但不绝对收敛，所以条件收敛。选 D。
\end{solution}

\par\medskip
\noindent\textbf{三、计算题（共 4 题，每题 7 分，共 28 分）}\par\nobreak\smallskip
\begin{exercise}
\par
设函数 \(\displaystyle u(x,y)=y f\left(\frac{x}{y}\right)+x g\left(\frac{y}{x}\right)\)，其中 $f$ 和 $g$ 具有连续的二阶导数，计算 \(\displaystyle x\frac{\partial^2u}{\partial x^2}+y\frac{\partial^2u}{\partial x\partial y}\)。
\end{exercise}

\begin{solution}
\par
函数
\[
  u(x,y)=y f\left(\frac{x}{y}\right)+x g\left(\frac{y}{x}\right)
\]
是关于 $x,y$ 的一阶齐次函数，因为对任意 $\lambda>0$，
\[
  u(\lambda x,\lambda y)=\lambda u(x,y).
\]
由 Euler 齐次函数公式，
\[
  xu_x+yu_y=u.
\]
对该恒等式关于 $x$ 求偏导，注意 $y$ 看作常数：
\[
  u_x+xu_{xx}+yu_{yx}=u_x.
\]
两边消去 $u_x$，并用 $u_{yx}=u_{xy}$，得
\[
  xu_{xx}+yu_{xy}=0.
\]
\end{solution}
\begin{exercise}
\par
计算二重积分 \(\displaystyle \iint_D \frac{1-x^2-y^2}{1+x^2+y^2}\,\mathrm{d}x\,\mathrm{d}y\)，其中 $D$ 是区域 $x^2+y^2\le1,\ x\ge0,\ y\ge0$。
\end{exercise}

\begin{solution}
\par
区域 $D$ 是第一象限内的单位四分之一圆盘，用极坐标
\[
  x=r\cos\theta,\qquad y=r\sin\theta,
\]
则
\[
  0\le r\le1,\qquad 0\le\theta\le\frac{\pi}{2}.
\]
被积函数化为
\[
  \frac{1-x^2-y^2}{1+x^2+y^2}
  =
  \frac{1-r^2}{1+r^2},
\]
且 $\mathrm{d}x\,\mathrm{d}y=r\,\mathrm{d}r\,\mathrm{d}\theta$。因此
\[
  I=\int_0^{\pi/2}\int_0^1\frac{1-r^2}{1+r^2}r\,\mathrm{d}r\,\mathrm{d}\theta.
\]
径向积分中
\[
  \frac{1-r^2}{1+r^2}r
  =
  \frac{2r}{1+r^2}-r.
\]
所以
\[
  \int_0^1\frac{1-r^2}{1+r^2}r\,\mathrm{d}r
  =
  \left[\ln(1+r^2)-\frac{r^2}{2}\right]_0^1
  =
  \ln2-\frac12.
\]
再乘角度长度 $\dfrac{\pi}{2}$，得
\[
  I=\frac{\pi}{2}\left(\ln2-\frac12\right)
  =
  \frac{\pi}{4}(2\ln2-1).
\]
\end{solution}
\begin{exercise}
\par
计算三重积分 \(\displaystyle \iiint_\Omega z^2\,\mathrm{d}x\,\mathrm{d}y\,\mathrm{d}z\)，
  其中 $\Omega$ 是平面 $x+y+z=1$ 与三个坐标面围成的区域。
\end{exercise}

\begin{solution}
\par
区域是第一卦限内由平面
\[
  x+y+z=1
\]
与三个坐标面围成的四面体。固定 $z$，有
\[
  0\le z\le1,
\]
截面为
\[
  x\ge0,\qquad y\ge0,\qquad x+y\le1-z.
\]
这是直角三角形，面积为
\[
  \frac12(1-z)^2.
\]
因此
\[
  \iiint_\Omega z^2\,\mathrm{d}x\,\mathrm{d}y\,\mathrm{d}z
  =
  \int_0^1 z^2\cdot\frac12(1-z)^2\,\mathrm{d}z.
\]
展开计算：
\[
  \frac12\int_0^1 z^2(1-2z+z^2)\,\mathrm{d}z
  =
  \frac12\left(\frac13-\frac12+\frac15\right)
  =
  \frac1{60}.
\]
\end{solution}
\begin{exercise}
\par
计算第一类曲面积分 \(\displaystyle \iint_\Sigma (x+y+z)\,\mathrm{d}S\)，
  其中 $\Sigma$ 为平面 $y+z=5$ 被柱面 $x^2+y^2=25$ 所截得的部分。
\end{exercise}

\begin{solution}
\par
平面 $y+z=5$ 可写为
\[
  z=5-y.
\]
于是
\[
  z_x=0,\qquad z_y=-1,
\]
所以
\[
  \mathrm{d}S=\sqrt{1+z_x^2+z_y^2}\,\mathrm{d}x\,\mathrm{d}y=\sqrt2\,\mathrm{d}x\,\mathrm{d}y.
\]
投影区域由柱面给出：
\[
  D:x^2+y^2\le25.
\]
在平面上
\[
  x+y+z=x+y+5-y=x+5.
\]
因此
\[
  I=\sqrt2\iint_D(x+5)\,\mathrm{d}x\,\mathrm{d}y.
\]
圆盘 $D$ 关于 $y$ 轴对称，故
\[
  \iint_Dx\,\mathrm{d}x\,\mathrm{d}y=0.
\]
又 $D$ 面积为 $25\pi$，所以
\[
  I=\sqrt2\cdot5\cdot25\pi=125\sqrt2\,\pi.
\]
\end{solution}

\par\medskip
\noindent\textbf{四、解答题（共 4 题，每题 8 分，共 32 分）}\par\nobreak\smallskip
\begin{exercise}
\par
求方程 \(\displaystyle y''+4y'+3y=0\) 的一个解 $y=y(x)$，使其在点 $(0,2)$ 处与直线 $x-y+2=0$ 相切。
\end{exercise}

\begin{solution}
\par
先解微分方程。特征方程为
\[
  \lambda^2+4\lambda+3=0,
\]
即
\[
  (\lambda+1)(\lambda+3)=0.
\]
所以通解为
\[
  y=C_1e^{-3x}+C_2e^{-x}.
\]
题目要求曲线在点 $(0,2)$ 处与直线
\[
  x-y+2=0
\]
相切。该直线可写为 $y=x+2$，斜率为 $1$。因此需要
\[
  y(0)=2,\qquad y'(0)=1.
\]
由通解得
\[
  y(0)=C_1+C_2=2.
\]
又
\[
  y'=-3C_1e^{-3x}-C_2e^{-x},
\]
所以
\[
  y'(0)=-3C_1-C_2=1.
\]
联立解得
\[
  C_1=-\frac32,\qquad C_2=\frac72.
\]
故所求解为
\[
  y=-\frac32e^{-3x}+\frac72e^{-x}.
\]
\end{solution}
\begin{exercise}
\par
求圆柱面 $x^2+y^2=4,\ x\ge0,\ y\ge0$ 介于 $xOy$ 平面和曲面 $z=xy$ 之间的部分的面积。
\end{exercise}

\begin{solution}
\par
圆柱面 $x^2+y^2=4$ 在第一卦限的部分可以参数化为
\[
  x=2\cos t,\qquad y=2\sin t,\qquad 0\le t\le\frac{\pi}{2}.
\]
曲面介于 $z=0$ 与 $z=xy$ 之间，因此对每个 $t$，竖直高度为
\[
  xy=4\sin t\cos t.
\]
底线弧长微元为
\[
  \mathrm{d}s=\sqrt{(-2\sin t)^2+(2\cos t)^2}\,\mathrm{d}t=2\,\mathrm{d}t.
\]
柱面侧面积等于沿底线积分“高度乘弧长微元”：
\[
  S=\int_\Gamma xy\,\mathrm{d}s
  =
  \int_0^{\pi/2}4\sin t\cos t\cdot2\,\mathrm{d}t.
\]
所以
\[
  S=8\int_0^{\pi/2}\sin t\cos t\,\mathrm{d}t=4.
\]
\end{solution}
\begin{exercise}
\par
设 $\Sigma$ 是锥面 $z=\sqrt{x^2+y^2}$ 被平面 $z=0$ 及 $z=1$ 截下的部分的下侧，计算第二类曲面积分 \(\displaystyle \iint_\Sigma x\,\mathrm{d}y\,\mathrm{d}z+y\,\mathrm{d}z\,\mathrm{d}x+(z^2-2z)\,\mathrm{d}x\,\mathrm{d}y\)。
\end{exercise}

\begin{solution}
\par
记向量场
\[
  \mathbf F=(P,Q,R)=(x,y,z^2-2z).
\]
锥面取下侧。对实心圆锥
\[
  \Omega=\{(x,y,z):x^2+y^2\le z^2,\ 0\le z\le1\}
\]
来说，锥面的外侧正是下侧。为了使用 Gauss 公式，补上顶面圆盘
\[
  \Sigma_1:\ z=1,\quad x^2+y^2\le1,
\]
并取上侧。散度为
\[
  \nabla\cdot\mathbf F=1+1+(2z-2)=2z.
\]
由 Gauss 公式，
\[
  I+I_{\Sigma_1}=\iiint_\Omega 2z\,\mathrm{d}V.
\]
固定 $z$ 时截面为半径 $z$ 的圆盘，面积 $\pi z^2$，所以
\[
  \iiint_\Omega 2z\,\mathrm{d}V
  =
  \int_0^1 2z\cdot\pi z^2\,\mathrm{d}z
  =
  \frac{\pi}{2}.
\]
顶面 $\Sigma_1$ 上 $z=1$，上侧法向量为 $(0,0,1)$，通量只取第三分量：
\[
  R=z^2-2z=1-2=-1.
\]
故
\[
  I_{\Sigma_1}
  =
  \iint_{x^2+y^2\le1}(-1)\,\mathrm{d}x\,\mathrm{d}y
  =
  -\pi.
\]
因此
\[
  I=\frac{\pi}{2}-(-\pi)=\frac{3\pi}{2}.
\]
\end{solution}
\begin{exercise}
\par
求幂级数 \(\displaystyle \sum_{n=1}^{\infty}\frac{x^n}{n}\) 的收敛域，并在收敛域上求其和函数。
\end{exercise}

\begin{solution}
\par
对幂级数
\[
  \sum_{n=1}^{\infty}\frac{x^n}{n},
\]
由根值或比值判别法可得收敛半径为 $1$。端点另判：当 $x=1$ 时，
\[
  \sum_{n=1}^{\infty}\frac1n
\]
发散；当 $x=-1$ 时，
\[
  \sum_{n=1}^{\infty}\frac{(-1)^n}{n}
\]
为交错调和级数，收敛。因此收敛域为
\[
  [-1,1).
\]
在 $(-1,1)$ 内可逐项求导。设
\[
  f(x)=\sum_{n=1}^{\infty}\frac{x^n}{n}.
\]
则
\[
  f'(x)=\sum_{n=1}^{\infty}x^{n-1}=\frac1{1-x}.
\]
又 $f(0)=0$，所以
\[
  f(x)=\int_0^x\frac{\mathrm{d}t}{1-t}
  =
  -\ln(1-x).
\]
在 $x=-1$ 处，级数收敛且右端也有意义，故和函数可写为
\[
  f(x)=-\ln(1-x),\qquad -1\le x<1.
\]
\end{solution}

\par\medskip
\noindent\textbf{五、证明题（共 2 题，每题 6 分，共 12 分）}\par\nobreak\smallskip
\begin{exercise}
\par
证明：曲面 \(\displaystyle z=xe^{y/x}\) 上所有点处的切平面都过一定点。
\end{exercise}

\begin{solution}
\par
设曲面上一点为 $(x_0,y_0,z_0)$，则
\[
  z_0=x_0e^{y_0/x_0}.
\]
先求偏导：
\[
  z_x=e^{y/x}+x e^{y/x}\left(-\frac{y}{x^2}\right)
  =
  \left(1-\frac yx\right)e^{y/x},
\]
\[
  z_y=x e^{y/x}\cdot\frac1x=e^{y/x}.
\]
曲面 $z=z(x,y)$ 在 $(x_0,y_0,z_0)$ 处的切平面为
\[
  z-z_0=z_x(x_0,y_0)(x-x_0)+z_y(x_0,y_0)(y-y_0).
\]
即
\[
  z-z_0
  =
  \left(1-\frac{y_0}{x_0}\right)e^{y_0/x_0}(x-x_0)
  +
  e^{y_0/x_0}(y-y_0).
\]
把点 $(0,0,0)$ 代入右边检验：
\[
  0-z_0
  =
  -\left(1-\frac{y_0}{x_0}\right)x_0e^{y_0/x_0}
  -
  y_0e^{y_0/x_0}.
\]
右端化简为
\[
  -x_0e^{y_0/x_0}.
\]
而 $z_0=x_0e^{y_0/x_0}$，所以等式成立。故所有切平面都通过定点
\[
  (0,0,0).
\]
\end{solution}
\begin{exercise}
\par
证明：函数项级数 \(\displaystyle \sum_{n=0}^{\infty}x^n\) 在 $(0,1)$ 上点态收敛，但不一致收敛。
\end{exercise}

\begin{solution}
\par
对任意固定的 $x\in(0,1)$，几何级数
\[
  \sum_{n=0}^{\infty}x^n
\]
收敛，和函数为
\[
  S(x)=\frac1{1-x}.
\]
所以它在 $(0,1)$ 上点态收敛。

第 $N$ 项部分和为
\[
  S_N(x)=\sum_{n=0}^{N}x^n=\frac{1-x^{N+1}}{1-x}.
\]
余项为
\[
  S(x)-S_N(x)=\frac{x^{N+1}}{1-x}.
\]
若一致收敛，则余项的上确界应趋于 $0$。取
\[
  x_N=1-\frac1N\in(0,1).
\]
则
\[
  |S(x_N)-S_N(x_N)|
  =
  \frac{x_N^{N+1}}{1-x_N}
  =
  N\left(1-\frac1N\right)^{N+1}.
\]
该量不趋于 $0$，事实上趋于无穷大。因此余项不能一致趋于 $0$，级数在 $(0,1)$ 上不一致收敛。
\end{solution}

\par\medskip
\noindent\textbf{六、应用题（共 1 题，共 8 分）}\par\nobreak\smallskip
\begin{exercise}
\par
从斜边长为 $l$ 的直角三角形中，求周长最大的直角三角形。
\end{exercise}

\begin{solution}
\par
设两条直角边为 $x,y>0$，斜边为 $l$，则
\[
  x^2+y^2=l^2.
\]
周长为
\[
  P=x+y+l.
\]
由于 $l$ 固定，只需最大化 $x+y$。由
\[
  (x-y)^2\ge0
\]
得
\[
  x^2+y^2\ge2xy.
\]
于是
\[
  (x+y)^2=x^2+y^2+2xy\le l^2+l^2=2l^2.
\]
所以
\[
  x+y\le\sqrt2\,l.
\]
等号成立当且仅当 $x=y$。代入约束得
\[
  x=y=\frac{l}{\sqrt2}.
\]
故周长最大的直角三角形是等腰直角三角形，最大周长为
\[
  P_{\max}=l+\sqrt2\,l=(1+\sqrt2)l.
\]
\end{solution}

\section{2017级工科数学分析下 A卷}

\par\medskip
\noindent\textbf{一、填空题：共 5 题，每题 2 分，共 10 分}\par\nobreak\smallskip
\begin{exercise}
\par
微分方程 $y''+y'-2y=1-2x$ 的通解为 \underline{\hspace{5cm}}。
\end{exercise}

\begin{solution}
\par
先解齐次方程 \(y''+y'-2y=0\)，特征方程为
\[
  r^2+r-2=0,
\]
得 \(r=1,-2\)，故齐次通解为
\[
  y_h=C_1e^x+C_2e^{-2x}.
\]
右端 \(1-2x\) 是一次多项式，设特解 \(y_p=ax+b\)。代入原方程：
\[
  0+a-2(ax+b)=1-2x.
\]
比较系数得 \(a=1,\ b=0\)。因此通解为
\[
  y=C_1e^{-2x}+C_2e^x+x.
\]
\end{solution}
\begin{exercise}
\par
设函数 $u=\ln(x^2+y^2+z^2)$，求 $\operatorname{div}(\operatorname{grad}u)=$ \underline{\hspace{4cm}}。
\end{exercise}

\begin{solution}
\par
设 \(r^2=x^2+y^2+z^2\)，则 \(u=\ln r^2\)。先求
\[
  u_x=\frac{2x}{r^2},\qquad
  u_{xx}=\frac{2}{r^2}-\frac{4x^2}{r^4}.
\]
同理，
\[
  u_{yy}=\frac{2}{r^2}-\frac{4y^2}{r^4},\qquad
  u_{zz}=\frac{2}{r^2}-\frac{4z^2}{r^4}.
\]
所以
\[
  \operatorname{div}(\operatorname{grad}u)=\Delta u
  =\frac{6}{r^2}-\frac{4(x^2+y^2+z^2)}{r^4}
  =\frac{2}{x^2+y^2+z^2}.
\]
\end{solution}
\begin{exercise}
\par
设 $\Gamma$ 是球面 $x^2+y^2+z^2=R^2$ 与平面 $x+y+z=0$ 的交线，则第一类曲线积分 \(\displaystyle \oint_{\Gamma} y^2\,\mathrm{d}s=\) \underline{\hspace{4cm}}。
\end{exercise}

\begin{solution}
\par
交线 \(\Gamma\) 是过球心的平面截球所得的大圆，半径为 \(R\)。由于平面
\(x+y+z=0\) 的法向量为 \((1,1,1)\)，大圆在三个坐标方向上完全对称，因此
\[
  \oint_\Gamma x^2\,\mathrm{d}s=\oint_\Gamma y^2\,\mathrm{d}s=\oint_\Gamma z^2\,\mathrm{d}s.
\]
又在 \(\Gamma\) 上 \(x^2+y^2+z^2=R^2\)，所以
\[
  3\oint_\Gamma y^2\,\mathrm{d}s
  =\oint_\Gamma (x^2+y^2+z^2)\,\mathrm{d}s
  =R^2\cdot 2\pi R.
\]
故
\[
  \oint_\Gamma y^2\,\mathrm{d}s=\frac{2}{3}\pi R^3.
\]
\end{solution}
\begin{exercise}
\par
参数曲线
  \[
    \begin{cases}
      x=t-\cos t,\\
      y=1-\sin t,\\
      z=t,
    \end{cases}
  \]
  在 $t=0$ 对应的点处的切线方程为 \underline{\hspace{6cm}}。
\end{exercise}

\begin{solution}
\par
当 \(t=0\) 时，对应点为 \((-1,1,0)\)。参数曲线的切向量为
\[
  (x'(t),y'(t),z'(t))=(1+\sin t,-\cos t,1),
\]
所以在 \(t=0\) 处切向量为 \((1,-1,1)\)。切线方程为
  \(\displaystyle \frac{x+1}{1}=\frac{y-1}{-1}=\frac{z}{1}\)。
也可写成参数式 \(x=-1+s,\ y=1-s,\ z=s\)。检查时只需确认直线过点 \((-1,1,0)\)，方向向量与曲线导向量一致。
\end{solution}
\begin{exercise}
\par
设周期为 $2\pi$ 的函数
  \[
    f(x)=
    \begin{cases}
      -1, & -\pi<x\le 0,\\
      1, & 0<x\le \pi,
    \end{cases}
  \]
  则 $f(x)$ 的傅里叶（Fourier）级数在 $x=\pi$ 处收敛于 \underline{\hspace{4cm}}。
\end{exercise}

\begin{solution}
\par
傅里叶级数在跳跃点或周期端点收敛于左右极限的平均值。在 \(x=\pi\) 处，
\[
  f(\pi-0)=1.
\]
由 \(2\pi\) 周期延拓，\(\pi+0\) 与 \(-\pi+0\) 对应，而
\[
  f(\pi+0)=f(-\pi+0)=-1.
\]
因此
\[
  \frac{f(\pi-0)+f(\pi+0)}{2}=\frac{1+(-1)}{2}=0.
\]
这是周期端点题，不是普通连续点代值题；端点右极限必须先用周期性折回到左端附近，再与左极限取平均。
\end{solution}

\par\medskip
\noindent\textbf{二、单选题：共 5 题，每题 2 分，共 10 分}\par\nobreak\smallskip
\begin{exercise}
\par
关于未知函数 $y$ 的微分方程
  \(\displaystyle (y-\sin x)\,\mathrm{d}x+\ln x\,\mathrm{d}y=0\)
  是（\quad）。
  \par\noindent\begin{tabular}{@{}p{0.48\linewidth}p{0.48\linewidth}@{}}
    A. 可分离变量方程； & B. 一阶线性非齐次方程；\\
    C. 一阶线性齐次方程； & D. 非线性方程。
  \end{tabular}
\end{exercise}

\begin{solution}
\par
原方程为
\[
  (y-\sin x)\,\mathrm{d}x+\ln x\,\mathrm{d}y=0.
\]
在 \(\ln x\ne0\) 的区间上除以 \(\ln x\,\mathrm{d}x\)，得
\[
  y'+\frac1{\ln x}y=\frac{\sin x}{\ln x}.
\]
它具有标准形式 \(y'+P(x)y=Q(x)\)，且 \(Q(x)\not\equiv0\)，所以是一阶线性非齐次方程，选 B。
这里要在 \(\ln x\ne0\) 的区间内化标准形，例如不跨过 \(x=1\)；但这种局部区间限制不影响题型判断。
\end{solution}
\begin{exercise}
\par
二元函数
  \[
    f(x,y)=
    \begin{cases}
      x\sin\dfrac{1}{y}, & xy\ne 0,\\
      0, & xy=0,
    \end{cases}
  \]
  则 $\displaystyle \lim_{(x,y)\to(0,0)}f(x,y)$（\quad）。
  \par\noindent\begin{tabular}{@{}p{0.48\linewidth}p{0.48\linewidth}@{}}
    A. 不存在； & B. 等于 1；\\
    C. 等于 0； & D. 等于 2。
  \end{tabular}
\end{exercise}

\begin{solution}
\par
当 \(xy\ne0\) 时，
\[
  |f(x,y)|=\left|x\sin\frac1y\right|\le |x|.
\]
当 \(xy=0\) 时，题设直接给出 \(f(x,y)=0\)，同样满足 \(|f(x,y)|\le |x|\)。因此当
\((x,y)\to(0,0)\) 时，由夹逼定理
\[
  |f(x,y)|\le |x|\to0,
\]
故极限为 \(0\)，选 C。
这个估计不依赖趋近路径，所以比单独沿直线或曲线代入更强；任意路径下 \(|x|\to0\) 都成立。
\end{solution}
\begin{exercise}
\par
设函数 $z=f(x,y)$ 在 $(a,b)$ 的某个邻域内有直到二阶的连续偏导数，且
  \(\displaystyle \frac{\partial z}{\partial x}(a,b)=0,\ \frac{\partial z}{\partial y}(a,b)=0\)。
  记 \(\displaystyle A=\frac{\partial^2 z}{\partial x^2}(a,b),\ B=\frac{\partial^2 z}{\partial x\partial y}(a,b),\ C=\frac{\partial^2 z}{\partial y^2}(a,b)\)。
  则函数 $z=f(x,y)$ 在点 $(a,b)$ 取极大值的充分条件是（\quad）。
  \par\noindent\begin{tabular}{@{}p{0.48\linewidth}p{0.48\linewidth}@{}}
    A. \(A>0,\ AC>B^2\); & B. \(A<0,\ AC>B^2\);\\
    C. \(A>0,\ AC<B^2\); & D. \(A<0,\ AC<B^2\).
  \end{tabular}
\end{exercise}

\begin{solution}
\par
二元函数在驻点处的二阶判别法为：令
\[
  D=AC-B^2.
\]
若 \(D>0\) 且 \(A>0\)，则取极小值；若 \(D>0\) 且 \(A<0\)，则取极大值；若 \(D<0\)，则不是极值点。因此取极大值的充分条件是
\[
  A<0,\qquad AC-B^2>0,
\]
也就是 \(A<0,\ AC>B^2\)，选 B。
这里要求的是“充分条件”，所以必须同时保证 Hessian 对应的二次型负定。单有 \(A<0\) 不够，因为还可能是鞍点；单有 \(AC>B^2\) 也不够，因为还要看 \(A\) 的符号决定极大还是极小。
\end{solution}
\begin{exercise}
\par
函数 $\ln(1-x)$ 在 $x=0$ 处的泰勒（Taylor）展开式正确的是（\quad）。
  \[
  \begin{array}{ll}
    \text{A. } \displaystyle \ln(1-x)=\sum_{n=1}^{\infty}\frac{(-1)^n}{n}x^n,\quad x\in(-1,1];&
    \text{B. } \displaystyle \ln(1-x)=\sum_{n=1}^{\infty}\frac{(-1)^{n-1}}{n}x^n,\quad x\in(-1,1];\\[1.2em]
    \text{C. } \displaystyle \ln(1-x)=-\sum_{n=1}^{\infty}\frac{1}{n}x^n,\quad x\in[-1,1);&
    \text{D. } \displaystyle \ln(1-x)=\sum_{n=1}^{\infty}\frac{1}{n}x^n,\quad x\in[-1,1).
  \end{array}
  \]
\end{exercise}

\begin{solution}
\par
由几何级数
\[
  \frac1{1-x}=\sum_{n=0}^{\infty}x^n,\qquad |x|<1,
\]
从 \(0\) 到 \(x\) 积分，得
\[
  -\ln(1-x)=\sum_{n=1}^{\infty}\frac{x^n}{n}.
\]
所以
\[
  \ln(1-x)=-\sum_{n=1}^{\infty}\frac{x^n}{n}.
\]
端点 \(x=-1\) 时为 \(\sum(-1)^{n+1}/n\)，收敛；端点 \(x=1\) 时为
\(-\sum1/n\)，发散。故收敛区间为 \([-1,1)\)，选 C。\textit{（校勘：答案页标为 A，与题面函数及端点范围不符。）}
\end{solution}
\begin{exercise}
\par
使得级数
  \(\displaystyle \sum_{n=1}^{\infty}\frac{(-1)^n\ln n}{n^p}\)
  条件收敛的常数 $p$ 的取值范围是（\quad）。
  \par\noindent\begin{tabular}{@{}p{0.48\linewidth}p{0.48\linewidth}@{}}
    A. \(p\le 0\); & B. \(0<p<1\);\\
    C. \(0<p\le 1\); & D. \(p>1\).
  \end{tabular}
\end{exercise}

\begin{solution}
\par
先看交错收敛。令
\(\displaystyle b_n=\frac{\ln n}{n^p}\)。当 \(p>0\) 时，\(b_n\to0\)，并且从某项起单调下降，故
\(\sum(-1)^n b_n\) 由交错级数判别法收敛；当 \(p\le0\) 时，\(b_n\) 不趋于 \(0\)，级数发散。再看绝对收敛：
\[
  \sum\frac{\ln n}{n^p}
\]
由积分判别法可知在 \(p>1\) 时收敛，在 \(p\le1\) 时发散。因此条件收敛范围为
\[
  0<p\le1.
\]
选 C。
\end{solution}

\par\medskip
\noindent\textbf{三、计算题：共 3 题，每题 10 分，共 30 分}\par\nobreak\smallskip
\begin{exercise}
\par
设 $u=f(x,y,z)$，其中函数 $f$ 有二阶连续的偏导数，且 $z=z(x,y)$ 由方程
  \(\displaystyle z^5-5xy+5z=1\)
  所确定，求 $\displaystyle \dfrac{\partial u}{\partial x}$ 和 $\displaystyle \dfrac{\partial^2u}{\partial x^2}$。
\end{exercise}

\begin{solution}
\par
对方程 $z^5-5xy+5z=1$ 两边求全微分，得 \(\displaystyle 5z^4\,\mathrm{d}z-5x\,\mathrm{d}y-5y\,\mathrm{d}x+5\,\mathrm{d}z=0\)。因此 \(\displaystyle \mathrm{d}z=\frac{y\,\mathrm{d}x+x\,\mathrm{d}y}{z^4+1}\)，
  即
  \(\displaystyle \frac{\partial z}{\partial x}=\frac{y}{z^4+1},\qquad
    \frac{\partial z}{\partial y}=\frac{x}{z^4+1}\)。
  由复合函数求导法则，
  \[
    \frac{\partial u}{\partial x}
    =f'_1(x,y,z)+f'_3(x,y,z)\frac{\partial z}{\partial x}
    =f'_1(x,y,z)+f'_3(x,y,z)\frac{y}{z^4+1}.
  \]
  进一步，
  \[
    \frac{\partial^2u}{\partial x^2}
    =f''_{11}(x,y,z)+2f''_{13}(x,y,z)\frac{y}{z^4+1}
    +f''_{33}(x,y,z)\left(\frac{y}{z^4+1}\right)^2
    -f'_3(x,y,z)\frac{4y^2z^3}{(z^4+1)^3}.
  \]
\end{solution}
\begin{exercise}
\par
计算累次积分
  \[
    \int_{\frac14}^{\frac12} \mathrm{d}x\int_{\frac12}^{\sqrt{x}} e^{x/y}\,\mathrm{d}y
    +\int_{\frac12}^{1} \mathrm{d}x\int_x^{\sqrt{x}} e^{x/y}\,\mathrm{d}y.
  \]
\end{exercise}

\begin{solution}
\par
记 $D$ 为由 $\displaystyle y=\dfrac12$、$y=\sqrt{x}$ 和 $y=x$ 围成的区域。原积分化为 \(\displaystyle \iint_D e^{x/y}\,\mathrm{d}x\,\mathrm{d}y=\int_{\frac12}^{1}\mathrm{d}y\int_{y^2}^{y}e^{x/y}\,\mathrm{d}x\)。
  于是
  \[
    \int_{\frac12}^{1}\mathrm{d}y\int_{y^2}^{y}e^{x/y}\,\mathrm{d}x
    =\int_{\frac12}^{1}\left(ey-ye^y\right)\,\mathrm{d}y
    =\frac{3}{8}e-\frac12e^{1/2}.
  \]
\end{solution}
\begin{exercise}
\par
计算球面 $x^2+y^2+z^2=4z$ 含在抛物面 $z=x^2+y^2$ 内的部分的面积。
\end{exercise}

\begin{solution}
\par
球面可写为 \(\displaystyle z=2+\sqrt{4-x^2-y^2}\)。其含在抛物面 $z=x^2+y^2$ 内的部分在 $xOy$ 平面上的投影区域为 \(\displaystyle D=\{(x,y)\mid x^2+y^2\le 3\}\)。曲面面积微元为 \(\displaystyle \mathrm{d}S=\sqrt{1+\left(\frac{\partial z}{\partial x}\right)^2+\left(\frac{\partial z}{\partial y}\right)^2}\,\mathrm{d}x\,\mathrm{d}y=\frac{2}{\sqrt{4-x^2-y^2}}\,\mathrm{d}x\,\mathrm{d}y\)。因此所求面积为 \(\displaystyle A=\iint_D\frac{2}{\sqrt{4-x^2-y^2}}\,\mathrm{d}x\,\mathrm{d}y=\int_0^{2\pi}\mathrm{d}\theta\int_0^{\sqrt3}\frac{2r}{\sqrt{4-r^2}}\,\mathrm{d}r=4\pi\)。
\end{solution}

\par\medskip
\noindent\textbf{四、解答题：共 3 题，每题 10 分，共 30 分}\par\nobreak\smallskip
\begin{exercise}
\par
设 $\Sigma$ 是锥面 $z=\sqrt{x^2+y^2}$ 被平面 $z=0$ 及 $z=1$ 截下的部分的下侧，计算第二类曲面积分 \(\displaystyle \iint_{\Sigma} xye^z\,\mathrm{d}y\,\mathrm{d}z+yz^2\,\mathrm{d}z\,\mathrm{d}x-ye^z\,\mathrm{d}x\,\mathrm{d}y\)。
\end{exercise}

\begin{solution}
\par
考虑平面 $\Sigma_1:z=1$ 在 $x^2+y^2\le 1$ 的部分，取上侧，则 $\Sigma\cup\Sigma_1$ 形成闭曲面，取外侧。记其围成的闭区域为 $\Omega$。由 Gauss 公式，
  \[
    \iint_{\Sigma\cup\Sigma_1} xye^z\,\mathrm{d}y\,\mathrm{d}z+yz^2\,\mathrm{d}z\,\mathrm{d}x-ye^z\,\mathrm{d}x\,\mathrm{d}y
    =\iiint_{\Omega}z^2\,\mathrm{d}x\,\mathrm{d}y\,\mathrm{d}z.
  \]
  而
  \[
    \iiint_{\Omega}z^2\,\mathrm{d}x\,\mathrm{d}y\,\mathrm{d}z
    =\int_0^1 \mathrm{d}z\iint_{x^2+y^2\le z^2}z^2\,\mathrm{d}x\,\mathrm{d}y
    =\int_0^1\pi z^4\,\mathrm{d}z=\frac{\pi}{5}.
  \]
  又 $\Sigma_1$ 的单位正法向量为 $\mathbf n=(0,0,1)$，故
  \[
    \iint_{\Sigma_1} xye^z\,\mathrm{d}y\,\mathrm{d}z+yz^2\,\mathrm{d}z\,\mathrm{d}x-ye^z\,\mathrm{d}x\,\mathrm{d}y
    =\iint_{x^2+y^2\le 1}(-ye)\,\mathrm{d}x\,\mathrm{d}y=0.
  \]
  因此 \(\displaystyle \iint_{\Sigma} xye^z\,\mathrm{d}y\,\mathrm{d}z+yz^2\,\mathrm{d}z\,\mathrm{d}x-ye^z\,\mathrm{d}x\,\mathrm{d}y=\frac{\pi}{5}\)。
\end{solution}
\begin{exercise}
\par
设曲线积分 \(\displaystyle \int_{\Gamma}\left(\sin x-f(x)\right)\frac{y}{x}\,\mathrm{d}x+f(x)\,\mathrm{d}y\) 与路径无关，其中 $f(x)$ 有一阶连续导数且 $f(\pi)=0$，求 $f(x)$ 并计算曲线积分 \(\displaystyle \int_{(1,0)}^{(\pi,\pi)}\left(\sin x-f(x)\right)\frac{y}{x}\,\mathrm{d}x+f(x)\,\mathrm{d}y\)。
\end{exercise}

\begin{solution}
\par
记 \(\displaystyle P=\left(\sin x-f(x)\right)\frac{y}{x},\ Q=f(x)\)。由曲线积分与路径无关，有 $\displaystyle \dfrac{\partial P}{\partial y}=\dfrac{\partial Q}{\partial x}$，即 \(\displaystyle \frac{\sin x}{x}-\frac{f(x)}{x}=f'(x)\)。
  这是一个一阶线性非齐次微分方程，积分因子为 $x$，从而 \(\displaystyle (xf(x))'=xf'(x)+f(x)=\sin x\)。
  其通解为 \(\displaystyle f(x)=-\frac{\cos x}{x}+\frac{C}{x}\)。
  由 $f(\pi)=0$ 得 $C=-1$。因此，\(\displaystyle f(x)=\frac{-\cos x-1}{x}\)。此时 \(\displaystyle P\,\mathrm{d}x+Q\,\mathrm{d}y=\mathrm{d}\left(\frac{(-\cos x-1)y}{x}\right)\)，因此 \(\displaystyle \int_{(1,0)}^{(\pi,\pi)}P\,\mathrm{d}x+Q\,\mathrm{d}y=0\)。
\end{solution}
\begin{exercise}
\par
求幂级数 \(\displaystyle \sum_{n=0}^{\infty}\frac{1}{2n+1}x^{2n+1}\) 的收敛域及和函数。
\end{exercise}

\begin{solution}
\par
考虑幂级数 \(\displaystyle \sum_{n=0}^{\infty}\frac{1}{2n+1}t^n\)。由于
  \begin{align*}
    \lim_{n\to\infty}
    \frac{\frac{1}{2(n+1)+1}}{\frac{1}{2n+1}}
    &=\lim_{n\to\infty}\frac{2n+1}{2n+3}=1,\\
    \sum_{n=0}^{\infty}\frac{1}{2n+1}x^{2n+1}
    &=x\sum_{n=0}^{\infty}\frac{1}{2n+1}(x^2)^n.
  \end{align*}
  因此辅助幂级数与原幂级数的收敛半径均为 $1$。当 $x=\pm1$ 时，级数为 $\displaystyle \pm\sum_{n=0}^{\infty}\dfrac{1}{2n+1}$，发散。因此收敛域为 $(-1,1)$。

  设其和函数为 \(\displaystyle f(x)=\sum_{n=0}^{\infty}\frac{1}{2n+1}x^{2n+1},\ x\in(-1,1)\)。则 \(\displaystyle f'(x)=\sum_{n=0}^{\infty}x^{2n}=\frac{1}{1-x^2}\)。
  因为 $f(0)=0$，所以 \(\displaystyle f(x)=\int_0^x\frac{1}{1-t^2}\,\mathrm{d}t=\frac12\ln\frac{1+x}{1-x},\qquad x\in(-1,1)\)。
\end{solution}

\par\medskip
\noindent\textbf{五、证明题：共 1 题，每题 10 分，共 10 分}\par\nobreak\smallskip
\begin{exercise}
\par
证明函数项级数 \(\displaystyle \sum_{n=0}^{\infty}x^2e^{-nx}\) 在 $[0,+\infty)$ 一致收敛。
\end{exercise}

\begin{solution}
\par
证明一：对任意 $x\in[0,+\infty)$，当 $n\ge 1$ 时，\(\displaystyle e^{nx}>1+nx+\frac12(nx)^2>\frac12(nx)^2\)。因此 \(\displaystyle |x^2e^{-nx}|=\frac{x^2}{e^{nx}}\le \frac{2}{n^2}\)。而数项级数 $\displaystyle \sum_{n=1}^{\infty}\frac{1}{n^2}$ 收敛。由优级数判别法，\(\displaystyle \sum_{n=0}^{\infty}x^2e^{-nx}\) 在 $[0,+\infty)$ 一致收敛。

证明二：该函数项级数的部分和为
\[
  S_n(x)=
  \begin{cases}
    \dfrac{x^2(1-e^{-(n+1)x})}{1-e^{-x}}, & x>0,\\
    0, & x=0.
  \end{cases}
\]
因此
\[
  S(x)=\lim_{n\to\infty}S_n(x)=
  \begin{cases}
    \dfrac{x^2}{1-e^{-x}}, & x>0,\\
    0, & x=0.
  \end{cases}
\]
下证 $S_n(x)$ 一致收敛于 $S(x)$。对任意 $\varepsilon>0$，取 $\displaystyle N>\dfrac1\varepsilon$，当 $n>N$ 时，对任意 $x\in[0,+\infty)$，
\[
  |S_n(x)-S(x)|
  \le \left|\frac{x^2e^{-(n+1)x}}{1-e^{-x}}\right|
  =\left|\frac{x}{e^x-1}\right|\left|\frac{x}{e^{nx}}\right|
  <\frac1n<\varepsilon.
\]
其中用到不等式 $e^x\ge x+1\ge x$。故该函数项级数在 $[0,+\infty)$ 一致收敛。
\end{solution}

\par\medskip
\noindent\textbf{六、应用题：共 1 题，每题 10 分，共 10 分}\par\nobreak\smallskip
\begin{exercise}
\par
将长度为 $2a$ 的铁丝分成三段，分别围成一个正方形、一个圆形和一个正三角形，求三个图形面积之和的最小值。
\end{exercise}

\begin{solution}
\par
设正方形的边长为 $x$，圆形的半径为 $y$，正三角形的边长为 $z$，则三个图形的周长之和为 \(\displaystyle 4x+2\pi y+3z=2a\)。
它们的面积之和为 \(\displaystyle f(x,y,z)=x^2+\pi y^2+\frac{\sqrt3}{4}z^2\)。
构造 Lagrange 函数
\[
  L(x,y,z,\lambda)
  =x^2+\pi y^2+\frac{\sqrt3}{4}z^2
  -\lambda(4x+2\pi y+3z-2a).
\]
极值点满足
\[
  \begin{cases}
    \dfrac{\partial L}{\partial x}=2x-4\lambda=0,\\
    \dfrac{\partial L}{\partial y}=2\pi y-2\pi\lambda=0,\\
    \dfrac{\partial L}{\partial z}=\dfrac{\sqrt3}{2}z-3\lambda=0,\\
    \dfrac{\partial L}{\partial \lambda}=-(4x+2\pi y+3z-2a)=0.
  \end{cases}
  \]
  解得
  \[
    \lambda=\frac{a}{4+\pi+3\sqrt3},\qquad
    x_0=\frac{2a}{4+\pi+3\sqrt3},\qquad
    y_0=\frac{a}{4+\pi+3\sqrt3},\qquad
    z_0=\frac{2\sqrt3a}{4+\pi+3\sqrt3}.
  \]
  所求最小面积为 \(\displaystyle f(x_0,y_0,z_0)=\frac{a^2}{4+\pi+3\sqrt3}\)。
\end{solution}

\section{2017级工科数学分析下 B卷}

\par\medskip
\noindent\textbf{一、填空题：共 5 题，每题 2 分，共 10 分。}\par\nobreak\smallskip
\begin{exercise}
\par
微分方程 $y''+2y'-3y=2-3x$ 的通解为\underline{\hspace{4cm}}。
\end{exercise}

\begin{solution}
\par
先解齐次方程 \(y''+2y'-3y=0\)。其特征方程为
\(\displaystyle r^2+2r-3=0\)，故 \(r=1,-3\)，齐次通解为
\(\displaystyle y_h=C_1e^x+C_2e^{-3x}\)。右端为一次多项式，设特解
\(\displaystyle y_p=ax+b\)。代入得
\(\displaystyle 2a-3(ax+b)=2-3x\)，比较系数可得 \(a=1,\ b=0\)。因此
\(\displaystyle y=C_1e^x+C_2e^{-3x}+x\)，其中 \(C_1,C_2\) 为任意常数。
\end{solution}
\begin{exercise}
\par
向量场 $2xy\vec{i}+e^z\sin y\vec{j}+(x^2+y^2+z^2)\vec{k}$ 的旋度为\underline{\hspace{4cm}}。
\end{exercise}

\begin{solution}
\par
记 \(\vec F=(P,Q,R)\)，其中
\(P=2xy,\ Q=e^z\sin y,\ R=x^2+y^2+z^2\)。旋度按
\[
  \nabla\times\vec F=(R_y-Q_z,\ P_z-R_x,\ Q_x-P_y)
\]
计算，故
  \[
    R_y-Q_z=2y-e^z\sin y,\qquad
    P_z-R_x=-2x,\qquad
    Q_x-P_y=-2x.
  \]
所以
\(\displaystyle \nabla\times\vec F=(2y-e^z\sin y)\vec i-2x\vec j-2x\vec k\)。
\end{solution}
\begin{exercise}
\par
设 $\Gamma$ 是平面上的下半圆周
  \(\displaystyle y=-\sqrt{1-x^2},\ -1\le x\le 1\)，
  则第一类曲线积分 \(\displaystyle \int_\Gamma (x^2+y^2)\,\mathrm{d}s=\underline{\hspace{4cm}}\)。
\end{exercise}

\begin{solution}
\par
曲线 \(\Gamma\) 是单位圆 \(x^2+y^2=1\) 的下半圆周，因此在整条曲线上都有
\(x^2+y^2=1\)。第一类曲线积分与方向无关，只需乘以弧长：
\[
  \int_\Gamma (x^2+y^2)\,\mathrm{d}s=\int_\Gamma 1\,\mathrm{d}s=\operatorname{length}(\Gamma)=\pi .
\]
这里用的是第一类曲线积分，积分元是 \(\mathrm{d}s\)，所以曲线方向不影响结果。若参数化，也可令 \(x=\cos t,y=-\sin t,\ 0\le t\le\pi\)，此时 \(\mathrm{d}s=\mathrm{d}t\)，积分同样等于 \(\int_0^\pi1\,\mathrm{d}t=\pi\)。
\end{solution}
\begin{exercise}
\par
幂级数 $\displaystyle \sum_{n=1}^{\infty}(-1)^n\dfrac{1}{n\cdot 2^n}x^n$ 的收敛域为\underline{\hspace{4cm}}。
\end{exercise}

\begin{solution}
\par
该幂级数的一般项为
\(\displaystyle \frac{(-1)^n}{n}\left(\frac{x}{2}\right)^n\)。由幂级数半径公式或比值法，收敛半径为
\(R=2\)。当 \(x=2\) 时，级数化为
\(\displaystyle \sum_{n=1}^{\infty}\frac{(-1)^n}{n}\)，由交错级数判别法收敛；当
\(x=-2\) 时，级数化为 \(\displaystyle \sum_{n=1}^{\infty}\frac1n\)，发散。因此收敛域为
\(\displaystyle (-2,2]\)。
\end{solution}
\begin{exercise}
\par
设周期为 $2\pi$ 的函数 $f(x)=|x|,\ -\pi<x\le\pi$，则 $f(x)$ 的傅里叶（Fourier）级数在 $x=\pi$ 处收敛于\underline{\hspace{4cm}}。
\end{exercise}

\begin{solution}
\par
按狄利克雷收敛定理，傅里叶级数在间断点或周期端点处收敛到左右极限的平均值。这里
\(f(\pi-0)=|\pi|=\pi\)，而 \(2\pi\) 周期延拓给出
\(f(\pi+0)=f(-\pi+0)=|-\pi|=\pi\)。所以
\[
  S(\pi)=\frac{f(\pi-0)+f(\pi+0)}2=\frac{\pi+\pi}{2}=\pi .
\]
虽然 \(x=\pi\) 是所给基本区间的端点，但周期延拓后的左右极限恰好相等，因此傅里叶级数在该点的和仍为 \(\pi\)。
\end{solution}

\par\medskip
\noindent\textbf{二、选择题：共 5 题，每题 2 分，共 10 分。}\par\nobreak\smallskip
\begin{exercise}
\par
下列微分方程中，属于二阶线性常微分方程的是（\quad）
  \par\noindent\begin{tabular}{@{}>{\raggedright\arraybackslash}p{0.46\linewidth}>{\raggedright\arraybackslash}p{0.46\linewidth}@{}}
  A. $(x+y^2)\,\mathrm{d}y+e^x\,\mathrm{d}x=0$； & B. $xy''+(\sin x)y'+(\ln x)y=\tan x$；\\
  C. $(y'')^2+y=2$； & D. $y''+\ln y=2$。
  \end{tabular}
\end{exercise}

\begin{solution}
\par
二阶线性常微分方程必须能写成
\(\displaystyle a_2(x)y''+a_1(x)y'+a_0(x)y=f(x)\)，未知函数 \(y\) 及其导数只能一次出现，系数只能依赖于自变量 \(x\)。A 是一阶微分形式；C 含有
\((y'')^2\)，不是线性；D 含有 \(\ln y\)，也不是线性。只有 B 满足定义，故选 B。判断线性时允许出现 \(\sin x,\ln x,\tan x\) 这类自变量函数作为系数或右端项，禁止的是未知函数 \(y\) 及其导数非一次出现。
\end{solution}
\begin{exercise}
\par
二元函数 $f(x,y)$ 在点 $(a,b)$ 的两个偏导数
  \(\displaystyle \frac{\partial f}{\partial x}(a,b),\ \frac{\partial f}{\partial y}(a,b)\)
  存在是 $f(x,y)$ 在该点连续的（\quad）
  \par\noindent\begin{tabular}{@{}>{\raggedright\arraybackslash}p{0.22\linewidth}>{\raggedright\arraybackslash}p{0.22\linewidth}>{\raggedright\arraybackslash}p{0.22\linewidth}>{\raggedright\arraybackslash}p{0.22\linewidth}@{}}
  A. 充分而非必要条件； & B. 必要而非充分条件； & C. 充分必要条件； & D. 既非充分也非必要条件。
  \end{tabular}
\end{exercise}

\begin{solution}
\par
偏导数存在只说明沿坐标轴方向的两个一元极限存在，不能控制任意路径趋近时的函数值，因此不能推出连续；反过来，函数在一点连续也不保证沿坐标轴方向的差商极限存在，例如 \(f(x,y)=|x|\) 在原点连续但
\(f_x(0,0)\) 不存在。故“两个偏导数存在”与“连续”既非充分也非必要，选 D。
这道题的易错点是把“一元函数可导推出连续”的结论照搬到多元函数。多元偏导只考察坐标轴方向，连续却要求任意路径同时趋近，因此二者没有充分必要关系。
\end{solution}
\begin{exercise}
\par
曲面 $3x^2+y^2+z^2=12$ 上点 $M(-1,0,3)$ 处的切平面与平面 $z=0$ 的夹角是（\quad）
  \par\noindent\begin{tabular}{@{}>{\raggedright\arraybackslash}p{0.22\linewidth}>{\raggedright\arraybackslash}p{0.22\linewidth}>{\raggedright\arraybackslash}p{0.22\linewidth}>{\raggedright\arraybackslash}p{0.22\linewidth}@{}}
  A. $\displaystyle \dfrac{\pi}{6}$； & B. $\displaystyle \dfrac{\pi}{4}$； & C. $\displaystyle \dfrac{\pi}{3}$； & D. $\displaystyle \dfrac{\pi}{2}$。
  \end{tabular}
\end{exercise}

\begin{solution}
\par
设 \(F(x,y,z)=3x^2+y^2+z^2-12\)，则曲面在一点的法向量为
\(\nabla F=(6x,2y,2z)\)。在 \(M(-1,0,3)\) 处，
\(\nabla F(M)=(-6,0,6)\)。平面 \(z=0\) 的法向量为 \((0,0,1)\)。两平面夹角等于其法向量夹角的锐角，故
\[
  \cos\theta=\frac{|(-6,0,6)\cdot(0,0,1)|}{\sqrt{72}\cdot1}
  =\frac1{\sqrt2},
\]
所以 \(\theta=\pi/4\)，选 B。
\end{solution}
\begin{exercise}
\par
设 $\Gamma$ 是平面曲线 \(\displaystyle \frac{x^2}{a^2}+\frac{y^2}{b^2}=1,\ a,b>0\)，取逆时针方向，则第二类曲线积分 \(\displaystyle \int_\Gamma x\,\mathrm{d}y-y\,\mathrm{d}x\) 等于（\quad）
  \par\noindent\begin{tabular}{@{}>{\raggedright\arraybackslash}p{0.22\linewidth}>{\raggedright\arraybackslash}p{0.22\linewidth}>{\raggedright\arraybackslash}p{0.22\linewidth}>{\raggedright\arraybackslash}p{0.22\linewidth}@{}}
  A. $\pi ab$； & B. $2\pi ab$； & C. $0$； & D. $-2\pi ab$。
  \end{tabular}
\end{exercise}

\begin{solution}
\par
令 \(P=-y,\ Q=x\)。椭圆取逆时针方向，满足 Green 公式正向，于是
\[
  \oint_\Gamma x\,\mathrm{d}y-y\,\mathrm{d}x
  =\iint_D(Q_x-P_y)\,\mathrm{d}A
  =\iint_D(1-(-1))\,\mathrm{d}A
  =2\,\operatorname{Area}(D).
\]
椭圆面积为 \(\pi ab\)，所以积分为 \(2\pi ab\)，选 B。
若题目把方向改成顺时针，结果就要变号；本题取逆时针，正好与 Green 公式的正向一致。
\end{solution}
\begin{exercise}
\par
下列级数发散的是（\quad）
  \par\noindent\begin{tabular}{@{}>{\raggedright\arraybackslash}p{0.46\linewidth}>{\raggedright\arraybackslash}p{0.46\linewidth}@{}}
  A. \(\displaystyle \sum_{n=1}^{\infty}2^n\sin\frac{\pi}{3^n}\)； & B. \(\displaystyle \sum_{n=2}^{\infty}\frac{1}{\ln n}\)；\\
  C. \(\displaystyle \sum_{n=2018}^{\infty}\frac{2^n}{n!}\)； & D. \(\displaystyle \sum_{n=1}^{\infty}\left(1-\frac1n\right)^{n^2}\)。
  \end{tabular}
\end{exercise}

\begin{solution}
\par
A 中
\(\displaystyle 2^n\sin\frac{\pi}{3^n}\sim \pi\left(\frac23\right)^n\)，故收敛；C 中
\(\displaystyle \sum 2^n/n!\) 收敛，去掉有限项仍收敛；D 中
\(\displaystyle \left(1-\frac1n\right)^{n^2}\sim e^{-n}\)，故也收敛。B 的通项虽然趋于 \(0\)，但当
\(n\) 足够大时 \(\ln n<n\)，从而 \(1/\ln n>1/n\)，由比较判别法发散。因此选 B。
\end{solution}

\par\medskip
\noindent\textbf{三、计算题：共 3 题，每题 10 分，共 30 分。}\par\nobreak\smallskip
\begin{exercise}
\par
设 \(\displaystyle z=xf\left(\frac{y}{x}\right)+g(x,xy)\)，其中 $f(t)$ 具有连续的二阶导数，$g(u,v)$ 具有连续的二阶偏导数，计算二阶偏导数 \(\displaystyle \frac{\partial^2 z}{\partial x\partial y}\)。
\end{exercise}

\begin{solution}
\par
先对 \(y\) 求偏导。第一项中 \(x\) 暂作常数，
\[
  \frac{\partial}{\partial y}\left[x f\left(\frac yx\right)\right]
  =x f'\left(\frac yx\right)\frac1x
  =f'\left(\frac yx\right).
\]
第二项 \(g(x,xy)\) 对 \(y\) 求导时只有第二个变量变化，故
\(\displaystyle \partial_y g(x,xy)=xg_2'(x,xy)\)。于是
\[
  z_y=f'\left(\frac yx\right)+xg_2'(x,xy).
\]
再对 \(x\) 求导：第一项给出
\(\displaystyle -\frac{y}{x^2}f''(y/x)\)；第二项用乘积法则和链式法则，得
  \[
    \frac{\partial^2z}{\partial x\partial y}
    =
    -\frac{y}{x^2}f''\left(\frac{y}{x}\right)
    +g_2'(x,xy)+xg_{21}''(x,xy)+xyg_{22}''(x,xy).
  \]
其中 \(g_2',g_{21}'',g_{22}''\) 均在 \((x,xy)\) 处取值。
\end{solution}
\begin{exercise}
\par
计算累次积分
  \[
    \int_{\frac14}^{\frac12}\mathrm{d}x\int_{\frac12}^{\sqrt{x}}\cos\left(\pi\frac{x}{y}\right)\,\mathrm{d}y
    +\int_{\frac12}^{1}\mathrm{d}x\int_x^{\sqrt{x}}\cos\left(\pi\frac{x}{y}\right)\,\mathrm{d}y.
  \]
\end{exercise}

\begin{solution}
\par
先把两个累次积分看成同一区域上的二重积分。第一个区域为
\(\displaystyle \frac14\le x\le\frac12,\ \frac12\le y\le\sqrt x\)；第二个区域为
\(\displaystyle \frac12\le x\le1,\ x\le y\le\sqrt x\)。合并后，区域 \(D\) 可写为
\[
  \frac12\le y\le1,\qquad y^2\le x\le y.
\]
因此原积分等于
  \begin{align*}
    &\int_{\frac14}^{\frac12}\mathrm{d}x\int_{\frac12}^{\sqrt{x}}\cos\left(\pi\frac{x}{y}\right)\,\mathrm{d}y
    +\int_{\frac12}^{1}\mathrm{d}x\int_x^{\sqrt{x}}\cos\left(\pi\frac{x}{y}\right)\,\mathrm{d}y\\
    &\qquad
    =\iint_D\cos\left(\pi\frac{x}{y}\right)\,\mathrm{d}x\,\mathrm{d}y
    =\int_{\frac12}^{1}\mathrm{d}y\int_{y^2}^{y}\cos\left(\pi\frac{x}{y}\right)\,\mathrm{d}x.
  \end{align*}
  对内层积分令 \(u=\pi x/y\)，或直接对 \(x\) 积分：
  \begin{align*}
    \int_{\frac12}^{1}\mathrm{d}y\int_{y^2}^{y}\cos\left(\pi\frac{x}{y}\right)\,\mathrm{d}x
    &=\int_{\frac12}^{1}\frac{y}{\pi}\left[\sin\left(\pi\frac{x}{y}\right)\right]_{x=y^2}^{x=y}\,\mathrm{d}y\\
    &=-\frac1\pi\int_{\frac12}^{1}y\sin(\pi y)\,\mathrm{d}y\\
    &=\frac{1-\pi}{\pi^3}.
  \end{align*}
\end{solution}
\begin{exercise}
\par
计算曲面 $z=\sqrt{x^2+y^2}$ 夹在两柱面 \(\displaystyle x^2+y^2=x,\ x^2+y^2=2x\) 之间的那一部分的面积。
\end{exercise}

\begin{solution}
\par
曲面可写为 \(z=\sqrt{x^2+y^2}\)。先求面积微元：
  \begin{align*}
    \mathrm{d}S
    &=\sqrt{1+\left(\frac{\partial z}{\partial x}\right)^2+\left(\frac{\partial z}{\partial y}\right)^2}\,\mathrm{d}x\,\mathrm{d}y\\
    &=\sqrt{1+\left(\frac{x}{\sqrt{x^2+y^2}}\right)^2+\left(\frac{y}{\sqrt{x^2+y^2}}\right)^2}\,\mathrm{d}x\,\mathrm{d}y\\
    &=\sqrt2\,\mathrm{d}x\,\mathrm{d}y.
  \end{align*}
两柱面在 \(Oxy\) 平面的投影为
\[
  x^2+y^2=x,\qquad x^2+y^2=2x,
\]
即半径分别为 \(1/2\)、\(1\)，圆心分别为 \((1/2,0)\)、\((1,0)\) 的两个圆。夹在两柱面之间的投影区域面积为
\(\displaystyle \pi\cdot1^2-\pi\cdot(1/2)^2=3\pi/4\)。故所求面积为
\[
  A=\iint_D\sqrt2\,\mathrm{d}x\,\mathrm{d}y
  =\sqrt2\cdot\frac{3\pi}{4}
  =\frac{3\sqrt2}{4}\pi .
\]
\end{solution}

\par\medskip
\noindent\textbf{四、解答题：共 3 题，每题 10 分，共 30 分。}\par\nobreak\smallskip
\begin{exercise}
\par
设 $\Sigma$ 是曲线
  \[
    \begin{cases}
      z=\sqrt{y-1},\\
      x=0,
    \end{cases}
    \qquad 1\le y\le 3
  \]
  绕 $y$ 轴旋转一周所成的曲面，其法向量与 $y$ 轴正向夹角恒大于 $\displaystyle \dfrac{\pi}{2}$。求曲面积分
  \[
    \iint_\Sigma (8y+1)x\,\mathrm{d}y\,\mathrm{d}z+2(1-y^2)\,\mathrm{d}z\,\mathrm{d}x+(z^2-2z)\,\mathrm{d}x\,\mathrm{d}y.
  \]
\end{exercise}

\begin{solution}
\par
旋转曲面满足 \(x^2+z^2=y-1\)，且 \(1\le y\le3\)。题设法向与 \(y\) 轴正向夹角大于
\(\pi/2\)，即侧面法向的 \(y\) 分量为负，这正是区域
\[
  \Omega:\quad 1+x^2+z^2\le y\le3
\]
的外侧方向。用平面
\(\displaystyle \Sigma_1:\ y=3,\ x^2+z^2\le2\)
补成闭曲面，\(\Sigma_1\) 取外侧法向 \((0,1,0)\)。令
\[
  P=(8y+1)x,\qquad Q=2(1-y^2),\qquad R=z^2-2z.
\]
由 Gauss 公式，
  \begin{align*}
    &\iint_{\Sigma\cup\Sigma_1}(8y+1)x\,\mathrm{d}y\,\mathrm{d}z+2(1-y^2)\,\mathrm{d}z\,\mathrm{d}x+(z^2-2z)\,\mathrm{d}x\,\mathrm{d}y\\
    &\qquad
    =\iiint_\Omega (4y+2z-1)\,\mathrm{d}x\,\mathrm{d}y\,\mathrm{d}z.
  \end{align*}
这里 \(\Omega\) 关于 \(z=0\) 对称，故 \(\iiint_\Omega z\,\mathrm{d}V=0\)。于是
  \begin{align*}
    \iiint_\Omega (4y+2z-1)\,\mathrm{d}x\,\mathrm{d}y\,\mathrm{d}z
    &=\iiint_\Omega (4y-1)\,\mathrm{d}x\,\mathrm{d}y\,\mathrm{d}z\\
    &=\int_1^3\mathrm{d}y\iint_{x^2+z^2\le y-1}(4y-1)\,\mathrm{d}x\,\mathrm{d}z\\
    &=\int_1^3(4y-1)\pi(y-1)\,\mathrm{d}y
    =\frac{50}{3}\pi.
  \end{align*}
在补面 \(\Sigma_1\) 上 \(y=3\)，且正向为 \(+y\) 方向，所以只有
\(Q\,\mathrm{d}z\,\mathrm{d}x\) 项贡献，\(Q=2(1-9)=-16\)。因此
  \begin{align*}
    \iint_{\Sigma_1}(8y+1)x\,\mathrm{d}y\,\mathrm{d}z+2(1-y^2)\,\mathrm{d}z\,\mathrm{d}x+(z^2-2z)\,\mathrm{d}x\,\mathrm{d}y
    &=\iint_{x^2+z^2\le2}(-16)\,\mathrm{d}x\,\mathrm{d}z=-32\pi,\\
    \iint_\Sigma (8y+1)x\,\mathrm{d}y\,\mathrm{d}z+2(1-y^2)\,\mathrm{d}z\,\mathrm{d}x+(z^2-2z)\,\mathrm{d}x\,\mathrm{d}y
    &=\frac{50}{3}\pi+32\pi
    =\frac{146}{3}\pi.
  \end{align*}
\end{solution}
\begin{exercise}
\par
设曲线积分 \(\displaystyle \int_\Gamma (f(x)-e^x)\sin\frac{y}{2}\,\mathrm{d}x-\frac12 f(x)\cos\frac{y}{2}\,\mathrm{d}y\)
  与路径无关，其中 $f(x)$ 有一阶连续导数且 $f(0)=0$。求 $f(x)$，并计算曲线积分
  \(\displaystyle \int_{(0,0)}^{(1,1)}(f(x)-e^x)\sin\frac{y}{2}\,\mathrm{d}x-\frac12 f(x)\cos\frac{y}{2}\,\mathrm{d}y\)。
\end{exercise}

\begin{solution}
\par
记
\(\displaystyle P=(f(x)-e^x)\sin\frac{y}{2},\ Q=-\frac12f(x)\cos\frac{y}{2}\)。曲线积分与路径无关等价于
\(P_y=Q_x\)，即
\[
  \frac12(f(x)-e^x)\cos\frac y2=-\frac12f'(x)\cos\frac y2.
\]
比较得 \(f'(x)+f(x)=e^x\)。解这个一阶线性方程：
\[
  (e^xf(x))'=e^{2x},
\]
于是
  \(\displaystyle f(x)=\frac12e^x+Ce^{-x}\)。
  由 $f(0)=0$，得 $\displaystyle C=-\dfrac12$，所以
  \(\displaystyle f(x)=\frac12e^x-\frac12e^{-x}\)。
  此时被积表达式为全微分：
  \begin{align*}
    P\,\mathrm{d}x+Q\,\mathrm{d}y
    &=-\frac12(e^x+e^{-x})\sin\frac{y}{2}\,\mathrm{d}x
      -\frac14(e^x-e^{-x})\cos\frac{y}{2}\,\mathrm{d}y\\
    &=\mathrm{d}\left[-\frac12(e^x-e^{-x})\sin\frac{y}{2}\right].
  \end{align*}
  因此积分只与端点有关：
  \begin{align*}
    \int_{(0,0)}^{(1,1)}P\,\mathrm{d}x+Q\,\mathrm{d}y
    &=
    \left[-\frac12(e^x-e^{-x})\sin\frac{y}{2}\right]_{(0,0)}^{(1,1)}\\
    &=-\frac12(e-e^{-1})\sin\frac12.
  \end{align*}
\end{solution}
\begin{exercise}
\par
将函数 $f(x)=\cos^2x$ 展开成 $x$ 的幂级数。
\end{exercise}

\begin{solution}
\par
利用恒等式 \(\displaystyle \cos^2x=\frac{1+\cos2x}{2}\)，再把 \(\cos2x\) 作麦克劳林展开：
  \[
  \cos^2x
  =\frac12+\frac12\sum_{n=0}^{\infty}\frac{(-1)^n}{(2n)!}(2x)^{2n}
  =1+\frac12\sum_{n=1}^{\infty}\frac{(-1)^n4^n}{(2n)!}x^{2n}.
  \]
也可写成
\(\displaystyle \frac12+\frac12\sum_{n=0}^{\infty}\frac{(-1)^n4^n}{(2n)!}x^{2n}\)。由余弦级数的收敛性可知收敛域为
\((-\infty,+\infty)\)。
\end{solution}

\par\medskip
\noindent\textbf{五、证明题：共 1 题，每题 10 分，共 10 分。}\par\nobreak\smallskip
\begin{exercise}
\par
设 $0<c<1$，证明函数项级数 \(\displaystyle \sum_{n=0}^{\infty}(1-x)^n\) 在 $[c,1]$ 一致收敛，但在 $(0,1)$ 不一致收敛。
\end{exercise}

\begin{solution}
\par
证明：当 \(x\in[c,1]\) 时，有 \(0\le1-x\le1-c<1\)，因此
\[
  0\le(1-x)^n\le(1-c)^n.
\]
而几何级数 \(\displaystyle \sum_{n=0}^{\infty}(1-c)^n\) 收敛，所以由 Weierstrass 判别法，
\(\displaystyle \sum_{n=0}^{\infty}(1-x)^n\) 在 \([c,1]\) 上一致收敛。

另一方面，在 \((0,1)\) 上它逐点收敛到几何级数和
\(\displaystyle S(x)=1/x\)。前 \(N+1\) 项部分和为
\[
  S_N(x)=\sum_{k=0}^{N}(1-x)^k
  =\frac{1-(1-x)^{N+1}}{x}.
\]
于是余项为
\[
  |S(x)-S_N(x)|=\frac{(1-x)^{N+1}}{x}.
\]
取 \(\displaystyle x_N=\frac1{N+1}\)，则 \(x_N\in(0,1)\)，并且
\[
  |S(x_N)-S_N(x_N)|
  =(N+1)\left(1-\frac1{N+1}\right)^{N+1}.
\]
该量不趋于 \(0\)，反而趋于无穷大，所以余项不能关于 \(x\in(0,1)\) 一致趋于 \(0\)。因此级数在
\((0,1)\) 上不一致收敛。
\end{solution}

\par\medskip
\noindent\textbf{六、应用题：共 1 题，每题 10 分，共 10 分。}\par\nobreak\smallskip
\begin{exercise}
\par
设椭圆 \(\displaystyle x^2+3y^2=12\) 的内接等腰三角形之底边平行于椭圆的长轴，求这样的三角形的最大面积。
\end{exercise}

\begin{solution}
\par
椭圆的长轴在 \(x\) 轴方向，上顶点为 \((0,2)\)。设等腰三角形的顶点为
\((0,2)\)，底边端点为 \((-x,y)\)、\((x,y)\)，其中 \(x>0\)。底边平行于长轴，且端点在椭圆上，所以
\[
  x^2+3y^2=12.
\]
三角形底长为 \(2x\)，高为 \(2-y\)，面积为
\[
  S(x,y)=\frac12\cdot2x(2-y)=x(2-y).
\]
于是问题化为在约束 \(x^2+3y^2=12\) 下求 \(S=x(2-y)\) 的最大值。设
\(\displaystyle L(x,y,\lambda)=x(2-y)-\lambda(x^2+3y^2-12)\)。
极值点满足
\[
  \begin{cases}
    \dfrac{\partial L}{\partial x}=2-y-2\lambda x=0,\\
    \dfrac{\partial L}{\partial y}=-x-6\lambda y=0,\\
    \dfrac{\partial L}{\partial \lambda}=-(x^2+3y^2-12)=0.
  \end{cases}
\]
解得 \(\displaystyle (x,y,\lambda)=\left(3,-1,\frac12\right),\quad
  \left(-3,-1,-\frac12\right),\quad
  (0,2,0)\)。
其中 \(x>0\) 且能形成非退化三角形的极大值点为 \((x,y)=(3,-1)\)。此时
\[
  S_{\max}=3\,[2-(-1)]=9.
\]
因此最大面积为 \(9\)。
\end{solution}

\section{2018级工科数学分析下 A卷}

\par\medskip
\noindent\textbf{一、填空题（每小题3分，共15分）}\par\nobreak\smallskip
\begin{exercise}
\par
微分方程 \(y''-y=2x\mathrm{e}^x\) 的特解形式为 \(y^*=x(ax+b)\mathrm{e}^x\)。
\end{exercise}

\begin{solution}
\par
对应齐次方程的特征方程为 \(r^2-1=0\)，特征根为 \(r=\pm1\)。右端
\(2x e^x\) 是“一次多项式乘 \(e^x\)”型，而 \(e^x\) 已经是齐次解的一部分，且 \(r=1\) 是单根，所以特解要在通常的
\((ax+b)e^x\) 前再乘一个 \(x\)。因此特解形式为
\(\displaystyle y^*=x(ax+b)e^x\)。
本题只要求“特解形式”，所以不需要继续把 \(a,b\) 代回方程求数值；关键是判断共振次数为一重。
\end{solution}
\begin{exercise}
\par
设函数 \(u(x,y,z)=\mathrm{e}^{x^2+y^2+z^2}\)，则 \(\operatorname{div}(\operatorname{grad}u)=\) \underline{\hspace{5cm}}。
\end{exercise}

\begin{solution}
\par
设 \(r^2=x^2+y^2+z^2\)，则 \(u=e^{r^2}\)。先求
\[
  u_x=2xe^{r^2},\qquad
  u_{xx}=(2+4x^2)e^{r^2}.
\]
同理
\(\displaystyle u_{yy}=(2+4y^2)e^{r^2},\ u_{zz}=(2+4z^2)e^{r^2}\)。因此
\[
  \operatorname{div}(\operatorname{grad}u)=\Delta u
  =u_{xx}+u_{yy}+u_{zz}
  =\bigl[6+4(x^2+y^2+z^2)\bigr]e^{x^2+y^2+z^2}.
\]
\end{solution}
\begin{exercise}
\par
设 \(L\) 是任意一条光滑的闭曲线，取逆时针方向，则
\[
  \oint_L(2xy+\cos x)\,\mathrm{d}x+(x^2+\sin y)\,\mathrm{d}y
  =\underline{\hspace{5cm}}.
\]
\end{exercise}

\begin{solution}
\par
令 \(P=2xy+\cos x,\ Q=x^2+\sin y\)。曲线 \(L\) 是逆时针闭曲线，若其围成区域为 \(D\)，由 Green 公式，
\[
  \oint_L P\,\mathrm{d}x+Q\,\mathrm{d}y
  =\iint_D(Q_x-P_y)\,\mathrm{d}A.
\]
这里 \(Q_x=2x,\ P_y=2x\)，所以 \(Q_x-P_y=0\)，从而积分等于 \(0\)。
也就是说被积微分形式在平面内的 Green 旋度恒为零；闭曲线具体形状不影响结果，只要方向正向且区域内偏导连续即可。
\end{solution}
\begin{exercise}
\par
若级数 \(\displaystyle \sum_{n=1}^{\infty}\frac{(-1)^n(n+1)}{n^p}\) 条件收敛，则常数 \(p\) 的取值范围为 \underline{\hspace{5cm}}。
\end{exercise}

\begin{solution}
\par
记正项部分
\(\displaystyle b_n=\frac{n+1}{n^p}\sim n^{1-p}\)。若要由交错级数判别法得到收敛，至少要有
\(b_n\to0\)，这等价于 \(p>1\)，并且此时 \(b_n\) 从某项起单调下降。再看绝对收敛：
\[
  \sum_{n=1}^{\infty}\left|\frac{(-1)^n(n+1)}{n^p}\right|
  \sim \sum_{n=1}^{\infty}n^{1-p},
\]
该 \(p\)-型级数在 \(1-p<-1\)，即 \(p>2\) 时收敛。因此条件收敛要求“交错收敛但不绝对收敛”，得到
\(\displaystyle 1<p\le2\)。
\end{solution}
\begin{exercise}
\par
设 \(\displaystyle f(x)=\dfrac{x^2}{2}\)，\(0\le x<\pi\)，\(S(x)\) 为 \(f(x)\) 的展成以 \(2\pi\) 为周期的正弦级数的和函数，则
  \(\displaystyle S\left(-\dfrac{\pi}{2}\right)=\) \underline{\hspace{5cm}}。
\end{exercise}

\begin{solution}
\par
在 \([0,\pi)\) 上展开正弦级数，等价于把 \(f(x)=x^2/2\) 作奇延拓后再作
\(2\pi\) 周期傅里叶展开。点 \(-\pi/2\) 是奇延拓函数的连续点，因此和函数取延拓后的函数值：
\[
  S\left(-\frac{\pi}{2}\right)
  =-f\left(\frac{\pi}{2}\right)
  =-\frac12\left(\frac{\pi}{2}\right)^2
  =-\frac{\pi^2}{8}.
\]
\end{solution}

\par\medskip
\noindent\textbf{二、计算题（每小题8分，共32分）}\par\nobreak\smallskip
\begin{exercise}
\par
设 \(\displaystyle z=f\left(x,\dfrac{y^2}{x}\right)\)，其中函数 \(f\) 有二阶连续的偏导数，求
  \(\displaystyle \dfrac{\partial^2 z}{\partial x\partial y}\)。
\end{exercise}

\begin{solution}
\par
记 \(u=x,\ v=y^2/x\)，则 \(z=f(u,v)\)。对 \(y\) 求偏导时 \(u_y=0\)，
\(v_y=2y/x\)，所以
\[
  z_y=f_1u_y+f_2v_y=\frac{2y}{x}f_2.
\]
再对 \(x\) 求导，既要对系数 \(2y/x\) 求导，也要对 \(f_2(u,v)\) 求复合函数导数：
\[
  (f_2)_x=f_{21}u_x+f_{22}v_x=f_{21}-\frac{y^2}{x^2}f_{22}.
\]
于是
  \[
    \frac{\partial^2z}{\partial x\partial y}
    =\frac{2y}{x}f_{12}-\frac{2y^3}{x^3}f_{22}-\frac{2y}{x^2}f_2,
  \]
  其中 \(f_2,f_{12},f_{22}\) 均在 \(\displaystyle \left(x,\dfrac{y^2}{x}\right)\) 处取值。
\end{solution}
\begin{exercise}
\par
计算三重积分 \(\displaystyle \iiint_\Omega \left(x+\frac y2-z\right)^2\,\mathrm{d}x\,\mathrm{d}y\,\mathrm{d}z\)，
  其中 \(\Omega\) 是曲面 \(\displaystyle x^2+\dfrac{y^2}{4}+z^2=1\) 所围成的区域。
\end{exercise}

\begin{solution}
\par
展开平方：
\[
  \left(x+\frac y2-z\right)^2=x^2+\frac{y^2}{4}+z^2+xy-2xz-yz.
\]
椭球区域关于三个坐标面对称，\(xy,xz,yz\) 都是关于某个变量的奇函数，积分为 \(0\)。因此原积分化为
  \[
    \iiint_\Omega\left(x^2+\frac{y^2}{4}+z^2\right)\,\mathrm{d}V.
  \]
  作广义球坐标变换
  \[
    x=r\sin\varphi\cos\theta,\quad
    y=2r\sin\varphi\sin\theta,\quad
    z=r\cos\varphi,\quad J=2r^2\sin\varphi.
  \]
在该变换下 \(x^2+y^2/4+z^2=r^2\)，且
\(0\le r\le1,\ 0\le\varphi\le\pi,\ 0\le\theta\le2\pi\)。因此
  \[
    I=\int_0^{2\pi}\!\!\int_0^\pi\!\!\int_0^1 r^2\cdot2r^2\sin\varphi\,\mathrm{d}r\,\mathrm{d}\varphi\,\mathrm{d}\theta
    =\frac{8\pi}{5}.
  \]
\end{solution}
\begin{exercise}
\par
求密度均匀的曲面 \(\Sigma:z=\sqrt{a^2-x^2-y^2}\) 的质心坐标。
\end{exercise}

\begin{solution}
\par
曲面是半径为 \(a\) 的上半球面，关于 \(x=0\)、\(y=0\) 对称，密度均匀，故
\(\bar x=\bar y=0\)。质心的 \(z\) 坐标为
\[
  \bar z=\frac{\iint_\Sigma z\,\mathrm{d}S}{\iint_\Sigma \mathrm{d}S}.
\]
上半球面面积为 \(2\pi a^2\)。取球面参数
\(\displaystyle x=a\sin\varphi\cos\theta,\ y=a\sin\varphi\sin\theta,\ z=a\cos\varphi\)，
其中 \(0\le\varphi\le\pi/2,\ 0\le\theta\le2\pi\)，且 \(\mathrm{d}S=a^2\sin\varphi\,\mathrm{d}\varphi \mathrm{d}\theta\)。于是
  \[
    \iint_\Sigma z\,\mathrm{d}S
    =\int_0^{2\pi}\!\!\int_0^{\pi/2}a\cos\varphi\cdot a^2\sin\varphi\,\mathrm{d}\varphi\,\mathrm{d}\theta
    =\pi a^3.
  \]
  故质心为 \(\displaystyle \left(0,0,\frac a2\right)\)。
\end{solution}
\begin{exercise}
\par
求微分方程 \(y'\cot x=y\ln y\) 的通解。
\end{exercise}

\begin{solution}
\par
因为方程中有 \(\ln y\)，先在 \(y>0\) 的区间上讨论。由
\[
  y'\cot x=y\ln y
\]
得
\[
  \frac{y'}{y\ln y}=\tan x.
\]
令 \(u=\ln y\)，则 \(u'=y'/y\)，方程化为
\(\displaystyle \frac{u'}u=\tan x\)。积分得
\[
  \ln|u|=\int\tan x\,\mathrm{d}x=-\ln|\cos x|+C,
\]
即 \(u=C\sec x\)。于是
\[
  \ln y=C\sec x,\qquad y=e^{C\sec x}.
\]
当 \(C=0\) 时得到常值解 \(y=1\)，已经包含在上式中。
\end{solution}

\par\medskip
\noindent\textbf{三、解答题（每小题10分，共30分）}\par\nobreak\smallskip
\begin{exercise}
\par
设曲线积分
  \(\displaystyle \int_L xy^2\,\mathrm{d}x+y\varphi(x)\,\mathrm{d}y\)
  与积分路径无关，其中 \(\varphi(x)\) 有一阶连续导数且 \(\varphi(0)=0\)。求 \(\varphi(x)\)，并计算曲线积分
  \(\displaystyle \int_{(0,0)}^{(1,1)}xy^2\,\mathrm{d}x+y\varphi(x)\,\mathrm{d}y\)。
\end{exercise}

\begin{solution}
\par
令 \(P=xy^2,\ Q=y\varphi(x)\)。曲线积分与路径无关等价于 \(P_y=Q_x\)，所以
\[
  2xy=y\varphi'(x).
\]
在恒等意义下得到 \(\varphi'(x)=2x\)。结合 \(\varphi(0)=0\)，得
\(\varphi(x)=x^2\)。此时
\[
  P\,\mathrm{d}x+Q\,\mathrm{d}y=xy^2\,\mathrm{d}x+x^2y\,\mathrm{d}y=\mathrm{d}\left(\frac12x^2y^2\right).
\]
故势函数可取 \(F(x,y)=\frac12x^2y^2\)，所求积分为
\[
  F(1,1)-F(0,0)=\frac12.
\]
\end{solution}
\begin{exercise}
\par
计算
  \[
    I=\oiint_\Sigma
    \frac{x\,\mathrm{d}y\,\mathrm{d}z+y\,\mathrm{d}z\,\mathrm{d}x+z\,\mathrm{d}x\,\mathrm{d}y}
    {(a^2x^2+b^2y^2+c^2z^2)^{3/2}},
  \]
  其中 \(a,b,c>0\)，\(\Sigma:x^2+y^2+z^2=1\)，取外侧。
\end{exercise}

\begin{solution}
\par
在单位球面上，外法向量为 \((x,y,z)\)，故 \(\displaystyle x\,\mathrm{d}y\,\mathrm{d}z+y\,\mathrm{d}z\,\mathrm{d}x+z\,\mathrm{d}x\,\mathrm{d}y=\mathrm{d}S\)。
  于是
  \[
    I=\iint_{x^2+y^2+z^2=1}
    \frac{\mathrm{d}S}{(a^2x^2+b^2y^2+c^2z^2)^{3/2}}.
  \]
  令 \(X=ax,\ Y=by,\ Z=cz\)。从原点看线性变换后的曲面，其立体角面积元满足
\[
  \mathrm{d}\omega=\frac{abc}{(a^2x^2+b^2y^2+c^2z^2)^{3/2}}\,\mathrm{d}S.
\]
这是线性变换对球面立体角的 Jacobian；也可由参数化单位球面后直接计算得到。单位球面经该线性变换后仍包围原点一次，所以总立体角为 \(4\pi\)。因此
\[
  abc\, I=\iint \mathrm{d}\omega=4\pi,\qquad I=\frac{4\pi}{abc}.
\]
\end{solution}
\begin{exercise}
\par
求幂级数 \(\displaystyle \sum_{n=1}^{\infty}\frac{2n-1}{n!}x^{2n}\) 的收敛域及和函数。
\end{exercise}

\begin{solution}
\par
令 \(t=x^2\)。因为 \(\displaystyle \sum_{n=1}^{\infty}\frac{nt^n}{n!}=t\mathrm{e}^t,\ \sum_{n=1}^{\infty}\frac{t^n}{n!}=\mathrm{e}^t-1\)，
  所以 \(\displaystyle \sum_{n=1}^{\infty}\frac{2n-1}{n!}x^{2n}=2x^2\mathrm{e}^{x^2}-(\mathrm{e}^{x^2}-1)=(2x^2-1)\mathrm{e}^{x^2}+1\)。
  收敛域为 \((-\infty,+\infty)\)。
\end{solution}

\par\medskip
\noindent\textbf{四、证明题（每小题7分，共14分）}\par\nobreak\smallskip
\begin{exercise}
\par
设 \(f(x)\in C[0,a]\)，证明：\(\displaystyle 2\int_0^a f(x)\,\mathrm{d}x\int_x^a f(y)\,\mathrm{d}y=\left[\int_0^a f(x)\,\mathrm{d}x\right]^2\)。
\end{exercise}

\begin{solution}
\par
令 \(\displaystyle F(x)=\int_x^a f(y)\,\mathrm{d}y\)。
  则 \(F'(x)=-f(x)\)，从而
  \[
    2\int_0^a f(x)\,\mathrm{d}x\int_x^a f(y)\,\mathrm{d}y
    =2\int_0^a f(x)F(x)\,\mathrm{d}x
    =-2\int_0^aF(x)F'(x)\,\mathrm{d}x.
  \]
  因此 \(\displaystyle -2\int_0^aF(x)F'(x)\,\mathrm{d}x=F(0)^2-F(a)^2=\left[\int_0^a f(x)\,\mathrm{d}x\right]^2\)。
\end{solution}
\begin{exercise}
\par
证明函数项级数
  \(\displaystyle \sum_{n=1}^{\infty}\frac{\arctan(nx)}{n^2}\)
  在 \((-\infty,+\infty)\) 上一致收敛。
\end{exercise}

\begin{solution}
\par
对任意 \(x\in(-\infty,+\infty)\)，
  \(\displaystyle \left|\frac{\arctan(nx)}{n^2}\right|\le \frac{\pi}{2n^2}\)。
  由于 \(\displaystyle \sum_{n=1}^{\infty}\dfrac{\pi}{2n^2}\) 收敛，由 Weierstrass 判别法，原函数项级数在 \((-\infty,+\infty)\) 上一致收敛。
\end{solution}

\par\medskip
\noindent\textbf{五、应用题（本题9分）}\par\nobreak\smallskip
\begin{exercise}
\par
在椭圆 \(x^2+4y^2=4\) 的第一象限部分上求一点，使其到直线 \(2x+3y-6=0\) 的距离最短。
\end{exercise}

\begin{solution}
\par
点到直线 \(2x+3y-6=0\) 的距离为
\[
  d=\frac{|2x+3y-6|}{\sqrt{13}}.
\]
在第一象限椭圆 \(x^2+4y^2=4\) 上，
\[
  2x+3y\le \sqrt{(1+9/4)(x^2+4y^2)}=\sqrt{13}<6,
\]
故绝对值内恒为负，距离最短等价于使 \(2x+3y\) 最大。设
\(F(x,y)=2x+3y\)，约束为 \(x^2+4y^2=4\)。由 Lagrange 乘子法，
\[
  2=2\lambda x,\qquad 3=8\lambda y,\qquad x^2+4y^2=4.
\]
前两式得 \(x=1/\lambda,\ y=3/(8\lambda)\)，代入约束解得
\(\lambda=5/8\)，故
\[
  x=\frac85,\qquad y=\frac35.
\]
因此所求点为 \(\displaystyle \left(\frac85,\frac35\right)\)。
\end{solution}

\section{2018级工科数学分析下 B卷}

\par\medskip
\noindent\textbf{一、填空题（每小题3分，共15分）}\par\nobreak\smallskip
\begin{exercise}
\par
设函数 \(u(x,y,z)=\sqrt{x^2+y^2+z^2}\)，则 \(\operatorname{div}(\operatorname{grad}u)=\) \underline{\hspace{5cm}}。
\end{exercise}

\begin{solution}
\par
当 \((x,y,z)\ne(0,0,0)\) 时，
设 \(r=\sqrt{x^2+y^2+z^2}\)。有
\(\displaystyle \nabla r=\frac{(x,y,z)}{r}\)，故
\[
  \Delta r
  =\frac{\partial}{\partial x}\frac{x}{r}
   +\frac{\partial}{\partial y}\frac{y}{r}
   +\frac{\partial}{\partial z}\frac{z}{r}
  =\left(\frac1r-\frac{x^2}{r^3}\right)
   +\left(\frac1r-\frac{y^2}{r^3}\right)
   +\left(\frac1r-\frac{z^2}{r^3}\right)
  =\frac2r.
\]
所以
\(\displaystyle \operatorname{div}(\operatorname{grad}u)=2/\sqrt{x^2+y^2+z^2}\)。
\end{solution}
\begin{exercise}
\par
交换积分次序，则 \(\displaystyle \int_0^2\mathrm{d}y\int_{y^2}^{4}f(x,y)\,\mathrm{d}x=\) \underline{\hspace{5cm}}。
\end{exercise}

\begin{solution}
\par
原积分区域为
\[
  D=\{(x,y):0\le y\le2,\ y^2\le x\le4\}.
\]
从 \(x\) 的角度看，最小 \(x\) 为 \(0\)，最大 \(x\) 为 \(4\)。固定 \(x\in[0,4]\) 时，由
\(y^2\le x\) 且 \(y\ge0\)，得 \(0\le y\le\sqrt x\)。因此
\[
  \int_0^2\mathrm{d}y\int_{y^2}^{4}f(x,y)\,\mathrm{d}x
  =\int_0^4\mathrm{d}x\int_0^{\sqrt x}f(x,y)\,\mathrm{d}y.
\]
\end{solution}
\begin{exercise}
\par
设曲线 \(\displaystyle L:\dfrac{x^2}{a^2}+\dfrac{y^2}{b^2}=1\)，取顺时针方向，则
  \[
    \oint_L(x-y)\,\mathrm{d}x+(x+y)\,\mathrm{d}y=\underline{\hspace{5cm}}.
  \]
\end{exercise}

\begin{solution}
\par
令 \(P=x-y,\ Q=x+y\)。若 \(L\) 取逆时针方向，则由 Green 公式
\[
  \oint_LP\,\mathrm{d}x+Q\,\mathrm{d}y
  =\iint_D(Q_x-P_y)\,\mathrm{d}A
  =\iint_D(1-(-1))\,\mathrm{d}A
  =2\pi ab.
\]
题中椭圆取顺时针方向，方向与 Green 公式正向相反，所以积分为
\(\displaystyle -2\pi ab\)。
这里先按逆时针正向算出 \(2\pi ab\)，最后再因方向相反取负号；方向账本是第二类曲线积分最容易出错的地方。
\end{solution}
\begin{exercise}
\par
若级数 \(\displaystyle \sum_{n=1}^{\infty}\frac{(-1)^n(n+2)}{n^p}\) 条件收敛，则常数 \(p\) 的取值范围为 \underline{\hspace{5cm}}。
\end{exercise}

\begin{solution}
\par
记
\[
  b_n=\frac{n+2}{n^p}.
\]
若要使交错级数收敛，必须有 \(b_n\to0\)，这等价于 \(p>1\)；并且此时 \(b_n\) 从某项起单调递减，可用交错级数判别法。再看绝对收敛：
\[
  \sum_{n=1}^{\infty}\left|\frac{(-1)^n(n+2)}{n^p}\right|
  \sim \sum_{n=1}^{\infty}n^{1-p},
\]
该级数在 \(p>2\) 时收敛。因此条件收敛要求 \(p>1\) 且不满足 \(p>2\)，即
\[
  1<p\le2.
\]
\end{solution}
\begin{exercise}
\par
设
  \[
    f(x)=
    \begin{cases}
      -1, & -1<x\le0,\\
      x^2+2, & 0<x\le1,
    \end{cases}
  \]
  设 \(S(x)\) 为 \(f(x)\) 展成以 \(2\) 为周期的傅里叶级数的和函数。
  则 \(S(0)=\) \underline{\hspace{5cm}}。
\end{exercise}

\begin{solution}
\par
傅里叶级数在跳跃间断点收敛于左右极限的平均值。这里
\[
  f(0-0)=-1,\qquad f(0+0)=0^2+2=2,
\]
因此
\[
  S(0)=\frac{f(0-0)+f(0+0)}2=\frac{-1+2}{2}=\frac12.
\]
这里 \(S(0)\) 表示傅里叶级数的和函数值，不一定等于原函数在 \(0\) 处被定义的值。遇到分段点先算左右极限，是傅里叶端点题最容易漏的一步。
所以本题的计算完全由左右极限决定，不需要展开出具体的正弦、余弦系数。
\end{solution}

\par\medskip
\noindent\textbf{二、计算题（每小题9分，共36分）}\par\nobreak\smallskip
\begin{exercise}
\par
设 \(u=f(x+y+z,x^2+y^2+z^2)\)，其中函数 \(f\) 有二阶连续的偏导数，求
  \(\displaystyle \frac{\partial^2u}{\partial x^2}+\frac{\partial^2u}{\partial y^2}+\frac{\partial^2u}{\partial z^2}\)。
\end{exercise}

\begin{solution}
\par
记 \(\displaystyle s=x+y+z,\ r=x^2+y^2+z^2\)。有 \(\displaystyle u_x=f_1+2xf_2\)，从而 \(\displaystyle u_{xx}=f_{11}+4xf_{12}+4x^2f_{22}+2f_2\)。
  同理可得 \(u_{yy}\)、\(u_{zz}\)，相加得到
  \[
    u_{xx}+u_{yy}+u_{zz}
    =3f_{11}+4(x+y+z)f_{12}+4(x^2+y^2+z^2)f_{22}+6f_2.
  \]
  各偏导数均在 \((x+y+z,x^2+y^2+z^2)\) 处取值。
\end{solution}
\begin{exercise}
\par
计算三重积分 \(\displaystyle \iiint_\Omega\left(\frac x2-y-z\right)^2\,\mathrm{d}x\,\mathrm{d}y\,\mathrm{d}z\)，其中 \(\Omega\) 是曲面 \(\displaystyle \dfrac{x^2}{4}+y^2+z^2=1\) 所围成的区域。
\end{exercise}

\begin{solution}
\par
展开平方后由椭球区域对称性，交叉项积分为 \(0\)，所以
  \[
    \iiint_\Omega\left(\frac x2-y-z\right)^2\,\mathrm{d}V
    =\iiint_\Omega\left(\frac{x^2}{4}+y^2+z^2\right)\,\mathrm{d}V.
  \]
  作广义球坐标变换
  \[
    x=2r\sin\varphi\cos\theta,\quad
    y=r\sin\varphi\sin\theta,\quad
    z=r\cos\varphi,\quad J=2r^2\sin\varphi.
  \]
  因此
  \[
    I=2\int_0^{2\pi}\int_0^\pi\int_0^1r^4\sin\varphi\,\mathrm{d}r\,\mathrm{d}\varphi\,\mathrm{d}\theta
    =\frac{8\pi}{5}.
  \]
\end{solution}
\begin{exercise}
\par
求抛物面 \(\Sigma:z=2-(x^2+y^2)\) 的质量（\(z\ge0\) 的部分），其密度函数为
  \(\rho(x,y,z)=x^2+y^2\)。
\end{exercise}

\begin{solution}
\par
抛物面为 \(z=2-(x^2+y^2)\)，由 \(z\ge0\) 得投影区域
\(D:x^2+y^2\le2\)。因为
\[
  z_x=-2x,\qquad z_y=-2y,
\]
所以
\[
  \mathrm{d}S=\sqrt{1+z_x^2+z_y^2}\,\mathrm{d}x\,\mathrm{d}y
  =\sqrt{1+4x^2+4y^2}\,\mathrm{d}x\,\mathrm{d}y.
\]
密度为 \(\rho=x^2+y^2\)，于是质量为
  \[
    M=\iint_D(x^2+y^2)\sqrt{1+4x^2+4y^2}\,\mathrm{d}x\,\mathrm{d}y
    =2\pi\int_0^{\sqrt2}r^3\sqrt{1+4r^2}\,\mathrm{d}r
    =\frac{149\pi}{30}.
  \]
\end{solution}
\begin{exercise}
\par
求微分方程
  \(\displaystyle x(2+y)\,\mathrm{d}x+y(1-x)\,\mathrm{d}y=0\)
  的通解。
\end{exercise}

\begin{solution}
\par
原方程分离为
  \(\displaystyle \frac{y\,\mathrm{d}y}{2+y}=\frac{x\,\mathrm{d}x}{x-1}\)。
  两边分别作除法：
\[
  \frac{y}{y+2}=1-\frac2{y+2},\qquad
  \frac{x}{x-1}=1+\frac1{x-1}.
\]
积分得
\[
  y-2\ln|y+2|=x+\ln|x-1|+C.
\]
移项并指数化，也可写为
\[
  (x-1)(y+2)^2=C e^{\,y-x}.
\]
\end{solution}

\par\medskip
\noindent\textbf{三、解答题（每小题10分，共30分）}\par\nobreak\smallskip
\begin{exercise}
\par
设曲线积分
\[
  \int_L\left(\cos x-\varphi(x)\right)\frac yx\,\mathrm{d}x+\varphi(x)\,\mathrm{d}y
\]
与积分路径无关，其中 \(\varphi(x)\) 有一阶连续导数且
  \(\displaystyle \varphi\left(\dfrac{\pi}{2}\right)=0\)，求 \(\varphi(x)\)，并计算曲线积分
\[
  \int_{(1,0)}^{(\pi,\pi)}
  \left(\cos x-\varphi(x)\right)\frac yx\,\mathrm{d}x+\varphi(x)\,\mathrm{d}y .
\]
\end{exercise}

\begin{solution}
\par
取右半平面作为单连通区域。令 \(\displaystyle P=\left(\cos x-\varphi(x)\right)\frac yx,\ Q=\varphi(x)\)。路径无关要求 \(P_y=Q_x\)，即 \(\displaystyle \varphi'(x)+\frac1x\varphi(x)=\frac{\cos x}{x}\)。
  因此 \(\displaystyle (x\varphi(x))'=\cos x,\ x\varphi(x)=\sin x+C\)。由 \(\varphi(\pi/2)=0\)，得 \(C=-1\)，故 \(\displaystyle \varphi(x)=\frac{\sin x-1}{x}\)。
  势函数可取 \(F(x,y)=y\varphi(x)\)，于是 \(\displaystyle \int_{(1,0)}^{(\pi,\pi)}P\,\mathrm{d}x+Q\,\mathrm{d}y=F(\pi,\pi)-F(1,0)=-1\)。
\end{solution}
\begin{exercise}
\par
计算
  \[
    I=\oiint_\Sigma
    \frac{x\,\mathrm{d}y\,\mathrm{d}z+y\,\mathrm{d}z\,\mathrm{d}x+z\,\mathrm{d}x\,\mathrm{d}y}
    {(x^2+y^2+z^2)^{3/2}},
  \]
  其中 \(\Sigma\) 为有界单连通区域的闭边界曲面，取外侧。
\end{exercise}

\begin{solution}
\par
记 \(\Sigma\) 围成的区域为 \(\Omega\)。若原点不在 \(\Omega\) 内，则向量场 \(\displaystyle \mathbf F=\frac{(x,y,z)}{(x^2+y^2+z^2)^{3/2}}\) 在 \(\Omega\) 内无奇点，且散度为 \(0\)，由 Gauss 公式得 \(\displaystyle I=0\)。
  若原点在 \(\Omega\) 内，则挖去小球 \(B_\varepsilon:x^2+y^2+z^2<\varepsilon^2\)。在
\(\Omega\setminus B_\varepsilon\) 上散度仍为 \(0\)，其边界由外边界 \(\Sigma\) 和内侧小球面组成；内侧小球面的外法向是相对于挖空区域的法向，方向指向球心。移项后，\(\Sigma\) 上通量等于小球面按通常外法向的通量：
  \[
    I=\frac1{\varepsilon^3}\oiint_{x^2+y^2+z^2=\varepsilon^2}
    x\,\mathrm{d}y\,\mathrm{d}z+y\,\mathrm{d}z\,\mathrm{d}x+z\,\mathrm{d}x\,\mathrm{d}y
    =4\pi.
  \]
因此
\[
  I=\begin{cases}
  0, & 0\notin\Omega,\\
  4\pi, & 0\in\Omega.
  \end{cases}
\]
\end{solution}
\begin{exercise}
\par
求幂级数 \(\displaystyle \sum_{n=0}^{\infty}\frac{3n+1}{n!}x^{3n}\) 的收敛域及和函数。
\end{exercise}

\begin{solution}
\par
记 \(\displaystyle s(x)=\sum_{n=0}^{\infty}\frac{3n+1}{n!}x^{3n}\)。因为 \(\displaystyle \lim_{n\to\infty}\frac{a_{n+1}}{a_n}=0\)，收敛半径为 \(+\infty\)，收敛域为 \((-\infty,+\infty)\)。又 \(\displaystyle s(x)=3\sum_{n=0}^{\infty}\frac{n(x^3)^n}{n!}+\sum_{n=0}^{\infty}\frac{(x^3)^n}{n!}=(3x^3+1)\mathrm{e}^{x^3}\)。
\end{solution}

\par\medskip
\noindent\textbf{四、证明题（本题10分）}\par\nobreak\smallskip
\begin{exercise}
\par
证明函数项级数 \(\displaystyle \sum_{n=1}^{\infty}\frac{\cos(nx)}{n^2}\) 在 \((-\infty,+\infty)\) 上一致收敛。
\end{exercise}

\begin{solution}
\par
对任意 \(x\in(-\infty,+\infty)\)，有 \(\displaystyle \left|\frac{\cos(nx)}{n^2}\right|\le\frac1{n^2}\)。又 \(\displaystyle \sum_{n=1}^{\infty}1/n^2\) 收敛，由 Weierstrass 判别法，原函数项级数在 \((-\infty,+\infty)\) 上一致收敛。
这里的控制级数与 \(x\) 无关，因此同一个尾项估计同时适用于全体实数 \(x\)，这正是一致收敛而不只是逐点收敛的关键。
\end{solution}

\par\medskip
\noindent\textbf{五、应用题（本题9分）}\par\nobreak\smallskip
\begin{exercise}
\par
欲造一无盖的长方体容器，已知底部造价为每平方米3元，侧面造价均为每平方米1元，现想用36元造一个容积最大的容器，求它的尺寸。
\end{exercise}

\begin{solution}
\par
设长方体的长、宽、高分别为 \(x,y,z\) 米。费用约束为 \(\displaystyle 3xy+2xz+2yz=36\)。最大化 \(V=xyz\)。令 \(\displaystyle F=\ln x+\ln y+\ln z+\lambda(3xy+2xz+2yz-36)\)。
由
\[
  \frac1x+3\lambda y+2\lambda z=0,\quad
  \frac1y+3\lambda x+2\lambda z=0,\quad
  \frac1z+2\lambda x+2\lambda y=0
\]
前两式相减得
\[
  \left(\frac1x-\frac1y\right)+3\lambda(y-x)=0,
\]
即 \((y-x)\left(3\lambda-\dfrac1{xy}\right)=0\)。若 \(3\lambda=1/(xy)\)，代回
\(\displaystyle \frac1x+\lambda(3y+2z)=0\) 会使左端为正，矛盾；故只能有 \(x=y\)。令 \(x=y=t\)，约束变为
\[
  3t^2+4tz=36,\qquad z=\frac{36-3t^2}{4t}.
\]
体积
\[
  V(t)=t^2z=\frac{36t-3t^3}{4}.
\]
令 \(V'(t)=\frac{36-9t^2}{4}=0\)，得 \(t=2\)，从而 \(z=3\)。实际问题中可行域内最大体积存在，故长、宽、高分别为
\(2\) 米、\(2\) 米、\(3\) 米。
\end{solution}

\section{2019级工科数学分析下 A卷}

\par\medskip
\noindent\textbf{一、填空题（共5小题，每题3分，共15分）}\par\nobreak\smallskip
\begin{exercise}
\par
函数 \(u=(2x^2+y^2)z^2\) 在点 \((0,1,1)\) 沿该点梯度方向的方向导数为 \underline{\hspace{5cm}}。
\end{exercise}

\begin{solution}
\par
方向导数在单位方向 \(\vec e\) 上等于 \(\nabla u\cdot\vec e\)。若方向取梯度方向，则单位方向为
\(\displaystyle \vec e=\nabla u/|\nabla u|\)，方向导数达到最大值 \(|\nabla u|\)。这里
\[
  \nabla u=(4xz^2,2yz^2,2(2x^2+y^2)z),
\]
故
\[
  \nabla u(0,1,1)=(0,2,2),\qquad
  |\nabla u(0,1,1)|=\sqrt{0^2+2^2+2^2}=2\sqrt2.
\]
\end{solution}
\begin{exercise}
\par
曲线 \(x=t+1,\ y=t^2+1,\ z=t^3+1\) 在点 \((1,1,1)\) 的切线方程为 \underline{\hspace{5cm}}。
\end{exercise}

\begin{solution}
\par
点 \((1,1,1)\) 对应参数 \(t=0\)。参数曲线的切向量为
\[
  (x'(t),y'(t),z'(t))=(1,2t,3t^2),
\]
在 \(t=0\) 处为 \((1,0,0)\)。因此切线过点 \((1,1,1)\)，方向平行于 \(x\) 轴，可写为
\[
  y=1,\qquad z=1.
\]
若写成对称式，则由于方向向量的 \(y,z\) 分量为 \(0\)，不能硬写成分母为 \(0\) 的比例式；保留 \(y=1,z=1\) 正是表示过该点且平行于 \(x\) 轴的直线。
\end{solution}
\begin{exercise}
\par
函数 \(z=x^3+y^3-3xy\) 的极小值为 \underline{\hspace{5cm}}。
\end{exercise}

\begin{solution}
\par
先求驻点：
\[
  z_x=3x^2-3y=0,\qquad z_y=3y^2-3x=0.
\]
由 \(y=x^2,\ x=y^2\) 得驻点 \((0,0)\)、\((1,1)\)。Hessian 矩阵为
\[
  H=\begin{pmatrix}6x&-3\\-3&6y\end{pmatrix}.
\]
在 \((0,0)\) 处，\(\det H=-9<0\)，不是极值点；在 \((1,1)\) 处，
\(\det H=36-9=27>0\)，且 \(z_{xx}=6>0\)，故为极小点。极小值为
\[
  z(1,1)=1+1-3=-1.
\]
\end{solution}
\begin{exercise}
\par
级数 \(\displaystyle \sum_{n=1}^{\infty}\frac{x^n}{1+x^{2n}}\) 的收敛域为 \underline{\hspace{5cm}}。
\end{exercise}

\begin{solution}
\par
当 \(|x|<1\) 时，有 \(x^{2n}\to0\)，且
\[
  \frac{x^n}{1+x^{2n}}\sim x^n,
\]
故与几何级数同敛散，收敛。当 \(|x|>1\) 时，分母主项为 \(x^{2n}\)，所以
\[
  \frac{x^n}{1+x^{2n}}\sim x^{-n},
\]
仍为收敛的几何型级数。当 \(x=1\) 时通项为 \(1/2\)，当 \(x=-1\) 时通项为
\((-1)^n/2\)，均不趋于 \(0\)，级数发散。故收敛域为
\[
  (-\infty,-1)\cup(-1,1)\cup(1,+\infty).
\]
\end{solution}
\begin{exercise}
\par
设 \(f(x)=\mathrm{e}^{x^2}\)，\(0\le x<\pi\)，\(S(x)\) 为 \(f(x)\) 的展成以 \(2\pi\) 为周期的正弦级数的和函数，则 \(S(0)=\) \underline{\hspace{5cm}}。
\end{exercise}

\begin{solution}
\par
在 \([0,\pi]\) 上展开正弦级数，对应先作奇延拓，再作 \(2\pi\) 周期延拓。奇延拓在
\(0\) 点左右极限分别为
\[
  S(0+)=f(0+)=1,\qquad S(0-)=-f(0+)=-1.
\]
傅里叶级数在端点取左右极限平均值，所以
\[
  S(0)=\frac{1+(-1)}2=0.
\]
正弦级数在端点处特别容易误代 \(f(0)=1\)。实际上正弦展开对应奇延拓，端点两侧极限相反，所以平均值为 \(0\)，这也是正弦级数端点常出现零值的原因。
\end{solution}

\par\medskip
\noindent\textbf{二、计算题（共5小题，每题9分，共45分）}\par\nobreak\smallskip
\begin{exercise}
\par
求微分方程
  \(\displaystyle y''+y'-2y=18x\mathrm{e}^x\)
  的通解。
\end{exercise}

\begin{solution}
\par
特征方程为 \(\displaystyle r^2+r-2=0\)，
  特征根为 \(r_1=-2,\ r_2=1\)，齐次解为 \(C_1\mathrm{e}^{-2x}+C_2\mathrm{e}^x\)。因 \(1\) 是特征根，设特解
  \(\displaystyle y^*=x\mathrm{e}^x(Ax+B)\)。将 \(y^*\) 代入算子
\(\displaystyle L[y]=y''+y'-2y\)，整理得
\[
  L[y^*]=e^x\bigl(6Ax+3B+2A\bigr).
\]
与右端 \(18xe^x\) 比较系数，得
\[
  6A=18,\qquad 3B+2A=0,
\]
所以 \(A=3,\ B=-2\)。通解为
\[
  y=C_1e^{-2x}+C_2e^x+x(3x-2)e^x.
\]
\end{solution}
\begin{exercise}
\par
设 \(z=f(x^2+y^2,x^2y^2)\)，其中 \(f\) 为任意阶可微函数，求
  \(\displaystyle \frac{\partial^2z}{\partial x^2},\ \frac{\partial^2z}{\partial y\partial x}\)。
\end{exercise}

\begin{solution}
\par
记 \(\displaystyle u=x^2+y^2,\ v=x^2y^2\)。先有 \(\displaystyle z_x=2xf_1+2xy^2f_2\)。
  继续求导得
  \begin{align*}
    z_{xx}&=2f_1+2y^2f_2+4x^2f_{11}+8x^2y^2f_{12}+4x^2y^4f_{22},\\
    z_{yx}&=4xyf_{11}+4xy(x^2+y^2)f_{12}+4xyf_2+4x^3y^3f_{22}.
  \end{align*}
  其中各偏导数均在 \((x^2+y^2,x^2y^2)\) 处取值。
\end{solution}
\begin{exercise}
\par
计算
  \(\displaystyle \iiint_\Omega (x+y+z)\,\mathrm{d}x\,\mathrm{d}y\,\mathrm{d}z\)，
  其中 \(\Omega\) 是曲面 \(z=\sqrt{x^2+y^2}\) 与 \(z=\sqrt{1-x^2-y^2}\) 所围成的区域。
\end{exercise}

\begin{solution}
\par
用球坐标 \(\displaystyle x=\rho\sin\varphi\cos\theta,\ y=\rho\sin\varphi\sin\theta,\ z=\rho\cos\varphi\)。两个边界分别为 \(\varphi=\pi/4\) 与 \(\rho=1\)，故 \(\displaystyle 0\le\theta\le2\pi,\ 0\le\varphi\le\frac{\pi}{4},\ 0\le\rho\le1\)。
由于 \(x,y\) 关于区域对称积分为 \(0\)，只需积 \(z\)：\(\displaystyle I=2\pi\int_0^{\pi/4}\int_0^1 \rho\cos\varphi\cdot\rho^2\sin\varphi\,\mathrm{d}\rho\,\mathrm{d}\varphi=\frac{\pi}{8}\)。
\end{solution}
\begin{exercise}
\par
计算第二类曲线积分
  \(\displaystyle \oint_L xy^2\,\mathrm{d}y-x^2y\,\mathrm{d}x\)，
  其中 \(\displaystyle L:\dfrac{x^2}{a^2}+\dfrac{y^2}{b^2}=1\)，\(a,b>0\)，取逆时针方向。
\end{exercise}

\begin{solution}
\par
令 \(\displaystyle x=a\cos t,\ y=b\sin t,\ 0\le t\le2\pi\)。
也可直接用 Green 公式。把积分写成
\[
  \oint_L P\,\mathrm{d}x+Q\,\mathrm{d}y,\qquad P=-x^2y,\quad Q=xy^2.
\]
椭圆取逆时针方向，所以
\[
  \oint_L xy^2\,\mathrm{d}y-x^2y\,\mathrm{d}x
  =\iint_D(Q_x-P_y)\,\mathrm{d}A
  =\iint_D(x^2+y^2)\,\mathrm{d}A.
\]
在椭圆 \(D:x^2/a^2+y^2/b^2\le1\) 中，
\(\displaystyle \iint_Dx^2\,\mathrm{d}A=\frac{\pi a^3b}{4}\)，
\(\displaystyle \iint_Dy^2\,\mathrm{d}A=\frac{\pi ab^3}{4}\)。因此
\[
  I=\frac{\pi}{4}(a^3b+ab^3).
\]
\end{solution}
\begin{exercise}
\par
计算
  \[
    I=\iint_\Sigma 2x(y-z)\,\mathrm{d}y\,\mathrm{d}z+(1-y^2)\,\mathrm{d}z\,\mathrm{d}x+(1+z^2)\,\mathrm{d}x\,\mathrm{d}y,
  \]
  其中 \(\Sigma\) 为柱面 \(x^2+y^2=1\) 夹在平面 \(z=0\) 和 \(z=3\) 部分的曲面，方向指向外侧。
\end{exercise}

\begin{solution}
\par
添加底面 \(\Sigma_1:z=0\) 和顶面 \(\Sigma_2:z=3\)，构成闭曲面。令 \(\displaystyle P=2x(y-z),\ Q=1-y^2,\ R=1+z^2\)。因为 \(\displaystyle P_x+Q_y+R_z=2(y-z)-2y+2z=0\)，
  闭曲面总通量为 \(0\)。底面 \(z=0\) 的外法向向下，只有 \(R\,\mathrm{d}x\,\mathrm{d}y\) 项贡献，通量为
\[
  I_{\Sigma_1}=-\iint_{x^2+y^2\le1}R\big|_{z=0}\,\mathrm{d}x\,\mathrm{d}y
  =-\iint_{x^2+y^2\le1}1\,\mathrm{d}x\,\mathrm{d}y=-\pi.
\]
顶面 \(z=3\) 的外法向向上，通量为
\[
  I_{\Sigma_2}=\iint_{x^2+y^2\le1}R\big|_{z=3}\,\mathrm{d}x\,\mathrm{d}y
  =\iint_{x^2+y^2\le1}10\,\mathrm{d}x\,\mathrm{d}y=10\pi.
\]
因此侧面通量为
\[
  I=0-I_{\Sigma_1}-I_{\Sigma_2}=0-(-\pi)-10\pi=-9\pi.
\]
\end{solution}

\par\medskip
\noindent\textbf{三、解答下列各题（共2小题，每题8分，共16分）}\par\nobreak\smallskip
\begin{exercise}
\par
试确定 \(\alpha\) 的值，使得函数
  \[
    f(x,y)=
    \begin{cases}
      (x^2+y^2)^\alpha\sin\dfrac1{x^2+y^2}, & (x,y)\ne(0,0),\\
      0, & (x,y)=(0,0)
    \end{cases}
  \]
  在点 \((0,0)\) 处可微。
\end{exercise}

\begin{solution}
\par
由
  \(\displaystyle |f(x,y)|\le (x^2+y^2)^\alpha=r^{2\alpha}\)
  可知若 \(\displaystyle \alpha>\dfrac12\)，则 \(\displaystyle \frac{|f(x,y)-f(0,0)|}{\sqrt{x^2+y^2}}\le r^{2\alpha-1}\to0\)。此时 \(f_x(0,0)=f_y(0,0)=0\)，且函数在原点可微。若 \(\displaystyle \alpha\le\dfrac12\)，上述商不趋于 \(0\)，不可微。因此 \(\displaystyle \alpha>\frac12\)。
\end{solution}
\begin{exercise}
\par
求正项级数 \(\displaystyle \sum_{n=0}^{\infty}\frac{2n+1}{n!}\) 的和。
\end{exercise}

\begin{solution}
\par
把分子拆开：
\[
  \sum_{n=0}^{\infty}\frac{2n+1}{n!}
  =2\sum_{n=0}^{\infty}\frac{n}{n!}
   +\sum_{n=0}^{\infty}\frac1{n!}.
\]
其中 \(n=0\) 时第一项为 \(0\)，且
\[
  \sum_{n=1}^{\infty}\frac{n}{n!}
  =\sum_{n=1}^{\infty}\frac1{(n-1)!}=e,\qquad
  \sum_{n=0}^{\infty}\frac1{n!}=e.
\]
故原级数的和为 \(2e+e=3e\)。
\end{solution}

\par\medskip
\noindent\textbf{四、证明题（共2小题，每题8分，共16分）}\par\nobreak\smallskip
\begin{exercise}
\par
设 \(x=x(y,z)\)，\(y=y(x,z)\)，\(z=z(x,y)\)，都是由方程 \(F(x,y,z)=0\) 所确定的具有连续偏导数的函数，证明
  \(\displaystyle \frac{\partial x}{\partial y}\cdot\frac{\partial y}{\partial z}\cdot\frac{\partial z}{\partial x}=-1\)。
\end{exercise}

\begin{solution}
\par
由隐函数求导公式，\(\displaystyle \frac{\partial x}{\partial y}=-\frac{F_y}{F_x},\quad \frac{\partial y}{\partial z}=-\frac{F_z}{F_y},\quad \frac{\partial z}{\partial x}=-\frac{F_x}{F_z}\)。三式相乘即得 \(\displaystyle \frac{\partial x}{\partial y}\cdot \frac{\partial y}{\partial z}\cdot \frac{\partial z}{\partial x}=-1\)。
\end{solution}
\begin{exercise}
\par
证明函数项级数
  \(\displaystyle \sum_{n=2}^{\infty}\frac{(-1)^n(1-\mathrm{e}^{-nx})}{n^2+\sin x}\)
  在 \([0,+\infty)\) 上一致收敛。
\end{exercise}

\begin{solution}
\par
对任意 \(x\in[0,+\infty)\) 且 \(n\ge2\)，有 \(\displaystyle \left|\frac{(-1)^n(1-\mathrm{e}^{-nx})}{n^2+\sin x}\right|\le\frac1{n^2-1}\)。而数项级数 \(\displaystyle \sum_{n=2}^{\infty}\dfrac1{n^2-1}\) 收敛，由 Weierstrass 判别法，原函数项级数在 \([0,+\infty)\) 上一致收敛。
\end{solution}

\par\medskip
\noindent\textbf{五、应用题（共1题，每题8分，共8分）}\par\nobreak\smallskip
\begin{exercise}
\par
将正数 \(9\) 分成三个正数 \(x,y,z\) 之和，使得 \(u=x^2y^3z^4\) 最大。
\end{exercise}

\begin{solution}
\par
在约束 \(x+y+z=9\) 下最大化 \(u=x^2y^3z^4\)。等价于最大化 \(\displaystyle \ln u=2\ln x+3\ln y+4\ln z\)。由 Lagrange 乘子法，\(\displaystyle \frac2x=\lambda,\ \frac3y=\lambda,\ \frac4z=\lambda\)。故 \(\displaystyle x:y:z=2:3:4\)。又 \(x+y+z=9\)，所以 \(\displaystyle x=2,\ y=3,\ z=4\)。最大值为 \(\displaystyle 2^2\cdot3^3\cdot4^4=27648\)。
\end{solution}

\section{2019级工科数学分析下 B卷补考}

\par\medskip
\noindent\textbf{一、填空题（共5小题，每题3分，共15分）}\par\nobreak\smallskip
\begin{exercise}
\par
函数 \(z=x^2+2y^2\) 在点 \((1,1)\) 沿该点梯度方向的方向导数为 \underline{\hspace{5cm}}。
\end{exercise}

\begin{solution}
\par
函数 \(z=x^2+2y^2\) 的梯度为
\[
  \nabla z=(2x,4y).
\]
在 \((1,1)\) 处，\(\nabla z(1,1)=(2,4)\)。沿梯度方向的方向导数等于梯度模长，因此
\[
  D_{\nabla z/|\nabla z|}z=\|\nabla z(1,1)\|
  =\sqrt{2^2+4^2}=2\sqrt5.
\]
方向导数公式为 \(D_{\mathbf e}z=\nabla z\cdot\mathbf e\)。当单位方向 \(\mathbf e\) 与梯度同向时取最大值，所以题目说“沿梯度方向”时可以直接取梯度模长。
\end{solution}
\begin{exercise}
\par
曲线 \(x=t,\ y=t^2,\ z=t^3\) 在点 \((1,1,1)\) 的法平面方程为 \underline{\hspace{5cm}}。
\end{exercise}

\begin{solution}
\par
点 \((1,1,1)\) 对应 \(t=1\)。曲线切向量为
\[
  (x'(t),y'(t),z'(t))=(1,2t,3t^2),
\]
在 \(t=1\) 处为 \((1,2,3)\)。法平面垂直于切向量，并经过 \((1,1,1)\)，所以
\[
  (1,2,3)\cdot(x-1,y-1,z-1)=0,
\]
即
  \(\displaystyle (x-1)+2(y-1)+3(z-1)=0\)，
  即 \(x+2y+3z-6=0\)。
\end{solution}
\begin{exercise}
\par
函数 \(z=xy\) 在条件 \(x^2+y^2=1\) 下的极大值为 \underline{\hspace{5cm}}。
\end{exercise}

\begin{solution}
\par
在约束 \(x^2+y^2=1\) 下，由
\[
  (x-y)^2\ge0
\]
得 \(2xy\le x^2+y^2=1\)，所以
\[
  xy\le\frac12.
\]
当 \(x=y=1/\sqrt2\) 时等号成立，故极大值为 \(\displaystyle \frac12\)。这个做法等价于在圆上令 \(x=\cos t,y=\sin t\)，此时 \(xy=\frac12\sin2t\)，最大值同样为 \(\frac12\)；等号点也说明该上界确实能取到。
\end{solution}
\begin{exercise}
\par
级数 \(\displaystyle \sum_{n=1}^{\infty}\frac{x^n}{n\cdot4^n}\) 的和函数为 \underline{\hspace{5cm}}。
\end{exercise}

\begin{solution}
\par
令 \(t=x/4\)。当 \(|t|<1\) 时，经典展开式
\[
  -\ln(1-t)=\sum_{n=1}^{\infty}\frac{t^n}{n}
\]
给出
\[
  \sum_{n=1}^{\infty}\frac{x^n}{n4^n}
  =-\ln\left(1-\frac x4\right).
\]
端点 \(x=-4\) 时为 \(\sum(-1)^n/n=-\ln2\)，与上式极限一致；端点
\(x=4\) 时为调和级数发散。因此和函数为
\(\displaystyle -\ln(1-x/4)\)，收敛域为 \([-4,4)\)。
\end{solution}
\begin{exercise}
\par
设 \(S(x)\) 为函数
\[
  f(x)=\mathrm{e}^{\cos x},\qquad x\in[0,\pi],
\]
展成的以 \(2\pi\) 为周期的余弦级数的和函数，则 \(S(0)=\) \underline{\hspace{5cm}}。
\end{exercise}

\begin{solution}
\par
在 \([0,\pi]\) 上展开余弦级数，等价于把函数作偶延拓后再作 \(2\pi\) 周期展开。点 \(x=0\) 是偶延拓后的连续点，因此傅里叶级数收敛到函数值本身：
\[
  S(0)=f(0)=e^{\cos0}=e.
\]
余弦级数对应偶延拓，故 \(0\) 点两侧极限相同；这与正弦级数对应奇延拓的端点取值规则不同。
因此本题虽然也问端点 \(x=0\)，但它不是跳跃平均成 \(0\) 的情形；偶延拓后左右极限都等于原函数右端极限 \(f(0)=e\)，所以和函数取 \(e\)。
\end{solution}

\par\medskip
\noindent\textbf{二、计算题（共5小题，每题9分，共45分）}\par\nobreak\smallskip
\begin{exercise}
\par
求微分方程
  \(\displaystyle \frac{\mathrm{d}y}{\mathrm{d}x}=\frac{2x-y+2}{x+2y-4}\)
  的通解。
\end{exercise}

\begin{solution}
\par
分子、分母中的两条直线 \(2x-y+2=0\)、\(x+2y-4=0\) 交于
\((0,2)\)，故作平移
\[
  X=x,\qquad Y=y-2.
\]
由于 \(\mathrm{d}Y/\mathrm{d}X=\mathrm{d}y/\mathrm{d}x\)，方程化为齐次方程
\[
  \frac{\mathrm{d}Y}{\mathrm{d}X}=\frac{2X-Y}{X+2Y}.
\]
写成微分形式：
\[
  (2X-Y)\,\mathrm{d}X-(X+2Y)\,\mathrm{d}Y=0.
\]
这里 \(M=2X-Y,\ N=-(X+2Y)\)，有 \(M_Y=-1=N_X\)，所以是恰当方程。取势函数
\(\Phi_X=2X-Y\)，得
\[
  \Phi=X^2-XY+\psi(Y).
\]
由 \(\Phi_Y=-X+\psi'(Y)=-(X+2Y)\)，得 \(\psi'(Y)=-2Y\)，于是
\(\psi(Y)=-Y^2\)。故
\[
  X^2-XY-Y^2=C.
\]
还原 \(X=x,\ Y=y-2\)，通解为
\[
  x^2-x(y-2)-(y-2)^2=C.
\]
\end{solution}
\begin{exercise}
\par
假设隐函数存在定理的条件都满足，且 \(y(t)\) 由方程组
  \[
    \begin{cases}
      y=f(x,t),\\
      x=g(y,t)
    \end{cases}
  \]
  确定，求 \(\displaystyle \dfrac{\mathrm{d}y}{\mathrm{d}t}\)。
\end{exercise}

\begin{solution}
\par
把 \(x,y\) 都看成 \(t\) 的函数。对
\[
  y=f(x,t),\qquad x=g(y,t)
\]
分别关于 \(t\) 求导，得
\[
  y'=f_xx'+f_t,\qquad x'=g_yy'+g_t.
\]
将第二式代入第一式：
\[
  y'=f_x(g_yy'+g_t)+f_t
  =f_xg_y\,y'+f_xg_t+f_t.
\]
因此
\[
  (1-f_xg_y)y'=f_t+f_xg_t,
\]
在隐函数存在定理保证分母不为 \(0\) 的情形下，
\[
  \frac{\mathrm{d}y}{\mathrm{d}t}=\frac{f_t+f_xg_t}{1-f_xg_y}.
\]
\end{solution}
\begin{exercise}
\par
计算
  \(\displaystyle \iiint_\Omega (x^2+y^2+z^2)\,\mathrm{d}x\,\mathrm{d}y\,\mathrm{d}z\)，
  其中 \(\Omega\) 是曲面 \(z=x^2+y^2\) 与 \(z=1-x^2-y^2\) 所围成的区域。
\end{exercise}

\begin{solution}
\par
用柱坐标。区域为 \(\displaystyle 0\le r\le \frac1{\sqrt2},\ r^2\le z\le 1-r^2\)，且 \(\displaystyle 0\le\theta\le2\pi\)。所以
  \[
  I=2\pi\int_0^{1/\sqrt2}\int_{r^2}^{1-r^2}(r^2+z^2)r\,\mathrm{d}z\,\mathrm{d}r.
  \]
先对 \(z\) 积分：
\[
  \int_{r^2}^{1-r^2}(r^2+z^2)\,\mathrm{d}z
  =r^2(1-2r^2)+\frac{(1-r^2)^3-r^6}{3}.
\]
再对 \(r\) 从 \(0\) 到 \(1/\sqrt2\) 积分，得
\[
  I=\frac{11\pi}{96}.
\]
\end{solution}
\begin{exercise}
\par
计算曲面积分
  \(\displaystyle \iint_\Sigma z\,\mathrm{d}S\)，
  其中 \(\Sigma\) 是球面 \(x^2+y^2+z^2=4\) 被平面 \(z=\sqrt2\) 截出的顶部。
\end{exercise}

\begin{solution}
\par
在球面 \(x^2+y^2+z^2=4\) 上取球坐标 \(z=2\cos\varphi\)，顶部对应
\(0\le\varphi\le\pi/4\)。球面半径为 \(2\)，故
\(\mathrm{d}S=4\sin\varphi\,\mathrm{d}\varphi \mathrm{d}\theta\)。于是
\[
  \iint_\Sigma z\,\mathrm{d}S
  =\int_0^{2\pi}\int_0^{\pi/4}2\cos\varphi\cdot4\sin\varphi\,\mathrm{d}\varphi \mathrm{d}\theta
  =16\pi\int_0^{\pi/4}\sin\varphi\cos\varphi\,\mathrm{d}\varphi
  =4\pi.
\]
\end{solution}
\begin{exercise}
\par
计算
  \(\displaystyle I=\iint_{\Sigma^+}\frac{\cos(\vec r,\vec n)}{r^2}\,\mathrm{d}S\)，
  其中 \(\Sigma^+\) 为不经过 \(P(a,b,c)\) 的光滑闭曲面，方向向外，\(\Sigma\) 不包含点 \(P(a,b,c)\)。
  \(\vec n\) 为曲面 \(\Sigma\) 上任一点 \(M(x,y,z)\) 处的外法向量的单位向量，\(\cos(\vec r,\vec n)\) 表示两向量 \(\vec r,\vec n\) 夹角的余弦，且
  \(\displaystyle \vec r=\overrightarrow{MP}=(a-x,b-y,c-z),\quad
    r=|\vec r|=\sqrt{(a-x)^2+(b-y)^2+(c-z)^2}\)。
\end{exercise}

\begin{solution}
\par
被积函数可看作向量场 \(\displaystyle \mathbf F(M)=\frac{\overrightarrow{MP}}{|\overrightarrow{MP}|^3}\) 通过 \(\Sigma\) 的通量。若点 \(P\) 在 \(\Sigma\) 所围区域外，则 \(\mathbf F\) 在该区域内无奇点且散度为 \(0\)，由 Gauss 公式可得 \(\displaystyle I=0\)。若点 \(P\) 在 \(\Sigma\) 所围区域内，由于 \(\vec r=\overrightarrow{MP}=P-M\) 与标准径向场方向相反，通量为 \(\displaystyle I=-4\pi\)。
\end{solution}

\par\medskip
\noindent\textbf{三、解答下列各题（共2小题，每题8分，共16分）}\par\nobreak\smallskip
\begin{exercise}
\par
试确定 \(\alpha\) 的值，使得函数
  \[
    f(x,y)=
    \begin{cases}
      \dfrac{x^\alpha y^\alpha}{x^2+y^2}, & (x,y)\ne(0,0),\\
      0, & (x,y)=(0,0)
    \end{cases}
  \]
  在点 \((0,0)\) 处可微。
\end{exercise}

\begin{solution}
\par
由极坐标估计，
  \(\displaystyle |f(x,y)|\le C r^{2\alpha-2}\)。
  若在原点可微且导数为零，需要 \(\displaystyle \frac{|f(x,y)|}{r}\le C r^{2\alpha-3}\to0\)。因此 \(\displaystyle \alpha>\frac32\)。
  此时 \(f_x(0,0)=f_y(0,0)=0\)，且 \(f(x,y)=o(\sqrt{x^2+y^2})\)，故函数在原点可微。
\end{solution}
\begin{exercise}
\par
将 \(\ln(1-3x+2x^2)\) 展成麦克劳林级数。
\end{exercise}

\begin{solution}
\par
因为 \(\displaystyle 1-3x+2x^2=(1-x)(1-2x)\)，
  所以
\[
  \ln(1-3x+2x^2)=\ln(1-x)+\ln(1-2x).
\]
利用
\[
  \ln(1-t)=-\sum_{n=1}^{\infty}\frac{t^n}{n},\qquad |t|<1,
\]
分别取 \(t=x\) 与 \(t=2x\)，得
\[
  \ln(1-3x+2x^2)
  =-\sum_{n=1}^{\infty}\frac{x^n}{n}
   -\sum_{n=1}^{\infty}\frac{(2x)^n}{n}
  =-\sum_{n=1}^{\infty}\frac{1+2^n}{n}x^n.
\]
两部分共同要求 \(|x|<1\) 且 \(|2x|<1\)，故展开区间为
\(|x|<1/2\)。
\end{solution}

\par\medskip
\noindent\textbf{四、证明题（共2小题，每题8分，共16分）}\par\nobreak\smallskip
\begin{exercise}
\par
设 \(f\) 为连续函数，\(C\) 为分段光滑的简单闭曲线，证明 \(\displaystyle \oint_C f(x^2-y^2)(x\,\mathrm{d}x-y\,\mathrm{d}y)=0\)。
\end{exercise}

\begin{solution}
\par
取原函数 \(F'(u)=f(u)\)。由于
  \(\displaystyle \mathrm{d}(x^2-y^2)=2x\,\mathrm{d}x-2y\,\mathrm{d}y\)，
  故
  \(\displaystyle f(x^2-y^2)(x\,\mathrm{d}x-y\,\mathrm{d}y)
    =\frac12\,\mathrm{d}F(x^2-y^2)\)。
  闭曲线上的全微分积分为 \(0\)，因此结论成立。
这里 \(f\) 连续保证原函数 \(F\) 存在；被积表达式已经化成某个函数的全微分，所以沿任意闭曲线一周的端点值相同，积分必为零。
\end{solution}
\begin{exercise}
\par
证明函数项级数
  \(\displaystyle \sum_{n=1}^{\infty}\frac{(-1)^n(1-\mathrm{e}^{-nx})}{n^2+x^2}\)
  在 \([0,+\infty)\) 上一致收敛。
\end{exercise}

\begin{solution}
\par
对 \(x\in[0,+\infty)\)，有
  \(\displaystyle \left|\frac{(-1)^n(1-\mathrm{e}^{-nx})}{n^2+x^2}\right|
    \le\frac1{n^2}\)。
  因为 \(\displaystyle \sum_{n=1}^{\infty}1/n^2\) 收敛，由 Weierstrass 判别法，原函数项级数在 \([0,+\infty)\) 上一致收敛。
不等式中用到了 \(0\le 1-e^{-nx}\le1\) 与 \(n^2+x^2\ge n^2\)。右端控制项不含 \(x\)，所以可直接得到整个区间上的一致收敛。
\end{solution}

\par\medskip
\noindent\textbf{五、应用题（共1题，每题8分，共8分）}\par\nobreak\smallskip
\begin{exercise}
\par
求点 \(\displaystyle \left(2,2,\dfrac12\right)\) 到抛物面 \(z=x^2+y^2\) 的最小距离。
\end{exercise}

\begin{solution}
\par
设抛物面上点为 \((x,y,x^2+y^2)\)。距离平方为 \(\displaystyle D=(x-2)^2+(y-2)^2+\left(x^2+y^2-\frac12\right)^2\)。
目标点 \((2,2,1/2)\) 与抛物面都关于平面 \(x=y\) 对称，最近点可取在该平面上，令
\(x=y=s\)。于是
\[
  D(s)=2(s-2)^2+\left(2s^2-\frac12\right)^2.
\]
求导：
\[
  D'(s)=4(s-2)+8s\left(2s^2-\frac12\right)
  =8(2s^3-1).
\]
令 \(D'(s)=0\)，得 \(s=2^{-1/3}\)。代回可得
\[
  D_{\min}^2=\frac{33}{4}-3\sqrt[3]{4}.
\]
因此最小距离为
\[
  \sqrt{\frac{33}{4}-3\sqrt[3]{4}}.
\]
\end{solution}

\section{2020级工科数学分析（二）期末考试 A 卷}

\par\medskip
\noindent\textbf{一、计算题（共 5 小题，每小题 8 分，共 40 分）}\par\nobreak\smallskip
\begin{exercise}
\par
设 $z=f(xy^2,x^2y)$，求 $\displaystyle \dfrac{\partial^2 z}{\partial y^2}$。
\end{exercise}

\begin{solution}
\par
令 $u=xy^2$，$v=x^2y$。记 $f_1,f_2$ 为 $f$ 对第一、第二个变量的偏导数，则 \(\displaystyle z_y=f_1u_y+f_2v_y=2xyf_1+x^2f_2\)。
  继续对 \(y\) 求导。因为 \(u_y=2xy,\ v_y=x^2\)，所以
\[
  \frac{\partial}{\partial y}f_1=f_{11}\,2xy+f_{12}\,x^2,\qquad
  \frac{\partial}{\partial y}f_2=f_{21}\,2xy+f_{22}\,x^2.
\]
于是
\[
  z_{yy}=2xf_1+2xy(2xyf_{11}+x^2f_{12})
  +x^2(2xyf_{21}+x^2f_{22}).
\]
由 \(f_{12}=f_{21}\)，得
\[
  \frac{\partial^2 z}{\partial y^2}
  =4x^2y^2f_{11}+4x^3yf_{12}+x^4f_{22}+2xf_1,
\]
其中各偏导数均在 \((xy^2,x^2y)\) 处取值。
\end{solution}
\begin{exercise}
\par
设
  \[
    f(x)=
    \begin{cases}
      -x, & |x|\le \dfrac{\pi}{2},\\
      x, & \dfrac{\pi}{2}<|x|\le \pi,
    \end{cases}
  \]
  写出 $f(x)$ 的以 $2\pi$ 为周期的 Fourier 级数的和函数 $S(x)$ 在 $[-\pi,\pi]$ 上的表达式。
\end{exercise}

\begin{solution}
\par
先看原函数在 \([-\pi,\pi]\) 内的连续点与跳跃点。除
\(\displaystyle x=\pm\frac{\pi}{2}\) 以及周期端点 \(\pm\pi\) 外，函数连续，傅里叶级数收敛到原函数值。 在
\(\displaystyle x=\frac{\pi}{2}\) 处，左右极限分别为
\(-\pi/2\) 与 \(\pi/2\)，平均值为 \(0\)；在
\(\displaystyle x=-\frac{\pi}{2}\) 处同理平均值也为 \(0\)。周期端点处左右极限为 \(\pi\) 与 \(-\pi\)，平均值为
\(0\)。故
  \[
    S(x)=
    \begin{cases}
      -x, & |x|<\dfrac{\pi}{2},\\
      x, & \dfrac{\pi}{2}<|x|<\pi,\\
      0, & x=\pm\dfrac{\pi}{2}\text{ 或 }x=\pm\pi.
    \end{cases}
  \]
\end{solution}
\begin{exercise}
\par
计算二次积分
  \(\displaystyle \int_0^2 \mathrm{d}x\int_{x/2}^1 e^{y^2}\,\mathrm{d}y\)。
\end{exercise}

\begin{solution}
\par
积分区域为 $0\le x\le 2$，$x/2\le y\le 1$。交换积分次序得 \(\displaystyle 0\le y\le 1,\ 0\le x\le 2y\)。
  因此 \(\displaystyle \int_0^2 \mathrm{d}x\int_{x/2}^1 e^{y^2}\,\mathrm{d}y=\int_0^1\mathrm{d}y\int_0^{2y}e^{y^2}\,\mathrm{d}x=\int_0^1 2y e^{y^2}\,\mathrm{d}y=e-1\)。
原积分若先对 \(y\) 积分没有初等原函数，所以换序是关键；换序后被积函数只含 \(y\)，先对 \(x\) 积分就产生因子 \(2y\)。
\end{solution}
\begin{exercise}
\par
计算曲线积分 \(\displaystyle I=\oint_L\frac{-y\,\mathrm{d}x+x\,\mathrm{d}y}{x^2+y^2}\)，其中 $\displaystyle L:\dfrac{x^2}{a^2}+\dfrac{y^2}{b^2}=1$，取逆时针方向。
\end{exercise}

\begin{solution}
\par
在极坐标中，极角微分满足
\[
  \mathrm{d}\theta=\frac{-y\,\mathrm{d}x+x\,\mathrm{d}y}{x^2+y^2}.
\]
题中椭圆 \(\dfrac{x^2}{a^2}+\dfrac{y^2}{b^2}=1\) 包围原点一次，并取逆时针方向。沿闭曲线走一周时，极角连续增加
\(2\pi\)。因此
\[
  I=\oint_L \mathrm{d}\theta=2\pi.
\]
这里不能直接用 Green 公式在椭圆内部计算旋度，因为微分形式在原点处有奇点，而原点又被椭圆包围。正确的判断是看曲线绕原点的圈数：逆时针绕一圈对应极角增加 \(2\pi\)。
\end{solution}
\begin{exercise}
\par
计算曲面积分 \(\displaystyle I=\iint_\Sigma (4xz+x)\,\mathrm{d}y\,\mathrm{d}z-2yz\,\mathrm{d}z\,\mathrm{d}x+(1-z^2)\,\mathrm{d}x\,\mathrm{d}y\)，其中 $\Sigma:z=x^2+y^2$，$0\le z\le 1$，取下侧。
\end{exercise}

\begin{solution}
\par
令 \(\displaystyle P=4xz+x,\ Q=-2yz,\ R=1-z^2\)。
  曲面 $\Sigma$ 与平面 $z=1$ 围成区域
  \(\displaystyle \Omega=\{(x,y,z):x^2+y^2\le z\le 1\}\)。$\Sigma$ 取下侧时正好是 $\Omega$ 的外侧。由高斯公式，\(\displaystyle \nabla\cdot(P,Q,R)=(4z+1)-2z-2z=1\)。
  上盖圆盘 \(z=1\) 的外法向向上，只有 \(R\,\mathrm{d}x\,\mathrm{d}y\) 项贡献；但此时 \(R=1-z^2=0\)，故上盖通量为
\(0\)。因此侧面通量就是闭曲面总通量：
\[
  I=\iiint_\Omega 1\,\mathrm{d}V.
\]
把体积积分投影到 \(Oxy\) 平面，得
\[
  I=\iint_{x^2+y^2\le1}\bigl(1-x^2-y^2\bigr)\,\mathrm{d}x\,\mathrm{d}y
  =2\pi\int_0^1(1-r^2)r\,\mathrm{d}r
  =\frac{\pi}{2}.
\]
\end{solution}

\par\medskip
\noindent\textbf{二、解答下列各题（共 5 小题，每小题 8 分，共 40 分）}\par\nobreak\smallskip
\begin{exercise}
\par
求曲线
  \[
    \Gamma:
    \begin{cases}
      x^2+y^2+z^2-3x=0,\\
      x+y-z=0
    \end{cases}
  \]
  在点 $(1,1,1)$ 处的单位切向量 $\vec t$，并求函数
  \(\displaystyle u(x,y,z)=\ln(x^2+y^2+z^2)\)
  在点 $(1,1,1)$ 处沿该单位切向量 $\vec t$ 的方向导数。
\end{exercise}

\begin{solution}
\par
设 \(\displaystyle F=x^2+y^2+z^2-3x,\ G=x+y-z\)。在 $(1,1,1)$ 处，\(\displaystyle \nabla F=(-1,2,2),\ \nabla G=(1,1,-1)\)。因此切向量可取 \(\displaystyle \nabla F\times\nabla G=(-4,1,-3)\)，
  单位切向量为
  \(\displaystyle \vec t=\pm\frac{(-4,1,-3)}{\sqrt{26}}\)。
  又 \(\displaystyle \nabla u(1,1,1)=\left(\frac23,\frac23,\frac23\right)\)。
  若取 $\vec t=(-4,1,-3)/\sqrt{26}$，方向导数为 \(\displaystyle D_{\vec t}u=\nabla u\cdot \vec t=-\frac4{\sqrt{26}}\)。
  若取相反方向，则方向导数为 $4/\sqrt{26}$。
\end{solution}
\begin{exercise}
\par
求曲面 $\Sigma:z=\sqrt{x^2+y^2}$ 夹在两曲面 $x^2+y^2=y$，$x^2+y^2=2y$ 之间的那部分面积。
\end{exercise}

\begin{solution}
\par
用柱坐标表示两圆柱为 $r=\sin\theta$ 与 $r=2\sin\theta$，其中 $0\le\theta\le\pi$。圆锥 $z=r$ 的面积元为 \(\displaystyle \mathrm{d}S=\sqrt{1+z_x^2+z_y^2}\,\mathrm{d}x\,\mathrm{d}y=\sqrt2\,r\,\mathrm{d}r\,\mathrm{d}\theta\)。故所求面积为 \(\displaystyle S=\frac{3\pi\sqrt2}{4}\)。
\[
  S=\sqrt2\int_0^\pi\int_{\sin\theta}^{2\sin\theta}r\,\mathrm{d}r\,\mathrm{d}\theta
  =\frac{\sqrt2}{2}\int_0^\pi 3\sin^2\theta\,\mathrm{d}\theta
  =\frac{3\pi\sqrt2}{4}.
\]
\end{solution}
\begin{exercise}
\par
求微分方程 $y''-6y'+9y=0$ 的通解，并写出 $y''-6y'+9y=(x+1)e^{3x}$ 的特解形式。
\end{exercise}

\begin{solution}
\par
齐次方程的特征方程为 \(\displaystyle \lambda^2-6\lambda+9=(\lambda-3)^2=0\)。因此通解为 \(\displaystyle y=(C_1+C_2x)e^{3x}\)。对方程 $y''-6y'+9y=(x+1)e^{3x}$，由于 $3$ 是二重特征根，特解可设为 \(\displaystyle y_p=x^2(Ax+B)e^{3x}\)。
这里右端是一次多项式乘 \(e^{3x}\)，通常特解为 \(e^{3x}(Ax+B)\)；但 \(3\) 是二重特征根，所以必须再乘 \(x^2\)。
\end{solution}
\begin{exercise}
\par
判断方程 $(x^2-y)\,\mathrm{d}x-x\,\mathrm{d}y=0$ 是否是全微分方程，如果是，求其通解。
\end{exercise}

\begin{solution}
\par
记 $P=x^2-y$，$Q=-x$。因为 \(\displaystyle P_y=-1,\ Q_x=-1\)，所以这是全微分方程。设势函数为 $\Phi$，则 \(\displaystyle \Phi_x=x^2-y\)，从而 \(\displaystyle \Phi=\frac{x^3}{3}-xy+\varphi(y)\)。
  由 $\Phi_y=-x+\varphi'(y)=Q=-x$，得 $\varphi'(y)=0$。通解为
  \(\displaystyle \frac{x^3}{3}-xy=C\)。
\end{solution}
\begin{exercise}
\par
求幂级数
  \(\displaystyle \sum_{n=0}^{\infty}(-1)^n\frac{2n^2+1}{(2n+1)!}x^{2n+1}\)
  的和函数，并指出其收敛域。
\end{exercise}

\begin{solution}
\par
设 \(\displaystyle S(x)=\sum_{n=0}^{\infty}(-1)^n\frac{2n^2+1}{(2n+1)!}x^{2n+1}\)。对 $\displaystyle \sin x=\sum_{n=0}^\infty(-1)^n\dfrac{x^{2n+1}}{(2n+1)!}$ 使用算子 $\displaystyle D=x\dfrac{\mathrm{d}}{\mathrm{d}x}$。因 $2n+1=m$ 时
  \(\displaystyle 2n^2+1=\frac{m^2-2m+3}{2}\)，所以 \(\displaystyle S(x)=\frac12(D^2-2D+3)\sin x\)。又 \(\displaystyle D\sin x=x\cos x,\ D^2\sin x=x\cos x-x^2\sin x\)。
  故
\[
  S(x)=\frac12\left[(x\cos x-x^2\sin x)-2x\cos x+3\sin x\right]
  =\frac{(3-x^2)\sin x-x\cos x}{2}.
\]
该幂级数由正弦级数逐项微分和线性组合得到，收敛域为
\((-\infty,+\infty)\)。
\end{solution}

\par\medskip
\noindent\textbf{三、证明下列各题（共 2 小题，每小题 6 分，共 12 分）}\par\nobreak\smallskip
\begin{exercise}
\par
设幂级数 $\displaystyle \sum_{n=0}^\infty a_nx^n$ 的收敛半径为 $R_1$，幂级数 $\displaystyle \sum_{n=0}^\infty b_nx^n$ 的收敛半径为 $R_2$，且 $R_1<R_2$。证明幂级数
  \(\displaystyle \sum_{n=0}^\infty (a_n+b_n)x^n\)
  的收敛半径为 $R_1$。
\end{exercise}

\begin{solution}
\par
当 $|x|<R_1$ 时，两个幂级数都收敛，所以 $\displaystyle \sum(a_n+b_n)x^n$ 收敛，故其收敛半径至少为 $R_1$。若其收敛半径大于 $R_1$，取 $x_0$ 使
  \(\displaystyle R_1<|x_0|<R_2\)。
  则 $\displaystyle \sum b_nx_0^n$ 收敛，而 $\displaystyle \sum a_nx_0^n$ 发散。若 $\displaystyle \sum(a_n+b_n)x_0^n$ 收敛，则 \(\displaystyle \sum a_nx_0^n=\sum(a_n+b_n)x_0^n-\sum b_nx_0^n\) 也应收敛，矛盾。因此 $\displaystyle \sum(a_n+b_n)x^n$ 的收敛半径为 $R_1$。
\end{solution}
\begin{exercise}
\par
证明函数项级数
  \(\displaystyle \sum_{n=0}^\infty \frac{\cos(nx)}{2^n}\)
  在 $(-\infty,+\infty)$ 上一致收敛。
\end{exercise}

\begin{solution}
\par
对一切实数 $x$，\(\displaystyle \left|\frac{\cos(nx)}{2^n}\right|\le \frac1{2^n}\)。而 \(\displaystyle \sum_{n=0}^{\infty}\frac1{2^n}\) 收敛，所以由 Weierstrass 判别法，原函数项级数在 $(-\infty,+\infty)$ 上一致收敛。
这里控制项 \(\frac1{2^n}\) 与 \(x\) 无关，正好满足一致收敛判别法的要求；若控制项还依赖 \(x\)，就不能直接推出一致收敛。
\end{solution}

\par\medskip
\noindent\textbf{四、应用题（本题 8 分）}\par\nobreak\smallskip
\begin{exercise}
\par
在第一象限内作椭球面 \(\displaystyle \frac{x^2}{a^2}+\frac{y^2}{b^2}+\frac{z^2}{c^2}=1\) 的切平面，使得切平面与三个坐标面所围成的四面体体积最小，求切点坐标。
\end{exercise}

\begin{solution}
\par
设切点为 $(x_0,y_0,z_0)$，其中 $x_0,y_0,z_0>0$。椭球面在该点的切平面为 \(\displaystyle \frac{x_0x}{a^2}+\frac{y_0y}{b^2}+\frac{z_0z}{c^2}=1\)。
它在三个坐标轴上的截距分别为 \(\displaystyle \frac{a^2}{x_0},\ \frac{b^2}{y_0},\ \frac{c^2}{z_0}\)。
故四面体体积为 \(\displaystyle V_T=\frac16\frac{a^2}{x_0}\frac{b^2}{y_0}\frac{c^2}{z_0}=\frac{a^2b^2c^2}{6x_0y_0z_0}\)。
因此要使 $V_T$ 最小，只需在约束 \(\displaystyle \frac{x_0^2}{a^2}+\frac{y_0^2}{b^2}+\frac{z_0^2}{c^2}=1\) 下使 $x_0y_0z_0$ 最大。令 \(\displaystyle X=\frac{x_0}{a},\ Y=\frac{y_0}{b},\ Z=\frac{z_0}{c}\)，
则 $X^2+Y^2+Z^2=1$。由对称性或拉格朗日乘数法，最大值在
\(\displaystyle X=Y=Z=\frac1{\sqrt3}\)
处取得。所以所求切点为 \(\displaystyle \left(\frac{a}{\sqrt3},\frac{b}{\sqrt3},\frac{c}{\sqrt3}\right)\)。
\end{solution}

\section{2020级工科数学分析（二）期末考试 B 卷}

\par\medskip
\noindent\textbf{一、计算题（共 5 小题，每小题 8 分，共 40 分）}\par\nobreak\smallskip
\begin{exercise}
\par
设 $z=f(x+y,xy)$，求 $\displaystyle \dfrac{\partial^2z}{\partial x\partial y}$。
\end{exercise}

\begin{solution}
\par
令 $u=x+y$，$v=xy$。则 \(\displaystyle z_x=f_1+ yf_2\)。
  对 \(y\) 求导，注意 \(u_y=1,\ v_y=x\)。先有
\[
  \frac{\partial f_1}{\partial y}=f_{11}+xf_{12},\qquad
  \frac{\partial f_2}{\partial y}=f_{21}+xf_{22}.
\]
因此
\[
  z_{xy}=f_{11}+xf_{12}+f_2+y(f_{21}+xf_{22}).
\]
因 \(f_{12}=f_{21}\)，可写为
\[
  z_{xy}=f_{11}+(x+y)f_{12}+xyf_{22}+f_2,
\]
其中各偏导数均在 \((x+y,xy)\) 处取值。
\end{solution}
\begin{exercise}
\par
设
  \[
    f(x)=
    \begin{cases}
      x+1, & -\pi\le x<0,\\
      x^2, & 0\le x<\pi,
    \end{cases}
  \]
  写出 $f(x)$ 的以 $2\pi$ 为周期的 Fourier 级数的和函数 $S(x)$ 在 $[-\pi,\pi]$ 上的表达式。
\end{exercise}

\begin{solution}
\par
傅里叶级数在连续点处收敛到函数值，在跳跃间断点处收敛到左右极限的平均值。故
  \[
    S(x)=
    \begin{cases}
      x+1, & -\pi<x<0,\\
      x^2, & 0<x<\pi,\\
      \dfrac12, & x=0,\\
      \dfrac{\pi^2-\pi+1}{2}, & x=\pm\pi.
    \end{cases}
  \]
其中 \(x=0\) 是内部跳跃点；\(x=\pm\pi\) 是周期端点，要把 \(\pi-0\) 的值与由周期性折回的 \(-\pi+0\) 的值取平均，所以端点值为 \((\pi^2+(-\pi+1))/2\)。
\end{solution}
\begin{exercise}
\par
计算二次积分
  \(\displaystyle \int_0^1\mathrm{d}y\int_{y^2}^{y}y\sin(x^2)\,\mathrm{d}x\)。
\end{exercise}

\begin{solution}
\par
积分区域为 $0\le y\le 1$，$y^2\le x\le y$。交换积分次序得
\[
  0\le x\le 1,\qquad x\le y\le \sqrt{x}.
\]
因此
\begin{align*}
  \int_0^1\mathrm{d}y\int_{y^2}^{y}y\sin(x^2)\,\mathrm{d}x
  &=\int_0^1\sin(x^2)\,\mathrm{d}x\int_x^{\sqrt{x}}y\,\mathrm{d}y\\
  &=\frac12\int_0^1(x-x^2)\sin(x^2)\,\mathrm{d}x .
\end{align*}
进一步化简，
\[
  \frac12\int_0^1x\sin(x^2)\,\mathrm{d}x=\frac{1-\cos1}{4},
  \qquad
  \int_0^1x^2\sin(x^2)\,\mathrm{d}x=-\frac{\cos1}{2}+\frac12\int_0^1\cos(x^2)\,\mathrm{d}x .
\]
所以 \(\displaystyle I=\frac14\left(1-\int_0^1\cos(x^2)\,\mathrm{d}x\right)\)。
其中 $\displaystyle \int_0^1\cos(x^2)\,\mathrm{d}x$ 可用 Fresnel 积分表示。
\end{solution}
\begin{exercise}
\par
计算曲线积分 \(\displaystyle I=\oint_L\frac{-y\,\mathrm{d}x+x\,\mathrm{d}y}{4x^2+y^2}\)，其中 $L:(x-1)^2+y^2=4$，取逆时针方向。
\end{exercise}

\begin{solution}
\par
作变量变换 \(u=2x,\ v=y\)，则 \(\mathrm{d}x=\mathrm{d}u/2,\ \mathrm{d}y=\mathrm{d}v\)，并且
\[
  \frac{-y\,\mathrm{d}x+x\,\mathrm{d}y}{4x^2+y^2}
  =\frac{-v\,\mathrm{d}u/2+(u/2)\,\mathrm{d}v}{u^2+v^2}
  =\frac12\frac{-v\,\mathrm{d}u+u\,\mathrm{d}v}{u^2+v^2}.
\]
曲线 \(L\) 在 \((u,v)\) 平面中的像仍是一条逆时针围绕原点一周的椭圆。被积式
\[
  \frac{-v\,\mathrm{d}u+u\,\mathrm{d}v}{u^2+v^2}
\]
是极角微分 \(\mathrm{d}\theta\)，沿该闭曲线积分为 \(2\pi\)。因此
\[
  I=\frac12\cdot2\pi=\pi.
\]
\end{solution}
\begin{exercise}
\par
计算曲面积分 \(\displaystyle I=\iint_\Sigma \frac{x^3\,\mathrm{d}y\,\mathrm{d}z+y^3\,\mathrm{d}z\,\mathrm{d}x+z^3\,\mathrm{d}x\,\mathrm{d}y}{\sqrt{x^2+y^2+z^2}}\)，其中 $\Sigma:z=\sqrt{a^2-x^2-y^2}$，取下侧。
\end{exercise}

\begin{solution}
\par
在上半球面上 \(\sqrt{x^2+y^2+z^2}=a\)。球面外侧单位法向量为
\((x,y,z)/a\)，而题设取下侧，即取相反方向。因此第二类曲面积分可写为
\[
  I=-\iint_\Sigma
  \frac{x^3x+y^3y+z^3z}{a^2}\,\mathrm{d}S
  =-\frac1{a^2}\iint_\Sigma(x^4+y^4+z^4)\,\mathrm{d}S.
\]
由球面对称性，上半球面上
\[
  \iint_\Sigma x^4\,\mathrm{d}S=\iint_\Sigma y^4\,\mathrm{d}S=\iint_\Sigma z^4\,\mathrm{d}S
  =\frac{2\pi a^6}{5}.
\]
所以
\[
  I=-\frac1{a^2}\cdot3\cdot\frac{2\pi a^6}{5}
  =-\frac{6\pi a^4}{5}.
\]
\end{solution}

\par\medskip
\noindent\textbf{二、解答下列各题（共 5 小题，每小题 8 分，共 40 分）}\par\nobreak\smallskip
\begin{exercise}
\par
求曲线
  \[
    \Gamma:
    \begin{cases}
      x^2+y^2+z^2=6,\\
      x+y+z=0
    \end{cases}
  \]
  在点 $(1,-2,1)$ 处的单位切向量 $\vec t$，并求函数
  \(\displaystyle u(x,y,z)=x^2+y^2+z^2\)
  在点 $(1,-2,1)$ 处沿该单位切向量 $\vec t$ 的方向导数。
\end{exercise}

\begin{solution}
\par
令 \(\displaystyle F=x^2+y^2+z^2,\ G=x+y+z\)。在 $(1,-2,1)$ 处，\(\displaystyle \nabla F=(2,-4,2),\ \nabla G=(1,1,1)\)。切向量可取 \(\displaystyle \nabla F\times\nabla G=(-6,0,6)\)，
  故单位切向量为
  \(\displaystyle \vec t=\pm\frac{(-1,0,1)}{\sqrt2}\)。
  又 $\nabla u(1,-2,1)=(2,-4,2)$，所以
  \(\displaystyle D_{\vec t}u=\nabla u\cdot\vec t=0\)。
\end{solution}
\begin{exercise}
\par
求密度 $\rho=1$ 的均匀上半球面 $\Sigma:z=\sqrt{a^2-x^2-y^2}$ 对 $z$ 轴的转动惯量。
\end{exercise}

\begin{solution}
\par
对 $z$ 轴的转动惯量为
  \(\displaystyle J_z=\iint_\Sigma (x^2+y^2)\rho\,\mathrm{d}S\)。
  用球坐标参数化上半球面，$x^2+y^2=a^2\sin^2\varphi$，$\mathrm{d}S=a^2\sin\varphi\,\mathrm{d}\varphi\,\mathrm{d}\theta$，其中 $0\le\varphi\le\pi/2$，$0\le\theta\le2\pi$。于是
  \(\displaystyle J_z=\int_0^{2\pi}\int_0^{\pi/2}a^2\sin^2\varphi\cdot a^2\sin\varphi\,\mathrm{d}\varphi\,\mathrm{d}\theta=2\pi a^4\cdot\frac23=\frac{4\pi a^4}{3}\)。
\end{solution}
\begin{exercise}
\par
求微分方程 $y'+y\tan x=\cos x$ 的通解。
\end{exercise}

\begin{solution}
\par
这是一阶线性方程
\[
  y'+(\tan x)y=\cos x.
\]
积分因子为
\[
  \mu(x)=e^{\int\tan x\,\mathrm{d}x}=e^{-\ln|\cos x|}=\sec x
\]
（在不跨越 \(\cos x=0\) 的区间上讨论）。方程两边乘以 \(\sec x\)，得
\[
  (y\sec x)'=1.
\]
积分得 \(y\sec x=x+C\)，即
\[
  y=(x+C)\cos x.
\]
这个通解在任一不跨越 \(\cos x=0\) 的区间上成立；积分因子中把 \(e^{-\ln|\cos x|}\) 写成 \(\sec x\) 时，本质上已经固定了这样的局部区间。
\end{solution}
\begin{exercise}
\par
求微分方程 $y''-y=0$ 的通解，并写出 $y''-y=4xe^{-x}$ 的特解形式。
\end{exercise}

\begin{solution}
\par
齐次方程 $y''-y=0$ 的特征根为 $\lambda=\pm1$，故通解为
  \(\displaystyle y=C_1e^x+C_2e^{-x}\)。
  对非齐次方程 \(y''-y=4xe^{-x}\)，右端为“一次多项式乘 \(e^{-x}\)”。通常应设
\((Ax+B)e^{-x}\)，但 \(-1\) 是齐次方程的单特征根，发生共振，所以需再乘一个
\(x\)。因此特解可设为
  \(\displaystyle y_p=x(Ax+B)e^{-x}\)。
\end{solution}
\begin{exercise}
\par
求幂级数
  \(\displaystyle \sum_{n=1}^{\infty}\frac{n^2+1}{2^n n!}x^n\)
  的和函数，并指出其收敛域。
\end{exercise}

\begin{solution}
\par
令 $t=x/2$，则 \(\displaystyle \sum_{n=1}^{\infty}\frac{n^2+1}{2^n n!}x^n=\sum_{n=1}^{\infty}\frac{(n^2+1)t^n}{n!}\)。利用 \(\displaystyle \sum_{n=1}^{\infty}\frac{n^2t^n}{n!}=(t^2+t)e^t,\quad \sum_{n=1}^{\infty}\frac{t^n}{n!}=e^t-1\)，得和函数 \(\displaystyle S(x)=\left(\frac{x^2}{4}+\frac{x}{2}+1\right)e^{x/2}-1\)。该幂级数收敛域为 $(-\infty,+\infty)$。
\end{solution}

\par\medskip
\noindent\textbf{三、证明下列各题（共 2 小题，每小题 6 分，共 12 分）}\par\nobreak\smallskip
\begin{exercise}
\par
设正项级数 $\displaystyle \sum_{n=0}^{\infty}a_n$ 收敛，证明正项级数
  \(\displaystyle \sum_{n=1}^{\infty}\frac{\sqrt{a_n}}{n}\)
  也收敛。【原卷第二个求和下限印作 $n=0$，但含分母 $n$，本稿按 $n=1$ 校正。】
\end{exercise}

\begin{solution}
\par
对任意 $N$，由 Cauchy--Schwarz 不等式，\(\displaystyle \sum_{n=1}^{N}\frac{\sqrt{a_n}}{n}\le\left(\sum_{n=1}^{N}a_n\right)^{1/2}\left(\sum_{n=1}^{N}\frac1{n^2}\right)^{1/2}\)。
  由于 $\displaystyle \sum a_n$ 与 $\displaystyle \sum 1/n^2$ 均收敛，右侧对 $N$ 有一致上界。又原级数为正项级数，其部分和单调递增且有上界，故 \(\displaystyle \sum_{n=1}^{\infty}\frac{\sqrt{a_n}}{n}\) 收敛。
\end{solution}
\begin{exercise}
\par
设正项级数 $\displaystyle \sum_{n=0}^{\infty}a_n$ 收敛，证明函数
  \(\displaystyle f(x)=\sum_{n=0}^{\infty}a_nx^n\)
  在 $(-1,1)$ 上连续。
\end{exercise}

\begin{solution}
\par
对 $x\in(-1,1)$，每一项 $a_nx^n$ 都连续。并且在 $[-1,1]$ 上 \(\displaystyle |a_nx^n|\le a_n\)。
  因为 $\displaystyle \sum a_n$ 收敛，由 Weierstrass 判别法，$\displaystyle \sum a_nx^n$ 在 $[-1,1]$ 上一致收敛，从而其和函数在 $[-1,1]$ 上连续，特别地在 $(-1,1)$ 上连续。
\end{solution}

\par\medskip
\noindent\textbf{四、应用题（本题 8 分）}\par\nobreak\smallskip
\begin{exercise}
\par
要造一容积为 $V$ 的长方体无盖铝盒，请问怎样设计尺寸，使它的表面积最小？
\end{exercise}

\begin{solution}
\par
设无盖铝盒底面边长分别为 $x,y$，高为 $h$。约束条件为 \(\displaystyle xyh=V\)。
表面积为 \(\displaystyle S=xy+2xh+2yh=xy+\frac{2V}{y}+\frac{2V}{x}\)。
求偏导并令其为零：\(\displaystyle S_x=y-\frac{2V}{x^2}=0,\qquad S_y=x-\frac{2V}{y^2}=0\)。由此得 $x=y$。设 $x=y=t$，则 \(\displaystyle t^3=2V,\ h=\frac{V}{t^2}=\frac{t}{2}\)。所以当底面为正方形，边长 \(\displaystyle t=\sqrt[3]{2V}\)，高为 \(\displaystyle h=\frac12\sqrt[3]{2V}\) 时，表面积最小。
\end{solution}

\section{2021级工科数学分析（二）期末考试 A 卷}

\par\medskip
\noindent\textbf{一、叙述定义（共 5 小题，每小题 2 分，共 10 分）}\par\nobreak\smallskip
\begin{exercise}
\par
二元函数的偏导数。
\end{exercise}

\begin{solution}
\par
若极限 \(\displaystyle \lim_{\Delta x\to 0}\frac{f(x_0+\Delta x,y_0)-f(x_0,y_0)}{\Delta x}\) 存在，则称其为 $f$ 在 $(x_0,y_0)$ 处关于 $x$ 的偏导数，记为 $f_x(x_0,y_0)$；关于 $y$ 的偏导数类似。
定义的关键是只让一个变量变化、另一个变量固定。求 \(x\) 偏导时固定 \(y=y_0\)，求 \(y\) 偏导时固定 \(x=x_0\)。
\end{solution}
\begin{exercise}
\par
高斯公式。
\end{exercise}

\begin{solution}
\par
若闭区域 $V$ 的边界为分片光滑闭曲面 $\Sigma$，取外侧，且 $P,Q,R$ 在 $V$ 上有连续一阶偏导数，则
  \[
    \iiint_V\left(\frac{\partial P}{\partial x}+\frac{\partial Q}{\partial y}+\frac{\partial R}{\partial z}\right)\,\mathrm{d}V
    =
    \iint_\Sigma P\,\mathrm{d}y\,\mathrm{d}z+Q\,\mathrm{d}z\,\mathrm{d}x+R\,\mathrm{d}x\,\mathrm{d}y .
  \]
\end{solution}
\begin{exercise}
\par
条件收敛级数。
\end{exercise}

\begin{solution}
\par
若 $\displaystyle \sum a_n$ 收敛而 $\displaystyle \sum |a_n|$ 发散，则称 $\displaystyle \sum a_n$ 条件收敛。
也就是说，原级数靠正负项相互抵消后能够收敛，但取绝对值后相当于去掉抵消，级数就发散。典型例子是交错调和级数
\(\displaystyle \sum_{n=1}^{\infty}(-1)^{n-1}/n\)。
因此判断条件收敛时要分两步：先证明原级数收敛，再证明绝对值级数发散；只验证其中一步是不够的。
\end{solution}
\begin{exercise}
\par
函数列的一致收敛。
\end{exercise}

\begin{solution}
\par
若函数列 $\{f_n(x)\}$ 在 $I$ 上逐点收敛于 $f(x)$，且对任意 $\varepsilon>0$，存在 $N$，使当 $n>N$ 时，对一切 $x\in I$ 都有
  \(\displaystyle |f_n(x)-f(x)|<\varepsilon\)，
  则称 $\{f_n\}$ 在 $I$ 上一致收敛于 $f$。
这里的 \(N\) 只能依赖于 \(\varepsilon\)，不能随着 \(x\) 改变；这正是一致收敛比逐点收敛更强的地方。
\end{solution}
\begin{exercise}
\par
二阶线性齐次常微分方程 $L[y]=0$ 的基础解系。
\end{exercise}

\begin{solution}
\par
若二阶线性齐次方程 $L[y]=0$ 的两个解 $\varphi_1,\varphi_2$ 线性无关，则 $\{\varphi_1,\varphi_2\}$ 称为该方程的一个基础解系。
此时该方程的通解可写成
\[
  y=C_1\varphi_1+C_2\varphi_2,
\]
其中 \(C_1,C_2\) 为任意常数。二阶方程需要两个线性无关解，正是因为通解中应含有两个独立任意常数。
若需要检验两个解是否线性无关，可看它们的 Wronski 行列式是否不恒为零；二阶基础解系必须恰好提供两个独立方向。
\end{solution}

\par\medskip
\noindent\textbf{二、计算题（共 4 小题，每小题 10 分，共 40 分）}\par\nobreak\smallskip
\begin{exercise}
\par
设 $u=z^2+2z$，且 $z=z(x,y)$ 由方程 $xe^x-ye^y=ze^z$（$z\ne 1$）所确定，求 $\mathrm{d}u$ 和 $\displaystyle \dfrac{\partial^2u}{\partial x^2}$。
\end{exercise}

\begin{solution}
\par
令 $F(x,y,z)=xe^x-ye^y-ze^z$。由隐函数求导，\(\displaystyle z_x=\frac{x+1}{z+1}e^{x-z},\qquad z_y=-\frac{y+1}{z+1}e^{y-z}\)。题面括号中写作 $z\ne 1$，而隐函数求导实际要求 $F_z=-(z+1)e^z\ne0$，即 $z\ne-1$；以下按该条件计算。
  因为 \(u=z^2+2z\)，所以
\[
  \mathrm{d}u=(2z+2)\,\mathrm{d}z=2(z+1)(z_x\,\mathrm{d}x+z_y\,\mathrm{d}y).
\]
代入 \(z_x,z_y\) 得
\[
  \mathrm{d}u=2(x+1)e^{x-z}\,\mathrm{d}x-2(y+1)e^{y-z}\,\mathrm{d}y.
\]
再由
\[
  u_x=2(x+1)e^{x-z}
\]
对 \(x\) 求导：
\[
  u_{xx}=2e^{x-z}+2(x+1)e^{x-z}(1-z_x)
  =2e^{x-z}\left[(x+2)-\frac{(x+1)^2e^{x-z}}{z+1}\right].
\]
即
\[
  \frac{\partial^2u}{\partial x^2}
  =\frac{2e^{x-z}\bigl[(z+1)(x+2)-(x+1)^2e^{x-z}\bigr]}{z+1}.
\]
\end{solution}
\begin{exercise}
\par
计算
  \[
    \int_1^2 \mathrm{d}x\int_{\sqrt{x}}^x \sin\frac{\pi x}{2y}\,\mathrm{d}y
    +\int_2^4 \mathrm{d}x\int_{\sqrt{x}}^2 \sin\frac{\pi x}{2y}\,\mathrm{d}y .
  \]
\end{exercise}

\begin{solution}
\par
交换积分次序，积分区域可写为 \(1\le y\le 2,\ y\le x\le y^2\)，故
  \[
  \int_1^2 \mathrm{d}y\int_y^{y^2}\sin\frac{\pi x}{2y}\,\mathrm{d}x
  =\int_1^2\frac{2y}{\pi}
  \left[\!-\,\cos\frac{\pi x}{2y}\right]_{x=y}^{x=y^2}\,\mathrm{d}y
  =-\frac{2}{\pi}\int_1^2 y\cos\frac{\pi y}{2}\,\mathrm{d}y
  =\frac{4(\pi+2)}{\pi^3}.
  \]
\end{solution}
\begin{exercise}
\par
计算 $\displaystyle \iint_\Sigma x\,\mathrm{d}y\,\mathrm{d}z$，其中 $\Sigma$ 为 $x^2+y^2=1$，$z=x+2$ 和 $z=0$ 所围立体的表面，取外侧。
\end{exercise}

\begin{solution}
\par
由高斯公式，取 \(P=x,Q=R=0\)，则 \(\displaystyle \iint_\Sigma x\,\mathrm{d}y\,\mathrm{d}z=\iiint_\Omega 1\,\mathrm{d}V\)。又
  \[
  \iiint_\Omega 1\,\mathrm{d}V
  =\iint_{x^2+y^2\le 1}(x+2)\,\mathrm{d}x\,\mathrm{d}y=2\pi .
  \]
  这里积分区域在 \(xy\) 平面上的投影是单位圆盘，顶面高度为 \(x+2\)，所以体积为 \(2\pi\)。因此该闭曲面的外向通量为 \(2\pi\)。
\end{solution}
\begin{exercise}
\par
计算
  \(\displaystyle \int_L e^{\sqrt{x^2+y^2}}\,\mathrm{d}s\)，
  其中 $L$ 为圆周 $x^2+y^2=a^2$（$a>0$）、直线 $y=x$ 与 $x$ 轴在第一象限内所围成的扇形的整个边界。
\end{exercise}

\begin{solution}
\par
边界由三段组成：$L_1:y=0,0\le x\le a$；$L_2:y=x,0\le x\le a/\sqrt2$；$L_3:x=a\cos\theta,y=a\sin\theta,0\le\theta\le\pi/4$。于是 \(L_1,L_2\) 上的积分均为 \(e^a-1\)，而 \(\displaystyle \int_{L_3}e^{\sqrt{x^2+y^2}}\,\mathrm{d}s=\frac{\pi a}{4}e^a\)。因而 \(\displaystyle \int_L e^{\sqrt{x^2+y^2}}\,\mathrm{d}s=2(e^a-1)+\frac{\pi a}{4}e^a\)。
\[
  \int_{L_1}e^{\sqrt{x^2+y^2}}\,\mathrm{d}s=\int_0^ae^x\,\mathrm{d}x=e^a-1.
\]
在 \(L_2\) 上令 \(x=t,\ y=t\)，则 \(0\le t\le a/\sqrt2\)，
\(\sqrt{x^2+y^2}=\sqrt2t,\ \mathrm{d}s=\sqrt2\,\mathrm{d}t\)，故
\[
  \int_{L_2}e^{\sqrt{x^2+y^2}}\,\mathrm{d}s
  =\int_0^{a/\sqrt2}e^{\sqrt2t}\sqrt2\,\mathrm{d}t=e^a-1.
\]
圆弧 \(L_3\) 上半径恒为 \(a\)，弧长为 \(a\cdot\pi/4\)，所以圆弧贡献为
\(\displaystyle \frac{\pi a}{4}e^a\)。三段相加即得上述结果。
\end{solution}

\par\medskip
\noindent\textbf{三、解答题（共 3 小题，每小题 10 分，共 30 分）}\par\nobreak\smallskip
\begin{exercise}
\par
设曲线积分
  \(\displaystyle \int_L xy^2\,\mathrm{d}x+y\varphi(x)\,\mathrm{d}y\)
  与路径无关，其中 $\varphi(x)$ 具有连续的导数，且 $\varphi(0)=0$，计算 \(\displaystyle \int_{(0,0)}^{(1,1)} xy^2\,\mathrm{d}x+y\varphi(x)\,\mathrm{d}y\)。
\end{exercise}

\begin{solution}
\par
设 $P=xy^2,Q=y\varphi(x)$。路径无关给出 \(\displaystyle P_y=Q_x,\ 2xy=y\varphi'(x)\)，
  从而 $\varphi'(x)=2x$，又 $\varphi(0)=0$，故 $\varphi(x)=x^2$。沿折线 $(0,0)\to(1,0)\to(1,1)$ 计算，
  第一段 \(y=0\) 上 \(\mathrm{d}x\) 项和 \(\mathrm{d}y\) 项都为 \(0\)；第二段 \(x=1,\ \mathrm{d}x=0\)，被积式化为
\(y\,\mathrm{d}y\)。因此
\[
  \int_{(0,0)}^{(1,1)} xy^2\,\mathrm{d}x+y\varphi(x)\,\mathrm{d}y
  =\int_0^1 y\,\mathrm{d}y=\frac12.
\]
\end{solution}
\begin{exercise}
\par
求级数 $\displaystyle \sum_{n=1}^\infty \frac{n(n+1)}{2^n}$ 的和。
\end{exercise}

\begin{solution}
\par
对 $|x|<1$，设 \(\displaystyle S(x)=\sum_{n=1}^\infty n(n+1)x^n\)。由 $\displaystyle \sum_{n=1}^\infty nx^{n-1}=1/(1-x)^2$ 可得 \(\displaystyle S(x)=\frac{2x}{(1-x)^3}\)。
具体地，
\[
  \sum_{n=1}^{\infty}nx^n=\frac{x}{(1-x)^2},
\]
再对两边求导并乘以 \(x\)，得
\[
  \sum_{n=1}^{\infty}n^2x^n
  =x\left(\frac{x}{(1-x)^2}\right)'
  =\frac{x(1+x)}{(1-x)^3}.
\]
于是
\[
  \sum_{n=1}^{\infty}n(n+1)x^n
  =\sum n^2x^n+\sum nx^n
  =\frac{2x}{(1-x)^3}.
\]
取 \(x=1/2\)，得所求和为 \(8\)。
\end{solution}
\begin{exercise}
\par
求方程 $y''+4y'+4y=\cos 2x$ 的通解。
\end{exercise}

\begin{solution}
\par
齐次方程特征根为 $r=-2$（二重根），故 \(\displaystyle y_h=(C_1+C_2x)e^{-2x}\)。
  设特解 \(y_p=A\sin2x+B\cos2x\)。代入 \(L[y]=y''+4y'+4y\) 后，
\[
  L[A\sin2x+B\cos2x]=8A\cos2x-8B\sin2x.
\]
与右端 \(\cos2x\) 比较，得 \(8A=1,\ -8B=0\)，即 \(A=1/8,\ B=0\)。所以
  \(\displaystyle y=(C_1+C_2x)e^{-2x}+\frac18\sin 2x\)。
\end{solution}

\par\medskip
\noindent\textbf{四、应用题（10 分）}\par\nobreak\smallskip
\begin{exercise}
\par
造一容积为 $V$ 的长方形无盖铝盆，怎样设计尺寸，使它的表面积最小？
\end{exercise}

\begin{solution}
\par
设长方形无盖铝盆底边为 $x,y$，高为 $h$。约束为 $xyh=V$，表面积
\(\displaystyle S=xy+2xh+2yh=xy+\frac{2V}{y}+\frac{2V}{x}\)。
由对称性和驻点方程可得 $x=y$，令 $x=y=t$，则 \(\displaystyle S(t)=t^2+\frac{4V}{t},\ S'(t)=2t-\frac{4V}{t^2}=0\)。故 $t^3=2V$，$h=V/t^2=t/2$。即底面为边长 $\sqrt[3]{2V}$ 的正方形，高为 $\sqrt[3]{2V}/2$ 时表面积最小。
\end{solution}

\par\medskip
\noindent\textbf{五、证明题（10 分）}\par\nobreak\smallskip
\begin{exercise}
\par
证明级数 $\displaystyle \sum_{n=1}^\infty x^n(1-x)^2$ 在 $[0,1]$ 上一致收敛。
\end{exercise}

\begin{solution}
\par
记 $u_n(x)=x^n(1-x)^2$。在 $[0,1]$ 上，\(\displaystyle u_n'(x)=nx^{n-1}(1-x)^2-2x^n(1-x)\)，其最大值在 $x=n/(n+2)$ 处取得，且
  \(\displaystyle 0\le u_n(x)\le \left(\frac{n}{n+2}\right)^n\left(\frac{2}{n+2}\right)^2\le \frac{4}{n^2}\)。
由于 $\displaystyle \sum 4/n^2$ 收敛，由 Weierstrass 判别法知 \(\displaystyle \sum_{n=1}^\infty x^n(1-x)^2\) 在 $[0,1]$ 上一致收敛。
\end{solution}

\section{2021级工科数学分析（二）期末考试 B 卷}

\par\medskip
\noindent\textbf{一、叙述定义（共 5 小题，每小题 2 分，共 10 分）}\par\nobreak\smallskip
\begin{exercise}
\par
二元函数的方向导数。
\end{exercise}

\begin{solution}
\par
若极限 \(\displaystyle \lim_{\rho\to0^+}\frac{f(x_0+\rho\cos\alpha,y_0+\rho\sin\alpha)-f(x_0,y_0)}{\rho}\)
  存在，则称其为 $f$ 在 $(x_0,y_0)$ 处沿方向 $\vec l=(\cos\alpha,\sin\alpha)$ 的方向导数。
其中方向向量 \(\vec l\) 应为单位向量；若 \(f\) 在该点可微，则方向导数还可写成 \(\nabla f(x_0,y_0)\cdot\vec l\)。
\end{solution}
\begin{exercise}
\par
斯托克斯公式。
\end{exercise}

\begin{solution}
\par
若 $\Sigma$ 是分片光滑有向曲面，边界为正向闭曲线 $L$，且 $P,Q,R$ 有连续一阶偏导数，则
  \[
    \oint_L P\,\mathrm{d}x+Q\,\mathrm{d}y+R\,\mathrm{d}z
    =
    \iint_\Sigma
    \begin{vmatrix}
      \mathrm{d}y\,\mathrm{d}z & \mathrm{d}z\,\mathrm{d}x & \mathrm{d}x\,\mathrm{d}y\\
      \partial_x & \partial_y & \partial_z\\
      P & Q & R
    \end{vmatrix}.
  \]
其中曲线方向与曲面法向必须满足右手法则：拇指指向曲面正法向时，四指弯曲方向就是边界曲线的正向。
\end{solution}
\begin{exercise}
\par
绝对收敛级数。
\end{exercise}

\begin{solution}
\par
若级数 $\displaystyle \sum a_n$ 满足 $\displaystyle \sum |a_n|$ 收敛，则称 $\displaystyle \sum a_n$ 绝对收敛。
绝对收敛一定推出原级数收敛，因为
\[
  0\le a_n^+\le |a_n|,\qquad 0\le a_n^-\le |a_n|
\]
可分别控制正项部分和负项部分。考试中判断一般级数时，通常先查绝对收敛；若绝对收敛失败，再考虑条件收敛。
\end{solution}
\begin{exercise}
\par
函数项级数的一致收敛。
\end{exercise}

\begin{solution}
\par
若函数项级数 $\displaystyle \sum u_n(x)$ 的部分和 $S_n(x)$ 在区间 $I$ 上一致收敛于 $S(x)$，即对任意 $\varepsilon>0$，存在 $N$，使 $n>N$ 时对一切 $x\in I$ 有
  \(\displaystyle |S_n(x)-S(x)|<\varepsilon\)，
  则称该级数在 $I$ 上一致收敛。
定义中的关键是同一个 \(N\) 要同时适用于所有 \(x\in I\)。若 \(N\) 可以随 \(x\) 改变，那只是逐点收敛，不是一致收敛。
\end{solution}
\begin{exercise}
\par
$n$ 阶线性齐次常微分方程 $L[y]=0$ 的基础解系。
\end{exercise}

\begin{solution}
\par
$n$ 阶线性齐次常微分方程 $L[y]=0$ 的 $n$ 个线性无关解 $y_1,\ldots,y_n$ 构成的解组称为一个基础解系。
有了基础解系后，齐次方程的通解为
\[
  y=C_1y_1+C_2y_2+\cdots+C_ny_n.
\]
因此“基础解系”的核心要求是线性无关；若解的个数不足或线性相关，就不能表示全部通解。
实际判定时常用 Wronski 行列式：若在某点 \(W(y_1,\ldots,y_n)\ne0\)，则这些解线性无关，从而可以作为基础解系。
\end{solution}

\par\medskip
\noindent\textbf{二、计算题（共 4 小题，每小题 10 分，共 40 分）}\par\nobreak\smallskip
\begin{exercise}
\par
设 $z=\ln(u^2+v^2)$，而 $u=e^{x^2+y^2}$，$v=x^2+y$，求 $\displaystyle \dfrac{\partial z}{\partial x}$ 和 $\displaystyle \dfrac{\partial^2z}{\partial x\partial y}$。
\end{exercise}

\begin{solution}
\par
令 \(\displaystyle u=e^{x^2+y^2},\ v=x^2+y,\ A=u^2+v^2\)。则
  \(\displaystyle z_x=\frac{4x(u^2+v)}{A}\)。继续求导得
  \[
  z_{xy}=\frac{4x\bigl[(4yu^2+1)A-(u^2+v)(4yu^2+2v)\bigr]}{A^2}.
  \]
这里用到
\[
  (u^2)_y=2u\,u_y=4yu^2,\qquad v_y=1,\qquad A_y=4yu^2+2v.
\]
也就是说，先把 \(z_x\) 看成
\(\displaystyle 4x\frac{u^2+v}{A}\)，再对 \(y\) 用商法则，就得到上式。
\end{solution}
\begin{exercise}
\par
计算
  \[
    \int_{-4}^0 \mathrm{d}x\int_0^{2x+8}(8-y)^{1/3}\,\mathrm{d}y
    +\int_0^4 \mathrm{d}x\int_0^{-2x+8}(8-y)^{1/3}\,\mathrm{d}y .
  \]
\end{exercise}

\begin{solution}
\par
积分区域可改写为 \(\displaystyle 0\le y\le8,\ \frac{y-8}{2}\le x\le\frac{8-y}{2}\)。对固定的 \(y\)，横向长度为
\[
  \frac{8-y}{2}-\frac{y-8}{2}=8-y.
\]
因此原式为
\[
  \int_0^8(8-y)^{1/3}(8-y)\,\mathrm{d}y
  =\int_0^8(8-y)^{4/3}\,\mathrm{d}y
  =\frac{3}{7}\,8^{7/3}
  =\frac{384}{7}.
\]
换序后被积函数只含 \(y\)，所以先把 \(x\) 方向长度乘进去即可。
\end{solution}
\begin{exercise}
\par
计算
  \(\displaystyle \iint_\Sigma (z^2+x)\,\mathrm{d}y\,\mathrm{d}z-z\,\mathrm{d}x\,\mathrm{d}y\)，
  其中 $\Sigma$ 为旋转抛物面 $\displaystyle z=\frac12(x^2+y^2)$ 介于平面 $z=0$ 与 $z=2$ 之间的部分的下侧。
\end{exercise}

\begin{solution}
\par
对图形 $\displaystyle z=\frac12(x^2+y^2)$ 取下侧，其向量面积元为 $(x,y,-1)\,\mathrm{d}x\,\mathrm{d}y$，投影区域为 $x^2+y^2\le4$。故
  \(\displaystyle \iint_\Sigma (z^2+x)\,\mathrm{d}y\,\mathrm{d}z-z\,\mathrm{d}x\,\mathrm{d}y\) 化为 \(\displaystyle \iint_{x^2+y^2\le4}(xz^2+x^2+z)\,\mathrm{d}x\,\mathrm{d}y\)。
  奇函数项积分为 $0$，且 $z=(x^2+y^2)/2$，所以
  \(\displaystyle I=\iint_{x^2+y^2\le4}\left(x^2+\frac{x^2+y^2}{2}\right)\,\mathrm{d}x\,\mathrm{d}y=8\pi\)。
\end{solution}
\begin{exercise}
\par
计算 \(\displaystyle I=\int_C (z-y)\,\mathrm{d}x+(x-z)\,\mathrm{d}y+(x-y)\,\mathrm{d}z\)，其中 $C$ 是曲线
  \[
    \begin{cases}
      x^2+y^2=1,\\
      x-y+z=2,
    \end{cases}
  \]
  从 $z$ 轴正向往 $z$ 轴负向看 $C$ 的方向是顺时针的。
\end{exercise}

\begin{solution}
\par
令 $\vec F=(z-y,x-z,x-y)$，则 \(\displaystyle \nabla\times\vec F=(0,0,2)\)。
  题设方向从 $z$ 轴正向看为顺时针，对应法向量的 $z$ 分量为负。由 Stokes 公式，
  \(\displaystyle I=\iint_D 2(-1)\,\mathrm{d}x\,\mathrm{d}y=-2\pi\)，
  其中 $D:x^2+y^2\le1$。
这里把曲线换成其所在平面内的圆盘投影来算 Stokes 面积分；题设“从 \(z\) 轴正向看为顺时针”说明所取法向量的 \(z\) 分量为负，所以面积分前有负号。
\end{solution}

\par\medskip
\noindent\textbf{三、解答题（共 3 小题，每小题 10 分，共 30 分）}\par\nobreak\smallskip
\begin{exercise}
\par
设曲线积分
  \(\displaystyle \int_L x\varphi(y)\,\mathrm{d}x+x^2y\,\mathrm{d}y\)
  与路径无关，其中 $\varphi(y)$ 具有连续的导数，且 $\varphi(0)=1$，计算 \(\displaystyle \int_{(0,0)}^{(1,1)}x\varphi(y)\,\mathrm{d}x+x^2y\,\mathrm{d}y\)。
\end{exercise}

\begin{solution}
\par
记 $P=x\varphi(y)$，$Q=x^2y$。路径无关给出 \(\displaystyle P_y=Q_x,\ x\varphi'(y)=2xy\)。
  故 $\varphi'(y)=2y$，由 $\varphi(0)=1$ 得 $\varphi(y)=y^2+1$。势函数可取
  \(\displaystyle \Phi(x,y)=\frac{x^2}{2}(y^2+1)\)，
  所求积分为 $\Phi(1,1)-\Phi(0,0)=1$。
\end{solution}
\begin{exercise}
\par
求级数
  \(\displaystyle \sum_{n=2}^\infty (-1)^n\frac{x^n}{n(n-1)}\)
  的收敛域，并求和函数。
\end{exercise}

\begin{solution}
\par
令 $t=-x$，则 \(\displaystyle \sum_{n=2}^\infty(-1)^n\frac{x^n}{n(n-1)}=\sum_{n=2}^\infty\frac{t^n}{n(n-1)}\)。
  当 $|t|<1$ 时，\(\displaystyle \sum_{n=2}^\infty\frac{t^n}{n(n-1)}=(1-t)\ln(1-t)+t\)。因而 \(\displaystyle S(x)=(1+x)\ln(1+x)-x\)。
  端点 $x=1$ 绝对收敛，$x=-1$ 也收敛，故收敛域为 $[-1,1]$。
\end{solution}
\begin{exercise}
\par
求方程 $y''+y'-2y=-2\sin x$ 的通解。
\end{exercise}

\begin{solution}
\par
齐次方程 \(y''+y'-2y=0\) 的特征方程为
\[
  r^2+r-2=0,
\]
故特征根为 \(1,-2\)，齐次解为 \(C_1e^x+C_2e^{-2x}\)。对右端 \(-2\sin x\)，设特解
\[
  y_p=A\sin x+B\cos x.
\]
代入左端：
\[
  y_p''+y_p'-2y_p=(-3A-B)\sin x+(A-3B)\cos x.
\]
比较系数
\[
  -3A-B=-2,\qquad A-3B=0,
\]
得 \(A=3/5,\ B=1/5\)。
  通解为 \(\displaystyle y=C_1e^x+C_2e^{-2x}+\frac35\sin x+\frac15\cos x\)。
\end{solution}

\par\medskip
\noindent\textbf{四、应用题（10 分）}\par\nobreak\smallskip
\begin{exercise}
\par
在椭圆 $x^2+4y^2=4$ 上求一点，使其到直线 $2x+3y-6=0$ 的距离最短。
\end{exercise}

\begin{solution}
\par
椭圆可参数化为 $x=2\cos t,\ y=\sin t$。点到直线距离为 \(\displaystyle d=\frac{|2x+3y-6|}{\sqrt{13}}\)。在椭圆上有 $2x+3y=4\cos t+3\sin t\le5<6$，故距离最小时需使 $4\cos t+3\sin t$ 最大。取 \(\displaystyle \cos t=\frac45,\ \sin t=\frac35\)，得所求点为 \(\displaystyle \left(\frac85,\frac35\right)\)。
\end{solution}

\par\medskip
\noindent\textbf{五、证明题（10 分）}\par\nobreak\smallskip
\begin{exercise}
\par
证明级数 $\displaystyle \sum_{n=1}^\infty x^2e^{-nx}$ 在 $(0,+\infty)$ 内一致收敛。
\end{exercise}

\begin{solution}
\par
设余项 \(\displaystyle R_N(x)=\sum_{n=N}^\infty x^2e^{-nx}
  =\frac{x^2e^{-Nx}}{1-e^{-x}},\qquad x>0\)。
函数 $h(x)=x^2/(e^x-1)$ 在 $(0,+\infty)$ 上有界，且当 $x\to0^+$ 与 $x\to+\infty$ 时均趋于 $0$。
并且
\[
  R_N(x)=\frac{x^2e^{-Nx}}{1-e^{-x}}
  =\frac{x^2}{e^x-1}\,e^{-(N-1)x}
  =h(x)e^{-(N-1)x}.
\]
给定 \(\varepsilon>0\)，先取 \(\delta>0\)，使 \(0<x<\delta\) 时 \(h(x)<\varepsilon\)，于是
\(R_N(x)\le h(x)<\varepsilon\)。对 \(x\ge\delta\)，由于 \(h\) 有界，设 \(h(x)\le M\)，则
\[
  R_N(x)\le M e^{-(N-1)\delta}.
\]
取 \(N\) 足够大，使右端小于 \(\varepsilon\)。两段合并，得到
\(\sup_{x>0}R_N(x)\to0\)，故原级数在 \((0,+\infty)\) 内一致收敛。
\end{solution}

\section{2022级工科数学分析下 A卷}

\par\medskip
\noindent\textbf{一、填空题（共 5 题，每题 2 分，共 10 分）}\par\nobreak\smallskip
\begin{exercise}
\par
微分方程 $y''+2y'+6y=0$ 的通解为 \underline{\hspace{4cm}}。
\end{exercise}

\begin{solution}
\par
特征方程为
\[
  r^2+2r+6=0.
\]
解得
\[
  r=\frac{-2\pm\sqrt{4-24}}2=-1\pm\sqrt5\,i.
\]
因此齐次方程的实通解为
\[
  y=e^{-x}\left(C_1\cos\sqrt5\,x+C_2\sin\sqrt5\,x\right).
\]
这里使用的是常系数齐次方程的复根规则：若特征根为 \(\alpha\pm\beta i\)，对应实解为 \(e^{\alpha x}\cos\beta x\) 与 \(e^{\alpha x}\sin\beta x\)。
\end{solution}
\begin{exercise}
\par
设函数 $u=2x^2+2y^2+3z^2$，则 $\operatorname{div}(\operatorname{grad}u)=$ \underline{\hspace{3cm}}。
\end{exercise}

\begin{solution}
\par
因为
\[
  \operatorname{div}(\operatorname{grad}u)=\Delta u=u_{xx}+u_{yy}+u_{zz},
\]
而 \(u=2x^2+2y^2+3z^2\)，所以
\[
  u_{xx}=4,\qquad u_{yy}=4,\qquad u_{zz}=6.
\]
故 \(\operatorname{div}(\operatorname{grad}u)=4+4+6=14\)。
也就是说，本题求的是拉普拉斯算子 \(\Delta u\)，只需分别求三个二阶偏导后相加。
\end{solution}
\begin{exercise}
\par
设 $\Gamma$ 为 $y^2=2x$ 上从原点到 $(2,2)$ 的一段，则第一类曲线积分 \(\displaystyle \int_\Gamma y\,\mathrm{d}s=\) \underline{\hspace{3cm}}。
\end{exercise}

\begin{solution}
\par
令 \(y=t,\ x=t^2/2\)，则 \(0\le t\le2\)，且
\[
  \mathrm{d}s=\sqrt{\left(\frac{\mathrm{d}x}{\mathrm{d}t}\right)^2+\left(\frac{\mathrm{d}y}{\mathrm{d}t}\right)^2}\,\mathrm{d}t
  =\sqrt{t^2+1}\,\mathrm{d}t.
\]
因此
\[
  \int_\Gamma y\,\mathrm{d}s
  =\int_0^2 t\sqrt{1+t^2}\,\mathrm{d}t
  =\frac13(1+t^2)^{3/2}\Big|_0^2
  =\frac{5\sqrt5-1}{3}.
\]
\end{solution}
\begin{exercise}
\par
参数曲线
  \[
    \begin{cases}
      x=t-\cos t,\\
      y=\sin t,\\
      z=2t
    \end{cases}
  \]
  在 $t=0$ 对应的点处的切线方程为 \underline{\hspace{5cm}}。
\end{exercise}

\begin{solution}
\par
当 \(t=0\) 时，点为 \((-1,0,0)\)。参数曲线的切向量为
\[
  (x'(t),y'(t),z'(t))=(1+\sin t,\cos t,2),
\]
故在 \(t=0\) 处切向量为 \((1,1,2)\)。切线方程为
  \(\displaystyle \frac{x+1}{1}=\frac{y}{1}=\frac{z}{2}\)。
检查时把 \(t=0\) 代入原曲线得到点，把导数代入得到方向向量；切线只由“过点 + 方向”确定。
\end{solution}
\begin{exercise}
\par
设周期为 $2\pi$ 的函数
  \[
    f(x)=
    \begin{cases}
      -1, & -\pi<x\le 0,\\
      1+x, & 0<x\le \pi,
    \end{cases}
  \]
  则 $f(x)$ 的傅里叶级数在 $x=\pi$ 处收敛于 \underline{\hspace{3cm}}。
\end{exercise}

\begin{solution}
\par
傅里叶级数在周期端点收敛于周期延拓后左右极限的平均值。这里
\[
  f(\pi-0)=1+\pi.
\]
而 \(\pi+0\) 由周期性对应 \(-\pi+0\)，此时函数值为 \(-1\)，所以
\[
  f(\pi+0)=f(-\pi+0)=-1.
\]
因此收敛值为
\[
  \frac{f(\pi-0)+f(\pi+0)}2
  =\frac{1+\pi-1}{2}
  =\frac{\pi}{2}.
\]
这里 \(x=\pi\) 是基本区间端点，不是普通连续点；右极限必须借助 \(2\pi\) 周期性转到 \(-\pi+0\) 来取。
\end{solution}

\par\medskip
\noindent\textbf{二、选择题（共 5 题，每题 2 分，共 10 分）}\par\nobreak\smallskip
\begin{exercise}
\par
关于未知函数 $y$ 的微分方程
  \(\displaystyle (y-\cos^2 x)\,\mathrm{d}x+e^x\,\mathrm{d}y=0\)
  是（\quad）。

  \begin{tabular}{llll}
    A. 可分离变量方程； & B. 一阶线性非齐次方程； &
    C. 一阶线性齐次方程； & D. 非线性方程。
  \end{tabular}
\end{exercise}

\begin{solution}
\par
原方程展开为
\[
  e^x\,\mathrm{d}y=-(y-\cos^2x)\,\mathrm{d}x.
\]
两边除以 \(e^x\,\mathrm{d}x\)，得
\[
  y'+e^{-x}y=e^{-x}\cos^2x.
\]
未知函数 \(y\) 及其导数都是一次，且右端不为 \(0\)，所以它是一阶线性非齐次方程，选 B。这里 \(e^{-x}\cos^2x\) 只是自变量函数，不能因为出现三角函数就误判为非线性；非线性只看 \(y\) 是否以平方、乘积、复合等形式出现。
\end{solution}
\begin{exercise}
\par
若二元函数 $f(x,y)$ 满足
  \(\displaystyle \lim_{(x,y)\to(0,0)}\frac{f(x,y)-f(0,0)+3x-5y}{\sqrt{x^2+y^2}}=0\)，
  则（\quad）。

  \begin{tabular}{ll}
    A. $\mathrm{d}f(0,0)=0$； & B. $\mathrm{d}f(0,0)=3\mathrm{d}x-5\mathrm{d}y$；\\
    C. $\mathrm{d}f(0,0)=-3\mathrm{d}x+5\mathrm{d}y$； & D. $\mathrm{d}f(0,0)$ 不存在。
  \end{tabular}
\end{exercise}

\begin{solution}
\par
可微定义为
\[
  f(x,y)-f(0,0)=A x+B y+o\left(\sqrt{x^2+y^2}\right).
\]
题设极限等于 \(0\)，即
\[
  f(x,y)-f(0,0)+3x-5y=o\left(\sqrt{x^2+y^2}\right),
\]
所以线性主部为 \(-3x+5y\)。因此
\[
  \mathrm{d}f(0,0)=-3\,\mathrm{d}x+5\,\mathrm{d}y,
\]
选 C。
本题的核心是把极限中的小 \(o\) 项与可微定义对齐：线性项 \(Ax+By\) 必须与 \(-3x+5y\) 相同，所以 \(A=f_x(0,0)=-3,\ B=f_y(0,0)=5\)。
\end{solution}
\begin{exercise}
\par
设函数 $f(x,y)$ 与 $\varphi(x,y)$ 均可微，且 $\varphi_y(x,y)\ne0$。若 $(x_0,y_0)$ 是 $f(x,y)$ 在约束条件
  $\varphi(x,y)=0$ 下的一个极值点，则下列选项正确的是（\quad）。

  \begin{tabular}{ll}
    A. 若 $f_x(x_0,y_0)=0$，则 $f_y(x_0,y_0)=0$； &
    B. 若 $f_x(x_0,y_0)=0$，则 $f_y(x_0,y_0)\ne0$；\\
    C. 若 $f_x(x_0,y_0)\ne0$，则 $f_y(x_0,y_0)=0$； &
    D. 若 $f_x(x_0,y_0)\ne0$，则 $f_y(x_0,y_0)\ne0$。
  \end{tabular}
\end{exercise}

\begin{solution}
\par
因为 \(\varphi_y(x_0,y_0)\ne0\)，约束方程 \(\varphi(x,y)=0\) 可在点附近解出
\(y=y(x)\)。约束极值的一阶必要条件等价于
\[
  f_x\varphi_y-f_y\varphi_x=0.
\]
若 \(f_x(x_0,y_0)\ne0\) 而 \(f_y(x_0,y_0)=0\)，上式会变成
\(f_x\varphi_y=0\)，与 \(f_x\ne0,\ \varphi_y\ne0\) 矛盾。因此 \(f_x\ne0\) 时必有 \(f_y\ne0\)，选 D。
\end{solution}
\begin{exercise}
\par
函数 $\ln(1+x)$ 在 $x=0$ 处的泰勒展开式正确的是（\quad）。

  \begin{tabular}{ll}
    A. $\displaystyle \ln(1+x)=\sum_{n=1}^{\infty}\frac{(-1)^n}{n}x^n,\ x\in(-1,1]$； &
    B. $\displaystyle \ln(1+x)=\sum_{n=1}^{\infty}\frac{(-1)^{n-1}}{n}x^n,\ x\in(-1,1]$；\\
    C. $\displaystyle \ln(1+x)=-\sum_{n=1}^{\infty}\frac{1}{n}x^n,\ x\in[-1,1)$； &
    D. $\displaystyle \ln(1+x)=\sum_{n=1}^{\infty}\frac{1}{n}x^n,\ x\in[-1,1)$。
  \end{tabular}
\end{exercise}

\begin{solution}
\par
由
\[
  \frac1{1+x}=1-x+x^2-x^3+\cdots,\qquad |x|<1,
\]
从 \(0\) 到 \(x\) 逐项积分，得
\[
  \ln(1+x)=\sum_{n=1}^{\infty}\frac{(-1)^{n-1}}{n}x^n.
\]
端点 \(x=1\) 为交错调和级数收敛，\(x=-1\) 为 \(-\sum1/n\) 发散，所以收敛区间为
\((-1,1]\)。选 B。
因此不仅符号要选 \((-1)^{n-1}\)，端点范围也要同时核对；只写 \(|x|<1\) 会漏掉 \(x=1\) 这个收敛端点。
\end{solution}
\begin{exercise}
\par
使得级数
  \(\displaystyle \sum_{n=1}^{\infty}\frac{(-1)^n}{n^p}\)
  条件收敛的常数 $p$ 的取值范围是（\quad）。

  \begin{tabular}{llll}
    A. $p\le0$； & B. $0<p\le1$； & C. $0<p<1$； & D. $p>1$。
  \end{tabular}
\end{exercise}

\begin{solution}
\par
若 \(p\le0\)，则 \(1/n^p\) 不趋于 \(0\)，级数发散。若 \(p>0\)，
\(\{1/n^p\}\) 单调趋于 \(0\)，由交错级数判别法收敛。再看绝对收敛：
\[
  \sum\left|\frac{(-1)^n}{n^p}\right|=\sum\frac1{n^p}
\]
在 \(p>1\) 时收敛。因此条件收敛要求 \(p>0\) 且 \(p\le1\)，即
\(\displaystyle 0<p\le1\)，选 B。
\end{solution}

\par\medskip
\noindent\textbf{三、计算题（共 3 题，每题 10 分，共 30 分）}\par\nobreak\smallskip
\begin{exercise}
\par
设
  \(\displaystyle z=\frac{1}{x}f(x+y,x-y)+g(xy)\)，
  其中函数 $f$ 有二阶连续的偏导数，且 $g$ 二阶可导，计算 $\displaystyle \dfrac{\partial z}{\partial x}$ 和
  $\displaystyle \dfrac{\partial^2z}{\partial x\partial y}$。
\end{exercise}

\begin{solution}
\par
记 $u=x+y,\ v=x-y$，并把 $f_1,f_2,f_{11},f_{22}$ 均理解为在 $(u,v)$ 处的取值，把 $g',g''$
  理解为在 $xy$ 处的取值。先对 $x$ 求导：\(\displaystyle \frac{\partial z}{\partial x}=-\frac{f}{x^2}+\frac{f_1+f_2}{x}+y\,g'(xy)\)。
  再对 \(y\) 求导。对 \(-f/x^2\) 求导得到
\(-\,(f_1-f_2)/x^2\)；对 \((f_1+f_2)/x\) 求导，由
\[
  (f_1+f_2)_y=(f_{11}+f_{21})-(f_{12}+f_{22})=f_{11}-f_{22}
\]
得到 \((f_{11}-f_{22})/x\)；最后
\(\partial_y[y g'(xy)]=g'(xy)+xyg''(xy)\)。故
  \(\displaystyle \frac{\partial^2z}{\partial x\partial y}=-\frac{f_1-f_2}{x^2}+\frac{f_{11}-f_{22}}{x}+g'(xy)+xy\,g''(xy)\)。
\end{solution}
\begin{exercise}
\par
计算累次积分
  \[
    \int_{-\sqrt2}^{0}\mathrm{d}x\int_{-x}^{\sqrt{4-x^2}}\sqrt{x^2+y^2}\,\mathrm{d}y
    +\int_{0}^{2}\mathrm{d}x\int_{\sqrt{2x-x^2}}^{\sqrt{4-x^2}}\sqrt{x^2+y^2}\,\mathrm{d}y.
  \]
\end{exercise}

\begin{solution}
\par
第一部分区域为极坐标下 \(\displaystyle \frac{\pi}{2}\le\theta\le\frac{3\pi}{4},\qquad 0\le r\le2\)，第二部分区域为 \(\displaystyle 0\le\theta\le\frac{\pi}{2},\ 2\cos\theta\le r\le2\)。因 $\sqrt{x^2+y^2}=r$，面积元为 $r\,\mathrm{d}r\,\mathrm{d}\theta$，所以原积分为
  \[
  \int_0^{\pi/2}\int_{2\cos\theta}^{2}r^2\,\mathrm{d}r\,\mathrm{d}\theta
  +\int_{\pi/2}^{3\pi/4}\int_0^2r^2\,\mathrm{d}r\,\mathrm{d}\theta.
  \]
计算得
\[
  \frac83\int_0^{\pi/2}(1-\cos^3\theta)\,\mathrm{d}\theta
  +\frac83\left(\frac{3\pi}{4}-\frac{\pi}{2}\right)
  =2\pi-\frac{16}{9}.
\]
\end{solution}
\begin{exercise}
\par
计算锥面 $z=\sqrt{x^2+y^2}$ 被柱面 $x^2+y^2=4x$ 截下的部分锥面面积。
\end{exercise}

\begin{solution}
\par
锥面写成 $z=\sqrt{x^2+y^2}$ 时，在非原点处有
\[
  z_x=\frac{x}{\sqrt{x^2+y^2}},\qquad z_y=\frac{y}{\sqrt{x^2+y^2}},
\]
所以
\[
  \mathrm{d}S=\sqrt{1+z_x^2+z_y^2}\,\mathrm{d}x\,\mathrm{d}y=\sqrt2\,\mathrm{d}x\,\mathrm{d}y.
\]
投影区域为 \(x^2+y^2\le4x\)，即 \((x-2)^2+y^2\le4\)，是半径为 \(2\) 的圆盘，面积为 \(4\pi\)。故锥面面积为
\[
  S=\sqrt2\cdot4\pi=4\sqrt2\,\pi.
\]
\end{solution}

\par\medskip
\noindent\textbf{四、综合题（共 3 题，每题 10 分，共 30 分）}\par\nobreak\smallskip
\begin{exercise}
\par
设 $\Sigma$ 是 \(\displaystyle \frac{x^2}{a^2}+\frac{y^2}{b^2}+\frac{z^2}{c^2}=1\quad (z\ge0)\) 的上侧，且 $a,b,c$ 都是正实数。计算第二类曲面积分 \(\displaystyle \iint_\Sigma (x^2-y\sin x)\,\mathrm{d}y\,\mathrm{d}z+(y^2-z^2)\,\mathrm{d}z\,\mathrm{d}x+(z^2-x^2)\,\mathrm{d}x\,\mathrm{d}y\)。
\end{exercise}

\begin{solution}
\par
记 \(\displaystyle P=x^2-y\sin x,\ Q=y^2-z^2,\ R=z^2-x^2\)。用底面 \(\displaystyle \Sigma_0:\ z=0,\quad \frac{x^2}{a^2}+\frac{y^2}{b^2}\le1\) 补成上半椭球体 $\Omega$ 的外侧闭曲面，底面取向向下。由高斯公式，
  \(\displaystyle I_\Sigma+I_{\Sigma_0}=\iiint_\Omega (P_x+Q_y+R_z)\,\mathrm{d}V=\iiint_\Omega (2x-y\cos x+2y+2z)\,\mathrm{d}V\)。
  上半椭球关于 $x,y$ 对称，含 $x,y$ 的奇函数项积分为 $0$，故
  \(\displaystyle I_\Sigma+I_{\Sigma_0}=\iiint_\Omega 2z\,\mathrm{d}V=\frac{\pi ab c^2}{2}\)。
  底面向下时
  \(\displaystyle I_{\Sigma_0}=\iint_{\frac{x^2}{a^2}+\frac{y^2}{b^2}\le1}x^2\,\mathrm{d}x\,\mathrm{d}y=\frac{\pi a^3b}{4}\)。
  因而 \(\displaystyle I_\Sigma=\frac{\pi ab c^2}{2}-\frac{\pi a^3b}{4}\)。
\end{solution}
\begin{exercise}
\par
计算曲线积分 \(\displaystyle \int_\Gamma \frac{(x+y)\,\mathrm{d}x-(x-y)\,\mathrm{d}y}{x^2+y^2}\),
  其中曲线 $\Gamma$ 是从点 $A(-1,0)$ 到点 $B(1,0)$ 的一条不经过原点的光滑曲线 $y=f(x)$，
  $-1\le x\le1$。
\end{exercise}

\begin{solution}
\par
积分可拆成
  \(\displaystyle \frac{(x+y)\,\mathrm{d}x-(x-y)\,\mathrm{d}y}{x^2+y^2}=\frac{x\,\mathrm{d}x+y\,\mathrm{d}y}{x^2+y^2}+\frac{y\,\mathrm{d}x-x\,\mathrm{d}y}{x^2+y^2}=\mathrm{d}\ln r-\mathrm{d}\theta\)。
  两端点到原点的距离都为 $1$，故 $\mathrm{d}\ln r$ 的积分为 $0$。剩下的值由曲线绕过原点的方向决定。若
  $f(0)>0$，曲线从上半平面绕过原点，$\theta$ 可由 $\pi$ 连续变到 $0$，积分为 $\pi$；若
  $f(0)<0$，曲线从下半平面绕过原点，积分为 $-\pi$。因此
  \(\displaystyle \int_\Gamma \frac{(x+y)\,\mathrm{d}x-(x-y)\,\mathrm{d}y}{x^2+y^2}=\begin{cases}\pi,&f(0)>0,\\-\pi,&f(0)<0.\end{cases}\)
\end{solution}
\begin{exercise}
\par
求幂级数
  \(\displaystyle \sum_{n=0}^{\infty}(n+1)x^{2n+1}\)
  的收敛域及和函数。
\end{exercise}

\begin{solution}
\par
令 $t=x^2$，则 \(\displaystyle \sum_{n=0}^{\infty}(n+1)x^{2n+1}=x\sum_{n=0}^{\infty}(n+1)t^n\)。
  当 $|t|<1$ 即 $|x|<1$ 时，\(\displaystyle \sum_{n=0}^{\infty}(n+1)t^n=\frac{1}{(1-t)^2}\)，所以和函数为 \(\displaystyle S(x)=\frac{x}{(1-x^2)^2},\ |x|<1\)。当 $x=\pm1$ 时通项不趋于 $0$，故发散。因此收敛域为 $(-1,1)$。
\end{solution}

\par\medskip
\noindent\textbf{五、证明题（共 1 题，每题 10 分，共 10 分）}\par\nobreak\smallskip
\begin{exercise}
\par
证明函数项级数 \(\displaystyle \sum_{n=1}^{\infty}\sin\frac{n!x^n}{x^2+n^n}\) 在区间 $[-2,2]$ 上一致收敛。
\end{exercise}

\begin{solution}
\par
对任意 $x\in[-2,2]$，有 \(\displaystyle \left|\sin\frac{n!x^n}{x^2+n^n}\right|\le\left|\frac{n!x^n}{x^2+n^n}\right|\le\frac{n!\,2^n}{n^n}\)。
令 $\displaystyle M_n=\dfrac{n!\,2^n}{n^n}$，则 \(\displaystyle \frac{M_{n+1}}{M_n}=2\left(\frac{n}{n+1}\right)^n\longrightarrow \frac{2}{e}<1\)。故 $\displaystyle \sum M_n$ 收敛。由魏尔斯特拉斯判别法，原函数项级数在 $[-2,2]$ 上一致收敛。
\end{solution}

\par\medskip
\noindent\textbf{六、应用题（共 1 题，每题 10 分，共 10 分）}\par\nobreak\smallskip
\begin{exercise}
\par
制作一个体积为 $2$ 的长方体盒子，利用拉格朗日乘数法求出长、宽和高分别为多少时才能使其表面积最小。
\end{exercise}

\begin{solution}
\par
设长方体长、宽、高分别为 $x,y,z>0$，体积约束为 $xyz=2$，闭盒子的表面积为 \(\displaystyle S=2(xy+xz+yz)\)。设 \(\displaystyle L=2(xy+xz+yz)+\lambda(xyz-2)\)。
由拉格朗日乘数法，\(\displaystyle 2(y+z)+\lambda yz=0,\quad 2(x+z)+\lambda xz=0,\quad 2(x+y)+\lambda xy=0\)。两两相减得 $x=y=z$。再由 $xyz=2$，得 \(\displaystyle x=y=z=\sqrt[3]{2}\)。
因此长、宽和高都等于 $\sqrt[3]{2}$ 时表面积最小。
\end{solution}

\section{2022级工科数学分析下 B卷}

\par\medskip
\noindent\textbf{一、填空题：共 5 题，每题 2 分，共 10 分。}\par\nobreak\smallskip
\begin{exercise}
\par
微分方程 $y''+3y'+2y=0$ 的通解为\underline{\hspace{4cm}}。
\end{exercise}

\begin{solution}
\par
特征方程为
\[
  r^2+3r+2=0=(r+1)(r+2).
\]
故特征根为 \(r=-1,-2\)，齐次方程的通解为
\[
  y=C_1e^{-x}+C_2e^{-2x}.
\]
两个特征根互异，所以直接对应两个线性无关的指数解 \(e^{-x}\) 与 \(e^{-2x}\)，不需要再乘 \(x\)。
也就是说，互异实根 \(r_1,r_2\) 的标准通解就是 \(C_1e^{r_1x}+C_2e^{r_2x}\)；只有重根时才需要补 \(xe^{rx}\)。
\end{solution}
\begin{exercise}
\par
函数 $u=2xy-z^2$ 在点 $(2,-1,-1)$ 处方向导数的最大值为\underline{\hspace{4cm}}。
\end{exercise}

\begin{solution}
\par
方向导数的最大值等于梯度的模。这里
\[
  \nabla u=(2y,2x,-2z).
\]
在点 \((2,-1,-1)\) 处，
\[
  \nabla u=(-2,4,2),
\]
因此最大方向导数为
\[
  |\nabla u|=\sqrt{(-2)^2+4^2+2^2}=2\sqrt6.
\]
方向导数公式为 \(D_{\mathbf e}u=\nabla u\cdot\mathbf e\)，其中 \(\mathbf e\) 是单位方向向量。由柯西不等式，最大值在 \(\mathbf e\) 与梯度同向时取得，所以最大方向导数等于梯度模长。
\end{solution}
\begin{exercise}
\par
设 \(\Gamma\) 是圆周 \(x^2+y^2=1\) 在第一象限的部分。则第一类曲线积分
  \[
    \int_\Gamma (x+y)^2\,\mathrm{d}s=\underline{\hspace{4cm}}.
  \]
\end{exercise}

\begin{solution}
\par
令
\[
  x=\cos t,\qquad y=\sin t,\qquad 0\le t\le\frac{\pi}{2}.
\]
这是单位圆第一象限弧，且 \(\mathrm{d}s=\mathrm{d}t\)。于是
\[
  \int_\Gamma (x+y)^2\,\mathrm{d}s
  =\int_0^{\pi/2}(\cos t+\sin t)^2\,\mathrm{d}t
  =\int_0^{\pi/2}(1+\sin2t)\,\mathrm{d}t
  =\frac{\pi}{2}+1.
\]
这里是第一类曲线积分，方向不影响结果；参数 \(t\) 只负责描出第一象限圆弧并给出弧长元 \(\mathrm{d}s=\mathrm{d}t\)。
\end{solution}
\begin{exercise}
\par
级数 $\displaystyle \sum_{n=0}^{\infty}(-1)^n\dfrac{2n+3}{(2n+1)!}$ 的和为\underline{\hspace{4cm}}。
\end{exercise}

\begin{solution}
\par
把分子拆成 \(2n+3=(2n+1)+2\)，则
\[
  \frac{2n+3}{(2n+1)!}
  =\frac{1}{(2n)!}+\frac{2}{(2n+1)!}.
\]
因此
\[
  \sum_{n=0}^{\infty}(-1)^n\frac{2n+3}{(2n+1)!}
  =\sum_{n=0}^{\infty}\frac{(-1)^n}{(2n)!}
   +2\sum_{n=0}^{\infty}\frac{(-1)^n}{(2n+1)!}
  =\cos1+2\sin1.
\]
\end{solution}
\begin{exercise}
\par
设周期为 $2\pi$ 的函数
  \[
    f(x)=
    \begin{cases}
      x^2, & -\pi<x\le 0,\\
      0, & 0<x\le \pi,
    \end{cases}
  \]
  则 $f(x)$ 的傅里叶（Fourier）级数在 $x=-\pi$ 处收敛于\underline{\hspace{4cm}}。
\end{exercise}

\begin{solution}
\par
傅里叶级数在周期端点处取周期延拓后左右极限的平均值。由题设，
\[
  f(-\pi+0)=(-\pi)^2=\pi^2.
\]
而 \(x=-\pi-0\) 由周期性对应 \(x=\pi-0\)，此时右半段函数值为 \(0\)，所以
\[
  f(-\pi-0)=f(\pi-0)=0.
\]
因此收敛值为
\[
  \frac{f(-\pi-0)+f(-\pi+0)}2=\frac{0+\pi^2}{2}=\frac{\pi^2}{2}.
\]
\end{solution}

\par\medskip
\noindent\textbf{二、选择题：共 5 题，每题 2 分，共 10 分。}\par\nobreak\smallskip
\begin{exercise}
\par
下列微分方程中，属于二阶线性常微分方程的是（\quad）
  \par\noindent\begin{tabular}{@{}p{0.48\linewidth}p{0.48\linewidth}@{}}
    1. $y''+x^3y'+(\ln x)y=\tan x$； & 2. $(x^2+y^2)\,\mathrm{d}y+y^2\,\mathrm{d}x=0$；\\
    3. $(y'')^2+x^2y=1$； & 4. $y''+2\ln y=2x$。
  \end{tabular}
\end{exercise}

\begin{solution}
\par
二阶线性常微分方程中，未知函数 \(y\)、\(y'\)、\(y''\) 只能一次出现，系数只能依赖于自变量。第 1 项
\[
  y''+x^3y'+(\ln x)y=\tan x
\]
满足这个结构，是二阶线性非齐次方程。第 2 项是一阶微分形式；第 3 项含
\((y'')^2\)，非线性；第 4 项含 \(\ln y\)，非线性。因此选第 1 项。
这里第 1 项虽然含有 \(x^3,\ln x,\tan x\)，但这些都只是自变量 \(x\) 的函数；线性性只限制未知函数 \(y\)、\(y'\)、\(y''\) 的出现方式。
\end{solution}
\begin{exercise}
\par
已知函数 $f(x)$ 具有二阶连续导函数，且 $f(x)>0,\ f'(0)=0$（$f$ 在 $0$ 处的导数为 $0$）。则函数 $z=f(x)\ln f(y)$ 在点 $(0,0)$ 处取得极小值的一个充分条件是（\quad）
  \par\noindent\begin{tabular}{@{}p{0.48\linewidth}p{0.48\linewidth}@{}}
    1. $f(0)>1,\ f''(0)>0$； & 2. $f(0)>1,\ f''(0)<0$；\\
    3. $f(0)<1,\ f''(0)>0$； & 4. $f(0)<1,\ f''(0)<0$。
  \end{tabular}
\end{exercise}

\begin{solution}
\par
设 \(a=f(0)>0\)。因为 \(f'(0)=0\)，有
\[
  z_x(0,0)=f'(0)\ln f(0)=0,\qquad
  z_y(0,0)=f(0)\frac{f'(0)}{f(0)}=0,
\]
所以 \((0,0)\) 是驻点。二阶偏导在该点为
\[
  z_{xx}(0,0)=f''(0)\ln f(0),\qquad
  z_{yy}(0,0)=f''(0),\qquad z_{xy}(0,0)=0.
\]
若 \(f(0)>1\)，则 \(\ln f(0)>0\)；再有 \(f''(0)>0\)，Hessian 矩阵正定，故取得极小值。因此选第 1 项。
\end{solution}
\begin{exercise}
\par
函数 $f(x,y)=1+x+y$ 在区域 $\{(x,y)\mid x^2+y^2\le 1\}$ 上的最大值与最小值之积是（\quad）
  \par\noindent\begin{tabular}{@{}p{0.48\linewidth}p{0.48\linewidth}@{}}
    1. $-1$； & 2. $1$；\\
    3. $3-\sqrt2$； & 4. $1+\sqrt2$。
  \end{tabular}
\end{exercise}

\begin{solution}
\par
在线性函数 \(1+x+y\) 中，只需看 \(x+y\) 在单位圆盘上的取值范围。由柯西不等式，
\[
  x+y\le \sqrt2\sqrt{x^2+y^2}\le\sqrt2,
\]
且下界为 \(-\sqrt2\)。因此最大值为 \(1+\sqrt2\)，最小值为 \(1-\sqrt2\)，二者乘积为
\[
  (1+\sqrt2)(1-\sqrt2)=-1.
\]
选第 1 项。
等号在 \((x,y)=\pm(1/\sqrt2,1/\sqrt2)\) 方向上取得，所以最大值和最小值都确实落在闭圆盘边界上，不只是估计上界。
\end{solution}
\begin{exercise}
\par
设 \(\Gamma\) 是以 \((1,1),(-1,1),(-1,-1),(1,-1)\) 为顶点的正方形边界。则第一类曲线积分
  \[
    \oint_\Gamma \frac{x+y+1}{|x|+|y|}\,\mathrm{d}s
  \]
  等于（\quad）
  \par\noindent\begin{tabular}{@{}p{0.48\linewidth}p{0.48\linewidth}@{}}
    1. $0$； & 2. $8$；\\
    3. $4\ln 2$； & 4. $8\ln 2$。
  \end{tabular}
\end{exercise}

\begin{solution}
\par
第一类曲线积分与方向无关，沿四条边分段计算。上边 \(y=1,\ -1\le x\le1\) 上，
\[
  \int_{-1}^1\frac{x+2}{|x|+1}\,\mathrm{d}x=4\ln2.
\]
右边 \(x=1,\ -1\le y\le1\) 上同理也为 \(4\ln2\)。下边 \(y=-1\) 上，被积函数为
\(\displaystyle x/(|x|+1)\)，在 \([-1,1]\) 上为奇函数，积分为 \(0\)；左边 \(x=-1\) 上同理也为
\(0\)。故总积分为 \(8\ln2\)，选第 4 项。
\end{solution}
\begin{exercise}
\par
设 $a$ 为常数，则级数 \(\displaystyle \sum_{n=1}^{\infty}\left(\frac{\sin(an)}{n^2}-\frac1{\sqrt n}\right)\)（\quad）
  \par\noindent\begin{tabular}{@{}p{0.48\linewidth}p{0.48\linewidth}@{}}
    1. 发散； & 2. 绝对收敛；\\
    3. 条件收敛； & 4. 敛散性与 $a$ 的取值相关。
  \end{tabular}
\end{exercise}

\begin{solution}
\par
对任意常数 \(a\)，有 \(|\sin(an)|\le1\)，所以
\[
  \sum_{n=1}^{\infty}\left|\frac{\sin(an)}{n^2}\right|
  \le\sum_{n=1}^{\infty}\frac1{n^2}
\]
绝对收敛。而 \(\displaystyle \sum 1/\sqrt n\) 是 \(p=1/2\) 的 \(p\)-级数，发散。原级数等于一个收敛级数减去一个发散到 \(+\infty\) 的正项级数，因此发散，选第 1 项。
\end{solution}

\par\medskip
\noindent\textbf{三、计算题：共 3 题，每题 10 分，共 30 分。}\par\nobreak\smallskip
\begin{exercise}
\par
设
  \(\displaystyle g(x,y)=f\left(xy,\frac12(x^2-y^2)\right)\)，
  其中 $f(u,v)$ 具有连续的二阶偏导数，且满足
  \(\displaystyle \frac{\partial^2 f}{\partial u^2}+\frac{\partial^2 f}{\partial v^2}=1\)。
  计算二阶偏导数 $\displaystyle \frac{\partial^2 g}{\partial x^2}+\frac{\partial^2 g}{\partial y^2}$。
\end{exercise}

\begin{solution}
\par
记 $\displaystyle u=xy,\ v=\dfrac12(x^2-y^2)$，并用 $f_1,f_2$ 表示 $f$ 对第一、第二个变量的偏导数。则 \(\displaystyle g_x=yf_1+xf_2,\ g_y=xf_1-yf_2\)。
  继续求导得 \(\displaystyle g_{xx}=y^2f_{11}+2xyf_{12}+x^2f_{22}+f_2\)，\(\displaystyle g_{yy}=x^2f_{11}-2xyf_{12}+y^2f_{22}-f_2\)。因此 \(\displaystyle g_{xx}+g_{yy}=(x^2+y^2)(f_{11}+f_{22})=x^2+y^2\)。
\end{solution}
\begin{exercise}
\par
计算累次积分 \(\displaystyle \int_0^1 \mathrm{d}x\int_0^x \mathrm{d}y\int_0^y \frac{\sin z}{(1-z)^2}\,\mathrm{d}z\)。
\end{exercise}

\begin{solution}
\par
积分区域为 $0\le z\le y\le x\le 1$。交换积分次序得 \(\displaystyle \int_0^1\frac{\sin z}{(1-z)^2}\left(\int_z^1\mathrm{d}y\int_y^1\mathrm{d}x\right)\mathrm{d}z=\int_0^1\frac{\sin z}{(1-z)^2}\cdot\frac{(1-z)^2}{2}\,\mathrm{d}z\)。
  故原式为 \(\displaystyle \frac12\int_0^1\sin z\,\mathrm{d}z=\frac{1-\cos 1}{2}\)。
\end{solution}
\begin{exercise}
\par
计算二重积分
  \(\displaystyle \iint_D \sqrt{x^2+y^2}\,\mathrm{d}x\,\mathrm{d}y\)，
  其中 $D$ 是由 $x=\sqrt{2y-y^2}$ 与 $y=x$ 所围成的闭区域。
\end{exercise}

\begin{solution}
\par
曲线 $x=\sqrt{2y-y^2}$ 等价于 $x^2+(y-1)^2=1$，在极坐标下边界为 $r=2\sin\theta$；直线 $y=x$ 为 $\displaystyle \theta=\dfrac{\pi}{4}$。故 \(\displaystyle D:\quad 0\le \theta\le \frac{\pi}{4},\quad 0\le r\le 2\sin\theta\)。于是
  \(\displaystyle \iint_D\sqrt{x^2+y^2}\,\mathrm{d}x\,\mathrm{d}y=\int_0^{\pi/4}\int_0^{2\sin\theta}r^2\,\mathrm{d}r\,\mathrm{d}\theta=\frac{16-10\sqrt2}{9}\)。
\end{solution}

\par\medskip
\noindent\textbf{四、综合题：共 3 题，每题 10 分，共 30 分。}\par\nobreak\smallskip
\begin{exercise}
\par
设 $\Sigma$ 是曲面 $z=1-x^2-y^2\ (z\ge 0)$ 的上侧，计算第二型曲面积分 \(\displaystyle \iint_\Sigma 2x^3\,\mathrm{d}y\,\mathrm{d}z+2y^3\,\mathrm{d}z\,\mathrm{d}x+3(z^2-1)\,\mathrm{d}x\,\mathrm{d}y\)。
\end{exercise}

\begin{solution}
\par
用平面 \(\Sigma_1:z=0,\ x^2+y^2\le 1\) 补成闭曲面，取外侧方向。由高斯公式，
  \[
  \iint_{\Sigma+\Sigma_1}2x^3\,\mathrm{d}y\,\mathrm{d}z+2y^3\,\mathrm{d}z\,\mathrm{d}x+3(z^2-1)\,\mathrm{d}x\,\mathrm{d}y
  =\iiint_\Omega (6x^2+6y^2+6z)\,\mathrm{d}V=2\pi .
  \]
  在补面 \(\Sigma_1\) 上取下侧，积分为 \(\displaystyle I_{\Sigma_1}=-\iint_{x^2+y^2\le 1}3(z^2-1)\,\mathrm{d}x\,\mathrm{d}y=3\pi\)。因此所求曲面积分为
  \(\displaystyle 2\pi-3\pi=-\pi\)。
\end{solution}
\begin{exercise}
\par
计算曲线积分
  \(\displaystyle \oint_L yz\,\mathrm{d}x+3zx\,\mathrm{d}y-xy\,\mathrm{d}z\)，
  其中 $L$ 是曲线
  \[
    \begin{cases}
      x^2+y^2=2y,\\
      y-z=1,
    \end{cases}
  \]
  其方向从 $z$ 轴正向往 $z$ 轴负向看去为逆时针方向。
\end{exercise}

\begin{solution}
\par
按题设方向可取参数
  \(\displaystyle x=\cos t,\ y=1+\sin t,\ z=\sin t,\ 0\le t\le 2\pi\)。
  代入得
  \[
  \oint_L yz\,\mathrm{d}x+3zx\,\mathrm{d}y-xy\,\mathrm{d}z
  =\int_0^{2\pi}\left[-(1+\sin t)\sin^2t+3\sin t\cos^2t-(1+\sin t)\cos^2t\right]\mathrm{d}t .
  \]
  利用一个周期内奇函数项积分为 $0$，得 \(\displaystyle -\int_0^{2\pi}\sin^2t\,\mathrm{d}t-\int_0^{2\pi}\cos^2t\,\mathrm{d}t=-2\pi\)。
\end{solution}
\begin{exercise}
\par
将函数
  \(\displaystyle f(x)=\frac{x}{x^2-5x+6}\)
  展开成 $x-1$ 的幂级数，并指出其收敛域。
\end{exercise}

\begin{solution}
\par
分解为 \(\displaystyle \frac{x}{x^2-5x+6}=\frac{3}{x-3}-\frac{2}{x-2}\)。令 $t=x-1$，则
  \begin{align*}
    \frac{3}{x-3}&=-\frac32\cdot\frac1{1-t/2}
    =-\frac32\sum_{n=0}^{\infty}\left(\frac{t}{2}\right)^n,\quad |t|<2,\\
    -\frac{2}{x-2}&=\frac2{1-t}
    =2\sum_{n=0}^{\infty}t^n,\quad |t|<1.
  \end{align*}
  因此 \(\displaystyle f(x)=\sum_{n=0}^{\infty}\left(2-\frac{3}{2^{n+1}}\right)(x-1)^n,\qquad |x-1|<1\)。
  收敛域为 $(0,2)$。
\end{solution}

\par\medskip
\noindent\textbf{五、证明题：共 1 题，每题 10 分，共 10 分。}\par\nobreak\smallskip
\begin{exercise}
\par
证明函数项级数 \(\displaystyle \sum_{n=1}^{\infty}\arctan\frac{2x}{x^2+n^3}\) 在区间 $(-\infty,+\infty)$ 上一致收敛。
\end{exercise}

\begin{solution}
\par
对任意 $x\in\mathbb R$，由 $|\arctan u|\le |u|$ 以及
\(\displaystyle x^2+n^3\ge 2|x|n^{3/2}\)
可得 \(\displaystyle \left|\arctan\frac{2x}{x^2+n^3}\right|
  \le
  \frac{2|x|}{x^2+n^3}
  \le
  \frac1{n^{3/2}}\)。
而级数 $\displaystyle \sum_{n=1}^{\infty}n^{-3/2}$ 收敛，故由 Weierstrass 判别法知 \(\displaystyle \sum_{n=1}^{\infty}\arctan\frac{2x}{x^2+n^3}\) 在 $(-\infty,+\infty)$ 上一致收敛。
\end{solution}

\par\medskip
\noindent\textbf{六、解答题：共 1 题，每题 10 分，共 10 分。}\par\nobreak\smallskip
\begin{exercise}
\par
限制点 $(x,y)$ 在圆周 $(x-1)^2+y^2=1$ 上变化时，求函数 $f(x,y)=xy$ 的最小值与最大值。
\end{exercise}

\begin{solution}
\par
令约束函数
\(\displaystyle \varphi(x,y)=(x-1)^2+y^2-1=0\)。
对 $F(x,y,\lambda)=xy+\lambda\varphi(x,y)$，有
\(\displaystyle y+2\lambda(x-1)=0,\qquad x+2\lambda y=0\)。
由约束得 $x^2+y^2=2x$。当 $\lambda=0$ 时得点 $(0,0)$，函数值为 $0$，不是最值点。设 $y\ne0$，由第二个方程得
\(\displaystyle \lambda=-\frac{x}{2y}\)。
代入第一个方程，得
\(\displaystyle y^2=x(x-1)\)。
再与 $y^2=2x-x^2$ 联立，得 \(\displaystyle x(x-1)=2x-x^2,\qquad x=\frac32,\qquad y=\pm\frac{\sqrt3}{2}\)。
因此 \(\displaystyle f\left(\frac32,\frac{\sqrt3}{2}\right)=\frac{3\sqrt3}{4}\) 为最大值，\(\displaystyle f\left(\frac32,-\frac{\sqrt3}{2}\right)=-\frac{3\sqrt3}{4}\) 为最小值。
\end{solution}

\section{2023级工科数学分析（二）期末考试 A 卷}

\par\medskip
\noindent\textbf{一、填空题（共 5 小题，每小题 2 分，共 10 分）}\par\nobreak\smallskip
\begin{exercise}
\par
方程 $\displaystyle y'-\dfrac{y}{x}=x^2$ 的通解为 \underline{\hspace{4cm}}。
\end{exercise}

\begin{solution}
\par
这是线性方程
\[
  y'-\frac1x y=x^2.
\]
积分因子为 \(\displaystyle \mu(x)=e^{\int -1/x\,\mathrm{d}x}=1/x\)。两边乘以 \(1/x\)，得
\[
  \left(\frac yx\right)'=x.
\]
积分得 \(\displaystyle y/x=x^2/2+C\)，故
\[
  y=Cx+\frac{x^3}{2}.
\]
由于标准形中含有 \(1/x\)，通解应理解为在不跨过 \(x=0\) 的区间上成立；积分因子法本身就是局部区间上的解法。
\end{solution}
\begin{exercise}
\par
设 $z=\ln(x^2+y^2)$，$\displaystyle \vec n=\left(\cos\dfrac{\pi}{3},\sin\dfrac{\pi}{3}\right)$。则 \(\displaystyle \left.\frac{\partial z}{\partial \vec n}\right|_{(1,1)}=\underline{\hspace{4cm}}\)。
\end{exercise}

\begin{solution}
\par
先求梯度：
\[
  \nabla z=\left(\frac{2x}{x^2+y^2},\frac{2y}{x^2+y^2}\right).
\]
在 \((1,1)\) 处，\(\nabla z(1,1)=(1,1)\)。方向
\(\displaystyle \vec n=(\cos\pi/3,\sin\pi/3)=\left(\frac12,\frac{\sqrt3}{2}\right)\)
已经是单位向量，因此
\[
  \left.\frac{\partial z}{\partial \vec n}\right|_{(1,1)}
  =\nabla z(1,1)\cdot\vec n
  =\frac12+\frac{\sqrt3}{2}
  =\frac{1+\sqrt3}{2}.
\]
\end{solution}
\begin{exercise}
\par
设 $\Omega=\{(x,y,z):x^2+y^2+z^2\le 1\}$，则 \(\displaystyle \iiint_\Omega (x^2+y^2+z^2)\,\mathrm{d}V=\underline{\hspace{4cm}}\)。
\end{exercise}

\begin{solution}
\par
在球坐标中 \(x^2+y^2+z^2=r^2\)，\(\mathrm{d}V=r^2\sin\varphi\,\mathrm{d}r\,\mathrm{d}\varphi\,\mathrm{d}\theta\)。单位球区域为
\[
  0\le r\le1,\qquad 0\le\varphi\le\pi,\qquad 0\le\theta\le2\pi.
\]
故
\[
  \iiint_\Omega (x^2+y^2+z^2)\,\mathrm{d}V
  =\int_0^{2\pi}\int_0^\pi\int_0^1 r^4\sin\varphi\,\mathrm{d}r\,\mathrm{d}\varphi\,\mathrm{d}\theta
  =4\pi\cdot\frac15
  =\frac{4\pi}{5}.
\]
\end{solution}
\begin{exercise}
\par
设 \(C\) 是从点 \(A=(3,2,1)\) 到点 \(B=(0,0,0)\) 的直线段。则
  \[
    \int_C x^3\,\mathrm{d}x+y^3\,\mathrm{d}y+z^3\,\mathrm{d}z=\underline{\hspace{4cm}}.
  \]
\end{exercise}

\begin{solution}
\par
被积式是全微分：
\[
  x^3\,\mathrm{d}x+y^3\,\mathrm{d}y+z^3\,\mathrm{d}z
  =\mathrm{d}\left(\frac{x^4+y^4+z^4}{4}\right).
\]
因此积分与路径无关，只与端点有关：
\[
  \int_C x^3\,\mathrm{d}x+y^3\,\mathrm{d}y+z^3\,\mathrm{d}z
  =\left[\frac{x^4+y^4+z^4}{4}\right]_{(3,2,1)}^{(0,0,0)}
  =-\frac{81+16+1}{4}
  =-\frac{49}{2}.
\]
\end{solution}
\begin{exercise}
\par
设 $f(x)$ 是以 $2\pi$ 为周期的函数，它在 $[-\pi,\pi)$ 上的表达式是
  \[
    f(x)=
    \begin{cases}
      x+1, & -\pi\le x<0,\\
      x^2, & 0\le x<\pi.
    \end{cases}
  \]
  若它的 Fourier 级数的和函数是 $S(x)$，则 $S(-\pi)=\underline{\hspace{4cm}}$。
\end{exercise}

\begin{solution}
\par
傅里叶级数在周期端点处收敛到周期延拓后左右极限的平均值。对 \(x=-\pi\)，右极限来自
\(-\pi<x<0\) 上的表达式：
\[
  f(-\pi+0)=-\pi+1.
\]
左极限由周期性等于 \(x=\pi-0\) 处的极限，而 \(0\le x<\pi\) 时 \(f(x)=x^2\)，所以
\[
  f(-\pi-0)=f(\pi-0)=\pi^2.
\]
故
\[
  S(-\pi)=\frac{f(-\pi-0)+f(-\pi+0)}2
  =\frac{\pi^2-\pi+1}{2}.
\]
\end{solution}

\par\medskip
\noindent\textbf{二、计算题（共 4 小题，每小题 9 分，共 36 分）}\par\nobreak\smallskip
\begin{exercise}
\par
求方程 $y''-y=\sin^2x$ 的通解。
\end{exercise}

\begin{solution}
\par
先解齐次方程 \(y''-y=0\)，特征根为 \(r=\pm1\)，故
\[
  y_h=C_1e^x+C_2e^{-x}.
\]
又
\[
  \sin^2x=\frac{1-\cos2x}{2}.
\]
对常数项 \(1/2\)，设特解为常数 \(A\)，代入 \(y''-y=1/2\) 得 \(-A=1/2\)，故
\(A=-1/2\)。对 \(-\frac12\cos2x\)，设特解为 \(B\cos2x\)，则
\[
  (B\cos2x)''-B\cos2x=-5B\cos2x.
\]
令 \(-5B=-1/2\)，得 \(B=1/10\)。因此通解为
\[
  y=C_1e^x+C_2e^{-x}-\frac12+\frac1{10}\cos2x.
\]
\end{solution}
\begin{exercise}
\par
设 $z=\arctan(x+y)$，$x=2u-v^2$，$y=u^2v$，求 $\displaystyle \dfrac{\partial z}{\partial u}$，$\displaystyle \dfrac{\partial z}{\partial v}$。
\end{exercise}

\begin{solution}
\par
令
\[
  s=x+y=2u-v^2+u^2v,
\qquad z=\arctan s.
\]
由链式法则，
\[
  z_u=\frac{s_u}{1+s^2},\qquad z_v=\frac{s_v}{1+s^2}.
\]
而
\[
  s_u=2+2uv,\qquad s_v=-2v+u^2.
\]
所以
\[
  \frac{\partial z}{\partial u}=\frac{2+2uv}{1+(2u-v^2+u^2v)^2},\qquad
  \frac{\partial z}{\partial v}=\frac{-2v+u^2}{1+(2u-v^2+u^2v)^2}.
\]
\end{solution}
\begin{exercise}
\par
设 $D$ 是以 $(0,0)$，$(\pi,-\pi)$，$(\pi,\pi)$ 为顶点的三角形区域，求
  \(\displaystyle \iint_D \cos(x-y)\sin(x+y)\,\mathrm{d}x\,\mathrm{d}y\)。
\end{exercise}

\begin{solution}
\par
作变换 $u=x-y,\ v=x+y$，则 $\displaystyle \mathrm{d}x\,\mathrm{d}y=\frac12\,\mathrm{d}u\,\mathrm{d}v$，区域变为
  \(\displaystyle u\ge0,\quad v\ge0,\quad u+v\le2\pi\)。
  因此
\[
  I=\frac12\int_0^{2\pi}\int_0^{2\pi-u}\cos u\sin v\,\mathrm{d}v\,\mathrm{d}u.
\]
先对 \(v\) 积分：
\[
  \int_0^{2\pi-u}\sin v\,\mathrm{d}v
  =1-\cos(2\pi-u)=1-\cos u.
\]
于是
\[
  I=\frac12\int_0^{2\pi}\cos u(1-\cos u)\,\mathrm{d}u
  =-\frac12\int_0^{2\pi}\cos^2u\,\mathrm{d}u
  =-\frac{\pi}{2}.
\]
\end{solution}
\begin{exercise}
\par
设 $\Omega\subset\mathbb R^3$ 是由曲面 $z=5-x^2-y^2$ 和 $z+1=\sqrt{x^2+y^2}$ 所围成的空间立体，且 $\Omega$ 内布满某种流体，其流速为 \(\displaystyle \vec F=\frac{\vec r}{\|\vec r\|^3},\qquad \vec r=(x,y,z)\)。求单位时间内流体通过 $\Omega$ 的表面 $\partial\Omega$ 的流量。
\end{exercise}

\begin{solution}
\par
两曲面分别为 \(z=5-r^2\) 与 \(z=r-1\)，其中 \(r=\sqrt{x^2+y^2}\)。当 \(r=0\) 时上下界为
\(-1\) 与 \(5\)，所以原点在 \(\Omega\) 内。向量场
\[
  \vec F=\frac{\vec r}{\|\vec r\|^3}
\]
在原点有孤立奇点，在原点以外散度为 \(0\)。挖去原点附近小球
\(B_\varepsilon\)，在 \(\Omega\setminus B_\varepsilon\) 上用高斯公式，外边界通量等于小球面按通常外法向的通量：
\[
  \iint_{\partial B_\varepsilon}\frac{\vec r}{\|\vec r\|^3}\cdot\vec n\,\mathrm{d}S
  =\iint_{\partial B_\varepsilon}\frac1{\varepsilon^2}\,\mathrm{d}S
  =4\pi.
\]
故通过 \(\partial\Omega\) 的流量为 \(4\pi\)。
\end{solution}

\par\medskip
\noindent\textbf{三、解答题（共 3 小题，每小题 9 分，共 27 分）}\par\nobreak\smallskip
\begin{exercise}
\par
求级数 $\displaystyle \sum_{n=1}^\infty \frac{2n-1}{2^n}$ 的和。
\end{exercise}

\begin{solution}
\par
利用
\[
  \sum_{n=1}^{\infty}nx^n=\frac{x}{(1-x)^2},\qquad |x|<1.
\]
取 \(x=1/2\)，得 \(\displaystyle \sum_{n=1}^{\infty}\frac{n}{2^n}=2\)。同时
\(\displaystyle \sum_{n=1}^{\infty}\frac1{2^n}=1\)。因此
\[
  \sum_{n=1}^{\infty}\frac{2n-1}{2^n}
  =2\sum_{n=1}^{\infty}\frac n{2^n}
   -\sum_{n=1}^{\infty}\frac1{2^n}
  =4-1=3.
\]
\end{solution}
\begin{exercise}
\par
求曲线积分
  \(\displaystyle \oint_L y\,\mathrm{d}x+z\,\mathrm{d}y+x\,\mathrm{d}z\)，
  其中 $L$ 为圆周
  \[
    \begin{cases}
      x^2+y^2+z^2=a^2,\quad a>0,\\
      x+y+z=0,
    \end{cases}
  \]
  且从 $x$ 轴正向看去，这个圆周取逆时针方向。
\end{exercise}

\begin{solution}
\par
令 $\vec F=(y,z,x)$，则 $\nabla\times\vec F=(-1,-1,-1)$。按题设方向取单位法向
  \(\displaystyle \vec n=\frac{(1,1,1)}{\sqrt3}\)。
  交线是过原点平面 \(x+y+z=0\) 截球 \(x^2+y^2+z^2=a^2\) 所得圆，半径为 \(a\)。由 Stokes 公式，
\[
  \oint_L y\,\mathrm{d}x+z\,\mathrm{d}y+x\,\mathrm{d}z
  =\iint_{\Sigma}(\nabla\times\vec F)\cdot\vec n\,\mathrm{d}S
  =(-1,-1,-1)\cdot\frac{(1,1,1)}{\sqrt3}\,\pi a^2
  =-\sqrt3\,\pi a^2.
\]
\end{solution}
\begin{exercise}
\par
给定 \(\displaystyle \mathrm{d}u=\frac{(x+2y)\,\mathrm{d}x+y\,\mathrm{d}y}{(x+y)^2}\)，试验证势函数 $u$ 的存在性，并求出 $u$。
\end{exercise}

\begin{solution}
\par
记 \(\displaystyle P=\frac{x+2y}{(x+y)^2},\ Q=\frac{y}{(x+y)^2}\)。有 $P_y=Q_x=-2y/(x+y)^3$，故势函数存在。对 $P$ 关于 $x$ 积分，得 \(\displaystyle u=\ln|x+y|-\frac{y}{x+y}+C\)。
\[
  \int \frac{x+2y}{(x+y)^2}\,\mathrm{d}x
  =\int\left(\frac1{x+y}+\frac{y}{(x+y)^2}\right)\,\mathrm{d}x
  =\ln|x+y|-\frac{y}{x+y}+\phi(y).
\]
再对 \(y\) 求导并与 \(Q\) 比较，可得 \(\phi'(y)=0\)。所以势函数为
\[
  u=\ln|x+y|-\frac{y}{x+y}+C.
\]
\end{solution}

\par\medskip
\noindent\textbf{四、应用题（9 分）}\par\nobreak\smallskip
\begin{exercise}
\par
求曲线
\[
  \begin{cases}
    z=x^2+y^2,\\
    y=\dfrac1x
  \end{cases}
\]
上到 $Oxy$ 平面距离最近的点。
\end{exercise}

\begin{solution}
\par
到 \(Oxy\) 平面的距离为 \(|z|\)。曲线上 \(y=1/x\)，所以
\[
  z=x^2+y^2=x^2+\frac1{x^2}.
\]
令 \(t=x^2>0\)，则
\[
  z=t+\frac1t\ge2,
\]
等号当且仅当 \(t=1\)，即 \(x=\pm1\) 时成立。对应 \(y=1/x\)，所以最近点为
\[
  (1,1,2),\qquad (-1,-1,2).
\]
由于 \(z=x^2+1/x^2>0\)，距离 \(|z|\) 就等于 \(z\)；上面的不等式给出的是全局最小值，而不是局部驻点判断。
\end{solution}

\par\medskip
\noindent\textbf{五、证明题（共 2 小题，每小题 9 分，共 18 分）}\par\nobreak\smallskip
\begin{exercise}
\par
证明曲线
  \[
    \begin{cases}
      x^2+y^2+z^2=6,\\
      x+y+z=0
    \end{cases}
  \]
  在点 $(1,-2,1)$ 处的法平面方程为 $x-z=0$。
\end{exercise}

\begin{solution}
\par
设 \(\displaystyle F=x^2+y^2+z^2-6,\ G=x+y+z\)。在 $(1,-2,1)$ 处，\(\displaystyle \nabla F=(2,-4,2),\ \nabla G=(1,1,1)\)，
  切向量可取 $\nabla F\times\nabla G=(-6,0,6)$，即平行于 $(1,0,-1)$。法平面垂直于切向量，故
  \(\displaystyle (1,0,-1)\cdot(x-1,y+2,z-1)=0\)，
  即 $x-z=0$。
\end{solution}
\begin{exercise}
\par
证明级数
  \(\displaystyle \sum_{n=3}^\infty \ln\left(1+\frac{x}{n\ln^2 n}\right)\) 在 $[-2,2]$ 内一致收敛。
\end{exercise}

\begin{solution}
\par
当 $|x|\le2$ 时，各项均有定义。对充分大的 $n$，\(\displaystyle \left|\ln\left(1+\frac{x}{n\ln^2n}\right)\right|\le C\frac{|x|}{n\ln^2n}\le \frac{2C}{n\ln^2n}\)。
  因 $\displaystyle \sum_{n=3}^\infty 1/(n\ln^2n)$ 收敛，有限个初始项不影响一致收敛，由 Weierstrass 判别法知原级数在 $[-2,2]$ 上一致收敛。
\end{solution}

\section{2023级工科数学分析（二）期末考试 B 卷}

\par\medskip
\noindent\textbf{一、填空题（共 5 小题，每小题 2 分，共 10 分）}\par\nobreak\smallskip
\begin{exercise}
\par
方程 $y'+y\tan x=\cos x$ 的通解为 \underline{\hspace{4cm}}。
\end{exercise}

\begin{solution}
\par
这是线性方程
\[
  y'+(\tan x)y=\cos x.
\]
积分因子为 \(\displaystyle e^{\int\tan x\,\mathrm{d}x}=\sec x\)。两边乘以 \(\sec x\)，得
\[
  (y\sec x)'=1.
\]
积分得 \(y\sec x=x+C\)，故
\[
  y=(x+C)\cos x.
\]
这个通解在任一不跨越 \(\cos x=0\) 的区间上成立；积分因子中把 \(e^{-\ln|\cos x|}\) 写成 \(\sec x\) 时，本质上已经固定了这样的局部区间。
\end{solution}
\begin{exercise}
\par
设 $z=e^{xy}$，$\displaystyle \vec n=\left(\cos\dfrac{\pi}{4},\sin\dfrac{\pi}{4}\right)$。则 \(\displaystyle \left.\frac{\partial z}{\partial \vec n}\right|_{(1,1)}=\underline{\hspace{4cm}}\)。
\end{exercise}

\begin{solution}
\par
梯度为
\[
  \nabla z=(ye^{xy},xe^{xy}).
\]
在 \((1,1)\) 处，\(\nabla z(1,1)=(e,e)\)。方向
\(\displaystyle \vec n=(\cos\pi/4,\sin\pi/4)=\left(\frac{\sqrt2}{2},\frac{\sqrt2}{2}\right)\)，故
\[
  \left.\frac{\partial z}{\partial\vec n}\right|_{(1,1)}
  =e\frac{\sqrt2}{2}+e\frac{\sqrt2}{2}
  =\sqrt2\,e.
\]
\end{solution}
\begin{exercise}
\par
\(\displaystyle \int_0^1 \mathrm{d}x\int_0^{\sqrt{1-x^2}}x^3\,\mathrm{d}y=\underline{\hspace{4cm}}\)。
\end{exercise}

\begin{solution}
\par
先对 \(y\) 积分，内层长度为 \(\sqrt{1-x^2}\)，所以
\[
  \int_0^1\mathrm{d}x\int_0^{\sqrt{1-x^2}}x^3\,\mathrm{d}y
  =\int_0^1x^3\sqrt{1-x^2}\,\mathrm{d}x.
\]
令 \(u=1-x^2\)，则 \(x^3\,\mathrm{d}x=x^2\cdot x\,\mathrm{d}x=(1-u)(-\mathrm{d}u/2)\)，得
\[
  \int_0^1x^3\sqrt{1-x^2}\,\mathrm{d}x
  =\frac12\int_0^1(1-u)u^{1/2}\,\mathrm{d}u
  =\frac12\left(\frac23-\frac25\right)
  =\frac{2}{15}.
\]
\end{solution}
\begin{exercise}
\par
设曲线的参数方程为 \(\displaystyle x=at,\qquad y=\frac a2t^2,\qquad z=\frac a3t^3,\qquad 0\le t\le 1\)，且其密度函数为 $\sqrt{2y/a}$，则曲线的质量为 \underline{\hspace{4cm}}。
\end{exercise}

\begin{solution}
\par
由参数方程，
\[
  (x',y',z')=(a,at,at^2),
\]
所以
\[
  \mathrm{d}s=a\sqrt{1+t^2+t^4}\,\mathrm{d}t.
\]
又 \(\rho=\sqrt{2y/a}=\sqrt{t^2}=t\)（因 \(0\le t\le1\)）。质量为
\[
  M=a\int_0^1t\sqrt{t^4+t^2+1}\,\mathrm{d}t.
\]
令 \(s=t^2\)，则
\[
  M=\frac a2\int_0^1\sqrt{s^2+s+1}\,\mathrm{d}s
  =\frac a2\int_0^1\sqrt{\left(s+\frac12\right)^2+\frac34}\,\mathrm{d}s.
\]
使用公式 \(\displaystyle \int\sqrt{w^2+\alpha^2}\,\mathrm{d}w=\frac12\left(w\sqrt{w^2+\alpha^2}+\alpha^2\ln|w+\sqrt{w^2+\alpha^2}|\right)\)，代入上下限得
\[
  M=\frac{a(3\sqrt3-1)}{8}
    +\frac{3a}{16}\ln\frac{3+2\sqrt3}{3}.
\]
\end{solution}
\begin{exercise}
\par
设 $f(x)$ 是以 $2\pi$ 为周期的函数，它在 $[-\pi,\pi)$ 上的表达式是 $f(x)=x$。则 $f(x)$ 的 Fourier 级数中 $\sin 3x$ 的系数是 \underline{\hspace{4cm}}。
\end{exercise}

\begin{solution}
\par
因为 \(f(x)=x\) 是奇函数，傅里叶级数只含正弦项。正弦系数为
\[
  b_n=\frac1\pi\int_{-\pi}^{\pi}x\sin nx\,\mathrm{d}x
  =\frac2\pi\int_0^\pi x\sin nx\,\mathrm{d}x.
\]
分部积分得
\[
  \int_0^\pi x\sin nx\,\mathrm{d}x
  =\left[-\frac{x\cos nx}{n}\right]_0^\pi
  =-\frac{\pi(-1)^n}{n}.
\]
所以
\[
  b_n=\frac{2(-1)^{n+1}}n.
\]
取 \(n=3\)，得 \(\sin3x\) 的系数为 \(b_3=2/3\)。
\end{solution}

\par\medskip
\noindent\textbf{二、计算题（共 4 小题，每小题 9 分，共 36 分）}\par\nobreak\smallskip
\begin{exercise}
\par
求方程 $y''+y=3xe^x$ 的通解。
\end{exercise}

\begin{solution}
\par
齐次方程 \(y''+y=0\) 的通解为
\[
  y_h=C_1\cos x+C_2\sin x.
\]
右端为 \(3xe^x\)，设特解
\[
  y_p=e^x(Ax+B).
\]
代入 \(y''+y\)，得
\[
  y_p''+y_p=e^x(2Ax+2A+2B).
\]
与 \(3xe^x\) 比较系数，得 \(2A=3,\ 2A+2B=0\)，故
\(A=3/2,\ B=-3/2\)。因此
\[
  y=C_1\cos x+C_2\sin x+\frac32(x-1)e^x.
\]
\end{solution}
\begin{exercise}
\par
设 $\displaystyle z=\sin\left(xy+\dfrac{x}{y}\right)$，求 $\displaystyle \dfrac{\partial^2z}{\partial x\partial y}$。
\end{exercise}

\begin{solution}
\par
令 \(\displaystyle w=xy+\frac{x}{y}\)，则
\[
  w_x=y+\frac1y,\qquad w_y=x-\frac{x}{y^2}.
\]
先对 \(x\) 求偏导：
\[
  z_x=\cos w\cdot w_x=\left(y+\frac1y\right)\cos w.
\]
再对 \(y\) 求导，前一因子的导数为 \(1-\dfrac1{y^2}\)，后一因子的导数为 \(-\sin w\cdot w_y\)。因此
\[
  z_{xy}
  =
  \left(1-\frac1{y^2}\right)\cos w
  -\left(y+\frac1y\right)\left(x-\frac{x}{y^2}\right)\sin w.
\]
最后把 \(w=xy+\dfrac{x}{y}\) 代回即可。
\end{solution}
\begin{exercise}
\par
设 $D$ 是由 $x^2+y^2=9$，$x^2+y^2=16$ 以及 $y^2-x^2=1$，$y^2-x^2=9$ 围成的位于第一象限的区域，求
  \(\displaystyle \iint_D xy\,\mathrm{d}x\,\mathrm{d}y\)。
\end{exercise}

\begin{solution}
\par
作变换 \(\displaystyle u=x^2+y^2,\ v=y^2-x^2\)。
  其 Jacobian 为
\[
  \frac{\partial(u,v)}{\partial(x,y)}
  =
  \begin{vmatrix}
  2x&2y\\
  -2x&2y
  \end{vmatrix}
  =8xy.
\]
在第一象限 \(xy>0\)，故 \(\mathrm{d}u\,\mathrm{d}v=8xy\,\mathrm{d}x\,\mathrm{d}y\)。四条边界分别化为
\[
  9\le u\le16,\qquad 1\le v\le9.
\]
于是
\[
  \iint_Dxy\,\mathrm{d}x\,\mathrm{d}y
  =\frac18\int_9^{16}\int_1^9\mathrm{d}u\,\mathrm{d}v
  =\frac18(16-9)(9-1)=7.
\]
\end{solution}
\begin{exercise}
\par
设 $\Omega\subset\mathbb R^3$ 以原点为内点，且边界曲面 $\partial\Omega$ 光滑，又设边界曲面单位外法向为 $\vec n$，向量场 \(\displaystyle \vec F=\frac{\vec r}{\|\vec r\|^3},\qquad \vec r=(x,y,z)\)。求
  \(\displaystyle \iint_{\partial\Omega}\vec F\cdot\vec n\,\mathrm{d}s\)。
\end{exercise}

\begin{solution}
\par
因 $\Omega$ 以原点为内点，而 $\vec F=\vec r/\|\vec r\|^3$ 在原点有孤立奇点，去掉原点附近小球并用高斯公式可得
  \[
  \iint_{\partial\Omega}\vec F\cdot\vec n\,\mathrm{d}S
  =\iint_{\partial B_\varepsilon}\frac{\vec r}{\|\vec r\|^3}\cdot\vec n\,\mathrm{d}S
  =\iint_{\partial B_\varepsilon}\frac1{\varepsilon^2}\,\mathrm{d}S
  =4\pi.
  \]
\end{solution}

\par\medskip
\noindent\textbf{三、解答题（共 3 小题，每小题 9 分，共 27 分）}\par\nobreak\smallskip
\begin{exercise}
\par
求级数 $\displaystyle \sum_{n=1}^\infty \frac1{4n(4n+2)}$ 的和。
\end{exercise}

\begin{solution}
\par
\(\displaystyle \sum_{n=1}^\infty\frac1{4n(4n+2)}
    =\frac18\sum_{n=1}^\infty\frac1{n(2n+1)}
    \)。
分解
\[
  \frac1{n(2n+1)}=\frac1n-\frac2{2n+1}.
\]
于是部分和为
\[
  \frac18\sum_{n=1}^{N}\left(\frac1n-\frac2{2n+1}\right).
\]
利用调和级数奇偶项拆分，令 \(N\to\infty\) 可得
\[
  \sum_{n=1}^{\infty}\frac1{n(2n+1)}=2(1-\ln2).
\]
因此原级数的和为
\[
  \frac18\cdot2(1-\ln2)=\frac{1-\ln2}{4}.
\]
\end{solution}
\begin{exercise}
\par
原卷配图表示第一卦限内的平面三角形及其法向；本题条件如下。设有向曲面 \(\displaystyle \Sigma=\{(x,y,z):x+y+z=1,\ x\ge 0,\ y\ge 0,\ z\ge 0\}\)，且其方向与 $(1,1,1)$ 同向。求力场 $\vec F=(y^2,z^2,x^2)$ 绕 $\Sigma$ 正向边界 $\partial\Sigma$ 一周所做的功。
\end{exercise}

\begin{solution}
\par
令 $\vec F=(y^2,z^2,x^2)$，则 \(\displaystyle \nabla\times\vec F=(-2z,-2x,-2y)\)。
  在平面 $x+y+z=1$ 上取与 $(1,1,1)$ 同向的单位法向量，则
  \(\displaystyle (\nabla\times\vec F)\cdot\vec n=-\frac2{\sqrt3}\)。
  三角形面积为 $\sqrt3/2$，由 Stokes 公式得所做的功为 $-1$。本题原卷有图，当前转写按题干给出的正向约定处理。
\end{solution}
\begin{exercise}
\par
已知 $\phi(\pi)=1$，试确定 $\phi(x)$，使得曲线积分 \(\displaystyle I=\int_A^B \left[\sin x-\phi(x)\right]\frac{y}{x}\,\mathrm{d}x+\phi(x)\,\mathrm{d}y\) 与路径无关，并求当 $A,B$ 两点分别为 $(1,0)$，$(\pi,\pi)$ 时积分 $I$ 的值。
\end{exercise}

\begin{solution}
\par
记
  \(\displaystyle P=\left[\sin x-\phi(x)\right]\frac{y}{x},\qquad Q=\phi(x)\)。
  路径无关要求
  \(\displaystyle P_y=Q_x,\qquad \phi'(x)+\frac{\phi(x)}{x}=\frac{\sin x}{x}\)。
  因而 $(x\phi(x))'=\sin x$。由 $\phi(\pi)=1$ 得 \(\displaystyle \phi(x)=\frac{\pi-1-\cos x}{x}\)。
  势函数可取 $U(x,y)=\phi(x)y$，故
  \(\displaystyle I=U(\pi,\pi)-U(1,0)=\pi\)。
\end{solution}

\par\medskip
\noindent\textbf{四、应用题（9 分）}\par\nobreak\smallskip
\begin{exercise}
\par
在力场 $\vec F=(yz,zx,xy)$ 的作用下，质点由原点沿直线运动到椭球面 \(\displaystyle \frac{x^2}{a^2}+\frac{y^2}{b^2}+\frac{z^2}{c^2}=1\) 上第一卦限的点 $M=(\xi,\eta,\zeta)$。问当 $\xi,\eta,\zeta$ 取何值时力场 $\vec F$ 做的功最大？最大值是多少？
\end{exercise}

\begin{solution}
\par
沿直线从原点到 $M=(\xi,\eta,\zeta)$，取参数 $\vec r(t)=tM$，$0\le t\le1$。则 \(\displaystyle W=\int_0^1 3t^2\xi\eta\zeta\,\mathrm{d}t=\xi\eta\zeta\)。
在约束 \(\displaystyle \frac{\xi^2}{a^2}+\frac{\eta^2}{b^2}+\frac{\zeta^2}{c^2}=1\) 下，令 $X=\xi/a,\ Y=\eta/b,\ Z=\zeta/c$，由对称性或拉格朗日乘数法知 $XYZ$ 在 \(\displaystyle X=Y=Z=\frac1{\sqrt3}\)
处最大。故
  \(\displaystyle M=\left(\frac a{\sqrt3},\frac b{\sqrt3},\frac c{\sqrt3}\right),\qquad W_{\max}=\frac{abc}{3\sqrt3}\)。
\end{solution}

\par\medskip
\noindent\textbf{五、证明题（共 2 小题，每小题 9 分，共 18 分）}\par\nobreak\smallskip
\begin{exercise}
\par
证明曲线
  \[
    \begin{cases}
      x^2+y^2+z^2-3x=0,\\
      2x-3y+5z-4=0
    \end{cases}
  \]
  在点 $(1,1,1)$ 处的法平面方程为 $16x+9y-z-24=0$。
\end{exercise}

\begin{solution}
\par
设 \(\displaystyle F=x^2+y^2+z^2-3x,\ G=2x-3y+5z-4\)。在 $(1,1,1)$ 处，\(\displaystyle \nabla F=(-1,2,2),\ \nabla G=(2,-3,5)\)。切向量可取 \(\displaystyle \nabla F\times\nabla G=(16,9,-1)\)。
  法平面垂直于切向量，故
  \(\displaystyle 16(x-1)+9(y-1)-(z-1)=0\)，
  即 $16x+9y-z-24=0$。
\end{solution}
\begin{exercise}
\par
证明级数 \(\displaystyle \sum_{n=1}^\infty \sin\frac{\ln n^x}{n^3+x^2}\) 在 $(-\infty,+\infty)$ 内一致收敛。
\end{exercise}

\begin{solution}
\par
因 $\ln n^x=x\ln n$，对任意 $x\in\mathbb R$ 有 \(\displaystyle \left|\sin\frac{\ln n^x}{n^3+x^2}\right|\le \frac{|x|\ln n}{n^3+x^2}\le \frac{\ln n}{2n^{3/2}}\)。
  而 $\displaystyle \sum_{n=1}^\infty \ln n/n^{3/2}$ 收敛，故由 Weierstrass 判别法，原级数在 $(-\infty,+\infty)$ 上一致收敛。
\end{solution}

\section{2024级工科数学分析（二）期末考试 A 卷}

\par\medskip
\noindent\textbf{一、填空题（共 5 小题，每小题 3 分，共 15 分）}\par\nobreak\smallskip
\begin{exercise}
\par
微分方程 \(y''-3y'+2y=3x-2e^x\) 的特解形式为 \underline{\hspace{4cm}}。
\end{exercise}

\begin{solution}
\par
齐次方程的特征根为 \(1,2\)。右端 \(3x\) 对应特解可设为 \(Ax+B\)；右端 \(-2e^x\) 中 \(1\) 是一重特征根，所以对应特解要乘以 \(x\)。因此特解形式为
\[
  y_p=Ax+B+Cxe^x.
\]
\end{solution}
\begin{exercise}
\par
函数 \(z=2x^2+y^2\) 在点 \((1,1)\) 沿着梯度方向的方向导数为 \underline{\hspace{4cm}}。
\end{exercise}

\begin{solution}
\par
\[
  \nabla z=(4x,2y),\qquad \nabla z(1,1)=(4,2).
\]
沿梯度方向的方向导数等于梯度的模，故所求为
\[
  |\nabla z(1,1)|=\sqrt{4^2+2^2}=2\sqrt5.
\]
\end{solution}
\begin{exercise}
\par
设
\[
  I=\int_0^1\mathrm{d}y\int_{-y}^{\sqrt{2y-y^2}} f(x,y)\,\mathrm{d}x,
\]
则交换积分次序后，\(I=\)\underline{\hspace{4cm}}。
\end{exercise}

\begin{solution}
\par
区域左边界为 \(x=-y\)，右边界满足
\[
  x^2+(y-1)^2=1,\qquad x\ge0.
\]
按 \(x\) 分段可得
\[
  I=\int_{-1}^{0}\mathrm{d}x\int_{-x}^{1}f(x,y)\,\mathrm{d}y
    +\int_{0}^{1}\mathrm{d}x\int_{1-\sqrt{1-x^2}}^{1}f(x,y)\,\mathrm{d}y.
\]
\end{solution}
\begin{exercise}
\par
设幂级数 \(\displaystyle \sum_{n=1}^{\infty}a_nx^n\) 的收敛半径为 \(3\)，则幂级数
\[
  \sum_{n=1}^{\infty}n a_n(x-1)^{n-1}
\]
的收敛区间为 \underline{\hspace{4cm}}。
\end{exercise}

\begin{solution}
\par
导数幂级数与原幂级数收敛半径相同。令 \(t=x-1\)，则该级数关于 \(t\) 的收敛半径为 \(3\)，即
\[
  |x-1|<3.
\]
所以收敛区间为 \((-2,4)\)。仅由题设不能确定端点 \(x=-2,4\) 处的收敛性。
\end{solution}
\begin{exercise}
\par
设 \(f(x)\) 是以 \(2\) 为周期的函数，它在区间 \((-1,1]\) 上的表达式为
\[
  f(x)=
  \begin{cases}
    2, & -1<x\le 0,\\
    x^3, & 0<x\le 1,
  \end{cases}
\]
则 \(f(x)\) 的 Fourier 级数在 \(x=1\) 处收敛于 \underline{\hspace{4cm}}。
\end{exercise}

\begin{solution}
\par
Fourier 级数在跳跃点收敛于左右极限的平均值。这里
\[
  f(1-0)=1,\qquad f(1+0)=f(-1+0)=2.
\]
故所求为
\[
  \frac{1+2}{2}=\frac32.
\]
\end{solution}

\par\medskip
\noindent\textbf{二、计算下列各题（共 8 小题，共 69 分）}\par\nobreak\smallskip
\begin{exercise}
\par
（本小题 8 分）求微分方程
\[
  y\ln y\,\mathrm{d}x+(x-\ln y)\,\mathrm{d}y=0
\]
的通解。
\end{exercise}

\begin{solution}
\par
把 \(x\) 看成 \(y\) 的函数，则
\[
  y\ln y\,\frac{\mathrm{d}x}{\mathrm{d}y}+x-\ln y=0,
\]
即
\[
  \frac{\mathrm{d}x}{\mathrm{d}y}+\frac{x}{y\ln y}=\frac1y.
\]
积分因子为
\[
  e^{\int\frac{\mathrm{d}y}{y\ln y}}=\ln y.
\]
因此
\[
  (x\ln y)'=\frac{\ln y}{y}.
\]
积分得通解
\[
  x\ln y-\frac12(\ln y)^2=C.
\]
\end{solution}
\begin{exercise}
\par
（本小题 8 分）设物体是曲面 \(\Sigma:z=\sqrt{x^2+y^2}\) 被柱面 \(z^2=2x\) 所截下的部分，
\(\Sigma\) 上任意一点处的密度为
\(\mu(x,y,z)=\sqrt{x^2+y^2+z^2}\)，求物体 \(\Sigma\) 的质量。
\end{exercise}

\begin{solution}
\par
在柱坐标中，曲面为 \(z=r\)，柱面 \(z^2=2x\) 化为
\[
  r^2=2r\cos\theta.
\]
故
\[
  -\frac{\pi}{2}\le\theta\le\frac{\pi}{2},\qquad 0\le r\le2\cos\theta.
\]
圆锥面 \(z=r\) 的面积元为
\[
  \mathrm{d}S=\sqrt2\,r\,\mathrm{d}r\,\mathrm{d}\theta,
\]
密度为
\[
  \mu=\sqrt{r^2+z^2}=\sqrt2\,r.
\]
所以质量为
\[
  M=\int_{-\pi/2}^{\pi/2}\int_0^{2\cos\theta}2r^2\,\mathrm{d}r\,\mathrm{d}\theta
  =\frac{16}{3}\int_{-\pi/2}^{\pi/2}\cos^3\theta\,\mathrm{d}\theta
  =\frac{64}{9}.
\]
\end{solution}
\begin{exercise}
\par
（本小题 8 分）计算
\[
  I=\oint_L\frac{(y+3x)\,\mathrm{d}x+(y-x)\,\mathrm{d}y}{3x^2+y^2},
\]
其中 \(L\) 为单位圆周，取逆时针方向。
\end{exercise}

\begin{solution}
\par
令
\[
  x=\cos t,\qquad y=\sin t,\qquad 0\le t\le2\pi.
\]
则
\[
  (y+3x)\,\mathrm{d}x+(y-x)\,\mathrm{d}y=-(1+\sin2t)\,\mathrm{d}t,
\]
\[
  3x^2+y^2=1+2\cos^2t.
\]
于是
\[
  I=-\int_0^{2\pi}\frac{\mathrm{d}t}{1+2\cos^2t}
    -\int_0^{2\pi}\frac{\sin2t}{1+2\cos^2t}\,\mathrm{d}t.
\]
后一项在整周期上为 \(0\)，而
\[
  \int_0^{2\pi}\frac{\mathrm{d}t}{1+2\cos^2t}=\frac{2\pi}{\sqrt3}.
\]
因此
\[
  I=-\frac{2\pi}{\sqrt3}.
\]
\end{solution}
\begin{exercise}
\par
（本小题 8 分）计算曲面积分
\[
  I=\iint_\Sigma x^3\,\mathrm{d}y\,\mathrm{d}z+y^3\,\mathrm{d}z\,\mathrm{d}x+z^3\,\mathrm{d}x\,\mathrm{d}y,
\]
其中 \(\Sigma\) 是上半球面
\[
  z=\sqrt{a^2-x^2-y^2}\qquad(a>0),
\]
取下侧。
\end{exercise}

\begin{solution}
\par
记 \(\vec F=(x^3,y^3,z^3)\)。上半球面的外法向为
\[
  \vec n=\frac{(x,y,z)}{a},
\]
而题设取下侧，即取 \(-\vec n\)。故
\[
  I=-\iint_\Sigma \frac{x^4+y^4+z^4}{a}\,\mathrm{d}S.
\]
由球面对称性，
\[
  \iint_\Sigma x^4\,\mathrm{d}S
  =\iint_\Sigma y^4\,\mathrm{d}S
  =\iint_\Sigma z^4\,\mathrm{d}S
  =\frac{2\pi a^6}{5}.
\]
所以
\[
  I=-\frac1a\cdot3\cdot\frac{2\pi a^6}{5}
  =-\frac{6\pi a^5}{5}.
\]
\end{solution}
\begin{exercise}
\par
（本小题 9 分）设 \(u=f(x^2-y^2,2xy)\)，其中 \(f(\xi,\eta)\) 具有二阶连续的偏导数，
且满足
\[
  \frac{\partial^2 f}{\partial \xi^2}
  +\frac{\partial^2 f}{\partial \eta^2}=0.
\]
求
\[
  \frac{\partial^2 u}{\partial x^2}
  +\frac{\partial^2 u}{\partial y^2}.
\]
\end{exercise}

\begin{solution}
\par
令
\[
  \xi=x^2-y^2,\qquad \eta=2xy.
\]
则
\[
  \xi_x=2x,\quad \xi_y=-2y,\quad \eta_x=2y,\quad \eta_y=2x.
\]
并且
\[
  \xi_x^2+\xi_y^2=\eta_x^2+\eta_y^2=4(x^2+y^2),
  \qquad
  \xi_x\eta_x+\xi_y\eta_y=0.
\]
又有
\[
  \xi_{xx}+\xi_{yy}=0,\qquad \eta_{xx}+\eta_{yy}=0.
\]
由链式法则，
\[
  u_{xx}+u_{yy}
  =4(x^2+y^2)\left(f_{\xi\xi}+f_{\eta\eta}\right)=0.
\]
故所求为 \(0\)。
\end{solution}
\begin{exercise}
\par
（本小题 9 分）设 \(f(x)\) 在 \((-\infty,+\infty)\) 内有连续导数，\(L\) 是从点
\(\displaystyle A\left(3,\frac23\right)\) 到点 \(B(1,2)\) 的直线段，计算曲线积分
\[
  I=\int_L
  \frac{1+y^2f(xy)}{y}\,\mathrm{d}x
  +\frac{x}{y^2}\bigl[y^2f(xy)-1\bigr]\,\mathrm{d}y.
\]
\end{exercise}

\begin{solution}
\par
设 \(\Phi'(t)=f(t)\)。被积表达式为
\[
  \left(\frac1y+yf(xy)\right)\mathrm{d}x
  +\left(xf(xy)-\frac{x}{y^2}\right)\mathrm{d}y
  =\mathrm{d}\left(\frac{x}{y}+\Phi(xy)\right).
\]
因此积分与路径无关。端点处
\[
  A:\quad \frac{x}{y}=\frac92,\quad xy=2,
  \qquad
  B:\quad \frac{x}{y}=\frac12,\quad xy=2.
\]
所以
\[
  I=\left(\frac12+\Phi(2)\right)-\left(\frac92+\Phi(2)\right)=-4.
\]
\end{solution}
\begin{exercise}
\par
（本小题 9 分）求幂级数
\[
  \sum_{n=1}^{\infty}\frac{n^2+1}{2^n n!}x^n
\]
的收敛域及和函数。
\end{exercise}

\begin{solution}
\par
令 \(t=x/2\)，则
\[
  S(x)=\sum_{n=1}^{\infty}\frac{n^2+1}{n!}t^n.
\]
利用指数级数，有
\[
  \sum_{n=1}^{\infty}\frac{nt^n}{n!}=te^t,\qquad
  \sum_{n=2}^{\infty}\frac{n(n-1)t^n}{n!}=t^2e^t.
\]
故
\[
  \sum_{n=1}^{\infty}\frac{n^2t^n}{n!}=(t^2+t)e^t,
  \qquad
  \sum_{n=1}^{\infty}\frac{t^n}{n!}=e^t-1.
\]
于是
\[
  S(x)=\left(t^2+t+1\right)e^t-1
  =\left(\frac{x^2}{4}+\frac{x}{2}+1\right)e^{x/2}-1.
\]
收敛域为 \((-\infty,+\infty)\)。
\end{solution}
\begin{exercise}
\par
（本小题 10 分）设在第一象限内从椭球面
\[
  \frac{x^2}{a^2}+\frac{y^2}{b^2}+\frac{z^2}{c^2}=1
\]
上任取一点 \(P(x,y,z)\) 作切平面，设切平面在三条坐标轴上的截距分别为 \(A,B,C\)。
问点 \(P\) 在何位置时能使 \(A+B+C\) 达到极小值。
\end{exercise}

\begin{solution}
\par
设切点为 \(P(x_0,y_0,z_0)\)，其中 \(x_0,y_0,z_0>0\)。切平面为
\[
  \frac{x_0x}{a^2}+\frac{y_0y}{b^2}+\frac{z_0z}{c^2}=1.
\]
三个坐标轴截距为
\[
  A=\frac{a^2}{x_0},\qquad B=\frac{b^2}{y_0},\qquad C=\frac{c^2}{z_0}.
\]
令
\[
  X=\frac{x_0}{a},\qquad Y=\frac{y_0}{b},\qquad Z=\frac{z_0}{c},
\]
则 \(X^2+Y^2+Z^2=1\)，且
\[
  A+B+C=\frac{a}{X}+\frac{b}{Y}+\frac{c}{Z}.
\]
用拉格朗日乘数法可得
\[
  -\frac{a}{X^2}+2\lambda X=0,\quad
  -\frac{b}{Y^2}+2\lambda Y=0,\quad
  -\frac{c}{Z^2}+2\lambda Z=0.
\]
于是
\[
  X:Y:Z=a^{1/3}:b^{1/3}:c^{1/3}.
\]
记 \(s=a^{2/3}+b^{2/3}+c^{2/3}\)，则
\[
  X=\frac{a^{1/3}}{\sqrt{s}},\qquad
  Y=\frac{b^{1/3}}{\sqrt{s}},\qquad
  Z=\frac{c^{1/3}}{\sqrt{s}}.
\]
故所求点为
\[
  P\left(
  \frac{a^{4/3}}{\sqrt{s}},
  \frac{b^{4/3}}{\sqrt{s}},
  \frac{c^{4/3}}{\sqrt{s}}
  \right),
  \qquad s=a^{2/3}+b^{2/3}+c^{2/3}.
\]
\end{solution}

\par\medskip
\noindent\textbf{三、证明题（共 2 小题，每小题 8 分，共 16 分）}\par\nobreak\smallskip
\begin{exercise}
\par
（本小题 8 分）设函数 \(f(x)\) 在 \(x=0\) 的邻域内有二阶连续的导数，且
\[
  \lim_{x\to0}\frac{f(x)}{x}=0,
\]
证明：级数
\[
  \sum_{n=1}^{\infty}f\left(\frac1n\right)
\]
绝对收敛。
\end{exercise}

\begin{solution}
\par
由
\[
  \lim_{x\to0}\frac{f(x)}{x}=0
\]
可得 \(f(0)=0\)，并且
\[
  f'(0)=\lim_{x\to0}\frac{f(x)-f(0)}{x}=0.
\]
由于 \(f\) 在 \(0\) 的邻域内二阶导数连续，由 Taylor 公式，充分小的 \(x\) 满足
\[
  f(x)=\frac12 f''(\xi_x)x^2,
\]
其中 \(\xi_x\) 介于 \(0\) 与 \(x\) 之间。因此存在常数 \(M>0\)，使得充分大的 \(n\) 有
\[
  \left|f\left(\frac1n\right)\right|\le \frac{M}{n^2}.
\]
由比较判别法，\(\sum f(1/n)\) 绝对收敛。
\end{solution}
\begin{exercise}
\par
（本小题 8 分）证明：
\[
  \sum_{n=1}^{\infty}\frac{\ln(1+nx)}{nx^n}
\]
在 \((1,+\infty)\) 上连续。
\end{exercise}

\begin{solution}
\par
记
\[
  u_n(x)=\frac{\ln(1+nx)}{nx^n}.
\]
每个 \(u_n(x)\) 在 \((1,+\infty)\) 上连续。任取闭区间 \([a,b]\subset(1,+\infty)\)，有 \(a>1\)，且对 \(x\in[a,b]\)，
\[
  0\le u_n(x)\le \frac{\ln(1+nb)}{na^n}.
\]
而
\[
  \sum_{n=1}^{\infty}\frac{\ln(1+nb)}{na^n}
\]
收敛，所以原函数项级数在 \([a,b]\) 上一致收敛。由连续函数项级数的一致收敛定理，和函数在 \([a,b]\) 上连续。
由于 \([a,b]\) 可任取，故该级数在 \((1,+\infty)\) 上连续。
\end{solution}

\section{2024级工科数学分析（二）期末考试 B 卷}

\par\medskip
\noindent\textbf{一、填空题（共 5 小题，每小题 3 分，共 15 分）}\par\nobreak\smallskip
\begin{exercise}
\par
已知线性齐次微分方程的基础解系为 \(y_1=e^x,\ y_2=xe^x\)，则该微分方程为 \underline{\hspace{4cm}}。
\end{exercise}

\begin{solution}
\par
基础解系 \(e^x,xe^x\) 对应特征根 \(r=1\) 为二重根，因此特征方程为
\[
  (r-1)^2=0.
\]
所求微分方程可写为
\[
  y''-2y'+y=0.
\]
\end{solution}
\begin{exercise}
\par
设数量场 \(u=\ln\sqrt{x^2+y^2+z^2}\)，则
\(\operatorname{div}(\operatorname{grad}u)=\)\underline{\hspace{4cm}}。
\end{exercise}

\begin{solution}
\par
令 \(r=\sqrt{x^2+y^2+z^2}\)，则 \(u=\ln r\)。三维径向函数满足
\[
  \Delta u=u''(r)+\frac2r u'(r).
\]
因为 \(u'(r)=1/r,\ u''(r)=-1/r^2\)，所以
\[
  \operatorname{div}(\operatorname{grad}u)=\Delta u
  =-\frac1{r^2}+\frac2{r^2}
  =\frac1{x^2+y^2+z^2}.
\]
\end{solution}
\begin{exercise}
\par
设
\[
  I=\int_1^2\mathrm{d}x\int_{\sqrt{x}}^x f(x,y)\,\mathrm{d}y
    +\int_2^4\mathrm{d}x\int_{\sqrt{x}}^2 f(x,y)\,\mathrm{d}y,
\]
则交换积分次序后，\(I=\)\underline{\hspace{4cm}}。
\end{exercise}

\begin{solution}
\par
第一部分对应区域
\[
  1\le x\le2,\qquad \sqrt{x}\le y\le x,
\]
即 \(y\le x\le y^2\)。第二部分对应区域
\[
  2\le x\le4,\qquad \sqrt{x}\le y\le2,
\]
即 \(2\le x\le y^2\)。两部分合并后为
\[
  1\le y\le2,\qquad y\le x\le y^2.
\]
因此
\[
  I=\int_1^2\mathrm{d}y\int_y^{y^2}f(x,y)\,\mathrm{d}x.
\]
\end{solution}
\begin{exercise}
\par
设幂级数 \(\displaystyle \sum_{n=1}^{\infty}a_nx^n\) 在 \(x=-3\) 条件收敛，则幂级数
\(\displaystyle \sum_{n=1}^{\infty}a_nx^n\) 的收敛半径为 \underline{\hspace{4cm}}。
\end{exercise}

\begin{solution}
\par
幂级数在 \(x=-3\) 收敛，说明收敛半径 \(R\ge3\)。又它在 \(x=-3\) 是条件收敛而不是绝对收敛；若 \(R>3\)，则 \(|x|=3\) 是收敛圆内部，必绝对收敛，矛盾。因此
\[
  R=3.
\]
\end{solution}
\begin{exercise}
\par
设 \(f(x)=x^2\ (0\le x<1)\)，而
\[
  S(x)=\sum_{n=1}^{\infty}b_n\sin n\pi x
  \qquad(-\infty<x<+\infty),
\]
其中
\[
  b_n=2\int_0^1 f(x)\sin n\pi x\,\mathrm{d}x
  \qquad(n=1,2,\cdots),
\]
则 \(\displaystyle S\left(-\frac12\right)=\)\underline{\hspace{4cm}}。
\end{exercise}

\begin{solution}
\par
这是 \(f(x)=x^2\) 在 \((0,1)\) 上的正弦级数，对应奇延拓且周期为 \(2\)。在连续点 \(x=-1/2\)，和函数等于奇延拓函数值：
\[
  S\left(-\frac12\right)=-f\left(\frac12\right)=-\frac14.
\]
\end{solution}

\par\medskip
\noindent\textbf{二、计算下列各题（共 8 小题，共 69 分）}\par\nobreak\smallskip
\begin{exercise}
\par
（本小题 8 分）求微分方程
\[
  x\sin(2y)\,y'+\sin^2y=\sin x
\]
的通解。
\end{exercise}

\begin{solution}
\par
令
\[
  u=\sin^2y.
\]
则
\[
  u'=\sin(2y)y'.
\]
原方程化为
\[
  xu'+u=\sin x,
\]
即
\[
  (xu)'=\sin x.
\]
积分得
\[
  x\sin^2y=-\cos x+C,
\]
故通解为
\[
  x\sin^2y+\cos x=C.
\]
\end{solution}
\begin{exercise}
\par
（本小题 8 分）求密度均匀的曲面
\[
  z=\sqrt{a^2-x^2-y^2}
\]
的质心坐标。
\end{exercise}

\begin{solution}
\par
该曲面是上半球面。由对称性，
\[
  \bar x=\bar y=0.
\]
设密度为常数，质心与密度大小无关。上半球面面积为
\[
  S=2\pi a^2.
\]
用球坐标参数化，\(z=a\cos\varphi\)，\(\mathrm{d}S=a^2\sin\varphi\,\mathrm{d}\varphi\,\mathrm{d}\theta\)，其中
\[
  0\le\varphi\le\frac{\pi}{2},\qquad 0\le\theta\le2\pi.
\]
于是
\[
  \iint_\Sigma z\,\mathrm{d}S
  =\int_0^{2\pi}\int_0^{\pi/2}a\cos\varphi\cdot a^2\sin\varphi\,\mathrm{d}\varphi\,\mathrm{d}\theta
  =\pi a^3.
\]
所以
\[
  \bar z=\frac{\pi a^3}{2\pi a^2}=\frac a2.
\]
质心坐标为
\[
  \left(0,0,\frac a2\right).
\]
\end{solution}
\begin{exercise}
\par
（本小题 8 分）计算
\[
  I=\int_L(e^x\sin y-y^3)\,\mathrm{d}x+(e^x\cos y+x^3)\,\mathrm{d}y,
\]
其中 \(L\) 为由 \(A(a,0)\) 到 \(O(0,0)\) 的上半圆周
\[
  x^2+y^2=ax.
\]
\end{exercise}

\begin{solution}
\par
被积式可拆成
\[
  \mathrm{d}(e^x\sin y)+(-y^3\,\mathrm{d}x+x^3\,\mathrm{d}y).
\]
在端点 \(A(a,0)\) 与 \(O(0,0)\) 处，\(e^x\sin y\) 均为 \(0\)，故只需计算第二部分。
上半圆周可参数化为
\[
  x=\frac a2(1+\cos t),\qquad y=\frac a2\sin t,\qquad 0\le t\le\pi.
\]
于是
\[
  \mathrm{d}x=-\frac a2\sin t\,\mathrm{d}t,\qquad \mathrm{d}y=\frac a2\cos t\,\mathrm{d}t.
\]
因此
\[
\begin{aligned}
  I
  &=\frac{a^4}{16}\int_0^\pi
    \left[\sin^4t+(1+\cos t)^3\cos t\right]\,\mathrm{d}t\\
  &=\frac{a^4}{16}\left(\frac{3\pi}{8}+\frac{3\pi}{2}+\frac{3\pi}{8}\right)
   =\frac{9\pi a^4}{64}.
\end{aligned}
\]
\end{solution}
\begin{exercise}
\par
（本小题 8 分）计算
\[
  I=\iint_\Sigma
  \frac{x^3\,\mathrm{d}y\,\mathrm{d}z+y^3\,\mathrm{d}z\,\mathrm{d}x+z^3\,\mathrm{d}x\,\mathrm{d}y}
       {\sqrt{x^2+y^2+z^2}},
\]
其中 \(\Sigma\) 为球面
\[
  x^2+y^2+z^2=a^2
\]
取外侧。
\end{exercise}

\begin{solution}
\par
球面上 \(\sqrt{x^2+y^2+z^2}=a\)，外法向为
\[
  \vec n=\frac{(x,y,z)}a.
\]
故
\[
  I=\iint_\Sigma \frac{x^4+y^4+z^4}{a^2}\,\mathrm{d}S.
\]
在整个球面上，由对称性
\[
  \iint_\Sigma x^4\,\mathrm{d}S
  =\iint_\Sigma y^4\,\mathrm{d}S
  =\iint_\Sigma z^4\,\mathrm{d}S
  =\frac{4\pi a^6}{5}.
\]
所以
\[
  I=\frac1{a^2}\cdot3\cdot\frac{4\pi a^6}{5}
  =\frac{12\pi a^4}{5}.
\]
\end{solution}
\begin{exercise}
\par
（本小题 9 分）设
\[
  u=yf\left(\frac{x}{y}\right)+xg\left(\frac{y}{x}\right),
\]
其中 \(f,g\) 具有二阶连续导数，求
\[
  x\frac{\partial^2u}{\partial x^2}
  +y\frac{\partial^2u}{\partial x\partial y}.
\]
\end{exercise}

\begin{solution}
\par
先看 \(u_1=yf(x/y)\)。令 \(s=x/y\)，则
\[
  (u_1)_x=f'(s),\qquad
  (u_1)_{xx}=\frac1y f''(s),\qquad
  (u_1)_{xy}=-\frac{x}{y^2}f''(s).
\]
因此
\[
  x(u_1)_{xx}+y(u_1)_{xy}
  =\frac{x}{y}f''(s)-\frac{x}{y}f''(s)=0.
\]
再看 \(u_2=xg(y/x)\)。令 \(t=y/x\)，则
\[
  (u_2)_x=g(t)-tg'(t),
\]
从而
\[
  (u_2)_{xx}=\frac{t^2}{x}g''(t),\qquad
  (u_2)_{xy}=-\frac{t}{x}g''(t).
\]
于是
\[
  x(u_2)_{xx}+y(u_2)_{xy}
  =t^2g''(t)-t\frac{y}{x}g''(t)=0.
\]
故所求为
\[
  x\frac{\partial^2u}{\partial x^2}
  +y\frac{\partial^2u}{\partial x\partial y}=0.
\]
\end{solution}
\begin{exercise}
\par
（本小题 9 分）设 \(L\) 是从点 \(A(-a,0)\) 经上半椭圆
\[
  \frac{x^2}{a^2}+\frac{y^2}{b^2}=1\qquad(y\ge0)
\]
到点 \(B(a,0)\) 的弧段，求
\[
  I=\int_L\frac{(x-y)\,\mathrm{d}x+(x+y)\,\mathrm{d}y}{x^2+y^2}.
\]
\end{exercise}

\begin{solution}
\par
把被积式分成两部分：
\[
  \frac{(x-y)\,\mathrm{d}x+(x+y)\,\mathrm{d}y}{x^2+y^2}
  =
  \frac{x\,\mathrm{d}x+y\,\mathrm{d}y}{x^2+y^2}
  +\frac{-y\,\mathrm{d}x+x\,\mathrm{d}y}{x^2+y^2}.
\]
第一项为
\[
  \mathrm{d}\left(\frac12\ln(x^2+y^2)\right),
\]
在两端点 \(A(-a,0)\)、\(B(a,0)\) 的取值相同，积分为 \(0\)。第二项为极角微分 \(\mathrm{d}\theta\)。
沿上半椭圆从 \((-a,0)\) 到 \((a,0)\) 时，极角由 \(\pi\) 连续减少到 \(0\)，故
\[
  I=0+(0-\pi)=-\pi.
\]
\end{solution}
\begin{exercise}
\par
（本小题 9 分）求函数 \(f(x)=\ln x\) 按分式
\[
  \frac{x-1}{x+1}
\]
的正整数次幂的展开式。
\end{exercise}

\begin{solution}
\par
令
\[
  t=\frac{x-1}{x+1},
\]
则
\[
  x=\frac{1+t}{1-t}.
\]
所以
\[
  \ln x=\ln(1+t)-\ln(1-t).
\]
当 \(|t|<1\) 时，
\[
  \ln(1+t)-\ln(1-t)
  =2\left(t+\frac{t^3}{3}+\frac{t^5}{5}+\cdots\right).
\]
代回 \(t=(x-1)/(x+1)\)，得
\[
  \ln x
  =2\sum_{k=0}^{\infty}\frac1{2k+1}
  \left(\frac{x-1}{x+1}\right)^{2k+1},
  \qquad x>0.
\]
\end{solution}
\begin{exercise}
\par
（本小题 10 分）试说明二次型
\[
  f(x,y,z)=Ax^2+By^2+Cz^2+2Dyz+2Ezx+2Fxy
\]
在单位球面
\[
  x^2+y^2+z^2=1
\]
上的最大值和最小值恰好是矩阵
\[
  M=
  \begin{bmatrix}
    A&F&E\\
    F&B&D\\
    E&D&C
  \end{bmatrix}
\]
的最大特征值和最小特征值。
\end{exercise}

\begin{solution}
\par
记
\[
  \mathbf x=
  \begin{bmatrix}
    x\\y\\z
  \end{bmatrix}.
\]
则二次型可写成
\[
  f(x,y,z)=\mathbf x^T M\mathbf x.
\]
因为 \(M\) 是实对称矩阵，存在正交矩阵 \(Q\)，使
\[
  Q^TMQ=\operatorname{diag}(\lambda_1,\lambda_2,\lambda_3),
\]
其中 \(\lambda_1,\lambda_2,\lambda_3\) 是 \(M\) 的特征值。令 \(\mathbf x=Q\mathbf y\)，则
\[
  x^2+y^2+z^2=\mathbf x^T\mathbf x=\mathbf y^T\mathbf y.
\]
在单位球面上，
\[
  f=\lambda_1y_1^2+\lambda_2y_2^2+\lambda_3y_3^2,
  \qquad
  y_1^2+y_2^2+y_3^2=1.
\]
设
\[
  \lambda_{\min}\le \lambda_i\le \lambda_{\max},
\]
则
\[
  \lambda_{\min}\le f\le \lambda_{\max}.
\]
当 \(\mathbf y\) 取最大特征值或最小特征值对应的单位特征向量时，上下界可取到。因此二次型在单位球面上的最大值和最小值分别为 \(M\) 的最大特征值和最小特征值。
\end{solution}

\par\medskip
\noindent\textbf{三、证明题（共 2 小题，每小题 8 分，共 16 分）}\par\nobreak\smallskip
\begin{exercise}
\par
（本小题 8 分）设偶函数 \(f(x)\) 在 \(x=0\) 的邻域内二阶导数连续，且
\[
  f(0)=1,\qquad f''(0)=2,
\]
证明：级数
\[
  \sum_{n=1}^{\infty}\left[f\left(\frac1n\right)-1\right]
\]
绝对收敛。
\end{exercise}

\begin{solution}
\par
因为 \(f(x)\) 是偶函数且在 \(0\) 的邻域内可导，所以
\[
  f'(0)=0.
\]
由 Taylor 公式，
\[
  f(x)=f(0)+f'(0)x+\frac12f''(0)x^2+o(x^2)
  =1+x^2+o(x^2).
\]
因此当 \(x\to0\) 时，
\[
  f(x)-1=O(x^2).
\]
取 \(x=1/n\)，有
\[
  f\left(\frac1n\right)-1=O\left(\frac1{n^2}\right).
\]
由比较判别法，
\[
  \sum_{n=1}^{\infty}\left[f\left(\frac1n\right)-1\right]
\]
绝对收敛。
\end{solution}
\begin{exercise}
\par
（本小题 8 分）设
\[
  F(x)=\sum_{n=1}^{\infty}\frac{\cos(nx)}{\sqrt{n^3+n}},
\]
求证：
\begin{enumerate}
  \item \(F(x)\) 在 \((-\infty,+\infty)\) 上连续；
  \item 若 \(g(x)\) 是 \(F(x)\) 的原函数，且 \(g(0)=0\)，则
  \[
    \frac1{\sqrt2}-\frac1{15}
    <g\left(\frac{\pi}{2}\right)
    <\frac1{\sqrt2}.
\]
\end{enumerate}
\end{exercise}

\begin{solution}
\par
对任意 \(x\)，有
\[
  \left|\frac{\cos(nx)}{\sqrt{n^3+n}}\right|
  \le \frac1{\sqrt{n^3+n}}
  \le \frac1{n^{3/2}}.
\]
而 \(\sum n^{-3/2}\) 收敛，所以由 Weierstrass 判别法，定义 \(F(x)\) 的函数项级数在 \((-\infty,+\infty)\) 上一致收敛。每一项连续，故 \(F(x)\) 在 \((-\infty,+\infty)\) 上连续。

又 \(g(0)=0\) 且 \(g'(x)=F(x)\)，所以
\[
  g\left(\frac{\pi}{2}\right)=\int_0^{\pi/2}F(x)\,\mathrm{d}x.
\]
由一致收敛可逐项积分，得
\[
  g\left(\frac{\pi}{2}\right)
  =\sum_{n=1}^{\infty}
  \frac{\sin(n\pi/2)}{n\sqrt{n^3+n}}.
\]
偶数项为 \(0\)，奇数项交错，故
\[
  g\left(\frac{\pi}{2}\right)
  =\frac1{\sqrt2}
  -\frac1{3\sqrt{30}}+\frac1{5\sqrt{130}}-\cdots .
\]
这是绝对值递减的交错级数，因此
\[
  \frac1{\sqrt2}-\frac1{3\sqrt{30}}
  <g\left(\frac{\pi}{2}\right)
  <\frac1{\sqrt2}.
\]
又
\[
  \frac1{3\sqrt{30}}<\frac1{15},
\]
于是
\[
  \frac1{\sqrt2}-\frac1{15}
  <g\left(\frac{\pi}{2}\right)
  <\frac1{\sqrt2}.
\]
\end{solution}

\end{document}

